\documentclass{situationist-preprint}
\usepackage{longtable}
\usepackage{array}

\title{On the Reality of Time}
\msid{TSM-01-a}
\msfamily{TSM-01}
\msdiscipline{Time, Space and Modality}
\mstype{Principal manuscript}
\mslean{lean/Philosophy/TSM/TSM\_01.lean}
\mskeywords{McTaggart; A-series; B-series; reality of time; tenseless theory; branching time; settledness; growing block; causal structure of spacetime}

% generated by gen/generate.py -- do not edit
\newcommand{\NFrames}{3047}
\newcommand{\NFramesDetail}{1205}
\newcommand{\NFramesEleven}{1842}
\newcommand{\NMoments}{11210}
\newcommand{\NMonoChecks}{129758668}
\newcommand{\NResolutions}{51747976}
\newcommand{\NReversals}{40535134}
\newcommand{\NRigid}{463}
\newcommand{\NNoFixH}{464}
\newcommand{\NHist}{14595}
\newcommand{\NHistRows}{1891}
\newcommand{\NLabelled}{409114}


\begin{document}
\maketitle

\begin{abstract}
McTaggart argued that time is unreal: there is no time without change, there is no change without the A-series of past, present and future, and the A-series is contradictory. We argue that the argument establishes less and more than he thought. It establishes that the A-series is incoherent: no world contains a changing A-series, and the only coherent way of making the A-determinations vary, relativizing them to moments, returns the order of earlier and later under another description. It does not establish that time is unreal, because its bridge premise, that time requires a changing A-series, owes its credibility to an equivocation between change in things over time and the change of A-determinations. Run in reverse from the reality of time, the argument refutes the bridge. Time is real and is constituted by the B-series, which we identify with the causal and metrical structure of spacetime. For an indeterministic world we add branching histories and moment-relative settledness, prove that settledness is monotone along every history, and exclude the growing block. The structural results are verified on all \NFrames{} branching frames with at most eleven moments, and the logical skeleton is formalized in Lean.
\end{abstract}

\section{Introduction}\label{s:intro}

Is time real? The question sounds idle until one meets the argument that it is not. In \citeyear{McTaggart1908} J.~M.~E. McTaggart argued that the positions of events in time are ordered in two ways: as past, present and future, which he called the A-series, and as earlier and later, which he called the B-series. He argued that time requires change, that change requires the A-series, and that the A-series is contradictory. It follows that nothing is in time \citep{McTaggart1908,McTaggart1927}.

Almost no one has accepted the conclusion, and almost everyone has accepted something from the argument. The literature since McTaggart can be organized by the premise each party rejects. Those who reject the third premise, that the A-series is contradictory, are A-theorists: they hold that one time is objectively present and that the present moves, and they owe a statement of the A-series that escapes McTaggart's regress \citep{Broad1938,Prior1967,Zimmerman2005,Cameron2015}. Those who reject the second premise, that change requires the A-series, are B-theorists: they hold that change is the variation of things across times and that the A-series is either a feature of our representation of time or a mistake \citep{Russell1903,Russell1915,Smart1963,Mellor1981,Mellor1998}. A third group accepts that the argument shows something about reality as a whole, that reality cannot be completely described from no temporal standpoint \citep{Dummett1960} or that it is not coherent as a whole \citep{Fine2005}.

We defend a version of the second response and give it a form that clarifies what the argument proves. Our thesis has two parts. First, McTaggart was right about the A-series: it is incoherent, and no world contains a changing A-series (Theorem~\ref{t:nochange}). Second, he was wrong about time: time is real and is constituted by the B-series, the order of events by earlier and later, which in our world is the causal and metrical structure of spacetime (Principle~\ref{p:constitution}). The argument is valid, and what its validity shows depends on the direction in which it is run. Run forward from the claim that time requires a changing A-series, it denies time. Run backward from the reality of time, for which our evidence is as good as for anything in physics, it denies that time requires a changing A-series (Theorem~\ref{t:inversion}). The backward direction is the rational one, because the bridge premise owes its appeal to an equivocation on \enquote{change} (Proposition~\ref{p:equivocation}).

A B-theory must also accommodate the open future. If the world is indeterministic, as current physics and our practice as agents both suggest, the B-series may branch toward the later: one past for each moment and several possible continuations. We develop the view for such a world. Relative to each moment some matters are settled and some open, and we prove that settledness is monotone along every history: what is settled at a moment remains settled, with the same value, at every later moment (Theorem~\ref{t:monotone}). Settledness is a relation between a moment and a proposition. No moment has it absolutely, and so the open future requires no privileged present and no growing block (Theorem~\ref{t:noblock}).

The structural results lend themselves to exhaustive verification. In the appendices we enumerate every finite branching frame with at most eleven moments, one per isomorphism class, \NFrames{} in all, and verify for each frame and each of its valuations the monotonicity of settledness, the bivalence of truth relative to a history together with the failure of bivalence for settledness, and the McTaggart-style facts about A-assignments. The tables record, for each of the \NFramesDetail{} frames with at most ten moments, the computed values at every moment. Section~\ref{s:computation} discusses particular entries.

The plan is as follows. Section~\ref{s:history} reviews the argument and its reception. Section~\ref{s:framework} sets out the framework. Section~\ref{s:incoherence} proves the incoherence result, Section~\ref{s:inversion} the inversion, and Section~\ref{s:branching} develops branching and settledness. Section~\ref{s:spacetime} defends the identification of the B-series with the causal and metrical structure of spacetime, and Section~\ref{s:block} excludes the growing block and the other A-theories. Section~\ref{s:objections} answers objections, Section~\ref{s:computation} reports the computation, Section~\ref{s:limits} records what the result does not establish, and Section~\ref{s:formal} describes the formalization.

\section{The argument and its reception}\label{s:history}

\subsection{Before McTaggart}

The question whether time is real is older than McTaggart's argument, and two strands of the earlier discussion bear directly on it. The first concerns the relation of time to change. Aristotle held that time is not change but is not without change: we perceive time when we perceive change, and time is the number of change with respect to before and after (\emph{Physics} IV.11). He also raised the puzzles about the now that every later A-theory inherits: if the now is always different, it seems that a now must perish, and it cannot perish in itself nor in another now (\emph{Physics} IV.10). The Aristotelian claim that time requires change is the ancestor of McTaggart's first premise, and Aristotle's own gloss shows what it concerns: change in things, the motion that time numbers, which is what we call B-change.

The second strand concerns the future. In \emph{De Interpretatione} 9 Aristotle considered whether it is now true or false that there will be a sea battle tomorrow, and argued that if every statement about the future were already true or false, everything would happen of necessity. The passage is the source of the problem of future contingents and of the branching treatment of it, on which the future is open at a moment because the moment has several possible continuations. Augustine's reflections in \emph{Confessions} XI add the epistemic and psychological side: the past is no longer and the future not yet, and the present, if it had extension, would divide into past and future; Augustine concluded that time is in some way an extension of the mind, which measures in memory and expectation what it cannot measure in things (\emph{Confessions} XI.14--28). Augustine's puzzles about the present are the ancestors of the presentist's view that only the present exists and of the A-theorist's sense that the present is special; his solution, placing the measure of time in the mind, is an early statement of the view that the experience of passage is a feature of experience (TSM-02).

In the modern period the reality of time was argued over in the dispute between Newton's absolute time, which flows equably without relation to anything external (\emph{Principia}, Scholium to the Definitions), and Leibniz's relational time, an order of successions (Correspondence with Clarke, Third Paper, §4). Kant held that time is the form of inner sense and is transcendentally ideal (\emph{Critique of Pure Reason}, Transcendental Aesthetic, §§4--7), which McTaggart cites as one of the predecessors of his conclusion. Our view takes from Leibniz the claim that time is an order, and from relativity, rather than from Newton, the claim that the order has a physical measure; it takes from neither the idea of a flow.

\subsection{McTaggart, 1908 and 1927}

McTaggart's paper of 1908 opens by remarking that the unreality of time is a doctrine that has attracted philosophers of very different schools, and he cites Spinoza, Kant, Hegel and Schopenhauer as well as the Buddhist tradition, before offering an argument of his own. The argument was restated, with minor changes, as chapter 33 of the second volume of \emph{The Nature of Existence}, published posthumously \citep{McTaggart1927}. We follow the later version where the two differ, and we use McTaggart's terminology throughout.

Positions in time, McTaggart says, appear to us in two ways. Each position is earlier than some and later than others; the series of positions so ordered is the B-series. Each position is also past, present or future; the series of positions running from the far past through the present to the far future is the A-series. The distinctions of the B-series are permanent: if an event is ever earlier than another, it is always earlier. The distinctions of the A-series are not: an event that is future becomes present and then past. McTaggart also introduces a third series, the C-series, an order of positions that is not temporal at all but which, combined with the A-series, gives rise to the B-series. The C-series has an order but no direction; the A-series supplies the direction and the change \citep[ch.~33]{McTaggart1927}.

The first stage of the argument claims that the A-series is essential to time. McTaggart's ground is that time involves change, and that the only change there can be is change in the A-determinations of events. An event is a content at a position in the B-series; its B-relations to other events cannot change, and its non-temporal characteristics cannot change either, since a change in them would make it a different event. A poker that is hot on Monday and cold on Tuesday does not, he says, change in the required sense, since it is always true that it is hot on Monday and cold on Tuesday. What does change is that the event of the poker's being hot, once future, becomes present and then past. If there were no A-series, then, there would be no change, and so no time. The B-series, which presupposes time, would not exist either.

The second stage claims that the A-series is contradictory. Past, present and future are incompatible determinations: no event can have more than one. Yet every event has all three, since every event is past, present and future. The natural reply is that no event has them all at once: an event is future, will be present, and will be past, or it is past, was present and was future. McTaggart argues that this reply is circular or regressive. To say that an event \emph{is} past, \emph{was} present and \emph{will be} future is to say that it is past at a present moment, present at a past moment and future at a future moment. But the moments at which it has these determinations themselves have all three A-determinations, and the same contradiction recurs at the second level, and so on. The regress is vicious, McTaggart holds, because each stage is needed to remove the contradiction of the previous one, and none succeeds without the next.

The conclusion follows. Time requires the A-series; the A-series is contradictory and so not real; therefore nothing is in time. McTaggart did not deny that there is some series in reality that we misperceive as temporal: the C-series, in his metaphysics, is the order of a timeless reality of perceptions, which appears to us as a series of events in time \citep{McTaggart1927}. He denied that this series is a temporal one.

\subsection{McTaggart's positive view}

It is easy to forget that McTaggart's argument was part of a system. In the second volume of \emph{The Nature of Existence} the unreality of time is one of several results, alongside the unreality of matter and of sense data, on the way to the conclusion that reality consists of selves and their perceptions, related by love \citep{McTaggart1927}. The C-series has a place in this system. It is the real, non-temporal order of a self's perceptions, ordered by their degree of inclusiveness, and the appearance of time is the misperception of this order as a series of events in which some are past and some future. The series runs, McTaggart held, from a least inclusive term to a most inclusive one, and the appearance of temporal passage is the appearance of a self moving along it.

The system explains why McTaggart did not regard the B-series as a fallback. On his view the B-series is constituted by the A-series together with the C-series: an event is earlier than another just if it is past when the other is present, or present when the other is future. Take away the A-series and what is left is the C-series, which is real but not temporal. Our disagreement with McTaggart is located precisely here. We hold that the B-relation is not defined from the A-series and the C-series but is a physical relation in its own right, the causal order of events, and that it has a direction of its own (Section~\ref{s:spacetime}). McTaggart's reduction of the B-series to the A-series is the bridge in another form, and its failure is the subject of Section~\ref{s:inversion}.

\subsection{Russell and Broad}

The argument met an immediate response from the view it was designed to refute. \citet[§442]{Russell1903} had already defined change as the difference, in respect of truth or falsity, between a proposition concerning an entity and a time and a proposition concerning the same entity and another time. On this definition the poker changes: the proposition that it is hot at Monday is true, and the proposition that it is hot at Tuesday is false. McTaggart's reply, that these propositions are both always true or always false and so do not change, assumes that change must be the change of a proposition's truth-value over time, which on Russell's view is not what change is. \citet{Russell1915} added that the A-determinations are relations to a subject: \enquote{past} means earlier than the time of the present experience, so that the A-series is the B-series viewed from a position within it. This is the ancestor of the indexical treatment of tense.

C.~D. Broad responded twice. In \emph{Scientific Thought} he proposed the view now called the growing block: past and present events exist and future events do not, and the sum total of existence increases as events come into being \citep{Broad1923}. The present is the edge of the block, and becoming is the addition of new slices. In his examination of McTaggart's philosophy he gave a criticism of the regress that has been influential among A-theorists \citep{Broad1938}. He argued that the regress arises only if one reads \enquote{was}, \enquote{is} and \enquote{will be} as a timeless copula followed by a temporal predicate. Read as irreducibly tensed copulas, the statements that an event was future, is present and will be past are consistent and do not need to be paraphrased in terms of moments. Broad thus located McTaggart's error in his translation of tensed statements into tenseless ones. We return to this reply in Section~\ref{s:incoherence}, where we argue that it secures consistency only by giving up the change it was meant to protect.

\subsection{Dummett}

Michael \citet{Dummett1960} argued that McTaggart's argument is better than its critics allow. The key move is to read it as an argument about descriptions. The A-series is essential to time because a description of a temporal reality from no temporal standpoint omits something: it cannot say which events are past and which future, and it cannot be supplemented to say this without adopting a standpoint. McTaggart's assumption, on this reading, is that there must be a complete description of reality, one that is not given from any standpoint. Since a temporal reality has no such description, it is not real. Dummett did not endorse the conclusion. He held that the argument shows that we must give up either the reality of time or the idea that reality has a complete, standpoint-free description, and he inclined to the latter.

Dummett's reading has proved fertile. It locates the issue in the relation between perspective and reality and is the ancestor of Fine's reconstruction. We think it also shows where the argument goes wrong. The ineliminability of tensed thought is that of indexical thought in general, and the same considerations show that no description from no standpoint says which place is here or which person is me. That thought has a perspective does not show that reality has one (Objection~\ref{o:dummett}).

\subsection{Prior and the logic of tense}

A.~N. Prior made tense the subject of a logic \citep{Prior1957,Prior1967}. Tense operators $P$ (\enquote{it was the case that}) and $F$ (\enquote{it will be the case that}) are prefixed to present-tensed sentences, and the logic of their combinations is studied in the same way as the logic of necessity and possibility. Prior was an A-theorist and in his later work a presentist: he held that tensed sentences are the fundamental bearers of truth, and that the truths about past and future are tensed truths about the present \citep{Prior1968}. Against the B-theory he offered his famous argument from relief: when a painful experience is over, one says \enquote{thank goodness that's over}, and what one is thankful for is not the tenseless fact that the experience ends before the utterance, which obtains at every time, but the tensed fact that it is now over \citep{Prior1959}.

Prior's logic is neutral between metaphysical views, and the B-theorist can use it. Two of his innovations are central to what follows. The first is his semantics for an indeterministic world, in which the moments of time form a tree branching toward the future and future-tensed sentences are evaluated relative to histories through a moment \citep{Prior1967}. Prior distinguished a Peircean reading, on which \enquote{it will be that $p$} means that $p$ on every continuation, from an Ockhamist reading, on which it is evaluated relative to a history. \citet{Thomason1970} combined the Ockhamist reading with supervaluation to give a semantics on which future contingents are neither true nor false at a moment, and the framework was developed by \citet{Belnap2001} into a general theory of agency in indeterministic time. The second innovation is the observation that the semantics of a tensed language may be stated in a tenseless metalanguage: the truth conditions of \enquote{it will be that $p$} at a moment are given by the existence of a later moment at which $p$. We adopt both.

\subsection{Smart, Mellor and the new B-theory}

The B-theory received its classical modern statement in \citet{Smart1963} and \citet{Mellor1981}. Smart held that tensed expressions are token-reflexive or indexical, so that \enquote{it is raining now} uttered at $t$ is true just if it rains at $t$, and that the A-series has no place in the scientific image. Mellor's \emph{Real Time} and its successor \citep{Mellor1998} gave a version of McTaggart's argument that the B-theorist can accept. Mellor's version concerns tensed sentence tokens and their truth conditions. Every A-determination is had by an event only relative to a time, since an event is past at some times and future at others. If we try to make A-facts absolute by saying that an event is past at some time and future at another, we must specify those times by their A-determinations or by their B-positions. In the first case the problem recurs; in the second, the A-facts are B-facts. Mellor concludes that there are no A-facts and that tensed sentences have tenseless truth conditions: a token of \enquote{$e$ is past} is true just if it occurs later than $e$.

Mellor's reconstruction drew a sharp debate. \citet{Lowe1987} charged it with an \enquote{indexical fallacy}: McTaggart's argument, Lowe held, depends on treating tensed sentences as having tenseless truth conditions, and so begs the question against the A-theorist. \citet{LePoidevinMellor1987} replied that the argument requires only that tensed tokens have truth conditions that a tenseless theory can state, and that the A-theorist who denies this owes an account of what makes a tensed token true. We agree with the reply and offer a version of the argument in Section~\ref{s:incoherence} that does not require the premise Lowe objected to. Its central claim is that an absolute present, a monadic property that one moment has simpliciter, cannot change without being had by more than one moment, and that the only way to let it change coherently is to relativize it to moments.

The new B-theory added an account of tensed belief. \citet{Perry1979} and \citet{Lewis1979} showed that the content of indexical belief cannot be captured by a proposition true or false absolutely: the belief that I am making a mess, or that it is now noon, has self-locating content. Mellor and others concluded that the irreducibility of tensed belief, which Prior's relief argument exploits, is an instance of the irreducibility of indexical belief in general and is no evidence of tensed facts \citep{Mellor1998}. The same conclusion governs the treatment of \enquote{now} in EPI-18-c, where Beauty's uncertainty whether it is now Monday is uncertainty about her location in a world whose history she knows, and in REL-19, where God's timelessness entails a lack of indexical knowledge of what time it is.

\subsection{Fine's reconstruction}

Kit \citet{Fine2005} reconstructed McTaggart's argument as a conflict between four assumptions. \emph{Realism}: reality is constituted, at least in part, by tensed facts. \emph{Neutrality}: no time is privileged; the tensed facts that constitute reality are not oriented toward one time as opposed to another. \emph{Absolutism}: the constitution of reality is an absolute matter, not relative to a time or other standpoint. \emph{Coherence}: reality is not contradictory; it is not constituted by facts with incompatible content. Fine argued that the four are inconsistent, since by Realism and Neutrality reality contains tensed facts oriented toward each time, and these are incompatible, so that by Absolutism reality contains incompatible facts, contrary to Coherence. The standard A-theorist rejects Neutrality: one time, the present, is privileged. The B-theorist rejects Realism. Fine himself considered two non-standard options: \emph{external relativism}, which gives up Absolutism and holds that reality is constituted relative to a time, and \emph{fragmentalism}, which gives up Coherence for reality as a whole while holding that each fragment, the facts oriented toward a single time, is coherent.

Fine's reconstruction is valuable because it shows how much the A-theorist must give up. The standard A-theorist must privilege one time, and so must explain why reality favors this time over every other, and how the favor can move. The non-standard A-theorist preserves Neutrality only by giving up the absoluteness or the coherence of reality. Our view rejects Realism, and in Section~\ref{s:incoherence} we show that the relativist's \enquote{reality relative to a time} is, so far as its temporal content goes, the B-series under another description (Proposition~\ref{p:stage}). Fragmentalism we discuss in Objection~\ref{o:fine}.

\subsection{Recent A-theories}

The last two decades have seen a revival of the A-theory in several forms. Presentism, the view that only present things exist, has been defended against the objection from relativity and the objection that it lacks truthmakers for truths about the past \citep{Markosian2004}. The growing block has found defenders who argue that it captures the asymmetry between the fixed past and the open future \citep{Tooley1997,CorreiaRosenkranz2018}, and it has been attacked on the ground that its inhabitants could not know whether they were at the edge of the block \citep{Bourne2002,BraddonMitchell2004}; \citet{Forrest2004} replied that only the present is conscious, a reply whose costs we discuss in Section~\ref{s:block}. The moving spotlight, on which all times exist but one is objectively present and the present moves, has been given a careful defense by \citet{Cameron2015}, who treats the past and future as existing but having their temporal distance from the present as a fundamental property, and by \citet{Skow2015}, who examines what objective becoming would have to be and argues that the best version of the A-theory is a moving spotlight with a privileged present. \citet{Zimmerman2005} argued that the debate between A- and B-theorists is substantive, against those who would dissolve it, and that \enquote{taking tense seriously} commits one to a distinction in what exists or in what is fundamental.

These views share a feature that is decisive for our argument: each privileges one moment, the present or the edge, absolutely. In Section~\ref{s:block} we show that this feature is what McTaggart's argument refutes, and that a view on which the openness of the future and the fixity of the past are moment-relative obtains every asymmetry the A-theorist wants without it.

\subsection{Where we stand}

Our position combines elements of the tradition. With McTaggart and Mellor, we hold that the A-series is incoherent. With Russell, Smart and Mellor, we hold that change is variation across times and that tensed thought is indexical. With Prior, Thomason and Belnap, we hold that the future may be open, and we model openness by branching. Against McTaggart, we hold that time is real, and against Broad and Prior, that the open future does not require an A-series. The novelty of the treatment, if it has one, lies in three things: in presenting the argument as an inversion rather than a selective rejection of premises, in following TSM-01-b; in proving that the coherent relativization of the A-series is the B-series (Proposition~\ref{p:stage}); and in showing, by proof and by exhaustive computation, that moment-relative settledness has the formal properties the open future requires and none of those the growing block attributes to it.

\section{Framework}\label{s:framework}

We begin with the two series and the two notions of change.

\begin{definition}[B-series]\label{d:bseries}
The \emph{earlier-than} relation is a strict partial order on events (or on the moments at which they occur): irreflexive and transitive. The \emph{B-series} is the set of events so ordered, together with the metrical relations of temporal separation between them. A \emph{B-fact} is a fact about what occurs at which moments and about the relations of earlier, later, simultaneous and temporal separation among them.
\end{definition}

We do not assume that the order is linear. In a relativistic spacetime it is the causal order, under which spacelike separated events are incomparable, and in an indeterministic world it may branch toward the later (Section~\ref{s:branching}).

\begin{definition}[A-determinations]\label{d:adet}
The \emph{A-determinations} are past, present and future, taken as the A-theorist takes them: monadic properties that events or moments have absolutely, not relative to a time or to a thinker. The \emph{A-series} is the series of positions ordered by these determinations. In particular, \emph{absolute presentness} is the property a moment has if it is present simpliciter.
\end{definition}

\begin{definition}[B-change and A-change]\label{d:change}
An object undergoes \emph{B-change} if it has a property at one moment and an incompatible property at a later moment. An event undergoes \emph{A-change} if it passes from future to present to past, that is, if its A-determination changes.
\end{definition}

McTaggart's premises, in the form in which TSM-01-b states them, are these.

\begin{assumption}[Time requires change]\label{a:m1}
There is no time without change.
\end{assumption}

\begin{assumption}[Change requires the A-series]\label{a:m2}
The only change there could be is change in A-determinations: an event's ceasing to be future, becoming present, and becoming past.
\end{assumption}

\begin{assumption}[Incoherence]\label{a:incoherence}
The A-series is incoherent: every event would have to be past, present and future, which are incompatible, and the attempt to remove the contradiction by tensing the copula generates a regress in which it recurs. In its formal core: the A-determinations are incompatible, so that absolute presentness is had by at most one moment of a world.
\end{assumption}

The formal core requires comment. If two distinct moments were both present simpliciter, then, since distinct moments of one history are related as earlier and later, one would be earlier than the other, and the earlier would be both present and, relative to the later, past. Absolute presentness would then be had by a moment that also has an incompatible determination. Where the moments are incomparable, as on different branches or at spacelike separation, the A-theorist who holds that both are present simpliciter has given up the uniqueness of the present that distinguishes the A-series from a mere labeling of moments; we return to this in Section~\ref{s:spacetime}. The formal core is therefore what remains of McTaggart's premise once the contradiction is stated without his regress; the regress itself is the subject of Corollary~\ref{c:dilemma}.

We accept Assumption~\ref{a:incoherence}. We reject the conjunction of Assumptions~\ref{a:m1} and~\ref{a:m2}. Our positive view is stated by two principles.

\begin{principle}[Constitution]\label{p:constitution}
Time is constituted by the B-series. For there to be time is for there to be events ordered as earlier and later; temporal facts are B-facts.
\end{principle}

\begin{principle}[Physical identification]\label{p:physical}
In our world the B-series is the causal and metrical structure of spacetime: the earlier-than relation is the causal order of events, and temporal separation is the proper time along timelike curves.
\end{principle}

Principle~\ref{p:constitution} is the thesis that time is the B-series. Principle~\ref{p:physical} says what the B-series is in our world, and it is the identification that TSM-14 uses when it asks whether time had a beginning. We defend it in Section~\ref{s:spacetime}. Two consequences may be noted now. First, Principle~\ref{p:constitution} does not say that time requires change. Whether a stretch of time could pass with no change in anything is a separate question, which TSM-17 answers affirmatively with Shoemaker's argument: the B-series is an order of times, and it needs no change to be populated. Second, Principle~\ref{p:constitution} makes tensed thought indexical. \enquote{Now} refers to the time of the thought or utterance (TSM-18), the present is the time at which a thought or utterance occurs (TSM-15), and to know what time it is is to have self-locating information about one's position in the B-series (EPI-18-c).

\section{The incoherence of the A-series}\label{s:incoherence}

McTaggart's contradiction is often presented as a puzzle about tensed sentences. We present it as a result about assignments of A-determinations to moments. This has two advantages. It makes clear exactly which feature of the A-series generates the contradiction, namely its absoluteness, and it shows that the only coherent repair, relativization to moments, yields the B-series.

Let a world have a set of moments ordered by earlier-than. An A-determination of a moment $x$ relative to a moment $m$ is one of four values: \emph{past} if $x$ is earlier than $m$, \emph{present} if $x = m$, \emph{future} if $x$ is later than $m$, and \emph{other} if $x$ and $m$ are incomparable (which can happen only if the order branches or is relativistic). We write $P$, $N$, $F$, $O$ for these values.

\begin{definition}[Monadic A-series]\label{d:monadic}
A \emph{monadic A-series} for a world is an assignment of A-determinations to its moments that the world has absolutely: not relative to a time, a thinker or a standpoint. It is \emph{coherent} if it gives each moment exactly one A-determination and its determinations agree with the order: if $x$ is present, every moment earlier than $x$ is past and every moment later than $x$ is future.
\end{definition}

\begin{lemma}\label{l:onepresent}
By Assumption~\ref{a:incoherence}, a monadic A-series has at most one present moment. A coherent monadic A-series with a present moment $m$ gives every moment the determination it has relative to $m$.
\end{lemma}

\begin{proof}
The first claim is the formal core of Assumption~\ref{a:incoherence}. For the second, let $m$ be the present moment. By coherence, the moments earlier than $m$ are past and the moments later than $m$ are future. No other moment is present, by the first claim, and a moment incomparable with $m$ is neither earlier nor later than the present and so is assigned the remaining value. The assignment is therefore the one relative to $m$.
\end{proof}

\begin{definition}[Changing A-series]\label{d:changing}
A world contains a \emph{changing A-series} if the present is at more than one of its moments: there are distinct moments $x$ and $y$ each of which is present.
\end{definition}

Definition~\ref{d:changing} is McTaggart's requirement that every event be present: an A-series that changes is one in which the present is not fixed at one moment but is had, in turn, by many. An A-series whose present is at only one moment does not change. It is a single labeling of the moments, and the events it labels are each past, present or future for good. Such a series is no rival to the B-series as the source of change; it is a further, static feature of the world that contributes no change at all.

\begin{theorem}[No changing A-series]\label{t:nochange}
No world contains a changing A-series.
\end{theorem}

\begin{proof}
Suppose a world contains a changing A-series, so that distinct moments $x$ and $y$ are each present. By Lemma~\ref{l:onepresent}, a monadic A-series has at most one present moment, so $x = y$, which contradicts their distinctness.
\end{proof}

The proof is short because the work has been done by isolating the formal core of McTaggart's premise. The A-theorist will reply that the present is at different moments \emph{at different times}: $x$ is present at $x$, and $y$ is present at $y$. This is the move to the tensed copula, and the next results examine where it leads.

\begin{definition}[Moment-relative A-assignment]\label{d:relative}
For each moment $m$, the \emph{A-assignment relative to $m$}, $A_m$, gives each moment $x$ the value $P$ if $x$ is earlier than $m$, $N$ if $x = m$, $F$ if $x$ is later than $m$, and $O$ otherwise.
\end{definition}

The family $\{A_m\}$ is coherent and changing in the only sense the A-theorist's reply can give these words: each $A_m$ is coherent, and every moment is present according to one of them, namely its own. The question is what this family is.

\begin{proposition}[The relativized A-series is the B-series]\label{p:stage}
Let $\{A_m\}$ be the family of moment-relative A-assignments of a frame.
\begin{enumerate}[label=(\alph*)]
\item $A_m(x) = P$ if and only if $x$ is earlier than $m$. Hence the earlier-than order can be read off the family, and the family can be read off the order.
\item $A_m(x) = F$ if and only if $A_x(m) = P$.
\item The family commutes with every automorphism $\sigma$ of the order: $A_{\sigma(m)}(\sigma(x)) = A_m(x)$.
\item No single member of the family gives every moment the value $N$ unless the frame has exactly one moment.
\end{enumerate}
\end{proposition}

\begin{proof}
(a) If $x$ is earlier than $m$, then $x \ne m$ by irreflexivity, so $A_m(x) = P$. Conversely, if $A_m(x) = P$ then by definition $x \ne m$ and $x$ is earlier than $m$. (b) $A_m(x) = F$ just if $x \ne m$ and $m$ is earlier than $x$, which by (a) holds just if $A_x(m) = P$. (c) Automorphisms preserve identity and the order, and $A_m(x)$ is defined from these alone. (d) $A_m$ gives $N$ only to $m$.
\end{proof}

The family $\{A_m\}$ contains no more information than the order, by (a), and treats moments only by their position in it, by (c). It is the B-series described in the vocabulary of tense. Clause (d) is the computational shadow of Theorem~\ref{t:nochange}: the family makes every moment present only by having many members, one for each moment, none of which is the A-series of the world.

An example makes the point concrete. In the three-moment chain T3.1, with moments $0$, $1$, $2$ in that order, and in the fork T3.2, in which $0$ precedes the incomparable moments $1$ and $2$, the family $\{A_m\}$ is given by the following matrices, whose row $m$ lists $A_m(0)$, $A_m(1)$, $A_m(2)$.
\begin{center}
\begin{tabular}{c@{\qquad}ccc@{\qquad\qquad}c@{\qquad}ccc}
\toprule
\multicolumn{4}{c}{T3.1 (chain)} & \multicolumn{4}{c}{T3.2 (fork)}\\
$m$ & $x=0$ & $x=1$ & $x=2$ & $m$ & $x=0$ & $x=1$ & $x=2$\\
\midrule
0 & $N$ & $F$ & $F$ & 0 & $N$ & $F$ & $F$\\
1 & $P$ & $N$ & $F$ & 1 & $P$ & $N$ & $O$\\
2 & $P$ & $P$ & $N$ & 2 & $P$ & $O$ & $N$\\
\bottomrule
\end{tabular}
\end{center}
Each matrix is the order written out: the entries $P$ below the diagonal of the chain are its pairs $x < m$, the entries $F$ above it the converse pairs, and in the fork the entries $O$ mark the incomparable pair. Each row is coherent, every moment is $N$ in exactly one row, and no row is $N$ throughout, as clause (d) requires. In the fork the automorphism exchanging $1$ and $2$ exchanges rows $1$ and $2$ and columns $1$ and $2$ together and leaves the matrix unchanged, as clause (c) requires. The A-theorist who wants a changing A-series wants the matrix to be a single row that changes. What the B-series provides is the whole matrix, which does not change and does not need to.


\begin{corollary}[The dilemma of the tensed copula]\label{c:dilemma}
An A-theorist who answers McTaggart by tensing the copula must choose. If \enquote{$x$ is present at $m$} means that $x$ is present relative to $m$, the A-series so described is the family $\{A_m\}$, which is the B-series under another description (Proposition~\ref{p:stage}). If it means that, at the stage at which $m$ is absolutely present, $x$ is absolutely present, then the stages are ordered, and either their order is the B-order of moments, and the first horn applies, or the stages are a second series whose own A-determinations generate the problem again.
\end{corollary}

\begin{proof}
The first horn is Proposition~\ref{p:stage}. For the second, the stages at which the various moments are absolutely present are distinct, since a moment is absolutely present at one stage and past at another, and they are ordered, since one moment's presentness precedes another's. If the stages are ordered as their moments are, the stage order is the earlier-than order, and the absolute presentness of a moment \enquote{at a stage} is its presentness relative to a moment, which is the first horn. If they are not, the stages form a series of their own; for it to be a series in which the present changes, a stage must be present absolutely, and by Theorem~\ref{t:nochange} applied to the series of stages, at most one is. This is McTaggart's regress.
\end{proof}

Corollary~\ref{c:dilemma} is our version of McTaggart's regress. It answers Broad's reply \citep{Broad1938}. Broad held that the irreducibly tensed copula renders the A-theorist's claims consistent, and so it does: \enquote{$e$ was future, is present and will be past} is consistent if \enquote{was}, \enquote{is} and \enquote{will be} are read as tense operators evaluated at a moment. But evaluation at a moment is what the first horn describes. The consistent reading of tensed claims is the reading on which their truth at a moment is fixed by the moment's position in the B-series, and this reading supplies no absolute present that changes. The A-theorist who wants more must say what it is for the present to move other than for different moments to be present relative to themselves, and every answer either repeats the B-series or starts a regress.

Lowe's charge of an indexical fallacy \citep{Lowe1987} is not engaged by this argument, since we have not assumed that tensed sentences have tenseless truth conditions. We have assumed only the formal core of incompatibility, which the A-theorist accepts: the present is one. The argument shows that a unique absolute present cannot change, and that a changing present is not absolute.

\subsection{The regress and its readers}

McTaggart's regress has been read in three ways, and the dilemma of Corollary~\ref{c:dilemma} shows what is right in each. On the first reading, the regress is a regress of times: to remove the contradiction we say that an event is present at one moment and past at another, the moments themselves must be past, present and future, and so on. Broad's reply was directed at this reading, and it is correct as far as it goes: if \enquote{at a moment} is read as a tense operator, no further moments are introduced, and nothing regresses \citep{Broad1938}. On the second reading, which \citet{Dummett1960} favored, the regress is a regress of descriptions: each attempt to describe the A-series consistently is given from a standpoint, and to describe the standpoints we must adopt another. On the third reading, Mellor's, the regress is a regress of truth conditions: a tensed token is true or false at a time, and if the A-facts are to be absolute, the truth of a token must not depend on when it is assessed, which can be secured only by giving tokens tenseless truth conditions \citep{Mellor1981,Mellor1998}.

The dilemma accommodates all three. On the first horn, the tensed copula is evaluated at moments, and Broad is right that no regress begins; but what is consistently described is the family $\{A_m\}$, whose temporal content is the B-order (Proposition~\ref{p:stage}). Mellor's conclusion, that the truth conditions of tensed tokens are tenseless, is then a consequence rather than a premise, which is why the argument is not open to Lowe's charge. On the second horn, stages are introduced over and above moments, and Dummett's regress of standpoints is the regress of series of stages, each requiring a present of its own. The three readers disagree about where McTaggart's argument has its force; they agree, and Corollary~\ref{c:dilemma} confirms, that the only consistent A-series is one whose determinations are relative to moments, and that such a series does not change in the A-theorist's sense.

\subsection{The present as a further fact}

One more feature of the moment-relative family deserves statement, because it bears on the A-theorist's claim that the present is a further fact over and above the B-facts.

\begin{proposition}[The B-facts do not select a present]\label{p:select}
Let a frame have more than one moment.
\begin{enumerate}[label=(\alph*)]
\item The B-facts of the frame are the same whichever moment is taken to be present; hence they do not determine which moment is present.
\item A moment that is singled out by the B-facts alone is fixed by every automorphism of the frame. In a finite branching frame the least moment is always so fixed.
\item If the frame has a non-trivial automorphism, there are distinct moments $x$ and $y$ such that the world with $x$ present and the world with $y$ present differ in no B-fact and in no respect of the position of the present in the B-series.
\end{enumerate}
\end{proposition}

\begin{proof}
(a) The B-facts are the facts of occurrence and order, and Definition~\ref{d:adet} makes absolute presentness a further, monadic property; replacing one present moment by another leaves the order and the occurrences untouched. (b) If a moment is the unique moment satisfying a condition stated in terms of the order, an automorphism, which preserves the order, maps it to a moment satisfying the same condition, hence to itself. The least moment of a finite branching frame is the unique moment earlier than every other. (c) A non-trivial automorphism $\sigma$ moves some moment $x$ to $y = \sigma(x) \ne x$. The pointed frames $(F, x)$ and $(F, y)$ are isomorphic by $\sigma$, so every condition on the position of the present in the order that holds of one holds of the other.
\end{proof}

Proposition~\ref{p:select} shows that a monadic present is, by the A-theorist's own lights, a fact that the B-facts leave open, and in frames with symmetry it distinguishes worlds that agree in everything else, including in where the present stands relative to every event. The computation in Appendix~\ref{app:detail} records, for each frame, the orbits of its automorphism group: in each of the frames that are not rigid, which number $\NFrames - \NRigid$ of the \NFrames{} frames, some moments share an orbit, and the A-theorist's present would there distinguish worlds that differ only in which of two indiscernible moments has it. In the \NRigid{} rigid frames every moment is singled out by the B-facts, so that the B-facts single out too many moments rather than none, and they still do not say which one is present.

\section{The inversion}\label{s:inversion}

The results of the previous section concern the A-series. McTaggart's conclusion concerns time. To reach it he needs a premise linking the two, and his first two premises supply it.

\begin{assumption}[Reality of time]\label{a:real}
Time is real in the actual world: there are events ordered as earlier and later, at temporal separations.
\end{assumption}

Our evidence for Assumption~\ref{a:real} is the evidence for the temporal order and duration of events: clocks agree with each other and with the laws that predict their readings; records and memories are ordered consistently with the causal structure of the processes that produce them; physical theory, from mechanics to general relativity, describes a world whose events stand in a causal and metrical order and is confirmed in that description to extraordinary precision. Nothing in this evidence concerns an absolute present. It is evidence for the B-series.

\begin{definition}[The bridge]\label{d:bridge}
The \emph{bridge} is the claim that there is no time without a changing A-series. It is the conjunction of Assumptions~\ref{a:m1} and~\ref{a:m2}.
\end{definition}

TSM-01-b shows that the argument needs exactly the bridge and that the conjunction of Assumptions~\ref{a:m1} and~\ref{a:m2} is equivalent to it. McTaggart's argument has the form: bridge, incoherence, so no time. Its inversion has the form: time, incoherence, so not the bridge.

\begin{theorem}[The inversion]\label{t:inversion}
If time is real in some world, the bridge is false. Given Assumption~\ref{a:real}, the bridge is false.
\end{theorem}

\begin{proof}
Suppose time is real in $w$ and the bridge holds. Then $w$ contains a changing A-series, which contradicts Theorem~\ref{t:nochange}. Hence if time is real in some world, the bridge is false, and by Assumption~\ref{a:real} time is real in the actual world.
\end{proof}

The theorem does not by itself say which direction is rational. A valid argument with premises $A$ and $B$ and conclusion $C$ supports $\neg C$ from $A$ and $B$ no more and no less than it supports $\neg B$ from $A$ and $C$. The choice depends on the comparative credibility of the bridge and of the reality of time. The reality of time has the support just described. The bridge, we now show, has none that survives disambiguation.

\begin{proposition}[The bridge equivocates]\label{p:equivocation}
Assumption~\ref{a:m1} is credible only when \enquote{change} means B-change, and Assumption~\ref{a:m2} only when it means A-change. Hence the bridge has no reading on which both of its conjuncts are credible.
\end{proposition}

\begin{proof}
The considerations that recommend Assumption~\ref{a:m1} are the Aristotelian thought that time is the number of motion with respect to before and after (\emph{Physics} IV.11), the verificationist thought that a time at which nothing differs from any other could not be detected, and the relationist thought that times are constructed from the states of things. Each concerns variation in things across times, which is B-change; none concerns events' passing from future to past. The consideration that recommends Assumption~\ref{a:m2} is McTaggart's argument from the poker: the poker's being hot on Monday and cold on Tuesday is a permanent fact, so the poker does not change. This presupposes that change must be a change in the truth-value of a fixed proposition, which is the definition of A-change. Read as a claim about B-change, Assumption~\ref{a:m2} is false, since the poker is hot at one time and cold at a later one, which is all that B-change is \citep[§442]{Russell1903}. Read as a claim about A-change, Assumption~\ref{a:m1} has no independent support.
\end{proof}

Proposition~\ref{p:equivocation} reproduces the diagnosis of TSM-01-b, which we accept in full: the credibility of the bridge derives from an equivocation between B-change and A-change, and the reality of time does not. Geach's term for the poker's variation, a \enquote{Cambridge change}, was intended to mark it as a degenerate case \citep{Geach1969}; what the term marks is only that the variation is not A-change. Russell's definition is not a stipulation that changes the subject. It describes the change we observe and measure: the poker cooling, the planet moving, the cell dividing. An argument that such things do not change unless their events also change their A-determinations rests on a premise that has no support of its own.

\begin{corollary}[What survives]\label{c:survives}
Time is an order of events by earlier and later. McTaggart's premise that time requires change, insofar as it supports his conclusion, concerns the change of A-determinations, and in that sense it is false; in the sense in which it is plausible, it concerns B-change, and it neither requires the A-series nor entails that time is unreal.
\end{corollary}

\begin{proof}
By Theorem~\ref{t:inversion} and Assumption~\ref{a:real}, time is real and does not require a changing A-series; by Theorem~\ref{t:nochange} there is none. What there is, in every world with time, is an order of events, which is the B-series of Definition~\ref{d:bseries}. The second sentence restates Proposition~\ref{p:equivocation}.
\end{proof}

Whether time requires B-change is a further question. We need not answer it here, since nothing in the argument depends on it, and TSM-17 argues that it does not: time without change is possible and could be known to occur. On the B-theory a total freeze is a stretch of the B-series at whose times the world's states are qualitatively alike, and that is a possibility the B-series leaves open.

\subsection{The evidence for the reality of time}

The inversion is rational only if the reality of time is more credible than the bridge. Proposition~\ref{p:equivocation} shows that the bridge has no credibility of its own; it remains to say why the reality of time has a great deal. Three kinds of evidence converge on it.

The first is the agreement of clocks. Physical processes of very different kinds, the swing of a pendulum, the decay of an isotope, the oscillation of a caesium atom, the rotation of the earth, assign durations to the same intervals between the same events, and they agree with one another to within their known errors and with the laws that predict their readings. Where they disagree, as moving and stationary clocks do, the disagreement is predicted by the metrical structure of spacetime. Agreement of this kind among independent measures is the standard evidence for the reality of a magnitude, and the magnitude here is temporal separation, a B-relation.

The second is the order of records. Fossils, sediments, tree rings, light from distant galaxies and memories record events in an order that is consistent across kinds of record and consistent with the causal processes that produce them. A record is formed after what it records, never before; and the orders inferred from different records coincide. This is evidence for a causal order of events, which is the earlier-than relation of Principle~\ref{p:physical}.

The third is the success of theories that presuppose the B-series. Every dynamical theory, from Newtonian mechanics to general relativity and quantum theory, describes how states at earlier times are related to states at later ones, and every one is confirmed in that description. None contains a quantity that measures the motion of a present, and none needs one. The A-theorist may say that physics is silent about the A-series, not that it denies it. That is correct, and it is why the evidence for the reality of time supports the B-series and not the bridge: what the evidence shows to be real is an order and a measure of events, and the bridge says that there is no such order without a changing A-series. The evidence, in other words, is the evidence of physics for the B-series, and the bridge would require that this evidence be evidence for nothing.


\section{Branching and settledness}\label{s:branching}

A B-theory must be stated for an indeterministic world as well as for a deterministic one. If the laws and the past do not fix the future, there are several ways the world could go on from a given moment, and the B-series, which contains the events that occur, has several continuations. The tradition of Prior, Thomason and Belnap models this by letting the order of moments branch toward the later \citep{Prior1967,Thomason1970,Belnap2001}. We adopt the model and show that it gives the open future everything it needs without an A-series.

\begin{definition}[Branching frame]\label{d:frame}
A \emph{branching frame} is a set of moments with a strict partial order, earlier-than, in which the moments earlier than any given moment are linearly ordered, and any two moments have a common earlier-or-equal moment. A \emph{finite} branching frame is a finite tree: it has a least moment, and every other moment has a unique immediate predecessor.
\end{definition}

The frame branches toward the later and not toward the earlier: each moment has one past and possibly several futures. The structure exists tenselessly. Every moment of every branch is as real as every other; the branches are the ways things go on from each moment, and branching is a feature of the B-series, not an addition to it.

\begin{definition}[History]\label{d:history}
A \emph{history} is a maximal chain of moments: a set of pairwise comparable moments to which no further moment can be added without loss of comparability. The histories \emph{through} $m$ are those containing $m$.
\end{definition}

In a finite branching frame the histories correspond one to one to the maximal moments: each history is the chain from the least moment to one maximal moment. We write $h(m)$ for the number of histories through $m$.

\begin{definition}[Tenseless truth relative to a history]\label{d:truth}
A \emph{valuation} assigns to an atomic sentence $p$ the set of moments at which it holds. A tenseless proposition is anything that is true or false relative to each history. For each moment $m$, the proposition $F_m p$, \enquote{at some moment later than $m$, $p$}, is true relative to a history $h$ if $h$ contains a moment later than $m$ at which $p$ holds, and false relative to $h$ otherwise. \enquote{It will be that $p$}, uttered or thought at $m$, expresses $F_m p$.
\end{definition}

Definition~\ref{d:truth} is the Ockhamist clause of \citet{Prior1967}, stated tenselessly: the index $m$ is a moment of the B-series and not a privileged present, and the tensed sentence expresses, at each moment, a proposition indexed to that moment, in the way that \enquote{it is raining here} expresses at each place a proposition about that place.

\begin{definition}[Settledness]\label{d:settled}
A tenseless proposition $\varphi$ is \emph{settled true at} $m$ if it is true relative to every history through $m$, \emph{settled false at} $m$ if it is false relative to every history through $m$, and \emph{open at} $m$ otherwise.
\end{definition}

Settledness is a relation between a moment and a proposition, like \enquote{north of} between a place and a thing. Every moment stands in it to some propositions. The definition agrees with Definition~2.2 of TSM-12-a, where a sentence is settled at a moment if it has the same truth value on every history through the moment, and with Definition~2.1 of REL-05-b; TSM-12-a adds the thesis that a future contingent, a sentence open at $m$, is neither true nor false at $m$, which we accept.

\begin{lemma}[Histories are closed under the past]\label{l:pastclosed}
In a branching frame, every history through a moment $m'$ passes through every moment earlier than $m'$.
\end{lemma}

\begin{proof}
Let $h$ pass through $m'$ and let $m$ be earlier than $m'$. We show that $m$ is comparable with every moment $x$ of $h$, so that, by maximality, $m$ belongs to $h$. Since $h$ is a chain, $x$ is comparable with $m'$. If $m'$ is earlier than $x$ or equal to it, then $m$ is earlier than $x$ by transitivity. If $x$ is earlier than $m'$, then $m$ and $x$ are both earlier than $m'$, and the moments earlier than $m'$ are linearly ordered, so $m$ and $x$ are comparable.
\end{proof}

\begin{theorem}[Monotonicity of settledness]\label{t:monotone}
For every tenseless proposition $\varphi$ and moments $m$ earlier than $m'$: if $\varphi$ is settled true at $m$, it is settled true at $m'$, and if it is settled false at $m$, it is settled false at $m'$.
\end{theorem}

\begin{proof}
By Lemma~\ref{l:pastclosed}, the histories through $m'$ are among the histories through $m$. A proposition true (false) relative to all of the latter is true (false) relative to all of the former.
\end{proof}

Theorem~\ref{t:monotone} is the fixity of the past in its general form, the form REL-05-b uses when it states that every proposition settled at a moment is settled at every later moment. It holds for every proposition, not only for propositions about the past. In particular it holds for $F_m p$: if at $m$ it is settled that $p$ will hold at some later moment, then at every moment later than $m$ it is settled that $p$ holds at some moment later than $m$. The computation verifies the theorem for $F_m p$ under every valuation of every frame with at most eleven moments, \NMonoChecks{} instances in all (Appendix~\ref{app:detail}, column M).

\begin{proposition}[Bivalence and its failure]\label{p:bivalence}
\begin{enumerate}[label=(\alph*)]
\item Truth relative to a history is bivalent: every tenseless proposition is true or false, and not both, relative to each history.
\item Settledness is not bivalent: in a frame where $m$ has more than one history through it, there is a valuation for which $F_m p$ is neither settled true nor settled false at $m$.
\item Conversely, if exactly one history passes through $m$, every tenseless proposition is settled at $m$. Hence $F_m p$ is open at $m$ for some valuation if and only if $h(m) \geq 2$.
\end{enumerate}
\end{proposition}

\begin{proof}
(a) is Definition~\ref{d:truth}. (b) Let $h_1$ and $h_2$ be distinct histories through $m$. They diverge at some moment not earlier than $m$, after which each contains a moment, $x_1$ and $x_2$, that is later than $m$ and not contained in the other. Let $p$ hold at $x_1$ only. Then $F_m p$ is true relative to $h_1$ and false relative to $h_2$, and so open at $m$. (c) If $h$ is the only history through $m$, every proposition has, relative to all histories through $m$, the value it has relative to $h$.
\end{proof}

Proposition~\ref{p:bivalence} is the formal content of the open future: the openness of a proposition at a moment is a matter of the shape of the frame above that moment, and nothing else. Combined with the thesis of TSM-12-a, it yields the conclusion that future contingents at $m$ lack truth values at $m$, without any appeal to a present that is privileged absolutely.

\begin{corollary}[Profiles and settling moments]\label{c:profile}
Along every history the number $h(m)$ of histories through its successive moments is non-increasing, and it reaches $1$ at a first moment, the \emph{settling moment} of the history, at and after which every tenseless proposition is settled.
\end{corollary}

\begin{proof}
By Lemma~\ref{l:pastclosed}, if $m$ is earlier than $m'$ then every history through $m'$ is a history through $m$, so $h(m') \leq h(m)$. The maximal moment of a history in a finite frame has exactly one history through it, so the value $1$ is reached; by Proposition~\ref{p:bivalence}(c) every proposition is settled from that moment on.
\end{proof}

Appendix~\ref{app:histories} lists, for every frame with at most nine moments, every history with its profile and its settling moment.

A point about indexicality must now be made explicit, because it is where the equivocation of Proposition~\ref{p:equivocation} returns in the setting of branching. The sentence \enquote{it will be that $p$}, evaluated at successive moments, expresses different propositions: at $m$ it expresses $F_m p$, at $m'$ it expresses $F_{m'} p$. Theorem~\ref{t:monotone} concerns a fixed proposition, and it says nothing about the sequence of propositions that a fixed sentence expresses at successive moments.

\begin{proposition}[Indexical reversal]\label{p:reversal}
There are frames, valuations and moments $m$ earlier than $m'$ such that the sentence \enquote{it will be that $p$} is settled true at $m$ (that is, $F_m p$ is settled true at $m$) and settled false at $m'$ (that is, $F_{m'} p$ is settled false at $m'$).
\end{proposition}

\begin{proof}
In the two-moment frame in which $m$ is earlier than $m'$, let $p$ hold at $m'$ only. Then $F_m p$ is true relative to the only history, and $F_{m'} p$ is false relative to it.
\end{proof}

The computation counts such reversals: there are \NReversals{} of them, summed over all frames, valuations and pairs of moments (column rev). None of them is a counterexample to Theorem~\ref{t:monotone}, and each is a case of B-variation: the sentence is true at one moment and false at a later one, as the poker is hot at one time and cold at a later one. The A-theorist who reads the reversal as the future's \enquote{becoming past} is reading a B-fact about the truth of a sentence at different moments as a change in a single proposition. That is the equivocation of Proposition~\ref{p:equivocation}, made in the vocabulary of tense logic. The same holds of the \NResolutions{} \emph{resolutions} recorded in column res: pairs of a proposition $F_m p$ open at $m$ and a later moment $m'$ at which it is settled. That a proposition is open at one moment and settled at a later one is a difference between two moments, not a movement of settledness along the series. As TSM-01-c puts it, to say that a contingent \enquote{has now become} settled is to say that it is settled at the moment of utterance.

\begin{proposition}[No thin red line from the B-facts]\label{p:trl}
A history singled out by the B-facts of a frame alone is fixed by every automorphism of the frame. There are frames in which no history is so fixed.
\end{proposition}

\begin{proof}
The first claim is proved as Proposition~\ref{p:select}(b). For the second, in the frame with three moments in which one moment precedes two incomparable ones, the automorphism exchanging the two later moments exchanges the two histories.
\end{proof}

Proposition~\ref{p:trl} supports, from the side of the B-facts, Principle P2 of TSM-12-a, that nothing relative to a moment distinguishes one history through it as the one that will be followed. A \enquote{thin red line}, a privileged actual history relative to which future contingents are true at earlier moments, would be a further fact over and above the B-structure, as an absolute present would be a further fact over and above the B-order. Of the \NFrames{} frames enumerated, \NNoFixH{} have no history fixed by every automorphism (column h-fix), so that in them a thin red line could not even be specified by the shape of the frame. We do not rest the rejection of the thin red line on this count, since a privileged history might be specified by the events on it. The point of the proposition is the parallel: the thin red line and the absolute present are the same kind of addition, one along the branches and the other along the order, and the B-theory needs neither \citep{BelnapGreen1994,MacFarlane2003}.

Theorem~\ref{t:monotone} depends on one structural feature of branching frames: the moments earlier than any moment are linearly ordered, so that the order branches toward the later and not toward the earlier. It is worth seeing that this feature is not an incidental assumption but exactly what the fixity of the past requires.

\begin{proposition}[Fixity of the past and the uniqueness of the past]\label{p:converse}
Let a strict partial order have histories as in Definition~\ref{d:history}. Settledness is monotone along the order for every tenseless proposition (as in Theorem~\ref{t:monotone}) if and only if every history through a moment passes through every moment earlier than it. In particular, if some moment $m'$ has a history through it that omits a moment $m$ earlier than $m'$, then the proposition that the history passes through $m$ is settled true at $m$ and not settled at $m'$.
\end{proposition}

\begin{proof}
If every history through $m'$ passes through every earlier $m$, monotonicity follows as in the proof of Theorem~\ref{t:monotone}. Conversely, suppose $h$ passes through $m'$ but not through $m$, where $m$ is earlier than $m'$. Let $\varphi$ be the proposition true relative to a history just if the history passes through $m$. Then $\varphi$ is true relative to every history through $m$, and so settled true at $m$; but it is false relative to $h$, which passes through $m'$, so it is not settled true at $m'$. Since $m$ itself lies on some history through $m'$ (extend the chain $\{m, m'\}$ to a maximal one), $\varphi$ is not settled false at $m'$ either.
\end{proof}

In a frame that branches toward the earlier, for instance one in which two incomparable moments $a$ and $b$ both precede a moment $c$, the histories through $c$ are $\{a, c\}$ and $\{b, c\}$, and at $c$ it is open whether $a$ occurred, although at $a$ it was settled. The fixity of the past is therefore equivalent, given the definition of histories, to the uniqueness of the past: a moment's past is fixed at it just because it has one past. This is the formal content of the asymmetry that TSM-12-a describes when it says that, relative to every moment, the past is fixed and the contingent future open, and that the asymmetry lies in the shape of the structure, not in a privileged present. It is also a reason why the asymmetry is a B-fact: it is a fact about the shape of the earlier-than order, not about where a present is.

Why should the order branch toward the later and not the earlier? The branching frame is a model of an indeterministic world in which the laws together with the course of events up to a moment do not fix what comes after it. The laws that do this in our world are time-asymmetric in their application, not in their form: the asymmetry of records, of influence and of the fixity of the past derives from the thermodynamic asymmetry and so, by the argument of TSM-05, from the low entropy of the early universe. Branching toward the later is the representation of that asymmetry in the modal structure of the world, and it requires nothing that the B-series lacks.

Two further propositions relate the framework to the semantics of tense. The first concerns the logic of future contingents at a moment, the second the alternative reading of the future tense that Prior called Peircean.

\begin{proposition}[Excluded middle without bivalence]\label{p:superval}
For every moment $m$ and every tenseless proposition $\varphi$, the disjunction $\varphi \lor \neg\varphi$ is settled true at $m$, although, if $\varphi$ is open at $m$, neither disjunct is settled true at $m$.
\end{proposition}

\begin{proof}
By Proposition~\ref{p:bivalence}(a), $\varphi$ or $\neg\varphi$ is true relative to each history, so the disjunction is true relative to every history through $m$. If $\varphi$ is open at $m$, it is false relative to some history through $m$ and true relative to another, so neither it nor its negation is true relative to all.
\end{proof}

Proposition~\ref{p:superval} is the supervaluational treatment of future contingents \citep{Thomason1970}: settled truth at a moment is supertruth over the histories through it, and it keeps classical logic while denying bivalence. TSM-12-b argues that the gap it leaves is a gap in truth at the moment, not a gap in what can be known about a sharp fact, since relative to a moment nothing grounds the choice of one history. Our framework agrees. The disjunction \enquote{there will be a sea battle or there will not} is settled at the moment of utterance; neither disjunct is; and nothing in the B-series relative to that moment selects one.

\begin{proposition}[The Peircean future]\label{p:peirce}
Read \enquote{it will be that $p$} at $m$ in Prior's Peircean way, as true at $m$ just if $F_m p$ is settled true at $m$, and \enquote{it will never be that $p$} as true at $m$ just if $F_m p$ is settled false at $m$. Then at a moment through which more than one history passes there is a valuation for which neither is true, and at a moment through which only one history passes, exactly one of them is true for every valuation.
\end{proposition}

\begin{proof}
Immediate from Definition~\ref{d:settled} and Proposition~\ref{p:bivalence}(b) and~(c).
\end{proof}

The Peircean values are what the tables record: in each row of Appendix~\ref{app:detail}, the column \emph{st} counts the valuations under which the Peircean \enquote{it will be that $p$} is true at the moment, \emph{sf} those under which \enquote{it will never be that $p$} is true, and \emph{op} those under which neither is. The Ockhamist and Peircean readings thus differ in a way the B-theory makes precise: the Ockhamist sentence expresses at $m$ a proposition true or false relative to each history, and the Peircean sentence expresses at $m$ the settledness of that proposition at $m$. Neither reading requires an absolute present, and the choice between them is a choice of which tenseless fact the future tense reports. \citet{MacFarlane2003} has proposed a third option, on which the truth of a future contingent is relative to a context of assessment as well as a context of utterance, so that a prediction made at $m$ may be neither true nor false as assessed from $m$ and true as assessed from a later moment on which the predicted event occurs. That view, too, is stated in a tenseless metalanguage: it relativizes truth to pairs of moments, and the relation it uses is the monotone settledness of Theorem~\ref{t:monotone}.

\section{The B-series as the structure of spacetime}\label{s:spacetime}

Principle~\ref{p:constitution} says that time is the B-series; Principle~\ref{p:physical} says what the B-series is in our world. The second principle is not needed to answer McTaggart, since the inversion goes through for any world in which events are ordered as earlier and later. It is needed to answer the question of the family, since a claim that time is real is a claim about our world, and it is the identification that TSM-14 relies on when it argues that time had a beginning. We defend it in four steps.

\subsection{Order and measure}

The relation of being earlier than, between events in a relativistic spacetime, is the causal order: $e$ is earlier than $e'$ if a future-directed causal curve runs from $e$ to $e'$, so that a signal sent at $e$ could reach $e'$. This relation is a strict partial order, and two events are incomparable under it just if they are spacelike separated. That the causal order carries nearly all the geometry of spacetime was shown early. \citet{Robb1914} gave axioms for a relation of \enquote{after} among the points of a four-dimensional manifold from which the whole of Minkowski geometry can be recovered, so that the causal order of special relativity determines its metric up to a constant factor. For general relativity, \citet{HawkingKingMcCarthy1976} and \citet{Malament1977} showed that in spacetimes without causal pathologies the causal order determines the topology and the conformal structure, that is, the metric up to a conformal factor. What the causal order leaves undetermined is a measure of volume: given the causal order and the volume element, the metric is fixed. This is the basis of the slogan of the causal set program, \enquote{order plus number equals geometry} \citep{Bombelli1987}.

The B-series of Definition~\ref{d:bseries} has exactly these two components: an order of earlier and later, and a measure of temporal separation. In a relativistic spacetime the second is the proper time along timelike curves, which is fixed by the metric and recorded by clocks. Principle~\ref{p:physical} therefore identifies the B-series with structure that physics already contains and that our best theories describe in detail. Nothing in that structure corresponds to an absolute present. A clock measures the proper time between two events on its worldline; it does not measure the rate at which the present moves along it, and the theory in which clocks are described has no quantity that could play that role.

\subsection{Direction}

McTaggart held that without the A-series the B-series reduces to the C-series, an order without direction, and that an order without direction is not temporal. The causal order is a strict partial order, and so has a direction in the formal sense: it distinguishes $e$'s being earlier than $e'$ from $e'$'s being earlier than $e$. But a critic may press that the formal direction is a matter of convention: the converse relation is equally a strict partial order, and nothing in the structure says which of the two is \enquote{earlier than} and which \enquote{later than}.

The reply is that the two orientations of the causal order are physically distinguished. In a time-orientable spacetime the light cones at each point divide into two halves, and a continuous choice of one half as future is a time orientation. Which orientation is the future one is fixed by the asymmetries of physical processes along it: the increase of entropy, the direction in which records and memories are formed, the direction of causal influence as agents exploit it. TSM-05 argues that these asymmetries derive from the low-entropy state of the early universe, which itself requires no further explanation. The direction of time is therefore the orientation of the physical asymmetries along the B-series; it is fixed by what occurs at each time, and it is not supplied by a moving present \citep{Price1996,Maudlin2007}. An order whose direction is grounded in this way is not a C-series. McTaggart's assumption that only the A-series could direct the order is the bridge once again: the claim that what makes time temporal is a change in A-determinations.

\subsection{Relativity and the present}

In special relativity there is no frame-independent relation of simultaneity between spacelike separated events. A global present, a set of events simultaneous with a given one, can be defined only relative to a frame, and different frames give different presents. \citet{Putnam1967} argued that this refutes the view that only present things are real, since by the relativity of simultaneity and the transitivity of \enquote{being real for}, every event would be real for every observer. \citet{Godel1949} argued, more radically, that relativity shows the \enquote{objective lapse of time} to be an illusion, and that in some solutions of the field equations there is not even a global time function. Gödel's conclusion resembles McTaggart's, and our response to it is the same. What relativity shows is the absence of a frame-independent global present, and hence the absence of anything that could be the A-series of the world. It does not show that time is unreal, since the causal and metrical structure that relativity describes is the B-series, and that structure is not frame-relative. TSM-16 draws the conclusion that the present is a local rather than a global notion.

\citet{Stein1968} argued that, in Minkowski spacetime, the only relation of \enquote{having become} that respects the structure is relative to a point: the events that have become, relative to a point-event, are those in or on its past light cone. Stein's relativized becoming is, in our terms, a moment-relative notion like settledness, and it has the formal property that makes settledness innocuous: it is defined at every point, and no point is privileged. We accept Stein's notion as a tenseless relation and reject the inference, sometimes drawn from it, that it vindicates an A-theory. A relation that every event bears to the events in its past cone is a B-relation.

The formal core of Assumption~\ref{a:incoherence} uses the fact that two distinct moments of one history are comparable. In a relativistic setting, an A-theorist might try to escape it by making many spacelike separated events present together, the present being a spacelike hypersurface. This saves uniqueness at the cost of choosing a hypersurface, which, in the absence of a privileged frame, is a choice that the B-facts do not make. The A-theorist must then hold that reality contains a privileged foliation which physics does not detect. Such a view is consistent, but it adds to the B-series a structure for which there is no evidence and which does no physical work, and it remains subject to Theorem~\ref{t:nochange}, since the privileged hypersurface could change only by being had by more than one of the hypersurfaces of a single history of foliations.

\subsection{Branching in spacetime}

The branching frames of Section~\ref{s:branching} have moments as their points: global states of the world that are linearly ordered on each history. In a relativistic setting, the points are better taken to be point-events, and the branching structure to be a branching of spacetimes. \citet{Belnap1992} developed a theory of branching space-times in which histories are maximal directed sets of point-events and the past of a point-event need not be linearly ordered. The results we need carry over. What Lemma~\ref{l:pastclosed} requires is that a history containing an event contain everything in that event's causal past, and in Belnap's theory histories are closed downward in this sense. Settledness at a point-event is then defined as in Definition~\ref{d:settled}, and Theorem~\ref{t:monotone} holds with the same proof. We work with linear moments in the body of the paper and in the computation because they suffice for the questions McTaggart raises, all of which concern the order of times along a history.

\subsection{Realism about the B-series}

An idealist may grant the inversion and hold that it shows only that the B-series is the order in which we represent events: Kant held that time is the form of inner sense, and McTaggart that the B-series is the appearance of a non-temporal C-series. On either view the B-series is real as a representation and not as a feature of things in themselves. This raises a different question from McTaggart's, that of realism about the physical world. Principle~\ref{p:physical} identifies the B-series with the causal and metrical structure of spacetime, and our grounds for thinking that this structure is mind-independent are our grounds for thinking that spacetime and its contents are: the success of theories that describe it, the agreement of independent measurements, and the absence of any account of why a merely represented order would display the structure it does. An idealist who holds that the whole physical world is appearance will hold that time is too, and we do not argue against general idealism here. What the idealist cannot do is use McTaggart's argument to single out time: the argument shows that the A-series is incoherent, and it gives no reason to treat the B-series differently from any other physical structure. McTaggart's own reason for doing so was his reduction of the B-series to the A-series and the C-series, which presupposes the bridge.

\subsection{The beginning of time}

Principle~\ref{p:physical} has a consequence that TSM-14 develops. If time is the causal and metrical structure of spacetime, then whether time had a beginning is the question whether that structure has a finite past: whether there is a bound on the duration of every past-directed history. The singularity theorems of general relativity give evidence that our universe has a finite past, and on the identification of the B-series with that structure, a question about what preceded every time presupposes a time outside the series. In the finite frames of the appendices every frame has a least moment, and every history begins there; nothing in a frame is earlier than its least moment, and nothing in the definitions invites the question what is.

\section{Against the growing block and the other A-theories}\label{s:block}

The open future is often thought to favor the growing block. If the future is open and the past fixed, it seems that the past exists and the future does not, and that the boundary between them, the present, moves as the future becomes past \citep{Broad1923,Tooley1997,CorreiaRosenkranz2018}. Branching is then naturally pictured as a tree that grows: at each stage the tree consists of the past and present, and its growth selects one continuation and discards the others.

\begin{definition}[Growing block]\label{d:block}
The \emph{growing block} is the view that the totality of existing moments is the set of moments up to and including an absolute present, its \emph{edge}, and that the edge moves, so that the totality of what exists is different at different stages.
\end{definition}

\begin{theorem}[The growing block is excluded]\label{t:noblock}
The growing block requires a changing A-series and is therefore excluded. Its static counterpart, on which every moment exists and settledness is relative to moments, is the branching B-theory of Section~\ref{s:branching}.
\end{theorem}

\begin{proof}
The edge of the block is an absolute present in the sense of Definition~\ref{d:adet}: it is the moment that is present simpliciter, and what exists is fixed by it. For the edge to move is for the present to be at more than one moment, which by Definition~\ref{d:changing} is for the world to contain a changing A-series. By Theorem~\ref{t:nochange} no world does. If the edge is relativized to moments, so that relative to each moment the block extends up to that moment, then by Corollary~\ref{c:dilemma} the relativized edge is the moment-relative A-assignment of Definition~\ref{d:relative}, and the view is a description of the whole frame, every moment of which exists, together with the relation \enquote{earlier than or equal to}. What varies between moments is then which moments are earlier, and which propositions are settled; by Theorem~\ref{t:monotone} the latter is a structural fact about the frame.
\end{proof}

The growing block has a further cost that TSM-12-a, following \citet{Bourne2002} and \citet{BraddonMitchell2004}, presses against it. On the block view, the past is as real as the present; past thinkers exist, and many of them believe, as we do, that they are at the edge. They are wrong. Since their evidence is the same as ours, we have no reason to think that we are not in their position, and the growing-blocker cannot know that it is now now. \citet{Forrest2004} replied that only the present is conscious: the past is real but dead. The reply concedes that existence does not mark the present, since the dead past exists, and makes consciousness do the marking instead, at the cost of holding that the conscious lives of our predecessors are not conscious. On our view the argument does not arise. The present is the time of the thought, and every thinker who believes that it is now the time of her thought is right \citep{Mellor1998}.

Settledness is the notion the growing-blocker wanted, and it is available without the block. The past of a moment is settled at that moment, the future of a moment may be open at it, and every history through a moment contains its whole past (Lemma~\ref{l:pastclosed}). These are relations of moments to propositions. They are moment-relative in the way that TSM-12-a requires: existence is not relativized to moments, only settledness is. TSM-03 adds that the past exists in the way distant places do, as elsewhere in time; we add that the branches above a moment exist in the same way, as other ways in which the world goes on from that moment, none of which is \enquote{the one that will come to exist}, since there is no moving present for it to come to exist at.

The other A-theories fare no better. The moving spotlight holds that all times exist and that one of them is absolutely present, the present moving from earlier to later \citep{Cameron2015,Skow2015}. The spotlight is an absolute present, and its motion is a changing A-series; Theorem~\ref{t:nochange} applies directly. Cameron's version replaces the monadic present by a fundamental property of temporal distance from the present, had by each time in some degree; but for these properties to change is for each time to have different degrees at different stages, and Corollary~\ref{c:dilemma} applies to the stages. Presentism, on which only present things exist, makes the present absolute in the strongest way, since it is the boundary of existence itself, and it faces in addition the objection from relativity and the problem of grounding truths about the past \citep{Markosian2004}; these are the subject of TSM-03 and TSM-16, and we do not repeat them here. Each view privileges one moment absolutely and has it move, and that combination is what Theorem~\ref{t:nochange} excludes.

Fine's non-standard views are different, and we treat them as objections.

\section{Cases and variants}\label{s:cases}

The principles of Section~\ref{s:framework} should give a verdict on the cases that have been used to test theories of time. We consider six.

\paragraph{The frozen world.} \citet{Shoemaker1969} described a world divided into three regions, each of which undergoes a total freeze at regular intervals, every third, fourth and fifth year respectively, for a year at a time. The inhabitants observe the local freezes from outside and infer, by the simplest law that fits their observations, that every sixtieth year all three regions freeze together. If they are right, a year passes in which nothing changes anywhere. TSM-17 argues that this is possible and could be known to occur. On our view the case is straightforward. A total freeze is a stretch of the B-series whose times are qualitatively alike: the order and the duration are there, and the change is not. Assumption~\ref{a:m1}, in its B-reading, is false in the frozen world, and nothing in the inversion depended on it. The case also confirms Proposition~\ref{p:equivocation} from the other side. The considerations that recommend Assumption~\ref{a:m1}, that idle time could not be detected and that times are constructed from the states of things, are exactly the considerations Shoemaker's case is designed to answer, and they concern B-change. McTaggart's A-reading of the premise is not engaged by the case at all.

\paragraph{The one-moment world.} A world with a single moment, the frame T1.1, is the limiting case. Its B-series is a one-element order with no pairs in the earlier-than relation and no temporal separations. On Principle~\ref{p:constitution} there is time in such a world only in a degenerate sense: there is a moment, but nothing is earlier or later than anything. It is also the only world in which a single A-assignment makes every moment present (Proposition~\ref{p:stage}(d)), and so the only world in which an A-series is complete in McTaggart's sense. This is no comfort to the A-theorist. The world in which the A-series is complete is the world in which nothing passes, since the only event is present and has never been future and will never be past. McTaggart's requirement that every event be past, present and future is unsatisfiable in every world with more than one moment and vacuous in the one world where it can be met.

\paragraph{The twins.} Two clocks are synchronized at an event $e$; one remains inertial and the other travels and returns, and at their reunion $e'$ the traveling clock records less elapsed time. The B-theory describes the case directly. The causal order puts $e$ earlier than $e'$; the metrical structure assigns different proper times to the two worldlines between them; each clock correctly measures the proper time along its own worldline. Principle~\ref{p:physical} identifies the temporal separation of events along a curve with this proper time. An A-theorist who holds that the present advances must say how fast it advances along each worldline. If it advances at one proper second per second along each, the \enquote{rate} is the B-fact of proper time under another description \citep{Markosian1993}, which is the topic of TSM-02-b; if it advances along a privileged global time, the theory adds a structure that the clocks do not record.

\paragraph{Time travel and closed curves.} TSM-04 argues that backward time travel is consistent provided the past is not changed, and that the grandfather paradox shows only that the traveler fails. The B-theory explains why the past cannot be changed. To change the past would be for an event to occupy a position in the B-series and later not to occupy it, which would be for its location in the order to vary across stages; and a variation across stages of what is at each position is a changing A-series in the sense of Corollary~\ref{c:dilemma}, now applied to occupancy rather than presentness. On the B-theory each event is where it is in the order, once; the traveler's attempt and its failure are both in the order, and the failure is a B-fact about the attempt. Spacetimes with closed timelike curves, such as the solution described by \citet{Godel1949}, raise a different issue: in them the causal order is not irreflexive, since an event can be causally earlier than itself. Principle~\ref{p:physical} must then be read locally. In any sufficiently small region the causal order is a strict partial order and is the B-series of that region; globally, the structure is a causal structure that fails to be an order. We do not think this threatens Principle~\ref{p:constitution}, since a causally pathological spacetime still has a causal and metrical structure, and the temporal facts in it are facts about that structure; but we record it as a case in which the identification of time with an order is local rather than global.

\paragraph{Foreknowledge.} REL-05-b argues that Ockhamism, in branching time with moment-relative settledness, is a confused statement of the view that God's knowledge is timeless. The argument uses the monotonicity of settledness in exactly the form of Theorem~\ref{t:monotone}: every proposition settled at a moment is settled at every later moment, so that a divine belief at an earlier moment about a contingent choice would have to be settled before the choice was, which branching excludes. The soft fact, correctly described, is a tenseless truth relative to a history. Our framework supplies the distinction REL-05-b needs: between truth relative to a history, which is bivalent (Proposition~\ref{p:bivalence}(a)), and settledness at a moment, which is not (Proposition~\ref{p:bivalence}(b)). A timeless knower knows the former and has no position in the order relative to which the latter could be assessed; REL-19 draws the consequence that such a knower lacks indexical knowledge of what time it is now.

\paragraph{Self-location.} In the Sleeping Beauty problem, Beauty wakes not knowing whether it is Monday or Tuesday. EPI-18-c argues that her uncertainty is uncertainty about her location in a world whose history she knows apart from the coin, and that her credences are defined on centered worlds. On our view this is the general form of tensed uncertainty. Every question of the form \enquote{is it now $t$?} is a question about where the thought is in the B-series. The B-series itself is not uncertain for her in that respect: she knows that there is an awakening on Monday and, on tails, one on Tuesday. What she lacks is self-locating information, and that is all that knowing what time it is ever comes to.

\section{Consequences for neighboring questions}\label{s:neighbors}

The answer to the question whether time is real constrains the answers to several neighboring questions, and the families that address them presuppose it. We indicate how.

\paragraph{Passage (TSM-02).} If the A-series is incoherent, nothing passes, since passage is the change of A-determinations, the present moving from earlier to later times. TSM-02-a takes Theorem~\ref{t:nochange} as a premise and adds an independent argument that the impression of passage is fixed by the B-facts about experience. The two arguments are complementary: ours shows that there is no passage for the impression to represent; TSM-02-a's shows that the impression would be had whether or not there were. The impression itself is not thereby an illusion. Its source lies in the structure of experience, in the succession of experiences each of which represents its predecessors as just past \citep{Williams1951}; what is mistaken, if anything, is the inference from it to a moving present, and the view requires only that a description of experience as of passage not be read as a description of time.

\paragraph{The past and the present (TSM-03, TSM-15, TSM-16, TSM-18).} On the B-theory every moment is as real as every other, and the past is elsewhere in time, as distant places are elsewhere in space (TSM-03). The present has no metaphysical privilege; it is the time at which a thought or utterance occurs (TSM-15), and \enquote{now} is an indexical (TSM-18). Proposition~\ref{p:select} shows from our side why a privileged present would have to be a further fact not fixed by the B-facts, and TSM-16 shows that relativity leaves no frame-independent global present for such a fact to concern.

\paragraph{Direction (TSM-05).} The B-series has a direction, and the direction is not supplied by a moving present. TSM-05 grounds it in the low-entropy initial state. Section~\ref{s:spacetime} uses that result to answer McTaggart's claim that without the A-series the B-series is a C-series.

\paragraph{The open future (TSM-12, ACT-17).} TSM-12-a argues that future contingents are neither true nor false until they are settled, and ACT-17 that fatalism is therefore unsupported. Section~\ref{s:branching} supplies the B-theoretic setting: settledness is moment-relative and monotone, the past of each moment is fixed at it because it has one past (Proposition~\ref{p:converse}), and the open future at a moment is the plurality of histories through it (Proposition~\ref{p:bivalence}). No moment is the edge of a growing block (Theorem~\ref{t:noblock}). It may be objected that if every moment is as real as every other, including the moments later than this one, it is already true what will happen, so that the B-theory entails fatalism. The objection runs together two senses in which the future might be \enquote{real}. In a branching frame every moment exists tenselessly, including every moment on every branch above this one. What does not exist is a fact, relative to this moment, about which branch is the one that will be followed: by Proposition~\ref{p:trl} nothing in the B-structure singles one out, and Principle P2 of TSM-12-a denies that anything else does. The reality of later moments is therefore compatible with the openness of what happens later, since what is real above a moment is several ways the world goes on, and a future contingent at the moment is true on some of them and false on others. That is why, as TSM-12-a and ACT-17 conclude, future contingents lack a truth value at the moment and fatalism is unsupported. In a deterministic world, the frame is a chain and every future-tensed proposition is settled; but then fatalism follows from determinism, not from the reality of the future.

\paragraph{The beginning of time (TSM-14).} TSM-14-a identifies time with the causal and metrical structure of spacetime, as Principle~\ref{p:physical} does, and argues that this structure has a finite past and that the question what preceded it is ill-formed. The identification is what makes the question of a beginning a question about the structure of spacetime rather than about a time outside it.

\paragraph{Time without change (TSM-17).} TSM-17-a argues, with Shoemaker, that time without change is possible. On our view this is the falsity of Assumption~\ref{a:m1} in its B-reading, which the inversion does not need. McTaggart's premise that time requires change, insofar as it was a premise of his argument, concerned the change of A-determinations (Corollary~\ref{c:survives}).

\paragraph{What time was asked to do.} TSM-01-c lists the roles standardly assigned to time: to order events and measure their separation, to make change possible, to give the order a direction, to make tensed sentences true or false, to leave the future open while the past is fixed, and to present itself as passing. It derives that the B-series, with the location of thinkers in it, suffices for each. The results of this manuscript supply the parts of that derivation that concern the order, change, the open future and the fixity of the past. The tenseless view, branching and with moment-relative settledness, gives up none of what time was asked to do.

\section{Objections and replies}\label{s:objections}

\begin{objection}[The C-series]\label{o:cseries}
Without the A-series, the B-series is a mere order, McTaggart's C-series. An order of positions is no more temporal than an order of points on a line or of numbers by magnitude. What makes an order temporal is that its members pass from future to past. Hence time is not constituted by the B-series.
\end{objection}

\begin{reply}
The objection assumes that the B-series is distinguished from other orders only by its form. It is not. The B-series is the order of events by the relation of being earlier than, and this relation is distinguished from spatial and numerical orders by its role in causation, in the laws of motion, in the formation of records and in the orientation of the thermodynamic asymmetry (Section~\ref{s:spacetime}). An order of events in which causes precede their effects, entropy increases toward the later, and memories record the earlier and not the later is a temporal order, and nothing in it is a passage from future to past. The objection's premise that only passage makes an order temporal is an instance of the bridge, and Proposition~\ref{p:equivocation} has shown that the bridge has no support of its own.
\end{reply}

\begin{objection}[Dummett: the completeness of description]\label{o:dummett}
A complete description of a temporal world must say which events are past, which present and which future. No description from no temporal standpoint can say this, and any description from a standpoint is incomplete, since it omits how things are from other standpoints. Hence either reality has no complete description or the B-theory omits something real \citep{Dummett1960}.
\end{objection}

\begin{reply}
The first premise trades on an ambiguity in \enquote{say which events are past}. A description that says, of each event, at which moments it is past (all those later than it) says which events are past relative to each moment, and this description is complete and from no standpoint: it is the moment-relative A-assignment of Definition~\ref{d:relative}, which Proposition~\ref{p:stage} shows to be equivalent to the B-series. What it does not say is which events are past relative to the moment at which the description is read. That is self-locating information, and the parallel with space and with persons shows that its absence is no incompleteness in the description of the world. A complete atlas does not say which place is here, and a complete biography of everyone does not say which of them I am \citep{Perry1979,Lewis1979}. What Dummett's point establishes is that thought has a perspective, not that time does. EPI-18-c gives the same treatment of Beauty's uncertainty whether it is now Monday, and REL-19 of the indexical knowledge that a timeless God lacks without defect in omniscience.
\end{reply}

\begin{objection}[Prior: thank goodness that's over]\label{o:prior}
When a headache ends, I am glad that it is over. What I am glad about is not that its end is earlier than my utterance, a tenseless fact I could have known in advance and to which I am indifferent, but that it is now over. The B-theory has no fact for my gladness to be about \citep{Prior1959}.
\end{objection}

\begin{reply}
The content of the relief is self-locating: that the headache ends before \emph{this} time, the time of the thought. Its content is true at some centers and false at others, and it is the content one has reason to welcome, because the location of one's present thinking relative to the headache is what determines whether one is in pain. That the B-theory has no non-indexical fact that one could be glad about in advance is correct and is no objection, since one could not, in advance, be glad of a content whose truth depends on the time of the gladness. The reply is Mellor's \citep{Mellor1998}, and TSM-01-c shows that it extends to every role tense is asked to play. On our view the facts that make the relief appropriate are B-facts: the relief occurs later than the headache, and its subject is not in pain at its time.

The same holds of a related objection, that action requires tensed facts. An agent who must leave at noon acts when she believes that it is now noon, and no belief with tenseless content, such as the belief that her leaving is simultaneous with noon, explains why she acts now rather than at some other time. The argument is Perry's argument for the essential indexical, applied to time, and Perry's conclusion is that indexical belief is essential to the explanation of action, not that there are indexical facts \citep{Perry1979}. The agent who believes that it is now noon believes that her thought occurs at noon; the content is self-locating, and it explains her acting at the time of the thought because the belief and the action are at the same time. The same structure explains why a person who believes \enquote{the bear is attacking me} runs, while one who believes \enquote{the bear is attacking the person in the mirror} need not. In both cases what explains action is the location of the believer, and in both cases the facts that make the belief true are non-indexical: that the thought occurs at noon, that the bear attacks the believer. The timing of action is therefore explained by B-facts about where the believer and her thoughts are located \citep{Mellor1998}.
\end{reply}

\begin{objection}[Fine: relativism and fragmentalism]\label{o:fine}
The B-theorist rejects Realism, the claim that reality is constituted in part by tensed facts. But there are A-theories that keep Realism and Neutrality and give up Absolutism or Coherence instead \citep{Fine2005}. External relativism holds that reality is constituted only relative to a time; fragmentalism holds that reality consists of fragments, each coherent and oriented toward one time, which do not cohere with one another. Neither privileges a present, and neither is touched by Theorem~\ref{t:nochange}.
\end{objection}

\begin{reply}
We take the two views in turn. External relativism says that reality relative to $t$ contains the tensed facts oriented toward $t$. The temporal content of this reality is the family $\{A_t\}$ of Definition~\ref{d:relative}: relative to $t$, the events earlier than $t$ are past, $t$ is present and the later events future. By Proposition~\ref{p:stage}, this family is equivalent to the B-order, and the relativist's \enquote{reality relative to $t$} is the B-series seen from $t$. If the relativist adds that one of these relativized realities is the one that obtains, Absolutism returns and with it Theorem~\ref{t:nochange}; if not, the view differs from ours in the vocabulary in which it describes the same structure. Fragmentalism grants what TSM-01-b notes the inversion needs: a reality whose tensed facts do not form a coherent whole is one in which the A-series is not coherent. Our use of Assumption~\ref{a:incoherence} is only to conclude that time does not consist in an A-series that is coherent and absolute. The fragmentalist may say that time consists in an incoherent totality of coherent fragments. We reply that the fragments, too, are ordered: the fragment oriented toward $t$ and the fragment oriented toward $t'$ stand in the relation that $t$ and $t'$ do, and the order of fragments is again the B-series. Fragmentalism adds to the B-series a refusal to combine the fragments, and that refusal does no temporal work: everything time is asked to do is done by the order of the fragments and their content, as TSM-01-c argues.
\end{reply}

\begin{objection}[Becoming in a branching world]\label{o:becoming}
In an indeterministic world the future becomes settled. At one moment it is open whether the coin lands heads; at a later moment it is settled. Becoming settled is a change, and it is not B-change in any object: it is the change of a proposition from open to settled. Hence the open future requires A-change after all.
\end{objection}

\begin{reply}
That a proposition is open at $m$ and settled at a later $m'$ is a difference between two moments in their relation to the proposition: $m$ stands in the relation \enquote{settles} to it and $m'$ does not. It is a B-fact about the branching structure, of the same kind as the poker's being hot at one time and cold at another. Theorem~\ref{t:monotone} says what is invariant along a history: what is settled stays settled. Proposition~\ref{p:reversal} says what is not: the proposition expressed by a tensed sentence differs between moments. Neither requires that settledness move along the frame. The objector's \enquote{change of a proposition from open to settled} is, on inspection, a pair of facts about two moments described as if it were a single fact about one proposition across stages; that is the equivocation of Proposition~\ref{p:equivocation}, and the computation counts \NResolutions{} such pairs among the frames enumerated without finding any that is anything else. To say that the coin's landing \enquote{has now become} settled is to say that it is settled at the moment of utterance (TSM-12, TSM-18).
\end{reply}





\begin{objection}[Lowe: the indexical fallacy]\label{o:lowe}
Arguments of McTaggart's kind assume that tensed sentences have tenseless truth conditions, so that \enquote{$e$ is past} uttered at $t$ is true just if $e$ is earlier than $t$. The A-theorist denies this, and so the argument begs the question \citep{Lowe1987}.
\end{objection}

\begin{reply}
Our argument does not assume it. Theorem~\ref{t:nochange} uses only the formal core of the incoherence premise, that absolute presentness is had by at most one moment, which the A-theorist accepts, since it is the uniqueness of the present. Corollary~\ref{c:dilemma} uses only the claim that stages, if they are distinct from moments, are ordered. The tenseless truth conditions of tensed tokens are a conclusion of the argument: the consistent reading of the tensed copula is evaluation at a moment, and evaluation at a moment is the family $\{A_m\}$, whose content is the B-order. As \citet{LePoidevinMellor1987} replied to Lowe, an A-theorist who rejects the conclusion owes an account of what makes a tensed token true, and the dilemma shows what such an account must either repeat or regress through.
\end{reply}


\section{The computation}\label{s:computation}

Several of the results of Sections~\ref{s:incoherence} and~\ref{s:branching} are general theorems with short proofs. We have also verified them exhaustively on all finite branching frames up to a bound, for two reasons. The first is the usual one: a mechanical check of a general claim on all small cases guards against errors in the proofs and in the definitions, especially in definitions of maximal chains and of settledness, where small slips produce false theorems. The second is particular to the subject. The claims of the A-theorist and the B-theorist about branching time are often made with diagrams of one or two forks, and it is useful to have, for every frame up to a bound, the exact pattern of settledness at every moment under every valuation, together with the symmetry facts on which Propositions~\ref{p:select} and~\ref{p:trl} depend.

\subsection{What is enumerated and checked}

A finite branching frame (Definition~\ref{d:frame}) is a finite rooted tree, the order being the ancestor relation. We enumerate every frame with at most eleven moments, one per isomorphism class. Frames are generated by adding a maximal moment in every possible position to every frame with one moment fewer and reducing each result to a canonical form; the canonical form is the level sequence of the tree in which the subtrees of every moment are arranged in decreasing order of their own canonical sequences, which is a complete isomorphism invariant for rooted trees \citep{AhoHopcroftUllman1974,BeyerHedetniemi1980}. The numbers of classes obtained for $n = 1, \dots, 11$ moments are $1, 1, 2, 4, 9, 20, 48, 115, 286, 719, 1842$, which are Otter's numbers of rooted trees \citep{Otter1948,Cayley1857}. As an independent check of exhaustiveness, for $n \leq 10$ the program also canonicalizes every labeled tree whose moments are numbered so that each moment's predecessor has a smaller number, \NLabelled{} labeled trees in all, and confirms that they yield exactly the same classes.

For each frame and for every one of the $2^n$ valuations of an atomic sentence $p$ the program computes, at every moment $m$, whether $F_m p$ is settled true, settled false or open, and checks the following.
\begin{enumerate}[label=(\roman*)]
\item \emph{Monotonicity} (Lemma~\ref{l:pastclosed}, Theorem~\ref{t:monotone}): every history through a later moment passes through every earlier one, and for every pair of moments $m$ earlier than $m'$, if $F_m p$ is settled at $m$ it is settled with the same value at $m'$.
\item \emph{Bivalence and its failure} (Proposition~\ref{p:bivalence}): truth relative to a history partitions the valuations, while $F_m p$ is open at $m$ for some valuation exactly when $h(m) \geq 2$.
\item \emph{Profiles} (Corollary~\ref{c:profile}): along every history the profile of $h(m)$ is non-increasing from the number of histories to $1$.
\item \emph{A-assignments} (Proposition~\ref{p:stage}): the moment-relative assignments recover the earlier-than order and its converse, commute with a generating set of automorphisms, and no single assignment makes every moment present unless $n = 1$.
\end{enumerate}
It also computes the order of the automorphism group, its orbits on moments and on histories, and the counts of resolutions and reversals defined in Section~\ref{s:branching}. Every check holds in every case. Table~\ref{tab:stats} summarizes the results; Appendix~\ref{app:detail} gives the per-moment values for the \NFramesDetail{} frames with at most ten moments, Appendix~\ref{app:histories} the histories of the frames with at most nine moments, and Appendix~\ref{app:eleven} one line for each of the \NFramesEleven{} frames with eleven moments.

\begin{table}[ht]
\centering
% generated by gen/generate.py -- do not edit
{\small\setlength{\tabcolsep}{4pt}
\begin{tabular}{@{}rrrrrrrrrr@{}}
\toprule
$n$ & frames & labelled & $\sum|H|$ & linear & rigid & no fixed $h$ & M checks & resolutions & reversals\\
\midrule
1 & 1 & 1 & 1 & 1 & 1 & 0 & 0 & 0 & 0\\
2 & 1 & 1 & 1 & 1 & 1 & 0 & 4 & 0 & 2\\
3 & 2 & 2 & 3 & 1 & 1 & 1 & 40 & 8 & 16\\
4 & 4 & 6 & 8 & 1 & 2 & 2 & 288 & 84 & 104\\
5 & 9 & 24 & 22 & 1 & 3 & 4 & 1984 & 660 & 688\\
6 & 20 & 120 & 58 & 1 & 6 & 6 & 12480 & 4476 & 4184\\
7 & 48 & 720 & 160 & 1 & 12 & 14 & 79104 & 29460 & 25992\\
8 & 115 & 5040 & 434 & 1 & 25 & 26 & 483584 & 185144 & 156082\\
9 & 286 & 40320 & 1204 & 1 & 52 & 56 & 2960896 & 1154152 & 942788\\
10 & 719 & 362880 & 3341 & 1 & 113 & 113 & 17920000 & 7084856 & 5640578\\
11 & 1842 & -- & 9363 & 1 & 247 & 242 & 108300288 & 43289136 & 33764700\\
\midrule
 all & 3047 & & 14595 & 11 & 463 & 464 & 129758668 & 51747976 & 40535134\\
\bottomrule
\end{tabular}}

\caption{Summary of the enumeration by number of moments $n$. \emph{Frames}: isomorphism classes; \emph{labeled}: labeled trees canonicalized in the independent check; $\sum|H|$: total number of histories; \emph{linear}: frames with one history; \emph{rigid}: frames with trivial automorphism group; \emph{no fixed $h$}: frames in which no history is fixed by every automorphism; \emph{M checks}: instances of Theorem~\ref{t:monotone} checked (pairs $m < m'$ times valuations); \emph{resolutions} and \emph{reversals} as in Section~\ref{s:branching}.}
\label{tab:stats}
\end{table}

\subsection{A row reproduced by hand}

The values in the tables can be checked by hand for small frames, and doing so for one row shows what they mean. Take the root of T4.3, whose moments are $0$ (the root), $1$ and $3$ (its successors) and $2$ (the successor of $1$). The histories through the root are $\{0, 1, 2\}$ and $\{0, 3\}$. The proposition $F_0 p$ is true relative to the first if $p$ holds at $1$ or $2$, and relative to the second if $p$ holds at $3$. It is settled true at the root just if both conditions hold: $p$ at $3$, and $p$ at $1$ or at $2$. Of the $2^4 = 16$ valuations, those satisfying both are fixed at moment $3$ and allow three of the four combinations at moments $1$ and $2$, with moment $0$ free, which gives $3 \times 2 = 6$; this is the entry \emph{st}. It is settled false just if $p$ holds at none of $1$, $2$, $3$, which leaves moment $0$ free and gives $2$, the entry \emph{sf}. The remaining $16 - 6 - 2 = 8$ valuations make it open, the entry \emph{op}. At moment $1$, only one history passes, so nothing is open; $F_1 p$ is settled true just if $p$ holds at $2$, which gives $8$ valuations. Every valuation under which $F_0 p$ is settled true at the root remains one under which $F_0 p$ is settled true at $1$, since the history through $1$ is one of the two through the root; that is the instance of Theorem~\ref{t:monotone} the program checks. The program performs this computation for every moment of every frame, for every valuation, and confirms by direct evaluation for all frames with at most six moments that the bit-string method gives the same counts.

\subsection{Patterns in the tables}

Three patterns in Table~\ref{tab:stats} are worth recording. First, the number of rigid frames, those with no non-trivial automorphism, grows with $n$, from two of the four frames at $n = 4$ to 247 of the 1842 at $n = 11$, while their proportion falls; the number of frames in which no history is fixed by every automorphism grows in step with it. So the two situations that Proposition~\ref{p:select} distinguishes, that in which the B-facts single out every moment and that in which they leave pairs of moments indiscernible, both persist as frames grow, and neither offers the A-theorist a present selected by the B-facts. Second, resolutions outnumber reversals from $n = 6$ on, and the ratio increases steadily with $n$, from about 1.07 at $n = 6$ to about 1.28 at $n = 11$, as the number of branching moments, and so of open future contingents to be resolved by later moments, grows. Third, the chains, which are the deterministic frames, are the only frames without resolutions, while the stars have, for each $n$, the fewest reversals: in T11.1842 there are 20, against 20,440 resolutions among its 20,480 instances of monotonicity. In a deterministic world all variation across moments in the truth of tensed sentences is reversal; in a maximally branching one nearly all is resolution. Neither kind of variation is a change in a single proposition.

\subsection{Particular entries}

Some entries illustrate the results with special clarity.

\paragraph{T1.1.} The frame with one moment is the only frame in which a single A-assignment makes every moment present (Proposition~\ref{p:stage}(d)). It is also the only frame in which the A-theorist's present could be fixed by the B-facts without remainder, since there is only one moment for it to be at. It is the degenerate case in which the A-series and the B-series coincide, because there is no order for them to disagree about. With one moment, $F_m p$ is settled false under both valuations, since nothing is later.

\paragraph{T2.1 and the chains.} The two-moment chain is the frame used in the proof of Proposition~\ref{p:reversal}: under the two valuations in which $p$ holds at the later moment, \enquote{it will be that $p$} is settled true at the earlier moment and settled false at the later. Its entry records two reversals. The chains T$n$.1 are the deterministic frames: one history, all moments fixed by the (trivial) automorphism group, no open entries and no resolutions. T10.001 records 46,080 instances of monotonicity, all trivial since every proposition is settled at every moment, and 14,362 reversals; T11.0001 in Appendix~\ref{app:eleven} records 32,796 reversals and no resolutions. In a deterministic world, then, every tensed sentence is settled at every moment, and every variation in the truth of \enquote{it will be that $p$} is a reversal of the kind the A-theorist misreads as A-change. The linear B-series is the special case of the branching one in which the open future is empty.

\paragraph{T3.2.} The fork, a moment with two incomparable successors, is the smallest frame in which settledness fails to be bivalent and the frame used in the proof of Proposition~\ref{p:trl}. At the root $F_0 p$ is settled true under 2 valuations (those in which $p$ holds at both later moments), settled false under 2 (those in which it holds at neither), and open under 4. The single-point vector at the root, \texttt{-??}, shows that $p$ at either later moment alone yields an open future contingent. The automorphism group has order 2 and exchanges the two histories, so neither is fixed: no thin red line can be specified by the shape of the frame. It is the frame formalized in the Lean module.

\paragraph{T4.3.} The smallest frame that branches and is rigid: the root has one successor with a further successor and one maximal successor. Every moment and every history is fixed by the trivial automorphism group, so every moment is singled out by the B-facts, and Proposition~\ref{p:select}(b) offers the A-theorist no help in choosing a present: the B-facts single out each of the four moments equally. At the root, $F_0 p$ is open under 8 of the 16 valuations, and its single-point vector \texttt{-???} shows that $p$ at moment 2, which lies on only one of the two histories, is open at the root but, as the row for moment 1 shows (\texttt{--+-}), settled true at moment 1. This is Theorem~\ref{t:monotone} in its simplest non-trivial instance.

\paragraph{T7.42.} A frame in which the root has two isomorphic successors, each of which is a fork. Its automorphism group has order 8, and of its three orbits on moments only the root's is a singleton. The root is thus the only moment singled out by the B-facts, and every other moment has an indiscernible partner, so that the A-theorist's present, placed at any moment but the root, distinguishes worlds that differ in nothing else (Proposition~\ref{p:select}(c)). None of its four histories is fixed.

\paragraph{T10.719 and T11.1842.} The stars, one moment with nine or ten maximal successors, maximize symmetry: their automorphism groups have orders $9! = 362{,}880$ and $10! = 3{,}628{,}800$. At the root of T10.719, $F_0 p$ is settled true under 2 of the 1024 valuations and settled false under 2, and open under the remaining 1,020; nearly every proposition about the immediate future is open, and every one of them is settled at the next moment, whichever it is. The star is the frame in which the open future is as open as the number of moments allows.

\medskip
In aggregate, of the \NFrames{} frames, \NRigid{} are rigid and \NNoFixH{} have no history fixed by every automorphism. The \NMonoChecks{} instances of Theorem~\ref{t:monotone} hold without exception, as the proof requires; the \NReversals{} reversals and \NResolutions{} resolutions show how much variation across moments there is in the truth and settledness of tensed sentences, none of which is a change in a single proposition.

\section{What the result does not establish}\label{s:limits}

The argument establishes that the A-series is incoherent in the sense of Theorem~\ref{t:nochange} and Corollary~\ref{c:dilemma}, and that, given the reality of time, the bridge premise of McTaggart's argument is false. It does not establish the incoherence premise from premises the A-theorist must accept. We have stated its formal core, that absolute presentness is had by at most one moment, and argued that the A-theorist who denies it gives up the uniqueness of the present that distinguishes the A-series from a labeling of moments. An A-theorist who holds that absolute presentness is had by many moments, but by each at a different stage, has been answered by Corollary~\ref{c:dilemma}, whose proof relies on the claim that the stages are ordered. A view on which they are not ordered at all would not be a theory of time, and we have not addressed it.

The argument does not establish Principle~\ref{p:physical} from the philosophy of time alone. It rests on physics, and in particular on the theorems relating causal order to the geometry of spacetime. If a future physics replaced spacetime by a structure without a causal order, Principle~\ref{p:constitution} would require that the B-series be identified with whatever ordering of events that physics retains; we have not considered whether every such physics retains one.

The argument does not decide whether time requires B-change. That is the question of TSM-17, which answers it negatively. It does not decide whether the world is indeterministic. The branching framework is a framework for indeterminism, and Section~\ref{s:branching} shows only that a B-theorist can accept indeterminism and the open future without the A-series. If the world is deterministic, the frame is a chain and every proposition is settled at every moment, as the entries T$n$.1 show.

Finally, the computation verifies the results on finite frames only. The proofs are general, and the module formalizes Lemma~\ref{l:pastclosed} and Theorem~\ref{t:monotone} for arbitrary branching frames, finite or infinite; but the counts in the tables concern the finite frames up to the stated bound, and the symmetry facts of Propositions~\ref{p:select} and~\ref{p:trl} are illustrated, not exhausted, by them.

\section{Formalization}\label{s:formal}

The module \texttt{lean/Philosophy/TSM/TSM\_01.lean} formalizes the logical skeleton of the argument in core Lean~4. Part I treats McTaggart's argument and its inversion. Worlds, their moments, absolute presentness and the reality of time are primitives. The formal core of Assumption~\ref{a:incoherence} is the axiom \texttt{present\_unique}; Definition~\ref{d:changing} is the definition \texttt{ChangingASeries}; the bridge of Definition~\ref{d:bridge} is \texttt{Bridge}; Assumption~\ref{a:real} is the axiom \texttt{time\_real\_actual}. The main results are these.

\begin{lstlisting}
/-- Theorem 4.4. No world contains a changing A-series. -/
theorem no_changing_A_series (w : World) : ¬ ChangingASeries w

/-- Theorem 5.3 (the inversion). -/
theorem bridge_fails : ¬ Bridge
\end{lstlisting}

Part II formalizes branching frames (Definition~\ref{d:frame}), histories as maximal chains (Definition~\ref{d:history}) and settledness (Definition~\ref{d:settled}) for arbitrary frames, proves Lemma~\ref{l:pastclosed}, Theorem~\ref{t:monotone} and the converse direction of Proposition~\ref{p:converse}, proves that settledness is not bivalent by exhibiting the fork T3.2 with an open future contingent at its root (Proposition~\ref{p:bivalence}), and proves clauses (a) and (b) of Proposition~\ref{p:stage} for the moment-relative A-assignment of Definition~\ref{d:relative}.

\begin{lstlisting}
/-- Theorem 6.6 (monotonicity of settledness). -/
theorem settled_monotone (phi : History F → Prop) {m m' : F.M}
    (hlt : F.lt m m') :
    (SettledTrue phi m → SettledTrue phi m') ∧
    (SettledFalse phi m → SettledFalse phi m')
\end{lstlisting}

The checker reports the following dependencies.
\begin{lstlisting}
'Philosophy.TSM_01.no_changing_A_series' depends on axioms:
  [Philosophy.TSM_01.Mom, Philosophy.TSM_01.Present,
   Philosophy.TSM_01.World, Philosophy.TSM_01.present_unique]
'Philosophy.TSM_01.bridge_fails' depends on axioms:
  [Philosophy.TSM_01.Mom, Philosophy.TSM_01.Present,
   Philosophy.TSM_01.TimeReal, Philosophy.TSM_01.World,
   Philosophy.TSM_01.actual, Philosophy.TSM_01.present_unique,
   Philosophy.TSM_01.time_real_actual]
'Philosophy.TSM_01.settled_monotone' does not depend on any axioms
'Philosophy.TSM_01.fixity_requires_unique_past' does not depend
  on any axioms
'Philosophy.TSM_01.fork_open' depends on axioms: [propext]
'Philosophy.TSM_01.stage_order' depends on axioms:
  [propext, Classical.choice, Quot.sound]
\end{lstlisting}

\texttt{World}, \texttt{Mom}, \texttt{Present}, \texttt{TimeReal} and \texttt{actual} are primitives. The two substantive axioms are \texttt{present\_unique}, the formal core of the incoherence premise, and \texttt{time\_real\_actual}, the reality of time; the inversion uses both, and Theorem~\ref{t:nochange} uses only the first. The monotonicity theorem and its converse depend on no axioms at all: it is a theorem about branching frames, and it holds whatever one thinks about the A-series. The remaining dependencies are the standard axioms of Lean's logic. The module does not establish the premises, does not formalize Principle~\ref{p:physical}, and does not reproduce the enumeration, which is carried out by the generator \texttt{gen/generate.py} distributed with the manuscript.

\section{Conclusion}\label{s:conclusion}

McTaggart's argument is valid, and its incoherence premise is true: an absolute present cannot change without being had by more than one moment, and a present relativized to moments is the B-series under another name. Its bridge premise is false, because it is credible only by equivocating between the change of things across times and the change of A-determinations. Run in reverse from the reality of time, the argument refutes the bridge and leaves time to the B-series, which in our world is the causal and metrical structure of spacetime. In an indeterministic world the B-series branches, settledness is a relation between moments and propositions, and what is settled at a moment stays settled at every later one; the open future and the fixed past are features of the shape of the series, not of a moving present, and the growing block is excluded with the A-series. Time is real, and it is the B-series.

\clearpage
\appendix

\section{Volume of frames: settledness at every moment}\label{app:detail}

\paragraph{What is enumerated.} This appendix lists every finite branching frame with at most ten moments, one representative of each isomorphism class, \NFramesDetail{} frames and \NMoments{} moments in all. A frame is a finite rooted tree of moments under the relation earlier-than (Definition~\ref{d:frame}); its histories are the maximal chains (Definition~\ref{d:history}), one for each maximal moment.

\paragraph{Encoding.} Each frame is given by its canonical level sequence: the moments are listed in preorder, the subtrees of each moment arranged in decreasing order of their own canonical sequences, and the $i$-th digit is the number of moments earlier than the $i$-th moment (the digits $a$ and $b$, used in Appendix~\ref{app:eleven}, stand for 10 and 11). Two frames are isomorphic if and only if their canonical sequences coincide, so each sequence names exactly one class. Moments are numbered $0, \dots, n-1$ in the order of the sequence; moment 0 is the least moment, and the predecessor of each moment has a smaller number. The frame identifier T$n$.$k$ gives the number of moments and the position of the frame in the list for $n$, the frames being listed in decreasing lexicographic order of their sequences, so that T$n$.1 is the chain and the last frame for each $n$ is the star.

\paragraph{The program.} The generator \texttt{gen/generate.py} (Python~3, standard library only) produces the frames by adding a maximal moment in every possible position to every frame with one moment fewer, canonicalizing, and removing duplicates. Exhaustiveness rests on the fact that removing any maximal moment from a frame with $n$ moments leaves a frame with $n-1$. The counts agree with Otter's numbers, and the program confirms them independently for $n \leq 10$ by canonicalizing all \NLabelled{} increasingly labeled trees (Section~\ref{s:computation}). Settledness is computed for all $2^n$ valuations simultaneously, by representing, for each moment $m$ and history $h$ through it, the set of valuations under which $F_m p$ is true relative to $h$ as a bit string of length $2^n$. For $n \leq 6$ the program recomputes every settledness count by direct evaluation, valuation by valuation, and checks agreement. The output is deterministic, and the program runs in a few seconds.

\paragraph{Reading an entry.} Each frame has a header line:
the identifier; the level sequence; $n$, the number of moments; $|H|$, the number of histories; \emph{len}, the number of moments on a longest history; \emph{br}, the number of moments with more than one immediate successor; $|\mathrm{Aut}|$, the order of the automorphism group; \emph{orb} $a/b$, the number $a$ of orbits of the group on moments and the number $b$ of moments fixed by every automorphism; \emph{hist-orb} $c/d$, the same for histories; \emph{M}, the number of instances of Theorem~\ref{t:monotone} checked, that is, the number of pairs $m < m'$ times $2^n$; \emph{res}, the number of resolutions, triples of a valuation and moments $m < m'$ such that $F_m p$ is open at $m$ and settled at $m'$; \emph{rev}, the number of reversals of Proposition~\ref{p:reversal}, triples such that $F_m p$ is settled true at $m$ and $F_{m'} p$ settled false at $m'$; and the verdict \enquote{ok}, which means that checks (i)--(iv) of Section~\ref{s:computation} hold for the frame.

One row follows for each moment $m$: its number; \emph{par}, its immediate predecessor; \emph{lvl}, the number of moments earlier than it; $h(m)$, the number of histories through it; \emph{orb}, the size of its orbit under the automorphism group (1 if it is singled out by the B-facts in the sense of Proposition~\ref{p:select}(b)); \emph{st}, \emph{sf} and \emph{op}, the numbers of valuations, out of $2^n$, under which $F_m p$ is settled true, settled false and open at $m$; and the \emph{single-point vector}, whose $x$-th character gives the status at $m$ of $F_m p$ when $p$ holds at moment $x$ alone: \texttt{+} settled true, \texttt{-} settled false, \texttt{?} open.

Several identities can be checked by eye. In every row $\mathit{st} + \mathit{sf} + \mathit{op} = 2^n$. The value of \emph{sf} is $2^{n-u}$, where $u$ is the number of moments later than $m$, since $F_m p$ is settled false at $m$ just if $p$ holds at no later moment. The value of \emph{op} is $0$ exactly when $h(m) = 1$ (Proposition~\ref{p:bivalence}). The single-point vector has \texttt{-} at $x$ unless $x$ is later than $m$, \texttt{+} if $x$ is later than $m$ and lies on every history through $m$, and \texttt{?} otherwise; along every history, a \texttt{+} at $x$ in the row of $m$ remains a \texttt{+} in the row of every later moment that is earlier than $x$ (Theorem~\ref{t:monotone}). The single-point vectors of maximal moments consist of \texttt{-} only.

\paragraph{Summary.} All checks hold in all \NFramesDetail{} frames. The cases of philosophical interest are discussed in Section~\ref{s:computation}: T1.1, the only frame in which one A-assignment makes every moment present; T2.1 and the chains, in which nothing is open and every variation is a reversal; T3.2, the fork, the smallest frame with an open future contingent and with no history singled out by the B-facts; T4.3, the smallest rigid branching frame; T7.42, in which the root is the only moment singled out by the B-facts; and T10.719, the star.

% generated by gen/generate.py -- do not edit
{\scriptsize\setlength{\tabcolsep}{4pt}
\begin{longtable}{@{}rrrrrrrrl@{}}
\toprule
$m$ & par & lvl & $h(m)$ & orb & st & sf & op & single-point vector $x=0,\dots,n-1$\\
\midrule
\endfirsthead
\toprule
$m$ & par & lvl & $h(m)$ & orb & st & sf & op & single-point vector $x=0,\dots,n-1$\\
\midrule
\endhead
\multicolumn{9}{@{}l@{}}{\rule{0pt}{3.4ex}\normalsize\bfseries Frames with 1 moment (1 classes)}\\*
\multicolumn{9}{@{}l@{}}{\rule{0pt}{2.6ex}\textbf{T1.1}\enspace \texttt{0}\enspace $n{=}1$, $|H|{=}1$, len 1, br 0, $|\mathrm{Aut}|{=}1$, orb 1/1, hist-orb 1/1; M 0, res 0, rev 0: ok}\\*
0 & -- & 0 & 1 & 1 & 0 & 2 & 0 & \texttt{-}\\
\multicolumn{9}{@{}l@{}}{\rule{0pt}{3.4ex}\normalsize\bfseries Frames with 2 moments (1 classes)}\\*
\multicolumn{9}{@{}l@{}}{\rule{0pt}{2.6ex}\textbf{T2.1}\enspace \texttt{01}\enspace $n{=}2$, $|H|{=}1$, len 2, br 0, $|\mathrm{Aut}|{=}1$, orb 2/2, hist-orb 1/1; M 4, res 0, rev 2: ok}\\*
0 & -- & 0 & 1 & 1 & 2 & 2 & 0 & \texttt{-+}\\*
1 & 0 & 1 & 1 & 1 & 0 & 4 & 0 & \texttt{--}\\
\multicolumn{9}{@{}l@{}}{\rule{0pt}{3.4ex}\normalsize\bfseries Frames with 3 moments (2 classes)}\\*
\multicolumn{9}{@{}l@{}}{\rule{0pt}{2.6ex}\textbf{T3.1}\enspace \texttt{012}\enspace $n{=}3$, $|H|{=}1$, len 3, br 0, $|\mathrm{Aut}|{=}1$, orb 3/3, hist-orb 1/1; M 24, res 0, rev 12: ok}\\*
0 & -- & 0 & 1 & 1 & 6 & 2 & 0 & \texttt{-++}\\*
1 & 0 & 1 & 1 & 1 & 4 & 4 & 0 & \texttt{--+}\\
2 & 1 & 2 & 1 & 1 & 0 & 8 & 0 & \texttt{---}\\
\multicolumn{9}{@{}l@{}}{\rule{0pt}{2.6ex}\textbf{T3.2}\enspace \texttt{011}\enspace $n{=}3$, $|H|{=}2$, len 2, br 1, $|\mathrm{Aut}|{=}2$, orb 2/1, hist-orb 1/0; M 16, res 8, rev 4: ok}\\*
0 & -- & 0 & 2 & 1 & 2 & 2 & 4 & \texttt{-??}\\*
1 & 0 & 1 & 1 & 2 & 0 & 8 & 0 & \texttt{---}\\
2 & 0 & 1 & 1 & 2 & 0 & 8 & 0 & \texttt{---}\\
\multicolumn{9}{@{}l@{}}{\rule{0pt}{3.4ex}\normalsize\bfseries Frames with 4 moments (4 classes)}\\*
\multicolumn{9}{@{}l@{}}{\rule{0pt}{2.6ex}\textbf{T4.1}\enspace \texttt{0123}\enspace $n{=}4$, $|H|{=}1$, len 4, br 0, $|\mathrm{Aut}|{=}1$, orb 4/4, hist-orb 1/1; M 96, res 0, rev 46: ok}\\*
0 & -- & 0 & 1 & 1 & 14 & 2 & 0 & \texttt{-+++}\\*
1 & 0 & 1 & 1 & 1 & 12 & 4 & 0 & \texttt{--++}\\
2 & 1 & 2 & 1 & 1 & 8 & 8 & 0 & \texttt{---+}\\
3 & 2 & 3 & 1 & 1 & 0 & 16 & 0 & \texttt{----}\\
\multicolumn{9}{@{}l@{}}{\rule{0pt}{2.6ex}\textbf{T4.2}\enspace \texttt{0122}\enspace $n{=}4$, $|H|{=}2$, len 3, br 1, $|\mathrm{Aut}|{=}2$, orb 3/2, hist-orb 1/0; M 80, res 24, rev 30: ok}\\*
0 & -- & 0 & 2 & 1 & 10 & 2 & 4 & \texttt{-+??}\\*
1 & 0 & 1 & 2 & 1 & 4 & 4 & 8 & \texttt{--??}\\
2 & 1 & 2 & 1 & 2 & 0 & 16 & 0 & \texttt{----}\\
3 & 1 & 2 & 1 & 2 & 0 & 16 & 0 & \texttt{----}\\
\multicolumn{9}{@{}l@{}}{\rule{0pt}{2.6ex}\textbf{T4.3}\enspace \texttt{0121}\enspace $n{=}4$, $|H|{=}2$, len 3, br 1, $|\mathrm{Aut}|{=}1$, orb 4/4, hist-orb 2/2; M 64, res 24, rev 22: ok}\\*
0 & -- & 0 & 2 & 1 & 6 & 2 & 8 & \texttt{-???}\\*
1 & 0 & 1 & 1 & 1 & 8 & 8 & 0 & \texttt{--+-}\\
2 & 1 & 2 & 1 & 1 & 0 & 16 & 0 & \texttt{----}\\
3 & 0 & 1 & 1 & 1 & 0 & 16 & 0 & \texttt{----}\\
\multicolumn{9}{@{}l@{}}{\rule{0pt}{2.6ex}\textbf{T4.4}\enspace \texttt{0111}\enspace $n{=}4$, $|H|{=}3$, len 2, br 1, $|\mathrm{Aut}|{=}6$, orb 2/1, hist-orb 1/0; M 48, res 36, rev 6: ok}\\*
0 & -- & 0 & 3 & 1 & 2 & 2 & 12 & \texttt{-???}\\*
1 & 0 & 1 & 1 & 3 & 0 & 16 & 0 & \texttt{----}\\
2 & 0 & 1 & 1 & 3 & 0 & 16 & 0 & \texttt{----}\\
3 & 0 & 1 & 1 & 3 & 0 & 16 & 0 & \texttt{----}\\
\multicolumn{9}{@{}l@{}}{\rule{0pt}{3.4ex}\normalsize\bfseries Frames with 5 moments (9 classes)}\\*
\multicolumn{9}{@{}l@{}}{\rule{0pt}{2.6ex}\textbf{T5.1}\enspace \texttt{01234}\enspace $n{=}5$, $|H|{=}1$, len 5, br 0, $|\mathrm{Aut}|{=}1$, orb 5/5, hist-orb 1/1; M 320, res 0, rev 144: ok}\\*
0 & -- & 0 & 1 & 1 & 30 & 2 & 0 & \texttt{-++++}\\*
1 & 0 & 1 & 1 & 1 & 28 & 4 & 0 & \texttt{--+++}\\
2 & 1 & 2 & 1 & 1 & 24 & 8 & 0 & \texttt{---++}\\
3 & 2 & 3 & 1 & 1 & 16 & 16 & 0 & \texttt{----+}\\
4 & 3 & 4 & 1 & 1 & 0 & 32 & 0 & \texttt{-----}\\
\multicolumn{9}{@{}l@{}}{\rule{0pt}{2.6ex}\textbf{T5.2}\enspace \texttt{01233}\enspace $n{=}5$, $|H|{=}2$, len 4, br 1, $|\mathrm{Aut}|{=}2$, orb 4/3, hist-orb 1/0; M 288, res 56, rev 120: ok}\\*
0 & -- & 0 & 2 & 1 & 26 & 2 & 4 & \texttt{-++??}\\*
1 & 0 & 1 & 2 & 1 & 20 & 4 & 8 & \texttt{--+??}\\
2 & 1 & 2 & 2 & 1 & 8 & 8 & 16 & \texttt{---??}\\
3 & 2 & 3 & 1 & 2 & 0 & 32 & 0 & \texttt{-----}\\
4 & 2 & 3 & 1 & 2 & 0 & 32 & 0 & \texttt{-----}\\
\multicolumn{9}{@{}l@{}}{\rule{0pt}{2.6ex}\textbf{T5.3}\enspace \texttt{01232}\enspace $n{=}5$, $|H|{=}2$, len 4, br 1, $|\mathrm{Aut}|{=}1$, orb 5/5, hist-orb 2/2; M 256, res 72, rev 100: ok}\\*
0 & -- & 0 & 2 & 1 & 22 & 2 & 8 & \texttt{-+???}\\*
1 & 0 & 1 & 2 & 1 & 12 & 4 & 16 & \texttt{--???}\\
2 & 1 & 2 & 1 & 1 & 16 & 16 & 0 & \texttt{---+-}\\
3 & 2 & 3 & 1 & 1 & 0 & 32 & 0 & \texttt{-----}\\
4 & 1 & 2 & 1 & 1 & 0 & 32 & 0 & \texttt{-----}\\
\multicolumn{9}{@{}l@{}}{\rule{0pt}{2.6ex}\textbf{T5.4}\enspace \texttt{01231}\enspace $n{=}5$, $|H|{=}2$, len 4, br 1, $|\mathrm{Aut}|{=}1$, orb 5/5, hist-orb 2/2; M 224, res 64, rev 84: ok}\\*
0 & -- & 0 & 2 & 1 & 14 & 2 & 16 & \texttt{-????}\\*
1 & 0 & 1 & 1 & 1 & 24 & 8 & 0 & \texttt{--++-}\\
2 & 1 & 2 & 1 & 1 & 16 & 16 & 0 & \texttt{---+-}\\
3 & 2 & 3 & 1 & 1 & 0 & 32 & 0 & \texttt{-----}\\
4 & 0 & 1 & 1 & 1 & 0 & 32 & 0 & \texttt{-----}\\
\multicolumn{9}{@{}l@{}}{\rule{0pt}{2.6ex}\textbf{T5.5}\enspace \texttt{01222}\enspace $n{=}5$, $|H|{=}3$, len 3, br 1, $|\mathrm{Aut}|{=}6$, orb 3/2, hist-orb 1/0; M 224, res 108, rev 68: ok}\\*
0 & -- & 0 & 3 & 1 & 18 & 2 & 12 & \texttt{-+???}\\*
1 & 0 & 1 & 3 & 1 & 4 & 4 & 24 & \texttt{--???}\\
2 & 1 & 2 & 1 & 3 & 0 & 32 & 0 & \texttt{-----}\\
3 & 1 & 2 & 1 & 3 & 0 & 32 & 0 & \texttt{-----}\\
4 & 1 & 2 & 1 & 3 & 0 & 32 & 0 & \texttt{-----}\\
\multicolumn{9}{@{}l@{}}{\rule{0pt}{2.6ex}\textbf{T5.6}\enspace \texttt{01221}\enspace $n{=}5$, $|H|{=}3$, len 3, br 2, $|\mathrm{Aut}|{=}2$, orb 4/3, hist-orb 2/1; M 192, res 104, rev 48: ok}\\*
0 & -- & 0 & 3 & 1 & 10 & 2 & 20 & \texttt{-????}\\*
1 & 0 & 1 & 2 & 1 & 8 & 8 & 16 & \texttt{--??-}\\
2 & 1 & 2 & 1 & 2 & 0 & 32 & 0 & \texttt{-----}\\
3 & 1 & 2 & 1 & 2 & 0 & 32 & 0 & \texttt{-----}\\
4 & 0 & 1 & 1 & 1 & 0 & 32 & 0 & \texttt{-----}\\
\multicolumn{9}{@{}l@{}}{\rule{0pt}{2.6ex}\textbf{T5.7}\enspace \texttt{01212}\enspace $n{=}5$, $|H|{=}2$, len 3, br 1, $|\mathrm{Aut}|{=}2$, orb 3/1, hist-orb 1/0; M 192, res 48, rev 80: ok}\\*
0 & -- & 0 & 2 & 1 & 18 & 2 & 12 & \texttt{-????}\\*
1 & 0 & 1 & 1 & 2 & 16 & 16 & 0 & \texttt{--+--}\\
2 & 1 & 2 & 1 & 2 & 0 & 32 & 0 & \texttt{-----}\\
3 & 0 & 1 & 1 & 2 & 16 & 16 & 0 & \texttt{----+}\\
4 & 3 & 2 & 1 & 2 & 0 & 32 & 0 & \texttt{-----}\\
\multicolumn{9}{@{}l@{}}{\rule{0pt}{2.6ex}\textbf{T5.8}\enspace \texttt{01211}\enspace $n{=}5$, $|H|{=}3$, len 3, br 1, $|\mathrm{Aut}|{=}2$, orb 4/3, hist-orb 2/1; M 160, res 96, rev 36: ok}\\*
0 & -- & 0 & 3 & 1 & 6 & 2 & 24 & \texttt{-????}\\*
1 & 0 & 1 & 1 & 1 & 16 & 16 & 0 & \texttt{--+--}\\
2 & 1 & 2 & 1 & 1 & 0 & 32 & 0 & \texttt{-----}\\
3 & 0 & 1 & 1 & 2 & 0 & 32 & 0 & \texttt{-----}\\
4 & 0 & 1 & 1 & 2 & 0 & 32 & 0 & \texttt{-----}\\
\multicolumn{9}{@{}l@{}}{\rule{0pt}{2.6ex}\textbf{T5.9}\enspace \texttt{01111}\enspace $n{=}5$, $|H|{=}4$, len 2, br 1, $|\mathrm{Aut}|{=}24$, orb 2/1, hist-orb 1/0; M 128, res 112, rev 8: ok}\\*
0 & -- & 0 & 4 & 1 & 2 & 2 & 28 & \texttt{-????}\\*
1 & 0 & 1 & 1 & 4 & 0 & 32 & 0 & \texttt{-----}\\
2 & 0 & 1 & 1 & 4 & 0 & 32 & 0 & \texttt{-----}\\
3 & 0 & 1 & 1 & 4 & 0 & 32 & 0 & \texttt{-----}\\
4 & 0 & 1 & 1 & 4 & 0 & 32 & 0 & \texttt{-----}\\
\multicolumn{9}{@{}l@{}}{\rule{0pt}{3.4ex}\normalsize\bfseries Frames with 6 moments (20 classes)}\\*
\multicolumn{9}{@{}l@{}}{\rule{0pt}{2.6ex}\textbf{T6.01}\enspace \texttt{012345}\enspace $n{=}6$, $|H|{=}1$, len 6, br 0, $|\mathrm{Aut}|{=}1$, orb 6/6, hist-orb 1/1; M 960, res 0, rev 402: ok}\\*
0 & -- & 0 & 1 & 1 & 62 & 2 & 0 & \texttt{-+++++}\\*
1 & 0 & 1 & 1 & 1 & 60 & 4 & 0 & \texttt{--++++}\\
2 & 1 & 2 & 1 & 1 & 56 & 8 & 0 & \texttt{---+++}\\
3 & 2 & 3 & 1 & 1 & 48 & 16 & 0 & \texttt{----++}\\
4 & 3 & 4 & 1 & 1 & 32 & 32 & 0 & \texttt{-----+}\\
5 & 4 & 5 & 1 & 1 & 0 & 64 & 0 & \texttt{------}\\
\multicolumn{9}{@{}l@{}}{\rule{0pt}{2.6ex}\textbf{T6.02}\enspace \texttt{012344}\enspace $n{=}6$, $|H|{=}2$, len 5, br 1, $|\mathrm{Aut}|{=}2$, orb 5/4, hist-orb 1/0; M 896, res 120, rev 378: ok}\\*
0 & -- & 0 & 2 & 1 & 58 & 2 & 4 & \texttt{-+++??}\\*
1 & 0 & 1 & 2 & 1 & 52 & 4 & 8 & \texttt{--++??}\\
2 & 1 & 2 & 2 & 1 & 40 & 8 & 16 & \texttt{---+??}\\
3 & 2 & 3 & 2 & 1 & 16 & 16 & 32 & \texttt{----??}\\
4 & 3 & 4 & 1 & 2 & 0 & 64 & 0 & \texttt{------}\\
5 & 3 & 4 & 1 & 2 & 0 & 64 & 0 & \texttt{------}\\
\multicolumn{9}{@{}l@{}}{\rule{0pt}{2.6ex}\textbf{T6.03}\enspace \texttt{012343}\enspace $n{=}6$, $|H|{=}2$, len 5, br 1, $|\mathrm{Aut}|{=}1$, orb 6/6, hist-orb 2/2; M 832, res 168, rev 342: ok}\\*
0 & -- & 0 & 2 & 1 & 54 & 2 & 8 & \texttt{-++???}\\*
1 & 0 & 1 & 2 & 1 & 44 & 4 & 16 & \texttt{--+???}\\
2 & 1 & 2 & 2 & 1 & 24 & 8 & 32 & \texttt{---???}\\
3 & 2 & 3 & 1 & 1 & 32 & 32 & 0 & \texttt{----+-}\\
4 & 3 & 4 & 1 & 1 & 0 & 64 & 0 & \texttt{------}\\
5 & 2 & 3 & 1 & 1 & 0 & 64 & 0 & \texttt{------}\\
\multicolumn{9}{@{}l@{}}{\rule{0pt}{2.6ex}\textbf{T6.04}\enspace \texttt{012342}\enspace $n{=}6$, $|H|{=}2$, len 5, br 1, $|\mathrm{Aut}|{=}1$, orb 6/6, hist-orb 2/2; M 768, res 192, rev 294: ok}\\*
0 & -- & 0 & 2 & 1 & 46 & 2 & 16 & \texttt{-+????}\\*
1 & 0 & 1 & 2 & 1 & 28 & 4 & 32 & \texttt{--????}\\
2 & 1 & 2 & 1 & 1 & 48 & 16 & 0 & \texttt{---++-}\\
3 & 2 & 3 & 1 & 1 & 32 & 32 & 0 & \texttt{----+-}\\
4 & 3 & 4 & 1 & 1 & 0 & 64 & 0 & \texttt{------}\\
5 & 1 & 2 & 1 & 1 & 0 & 64 & 0 & \texttt{------}\\
\multicolumn{9}{@{}l@{}}{\rule{0pt}{2.6ex}\textbf{T6.05}\enspace \texttt{012341}\enspace $n{=}6$, $|H|{=}2$, len 5, br 1, $|\mathrm{Aut}|{=}1$, orb 6/6, hist-orb 2/2; M 704, res 160, rev 266: ok}\\*
0 & -- & 0 & 2 & 1 & 30 & 2 & 32 & \texttt{-?????}\\*
1 & 0 & 1 & 1 & 1 & 56 & 8 & 0 & \texttt{--+++-}\\
2 & 1 & 2 & 1 & 1 & 48 & 16 & 0 & \texttt{---++-}\\
3 & 2 & 3 & 1 & 1 & 32 & 32 & 0 & \texttt{----+-}\\
4 & 3 & 4 & 1 & 1 & 0 & 64 & 0 & \texttt{------}\\
5 & 0 & 1 & 1 & 1 & 0 & 64 & 0 & \texttt{------}\\
\multicolumn{9}{@{}l@{}}{\rule{0pt}{2.6ex}\textbf{T6.06}\enspace \texttt{012333}\enspace $n{=}6$, $|H|{=}3$, len 4, br 1, $|\mathrm{Aut}|{=}6$, orb 4/3, hist-orb 1/0; M 768, res 252, rev 294: ok}\\*
0 & -- & 0 & 3 & 1 & 50 & 2 & 12 & \texttt{-++???}\\*
1 & 0 & 1 & 3 & 1 & 36 & 4 & 24 & \texttt{--+???}\\
2 & 1 & 2 & 3 & 1 & 8 & 8 & 48 & \texttt{---???}\\
3 & 2 & 3 & 1 & 3 & 0 & 64 & 0 & \texttt{------}\\
4 & 2 & 3 & 1 & 3 & 0 & 64 & 0 & \texttt{------}\\
5 & 2 & 3 & 1 & 3 & 0 & 64 & 0 & \texttt{------}\\
\multicolumn{9}{@{}l@{}}{\rule{0pt}{2.6ex}\textbf{T6.07}\enspace \texttt{012332}\enspace $n{=}6$, $|H|{=}3$, len 4, br 2, $|\mathrm{Aut}|{=}2$, orb 5/4, hist-orb 2/1; M 704, res 280, rev 234: ok}\\*
0 & -- & 0 & 3 & 1 & 42 & 2 & 20 & \texttt{-+????}\\*
1 & 0 & 1 & 3 & 1 & 20 & 4 & 40 & \texttt{--????}\\
2 & 1 & 2 & 2 & 1 & 16 & 16 & 32 & \texttt{---??-}\\
3 & 2 & 3 & 1 & 2 & 0 & 64 & 0 & \texttt{------}\\
4 & 2 & 3 & 1 & 2 & 0 & 64 & 0 & \texttt{------}\\
5 & 1 & 2 & 1 & 1 & 0 & 64 & 0 & \texttt{------}\\
\multicolumn{9}{@{}l@{}}{\rule{0pt}{2.6ex}\textbf{T6.08}\enspace \texttt{012331}\enspace $n{=}6$, $|H|{=}3$, len 4, br 2, $|\mathrm{Aut}|{=}2$, orb 5/4, hist-orb 2/1; M 640, res 260, rev 206: ok}\\*
0 & -- & 0 & 3 & 1 & 26 & 2 & 36 & \texttt{-?????}\\*
1 & 0 & 1 & 2 & 1 & 40 & 8 & 16 & \texttt{--+??-}\\
2 & 1 & 2 & 2 & 1 & 16 & 16 & 32 & \texttt{---??-}\\
3 & 2 & 3 & 1 & 2 & 0 & 64 & 0 & \texttt{------}\\
4 & 2 & 3 & 1 & 2 & 0 & 64 & 0 & \texttt{------}\\
5 & 0 & 1 & 1 & 1 & 0 & 64 & 0 & \texttt{------}\\
\multicolumn{9}{@{}l@{}}{\rule{0pt}{2.6ex}\textbf{T6.09}\enspace \texttt{012323}\enspace $n{=}6$, $|H|{=}2$, len 4, br 1, $|\mathrm{Aut}|{=}2$, orb 4/2, hist-orb 1/0; M 704, res 144, rev 306: ok}\\*
0 & -- & 0 & 2 & 1 & 50 & 2 & 12 & \texttt{-+????}\\*
1 & 0 & 1 & 2 & 1 & 36 & 4 & 24 & \texttt{--????}\\
2 & 1 & 2 & 1 & 2 & 32 & 32 & 0 & \texttt{---+--}\\
3 & 2 & 3 & 1 & 2 & 0 & 64 & 0 & \texttt{------}\\
4 & 1 & 2 & 1 & 2 & 32 & 32 & 0 & \texttt{-----+}\\
5 & 4 & 3 & 1 & 2 & 0 & 64 & 0 & \texttt{------}\\
\multicolumn{9}{@{}l@{}}{\rule{0pt}{2.6ex}\textbf{T6.10}\enspace \texttt{012322}\enspace $n{=}6$, $|H|{=}3$, len 4, br 1, $|\mathrm{Aut}|{=}2$, orb 5/4, hist-orb 2/1; M 640, res 288, rev 206: ok}\\*
0 & -- & 0 & 3 & 1 & 38 & 2 & 24 & \texttt{-+????}\\*
1 & 0 & 1 & 3 & 1 & 12 & 4 & 48 & \texttt{--????}\\
2 & 1 & 2 & 1 & 1 & 32 & 32 & 0 & \texttt{---+--}\\
3 & 2 & 3 & 1 & 1 & 0 & 64 & 0 & \texttt{------}\\
4 & 1 & 2 & 1 & 2 & 0 & 64 & 0 & \texttt{------}\\
5 & 1 & 2 & 1 & 2 & 0 & 64 & 0 & \texttt{------}\\
\multicolumn{9}{@{}l@{}}{\rule{0pt}{2.6ex}\textbf{T6.11}\enspace \texttt{012321}\enspace $n{=}6$, $|H|{=}3$, len 4, br 2, $|\mathrm{Aut}|{=}1$, orb 6/6, hist-orb 3/3; M 576, res 280, rev 166: ok}\\*
0 & -- & 0 & 3 & 1 & 22 & 2 & 40 & \texttt{-?????}\\*
1 & 0 & 1 & 2 & 1 & 24 & 8 & 32 & \texttt{--???-}\\
2 & 1 & 2 & 1 & 1 & 32 & 32 & 0 & \texttt{---+--}\\
3 & 2 & 3 & 1 & 1 & 0 & 64 & 0 & \texttt{------}\\
4 & 1 & 2 & 1 & 1 & 0 & 64 & 0 & \texttt{------}\\
5 & 0 & 1 & 1 & 1 & 0 & 64 & 0 & \texttt{------}\\
\multicolumn{9}{@{}l@{}}{\rule{0pt}{2.6ex}\textbf{T6.12}\enspace \texttt{012312}\enspace $n{=}6$, $|H|{=}2$, len 4, br 1, $|\mathrm{Aut}|{=}1$, orb 6/6, hist-orb 2/2; M 576, res 100, rev 250: ok}\\*
0 & -- & 0 & 2 & 1 & 42 & 2 & 20 & \texttt{-?????}\\*
1 & 0 & 1 & 1 & 1 & 48 & 16 & 0 & \texttt{--++--}\\
2 & 1 & 2 & 1 & 1 & 32 & 32 & 0 & \texttt{---+--}\\
3 & 2 & 3 & 1 & 1 & 0 & 64 & 0 & \texttt{------}\\
4 & 0 & 1 & 1 & 1 & 32 & 32 & 0 & \texttt{-----+}\\
5 & 4 & 2 & 1 & 1 & 0 & 64 & 0 & \texttt{------}\\
\multicolumn{9}{@{}l@{}}{\rule{0pt}{2.6ex}\textbf{T6.13}\enspace \texttt{012311}\enspace $n{=}6$, $|H|{=}3$, len 4, br 1, $|\mathrm{Aut}|{=}2$, orb 5/4, hist-orb 2/1; M 512, res 240, rev 146: ok}\\*
0 & -- & 0 & 3 & 1 & 14 & 2 & 48 & \texttt{-?????}\\*
1 & 0 & 1 & 1 & 1 & 48 & 16 & 0 & \texttt{--++--}\\
2 & 1 & 2 & 1 & 1 & 32 & 32 & 0 & \texttt{---+--}\\
3 & 2 & 3 & 1 & 1 & 0 & 64 & 0 & \texttt{------}\\
4 & 0 & 1 & 1 & 2 & 0 & 64 & 0 & \texttt{------}\\
5 & 0 & 1 & 1 & 2 & 0 & 64 & 0 & \texttt{------}\\
\multicolumn{9}{@{}l@{}}{\rule{0pt}{2.6ex}\textbf{T6.14}\enspace \texttt{012222}\enspace $n{=}6$, $|H|{=}4$, len 3, br 1, $|\mathrm{Aut}|{=}24$, orb 3/2, hist-orb 1/0; M 576, res 336, rev 154: ok}\\*
0 & -- & 0 & 4 & 1 & 34 & 2 & 28 & \texttt{-+????}\\*
1 & 0 & 1 & 4 & 1 & 4 & 4 & 56 & \texttt{--????}\\
2 & 1 & 2 & 1 & 4 & 0 & 64 & 0 & \texttt{------}\\
3 & 1 & 2 & 1 & 4 & 0 & 64 & 0 & \texttt{------}\\
4 & 1 & 2 & 1 & 4 & 0 & 64 & 0 & \texttt{------}\\
5 & 1 & 2 & 1 & 4 & 0 & 64 & 0 & \texttt{------}\\
\multicolumn{9}{@{}l@{}}{\rule{0pt}{2.6ex}\textbf{T6.15}\enspace \texttt{012221}\enspace $n{=}6$, $|H|{=}4$, len 3, br 2, $|\mathrm{Aut}|{=}6$, orb 4/3, hist-orb 2/1; M 512, res 340, rev 98: ok}\\*
0 & -- & 0 & 4 & 1 & 18 & 2 & 44 & \texttt{-?????}\\*
1 & 0 & 1 & 3 & 1 & 8 & 8 & 48 & \texttt{--???-}\\
2 & 1 & 2 & 1 & 3 & 0 & 64 & 0 & \texttt{------}\\
3 & 1 & 2 & 1 & 3 & 0 & 64 & 0 & \texttt{------}\\
4 & 1 & 2 & 1 & 3 & 0 & 64 & 0 & \texttt{------}\\
5 & 0 & 1 & 1 & 1 & 0 & 64 & 0 & \texttt{------}\\
\multicolumn{9}{@{}l@{}}{\rule{0pt}{2.6ex}\textbf{T6.16}\enspace \texttt{012212}\enspace $n{=}6$, $|H|{=}3$, len 3, br 2, $|\mathrm{Aut}|{=}2$, orb 5/4, hist-orb 2/1; M 512, res 208, rev 170: ok}\\*
0 & -- & 0 & 3 & 1 & 30 & 2 & 32 & \texttt{-?????}\\*
1 & 0 & 1 & 2 & 1 & 16 & 16 & 32 & \texttt{--??--}\\
2 & 1 & 2 & 1 & 2 & 0 & 64 & 0 & \texttt{------}\\
3 & 1 & 2 & 1 & 2 & 0 & 64 & 0 & \texttt{------}\\
4 & 0 & 1 & 1 & 1 & 32 & 32 & 0 & \texttt{-----+}\\
5 & 4 & 2 & 1 & 1 & 0 & 64 & 0 & \texttt{------}\\
\multicolumn{9}{@{}l@{}}{\rule{0pt}{2.6ex}\textbf{T6.17}\enspace \texttt{012211}\enspace $n{=}6$, $|H|{=}4$, len 3, br 2, $|\mathrm{Aut}|{=}4$, orb 4/2, hist-orb 2/0; M 448, res 308, rev 74: ok}\\*
0 & -- & 0 & 4 & 1 & 10 & 2 & 52 & \texttt{-?????}\\*
1 & 0 & 1 & 2 & 1 & 16 & 16 & 32 & \texttt{--??--}\\
2 & 1 & 2 & 1 & 2 & 0 & 64 & 0 & \texttt{------}\\
3 & 1 & 2 & 1 & 2 & 0 & 64 & 0 & \texttt{------}\\
4 & 0 & 1 & 1 & 2 & 0 & 64 & 0 & \texttt{------}\\
5 & 0 & 1 & 1 & 2 & 0 & 64 & 0 & \texttt{------}\\
\multicolumn{9}{@{}l@{}}{\rule{0pt}{2.6ex}\textbf{T6.18}\enspace \texttt{012121}\enspace $n{=}6$, $|H|{=}3$, len 3, br 1, $|\mathrm{Aut}|{=}2$, orb 4/2, hist-orb 2/1; M 448, res 220, rev 130: ok}\\*
0 & -- & 0 & 3 & 1 & 18 & 2 & 44 & \texttt{-?????}\\*
1 & 0 & 1 & 1 & 2 & 32 & 32 & 0 & \texttt{--+---}\\
2 & 1 & 2 & 1 & 2 & 0 & 64 & 0 & \texttt{------}\\
3 & 0 & 1 & 1 & 2 & 32 & 32 & 0 & \texttt{----+-}\\
4 & 3 & 2 & 1 & 2 & 0 & 64 & 0 & \texttt{------}\\
5 & 0 & 1 & 1 & 1 & 0 & 64 & 0 & \texttt{------}\\
\multicolumn{9}{@{}l@{}}{\rule{0pt}{2.6ex}\textbf{T6.19}\enspace \texttt{012111}\enspace $n{=}6$, $|H|{=}4$, len 3, br 1, $|\mathrm{Aut}|{=}6$, orb 4/3, hist-orb 2/1; M 384, res 280, rev 58: ok}\\*
0 & -- & 0 & 4 & 1 & 6 & 2 & 56 & \texttt{-?????}\\*
1 & 0 & 1 & 1 & 1 & 32 & 32 & 0 & \texttt{--+---}\\
2 & 1 & 2 & 1 & 1 & 0 & 64 & 0 & \texttt{------}\\
3 & 0 & 1 & 1 & 3 & 0 & 64 & 0 & \texttt{------}\\
4 & 0 & 1 & 1 & 3 & 0 & 64 & 0 & \texttt{------}\\
5 & 0 & 1 & 1 & 3 & 0 & 64 & 0 & \texttt{------}\\
\multicolumn{9}{@{}l@{}}{\rule{0pt}{2.6ex}\textbf{T6.20}\enspace \texttt{011111}\enspace $n{=}6$, $|H|{=}5$, len 2, br 1, $|\mathrm{Aut}|{=}120$, orb 2/1, hist-orb 1/0; M 320, res 300, rev 10: ok}\\*
0 & -- & 0 & 5 & 1 & 2 & 2 & 60 & \texttt{-?????}\\*
1 & 0 & 1 & 1 & 5 & 0 & 64 & 0 & \texttt{------}\\
2 & 0 & 1 & 1 & 5 & 0 & 64 & 0 & \texttt{------}\\
3 & 0 & 1 & 1 & 5 & 0 & 64 & 0 & \texttt{------}\\
4 & 0 & 1 & 1 & 5 & 0 & 64 & 0 & \texttt{------}\\
5 & 0 & 1 & 1 & 5 & 0 & 64 & 0 & \texttt{------}\\
\multicolumn{9}{@{}l@{}}{\rule{0pt}{3.4ex}\normalsize\bfseries Frames with 7 moments (48 classes)}\\*
\multicolumn{9}{@{}l@{}}{\rule{0pt}{2.6ex}\textbf{T7.01}\enspace \texttt{0123456}\enspace $n{=}7$, $|H|{=}1$, len 7, br 0, $|\mathrm{Aut}|{=}1$, orb 7/7, hist-orb 1/1; M 2688, res 0, rev 1044: ok}\\*
0 & -- & 0 & 1 & 1 & 126 & 2 & 0 & \texttt{-++++++}\\*
1 & 0 & 1 & 1 & 1 & 124 & 4 & 0 & \texttt{--+++++}\\
2 & 1 & 2 & 1 & 1 & 120 & 8 & 0 & \texttt{---++++}\\
3 & 2 & 3 & 1 & 1 & 112 & 16 & 0 & \texttt{----+++}\\
4 & 3 & 4 & 1 & 1 & 96 & 32 & 0 & \texttt{-----++}\\
5 & 4 & 5 & 1 & 1 & 64 & 64 & 0 & \texttt{------+}\\
6 & 5 & 6 & 1 & 1 & 0 & 128 & 0 & \texttt{-------}\\
\multicolumn{9}{@{}l@{}}{\rule{0pt}{2.6ex}\textbf{T7.02}\enspace \texttt{0123455}\enspace $n{=}7$, $|H|{=}2$, len 6, br 1, $|\mathrm{Aut}|{=}2$, orb 6/5, hist-orb 1/0; M 2560, res 248, rev 1052: ok}\\*
0 & -- & 0 & 2 & 1 & 122 & 2 & 4 & \texttt{-++++??}\\*
1 & 0 & 1 & 2 & 1 & 116 & 4 & 8 & \texttt{--+++??}\\
2 & 1 & 2 & 2 & 1 & 104 & 8 & 16 & \texttt{---++??}\\
3 & 2 & 3 & 2 & 1 & 80 & 16 & 32 & \texttt{----+??}\\
4 & 3 & 4 & 2 & 1 & 32 & 32 & 64 & \texttt{-----??}\\
5 & 4 & 5 & 1 & 2 & 0 & 128 & 0 & \texttt{-------}\\
6 & 4 & 5 & 1 & 2 & 0 & 128 & 0 & \texttt{-------}\\
\multicolumn{9}{@{}l@{}}{\rule{0pt}{2.6ex}\textbf{T7.03}\enspace \texttt{0123454}\enspace $n{=}7$, $|H|{=}2$, len 6, br 1, $|\mathrm{Aut}|{=}1$, orb 7/7, hist-orb 2/2; M 2432, res 360, rev 1000: ok}\\*
0 & -- & 0 & 2 & 1 & 118 & 2 & 8 & \texttt{-+++???}\\*
1 & 0 & 1 & 2 & 1 & 108 & 4 & 16 & \texttt{--++???}\\
2 & 1 & 2 & 2 & 1 & 88 & 8 & 32 & \texttt{---+???}\\
3 & 2 & 3 & 2 & 1 & 48 & 16 & 64 & \texttt{----???}\\
4 & 3 & 4 & 1 & 1 & 64 & 64 & 0 & \texttt{-----+-}\\
5 & 4 & 5 & 1 & 1 & 0 & 128 & 0 & \texttt{-------}\\
6 & 3 & 4 & 1 & 1 & 0 & 128 & 0 & \texttt{-------}\\
\multicolumn{9}{@{}l@{}}{\rule{0pt}{2.6ex}\textbf{T7.04}\enspace \texttt{0123453}\enspace $n{=}7$, $|H|{=}2$, len 6, br 1, $|\mathrm{Aut}|{=}1$, orb 7/7, hist-orb 2/2; M 2304, res 448, rev 896: ok}\\*
0 & -- & 0 & 2 & 1 & 110 & 2 & 16 & \texttt{-++????}\\*
1 & 0 & 1 & 2 & 1 & 92 & 4 & 32 & \texttt{--+????}\\
2 & 1 & 2 & 2 & 1 & 56 & 8 & 64 & \texttt{---????}\\
3 & 2 & 3 & 1 & 1 & 96 & 32 & 0 & \texttt{----++-}\\
4 & 3 & 4 & 1 & 1 & 64 & 64 & 0 & \texttt{-----+-}\\
5 & 4 & 5 & 1 & 1 & 0 & 128 & 0 & \texttt{-------}\\
6 & 2 & 3 & 1 & 1 & 0 & 128 & 0 & \texttt{-------}\\
\multicolumn{9}{@{}l@{}}{\rule{0pt}{2.6ex}\textbf{T7.05}\enspace \texttt{0123452}\enspace $n{=}7$, $|H|{=}2$, len 6, br 1, $|\mathrm{Aut}|{=}1$, orb 7/7, hist-orb 2/2; M 2176, res 480, rev 800: ok}\\*
0 & -- & 0 & 2 & 1 & 94 & 2 & 32 & \texttt{-+?????}\\*
1 & 0 & 1 & 2 & 1 & 60 & 4 & 64 & \texttt{--?????}\\
2 & 1 & 2 & 1 & 1 & 112 & 16 & 0 & \texttt{---+++-}\\
3 & 2 & 3 & 1 & 1 & 96 & 32 & 0 & \texttt{----++-}\\
4 & 3 & 4 & 1 & 1 & 64 & 64 & 0 & \texttt{-----+-}\\
5 & 4 & 5 & 1 & 1 & 0 & 128 & 0 & \texttt{-------}\\
6 & 1 & 2 & 1 & 1 & 0 & 128 & 0 & \texttt{-------}\\
\multicolumn{9}{@{}l@{}}{\rule{0pt}{2.6ex}\textbf{T7.06}\enspace \texttt{0123451}\enspace $n{=}7$, $|H|{=}2$, len 6, br 1, $|\mathrm{Aut}|{=}1$, orb 7/7, hist-orb 2/2; M 2048, res 384, rev 752: ok}\\*
0 & -- & 0 & 2 & 1 & 62 & 2 & 64 & \texttt{-??????}\\*
1 & 0 & 1 & 1 & 1 & 120 & 8 & 0 & \texttt{--++++-}\\
2 & 1 & 2 & 1 & 1 & 112 & 16 & 0 & \texttt{---+++-}\\
3 & 2 & 3 & 1 & 1 & 96 & 32 & 0 & \texttt{----++-}\\
4 & 3 & 4 & 1 & 1 & 64 & 64 & 0 & \texttt{-----+-}\\
5 & 4 & 5 & 1 & 1 & 0 & 128 & 0 & \texttt{-------}\\
6 & 0 & 1 & 1 & 1 & 0 & 128 & 0 & \texttt{-------}\\
\multicolumn{9}{@{}l@{}}{\rule{0pt}{2.6ex}\textbf{T7.07}\enspace \texttt{0123444}\enspace $n{=}7$, $|H|{=}3$, len 5, br 1, $|\mathrm{Aut}|{=}6$, orb 5/4, hist-orb 1/0; M 2304, res 540, rev 952: ok}\\*
0 & -- & 0 & 3 & 1 & 114 & 2 & 12 & \texttt{-+++???}\\*
1 & 0 & 1 & 3 & 1 & 100 & 4 & 24 & \texttt{--++???}\\
2 & 1 & 2 & 3 & 1 & 72 & 8 & 48 & \texttt{---+???}\\
3 & 2 & 3 & 3 & 1 & 16 & 16 & 96 & \texttt{----???}\\
4 & 3 & 4 & 1 & 3 & 0 & 128 & 0 & \texttt{-------}\\
5 & 3 & 4 & 1 & 3 & 0 & 128 & 0 & \texttt{-------}\\
6 & 3 & 4 & 1 & 3 & 0 & 128 & 0 & \texttt{-------}\\
\multicolumn{9}{@{}l@{}}{\rule{0pt}{2.6ex}\textbf{T7.08}\enspace \texttt{0123443}\enspace $n{=}7$, $|H|{=}3$, len 5, br 2, $|\mathrm{Aut}|{=}2$, orb 6/5, hist-orb 2/1; M 2176, res 632, rev 820: ok}\\*
0 & -- & 0 & 3 & 1 & 106 & 2 & 20 & \texttt{-++????}\\*
1 & 0 & 1 & 3 & 1 & 84 & 4 & 40 & \texttt{--+????}\\
2 & 1 & 2 & 3 & 1 & 40 & 8 & 80 & \texttt{---????}\\
3 & 2 & 3 & 2 & 1 & 32 & 32 & 64 & \texttt{----??-}\\
4 & 3 & 4 & 1 & 2 & 0 & 128 & 0 & \texttt{-------}\\
5 & 3 & 4 & 1 & 2 & 0 & 128 & 0 & \texttt{-------}\\
6 & 2 & 3 & 1 & 1 & 0 & 128 & 0 & \texttt{-------}\\
\multicolumn{9}{@{}l@{}}{\rule{0pt}{2.6ex}\textbf{T7.09}\enspace \texttt{0123442}\enspace $n{=}7$, $|H|{=}3$, len 5, br 2, $|\mathrm{Aut}|{=}2$, orb 6/5, hist-orb 2/1; M 2048, res 684, rev 716: ok}\\*
0 & -- & 0 & 3 & 1 & 90 & 2 & 36 & \texttt{-+?????}\\*
1 & 0 & 1 & 3 & 1 & 52 & 4 & 72 & \texttt{--?????}\\
2 & 1 & 2 & 2 & 1 & 80 & 16 & 32 & \texttt{---+??-}\\
3 & 2 & 3 & 2 & 1 & 32 & 32 & 64 & \texttt{----??-}\\
4 & 3 & 4 & 1 & 2 & 0 & 128 & 0 & \texttt{-------}\\
5 & 3 & 4 & 1 & 2 & 0 & 128 & 0 & \texttt{-------}\\
6 & 1 & 2 & 1 & 1 & 0 & 128 & 0 & \texttt{-------}\\
\multicolumn{9}{@{}l@{}}{\rule{0pt}{2.6ex}\textbf{T7.10}\enspace \texttt{0123441}\enspace $n{=}7$, $|H|{=}3$, len 5, br 2, $|\mathrm{Aut}|{=}2$, orb 6/5, hist-orb 2/1; M 1920, res 608, rev 676: ok}\\*
0 & -- & 0 & 3 & 1 & 58 & 2 & 68 & \texttt{-??????}\\*
1 & 0 & 1 & 2 & 1 & 104 & 8 & 16 & \texttt{--++??-}\\
2 & 1 & 2 & 2 & 1 & 80 & 16 & 32 & \texttt{---+??-}\\
3 & 2 & 3 & 2 & 1 & 32 & 32 & 64 & \texttt{----??-}\\
4 & 3 & 4 & 1 & 2 & 0 & 128 & 0 & \texttt{-------}\\
5 & 3 & 4 & 1 & 2 & 0 & 128 & 0 & \texttt{-------}\\
6 & 0 & 1 & 1 & 1 & 0 & 128 & 0 & \texttt{-------}\\
\multicolumn{9}{@{}l@{}}{\rule{0pt}{2.6ex}\textbf{T7.11}\enspace \texttt{0123434}\enspace $n{=}7$, $|H|{=}2$, len 5, br 1, $|\mathrm{Aut}|{=}2$, orb 5/3, hist-orb 1/0; M 2176, res 336, rev 956: ok}\\*
0 & -- & 0 & 2 & 1 & 114 & 2 & 12 & \texttt{-++????}\\*
1 & 0 & 1 & 2 & 1 & 100 & 4 & 24 & \texttt{--+????}\\
2 & 1 & 2 & 2 & 1 & 72 & 8 & 48 & \texttt{---????}\\
3 & 2 & 3 & 1 & 2 & 64 & 64 & 0 & \texttt{----+--}\\
4 & 3 & 4 & 1 & 2 & 0 & 128 & 0 & \texttt{-------}\\
5 & 2 & 3 & 1 & 2 & 64 & 64 & 0 & \texttt{------+}\\
6 & 5 & 4 & 1 & 2 & 0 & 128 & 0 & \texttt{-------}\\
\multicolumn{9}{@{}l@{}}{\rule{0pt}{2.6ex}\textbf{T7.12}\enspace \texttt{0123433}\enspace $n{=}7$, $|H|{=}3$, len 5, br 1, $|\mathrm{Aut}|{=}2$, orb 6/5, hist-orb 2/1; M 2048, res 672, rev 776: ok}\\*
0 & -- & 0 & 3 & 1 & 102 & 2 & 24 & \texttt{-++????}\\*
1 & 0 & 1 & 3 & 1 & 76 & 4 & 48 & \texttt{--+????}\\
2 & 1 & 2 & 3 & 1 & 24 & 8 & 96 & \texttt{---????}\\
3 & 2 & 3 & 1 & 1 & 64 & 64 & 0 & \texttt{----+--}\\
4 & 3 & 4 & 1 & 1 & 0 & 128 & 0 & \texttt{-------}\\
5 & 2 & 3 & 1 & 2 & 0 & 128 & 0 & \texttt{-------}\\
6 & 2 & 3 & 1 & 2 & 0 & 128 & 0 & \texttt{-------}\\
\multicolumn{9}{@{}l@{}}{\rule{0pt}{2.6ex}\textbf{T7.13}\enspace \texttt{0123432}\enspace $n{=}7$, $|H|{=}3$, len 5, br 2, $|\mathrm{Aut}|{=}1$, orb 7/7, hist-orb 3/3; M 1920, res 744, rev 644: ok}\\*
0 & -- & 0 & 3 & 1 & 86 & 2 & 40 & \texttt{-+?????}\\*
1 & 0 & 1 & 3 & 1 & 44 & 4 & 80 & \texttt{--?????}\\
2 & 1 & 2 & 2 & 1 & 48 & 16 & 64 & \texttt{---???-}\\
3 & 2 & 3 & 1 & 1 & 64 & 64 & 0 & \texttt{----+--}\\
4 & 3 & 4 & 1 & 1 & 0 & 128 & 0 & \texttt{-------}\\
5 & 2 & 3 & 1 & 1 & 0 & 128 & 0 & \texttt{-------}\\
6 & 1 & 2 & 1 & 1 & 0 & 128 & 0 & \texttt{-------}\\
\multicolumn{9}{@{}l@{}}{\rule{0pt}{2.6ex}\textbf{T7.14}\enspace \texttt{0123431}\enspace $n{=}7$, $|H|{=}3$, len 5, br 2, $|\mathrm{Aut}|{=}1$, orb 7/7, hist-orb 3/3; M 1792, res 688, rev 596: ok}\\*
0 & -- & 0 & 3 & 1 & 54 & 2 & 72 & \texttt{-??????}\\*
1 & 0 & 1 & 2 & 1 & 88 & 8 & 32 & \texttt{--+???-}\\
2 & 1 & 2 & 2 & 1 & 48 & 16 & 64 & \texttt{---???-}\\
3 & 2 & 3 & 1 & 1 & 64 & 64 & 0 & \texttt{----+--}\\
4 & 3 & 4 & 1 & 1 & 0 & 128 & 0 & \texttt{-------}\\
5 & 2 & 3 & 1 & 1 & 0 & 128 & 0 & \texttt{-------}\\
6 & 0 & 1 & 1 & 1 & 0 & 128 & 0 & \texttt{-------}\\
\multicolumn{9}{@{}l@{}}{\rule{0pt}{2.6ex}\textbf{T7.15}\enspace \texttt{0123423}\enspace $n{=}7$, $|H|{=}2$, len 5, br 1, $|\mathrm{Aut}|{=}1$, orb 7/7, hist-orb 2/2; M 1920, res 300, rev 832: ok}\\*
0 & -- & 0 & 2 & 1 & 106 & 2 & 20 & \texttt{-+?????}\\*
1 & 0 & 1 & 2 & 1 & 84 & 4 & 40 & \texttt{--?????}\\
2 & 1 & 2 & 1 & 1 & 96 & 32 & 0 & \texttt{---++--}\\
3 & 2 & 3 & 1 & 1 & 64 & 64 & 0 & \texttt{----+--}\\
4 & 3 & 4 & 1 & 1 & 0 & 128 & 0 & \texttt{-------}\\
5 & 1 & 2 & 1 & 1 & 64 & 64 & 0 & \texttt{------+}\\
6 & 5 & 3 & 1 & 1 & 0 & 128 & 0 & \texttt{-------}\\
\multicolumn{9}{@{}l@{}}{\rule{0pt}{2.6ex}\textbf{T7.16}\enspace \texttt{0123422}\enspace $n{=}7$, $|H|{=}3$, len 5, br 1, $|\mathrm{Aut}|{=}2$, orb 6/5, hist-orb 2/1; M 1792, res 720, rev 584: ok}\\*
0 & -- & 0 & 3 & 1 & 78 & 2 & 48 & \texttt{-+?????}\\*
1 & 0 & 1 & 3 & 1 & 28 & 4 & 96 & \texttt{--?????}\\
2 & 1 & 2 & 1 & 1 & 96 & 32 & 0 & \texttt{---++--}\\
3 & 2 & 3 & 1 & 1 & 64 & 64 & 0 & \texttt{----+--}\\
4 & 3 & 4 & 1 & 1 & 0 & 128 & 0 & \texttt{-------}\\
5 & 1 & 2 & 1 & 2 & 0 & 128 & 0 & \texttt{-------}\\
6 & 1 & 2 & 1 & 2 & 0 & 128 & 0 & \texttt{-------}\\
\multicolumn{9}{@{}l@{}}{\rule{0pt}{2.6ex}\textbf{T7.17}\enspace \texttt{0123421}\enspace $n{=}7$, $|H|{=}3$, len 5, br 2, $|\mathrm{Aut}|{=}1$, orb 7/7, hist-orb 3/3; M 1664, res 704, rev 508: ok}\\*
0 & -- & 0 & 3 & 1 & 46 & 2 & 80 & \texttt{-??????}\\*
1 & 0 & 1 & 2 & 1 & 56 & 8 & 64 & \texttt{--????-}\\
2 & 1 & 2 & 1 & 1 & 96 & 32 & 0 & \texttt{---++--}\\
3 & 2 & 3 & 1 & 1 & 64 & 64 & 0 & \texttt{----+--}\\
4 & 3 & 4 & 1 & 1 & 0 & 128 & 0 & \texttt{-------}\\
5 & 1 & 2 & 1 & 1 & 0 & 128 & 0 & \texttt{-------}\\
6 & 0 & 1 & 1 & 1 & 0 & 128 & 0 & \texttt{-------}\\
\multicolumn{9}{@{}l@{}}{\rule{0pt}{2.6ex}\textbf{T7.18}\enspace \texttt{0123412}\enspace $n{=}7$, $|H|{=}2$, len 5, br 1, $|\mathrm{Aut}|{=}1$, orb 7/7, hist-orb 2/2; M 1664, res 216, rev 708: ok}\\*
0 & -- & 0 & 2 & 1 & 90 & 2 & 36 & \texttt{-??????}\\*
1 & 0 & 1 & 1 & 1 & 112 & 16 & 0 & \texttt{--+++--}\\
2 & 1 & 2 & 1 & 1 & 96 & 32 & 0 & \texttt{---++--}\\
3 & 2 & 3 & 1 & 1 & 64 & 64 & 0 & \texttt{----+--}\\
4 & 3 & 4 & 1 & 1 & 0 & 128 & 0 & \texttt{-------}\\
5 & 0 & 1 & 1 & 1 & 64 & 64 & 0 & \texttt{------+}\\
6 & 5 & 2 & 1 & 1 & 0 & 128 & 0 & \texttt{-------}\\
\multicolumn{9}{@{}l@{}}{\rule{0pt}{2.6ex}\textbf{T7.19}\enspace \texttt{0123411}\enspace $n{=}7$, $|H|{=}3$, len 5, br 1, $|\mathrm{Aut}|{=}2$, orb 6/5, hist-orb 2/1; M 1536, res 576, rev 480: ok}\\*
0 & -- & 0 & 3 & 1 & 30 & 2 & 96 & \texttt{-??????}\\*
1 & 0 & 1 & 1 & 1 & 112 & 16 & 0 & \texttt{--+++--}\\
2 & 1 & 2 & 1 & 1 & 96 & 32 & 0 & \texttt{---++--}\\
3 & 2 & 3 & 1 & 1 & 64 & 64 & 0 & \texttt{----+--}\\
4 & 3 & 4 & 1 & 1 & 0 & 128 & 0 & \texttt{-------}\\
5 & 0 & 1 & 1 & 2 & 0 & 128 & 0 & \texttt{-------}\\
6 & 0 & 1 & 1 & 2 & 0 & 128 & 0 & \texttt{-------}\\
\multicolumn{9}{@{}l@{}}{\rule{0pt}{2.6ex}\textbf{T7.20}\enspace \texttt{0123333}\enspace $n{=}7$, $|H|{=}4$, len 4, br 1, $|\mathrm{Aut}|{=}24$, orb 4/3, hist-orb 1/0; M 1920, res 784, rev 708: ok}\\*
0 & -- & 0 & 4 & 1 & 98 & 2 & 28 & \texttt{-++????}\\*
1 & 0 & 1 & 4 & 1 & 68 & 4 & 56 & \texttt{--+????}\\
2 & 1 & 2 & 4 & 1 & 8 & 8 & 112 & \texttt{---????}\\
3 & 2 & 3 & 1 & 4 & 0 & 128 & 0 & \texttt{-------}\\
4 & 2 & 3 & 1 & 4 & 0 & 128 & 0 & \texttt{-------}\\
5 & 2 & 3 & 1 & 4 & 0 & 128 & 0 & \texttt{-------}\\
6 & 2 & 3 & 1 & 4 & 0 & 128 & 0 & \texttt{-------}\\
\multicolumn{9}{@{}l@{}}{\rule{0pt}{2.6ex}\textbf{T7.21}\enspace \texttt{0123332}\enspace $n{=}7$, $|H|{=}4$, len 4, br 2, $|\mathrm{Aut}|{=}6$, orb 5/4, hist-orb 2/1; M 1792, res 876, rev 536: ok}\\*
0 & -- & 0 & 4 & 1 & 82 & 2 & 44 & \texttt{-+?????}\\*
1 & 0 & 1 & 4 & 1 & 36 & 4 & 88 & \texttt{--?????}\\
2 & 1 & 2 & 3 & 1 & 16 & 16 & 96 & \texttt{---???-}\\
3 & 2 & 3 & 1 & 3 & 0 & 128 & 0 & \texttt{-------}\\
4 & 2 & 3 & 1 & 3 & 0 & 128 & 0 & \texttt{-------}\\
5 & 2 & 3 & 1 & 3 & 0 & 128 & 0 & \texttt{-------}\\
6 & 1 & 2 & 1 & 1 & 0 & 128 & 0 & \texttt{-------}\\
\multicolumn{9}{@{}l@{}}{\rule{0pt}{2.6ex}\textbf{T7.22}\enspace \texttt{0123331}\enspace $n{=}7$, $|H|{=}4$, len 4, br 2, $|\mathrm{Aut}|{=}6$, orb 5/4, hist-orb 2/1; M 1664, res 840, rev 480: ok}\\*
0 & -- & 0 & 4 & 1 & 50 & 2 & 76 & \texttt{-??????}\\*
1 & 0 & 1 & 3 & 1 & 72 & 8 & 48 & \texttt{--+???-}\\
2 & 1 & 2 & 3 & 1 & 16 & 16 & 96 & \texttt{---???-}\\
3 & 2 & 3 & 1 & 3 & 0 & 128 & 0 & \texttt{-------}\\
4 & 2 & 3 & 1 & 3 & 0 & 128 & 0 & \texttt{-------}\\
5 & 2 & 3 & 1 & 3 & 0 & 128 & 0 & \texttt{-------}\\
6 & 0 & 1 & 1 & 1 & 0 & 128 & 0 & \texttt{-------}\\
\multicolumn{9}{@{}l@{}}{\rule{0pt}{2.6ex}\textbf{T7.23}\enspace \texttt{0123323}\enspace $n{=}7$, $|H|{=}3$, len 4, br 2, $|\mathrm{Aut}|{=}2$, orb 6/5, hist-orb 2/1; M 1792, res 560, rev 688: ok}\\*
0 & -- & 0 & 3 & 1 & 94 & 2 & 32 & \texttt{-+?????}\\*
1 & 0 & 1 & 3 & 1 & 60 & 4 & 64 & \texttt{--?????}\\
2 & 1 & 2 & 2 & 1 & 32 & 32 & 64 & \texttt{---??--}\\
3 & 2 & 3 & 1 & 2 & 0 & 128 & 0 & \texttt{-------}\\
4 & 2 & 3 & 1 & 2 & 0 & 128 & 0 & \texttt{-------}\\
5 & 1 & 2 & 1 & 1 & 64 & 64 & 0 & \texttt{------+}\\
6 & 5 & 3 & 1 & 1 & 0 & 128 & 0 & \texttt{-------}\\
\multicolumn{9}{@{}l@{}}{\rule{0pt}{2.6ex}\textbf{T7.24}\enspace \texttt{0123322}\enspace $n{=}7$, $|H|{=}4$, len 4, br 2, $|\mathrm{Aut}|{=}4$, orb 5/3, hist-orb 2/0; M 1664, res 860, rev 464: ok}\\*
0 & -- & 0 & 4 & 1 & 74 & 2 & 52 & \texttt{-+?????}\\*
1 & 0 & 1 & 4 & 1 & 20 & 4 & 104 & \texttt{--?????}\\
2 & 1 & 2 & 2 & 1 & 32 & 32 & 64 & \texttt{---??--}\\
3 & 2 & 3 & 1 & 2 & 0 & 128 & 0 & \texttt{-------}\\
4 & 2 & 3 & 1 & 2 & 0 & 128 & 0 & \texttt{-------}\\
5 & 1 & 2 & 1 & 2 & 0 & 128 & 0 & \texttt{-------}\\
6 & 1 & 2 & 1 & 2 & 0 & 128 & 0 & \texttt{-------}\\
\multicolumn{9}{@{}l@{}}{\rule{0pt}{2.6ex}\textbf{T7.25}\enspace \texttt{0123321}\enspace $n{=}7$, $|H|{=}4$, len 4, br 3, $|\mathrm{Aut}|{=}2$, orb 6/5, hist-orb 3/2; M 1536, res 864, rev 372: ok}\\*
0 & -- & 0 & 4 & 1 & 42 & 2 & 84 & \texttt{-??????}\\*
1 & 0 & 1 & 3 & 1 & 40 & 8 & 80 & \texttt{--????-}\\
2 & 1 & 2 & 2 & 1 & 32 & 32 & 64 & \texttt{---??--}\\
3 & 2 & 3 & 1 & 2 & 0 & 128 & 0 & \texttt{-------}\\
4 & 2 & 3 & 1 & 2 & 0 & 128 & 0 & \texttt{-------}\\
5 & 1 & 2 & 1 & 1 & 0 & 128 & 0 & \texttt{-------}\\
6 & 0 & 1 & 1 & 1 & 0 & 128 & 0 & \texttt{-------}\\
\multicolumn{9}{@{}l@{}}{\rule{0pt}{2.6ex}\textbf{T7.26}\enspace \texttt{0123312}\enspace $n{=}7$, $|H|{=}3$, len 4, br 2, $|\mathrm{Aut}|{=}2$, orb 6/5, hist-orb 2/1; M 1536, res 448, rev 588: ok}\\*
0 & -- & 0 & 3 & 1 & 78 & 2 & 48 & \texttt{-??????}\\*
1 & 0 & 1 & 2 & 1 & 80 & 16 & 32 & \texttt{--+??--}\\
2 & 1 & 2 & 2 & 1 & 32 & 32 & 64 & \texttt{---??--}\\
3 & 2 & 3 & 1 & 2 & 0 & 128 & 0 & \texttt{-------}\\
4 & 2 & 3 & 1 & 2 & 0 & 128 & 0 & \texttt{-------}\\
5 & 0 & 1 & 1 & 1 & 64 & 64 & 0 & \texttt{------+}\\
6 & 5 & 2 & 1 & 1 & 0 & 128 & 0 & \texttt{-------}\\
\multicolumn{9}{@{}l@{}}{\rule{0pt}{2.6ex}\textbf{T7.27}\enspace \texttt{0123311}\enspace $n{=}7$, $|H|{=}4$, len 4, br 2, $|\mathrm{Aut}|{=}4$, orb 5/3, hist-orb 2/0; M 1408, res 760, rev 352: ok}\\*
0 & -- & 0 & 4 & 1 & 26 & 2 & 100 & \texttt{-??????}\\*
1 & 0 & 1 & 2 & 1 & 80 & 16 & 32 & \texttt{--+??--}\\
2 & 1 & 2 & 2 & 1 & 32 & 32 & 64 & \texttt{---??--}\\
3 & 2 & 3 & 1 & 2 & 0 & 128 & 0 & \texttt{-------}\\
4 & 2 & 3 & 1 & 2 & 0 & 128 & 0 & \texttt{-------}\\
5 & 0 & 1 & 1 & 2 & 0 & 128 & 0 & \texttt{-------}\\
6 & 0 & 1 & 1 & 2 & 0 & 128 & 0 & \texttt{-------}\\
\multicolumn{9}{@{}l@{}}{\rule{0pt}{2.6ex}\textbf{T7.28}\enspace \texttt{0123232}\enspace $n{=}7$, $|H|{=}3$, len 4, br 1, $|\mathrm{Aut}|{=}2$, orb 5/3, hist-orb 2/1; M 1664, res 660, rev 584: ok}\\*
0 & -- & 0 & 3 & 1 & 82 & 2 & 44 & \texttt{-+?????}\\*
1 & 0 & 1 & 3 & 1 & 36 & 4 & 88 & \texttt{--?????}\\
2 & 1 & 2 & 1 & 2 & 64 & 64 & 0 & \texttt{---+---}\\
3 & 2 & 3 & 1 & 2 & 0 & 128 & 0 & \texttt{-------}\\
4 & 1 & 2 & 1 & 2 & 64 & 64 & 0 & \texttt{-----+-}\\
5 & 4 & 3 & 1 & 2 & 0 & 128 & 0 & \texttt{-------}\\
6 & 1 & 2 & 1 & 1 & 0 & 128 & 0 & \texttt{-------}\\
\multicolumn{9}{@{}l@{}}{\rule{0pt}{2.6ex}\textbf{T7.29}\enspace \texttt{0123231}\enspace $n{=}7$, $|H|{=}3$, len 4, br 2, $|\mathrm{Aut}|{=}2$, orb 5/3, hist-orb 2/1; M 1536, res 624, rev 516: ok}\\*
0 & -- & 0 & 3 & 1 & 50 & 2 & 76 & \texttt{-??????}\\*
1 & 0 & 1 & 2 & 1 & 72 & 8 & 48 & \texttt{--????-}\\
2 & 1 & 2 & 1 & 2 & 64 & 64 & 0 & \texttt{---+---}\\
3 & 2 & 3 & 1 & 2 & 0 & 128 & 0 & \texttt{-------}\\
4 & 1 & 2 & 1 & 2 & 64 & 64 & 0 & \texttt{-----+-}\\
5 & 4 & 3 & 1 & 2 & 0 & 128 & 0 & \texttt{-------}\\
6 & 0 & 1 & 1 & 1 & 0 & 128 & 0 & \texttt{-------}\\
\multicolumn{9}{@{}l@{}}{\rule{0pt}{2.6ex}\textbf{T7.30}\enspace \texttt{0123222}\enspace $n{=}7$, $|H|{=}4$, len 4, br 1, $|\mathrm{Aut}|{=}6$, orb 5/4, hist-orb 2/1; M 1536, res 840, rev 432: ok}\\*
0 & -- & 0 & 4 & 1 & 70 & 2 & 56 & \texttt{-+?????}\\*
1 & 0 & 1 & 4 & 1 & 12 & 4 & 112 & \texttt{--?????}\\
2 & 1 & 2 & 1 & 1 & 64 & 64 & 0 & \texttt{---+---}\\
3 & 2 & 3 & 1 & 1 & 0 & 128 & 0 & \texttt{-------}\\
4 & 1 & 2 & 1 & 3 & 0 & 128 & 0 & \texttt{-------}\\
5 & 1 & 2 & 1 & 3 & 0 & 128 & 0 & \texttt{-------}\\
6 & 1 & 2 & 1 & 3 & 0 & 128 & 0 & \texttt{-------}\\
\multicolumn{9}{@{}l@{}}{\rule{0pt}{2.6ex}\textbf{T7.31}\enspace \texttt{0123221}\enspace $n{=}7$, $|H|{=}4$, len 4, br 2, $|\mathrm{Aut}|{=}2$, orb 6/5, hist-orb 3/2; M 1408, res 864, rev 316: ok}\\*
0 & -- & 0 & 4 & 1 & 38 & 2 & 88 & \texttt{-??????}\\*
1 & 0 & 1 & 3 & 1 & 24 & 8 & 96 & \texttt{--????-}\\
2 & 1 & 2 & 1 & 1 & 64 & 64 & 0 & \texttt{---+---}\\
3 & 2 & 3 & 1 & 1 & 0 & 128 & 0 & \texttt{-------}\\
4 & 1 & 2 & 1 & 2 & 0 & 128 & 0 & \texttt{-------}\\
5 & 1 & 2 & 1 & 2 & 0 & 128 & 0 & \texttt{-------}\\
6 & 0 & 1 & 1 & 1 & 0 & 128 & 0 & \texttt{-------}\\
\multicolumn{9}{@{}l@{}}{\rule{0pt}{2.6ex}\textbf{T7.32}\enspace \texttt{0123212}\enspace $n{=}7$, $|H|{=}3$, len 4, br 2, $|\mathrm{Aut}|{=}1$, orb 7/7, hist-orb 3/3; M 1408, res 520, rev 496: ok}\\*
0 & -- & 0 & 3 & 1 & 66 & 2 & 60 & \texttt{-??????}\\*
1 & 0 & 1 & 2 & 1 & 48 & 16 & 64 & \texttt{--???--}\\
2 & 1 & 2 & 1 & 1 & 64 & 64 & 0 & \texttt{---+---}\\
3 & 2 & 3 & 1 & 1 & 0 & 128 & 0 & \texttt{-------}\\
4 & 1 & 2 & 1 & 1 & 0 & 128 & 0 & \texttt{-------}\\
5 & 0 & 1 & 1 & 1 & 64 & 64 & 0 & \texttt{------+}\\
6 & 5 & 2 & 1 & 1 & 0 & 128 & 0 & \texttt{-------}\\
\multicolumn{9}{@{}l@{}}{\rule{0pt}{2.6ex}\textbf{T7.33}\enspace \texttt{0123211}\enspace $n{=}7$, $|H|{=}4$, len 4, br 2, $|\mathrm{Aut}|{=}2$, orb 6/5, hist-orb 3/2; M 1280, res 784, rev 276: ok}\\*
0 & -- & 0 & 4 & 1 & 22 & 2 & 104 & \texttt{-??????}\\*
1 & 0 & 1 & 2 & 1 & 48 & 16 & 64 & \texttt{--???--}\\
2 & 1 & 2 & 1 & 1 & 64 & 64 & 0 & \texttt{---+---}\\
3 & 2 & 3 & 1 & 1 & 0 & 128 & 0 & \texttt{-------}\\
4 & 1 & 2 & 1 & 1 & 0 & 128 & 0 & \texttt{-------}\\
5 & 0 & 1 & 1 & 2 & 0 & 128 & 0 & \texttt{-------}\\
6 & 0 & 1 & 1 & 2 & 0 & 128 & 0 & \texttt{-------}\\
\multicolumn{9}{@{}l@{}}{\rule{0pt}{2.6ex}\textbf{T7.34}\enspace \texttt{0123123}\enspace $n{=}7$, $|H|{=}2$, len 4, br 1, $|\mathrm{Aut}|{=}2$, orb 4/1, hist-orb 1/0; M 1536, res 168, rev 692: ok}\\*
0 & -- & 0 & 2 & 1 & 98 & 2 & 28 & \texttt{-??????}\\*
1 & 0 & 1 & 1 & 2 & 96 & 32 & 0 & \texttt{--++---}\\
2 & 1 & 2 & 1 & 2 & 64 & 64 & 0 & \texttt{---+---}\\
3 & 2 & 3 & 1 & 2 & 0 & 128 & 0 & \texttt{-------}\\
4 & 0 & 1 & 1 & 2 & 96 & 32 & 0 & \texttt{-----++}\\
5 & 4 & 2 & 1 & 2 & 64 & 64 & 0 & \texttt{------+}\\
6 & 5 & 3 & 1 & 2 & 0 & 128 & 0 & \texttt{-------}\\
\multicolumn{9}{@{}l@{}}{\rule{0pt}{2.6ex}\textbf{T7.35}\enspace \texttt{0123122}\enspace $n{=}7$, $|H|{=}3$, len 4, br 2, $|\mathrm{Aut}|{=}2$, orb 6/5, hist-orb 2/1; M 1408, res 432, rev 520: ok}\\*
0 & -- & 0 & 3 & 1 & 70 & 2 & 56 & \texttt{-??????}\\*
1 & 0 & 1 & 1 & 1 & 96 & 32 & 0 & \texttt{--++---}\\
2 & 1 & 2 & 1 & 1 & 64 & 64 & 0 & \texttt{---+---}\\
3 & 2 & 3 & 1 & 1 & 0 & 128 & 0 & \texttt{-------}\\
4 & 0 & 1 & 2 & 1 & 32 & 32 & 64 & \texttt{-----??}\\
5 & 4 & 2 & 1 & 2 & 0 & 128 & 0 & \texttt{-------}\\
6 & 4 & 2 & 1 & 2 & 0 & 128 & 0 & \texttt{-------}\\
\multicolumn{9}{@{}l@{}}{\rule{0pt}{2.6ex}\textbf{T7.36}\enspace \texttt{0123121}\enspace $n{=}7$, $|H|{=}3$, len 4, br 1, $|\mathrm{Aut}|{=}1$, orb 7/7, hist-orb 3/3; M 1280, res 504, rev 420: ok}\\*
0 & -- & 0 & 3 & 1 & 42 & 2 & 84 & \texttt{-??????}\\*
1 & 0 & 1 & 1 & 1 & 96 & 32 & 0 & \texttt{--++---}\\
2 & 1 & 2 & 1 & 1 & 64 & 64 & 0 & \texttt{---+---}\\
3 & 2 & 3 & 1 & 1 & 0 & 128 & 0 & \texttt{-------}\\
4 & 0 & 1 & 1 & 1 & 64 & 64 & 0 & \texttt{-----+-}\\
5 & 4 & 2 & 1 & 1 & 0 & 128 & 0 & \texttt{-------}\\
6 & 0 & 1 & 1 & 1 & 0 & 128 & 0 & \texttt{-------}\\
\multicolumn{9}{@{}l@{}}{\rule{0pt}{2.6ex}\textbf{T7.37}\enspace \texttt{0123111}\enspace $n{=}7$, $|H|{=}4$, len 4, br 1, $|\mathrm{Aut}|{=}6$, orb 5/4, hist-orb 2/1; M 1152, res 672, rev 256: ok}\\*
0 & -- & 0 & 4 & 1 & 14 & 2 & 112 & \texttt{-??????}\\*
1 & 0 & 1 & 1 & 1 & 96 & 32 & 0 & \texttt{--++---}\\
2 & 1 & 2 & 1 & 1 & 64 & 64 & 0 & \texttt{---+---}\\
3 & 2 & 3 & 1 & 1 & 0 & 128 & 0 & \texttt{-------}\\
4 & 0 & 1 & 1 & 3 & 0 & 128 & 0 & \texttt{-------}\\
5 & 0 & 1 & 1 & 3 & 0 & 128 & 0 & \texttt{-------}\\
6 & 0 & 1 & 1 & 3 & 0 & 128 & 0 & \texttt{-------}\\
\multicolumn{9}{@{}l@{}}{\rule{0pt}{2.6ex}\textbf{T7.38}\enspace \texttt{0122222}\enspace $n{=}7$, $|H|{=}5$, len 3, br 1, $|\mathrm{Aut}|{=}120$, orb 3/2, hist-orb 1/0; M 1408, res 900, rev 352: ok}\\*
0 & -- & 0 & 5 & 1 & 66 & 2 & 60 & \texttt{-+?????}\\*
1 & 0 & 1 & 5 & 1 & 4 & 4 & 120 & \texttt{--?????}\\
2 & 1 & 2 & 1 & 5 & 0 & 128 & 0 & \texttt{-------}\\
3 & 1 & 2 & 1 & 5 & 0 & 128 & 0 & \texttt{-------}\\
4 & 1 & 2 & 1 & 5 & 0 & 128 & 0 & \texttt{-------}\\
5 & 1 & 2 & 1 & 5 & 0 & 128 & 0 & \texttt{-------}\\
6 & 1 & 2 & 1 & 5 & 0 & 128 & 0 & \texttt{-------}\\
\multicolumn{9}{@{}l@{}}{\rule{0pt}{2.6ex}\textbf{T7.39}\enspace \texttt{0122221}\enspace $n{=}7$, $|H|{=}5$, len 3, br 2, $|\mathrm{Aut}|{=}24$, orb 4/3, hist-orb 2/1; M 1280, res 944, rev 204: ok}\\*
0 & -- & 0 & 5 & 1 & 34 & 2 & 92 & \texttt{-??????}\\*
1 & 0 & 1 & 4 & 1 & 8 & 8 & 112 & \texttt{--????-}\\
2 & 1 & 2 & 1 & 4 & 0 & 128 & 0 & \texttt{-------}\\
3 & 1 & 2 & 1 & 4 & 0 & 128 & 0 & \texttt{-------}\\
4 & 1 & 2 & 1 & 4 & 0 & 128 & 0 & \texttt{-------}\\
5 & 1 & 2 & 1 & 4 & 0 & 128 & 0 & \texttt{-------}\\
6 & 0 & 1 & 1 & 1 & 0 & 128 & 0 & \texttt{-------}\\
\multicolumn{9}{@{}l@{}}{\rule{0pt}{2.6ex}\textbf{T7.40}\enspace \texttt{0122212}\enspace $n{=}7$, $|H|{=}4$, len 3, br 2, $|\mathrm{Aut}|{=}6$, orb 5/4, hist-orb 2/1; M 1280, res 672, rev 352: ok}\\*
0 & -- & 0 & 4 & 1 & 54 & 2 & 72 & \texttt{-??????}\\*
1 & 0 & 1 & 3 & 1 & 16 & 16 & 96 & \texttt{--???--}\\
2 & 1 & 2 & 1 & 3 & 0 & 128 & 0 & \texttt{-------}\\
3 & 1 & 2 & 1 & 3 & 0 & 128 & 0 & \texttt{-------}\\
4 & 1 & 2 & 1 & 3 & 0 & 128 & 0 & \texttt{-------}\\
5 & 0 & 1 & 1 & 1 & 64 & 64 & 0 & \texttt{------+}\\
6 & 5 & 2 & 1 & 1 & 0 & 128 & 0 & \texttt{-------}\\
\multicolumn{9}{@{}l@{}}{\rule{0pt}{2.6ex}\textbf{T7.41}\enspace \texttt{0122211}\enspace $n{=}7$, $|H|{=}5$, len 3, br 2, $|\mathrm{Aut}|{=}12$, orb 4/2, hist-orb 2/0; M 1152, res 888, rev 140: ok}\\*
0 & -- & 0 & 5 & 1 & 18 & 2 & 108 & \texttt{-??????}\\*
1 & 0 & 1 & 3 & 1 & 16 & 16 & 96 & \texttt{--???--}\\
2 & 1 & 2 & 1 & 3 & 0 & 128 & 0 & \texttt{-------}\\
3 & 1 & 2 & 1 & 3 & 0 & 128 & 0 & \texttt{-------}\\
4 & 1 & 2 & 1 & 3 & 0 & 128 & 0 & \texttt{-------}\\
5 & 0 & 1 & 1 & 2 & 0 & 128 & 0 & \texttt{-------}\\
6 & 0 & 1 & 1 & 2 & 0 & 128 & 0 & \texttt{-------}\\
\multicolumn{9}{@{}l@{}}{\rule{0pt}{2.6ex}\textbf{T7.42}\enspace \texttt{0122122}\enspace $n{=}7$, $|H|{=}4$, len 3, br 3, $|\mathrm{Aut}|{=}8$, orb 3/1, hist-orb 1/0; M 1280, res 648, rev 348: ok}\\*
0 & -- & 0 & 4 & 1 & 50 & 2 & 76 & \texttt{-??????}\\*
1 & 0 & 1 & 2 & 2 & 32 & 32 & 64 & \texttt{--??---}\\
2 & 1 & 2 & 1 & 4 & 0 & 128 & 0 & \texttt{-------}\\
3 & 1 & 2 & 1 & 4 & 0 & 128 & 0 & \texttt{-------}\\
4 & 0 & 1 & 2 & 2 & 32 & 32 & 64 & \texttt{-----??}\\
5 & 4 & 2 & 1 & 4 & 0 & 128 & 0 & \texttt{-------}\\
6 & 4 & 2 & 1 & 4 & 0 & 128 & 0 & \texttt{-------}\\
\multicolumn{9}{@{}l@{}}{\rule{0pt}{2.6ex}\textbf{T7.43}\enspace \texttt{0122121}\enspace $n{=}7$, $|H|{=}4$, len 3, br 2, $|\mathrm{Aut}|{=}2$, orb 6/5, hist-orb 3/2; M 1152, res 672, rev 264: ok}\\*
0 & -- & 0 & 4 & 1 & 30 & 2 & 96 & \texttt{-??????}\\*
1 & 0 & 1 & 2 & 1 & 32 & 32 & 64 & \texttt{--??---}\\
2 & 1 & 2 & 1 & 2 & 0 & 128 & 0 & \texttt{-------}\\
3 & 1 & 2 & 1 & 2 & 0 & 128 & 0 & \texttt{-------}\\
4 & 0 & 1 & 1 & 1 & 64 & 64 & 0 & \texttt{-----+-}\\
5 & 4 & 2 & 1 & 1 & 0 & 128 & 0 & \texttt{-------}\\
6 & 0 & 1 & 1 & 1 & 0 & 128 & 0 & \texttt{-------}\\
\multicolumn{9}{@{}l@{}}{\rule{0pt}{2.6ex}\textbf{T7.44}\enspace \texttt{0122111}\enspace $n{=}7$, $|H|{=}5$, len 3, br 2, $|\mathrm{Aut}|{=}12$, orb 4/2, hist-orb 2/0; M 1024, res 792, rev 116: ok}\\*
0 & -- & 0 & 5 & 1 & 10 & 2 & 116 & \texttt{-??????}\\*
1 & 0 & 1 & 2 & 1 & 32 & 32 & 64 & \texttt{--??---}\\
2 & 1 & 2 & 1 & 2 & 0 & 128 & 0 & \texttt{-------}\\
3 & 1 & 2 & 1 & 2 & 0 & 128 & 0 & \texttt{-------}\\
4 & 0 & 1 & 1 & 3 & 0 & 128 & 0 & \texttt{-------}\\
5 & 0 & 1 & 1 & 3 & 0 & 128 & 0 & \texttt{-------}\\
6 & 0 & 1 & 1 & 3 & 0 & 128 & 0 & \texttt{-------}\\
\multicolumn{9}{@{}l@{}}{\rule{0pt}{2.6ex}\textbf{T7.45}\enspace \texttt{0121212}\enspace $n{=}7$, $|H|{=}3$, len 3, br 1, $|\mathrm{Aut}|{=}6$, orb 3/1, hist-orb 1/0; M 1152, res 432, rev 408: ok}\\*
0 & -- & 0 & 3 & 1 & 54 & 2 & 72 & \texttt{-??????}\\*
1 & 0 & 1 & 1 & 3 & 64 & 64 & 0 & \texttt{--+----}\\
2 & 1 & 2 & 1 & 3 & 0 & 128 & 0 & \texttt{-------}\\
3 & 0 & 1 & 1 & 3 & 64 & 64 & 0 & \texttt{----+--}\\
4 & 3 & 2 & 1 & 3 & 0 & 128 & 0 & \texttt{-------}\\
5 & 0 & 1 & 1 & 3 & 64 & 64 & 0 & \texttt{------+}\\
6 & 5 & 2 & 1 & 3 & 0 & 128 & 0 & \texttt{-------}\\
\multicolumn{9}{@{}l@{}}{\rule{0pt}{2.6ex}\textbf{T7.46}\enspace \texttt{0121211}\enspace $n{=}7$, $|H|{=}4$, len 3, br 1, $|\mathrm{Aut}|{=}4$, orb 4/1, hist-orb 2/0; M 1024, res 648, rev 212: ok}\\*
0 & -- & 0 & 4 & 1 & 18 & 2 & 108 & \texttt{-??????}\\*
1 & 0 & 1 & 1 & 2 & 64 & 64 & 0 & \texttt{--+----}\\
2 & 1 & 2 & 1 & 2 & 0 & 128 & 0 & \texttt{-------}\\
3 & 0 & 1 & 1 & 2 & 64 & 64 & 0 & \texttt{----+--}\\
4 & 3 & 2 & 1 & 2 & 0 & 128 & 0 & \texttt{-------}\\
5 & 0 & 1 & 1 & 2 & 0 & 128 & 0 & \texttt{-------}\\
6 & 0 & 1 & 1 & 2 & 0 & 128 & 0 & \texttt{-------}\\
\multicolumn{9}{@{}l@{}}{\rule{0pt}{2.6ex}\textbf{T7.47}\enspace \texttt{0121111}\enspace $n{=}7$, $|H|{=}5$, len 3, br 1, $|\mathrm{Aut}|{=}24$, orb 4/3, hist-orb 2/1; M 896, res 720, rev 96: ok}\\*
0 & -- & 0 & 5 & 1 & 6 & 2 & 120 & \texttt{-??????}\\*
1 & 0 & 1 & 1 & 1 & 64 & 64 & 0 & \texttt{--+----}\\
2 & 1 & 2 & 1 & 1 & 0 & 128 & 0 & \texttt{-------}\\
3 & 0 & 1 & 1 & 4 & 0 & 128 & 0 & \texttt{-------}\\
4 & 0 & 1 & 1 & 4 & 0 & 128 & 0 & \texttt{-------}\\
5 & 0 & 1 & 1 & 4 & 0 & 128 & 0 & \texttt{-------}\\
6 & 0 & 1 & 1 & 4 & 0 & 128 & 0 & \texttt{-------}\\
\multicolumn{9}{@{}l@{}}{\rule{0pt}{2.6ex}\textbf{T7.48}\enspace \texttt{0111111}\enspace $n{=}7$, $|H|{=}6$, len 2, br 1, $|\mathrm{Aut}|{=}720$, orb 2/1, hist-orb 1/0; M 768, res 744, rev 12: ok}\\*
0 & -- & 0 & 6 & 1 & 2 & 2 & 124 & \texttt{-??????}\\*
1 & 0 & 1 & 1 & 6 & 0 & 128 & 0 & \texttt{-------}\\
2 & 0 & 1 & 1 & 6 & 0 & 128 & 0 & \texttt{-------}\\
3 & 0 & 1 & 1 & 6 & 0 & 128 & 0 & \texttt{-------}\\
4 & 0 & 1 & 1 & 6 & 0 & 128 & 0 & \texttt{-------}\\
5 & 0 & 1 & 1 & 6 & 0 & 128 & 0 & \texttt{-------}\\
6 & 0 & 1 & 1 & 6 & 0 & 128 & 0 & \texttt{-------}\\
\multicolumn{9}{@{}l@{}}{\rule{0pt}{3.4ex}\normalsize\bfseries Frames with 8 moments (115 classes)}\\*
\multicolumn{9}{@{}l@{}}{\rule{0pt}{2.6ex}\textbf{T8.001}\enspace \texttt{01234567}\enspace $n{=}8$, $|H|{=}1$, len 8, br 0, $|\mathrm{Aut}|{=}1$, orb 8/8, hist-orb 1/1; M 7168, res 0, rev 2582: ok}\\*
0 & -- & 0 & 1 & 1 & 254 & 2 & 0 & \texttt{-+++++++}\\*
1 & 0 & 1 & 1 & 1 & 252 & 4 & 0 & \texttt{--++++++}\\
2 & 1 & 2 & 1 & 1 & 248 & 8 & 0 & \texttt{---+++++}\\
3 & 2 & 3 & 1 & 1 & 240 & 16 & 0 & \texttt{----++++}\\
4 & 3 & 4 & 1 & 1 & 224 & 32 & 0 & \texttt{-----+++}\\
5 & 4 & 5 & 1 & 1 & 192 & 64 & 0 & \texttt{------++}\\
6 & 5 & 6 & 1 & 1 & 128 & 128 & 0 & \texttt{-------+}\\
7 & 6 & 7 & 1 & 1 & 0 & 256 & 0 & \texttt{--------}\\
\multicolumn{9}{@{}l@{}}{\rule{0pt}{2.6ex}\textbf{T8.002}\enspace \texttt{01234566}\enspace $n{=}8$, $|H|{=}2$, len 7, br 1, $|\mathrm{Aut}|{=}2$, orb 7/6, hist-orb 1/0; M 6912, res 504, rev 2718: ok}\\*
0 & -- & 0 & 2 & 1 & 250 & 2 & 4 & \texttt{-+++++??}\\*
1 & 0 & 1 & 2 & 1 & 244 & 4 & 8 & \texttt{--++++??}\\
2 & 1 & 2 & 2 & 1 & 232 & 8 & 16 & \texttt{---+++??}\\
3 & 2 & 3 & 2 & 1 & 208 & 16 & 32 & \texttt{----++??}\\
4 & 3 & 4 & 2 & 1 & 160 & 32 & 64 & \texttt{-----+??}\\
5 & 4 & 5 & 2 & 1 & 64 & 64 & 128 & \texttt{------??}\\
6 & 5 & 6 & 1 & 2 & 0 & 256 & 0 & \texttt{--------}\\
7 & 5 & 6 & 1 & 2 & 0 & 256 & 0 & \texttt{--------}\\
\multicolumn{9}{@{}l@{}}{\rule{0pt}{2.6ex}\textbf{T8.003}\enspace \texttt{01234565}\enspace $n{=}8$, $|H|{=}2$, len 7, br 1, $|\mathrm{Aut}|{=}1$, orb 8/8, hist-orb 2/2; M 6656, res 744, rev 2666: ok}\\*
0 & -- & 0 & 2 & 1 & 246 & 2 & 8 & \texttt{-++++???}\\*
1 & 0 & 1 & 2 & 1 & 236 & 4 & 16 & \texttt{--+++???}\\
2 & 1 & 2 & 2 & 1 & 216 & 8 & 32 & \texttt{---++???}\\
3 & 2 & 3 & 2 & 1 & 176 & 16 & 64 & \texttt{----+???}\\
4 & 3 & 4 & 2 & 1 & 96 & 32 & 128 & \texttt{-----???}\\
5 & 4 & 5 & 1 & 1 & 128 & 128 & 0 & \texttt{------+-}\\
6 & 5 & 6 & 1 & 1 & 0 & 256 & 0 & \texttt{--------}\\
7 & 4 & 5 & 1 & 1 & 0 & 256 & 0 & \texttt{--------}\\
\multicolumn{9}{@{}l@{}}{\rule{0pt}{2.6ex}\textbf{T8.004}\enspace \texttt{01234564}\enspace $n{=}8$, $|H|{=}2$, len 7, br 1, $|\mathrm{Aut}|{=}1$, orb 8/8, hist-orb 2/2; M 6400, res 960, rev 2466: ok}\\*
0 & -- & 0 & 2 & 1 & 238 & 2 & 16 & \texttt{-+++????}\\*
1 & 0 & 1 & 2 & 1 & 220 & 4 & 32 & \texttt{--++????}\\
2 & 1 & 2 & 2 & 1 & 184 & 8 & 64 & \texttt{---+????}\\
3 & 2 & 3 & 2 & 1 & 112 & 16 & 128 & \texttt{----????}\\
4 & 3 & 4 & 1 & 1 & 192 & 64 & 0 & \texttt{-----++-}\\
5 & 4 & 5 & 1 & 1 & 128 & 128 & 0 & \texttt{------+-}\\
6 & 5 & 6 & 1 & 1 & 0 & 256 & 0 & \texttt{--------}\\
7 & 3 & 4 & 1 & 1 & 0 & 256 & 0 & \texttt{--------}\\
\multicolumn{9}{@{}l@{}}{\rule{0pt}{2.6ex}\textbf{T8.005}\enspace \texttt{01234563}\enspace $n{=}8$, $|H|{=}2$, len 7, br 1, $|\mathrm{Aut}|{=}1$, orb 8/8, hist-orb 2/2; M 6144, res 1120, rev 2242: ok}\\*
0 & -- & 0 & 2 & 1 & 222 & 2 & 32 & \texttt{-++?????}\\*
1 & 0 & 1 & 2 & 1 & 188 & 4 & 64 & \texttt{--+?????}\\
2 & 1 & 2 & 2 & 1 & 120 & 8 & 128 & \texttt{---?????}\\
3 & 2 & 3 & 1 & 1 & 224 & 32 & 0 & \texttt{----+++-}\\
4 & 3 & 4 & 1 & 1 & 192 & 64 & 0 & \texttt{-----++-}\\
5 & 4 & 5 & 1 & 1 & 128 & 128 & 0 & \texttt{------+-}\\
6 & 5 & 6 & 1 & 1 & 0 & 256 & 0 & \texttt{--------}\\
7 & 2 & 3 & 1 & 1 & 0 & 256 & 0 & \texttt{--------}\\
\multicolumn{9}{@{}l@{}}{\rule{0pt}{2.6ex}\textbf{T8.006}\enspace \texttt{01234562}\enspace $n{=}8$, $|H|{=}2$, len 7, br 1, $|\mathrm{Aut}|{=}1$, orb 8/8, hist-orb 2/2; M 5888, res 1152, rev 2058: ok}\\*
0 & -- & 0 & 2 & 1 & 190 & 2 & 64 & \texttt{-+??????}\\*
1 & 0 & 1 & 2 & 1 & 124 & 4 & 128 & \texttt{--??????}\\
2 & 1 & 2 & 1 & 1 & 240 & 16 & 0 & \texttt{---++++-}\\
3 & 2 & 3 & 1 & 1 & 224 & 32 & 0 & \texttt{----+++-}\\
4 & 3 & 4 & 1 & 1 & 192 & 64 & 0 & \texttt{-----++-}\\
5 & 4 & 5 & 1 & 1 & 128 & 128 & 0 & \texttt{------+-}\\
6 & 5 & 6 & 1 & 1 & 0 & 256 & 0 & \texttt{--------}\\
7 & 1 & 2 & 1 & 1 & 0 & 256 & 0 & \texttt{--------}\\
\multicolumn{9}{@{}l@{}}{\rule{0pt}{2.6ex}\textbf{T8.007}\enspace \texttt{01234561}\enspace $n{=}8$, $|H|{=}2$, len 7, br 1, $|\mathrm{Aut}|{=}1$, orb 8/8, hist-orb 2/2; M 5632, res 896, rev 1974: ok}\\*
0 & -- & 0 & 2 & 1 & 126 & 2 & 128 & \texttt{-???????}\\*
1 & 0 & 1 & 1 & 1 & 248 & 8 & 0 & \texttt{--+++++-}\\
2 & 1 & 2 & 1 & 1 & 240 & 16 & 0 & \texttt{---++++-}\\
3 & 2 & 3 & 1 & 1 & 224 & 32 & 0 & \texttt{----+++-}\\
4 & 3 & 4 & 1 & 1 & 192 & 64 & 0 & \texttt{-----++-}\\
5 & 4 & 5 & 1 & 1 & 128 & 128 & 0 & \texttt{------+-}\\
6 & 5 & 6 & 1 & 1 & 0 & 256 & 0 & \texttt{--------}\\
7 & 0 & 1 & 1 & 1 & 0 & 256 & 0 & \texttt{--------}\\
\multicolumn{9}{@{}l@{}}{\rule{0pt}{2.6ex}\textbf{T8.008}\enspace \texttt{01234555}\enspace $n{=}8$, $|H|{=}3$, len 6, br 1, $|\mathrm{Aut}|{=}6$, orb 6/5, hist-orb 1/0; M 6400, res 1116, rev 2682: ok}\\*
0 & -- & 0 & 3 & 1 & 242 & 2 & 12 & \texttt{-++++???}\\*
1 & 0 & 1 & 3 & 1 & 228 & 4 & 24 & \texttt{--+++???}\\
2 & 1 & 2 & 3 & 1 & 200 & 8 & 48 & \texttt{---++???}\\
3 & 2 & 3 & 3 & 1 & 144 & 16 & 96 & \texttt{----+???}\\
4 & 3 & 4 & 3 & 1 & 32 & 32 & 192 & \texttt{-----???}\\
5 & 4 & 5 & 1 & 3 & 0 & 256 & 0 & \texttt{--------}\\
6 & 4 & 5 & 1 & 3 & 0 & 256 & 0 & \texttt{--------}\\
7 & 4 & 5 & 1 & 3 & 0 & 256 & 0 & \texttt{--------}\\
\multicolumn{9}{@{}l@{}}{\rule{0pt}{2.6ex}\textbf{T8.009}\enspace \texttt{01234554}\enspace $n{=}8$, $|H|{=}3$, len 6, br 2, $|\mathrm{Aut}|{=}2$, orb 7/6, hist-orb 2/1; M 6144, res 1336, rev 2422: ok}\\*
0 & -- & 0 & 3 & 1 & 234 & 2 & 20 & \texttt{-+++????}\\*
1 & 0 & 1 & 3 & 1 & 212 & 4 & 40 & \texttt{--++????}\\
2 & 1 & 2 & 3 & 1 & 168 & 8 & 80 & \texttt{---+????}\\
3 & 2 & 3 & 3 & 1 & 80 & 16 & 160 & \texttt{----????}\\
4 & 3 & 4 & 2 & 1 & 64 & 64 & 128 & \texttt{-----??-}\\
5 & 4 & 5 & 1 & 2 & 0 & 256 & 0 & \texttt{--------}\\
6 & 4 & 5 & 1 & 2 & 0 & 256 & 0 & \texttt{--------}\\
7 & 3 & 4 & 1 & 1 & 0 & 256 & 0 & \texttt{--------}\\
\multicolumn{9}{@{}l@{}}{\rule{0pt}{2.6ex}\textbf{T8.010}\enspace \texttt{01234553}\enspace $n{=}8$, $|H|{=}3$, len 6, br 2, $|\mathrm{Aut}|{=}2$, orb 7/6, hist-orb 2/1; M 5888, res 1532, rev 2174: ok}\\*
0 & -- & 0 & 3 & 1 & 218 & 2 & 36 & \texttt{-++?????}\\*
1 & 0 & 1 & 3 & 1 & 180 & 4 & 72 & \texttt{--+?????}\\
2 & 1 & 2 & 3 & 1 & 104 & 8 & 144 & \texttt{---?????}\\
3 & 2 & 3 & 2 & 1 & 160 & 32 & 64 & \texttt{----+??-}\\
4 & 3 & 4 & 2 & 1 & 64 & 64 & 128 & \texttt{-----??-}\\
5 & 4 & 5 & 1 & 2 & 0 & 256 & 0 & \texttt{--------}\\
6 & 4 & 5 & 1 & 2 & 0 & 256 & 0 & \texttt{--------}\\
7 & 2 & 3 & 1 & 1 & 0 & 256 & 0 & \texttt{--------}\\
\multicolumn{9}{@{}l@{}}{\rule{0pt}{2.6ex}\textbf{T8.011}\enspace \texttt{01234552}\enspace $n{=}8$, $|H|{=}3$, len 6, br 2, $|\mathrm{Aut}|{=}2$, orb 7/6, hist-orb 2/1; M 5632, res 1600, rev 1990: ok}\\*
0 & -- & 0 & 3 & 1 & 186 & 2 & 68 & \texttt{-+??????}\\*
1 & 0 & 1 & 3 & 1 & 116 & 4 & 136 & \texttt{--??????}\\
2 & 1 & 2 & 2 & 1 & 208 & 16 & 32 & \texttt{---++??-}\\
3 & 2 & 3 & 2 & 1 & 160 & 32 & 64 & \texttt{----+??-}\\
4 & 3 & 4 & 2 & 1 & 64 & 64 & 128 & \texttt{-----??-}\\
5 & 4 & 5 & 1 & 2 & 0 & 256 & 0 & \texttt{--------}\\
6 & 4 & 5 & 1 & 2 & 0 & 256 & 0 & \texttt{--------}\\
7 & 1 & 2 & 1 & 1 & 0 & 256 & 0 & \texttt{--------}\\
\multicolumn{9}{@{}l@{}}{\rule{0pt}{2.6ex}\textbf{T8.012}\enspace \texttt{01234551}\enspace $n{=}8$, $|H|{=}3$, len 6, br 2, $|\mathrm{Aut}|{=}2$, orb 7/6, hist-orb 2/1; M 5376, res 1372, rev 1930: ok}\\*
0 & -- & 0 & 3 & 1 & 122 & 2 & 132 & \texttt{-???????}\\*
1 & 0 & 1 & 2 & 1 & 232 & 8 & 16 & \texttt{--+++??-}\\
2 & 1 & 2 & 2 & 1 & 208 & 16 & 32 & \texttt{---++??-}\\
3 & 2 & 3 & 2 & 1 & 160 & 32 & 64 & \texttt{----+??-}\\
4 & 3 & 4 & 2 & 1 & 64 & 64 & 128 & \texttt{-----??-}\\
5 & 4 & 5 & 1 & 2 & 0 & 256 & 0 & \texttt{--------}\\
6 & 4 & 5 & 1 & 2 & 0 & 256 & 0 & \texttt{--------}\\
7 & 0 & 1 & 1 & 1 & 0 & 256 & 0 & \texttt{--------}\\
\multicolumn{9}{@{}l@{}}{\rule{0pt}{2.6ex}\textbf{T8.013}\enspace \texttt{01234545}\enspace $n{=}8$, $|H|{=}2$, len 6, br 1, $|\mathrm{Aut}|{=}2$, orb 6/4, hist-orb 1/0; M 6144, res 720, rev 2654: ok}\\*
0 & -- & 0 & 2 & 1 & 242 & 2 & 12 & \texttt{-+++????}\\*
1 & 0 & 1 & 2 & 1 & 228 & 4 & 24 & \texttt{--++????}\\
2 & 1 & 2 & 2 & 1 & 200 & 8 & 48 & \texttt{---+????}\\
3 & 2 & 3 & 2 & 1 & 144 & 16 & 96 & \texttt{----????}\\
4 & 3 & 4 & 1 & 2 & 128 & 128 & 0 & \texttt{-----+--}\\
5 & 4 & 5 & 1 & 2 & 0 & 256 & 0 & \texttt{--------}\\
6 & 3 & 4 & 1 & 2 & 128 & 128 & 0 & \texttt{-------+}\\
7 & 6 & 5 & 1 & 2 & 0 & 256 & 0 & \texttt{--------}\\
\multicolumn{9}{@{}l@{}}{\rule{0pt}{2.6ex}\textbf{T8.014}\enspace \texttt{01234544}\enspace $n{=}8$, $|H|{=}3$, len 6, br 1, $|\mathrm{Aut}|{=}2$, orb 7/6, hist-orb 2/1; M 5888, res 1440, rev 2378: ok}\\*
0 & -- & 0 & 3 & 1 & 230 & 2 & 24 & \texttt{-+++????}\\*
1 & 0 & 1 & 3 & 1 & 204 & 4 & 48 & \texttt{--++????}\\
2 & 1 & 2 & 3 & 1 & 152 & 8 & 96 & \texttt{---+????}\\
3 & 2 & 3 & 3 & 1 & 48 & 16 & 192 & \texttt{----????}\\
4 & 3 & 4 & 1 & 1 & 128 & 128 & 0 & \texttt{-----+--}\\
5 & 4 & 5 & 1 & 1 & 0 & 256 & 0 & \texttt{--------}\\
6 & 3 & 4 & 1 & 2 & 0 & 256 & 0 & \texttt{--------}\\
7 & 3 & 4 & 1 & 2 & 0 & 256 & 0 & \texttt{--------}\\
\multicolumn{9}{@{}l@{}}{\rule{0pt}{2.6ex}\textbf{T8.015}\enspace \texttt{01234543}\enspace $n{=}8$, $|H|{=}3$, len 6, br 2, $|\mathrm{Aut}|{=}1$, orb 8/8, hist-orb 3/3; M 5632, res 1672, rev 2070: ok}\\*
0 & -- & 0 & 3 & 1 & 214 & 2 & 40 & \texttt{-++?????}\\*
1 & 0 & 1 & 3 & 1 & 172 & 4 & 80 & \texttt{--+?????}\\
2 & 1 & 2 & 3 & 1 & 88 & 8 & 160 & \texttt{---?????}\\
3 & 2 & 3 & 2 & 1 & 96 & 32 & 128 & \texttt{----???-}\\
4 & 3 & 4 & 1 & 1 & 128 & 128 & 0 & \texttt{-----+--}\\
5 & 4 & 5 & 1 & 1 & 0 & 256 & 0 & \texttt{--------}\\
6 & 3 & 4 & 1 & 1 & 0 & 256 & 0 & \texttt{--------}\\
7 & 2 & 3 & 1 & 1 & 0 & 256 & 0 & \texttt{--------}\\
\multicolumn{9}{@{}l@{}}{\rule{0pt}{2.6ex}\textbf{T8.016}\enspace \texttt{01234542}\enspace $n{=}8$, $|H|{=}3$, len 6, br 2, $|\mathrm{Aut}|{=}1$, orb 8/8, hist-orb 3/3; M 5376, res 1776, rev 1862: ok}\\*
0 & -- & 0 & 3 & 1 & 182 & 2 & 72 & \texttt{-+??????}\\*
1 & 0 & 1 & 3 & 1 & 108 & 4 & 144 & \texttt{--??????}\\
2 & 1 & 2 & 2 & 1 & 176 & 16 & 64 & \texttt{---+???-}\\
3 & 2 & 3 & 2 & 1 & 96 & 32 & 128 & \texttt{----???-}\\
4 & 3 & 4 & 1 & 1 & 128 & 128 & 0 & \texttt{-----+--}\\
5 & 4 & 5 & 1 & 1 & 0 & 256 & 0 & \texttt{--------}\\
6 & 3 & 4 & 1 & 1 & 0 & 256 & 0 & \texttt{--------}\\
7 & 1 & 2 & 1 & 1 & 0 & 256 & 0 & \texttt{--------}\\
\multicolumn{9}{@{}l@{}}{\rule{0pt}{2.6ex}\textbf{T8.017}\enspace \texttt{01234541}\enspace $n{=}8$, $|H|{=}3$, len 6, br 2, $|\mathrm{Aut}|{=}1$, orb 8/8, hist-orb 3/3; M 5120, res 1576, rev 1802: ok}\\*
0 & -- & 0 & 3 & 1 & 118 & 2 & 136 & \texttt{-???????}\\*
1 & 0 & 1 & 2 & 1 & 216 & 8 & 32 & \texttt{--++???-}\\
2 & 1 & 2 & 2 & 1 & 176 & 16 & 64 & \texttt{---+???-}\\
3 & 2 & 3 & 2 & 1 & 96 & 32 & 128 & \texttt{----???-}\\
4 & 3 & 4 & 1 & 1 & 128 & 128 & 0 & \texttt{-----+--}\\
5 & 4 & 5 & 1 & 1 & 0 & 256 & 0 & \texttt{--------}\\
6 & 3 & 4 & 1 & 1 & 0 & 256 & 0 & \texttt{--------}\\
7 & 0 & 1 & 1 & 1 & 0 & 256 & 0 & \texttt{--------}\\
\multicolumn{9}{@{}l@{}}{\rule{0pt}{2.6ex}\textbf{T8.018}\enspace \texttt{01234534}\enspace $n{=}8$, $|H|{=}2$, len 6, br 1, $|\mathrm{Aut}|{=}1$, orb 8/8, hist-orb 2/2; M 5632, res 700, rev 2418: ok}\\*
0 & -- & 0 & 2 & 1 & 234 & 2 & 20 & \texttt{-++?????}\\*
1 & 0 & 1 & 2 & 1 & 212 & 4 & 40 & \texttt{--+?????}\\
2 & 1 & 2 & 2 & 1 & 168 & 8 & 80 & \texttt{---?????}\\
3 & 2 & 3 & 1 & 1 & 192 & 64 & 0 & \texttt{----++--}\\
4 & 3 & 4 & 1 & 1 & 128 & 128 & 0 & \texttt{-----+--}\\
5 & 4 & 5 & 1 & 1 & 0 & 256 & 0 & \texttt{--------}\\
6 & 2 & 3 & 1 & 1 & 128 & 128 & 0 & \texttt{-------+}\\
7 & 6 & 4 & 1 & 1 & 0 & 256 & 0 & \texttt{--------}\\
\multicolumn{9}{@{}l@{}}{\rule{0pt}{2.6ex}\textbf{T8.019}\enspace \texttt{01234533}\enspace $n{=}8$, $|H|{=}3$, len 6, br 1, $|\mathrm{Aut}|{=}2$, orb 7/6, hist-orb 2/1; M 5376, res 1680, rev 1946: ok}\\*
0 & -- & 0 & 3 & 1 & 206 & 2 & 48 & \texttt{-++?????}\\*
1 & 0 & 1 & 3 & 1 & 156 & 4 & 96 & \texttt{--+?????}\\
2 & 1 & 2 & 3 & 1 & 56 & 8 & 192 & \texttt{---?????}\\
3 & 2 & 3 & 1 & 1 & 192 & 64 & 0 & \texttt{----++--}\\
4 & 3 & 4 & 1 & 1 & 128 & 128 & 0 & \texttt{-----+--}\\
5 & 4 & 5 & 1 & 1 & 0 & 256 & 0 & \texttt{--------}\\
6 & 2 & 3 & 1 & 2 & 0 & 256 & 0 & \texttt{--------}\\
7 & 2 & 3 & 1 & 2 & 0 & 256 & 0 & \texttt{--------}\\
\multicolumn{9}{@{}l@{}}{\rule{0pt}{2.6ex}\textbf{T8.020}\enspace \texttt{01234532}\enspace $n{=}8$, $|H|{=}3$, len 6, br 2, $|\mathrm{Aut}|{=}1$, orb 8/8, hist-orb 3/3; M 5120, res 1856, rev 1678: ok}\\*
0 & -- & 0 & 3 & 1 & 174 & 2 & 80 & \texttt{-+??????}\\*
1 & 0 & 1 & 3 & 1 & 92 & 4 & 160 & \texttt{--??????}\\
2 & 1 & 2 & 2 & 1 & 112 & 16 & 128 & \texttt{---????-}\\
3 & 2 & 3 & 1 & 1 & 192 & 64 & 0 & \texttt{----++--}\\
4 & 3 & 4 & 1 & 1 & 128 & 128 & 0 & \texttt{-----+--}\\
5 & 4 & 5 & 1 & 1 & 0 & 256 & 0 & \texttt{--------}\\
6 & 2 & 3 & 1 & 1 & 0 & 256 & 0 & \texttt{--------}\\
7 & 1 & 2 & 1 & 1 & 0 & 256 & 0 & \texttt{--------}\\
\multicolumn{9}{@{}l@{}}{\rule{0pt}{2.6ex}\textbf{T8.021}\enspace \texttt{01234531}\enspace $n{=}8$, $|H|{=}3$, len 6, br 2, $|\mathrm{Aut}|{=}1$, orb 8/8, hist-orb 3/3; M 4864, res 1712, rev 1594: ok}\\*
0 & -- & 0 & 3 & 1 & 110 & 2 & 144 & \texttt{-???????}\\*
1 & 0 & 1 & 2 & 1 & 184 & 8 & 64 & \texttt{--+????-}\\
2 & 1 & 2 & 2 & 1 & 112 & 16 & 128 & \texttt{---????-}\\
3 & 2 & 3 & 1 & 1 & 192 & 64 & 0 & \texttt{----++--}\\
4 & 3 & 4 & 1 & 1 & 128 & 128 & 0 & \texttt{-----+--}\\
5 & 4 & 5 & 1 & 1 & 0 & 256 & 0 & \texttt{--------}\\
6 & 2 & 3 & 1 & 1 & 0 & 256 & 0 & \texttt{--------}\\
7 & 0 & 1 & 1 & 1 & 0 & 256 & 0 & \texttt{--------}\\
\multicolumn{9}{@{}l@{}}{\rule{0pt}{2.6ex}\textbf{T8.022}\enspace \texttt{01234523}\enspace $n{=}8$, $|H|{=}2$, len 6, br 1, $|\mathrm{Aut}|{=}1$, orb 8/8, hist-orb 2/2; M 5120, res 648, rev 2126: ok}\\*
0 & -- & 0 & 2 & 1 & 218 & 2 & 36 & \texttt{-+??????}\\*
1 & 0 & 1 & 2 & 1 & 180 & 4 & 72 & \texttt{--??????}\\
2 & 1 & 2 & 1 & 1 & 224 & 32 & 0 & \texttt{---+++--}\\
3 & 2 & 3 & 1 & 1 & 192 & 64 & 0 & \texttt{----++--}\\
4 & 3 & 4 & 1 & 1 & 128 & 128 & 0 & \texttt{-----+--}\\
5 & 4 & 5 & 1 & 1 & 0 & 256 & 0 & \texttt{--------}\\
6 & 1 & 2 & 1 & 1 & 128 & 128 & 0 & \texttt{-------+}\\
7 & 6 & 3 & 1 & 1 & 0 & 256 & 0 & \texttt{--------}\\
\multicolumn{9}{@{}l@{}}{\rule{0pt}{2.6ex}\textbf{T8.023}\enspace \texttt{01234522}\enspace $n{=}8$, $|H|{=}3$, len 6, br 1, $|\mathrm{Aut}|{=}2$, orb 7/6, hist-orb 2/1; M 4864, res 1728, rev 1570: ok}\\*
0 & -- & 0 & 3 & 1 & 158 & 2 & 96 & \texttt{-+??????}\\*
1 & 0 & 1 & 3 & 1 & 60 & 4 & 192 & \texttt{--??????}\\
2 & 1 & 2 & 1 & 1 & 224 & 32 & 0 & \texttt{---+++--}\\
3 & 2 & 3 & 1 & 1 & 192 & 64 & 0 & \texttt{----++--}\\
4 & 3 & 4 & 1 & 1 & 128 & 128 & 0 & \texttt{-----+--}\\
5 & 4 & 5 & 1 & 1 & 0 & 256 & 0 & \texttt{--------}\\
6 & 1 & 2 & 1 & 2 & 0 & 256 & 0 & \texttt{--------}\\
7 & 1 & 2 & 1 & 2 & 0 & 256 & 0 & \texttt{--------}\\
\multicolumn{9}{@{}l@{}}{\rule{0pt}{2.6ex}\textbf{T8.024}\enspace \texttt{01234521}\enspace $n{=}8$, $|H|{=}3$, len 6, br 2, $|\mathrm{Aut}|{=}1$, orb 8/8, hist-orb 3/3; M 4608, res 1696, rev 1426: ok}\\*
0 & -- & 0 & 3 & 1 & 94 & 2 & 160 & \texttt{-???????}\\*
1 & 0 & 1 & 2 & 1 & 120 & 8 & 128 & \texttt{--?????-}\\
2 & 1 & 2 & 1 & 1 & 224 & 32 & 0 & \texttt{---+++--}\\
3 & 2 & 3 & 1 & 1 & 192 & 64 & 0 & \texttt{----++--}\\
4 & 3 & 4 & 1 & 1 & 128 & 128 & 0 & \texttt{-----+--}\\
5 & 4 & 5 & 1 & 1 & 0 & 256 & 0 & \texttt{--------}\\
6 & 1 & 2 & 1 & 1 & 0 & 256 & 0 & \texttt{--------}\\
7 & 0 & 1 & 1 & 1 & 0 & 256 & 0 & \texttt{--------}\\
\multicolumn{9}{@{}l@{}}{\rule{0pt}{2.6ex}\textbf{T8.025}\enspace \texttt{01234512}\enspace $n{=}8$, $|H|{=}2$, len 6, br 1, $|\mathrm{Aut}|{=}1$, orb 8/8, hist-orb 2/2; M 4608, res 476, rev 1870: ok}\\*
0 & -- & 0 & 2 & 1 & 186 & 2 & 68 & \texttt{-???????}\\*
1 & 0 & 1 & 1 & 1 & 240 & 16 & 0 & \texttt{--++++--}\\
2 & 1 & 2 & 1 & 1 & 224 & 32 & 0 & \texttt{---+++--}\\
3 & 2 & 3 & 1 & 1 & 192 & 64 & 0 & \texttt{----++--}\\
4 & 3 & 4 & 1 & 1 & 128 & 128 & 0 & \texttt{-----+--}\\
5 & 4 & 5 & 1 & 1 & 0 & 256 & 0 & \texttt{--------}\\
6 & 0 & 1 & 1 & 1 & 128 & 128 & 0 & \texttt{-------+}\\
7 & 6 & 2 & 1 & 1 & 0 & 256 & 0 & \texttt{--------}\\
\multicolumn{9}{@{}l@{}}{\rule{0pt}{2.6ex}\textbf{T8.026}\enspace \texttt{01234511}\enspace $n{=}8$, $|H|{=}3$, len 6, br 1, $|\mathrm{Aut}|{=}2$, orb 7/6, hist-orb 2/1; M 4352, res 1344, rev 1390: ok}\\*
0 & -- & 0 & 3 & 1 & 62 & 2 & 192 & \texttt{-???????}\\*
1 & 0 & 1 & 1 & 1 & 240 & 16 & 0 & \texttt{--++++--}\\
2 & 1 & 2 & 1 & 1 & 224 & 32 & 0 & \texttt{---+++--}\\
3 & 2 & 3 & 1 & 1 & 192 & 64 & 0 & \texttt{----++--}\\
4 & 3 & 4 & 1 & 1 & 128 & 128 & 0 & \texttt{-----+--}\\
5 & 4 & 5 & 1 & 1 & 0 & 256 & 0 & \texttt{--------}\\
6 & 0 & 1 & 1 & 2 & 0 & 256 & 0 & \texttt{--------}\\
7 & 0 & 1 & 1 & 2 & 0 & 256 & 0 & \texttt{--------}\\
\multicolumn{9}{@{}l@{}}{\rule{0pt}{2.6ex}\textbf{T8.027}\enspace \texttt{01234444}\enspace $n{=}8$, $|H|{=}4$, len 5, br 1, $|\mathrm{Aut}|{=}24$, orb 5/4, hist-orb 1/0; M 5632, res 1680, rev 2342: ok}\\*
0 & -- & 0 & 4 & 1 & 226 & 2 & 28 & \texttt{-+++????}\\*
1 & 0 & 1 & 4 & 1 & 196 & 4 & 56 & \texttt{--++????}\\
2 & 1 & 2 & 4 & 1 & 136 & 8 & 112 & \texttt{---+????}\\
3 & 2 & 3 & 4 & 1 & 16 & 16 & 224 & \texttt{----????}\\
4 & 3 & 4 & 1 & 4 & 0 & 256 & 0 & \texttt{--------}\\
5 & 3 & 4 & 1 & 4 & 0 & 256 & 0 & \texttt{--------}\\
6 & 3 & 4 & 1 & 4 & 0 & 256 & 0 & \texttt{--------}\\
7 & 3 & 4 & 1 & 4 & 0 & 256 & 0 & \texttt{--------}\\
\multicolumn{9}{@{}l@{}}{\rule{0pt}{2.6ex}\textbf{T8.028}\enspace \texttt{01234443}\enspace $n{=}8$, $|H|{=}4$, len 5, br 2, $|\mathrm{Aut}|{=}6$, orb 6/5, hist-orb 2/1; M 5376, res 1948, rev 1946: ok}\\*
0 & -- & 0 & 4 & 1 & 210 & 2 & 44 & \texttt{-++?????}\\*
1 & 0 & 1 & 4 & 1 & 164 & 4 & 88 & \texttt{--+?????}\\
2 & 1 & 2 & 4 & 1 & 72 & 8 & 176 & \texttt{---?????}\\
3 & 2 & 3 & 3 & 1 & 32 & 32 & 192 & \texttt{----???-}\\
4 & 3 & 4 & 1 & 3 & 0 & 256 & 0 & \texttt{--------}\\
5 & 3 & 4 & 1 & 3 & 0 & 256 & 0 & \texttt{--------}\\
6 & 3 & 4 & 1 & 3 & 0 & 256 & 0 & \texttt{--------}\\
7 & 2 & 3 & 1 & 1 & 0 & 256 & 0 & \texttt{--------}\\
\multicolumn{9}{@{}l@{}}{\rule{0pt}{2.6ex}\textbf{T8.029}\enspace \texttt{01234442}\enspace $n{=}8$, $|H|{=}4$, len 5, br 2, $|\mathrm{Aut}|{=}6$, orb 6/5, hist-orb 2/1; M 5120, res 2088, rev 1706: ok}\\*
0 & -- & 0 & 4 & 1 & 178 & 2 & 76 & \texttt{-+??????}\\*
1 & 0 & 1 & 4 & 1 & 100 & 4 & 152 & \texttt{--??????}\\
2 & 1 & 2 & 3 & 1 & 144 & 16 & 96 & \texttt{---+???-}\\
3 & 2 & 3 & 3 & 1 & 32 & 32 & 192 & \texttt{----???-}\\
4 & 3 & 4 & 1 & 3 & 0 & 256 & 0 & \texttt{--------}\\
5 & 3 & 4 & 1 & 3 & 0 & 256 & 0 & \texttt{--------}\\
6 & 3 & 4 & 1 & 3 & 0 & 256 & 0 & \texttt{--------}\\
7 & 1 & 2 & 1 & 1 & 0 & 256 & 0 & \texttt{--------}\\
\multicolumn{9}{@{}l@{}}{\rule{0pt}{2.6ex}\textbf{T8.030}\enspace \texttt{01234441}\enspace $n{=}8$, $|H|{=}4$, len 5, br 2, $|\mathrm{Aut}|{=}6$, orb 6/5, hist-orb 2/1; M 4864, res 1916, rev 1654: ok}\\*
0 & -- & 0 & 4 & 1 & 114 & 2 & 140 & \texttt{-???????}\\*
1 & 0 & 1 & 3 & 1 & 200 & 8 & 48 & \texttt{--++???-}\\
2 & 1 & 2 & 3 & 1 & 144 & 16 & 96 & \texttt{---+???-}\\
3 & 2 & 3 & 3 & 1 & 32 & 32 & 192 & \texttt{----???-}\\
4 & 3 & 4 & 1 & 3 & 0 & 256 & 0 & \texttt{--------}\\
5 & 3 & 4 & 1 & 3 & 0 & 256 & 0 & \texttt{--------}\\
6 & 3 & 4 & 1 & 3 & 0 & 256 & 0 & \texttt{--------}\\
7 & 0 & 1 & 1 & 1 & 0 & 256 & 0 & \texttt{--------}\\
\multicolumn{9}{@{}l@{}}{\rule{0pt}{2.6ex}\textbf{T8.031}\enspace \texttt{01234434}\enspace $n{=}8$, $|H|{=}3$, len 5, br 2, $|\mathrm{Aut}|{=}2$, orb 7/6, hist-orb 2/1; M 5376, res 1264, rev 2210: ok}\\*
0 & -- & 0 & 3 & 1 & 222 & 2 & 32 & \texttt{-++?????}\\*
1 & 0 & 1 & 3 & 1 & 188 & 4 & 64 & \texttt{--+?????}\\
2 & 1 & 2 & 3 & 1 & 120 & 8 & 128 & \texttt{---?????}\\
3 & 2 & 3 & 2 & 1 & 64 & 64 & 128 & \texttt{----??--}\\
4 & 3 & 4 & 1 & 2 & 0 & 256 & 0 & \texttt{--------}\\
5 & 3 & 4 & 1 & 2 & 0 & 256 & 0 & \texttt{--------}\\
6 & 2 & 3 & 1 & 1 & 128 & 128 & 0 & \texttt{-------+}\\
7 & 6 & 4 & 1 & 1 & 0 & 256 & 0 & \texttt{--------}\\
\multicolumn{9}{@{}l@{}}{\rule{0pt}{2.6ex}\textbf{T8.032}\enspace \texttt{01234433}\enspace $n{=}8$, $|H|{=}4$, len 5, br 2, $|\mathrm{Aut}|{=}4$, orb 6/4, hist-orb 2/0; M 5120, res 1964, rev 1794: ok}\\*
0 & -- & 0 & 4 & 1 & 202 & 2 & 52 & \texttt{-++?????}\\*
1 & 0 & 1 & 4 & 1 & 148 & 4 & 104 & \texttt{--+?????}\\
2 & 1 & 2 & 4 & 1 & 40 & 8 & 208 & \texttt{---?????}\\
3 & 2 & 3 & 2 & 1 & 64 & 64 & 128 & \texttt{----??--}\\
4 & 3 & 4 & 1 & 2 & 0 & 256 & 0 & \texttt{--------}\\
5 & 3 & 4 & 1 & 2 & 0 & 256 & 0 & \texttt{--------}\\
6 & 2 & 3 & 1 & 2 & 0 & 256 & 0 & \texttt{--------}\\
7 & 2 & 3 & 1 & 2 & 0 & 256 & 0 & \texttt{--------}\\
\multicolumn{9}{@{}l@{}}{\rule{0pt}{2.6ex}\textbf{T8.033}\enspace \texttt{01234432}\enspace $n{=}8$, $|H|{=}4$, len 5, br 3, $|\mathrm{Aut}|{=}2$, orb 7/6, hist-orb 3/2; M 4864, res 2176, rev 1478: ok}\\*
0 & -- & 0 & 4 & 1 & 170 & 2 & 84 & \texttt{-+??????}\\*
1 & 0 & 1 & 4 & 1 & 84 & 4 & 168 & \texttt{--??????}\\
2 & 1 & 2 & 3 & 1 & 80 & 16 & 160 & \texttt{---????-}\\
3 & 2 & 3 & 2 & 1 & 64 & 64 & 128 & \texttt{----??--}\\
4 & 3 & 4 & 1 & 2 & 0 & 256 & 0 & \texttt{--------}\\
5 & 3 & 4 & 1 & 2 & 0 & 256 & 0 & \texttt{--------}\\
6 & 2 & 3 & 1 & 1 & 0 & 256 & 0 & \texttt{--------}\\
7 & 1 & 2 & 1 & 1 & 0 & 256 & 0 & \texttt{--------}\\
\multicolumn{9}{@{}l@{}}{\rule{0pt}{2.6ex}\textbf{T8.034}\enspace \texttt{01234431}\enspace $n{=}8$, $|H|{=}4$, len 5, br 3, $|\mathrm{Aut}|{=}2$, orb 7/6, hist-orb 3/2; M 4608, res 2060, rev 1394: ok}\\*
0 & -- & 0 & 4 & 1 & 106 & 2 & 148 & \texttt{-???????}\\*
1 & 0 & 1 & 3 & 1 & 168 & 8 & 80 & \texttt{--+????-}\\
2 & 1 & 2 & 3 & 1 & 80 & 16 & 160 & \texttt{---????-}\\
3 & 2 & 3 & 2 & 1 & 64 & 64 & 128 & \texttt{----??--}\\
4 & 3 & 4 & 1 & 2 & 0 & 256 & 0 & \texttt{--------}\\
5 & 3 & 4 & 1 & 2 & 0 & 256 & 0 & \texttt{--------}\\
6 & 2 & 3 & 1 & 1 & 0 & 256 & 0 & \texttt{--------}\\
7 & 0 & 1 & 1 & 1 & 0 & 256 & 0 & \texttt{--------}\\
\multicolumn{9}{@{}l@{}}{\rule{0pt}{2.6ex}\textbf{T8.035}\enspace \texttt{01234423}\enspace $n{=}8$, $|H|{=}3$, len 5, br 2, $|\mathrm{Aut}|{=}2$, orb 7/6, hist-orb 2/1; M 4864, res 1152, rev 1958: ok}\\*
0 & -- & 0 & 3 & 1 & 206 & 2 & 48 & \texttt{-+??????}\\*
1 & 0 & 1 & 3 & 1 & 156 & 4 & 96 & \texttt{--??????}\\
2 & 1 & 2 & 2 & 1 & 160 & 32 & 64 & \texttt{---+??--}\\
3 & 2 & 3 & 2 & 1 & 64 & 64 & 128 & \texttt{----??--}\\
4 & 3 & 4 & 1 & 2 & 0 & 256 & 0 & \texttt{--------}\\
5 & 3 & 4 & 1 & 2 & 0 & 256 & 0 & \texttt{--------}\\
6 & 1 & 2 & 1 & 1 & 128 & 128 & 0 & \texttt{-------+}\\
7 & 6 & 3 & 1 & 1 & 0 & 256 & 0 & \texttt{--------}\\
\multicolumn{9}{@{}l@{}}{\rule{0pt}{2.6ex}\textbf{T8.036}\enspace \texttt{01234422}\enspace $n{=}8$, $|H|{=}4$, len 5, br 2, $|\mathrm{Aut}|{=}4$, orb 6/4, hist-orb 2/0; M 4608, res 2088, rev 1378: ok}\\*
0 & -- & 0 & 4 & 1 & 154 & 2 & 100 & \texttt{-+??????}\\*
1 & 0 & 1 & 4 & 1 & 52 & 4 & 200 & \texttt{--??????}\\
2 & 1 & 2 & 2 & 1 & 160 & 32 & 64 & \texttt{---+??--}\\
3 & 2 & 3 & 2 & 1 & 64 & 64 & 128 & \texttt{----??--}\\
4 & 3 & 4 & 1 & 2 & 0 & 256 & 0 & \texttt{--------}\\
5 & 3 & 4 & 1 & 2 & 0 & 256 & 0 & \texttt{--------}\\
6 & 1 & 2 & 1 & 2 & 0 & 256 & 0 & \texttt{--------}\\
7 & 1 & 2 & 1 & 2 & 0 & 256 & 0 & \texttt{--------}\\
\multicolumn{9}{@{}l@{}}{\rule{0pt}{2.6ex}\textbf{T8.037}\enspace \texttt{01234421}\enspace $n{=}8$, $|H|{=}4$, len 5, br 3, $|\mathrm{Aut}|{=}2$, orb 7/6, hist-orb 3/2; M 4352, res 2084, rev 1218: ok}\\*
0 & -- & 0 & 4 & 1 & 90 & 2 & 164 & \texttt{-???????}\\*
1 & 0 & 1 & 3 & 1 & 104 & 8 & 144 & \texttt{--?????-}\\
2 & 1 & 2 & 2 & 1 & 160 & 32 & 64 & \texttt{---+??--}\\
3 & 2 & 3 & 2 & 1 & 64 & 64 & 128 & \texttt{----??--}\\
4 & 3 & 4 & 1 & 2 & 0 & 256 & 0 & \texttt{--------}\\
5 & 3 & 4 & 1 & 2 & 0 & 256 & 0 & \texttt{--------}\\
6 & 1 & 2 & 1 & 1 & 0 & 256 & 0 & \texttt{--------}\\
7 & 0 & 1 & 1 & 1 & 0 & 256 & 0 & \texttt{--------}\\
\multicolumn{9}{@{}l@{}}{\rule{0pt}{2.6ex}\textbf{T8.038}\enspace \texttt{01234412}\enspace $n{=}8$, $|H|{=}3$, len 5, br 2, $|\mathrm{Aut}|{=}2$, orb 7/6, hist-orb 2/1; M 4352, res 960, rev 1734: ok}\\*
0 & -- & 0 & 3 & 1 & 174 & 2 & 80 & \texttt{-???????}\\*
1 & 0 & 1 & 2 & 1 & 208 & 16 & 32 & \texttt{--++??--}\\
2 & 1 & 2 & 2 & 1 & 160 & 32 & 64 & \texttt{---+??--}\\
3 & 2 & 3 & 2 & 1 & 64 & 64 & 128 & \texttt{----??--}\\
4 & 3 & 4 & 1 & 2 & 0 & 256 & 0 & \texttt{--------}\\
5 & 3 & 4 & 1 & 2 & 0 & 256 & 0 & \texttt{--------}\\
6 & 0 & 1 & 1 & 1 & 128 & 128 & 0 & \texttt{-------+}\\
7 & 6 & 2 & 1 & 1 & 0 & 256 & 0 & \texttt{--------}\\
\multicolumn{9}{@{}l@{}}{\rule{0pt}{2.6ex}\textbf{T8.039}\enspace \texttt{01234411}\enspace $n{=}8$, $|H|{=}4$, len 5, br 2, $|\mathrm{Aut}|{=}4$, orb 6/4, hist-orb 2/0; M 4096, res 1772, rev 1214: ok}\\*
0 & -- & 0 & 4 & 1 & 58 & 2 & 196 & \texttt{-???????}\\*
1 & 0 & 1 & 2 & 1 & 208 & 16 & 32 & \texttt{--++??--}\\
2 & 1 & 2 & 2 & 1 & 160 & 32 & 64 & \texttt{---+??--}\\
3 & 2 & 3 & 2 & 1 & 64 & 64 & 128 & \texttt{----??--}\\
4 & 3 & 4 & 1 & 2 & 0 & 256 & 0 & \texttt{--------}\\
5 & 3 & 4 & 1 & 2 & 0 & 256 & 0 & \texttt{--------}\\
6 & 0 & 1 & 1 & 2 & 0 & 256 & 0 & \texttt{--------}\\
7 & 0 & 1 & 1 & 2 & 0 & 256 & 0 & \texttt{--------}\\
\multicolumn{9}{@{}l@{}}{\rule{0pt}{2.6ex}\textbf{T8.040}\enspace \texttt{01234343}\enspace $n{=}8$, $|H|{=}3$, len 5, br 1, $|\mathrm{Aut}|{=}2$, orb 6/4, hist-orb 2/1; M 5120, res 1540, rev 2010: ok}\\*
0 & -- & 0 & 3 & 1 & 210 & 2 & 44 & \texttt{-++?????}\\*
1 & 0 & 1 & 3 & 1 & 164 & 4 & 88 & \texttt{--+?????}\\
2 & 1 & 2 & 3 & 1 & 72 & 8 & 176 & \texttt{---?????}\\
3 & 2 & 3 & 1 & 2 & 128 & 128 & 0 & \texttt{----+---}\\
4 & 3 & 4 & 1 & 2 & 0 & 256 & 0 & \texttt{--------}\\
5 & 2 & 3 & 1 & 2 & 128 & 128 & 0 & \texttt{------+-}\\
6 & 5 & 4 & 1 & 2 & 0 & 256 & 0 & \texttt{--------}\\
7 & 2 & 3 & 1 & 1 & 0 & 256 & 0 & \texttt{--------}\\
\multicolumn{9}{@{}l@{}}{\rule{0pt}{2.6ex}\textbf{T8.041}\enspace \texttt{01234342}\enspace $n{=}8$, $|H|{=}3$, len 5, br 2, $|\mathrm{Aut}|{=}2$, orb 6/4, hist-orb 2/1; M 4864, res 1680, rev 1750: ok}\\*
0 & -- & 0 & 3 & 1 & 178 & 2 & 76 & \texttt{-+??????}\\*
1 & 0 & 1 & 3 & 1 & 100 & 4 & 152 & \texttt{--??????}\\
2 & 1 & 2 & 2 & 1 & 144 & 16 & 96 & \texttt{---????-}\\
3 & 2 & 3 & 1 & 2 & 128 & 128 & 0 & \texttt{----+---}\\
4 & 3 & 4 & 1 & 2 & 0 & 256 & 0 & \texttt{--------}\\
5 & 2 & 3 & 1 & 2 & 128 & 128 & 0 & \texttt{------+-}\\
6 & 5 & 4 & 1 & 2 & 0 & 256 & 0 & \texttt{--------}\\
7 & 1 & 2 & 1 & 1 & 0 & 256 & 0 & \texttt{--------}\\
\multicolumn{9}{@{}l@{}}{\rule{0pt}{2.6ex}\textbf{T8.042}\enspace \texttt{01234341}\enspace $n{=}8$, $|H|{=}3$, len 5, br 2, $|\mathrm{Aut}|{=}2$, orb 6/4, hist-orb 2/1; M 4608, res 1508, rev 1682: ok}\\*
0 & -- & 0 & 3 & 1 & 114 & 2 & 140 & \texttt{-???????}\\*
1 & 0 & 1 & 2 & 1 & 200 & 8 & 48 & \texttt{--+????-}\\
2 & 1 & 2 & 2 & 1 & 144 & 16 & 96 & \texttt{---????-}\\
3 & 2 & 3 & 1 & 2 & 128 & 128 & 0 & \texttt{----+---}\\
4 & 3 & 4 & 1 & 2 & 0 & 256 & 0 & \texttt{--------}\\
5 & 2 & 3 & 1 & 2 & 128 & 128 & 0 & \texttt{------+-}\\
6 & 5 & 4 & 1 & 2 & 0 & 256 & 0 & \texttt{--------}\\
7 & 0 & 1 & 1 & 1 & 0 & 256 & 0 & \texttt{--------}\\
\multicolumn{9}{@{}l@{}}{\rule{0pt}{2.6ex}\textbf{T8.043}\enspace \texttt{01234333}\enspace $n{=}8$, $|H|{=}4$, len 5, br 1, $|\mathrm{Aut}|{=}6$, orb 6/5, hist-orb 2/1; M 4864, res 1960, rev 1762: ok}\\*
0 & -- & 0 & 4 & 1 & 198 & 2 & 56 & \texttt{-++?????}\\*
1 & 0 & 1 & 4 & 1 & 140 & 4 & 112 & \texttt{--+?????}\\
2 & 1 & 2 & 4 & 1 & 24 & 8 & 224 & \texttt{---?????}\\
3 & 2 & 3 & 1 & 1 & 128 & 128 & 0 & \texttt{----+---}\\
4 & 3 & 4 & 1 & 1 & 0 & 256 & 0 & \texttt{--------}\\
5 & 2 & 3 & 1 & 3 & 0 & 256 & 0 & \texttt{--------}\\
6 & 2 & 3 & 1 & 3 & 0 & 256 & 0 & \texttt{--------}\\
7 & 2 & 3 & 1 & 3 & 0 & 256 & 0 & \texttt{--------}\\
\multicolumn{9}{@{}l@{}}{\rule{0pt}{2.6ex}\textbf{T8.044}\enspace \texttt{01234332}\enspace $n{=}8$, $|H|{=}4$, len 5, br 2, $|\mathrm{Aut}|{=}2$, orb 7/6, hist-orb 3/2; M 4608, res 2208, rev 1390: ok}\\*
0 & -- & 0 & 4 & 1 & 166 & 2 & 88 & \texttt{-+??????}\\*
1 & 0 & 1 & 4 & 1 & 76 & 4 & 176 & \texttt{--??????}\\
2 & 1 & 2 & 3 & 1 & 48 & 16 & 192 & \texttt{---????-}\\
3 & 2 & 3 & 1 & 1 & 128 & 128 & 0 & \texttt{----+---}\\
4 & 3 & 4 & 1 & 1 & 0 & 256 & 0 & \texttt{--------}\\
5 & 2 & 3 & 1 & 2 & 0 & 256 & 0 & \texttt{--------}\\
6 & 2 & 3 & 1 & 2 & 0 & 256 & 0 & \texttt{--------}\\
7 & 1 & 2 & 1 & 1 & 0 & 256 & 0 & \texttt{--------}\\
\multicolumn{9}{@{}l@{}}{\rule{0pt}{2.6ex}\textbf{T8.045}\enspace \texttt{01234331}\enspace $n{=}8$, $|H|{=}4$, len 5, br 2, $|\mathrm{Aut}|{=}2$, orb 7/6, hist-orb 3/2; M 4352, res 2120, rev 1290: ok}\\*
0 & -- & 0 & 4 & 1 & 102 & 2 & 152 & \texttt{-???????}\\*
1 & 0 & 1 & 3 & 1 & 152 & 8 & 96 & \texttt{--+????-}\\
2 & 1 & 2 & 3 & 1 & 48 & 16 & 192 & \texttt{---????-}\\
3 & 2 & 3 & 1 & 1 & 128 & 128 & 0 & \texttt{----+---}\\
4 & 3 & 4 & 1 & 1 & 0 & 256 & 0 & \texttt{--------}\\
5 & 2 & 3 & 1 & 2 & 0 & 256 & 0 & \texttt{--------}\\
6 & 2 & 3 & 1 & 2 & 0 & 256 & 0 & \texttt{--------}\\
7 & 0 & 1 & 1 & 1 & 0 & 256 & 0 & \texttt{--------}\\
\multicolumn{9}{@{}l@{}}{\rule{0pt}{2.6ex}\textbf{T8.046}\enspace \texttt{01234323}\enspace $n{=}8$, $|H|{=}3$, len 5, br 2, $|\mathrm{Aut}|{=}1$, orb 8/8, hist-orb 3/3; M 4608, res 1368, rev 1778: ok}\\*
0 & -- & 0 & 3 & 1 & 194 & 2 & 60 & \texttt{-+??????}\\*
1 & 0 & 1 & 3 & 1 & 132 & 4 & 120 & \texttt{--??????}\\
2 & 1 & 2 & 2 & 1 & 96 & 32 & 128 & \texttt{---???--}\\
3 & 2 & 3 & 1 & 1 & 128 & 128 & 0 & \texttt{----+---}\\
4 & 3 & 4 & 1 & 1 & 0 & 256 & 0 & \texttt{--------}\\
5 & 2 & 3 & 1 & 1 & 0 & 256 & 0 & \texttt{--------}\\
6 & 1 & 2 & 1 & 1 & 128 & 128 & 0 & \texttt{-------+}\\
7 & 6 & 3 & 1 & 1 & 0 & 256 & 0 & \texttt{--------}\\
\multicolumn{9}{@{}l@{}}{\rule{0pt}{2.6ex}\textbf{T8.047}\enspace \texttt{01234322}\enspace $n{=}8$, $|H|{=}4$, len 5, br 2, $|\mathrm{Aut}|{=}2$, orb 7/6, hist-orb 3/2; M 4352, res 2160, rev 1246: ok}\\*
0 & -- & 0 & 4 & 1 & 150 & 2 & 104 & \texttt{-+??????}\\*
1 & 0 & 1 & 4 & 1 & 44 & 4 & 208 & \texttt{--??????}\\
2 & 1 & 2 & 2 & 1 & 96 & 32 & 128 & \texttt{---???--}\\
3 & 2 & 3 & 1 & 1 & 128 & 128 & 0 & \texttt{----+---}\\
4 & 3 & 4 & 1 & 1 & 0 & 256 & 0 & \texttt{--------}\\
5 & 2 & 3 & 1 & 1 & 0 & 256 & 0 & \texttt{--------}\\
6 & 1 & 2 & 1 & 2 & 0 & 256 & 0 & \texttt{--------}\\
7 & 1 & 2 & 1 & 2 & 0 & 256 & 0 & \texttt{--------}\\
\multicolumn{9}{@{}l@{}}{\rule{0pt}{2.6ex}\textbf{T8.048}\enspace \texttt{01234321}\enspace $n{=}8$, $|H|{=}4$, len 5, br 3, $|\mathrm{Aut}|{=}1$, orb 8/8, hist-orb 4/4; M 4096, res 2184, rev 1062: ok}\\*
0 & -- & 0 & 4 & 1 & 86 & 2 & 168 & \texttt{-???????}\\*
1 & 0 & 1 & 3 & 1 & 88 & 8 & 160 & \texttt{--?????-}\\
2 & 1 & 2 & 2 & 1 & 96 & 32 & 128 & \texttt{---???--}\\
3 & 2 & 3 & 1 & 1 & 128 & 128 & 0 & \texttt{----+---}\\
4 & 3 & 4 & 1 & 1 & 0 & 256 & 0 & \texttt{--------}\\
5 & 2 & 3 & 1 & 1 & 0 & 256 & 0 & \texttt{--------}\\
6 & 1 & 2 & 1 & 1 & 0 & 256 & 0 & \texttt{--------}\\
7 & 0 & 1 & 1 & 1 & 0 & 256 & 0 & \texttt{--------}\\
\multicolumn{9}{@{}l@{}}{\rule{0pt}{2.6ex}\textbf{T8.049}\enspace \texttt{01234312}\enspace $n{=}8$, $|H|{=}3$, len 5, br 2, $|\mathrm{Aut}|{=}1$, orb 8/8, hist-orb 3/3; M 4096, res 1156, rev 1570: ok}\\*
0 & -- & 0 & 3 & 1 & 162 & 2 & 92 & \texttt{-???????}\\*
1 & 0 & 1 & 2 & 1 & 176 & 16 & 64 & \texttt{--+???--}\\
2 & 1 & 2 & 2 & 1 & 96 & 32 & 128 & \texttt{---???--}\\
3 & 2 & 3 & 1 & 1 & 128 & 128 & 0 & \texttt{----+---}\\
4 & 3 & 4 & 1 & 1 & 0 & 256 & 0 & \texttt{--------}\\
5 & 2 & 3 & 1 & 1 & 0 & 256 & 0 & \texttt{--------}\\
6 & 0 & 1 & 1 & 1 & 128 & 128 & 0 & \texttt{-------+}\\
7 & 6 & 2 & 1 & 1 & 0 & 256 & 0 & \texttt{--------}\\
\multicolumn{9}{@{}l@{}}{\rule{0pt}{2.6ex}\textbf{T8.050}\enspace \texttt{01234311}\enspace $n{=}8$, $|H|{=}4$, len 5, br 2, $|\mathrm{Aut}|{=}2$, orb 7/6, hist-orb 3/2; M 3840, res 1912, rev 1050: ok}\\*
0 & -- & 0 & 4 & 1 & 54 & 2 & 200 & \texttt{-???????}\\*
1 & 0 & 1 & 2 & 1 & 176 & 16 & 64 & \texttt{--+???--}\\
2 & 1 & 2 & 2 & 1 & 96 & 32 & 128 & \texttt{---???--}\\
3 & 2 & 3 & 1 & 1 & 128 & 128 & 0 & \texttt{----+---}\\
4 & 3 & 4 & 1 & 1 & 0 & 256 & 0 & \texttt{--------}\\
5 & 2 & 3 & 1 & 1 & 0 & 256 & 0 & \texttt{--------}\\
6 & 0 & 1 & 1 & 2 & 0 & 256 & 0 & \texttt{--------}\\
7 & 0 & 1 & 1 & 2 & 0 & 256 & 0 & \texttt{--------}\\
\multicolumn{9}{@{}l@{}}{\rule{0pt}{2.6ex}\textbf{T8.051}\enspace \texttt{01234234}\enspace $n{=}8$, $|H|{=}2$, len 5, br 1, $|\mathrm{Aut}|{=}2$, orb 5/2, hist-orb 1/0; M 4864, res 504, rev 2142: ok}\\*
0 & -- & 0 & 2 & 1 & 226 & 2 & 28 & \texttt{-+??????}\\*
1 & 0 & 1 & 2 & 1 & 196 & 4 & 56 & \texttt{--??????}\\
2 & 1 & 2 & 1 & 2 & 192 & 64 & 0 & \texttt{---++---}\\
3 & 2 & 3 & 1 & 2 & 128 & 128 & 0 & \texttt{----+---}\\
4 & 3 & 4 & 1 & 2 & 0 & 256 & 0 & \texttt{--------}\\
5 & 1 & 2 & 1 & 2 & 192 & 64 & 0 & \texttt{------++}\\
6 & 5 & 3 & 1 & 2 & 128 & 128 & 0 & \texttt{-------+}\\
7 & 6 & 4 & 1 & 2 & 0 & 256 & 0 & \texttt{--------}\\
\multicolumn{9}{@{}l@{}}{\rule{0pt}{2.6ex}\textbf{T8.052}\enspace \texttt{01234233}\enspace $n{=}8$, $|H|{=}3$, len 5, br 2, $|\mathrm{Aut}|{=}2$, orb 7/6, hist-orb 2/1; M 4608, res 1168, rev 1818: ok}\\*
0 & -- & 0 & 3 & 1 & 198 & 2 & 56 & \texttt{-+??????}\\*
1 & 0 & 1 & 3 & 1 & 140 & 4 & 112 & \texttt{--??????}\\
2 & 1 & 2 & 1 & 1 & 192 & 64 & 0 & \texttt{---++---}\\
3 & 2 & 3 & 1 & 1 & 128 & 128 & 0 & \texttt{----+---}\\
4 & 3 & 4 & 1 & 1 & 0 & 256 & 0 & \texttt{--------}\\
5 & 1 & 2 & 2 & 1 & 64 & 64 & 128 & \texttt{------??}\\
6 & 5 & 3 & 1 & 2 & 0 & 256 & 0 & \texttt{--------}\\
7 & 5 & 3 & 1 & 2 & 0 & 256 & 0 & \texttt{--------}\\
\multicolumn{9}{@{}l@{}}{\rule{0pt}{2.6ex}\textbf{T8.053}\enspace \texttt{01234232}\enspace $n{=}8$, $|H|{=}3$, len 5, br 1, $|\mathrm{Aut}|{=}1$, orb 8/8, hist-orb 3/3; M 4352, res 1512, rev 1550: ok}\\*
0 & -- & 0 & 3 & 1 & 170 & 2 & 84 & \texttt{-+??????}\\*
1 & 0 & 1 & 3 & 1 & 84 & 4 & 168 & \texttt{--??????}\\
2 & 1 & 2 & 1 & 1 & 192 & 64 & 0 & \texttt{---++---}\\
3 & 2 & 3 & 1 & 1 & 128 & 128 & 0 & \texttt{----+---}\\
4 & 3 & 4 & 1 & 1 & 0 & 256 & 0 & \texttt{--------}\\
5 & 1 & 2 & 1 & 1 & 128 & 128 & 0 & \texttt{------+-}\\
6 & 5 & 3 & 1 & 1 & 0 & 256 & 0 & \texttt{--------}\\
7 & 1 & 2 & 1 & 1 & 0 & 256 & 0 & \texttt{--------}\\
\multicolumn{9}{@{}l@{}}{\rule{0pt}{2.6ex}\textbf{T8.054}\enspace \texttt{01234231}\enspace $n{=}8$, $|H|{=}3$, len 5, br 2, $|\mathrm{Aut}|{=}1$, orb 8/8, hist-orb 3/3; M 4096, res 1396, rev 1438: ok}\\*
0 & -- & 0 & 3 & 1 & 106 & 2 & 148 & \texttt{-???????}\\*
1 & 0 & 1 & 2 & 1 & 168 & 8 & 80 & \texttt{--?????-}\\
2 & 1 & 2 & 1 & 1 & 192 & 64 & 0 & \texttt{---++---}\\
3 & 2 & 3 & 1 & 1 & 128 & 128 & 0 & \texttt{----+---}\\
4 & 3 & 4 & 1 & 1 & 0 & 256 & 0 & \texttt{--------}\\
5 & 1 & 2 & 1 & 1 & 128 & 128 & 0 & \texttt{------+-}\\
6 & 5 & 3 & 1 & 1 & 0 & 256 & 0 & \texttt{--------}\\
7 & 0 & 1 & 1 & 1 & 0 & 256 & 0 & \texttt{--------}\\
\multicolumn{9}{@{}l@{}}{\rule{0pt}{2.6ex}\textbf{T8.055}\enspace \texttt{01234222}\enspace $n{=}8$, $|H|{=}4$, len 5, br 1, $|\mathrm{Aut}|{=}6$, orb 6/5, hist-orb 2/1; M 4096, res 2016, rev 1186: ok}\\*
0 & -- & 0 & 4 & 1 & 142 & 2 & 112 & \texttt{-+??????}\\*
1 & 0 & 1 & 4 & 1 & 28 & 4 & 224 & \texttt{--??????}\\
2 & 1 & 2 & 1 & 1 & 192 & 64 & 0 & \texttt{---++---}\\
3 & 2 & 3 & 1 & 1 & 128 & 128 & 0 & \texttt{----+---}\\
4 & 3 & 4 & 1 & 1 & 0 & 256 & 0 & \texttt{--------}\\
5 & 1 & 2 & 1 & 3 & 0 & 256 & 0 & \texttt{--------}\\
6 & 1 & 2 & 1 & 3 & 0 & 256 & 0 & \texttt{--------}\\
7 & 1 & 2 & 1 & 3 & 0 & 256 & 0 & \texttt{--------}\\
\multicolumn{9}{@{}l@{}}{\rule{0pt}{2.6ex}\textbf{T8.056}\enspace \texttt{01234221}\enspace $n{=}8$, $|H|{=}4$, len 5, br 2, $|\mathrm{Aut}|{=}2$, orb 7/6, hist-orb 3/2; M 3840, res 2096, rev 954: ok}\\*
0 & -- & 0 & 4 & 1 & 78 & 2 & 176 & \texttt{-???????}\\*
1 & 0 & 1 & 3 & 1 & 56 & 8 & 192 & \texttt{--?????-}\\
2 & 1 & 2 & 1 & 1 & 192 & 64 & 0 & \texttt{---++---}\\
3 & 2 & 3 & 1 & 1 & 128 & 128 & 0 & \texttt{----+---}\\
4 & 3 & 4 & 1 & 1 & 0 & 256 & 0 & \texttt{--------}\\
5 & 1 & 2 & 1 & 2 & 0 & 256 & 0 & \texttt{--------}\\
6 & 1 & 2 & 1 & 2 & 0 & 256 & 0 & \texttt{--------}\\
7 & 0 & 1 & 1 & 1 & 0 & 256 & 0 & \texttt{--------}\\
\multicolumn{9}{@{}l@{}}{\rule{0pt}{2.6ex}\textbf{T8.057}\enspace \texttt{01234212}\enspace $n{=}8$, $|H|{=}3$, len 5, br 2, $|\mathrm{Aut}|{=}1$, orb 8/8, hist-orb 3/3; M 3840, res 1260, rev 1362: ok}\\*
0 & -- & 0 & 3 & 1 & 138 & 2 & 116 & \texttt{-???????}\\*
1 & 0 & 1 & 2 & 1 & 112 & 16 & 128 & \texttt{--????--}\\
2 & 1 & 2 & 1 & 1 & 192 & 64 & 0 & \texttt{---++---}\\
3 & 2 & 3 & 1 & 1 & 128 & 128 & 0 & \texttt{----+---}\\
4 & 3 & 4 & 1 & 1 & 0 & 256 & 0 & \texttt{--------}\\
5 & 1 & 2 & 1 & 1 & 0 & 256 & 0 & \texttt{--------}\\
6 & 0 & 1 & 1 & 1 & 128 & 128 & 0 & \texttt{-------+}\\
7 & 6 & 2 & 1 & 1 & 0 & 256 & 0 & \texttt{--------}\\
\multicolumn{9}{@{}l@{}}{\rule{0pt}{2.6ex}\textbf{T8.058}\enspace \texttt{01234211}\enspace $n{=}8$, $|H|{=}4$, len 5, br 2, $|\mathrm{Aut}|{=}2$, orb 7/6, hist-orb 3/2; M 3584, res 1904, rev 890: ok}\\*
0 & -- & 0 & 4 & 1 & 46 & 2 & 208 & \texttt{-???????}\\*
1 & 0 & 1 & 2 & 1 & 112 & 16 & 128 & \texttt{--????--}\\
2 & 1 & 2 & 1 & 1 & 192 & 64 & 0 & \texttt{---++---}\\
3 & 2 & 3 & 1 & 1 & 128 & 128 & 0 & \texttt{----+---}\\
4 & 3 & 4 & 1 & 1 & 0 & 256 & 0 & \texttt{--------}\\
5 & 1 & 2 & 1 & 1 & 0 & 256 & 0 & \texttt{--------}\\
6 & 0 & 1 & 1 & 2 & 0 & 256 & 0 & \texttt{--------}\\
7 & 0 & 1 & 1 & 2 & 0 & 256 & 0 & \texttt{--------}\\
\multicolumn{9}{@{}l@{}}{\rule{0pt}{2.6ex}\textbf{T8.059}\enspace \texttt{01234123}\enspace $n{=}8$, $|H|{=}2$, len 5, br 1, $|\mathrm{Aut}|{=}1$, orb 8/8, hist-orb 2/2; M 4096, res 308, rev 1814: ok}\\*
0 & -- & 0 & 2 & 1 & 210 & 2 & 44 & \texttt{-???????}\\*
1 & 0 & 1 & 1 & 1 & 224 & 32 & 0 & \texttt{--+++---}\\
2 & 1 & 2 & 1 & 1 & 192 & 64 & 0 & \texttt{---++---}\\
3 & 2 & 3 & 1 & 1 & 128 & 128 & 0 & \texttt{----+---}\\
4 & 3 & 4 & 1 & 1 & 0 & 256 & 0 & \texttt{--------}\\
5 & 0 & 1 & 1 & 1 & 192 & 64 & 0 & \texttt{------++}\\
6 & 5 & 2 & 1 & 1 & 128 & 128 & 0 & \texttt{-------+}\\
7 & 6 & 3 & 1 & 1 & 0 & 256 & 0 & \texttt{--------}\\
\multicolumn{9}{@{}l@{}}{\rule{0pt}{2.6ex}\textbf{T8.060}\enspace \texttt{01234122}\enspace $n{=}8$, $|H|{=}3$, len 5, br 2, $|\mathrm{Aut}|{=}2$, orb 7/6, hist-orb 2/1; M 3840, res 920, rev 1454: ok}\\*
0 & -- & 0 & 3 & 1 & 150 & 2 & 104 & \texttt{-???????}\\*
1 & 0 & 1 & 1 & 1 & 224 & 32 & 0 & \texttt{--+++---}\\
2 & 1 & 2 & 1 & 1 & 192 & 64 & 0 & \texttt{---++---}\\
3 & 2 & 3 & 1 & 1 & 128 & 128 & 0 & \texttt{----+---}\\
4 & 3 & 4 & 1 & 1 & 0 & 256 & 0 & \texttt{--------}\\
5 & 0 & 1 & 2 & 1 & 64 & 64 & 128 & \texttt{------??}\\
6 & 5 & 2 & 1 & 2 & 0 & 256 & 0 & \texttt{--------}\\
7 & 5 & 2 & 1 & 2 & 0 & 256 & 0 & \texttt{--------}\\
\multicolumn{9}{@{}l@{}}{\rule{0pt}{2.6ex}\textbf{T8.061}\enspace \texttt{01234121}\enspace $n{=}8$, $|H|{=}3$, len 5, br 1, $|\mathrm{Aut}|{=}1$, orb 8/8, hist-orb 3/3; M 3584, res 1148, rev 1230: ok}\\*
0 & -- & 0 & 3 & 1 & 90 & 2 & 164 & \texttt{-???????}\\*
1 & 0 & 1 & 1 & 1 & 224 & 32 & 0 & \texttt{--+++---}\\
2 & 1 & 2 & 1 & 1 & 192 & 64 & 0 & \texttt{---++---}\\
3 & 2 & 3 & 1 & 1 & 128 & 128 & 0 & \texttt{----+---}\\
4 & 3 & 4 & 1 & 1 & 0 & 256 & 0 & \texttt{--------}\\
5 & 0 & 1 & 1 & 1 & 128 & 128 & 0 & \texttt{------+-}\\
6 & 5 & 2 & 1 & 1 & 0 & 256 & 0 & \texttt{--------}\\
7 & 0 & 1 & 1 & 1 & 0 & 256 & 0 & \texttt{--------}\\
\multicolumn{9}{@{}l@{}}{\rule{0pt}{2.6ex}\textbf{T8.062}\enspace \texttt{01234111}\enspace $n{=}8$, $|H|{=}4$, len 5, br 1, $|\mathrm{Aut}|{=}6$, orb 6/5, hist-orb 2/1; M 3328, res 1568, rev 878: ok}\\*
0 & -- & 0 & 4 & 1 & 30 & 2 & 224 & \texttt{-???????}\\*
1 & 0 & 1 & 1 & 1 & 224 & 32 & 0 & \texttt{--+++---}\\
2 & 1 & 2 & 1 & 1 & 192 & 64 & 0 & \texttt{---++---}\\
3 & 2 & 3 & 1 & 1 & 128 & 128 & 0 & \texttt{----+---}\\
4 & 3 & 4 & 1 & 1 & 0 & 256 & 0 & \texttt{--------}\\
5 & 0 & 1 & 1 & 3 & 0 & 256 & 0 & \texttt{--------}\\
6 & 0 & 1 & 1 & 3 & 0 & 256 & 0 & \texttt{--------}\\
7 & 0 & 1 & 1 & 3 & 0 & 256 & 0 & \texttt{--------}\\
\multicolumn{9}{@{}l@{}}{\rule{0pt}{2.6ex}\textbf{T8.063}\enspace \texttt{01233333}\enspace $n{=}8$, $|H|{=}5$, len 4, br 1, $|\mathrm{Aut}|{=}120$, orb 4/3, hist-orb 1/0; M 4608, res 2100, rev 1682: ok}\\*
0 & -- & 0 & 5 & 1 & 194 & 2 & 60 & \texttt{-++?????}\\*
1 & 0 & 1 & 5 & 1 & 132 & 4 & 120 & \texttt{--+?????}\\
2 & 1 & 2 & 5 & 1 & 8 & 8 & 240 & \texttt{---?????}\\
3 & 2 & 3 & 1 & 5 & 0 & 256 & 0 & \texttt{--------}\\
4 & 2 & 3 & 1 & 5 & 0 & 256 & 0 & \texttt{--------}\\
5 & 2 & 3 & 1 & 5 & 0 & 256 & 0 & \texttt{--------}\\
6 & 2 & 3 & 1 & 5 & 0 & 256 & 0 & \texttt{--------}\\
7 & 2 & 3 & 1 & 5 & 0 & 256 & 0 & \texttt{--------}\\
\multicolumn{9}{@{}l@{}}{\rule{0pt}{2.6ex}\textbf{T8.064}\enspace \texttt{01233332}\enspace $n{=}8$, $|H|{=}5$, len 4, br 2, $|\mathrm{Aut}|{=}24$, orb 5/4, hist-orb 2/1; M 4352, res 2384, rev 1230: ok}\\*
0 & -- & 0 & 5 & 1 & 162 & 2 & 92 & \texttt{-+??????}\\*
1 & 0 & 1 & 5 & 1 & 68 & 4 & 184 & \texttt{--??????}\\
2 & 1 & 2 & 4 & 1 & 16 & 16 & 224 & \texttt{---????-}\\
3 & 2 & 3 & 1 & 4 & 0 & 256 & 0 & \texttt{--------}\\
4 & 2 & 3 & 1 & 4 & 0 & 256 & 0 & \texttt{--------}\\
5 & 2 & 3 & 1 & 4 & 0 & 256 & 0 & \texttt{--------}\\
6 & 2 & 3 & 1 & 4 & 0 & 256 & 0 & \texttt{--------}\\
7 & 1 & 2 & 1 & 1 & 0 & 256 & 0 & \texttt{--------}\\
\multicolumn{9}{@{}l@{}}{\rule{0pt}{2.6ex}\textbf{T8.065}\enspace \texttt{01233331}\enspace $n{=}8$, $|H|{=}5$, len 4, br 2, $|\mathrm{Aut}|{=}24$, orb 5/4, hist-orb 2/1; M 4096, res 2324, rev 1114: ok}\\*
0 & -- & 0 & 5 & 1 & 98 & 2 & 156 & \texttt{-???????}\\*
1 & 0 & 1 & 4 & 1 & 136 & 8 & 112 & \texttt{--+????-}\\
2 & 1 & 2 & 4 & 1 & 16 & 16 & 224 & \texttt{---????-}\\
3 & 2 & 3 & 1 & 4 & 0 & 256 & 0 & \texttt{--------}\\
4 & 2 & 3 & 1 & 4 & 0 & 256 & 0 & \texttt{--------}\\
5 & 2 & 3 & 1 & 4 & 0 & 256 & 0 & \texttt{--------}\\
6 & 2 & 3 & 1 & 4 & 0 & 256 & 0 & \texttt{--------}\\
7 & 0 & 1 & 1 & 1 & 0 & 256 & 0 & \texttt{--------}\\
\multicolumn{9}{@{}l@{}}{\rule{0pt}{2.6ex}\textbf{T8.066}\enspace \texttt{01233323}\enspace $n{=}8$, $|H|{=}4$, len 4, br 2, $|\mathrm{Aut}|{=}6$, orb 6/5, hist-orb 2/1; M 4352, res 1728, rev 1538: ok}\\*
0 & -- & 0 & 4 & 1 & 182 & 2 & 72 & \texttt{-+??????}\\*
1 & 0 & 1 & 4 & 1 & 108 & 4 & 144 & \texttt{--??????}\\
2 & 1 & 2 & 3 & 1 & 32 & 32 & 192 & \texttt{---???--}\\
3 & 2 & 3 & 1 & 3 & 0 & 256 & 0 & \texttt{--------}\\
4 & 2 & 3 & 1 & 3 & 0 & 256 & 0 & \texttt{--------}\\
5 & 2 & 3 & 1 & 3 & 0 & 256 & 0 & \texttt{--------}\\
6 & 1 & 2 & 1 & 1 & 128 & 128 & 0 & \texttt{-------+}\\
7 & 6 & 3 & 1 & 1 & 0 & 256 & 0 & \texttt{--------}\\
\multicolumn{9}{@{}l@{}}{\rule{0pt}{2.6ex}\textbf{T8.067}\enspace \texttt{01233322}\enspace $n{=}8$, $|H|{=}5$, len 4, br 2, $|\mathrm{Aut}|{=}12$, orb 5/3, hist-orb 2/0; M 4096, res 2376, rev 1030: ok}\\*
0 & -- & 0 & 5 & 1 & 146 & 2 & 108 & \texttt{-+??????}\\*
1 & 0 & 1 & 5 & 1 & 36 & 4 & 216 & \texttt{--??????}\\
2 & 1 & 2 & 3 & 1 & 32 & 32 & 192 & \texttt{---???--}\\
3 & 2 & 3 & 1 & 3 & 0 & 256 & 0 & \texttt{--------}\\
4 & 2 & 3 & 1 & 3 & 0 & 256 & 0 & \texttt{--------}\\
5 & 2 & 3 & 1 & 3 & 0 & 256 & 0 & \texttt{--------}\\
6 & 1 & 2 & 1 & 2 & 0 & 256 & 0 & \texttt{--------}\\
7 & 1 & 2 & 1 & 2 & 0 & 256 & 0 & \texttt{--------}\\
\multicolumn{9}{@{}l@{}}{\rule{0pt}{2.6ex}\textbf{T8.068}\enspace \texttt{01233321}\enspace $n{=}8$, $|H|{=}5$, len 4, br 3, $|\mathrm{Aut}|{=}6$, orb 6/5, hist-orb 3/2; M 3840, res 2428, rev 814: ok}\\*
0 & -- & 0 & 5 & 1 & 82 & 2 & 172 & \texttt{-???????}\\*
1 & 0 & 1 & 4 & 1 & 72 & 8 & 176 & \texttt{--?????-}\\
2 & 1 & 2 & 3 & 1 & 32 & 32 & 192 & \texttt{---???--}\\
3 & 2 & 3 & 1 & 3 & 0 & 256 & 0 & \texttt{--------}\\
4 & 2 & 3 & 1 & 3 & 0 & 256 & 0 & \texttt{--------}\\
5 & 2 & 3 & 1 & 3 & 0 & 256 & 0 & \texttt{--------}\\
6 & 1 & 2 & 1 & 1 & 0 & 256 & 0 & \texttt{--------}\\
7 & 0 & 1 & 1 & 1 & 0 & 256 & 0 & \texttt{--------}\\
\multicolumn{9}{@{}l@{}}{\rule{0pt}{2.6ex}\textbf{T8.069}\enspace \texttt{01233312}\enspace $n{=}8$, $|H|{=}4$, len 4, br 2, $|\mathrm{Aut}|{=}6$, orb 6/5, hist-orb 2/1; M 3840, res 1496, rev 1346: ok}\\*
0 & -- & 0 & 4 & 1 & 150 & 2 & 104 & \texttt{-???????}\\*
1 & 0 & 1 & 3 & 1 & 144 & 16 & 96 & \texttt{--+???--}\\
2 & 1 & 2 & 3 & 1 & 32 & 32 & 192 & \texttt{---???--}\\
3 & 2 & 3 & 1 & 3 & 0 & 256 & 0 & \texttt{--------}\\
4 & 2 & 3 & 1 & 3 & 0 & 256 & 0 & \texttt{--------}\\
5 & 2 & 3 & 1 & 3 & 0 & 256 & 0 & \texttt{--------}\\
6 & 0 & 1 & 1 & 1 & 128 & 128 & 0 & \texttt{-------+}\\
7 & 6 & 2 & 1 & 1 & 0 & 256 & 0 & \texttt{--------}\\
\multicolumn{9}{@{}l@{}}{\rule{0pt}{2.6ex}\textbf{T8.070}\enspace \texttt{01233311}\enspace $n{=}8$, $|H|{=}5$, len 4, br 2, $|\mathrm{Aut}|{=}12$, orb 5/3, hist-orb 2/0; M 3584, res 2196, rev 802: ok}\\*
0 & -- & 0 & 5 & 1 & 50 & 2 & 204 & \texttt{-???????}\\*
1 & 0 & 1 & 3 & 1 & 144 & 16 & 96 & \texttt{--+???--}\\
2 & 1 & 2 & 3 & 1 & 32 & 32 & 192 & \texttt{---???--}\\
3 & 2 & 3 & 1 & 3 & 0 & 256 & 0 & \texttt{--------}\\
4 & 2 & 3 & 1 & 3 & 0 & 256 & 0 & \texttt{--------}\\
5 & 2 & 3 & 1 & 3 & 0 & 256 & 0 & \texttt{--------}\\
6 & 0 & 1 & 1 & 2 & 0 & 256 & 0 & \texttt{--------}\\
7 & 0 & 1 & 1 & 2 & 0 & 256 & 0 & \texttt{--------}\\
\multicolumn{9}{@{}l@{}}{\rule{0pt}{2.6ex}\textbf{T8.071}\enspace \texttt{01233233}\enspace $n{=}8$, $|H|{=}4$, len 4, br 3, $|\mathrm{Aut}|{=}8$, orb 4/2, hist-orb 1/0; M 4352, res 1688, rev 1494: ok}\\*
0 & -- & 0 & 4 & 1 & 178 & 2 & 76 & \texttt{-+??????}\\*
1 & 0 & 1 & 4 & 1 & 100 & 4 & 152 & \texttt{--??????}\\
2 & 1 & 2 & 2 & 2 & 64 & 64 & 128 & \texttt{---??---}\\
3 & 2 & 3 & 1 & 4 & 0 & 256 & 0 & \texttt{--------}\\
4 & 2 & 3 & 1 & 4 & 0 & 256 & 0 & \texttt{--------}\\
5 & 1 & 2 & 2 & 2 & 64 & 64 & 128 & \texttt{------??}\\
6 & 5 & 3 & 1 & 4 & 0 & 256 & 0 & \texttt{--------}\\
7 & 5 & 3 & 1 & 4 & 0 & 256 & 0 & \texttt{--------}\\
\multicolumn{9}{@{}l@{}}{\rule{0pt}{2.6ex}\textbf{T8.072}\enspace \texttt{01233232}\enspace $n{=}8$, $|H|{=}4$, len 4, br 2, $|\mathrm{Aut}|{=}2$, orb 7/6, hist-orb 3/2; M 4096, res 1888, rev 1274: ok}\\*
0 & -- & 0 & 4 & 1 & 158 & 2 & 96 & \texttt{-+??????}\\*
1 & 0 & 1 & 4 & 1 & 60 & 4 & 192 & \texttt{--??????}\\
2 & 1 & 2 & 2 & 1 & 64 & 64 & 128 & \texttt{---??---}\\
3 & 2 & 3 & 1 & 2 & 0 & 256 & 0 & \texttt{--------}\\
4 & 2 & 3 & 1 & 2 & 0 & 256 & 0 & \texttt{--------}\\
5 & 1 & 2 & 1 & 1 & 128 & 128 & 0 & \texttt{------+-}\\
6 & 5 & 3 & 1 & 1 & 0 & 256 & 0 & \texttt{--------}\\
7 & 1 & 2 & 1 & 1 & 0 & 256 & 0 & \texttt{--------}\\
\multicolumn{9}{@{}l@{}}{\rule{0pt}{2.6ex}\textbf{T8.073}\enspace \texttt{01233231}\enspace $n{=}8$, $|H|{=}4$, len 4, br 3, $|\mathrm{Aut}|{=}2$, orb 7/6, hist-orb 3/2; M 3840, res 1856, rev 1122: ok}\\*
0 & -- & 0 & 4 & 1 & 94 & 2 & 160 & \texttt{-???????}\\*
1 & 0 & 1 & 3 & 1 & 120 & 8 & 128 & \texttt{--?????-}\\
2 & 1 & 2 & 2 & 1 & 64 & 64 & 128 & \texttt{---??---}\\
3 & 2 & 3 & 1 & 2 & 0 & 256 & 0 & \texttt{--------}\\
4 & 2 & 3 & 1 & 2 & 0 & 256 & 0 & \texttt{--------}\\
5 & 1 & 2 & 1 & 1 & 128 & 128 & 0 & \texttt{------+-}\\
6 & 5 & 3 & 1 & 1 & 0 & 256 & 0 & \texttt{--------}\\
7 & 0 & 1 & 1 & 1 & 0 & 256 & 0 & \texttt{--------}\\
\multicolumn{9}{@{}l@{}}{\rule{0pt}{2.6ex}\textbf{T8.074}\enspace \texttt{01233222}\enspace $n{=}8$, $|H|{=}5$, len 4, br 2, $|\mathrm{Aut}|{=}12$, orb 5/3, hist-orb 2/0; M 3840, res 2248, rev 958: ok}\\*
0 & -- & 0 & 5 & 1 & 138 & 2 & 116 & \texttt{-+??????}\\*
1 & 0 & 1 & 5 & 1 & 20 & 4 & 232 & \texttt{--??????}\\
2 & 1 & 2 & 2 & 1 & 64 & 64 & 128 & \texttt{---??---}\\
3 & 2 & 3 & 1 & 2 & 0 & 256 & 0 & \texttt{--------}\\
4 & 2 & 3 & 1 & 2 & 0 & 256 & 0 & \texttt{--------}\\
5 & 1 & 2 & 1 & 3 & 0 & 256 & 0 & \texttt{--------}\\
6 & 1 & 2 & 1 & 3 & 0 & 256 & 0 & \texttt{--------}\\
7 & 1 & 2 & 1 & 3 & 0 & 256 & 0 & \texttt{--------}\\
\multicolumn{9}{@{}l@{}}{\rule{0pt}{2.6ex}\textbf{T8.075}\enspace \texttt{01233221}\enspace $n{=}8$, $|H|{=}5$, len 4, br 3, $|\mathrm{Aut}|{=}4$, orb 6/4, hist-orb 3/1; M 3584, res 2356, rev 686: ok}\\*
0 & -- & 0 & 5 & 1 & 74 & 2 & 180 & \texttt{-???????}\\*
1 & 0 & 1 & 4 & 1 & 40 & 8 & 208 & \texttt{--?????-}\\
2 & 1 & 2 & 2 & 1 & 64 & 64 & 128 & \texttt{---??---}\\
3 & 2 & 3 & 1 & 2 & 0 & 256 & 0 & \texttt{--------}\\
4 & 2 & 3 & 1 & 2 & 0 & 256 & 0 & \texttt{--------}\\
5 & 1 & 2 & 1 & 2 & 0 & 256 & 0 & \texttt{--------}\\
6 & 1 & 2 & 1 & 2 & 0 & 256 & 0 & \texttt{--------}\\
7 & 0 & 1 & 1 & 1 & 0 & 256 & 0 & \texttt{--------}\\
\multicolumn{9}{@{}l@{}}{\rule{0pt}{2.6ex}\textbf{T8.076}\enspace \texttt{01233212}\enspace $n{=}8$, $|H|{=}4$, len 4, br 3, $|\mathrm{Aut}|{=}2$, orb 7/6, hist-orb 3/2; M 3584, res 1616, rev 1094: ok}\\*
0 & -- & 0 & 4 & 1 & 126 & 2 & 128 & \texttt{-???????}\\*
1 & 0 & 1 & 3 & 1 & 80 & 16 & 160 & \texttt{--????--}\\
2 & 1 & 2 & 2 & 1 & 64 & 64 & 128 & \texttt{---??---}\\
3 & 2 & 3 & 1 & 2 & 0 & 256 & 0 & \texttt{--------}\\
4 & 2 & 3 & 1 & 2 & 0 & 256 & 0 & \texttt{--------}\\
5 & 1 & 2 & 1 & 1 & 0 & 256 & 0 & \texttt{--------}\\
6 & 0 & 1 & 1 & 1 & 128 & 128 & 0 & \texttt{-------+}\\
7 & 6 & 2 & 1 & 1 & 0 & 256 & 0 & \texttt{--------}\\
\multicolumn{9}{@{}l@{}}{\rule{0pt}{2.6ex}\textbf{T8.077}\enspace \texttt{01233211}\enspace $n{=}8$, $|H|{=}5$, len 4, br 3, $|\mathrm{Aut}|{=}4$, orb 6/4, hist-orb 3/1; M 3328, res 2204, rev 606: ok}\\*
0 & -- & 0 & 5 & 1 & 42 & 2 & 212 & \texttt{-???????}\\*
1 & 0 & 1 & 3 & 1 & 80 & 16 & 160 & \texttt{--????--}\\
2 & 1 & 2 & 2 & 1 & 64 & 64 & 128 & \texttt{---??---}\\
3 & 2 & 3 & 1 & 2 & 0 & 256 & 0 & \texttt{--------}\\
4 & 2 & 3 & 1 & 2 & 0 & 256 & 0 & \texttt{--------}\\
5 & 1 & 2 & 1 & 1 & 0 & 256 & 0 & \texttt{--------}\\
6 & 0 & 1 & 1 & 2 & 0 & 256 & 0 & \texttt{--------}\\
7 & 0 & 1 & 1 & 2 & 0 & 256 & 0 & \texttt{--------}\\
\multicolumn{9}{@{}l@{}}{\rule{0pt}{2.6ex}\textbf{T8.078}\enspace \texttt{01233123}\enspace $n{=}8$, $|H|{=}3$, len 4, br 2, $|\mathrm{Aut}|{=}2$, orb 7/6, hist-orb 2/1; M 3840, res 824, rev 1570: ok}\\*
0 & -- & 0 & 3 & 1 & 182 & 2 & 72 & \texttt{-???????}\\*
1 & 0 & 1 & 2 & 1 & 160 & 32 & 64 & \texttt{--+??---}\\
2 & 1 & 2 & 2 & 1 & 64 & 64 & 128 & \texttt{---??---}\\
3 & 2 & 3 & 1 & 2 & 0 & 256 & 0 & \texttt{--------}\\
4 & 2 & 3 & 1 & 2 & 0 & 256 & 0 & \texttt{--------}\\
5 & 0 & 1 & 1 & 1 & 192 & 64 & 0 & \texttt{------++}\\
6 & 5 & 2 & 1 & 1 & 128 & 128 & 0 & \texttt{-------+}\\
7 & 6 & 3 & 1 & 1 & 0 & 256 & 0 & \texttt{--------}\\
\multicolumn{9}{@{}l@{}}{\rule{0pt}{2.6ex}\textbf{T8.079}\enspace \texttt{01233122}\enspace $n{=}8$, $|H|{=}4$, len 4, br 3, $|\mathrm{Aut}|{=}4$, orb 6/4, hist-orb 2/0; M 3584, res 1380, rev 1194: ok}\\*
0 & -- & 0 & 4 & 1 & 130 & 2 & 124 & \texttt{-???????}\\*
1 & 0 & 1 & 2 & 1 & 160 & 32 & 64 & \texttt{--+??---}\\
2 & 1 & 2 & 2 & 1 & 64 & 64 & 128 & \texttt{---??---}\\
3 & 2 & 3 & 1 & 2 & 0 & 256 & 0 & \texttt{--------}\\
4 & 2 & 3 & 1 & 2 & 0 & 256 & 0 & \texttt{--------}\\
5 & 0 & 1 & 2 & 1 & 64 & 64 & 128 & \texttt{------??}\\
6 & 5 & 2 & 1 & 2 & 0 & 256 & 0 & \texttt{--------}\\
7 & 5 & 2 & 1 & 2 & 0 & 256 & 0 & \texttt{--------}\\
\multicolumn{9}{@{}l@{}}{\rule{0pt}{2.6ex}\textbf{T8.080}\enspace \texttt{01233121}\enspace $n{=}8$, $|H|{=}4$, len 4, br 2, $|\mathrm{Aut}|{=}2$, orb 7/6, hist-orb 3/2; M 3328, res 1552, rev 970: ok}\\*
0 & -- & 0 & 4 & 1 & 78 & 2 & 176 & \texttt{-???????}\\*
1 & 0 & 1 & 2 & 1 & 160 & 32 & 64 & \texttt{--+??---}\\
2 & 1 & 2 & 2 & 1 & 64 & 64 & 128 & \texttt{---??---}\\
3 & 2 & 3 & 1 & 2 & 0 & 256 & 0 & \texttt{--------}\\
4 & 2 & 3 & 1 & 2 & 0 & 256 & 0 & \texttt{--------}\\
5 & 0 & 1 & 1 & 1 & 128 & 128 & 0 & \texttt{------+-}\\
6 & 5 & 2 & 1 & 1 & 0 & 256 & 0 & \texttt{--------}\\
7 & 0 & 1 & 1 & 1 & 0 & 256 & 0 & \texttt{--------}\\
\multicolumn{9}{@{}l@{}}{\rule{0pt}{2.6ex}\textbf{T8.081}\enspace \texttt{01233111}\enspace $n{=}8$, $|H|{=}5$, len 4, br 2, $|\mathrm{Aut}|{=}12$, orb 5/3, hist-orb 2/0; M 3072, res 1916, rev 618: ok}\\*
0 & -- & 0 & 5 & 1 & 26 & 2 & 228 & \texttt{-???????}\\*
1 & 0 & 1 & 2 & 1 & 160 & 32 & 64 & \texttt{--+??---}\\
2 & 1 & 2 & 2 & 1 & 64 & 64 & 128 & \texttt{---??---}\\
3 & 2 & 3 & 1 & 2 & 0 & 256 & 0 & \texttt{--------}\\
4 & 2 & 3 & 1 & 2 & 0 & 256 & 0 & \texttt{--------}\\
5 & 0 & 1 & 1 & 3 & 0 & 256 & 0 & \texttt{--------}\\
6 & 0 & 1 & 1 & 3 & 0 & 256 & 0 & \texttt{--------}\\
7 & 0 & 1 & 1 & 3 & 0 & 256 & 0 & \texttt{--------}\\
\multicolumn{9}{@{}l@{}}{\rule{0pt}{2.6ex}\textbf{T8.082}\enspace \texttt{01232323}\enspace $n{=}8$, $|H|{=}3$, len 4, br 1, $|\mathrm{Aut}|{=}6$, orb 4/2, hist-orb 1/0; M 4096, res 1296, rev 1610: ok}\\*
0 & -- & 0 & 3 & 1 & 182 & 2 & 72 & \texttt{-+??????}\\*
1 & 0 & 1 & 3 & 1 & 108 & 4 & 144 & \texttt{--??????}\\
2 & 1 & 2 & 1 & 3 & 128 & 128 & 0 & \texttt{---+----}\\
3 & 2 & 3 & 1 & 3 & 0 & 256 & 0 & \texttt{--------}\\
4 & 1 & 2 & 1 & 3 & 128 & 128 & 0 & \texttt{-----+--}\\
5 & 4 & 3 & 1 & 3 & 0 & 256 & 0 & \texttt{--------}\\
6 & 1 & 2 & 1 & 3 & 128 & 128 & 0 & \texttt{-------+}\\
7 & 6 & 3 & 1 & 3 & 0 & 256 & 0 & \texttt{--------}\\
\multicolumn{9}{@{}l@{}}{\rule{0pt}{2.6ex}\textbf{T8.083}\enspace \texttt{01232322}\enspace $n{=}8$, $|H|{=}4$, len 4, br 1, $|\mathrm{Aut}|{=}4$, orb 5/2, hist-orb 2/0; M 3840, res 1944, rev 1150: ok}\\*
0 & -- & 0 & 4 & 1 & 146 & 2 & 108 & \texttt{-+??????}\\*
1 & 0 & 1 & 4 & 1 & 36 & 4 & 216 & \texttt{--??????}\\
2 & 1 & 2 & 1 & 2 & 128 & 128 & 0 & \texttt{---+----}\\
3 & 2 & 3 & 1 & 2 & 0 & 256 & 0 & \texttt{--------}\\
4 & 1 & 2 & 1 & 2 & 128 & 128 & 0 & \texttt{-----+--}\\
5 & 4 & 3 & 1 & 2 & 0 & 256 & 0 & \texttt{--------}\\
6 & 1 & 2 & 1 & 2 & 0 & 256 & 0 & \texttt{--------}\\
7 & 1 & 2 & 1 & 2 & 0 & 256 & 0 & \texttt{--------}\\
\multicolumn{9}{@{}l@{}}{\rule{0pt}{2.6ex}\textbf{T8.084}\enspace \texttt{01232321}\enspace $n{=}8$, $|H|{=}4$, len 4, br 2, $|\mathrm{Aut}|{=}2$, orb 6/4, hist-orb 3/2; M 3584, res 1996, rev 926: ok}\\*
0 & -- & 0 & 4 & 1 & 82 & 2 & 172 & \texttt{-???????}\\*
1 & 0 & 1 & 3 & 1 & 72 & 8 & 176 & \texttt{--?????-}\\
2 & 1 & 2 & 1 & 2 & 128 & 128 & 0 & \texttt{---+----}\\
3 & 2 & 3 & 1 & 2 & 0 & 256 & 0 & \texttt{--------}\\
4 & 1 & 2 & 1 & 2 & 128 & 128 & 0 & \texttt{-----+--}\\
5 & 4 & 3 & 1 & 2 & 0 & 256 & 0 & \texttt{--------}\\
6 & 1 & 2 & 1 & 1 & 0 & 256 & 0 & \texttt{--------}\\
7 & 0 & 1 & 1 & 1 & 0 & 256 & 0 & \texttt{--------}\\
\multicolumn{9}{@{}l@{}}{\rule{0pt}{2.6ex}\textbf{T8.085}\enspace \texttt{01232312}\enspace $n{=}8$, $|H|{=}3$, len 4, br 2, $|\mathrm{Aut}|{=}2$, orb 6/4, hist-orb 2/1; M 3584, res 1064, rev 1406: ok}\\*
0 & -- & 0 & 3 & 1 & 150 & 2 & 104 & \texttt{-???????}\\*
1 & 0 & 1 & 2 & 1 & 144 & 16 & 96 & \texttt{--????--}\\
2 & 1 & 2 & 1 & 2 & 128 & 128 & 0 & \texttt{---+----}\\
3 & 2 & 3 & 1 & 2 & 0 & 256 & 0 & \texttt{--------}\\
4 & 1 & 2 & 1 & 2 & 128 & 128 & 0 & \texttt{-----+--}\\
5 & 4 & 3 & 1 & 2 & 0 & 256 & 0 & \texttt{--------}\\
6 & 0 & 1 & 1 & 1 & 128 & 128 & 0 & \texttt{-------+}\\
7 & 6 & 2 & 1 & 1 & 0 & 256 & 0 & \texttt{--------}\\
\multicolumn{9}{@{}l@{}}{\rule{0pt}{2.6ex}\textbf{T8.086}\enspace \texttt{01232311}\enspace $n{=}8$, $|H|{=}4$, len 4, br 2, $|\mathrm{Aut}|{=}4$, orb 5/2, hist-orb 2/0; M 3328, res 1764, rev 886: ok}\\*
0 & -- & 0 & 4 & 1 & 50 & 2 & 204 & \texttt{-???????}\\*
1 & 0 & 1 & 2 & 1 & 144 & 16 & 96 & \texttt{--????--}\\
2 & 1 & 2 & 1 & 2 & 128 & 128 & 0 & \texttt{---+----}\\
3 & 2 & 3 & 1 & 2 & 0 & 256 & 0 & \texttt{--------}\\
4 & 1 & 2 & 1 & 2 & 128 & 128 & 0 & \texttt{-----+--}\\
5 & 4 & 3 & 1 & 2 & 0 & 256 & 0 & \texttt{--------}\\
6 & 0 & 1 & 1 & 2 & 0 & 256 & 0 & \texttt{--------}\\
7 & 0 & 1 & 1 & 2 & 0 & 256 & 0 & \texttt{--------}\\
\multicolumn{9}{@{}l@{}}{\rule{0pt}{2.6ex}\textbf{T8.087}\enspace \texttt{01232222}\enspace $n{=}8$, $|H|{=}5$, len 4, br 1, $|\mathrm{Aut}|{=}24$, orb 5/4, hist-orb 2/1; M 3584, res 2160, rev 930: ok}\\*
0 & -- & 0 & 5 & 1 & 134 & 2 & 120 & \texttt{-+??????}\\*
1 & 0 & 1 & 5 & 1 & 12 & 4 & 240 & \texttt{--??????}\\
2 & 1 & 2 & 1 & 1 & 128 & 128 & 0 & \texttt{---+----}\\
3 & 2 & 3 & 1 & 1 & 0 & 256 & 0 & \texttt{--------}\\
4 & 1 & 2 & 1 & 4 & 0 & 256 & 0 & \texttt{--------}\\
5 & 1 & 2 & 1 & 4 & 0 & 256 & 0 & \texttt{--------}\\
6 & 1 & 2 & 1 & 4 & 0 & 256 & 0 & \texttt{--------}\\
7 & 1 & 2 & 1 & 4 & 0 & 256 & 0 & \texttt{--------}\\
\multicolumn{9}{@{}l@{}}{\rule{0pt}{2.6ex}\textbf{T8.088}\enspace \texttt{01232221}\enspace $n{=}8$, $|H|{=}5$, len 4, br 2, $|\mathrm{Aut}|{=}6$, orb 6/5, hist-orb 3/2; M 3328, res 2296, rev 618: ok}\\*
0 & -- & 0 & 5 & 1 & 70 & 2 & 184 & \texttt{-???????}\\*
1 & 0 & 1 & 4 & 1 & 24 & 8 & 224 & \texttt{--?????-}\\
2 & 1 & 2 & 1 & 1 & 128 & 128 & 0 & \texttt{---+----}\\
3 & 2 & 3 & 1 & 1 & 0 & 256 & 0 & \texttt{--------}\\
4 & 1 & 2 & 1 & 3 & 0 & 256 & 0 & \texttt{--------}\\
5 & 1 & 2 & 1 & 3 & 0 & 256 & 0 & \texttt{--------}\\
6 & 1 & 2 & 1 & 3 & 0 & 256 & 0 & \texttt{--------}\\
7 & 0 & 1 & 1 & 1 & 0 & 256 & 0 & \texttt{--------}\\
\multicolumn{9}{@{}l@{}}{\rule{0pt}{2.6ex}\textbf{T8.089}\enspace \texttt{01232212}\enspace $n{=}8$, $|H|{=}4$, len 4, br 2, $|\mathrm{Aut}|{=}2$, orb 7/6, hist-orb 3/2; M 3328, res 1652, rev 970: ok}\\*
0 & -- & 0 & 4 & 1 & 114 & 2 & 140 & \texttt{-???????}\\*
1 & 0 & 1 & 3 & 1 & 48 & 16 & 192 & \texttt{--????--}\\
2 & 1 & 2 & 1 & 1 & 128 & 128 & 0 & \texttt{---+----}\\
3 & 2 & 3 & 1 & 1 & 0 & 256 & 0 & \texttt{--------}\\
4 & 1 & 2 & 1 & 2 & 0 & 256 & 0 & \texttt{--------}\\
5 & 1 & 2 & 1 & 2 & 0 & 256 & 0 & \texttt{--------}\\
6 & 0 & 1 & 1 & 1 & 128 & 128 & 0 & \texttt{-------+}\\
7 & 6 & 2 & 1 & 1 & 0 & 256 & 0 & \texttt{--------}\\
\multicolumn{9}{@{}l@{}}{\rule{0pt}{2.6ex}\textbf{T8.090}\enspace \texttt{01232211}\enspace $n{=}8$, $|H|{=}5$, len 4, br 2, $|\mathrm{Aut}|{=}4$, orb 6/4, hist-orb 3/1; M 3072, res 2184, rev 498: ok}\\*
0 & -- & 0 & 5 & 1 & 38 & 2 & 216 & \texttt{-???????}\\*
1 & 0 & 1 & 3 & 1 & 48 & 16 & 192 & \texttt{--????--}\\
2 & 1 & 2 & 1 & 1 & 128 & 128 & 0 & \texttt{---+----}\\
3 & 2 & 3 & 1 & 1 & 0 & 256 & 0 & \texttt{--------}\\
4 & 1 & 2 & 1 & 2 & 0 & 256 & 0 & \texttt{--------}\\
5 & 1 & 2 & 1 & 2 & 0 & 256 & 0 & \texttt{--------}\\
6 & 0 & 1 & 1 & 2 & 0 & 256 & 0 & \texttt{--------}\\
7 & 0 & 1 & 1 & 2 & 0 & 256 & 0 & \texttt{--------}\\
\multicolumn{9}{@{}l@{}}{\rule{0pt}{2.6ex}\textbf{T8.091}\enspace \texttt{01232123}\enspace $n{=}8$, $|H|{=}3$, len 4, br 2, $|\mathrm{Aut}|{=}1$, orb 8/8, hist-orb 3/3; M 3584, res 1020, rev 1370: ok}\\*
0 & -- & 0 & 3 & 1 & 154 & 2 & 100 & \texttt{-???????}\\*
1 & 0 & 1 & 2 & 1 & 96 & 32 & 128 & \texttt{--???---}\\
2 & 1 & 2 & 1 & 1 & 128 & 128 & 0 & \texttt{---+----}\\
3 & 2 & 3 & 1 & 1 & 0 & 256 & 0 & \texttt{--------}\\
4 & 1 & 2 & 1 & 1 & 0 & 256 & 0 & \texttt{--------}\\
5 & 0 & 1 & 1 & 1 & 192 & 64 & 0 & \texttt{------++}\\
6 & 5 & 2 & 1 & 1 & 128 & 128 & 0 & \texttt{-------+}\\
7 & 6 & 3 & 1 & 1 & 0 & 256 & 0 & \texttt{--------}\\
\multicolumn{9}{@{}l@{}}{\rule{0pt}{2.6ex}\textbf{T8.092}\enspace \texttt{01232122}\enspace $n{=}8$, $|H|{=}4$, len 4, br 3, $|\mathrm{Aut}|{=}2$, orb 7/6, hist-orb 3/2; M 3328, res 1520, rev 1002: ok}\\*
0 & -- & 0 & 4 & 1 & 110 & 2 & 144 & \texttt{-???????}\\*
1 & 0 & 1 & 2 & 1 & 96 & 32 & 128 & \texttt{--???---}\\
2 & 1 & 2 & 1 & 1 & 128 & 128 & 0 & \texttt{---+----}\\
3 & 2 & 3 & 1 & 1 & 0 & 256 & 0 & \texttt{--------}\\
4 & 1 & 2 & 1 & 1 & 0 & 256 & 0 & \texttt{--------}\\
5 & 0 & 1 & 2 & 1 & 64 & 64 & 128 & \texttt{------??}\\
6 & 5 & 2 & 1 & 2 & 0 & 256 & 0 & \texttt{--------}\\
7 & 5 & 2 & 1 & 2 & 0 & 256 & 0 & \texttt{--------}\\
\multicolumn{9}{@{}l@{}}{\rule{0pt}{2.6ex}\textbf{T8.093}\enspace \texttt{01232121}\enspace $n{=}8$, $|H|{=}4$, len 4, br 2, $|\mathrm{Aut}|{=}1$, orb 8/8, hist-orb 4/4; M 3072, res 1636, rev 802: ok}\\*
0 & -- & 0 & 4 & 1 & 66 & 2 & 188 & \texttt{-???????}\\*
1 & 0 & 1 & 2 & 1 & 96 & 32 & 128 & \texttt{--???---}\\
2 & 1 & 2 & 1 & 1 & 128 & 128 & 0 & \texttt{---+----}\\
3 & 2 & 3 & 1 & 1 & 0 & 256 & 0 & \texttt{--------}\\
4 & 1 & 2 & 1 & 1 & 0 & 256 & 0 & \texttt{--------}\\
5 & 0 & 1 & 1 & 1 & 128 & 128 & 0 & \texttt{------+-}\\
6 & 5 & 2 & 1 & 1 & 0 & 256 & 0 & \texttt{--------}\\
7 & 0 & 1 & 1 & 1 & 0 & 256 & 0 & \texttt{--------}\\
\multicolumn{9}{@{}l@{}}{\rule{0pt}{2.6ex}\textbf{T8.094}\enspace \texttt{01232111}\enspace $n{=}8$, $|H|{=}5$, len 4, br 2, $|\mathrm{Aut}|{=}6$, orb 6/5, hist-orb 3/2; M 2816, res 1944, rev 474: ok}\\*
0 & -- & 0 & 5 & 1 & 22 & 2 & 232 & \texttt{-???????}\\*
1 & 0 & 1 & 2 & 1 & 96 & 32 & 128 & \texttt{--???---}\\
2 & 1 & 2 & 1 & 1 & 128 & 128 & 0 & \texttt{---+----}\\
3 & 2 & 3 & 1 & 1 & 0 & 256 & 0 & \texttt{--------}\\
4 & 1 & 2 & 1 & 1 & 0 & 256 & 0 & \texttt{--------}\\
5 & 0 & 1 & 1 & 3 & 0 & 256 & 0 & \texttt{--------}\\
6 & 0 & 1 & 1 & 3 & 0 & 256 & 0 & \texttt{--------}\\
7 & 0 & 1 & 1 & 3 & 0 & 256 & 0 & \texttt{--------}\\
\multicolumn{9}{@{}l@{}}{\rule{0pt}{2.6ex}\textbf{T8.095}\enspace \texttt{01231231}\enspace $n{=}8$, $|H|{=}3$, len 4, br 1, $|\mathrm{Aut}|{=}2$, orb 5/2, hist-orb 2/1; M 3328, res 1092, rev 1174: ok}\\*
0 & -- & 0 & 3 & 1 & 98 & 2 & 156 & \texttt{-???????}\\*
1 & 0 & 1 & 1 & 2 & 192 & 64 & 0 & \texttt{--++----}\\
2 & 1 & 2 & 1 & 2 & 128 & 128 & 0 & \texttt{---+----}\\
3 & 2 & 3 & 1 & 2 & 0 & 256 & 0 & \texttt{--------}\\
4 & 0 & 1 & 1 & 2 & 192 & 64 & 0 & \texttt{-----++-}\\
5 & 4 & 2 & 1 & 2 & 128 & 128 & 0 & \texttt{------+-}\\
6 & 5 & 3 & 1 & 2 & 0 & 256 & 0 & \texttt{--------}\\
7 & 0 & 1 & 1 & 1 & 0 & 256 & 0 & \texttt{--------}\\
\multicolumn{9}{@{}l@{}}{\rule{0pt}{2.6ex}\textbf{T8.096}\enspace \texttt{01231222}\enspace $n{=}8$, $|H|{=}4$, len 4, br 2, $|\mathrm{Aut}|{=}6$, orb 6/5, hist-orb 2/1; M 3328, res 1376, rev 1070: ok}\\*
0 & -- & 0 & 4 & 1 & 126 & 2 & 128 & \texttt{-???????}\\*
1 & 0 & 1 & 1 & 1 & 192 & 64 & 0 & \texttt{--++----}\\
2 & 1 & 2 & 1 & 1 & 128 & 128 & 0 & \texttt{---+----}\\
3 & 2 & 3 & 1 & 1 & 0 & 256 & 0 & \texttt{--------}\\
4 & 0 & 1 & 3 & 1 & 32 & 32 & 192 & \texttt{-----???}\\
5 & 4 & 2 & 1 & 3 & 0 & 256 & 0 & \texttt{--------}\\
6 & 4 & 2 & 1 & 3 & 0 & 256 & 0 & \texttt{--------}\\
7 & 4 & 2 & 1 & 3 & 0 & 256 & 0 & \texttt{--------}\\
\multicolumn{9}{@{}l@{}}{\rule{0pt}{2.6ex}\textbf{T8.097}\enspace \texttt{01231221}\enspace $n{=}8$, $|H|{=}4$, len 4, br 2, $|\mathrm{Aut}|{=}2$, orb 7/6, hist-orb 3/2; M 3072, res 1480, rev 846: ok}\\*
0 & -- & 0 & 4 & 1 & 70 & 2 & 184 & \texttt{-???????}\\*
1 & 0 & 1 & 1 & 1 & 192 & 64 & 0 & \texttt{--++----}\\
2 & 1 & 2 & 1 & 1 & 128 & 128 & 0 & \texttt{---+----}\\
3 & 2 & 3 & 1 & 1 & 0 & 256 & 0 & \texttt{--------}\\
4 & 0 & 1 & 2 & 1 & 64 & 64 & 128 & \texttt{-----??-}\\
5 & 4 & 2 & 1 & 2 & 0 & 256 & 0 & \texttt{--------}\\
6 & 4 & 2 & 1 & 2 & 0 & 256 & 0 & \texttt{--------}\\
7 & 0 & 1 & 1 & 1 & 0 & 256 & 0 & \texttt{--------}\\
\multicolumn{9}{@{}l@{}}{\rule{0pt}{2.6ex}\textbf{T8.098}\enspace \texttt{01231212}\enspace $n{=}8$, $|H|{=}3$, len 4, br 1, $|\mathrm{Aut}|{=}2$, orb 6/4, hist-orb 2/1; M 3072, res 896, rev 1174: ok}\\*
0 & -- & 0 & 3 & 1 & 126 & 2 & 128 & \texttt{-???????}\\*
1 & 0 & 1 & 1 & 1 & 192 & 64 & 0 & \texttt{--++----}\\
2 & 1 & 2 & 1 & 1 & 128 & 128 & 0 & \texttt{---+----}\\
3 & 2 & 3 & 1 & 1 & 0 & 256 & 0 & \texttt{--------}\\
4 & 0 & 1 & 1 & 2 & 128 & 128 & 0 & \texttt{-----+--}\\
5 & 4 & 2 & 1 & 2 & 0 & 256 & 0 & \texttt{--------}\\
6 & 0 & 1 & 1 & 2 & 128 & 128 & 0 & \texttt{-------+}\\
7 & 6 & 2 & 1 & 2 & 0 & 256 & 0 & \texttt{--------}\\
\multicolumn{9}{@{}l@{}}{\rule{0pt}{2.6ex}\textbf{T8.099}\enspace \texttt{01231211}\enspace $n{=}8$, $|H|{=}4$, len 4, br 1, $|\mathrm{Aut}|{=}2$, orb 7/6, hist-orb 3/2; M 2816, res 1484, rev 718: ok}\\*
0 & -- & 0 & 4 & 1 & 42 & 2 & 212 & \texttt{-???????}\\*
1 & 0 & 1 & 1 & 1 & 192 & 64 & 0 & \texttt{--++----}\\
2 & 1 & 2 & 1 & 1 & 128 & 128 & 0 & \texttt{---+----}\\
3 & 2 & 3 & 1 & 1 & 0 & 256 & 0 & \texttt{--------}\\
4 & 0 & 1 & 1 & 1 & 128 & 128 & 0 & \texttt{-----+--}\\
5 & 4 & 2 & 1 & 1 & 0 & 256 & 0 & \texttt{--------}\\
6 & 0 & 1 & 1 & 2 & 0 & 256 & 0 & \texttt{--------}\\
7 & 0 & 1 & 1 & 2 & 0 & 256 & 0 & \texttt{--------}\\
\multicolumn{9}{@{}l@{}}{\rule{0pt}{2.6ex}\textbf{T8.100}\enspace \texttt{01231111}\enspace $n{=}8$, $|H|{=}5$, len 4, br 1, $|\mathrm{Aut}|{=}24$, orb 5/4, hist-orb 2/1; M 2560, res 1680, rev 462: ok}\\*
0 & -- & 0 & 5 & 1 & 14 & 2 & 240 & \texttt{-???????}\\*
1 & 0 & 1 & 1 & 1 & 192 & 64 & 0 & \texttt{--++----}\\
2 & 1 & 2 & 1 & 1 & 128 & 128 & 0 & \texttt{---+----}\\
3 & 2 & 3 & 1 & 1 & 0 & 256 & 0 & \texttt{--------}\\
4 & 0 & 1 & 1 & 4 & 0 & 256 & 0 & \texttt{--------}\\
5 & 0 & 1 & 1 & 4 & 0 & 256 & 0 & \texttt{--------}\\
6 & 0 & 1 & 1 & 4 & 0 & 256 & 0 & \texttt{--------}\\
7 & 0 & 1 & 1 & 4 & 0 & 256 & 0 & \texttt{--------}\\
\multicolumn{9}{@{}l@{}}{\rule{0pt}{2.6ex}\textbf{T8.101}\enspace \texttt{01222222}\enspace $n{=}8$, $|H|{=}6$, len 3, br 1, $|\mathrm{Aut}|{=}720$, orb 3/2, hist-orb 1/0; M 3328, res 2232, rev 806: ok}\\*
0 & -- & 0 & 6 & 1 & 130 & 2 & 124 & \texttt{-+??????}\\*
1 & 0 & 1 & 6 & 1 & 4 & 4 & 248 & \texttt{--??????}\\
2 & 1 & 2 & 1 & 6 & 0 & 256 & 0 & \texttt{--------}\\
3 & 1 & 2 & 1 & 6 & 0 & 256 & 0 & \texttt{--------}\\
4 & 1 & 2 & 1 & 6 & 0 & 256 & 0 & \texttt{--------}\\
5 & 1 & 2 & 1 & 6 & 0 & 256 & 0 & \texttt{--------}\\
6 & 1 & 2 & 1 & 6 & 0 & 256 & 0 & \texttt{--------}\\
7 & 1 & 2 & 1 & 6 & 0 & 256 & 0 & \texttt{--------}\\
\multicolumn{9}{@{}l@{}}{\rule{0pt}{2.6ex}\textbf{T8.102}\enspace \texttt{01222221}\enspace $n{=}8$, $|H|{=}6$, len 3, br 2, $|\mathrm{Aut}|{=}120$, orb 4/3, hist-orb 2/1; M 3072, res 2396, rev 438: ok}\\*
0 & -- & 0 & 6 & 1 & 66 & 2 & 188 & \texttt{-???????}\\*
1 & 0 & 1 & 5 & 1 & 8 & 8 & 240 & \texttt{--?????-}\\
2 & 1 & 2 & 1 & 5 & 0 & 256 & 0 & \texttt{--------}\\
3 & 1 & 2 & 1 & 5 & 0 & 256 & 0 & \texttt{--------}\\
4 & 1 & 2 & 1 & 5 & 0 & 256 & 0 & \texttt{--------}\\
5 & 1 & 2 & 1 & 5 & 0 & 256 & 0 & \texttt{--------}\\
6 & 1 & 2 & 1 & 5 & 0 & 256 & 0 & \texttt{--------}\\
7 & 0 & 1 & 1 & 1 & 0 & 256 & 0 & \texttt{--------}\\
\multicolumn{9}{@{}l@{}}{\rule{0pt}{2.6ex}\textbf{T8.103}\enspace \texttt{01222212}\enspace $n{=}8$, $|H|{=}5$, len 3, br 2, $|\mathrm{Aut}|{=}24$, orb 5/4, hist-orb 2/1; M 3072, res 1848, rev 742: ok}\\*
0 & -- & 0 & 5 & 1 & 102 & 2 & 152 & \texttt{-???????}\\*
1 & 0 & 1 & 4 & 1 & 16 & 16 & 224 & \texttt{--????--}\\
2 & 1 & 2 & 1 & 4 & 0 & 256 & 0 & \texttt{--------}\\
3 & 1 & 2 & 1 & 4 & 0 & 256 & 0 & \texttt{--------}\\
4 & 1 & 2 & 1 & 4 & 0 & 256 & 0 & \texttt{--------}\\
5 & 1 & 2 & 1 & 4 & 0 & 256 & 0 & \texttt{--------}\\
6 & 0 & 1 & 1 & 1 & 128 & 128 & 0 & \texttt{-------+}\\
7 & 6 & 2 & 1 & 1 & 0 & 256 & 0 & \texttt{--------}\\
\multicolumn{9}{@{}l@{}}{\rule{0pt}{2.6ex}\textbf{T8.104}\enspace \texttt{01222211}\enspace $n{=}8$, $|H|{=}6$, len 3, br 2, $|\mathrm{Aut}|{=}48$, orb 4/2, hist-orb 2/0; M 2816, res 2324, rev 270: ok}\\*
0 & -- & 0 & 6 & 1 & 34 & 2 & 220 & \texttt{-???????}\\*
1 & 0 & 1 & 4 & 1 & 16 & 16 & 224 & \texttt{--????--}\\
2 & 1 & 2 & 1 & 4 & 0 & 256 & 0 & \texttt{--------}\\
3 & 1 & 2 & 1 & 4 & 0 & 256 & 0 & \texttt{--------}\\
4 & 1 & 2 & 1 & 4 & 0 & 256 & 0 & \texttt{--------}\\
5 & 1 & 2 & 1 & 4 & 0 & 256 & 0 & \texttt{--------}\\
6 & 0 & 1 & 1 & 2 & 0 & 256 & 0 & \texttt{--------}\\
7 & 0 & 1 & 1 & 2 & 0 & 256 & 0 & \texttt{--------}\\
\multicolumn{9}{@{}l@{}}{\rule{0pt}{2.6ex}\textbf{T8.105}\enspace \texttt{01222122}\enspace $n{=}8$, $|H|{=}5$, len 3, br 3, $|\mathrm{Aut}|{=}12$, orb 5/3, hist-orb 2/0; M 3072, res 1820, rev 702: ok}\\*
0 & -- & 0 & 5 & 1 & 90 & 2 & 164 & \texttt{-???????}\\*
1 & 0 & 1 & 3 & 1 & 32 & 32 & 192 & \texttt{--???---}\\
2 & 1 & 2 & 1 & 3 & 0 & 256 & 0 & \texttt{--------}\\
3 & 1 & 2 & 1 & 3 & 0 & 256 & 0 & \texttt{--------}\\
4 & 1 & 2 & 1 & 3 & 0 & 256 & 0 & \texttt{--------}\\
5 & 0 & 1 & 2 & 1 & 64 & 64 & 128 & \texttt{------??}\\
6 & 5 & 2 & 1 & 2 & 0 & 256 & 0 & \texttt{--------}\\
7 & 5 & 2 & 1 & 2 & 0 & 256 & 0 & \texttt{--------}\\
\multicolumn{9}{@{}l@{}}{\rule{0pt}{2.6ex}\textbf{T8.106}\enspace \texttt{01222121}\enspace $n{=}8$, $|H|{=}5$, len 3, br 2, $|\mathrm{Aut}|{=}6$, orb 6/5, hist-orb 3/2; M 2816, res 1880, rev 518: ok}\\*
0 & -- & 0 & 5 & 1 & 54 & 2 & 200 & \texttt{-???????}\\*
1 & 0 & 1 & 3 & 1 & 32 & 32 & 192 & \texttt{--???---}\\
2 & 1 & 2 & 1 & 3 & 0 & 256 & 0 & \texttt{--------}\\
3 & 1 & 2 & 1 & 3 & 0 & 256 & 0 & \texttt{--------}\\
4 & 1 & 2 & 1 & 3 & 0 & 256 & 0 & \texttt{--------}\\
5 & 0 & 1 & 1 & 1 & 128 & 128 & 0 & \texttt{------+-}\\
6 & 5 & 2 & 1 & 1 & 0 & 256 & 0 & \texttt{--------}\\
7 & 0 & 1 & 1 & 1 & 0 & 256 & 0 & \texttt{--------}\\
\multicolumn{9}{@{}l@{}}{\rule{0pt}{2.6ex}\textbf{T8.107}\enspace \texttt{01222111}\enspace $n{=}8$, $|H|{=}6$, len 3, br 2, $|\mathrm{Aut}|{=}36$, orb 4/2, hist-orb 2/0; M 2560, res 2132, rev 206: ok}\\*
0 & -- & 0 & 6 & 1 & 18 & 2 & 236 & \texttt{-???????}\\*
1 & 0 & 1 & 3 & 1 & 32 & 32 & 192 & \texttt{--???---}\\
2 & 1 & 2 & 1 & 3 & 0 & 256 & 0 & \texttt{--------}\\
3 & 1 & 2 & 1 & 3 & 0 & 256 & 0 & \texttt{--------}\\
4 & 1 & 2 & 1 & 3 & 0 & 256 & 0 & \texttt{--------}\\
5 & 0 & 1 & 1 & 3 & 0 & 256 & 0 & \texttt{--------}\\
6 & 0 & 1 & 1 & 3 & 0 & 256 & 0 & \texttt{--------}\\
7 & 0 & 1 & 1 & 3 & 0 & 256 & 0 & \texttt{--------}\\
\multicolumn{9}{@{}l@{}}{\rule{0pt}{2.6ex}\textbf{T8.108}\enspace \texttt{01221221}\enspace $n{=}8$, $|H|{=}5$, len 3, br 3, $|\mathrm{Aut}|{=}8$, orb 4/2, hist-orb 2/1; M 2816, res 1812, rev 526: ok}\\*
0 & -- & 0 & 5 & 1 & 50 & 2 & 204 & \texttt{-???????}\\*
1 & 0 & 1 & 2 & 2 & 64 & 64 & 128 & \texttt{--??----}\\
2 & 1 & 2 & 1 & 4 & 0 & 256 & 0 & \texttt{--------}\\
3 & 1 & 2 & 1 & 4 & 0 & 256 & 0 & \texttt{--------}\\
4 & 0 & 1 & 2 & 2 & 64 & 64 & 128 & \texttt{-----??-}\\
5 & 4 & 2 & 1 & 4 & 0 & 256 & 0 & \texttt{--------}\\
6 & 4 & 2 & 1 & 4 & 0 & 256 & 0 & \texttt{--------}\\
7 & 0 & 1 & 1 & 1 & 0 & 256 & 0 & \texttt{--------}\\
\multicolumn{9}{@{}l@{}}{\rule{0pt}{2.6ex}\textbf{T8.109}\enspace \texttt{01221212}\enspace $n{=}8$, $|H|{=}4$, len 3, br 2, $|\mathrm{Aut}|{=}4$, orb 5/2, hist-orb 2/0; M 2816, res 1340, rev 822: ok}\\*
0 & -- & 0 & 4 & 1 & 90 & 2 & 164 & \texttt{-???????}\\*
1 & 0 & 1 & 2 & 1 & 64 & 64 & 128 & \texttt{--??----}\\
2 & 1 & 2 & 1 & 2 & 0 & 256 & 0 & \texttt{--------}\\
3 & 1 & 2 & 1 & 2 & 0 & 256 & 0 & \texttt{--------}\\
4 & 0 & 1 & 1 & 2 & 128 & 128 & 0 & \texttt{-----+--}\\
5 & 4 & 2 & 1 & 2 & 0 & 256 & 0 & \texttt{--------}\\
6 & 0 & 1 & 1 & 2 & 128 & 128 & 0 & \texttt{-------+}\\
7 & 6 & 2 & 1 & 2 & 0 & 256 & 0 & \texttt{--------}\\
\multicolumn{9}{@{}l@{}}{\rule{0pt}{2.6ex}\textbf{T8.110}\enspace \texttt{01221211}\enspace $n{=}8$, $|H|{=}5$, len 3, br 2, $|\mathrm{Aut}|{=}4$, orb 6/4, hist-orb 3/1; M 2560, res 1760, rev 422: ok}\\*
0 & -- & 0 & 5 & 1 & 30 & 2 & 224 & \texttt{-???????}\\*
1 & 0 & 1 & 2 & 1 & 64 & 64 & 128 & \texttt{--??----}\\
2 & 1 & 2 & 1 & 2 & 0 & 256 & 0 & \texttt{--------}\\
3 & 1 & 2 & 1 & 2 & 0 & 256 & 0 & \texttt{--------}\\
4 & 0 & 1 & 1 & 1 & 128 & 128 & 0 & \texttt{-----+--}\\
5 & 4 & 2 & 1 & 1 & 0 & 256 & 0 & \texttt{--------}\\
6 & 0 & 1 & 1 & 2 & 0 & 256 & 0 & \texttt{--------}\\
7 & 0 & 1 & 1 & 2 & 0 & 256 & 0 & \texttt{--------}\\
\multicolumn{9}{@{}l@{}}{\rule{0pt}{2.6ex}\textbf{T8.111}\enspace \texttt{01221111}\enspace $n{=}8$, $|H|{=}6$, len 3, br 2, $|\mathrm{Aut}|{=}48$, orb 4/2, hist-orb 2/0; M 2304, res 1900, rev 190: ok}\\*
0 & -- & 0 & 6 & 1 & 10 & 2 & 244 & \texttt{-???????}\\*
1 & 0 & 1 & 2 & 1 & 64 & 64 & 128 & \texttt{--??----}\\
2 & 1 & 2 & 1 & 2 & 0 & 256 & 0 & \texttt{--------}\\
3 & 1 & 2 & 1 & 2 & 0 & 256 & 0 & \texttt{--------}\\
4 & 0 & 1 & 1 & 4 & 0 & 256 & 0 & \texttt{--------}\\
5 & 0 & 1 & 1 & 4 & 0 & 256 & 0 & \texttt{--------}\\
6 & 0 & 1 & 1 & 4 & 0 & 256 & 0 & \texttt{--------}\\
7 & 0 & 1 & 1 & 4 & 0 & 256 & 0 & \texttt{--------}\\
\multicolumn{9}{@{}l@{}}{\rule{0pt}{2.6ex}\textbf{T8.112}\enspace \texttt{01212121}\enspace $n{=}8$, $|H|{=}4$, len 3, br 1, $|\mathrm{Aut}|{=}6$, orb 4/2, hist-orb 2/1; M 2560, res 1400, rev 654: ok}\\*
0 & -- & 0 & 4 & 1 & 54 & 2 & 200 & \texttt{-???????}\\*
1 & 0 & 1 & 1 & 3 & 128 & 128 & 0 & \texttt{--+-----}\\
2 & 1 & 2 & 1 & 3 & 0 & 256 & 0 & \texttt{--------}\\
3 & 0 & 1 & 1 & 3 & 128 & 128 & 0 & \texttt{----+---}\\
4 & 3 & 2 & 1 & 3 & 0 & 256 & 0 & \texttt{--------}\\
5 & 0 & 1 & 1 & 3 & 128 & 128 & 0 & \texttt{------+-}\\
6 & 5 & 2 & 1 & 3 & 0 & 256 & 0 & \texttt{--------}\\
7 & 0 & 1 & 1 & 1 & 0 & 256 & 0 & \texttt{--------}\\
\multicolumn{9}{@{}l@{}}{\rule{0pt}{2.6ex}\textbf{T8.113}\enspace \texttt{01212111}\enspace $n{=}8$, $|H|{=}5$, len 3, br 1, $|\mathrm{Aut}|{=}12$, orb 4/1, hist-orb 2/0; M 2304, res 1652, rev 358: ok}\\*
0 & -- & 0 & 5 & 1 & 18 & 2 & 236 & \texttt{-???????}\\*
1 & 0 & 1 & 1 & 2 & 128 & 128 & 0 & \texttt{--+-----}\\
2 & 1 & 2 & 1 & 2 & 0 & 256 & 0 & \texttt{--------}\\
3 & 0 & 1 & 1 & 2 & 128 & 128 & 0 & \texttt{----+---}\\
4 & 3 & 2 & 1 & 2 & 0 & 256 & 0 & \texttt{--------}\\
5 & 0 & 1 & 1 & 3 & 0 & 256 & 0 & \texttt{--------}\\
6 & 0 & 1 & 1 & 3 & 0 & 256 & 0 & \texttt{--------}\\
7 & 0 & 1 & 1 & 3 & 0 & 256 & 0 & \texttt{--------}\\
\multicolumn{9}{@{}l@{}}{\rule{0pt}{2.6ex}\textbf{T8.114}\enspace \texttt{01211111}\enspace $n{=}8$, $|H|{=}6$, len 3, br 1, $|\mathrm{Aut}|{=}120$, orb 4/3, hist-orb 2/1; M 2048, res 1736, rev 166: ok}\\*
0 & -- & 0 & 6 & 1 & 6 & 2 & 248 & \texttt{-???????}\\*
1 & 0 & 1 & 1 & 1 & 128 & 128 & 0 & \texttt{--+-----}\\
2 & 1 & 2 & 1 & 1 & 0 & 256 & 0 & \texttt{--------}\\
3 & 0 & 1 & 1 & 5 & 0 & 256 & 0 & \texttt{--------}\\
4 & 0 & 1 & 1 & 5 & 0 & 256 & 0 & \texttt{--------}\\
5 & 0 & 1 & 1 & 5 & 0 & 256 & 0 & \texttt{--------}\\
6 & 0 & 1 & 1 & 5 & 0 & 256 & 0 & \texttt{--------}\\
7 & 0 & 1 & 1 & 5 & 0 & 256 & 0 & \texttt{--------}\\
\multicolumn{9}{@{}l@{}}{\rule{0pt}{2.6ex}\textbf{T8.115}\enspace \texttt{01111111}\enspace $n{=}8$, $|H|{=}7$, len 2, br 1, $|\mathrm{Aut}|{=}5040$, orb 2/1, hist-orb 1/0; M 1792, res 1764, rev 14: ok}\\*
0 & -- & 0 & 7 & 1 & 2 & 2 & 252 & \texttt{-???????}\\*
1 & 0 & 1 & 1 & 7 & 0 & 256 & 0 & \texttt{--------}\\
2 & 0 & 1 & 1 & 7 & 0 & 256 & 0 & \texttt{--------}\\
3 & 0 & 1 & 1 & 7 & 0 & 256 & 0 & \texttt{--------}\\
4 & 0 & 1 & 1 & 7 & 0 & 256 & 0 & \texttt{--------}\\
5 & 0 & 1 & 1 & 7 & 0 & 256 & 0 & \texttt{--------}\\
6 & 0 & 1 & 1 & 7 & 0 & 256 & 0 & \texttt{--------}\\
7 & 0 & 1 & 1 & 7 & 0 & 256 & 0 & \texttt{--------}\\
\multicolumn{9}{@{}l@{}}{\rule{0pt}{3.4ex}\normalsize\bfseries Frames with 9 moments (286 classes)}\\*
\multicolumn{9}{@{}l@{}}{\rule{0pt}{2.6ex}\textbf{T9.001}\enspace \texttt{012345678}\enspace $n{=}9$, $|H|{=}1$, len 9, br 0, $|\mathrm{Aut}|{=}1$, orb 9/9, hist-orb 1/1; M 18432, res 0, rev 6168: ok}\\*
0 & -- & 0 & 1 & 1 & 510 & 2 & 0 & \texttt{-++++++++}\\*
1 & 0 & 1 & 1 & 1 & 508 & 4 & 0 & \texttt{--+++++++}\\
2 & 1 & 2 & 1 & 1 & 504 & 8 & 0 & \texttt{---++++++}\\
3 & 2 & 3 & 1 & 1 & 496 & 16 & 0 & \texttt{----+++++}\\
4 & 3 & 4 & 1 & 1 & 480 & 32 & 0 & \texttt{-----++++}\\
5 & 4 & 5 & 1 & 1 & 448 & 64 & 0 & \texttt{------+++}\\
6 & 5 & 6 & 1 & 1 & 384 & 128 & 0 & \texttt{-------++}\\
7 & 6 & 7 & 1 & 1 & 256 & 256 & 0 & \texttt{--------+}\\
8 & 7 & 8 & 1 & 1 & 0 & 512 & 0 & \texttt{---------}\\
\multicolumn{9}{@{}l@{}}{\rule{0pt}{2.6ex}\textbf{T9.002}\enspace \texttt{012345677}\enspace $n{=}9$, $|H|{=}2$, len 8, br 1, $|\mathrm{Aut}|{=}2$, orb 8/7, hist-orb 1/0; M 17920, res 1016, rev 6688: ok}\\*
0 & -- & 0 & 2 & 1 & 506 & 2 & 4 & \texttt{-++++++??}\\*
1 & 0 & 1 & 2 & 1 & 500 & 4 & 8 & \texttt{--+++++??}\\
2 & 1 & 2 & 2 & 1 & 488 & 8 & 16 & \texttt{---++++??}\\
3 & 2 & 3 & 2 & 1 & 464 & 16 & 32 & \texttt{----+++??}\\
4 & 3 & 4 & 2 & 1 & 416 & 32 & 64 & \texttt{-----++??}\\
5 & 4 & 5 & 2 & 1 & 320 & 64 & 128 & \texttt{------+??}\\
6 & 5 & 6 & 2 & 1 & 128 & 128 & 256 & \texttt{-------??}\\
7 & 6 & 7 & 1 & 2 & 0 & 512 & 0 & \texttt{---------}\\
8 & 6 & 7 & 1 & 2 & 0 & 512 & 0 & \texttt{---------}\\
\multicolumn{9}{@{}l@{}}{\rule{0pt}{2.6ex}\textbf{T9.003}\enspace \texttt{012345676}\enspace $n{=}9$, $|H|{=}2$, len 8, br 1, $|\mathrm{Aut}|{=}1$, orb 9/9, hist-orb 2/2; M 17408, res 1512, rev 6700: ok}\\*
0 & -- & 0 & 2 & 1 & 502 & 2 & 8 & \texttt{-+++++???}\\*
1 & 0 & 1 & 2 & 1 & 492 & 4 & 16 & \texttt{--++++???}\\
2 & 1 & 2 & 2 & 1 & 472 & 8 & 32 & \texttt{---+++???}\\
3 & 2 & 3 & 2 & 1 & 432 & 16 & 64 & \texttt{----++???}\\
4 & 3 & 4 & 2 & 1 & 352 & 32 & 128 & \texttt{-----+???}\\
5 & 4 & 5 & 2 & 1 & 192 & 64 & 256 & \texttt{------???}\\
6 & 5 & 6 & 1 & 1 & 256 & 256 & 0 & \texttt{-------+-}\\
7 & 6 & 7 & 1 & 1 & 0 & 512 & 0 & \texttt{---------}\\
8 & 5 & 6 & 1 & 1 & 0 & 512 & 0 & \texttt{---------}\\
\multicolumn{9}{@{}l@{}}{\rule{0pt}{2.6ex}\textbf{T9.004}\enspace \texttt{012345675}\enspace $n{=}9$, $|H|{=}2$, len 8, br 1, $|\mathrm{Aut}|{=}1$, orb 9/9, hist-orb 2/2; M 16896, res 1984, rev 6340: ok}\\*
0 & -- & 0 & 2 & 1 & 494 & 2 & 16 & \texttt{-++++????}\\*
1 & 0 & 1 & 2 & 1 & 476 & 4 & 32 & \texttt{--+++????}\\
2 & 1 & 2 & 2 & 1 & 440 & 8 & 64 & \texttt{---++????}\\
3 & 2 & 3 & 2 & 1 & 368 & 16 & 128 & \texttt{----+????}\\
4 & 3 & 4 & 2 & 1 & 224 & 32 & 256 & \texttt{-----????}\\
5 & 4 & 5 & 1 & 1 & 384 & 128 & 0 & \texttt{------++-}\\
6 & 5 & 6 & 1 & 1 & 256 & 256 & 0 & \texttt{-------+-}\\
7 & 6 & 7 & 1 & 1 & 0 & 512 & 0 & \texttt{---------}\\
8 & 4 & 5 & 1 & 1 & 0 & 512 & 0 & \texttt{---------}\\
\multicolumn{9}{@{}l@{}}{\rule{0pt}{2.6ex}\textbf{T9.005}\enspace \texttt{012345674}\enspace $n{=}9$, $|H|{=}2$, len 8, br 1, $|\mathrm{Aut}|{=}1$, orb 9/9, hist-orb 2/2; M 16384, res 2400, rev 5876: ok}\\*
0 & -- & 0 & 2 & 1 & 478 & 2 & 32 & \texttt{-+++?????}\\*
1 & 0 & 1 & 2 & 1 & 444 & 4 & 64 & \texttt{--++?????}\\
2 & 1 & 2 & 2 & 1 & 376 & 8 & 128 & \texttt{---+?????}\\
3 & 2 & 3 & 2 & 1 & 240 & 16 & 256 & \texttt{----?????}\\
4 & 3 & 4 & 1 & 1 & 448 & 64 & 0 & \texttt{-----+++-}\\
5 & 4 & 5 & 1 & 1 & 384 & 128 & 0 & \texttt{------++-}\\
6 & 5 & 6 & 1 & 1 & 256 & 256 & 0 & \texttt{-------+-}\\
7 & 6 & 7 & 1 & 1 & 0 & 512 & 0 & \texttt{---------}\\
8 & 3 & 4 & 1 & 1 & 0 & 512 & 0 & \texttt{---------}\\
\multicolumn{9}{@{}l@{}}{\rule{0pt}{2.6ex}\textbf{T9.006}\enspace \texttt{012345673}\enspace $n{=}9$, $|H|{=}2$, len 8, br 1, $|\mathrm{Aut}|{=}1$, orb 9/9, hist-orb 2/2; M 15872, res 2688, rev 5428: ok}\\*
0 & -- & 0 & 2 & 1 & 446 & 2 & 64 & \texttt{-++??????}\\*
1 & 0 & 1 & 2 & 1 & 380 & 4 & 128 & \texttt{--+??????}\\
2 & 1 & 2 & 2 & 1 & 248 & 8 & 256 & \texttt{---??????}\\
3 & 2 & 3 & 1 & 1 & 480 & 32 & 0 & \texttt{----++++-}\\
4 & 3 & 4 & 1 & 1 & 448 & 64 & 0 & \texttt{-----+++-}\\
5 & 4 & 5 & 1 & 1 & 384 & 128 & 0 & \texttt{------++-}\\
6 & 5 & 6 & 1 & 1 & 256 & 256 & 0 & \texttt{-------+-}\\
7 & 6 & 7 & 1 & 1 & 0 & 512 & 0 & \texttt{---------}\\
8 & 2 & 3 & 1 & 1 & 0 & 512 & 0 & \texttt{---------}\\
\multicolumn{9}{@{}l@{}}{\rule{0pt}{2.6ex}\textbf{T9.007}\enspace \texttt{012345672}\enspace $n{=}9$, $|H|{=}2$, len 8, br 1, $|\mathrm{Aut}|{=}1$, orb 9/9, hist-orb 2/2; M 15360, res 2688, rev 5076: ok}\\*
0 & -- & 0 & 2 & 1 & 382 & 2 & 128 & \texttt{-+???????}\\*
1 & 0 & 1 & 2 & 1 & 252 & 4 & 256 & \texttt{--???????}\\
2 & 1 & 2 & 1 & 1 & 496 & 16 & 0 & \texttt{---+++++-}\\
3 & 2 & 3 & 1 & 1 & 480 & 32 & 0 & \texttt{----++++-}\\
4 & 3 & 4 & 1 & 1 & 448 & 64 & 0 & \texttt{-----+++-}\\
5 & 4 & 5 & 1 & 1 & 384 & 128 & 0 & \texttt{------++-}\\
6 & 5 & 6 & 1 & 1 & 256 & 256 & 0 & \texttt{-------+-}\\
7 & 6 & 7 & 1 & 1 & 0 & 512 & 0 & \texttt{---------}\\
8 & 1 & 2 & 1 & 1 & 0 & 512 & 0 & \texttt{---------}\\
\multicolumn{9}{@{}l@{}}{\rule{0pt}{2.6ex}\textbf{T9.008}\enspace \texttt{012345671}\enspace $n{=}9$, $|H|{=}2$, len 8, br 1, $|\mathrm{Aut}|{=}1$, orb 9/9, hist-orb 2/2; M 14848, res 2048, rev 4924: ok}\\*
0 & -- & 0 & 2 & 1 & 254 & 2 & 256 & \texttt{-????????}\\*
1 & 0 & 1 & 1 & 1 & 504 & 8 & 0 & \texttt{--++++++-}\\
2 & 1 & 2 & 1 & 1 & 496 & 16 & 0 & \texttt{---+++++-}\\
3 & 2 & 3 & 1 & 1 & 480 & 32 & 0 & \texttt{----++++-}\\
4 & 3 & 4 & 1 & 1 & 448 & 64 & 0 & \texttt{-----+++-}\\
5 & 4 & 5 & 1 & 1 & 384 & 128 & 0 & \texttt{------++-}\\
6 & 5 & 6 & 1 & 1 & 256 & 256 & 0 & \texttt{-------+-}\\
7 & 6 & 7 & 1 & 1 & 0 & 512 & 0 & \texttt{---------}\\
8 & 0 & 1 & 1 & 1 & 0 & 512 & 0 & \texttt{---------}\\
\multicolumn{9}{@{}l@{}}{\rule{0pt}{2.6ex}\textbf{T9.009}\enspace \texttt{012345666}\enspace $n{=}9$, $|H|{=}3$, len 7, br 1, $|\mathrm{Aut}|{=}6$, orb 7/6, hist-orb 1/0; M 16896, res 2268, rev 6972: ok}\\*
0 & -- & 0 & 3 & 1 & 498 & 2 & 12 & \texttt{-+++++???}\\*
1 & 0 & 1 & 3 & 1 & 484 & 4 & 24 & \texttt{--++++???}\\
2 & 1 & 2 & 3 & 1 & 456 & 8 & 48 & \texttt{---+++???}\\
3 & 2 & 3 & 3 & 1 & 400 & 16 & 96 & \texttt{----++???}\\
4 & 3 & 4 & 3 & 1 & 288 & 32 & 192 & \texttt{-----+???}\\
5 & 4 & 5 & 3 & 1 & 64 & 64 & 384 & \texttt{------???}\\
6 & 5 & 6 & 1 & 3 & 0 & 512 & 0 & \texttt{---------}\\
7 & 5 & 6 & 1 & 3 & 0 & 512 & 0 & \texttt{---------}\\
8 & 5 & 6 & 1 & 3 & 0 & 512 & 0 & \texttt{---------}\\
\multicolumn{9}{@{}l@{}}{\rule{0pt}{2.6ex}\textbf{T9.010}\enspace \texttt{012345665}\enspace $n{=}9$, $|H|{=}3$, len 7, br 2, $|\mathrm{Aut}|{=}2$, orb 8/7, hist-orb 2/1; M 16384, res 2744, rev 6488: ok}\\*
0 & -- & 0 & 3 & 1 & 490 & 2 & 20 & \texttt{-++++????}\\*
1 & 0 & 1 & 3 & 1 & 468 & 4 & 40 & \texttt{--+++????}\\
2 & 1 & 2 & 3 & 1 & 424 & 8 & 80 & \texttt{---++????}\\
3 & 2 & 3 & 3 & 1 & 336 & 16 & 160 & \texttt{----+????}\\
4 & 3 & 4 & 3 & 1 & 160 & 32 & 320 & \texttt{-----????}\\
5 & 4 & 5 & 2 & 1 & 128 & 128 & 256 & \texttt{------??-}\\
6 & 5 & 6 & 1 & 2 & 0 & 512 & 0 & \texttt{---------}\\
7 & 5 & 6 & 1 & 2 & 0 & 512 & 0 & \texttt{---------}\\
8 & 4 & 5 & 1 & 1 & 0 & 512 & 0 & \texttt{---------}\\
\multicolumn{9}{@{}l@{}}{\rule{0pt}{2.6ex}\textbf{T9.011}\enspace \texttt{012345664}\enspace $n{=}9$, $|H|{=}3$, len 7, br 2, $|\mathrm{Aut}|{=}2$, orb 8/7, hist-orb 2/1; M 15872, res 3228, rev 5968: ok}\\*
0 & -- & 0 & 3 & 1 & 474 & 2 & 36 & \texttt{-+++?????}\\*
1 & 0 & 1 & 3 & 1 & 436 & 4 & 72 & \texttt{--++?????}\\
2 & 1 & 2 & 3 & 1 & 360 & 8 & 144 & \texttt{---+?????}\\
3 & 2 & 3 & 3 & 1 & 208 & 16 & 288 & \texttt{----?????}\\
4 & 3 & 4 & 2 & 1 & 320 & 64 & 128 & \texttt{-----+??-}\\
5 & 4 & 5 & 2 & 1 & 128 & 128 & 256 & \texttt{------??-}\\
6 & 5 & 6 & 1 & 2 & 0 & 512 & 0 & \texttt{---------}\\
7 & 5 & 6 & 1 & 2 & 0 & 512 & 0 & \texttt{---------}\\
8 & 3 & 4 & 1 & 1 & 0 & 512 & 0 & \texttt{---------}\\
\multicolumn{9}{@{}l@{}}{\rule{0pt}{2.6ex}\textbf{T9.012}\enspace \texttt{012345663}\enspace $n{=}9$, $|H|{=}3$, len 7, br 2, $|\mathrm{Aut}|{=}2$, orb 8/7, hist-orb 2/1; M 15360, res 3584, rev 5504: ok}\\*
0 & -- & 0 & 3 & 1 & 442 & 2 & 68 & \texttt{-++??????}\\*
1 & 0 & 1 & 3 & 1 & 372 & 4 & 136 & \texttt{--+??????}\\
2 & 1 & 2 & 3 & 1 & 232 & 8 & 272 & \texttt{---??????}\\
3 & 2 & 3 & 2 & 1 & 416 & 32 & 64 & \texttt{----++??-}\\
4 & 3 & 4 & 2 & 1 & 320 & 64 & 128 & \texttt{-----+??-}\\
5 & 4 & 5 & 2 & 1 & 128 & 128 & 256 & \texttt{------??-}\\
6 & 5 & 6 & 1 & 2 & 0 & 512 & 0 & \texttt{---------}\\
7 & 5 & 6 & 1 & 2 & 0 & 512 & 0 & \texttt{---------}\\
8 & 2 & 3 & 1 & 1 & 0 & 512 & 0 & \texttt{---------}\\
\multicolumn{9}{@{}l@{}}{\rule{0pt}{2.6ex}\textbf{T9.013}\enspace \texttt{012345662}\enspace $n{=}9$, $|H|{=}3$, len 7, br 2, $|\mathrm{Aut}|{=}2$, orb 8/7, hist-orb 2/1; M 14848, res 3636, rev 5168: ok}\\*
0 & -- & 0 & 3 & 1 & 378 & 2 & 132 & \texttt{-+???????}\\*
1 & 0 & 1 & 3 & 1 & 244 & 4 & 264 & \texttt{--???????}\\
2 & 1 & 2 & 2 & 1 & 464 & 16 & 32 & \texttt{---+++??-}\\
3 & 2 & 3 & 2 & 1 & 416 & 32 & 64 & \texttt{----++??-}\\
4 & 3 & 4 & 2 & 1 & 320 & 64 & 128 & \texttt{-----+??-}\\
5 & 4 & 5 & 2 & 1 & 128 & 128 & 256 & \texttt{------??-}\\
6 & 5 & 6 & 1 & 2 & 0 & 512 & 0 & \texttt{---------}\\
7 & 5 & 6 & 1 & 2 & 0 & 512 & 0 & \texttt{---------}\\
8 & 1 & 2 & 1 & 1 & 0 & 512 & 0 & \texttt{---------}\\
\multicolumn{9}{@{}l@{}}{\rule{0pt}{2.6ex}\textbf{T9.014}\enspace \texttt{012345661}\enspace $n{=}9$, $|H|{=}3$, len 7, br 2, $|\mathrm{Aut}|{=}2$, orb 8/7, hist-orb 2/1; M 14336, res 3032, rev 5072: ok}\\*
0 & -- & 0 & 3 & 1 & 250 & 2 & 260 & \texttt{-????????}\\*
1 & 0 & 1 & 2 & 1 & 488 & 8 & 16 & \texttt{--++++??-}\\
2 & 1 & 2 & 2 & 1 & 464 & 16 & 32 & \texttt{---+++??-}\\
3 & 2 & 3 & 2 & 1 & 416 & 32 & 64 & \texttt{----++??-}\\
4 & 3 & 4 & 2 & 1 & 320 & 64 & 128 & \texttt{-----+??-}\\
5 & 4 & 5 & 2 & 1 & 128 & 128 & 256 & \texttt{------??-}\\
6 & 5 & 6 & 1 & 2 & 0 & 512 & 0 & \texttt{---------}\\
7 & 5 & 6 & 1 & 2 & 0 & 512 & 0 & \texttt{---------}\\
8 & 0 & 1 & 1 & 1 & 0 & 512 & 0 & \texttt{---------}\\
\multicolumn{9}{@{}l@{}}{\rule{0pt}{2.6ex}\textbf{T9.015}\enspace \texttt{012345656}\enspace $n{=}9$, $|H|{=}2$, len 7, br 1, $|\mathrm{Aut}|{=}2$, orb 7/5, hist-orb 1/0; M 16384, res 1488, rev 6848: ok}\\*
0 & -- & 0 & 2 & 1 & 498 & 2 & 12 & \texttt{-++++????}\\*
1 & 0 & 1 & 2 & 1 & 484 & 4 & 24 & \texttt{--+++????}\\
2 & 1 & 2 & 2 & 1 & 456 & 8 & 48 & \texttt{---++????}\\
3 & 2 & 3 & 2 & 1 & 400 & 16 & 96 & \texttt{----+????}\\
4 & 3 & 4 & 2 & 1 & 288 & 32 & 192 & \texttt{-----????}\\
5 & 4 & 5 & 1 & 2 & 256 & 256 & 0 & \texttt{------+--}\\
6 & 5 & 6 & 1 & 2 & 0 & 512 & 0 & \texttt{---------}\\
7 & 4 & 5 & 1 & 2 & 256 & 256 & 0 & \texttt{--------+}\\
8 & 7 & 6 & 1 & 2 & 0 & 512 & 0 & \texttt{---------}\\
\multicolumn{9}{@{}l@{}}{\rule{0pt}{2.6ex}\textbf{T9.016}\enspace \texttt{012345655}\enspace $n{=}9$, $|H|{=}3$, len 7, br 1, $|\mathrm{Aut}|{=}2$, orb 8/7, hist-orb 2/1; M 15872, res 2976, rev 6508: ok}\\*
0 & -- & 0 & 3 & 1 & 486 & 2 & 24 & \texttt{-++++????}\\*
1 & 0 & 1 & 3 & 1 & 460 & 4 & 48 & \texttt{--+++????}\\
2 & 1 & 2 & 3 & 1 & 408 & 8 & 96 & \texttt{---++????}\\
3 & 2 & 3 & 3 & 1 & 304 & 16 & 192 & \texttt{----+????}\\
4 & 3 & 4 & 3 & 1 & 96 & 32 & 384 & \texttt{-----????}\\
5 & 4 & 5 & 1 & 1 & 256 & 256 & 0 & \texttt{------+--}\\
6 & 5 & 6 & 1 & 1 & 0 & 512 & 0 & \texttt{---------}\\
7 & 4 & 5 & 1 & 2 & 0 & 512 & 0 & \texttt{---------}\\
8 & 4 & 5 & 1 & 2 & 0 & 512 & 0 & \texttt{---------}\\
\multicolumn{9}{@{}l@{}}{\rule{0pt}{2.6ex}\textbf{T9.017}\enspace \texttt{012345654}\enspace $n{=}9$, $|H|{=}3$, len 7, br 2, $|\mathrm{Aut}|{=}1$, orb 9/9, hist-orb 3/3; M 15360, res 3528, rev 5864: ok}\\*
0 & -- & 0 & 3 & 1 & 470 & 2 & 40 & \texttt{-+++?????}\\*
1 & 0 & 1 & 3 & 1 & 428 & 4 & 80 & \texttt{--++?????}\\
2 & 1 & 2 & 3 & 1 & 344 & 8 & 160 & \texttt{---+?????}\\
3 & 2 & 3 & 3 & 1 & 176 & 16 & 320 & \texttt{----?????}\\
4 & 3 & 4 & 2 & 1 & 192 & 64 & 256 & \texttt{-----???-}\\
5 & 4 & 5 & 1 & 1 & 256 & 256 & 0 & \texttt{------+--}\\
6 & 5 & 6 & 1 & 1 & 0 & 512 & 0 & \texttt{---------}\\
7 & 4 & 5 & 1 & 1 & 0 & 512 & 0 & \texttt{---------}\\
8 & 3 & 4 & 1 & 1 & 0 & 512 & 0 & \texttt{---------}\\
\multicolumn{9}{@{}l@{}}{\rule{0pt}{2.6ex}\textbf{T9.018}\enspace \texttt{012345653}\enspace $n{=}9$, $|H|{=}3$, len 7, br 2, $|\mathrm{Aut}|{=}1$, orb 9/9, hist-orb 3/3; M 14848, res 3952, rev 5344: ok}\\*
0 & -- & 0 & 3 & 1 & 438 & 2 & 72 & \texttt{-++??????}\\*
1 & 0 & 1 & 3 & 1 & 364 & 4 & 144 & \texttt{--+??????}\\
2 & 1 & 2 & 3 & 1 & 216 & 8 & 288 & \texttt{---??????}\\
3 & 2 & 3 & 2 & 1 & 352 & 32 & 128 & \texttt{----+???-}\\
4 & 3 & 4 & 2 & 1 & 192 & 64 & 256 & \texttt{-----???-}\\
5 & 4 & 5 & 1 & 1 & 256 & 256 & 0 & \texttt{------+--}\\
6 & 5 & 6 & 1 & 1 & 0 & 512 & 0 & \texttt{---------}\\
7 & 4 & 5 & 1 & 1 & 0 & 512 & 0 & \texttt{---------}\\
8 & 2 & 3 & 1 & 1 & 0 & 512 & 0 & \texttt{---------}\\
\multicolumn{9}{@{}l@{}}{\rule{0pt}{2.6ex}\textbf{T9.019}\enspace \texttt{012345652}\enspace $n{=}9$, $|H|{=}3$, len 7, br 2, $|\mathrm{Aut}|{=}1$, orb 9/9, hist-orb 3/3; M 14336, res 4056, rev 4992: ok}\\*
0 & -- & 0 & 3 & 1 & 374 & 2 & 136 & \texttt{-+???????}\\*
1 & 0 & 1 & 3 & 1 & 236 & 4 & 272 & \texttt{--???????}\\
2 & 1 & 2 & 2 & 1 & 432 & 16 & 64 & \texttt{---++???-}\\
3 & 2 & 3 & 2 & 1 & 352 & 32 & 128 & \texttt{----+???-}\\
4 & 3 & 4 & 2 & 1 & 192 & 64 & 256 & \texttt{-----???-}\\
5 & 4 & 5 & 1 & 1 & 256 & 256 & 0 & \texttt{------+--}\\
6 & 5 & 6 & 1 & 1 & 0 & 512 & 0 & \texttt{---------}\\
7 & 4 & 5 & 1 & 1 & 0 & 512 & 0 & \texttt{---------}\\
8 & 1 & 2 & 1 & 1 & 0 & 512 & 0 & \texttt{---------}\\
\multicolumn{9}{@{}l@{}}{\rule{0pt}{2.6ex}\textbf{T9.020}\enspace \texttt{012345651}\enspace $n{=}9$, $|H|{=}3$, len 7, br 2, $|\mathrm{Aut}|{=}1$, orb 9/9, hist-orb 3/3; M 13824, res 3488, rev 4912: ok}\\*
0 & -- & 0 & 3 & 1 & 246 & 2 & 264 & \texttt{-????????}\\*
1 & 0 & 1 & 2 & 1 & 472 & 8 & 32 & \texttt{--+++???-}\\
2 & 1 & 2 & 2 & 1 & 432 & 16 & 64 & \texttt{---++???-}\\
3 & 2 & 3 & 2 & 1 & 352 & 32 & 128 & \texttt{----+???-}\\
4 & 3 & 4 & 2 & 1 & 192 & 64 & 256 & \texttt{-----???-}\\
5 & 4 & 5 & 1 & 1 & 256 & 256 & 0 & \texttt{------+--}\\
6 & 5 & 6 & 1 & 1 & 0 & 512 & 0 & \texttt{---------}\\
7 & 4 & 5 & 1 & 1 & 0 & 512 & 0 & \texttt{---------}\\
8 & 0 & 1 & 1 & 1 & 0 & 512 & 0 & \texttt{---------}\\
\multicolumn{9}{@{}l@{}}{\rule{0pt}{2.6ex}\textbf{T9.021}\enspace \texttt{012345645}\enspace $n{=}9$, $|H|{=}2$, len 7, br 1, $|\mathrm{Aut}|{=}1$, orb 9/9, hist-orb 2/2; M 15360, res 1500, rev 6436: ok}\\*
0 & -- & 0 & 2 & 1 & 490 & 2 & 20 & \texttt{-+++?????}\\*
1 & 0 & 1 & 2 & 1 & 468 & 4 & 40 & \texttt{--++?????}\\
2 & 1 & 2 & 2 & 1 & 424 & 8 & 80 & \texttt{---+?????}\\
3 & 2 & 3 & 2 & 1 & 336 & 16 & 160 & \texttt{----?????}\\
4 & 3 & 4 & 1 & 1 & 384 & 128 & 0 & \texttt{-----++--}\\
5 & 4 & 5 & 1 & 1 & 256 & 256 & 0 & \texttt{------+--}\\
6 & 5 & 6 & 1 & 1 & 0 & 512 & 0 & \texttt{---------}\\
7 & 3 & 4 & 1 & 1 & 256 & 256 & 0 & \texttt{--------+}\\
8 & 7 & 5 & 1 & 1 & 0 & 512 & 0 & \texttt{---------}\\
\multicolumn{9}{@{}l@{}}{\rule{0pt}{2.6ex}\textbf{T9.022}\enspace \texttt{012345644}\enspace $n{=}9$, $|H|{=}3$, len 7, br 1, $|\mathrm{Aut}|{=}2$, orb 8/7, hist-orb 2/1; M 14848, res 3600, rev 5644: ok}\\*
0 & -- & 0 & 3 & 1 & 462 & 2 & 48 & \texttt{-+++?????}\\*
1 & 0 & 1 & 3 & 1 & 412 & 4 & 96 & \texttt{--++?????}\\
2 & 1 & 2 & 3 & 1 & 312 & 8 & 192 & \texttt{---+?????}\\
3 & 2 & 3 & 3 & 1 & 112 & 16 & 384 & \texttt{----?????}\\
4 & 3 & 4 & 1 & 1 & 384 & 128 & 0 & \texttt{-----++--}\\
5 & 4 & 5 & 1 & 1 & 256 & 256 & 0 & \texttt{------+--}\\
6 & 5 & 6 & 1 & 1 & 0 & 512 & 0 & \texttt{---------}\\
7 & 3 & 4 & 1 & 2 & 0 & 512 & 0 & \texttt{---------}\\
8 & 3 & 4 & 1 & 2 & 0 & 512 & 0 & \texttt{---------}\\
\multicolumn{9}{@{}l@{}}{\rule{0pt}{2.6ex}\textbf{T9.023}\enspace \texttt{012345643}\enspace $n{=}9$, $|H|{=}3$, len 7, br 2, $|\mathrm{Aut}|{=}1$, orb 9/9, hist-orb 3/3; M 14336, res 4160, rev 5000: ok}\\*
0 & -- & 0 & 3 & 1 & 430 & 2 & 80 & \texttt{-++??????}\\*
1 & 0 & 1 & 3 & 1 & 348 & 4 & 160 & \texttt{--+??????}\\
2 & 1 & 2 & 3 & 1 & 184 & 8 & 320 & \texttt{---??????}\\
3 & 2 & 3 & 2 & 1 & 224 & 32 & 256 & \texttt{----????-}\\
4 & 3 & 4 & 1 & 1 & 384 & 128 & 0 & \texttt{-----++--}\\
5 & 4 & 5 & 1 & 1 & 256 & 256 & 0 & \texttt{------+--}\\
6 & 5 & 6 & 1 & 1 & 0 & 512 & 0 & \texttt{---------}\\
7 & 3 & 4 & 1 & 1 & 0 & 512 & 0 & \texttt{---------}\\
8 & 2 & 3 & 1 & 1 & 0 & 512 & 0 & \texttt{---------}\\
\multicolumn{9}{@{}l@{}}{\rule{0pt}{2.6ex}\textbf{T9.024}\enspace \texttt{012345642}\enspace $n{=}9$, $|H|{=}3$, len 7, br 2, $|\mathrm{Aut}|{=}1$, orb 9/9, hist-orb 3/3; M 13824, res 4368, rev 4592: ok}\\*
0 & -- & 0 & 3 & 1 & 366 & 2 & 144 & \texttt{-+???????}\\*
1 & 0 & 1 & 3 & 1 & 220 & 4 & 288 & \texttt{--???????}\\
2 & 1 & 2 & 2 & 1 & 368 & 16 & 128 & \texttt{---+????-}\\
3 & 2 & 3 & 2 & 1 & 224 & 32 & 256 & \texttt{----????-}\\
4 & 3 & 4 & 1 & 1 & 384 & 128 & 0 & \texttt{-----++--}\\
5 & 4 & 5 & 1 & 1 & 256 & 256 & 0 & \texttt{------+--}\\
6 & 5 & 6 & 1 & 1 & 0 & 512 & 0 & \texttt{---------}\\
7 & 3 & 4 & 1 & 1 & 0 & 512 & 0 & \texttt{---------}\\
8 & 1 & 2 & 1 & 1 & 0 & 512 & 0 & \texttt{---------}\\
\multicolumn{9}{@{}l@{}}{\rule{0pt}{2.6ex}\textbf{T9.025}\enspace \texttt{012345641}\enspace $n{=}9$, $|H|{=}3$, len 7, br 2, $|\mathrm{Aut}|{=}1$, orb 9/9, hist-orb 3/3; M 13312, res 3872, rev 4496: ok}\\*
0 & -- & 0 & 3 & 1 & 238 & 2 & 272 & \texttt{-????????}\\*
1 & 0 & 1 & 2 & 1 & 440 & 8 & 64 & \texttt{--++????-}\\
2 & 1 & 2 & 2 & 1 & 368 & 16 & 128 & \texttt{---+????-}\\
3 & 2 & 3 & 2 & 1 & 224 & 32 & 256 & \texttt{----????-}\\
4 & 3 & 4 & 1 & 1 & 384 & 128 & 0 & \texttt{-----++--}\\
5 & 4 & 5 & 1 & 1 & 256 & 256 & 0 & \texttt{------+--}\\
6 & 5 & 6 & 1 & 1 & 0 & 512 & 0 & \texttt{---------}\\
7 & 3 & 4 & 1 & 1 & 0 & 512 & 0 & \texttt{---------}\\
8 & 0 & 1 & 1 & 1 & 0 & 512 & 0 & \texttt{---------}\\
\multicolumn{9}{@{}l@{}}{\rule{0pt}{2.6ex}\textbf{T9.026}\enspace \texttt{012345634}\enspace $n{=}9$, $|H|{=}2$, len 7, br 1, $|\mathrm{Aut}|{=}1$, orb 9/9, hist-orb 2/2; M 14336, res 1512, rev 5832: ok}\\*
0 & -- & 0 & 2 & 1 & 474 & 2 & 36 & \texttt{-++??????}\\*
1 & 0 & 1 & 2 & 1 & 436 & 4 & 72 & \texttt{--+??????}\\
2 & 1 & 2 & 2 & 1 & 360 & 8 & 144 & \texttt{---??????}\\
3 & 2 & 3 & 1 & 1 & 448 & 64 & 0 & \texttt{----+++--}\\
4 & 3 & 4 & 1 & 1 & 384 & 128 & 0 & \texttt{-----++--}\\
5 & 4 & 5 & 1 & 1 & 256 & 256 & 0 & \texttt{------+--}\\
6 & 5 & 6 & 1 & 1 & 0 & 512 & 0 & \texttt{---------}\\
7 & 2 & 3 & 1 & 1 & 256 & 256 & 0 & \texttt{--------+}\\
8 & 7 & 4 & 1 & 1 & 0 & 512 & 0 & \texttt{---------}\\
\multicolumn{9}{@{}l@{}}{\rule{0pt}{2.6ex}\textbf{T9.027}\enspace \texttt{012345633}\enspace $n{=}9$, $|H|{=}3$, len 7, br 1, $|\mathrm{Aut}|{=}2$, orb 8/7, hist-orb 2/1; M 13824, res 4032, rev 4748: ok}\\*
0 & -- & 0 & 3 & 1 & 414 & 2 & 96 & \texttt{-++??????}\\*
1 & 0 & 1 & 3 & 1 & 316 & 4 & 192 & \texttt{--+??????}\\
2 & 1 & 2 & 3 & 1 & 120 & 8 & 384 & \texttt{---??????}\\
3 & 2 & 3 & 1 & 1 & 448 & 64 & 0 & \texttt{----+++--}\\
4 & 3 & 4 & 1 & 1 & 384 & 128 & 0 & \texttt{-----++--}\\
5 & 4 & 5 & 1 & 1 & 256 & 256 & 0 & \texttt{------+--}\\
6 & 5 & 6 & 1 & 1 & 0 & 512 & 0 & \texttt{---------}\\
7 & 2 & 3 & 1 & 2 & 0 & 512 & 0 & \texttt{---------}\\
8 & 2 & 3 & 1 & 2 & 0 & 512 & 0 & \texttt{---------}\\
\multicolumn{9}{@{}l@{}}{\rule{0pt}{2.6ex}\textbf{T9.028}\enspace \texttt{012345632}\enspace $n{=}9$, $|H|{=}3$, len 7, br 2, $|\mathrm{Aut}|{=}1$, orb 9/9, hist-orb 3/3; M 13312, res 4448, rev 4216: ok}\\*
0 & -- & 0 & 3 & 1 & 350 & 2 & 160 & \texttt{-+???????}\\*
1 & 0 & 1 & 3 & 1 & 188 & 4 & 320 & \texttt{--???????}\\
2 & 1 & 2 & 2 & 1 & 240 & 16 & 256 & \texttt{---?????-}\\
3 & 2 & 3 & 1 & 1 & 448 & 64 & 0 & \texttt{----+++--}\\
4 & 3 & 4 & 1 & 1 & 384 & 128 & 0 & \texttt{-----++--}\\
5 & 4 & 5 & 1 & 1 & 256 & 256 & 0 & \texttt{------+--}\\
6 & 5 & 6 & 1 & 1 & 0 & 512 & 0 & \texttt{---------}\\
7 & 2 & 3 & 1 & 1 & 0 & 512 & 0 & \texttt{---------}\\
8 & 1 & 2 & 1 & 1 & 0 & 512 & 0 & \texttt{---------}\\
\multicolumn{9}{@{}l@{}}{\rule{0pt}{2.6ex}\textbf{T9.029}\enspace \texttt{012345631}\enspace $n{=}9$, $|H|{=}3$, len 7, br 2, $|\mathrm{Aut}|{=}1$, orb 9/9, hist-orb 3/3; M 12800, res 4096, rev 4064: ok}\\*
0 & -- & 0 & 3 & 1 & 222 & 2 & 288 & \texttt{-????????}\\*
1 & 0 & 1 & 2 & 1 & 376 & 8 & 128 & \texttt{--+?????-}\\
2 & 1 & 2 & 2 & 1 & 240 & 16 & 256 & \texttt{---?????-}\\
3 & 2 & 3 & 1 & 1 & 448 & 64 & 0 & \texttt{----+++--}\\
4 & 3 & 4 & 1 & 1 & 384 & 128 & 0 & \texttt{-----++--}\\
5 & 4 & 5 & 1 & 1 & 256 & 256 & 0 & \texttt{------+--}\\
6 & 5 & 6 & 1 & 1 & 0 & 512 & 0 & \texttt{---------}\\
7 & 2 & 3 & 1 & 1 & 0 & 512 & 0 & \texttt{---------}\\
8 & 0 & 1 & 1 & 1 & 0 & 512 & 0 & \texttt{---------}\\
\multicolumn{9}{@{}l@{}}{\rule{0pt}{2.6ex}\textbf{T9.030}\enspace \texttt{012345623}\enspace $n{=}9$, $|H|{=}2$, len 7, br 1, $|\mathrm{Aut}|{=}1$, orb 9/9, hist-orb 2/2; M 13312, res 1428, rev 5212: ok}\\*
0 & -- & 0 & 2 & 1 & 442 & 2 & 68 & \texttt{-+???????}\\*
1 & 0 & 1 & 2 & 1 & 372 & 4 & 136 & \texttt{--???????}\\
2 & 1 & 2 & 1 & 1 & 480 & 32 & 0 & \texttt{---++++--}\\
3 & 2 & 3 & 1 & 1 & 448 & 64 & 0 & \texttt{----+++--}\\
4 & 3 & 4 & 1 & 1 & 384 & 128 & 0 & \texttt{-----++--}\\
5 & 4 & 5 & 1 & 1 & 256 & 256 & 0 & \texttt{------+--}\\
6 & 5 & 6 & 1 & 1 & 0 & 512 & 0 & \texttt{---------}\\
7 & 1 & 2 & 1 & 1 & 256 & 256 & 0 & \texttt{--------+}\\
8 & 7 & 3 & 1 & 1 & 0 & 512 & 0 & \texttt{---------}\\
\multicolumn{9}{@{}l@{}}{\rule{0pt}{2.6ex}\textbf{T9.031}\enspace \texttt{012345622}\enspace $n{=}9$, $|H|{=}3$, len 7, br 1, $|\mathrm{Aut}|{=}2$, orb 8/7, hist-orb 2/1; M 12800, res 4032, rev 4028: ok}\\*
0 & -- & 0 & 3 & 1 & 318 & 2 & 192 & \texttt{-+???????}\\*
1 & 0 & 1 & 3 & 1 & 124 & 4 & 384 & \texttt{--???????}\\
2 & 1 & 2 & 1 & 1 & 480 & 32 & 0 & \texttt{---++++--}\\
3 & 2 & 3 & 1 & 1 & 448 & 64 & 0 & \texttt{----+++--}\\
4 & 3 & 4 & 1 & 1 & 384 & 128 & 0 & \texttt{-----++--}\\
5 & 4 & 5 & 1 & 1 & 256 & 256 & 0 & \texttt{------+--}\\
6 & 5 & 6 & 1 & 1 & 0 & 512 & 0 & \texttt{---------}\\
7 & 1 & 2 & 1 & 2 & 0 & 512 & 0 & \texttt{---------}\\
8 & 1 & 2 & 1 & 2 & 0 & 512 & 0 & \texttt{---------}\\
\multicolumn{9}{@{}l@{}}{\rule{0pt}{2.6ex}\textbf{T9.032}\enspace \texttt{012345621}\enspace $n{=}9$, $|H|{=}3$, len 7, br 2, $|\mathrm{Aut}|{=}1$, orb 9/9, hist-orb 3/3; M 12288, res 3968, rev 3752: ok}\\*
0 & -- & 0 & 3 & 1 & 190 & 2 & 320 & \texttt{-????????}\\*
1 & 0 & 1 & 2 & 1 & 248 & 8 & 256 & \texttt{--??????-}\\
2 & 1 & 2 & 1 & 1 & 480 & 32 & 0 & \texttt{---++++--}\\
3 & 2 & 3 & 1 & 1 & 448 & 64 & 0 & \texttt{----+++--}\\
4 & 3 & 4 & 1 & 1 & 384 & 128 & 0 & \texttt{-----++--}\\
5 & 4 & 5 & 1 & 1 & 256 & 256 & 0 & \texttt{------+--}\\
6 & 5 & 6 & 1 & 1 & 0 & 512 & 0 & \texttt{---------}\\
7 & 1 & 2 & 1 & 1 & 0 & 512 & 0 & \texttt{---------}\\
8 & 0 & 1 & 1 & 1 & 0 & 512 & 0 & \texttt{---------}\\
\multicolumn{9}{@{}l@{}}{\rule{0pt}{2.6ex}\textbf{T9.033}\enspace \texttt{012345612}\enspace $n{=}9$, $|H|{=}2$, len 7, br 1, $|\mathrm{Aut}|{=}1$, orb 9/9, hist-orb 2/2; M 12288, res 1056, rev 4696: ok}\\*
0 & -- & 0 & 2 & 1 & 378 & 2 & 132 & \texttt{-????????}\\*
1 & 0 & 1 & 1 & 1 & 496 & 16 & 0 & \texttt{--+++++--}\\
2 & 1 & 2 & 1 & 1 & 480 & 32 & 0 & \texttt{---++++--}\\
3 & 2 & 3 & 1 & 1 & 448 & 64 & 0 & \texttt{----+++--}\\
4 & 3 & 4 & 1 & 1 & 384 & 128 & 0 & \texttt{-----++--}\\
5 & 4 & 5 & 1 & 1 & 256 & 256 & 0 & \texttt{------+--}\\
6 & 5 & 6 & 1 & 1 & 0 & 512 & 0 & \texttt{---------}\\
7 & 0 & 1 & 1 & 1 & 256 & 256 & 0 & \texttt{--------+}\\
8 & 7 & 2 & 1 & 1 & 0 & 512 & 0 & \texttt{---------}\\
\multicolumn{9}{@{}l@{}}{\rule{0pt}{2.6ex}\textbf{T9.034}\enspace \texttt{012345611}\enspace $n{=}9$, $|H|{=}3$, len 7, br 1, $|\mathrm{Aut}|{=}2$, orb 8/7, hist-orb 2/1; M 11776, res 3072, rev 3708: ok}\\*
0 & -- & 0 & 3 & 1 & 126 & 2 & 384 & \texttt{-????????}\\*
1 & 0 & 1 & 1 & 1 & 496 & 16 & 0 & \texttt{--+++++--}\\
2 & 1 & 2 & 1 & 1 & 480 & 32 & 0 & \texttt{---++++--}\\
3 & 2 & 3 & 1 & 1 & 448 & 64 & 0 & \texttt{----+++--}\\
4 & 3 & 4 & 1 & 1 & 384 & 128 & 0 & \texttt{-----++--}\\
5 & 4 & 5 & 1 & 1 & 256 & 256 & 0 & \texttt{------+--}\\
6 & 5 & 6 & 1 & 1 & 0 & 512 & 0 & \texttt{---------}\\
7 & 0 & 1 & 1 & 2 & 0 & 512 & 0 & \texttt{---------}\\
8 & 0 & 1 & 1 & 2 & 0 & 512 & 0 & \texttt{---------}\\
\multicolumn{9}{@{}l@{}}{\rule{0pt}{2.6ex}\textbf{T9.035}\enspace \texttt{012345555}\enspace $n{=}9$, $|H|{=}4$, len 6, br 1, $|\mathrm{Aut}|{=}24$, orb 6/5, hist-orb 1/0; M 15360, res 3472, rev 6664: ok}\\*
0 & -- & 0 & 4 & 1 & 482 & 2 & 28 & \texttt{-++++????}\\*
1 & 0 & 1 & 4 & 1 & 452 & 4 & 56 & \texttt{--+++????}\\
2 & 1 & 2 & 4 & 1 & 392 & 8 & 112 & \texttt{---++????}\\
3 & 2 & 3 & 4 & 1 & 272 & 16 & 224 & \texttt{----+????}\\
4 & 3 & 4 & 4 & 1 & 32 & 32 & 448 & \texttt{-----????}\\
5 & 4 & 5 & 1 & 4 & 0 & 512 & 0 & \texttt{---------}\\
6 & 4 & 5 & 1 & 4 & 0 & 512 & 0 & \texttt{---------}\\
7 & 4 & 5 & 1 & 4 & 0 & 512 & 0 & \texttt{---------}\\
8 & 4 & 5 & 1 & 4 & 0 & 512 & 0 & \texttt{---------}\\
\multicolumn{9}{@{}l@{}}{\rule{0pt}{2.6ex}\textbf{T9.036}\enspace \texttt{012345554}\enspace $n{=}9$, $|H|{=}4$, len 6, br 2, $|\mathrm{Aut}|{=}6$, orb 7/6, hist-orb 2/1; M 14848, res 4092, rev 5836: ok}\\*
0 & -- & 0 & 4 & 1 & 466 & 2 & 44 & \texttt{-+++?????}\\*
1 & 0 & 1 & 4 & 1 & 420 & 4 & 88 & \texttt{--++?????}\\
2 & 1 & 2 & 4 & 1 & 328 & 8 & 176 & \texttt{---+?????}\\
3 & 2 & 3 & 4 & 1 & 144 & 16 & 352 & \texttt{----?????}\\
4 & 3 & 4 & 3 & 1 & 64 & 64 & 384 & \texttt{-----???-}\\
5 & 4 & 5 & 1 & 3 & 0 & 512 & 0 & \texttt{---------}\\
6 & 4 & 5 & 1 & 3 & 0 & 512 & 0 & \texttt{---------}\\
7 & 4 & 5 & 1 & 3 & 0 & 512 & 0 & \texttt{---------}\\
8 & 3 & 4 & 1 & 1 & 0 & 512 & 0 & \texttt{---------}\\
\multicolumn{9}{@{}l@{}}{\rule{0pt}{2.6ex}\textbf{T9.037}\enspace \texttt{012345553}\enspace $n{=}9$, $|H|{=}4$, len 6, br 2, $|\mathrm{Aut}|{=}6$, orb 7/6, hist-orb 2/1; M 14336, res 4584, rev 5236: ok}\\*
0 & -- & 0 & 4 & 1 & 434 & 2 & 76 & \texttt{-++??????}\\*
1 & 0 & 1 & 4 & 1 & 356 & 4 & 152 & \texttt{--+??????}\\
2 & 1 & 2 & 4 & 1 & 200 & 8 & 304 & \texttt{---??????}\\
3 & 2 & 3 & 3 & 1 & 288 & 32 & 192 & \texttt{----+???-}\\
4 & 3 & 4 & 3 & 1 & 64 & 64 & 384 & \texttt{-----???-}\\
5 & 4 & 5 & 1 & 3 & 0 & 512 & 0 & \texttt{---------}\\
6 & 4 & 5 & 1 & 3 & 0 & 512 & 0 & \texttt{---------}\\
7 & 4 & 5 & 1 & 3 & 0 & 512 & 0 & \texttt{---------}\\
8 & 2 & 3 & 1 & 1 & 0 & 512 & 0 & \texttt{---------}\\
\multicolumn{9}{@{}l@{}}{\rule{0pt}{2.6ex}\textbf{T9.038}\enspace \texttt{012345552}\enspace $n{=}9$, $|H|{=}4$, len 6, br 2, $|\mathrm{Aut}|{=}6$, orb 7/6, hist-orb 2/1; M 13824, res 4740, rev 4868: ok}\\*
0 & -- & 0 & 4 & 1 & 370 & 2 & 140 & \texttt{-+???????}\\*
1 & 0 & 1 & 4 & 1 & 228 & 4 & 280 & \texttt{--???????}\\
2 & 1 & 2 & 3 & 1 & 400 & 16 & 96 & \texttt{---++???-}\\
3 & 2 & 3 & 3 & 1 & 288 & 32 & 192 & \texttt{----+???-}\\
4 & 3 & 4 & 3 & 1 & 64 & 64 & 384 & \texttt{-----???-}\\
5 & 4 & 5 & 1 & 3 & 0 & 512 & 0 & \texttt{---------}\\
6 & 4 & 5 & 1 & 3 & 0 & 512 & 0 & \texttt{---------}\\
7 & 4 & 5 & 1 & 3 & 0 & 512 & 0 & \texttt{---------}\\
8 & 1 & 2 & 1 & 1 & 0 & 512 & 0 & \texttt{---------}\\
\multicolumn{9}{@{}l@{}}{\rule{0pt}{2.6ex}\textbf{T9.039}\enspace \texttt{012345551}\enspace $n{=}9$, $|H|{=}4$, len 6, br 2, $|\mathrm{Aut}|{=}6$, orb 7/6, hist-orb 2/1; M 13312, res 4208, rev 4828: ok}\\*
0 & -- & 0 & 4 & 1 & 242 & 2 & 268 & \texttt{-????????}\\*
1 & 0 & 1 & 3 & 1 & 456 & 8 & 48 & \texttt{--+++???-}\\
2 & 1 & 2 & 3 & 1 & 400 & 16 & 96 & \texttt{---++???-}\\
3 & 2 & 3 & 3 & 1 & 288 & 32 & 192 & \texttt{----+???-}\\
4 & 3 & 4 & 3 & 1 & 64 & 64 & 384 & \texttt{-----???-}\\
5 & 4 & 5 & 1 & 3 & 0 & 512 & 0 & \texttt{---------}\\
6 & 4 & 5 & 1 & 3 & 0 & 512 & 0 & \texttt{---------}\\
7 & 4 & 5 & 1 & 3 & 0 & 512 & 0 & \texttt{---------}\\
8 & 0 & 1 & 1 & 1 & 0 & 512 & 0 & \texttt{---------}\\
\multicolumn{9}{@{}l@{}}{\rule{0pt}{2.6ex}\textbf{T9.040}\enspace \texttt{012345545}\enspace $n{=}9$, $|H|{=}3$, len 6, br 2, $|\mathrm{Aut}|{=}2$, orb 8/7, hist-orb 2/1; M 14848, res 2672, rev 6228: ok}\\*
0 & -- & 0 & 3 & 1 & 478 & 2 & 32 & \texttt{-+++?????}\\*
1 & 0 & 1 & 3 & 1 & 444 & 4 & 64 & \texttt{--++?????}\\
2 & 1 & 2 & 3 & 1 & 376 & 8 & 128 & \texttt{---+?????}\\
3 & 2 & 3 & 3 & 1 & 240 & 16 & 256 & \texttt{----?????}\\
4 & 3 & 4 & 2 & 1 & 128 & 128 & 256 & \texttt{-----??--}\\
5 & 4 & 5 & 1 & 2 & 0 & 512 & 0 & \texttt{---------}\\
6 & 4 & 5 & 1 & 2 & 0 & 512 & 0 & \texttt{---------}\\
7 & 3 & 4 & 1 & 1 & 256 & 256 & 0 & \texttt{--------+}\\
8 & 7 & 5 & 1 & 1 & 0 & 512 & 0 & \texttt{---------}\\
\multicolumn{9}{@{}l@{}}{\rule{0pt}{2.6ex}\textbf{T9.041}\enspace \texttt{012345544}\enspace $n{=}9$, $|H|{=}4$, len 6, br 2, $|\mathrm{Aut}|{=}4$, orb 7/5, hist-orb 2/0; M 14336, res 4172, rev 5556: ok}\\*
0 & -- & 0 & 4 & 1 & 458 & 2 & 52 & \texttt{-+++?????}\\*
1 & 0 & 1 & 4 & 1 & 404 & 4 & 104 & \texttt{--++?????}\\
2 & 1 & 2 & 4 & 1 & 296 & 8 & 208 & \texttt{---+?????}\\
3 & 2 & 3 & 4 & 1 & 80 & 16 & 416 & \texttt{----?????}\\
4 & 3 & 4 & 2 & 1 & 128 & 128 & 256 & \texttt{-----??--}\\
5 & 4 & 5 & 1 & 2 & 0 & 512 & 0 & \texttt{---------}\\
6 & 4 & 5 & 1 & 2 & 0 & 512 & 0 & \texttt{---------}\\
7 & 3 & 4 & 1 & 2 & 0 & 512 & 0 & \texttt{---------}\\
8 & 3 & 4 & 1 & 2 & 0 & 512 & 0 & \texttt{---------}\\
\multicolumn{9}{@{}l@{}}{\rule{0pt}{2.6ex}\textbf{T9.042}\enspace \texttt{012345543}\enspace $n{=}9$, $|H|{=}4$, len 6, br 3, $|\mathrm{Aut}|{=}2$, orb 8/7, hist-orb 3/2; M 13824, res 4800, rev 4800: ok}\\*
0 & -- & 0 & 4 & 1 & 426 & 2 & 84 & \texttt{-++??????}\\*
1 & 0 & 1 & 4 & 1 & 340 & 4 & 168 & \texttt{--+??????}\\
2 & 1 & 2 & 4 & 1 & 168 & 8 & 336 & \texttt{---??????}\\
3 & 2 & 3 & 3 & 1 & 160 & 32 & 320 & \texttt{----????-}\\
4 & 3 & 4 & 2 & 1 & 128 & 128 & 256 & \texttt{-----??--}\\
5 & 4 & 5 & 1 & 2 & 0 & 512 & 0 & \texttt{---------}\\
6 & 4 & 5 & 1 & 2 & 0 & 512 & 0 & \texttt{---------}\\
7 & 3 & 4 & 1 & 1 & 0 & 512 & 0 & \texttt{---------}\\
8 & 2 & 3 & 1 & 1 & 0 & 512 & 0 & \texttt{---------}\\
\multicolumn{9}{@{}l@{}}{\rule{0pt}{2.6ex}\textbf{T9.043}\enspace \texttt{012345542}\enspace $n{=}9$, $|H|{=}4$, len 6, br 3, $|\mathrm{Aut}|{=}2$, orb 8/7, hist-orb 3/2; M 13312, res 5060, rev 4360: ok}\\*
0 & -- & 0 & 4 & 1 & 362 & 2 & 148 & \texttt{-+???????}\\*
1 & 0 & 1 & 4 & 1 & 212 & 4 & 296 & \texttt{--???????}\\
2 & 1 & 2 & 3 & 1 & 336 & 16 & 160 & \texttt{---+????-}\\
3 & 2 & 3 & 3 & 1 & 160 & 32 & 320 & \texttt{----????-}\\
4 & 3 & 4 & 2 & 1 & 128 & 128 & 256 & \texttt{-----??--}\\
5 & 4 & 5 & 1 & 2 & 0 & 512 & 0 & \texttt{---------}\\
6 & 4 & 5 & 1 & 2 & 0 & 512 & 0 & \texttt{---------}\\
7 & 3 & 4 & 1 & 1 & 0 & 512 & 0 & \texttt{---------}\\
8 & 1 & 2 & 1 & 1 & 0 & 512 & 0 & \texttt{---------}\\
\multicolumn{9}{@{}l@{}}{\rule{0pt}{2.6ex}\textbf{T9.044}\enspace \texttt{012345541}\enspace $n{=}9$, $|H|{=}4$, len 6, br 3, $|\mathrm{Aut}|{=}2$, orb 8/7, hist-orb 3/2; M 12800, res 4600, rev 4296: ok}\\*
0 & -- & 0 & 4 & 1 & 234 & 2 & 276 & \texttt{-????????}\\*
1 & 0 & 1 & 3 & 1 & 424 & 8 & 80 & \texttt{--++????-}\\
2 & 1 & 2 & 3 & 1 & 336 & 16 & 160 & \texttt{---+????-}\\
3 & 2 & 3 & 3 & 1 & 160 & 32 & 320 & \texttt{----????-}\\
4 & 3 & 4 & 2 & 1 & 128 & 128 & 256 & \texttt{-----??--}\\
5 & 4 & 5 & 1 & 2 & 0 & 512 & 0 & \texttt{---------}\\
6 & 4 & 5 & 1 & 2 & 0 & 512 & 0 & \texttt{---------}\\
7 & 3 & 4 & 1 & 1 & 0 & 512 & 0 & \texttt{---------}\\
8 & 0 & 1 & 1 & 1 & 0 & 512 & 0 & \texttt{---------}\\
\multicolumn{9}{@{}l@{}}{\rule{0pt}{2.6ex}\textbf{T9.045}\enspace \texttt{012345534}\enspace $n{=}9$, $|H|{=}3$, len 6, br 2, $|\mathrm{Aut}|{=}2$, orb 8/7, hist-orb 2/1; M 13824, res 2560, rev 5696: ok}\\*
0 & -- & 0 & 3 & 1 & 462 & 2 & 48 & \texttt{-++??????}\\*
1 & 0 & 1 & 3 & 1 & 412 & 4 & 96 & \texttt{--+??????}\\
2 & 1 & 2 & 3 & 1 & 312 & 8 & 192 & \texttt{---??????}\\
3 & 2 & 3 & 2 & 1 & 320 & 64 & 128 & \texttt{----+??--}\\
4 & 3 & 4 & 2 & 1 & 128 & 128 & 256 & \texttt{-----??--}\\
5 & 4 & 5 & 1 & 2 & 0 & 512 & 0 & \texttt{---------}\\
6 & 4 & 5 & 1 & 2 & 0 & 512 & 0 & \texttt{---------}\\
7 & 2 & 3 & 1 & 1 & 256 & 256 & 0 & \texttt{--------+}\\
8 & 7 & 4 & 1 & 1 & 0 & 512 & 0 & \texttt{---------}\\
\multicolumn{9}{@{}l@{}}{\rule{0pt}{2.6ex}\textbf{T9.046}\enspace \texttt{012345533}\enspace $n{=}9$, $|H|{=}4$, len 6, br 2, $|\mathrm{Aut}|{=}4$, orb 7/5, hist-orb 2/0; M 13312, res 4744, rev 4556: ok}\\*
0 & -- & 0 & 4 & 1 & 410 & 2 & 100 & \texttt{-++??????}\\*
1 & 0 & 1 & 4 & 1 & 308 & 4 & 200 & \texttt{--+??????}\\
2 & 1 & 2 & 4 & 1 & 104 & 8 & 400 & \texttt{---??????}\\
3 & 2 & 3 & 2 & 1 & 320 & 64 & 128 & \texttt{----+??--}\\
4 & 3 & 4 & 2 & 1 & 128 & 128 & 256 & \texttt{-----??--}\\
5 & 4 & 5 & 1 & 2 & 0 & 512 & 0 & \texttt{---------}\\
6 & 4 & 5 & 1 & 2 & 0 & 512 & 0 & \texttt{---------}\\
7 & 2 & 3 & 1 & 2 & 0 & 512 & 0 & \texttt{---------}\\
8 & 2 & 3 & 1 & 2 & 0 & 512 & 0 & \texttt{---------}\\
\multicolumn{9}{@{}l@{}}{\rule{0pt}{2.6ex}\textbf{T9.047}\enspace \texttt{012345532}\enspace $n{=}9$, $|H|{=}4$, len 6, br 3, $|\mathrm{Aut}|{=}2$, orb 8/7, hist-orb 3/2; M 12800, res 5212, rev 3960: ok}\\*
0 & -- & 0 & 4 & 1 & 346 & 2 & 164 & \texttt{-+???????}\\*
1 & 0 & 1 & 4 & 1 & 180 & 4 & 328 & \texttt{--???????}\\
2 & 1 & 2 & 3 & 1 & 208 & 16 & 288 & \texttt{---?????-}\\
3 & 2 & 3 & 2 & 1 & 320 & 64 & 128 & \texttt{----+??--}\\
4 & 3 & 4 & 2 & 1 & 128 & 128 & 256 & \texttt{-----??--}\\
5 & 4 & 5 & 1 & 2 & 0 & 512 & 0 & \texttt{---------}\\
6 & 4 & 5 & 1 & 2 & 0 & 512 & 0 & \texttt{---------}\\
7 & 2 & 3 & 1 & 1 & 0 & 512 & 0 & \texttt{---------}\\
8 & 1 & 2 & 1 & 1 & 0 & 512 & 0 & \texttt{---------}\\
\multicolumn{9}{@{}l@{}}{\rule{0pt}{2.6ex}\textbf{T9.048}\enspace \texttt{012345531}\enspace $n{=}9$, $|H|{=}4$, len 6, br 3, $|\mathrm{Aut}|{=}2$, orb 8/7, hist-orb 3/2; M 12288, res 4896, rev 3824: ok}\\*
0 & -- & 0 & 4 & 1 & 218 & 2 & 292 & \texttt{-????????}\\*
1 & 0 & 1 & 3 & 1 & 360 & 8 & 144 & \texttt{--+?????-}\\
2 & 1 & 2 & 3 & 1 & 208 & 16 & 288 & \texttt{---?????-}\\
3 & 2 & 3 & 2 & 1 & 320 & 64 & 128 & \texttt{----+??--}\\
4 & 3 & 4 & 2 & 1 & 128 & 128 & 256 & \texttt{-----??--}\\
5 & 4 & 5 & 1 & 2 & 0 & 512 & 0 & \texttt{---------}\\
6 & 4 & 5 & 1 & 2 & 0 & 512 & 0 & \texttt{---------}\\
7 & 2 & 3 & 1 & 1 & 0 & 512 & 0 & \texttt{---------}\\
8 & 0 & 1 & 1 & 1 & 0 & 512 & 0 & \texttt{---------}\\
\multicolumn{9}{@{}l@{}}{\rule{0pt}{2.6ex}\textbf{T9.049}\enspace \texttt{012345523}\enspace $n{=}9$, $|H|{=}3$, len 6, br 2, $|\mathrm{Aut}|{=}2$, orb 8/7, hist-orb 2/1; M 12800, res 2432, rev 5124: ok}\\*
0 & -- & 0 & 3 & 1 & 430 & 2 & 80 & \texttt{-+???????}\\*
1 & 0 & 1 & 3 & 1 & 348 & 4 & 160 & \texttt{--???????}\\
2 & 1 & 2 & 2 & 1 & 416 & 32 & 64 & \texttt{---++??--}\\
3 & 2 & 3 & 2 & 1 & 320 & 64 & 128 & \texttt{----+??--}\\
4 & 3 & 4 & 2 & 1 & 128 & 128 & 256 & \texttt{-----??--}\\
5 & 4 & 5 & 1 & 2 & 0 & 512 & 0 & \texttt{---------}\\
6 & 4 & 5 & 1 & 2 & 0 & 512 & 0 & \texttt{---------}\\
7 & 1 & 2 & 1 & 1 & 256 & 256 & 0 & \texttt{--------+}\\
8 & 7 & 3 & 1 & 1 & 0 & 512 & 0 & \texttt{---------}\\
\multicolumn{9}{@{}l@{}}{\rule{0pt}{2.6ex}\textbf{T9.050}\enspace \texttt{012345522}\enspace $n{=}9$, $|H|{=}4$, len 6, br 2, $|\mathrm{Aut}|{=}4$, orb 7/5, hist-orb 2/0; M 12288, res 4868, rev 3820: ok}\\*
0 & -- & 0 & 4 & 1 & 314 & 2 & 196 & \texttt{-+???????}\\*
1 & 0 & 1 & 4 & 1 & 116 & 4 & 392 & \texttt{--???????}\\
2 & 1 & 2 & 2 & 1 & 416 & 32 & 64 & \texttt{---++??--}\\
3 & 2 & 3 & 2 & 1 & 320 & 64 & 128 & \texttt{----+??--}\\
4 & 3 & 4 & 2 & 1 & 128 & 128 & 256 & \texttt{-----??--}\\
5 & 4 & 5 & 1 & 2 & 0 & 512 & 0 & \texttt{---------}\\
6 & 4 & 5 & 1 & 2 & 0 & 512 & 0 & \texttt{---------}\\
7 & 1 & 2 & 1 & 2 & 0 & 512 & 0 & \texttt{---------}\\
8 & 1 & 2 & 1 & 2 & 0 & 512 & 0 & \texttt{---------}\\
\multicolumn{9}{@{}l@{}}{\rule{0pt}{2.6ex}\textbf{T9.051}\enspace \texttt{012345521}\enspace $n{=}9$, $|H|{=}4$, len 6, br 3, $|\mathrm{Aut}|{=}2$, orb 8/7, hist-orb 3/2; M 11776, res 4840, rev 3528: ok}\\*
0 & -- & 0 & 4 & 1 & 186 & 2 & 324 & \texttt{-????????}\\*
1 & 0 & 1 & 3 & 1 & 232 & 8 & 272 & \texttt{--??????-}\\
2 & 1 & 2 & 2 & 1 & 416 & 32 & 64 & \texttt{---++??--}\\
3 & 2 & 3 & 2 & 1 & 320 & 64 & 128 & \texttt{----+??--}\\
4 & 3 & 4 & 2 & 1 & 128 & 128 & 256 & \texttt{-----??--}\\
5 & 4 & 5 & 1 & 2 & 0 & 512 & 0 & \texttt{---------}\\
6 & 4 & 5 & 1 & 2 & 0 & 512 & 0 & \texttt{---------}\\
7 & 1 & 2 & 1 & 1 & 0 & 512 & 0 & \texttt{---------}\\
8 & 0 & 1 & 1 & 1 & 0 & 512 & 0 & \texttt{---------}\\
\multicolumn{9}{@{}l@{}}{\rule{0pt}{2.6ex}\textbf{T9.052}\enspace \texttt{012345512}\enspace $n{=}9$, $|H|{=}3$, len 6, br 2, $|\mathrm{Aut}|{=}2$, orb 8/7, hist-orb 2/1; M 11776, res 2048, rev 4656: ok}\\*
0 & -- & 0 & 3 & 1 & 366 & 2 & 144 & \texttt{-????????}\\*
1 & 0 & 1 & 2 & 1 & 464 & 16 & 32 & \texttt{--+++??--}\\
2 & 1 & 2 & 2 & 1 & 416 & 32 & 64 & \texttt{---++??--}\\
3 & 2 & 3 & 2 & 1 & 320 & 64 & 128 & \texttt{----+??--}\\
4 & 3 & 4 & 2 & 1 & 128 & 128 & 256 & \texttt{-----??--}\\
5 & 4 & 5 & 1 & 2 & 0 & 512 & 0 & \texttt{---------}\\
6 & 4 & 5 & 1 & 2 & 0 & 512 & 0 & \texttt{---------}\\
7 & 0 & 1 & 1 & 1 & 256 & 256 & 0 & \texttt{--------+}\\
8 & 7 & 2 & 1 & 1 & 0 & 512 & 0 & \texttt{---------}\\
\multicolumn{9}{@{}l@{}}{\rule{0pt}{2.6ex}\textbf{T9.053}\enspace \texttt{012345511}\enspace $n{=}9$, $|H|{=}4$, len 6, br 2, $|\mathrm{Aut}|{=}4$, orb 7/5, hist-orb 2/0; M 11264, res 4000, rev 3564: ok}\\*
0 & -- & 0 & 4 & 1 & 122 & 2 & 388 & \texttt{-????????}\\*
1 & 0 & 1 & 2 & 1 & 464 & 16 & 32 & \texttt{--+++??--}\\
2 & 1 & 2 & 2 & 1 & 416 & 32 & 64 & \texttt{---++??--}\\
3 & 2 & 3 & 2 & 1 & 320 & 64 & 128 & \texttt{----+??--}\\
4 & 3 & 4 & 2 & 1 & 128 & 128 & 256 & \texttt{-----??--}\\
5 & 4 & 5 & 1 & 2 & 0 & 512 & 0 & \texttt{---------}\\
6 & 4 & 5 & 1 & 2 & 0 & 512 & 0 & \texttt{---------}\\
7 & 0 & 1 & 1 & 2 & 0 & 512 & 0 & \texttt{---------}\\
8 & 0 & 1 & 1 & 2 & 0 & 512 & 0 & \texttt{---------}\\
\multicolumn{9}{@{}l@{}}{\rule{0pt}{2.6ex}\textbf{T9.054}\enspace \texttt{012345454}\enspace $n{=}9$, $|H|{=}3$, len 6, br 1, $|\mathrm{Aut}|{=}2$, orb 7/5, hist-orb 2/1; M 14336, res 3300, rev 5900: ok}\\*
0 & -- & 0 & 3 & 1 & 466 & 2 & 44 & \texttt{-+++?????}\\*
1 & 0 & 1 & 3 & 1 & 420 & 4 & 88 & \texttt{--++?????}\\
2 & 1 & 2 & 3 & 1 & 328 & 8 & 176 & \texttt{---+?????}\\
3 & 2 & 3 & 3 & 1 & 144 & 16 & 352 & \texttt{----?????}\\
4 & 3 & 4 & 1 & 2 & 256 & 256 & 0 & \texttt{-----+---}\\
5 & 4 & 5 & 1 & 2 & 0 & 512 & 0 & \texttt{---------}\\
6 & 3 & 4 & 1 & 2 & 256 & 256 & 0 & \texttt{-------+-}\\
7 & 6 & 5 & 1 & 2 & 0 & 512 & 0 & \texttt{---------}\\
8 & 3 & 4 & 1 & 1 & 0 & 512 & 0 & \texttt{---------}\\
\multicolumn{9}{@{}l@{}}{\rule{0pt}{2.6ex}\textbf{T9.055}\enspace \texttt{012345453}\enspace $n{=}9$, $|H|{=}3$, len 6, br 2, $|\mathrm{Aut}|{=}2$, orb 7/5, hist-orb 2/1; M 13824, res 3792, rev 5264: ok}\\*
0 & -- & 0 & 3 & 1 & 434 & 2 & 76 & \texttt{-++??????}\\*
1 & 0 & 1 & 3 & 1 & 356 & 4 & 152 & \texttt{--+??????}\\
2 & 1 & 2 & 3 & 1 & 200 & 8 & 304 & \texttt{---??????}\\
3 & 2 & 3 & 2 & 1 & 288 & 32 & 192 & \texttt{----????-}\\
4 & 3 & 4 & 1 & 2 & 256 & 256 & 0 & \texttt{-----+---}\\
5 & 4 & 5 & 1 & 2 & 0 & 512 & 0 & \texttt{---------}\\
6 & 3 & 4 & 1 & 2 & 256 & 256 & 0 & \texttt{-------+-}\\
7 & 6 & 5 & 1 & 2 & 0 & 512 & 0 & \texttt{---------}\\
8 & 2 & 3 & 1 & 1 & 0 & 512 & 0 & \texttt{---------}\\
\multicolumn{9}{@{}l@{}}{\rule{0pt}{2.6ex}\textbf{T9.056}\enspace \texttt{012345452}\enspace $n{=}9$, $|H|{=}3$, len 6, br 2, $|\mathrm{Aut}|{=}2$, orb 7/5, hist-orb 2/1; M 13312, res 3948, rev 4872: ok}\\*
0 & -- & 0 & 3 & 1 & 370 & 2 & 140 & \texttt{-+???????}\\*
1 & 0 & 1 & 3 & 1 & 228 & 4 & 280 & \texttt{--???????}\\
2 & 1 & 2 & 2 & 1 & 400 & 16 & 96 & \texttt{---+????-}\\
3 & 2 & 3 & 2 & 1 & 288 & 32 & 192 & \texttt{----????-}\\
4 & 3 & 4 & 1 & 2 & 256 & 256 & 0 & \texttt{-----+---}\\
5 & 4 & 5 & 1 & 2 & 0 & 512 & 0 & \texttt{---------}\\
6 & 3 & 4 & 1 & 2 & 256 & 256 & 0 & \texttt{-------+-}\\
7 & 6 & 5 & 1 & 2 & 0 & 512 & 0 & \texttt{---------}\\
8 & 1 & 2 & 1 & 1 & 0 & 512 & 0 & \texttt{---------}\\
\multicolumn{9}{@{}l@{}}{\rule{0pt}{2.6ex}\textbf{T9.057}\enspace \texttt{012345451}\enspace $n{=}9$, $|H|{=}3$, len 6, br 2, $|\mathrm{Aut}|{=}2$, orb 7/5, hist-orb 2/1; M 12800, res 3416, rev 4808: ok}\\*
0 & -- & 0 & 3 & 1 & 242 & 2 & 268 & \texttt{-????????}\\*
1 & 0 & 1 & 2 & 1 & 456 & 8 & 48 & \texttt{--++????-}\\
2 & 1 & 2 & 2 & 1 & 400 & 16 & 96 & \texttt{---+????-}\\
3 & 2 & 3 & 2 & 1 & 288 & 32 & 192 & \texttt{----????-}\\
4 & 3 & 4 & 1 & 2 & 256 & 256 & 0 & \texttt{-----+---}\\
5 & 4 & 5 & 1 & 2 & 0 & 512 & 0 & \texttt{---------}\\
6 & 3 & 4 & 1 & 2 & 256 & 256 & 0 & \texttt{-------+-}\\
7 & 6 & 5 & 1 & 2 & 0 & 512 & 0 & \texttt{---------}\\
8 & 0 & 1 & 1 & 1 & 0 & 512 & 0 & \texttt{---------}\\
\multicolumn{9}{@{}l@{}}{\rule{0pt}{2.6ex}\textbf{T9.058}\enspace \texttt{012345444}\enspace $n{=}9$, $|H|{=}4$, len 6, br 1, $|\mathrm{Aut}|{=}6$, orb 7/6, hist-orb 2/1; M 13824, res 4200, rev 5588: ok}\\*
0 & -- & 0 & 4 & 1 & 454 & 2 & 56 & \texttt{-+++?????}\\*
1 & 0 & 1 & 4 & 1 & 396 & 4 & 112 & \texttt{--++?????}\\
2 & 1 & 2 & 4 & 1 & 280 & 8 & 224 & \texttt{---+?????}\\
3 & 2 & 3 & 4 & 1 & 48 & 16 & 448 & \texttt{----?????}\\
4 & 3 & 4 & 1 & 1 & 256 & 256 & 0 & \texttt{-----+---}\\
5 & 4 & 5 & 1 & 1 & 0 & 512 & 0 & \texttt{---------}\\
6 & 3 & 4 & 1 & 3 & 0 & 512 & 0 & \texttt{---------}\\
7 & 3 & 4 & 1 & 3 & 0 & 512 & 0 & \texttt{---------}\\
8 & 3 & 4 & 1 & 3 & 0 & 512 & 0 & \texttt{---------}\\
\multicolumn{9}{@{}l@{}}{\rule{0pt}{2.6ex}\textbf{T9.059}\enspace \texttt{012345443}\enspace $n{=}9$, $|H|{=}4$, len 6, br 2, $|\mathrm{Aut}|{=}2$, orb 8/7, hist-orb 3/2; M 13312, res 4896, rev 4712: ok}\\*
0 & -- & 0 & 4 & 1 & 422 & 2 & 88 & \texttt{-++??????}\\*
1 & 0 & 1 & 4 & 1 & 332 & 4 & 176 & \texttt{--+??????}\\
2 & 1 & 2 & 4 & 1 & 152 & 8 & 352 & \texttt{---??????}\\
3 & 2 & 3 & 3 & 1 & 96 & 32 & 384 & \texttt{----????-}\\
4 & 3 & 4 & 1 & 1 & 256 & 256 & 0 & \texttt{-----+---}\\
5 & 4 & 5 & 1 & 1 & 0 & 512 & 0 & \texttt{---------}\\
6 & 3 & 4 & 1 & 2 & 0 & 512 & 0 & \texttt{---------}\\
7 & 3 & 4 & 1 & 2 & 0 & 512 & 0 & \texttt{---------}\\
8 & 2 & 3 & 1 & 1 & 0 & 512 & 0 & \texttt{---------}\\
\multicolumn{9}{@{}l@{}}{\rule{0pt}{2.6ex}\textbf{T9.060}\enspace \texttt{012345442}\enspace $n{=}9$, $|H|{=}4$, len 6, br 2, $|\mathrm{Aut}|{=}2$, orb 8/7, hist-orb 3/2; M 12800, res 5208, rev 4224: ok}\\*
0 & -- & 0 & 4 & 1 & 358 & 2 & 152 & \texttt{-+???????}\\*
1 & 0 & 1 & 4 & 1 & 204 & 4 & 304 & \texttt{--???????}\\
2 & 1 & 2 & 3 & 1 & 304 & 16 & 192 & \texttt{---+????-}\\
3 & 2 & 3 & 3 & 1 & 96 & 32 & 384 & \texttt{----????-}\\
4 & 3 & 4 & 1 & 1 & 256 & 256 & 0 & \texttt{-----+---}\\
5 & 4 & 5 & 1 & 1 & 0 & 512 & 0 & \texttt{---------}\\
6 & 3 & 4 & 1 & 2 & 0 & 512 & 0 & \texttt{---------}\\
7 & 3 & 4 & 1 & 2 & 0 & 512 & 0 & \texttt{---------}\\
8 & 1 & 2 & 1 & 1 & 0 & 512 & 0 & \texttt{---------}\\
\multicolumn{9}{@{}l@{}}{\rule{0pt}{2.6ex}\textbf{T9.061}\enspace \texttt{012345441}\enspace $n{=}9$, $|H|{=}4$, len 6, br 2, $|\mathrm{Aut}|{=}2$, orb 8/7, hist-orb 3/2; M 12288, res 4784, rev 4160: ok}\\*
0 & -- & 0 & 4 & 1 & 230 & 2 & 280 & \texttt{-????????}\\*
1 & 0 & 1 & 3 & 1 & 408 & 8 & 96 & \texttt{--++????-}\\
2 & 1 & 2 & 3 & 1 & 304 & 16 & 192 & \texttt{---+????-}\\
3 & 2 & 3 & 3 & 1 & 96 & 32 & 384 & \texttt{----????-}\\
4 & 3 & 4 & 1 & 1 & 256 & 256 & 0 & \texttt{-----+---}\\
5 & 4 & 5 & 1 & 1 & 0 & 512 & 0 & \texttt{---------}\\
6 & 3 & 4 & 1 & 2 & 0 & 512 & 0 & \texttt{---------}\\
7 & 3 & 4 & 1 & 2 & 0 & 512 & 0 & \texttt{---------}\\
8 & 0 & 1 & 1 & 1 & 0 & 512 & 0 & \texttt{---------}\\
\multicolumn{9}{@{}l@{}}{\rule{0pt}{2.6ex}\textbf{T9.062}\enspace \texttt{012345434}\enspace $n{=}9$, $|H|{=}3$, len 6, br 2, $|\mathrm{Aut}|{=}1$, orb 9/9, hist-orb 3/3; M 13312, res 3064, rev 5404: ok}\\*
0 & -- & 0 & 3 & 1 & 450 & 2 & 60 & \texttt{-++??????}\\*
1 & 0 & 1 & 3 & 1 & 388 & 4 & 120 & \texttt{--+??????}\\
2 & 1 & 2 & 3 & 1 & 264 & 8 & 240 & \texttt{---??????}\\
3 & 2 & 3 & 2 & 1 & 192 & 64 & 256 & \texttt{----???--}\\
4 & 3 & 4 & 1 & 1 & 256 & 256 & 0 & \texttt{-----+---}\\
5 & 4 & 5 & 1 & 1 & 0 & 512 & 0 & \texttt{---------}\\
6 & 3 & 4 & 1 & 1 & 0 & 512 & 0 & \texttt{---------}\\
7 & 2 & 3 & 1 & 1 & 256 & 256 & 0 & \texttt{--------+}\\
8 & 7 & 4 & 1 & 1 & 0 & 512 & 0 & \texttt{---------}\\
\multicolumn{9}{@{}l@{}}{\rule{0pt}{2.6ex}\textbf{T9.063}\enspace \texttt{012345433}\enspace $n{=}9$, $|H|{=}4$, len 6, br 2, $|\mathrm{Aut}|{=}2$, orb 8/7, hist-orb 3/2; M 12800, res 4912, rev 4376: ok}\\*
0 & -- & 0 & 4 & 1 & 406 & 2 & 104 & \texttt{-++??????}\\*
1 & 0 & 1 & 4 & 1 & 300 & 4 & 208 & \texttt{--+??????}\\
2 & 1 & 2 & 4 & 1 & 88 & 8 & 416 & \texttt{---??????}\\
3 & 2 & 3 & 2 & 1 & 192 & 64 & 256 & \texttt{----???--}\\
4 & 3 & 4 & 1 & 1 & 256 & 256 & 0 & \texttt{-----+---}\\
5 & 4 & 5 & 1 & 1 & 0 & 512 & 0 & \texttt{---------}\\
6 & 3 & 4 & 1 & 1 & 0 & 512 & 0 & \texttt{---------}\\
7 & 2 & 3 & 1 & 2 & 0 & 512 & 0 & \texttt{---------}\\
8 & 2 & 3 & 1 & 2 & 0 & 512 & 0 & \texttt{---------}\\
\multicolumn{9}{@{}l@{}}{\rule{0pt}{2.6ex}\textbf{T9.064}\enspace \texttt{012345432}\enspace $n{=}9$, $|H|{=}4$, len 6, br 3, $|\mathrm{Aut}|{=}1$, orb 9/9, hist-orb 4/4; M 12288, res 5432, rev 3716: ok}\\*
0 & -- & 0 & 4 & 1 & 342 & 2 & 168 & \texttt{-+???????}\\*
1 & 0 & 1 & 4 & 1 & 172 & 4 & 336 & \texttt{--???????}\\
2 & 1 & 2 & 3 & 1 & 176 & 16 & 320 & \texttt{---?????-}\\
3 & 2 & 3 & 2 & 1 & 192 & 64 & 256 & \texttt{----???--}\\
4 & 3 & 4 & 1 & 1 & 256 & 256 & 0 & \texttt{-----+---}\\
5 & 4 & 5 & 1 & 1 & 0 & 512 & 0 & \texttt{---------}\\
6 & 3 & 4 & 1 & 1 & 0 & 512 & 0 & \texttt{---------}\\
7 & 2 & 3 & 1 & 1 & 0 & 512 & 0 & \texttt{---------}\\
8 & 1 & 2 & 1 & 1 & 0 & 512 & 0 & \texttt{---------}\\
\multicolumn{9}{@{}l@{}}{\rule{0pt}{2.6ex}\textbf{T9.065}\enspace \texttt{012345431}\enspace $n{=}9$, $|H|{=}4$, len 6, br 3, $|\mathrm{Aut}|{=}1$, orb 9/9, hist-orb 4/4; M 11776, res 5152, rev 3572: ok}\\*
0 & -- & 0 & 4 & 1 & 214 & 2 & 296 & \texttt{-????????}\\*
1 & 0 & 1 & 3 & 1 & 344 & 8 & 160 & \texttt{--+?????-}\\
2 & 1 & 2 & 3 & 1 & 176 & 16 & 320 & \texttt{---?????-}\\
3 & 2 & 3 & 2 & 1 & 192 & 64 & 256 & \texttt{----???--}\\
4 & 3 & 4 & 1 & 1 & 256 & 256 & 0 & \texttt{-----+---}\\
5 & 4 & 5 & 1 & 1 & 0 & 512 & 0 & \texttt{---------}\\
6 & 3 & 4 & 1 & 1 & 0 & 512 & 0 & \texttt{---------}\\
7 & 2 & 3 & 1 & 1 & 0 & 512 & 0 & \texttt{---------}\\
8 & 0 & 1 & 1 & 1 & 0 & 512 & 0 & \texttt{---------}\\
\multicolumn{9}{@{}l@{}}{\rule{0pt}{2.6ex}\textbf{T9.066}\enspace \texttt{012345423}\enspace $n{=}9$, $|H|{=}3$, len 6, br 2, $|\mathrm{Aut}|{=}1$, orb 9/9, hist-orb 3/3; M 12288, res 2892, rev 4856: ok}\\*
0 & -- & 0 & 3 & 1 & 418 & 2 & 92 & \texttt{-+???????}\\*
1 & 0 & 1 & 3 & 1 & 324 & 4 & 184 & \texttt{--???????}\\
2 & 1 & 2 & 2 & 1 & 352 & 32 & 128 & \texttt{---+???--}\\
3 & 2 & 3 & 2 & 1 & 192 & 64 & 256 & \texttt{----???--}\\
4 & 3 & 4 & 1 & 1 & 256 & 256 & 0 & \texttt{-----+---}\\
5 & 4 & 5 & 1 & 1 & 0 & 512 & 0 & \texttt{---------}\\
6 & 3 & 4 & 1 & 1 & 0 & 512 & 0 & \texttt{---------}\\
7 & 1 & 2 & 1 & 1 & 256 & 256 & 0 & \texttt{--------+}\\
8 & 7 & 3 & 1 & 1 & 0 & 512 & 0 & \texttt{---------}\\
\multicolumn{9}{@{}l@{}}{\rule{0pt}{2.6ex}\textbf{T9.067}\enspace \texttt{012345422}\enspace $n{=}9$, $|H|{=}4$, len 6, br 2, $|\mathrm{Aut}|{=}2$, orb 8/7, hist-orb 3/2; M 11776, res 5160, rev 3552: ok}\\*
0 & -- & 0 & 4 & 1 & 310 & 2 & 200 & \texttt{-+???????}\\*
1 & 0 & 1 & 4 & 1 & 108 & 4 & 400 & \texttt{--???????}\\
2 & 1 & 2 & 2 & 1 & 352 & 32 & 128 & \texttt{---+???--}\\
3 & 2 & 3 & 2 & 1 & 192 & 64 & 256 & \texttt{----???--}\\
4 & 3 & 4 & 1 & 1 & 256 & 256 & 0 & \texttt{-----+---}\\
5 & 4 & 5 & 1 & 1 & 0 & 512 & 0 & \texttt{---------}\\
6 & 3 & 4 & 1 & 1 & 0 & 512 & 0 & \texttt{---------}\\
7 & 1 & 2 & 1 & 2 & 0 & 512 & 0 & \texttt{---------}\\
8 & 1 & 2 & 1 & 2 & 0 & 512 & 0 & \texttt{---------}\\
\multicolumn{9}{@{}l@{}}{\rule{0pt}{2.6ex}\textbf{T9.068}\enspace \texttt{012345421}\enspace $n{=}9$, $|H|{=}4$, len 6, br 3, $|\mathrm{Aut}|{=}1$, orb 9/9, hist-orb 4/4; M 11264, res 5168, rev 3236: ok}\\*
0 & -- & 0 & 4 & 1 & 182 & 2 & 328 & \texttt{-????????}\\*
1 & 0 & 1 & 3 & 1 & 216 & 8 & 288 & \texttt{--??????-}\\
2 & 1 & 2 & 2 & 1 & 352 & 32 & 128 & \texttt{---+???--}\\
3 & 2 & 3 & 2 & 1 & 192 & 64 & 256 & \texttt{----???--}\\
4 & 3 & 4 & 1 & 1 & 256 & 256 & 0 & \texttt{-----+---}\\
5 & 4 & 5 & 1 & 1 & 0 & 512 & 0 & \texttt{---------}\\
6 & 3 & 4 & 1 & 1 & 0 & 512 & 0 & \texttt{---------}\\
7 & 1 & 2 & 1 & 1 & 0 & 512 & 0 & \texttt{---------}\\
8 & 0 & 1 & 1 & 1 & 0 & 512 & 0 & \texttt{---------}\\
\multicolumn{9}{@{}l@{}}{\rule{0pt}{2.6ex}\textbf{T9.069}\enspace \texttt{012345412}\enspace $n{=}9$, $|H|{=}3$, len 6, br 2, $|\mathrm{Aut}|{=}1$, orb 9/9, hist-orb 3/3; M 11264, res 2496, rev 4412: ok}\\*
0 & -- & 0 & 3 & 1 & 354 & 2 & 156 & \texttt{-????????}\\*
1 & 0 & 1 & 2 & 1 & 432 & 16 & 64 & \texttt{--++???--}\\
2 & 1 & 2 & 2 & 1 & 352 & 32 & 128 & \texttt{---+???--}\\
3 & 2 & 3 & 2 & 1 & 192 & 64 & 256 & \texttt{----???--}\\
4 & 3 & 4 & 1 & 1 & 256 & 256 & 0 & \texttt{-----+---}\\
5 & 4 & 5 & 1 & 1 & 0 & 512 & 0 & \texttt{---------}\\
6 & 3 & 4 & 1 & 1 & 0 & 512 & 0 & \texttt{---------}\\
7 & 0 & 1 & 1 & 1 & 256 & 256 & 0 & \texttt{--------+}\\
8 & 7 & 2 & 1 & 1 & 0 & 512 & 0 & \texttt{---------}\\
\multicolumn{9}{@{}l@{}}{\rule{0pt}{2.6ex}\textbf{T9.070}\enspace \texttt{012345411}\enspace $n{=}9$, $|H|{=}4$, len 6, br 2, $|\mathrm{Aut}|{=}2$, orb 8/7, hist-orb 3/2; M 10752, res 4384, rev 3288: ok}\\*
0 & -- & 0 & 4 & 1 & 118 & 2 & 392 & \texttt{-????????}\\*
1 & 0 & 1 & 2 & 1 & 432 & 16 & 64 & \texttt{--++???--}\\
2 & 1 & 2 & 2 & 1 & 352 & 32 & 128 & \texttt{---+???--}\\
3 & 2 & 3 & 2 & 1 & 192 & 64 & 256 & \texttt{----???--}\\
4 & 3 & 4 & 1 & 1 & 256 & 256 & 0 & \texttt{-----+---}\\
5 & 4 & 5 & 1 & 1 & 0 & 512 & 0 & \texttt{---------}\\
6 & 3 & 4 & 1 & 1 & 0 & 512 & 0 & \texttt{---------}\\
7 & 0 & 1 & 1 & 2 & 0 & 512 & 0 & \texttt{---------}\\
8 & 0 & 1 & 1 & 2 & 0 & 512 & 0 & \texttt{---------}\\
\multicolumn{9}{@{}l@{}}{\rule{0pt}{2.6ex}\textbf{T9.071}\enspace \texttt{012345345}\enspace $n{=}9$, $|H|{=}2$, len 6, br 1, $|\mathrm{Aut}|{=}2$, orb 6/3, hist-orb 1/0; M 13824, res 1176, rev 5944: ok}\\*
0 & -- & 0 & 2 & 1 & 482 & 2 & 28 & \texttt{-++??????}\\*
1 & 0 & 1 & 2 & 1 & 452 & 4 & 56 & \texttt{--+??????}\\
2 & 1 & 2 & 2 & 1 & 392 & 8 & 112 & \texttt{---??????}\\
3 & 2 & 3 & 1 & 2 & 384 & 128 & 0 & \texttt{----++---}\\
4 & 3 & 4 & 1 & 2 & 256 & 256 & 0 & \texttt{-----+---}\\
5 & 4 & 5 & 1 & 2 & 0 & 512 & 0 & \texttt{---------}\\
6 & 2 & 3 & 1 & 2 & 384 & 128 & 0 & \texttt{-------++}\\
7 & 6 & 4 & 1 & 2 & 256 & 256 & 0 & \texttt{--------+}\\
8 & 7 & 5 & 1 & 2 & 0 & 512 & 0 & \texttt{---------}\\
\multicolumn{9}{@{}l@{}}{\rule{0pt}{2.6ex}\textbf{T9.072}\enspace \texttt{012345344}\enspace $n{=}9$, $|H|{=}3$, len 6, br 2, $|\mathrm{Aut}|{=}2$, orb 8/7, hist-orb 2/1; M 13312, res 2640, rev 5444: ok}\\*
0 & -- & 0 & 3 & 1 & 454 & 2 & 56 & \texttt{-++??????}\\*
1 & 0 & 1 & 3 & 1 & 396 & 4 & 112 & \texttt{--+??????}\\
2 & 1 & 2 & 3 & 1 & 280 & 8 & 224 & \texttt{---??????}\\
3 & 2 & 3 & 1 & 1 & 384 & 128 & 0 & \texttt{----++---}\\
4 & 3 & 4 & 1 & 1 & 256 & 256 & 0 & \texttt{-----+---}\\
5 & 4 & 5 & 1 & 1 & 0 & 512 & 0 & \texttt{---------}\\
6 & 2 & 3 & 2 & 1 & 128 & 128 & 256 & \texttt{-------??}\\
7 & 6 & 4 & 1 & 2 & 0 & 512 & 0 & \texttt{---------}\\
8 & 6 & 4 & 1 & 2 & 0 & 512 & 0 & \texttt{---------}\\
\multicolumn{9}{@{}l@{}}{\rule{0pt}{2.6ex}\textbf{T9.073}\enspace \texttt{012345343}\enspace $n{=}9$, $|H|{=}3$, len 6, br 1, $|\mathrm{Aut}|{=}1$, orb 9/9, hist-orb 3/3; M 12800, res 3528, rev 4904: ok}\\*
0 & -- & 0 & 3 & 1 & 426 & 2 & 84 & \texttt{-++??????}\\*
1 & 0 & 1 & 3 & 1 & 340 & 4 & 168 & \texttt{--+??????}\\
2 & 1 & 2 & 3 & 1 & 168 & 8 & 336 & \texttt{---??????}\\
3 & 2 & 3 & 1 & 1 & 384 & 128 & 0 & \texttt{----++---}\\
4 & 3 & 4 & 1 & 1 & 256 & 256 & 0 & \texttt{-----+---}\\
5 & 4 & 5 & 1 & 1 & 0 & 512 & 0 & \texttt{---------}\\
6 & 2 & 3 & 1 & 1 & 256 & 256 & 0 & \texttt{-------+-}\\
7 & 6 & 4 & 1 & 1 & 0 & 512 & 0 & \texttt{---------}\\
8 & 2 & 3 & 1 & 1 & 0 & 512 & 0 & \texttt{---------}\\
\multicolumn{9}{@{}l@{}}{\rule{0pt}{2.6ex}\textbf{T9.074}\enspace \texttt{012345342}\enspace $n{=}9$, $|H|{=}3$, len 6, br 2, $|\mathrm{Aut}|{=}1$, orb 9/9, hist-orb 3/3; M 12288, res 3788, rev 4412: ok}\\*
0 & -- & 0 & 3 & 1 & 362 & 2 & 148 & \texttt{-+???????}\\*
1 & 0 & 1 & 3 & 1 & 212 & 4 & 296 & \texttt{--???????}\\
2 & 1 & 2 & 2 & 1 & 336 & 16 & 160 & \texttt{---?????-}\\
3 & 2 & 3 & 1 & 1 & 384 & 128 & 0 & \texttt{----++---}\\
4 & 3 & 4 & 1 & 1 & 256 & 256 & 0 & \texttt{-----+---}\\
5 & 4 & 5 & 1 & 1 & 0 & 512 & 0 & \texttt{---------}\\
6 & 2 & 3 & 1 & 1 & 256 & 256 & 0 & \texttt{-------+-}\\
7 & 6 & 4 & 1 & 1 & 0 & 512 & 0 & \texttt{---------}\\
8 & 1 & 2 & 1 & 1 & 0 & 512 & 0 & \texttt{---------}\\
\multicolumn{9}{@{}l@{}}{\rule{0pt}{2.6ex}\textbf{T9.075}\enspace \texttt{012345341}\enspace $n{=}9$, $|H|{=}3$, len 6, br 2, $|\mathrm{Aut}|{=}1$, orb 9/9, hist-orb 3/3; M 11776, res 3328, rev 4316: ok}\\*
0 & -- & 0 & 3 & 1 & 234 & 2 & 276 & \texttt{-????????}\\*
1 & 0 & 1 & 2 & 1 & 424 & 8 & 80 & \texttt{--+?????-}\\
2 & 1 & 2 & 2 & 1 & 336 & 16 & 160 & \texttt{---?????-}\\
3 & 2 & 3 & 1 & 1 & 384 & 128 & 0 & \texttt{----++---}\\
4 & 3 & 4 & 1 & 1 & 256 & 256 & 0 & \texttt{-----+---}\\
5 & 4 & 5 & 1 & 1 & 0 & 512 & 0 & \texttt{---------}\\
6 & 2 & 3 & 1 & 1 & 256 & 256 & 0 & \texttt{-------+-}\\
7 & 6 & 4 & 1 & 1 & 0 & 512 & 0 & \texttt{---------}\\
8 & 0 & 1 & 1 & 1 & 0 & 512 & 0 & \texttt{---------}\\
\multicolumn{9}{@{}l@{}}{\rule{0pt}{2.6ex}\textbf{T9.076}\enspace \texttt{012345333}\enspace $n{=}9$, $|H|{=}4$, len 6, br 1, $|\mathrm{Aut}|{=}6$, orb 7/6, hist-orb 2/1; M 12288, res 4704, rev 4268: ok}\\*
0 & -- & 0 & 4 & 1 & 398 & 2 & 112 & \texttt{-++??????}\\*
1 & 0 & 1 & 4 & 1 & 284 & 4 & 224 & \texttt{--+??????}\\
2 & 1 & 2 & 4 & 1 & 56 & 8 & 448 & \texttt{---??????}\\
3 & 2 & 3 & 1 & 1 & 384 & 128 & 0 & \texttt{----++---}\\
4 & 3 & 4 & 1 & 1 & 256 & 256 & 0 & \texttt{-----+---}\\
5 & 4 & 5 & 1 & 1 & 0 & 512 & 0 & \texttt{---------}\\
6 & 2 & 3 & 1 & 3 & 0 & 512 & 0 & \texttt{---------}\\
7 & 2 & 3 & 1 & 3 & 0 & 512 & 0 & \texttt{---------}\\
8 & 2 & 3 & 1 & 3 & 0 & 512 & 0 & \texttt{---------}\\
\multicolumn{9}{@{}l@{}}{\rule{0pt}{2.6ex}\textbf{T9.077}\enspace \texttt{012345332}\enspace $n{=}9$, $|H|{=}4$, len 6, br 2, $|\mathrm{Aut}|{=}2$, orb 8/7, hist-orb 3/2; M 11776, res 5328, rev 3504: ok}\\*
0 & -- & 0 & 4 & 1 & 334 & 2 & 176 & \texttt{-+???????}\\*
1 & 0 & 1 & 4 & 1 & 156 & 4 & 352 & \texttt{--???????}\\
2 & 1 & 2 & 3 & 1 & 112 & 16 & 384 & \texttt{---?????-}\\
3 & 2 & 3 & 1 & 1 & 384 & 128 & 0 & \texttt{----++---}\\
4 & 3 & 4 & 1 & 1 & 256 & 256 & 0 & \texttt{-----+---}\\
5 & 4 & 5 & 1 & 1 & 0 & 512 & 0 & \texttt{---------}\\
6 & 2 & 3 & 1 & 2 & 0 & 512 & 0 & \texttt{---------}\\
7 & 2 & 3 & 1 & 2 & 0 & 512 & 0 & \texttt{---------}\\
8 & 1 & 2 & 1 & 1 & 0 & 512 & 0 & \texttt{---------}\\
\multicolumn{9}{@{}l@{}}{\rule{0pt}{2.6ex}\textbf{T9.078}\enspace \texttt{012345331}\enspace $n{=}9$, $|H|{=}4$, len 6, br 2, $|\mathrm{Aut}|{=}2$, orb 8/7, hist-orb 3/2; M 11264, res 5120, rev 3320: ok}\\*
0 & -- & 0 & 4 & 1 & 206 & 2 & 304 & \texttt{-????????}\\*
1 & 0 & 1 & 3 & 1 & 312 & 8 & 192 & \texttt{--+?????-}\\
2 & 1 & 2 & 3 & 1 & 112 & 16 & 384 & \texttt{---?????-}\\
3 & 2 & 3 & 1 & 1 & 384 & 128 & 0 & \texttt{----++---}\\
4 & 3 & 4 & 1 & 1 & 256 & 256 & 0 & \texttt{-----+---}\\
5 & 4 & 5 & 1 & 1 & 0 & 512 & 0 & \texttt{---------}\\
6 & 2 & 3 & 1 & 2 & 0 & 512 & 0 & \texttt{---------}\\
7 & 2 & 3 & 1 & 2 & 0 & 512 & 0 & \texttt{---------}\\
8 & 0 & 1 & 1 & 1 & 0 & 512 & 0 & \texttt{---------}\\
\multicolumn{9}{@{}l@{}}{\rule{0pt}{2.6ex}\textbf{T9.079}\enspace \texttt{012345323}\enspace $n{=}9$, $|H|{=}3$, len 6, br 2, $|\mathrm{Aut}|{=}1$, orb 9/9, hist-orb 3/3; M 11776, res 3268, rev 4392: ok}\\*
0 & -- & 0 & 3 & 1 & 394 & 2 & 116 & \texttt{-+???????}\\*
1 & 0 & 1 & 3 & 1 & 276 & 4 & 232 & \texttt{--???????}\\
2 & 1 & 2 & 2 & 1 & 224 & 32 & 256 & \texttt{---????--}\\
3 & 2 & 3 & 1 & 1 & 384 & 128 & 0 & \texttt{----++---}\\
4 & 3 & 4 & 1 & 1 & 256 & 256 & 0 & \texttt{-----+---}\\
5 & 4 & 5 & 1 & 1 & 0 & 512 & 0 & \texttt{---------}\\
6 & 2 & 3 & 1 & 1 & 0 & 512 & 0 & \texttt{---------}\\
7 & 1 & 2 & 1 & 1 & 256 & 256 & 0 & \texttt{--------+}\\
8 & 7 & 3 & 1 & 1 & 0 & 512 & 0 & \texttt{---------}\\
\multicolumn{9}{@{}l@{}}{\rule{0pt}{2.6ex}\textbf{T9.080}\enspace \texttt{012345322}\enspace $n{=}9$, $|H|{=}4$, len 6, br 2, $|\mathrm{Aut}|{=}2$, orb 8/7, hist-orb 3/2; M 11264, res 5200, rev 3232: ok}\\*
0 & -- & 0 & 4 & 1 & 302 & 2 & 208 & \texttt{-+???????}\\*
1 & 0 & 1 & 4 & 1 & 92 & 4 & 416 & \texttt{--???????}\\
2 & 1 & 2 & 2 & 1 & 224 & 32 & 256 & \texttt{---????--}\\
3 & 2 & 3 & 1 & 1 & 384 & 128 & 0 & \texttt{----++---}\\
4 & 3 & 4 & 1 & 1 & 256 & 256 & 0 & \texttt{-----+---}\\
5 & 4 & 5 & 1 & 1 & 0 & 512 & 0 & \texttt{---------}\\
6 & 2 & 3 & 1 & 1 & 0 & 512 & 0 & \texttt{---------}\\
7 & 1 & 2 & 1 & 2 & 0 & 512 & 0 & \texttt{---------}\\
8 & 1 & 2 & 1 & 2 & 0 & 512 & 0 & \texttt{---------}\\
\multicolumn{9}{@{}l@{}}{\rule{0pt}{2.6ex}\textbf{T9.081}\enspace \texttt{012345321}\enspace $n{=}9$, $|H|{=}4$, len 6, br 3, $|\mathrm{Aut}|{=}1$, orb 9/9, hist-orb 4/4; M 10752, res 5280, rev 2868: ok}\\*
0 & -- & 0 & 4 & 1 & 174 & 2 & 336 & \texttt{-????????}\\*
1 & 0 & 1 & 3 & 1 & 184 & 8 & 320 & \texttt{--??????-}\\
2 & 1 & 2 & 2 & 1 & 224 & 32 & 256 & \texttt{---????--}\\
3 & 2 & 3 & 1 & 1 & 384 & 128 & 0 & \texttt{----++---}\\
4 & 3 & 4 & 1 & 1 & 256 & 256 & 0 & \texttt{-----+---}\\
5 & 4 & 5 & 1 & 1 & 0 & 512 & 0 & \texttt{---------}\\
6 & 2 & 3 & 1 & 1 & 0 & 512 & 0 & \texttt{---------}\\
7 & 1 & 2 & 1 & 1 & 0 & 512 & 0 & \texttt{---------}\\
8 & 0 & 1 & 1 & 1 & 0 & 512 & 0 & \texttt{---------}\\
\multicolumn{9}{@{}l@{}}{\rule{0pt}{2.6ex}\textbf{T9.082}\enspace \texttt{012345312}\enspace $n{=}9$, $|H|{=}3$, len 6, br 2, $|\mathrm{Aut}|{=}1$, orb 9/9, hist-orb 3/3; M 10752, res 2848, rev 3972: ok}\\*
0 & -- & 0 & 3 & 1 & 330 & 2 & 180 & \texttt{-????????}\\*
1 & 0 & 1 & 2 & 1 & 368 & 16 & 128 & \texttt{--+????--}\\
2 & 1 & 2 & 2 & 1 & 224 & 32 & 256 & \texttt{---????--}\\
3 & 2 & 3 & 1 & 1 & 384 & 128 & 0 & \texttt{----++---}\\
4 & 3 & 4 & 1 & 1 & 256 & 256 & 0 & \texttt{-----+---}\\
5 & 4 & 5 & 1 & 1 & 0 & 512 & 0 & \texttt{---------}\\
6 & 2 & 3 & 1 & 1 & 0 & 512 & 0 & \texttt{---------}\\
7 & 0 & 1 & 1 & 1 & 256 & 256 & 0 & \texttt{--------+}\\
8 & 7 & 2 & 1 & 1 & 0 & 512 & 0 & \texttt{---------}\\
\multicolumn{9}{@{}l@{}}{\rule{0pt}{2.6ex}\textbf{T9.083}\enspace \texttt{012345311}\enspace $n{=}9$, $|H|{=}4$, len 6, br 2, $|\mathrm{Aut}|{=}2$, orb 8/7, hist-orb 3/2; M 10240, res 4608, rev 2880: ok}\\*
0 & -- & 0 & 4 & 1 & 110 & 2 & 400 & \texttt{-????????}\\*
1 & 0 & 1 & 2 & 1 & 368 & 16 & 128 & \texttt{--+????--}\\
2 & 1 & 2 & 2 & 1 & 224 & 32 & 256 & \texttt{---????--}\\
3 & 2 & 3 & 1 & 1 & 384 & 128 & 0 & \texttt{----++---}\\
4 & 3 & 4 & 1 & 1 & 256 & 256 & 0 & \texttt{-----+---}\\
5 & 4 & 5 & 1 & 1 & 0 & 512 & 0 & \texttt{---------}\\
6 & 2 & 3 & 1 & 1 & 0 & 512 & 0 & \texttt{---------}\\
7 & 0 & 1 & 1 & 2 & 0 & 512 & 0 & \texttt{---------}\\
8 & 0 & 1 & 1 & 2 & 0 & 512 & 0 & \texttt{---------}\\
\multicolumn{9}{@{}l@{}}{\rule{0pt}{2.6ex}\textbf{T9.084}\enspace \texttt{012345234}\enspace $n{=}9$, $|H|{=}2$, len 6, br 1, $|\mathrm{Aut}|{=}1$, orb 9/9, hist-orb 2/2; M 12288, res 924, rev 5252: ok}\\*
0 & -- & 0 & 2 & 1 & 466 & 2 & 44 & \texttt{-+???????}\\*
1 & 0 & 1 & 2 & 1 & 420 & 4 & 88 & \texttt{--???????}\\
2 & 1 & 2 & 1 & 1 & 448 & 64 & 0 & \texttt{---+++---}\\
3 & 2 & 3 & 1 & 1 & 384 & 128 & 0 & \texttt{----++---}\\
4 & 3 & 4 & 1 & 1 & 256 & 256 & 0 & \texttt{-----+---}\\
5 & 4 & 5 & 1 & 1 & 0 & 512 & 0 & \texttt{---------}\\
6 & 1 & 2 & 1 & 1 & 384 & 128 & 0 & \texttt{-------++}\\
7 & 6 & 3 & 1 & 1 & 256 & 256 & 0 & \texttt{--------+}\\
8 & 7 & 4 & 1 & 1 & 0 & 512 & 0 & \texttt{---------}\\
\multicolumn{9}{@{}l@{}}{\rule{0pt}{2.6ex}\textbf{T9.085}\enspace \texttt{012345233}\enspace $n{=}9$, $|H|{=}3$, len 6, br 2, $|\mathrm{Aut}|{=}2$, orb 8/7, hist-orb 2/1; M 11776, res 2504, rev 4556: ok}\\*
0 & -- & 0 & 3 & 1 & 406 & 2 & 104 & \texttt{-+???????}\\*
1 & 0 & 1 & 3 & 1 & 300 & 4 & 208 & \texttt{--???????}\\
2 & 1 & 2 & 1 & 1 & 448 & 64 & 0 & \texttt{---+++---}\\
3 & 2 & 3 & 1 & 1 & 384 & 128 & 0 & \texttt{----++---}\\
4 & 3 & 4 & 1 & 1 & 256 & 256 & 0 & \texttt{-----+---}\\
5 & 4 & 5 & 1 & 1 & 0 & 512 & 0 & \texttt{---------}\\
6 & 1 & 2 & 2 & 1 & 128 & 128 & 256 & \texttt{-------??}\\
7 & 6 & 3 & 1 & 2 & 0 & 512 & 0 & \texttt{---------}\\
8 & 6 & 3 & 1 & 2 & 0 & 512 & 0 & \texttt{---------}\\
\multicolumn{9}{@{}l@{}}{\rule{0pt}{2.6ex}\textbf{T9.086}\enspace \texttt{012345232}\enspace $n{=}9$, $|H|{=}3$, len 6, br 1, $|\mathrm{Aut}|{=}1$, orb 9/9, hist-orb 3/3; M 11264, res 3444, rev 3948: ok}\\*
0 & -- & 0 & 3 & 1 & 346 & 2 & 164 & \texttt{-+???????}\\*
1 & 0 & 1 & 3 & 1 & 180 & 4 & 328 & \texttt{--???????}\\
2 & 1 & 2 & 1 & 1 & 448 & 64 & 0 & \texttt{---+++---}\\
3 & 2 & 3 & 1 & 1 & 384 & 128 & 0 & \texttt{----++---}\\
4 & 3 & 4 & 1 & 1 & 256 & 256 & 0 & \texttt{-----+---}\\
5 & 4 & 5 & 1 & 1 & 0 & 512 & 0 & \texttt{---------}\\
6 & 1 & 2 & 1 & 1 & 256 & 256 & 0 & \texttt{-------+-}\\
7 & 6 & 3 & 1 & 1 & 0 & 512 & 0 & \texttt{---------}\\
8 & 1 & 2 & 1 & 1 & 0 & 512 & 0 & \texttt{---------}\\
\multicolumn{9}{@{}l@{}}{\rule{0pt}{2.6ex}\textbf{T9.087}\enspace \texttt{012345231}\enspace $n{=}9$, $|H|{=}3$, len 6, br 2, $|\mathrm{Aut}|{=}1$, orb 9/9, hist-orb 3/3; M 10752, res 3128, rev 3760: ok}\\*
0 & -- & 0 & 3 & 1 & 218 & 2 & 292 & \texttt{-????????}\\*
1 & 0 & 1 & 2 & 1 & 360 & 8 & 144 & \texttt{--??????-}\\
2 & 1 & 2 & 1 & 1 & 448 & 64 & 0 & \texttt{---+++---}\\
3 & 2 & 3 & 1 & 1 & 384 & 128 & 0 & \texttt{----++---}\\
4 & 3 & 4 & 1 & 1 & 256 & 256 & 0 & \texttt{-----+---}\\
5 & 4 & 5 & 1 & 1 & 0 & 512 & 0 & \texttt{---------}\\
6 & 1 & 2 & 1 & 1 & 256 & 256 & 0 & \texttt{-------+-}\\
7 & 6 & 3 & 1 & 1 & 0 & 512 & 0 & \texttt{---------}\\
8 & 0 & 1 & 1 & 1 & 0 & 512 & 0 & \texttt{---------}\\
\multicolumn{9}{@{}l@{}}{\rule{0pt}{2.6ex}\textbf{T9.088}\enspace \texttt{012345222}\enspace $n{=}9$, $|H|{=}4$, len 6, br 1, $|\mathrm{Aut}|{=}6$, orb 7/6, hist-orb 2/1; M 10752, res 4704, rev 3148: ok}\\*
0 & -- & 0 & 4 & 1 & 286 & 2 & 224 & \texttt{-+???????}\\*
1 & 0 & 1 & 4 & 1 & 60 & 4 & 448 & \texttt{--???????}\\
2 & 1 & 2 & 1 & 1 & 448 & 64 & 0 & \texttt{---+++---}\\
3 & 2 & 3 & 1 & 1 & 384 & 128 & 0 & \texttt{----++---}\\
4 & 3 & 4 & 1 & 1 & 256 & 256 & 0 & \texttt{-----+---}\\
5 & 4 & 5 & 1 & 1 & 0 & 512 & 0 & \texttt{---------}\\
6 & 1 & 2 & 1 & 3 & 0 & 512 & 0 & \texttt{---------}\\
7 & 1 & 2 & 1 & 3 & 0 & 512 & 0 & \texttt{---------}\\
8 & 1 & 2 & 1 & 3 & 0 & 512 & 0 & \texttt{---------}\\
\multicolumn{9}{@{}l@{}}{\rule{0pt}{2.6ex}\textbf{T9.089}\enspace \texttt{012345221}\enspace $n{=}9$, $|H|{=}4$, len 6, br 2, $|\mathrm{Aut}|{=}2$, orb 8/7, hist-orb 3/2; M 10240, res 4928, rev 2688: ok}\\*
0 & -- & 0 & 4 & 1 & 158 & 2 & 352 & \texttt{-????????}\\*
1 & 0 & 1 & 3 & 1 & 120 & 8 & 384 & \texttt{--??????-}\\
2 & 1 & 2 & 1 & 1 & 448 & 64 & 0 & \texttt{---+++---}\\
3 & 2 & 3 & 1 & 1 & 384 & 128 & 0 & \texttt{----++---}\\
4 & 3 & 4 & 1 & 1 & 256 & 256 & 0 & \texttt{-----+---}\\
5 & 4 & 5 & 1 & 1 & 0 & 512 & 0 & \texttt{---------}\\
6 & 1 & 2 & 1 & 2 & 0 & 512 & 0 & \texttt{---------}\\
7 & 1 & 2 & 1 & 2 & 0 & 512 & 0 & \texttt{---------}\\
8 & 0 & 1 & 1 & 1 & 0 & 512 & 0 & \texttt{---------}\\
\multicolumn{9}{@{}l@{}}{\rule{0pt}{2.6ex}\textbf{T9.090}\enspace \texttt{012345212}\enspace $n{=}9$, $|H|{=}3$, len 6, br 2, $|\mathrm{Aut}|{=}1$, orb 9/9, hist-orb 3/3; M 10240, res 2976, rev 3564: ok}\\*
0 & -- & 0 & 3 & 1 & 282 & 2 & 228 & \texttt{-????????}\\*
1 & 0 & 1 & 2 & 1 & 240 & 16 & 256 & \texttt{--?????--}\\
2 & 1 & 2 & 1 & 1 & 448 & 64 & 0 & \texttt{---+++---}\\
3 & 2 & 3 & 1 & 1 & 384 & 128 & 0 & \texttt{----++---}\\
4 & 3 & 4 & 1 & 1 & 256 & 256 & 0 & \texttt{-----+---}\\
5 & 4 & 5 & 1 & 1 & 0 & 512 & 0 & \texttt{---------}\\
6 & 1 & 2 & 1 & 1 & 0 & 512 & 0 & \texttt{---------}\\
7 & 0 & 1 & 1 & 1 & 256 & 256 & 0 & \texttt{--------+}\\
8 & 7 & 2 & 1 & 1 & 0 & 512 & 0 & \texttt{---------}\\
\multicolumn{9}{@{}l@{}}{\rule{0pt}{2.6ex}\textbf{T9.091}\enspace \texttt{012345211}\enspace $n{=}9$, $|H|{=}4$, len 6, br 2, $|\mathrm{Aut}|{=}2$, orb 8/7, hist-orb 3/2; M 9728, res 4480, rev 2584: ok}\\*
0 & -- & 0 & 4 & 1 & 94 & 2 & 416 & \texttt{-????????}\\*
1 & 0 & 1 & 2 & 1 & 240 & 16 & 256 & \texttt{--?????--}\\
2 & 1 & 2 & 1 & 1 & 448 & 64 & 0 & \texttt{---+++---}\\
3 & 2 & 3 & 1 & 1 & 384 & 128 & 0 & \texttt{----++---}\\
4 & 3 & 4 & 1 & 1 & 256 & 256 & 0 & \texttt{-----+---}\\
5 & 4 & 5 & 1 & 1 & 0 & 512 & 0 & \texttt{---------}\\
6 & 1 & 2 & 1 & 1 & 0 & 512 & 0 & \texttt{---------}\\
7 & 0 & 1 & 1 & 2 & 0 & 512 & 0 & \texttt{---------}\\
8 & 0 & 1 & 1 & 2 & 0 & 512 & 0 & \texttt{---------}\\
\multicolumn{9}{@{}l@{}}{\rule{0pt}{2.6ex}\textbf{T9.092}\enspace \texttt{012345123}\enspace $n{=}9$, $|H|{=}2$, len 6, br 1, $|\mathrm{Aut}|{=}1$, orb 9/9, hist-orb 2/2; M 10752, res 608, rev 4552: ok}\\*
0 & -- & 0 & 2 & 1 & 434 & 2 & 76 & \texttt{-????????}\\*
1 & 0 & 1 & 1 & 1 & 480 & 32 & 0 & \texttt{--++++---}\\
2 & 1 & 2 & 1 & 1 & 448 & 64 & 0 & \texttt{---+++---}\\
3 & 2 & 3 & 1 & 1 & 384 & 128 & 0 & \texttt{----++---}\\
4 & 3 & 4 & 1 & 1 & 256 & 256 & 0 & \texttt{-----+---}\\
5 & 4 & 5 & 1 & 1 & 0 & 512 & 0 & \texttt{---------}\\
6 & 0 & 1 & 1 & 1 & 384 & 128 & 0 & \texttt{-------++}\\
7 & 6 & 2 & 1 & 1 & 256 & 256 & 0 & \texttt{--------+}\\
8 & 7 & 3 & 1 & 1 & 0 & 512 & 0 & \texttt{---------}\\
\multicolumn{9}{@{}l@{}}{\rule{0pt}{2.6ex}\textbf{T9.093}\enspace \texttt{012345122}\enspace $n{=}9$, $|H|{=}3$, len 6, br 2, $|\mathrm{Aut}|{=}2$, orb 8/7, hist-orb 2/1; M 10240, res 1984, rev 3812: ok}\\*
0 & -- & 0 & 3 & 1 & 310 & 2 & 200 & \texttt{-????????}\\*
1 & 0 & 1 & 1 & 1 & 480 & 32 & 0 & \texttt{--++++---}\\
2 & 1 & 2 & 1 & 1 & 448 & 64 & 0 & \texttt{---+++---}\\
3 & 2 & 3 & 1 & 1 & 384 & 128 & 0 & \texttt{----++---}\\
4 & 3 & 4 & 1 & 1 & 256 & 256 & 0 & \texttt{-----+---}\\
5 & 4 & 5 & 1 & 1 & 0 & 512 & 0 & \texttt{---------}\\
6 & 0 & 1 & 2 & 1 & 128 & 128 & 256 & \texttt{-------??}\\
7 & 6 & 2 & 1 & 2 & 0 & 512 & 0 & \texttt{---------}\\
8 & 6 & 2 & 1 & 2 & 0 & 512 & 0 & \texttt{---------}\\
\multicolumn{9}{@{}l@{}}{\rule{0pt}{2.6ex}\textbf{T9.094}\enspace \texttt{012345121}\enspace $n{=}9$, $|H|{=}3$, len 6, br 1, $|\mathrm{Aut}|{=}1$, orb 9/9, hist-orb 3/3; M 9728, res 2592, rev 3336: ok}\\*
0 & -- & 0 & 3 & 1 & 186 & 2 & 324 & \texttt{-????????}\\*
1 & 0 & 1 & 1 & 1 & 480 & 32 & 0 & \texttt{--++++---}\\
2 & 1 & 2 & 1 & 1 & 448 & 64 & 0 & \texttt{---+++---}\\
3 & 2 & 3 & 1 & 1 & 384 & 128 & 0 & \texttt{----++---}\\
4 & 3 & 4 & 1 & 1 & 256 & 256 & 0 & \texttt{-----+---}\\
5 & 4 & 5 & 1 & 1 & 0 & 512 & 0 & \texttt{---------}\\
6 & 0 & 1 & 1 & 1 & 256 & 256 & 0 & \texttt{-------+-}\\
7 & 6 & 2 & 1 & 1 & 0 & 512 & 0 & \texttt{---------}\\
8 & 0 & 1 & 1 & 1 & 0 & 512 & 0 & \texttt{---------}\\
\multicolumn{9}{@{}l@{}}{\rule{0pt}{2.6ex}\textbf{T9.095}\enspace \texttt{012345111}\enspace $n{=}9$, $|H|{=}4$, len 6, br 1, $|\mathrm{Aut}|{=}6$, orb 7/6, hist-orb 2/1; M 9216, res 3584, rev 2604: ok}\\*
0 & -- & 0 & 4 & 1 & 62 & 2 & 448 & \texttt{-????????}\\*
1 & 0 & 1 & 1 & 1 & 480 & 32 & 0 & \texttt{--++++---}\\
2 & 1 & 2 & 1 & 1 & 448 & 64 & 0 & \texttt{---+++---}\\
3 & 2 & 3 & 1 & 1 & 384 & 128 & 0 & \texttt{----++---}\\
4 & 3 & 4 & 1 & 1 & 256 & 256 & 0 & \texttt{-----+---}\\
5 & 4 & 5 & 1 & 1 & 0 & 512 & 0 & \texttt{---------}\\
6 & 0 & 1 & 1 & 3 & 0 & 512 & 0 & \texttt{---------}\\
7 & 0 & 1 & 1 & 3 & 0 & 512 & 0 & \texttt{---------}\\
8 & 0 & 1 & 1 & 3 & 0 & 512 & 0 & \texttt{---------}\\
\multicolumn{9}{@{}l@{}}{\rule{0pt}{2.6ex}\textbf{T9.096}\enspace \texttt{012344444}\enspace $n{=}9$, $|H|{=}5$, len 5, br 1, $|\mathrm{Aut}|{=}120$, orb 5/4, hist-orb 1/0; M 13312, res 4500, rev 5636: ok}\\*
0 & -- & 0 & 5 & 1 & 450 & 2 & 60 & \texttt{-+++?????}\\*
1 & 0 & 1 & 5 & 1 & 388 & 4 & 120 & \texttt{--++?????}\\
2 & 1 & 2 & 5 & 1 & 264 & 8 & 240 & \texttt{---+?????}\\
3 & 2 & 3 & 5 & 1 & 16 & 16 & 480 & \texttt{----?????}\\
4 & 3 & 4 & 1 & 5 & 0 & 512 & 0 & \texttt{---------}\\
5 & 3 & 4 & 1 & 5 & 0 & 512 & 0 & \texttt{---------}\\
6 & 3 & 4 & 1 & 5 & 0 & 512 & 0 & \texttt{---------}\\
7 & 3 & 4 & 1 & 5 & 0 & 512 & 0 & \texttt{---------}\\
8 & 3 & 4 & 1 & 5 & 0 & 512 & 0 & \texttt{---------}\\
\multicolumn{9}{@{}l@{}}{\rule{0pt}{2.6ex}\textbf{T9.097}\enspace \texttt{012344443}\enspace $n{=}9$, $|H|{=}5$, len 5, br 2, $|\mathrm{Aut}|{=}24$, orb 6/5, hist-orb 2/1; M 12800, res 5264, rev 4584: ok}\\*
0 & -- & 0 & 5 & 1 & 418 & 2 & 92 & \texttt{-++??????}\\*
1 & 0 & 1 & 5 & 1 & 324 & 4 & 184 & \texttt{--+??????}\\
2 & 1 & 2 & 5 & 1 & 136 & 8 & 368 & \texttt{---??????}\\
3 & 2 & 3 & 4 & 1 & 32 & 32 & 448 & \texttt{----????-}\\
4 & 3 & 4 & 1 & 4 & 0 & 512 & 0 & \texttt{---------}\\
5 & 3 & 4 & 1 & 4 & 0 & 512 & 0 & \texttt{---------}\\
6 & 3 & 4 & 1 & 4 & 0 & 512 & 0 & \texttt{---------}\\
7 & 3 & 4 & 1 & 4 & 0 & 512 & 0 & \texttt{---------}\\
8 & 2 & 3 & 1 & 1 & 0 & 512 & 0 & \texttt{---------}\\
\multicolumn{9}{@{}l@{}}{\rule{0pt}{2.6ex}\textbf{T9.098}\enspace \texttt{012344442}\enspace $n{=}9$, $|H|{=}5$, len 5, br 2, $|\mathrm{Aut}|{=}24$, orb 6/5, hist-orb 2/1; M 12288, res 5628, rev 4032: ok}\\*
0 & -- & 0 & 5 & 1 & 354 & 2 & 156 & \texttt{-+???????}\\*
1 & 0 & 1 & 5 & 1 & 196 & 4 & 312 & \texttt{--???????}\\
2 & 1 & 2 & 4 & 1 & 272 & 16 & 224 & \texttt{---+????-}\\
3 & 2 & 3 & 4 & 1 & 32 & 32 & 448 & \texttt{----????-}\\
4 & 3 & 4 & 1 & 4 & 0 & 512 & 0 & \texttt{---------}\\
5 & 3 & 4 & 1 & 4 & 0 & 512 & 0 & \texttt{---------}\\
6 & 3 & 4 & 1 & 4 & 0 & 512 & 0 & \texttt{---------}\\
7 & 3 & 4 & 1 & 4 & 0 & 512 & 0 & \texttt{---------}\\
8 & 1 & 2 & 1 & 1 & 0 & 512 & 0 & \texttt{---------}\\
\multicolumn{9}{@{}l@{}}{\rule{0pt}{2.6ex}\textbf{T9.099}\enspace \texttt{012344441}\enspace $n{=}9$, $|H|{=}5$, len 5, br 2, $|\mathrm{Aut}|{=}24$, orb 6/5, hist-orb 2/1; M 11776, res 5240, rev 3984: ok}\\*
0 & -- & 0 & 5 & 1 & 226 & 2 & 284 & \texttt{-????????}\\*
1 & 0 & 1 & 4 & 1 & 392 & 8 & 112 & \texttt{--++????-}\\
2 & 1 & 2 & 4 & 1 & 272 & 16 & 224 & \texttt{---+????-}\\
3 & 2 & 3 & 4 & 1 & 32 & 32 & 448 & \texttt{----????-}\\
4 & 3 & 4 & 1 & 4 & 0 & 512 & 0 & \texttt{---------}\\
5 & 3 & 4 & 1 & 4 & 0 & 512 & 0 & \texttt{---------}\\
6 & 3 & 4 & 1 & 4 & 0 & 512 & 0 & \texttt{---------}\\
7 & 3 & 4 & 1 & 4 & 0 & 512 & 0 & \texttt{---------}\\
8 & 0 & 1 & 1 & 1 & 0 & 512 & 0 & \texttt{---------}\\
\multicolumn{9}{@{}l@{}}{\rule{0pt}{2.6ex}\textbf{T9.100}\enspace \texttt{012344434}\enspace $n{=}9$, $|H|{=}4$, len 5, br 2, $|\mathrm{Aut}|{=}6$, orb 7/6, hist-orb 2/1; M 12800, res 3840, rev 5100: ok}\\*
0 & -- & 0 & 4 & 1 & 438 & 2 & 72 & \texttt{-++??????}\\*
1 & 0 & 1 & 4 & 1 & 364 & 4 & 144 & \texttt{--+??????}\\
2 & 1 & 2 & 4 & 1 & 216 & 8 & 288 & \texttt{---??????}\\
3 & 2 & 3 & 3 & 1 & 64 & 64 & 384 & \texttt{----???--}\\
4 & 3 & 4 & 1 & 3 & 0 & 512 & 0 & \texttt{---------}\\
5 & 3 & 4 & 1 & 3 & 0 & 512 & 0 & \texttt{---------}\\
6 & 3 & 4 & 1 & 3 & 0 & 512 & 0 & \texttt{---------}\\
7 & 2 & 3 & 1 & 1 & 256 & 256 & 0 & \texttt{--------+}\\
8 & 7 & 4 & 1 & 1 & 0 & 512 & 0 & \texttt{---------}\\
\multicolumn{9}{@{}l@{}}{\rule{0pt}{2.6ex}\textbf{T9.101}\enspace \texttt{012344433}\enspace $n{=}9$, $|H|{=}5$, len 5, br 2, $|\mathrm{Aut}|{=}12$, orb 6/4, hist-orb 2/0; M 12288, res 5352, rev 4128: ok}\\*
0 & -- & 0 & 5 & 1 & 402 & 2 & 108 & \texttt{-++??????}\\*
1 & 0 & 1 & 5 & 1 & 292 & 4 & 216 & \texttt{--+??????}\\
2 & 1 & 2 & 5 & 1 & 72 & 8 & 432 & \texttt{---??????}\\
3 & 2 & 3 & 3 & 1 & 64 & 64 & 384 & \texttt{----???--}\\
4 & 3 & 4 & 1 & 3 & 0 & 512 & 0 & \texttt{---------}\\
5 & 3 & 4 & 1 & 3 & 0 & 512 & 0 & \texttt{---------}\\
6 & 3 & 4 & 1 & 3 & 0 & 512 & 0 & \texttt{---------}\\
7 & 2 & 3 & 1 & 2 & 0 & 512 & 0 & \texttt{---------}\\
8 & 2 & 3 & 1 & 2 & 0 & 512 & 0 & \texttt{---------}\\
\multicolumn{9}{@{}l@{}}{\rule{0pt}{2.6ex}\textbf{T9.102}\enspace \texttt{012344432}\enspace $n{=}9$, $|H|{=}5$, len 5, br 3, $|\mathrm{Aut}|{=}6$, orb 7/6, hist-orb 3/2; M 11776, res 5924, rev 3372: ok}\\*
0 & -- & 0 & 5 & 1 & 338 & 2 & 172 & \texttt{-+???????}\\*
1 & 0 & 1 & 5 & 1 & 164 & 4 & 344 & \texttt{--???????}\\
2 & 1 & 2 & 4 & 1 & 144 & 16 & 352 & \texttt{---?????-}\\
3 & 2 & 3 & 3 & 1 & 64 & 64 & 384 & \texttt{----???--}\\
4 & 3 & 4 & 1 & 3 & 0 & 512 & 0 & \texttt{---------}\\
5 & 3 & 4 & 1 & 3 & 0 & 512 & 0 & \texttt{---------}\\
6 & 3 & 4 & 1 & 3 & 0 & 512 & 0 & \texttt{---------}\\
7 & 2 & 3 & 1 & 1 & 0 & 512 & 0 & \texttt{---------}\\
8 & 1 & 2 & 1 & 1 & 0 & 512 & 0 & \texttt{---------}\\
\multicolumn{9}{@{}l@{}}{\rule{0pt}{2.6ex}\textbf{T9.103}\enspace \texttt{012344431}\enspace $n{=}9$, $|H|{=}5$, len 5, br 3, $|\mathrm{Aut}|{=}6$, orb 7/6, hist-orb 3/2; M 11264, res 5680, rev 3228: ok}\\*
0 & -- & 0 & 5 & 1 & 210 & 2 & 300 & \texttt{-????????}\\*
1 & 0 & 1 & 4 & 1 & 328 & 8 & 176 & \texttt{--+?????-}\\
2 & 1 & 2 & 4 & 1 & 144 & 16 & 352 & \texttt{---?????-}\\
3 & 2 & 3 & 3 & 1 & 64 & 64 & 384 & \texttt{----???--}\\
4 & 3 & 4 & 1 & 3 & 0 & 512 & 0 & \texttt{---------}\\
5 & 3 & 4 & 1 & 3 & 0 & 512 & 0 & \texttt{---------}\\
6 & 3 & 4 & 1 & 3 & 0 & 512 & 0 & \texttt{---------}\\
7 & 2 & 3 & 1 & 1 & 0 & 512 & 0 & \texttt{---------}\\
8 & 0 & 1 & 1 & 1 & 0 & 512 & 0 & \texttt{---------}\\
\multicolumn{9}{@{}l@{}}{\rule{0pt}{2.6ex}\textbf{T9.104}\enspace \texttt{012344423}\enspace $n{=}9$, $|H|{=}4$, len 5, br 2, $|\mathrm{Aut}|{=}6$, orb 7/6, hist-orb 2/1; M 11776, res 3624, rev 4568: ok}\\*
0 & -- & 0 & 4 & 1 & 406 & 2 & 104 & \texttt{-+???????}\\*
1 & 0 & 1 & 4 & 1 & 300 & 4 & 208 & \texttt{--???????}\\
2 & 1 & 2 & 3 & 1 & 288 & 32 & 192 & \texttt{---+???--}\\
3 & 2 & 3 & 3 & 1 & 64 & 64 & 384 & \texttt{----???--}\\
4 & 3 & 4 & 1 & 3 & 0 & 512 & 0 & \texttt{---------}\\
5 & 3 & 4 & 1 & 3 & 0 & 512 & 0 & \texttt{---------}\\
6 & 3 & 4 & 1 & 3 & 0 & 512 & 0 & \texttt{---------}\\
7 & 1 & 2 & 1 & 1 & 256 & 256 & 0 & \texttt{--------+}\\
8 & 7 & 3 & 1 & 1 & 0 & 512 & 0 & \texttt{---------}\\
\multicolumn{9}{@{}l@{}}{\rule{0pt}{2.6ex}\textbf{T9.105}\enspace \texttt{012344422}\enspace $n{=}9$, $|H|{=}5$, len 5, br 2, $|\mathrm{Aut}|{=}12$, orb 6/4, hist-orb 2/0; M 11264, res 5724, rev 3192: ok}\\*
0 & -- & 0 & 5 & 1 & 306 & 2 & 204 & \texttt{-+???????}\\*
1 & 0 & 1 & 5 & 1 & 100 & 4 & 408 & \texttt{--???????}\\
2 & 1 & 2 & 3 & 1 & 288 & 32 & 192 & \texttt{---+???--}\\
3 & 2 & 3 & 3 & 1 & 64 & 64 & 384 & \texttt{----???--}\\
4 & 3 & 4 & 1 & 3 & 0 & 512 & 0 & \texttt{---------}\\
5 & 3 & 4 & 1 & 3 & 0 & 512 & 0 & \texttt{---------}\\
6 & 3 & 4 & 1 & 3 & 0 & 512 & 0 & \texttt{---------}\\
7 & 1 & 2 & 1 & 2 & 0 & 512 & 0 & \texttt{---------}\\
8 & 1 & 2 & 1 & 2 & 0 & 512 & 0 & \texttt{---------}\\
\multicolumn{9}{@{}l@{}}{\rule{0pt}{2.6ex}\textbf{T9.106}\enspace \texttt{012344421}\enspace $n{=}9$, $|H|{=}5$, len 5, br 3, $|\mathrm{Aut}|{=}6$, orb 7/6, hist-orb 3/2; M 10752, res 5768, rev 2844: ok}\\*
0 & -- & 0 & 5 & 1 & 178 & 2 & 332 & \texttt{-????????}\\*
1 & 0 & 1 & 4 & 1 & 200 & 8 & 304 & \texttt{--??????-}\\
2 & 1 & 2 & 3 & 1 & 288 & 32 & 192 & \texttt{---+???--}\\
3 & 2 & 3 & 3 & 1 & 64 & 64 & 384 & \texttt{----???--}\\
4 & 3 & 4 & 1 & 3 & 0 & 512 & 0 & \texttt{---------}\\
5 & 3 & 4 & 1 & 3 & 0 & 512 & 0 & \texttt{---------}\\
6 & 3 & 4 & 1 & 3 & 0 & 512 & 0 & \texttt{---------}\\
7 & 1 & 2 & 1 & 1 & 0 & 512 & 0 & \texttt{---------}\\
8 & 0 & 1 & 1 & 1 & 0 & 512 & 0 & \texttt{---------}\\
\multicolumn{9}{@{}l@{}}{\rule{0pt}{2.6ex}\textbf{T9.107}\enspace \texttt{012344412}\enspace $n{=}9$, $|H|{=}4$, len 5, br 2, $|\mathrm{Aut}|{=}6$, orb 7/6, hist-orb 2/1; M 10752, res 3216, rev 4156: ok}\\*
0 & -- & 0 & 4 & 1 & 342 & 2 & 168 & \texttt{-????????}\\*
1 & 0 & 1 & 3 & 1 & 400 & 16 & 96 & \texttt{--++???--}\\
2 & 1 & 2 & 3 & 1 & 288 & 32 & 192 & \texttt{---+???--}\\
3 & 2 & 3 & 3 & 1 & 64 & 64 & 384 & \texttt{----???--}\\
4 & 3 & 4 & 1 & 3 & 0 & 512 & 0 & \texttt{---------}\\
5 & 3 & 4 & 1 & 3 & 0 & 512 & 0 & \texttt{---------}\\
6 & 3 & 4 & 1 & 3 & 0 & 512 & 0 & \texttt{---------}\\
7 & 0 & 1 & 1 & 1 & 256 & 256 & 0 & \texttt{--------+}\\
8 & 7 & 2 & 1 & 1 & 0 & 512 & 0 & \texttt{---------}\\
\multicolumn{9}{@{}l@{}}{\rule{0pt}{2.6ex}\textbf{T9.108}\enspace \texttt{012344411}\enspace $n{=}9$, $|H|{=}5$, len 5, br 2, $|\mathrm{Aut}|{=}12$, orb 6/4, hist-orb 2/0; M 10240, res 5040, rev 2944: ok}\\*
0 & -- & 0 & 5 & 1 & 114 & 2 & 396 & \texttt{-????????}\\*
1 & 0 & 1 & 3 & 1 & 400 & 16 & 96 & \texttt{--++???--}\\
2 & 1 & 2 & 3 & 1 & 288 & 32 & 192 & \texttt{---+???--}\\
3 & 2 & 3 & 3 & 1 & 64 & 64 & 384 & \texttt{----???--}\\
4 & 3 & 4 & 1 & 3 & 0 & 512 & 0 & \texttt{---------}\\
5 & 3 & 4 & 1 & 3 & 0 & 512 & 0 & \texttt{---------}\\
6 & 3 & 4 & 1 & 3 & 0 & 512 & 0 & \texttt{---------}\\
7 & 0 & 1 & 1 & 2 & 0 & 512 & 0 & \texttt{---------}\\
8 & 0 & 1 & 1 & 2 & 0 & 512 & 0 & \texttt{---------}\\
\multicolumn{9}{@{}l@{}}{\rule{0pt}{2.6ex}\textbf{T9.109}\enspace \texttt{012344344}\enspace $n{=}9$, $|H|{=}4$, len 5, br 3, $|\mathrm{Aut}|{=}8$, orb 5/3, hist-orb 1/0; M 12800, res 3768, rev 4944: ok}\\*
0 & -- & 0 & 4 & 1 & 434 & 2 & 76 & \texttt{-++??????}\\*
1 & 0 & 1 & 4 & 1 & 356 & 4 & 152 & \texttt{--+??????}\\
2 & 1 & 2 & 4 & 1 & 200 & 8 & 304 & \texttt{---??????}\\
3 & 2 & 3 & 2 & 2 & 128 & 128 & 256 & \texttt{----??---}\\
4 & 3 & 4 & 1 & 4 & 0 & 512 & 0 & \texttt{---------}\\
5 & 3 & 4 & 1 & 4 & 0 & 512 & 0 & \texttt{---------}\\
6 & 2 & 3 & 2 & 2 & 128 & 128 & 256 & \texttt{-------??}\\
7 & 6 & 4 & 1 & 4 & 0 & 512 & 0 & \texttt{---------}\\
8 & 6 & 4 & 1 & 4 & 0 & 512 & 0 & \texttt{---------}\\
\multicolumn{9}{@{}l@{}}{\rule{0pt}{2.6ex}\textbf{T9.110}\enspace \texttt{012344343}\enspace $n{=}9$, $|H|{=}4$, len 5, br 2, $|\mathrm{Aut}|{=}2$, orb 8/7, hist-orb 3/2; M 12288, res 4320, rev 4516: ok}\\*
0 & -- & 0 & 4 & 1 & 414 & 2 & 96 & \texttt{-++??????}\\*
1 & 0 & 1 & 4 & 1 & 316 & 4 & 192 & \texttt{--+??????}\\
2 & 1 & 2 & 4 & 1 & 120 & 8 & 384 & \texttt{---??????}\\
3 & 2 & 3 & 2 & 1 & 128 & 128 & 256 & \texttt{----??---}\\
4 & 3 & 4 & 1 & 2 & 0 & 512 & 0 & \texttt{---------}\\
5 & 3 & 4 & 1 & 2 & 0 & 512 & 0 & \texttt{---------}\\
6 & 2 & 3 & 1 & 1 & 256 & 256 & 0 & \texttt{-------+-}\\
7 & 6 & 4 & 1 & 1 & 0 & 512 & 0 & \texttt{---------}\\
8 & 2 & 3 & 1 & 1 & 0 & 512 & 0 & \texttt{---------}\\
\multicolumn{9}{@{}l@{}}{\rule{0pt}{2.6ex}\textbf{T9.111}\enspace \texttt{012344342}\enspace $n{=}9$, $|H|{=}4$, len 5, br 3, $|\mathrm{Aut}|{=}2$, orb 8/7, hist-orb 3/2; M 11776, res 4736, rev 3912: ok}\\*
0 & -- & 0 & 4 & 1 & 350 & 2 & 160 & \texttt{-+???????}\\*
1 & 0 & 1 & 4 & 1 & 188 & 4 & 320 & \texttt{--???????}\\
2 & 1 & 2 & 3 & 1 & 240 & 16 & 256 & \texttt{---?????-}\\
3 & 2 & 3 & 2 & 1 & 128 & 128 & 256 & \texttt{----??---}\\
4 & 3 & 4 & 1 & 2 & 0 & 512 & 0 & \texttt{---------}\\
5 & 3 & 4 & 1 & 2 & 0 & 512 & 0 & \texttt{---------}\\
6 & 2 & 3 & 1 & 1 & 256 & 256 & 0 & \texttt{-------+-}\\
7 & 6 & 4 & 1 & 1 & 0 & 512 & 0 & \texttt{---------}\\
8 & 1 & 2 & 1 & 1 & 0 & 512 & 0 & \texttt{---------}\\
\multicolumn{9}{@{}l@{}}{\rule{0pt}{2.6ex}\textbf{T9.112}\enspace \texttt{012344341}\enspace $n{=}9$, $|H|{=}4$, len 5, br 3, $|\mathrm{Aut}|{=}2$, orb 8/7, hist-orb 3/2; M 11264, res 4384, rev 3808: ok}\\*
0 & -- & 0 & 4 & 1 & 222 & 2 & 288 & \texttt{-????????}\\*
1 & 0 & 1 & 3 & 1 & 376 & 8 & 128 & \texttt{--+?????-}\\
2 & 1 & 2 & 3 & 1 & 240 & 16 & 256 & \texttt{---?????-}\\
3 & 2 & 3 & 2 & 1 & 128 & 128 & 256 & \texttt{----??---}\\
4 & 3 & 4 & 1 & 2 & 0 & 512 & 0 & \texttt{---------}\\
5 & 3 & 4 & 1 & 2 & 0 & 512 & 0 & \texttt{---------}\\
6 & 2 & 3 & 1 & 1 & 256 & 256 & 0 & \texttt{-------+-}\\
7 & 6 & 4 & 1 & 1 & 0 & 512 & 0 & \texttt{---------}\\
8 & 0 & 1 & 1 & 1 & 0 & 512 & 0 & \texttt{---------}\\
\multicolumn{9}{@{}l@{}}{\rule{0pt}{2.6ex}\textbf{T9.113}\enspace \texttt{012344333}\enspace $n{=}9$, $|H|{=}5$, len 5, br 2, $|\mathrm{Aut}|{=}12$, orb 6/4, hist-orb 2/0; M 11776, res 5160, rev 3992: ok}\\*
0 & -- & 0 & 5 & 1 & 394 & 2 & 116 & \texttt{-++??????}\\*
1 & 0 & 1 & 5 & 1 & 276 & 4 & 232 & \texttt{--+??????}\\
2 & 1 & 2 & 5 & 1 & 40 & 8 & 464 & \texttt{---??????}\\
3 & 2 & 3 & 2 & 1 & 128 & 128 & 256 & \texttt{----??---}\\
4 & 3 & 4 & 1 & 2 & 0 & 512 & 0 & \texttt{---------}\\
5 & 3 & 4 & 1 & 2 & 0 & 512 & 0 & \texttt{---------}\\
6 & 2 & 3 & 1 & 3 & 0 & 512 & 0 & \texttt{---------}\\
7 & 2 & 3 & 1 & 3 & 0 & 512 & 0 & \texttt{---------}\\
8 & 2 & 3 & 1 & 3 & 0 & 512 & 0 & \texttt{---------}\\
\multicolumn{9}{@{}l@{}}{\rule{0pt}{2.6ex}\textbf{T9.114}\enspace \texttt{012344332}\enspace $n{=}9$, $|H|{=}5$, len 5, br 3, $|\mathrm{Aut}|{=}4$, orb 7/5, hist-orb 3/1; M 11264, res 5836, rev 3116: ok}\\*
0 & -- & 0 & 5 & 1 & 330 & 2 & 180 & \texttt{-+???????}\\*
1 & 0 & 1 & 5 & 1 & 148 & 4 & 360 & \texttt{--???????}\\
2 & 1 & 2 & 4 & 1 & 80 & 16 & 416 & \texttt{---?????-}\\
3 & 2 & 3 & 2 & 1 & 128 & 128 & 256 & \texttt{----??---}\\
4 & 3 & 4 & 1 & 2 & 0 & 512 & 0 & \texttt{---------}\\
5 & 3 & 4 & 1 & 2 & 0 & 512 & 0 & \texttt{---------}\\
6 & 2 & 3 & 1 & 2 & 0 & 512 & 0 & \texttt{---------}\\
7 & 2 & 3 & 1 & 2 & 0 & 512 & 0 & \texttt{---------}\\
8 & 1 & 2 & 1 & 1 & 0 & 512 & 0 & \texttt{---------}\\
\multicolumn{9}{@{}l@{}}{\rule{0pt}{2.6ex}\textbf{T9.115}\enspace \texttt{012344331}\enspace $n{=}9$, $|H|{=}5$, len 5, br 3, $|\mathrm{Aut}|{=}4$, orb 7/5, hist-orb 3/1; M 10752, res 5664, rev 2924: ok}\\*
0 & -- & 0 & 5 & 1 & 202 & 2 & 308 & \texttt{-????????}\\*
1 & 0 & 1 & 4 & 1 & 296 & 8 & 208 & \texttt{--+?????-}\\
2 & 1 & 2 & 4 & 1 & 80 & 16 & 416 & \texttt{---?????-}\\
3 & 2 & 3 & 2 & 1 & 128 & 128 & 256 & \texttt{----??---}\\
4 & 3 & 4 & 1 & 2 & 0 & 512 & 0 & \texttt{---------}\\
5 & 3 & 4 & 1 & 2 & 0 & 512 & 0 & \texttt{---------}\\
6 & 2 & 3 & 1 & 2 & 0 & 512 & 0 & \texttt{---------}\\
7 & 2 & 3 & 1 & 2 & 0 & 512 & 0 & \texttt{---------}\\
8 & 0 & 1 & 1 & 1 & 0 & 512 & 0 & \texttt{---------}\\
\multicolumn{9}{@{}l@{}}{\rule{0pt}{2.6ex}\textbf{T9.116}\enspace \texttt{012344323}\enspace $n{=}9$, $|H|{=}4$, len 5, br 3, $|\mathrm{Aut}|{=}2$, orb 8/7, hist-orb 3/2; M 11264, res 4016, rev 4004: ok}\\*
0 & -- & 0 & 4 & 1 & 382 & 2 & 128 & \texttt{-+???????}\\*
1 & 0 & 1 & 4 & 1 & 252 & 4 & 256 & \texttt{--???????}\\
2 & 1 & 2 & 3 & 1 & 160 & 32 & 320 & \texttt{---????--}\\
3 & 2 & 3 & 2 & 1 & 128 & 128 & 256 & \texttt{----??---}\\
4 & 3 & 4 & 1 & 2 & 0 & 512 & 0 & \texttt{---------}\\
5 & 3 & 4 & 1 & 2 & 0 & 512 & 0 & \texttt{---------}\\
6 & 2 & 3 & 1 & 1 & 0 & 512 & 0 & \texttt{---------}\\
7 & 1 & 2 & 1 & 1 & 256 & 256 & 0 & \texttt{--------+}\\
8 & 7 & 3 & 1 & 1 & 0 & 512 & 0 & \texttt{---------}\\
\multicolumn{9}{@{}l@{}}{\rule{0pt}{2.6ex}\textbf{T9.117}\enspace \texttt{012344322}\enspace $n{=}9$, $|H|{=}5$, len 5, br 3, $|\mathrm{Aut}|{=}4$, orb 7/5, hist-orb 3/1; M 10752, res 5780, rev 2796: ok}\\*
0 & -- & 0 & 5 & 1 & 298 & 2 & 212 & \texttt{-+???????}\\*
1 & 0 & 1 & 5 & 1 & 84 & 4 & 424 & \texttt{--???????}\\
2 & 1 & 2 & 3 & 1 & 160 & 32 & 320 & \texttt{---????--}\\
3 & 2 & 3 & 2 & 1 & 128 & 128 & 256 & \texttt{----??---}\\
4 & 3 & 4 & 1 & 2 & 0 & 512 & 0 & \texttt{---------}\\
5 & 3 & 4 & 1 & 2 & 0 & 512 & 0 & \texttt{---------}\\
6 & 2 & 3 & 1 & 1 & 0 & 512 & 0 & \texttt{---------}\\
7 & 1 & 2 & 1 & 2 & 0 & 512 & 0 & \texttt{---------}\\
8 & 1 & 2 & 1 & 2 & 0 & 512 & 0 & \texttt{---------}\\
\multicolumn{9}{@{}l@{}}{\rule{0pt}{2.6ex}\textbf{T9.118}\enspace \texttt{012344321}\enspace $n{=}9$, $|H|{=}5$, len 5, br 4, $|\mathrm{Aut}|{=}2$, orb 8/7, hist-orb 4/3; M 10240, res 5896, rev 2392: ok}\\*
0 & -- & 0 & 5 & 1 & 170 & 2 & 340 & \texttt{-????????}\\*
1 & 0 & 1 & 4 & 1 & 168 & 8 & 336 & \texttt{--??????-}\\
2 & 1 & 2 & 3 & 1 & 160 & 32 & 320 & \texttt{---????--}\\
3 & 2 & 3 & 2 & 1 & 128 & 128 & 256 & \texttt{----??---}\\
4 & 3 & 4 & 1 & 2 & 0 & 512 & 0 & \texttt{---------}\\
5 & 3 & 4 & 1 & 2 & 0 & 512 & 0 & \texttt{---------}\\
6 & 2 & 3 & 1 & 1 & 0 & 512 & 0 & \texttt{---------}\\
7 & 1 & 2 & 1 & 1 & 0 & 512 & 0 & \texttt{---------}\\
8 & 0 & 1 & 1 & 1 & 0 & 512 & 0 & \texttt{---------}\\
\multicolumn{9}{@{}l@{}}{\rule{0pt}{2.6ex}\textbf{T9.119}\enspace \texttt{012344312}\enspace $n{=}9$, $|H|{=}4$, len 5, br 3, $|\mathrm{Aut}|{=}2$, orb 8/7, hist-orb 3/2; M 10240, res 3584, rev 3608: ok}\\*
0 & -- & 0 & 4 & 1 & 318 & 2 & 192 & \texttt{-????????}\\*
1 & 0 & 1 & 3 & 1 & 336 & 16 & 160 & \texttt{--+????--}\\
2 & 1 & 2 & 3 & 1 & 160 & 32 & 320 & \texttt{---????--}\\
3 & 2 & 3 & 2 & 1 & 128 & 128 & 256 & \texttt{----??---}\\
4 & 3 & 4 & 1 & 2 & 0 & 512 & 0 & \texttt{---------}\\
5 & 3 & 4 & 1 & 2 & 0 & 512 & 0 & \texttt{---------}\\
6 & 2 & 3 & 1 & 1 & 0 & 512 & 0 & \texttt{---------}\\
7 & 0 & 1 & 1 & 1 & 256 & 256 & 0 & \texttt{--------+}\\
8 & 7 & 2 & 1 & 1 & 0 & 512 & 0 & \texttt{---------}\\
\multicolumn{9}{@{}l@{}}{\rule{0pt}{2.6ex}\textbf{T9.120}\enspace \texttt{012344311}\enspace $n{=}9$, $|H|{=}5$, len 5, br 3, $|\mathrm{Aut}|{=}4$, orb 7/5, hist-orb 3/1; M 9728, res 5280, rev 2436: ok}\\*
0 & -- & 0 & 5 & 1 & 106 & 2 & 404 & \texttt{-????????}\\*
1 & 0 & 1 & 3 & 1 & 336 & 16 & 160 & \texttt{--+????--}\\
2 & 1 & 2 & 3 & 1 & 160 & 32 & 320 & \texttt{---????--}\\
3 & 2 & 3 & 2 & 1 & 128 & 128 & 256 & \texttt{----??---}\\
4 & 3 & 4 & 1 & 2 & 0 & 512 & 0 & \texttt{---------}\\
5 & 3 & 4 & 1 & 2 & 0 & 512 & 0 & \texttt{---------}\\
6 & 2 & 3 & 1 & 1 & 0 & 512 & 0 & \texttt{---------}\\
7 & 0 & 1 & 1 & 2 & 0 & 512 & 0 & \texttt{---------}\\
8 & 0 & 1 & 1 & 2 & 0 & 512 & 0 & \texttt{---------}\\
\multicolumn{9}{@{}l@{}}{\rule{0pt}{2.6ex}\textbf{T9.121}\enspace \texttt{012344234}\enspace $n{=}9$, $|H|{=}3$, len 5, br 2, $|\mathrm{Aut}|{=}2$, orb 8/7, hist-orb 2/1; M 11776, res 2088, rev 4904: ok}\\*
0 & -- & 0 & 3 & 1 & 438 & 2 & 72 & \texttt{-+???????}\\*
1 & 0 & 1 & 3 & 1 & 364 & 4 & 144 & \texttt{--???????}\\
2 & 1 & 2 & 2 & 1 & 320 & 64 & 128 & \texttt{---+??---}\\
3 & 2 & 3 & 2 & 1 & 128 & 128 & 256 & \texttt{----??---}\\
4 & 3 & 4 & 1 & 2 & 0 & 512 & 0 & \texttt{---------}\\
5 & 3 & 4 & 1 & 2 & 0 & 512 & 0 & \texttt{---------}\\
6 & 1 & 2 & 1 & 1 & 384 & 128 & 0 & \texttt{-------++}\\
7 & 6 & 3 & 1 & 1 & 256 & 256 & 0 & \texttt{--------+}\\
8 & 7 & 4 & 1 & 1 & 0 & 512 & 0 & \texttt{---------}\\
\multicolumn{9}{@{}l@{}}{\rule{0pt}{2.6ex}\textbf{T9.122}\enspace \texttt{012344233}\enspace $n{=}9$, $|H|{=}4$, len 5, br 3, $|\mathrm{Aut}|{=}4$, orb 7/5, hist-orb 2/0; M 11264, res 3500, rev 4160: ok}\\*
0 & -- & 0 & 4 & 1 & 386 & 2 & 124 & \texttt{-+???????}\\*
1 & 0 & 1 & 4 & 1 & 260 & 4 & 248 & \texttt{--???????}\\
2 & 1 & 2 & 2 & 1 & 320 & 64 & 128 & \texttt{---+??---}\\
3 & 2 & 3 & 2 & 1 & 128 & 128 & 256 & \texttt{----??---}\\
4 & 3 & 4 & 1 & 2 & 0 & 512 & 0 & \texttt{---------}\\
5 & 3 & 4 & 1 & 2 & 0 & 512 & 0 & \texttt{---------}\\
6 & 1 & 2 & 2 & 1 & 128 & 128 & 256 & \texttt{-------??}\\
7 & 6 & 3 & 1 & 2 & 0 & 512 & 0 & \texttt{---------}\\
8 & 6 & 3 & 1 & 2 & 0 & 512 & 0 & \texttt{---------}\\
\multicolumn{9}{@{}l@{}}{\rule{0pt}{2.6ex}\textbf{T9.123}\enspace \texttt{012344232}\enspace $n{=}9$, $|H|{=}4$, len 5, br 2, $|\mathrm{Aut}|{=}2$, orb 8/7, hist-orb 3/2; M 10752, res 4272, rev 3552: ok}\\*
0 & -- & 0 & 4 & 1 & 334 & 2 & 176 & \texttt{-+???????}\\*
1 & 0 & 1 & 4 & 1 & 156 & 4 & 352 & \texttt{--???????}\\
2 & 1 & 2 & 2 & 1 & 320 & 64 & 128 & \texttt{---+??---}\\
3 & 2 & 3 & 2 & 1 & 128 & 128 & 256 & \texttt{----??---}\\
4 & 3 & 4 & 1 & 2 & 0 & 512 & 0 & \texttt{---------}\\
5 & 3 & 4 & 1 & 2 & 0 & 512 & 0 & \texttt{---------}\\
6 & 1 & 2 & 1 & 1 & 256 & 256 & 0 & \texttt{-------+-}\\
7 & 6 & 3 & 1 & 1 & 0 & 512 & 0 & \texttt{---------}\\
8 & 1 & 2 & 1 & 1 & 0 & 512 & 0 & \texttt{---------}\\
\multicolumn{9}{@{}l@{}}{\rule{0pt}{2.6ex}\textbf{T9.124}\enspace \texttt{012344231}\enspace $n{=}9$, $|H|{=}4$, len 5, br 3, $|\mathrm{Aut}|{=}2$, orb 8/7, hist-orb 3/2; M 10240, res 4064, rev 3340: ok}\\*
0 & -- & 0 & 4 & 1 & 206 & 2 & 304 & \texttt{-????????}\\*
1 & 0 & 1 & 3 & 1 & 312 & 8 & 192 & \texttt{--??????-}\\
2 & 1 & 2 & 2 & 1 & 320 & 64 & 128 & \texttt{---+??---}\\
3 & 2 & 3 & 2 & 1 & 128 & 128 & 256 & \texttt{----??---}\\
4 & 3 & 4 & 1 & 2 & 0 & 512 & 0 & \texttt{---------}\\
5 & 3 & 4 & 1 & 2 & 0 & 512 & 0 & \texttt{---------}\\
6 & 1 & 2 & 1 & 1 & 256 & 256 & 0 & \texttt{-------+-}\\
7 & 6 & 3 & 1 & 1 & 0 & 512 & 0 & \texttt{---------}\\
8 & 0 & 1 & 1 & 1 & 0 & 512 & 0 & \texttt{---------}\\
\multicolumn{9}{@{}l@{}}{\rule{0pt}{2.6ex}\textbf{T9.125}\enspace \texttt{012344222}\enspace $n{=}9$, $|H|{=}5$, len 5, br 2, $|\mathrm{Aut}|{=}12$, orb 6/4, hist-orb 2/0; M 10240, res 5364, rev 2752: ok}\\*
0 & -- & 0 & 5 & 1 & 282 & 2 & 228 & \texttt{-+???????}\\*
1 & 0 & 1 & 5 & 1 & 52 & 4 & 456 & \texttt{--???????}\\
2 & 1 & 2 & 2 & 1 & 320 & 64 & 128 & \texttt{---+??---}\\
3 & 2 & 3 & 2 & 1 & 128 & 128 & 256 & \texttt{----??---}\\
4 & 3 & 4 & 1 & 2 & 0 & 512 & 0 & \texttt{---------}\\
5 & 3 & 4 & 1 & 2 & 0 & 512 & 0 & \texttt{---------}\\
6 & 1 & 2 & 1 & 3 & 0 & 512 & 0 & \texttt{---------}\\
7 & 1 & 2 & 1 & 3 & 0 & 512 & 0 & \texttt{---------}\\
8 & 1 & 2 & 1 & 3 & 0 & 512 & 0 & \texttt{---------}\\
\multicolumn{9}{@{}l@{}}{\rule{0pt}{2.6ex}\textbf{T9.126}\enspace \texttt{012344221}\enspace $n{=}9$, $|H|{=}5$, len 5, br 3, $|\mathrm{Aut}|{=}4$, orb 7/5, hist-orb 3/1; M 9728, res 5624, rev 2236: ok}\\*
0 & -- & 0 & 5 & 1 & 154 & 2 & 356 & \texttt{-????????}\\*
1 & 0 & 1 & 4 & 1 & 104 & 8 & 400 & \texttt{--??????-}\\
2 & 1 & 2 & 2 & 1 & 320 & 64 & 128 & \texttt{---+??---}\\
3 & 2 & 3 & 2 & 1 & 128 & 128 & 256 & \texttt{----??---}\\
4 & 3 & 4 & 1 & 2 & 0 & 512 & 0 & \texttt{---------}\\
5 & 3 & 4 & 1 & 2 & 0 & 512 & 0 & \texttt{---------}\\
6 & 1 & 2 & 1 & 2 & 0 & 512 & 0 & \texttt{---------}\\
7 & 1 & 2 & 1 & 2 & 0 & 512 & 0 & \texttt{---------}\\
8 & 0 & 1 & 1 & 1 & 0 & 512 & 0 & \texttt{---------}\\
\multicolumn{9}{@{}l@{}}{\rule{0pt}{2.6ex}\textbf{T9.127}\enspace \texttt{012344212}\enspace $n{=}9$, $|H|{=}4$, len 5, br 3, $|\mathrm{Aut}|{=}2$, orb 8/7, hist-orb 3/2; M 9728, res 3792, rev 3176: ok}\\*
0 & -- & 0 & 4 & 1 & 270 & 2 & 240 & \texttt{-????????}\\*
1 & 0 & 1 & 3 & 1 & 208 & 16 & 288 & \texttt{--?????--}\\
2 & 1 & 2 & 2 & 1 & 320 & 64 & 128 & \texttt{---+??---}\\
3 & 2 & 3 & 2 & 1 & 128 & 128 & 256 & \texttt{----??---}\\
4 & 3 & 4 & 1 & 2 & 0 & 512 & 0 & \texttt{---------}\\
5 & 3 & 4 & 1 & 2 & 0 & 512 & 0 & \texttt{---------}\\
6 & 1 & 2 & 1 & 1 & 0 & 512 & 0 & \texttt{---------}\\
7 & 0 & 1 & 1 & 1 & 256 & 256 & 0 & \texttt{--------+}\\
8 & 7 & 2 & 1 & 1 & 0 & 512 & 0 & \texttt{---------}\\
\multicolumn{9}{@{}l@{}}{\rule{0pt}{2.6ex}\textbf{T9.128}\enspace \texttt{012344211}\enspace $n{=}9$, $|H|{=}5$, len 5, br 3, $|\mathrm{Aut}|{=}4$, orb 7/5, hist-orb 3/1; M 9216, res 5232, rev 2132: ok}\\*
0 & -- & 0 & 5 & 1 & 90 & 2 & 420 & \texttt{-????????}\\*
1 & 0 & 1 & 3 & 1 & 208 & 16 & 288 & \texttt{--?????--}\\
2 & 1 & 2 & 2 & 1 & 320 & 64 & 128 & \texttt{---+??---}\\
3 & 2 & 3 & 2 & 1 & 128 & 128 & 256 & \texttt{----??---}\\
4 & 3 & 4 & 1 & 2 & 0 & 512 & 0 & \texttt{---------}\\
5 & 3 & 4 & 1 & 2 & 0 & 512 & 0 & \texttt{---------}\\
6 & 1 & 2 & 1 & 1 & 0 & 512 & 0 & \texttt{---------}\\
7 & 0 & 1 & 1 & 2 & 0 & 512 & 0 & \texttt{---------}\\
8 & 0 & 1 & 1 & 2 & 0 & 512 & 0 & \texttt{---------}\\
\multicolumn{9}{@{}l@{}}{\rule{0pt}{2.6ex}\textbf{T9.129}\enspace \texttt{012344123}\enspace $n{=}9$, $|H|{=}3$, len 5, br 2, $|\mathrm{Aut}|{=}2$, orb 8/7, hist-orb 2/1; M 10240, res 1632, rev 4292: ok}\\*
0 & -- & 0 & 3 & 1 & 406 & 2 & 104 & \texttt{-????????}\\*
1 & 0 & 1 & 2 & 1 & 416 & 32 & 64 & \texttt{--++??---}\\
2 & 1 & 2 & 2 & 1 & 320 & 64 & 128 & \texttt{---+??---}\\
3 & 2 & 3 & 2 & 1 & 128 & 128 & 256 & \texttt{----??---}\\
4 & 3 & 4 & 1 & 2 & 0 & 512 & 0 & \texttt{---------}\\
5 & 3 & 4 & 1 & 2 & 0 & 512 & 0 & \texttt{---------}\\
6 & 0 & 1 & 1 & 1 & 384 & 128 & 0 & \texttt{-------++}\\
7 & 6 & 2 & 1 & 1 & 256 & 256 & 0 & \texttt{--------+}\\
8 & 7 & 3 & 1 & 1 & 0 & 512 & 0 & \texttt{---------}\\
\multicolumn{9}{@{}l@{}}{\rule{0pt}{2.6ex}\textbf{T9.130}\enspace \texttt{012344122}\enspace $n{=}9$, $|H|{=}4$, len 5, br 3, $|\mathrm{Aut}|{=}4$, orb 7/5, hist-orb 2/0; M 9728, res 2944, rev 3504: ok}\\*
0 & -- & 0 & 4 & 1 & 290 & 2 & 220 & \texttt{-????????}\\*
1 & 0 & 1 & 2 & 1 & 416 & 32 & 64 & \texttt{--++??---}\\
2 & 1 & 2 & 2 & 1 & 320 & 64 & 128 & \texttt{---+??---}\\
3 & 2 & 3 & 2 & 1 & 128 & 128 & 256 & \texttt{----??---}\\
4 & 3 & 4 & 1 & 2 & 0 & 512 & 0 & \texttt{---------}\\
5 & 3 & 4 & 1 & 2 & 0 & 512 & 0 & \texttt{---------}\\
6 & 0 & 1 & 2 & 1 & 128 & 128 & 256 & \texttt{-------??}\\
7 & 6 & 2 & 1 & 2 & 0 & 512 & 0 & \texttt{---------}\\
8 & 6 & 2 & 1 & 2 & 0 & 512 & 0 & \texttt{---------}\\
\multicolumn{9}{@{}l@{}}{\rule{0pt}{2.6ex}\textbf{T9.131}\enspace \texttt{012344121}\enspace $n{=}9$, $|H|{=}4$, len 5, br 2, $|\mathrm{Aut}|{=}2$, orb 8/7, hist-orb 3/2; M 9216, res 3488, rev 2996: ok}\\*
0 & -- & 0 & 4 & 1 & 174 & 2 & 336 & \texttt{-????????}\\*
1 & 0 & 1 & 2 & 1 & 416 & 32 & 64 & \texttt{--++??---}\\
2 & 1 & 2 & 2 & 1 & 320 & 64 & 128 & \texttt{---+??---}\\
3 & 2 & 3 & 2 & 1 & 128 & 128 & 256 & \texttt{----??---}\\
4 & 3 & 4 & 1 & 2 & 0 & 512 & 0 & \texttt{---------}\\
5 & 3 & 4 & 1 & 2 & 0 & 512 & 0 & \texttt{---------}\\
6 & 0 & 1 & 1 & 1 & 256 & 256 & 0 & \texttt{-------+-}\\
7 & 6 & 2 & 1 & 1 & 0 & 512 & 0 & \texttt{---------}\\
8 & 0 & 1 & 1 & 1 & 0 & 512 & 0 & \texttt{---------}\\
\multicolumn{9}{@{}l@{}}{\rule{0pt}{2.6ex}\textbf{T9.132}\enspace \texttt{012344111}\enspace $n{=}9$, $|H|{=}5$, len 5, br 2, $|\mathrm{Aut}|{=}12$, orb 6/4, hist-orb 2/0; M 8704, res 4416, rev 2232: ok}\\*
0 & -- & 0 & 5 & 1 & 58 & 2 & 452 & \texttt{-????????}\\*
1 & 0 & 1 & 2 & 1 & 416 & 32 & 64 & \texttt{--++??---}\\
2 & 1 & 2 & 2 & 1 & 320 & 64 & 128 & \texttt{---+??---}\\
3 & 2 & 3 & 2 & 1 & 128 & 128 & 256 & \texttt{----??---}\\
4 & 3 & 4 & 1 & 2 & 0 & 512 & 0 & \texttt{---------}\\
5 & 3 & 4 & 1 & 2 & 0 & 512 & 0 & \texttt{---------}\\
6 & 0 & 1 & 1 & 3 & 0 & 512 & 0 & \texttt{---------}\\
7 & 0 & 1 & 1 & 3 & 0 & 512 & 0 & \texttt{---------}\\
8 & 0 & 1 & 1 & 3 & 0 & 512 & 0 & \texttt{---------}\\
\multicolumn{9}{@{}l@{}}{\rule{0pt}{2.6ex}\textbf{T9.133}\enspace \texttt{012343434}\enspace $n{=}9$, $|H|{=}3$, len 5, br 1, $|\mathrm{Aut}|{=}6$, orb 5/3, hist-orb 1/0; M 12288, res 3024, rev 5172: ok}\\*
0 & -- & 0 & 3 & 1 & 438 & 2 & 72 & \texttt{-++??????}\\*
1 & 0 & 1 & 3 & 1 & 364 & 4 & 144 & \texttt{--+??????}\\
2 & 1 & 2 & 3 & 1 & 216 & 8 & 288 & \texttt{---??????}\\
3 & 2 & 3 & 1 & 3 & 256 & 256 & 0 & \texttt{----+----}\\
4 & 3 & 4 & 1 & 3 & 0 & 512 & 0 & \texttt{---------}\\
5 & 2 & 3 & 1 & 3 & 256 & 256 & 0 & \texttt{------+--}\\
6 & 5 & 4 & 1 & 3 & 0 & 512 & 0 & \texttt{---------}\\
7 & 2 & 3 & 1 & 3 & 256 & 256 & 0 & \texttt{--------+}\\
8 & 7 & 4 & 1 & 3 & 0 & 512 & 0 & \texttt{---------}\\
\multicolumn{9}{@{}l@{}}{\rule{0pt}{2.6ex}\textbf{T9.134}\enspace \texttt{012343433}\enspace $n{=}9$, $|H|{=}4$, len 5, br 1, $|\mathrm{Aut}|{=}4$, orb 6/3, hist-orb 2/0; M 11776, res 4536, rev 4312: ok}\\*
0 & -- & 0 & 4 & 1 & 402 & 2 & 108 & \texttt{-++??????}\\*
1 & 0 & 1 & 4 & 1 & 292 & 4 & 216 & \texttt{--+??????}\\
2 & 1 & 2 & 4 & 1 & 72 & 8 & 432 & \texttt{---??????}\\
3 & 2 & 3 & 1 & 2 & 256 & 256 & 0 & \texttt{----+----}\\
4 & 3 & 4 & 1 & 2 & 0 & 512 & 0 & \texttt{---------}\\
5 & 2 & 3 & 1 & 2 & 256 & 256 & 0 & \texttt{------+--}\\
6 & 5 & 4 & 1 & 2 & 0 & 512 & 0 & \texttt{---------}\\
7 & 2 & 3 & 1 & 2 & 0 & 512 & 0 & \texttt{---------}\\
8 & 2 & 3 & 1 & 2 & 0 & 512 & 0 & \texttt{---------}\\
\multicolumn{9}{@{}l@{}}{\rule{0pt}{2.6ex}\textbf{T9.135}\enspace \texttt{012343432}\enspace $n{=}9$, $|H|{=}4$, len 5, br 2, $|\mathrm{Aut}|{=}2$, orb 7/5, hist-orb 3/2; M 11264, res 5108, rev 3548: ok}\\*
0 & -- & 0 & 4 & 1 & 338 & 2 & 172 & \texttt{-+???????}\\*
1 & 0 & 1 & 4 & 1 & 164 & 4 & 344 & \texttt{--???????}\\
2 & 1 & 2 & 3 & 1 & 144 & 16 & 352 & \texttt{---?????-}\\
3 & 2 & 3 & 1 & 2 & 256 & 256 & 0 & \texttt{----+----}\\
4 & 3 & 4 & 1 & 2 & 0 & 512 & 0 & \texttt{---------}\\
5 & 2 & 3 & 1 & 2 & 256 & 256 & 0 & \texttt{------+--}\\
6 & 5 & 4 & 1 & 2 & 0 & 512 & 0 & \texttt{---------}\\
7 & 2 & 3 & 1 & 1 & 0 & 512 & 0 & \texttt{---------}\\
8 & 1 & 2 & 1 & 1 & 0 & 512 & 0 & \texttt{---------}\\
\multicolumn{9}{@{}l@{}}{\rule{0pt}{2.6ex}\textbf{T9.136}\enspace \texttt{012343431}\enspace $n{=}9$, $|H|{=}4$, len 5, br 2, $|\mathrm{Aut}|{=}2$, orb 7/5, hist-orb 3/2; M 10752, res 4864, rev 3388: ok}\\*
0 & -- & 0 & 4 & 1 & 210 & 2 & 300 & \texttt{-????????}\\*
1 & 0 & 1 & 3 & 1 & 328 & 8 & 176 & \texttt{--+?????-}\\
2 & 1 & 2 & 3 & 1 & 144 & 16 & 352 & \texttt{---?????-}\\
3 & 2 & 3 & 1 & 2 & 256 & 256 & 0 & \texttt{----+----}\\
4 & 3 & 4 & 1 & 2 & 0 & 512 & 0 & \texttt{---------}\\
5 & 2 & 3 & 1 & 2 & 256 & 256 & 0 & \texttt{------+--}\\
6 & 5 & 4 & 1 & 2 & 0 & 512 & 0 & \texttt{---------}\\
7 & 2 & 3 & 1 & 1 & 0 & 512 & 0 & \texttt{---------}\\
8 & 0 & 1 & 1 & 1 & 0 & 512 & 0 & \texttt{---------}\\
\multicolumn{9}{@{}l@{}}{\rule{0pt}{2.6ex}\textbf{T9.137}\enspace \texttt{012343423}\enspace $n{=}9$, $|H|{=}3$, len 5, br 2, $|\mathrm{Aut}|{=}2$, orb 7/5, hist-orb 2/1; M 11264, res 2808, rev 4620: ok}\\*
0 & -- & 0 & 3 & 1 & 406 & 2 & 104 & \texttt{-+???????}\\*
1 & 0 & 1 & 3 & 1 & 300 & 4 & 208 & \texttt{--???????}\\
2 & 1 & 2 & 2 & 1 & 288 & 32 & 192 & \texttt{---????--}\\
3 & 2 & 3 & 1 & 2 & 256 & 256 & 0 & \texttt{----+----}\\
4 & 3 & 4 & 1 & 2 & 0 & 512 & 0 & \texttt{---------}\\
5 & 2 & 3 & 1 & 2 & 256 & 256 & 0 & \texttt{------+--}\\
6 & 5 & 4 & 1 & 2 & 0 & 512 & 0 & \texttt{---------}\\
7 & 1 & 2 & 1 & 1 & 256 & 256 & 0 & \texttt{--------+}\\
8 & 7 & 3 & 1 & 1 & 0 & 512 & 0 & \texttt{---------}\\
\multicolumn{9}{@{}l@{}}{\rule{0pt}{2.6ex}\textbf{T9.138}\enspace \texttt{012343422}\enspace $n{=}9$, $|H|{=}4$, len 5, br 2, $|\mathrm{Aut}|{=}4$, orb 6/3, hist-orb 2/0; M 10752, res 4908, rev 3316: ok}\\*
0 & -- & 0 & 4 & 1 & 306 & 2 & 204 & \texttt{-+???????}\\*
1 & 0 & 1 & 4 & 1 & 100 & 4 & 408 & \texttt{--???????}\\
2 & 1 & 2 & 2 & 1 & 288 & 32 & 192 & \texttt{---????--}\\
3 & 2 & 3 & 1 & 2 & 256 & 256 & 0 & \texttt{----+----}\\
4 & 3 & 4 & 1 & 2 & 0 & 512 & 0 & \texttt{---------}\\
5 & 2 & 3 & 1 & 2 & 256 & 256 & 0 & \texttt{------+--}\\
6 & 5 & 4 & 1 & 2 & 0 & 512 & 0 & \texttt{---------}\\
7 & 1 & 2 & 1 & 2 & 0 & 512 & 0 & \texttt{---------}\\
8 & 1 & 2 & 1 & 2 & 0 & 512 & 0 & \texttt{---------}\\
\multicolumn{9}{@{}l@{}}{\rule{0pt}{2.6ex}\textbf{T9.139}\enspace \texttt{012343421}\enspace $n{=}9$, $|H|{=}4$, len 5, br 3, $|\mathrm{Aut}|{=}2$, orb 7/5, hist-orb 3/2; M 10240, res 4952, rev 2960: ok}\\*
0 & -- & 0 & 4 & 1 & 178 & 2 & 332 & \texttt{-????????}\\*
1 & 0 & 1 & 3 & 1 & 200 & 8 & 304 & \texttt{--??????-}\\
2 & 1 & 2 & 2 & 1 & 288 & 32 & 192 & \texttt{---????--}\\
3 & 2 & 3 & 1 & 2 & 256 & 256 & 0 & \texttt{----+----}\\
4 & 3 & 4 & 1 & 2 & 0 & 512 & 0 & \texttt{---------}\\
5 & 2 & 3 & 1 & 2 & 256 & 256 & 0 & \texttt{------+--}\\
6 & 5 & 4 & 1 & 2 & 0 & 512 & 0 & \texttt{---------}\\
7 & 1 & 2 & 1 & 1 & 0 & 512 & 0 & \texttt{---------}\\
8 & 0 & 1 & 1 & 1 & 0 & 512 & 0 & \texttt{---------}\\
\multicolumn{9}{@{}l@{}}{\rule{0pt}{2.6ex}\textbf{T9.140}\enspace \texttt{012343412}\enspace $n{=}9$, $|H|{=}3$, len 5, br 2, $|\mathrm{Aut}|{=}2$, orb 7/5, hist-orb 2/1; M 10240, res 2400, rev 4192: ok}\\*
0 & -- & 0 & 3 & 1 & 342 & 2 & 168 & \texttt{-????????}\\*
1 & 0 & 1 & 2 & 1 & 400 & 16 & 96 & \texttt{--+????--}\\
2 & 1 & 2 & 2 & 1 & 288 & 32 & 192 & \texttt{---????--}\\
3 & 2 & 3 & 1 & 2 & 256 & 256 & 0 & \texttt{----+----}\\
4 & 3 & 4 & 1 & 2 & 0 & 512 & 0 & \texttt{---------}\\
5 & 2 & 3 & 1 & 2 & 256 & 256 & 0 & \texttt{------+--}\\
6 & 5 & 4 & 1 & 2 & 0 & 512 & 0 & \texttt{---------}\\
7 & 0 & 1 & 1 & 1 & 256 & 256 & 0 & \texttt{--------+}\\
8 & 7 & 2 & 1 & 1 & 0 & 512 & 0 & \texttt{---------}\\
\multicolumn{9}{@{}l@{}}{\rule{0pt}{2.6ex}\textbf{T9.141}\enspace \texttt{012343411}\enspace $n{=}9$, $|H|{=}4$, len 5, br 2, $|\mathrm{Aut}|{=}4$, orb 6/3, hist-orb 2/0; M 9728, res 4224, rev 3020: ok}\\*
0 & -- & 0 & 4 & 1 & 114 & 2 & 396 & \texttt{-????????}\\*
1 & 0 & 1 & 2 & 1 & 400 & 16 & 96 & \texttt{--+????--}\\
2 & 1 & 2 & 2 & 1 & 288 & 32 & 192 & \texttt{---????--}\\
3 & 2 & 3 & 1 & 2 & 256 & 256 & 0 & \texttt{----+----}\\
4 & 3 & 4 & 1 & 2 & 0 & 512 & 0 & \texttt{---------}\\
5 & 2 & 3 & 1 & 2 & 256 & 256 & 0 & \texttt{------+--}\\
6 & 5 & 4 & 1 & 2 & 0 & 512 & 0 & \texttt{---------}\\
7 & 0 & 1 & 1 & 2 & 0 & 512 & 0 & \texttt{---------}\\
8 & 0 & 1 & 1 & 2 & 0 & 512 & 0 & \texttt{---------}\\
\multicolumn{9}{@{}l@{}}{\rule{0pt}{2.6ex}\textbf{T9.142}\enspace \texttt{012343333}\enspace $n{=}9$, $|H|{=}5$, len 5, br 1, $|\mathrm{Aut}|{=}24$, orb 6/5, hist-orb 2/1; M 11264, res 5040, rev 4012: ok}\\*
0 & -- & 0 & 5 & 1 & 390 & 2 & 120 & \texttt{-++??????}\\*
1 & 0 & 1 & 5 & 1 & 268 & 4 & 240 & \texttt{--+??????}\\
2 & 1 & 2 & 5 & 1 & 24 & 8 & 480 & \texttt{---??????}\\
3 & 2 & 3 & 1 & 1 & 256 & 256 & 0 & \texttt{----+----}\\
4 & 3 & 4 & 1 & 1 & 0 & 512 & 0 & \texttt{---------}\\
5 & 2 & 3 & 1 & 4 & 0 & 512 & 0 & \texttt{---------}\\
6 & 2 & 3 & 1 & 4 & 0 & 512 & 0 & \texttt{---------}\\
7 & 2 & 3 & 1 & 4 & 0 & 512 & 0 & \texttt{---------}\\
8 & 2 & 3 & 1 & 4 & 0 & 512 & 0 & \texttt{---------}\\
\multicolumn{9}{@{}l@{}}{\rule{0pt}{2.6ex}\textbf{T9.143}\enspace \texttt{012343332}\enspace $n{=}9$, $|H|{=}5$, len 5, br 2, $|\mathrm{Aut}|{=}6$, orb 7/6, hist-orb 3/2; M 10752, res 5768, rev 3040: ok}\\*
0 & -- & 0 & 5 & 1 & 326 & 2 & 184 & \texttt{-+???????}\\*
1 & 0 & 1 & 5 & 1 & 140 & 4 & 368 & \texttt{--???????}\\
2 & 1 & 2 & 4 & 1 & 48 & 16 & 448 & \texttt{---?????-}\\
3 & 2 & 3 & 1 & 1 & 256 & 256 & 0 & \texttt{----+----}\\
4 & 3 & 4 & 1 & 1 & 0 & 512 & 0 & \texttt{---------}\\
5 & 2 & 3 & 1 & 3 & 0 & 512 & 0 & \texttt{---------}\\
6 & 2 & 3 & 1 & 3 & 0 & 512 & 0 & \texttt{---------}\\
7 & 2 & 3 & 1 & 3 & 0 & 512 & 0 & \texttt{---------}\\
8 & 1 & 2 & 1 & 1 & 0 & 512 & 0 & \texttt{---------}\\
\multicolumn{9}{@{}l@{}}{\rule{0pt}{2.6ex}\textbf{T9.144}\enspace \texttt{012343331}\enspace $n{=}9$, $|H|{=}5$, len 5, br 2, $|\mathrm{Aut}|{=}6$, orb 7/6, hist-orb 3/2; M 10240, res 5632, rev 2824: ok}\\*
0 & -- & 0 & 5 & 1 & 198 & 2 & 312 & \texttt{-????????}\\*
1 & 0 & 1 & 4 & 1 & 280 & 8 & 224 & \texttt{--+?????-}\\
2 & 1 & 2 & 4 & 1 & 48 & 16 & 448 & \texttt{---?????-}\\
3 & 2 & 3 & 1 & 1 & 256 & 256 & 0 & \texttt{----+----}\\
4 & 3 & 4 & 1 & 1 & 0 & 512 & 0 & \texttt{---------}\\
5 & 2 & 3 & 1 & 3 & 0 & 512 & 0 & \texttt{---------}\\
6 & 2 & 3 & 1 & 3 & 0 & 512 & 0 & \texttt{---------}\\
7 & 2 & 3 & 1 & 3 & 0 & 512 & 0 & \texttt{---------}\\
8 & 0 & 1 & 1 & 1 & 0 & 512 & 0 & \texttt{---------}\\
\multicolumn{9}{@{}l@{}}{\rule{0pt}{2.6ex}\textbf{T9.145}\enspace \texttt{012343323}\enspace $n{=}9$, $|H|{=}4$, len 5, br 2, $|\mathrm{Aut}|{=}2$, orb 8/7, hist-orb 3/2; M 10752, res 4188, rev 3792: ok}\\*
0 & -- & 0 & 4 & 1 & 370 & 2 & 140 & \texttt{-+???????}\\*
1 & 0 & 1 & 4 & 1 & 228 & 4 & 280 & \texttt{--???????}\\
2 & 1 & 2 & 3 & 1 & 96 & 32 & 384 & \texttt{---????--}\\
3 & 2 & 3 & 1 & 1 & 256 & 256 & 0 & \texttt{----+----}\\
4 & 3 & 4 & 1 & 1 & 0 & 512 & 0 & \texttt{---------}\\
5 & 2 & 3 & 1 & 2 & 0 & 512 & 0 & \texttt{---------}\\
6 & 2 & 3 & 1 & 2 & 0 & 512 & 0 & \texttt{---------}\\
7 & 1 & 2 & 1 & 1 & 256 & 256 & 0 & \texttt{--------+}\\
8 & 7 & 3 & 1 & 1 & 0 & 512 & 0 & \texttt{---------}\\
\multicolumn{9}{@{}l@{}}{\rule{0pt}{2.6ex}\textbf{T9.146}\enspace \texttt{012343322}\enspace $n{=}9$, $|H|{=}5$, len 5, br 2, $|\mathrm{Aut}|{=}4$, orb 7/5, hist-orb 3/1; M 10240, res 5784, rev 2632: ok}\\*
0 & -- & 0 & 5 & 1 & 294 & 2 & 216 & \texttt{-+???????}\\*
1 & 0 & 1 & 5 & 1 & 76 & 4 & 432 & \texttt{--???????}\\
2 & 1 & 2 & 3 & 1 & 96 & 32 & 384 & \texttt{---????--}\\
3 & 2 & 3 & 1 & 1 & 256 & 256 & 0 & \texttt{----+----}\\
4 & 3 & 4 & 1 & 1 & 0 & 512 & 0 & \texttt{---------}\\
5 & 2 & 3 & 1 & 2 & 0 & 512 & 0 & \texttt{---------}\\
6 & 2 & 3 & 1 & 2 & 0 & 512 & 0 & \texttt{---------}\\
7 & 1 & 2 & 1 & 2 & 0 & 512 & 0 & \texttt{---------}\\
8 & 1 & 2 & 1 & 2 & 0 & 512 & 0 & \texttt{---------}\\
\multicolumn{9}{@{}l@{}}{\rule{0pt}{2.6ex}\textbf{T9.147}\enspace \texttt{012343321}\enspace $n{=}9$, $|H|{=}5$, len 5, br 3, $|\mathrm{Aut}|{=}2$, orb 8/7, hist-orb 4/3; M 9728, res 5936, rev 2188: ok}\\*
0 & -- & 0 & 5 & 1 & 166 & 2 & 344 & \texttt{-????????}\\*
1 & 0 & 1 & 4 & 1 & 152 & 8 & 352 & \texttt{--??????-}\\
2 & 1 & 2 & 3 & 1 & 96 & 32 & 384 & \texttt{---????--}\\
3 & 2 & 3 & 1 & 1 & 256 & 256 & 0 & \texttt{----+----}\\
4 & 3 & 4 & 1 & 1 & 0 & 512 & 0 & \texttt{---------}\\
5 & 2 & 3 & 1 & 2 & 0 & 512 & 0 & \texttt{---------}\\
6 & 2 & 3 & 1 & 2 & 0 & 512 & 0 & \texttt{---------}\\
7 & 1 & 2 & 1 & 1 & 0 & 512 & 0 & \texttt{---------}\\
8 & 0 & 1 & 1 & 1 & 0 & 512 & 0 & \texttt{---------}\\
\multicolumn{9}{@{}l@{}}{\rule{0pt}{2.6ex}\textbf{T9.148}\enspace \texttt{012343312}\enspace $n{=}9$, $|H|{=}4$, len 5, br 2, $|\mathrm{Aut}|{=}2$, orb 8/7, hist-orb 3/2; M 9728, res 3744, rev 3404: ok}\\*
0 & -- & 0 & 4 & 1 & 306 & 2 & 204 & \texttt{-????????}\\*
1 & 0 & 1 & 3 & 1 & 304 & 16 & 192 & \texttt{--+????--}\\
2 & 1 & 2 & 3 & 1 & 96 & 32 & 384 & \texttt{---????--}\\
3 & 2 & 3 & 1 & 1 & 256 & 256 & 0 & \texttt{----+----}\\
4 & 3 & 4 & 1 & 1 & 0 & 512 & 0 & \texttt{---------}\\
5 & 2 & 3 & 1 & 2 & 0 & 512 & 0 & \texttt{---------}\\
6 & 2 & 3 & 1 & 2 & 0 & 512 & 0 & \texttt{---------}\\
7 & 0 & 1 & 1 & 1 & 256 & 256 & 0 & \texttt{--------+}\\
8 & 7 & 2 & 1 & 1 & 0 & 512 & 0 & \texttt{---------}\\
\multicolumn{9}{@{}l@{}}{\rule{0pt}{2.6ex}\textbf{T9.149}\enspace \texttt{012343311}\enspace $n{=}9$, $|H|{=}5$, len 5, br 2, $|\mathrm{Aut}|{=}4$, orb 7/5, hist-orb 3/1; M 9216, res 5376, rev 2216: ok}\\*
0 & -- & 0 & 5 & 1 & 102 & 2 & 408 & \texttt{-????????}\\*
1 & 0 & 1 & 3 & 1 & 304 & 16 & 192 & \texttt{--+????--}\\
2 & 1 & 2 & 3 & 1 & 96 & 32 & 384 & \texttt{---????--}\\
3 & 2 & 3 & 1 & 1 & 256 & 256 & 0 & \texttt{----+----}\\
4 & 3 & 4 & 1 & 1 & 0 & 512 & 0 & \texttt{---------}\\
5 & 2 & 3 & 1 & 2 & 0 & 512 & 0 & \texttt{---------}\\
6 & 2 & 3 & 1 & 2 & 0 & 512 & 0 & \texttt{---------}\\
7 & 0 & 1 & 1 & 2 & 0 & 512 & 0 & \texttt{---------}\\
8 & 0 & 1 & 1 & 2 & 0 & 512 & 0 & \texttt{---------}\\
\multicolumn{9}{@{}l@{}}{\rule{0pt}{2.6ex}\textbf{T9.150}\enspace \texttt{012343234}\enspace $n{=}9$, $|H|{=}3$, len 5, br 2, $|\mathrm{Aut}|{=}1$, orb 9/9, hist-orb 3/3; M 11264, res 2676, rev 4496: ok}\\*
0 & -- & 0 & 3 & 1 & 410 & 2 & 100 & \texttt{-+???????}\\*
1 & 0 & 1 & 3 & 1 & 308 & 4 & 200 & \texttt{--???????}\\
2 & 1 & 2 & 2 & 1 & 192 & 64 & 256 & \texttt{---???---}\\
3 & 2 & 3 & 1 & 1 & 256 & 256 & 0 & \texttt{----+----}\\
4 & 3 & 4 & 1 & 1 & 0 & 512 & 0 & \texttt{---------}\\
5 & 2 & 3 & 1 & 1 & 0 & 512 & 0 & \texttt{---------}\\
6 & 1 & 2 & 1 & 1 & 384 & 128 & 0 & \texttt{-------++}\\
7 & 6 & 3 & 1 & 1 & 256 & 256 & 0 & \texttt{--------+}\\
8 & 7 & 4 & 1 & 1 & 0 & 512 & 0 & \texttt{---------}\\
\multicolumn{9}{@{}l@{}}{\rule{0pt}{2.6ex}\textbf{T9.151}\enspace \texttt{012343233}\enspace $n{=}9$, $|H|{=}4$, len 5, br 3, $|\mathrm{Aut}|{=}2$, orb 8/7, hist-orb 3/2; M 10752, res 3920, rev 3776: ok}\\*
0 & -- & 0 & 4 & 1 & 366 & 2 & 144 & \texttt{-+???????}\\*
1 & 0 & 1 & 4 & 1 & 220 & 4 & 288 & \texttt{--???????}\\
2 & 1 & 2 & 2 & 1 & 192 & 64 & 256 & \texttt{---???---}\\
3 & 2 & 3 & 1 & 1 & 256 & 256 & 0 & \texttt{----+----}\\
4 & 3 & 4 & 1 & 1 & 0 & 512 & 0 & \texttt{---------}\\
5 & 2 & 3 & 1 & 1 & 0 & 512 & 0 & \texttt{---------}\\
6 & 1 & 2 & 2 & 1 & 128 & 128 & 256 & \texttt{-------??}\\
7 & 6 & 3 & 1 & 2 & 0 & 512 & 0 & \texttt{---------}\\
8 & 6 & 3 & 1 & 2 & 0 & 512 & 0 & \texttt{---------}\\
\multicolumn{9}{@{}l@{}}{\rule{0pt}{2.6ex}\textbf{T9.152}\enspace \texttt{012343232}\enspace $n{=}9$, $|H|{=}4$, len 5, br 2, $|\mathrm{Aut}|{=}1$, orb 9/9, hist-orb 4/4; M 10240, res 4524, rev 3240: ok}\\*
0 & -- & 0 & 4 & 1 & 322 & 2 & 188 & \texttt{-+???????}\\*
1 & 0 & 1 & 4 & 1 & 132 & 4 & 376 & \texttt{--???????}\\
2 & 1 & 2 & 2 & 1 & 192 & 64 & 256 & \texttt{---???---}\\
3 & 2 & 3 & 1 & 1 & 256 & 256 & 0 & \texttt{----+----}\\
4 & 3 & 4 & 1 & 1 & 0 & 512 & 0 & \texttt{---------}\\
5 & 2 & 3 & 1 & 1 & 0 & 512 & 0 & \texttt{---------}\\
6 & 1 & 2 & 1 & 1 & 256 & 256 & 0 & \texttt{-------+-}\\
7 & 6 & 3 & 1 & 1 & 0 & 512 & 0 & \texttt{---------}\\
8 & 1 & 2 & 1 & 1 & 0 & 512 & 0 & \texttt{---------}\\
\multicolumn{9}{@{}l@{}}{\rule{0pt}{2.6ex}\textbf{T9.153}\enspace \texttt{012343231}\enspace $n{=}9$, $|H|{=}4$, len 5, br 3, $|\mathrm{Aut}|{=}1$, orb 9/9, hist-orb 4/4; M 9728, res 4424, rev 2964: ok}\\*
0 & -- & 0 & 4 & 1 & 194 & 2 & 316 & \texttt{-????????}\\*
1 & 0 & 1 & 3 & 1 & 264 & 8 & 240 & \texttt{--??????-}\\
2 & 1 & 2 & 2 & 1 & 192 & 64 & 256 & \texttt{---???---}\\
3 & 2 & 3 & 1 & 1 & 256 & 256 & 0 & \texttt{----+----}\\
4 & 3 & 4 & 1 & 1 & 0 & 512 & 0 & \texttt{---------}\\
5 & 2 & 3 & 1 & 1 & 0 & 512 & 0 & \texttt{---------}\\
6 & 1 & 2 & 1 & 1 & 256 & 256 & 0 & \texttt{-------+-}\\
7 & 6 & 3 & 1 & 1 & 0 & 512 & 0 & \texttt{---------}\\
8 & 0 & 1 & 1 & 1 & 0 & 512 & 0 & \texttt{---------}\\
\multicolumn{9}{@{}l@{}}{\rule{0pt}{2.6ex}\textbf{T9.154}\enspace \texttt{012343222}\enspace $n{=}9$, $|H|{=}5$, len 5, br 2, $|\mathrm{Aut}|{=}6$, orb 7/6, hist-orb 3/2; M 9728, res 5448, rev 2512: ok}\\*
0 & -- & 0 & 5 & 1 & 278 & 2 & 232 & \texttt{-+???????}\\*
1 & 0 & 1 & 5 & 1 & 44 & 4 & 464 & \texttt{--???????}\\
2 & 1 & 2 & 2 & 1 & 192 & 64 & 256 & \texttt{---???---}\\
3 & 2 & 3 & 1 & 1 & 256 & 256 & 0 & \texttt{----+----}\\
4 & 3 & 4 & 1 & 1 & 0 & 512 & 0 & \texttt{---------}\\
5 & 2 & 3 & 1 & 1 & 0 & 512 & 0 & \texttt{---------}\\
6 & 1 & 2 & 1 & 3 & 0 & 512 & 0 & \texttt{---------}\\
7 & 1 & 2 & 1 & 3 & 0 & 512 & 0 & \texttt{---------}\\
8 & 1 & 2 & 1 & 3 & 0 & 512 & 0 & \texttt{---------}\\
\multicolumn{9}{@{}l@{}}{\rule{0pt}{2.6ex}\textbf{T9.155}\enspace \texttt{012343221}\enspace $n{=}9$, $|H|{=}5$, len 5, br 3, $|\mathrm{Aut}|{=}2$, orb 8/7, hist-orb 4/3; M 9216, res 5744, rev 1948: ok}\\*
0 & -- & 0 & 5 & 1 & 150 & 2 & 360 & \texttt{-????????}\\*
1 & 0 & 1 & 4 & 1 & 88 & 8 & 416 & \texttt{--??????-}\\
2 & 1 & 2 & 2 & 1 & 192 & 64 & 256 & \texttt{---???---}\\
3 & 2 & 3 & 1 & 1 & 256 & 256 & 0 & \texttt{----+----}\\
4 & 3 & 4 & 1 & 1 & 0 & 512 & 0 & \texttt{---------}\\
5 & 2 & 3 & 1 & 1 & 0 & 512 & 0 & \texttt{---------}\\
6 & 1 & 2 & 1 & 2 & 0 & 512 & 0 & \texttt{---------}\\
7 & 1 & 2 & 1 & 2 & 0 & 512 & 0 & \texttt{---------}\\
8 & 0 & 1 & 1 & 1 & 0 & 512 & 0 & \texttt{---------}\\
\multicolumn{9}{@{}l@{}}{\rule{0pt}{2.6ex}\textbf{T9.156}\enspace \texttt{012343212}\enspace $n{=}9$, $|H|{=}4$, len 5, br 3, $|\mathrm{Aut}|{=}1$, orb 9/9, hist-orb 4/4; M 9216, res 4032, rev 2864: ok}\\*
0 & -- & 0 & 4 & 1 & 258 & 2 & 252 & \texttt{-????????}\\*
1 & 0 & 1 & 3 & 1 & 176 & 16 & 320 & \texttt{--?????--}\\
2 & 1 & 2 & 2 & 1 & 192 & 64 & 256 & \texttt{---???---}\\
3 & 2 & 3 & 1 & 1 & 256 & 256 & 0 & \texttt{----+----}\\
4 & 3 & 4 & 1 & 1 & 0 & 512 & 0 & \texttt{---------}\\
5 & 2 & 3 & 1 & 1 & 0 & 512 & 0 & \texttt{---------}\\
6 & 1 & 2 & 1 & 1 & 0 & 512 & 0 & \texttt{---------}\\
7 & 0 & 1 & 1 & 1 & 256 & 256 & 0 & \texttt{--------+}\\
8 & 7 & 2 & 1 & 1 & 0 & 512 & 0 & \texttt{---------}\\
\multicolumn{9}{@{}l@{}}{\rule{0pt}{2.6ex}\textbf{T9.157}\enspace \texttt{012343211}\enspace $n{=}9$, $|H|{=}5$, len 5, br 3, $|\mathrm{Aut}|{=}2$, orb 8/7, hist-orb 4/3; M 8704, res 5408, rev 1812: ok}\\*
0 & -- & 0 & 5 & 1 & 86 & 2 & 424 & \texttt{-????????}\\*
1 & 0 & 1 & 3 & 1 & 176 & 16 & 320 & \texttt{--?????--}\\
2 & 1 & 2 & 2 & 1 & 192 & 64 & 256 & \texttt{---???---}\\
3 & 2 & 3 & 1 & 1 & 256 & 256 & 0 & \texttt{----+----}\\
4 & 3 & 4 & 1 & 1 & 0 & 512 & 0 & \texttt{---------}\\
5 & 2 & 3 & 1 & 1 & 0 & 512 & 0 & \texttt{---------}\\
6 & 1 & 2 & 1 & 1 & 0 & 512 & 0 & \texttt{---------}\\
7 & 0 & 1 & 1 & 2 & 0 & 512 & 0 & \texttt{---------}\\
8 & 0 & 1 & 1 & 2 & 0 & 512 & 0 & \texttt{---------}\\
\multicolumn{9}{@{}l@{}}{\rule{0pt}{2.6ex}\textbf{T9.158}\enspace \texttt{012343123}\enspace $n{=}9$, $|H|{=}3$, len 5, br 2, $|\mathrm{Aut}|{=}1$, orb 9/9, hist-orb 3/3; M 9728, res 2080, rev 3956: ok}\\*
0 & -- & 0 & 3 & 1 & 378 & 2 & 132 & \texttt{-????????}\\*
1 & 0 & 1 & 2 & 1 & 352 & 32 & 128 & \texttt{--+???---}\\
2 & 1 & 2 & 2 & 1 & 192 & 64 & 256 & \texttt{---???---}\\
3 & 2 & 3 & 1 & 1 & 256 & 256 & 0 & \texttt{----+----}\\
4 & 3 & 4 & 1 & 1 & 0 & 512 & 0 & \texttt{---------}\\
5 & 2 & 3 & 1 & 1 & 0 & 512 & 0 & \texttt{---------}\\
6 & 0 & 1 & 1 & 1 & 384 & 128 & 0 & \texttt{-------++}\\
7 & 6 & 2 & 1 & 1 & 256 & 256 & 0 & \texttt{--------+}\\
8 & 7 & 3 & 1 & 1 & 0 & 512 & 0 & \texttt{---------}\\
\multicolumn{9}{@{}l@{}}{\rule{0pt}{2.6ex}\textbf{T9.159}\enspace \texttt{012343122}\enspace $n{=}9$, $|H|{=}4$, len 5, br 3, $|\mathrm{Aut}|{=}2$, orb 8/7, hist-orb 3/2; M 9216, res 3328, rev 3160: ok}\\*
0 & -- & 0 & 4 & 1 & 270 & 2 & 240 & \texttt{-????????}\\*
1 & 0 & 1 & 2 & 1 & 352 & 32 & 128 & \texttt{--+???---}\\
2 & 1 & 2 & 2 & 1 & 192 & 64 & 256 & \texttt{---???---}\\
3 & 2 & 3 & 1 & 1 & 256 & 256 & 0 & \texttt{----+----}\\
4 & 3 & 4 & 1 & 1 & 0 & 512 & 0 & \texttt{---------}\\
5 & 2 & 3 & 1 & 1 & 0 & 512 & 0 & \texttt{---------}\\
6 & 0 & 1 & 2 & 1 & 128 & 128 & 256 & \texttt{-------??}\\
7 & 6 & 2 & 1 & 2 & 0 & 512 & 0 & \texttt{---------}\\
8 & 6 & 2 & 1 & 2 & 0 & 512 & 0 & \texttt{---------}\\
\multicolumn{9}{@{}l@{}}{\rule{0pt}{2.6ex}\textbf{T9.160}\enspace \texttt{012343121}\enspace $n{=}9$, $|H|{=}4$, len 5, br 2, $|\mathrm{Aut}|{=}1$, orb 9/9, hist-orb 4/4; M 8704, res 3808, rev 2660: ok}\\*
0 & -- & 0 & 4 & 1 & 162 & 2 & 348 & \texttt{-????????}\\*
1 & 0 & 1 & 2 & 1 & 352 & 32 & 128 & \texttt{--+???---}\\
2 & 1 & 2 & 2 & 1 & 192 & 64 & 256 & \texttt{---???---}\\
3 & 2 & 3 & 1 & 1 & 256 & 256 & 0 & \texttt{----+----}\\
4 & 3 & 4 & 1 & 1 & 0 & 512 & 0 & \texttt{---------}\\
5 & 2 & 3 & 1 & 1 & 0 & 512 & 0 & \texttt{---------}\\
6 & 0 & 1 & 1 & 1 & 256 & 256 & 0 & \texttt{-------+-}\\
7 & 6 & 2 & 1 & 1 & 0 & 512 & 0 & \texttt{---------}\\
8 & 0 & 1 & 1 & 1 & 0 & 512 & 0 & \texttt{---------}\\
\multicolumn{9}{@{}l@{}}{\rule{0pt}{2.6ex}\textbf{T9.161}\enspace \texttt{012343111}\enspace $n{=}9$, $|H|{=}5$, len 5, br 2, $|\mathrm{Aut}|{=}6$, orb 7/6, hist-orb 3/2; M 8192, res 4672, rev 1904: ok}\\*
0 & -- & 0 & 5 & 1 & 54 & 2 & 456 & \texttt{-????????}\\*
1 & 0 & 1 & 2 & 1 & 352 & 32 & 128 & \texttt{--+???---}\\
2 & 1 & 2 & 2 & 1 & 192 & 64 & 256 & \texttt{---???---}\\
3 & 2 & 3 & 1 & 1 & 256 & 256 & 0 & \texttt{----+----}\\
4 & 3 & 4 & 1 & 1 & 0 & 512 & 0 & \texttt{---------}\\
5 & 2 & 3 & 1 & 1 & 0 & 512 & 0 & \texttt{---------}\\
6 & 0 & 1 & 1 & 3 & 0 & 512 & 0 & \texttt{---------}\\
7 & 0 & 1 & 1 & 3 & 0 & 512 & 0 & \texttt{---------}\\
8 & 0 & 1 & 1 & 3 & 0 & 512 & 0 & \texttt{---------}\\
\multicolumn{9}{@{}l@{}}{\rule{0pt}{2.6ex}\textbf{T9.162}\enspace \texttt{012342342}\enspace $n{=}9$, $|H|{=}3$, len 5, br 1, $|\mathrm{Aut}|{=}2$, orb 6/3, hist-orb 2/1; M 10752, res 3276, rev 3908: ok}\\*
0 & -- & 0 & 3 & 1 & 354 & 2 & 156 & \texttt{-+???????}\\*
1 & 0 & 1 & 3 & 1 & 196 & 4 & 312 & \texttt{--???????}\\
2 & 1 & 2 & 1 & 2 & 384 & 128 & 0 & \texttt{---++----}\\
3 & 2 & 3 & 1 & 2 & 256 & 256 & 0 & \texttt{----+----}\\
4 & 3 & 4 & 1 & 2 & 0 & 512 & 0 & \texttt{---------}\\
5 & 1 & 2 & 1 & 2 & 384 & 128 & 0 & \texttt{------++-}\\
6 & 5 & 3 & 1 & 2 & 256 & 256 & 0 & \texttt{-------+-}\\
7 & 6 & 4 & 1 & 2 & 0 & 512 & 0 & \texttt{---------}\\
8 & 1 & 2 & 1 & 1 & 0 & 512 & 0 & \texttt{---------}\\
\multicolumn{9}{@{}l@{}}{\rule{0pt}{2.6ex}\textbf{T9.163}\enspace \texttt{012342341}\enspace $n{=}9$, $|H|{=}3$, len 5, br 2, $|\mathrm{Aut}|{=}2$, orb 6/3, hist-orb 2/1; M 10240, res 2888, rev 3752: ok}\\*
0 & -- & 0 & 3 & 1 & 226 & 2 & 284 & \texttt{-????????}\\*
1 & 0 & 1 & 2 & 1 & 392 & 8 & 112 & \texttt{--??????-}\\
2 & 1 & 2 & 1 & 2 & 384 & 128 & 0 & \texttt{---++----}\\
3 & 2 & 3 & 1 & 2 & 256 & 256 & 0 & \texttt{----+----}\\
4 & 3 & 4 & 1 & 2 & 0 & 512 & 0 & \texttt{---------}\\
5 & 1 & 2 & 1 & 2 & 384 & 128 & 0 & \texttt{------++-}\\
6 & 5 & 3 & 1 & 2 & 256 & 256 & 0 & \texttt{-------+-}\\
7 & 6 & 4 & 1 & 2 & 0 & 512 & 0 & \texttt{---------}\\
8 & 0 & 1 & 1 & 1 & 0 & 512 & 0 & \texttt{---------}\\
\multicolumn{9}{@{}l@{}}{\rule{0pt}{2.6ex}\textbf{T9.164}\enspace \texttt{012342333}\enspace $n{=}9$, $|H|{=}4$, len 5, br 2, $|\mathrm{Aut}|{=}6$, orb 7/6, hist-orb 2/1; M 10752, res 3552, rev 3980: ok}\\*
0 & -- & 0 & 4 & 1 & 382 & 2 & 128 & \texttt{-+???????}\\*
1 & 0 & 1 & 4 & 1 & 252 & 4 & 256 & \texttt{--???????}\\
2 & 1 & 2 & 1 & 1 & 384 & 128 & 0 & \texttt{---++----}\\
3 & 2 & 3 & 1 & 1 & 256 & 256 & 0 & \texttt{----+----}\\
4 & 3 & 4 & 1 & 1 & 0 & 512 & 0 & \texttt{---------}\\
5 & 1 & 2 & 3 & 1 & 64 & 64 & 384 & \texttt{------???}\\
6 & 5 & 3 & 1 & 3 & 0 & 512 & 0 & \texttt{---------}\\
7 & 5 & 3 & 1 & 3 & 0 & 512 & 0 & \texttt{---------}\\
8 & 5 & 3 & 1 & 3 & 0 & 512 & 0 & \texttt{---------}\\
\multicolumn{9}{@{}l@{}}{\rule{0pt}{2.6ex}\textbf{T9.165}\enspace \texttt{012342332}\enspace $n{=}9$, $|H|{=}4$, len 5, br 2, $|\mathrm{Aut}|{=}2$, orb 8/7, hist-orb 3/2; M 10240, res 4184, rev 3308: ok}\\*
0 & -- & 0 & 4 & 1 & 326 & 2 & 184 & \texttt{-+???????}\\*
1 & 0 & 1 & 4 & 1 & 140 & 4 & 368 & \texttt{--???????}\\
2 & 1 & 2 & 1 & 1 & 384 & 128 & 0 & \texttt{---++----}\\
3 & 2 & 3 & 1 & 1 & 256 & 256 & 0 & \texttt{----+----}\\
4 & 3 & 4 & 1 & 1 & 0 & 512 & 0 & \texttt{---------}\\
5 & 1 & 2 & 2 & 1 & 128 & 128 & 256 & \texttt{------??-}\\
6 & 5 & 3 & 1 & 2 & 0 & 512 & 0 & \texttt{---------}\\
7 & 5 & 3 & 1 & 2 & 0 & 512 & 0 & \texttt{---------}\\
8 & 1 & 2 & 1 & 1 & 0 & 512 & 0 & \texttt{---------}\\
\multicolumn{9}{@{}l@{}}{\rule{0pt}{2.6ex}\textbf{T9.166}\enspace \texttt{012342331}\enspace $n{=}9$, $|H|{=}4$, len 5, br 3, $|\mathrm{Aut}|{=}2$, orb 8/7, hist-orb 3/2; M 9728, res 4048, rev 3056: ok}\\*
0 & -- & 0 & 4 & 1 & 198 & 2 & 312 & \texttt{-????????}\\*
1 & 0 & 1 & 3 & 1 & 280 & 8 & 224 & \texttt{--??????-}\\
2 & 1 & 2 & 1 & 1 & 384 & 128 & 0 & \texttt{---++----}\\
3 & 2 & 3 & 1 & 1 & 256 & 256 & 0 & \texttt{----+----}\\
4 & 3 & 4 & 1 & 1 & 0 & 512 & 0 & \texttt{---------}\\
5 & 1 & 2 & 2 & 1 & 128 & 128 & 256 & \texttt{------??-}\\
6 & 5 & 3 & 1 & 2 & 0 & 512 & 0 & \texttt{---------}\\
7 & 5 & 3 & 1 & 2 & 0 & 512 & 0 & \texttt{---------}\\
8 & 0 & 1 & 1 & 1 & 0 & 512 & 0 & \texttt{---------}\\
\multicolumn{9}{@{}l@{}}{\rule{0pt}{2.6ex}\textbf{T9.167}\enspace \texttt{012342323}\enspace $n{=}9$, $|H|{=}3$, len 5, br 1, $|\mathrm{Aut}|{=}2$, orb 7/5, hist-orb 2/1; M 10240, res 2688, rev 4100: ok}\\*
0 & -- & 0 & 3 & 1 & 382 & 2 & 128 & \texttt{-+???????}\\*
1 & 0 & 1 & 3 & 1 & 252 & 4 & 256 & \texttt{--???????}\\
2 & 1 & 2 & 1 & 1 & 384 & 128 & 0 & \texttt{---++----}\\
3 & 2 & 3 & 1 & 1 & 256 & 256 & 0 & \texttt{----+----}\\
4 & 3 & 4 & 1 & 1 & 0 & 512 & 0 & \texttt{---------}\\
5 & 1 & 2 & 1 & 2 & 256 & 256 & 0 & \texttt{------+--}\\
6 & 5 & 3 & 1 & 2 & 0 & 512 & 0 & \texttt{---------}\\
7 & 1 & 2 & 1 & 2 & 256 & 256 & 0 & \texttt{--------+}\\
8 & 7 & 3 & 1 & 2 & 0 & 512 & 0 & \texttt{---------}\\
\multicolumn{9}{@{}l@{}}{\rule{0pt}{2.6ex}\textbf{T9.168}\enspace \texttt{012342322}\enspace $n{=}9$, $|H|{=}4$, len 5, br 1, $|\mathrm{Aut}|{=}2$, orb 8/7, hist-orb 3/2; M 9728, res 4452, rev 2988: ok}\\*
0 & -- & 0 & 4 & 1 & 298 & 2 & 212 & \texttt{-+???????}\\*
1 & 0 & 1 & 4 & 1 & 84 & 4 & 424 & \texttt{--???????}\\
2 & 1 & 2 & 1 & 1 & 384 & 128 & 0 & \texttt{---++----}\\
3 & 2 & 3 & 1 & 1 & 256 & 256 & 0 & \texttt{----+----}\\
4 & 3 & 4 & 1 & 1 & 0 & 512 & 0 & \texttt{---------}\\
5 & 1 & 2 & 1 & 1 & 256 & 256 & 0 & \texttt{------+--}\\
6 & 5 & 3 & 1 & 1 & 0 & 512 & 0 & \texttt{---------}\\
7 & 1 & 2 & 1 & 2 & 0 & 512 & 0 & \texttt{---------}\\
8 & 1 & 2 & 1 & 2 & 0 & 512 & 0 & \texttt{---------}\\
\multicolumn{9}{@{}l@{}}{\rule{0pt}{2.6ex}\textbf{T9.169}\enspace \texttt{012342321}\enspace $n{=}9$, $|H|{=}4$, len 5, br 2, $|\mathrm{Aut}|{=}1$, orb 9/9, hist-orb 4/4; M 9216, res 4568, rev 2560: ok}\\*
0 & -- & 0 & 4 & 1 & 170 & 2 & 340 & \texttt{-????????}\\*
1 & 0 & 1 & 3 & 1 & 168 & 8 & 336 & \texttt{--??????-}\\
2 & 1 & 2 & 1 & 1 & 384 & 128 & 0 & \texttt{---++----}\\
3 & 2 & 3 & 1 & 1 & 256 & 256 & 0 & \texttt{----+----}\\
4 & 3 & 4 & 1 & 1 & 0 & 512 & 0 & \texttt{---------}\\
5 & 1 & 2 & 1 & 1 & 256 & 256 & 0 & \texttt{------+--}\\
6 & 5 & 3 & 1 & 1 & 0 & 512 & 0 & \texttt{---------}\\
7 & 1 & 2 & 1 & 1 & 0 & 512 & 0 & \texttt{---------}\\
8 & 0 & 1 & 1 & 1 & 0 & 512 & 0 & \texttt{---------}\\
\multicolumn{9}{@{}l@{}}{\rule{0pt}{2.6ex}\textbf{T9.170}\enspace \texttt{012342312}\enspace $n{=}9$, $|H|{=}3$, len 5, br 2, $|\mathrm{Aut}|{=}1$, orb 9/9, hist-orb 3/3; M 9216, res 2256, rev 3676: ok}\\*
0 & -- & 0 & 3 & 1 & 318 & 2 & 192 & \texttt{-????????}\\*
1 & 0 & 1 & 2 & 1 & 336 & 16 & 160 & \texttt{--?????--}\\
2 & 1 & 2 & 1 & 1 & 384 & 128 & 0 & \texttt{---++----}\\
3 & 2 & 3 & 1 & 1 & 256 & 256 & 0 & \texttt{----+----}\\
4 & 3 & 4 & 1 & 1 & 0 & 512 & 0 & \texttt{---------}\\
5 & 1 & 2 & 1 & 1 & 256 & 256 & 0 & \texttt{------+--}\\
6 & 5 & 3 & 1 & 1 & 0 & 512 & 0 & \texttt{---------}\\
7 & 0 & 1 & 1 & 1 & 256 & 256 & 0 & \texttt{--------+}\\
8 & 7 & 2 & 1 & 1 & 0 & 512 & 0 & \texttt{---------}\\
\multicolumn{9}{@{}l@{}}{\rule{0pt}{2.6ex}\textbf{T9.171}\enspace \texttt{012342311}\enspace $n{=}9$, $|H|{=}4$, len 5, br 2, $|\mathrm{Aut}|{=}2$, orb 8/7, hist-orb 3/2; M 8704, res 3952, rev 2544: ok}\\*
0 & -- & 0 & 4 & 1 & 106 & 2 & 404 & \texttt{-????????}\\*
1 & 0 & 1 & 2 & 1 & 336 & 16 & 160 & \texttt{--?????--}\\
2 & 1 & 2 & 1 & 1 & 384 & 128 & 0 & \texttt{---++----}\\
3 & 2 & 3 & 1 & 1 & 256 & 256 & 0 & \texttt{----+----}\\
4 & 3 & 4 & 1 & 1 & 0 & 512 & 0 & \texttt{---------}\\
5 & 1 & 2 & 1 & 1 & 256 & 256 & 0 & \texttt{------+--}\\
6 & 5 & 3 & 1 & 1 & 0 & 512 & 0 & \texttt{---------}\\
7 & 0 & 1 & 1 & 2 & 0 & 512 & 0 & \texttt{---------}\\
8 & 0 & 1 & 1 & 2 & 0 & 512 & 0 & \texttt{---------}\\
\multicolumn{9}{@{}l@{}}{\rule{0pt}{2.6ex}\textbf{T9.172}\enspace \texttt{012342222}\enspace $n{=}9$, $|H|{=}5$, len 5, br 1, $|\mathrm{Aut}|{=}24$, orb 6/5, hist-orb 2/1; M 9216, res 5040, rev 2476: ok}\\*
0 & -- & 0 & 5 & 1 & 270 & 2 & 240 & \texttt{-+???????}\\*
1 & 0 & 1 & 5 & 1 & 28 & 4 & 480 & \texttt{--???????}\\
2 & 1 & 2 & 1 & 1 & 384 & 128 & 0 & \texttt{---++----}\\
3 & 2 & 3 & 1 & 1 & 256 & 256 & 0 & \texttt{----+----}\\
4 & 3 & 4 & 1 & 1 & 0 & 512 & 0 & \texttt{---------}\\
5 & 1 & 2 & 1 & 4 & 0 & 512 & 0 & \texttt{---------}\\
6 & 1 & 2 & 1 & 4 & 0 & 512 & 0 & \texttt{---------}\\
7 & 1 & 2 & 1 & 4 & 0 & 512 & 0 & \texttt{---------}\\
8 & 1 & 2 & 1 & 4 & 0 & 512 & 0 & \texttt{---------}\\
\multicolumn{9}{@{}l@{}}{\rule{0pt}{2.6ex}\textbf{T9.173}\enspace \texttt{012342221}\enspace $n{=}9$, $|H|{=}5$, len 5, br 2, $|\mathrm{Aut}|{=}6$, orb 7/6, hist-orb 3/2; M 8704, res 5408, rev 1840: ok}\\*
0 & -- & 0 & 5 & 1 & 142 & 2 & 368 & \texttt{-????????}\\*
1 & 0 & 1 & 4 & 1 & 56 & 8 & 448 & \texttt{--??????-}\\
2 & 1 & 2 & 1 & 1 & 384 & 128 & 0 & \texttt{---++----}\\
3 & 2 & 3 & 1 & 1 & 256 & 256 & 0 & \texttt{----+----}\\
4 & 3 & 4 & 1 & 1 & 0 & 512 & 0 & \texttt{---------}\\
5 & 1 & 2 & 1 & 3 & 0 & 512 & 0 & \texttt{---------}\\
6 & 1 & 2 & 1 & 3 & 0 & 512 & 0 & \texttt{---------}\\
7 & 1 & 2 & 1 & 3 & 0 & 512 & 0 & \texttt{---------}\\
8 & 0 & 1 & 1 & 1 & 0 & 512 & 0 & \texttt{---------}\\
\multicolumn{9}{@{}l@{}}{\rule{0pt}{2.6ex}\textbf{T9.174}\enspace \texttt{012342212}\enspace $n{=}9$, $|H|{=}4$, len 5, br 2, $|\mathrm{Aut}|{=}2$, orb 8/7, hist-orb 3/2; M 8704, res 3936, rev 2612: ok}\\*
0 & -- & 0 & 4 & 1 & 234 & 2 & 276 & \texttt{-????????}\\*
1 & 0 & 1 & 3 & 1 & 112 & 16 & 384 & \texttt{--?????--}\\
2 & 1 & 2 & 1 & 1 & 384 & 128 & 0 & \texttt{---++----}\\
3 & 2 & 3 & 1 & 1 & 256 & 256 & 0 & \texttt{----+----}\\
4 & 3 & 4 & 1 & 1 & 0 & 512 & 0 & \texttt{---------}\\
5 & 1 & 2 & 1 & 2 & 0 & 512 & 0 & \texttt{---------}\\
6 & 1 & 2 & 1 & 2 & 0 & 512 & 0 & \texttt{---------}\\
7 & 0 & 1 & 1 & 1 & 256 & 256 & 0 & \texttt{--------+}\\
8 & 7 & 2 & 1 & 1 & 0 & 512 & 0 & \texttt{---------}\\
\multicolumn{9}{@{}l@{}}{\rule{0pt}{2.6ex}\textbf{T9.175}\enspace \texttt{012342211}\enspace $n{=}9$, $|H|{=}5$, len 5, br 2, $|\mathrm{Aut}|{=}4$, orb 7/5, hist-orb 3/1; M 8192, res 5184, rev 1616: ok}\\*
0 & -- & 0 & 5 & 1 & 78 & 2 & 432 & \texttt{-????????}\\*
1 & 0 & 1 & 3 & 1 & 112 & 16 & 384 & \texttt{--?????--}\\
2 & 1 & 2 & 1 & 1 & 384 & 128 & 0 & \texttt{---++----}\\
3 & 2 & 3 & 1 & 1 & 256 & 256 & 0 & \texttt{----+----}\\
4 & 3 & 4 & 1 & 1 & 0 & 512 & 0 & \texttt{---------}\\
5 & 1 & 2 & 1 & 2 & 0 & 512 & 0 & \texttt{---------}\\
6 & 1 & 2 & 1 & 2 & 0 & 512 & 0 & \texttt{---------}\\
7 & 0 & 1 & 1 & 2 & 0 & 512 & 0 & \texttt{---------}\\
8 & 0 & 1 & 1 & 2 & 0 & 512 & 0 & \texttt{---------}\\
\multicolumn{9}{@{}l@{}}{\rule{0pt}{2.6ex}\textbf{T9.176}\enspace \texttt{012342123}\enspace $n{=}9$, $|H|{=}3$, len 5, br 2, $|\mathrm{Aut}|{=}1$, orb 9/9, hist-orb 3/3; M 9216, res 2400, rev 3500: ok}\\*
0 & -- & 0 & 3 & 1 & 322 & 2 & 188 & \texttt{-????????}\\*
1 & 0 & 1 & 2 & 1 & 224 & 32 & 256 & \texttt{--????---}\\
2 & 1 & 2 & 1 & 1 & 384 & 128 & 0 & \texttt{---++----}\\
3 & 2 & 3 & 1 & 1 & 256 & 256 & 0 & \texttt{----+----}\\
4 & 3 & 4 & 1 & 1 & 0 & 512 & 0 & \texttt{---------}\\
5 & 1 & 2 & 1 & 1 & 0 & 512 & 0 & \texttt{---------}\\
6 & 0 & 1 & 1 & 1 & 384 & 128 & 0 & \texttt{-------++}\\
7 & 6 & 2 & 1 & 1 & 256 & 256 & 0 & \texttt{--------+}\\
8 & 7 & 3 & 1 & 1 & 0 & 512 & 0 & \texttt{---------}\\
\multicolumn{9}{@{}l@{}}{\rule{0pt}{2.6ex}\textbf{T9.177}\enspace \texttt{012342122}\enspace $n{=}9$, $|H|{=}4$, len 5, br 3, $|\mathrm{Aut}|{=}2$, orb 8/7, hist-orb 3/2; M 8704, res 3520, rev 2736: ok}\\*
0 & -- & 0 & 4 & 1 & 230 & 2 & 280 & \texttt{-????????}\\*
1 & 0 & 1 & 2 & 1 & 224 & 32 & 256 & \texttt{--????---}\\
2 & 1 & 2 & 1 & 1 & 384 & 128 & 0 & \texttt{---++----}\\
3 & 2 & 3 & 1 & 1 & 256 & 256 & 0 & \texttt{----+----}\\
4 & 3 & 4 & 1 & 1 & 0 & 512 & 0 & \texttt{---------}\\
5 & 1 & 2 & 1 & 1 & 0 & 512 & 0 & \texttt{---------}\\
6 & 0 & 1 & 2 & 1 & 128 & 128 & 256 & \texttt{-------??}\\
7 & 6 & 2 & 1 & 2 & 0 & 512 & 0 & \texttt{---------}\\
8 & 6 & 2 & 1 & 2 & 0 & 512 & 0 & \texttt{---------}\\
\multicolumn{9}{@{}l@{}}{\rule{0pt}{2.6ex}\textbf{T9.178}\enspace \texttt{012342121}\enspace $n{=}9$, $|H|{=}4$, len 5, br 2, $|\mathrm{Aut}|{=}1$, orb 9/9, hist-orb 4/4; M 8192, res 3872, rev 2300: ok}\\*
0 & -- & 0 & 4 & 1 & 138 & 2 & 372 & \texttt{-????????}\\*
1 & 0 & 1 & 2 & 1 & 224 & 32 & 256 & \texttt{--????---}\\
2 & 1 & 2 & 1 & 1 & 384 & 128 & 0 & \texttt{---++----}\\
3 & 2 & 3 & 1 & 1 & 256 & 256 & 0 & \texttt{----+----}\\
4 & 3 & 4 & 1 & 1 & 0 & 512 & 0 & \texttt{---------}\\
5 & 1 & 2 & 1 & 1 & 0 & 512 & 0 & \texttt{---------}\\
6 & 0 & 1 & 1 & 1 & 256 & 256 & 0 & \texttt{-------+-}\\
7 & 6 & 2 & 1 & 1 & 0 & 512 & 0 & \texttt{---------}\\
8 & 0 & 1 & 1 & 1 & 0 & 512 & 0 & \texttt{---------}\\
\multicolumn{9}{@{}l@{}}{\rule{0pt}{2.6ex}\textbf{T9.179}\enspace \texttt{012342111}\enspace $n{=}9$, $|H|{=}5$, len 5, br 2, $|\mathrm{Aut}|{=}6$, orb 7/6, hist-orb 3/2; M 7680, res 4608, rev 1608: ok}\\*
0 & -- & 0 & 5 & 1 & 46 & 2 & 464 & \texttt{-????????}\\*
1 & 0 & 1 & 2 & 1 & 224 & 32 & 256 & \texttt{--????---}\\
2 & 1 & 2 & 1 & 1 & 384 & 128 & 0 & \texttt{---++----}\\
3 & 2 & 3 & 1 & 1 & 256 & 256 & 0 & \texttt{----+----}\\
4 & 3 & 4 & 1 & 1 & 0 & 512 & 0 & \texttt{---------}\\
5 & 1 & 2 & 1 & 1 & 0 & 512 & 0 & \texttt{---------}\\
6 & 0 & 1 & 1 & 3 & 0 & 512 & 0 & \texttt{---------}\\
7 & 0 & 1 & 1 & 3 & 0 & 512 & 0 & \texttt{---------}\\
8 & 0 & 1 & 1 & 3 & 0 & 512 & 0 & \texttt{---------}\\
\multicolumn{9}{@{}l@{}}{\rule{0pt}{2.6ex}\textbf{T9.180}\enspace \texttt{012341234}\enspace $n{=}9$, $|H|{=}2$, len 5, br 1, $|\mathrm{Aut}|{=}2$, orb 5/1, hist-orb 1/0; M 10240, res 480, rev 4504: ok}\\*
0 & -- & 0 & 2 & 1 & 450 & 2 & 60 & \texttt{-????????}\\*
1 & 0 & 1 & 1 & 2 & 448 & 64 & 0 & \texttt{--+++----}\\
2 & 1 & 2 & 1 & 2 & 384 & 128 & 0 & \texttt{---++----}\\
3 & 2 & 3 & 1 & 2 & 256 & 256 & 0 & \texttt{----+----}\\
4 & 3 & 4 & 1 & 2 & 0 & 512 & 0 & \texttt{---------}\\
5 & 0 & 1 & 1 & 2 & 448 & 64 & 0 & \texttt{------+++}\\
6 & 5 & 2 & 1 & 2 & 384 & 128 & 0 & \texttt{-------++}\\
7 & 6 & 3 & 1 & 2 & 256 & 256 & 0 & \texttt{--------+}\\
8 & 7 & 4 & 1 & 2 & 0 & 512 & 0 & \texttt{---------}\\
\multicolumn{9}{@{}l@{}}{\rule{0pt}{2.6ex}\textbf{T9.181}\enspace \texttt{012341233}\enspace $n{=}9$, $|H|{=}3$, len 5, br 2, $|\mathrm{Aut}|{=}2$, orb 8/7, hist-orb 2/1; M 9728, res 1600, rev 4008: ok}\\*
0 & -- & 0 & 3 & 1 & 390 & 2 & 120 & \texttt{-????????}\\*
1 & 0 & 1 & 1 & 1 & 448 & 64 & 0 & \texttt{--+++----}\\
2 & 1 & 2 & 1 & 1 & 384 & 128 & 0 & \texttt{---++----}\\
3 & 2 & 3 & 1 & 1 & 256 & 256 & 0 & \texttt{----+----}\\
4 & 3 & 4 & 1 & 1 & 0 & 512 & 0 & \texttt{---------}\\
5 & 0 & 1 & 2 & 1 & 320 & 64 & 128 & \texttt{------+??}\\
6 & 5 & 2 & 2 & 1 & 128 & 128 & 256 & \texttt{-------??}\\
7 & 6 & 3 & 1 & 2 & 0 & 512 & 0 & \texttt{---------}\\
8 & 6 & 3 & 1 & 2 & 0 & 512 & 0 & \texttt{---------}\\
\multicolumn{9}{@{}l@{}}{\rule{0pt}{2.6ex}\textbf{T9.182}\enspace \texttt{012341232}\enspace $n{=}9$, $|H|{=}3$, len 5, br 2, $|\mathrm{Aut}|{=}1$, orb 9/9, hist-orb 3/3; M 9216, res 2080, rev 3588: ok}\\*
0 & -- & 0 & 3 & 1 & 330 & 2 & 180 & \texttt{-????????}\\*
1 & 0 & 1 & 1 & 1 & 448 & 64 & 0 & \texttt{--+++----}\\
2 & 1 & 2 & 1 & 1 & 384 & 128 & 0 & \texttt{---++----}\\
3 & 2 & 3 & 1 & 1 & 256 & 256 & 0 & \texttt{----+----}\\
4 & 3 & 4 & 1 & 1 & 0 & 512 & 0 & \texttt{---------}\\
5 & 0 & 1 & 2 & 1 & 192 & 64 & 256 & \texttt{------???}\\
6 & 5 & 2 & 1 & 1 & 256 & 256 & 0 & \texttt{-------+-}\\
7 & 6 & 3 & 1 & 1 & 0 & 512 & 0 & \texttt{---------}\\
8 & 5 & 2 & 1 & 1 & 0 & 512 & 0 & \texttt{---------}\\
\multicolumn{9}{@{}l@{}}{\rule{0pt}{2.6ex}\textbf{T9.183}\enspace \texttt{012341231}\enspace $n{=}9$, $|H|{=}3$, len 5, br 1, $|\mathrm{Aut}|{=}1$, orb 9/9, hist-orb 3/3; M 8704, res 2400, rev 3144: ok}\\*
0 & -- & 0 & 3 & 1 & 210 & 2 & 300 & \texttt{-????????}\\*
1 & 0 & 1 & 1 & 1 & 448 & 64 & 0 & \texttt{--+++----}\\
2 & 1 & 2 & 1 & 1 & 384 & 128 & 0 & \texttt{---++----}\\
3 & 2 & 3 & 1 & 1 & 256 & 256 & 0 & \texttt{----+----}\\
4 & 3 & 4 & 1 & 1 & 0 & 512 & 0 & \texttt{---------}\\
5 & 0 & 1 & 1 & 1 & 384 & 128 & 0 & \texttt{------++-}\\
6 & 5 & 2 & 1 & 1 & 256 & 256 & 0 & \texttt{-------+-}\\
7 & 6 & 3 & 1 & 1 & 0 & 512 & 0 & \texttt{---------}\\
8 & 0 & 1 & 1 & 1 & 0 & 512 & 0 & \texttt{---------}\\
\multicolumn{9}{@{}l@{}}{\rule{0pt}{2.6ex}\textbf{T9.184}\enspace \texttt{012341222}\enspace $n{=}9$, $|H|{=}4$, len 5, br 2, $|\mathrm{Aut}|{=}6$, orb 7/6, hist-orb 2/1; M 8704, res 2880, rev 2972: ok}\\*
0 & -- & 0 & 4 & 1 & 270 & 2 & 240 & \texttt{-????????}\\*
1 & 0 & 1 & 1 & 1 & 448 & 64 & 0 & \texttt{--+++----}\\
2 & 1 & 2 & 1 & 1 & 384 & 128 & 0 & \texttt{---++----}\\
3 & 2 & 3 & 1 & 1 & 256 & 256 & 0 & \texttt{----+----}\\
4 & 3 & 4 & 1 & 1 & 0 & 512 & 0 & \texttt{---------}\\
5 & 0 & 1 & 3 & 1 & 64 & 64 & 384 & \texttt{------???}\\
6 & 5 & 2 & 1 & 3 & 0 & 512 & 0 & \texttt{---------}\\
7 & 5 & 2 & 1 & 3 & 0 & 512 & 0 & \texttt{---------}\\
8 & 5 & 2 & 1 & 3 & 0 & 512 & 0 & \texttt{---------}\\
\multicolumn{9}{@{}l@{}}{\rule{0pt}{2.6ex}\textbf{T9.185}\enspace \texttt{012341221}\enspace $n{=}9$, $|H|{=}4$, len 5, br 2, $|\mathrm{Aut}|{=}2$, orb 8/7, hist-orb 3/2; M 8192, res 3264, rev 2468: ok}\\*
0 & -- & 0 & 4 & 1 & 150 & 2 & 360 & \texttt{-????????}\\*
1 & 0 & 1 & 1 & 1 & 448 & 64 & 0 & \texttt{--+++----}\\
2 & 1 & 2 & 1 & 1 & 384 & 128 & 0 & \texttt{---++----}\\
3 & 2 & 3 & 1 & 1 & 256 & 256 & 0 & \texttt{----+----}\\
4 & 3 & 4 & 1 & 1 & 0 & 512 & 0 & \texttt{---------}\\
5 & 0 & 1 & 2 & 1 & 128 & 128 & 256 & \texttt{------??-}\\
6 & 5 & 2 & 1 & 2 & 0 & 512 & 0 & \texttt{---------}\\
7 & 5 & 2 & 1 & 2 & 0 & 512 & 0 & \texttt{---------}\\
8 & 0 & 1 & 1 & 1 & 0 & 512 & 0 & \texttt{---------}\\
\multicolumn{9}{@{}l@{}}{\rule{0pt}{2.6ex}\textbf{T9.186}\enspace \texttt{012341212}\enspace $n{=}9$, $|H|{=}3$, len 5, br 1, $|\mathrm{Aut}|{=}2$, orb 7/5, hist-orb 2/1; M 8192, res 1920, rev 3172: ok}\\*
0 & -- & 0 & 3 & 1 & 270 & 2 & 240 & \texttt{-????????}\\*
1 & 0 & 1 & 1 & 1 & 448 & 64 & 0 & \texttt{--+++----}\\
2 & 1 & 2 & 1 & 1 & 384 & 128 & 0 & \texttt{---++----}\\
3 & 2 & 3 & 1 & 1 & 256 & 256 & 0 & \texttt{----+----}\\
4 & 3 & 4 & 1 & 1 & 0 & 512 & 0 & \texttt{---------}\\
5 & 0 & 1 & 1 & 2 & 256 & 256 & 0 & \texttt{------+--}\\
6 & 5 & 2 & 1 & 2 & 0 & 512 & 0 & \texttt{---------}\\
7 & 0 & 1 & 1 & 2 & 256 & 256 & 0 & \texttt{--------+}\\
8 & 7 & 2 & 1 & 2 & 0 & 512 & 0 & \texttt{---------}\\
\multicolumn{9}{@{}l@{}}{\rule{0pt}{2.6ex}\textbf{T9.187}\enspace \texttt{012341211}\enspace $n{=}9$, $|H|{=}4$, len 5, br 1, $|\mathrm{Aut}|{=}2$, orb 8/7, hist-orb 3/2; M 7680, res 3360, rev 2184: ok}\\*
0 & -- & 0 & 4 & 1 & 90 & 2 & 420 & \texttt{-????????}\\*
1 & 0 & 1 & 1 & 1 & 448 & 64 & 0 & \texttt{--+++----}\\
2 & 1 & 2 & 1 & 1 & 384 & 128 & 0 & \texttt{---++----}\\
3 & 2 & 3 & 1 & 1 & 256 & 256 & 0 & \texttt{----+----}\\
4 & 3 & 4 & 1 & 1 & 0 & 512 & 0 & \texttt{---------}\\
5 & 0 & 1 & 1 & 1 & 256 & 256 & 0 & \texttt{------+--}\\
6 & 5 & 2 & 1 & 1 & 0 & 512 & 0 & \texttt{---------}\\
7 & 0 & 1 & 1 & 2 & 0 & 512 & 0 & \texttt{---------}\\
8 & 0 & 1 & 1 & 2 & 0 & 512 & 0 & \texttt{---------}\\
\multicolumn{9}{@{}l@{}}{\rule{0pt}{2.6ex}\textbf{T9.188}\enspace \texttt{012341111}\enspace $n{=}9$, $|H|{=}5$, len 5, br 1, $|\mathrm{Aut}|{=}24$, orb 6/5, hist-orb 2/1; M 7168, res 3840, rev 1644: ok}\\*
0 & -- & 0 & 5 & 1 & 30 & 2 & 480 & \texttt{-????????}\\*
1 & 0 & 1 & 1 & 1 & 448 & 64 & 0 & \texttt{--+++----}\\
2 & 1 & 2 & 1 & 1 & 384 & 128 & 0 & \texttt{---++----}\\
3 & 2 & 3 & 1 & 1 & 256 & 256 & 0 & \texttt{----+----}\\
4 & 3 & 4 & 1 & 1 & 0 & 512 & 0 & \texttt{---------}\\
5 & 0 & 1 & 1 & 4 & 0 & 512 & 0 & \texttt{---------}\\
6 & 0 & 1 & 1 & 4 & 0 & 512 & 0 & \texttt{---------}\\
7 & 0 & 1 & 1 & 4 & 0 & 512 & 0 & \texttt{---------}\\
8 & 0 & 1 & 1 & 4 & 0 & 512 & 0 & \texttt{---------}\\
\multicolumn{9}{@{}l@{}}{\rule{0pt}{2.6ex}\textbf{T9.189}\enspace \texttt{012333333}\enspace $n{=}9$, $|H|{=}6$, len 4, br 1, $|\mathrm{Aut}|{=}720$, orb 4/3, hist-orb 1/0; M 10752, res 5208, rev 3936: ok}\\*
0 & -- & 0 & 6 & 1 & 386 & 2 & 124 & \texttt{-++??????}\\*
1 & 0 & 1 & 6 & 1 & 260 & 4 & 248 & \texttt{--+??????}\\
2 & 1 & 2 & 6 & 1 & 8 & 8 & 496 & \texttt{---??????}\\
3 & 2 & 3 & 1 & 6 & 0 & 512 & 0 & \texttt{---------}\\
4 & 2 & 3 & 1 & 6 & 0 & 512 & 0 & \texttt{---------}\\
5 & 2 & 3 & 1 & 6 & 0 & 512 & 0 & \texttt{---------}\\
6 & 2 & 3 & 1 & 6 & 0 & 512 & 0 & \texttt{---------}\\
7 & 2 & 3 & 1 & 6 & 0 & 512 & 0 & \texttt{---------}\\
8 & 2 & 3 & 1 & 6 & 0 & 512 & 0 & \texttt{---------}\\
\multicolumn{9}{@{}l@{}}{\rule{0pt}{2.6ex}\textbf{T9.190}\enspace \texttt{012333332}\enspace $n{=}9$, $|H|{=}6$, len 4, br 2, $|\mathrm{Aut}|{=}120$, orb 5/4, hist-orb 2/1; M 10240, res 5988, rev 2820: ok}\\*
0 & -- & 0 & 6 & 1 & 322 & 2 & 188 & \texttt{-+???????}\\*
1 & 0 & 1 & 6 & 1 & 132 & 4 & 376 & \texttt{--???????}\\
2 & 1 & 2 & 5 & 1 & 16 & 16 & 480 & \texttt{---?????-}\\
3 & 2 & 3 & 1 & 5 & 0 & 512 & 0 & \texttt{---------}\\
4 & 2 & 3 & 1 & 5 & 0 & 512 & 0 & \texttt{---------}\\
5 & 2 & 3 & 1 & 5 & 0 & 512 & 0 & \texttt{---------}\\
6 & 2 & 3 & 1 & 5 & 0 & 512 & 0 & \texttt{---------}\\
7 & 2 & 3 & 1 & 5 & 0 & 512 & 0 & \texttt{---------}\\
8 & 1 & 2 & 1 & 1 & 0 & 512 & 0 & \texttt{---------}\\
\multicolumn{9}{@{}l@{}}{\rule{0pt}{2.6ex}\textbf{T9.191}\enspace \texttt{012333331}\enspace $n{=}9$, $|H|{=}6$, len 4, br 2, $|\mathrm{Aut}|{=}120$, orb 5/4, hist-orb 2/1; M 9728, res 5888, rev 2580: ok}\\*
0 & -- & 0 & 6 & 1 & 194 & 2 & 316 & \texttt{-????????}\\*
1 & 0 & 1 & 5 & 1 & 264 & 8 & 240 & \texttt{--+?????-}\\
2 & 1 & 2 & 5 & 1 & 16 & 16 & 480 & \texttt{---?????-}\\
3 & 2 & 3 & 1 & 5 & 0 & 512 & 0 & \texttt{---------}\\
4 & 2 & 3 & 1 & 5 & 0 & 512 & 0 & \texttt{---------}\\
5 & 2 & 3 & 1 & 5 & 0 & 512 & 0 & \texttt{---------}\\
6 & 2 & 3 & 1 & 5 & 0 & 512 & 0 & \texttt{---------}\\
7 & 2 & 3 & 1 & 5 & 0 & 512 & 0 & \texttt{---------}\\
8 & 0 & 1 & 1 & 1 & 0 & 512 & 0 & \texttt{---------}\\
\multicolumn{9}{@{}l@{}}{\rule{0pt}{2.6ex}\textbf{T9.192}\enspace \texttt{012333323}\enspace $n{=}9$, $|H|{=}5$, len 4, br 2, $|\mathrm{Aut}|{=}24$, orb 6/5, hist-orb 2/1; M 10240, res 4648, rev 3460: ok}\\*
0 & -- & 0 & 5 & 1 & 358 & 2 & 152 & \texttt{-+???????}\\*
1 & 0 & 1 & 5 & 1 & 204 & 4 & 304 & \texttt{--???????}\\
2 & 1 & 2 & 4 & 1 & 32 & 32 & 448 & \texttt{---????--}\\
3 & 2 & 3 & 1 & 4 & 0 & 512 & 0 & \texttt{---------}\\
4 & 2 & 3 & 1 & 4 & 0 & 512 & 0 & \texttt{---------}\\
5 & 2 & 3 & 1 & 4 & 0 & 512 & 0 & \texttt{---------}\\
6 & 2 & 3 & 1 & 4 & 0 & 512 & 0 & \texttt{---------}\\
7 & 1 & 2 & 1 & 1 & 256 & 256 & 0 & \texttt{--------+}\\
8 & 7 & 3 & 1 & 1 & 0 & 512 & 0 & \texttt{---------}\\
\multicolumn{9}{@{}l@{}}{\rule{0pt}{2.6ex}\textbf{T9.193}\enspace \texttt{012333322}\enspace $n{=}9$, $|H|{=}6$, len 4, br 2, $|\mathrm{Aut}|{=}48$, orb 5/3, hist-orb 2/0; M 9728, res 6076, rev 2300: ok}\\*
0 & -- & 0 & 6 & 1 & 290 & 2 & 220 & \texttt{-+???????}\\*
1 & 0 & 1 & 6 & 1 & 68 & 4 & 440 & \texttt{--???????}\\
2 & 1 & 2 & 4 & 1 & 32 & 32 & 448 & \texttt{---????--}\\
3 & 2 & 3 & 1 & 4 & 0 & 512 & 0 & \texttt{---------}\\
4 & 2 & 3 & 1 & 4 & 0 & 512 & 0 & \texttt{---------}\\
5 & 2 & 3 & 1 & 4 & 0 & 512 & 0 & \texttt{---------}\\
6 & 2 & 3 & 1 & 4 & 0 & 512 & 0 & \texttt{---------}\\
7 & 1 & 2 & 1 & 2 & 0 & 512 & 0 & \texttt{---------}\\
8 & 1 & 2 & 1 & 2 & 0 & 512 & 0 & \texttt{---------}\\
\multicolumn{9}{@{}l@{}}{\rule{0pt}{2.6ex}\textbf{T9.194}\enspace \texttt{012333321}\enspace $n{=}9$, $|H|{=}6$, len 4, br 3, $|\mathrm{Aut}|{=}24$, orb 6/5, hist-orb 3/2; M 9216, res 6264, rev 1800: ok}\\*
0 & -- & 0 & 6 & 1 & 162 & 2 & 348 & \texttt{-????????}\\*
1 & 0 & 1 & 5 & 1 & 136 & 8 & 368 & \texttt{--??????-}\\
2 & 1 & 2 & 4 & 1 & 32 & 32 & 448 & \texttt{---????--}\\
3 & 2 & 3 & 1 & 4 & 0 & 512 & 0 & \texttt{---------}\\
4 & 2 & 3 & 1 & 4 & 0 & 512 & 0 & \texttt{---------}\\
5 & 2 & 3 & 1 & 4 & 0 & 512 & 0 & \texttt{---------}\\
6 & 2 & 3 & 1 & 4 & 0 & 512 & 0 & \texttt{---------}\\
7 & 1 & 2 & 1 & 1 & 0 & 512 & 0 & \texttt{---------}\\
8 & 0 & 1 & 1 & 1 & 0 & 512 & 0 & \texttt{---------}\\
\multicolumn{9}{@{}l@{}}{\rule{0pt}{2.6ex}\textbf{T9.195}\enspace \texttt{012333312}\enspace $n{=}9$, $|H|{=}5$, len 4, br 2, $|\mathrm{Aut}|{=}24$, orb 6/5, hist-orb 2/1; M 9216, res 4192, rev 3080: ok}\\*
0 & -- & 0 & 5 & 1 & 294 & 2 & 216 & \texttt{-????????}\\*
1 & 0 & 1 & 4 & 1 & 272 & 16 & 224 & \texttt{--+????--}\\
2 & 1 & 2 & 4 & 1 & 32 & 32 & 448 & \texttt{---????--}\\
3 & 2 & 3 & 1 & 4 & 0 & 512 & 0 & \texttt{---------}\\
4 & 2 & 3 & 1 & 4 & 0 & 512 & 0 & \texttt{---------}\\
5 & 2 & 3 & 1 & 4 & 0 & 512 & 0 & \texttt{---------}\\
6 & 2 & 3 & 1 & 4 & 0 & 512 & 0 & \texttt{---------}\\
7 & 0 & 1 & 1 & 1 & 256 & 256 & 0 & \texttt{--------+}\\
8 & 7 & 2 & 1 & 1 & 0 & 512 & 0 & \texttt{---------}\\
\multicolumn{9}{@{}l@{}}{\rule{0pt}{2.6ex}\textbf{T9.196}\enspace \texttt{012333311}\enspace $n{=}9$, $|H|{=}6$, len 4, br 2, $|\mathrm{Aut}|{=}48$, orb 5/3, hist-orb 2/0; M 8704, res 5760, rev 1828: ok}\\*
0 & -- & 0 & 6 & 1 & 98 & 2 & 412 & \texttt{-????????}\\*
1 & 0 & 1 & 4 & 1 & 272 & 16 & 224 & \texttt{--+????--}\\
2 & 1 & 2 & 4 & 1 & 32 & 32 & 448 & \texttt{---????--}\\
3 & 2 & 3 & 1 & 4 & 0 & 512 & 0 & \texttt{---------}\\
4 & 2 & 3 & 1 & 4 & 0 & 512 & 0 & \texttt{---------}\\
5 & 2 & 3 & 1 & 4 & 0 & 512 & 0 & \texttt{---------}\\
6 & 2 & 3 & 1 & 4 & 0 & 512 & 0 & \texttt{---------}\\
7 & 0 & 1 & 1 & 2 & 0 & 512 & 0 & \texttt{---------}\\
8 & 0 & 1 & 1 & 2 & 0 & 512 & 0 & \texttt{---------}\\
\multicolumn{9}{@{}l@{}}{\rule{0pt}{2.6ex}\textbf{T9.197}\enspace \texttt{012333233}\enspace $n{=}9$, $|H|{=}5$, len 4, br 3, $|\mathrm{Aut}|{=}12$, orb 6/4, hist-orb 2/0; M 10240, res 4628, rev 3260: ok}\\*
0 & -- & 0 & 5 & 1 & 346 & 2 & 164 & \texttt{-+???????}\\*
1 & 0 & 1 & 5 & 1 & 180 & 4 & 328 & \texttt{--???????}\\
2 & 1 & 2 & 3 & 1 & 64 & 64 & 384 & \texttt{---???---}\\
3 & 2 & 3 & 1 & 3 & 0 & 512 & 0 & \texttt{---------}\\
4 & 2 & 3 & 1 & 3 & 0 & 512 & 0 & \texttt{---------}\\
5 & 2 & 3 & 1 & 3 & 0 & 512 & 0 & \texttt{---------}\\
6 & 1 & 2 & 2 & 1 & 128 & 128 & 256 & \texttt{-------??}\\
7 & 6 & 3 & 1 & 2 & 0 & 512 & 0 & \texttt{---------}\\
8 & 6 & 3 & 1 & 2 & 0 & 512 & 0 & \texttt{---------}\\
\multicolumn{9}{@{}l@{}}{\rule{0pt}{2.6ex}\textbf{T9.198}\enspace \texttt{012333232}\enspace $n{=}9$, $|H|{=}5$, len 4, br 2, $|\mathrm{Aut}|{=}6$, orb 7/6, hist-orb 3/2; M 9728, res 5064, rev 2772: ok}\\*
0 & -- & 0 & 5 & 1 & 310 & 2 & 200 & \texttt{-+???????}\\*
1 & 0 & 1 & 5 & 1 & 108 & 4 & 400 & \texttt{--???????}\\
2 & 1 & 2 & 3 & 1 & 64 & 64 & 384 & \texttt{---???---}\\
3 & 2 & 3 & 1 & 3 & 0 & 512 & 0 & \texttt{---------}\\
4 & 2 & 3 & 1 & 3 & 0 & 512 & 0 & \texttt{---------}\\
5 & 2 & 3 & 1 & 3 & 0 & 512 & 0 & \texttt{---------}\\
6 & 1 & 2 & 1 & 1 & 256 & 256 & 0 & \texttt{-------+-}\\
7 & 6 & 3 & 1 & 1 & 0 & 512 & 0 & \texttt{---------}\\
8 & 1 & 2 & 1 & 1 & 0 & 512 & 0 & \texttt{---------}\\
\multicolumn{9}{@{}l@{}}{\rule{0pt}{2.6ex}\textbf{T9.199}\enspace \texttt{012333231}\enspace $n{=}9$, $|H|{=}5$, len 4, br 3, $|\mathrm{Aut}|{=}6$, orb 7/6, hist-orb 3/2; M 9216, res 5072, rev 2424: ok}\\*
0 & -- & 0 & 5 & 1 & 182 & 2 & 328 & \texttt{-????????}\\*
1 & 0 & 1 & 4 & 1 & 216 & 8 & 288 & \texttt{--??????-}\\
2 & 1 & 2 & 3 & 1 & 64 & 64 & 384 & \texttt{---???---}\\
3 & 2 & 3 & 1 & 3 & 0 & 512 & 0 & \texttt{---------}\\
4 & 2 & 3 & 1 & 3 & 0 & 512 & 0 & \texttt{---------}\\
5 & 2 & 3 & 1 & 3 & 0 & 512 & 0 & \texttt{---------}\\
6 & 1 & 2 & 1 & 1 & 256 & 256 & 0 & \texttt{-------+-}\\
7 & 6 & 3 & 1 & 1 & 0 & 512 & 0 & \texttt{---------}\\
8 & 0 & 1 & 1 & 1 & 0 & 512 & 0 & \texttt{---------}\\
\multicolumn{9}{@{}l@{}}{\rule{0pt}{2.6ex}\textbf{T9.200}\enspace \texttt{012333222}\enspace $n{=}9$, $|H|{=}6$, len 4, br 2, $|\mathrm{Aut}|{=}36$, orb 5/3, hist-orb 2/0; M 9216, res 5820, rev 2092: ok}\\*
0 & -- & 0 & 6 & 1 & 274 & 2 & 236 & \texttt{-+???????}\\*
1 & 0 & 1 & 6 & 1 & 36 & 4 & 472 & \texttt{--???????}\\
2 & 1 & 2 & 3 & 1 & 64 & 64 & 384 & \texttt{---???---}\\
3 & 2 & 3 & 1 & 3 & 0 & 512 & 0 & \texttt{---------}\\
4 & 2 & 3 & 1 & 3 & 0 & 512 & 0 & \texttt{---------}\\
5 & 2 & 3 & 1 & 3 & 0 & 512 & 0 & \texttt{---------}\\
6 & 1 & 2 & 1 & 3 & 0 & 512 & 0 & \texttt{---------}\\
7 & 1 & 2 & 1 & 3 & 0 & 512 & 0 & \texttt{---------}\\
8 & 1 & 2 & 1 & 3 & 0 & 512 & 0 & \texttt{---------}\\
\multicolumn{9}{@{}l@{}}{\rule{0pt}{2.6ex}\textbf{T9.201}\enspace \texttt{012333221}\enspace $n{=}9$, $|H|{=}6$, len 4, br 3, $|\mathrm{Aut}|{=}12$, orb 6/4, hist-orb 3/1; M 8704, res 6152, rev 1456: ok}\\*
0 & -- & 0 & 6 & 1 & 146 & 2 & 364 & \texttt{-????????}\\*
1 & 0 & 1 & 5 & 1 & 72 & 8 & 432 & \texttt{--??????-}\\
2 & 1 & 2 & 3 & 1 & 64 & 64 & 384 & \texttt{---???---}\\
3 & 2 & 3 & 1 & 3 & 0 & 512 & 0 & \texttt{---------}\\
4 & 2 & 3 & 1 & 3 & 0 & 512 & 0 & \texttt{---------}\\
5 & 2 & 3 & 1 & 3 & 0 & 512 & 0 & \texttt{---------}\\
6 & 1 & 2 & 1 & 2 & 0 & 512 & 0 & \texttt{---------}\\
7 & 1 & 2 & 1 & 2 & 0 & 512 & 0 & \texttt{---------}\\
8 & 0 & 1 & 1 & 1 & 0 & 512 & 0 & \texttt{---------}\\
\multicolumn{9}{@{}l@{}}{\rule{0pt}{2.6ex}\textbf{T9.202}\enspace \texttt{012333212}\enspace $n{=}9$, $|H|{=}5$, len 4, br 3, $|\mathrm{Aut}|{=}6$, orb 7/6, hist-orb 3/2; M 8704, res 4560, rev 2388: ok}\\*
0 & -- & 0 & 5 & 1 & 246 & 2 & 264 & \texttt{-????????}\\*
1 & 0 & 1 & 4 & 1 & 144 & 16 & 352 & \texttt{--?????--}\\
2 & 1 & 2 & 3 & 1 & 64 & 64 & 384 & \texttt{---???---}\\
3 & 2 & 3 & 1 & 3 & 0 & 512 & 0 & \texttt{---------}\\
4 & 2 & 3 & 1 & 3 & 0 & 512 & 0 & \texttt{---------}\\
5 & 2 & 3 & 1 & 3 & 0 & 512 & 0 & \texttt{---------}\\
6 & 1 & 2 & 1 & 1 & 0 & 512 & 0 & \texttt{---------}\\
7 & 0 & 1 & 1 & 1 & 256 & 256 & 0 & \texttt{--------+}\\
8 & 7 & 2 & 1 & 1 & 0 & 512 & 0 & \texttt{---------}\\
\multicolumn{9}{@{}l@{}}{\rule{0pt}{2.6ex}\textbf{T9.203}\enspace \texttt{012333211}\enspace $n{=}9$, $|H|{=}6$, len 4, br 3, $|\mathrm{Aut}|{=}12$, orb 6/4, hist-orb 3/1; M 8192, res 5872, rev 1288: ok}\\*
0 & -- & 0 & 6 & 1 & 82 & 2 & 428 & \texttt{-????????}\\*
1 & 0 & 1 & 4 & 1 & 144 & 16 & 352 & \texttt{--?????--}\\
2 & 1 & 2 & 3 & 1 & 64 & 64 & 384 & \texttt{---???---}\\
3 & 2 & 3 & 1 & 3 & 0 & 512 & 0 & \texttt{---------}\\
4 & 2 & 3 & 1 & 3 & 0 & 512 & 0 & \texttt{---------}\\
5 & 2 & 3 & 1 & 3 & 0 & 512 & 0 & \texttt{---------}\\
6 & 1 & 2 & 1 & 1 & 0 & 512 & 0 & \texttt{---------}\\
7 & 0 & 1 & 1 & 2 & 0 & 512 & 0 & \texttt{---------}\\
8 & 0 & 1 & 1 & 2 & 0 & 512 & 0 & \texttt{---------}\\
\multicolumn{9}{@{}l@{}}{\rule{0pt}{2.6ex}\textbf{T9.204}\enspace \texttt{012333123}\enspace $n{=}9$, $|H|{=}4$, len 4, br 2, $|\mathrm{Aut}|{=}6$, orb 7/6, hist-orb 2/1; M 9216, res 2816, rev 3512: ok}\\*
0 & -- & 0 & 4 & 1 & 350 & 2 & 160 & \texttt{-????????}\\*
1 & 0 & 1 & 3 & 1 & 288 & 32 & 192 & \texttt{--+???---}\\
2 & 1 & 2 & 3 & 1 & 64 & 64 & 384 & \texttt{---???---}\\
3 & 2 & 3 & 1 & 3 & 0 & 512 & 0 & \texttt{---------}\\
4 & 2 & 3 & 1 & 3 & 0 & 512 & 0 & \texttt{---------}\\
5 & 2 & 3 & 1 & 3 & 0 & 512 & 0 & \texttt{---------}\\
6 & 0 & 1 & 1 & 1 & 384 & 128 & 0 & \texttt{-------++}\\
7 & 6 & 2 & 1 & 1 & 256 & 256 & 0 & \texttt{--------+}\\
8 & 7 & 3 & 1 & 1 & 0 & 512 & 0 & \texttt{---------}\\
\multicolumn{9}{@{}l@{}}{\rule{0pt}{2.6ex}\textbf{T9.205}\enspace \texttt{012333122}\enspace $n{=}9$, $|H|{=}5$, len 4, br 3, $|\mathrm{Aut}|{=}12$, orb 6/4, hist-orb 2/0; M 8704, res 4000, rev 2684: ok}\\*
0 & -- & 0 & 5 & 1 & 250 & 2 & 260 & \texttt{-????????}\\*
1 & 0 & 1 & 3 & 1 & 288 & 32 & 192 & \texttt{--+???---}\\
2 & 1 & 2 & 3 & 1 & 64 & 64 & 384 & \texttt{---???---}\\
3 & 2 & 3 & 1 & 3 & 0 & 512 & 0 & \texttt{---------}\\
4 & 2 & 3 & 1 & 3 & 0 & 512 & 0 & \texttt{---------}\\
5 & 2 & 3 & 1 & 3 & 0 & 512 & 0 & \texttt{---------}\\
6 & 0 & 1 & 2 & 1 & 128 & 128 & 256 & \texttt{-------??}\\
7 & 6 & 2 & 1 & 2 & 0 & 512 & 0 & \texttt{---------}\\
8 & 6 & 2 & 1 & 2 & 0 & 512 & 0 & \texttt{---------}\\
\multicolumn{9}{@{}l@{}}{\rule{0pt}{2.6ex}\textbf{T9.206}\enspace \texttt{012333121}\enspace $n{=}9$, $|H|{=}5$, len 4, br 2, $|\mathrm{Aut}|{=}6$, orb 7/6, hist-orb 3/2; M 8192, res 4416, rev 2168: ok}\\*
0 & -- & 0 & 5 & 1 & 150 & 2 & 360 & \texttt{-????????}\\*
1 & 0 & 1 & 3 & 1 & 288 & 32 & 192 & \texttt{--+???---}\\
2 & 1 & 2 & 3 & 1 & 64 & 64 & 384 & \texttt{---???---}\\
3 & 2 & 3 & 1 & 3 & 0 & 512 & 0 & \texttt{---------}\\
4 & 2 & 3 & 1 & 3 & 0 & 512 & 0 & \texttt{---------}\\
5 & 2 & 3 & 1 & 3 & 0 & 512 & 0 & \texttt{---------}\\
6 & 0 & 1 & 1 & 1 & 256 & 256 & 0 & \texttt{-------+-}\\
7 & 6 & 2 & 1 & 1 & 0 & 512 & 0 & \texttt{---------}\\
8 & 0 & 1 & 1 & 1 & 0 & 512 & 0 & \texttt{---------}\\
\multicolumn{9}{@{}l@{}}{\rule{0pt}{2.6ex}\textbf{T9.207}\enspace \texttt{012333111}\enspace $n{=}9$, $|H|{=}6$, len 4, br 2, $|\mathrm{Aut}|{=}36$, orb 5/3, hist-orb 2/0; M 7680, res 5216, rev 1396: ok}\\*
0 & -- & 0 & 6 & 1 & 50 & 2 & 460 & \texttt{-????????}\\*
1 & 0 & 1 & 3 & 1 & 288 & 32 & 192 & \texttt{--+???---}\\
2 & 1 & 2 & 3 & 1 & 64 & 64 & 384 & \texttt{---???---}\\
3 & 2 & 3 & 1 & 3 & 0 & 512 & 0 & \texttt{---------}\\
4 & 2 & 3 & 1 & 3 & 0 & 512 & 0 & \texttt{---------}\\
5 & 2 & 3 & 1 & 3 & 0 & 512 & 0 & \texttt{---------}\\
6 & 0 & 1 & 1 & 3 & 0 & 512 & 0 & \texttt{---------}\\
7 & 0 & 1 & 1 & 3 & 0 & 512 & 0 & \texttt{---------}\\
8 & 0 & 1 & 1 & 3 & 0 & 512 & 0 & \texttt{---------}\\
\multicolumn{9}{@{}l@{}}{\rule{0pt}{2.6ex}\textbf{T9.208}\enspace \texttt{012332332}\enspace $n{=}9$, $|H|{=}5$, len 4, br 3, $|\mathrm{Aut}|{=}8$, orb 5/3, hist-orb 2/1; M 9728, res 4924, rev 2732: ok}\\*
0 & -- & 0 & 5 & 1 & 306 & 2 & 204 & \texttt{-+???????}\\*
1 & 0 & 1 & 5 & 1 & 100 & 4 & 408 & \texttt{--???????}\\
2 & 1 & 2 & 2 & 2 & 128 & 128 & 256 & \texttt{---??----}\\
3 & 2 & 3 & 1 & 4 & 0 & 512 & 0 & \texttt{---------}\\
4 & 2 & 3 & 1 & 4 & 0 & 512 & 0 & \texttt{---------}\\
5 & 1 & 2 & 2 & 2 & 128 & 128 & 256 & \texttt{------??-}\\
6 & 5 & 3 & 1 & 4 & 0 & 512 & 0 & \texttt{---------}\\
7 & 5 & 3 & 1 & 4 & 0 & 512 & 0 & \texttt{---------}\\
8 & 1 & 2 & 1 & 1 & 0 & 512 & 0 & \texttt{---------}\\
\multicolumn{9}{@{}l@{}}{\rule{0pt}{2.6ex}\textbf{T9.209}\enspace \texttt{012332331}\enspace $n{=}9$, $|H|{=}5$, len 4, br 4, $|\mathrm{Aut}|{=}8$, orb 5/3, hist-orb 2/1; M 9216, res 4968, rev 2368: ok}\\*
0 & -- & 0 & 5 & 1 & 178 & 2 & 332 & \texttt{-????????}\\*
1 & 0 & 1 & 4 & 1 & 200 & 8 & 304 & \texttt{--??????-}\\
2 & 1 & 2 & 2 & 2 & 128 & 128 & 256 & \texttt{---??----}\\
3 & 2 & 3 & 1 & 4 & 0 & 512 & 0 & \texttt{---------}\\
4 & 2 & 3 & 1 & 4 & 0 & 512 & 0 & \texttt{---------}\\
5 & 1 & 2 & 2 & 2 & 128 & 128 & 256 & \texttt{------??-}\\
6 & 5 & 3 & 1 & 4 & 0 & 512 & 0 & \texttt{---------}\\
7 & 5 & 3 & 1 & 4 & 0 & 512 & 0 & \texttt{---------}\\
8 & 0 & 1 & 1 & 1 & 0 & 512 & 0 & \texttt{---------}\\
\multicolumn{9}{@{}l@{}}{\rule{0pt}{2.6ex}\textbf{T9.210}\enspace \texttt{012332323}\enspace $n{=}9$, $|H|{=}4$, len 4, br 2, $|\mathrm{Aut}|{=}4$, orb 6/3, hist-orb 2/0; M 9728, res 3764, rev 3428: ok}\\*
0 & -- & 0 & 4 & 1 & 346 & 2 & 164 & \texttt{-+???????}\\*
1 & 0 & 1 & 4 & 1 & 180 & 4 & 328 & \texttt{--???????}\\
2 & 1 & 2 & 2 & 1 & 128 & 128 & 256 & \texttt{---??----}\\
3 & 2 & 3 & 1 & 2 & 0 & 512 & 0 & \texttt{---------}\\
4 & 2 & 3 & 1 & 2 & 0 & 512 & 0 & \texttt{---------}\\
5 & 1 & 2 & 1 & 2 & 256 & 256 & 0 & \texttt{------+--}\\
6 & 5 & 3 & 1 & 2 & 0 & 512 & 0 & \texttt{---------}\\
7 & 1 & 2 & 1 & 2 & 256 & 256 & 0 & \texttt{--------+}\\
8 & 7 & 3 & 1 & 2 & 0 & 512 & 0 & \texttt{---------}\\
\multicolumn{9}{@{}l@{}}{\rule{0pt}{2.6ex}\textbf{T9.211}\enspace \texttt{012332322}\enspace $n{=}9$, $|H|{=}5$, len 4, br 2, $|\mathrm{Aut}|{=}4$, orb 7/5, hist-orb 3/1; M 9216, res 5024, rev 2484: ok}\\*
0 & -- & 0 & 5 & 1 & 286 & 2 & 224 & \texttt{-+???????}\\*
1 & 0 & 1 & 5 & 1 & 60 & 4 & 448 & \texttt{--???????}\\
2 & 1 & 2 & 2 & 1 & 128 & 128 & 256 & \texttt{---??----}\\
3 & 2 & 3 & 1 & 2 & 0 & 512 & 0 & \texttt{---------}\\
4 & 2 & 3 & 1 & 2 & 0 & 512 & 0 & \texttt{---------}\\
5 & 1 & 2 & 1 & 1 & 256 & 256 & 0 & \texttt{------+--}\\
6 & 5 & 3 & 1 & 1 & 0 & 512 & 0 & \texttt{---------}\\
7 & 1 & 2 & 1 & 2 & 0 & 512 & 0 & \texttt{---------}\\
8 & 1 & 2 & 1 & 2 & 0 & 512 & 0 & \texttt{---------}\\
\multicolumn{9}{@{}l@{}}{\rule{0pt}{2.6ex}\textbf{T9.212}\enspace \texttt{012332321}\enspace $n{=}9$, $|H|{=}5$, len 4, br 3, $|\mathrm{Aut}|{=}2$, orb 8/7, hist-orb 4/3; M 8704, res 5248, rev 1960: ok}\\*
0 & -- & 0 & 5 & 1 & 158 & 2 & 352 & \texttt{-????????}\\*
1 & 0 & 1 & 4 & 1 & 120 & 8 & 384 & \texttt{--??????-}\\
2 & 1 & 2 & 2 & 1 & 128 & 128 & 256 & \texttt{---??----}\\
3 & 2 & 3 & 1 & 2 & 0 & 512 & 0 & \texttt{---------}\\
4 & 2 & 3 & 1 & 2 & 0 & 512 & 0 & \texttt{---------}\\
5 & 1 & 2 & 1 & 1 & 256 & 256 & 0 & \texttt{------+--}\\
6 & 5 & 3 & 1 & 1 & 0 & 512 & 0 & \texttt{---------}\\
7 & 1 & 2 & 1 & 1 & 0 & 512 & 0 & \texttt{---------}\\
8 & 0 & 1 & 1 & 1 & 0 & 512 & 0 & \texttt{---------}\\
\multicolumn{9}{@{}l@{}}{\rule{0pt}{2.6ex}\textbf{T9.213}\enspace \texttt{012332312}\enspace $n{=}9$, $|H|{=}4$, len 4, br 3, $|\mathrm{Aut}|{=}2$, orb 8/7, hist-orb 3/2; M 8704, res 3296, rev 3036: ok}\\*
0 & -- & 0 & 4 & 1 & 282 & 2 & 228 & \texttt{-????????}\\*
1 & 0 & 1 & 3 & 1 & 240 & 16 & 256 & \texttt{--?????--}\\
2 & 1 & 2 & 2 & 1 & 128 & 128 & 256 & \texttt{---??----}\\
3 & 2 & 3 & 1 & 2 & 0 & 512 & 0 & \texttt{---------}\\
4 & 2 & 3 & 1 & 2 & 0 & 512 & 0 & \texttt{---------}\\
5 & 1 & 2 & 1 & 1 & 256 & 256 & 0 & \texttt{------+--}\\
6 & 5 & 3 & 1 & 1 & 0 & 512 & 0 & \texttt{---------}\\
7 & 0 & 1 & 1 & 1 & 256 & 256 & 0 & \texttt{--------+}\\
8 & 7 & 2 & 1 & 1 & 0 & 512 & 0 & \texttt{---------}\\
\multicolumn{9}{@{}l@{}}{\rule{0pt}{2.6ex}\textbf{T9.214}\enspace \texttt{012332311}\enspace $n{=}9$, $|H|{=}5$, len 4, br 3, $|\mathrm{Aut}|{=}4$, orb 7/5, hist-orb 3/1; M 8192, res 4800, rev 1896: ok}\\*
0 & -- & 0 & 5 & 1 & 94 & 2 & 416 & \texttt{-????????}\\*
1 & 0 & 1 & 3 & 1 & 240 & 16 & 256 & \texttt{--?????--}\\
2 & 1 & 2 & 2 & 1 & 128 & 128 & 256 & \texttt{---??----}\\
3 & 2 & 3 & 1 & 2 & 0 & 512 & 0 & \texttt{---------}\\
4 & 2 & 3 & 1 & 2 & 0 & 512 & 0 & \texttt{---------}\\
5 & 1 & 2 & 1 & 1 & 256 & 256 & 0 & \texttt{------+--}\\
6 & 5 & 3 & 1 & 1 & 0 & 512 & 0 & \texttt{---------}\\
7 & 0 & 1 & 1 & 2 & 0 & 512 & 0 & \texttt{---------}\\
8 & 0 & 1 & 1 & 2 & 0 & 512 & 0 & \texttt{---------}\\
\multicolumn{9}{@{}l@{}}{\rule{0pt}{2.6ex}\textbf{T9.215}\enspace \texttt{012332222}\enspace $n{=}9$, $|H|{=}6$, len 4, br 2, $|\mathrm{Aut}|{=}48$, orb 5/3, hist-orb 2/0; M 8704, res 5444, rev 2044: ok}\\*
0 & -- & 0 & 6 & 1 & 266 & 2 & 244 & \texttt{-+???????}\\*
1 & 0 & 1 & 6 & 1 & 20 & 4 & 488 & \texttt{--???????}\\
2 & 1 & 2 & 2 & 1 & 128 & 128 & 256 & \texttt{---??----}\\
3 & 2 & 3 & 1 & 2 & 0 & 512 & 0 & \texttt{---------}\\
4 & 2 & 3 & 1 & 2 & 0 & 512 & 0 & \texttt{---------}\\
5 & 1 & 2 & 1 & 4 & 0 & 512 & 0 & \texttt{---------}\\
6 & 1 & 2 & 1 & 4 & 0 & 512 & 0 & \texttt{---------}\\
7 & 1 & 2 & 1 & 4 & 0 & 512 & 0 & \texttt{---------}\\
8 & 1 & 2 & 1 & 4 & 0 & 512 & 0 & \texttt{---------}\\
\multicolumn{9}{@{}l@{}}{\rule{0pt}{2.6ex}\textbf{T9.216}\enspace \texttt{012332221}\enspace $n{=}9$, $|H|{=}6$, len 4, br 3, $|\mathrm{Aut}|{=}12$, orb 6/4, hist-orb 3/1; M 8192, res 5848, rev 1328: ok}\\*
0 & -- & 0 & 6 & 1 & 138 & 2 & 372 & \texttt{-????????}\\*
1 & 0 & 1 & 5 & 1 & 40 & 8 & 464 & \texttt{--??????-}\\
2 & 1 & 2 & 2 & 1 & 128 & 128 & 256 & \texttt{---??----}\\
3 & 2 & 3 & 1 & 2 & 0 & 512 & 0 & \texttt{---------}\\
4 & 2 & 3 & 1 & 2 & 0 & 512 & 0 & \texttt{---------}\\
5 & 1 & 2 & 1 & 3 & 0 & 512 & 0 & \texttt{---------}\\
6 & 1 & 2 & 1 & 3 & 0 & 512 & 0 & \texttt{---------}\\
7 & 1 & 2 & 1 & 3 & 0 & 512 & 0 & \texttt{---------}\\
8 & 0 & 1 & 1 & 1 & 0 & 512 & 0 & \texttt{---------}\\
\multicolumn{9}{@{}l@{}}{\rule{0pt}{2.6ex}\textbf{T9.217}\enspace \texttt{012332212}\enspace $n{=}9$, $|H|{=}5$, len 4, br 3, $|\mathrm{Aut}|{=}4$, orb 7/5, hist-orb 3/1; M 8192, res 4496, rev 2092: ok}\\*
0 & -- & 0 & 5 & 1 & 222 & 2 & 288 & \texttt{-????????}\\*
1 & 0 & 1 & 4 & 1 & 80 & 16 & 416 & \texttt{--?????--}\\
2 & 1 & 2 & 2 & 1 & 128 & 128 & 256 & \texttt{---??----}\\
3 & 2 & 3 & 1 & 2 & 0 & 512 & 0 & \texttt{---------}\\
4 & 2 & 3 & 1 & 2 & 0 & 512 & 0 & \texttt{---------}\\
5 & 1 & 2 & 1 & 2 & 0 & 512 & 0 & \texttt{---------}\\
6 & 1 & 2 & 1 & 2 & 0 & 512 & 0 & \texttt{---------}\\
7 & 0 & 1 & 1 & 1 & 256 & 256 & 0 & \texttt{--------+}\\
8 & 7 & 2 & 1 & 1 & 0 & 512 & 0 & \texttt{---------}\\
\multicolumn{9}{@{}l@{}}{\rule{0pt}{2.6ex}\textbf{T9.218}\enspace \texttt{012332211}\enspace $n{=}9$, $|H|{=}6$, len 4, br 3, $|\mathrm{Aut}|{=}8$, orb 6/3, hist-orb 3/0; M 7680, res 5680, rev 1056: ok}\\*
0 & -- & 0 & 6 & 1 & 74 & 2 & 436 & \texttt{-????????}\\*
1 & 0 & 1 & 4 & 1 & 80 & 16 & 416 & \texttt{--?????--}\\
2 & 1 & 2 & 2 & 1 & 128 & 128 & 256 & \texttt{---??----}\\
3 & 2 & 3 & 1 & 2 & 0 & 512 & 0 & \texttt{---------}\\
4 & 2 & 3 & 1 & 2 & 0 & 512 & 0 & \texttt{---------}\\
5 & 1 & 2 & 1 & 2 & 0 & 512 & 0 & \texttt{---------}\\
6 & 1 & 2 & 1 & 2 & 0 & 512 & 0 & \texttt{---------}\\
7 & 0 & 1 & 1 & 2 & 0 & 512 & 0 & \texttt{---------}\\
8 & 0 & 1 & 1 & 2 & 0 & 512 & 0 & \texttt{---------}\\
\multicolumn{9}{@{}l@{}}{\rule{0pt}{2.6ex}\textbf{T9.219}\enspace \texttt{012332123}\enspace $n{=}9$, $|H|{=}4$, len 4, br 3, $|\mathrm{Aut}|{=}2$, orb 8/7, hist-orb 3/2; M 8704, res 3168, rev 2964: ok}\\*
0 & -- & 0 & 4 & 1 & 294 & 2 & 216 & \texttt{-????????}\\*
1 & 0 & 1 & 3 & 1 & 160 & 32 & 320 & \texttt{--????---}\\
2 & 1 & 2 & 2 & 1 & 128 & 128 & 256 & \texttt{---??----}\\
3 & 2 & 3 & 1 & 2 & 0 & 512 & 0 & \texttt{---------}\\
4 & 2 & 3 & 1 & 2 & 0 & 512 & 0 & \texttt{---------}\\
5 & 1 & 2 & 1 & 1 & 0 & 512 & 0 & \texttt{---------}\\
6 & 0 & 1 & 1 & 1 & 384 & 128 & 0 & \texttt{-------++}\\
7 & 6 & 2 & 1 & 1 & 256 & 256 & 0 & \texttt{--------+}\\
8 & 7 & 3 & 1 & 1 & 0 & 512 & 0 & \texttt{---------}\\
\multicolumn{9}{@{}l@{}}{\rule{0pt}{2.6ex}\textbf{T9.220}\enspace \texttt{012332122}\enspace $n{=}9$, $|H|{=}5$, len 4, br 4, $|\mathrm{Aut}|{=}4$, orb 7/5, hist-orb 3/1; M 8192, res 4224, rev 2176: ok}\\*
0 & -- & 0 & 5 & 1 & 210 & 2 & 300 & \texttt{-????????}\\*
1 & 0 & 1 & 3 & 1 & 160 & 32 & 320 & \texttt{--????---}\\
2 & 1 & 2 & 2 & 1 & 128 & 128 & 256 & \texttt{---??----}\\
3 & 2 & 3 & 1 & 2 & 0 & 512 & 0 & \texttt{---------}\\
4 & 2 & 3 & 1 & 2 & 0 & 512 & 0 & \texttt{---------}\\
5 & 1 & 2 & 1 & 1 & 0 & 512 & 0 & \texttt{---------}\\
6 & 0 & 1 & 2 & 1 & 128 & 128 & 256 & \texttt{-------??}\\
7 & 6 & 2 & 1 & 2 & 0 & 512 & 0 & \texttt{---------}\\
8 & 6 & 2 & 1 & 2 & 0 & 512 & 0 & \texttt{---------}\\
\multicolumn{9}{@{}l@{}}{\rule{0pt}{2.6ex}\textbf{T9.221}\enspace \texttt{012332121}\enspace $n{=}9$, $|H|{=}5$, len 4, br 3, $|\mathrm{Aut}|{=}2$, orb 8/7, hist-orb 4/3; M 7680, res 4512, rev 1732: ok}\\*
0 & -- & 0 & 5 & 1 & 126 & 2 & 384 & \texttt{-????????}\\*
1 & 0 & 1 & 3 & 1 & 160 & 32 & 320 & \texttt{--????---}\\
2 & 1 & 2 & 2 & 1 & 128 & 128 & 256 & \texttt{---??----}\\
3 & 2 & 3 & 1 & 2 & 0 & 512 & 0 & \texttt{---------}\\
4 & 2 & 3 & 1 & 2 & 0 & 512 & 0 & \texttt{---------}\\
5 & 1 & 2 & 1 & 1 & 0 & 512 & 0 & \texttt{---------}\\
6 & 0 & 1 & 1 & 1 & 256 & 256 & 0 & \texttt{-------+-}\\
7 & 6 & 2 & 1 & 1 & 0 & 512 & 0 & \texttt{---------}\\
8 & 0 & 1 & 1 & 1 & 0 & 512 & 0 & \texttt{---------}\\
\multicolumn{9}{@{}l@{}}{\rule{0pt}{2.6ex}\textbf{T9.222}\enspace \texttt{012332111}\enspace $n{=}9$, $|H|{=}6$, len 4, br 3, $|\mathrm{Aut}|{=}12$, orb 6/4, hist-orb 3/1; M 7168, res 5184, rev 1032: ok}\\*
0 & -- & 0 & 6 & 1 & 42 & 2 & 468 & \texttt{-????????}\\*
1 & 0 & 1 & 3 & 1 & 160 & 32 & 320 & \texttt{--????---}\\
2 & 1 & 2 & 2 & 1 & 128 & 128 & 256 & \texttt{---??----}\\
3 & 2 & 3 & 1 & 2 & 0 & 512 & 0 & \texttt{---------}\\
4 & 2 & 3 & 1 & 2 & 0 & 512 & 0 & \texttt{---------}\\
5 & 1 & 2 & 1 & 1 & 0 & 512 & 0 & \texttt{---------}\\
6 & 0 & 1 & 1 & 3 & 0 & 512 & 0 & \texttt{---------}\\
7 & 0 & 1 & 1 & 3 & 0 & 512 & 0 & \texttt{---------}\\
8 & 0 & 1 & 1 & 3 & 0 & 512 & 0 & \texttt{---------}\\
\multicolumn{9}{@{}l@{}}{\rule{0pt}{2.6ex}\textbf{T9.223}\enspace \texttt{012331233}\enspace $n{=}9$, $|H|{=}4$, len 4, br 3, $|\mathrm{Aut}|{=}8$, orb 4/1, hist-orb 1/0; M 9216, res 2656, rev 3480: ok}\\*
0 & -- & 0 & 4 & 1 & 338 & 2 & 172 & \texttt{-????????}\\*
1 & 0 & 1 & 2 & 2 & 320 & 64 & 128 & \texttt{--+??----}\\
2 & 1 & 2 & 2 & 2 & 128 & 128 & 256 & \texttt{---??----}\\
3 & 2 & 3 & 1 & 4 & 0 & 512 & 0 & \texttt{---------}\\
4 & 2 & 3 & 1 & 4 & 0 & 512 & 0 & \texttt{---------}\\
5 & 0 & 1 & 2 & 2 & 320 & 64 & 128 & \texttt{------+??}\\
6 & 5 & 2 & 2 & 2 & 128 & 128 & 256 & \texttt{-------??}\\
7 & 6 & 3 & 1 & 4 & 0 & 512 & 0 & \texttt{---------}\\
8 & 6 & 3 & 1 & 4 & 0 & 512 & 0 & \texttt{---------}\\
\multicolumn{9}{@{}l@{}}{\rule{0pt}{2.6ex}\textbf{T9.224}\enspace \texttt{012331232}\enspace $n{=}9$, $|H|{=}4$, len 4, br 3, $|\mathrm{Aut}|{=}2$, orb 8/7, hist-orb 3/2; M 8704, res 3072, rev 3052: ok}\\*
0 & -- & 0 & 4 & 1 & 286 & 2 & 224 & \texttt{-????????}\\*
1 & 0 & 1 & 2 & 1 & 320 & 64 & 128 & \texttt{--+??----}\\
2 & 1 & 2 & 2 & 1 & 128 & 128 & 256 & \texttt{---??----}\\
3 & 2 & 3 & 1 & 2 & 0 & 512 & 0 & \texttt{---------}\\
4 & 2 & 3 & 1 & 2 & 0 & 512 & 0 & \texttt{---------}\\
5 & 0 & 1 & 2 & 1 & 192 & 64 & 256 & \texttt{------???}\\
6 & 5 & 2 & 1 & 1 & 256 & 256 & 0 & \texttt{-------+-}\\
7 & 6 & 3 & 1 & 1 & 0 & 512 & 0 & \texttt{---------}\\
8 & 5 & 2 & 1 & 1 & 0 & 512 & 0 & \texttt{---------}\\
\multicolumn{9}{@{}l@{}}{\rule{0pt}{2.6ex}\textbf{T9.225}\enspace \texttt{012331231}\enspace $n{=}9$, $|H|{=}4$, len 4, br 2, $|\mathrm{Aut}|{=}2$, orb 8/7, hist-orb 3/2; M 8192, res 3264, rev 2616: ok}\\*
0 & -- & 0 & 4 & 1 & 182 & 2 & 328 & \texttt{-????????}\\*
1 & 0 & 1 & 2 & 1 & 320 & 64 & 128 & \texttt{--+??----}\\
2 & 1 & 2 & 2 & 1 & 128 & 128 & 256 & \texttt{---??----}\\
3 & 2 & 3 & 1 & 2 & 0 & 512 & 0 & \texttt{---------}\\
4 & 2 & 3 & 1 & 2 & 0 & 512 & 0 & \texttt{---------}\\
5 & 0 & 1 & 1 & 1 & 384 & 128 & 0 & \texttt{------++-}\\
6 & 5 & 2 & 1 & 1 & 256 & 256 & 0 & \texttt{-------+-}\\
7 & 6 & 3 & 1 & 1 & 0 & 512 & 0 & \texttt{---------}\\
8 & 0 & 1 & 1 & 1 & 0 & 512 & 0 & \texttt{---------}\\
\multicolumn{9}{@{}l@{}}{\rule{0pt}{2.6ex}\textbf{T9.226}\enspace \texttt{012331222}\enspace $n{=}9$, $|H|{=}5$, len 4, br 3, $|\mathrm{Aut}|{=}12$, orb 6/4, hist-orb 2/0; M 8192, res 3808, rev 2420: ok}\\*
0 & -- & 0 & 5 & 1 & 234 & 2 & 276 & \texttt{-????????}\\*
1 & 0 & 1 & 2 & 1 & 320 & 64 & 128 & \texttt{--+??----}\\
2 & 1 & 2 & 2 & 1 & 128 & 128 & 256 & \texttt{---??----}\\
3 & 2 & 3 & 1 & 2 & 0 & 512 & 0 & \texttt{---------}\\
4 & 2 & 3 & 1 & 2 & 0 & 512 & 0 & \texttt{---------}\\
5 & 0 & 1 & 3 & 1 & 64 & 64 & 384 & \texttt{------???}\\
6 & 5 & 2 & 1 & 3 & 0 & 512 & 0 & \texttt{---------}\\
7 & 5 & 2 & 1 & 3 & 0 & 512 & 0 & \texttt{---------}\\
8 & 5 & 2 & 1 & 3 & 0 & 512 & 0 & \texttt{---------}\\
\multicolumn{9}{@{}l@{}}{\rule{0pt}{2.6ex}\textbf{T9.227}\enspace \texttt{012331221}\enspace $n{=}9$, $|H|{=}5$, len 4, br 3, $|\mathrm{Aut}|{=}4$, orb 7/5, hist-orb 3/1; M 7680, res 4064, rev 1932: ok}\\*
0 & -- & 0 & 5 & 1 & 130 & 2 & 380 & \texttt{-????????}\\*
1 & 0 & 1 & 2 & 1 & 320 & 64 & 128 & \texttt{--+??----}\\
2 & 1 & 2 & 2 & 1 & 128 & 128 & 256 & \texttt{---??----}\\
3 & 2 & 3 & 1 & 2 & 0 & 512 & 0 & \texttt{---------}\\
4 & 2 & 3 & 1 & 2 & 0 & 512 & 0 & \texttt{---------}\\
5 & 0 & 1 & 2 & 1 & 128 & 128 & 256 & \texttt{------??-}\\
6 & 5 & 2 & 1 & 2 & 0 & 512 & 0 & \texttt{---------}\\
7 & 5 & 2 & 1 & 2 & 0 & 512 & 0 & \texttt{---------}\\
8 & 0 & 1 & 1 & 1 & 0 & 512 & 0 & \texttt{---------}\\
\multicolumn{9}{@{}l@{}}{\rule{0pt}{2.6ex}\textbf{T9.228}\enspace \texttt{012331212}\enspace $n{=}9$, $|H|{=}4$, len 4, br 2, $|\mathrm{Aut}|{=}4$, orb 6/3, hist-orb 2/0; M 7680, res 2848, rev 2636: ok}\\*
0 & -- & 0 & 4 & 1 & 234 & 2 & 276 & \texttt{-????????}\\*
1 & 0 & 1 & 2 & 1 & 320 & 64 & 128 & \texttt{--+??----}\\
2 & 1 & 2 & 2 & 1 & 128 & 128 & 256 & \texttt{---??----}\\
3 & 2 & 3 & 1 & 2 & 0 & 512 & 0 & \texttt{---------}\\
4 & 2 & 3 & 1 & 2 & 0 & 512 & 0 & \texttt{---------}\\
5 & 0 & 1 & 1 & 2 & 256 & 256 & 0 & \texttt{------+--}\\
6 & 5 & 2 & 1 & 2 & 0 & 512 & 0 & \texttt{---------}\\
7 & 0 & 1 & 1 & 2 & 256 & 256 & 0 & \texttt{--------+}\\
8 & 7 & 2 & 1 & 2 & 0 & 512 & 0 & \texttt{---------}\\
\multicolumn{9}{@{}l@{}}{\rule{0pt}{2.6ex}\textbf{T9.229}\enspace \texttt{012331211}\enspace $n{=}9$, $|H|{=}5$, len 4, br 2, $|\mathrm{Aut}|{=}4$, orb 7/5, hist-orb 3/1; M 7168, res 4096, rev 1656: ok}\\*
0 & -- & 0 & 5 & 1 & 78 & 2 & 432 & \texttt{-????????}\\*
1 & 0 & 1 & 2 & 1 & 320 & 64 & 128 & \texttt{--+??----}\\
2 & 1 & 2 & 2 & 1 & 128 & 128 & 256 & \texttt{---??----}\\
3 & 2 & 3 & 1 & 2 & 0 & 512 & 0 & \texttt{---------}\\
4 & 2 & 3 & 1 & 2 & 0 & 512 & 0 & \texttt{---------}\\
5 & 0 & 1 & 1 & 1 & 256 & 256 & 0 & \texttt{------+--}\\
6 & 5 & 2 & 1 & 1 & 0 & 512 & 0 & \texttt{---------}\\
7 & 0 & 1 & 1 & 2 & 0 & 512 & 0 & \texttt{---------}\\
8 & 0 & 1 & 1 & 2 & 0 & 512 & 0 & \texttt{---------}\\
\multicolumn{9}{@{}l@{}}{\rule{0pt}{2.6ex}\textbf{T9.230}\enspace \texttt{012331111}\enspace $n{=}9$, $|H|{=}6$, len 4, br 2, $|\mathrm{Aut}|{=}48$, orb 5/3, hist-orb 2/0; M 6656, res 4512, rev 1124: ok}\\*
0 & -- & 0 & 6 & 1 & 26 & 2 & 484 & \texttt{-????????}\\*
1 & 0 & 1 & 2 & 1 & 320 & 64 & 128 & \texttt{--+??----}\\
2 & 1 & 2 & 2 & 1 & 128 & 128 & 256 & \texttt{---??----}\\
3 & 2 & 3 & 1 & 2 & 0 & 512 & 0 & \texttt{---------}\\
4 & 2 & 3 & 1 & 2 & 0 & 512 & 0 & \texttt{---------}\\
5 & 0 & 1 & 1 & 4 & 0 & 512 & 0 & \texttt{---------}\\
6 & 0 & 1 & 1 & 4 & 0 & 512 & 0 & \texttt{---------}\\
7 & 0 & 1 & 1 & 4 & 0 & 512 & 0 & \texttt{---------}\\
8 & 0 & 1 & 1 & 4 & 0 & 512 & 0 & \texttt{---------}\\
\multicolumn{9}{@{}l@{}}{\rule{0pt}{2.6ex}\textbf{T9.231}\enspace \texttt{012323232}\enspace $n{=}9$, $|H|{=}4$, len 4, br 1, $|\mathrm{Aut}|{=}6$, orb 5/3, hist-orb 2/1; M 9216, res 4200, rev 2988: ok}\\*
0 & -- & 0 & 4 & 1 & 310 & 2 & 200 & \texttt{-+???????}\\*
1 & 0 & 1 & 4 & 1 & 108 & 4 & 400 & \texttt{--???????}\\
2 & 1 & 2 & 1 & 3 & 256 & 256 & 0 & \texttt{---+-----}\\
3 & 2 & 3 & 1 & 3 & 0 & 512 & 0 & \texttt{---------}\\
4 & 1 & 2 & 1 & 3 & 256 & 256 & 0 & \texttt{-----+---}\\
5 & 4 & 3 & 1 & 3 & 0 & 512 & 0 & \texttt{---------}\\
6 & 1 & 2 & 1 & 3 & 256 & 256 & 0 & \texttt{-------+-}\\
7 & 6 & 3 & 1 & 3 & 0 & 512 & 0 & \texttt{---------}\\
8 & 1 & 2 & 1 & 1 & 0 & 512 & 0 & \texttt{---------}\\
\multicolumn{9}{@{}l@{}}{\rule{0pt}{2.6ex}\textbf{T9.232}\enspace \texttt{012323231}\enspace $n{=}9$, $|H|{=}4$, len 4, br 2, $|\mathrm{Aut}|{=}6$, orb 5/3, hist-orb 2/1; M 8704, res 4208, rev 2608: ok}\\*
0 & -- & 0 & 4 & 1 & 182 & 2 & 328 & \texttt{-????????}\\*
1 & 0 & 1 & 3 & 1 & 216 & 8 & 288 & \texttt{--??????-}\\
2 & 1 & 2 & 1 & 3 & 256 & 256 & 0 & \texttt{---+-----}\\
3 & 2 & 3 & 1 & 3 & 0 & 512 & 0 & \texttt{---------}\\
4 & 1 & 2 & 1 & 3 & 256 & 256 & 0 & \texttt{-----+---}\\
5 & 4 & 3 & 1 & 3 & 0 & 512 & 0 & \texttt{---------}\\
6 & 1 & 2 & 1 & 3 & 256 & 256 & 0 & \texttt{-------+-}\\
7 & 6 & 3 & 1 & 3 & 0 & 512 & 0 & \texttt{---------}\\
8 & 0 & 1 & 1 & 1 & 0 & 512 & 0 & \texttt{---------}\\
\multicolumn{9}{@{}l@{}}{\rule{0pt}{2.6ex}\textbf{T9.233}\enspace \texttt{012323222}\enspace $n{=}9$, $|H|{=}5$, len 4, br 1, $|\mathrm{Aut}|{=}12$, orb 5/2, hist-orb 2/0; M 8704, res 4956, rev 2356: ok}\\*
0 & -- & 0 & 5 & 1 & 274 & 2 & 236 & \texttt{-+???????}\\*
1 & 0 & 1 & 5 & 1 & 36 & 4 & 472 & \texttt{--???????}\\
2 & 1 & 2 & 1 & 2 & 256 & 256 & 0 & \texttt{---+-----}\\
3 & 2 & 3 & 1 & 2 & 0 & 512 & 0 & \texttt{---------}\\
4 & 1 & 2 & 1 & 2 & 256 & 256 & 0 & \texttt{-----+---}\\
5 & 4 & 3 & 1 & 2 & 0 & 512 & 0 & \texttt{---------}\\
6 & 1 & 2 & 1 & 3 & 0 & 512 & 0 & \texttt{---------}\\
7 & 1 & 2 & 1 & 3 & 0 & 512 & 0 & \texttt{---------}\\
8 & 1 & 2 & 1 & 3 & 0 & 512 & 0 & \texttt{---------}\\
\multicolumn{9}{@{}l@{}}{\rule{0pt}{2.6ex}\textbf{T9.234}\enspace \texttt{012323221}\enspace $n{=}9$, $|H|{=}5$, len 4, br 2, $|\mathrm{Aut}|{=}4$, orb 6/3, hist-orb 3/1; M 8192, res 5288, rev 1720: ok}\\*
0 & -- & 0 & 5 & 1 & 146 & 2 & 364 & \texttt{-????????}\\*
1 & 0 & 1 & 4 & 1 & 72 & 8 & 432 & \texttt{--??????-}\\
2 & 1 & 2 & 1 & 2 & 256 & 256 & 0 & \texttt{---+-----}\\
3 & 2 & 3 & 1 & 2 & 0 & 512 & 0 & \texttt{---------}\\
4 & 1 & 2 & 1 & 2 & 256 & 256 & 0 & \texttt{-----+---}\\
5 & 4 & 3 & 1 & 2 & 0 & 512 & 0 & \texttt{---------}\\
6 & 1 & 2 & 1 & 2 & 0 & 512 & 0 & \texttt{---------}\\
7 & 1 & 2 & 1 & 2 & 0 & 512 & 0 & \texttt{---------}\\
8 & 0 & 1 & 1 & 1 & 0 & 512 & 0 & \texttt{---------}\\
\multicolumn{9}{@{}l@{}}{\rule{0pt}{2.6ex}\textbf{T9.235}\enspace \texttt{012323212}\enspace $n{=}9$, $|H|{=}4$, len 4, br 2, $|\mathrm{Aut}|{=}2$, orb 7/5, hist-orb 3/2; M 8192, res 3696, rev 2596: ok}\\*
0 & -- & 0 & 4 & 1 & 246 & 2 & 264 & \texttt{-????????}\\*
1 & 0 & 1 & 3 & 1 & 144 & 16 & 352 & \texttt{--?????--}\\
2 & 1 & 2 & 1 & 2 & 256 & 256 & 0 & \texttt{---+-----}\\
3 & 2 & 3 & 1 & 2 & 0 & 512 & 0 & \texttt{---------}\\
4 & 1 & 2 & 1 & 2 & 256 & 256 & 0 & \texttt{-----+---}\\
5 & 4 & 3 & 1 & 2 & 0 & 512 & 0 & \texttt{---------}\\
6 & 1 & 2 & 1 & 1 & 0 & 512 & 0 & \texttt{---------}\\
7 & 0 & 1 & 1 & 1 & 256 & 256 & 0 & \texttt{--------+}\\
8 & 7 & 2 & 1 & 1 & 0 & 512 & 0 & \texttt{---------}\\
\multicolumn{9}{@{}l@{}}{\rule{0pt}{2.6ex}\textbf{T9.236}\enspace \texttt{012323211}\enspace $n{=}9$, $|H|{=}5$, len 4, br 2, $|\mathrm{Aut}|{=}4$, orb 6/3, hist-orb 3/1; M 7680, res 5008, rev 1528: ok}\\*
0 & -- & 0 & 5 & 1 & 82 & 2 & 428 & \texttt{-????????}\\*
1 & 0 & 1 & 3 & 1 & 144 & 16 & 352 & \texttt{--?????--}\\
2 & 1 & 2 & 1 & 2 & 256 & 256 & 0 & \texttt{---+-----}\\
3 & 2 & 3 & 1 & 2 & 0 & 512 & 0 & \texttt{---------}\\
4 & 1 & 2 & 1 & 2 & 256 & 256 & 0 & \texttt{-----+---}\\
5 & 4 & 3 & 1 & 2 & 0 & 512 & 0 & \texttt{---------}\\
6 & 1 & 2 & 1 & 1 & 0 & 512 & 0 & \texttt{---------}\\
7 & 0 & 1 & 1 & 2 & 0 & 512 & 0 & \texttt{---------}\\
8 & 0 & 1 & 1 & 2 & 0 & 512 & 0 & \texttt{---------}\\
\multicolumn{9}{@{}l@{}}{\rule{0pt}{2.6ex}\textbf{T9.237}\enspace \texttt{012323123}\enspace $n{=}9$, $|H|{=}3$, len 4, br 2, $|\mathrm{Aut}|{=}2$, orb 7/5, hist-orb 2/1; M 8704, res 1952, rev 3620: ok}\\*
0 & -- & 0 & 3 & 1 & 350 & 2 & 160 & \texttt{-????????}\\*
1 & 0 & 1 & 2 & 1 & 288 & 32 & 192 & \texttt{--????---}\\
2 & 1 & 2 & 1 & 2 & 256 & 256 & 0 & \texttt{---+-----}\\
3 & 2 & 3 & 1 & 2 & 0 & 512 & 0 & \texttt{---------}\\
4 & 1 & 2 & 1 & 2 & 256 & 256 & 0 & \texttt{-----+---}\\
5 & 4 & 3 & 1 & 2 & 0 & 512 & 0 & \texttt{---------}\\
6 & 0 & 1 & 1 & 1 & 384 & 128 & 0 & \texttt{-------++}\\
7 & 6 & 2 & 1 & 1 & 256 & 256 & 0 & \texttt{--------+}\\
8 & 7 & 3 & 1 & 1 & 0 & 512 & 0 & \texttt{---------}\\
\multicolumn{9}{@{}l@{}}{\rule{0pt}{2.6ex}\textbf{T9.238}\enspace \texttt{012323122}\enspace $n{=}9$, $|H|{=}4$, len 4, br 3, $|\mathrm{Aut}|{=}4$, orb 6/3, hist-orb 2/0; M 8192, res 3136, rev 2816: ok}\\*
0 & -- & 0 & 4 & 1 & 250 & 2 & 260 & \texttt{-????????}\\*
1 & 0 & 1 & 2 & 1 & 288 & 32 & 192 & \texttt{--????---}\\
2 & 1 & 2 & 1 & 2 & 256 & 256 & 0 & \texttt{---+-----}\\
3 & 2 & 3 & 1 & 2 & 0 & 512 & 0 & \texttt{---------}\\
4 & 1 & 2 & 1 & 2 & 256 & 256 & 0 & \texttt{-----+---}\\
5 & 4 & 3 & 1 & 2 & 0 & 512 & 0 & \texttt{---------}\\
6 & 0 & 1 & 2 & 1 & 128 & 128 & 256 & \texttt{-------??}\\
7 & 6 & 2 & 1 & 2 & 0 & 512 & 0 & \texttt{---------}\\
8 & 6 & 2 & 1 & 2 & 0 & 512 & 0 & \texttt{---------}\\
\multicolumn{9}{@{}l@{}}{\rule{0pt}{2.6ex}\textbf{T9.239}\enspace \texttt{012323121}\enspace $n{=}9$, $|H|{=}4$, len 4, br 2, $|\mathrm{Aut}|{=}2$, orb 7/5, hist-orb 3/2; M 7680, res 3552, rev 2324: ok}\\*
0 & -- & 0 & 4 & 1 & 150 & 2 & 360 & \texttt{-????????}\\*
1 & 0 & 1 & 2 & 1 & 288 & 32 & 192 & \texttt{--????---}\\
2 & 1 & 2 & 1 & 2 & 256 & 256 & 0 & \texttt{---+-----}\\
3 & 2 & 3 & 1 & 2 & 0 & 512 & 0 & \texttt{---------}\\
4 & 1 & 2 & 1 & 2 & 256 & 256 & 0 & \texttt{-----+---}\\
5 & 4 & 3 & 1 & 2 & 0 & 512 & 0 & \texttt{---------}\\
6 & 0 & 1 & 1 & 1 & 256 & 256 & 0 & \texttt{-------+-}\\
7 & 6 & 2 & 1 & 1 & 0 & 512 & 0 & \texttt{---------}\\
8 & 0 & 1 & 1 & 1 & 0 & 512 & 0 & \texttt{---------}\\
\multicolumn{9}{@{}l@{}}{\rule{0pt}{2.6ex}\textbf{T9.240}\enspace \texttt{012323111}\enspace $n{=}9$, $|H|{=}5$, len 4, br 2, $|\mathrm{Aut}|{=}12$, orb 5/2, hist-orb 2/0; M 7168, res 4352, rev 1576: ok}\\*
0 & -- & 0 & 5 & 1 & 50 & 2 & 460 & \texttt{-????????}\\*
1 & 0 & 1 & 2 & 1 & 288 & 32 & 192 & \texttt{--????---}\\
2 & 1 & 2 & 1 & 2 & 256 & 256 & 0 & \texttt{---+-----}\\
3 & 2 & 3 & 1 & 2 & 0 & 512 & 0 & \texttt{---------}\\
4 & 1 & 2 & 1 & 2 & 256 & 256 & 0 & \texttt{-----+---}\\
5 & 4 & 3 & 1 & 2 & 0 & 512 & 0 & \texttt{---------}\\
6 & 0 & 1 & 1 & 3 & 0 & 512 & 0 & \texttt{---------}\\
7 & 0 & 1 & 1 & 3 & 0 & 512 & 0 & \texttt{---------}\\
8 & 0 & 1 & 1 & 3 & 0 & 512 & 0 & \texttt{---------}\\
\multicolumn{9}{@{}l@{}}{\rule{0pt}{2.6ex}\textbf{T9.241}\enspace \texttt{012322222}\enspace $n{=}9$, $|H|{=}6$, len 4, br 1, $|\mathrm{Aut}|{=}120$, orb 5/4, hist-orb 2/1; M 8192, res 5208, rev 2036: ok}\\*
0 & -- & 0 & 6 & 1 & 262 & 2 & 248 & \texttt{-+???????}\\*
1 & 0 & 1 & 6 & 1 & 12 & 4 & 496 & \texttt{--???????}\\
2 & 1 & 2 & 1 & 1 & 256 & 256 & 0 & \texttt{---+-----}\\
3 & 2 & 3 & 1 & 1 & 0 & 512 & 0 & \texttt{---------}\\
4 & 1 & 2 & 1 & 5 & 0 & 512 & 0 & \texttt{---------}\\
5 & 1 & 2 & 1 & 5 & 0 & 512 & 0 & \texttt{---------}\\
6 & 1 & 2 & 1 & 5 & 0 & 512 & 0 & \texttt{---------}\\
7 & 1 & 2 & 1 & 5 & 0 & 512 & 0 & \texttt{---------}\\
8 & 1 & 2 & 1 & 5 & 0 & 512 & 0 & \texttt{---------}\\
\multicolumn{9}{@{}l@{}}{\rule{0pt}{2.6ex}\textbf{T9.242}\enspace \texttt{012322221}\enspace $n{=}9$, $|H|{=}6$, len 4, br 2, $|\mathrm{Aut}|{=}24$, orb 6/5, hist-orb 3/2; M 7680, res 5648, rev 1256: ok}\\*
0 & -- & 0 & 6 & 1 & 134 & 2 & 376 & \texttt{-????????}\\*
1 & 0 & 1 & 5 & 1 & 24 & 8 & 480 & \texttt{--??????-}\\
2 & 1 & 2 & 1 & 1 & 256 & 256 & 0 & \texttt{---+-----}\\
3 & 2 & 3 & 1 & 1 & 0 & 512 & 0 & \texttt{---------}\\
4 & 1 & 2 & 1 & 4 & 0 & 512 & 0 & \texttt{---------}\\
5 & 1 & 2 & 1 & 4 & 0 & 512 & 0 & \texttt{---------}\\
6 & 1 & 2 & 1 & 4 & 0 & 512 & 0 & \texttt{---------}\\
7 & 1 & 2 & 1 & 4 & 0 & 512 & 0 & \texttt{---------}\\
8 & 0 & 1 & 1 & 1 & 0 & 512 & 0 & \texttt{---------}\\
\multicolumn{9}{@{}l@{}}{\rule{0pt}{2.6ex}\textbf{T9.243}\enspace \texttt{012322212}\enspace $n{=}9$, $|H|{=}5$, len 4, br 2, $|\mathrm{Aut}|{=}6$, orb 7/6, hist-orb 3/2; M 7680, res 4416, rev 1948: ok}\\*
0 & -- & 0 & 5 & 1 & 210 & 2 & 300 & \texttt{-????????}\\*
1 & 0 & 1 & 4 & 1 & 48 & 16 & 448 & \texttt{--?????--}\\
2 & 1 & 2 & 1 & 1 & 256 & 256 & 0 & \texttt{---+-----}\\
3 & 2 & 3 & 1 & 1 & 0 & 512 & 0 & \texttt{---------}\\
4 & 1 & 2 & 1 & 3 & 0 & 512 & 0 & \texttt{---------}\\
5 & 1 & 2 & 1 & 3 & 0 & 512 & 0 & \texttt{---------}\\
6 & 1 & 2 & 1 & 3 & 0 & 512 & 0 & \texttt{---------}\\
7 & 0 & 1 & 1 & 1 & 256 & 256 & 0 & \texttt{--------+}\\
8 & 7 & 2 & 1 & 1 & 0 & 512 & 0 & \texttt{---------}\\
\multicolumn{9}{@{}l@{}}{\rule{0pt}{2.6ex}\textbf{T9.244}\enspace \texttt{012322211}\enspace $n{=}9$, $|H|{=}6$, len 4, br 2, $|\mathrm{Aut}|{=}12$, orb 6/4, hist-orb 3/1; M 7168, res 5536, rev 920: ok}\\*
0 & -- & 0 & 6 & 1 & 70 & 2 & 440 & \texttt{-????????}\\*
1 & 0 & 1 & 4 & 1 & 48 & 16 & 448 & \texttt{--?????--}\\
2 & 1 & 2 & 1 & 1 & 256 & 256 & 0 & \texttt{---+-----}\\
3 & 2 & 3 & 1 & 1 & 0 & 512 & 0 & \texttt{---------}\\
4 & 1 & 2 & 1 & 3 & 0 & 512 & 0 & \texttt{---------}\\
5 & 1 & 2 & 1 & 3 & 0 & 512 & 0 & \texttt{---------}\\
6 & 1 & 2 & 1 & 3 & 0 & 512 & 0 & \texttt{---------}\\
7 & 0 & 1 & 1 & 2 & 0 & 512 & 0 & \texttt{---------}\\
8 & 0 & 1 & 1 & 2 & 0 & 512 & 0 & \texttt{---------}\\
\multicolumn{9}{@{}l@{}}{\rule{0pt}{2.6ex}\textbf{T9.245}\enspace \texttt{012322123}\enspace $n{=}9$, $|H|{=}4$, len 4, br 2, $|\mathrm{Aut}|{=}2$, orb 8/7, hist-orb 3/2; M 8192, res 3296, rev 2700: ok}\\*
0 & -- & 0 & 4 & 1 & 266 & 2 & 244 & \texttt{-????????}\\*
1 & 0 & 1 & 3 & 1 & 96 & 32 & 384 & \texttt{--????---}\\
2 & 1 & 2 & 1 & 1 & 256 & 256 & 0 & \texttt{---+-----}\\
3 & 2 & 3 & 1 & 1 & 0 & 512 & 0 & \texttt{---------}\\
4 & 1 & 2 & 1 & 2 & 0 & 512 & 0 & \texttt{---------}\\
5 & 1 & 2 & 1 & 2 & 0 & 512 & 0 & \texttt{---------}\\
6 & 0 & 1 & 1 & 1 & 384 & 128 & 0 & \texttt{-------++}\\
7 & 6 & 2 & 1 & 1 & 256 & 256 & 0 & \texttt{--------+}\\
8 & 7 & 3 & 1 & 1 & 0 & 512 & 0 & \texttt{---------}\\
\multicolumn{9}{@{}l@{}}{\rule{0pt}{2.6ex}\textbf{T9.246}\enspace \texttt{012322122}\enspace $n{=}9$, $|H|{=}5$, len 4, br 3, $|\mathrm{Aut}|{=}4$, orb 7/5, hist-orb 3/1; M 7680, res 4288, rev 1920: ok}\\*
0 & -- & 0 & 5 & 1 & 190 & 2 & 320 & \texttt{-????????}\\*
1 & 0 & 1 & 3 & 1 & 96 & 32 & 384 & \texttt{--????---}\\
2 & 1 & 2 & 1 & 1 & 256 & 256 & 0 & \texttt{---+-----}\\
3 & 2 & 3 & 1 & 1 & 0 & 512 & 0 & \texttt{---------}\\
4 & 1 & 2 & 1 & 2 & 0 & 512 & 0 & \texttt{---------}\\
5 & 1 & 2 & 1 & 2 & 0 & 512 & 0 & \texttt{---------}\\
6 & 0 & 1 & 2 & 1 & 128 & 128 & 256 & \texttt{-------??}\\
7 & 6 & 2 & 1 & 2 & 0 & 512 & 0 & \texttt{---------}\\
8 & 6 & 2 & 1 & 2 & 0 & 512 & 0 & \texttt{---------}\\
\multicolumn{9}{@{}l@{}}{\rule{0pt}{2.6ex}\textbf{T9.247}\enspace \texttt{012322121}\enspace $n{=}9$, $|H|{=}5$, len 4, br 2, $|\mathrm{Aut}|{=}2$, orb 8/7, hist-orb 4/3; M 7168, res 4512, rev 1500: ok}\\*
0 & -- & 0 & 5 & 1 & 114 & 2 & 396 & \texttt{-????????}\\*
1 & 0 & 1 & 3 & 1 & 96 & 32 & 384 & \texttt{--????---}\\
2 & 1 & 2 & 1 & 1 & 256 & 256 & 0 & \texttt{---+-----}\\
3 & 2 & 3 & 1 & 1 & 0 & 512 & 0 & \texttt{---------}\\
4 & 1 & 2 & 1 & 2 & 0 & 512 & 0 & \texttt{---------}\\
5 & 1 & 2 & 1 & 2 & 0 & 512 & 0 & \texttt{---------}\\
6 & 0 & 1 & 1 & 1 & 256 & 256 & 0 & \texttt{-------+-}\\
7 & 6 & 2 & 1 & 1 & 0 & 512 & 0 & \texttt{---------}\\
8 & 0 & 1 & 1 & 1 & 0 & 512 & 0 & \texttt{---------}\\
\multicolumn{9}{@{}l@{}}{\rule{0pt}{2.6ex}\textbf{T9.248}\enspace \texttt{012322111}\enspace $n{=}9$, $|H|{=}6$, len 4, br 2, $|\mathrm{Aut}|{=}12$, orb 6/4, hist-orb 3/1; M 6656, res 5120, rev 824: ok}\\*
0 & -- & 0 & 6 & 1 & 38 & 2 & 472 & \texttt{-????????}\\*
1 & 0 & 1 & 3 & 1 & 96 & 32 & 384 & \texttt{--????---}\\
2 & 1 & 2 & 1 & 1 & 256 & 256 & 0 & \texttt{---+-----}\\
3 & 2 & 3 & 1 & 1 & 0 & 512 & 0 & \texttt{---------}\\
4 & 1 & 2 & 1 & 2 & 0 & 512 & 0 & \texttt{---------}\\
5 & 1 & 2 & 1 & 2 & 0 & 512 & 0 & \texttt{---------}\\
6 & 0 & 1 & 1 & 3 & 0 & 512 & 0 & \texttt{---------}\\
7 & 0 & 1 & 1 & 3 & 0 & 512 & 0 & \texttt{---------}\\
8 & 0 & 1 & 1 & 3 & 0 & 512 & 0 & \texttt{---------}\\
\multicolumn{9}{@{}l@{}}{\rule{0pt}{2.6ex}\textbf{T9.249}\enspace \texttt{012321232}\enspace $n{=}9$, $|H|{=}4$, len 4, br 3, $|\mathrm{Aut}|{=}2$, orb 5/1, hist-orb 2/0; M 8192, res 3424, rev 2640: ok}\\*
0 & -- & 0 & 4 & 1 & 242 & 2 & 268 & \texttt{-????????}\\*
1 & 0 & 1 & 2 & 2 & 192 & 64 & 256 & \texttt{--???----}\\
2 & 1 & 2 & 1 & 2 & 256 & 256 & 0 & \texttt{---+-----}\\
3 & 2 & 3 & 1 & 2 & 0 & 512 & 0 & \texttt{---------}\\
4 & 1 & 2 & 1 & 2 & 0 & 512 & 0 & \texttt{---------}\\
5 & 0 & 1 & 2 & 2 & 192 & 64 & 256 & \texttt{------???}\\
6 & 5 & 2 & 1 & 2 & 256 & 256 & 0 & \texttt{-------+-}\\
7 & 6 & 3 & 1 & 2 & 0 & 512 & 0 & \texttt{---------}\\
8 & 5 & 2 & 1 & 2 & 0 & 512 & 0 & \texttt{---------}\\
\multicolumn{9}{@{}l@{}}{\rule{0pt}{2.6ex}\textbf{T9.250}\enspace \texttt{012321231}\enspace $n{=}9$, $|H|{=}4$, len 4, br 2, $|\mathrm{Aut}|{=}1$, orb 9/9, hist-orb 4/4; M 7680, res 3488, rev 2260: ok}\\*
0 & -- & 0 & 4 & 1 & 154 & 2 & 356 & \texttt{-????????}\\*
1 & 0 & 1 & 2 & 1 & 192 & 64 & 256 & \texttt{--???----}\\
2 & 1 & 2 & 1 & 1 & 256 & 256 & 0 & \texttt{---+-----}\\
3 & 2 & 3 & 1 & 1 & 0 & 512 & 0 & \texttt{---------}\\
4 & 1 & 2 & 1 & 1 & 0 & 512 & 0 & \texttt{---------}\\
5 & 0 & 1 & 1 & 1 & 384 & 128 & 0 & \texttt{------++-}\\
6 & 5 & 2 & 1 & 1 & 256 & 256 & 0 & \texttt{-------+-}\\
7 & 6 & 3 & 1 & 1 & 0 & 512 & 0 & \texttt{---------}\\
8 & 0 & 1 & 1 & 1 & 0 & 512 & 0 & \texttt{---------}\\
\multicolumn{9}{@{}l@{}}{\rule{0pt}{2.6ex}\textbf{T9.251}\enspace \texttt{012321222}\enspace $n{=}9$, $|H|{=}5$, len 4, br 3, $|\mathrm{Aut}|{=}6$, orb 7/6, hist-orb 3/2; M 7680, res 4096, rev 2016: ok}\\*
0 & -- & 0 & 5 & 1 & 198 & 2 & 312 & \texttt{-????????}\\*
1 & 0 & 1 & 2 & 1 & 192 & 64 & 256 & \texttt{--???----}\\
2 & 1 & 2 & 1 & 1 & 256 & 256 & 0 & \texttt{---+-----}\\
3 & 2 & 3 & 1 & 1 & 0 & 512 & 0 & \texttt{---------}\\
4 & 1 & 2 & 1 & 1 & 0 & 512 & 0 & \texttt{---------}\\
5 & 0 & 1 & 3 & 1 & 64 & 64 & 384 & \texttt{------???}\\
6 & 5 & 2 & 1 & 3 & 0 & 512 & 0 & \texttt{---------}\\
7 & 5 & 2 & 1 & 3 & 0 & 512 & 0 & \texttt{---------}\\
8 & 5 & 2 & 1 & 3 & 0 & 512 & 0 & \texttt{---------}\\
\multicolumn{9}{@{}l@{}}{\rule{0pt}{2.6ex}\textbf{T9.252}\enspace \texttt{012321221}\enspace $n{=}9$, $|H|{=}5$, len 4, br 3, $|\mathrm{Aut}|{=}2$, orb 8/7, hist-orb 4/3; M 7168, res 4224, rev 1592: ok}\\*
0 & -- & 0 & 5 & 1 & 110 & 2 & 400 & \texttt{-????????}\\*
1 & 0 & 1 & 2 & 1 & 192 & 64 & 256 & \texttt{--???----}\\
2 & 1 & 2 & 1 & 1 & 256 & 256 & 0 & \texttt{---+-----}\\
3 & 2 & 3 & 1 & 1 & 0 & 512 & 0 & \texttt{---------}\\
4 & 1 & 2 & 1 & 1 & 0 & 512 & 0 & \texttt{---------}\\
5 & 0 & 1 & 2 & 1 & 128 & 128 & 256 & \texttt{------??-}\\
6 & 5 & 2 & 1 & 2 & 0 & 512 & 0 & \texttt{---------}\\
7 & 5 & 2 & 1 & 2 & 0 & 512 & 0 & \texttt{---------}\\
8 & 0 & 1 & 1 & 1 & 0 & 512 & 0 & \texttt{---------}\\
\multicolumn{9}{@{}l@{}}{\rule{0pt}{2.6ex}\textbf{T9.253}\enspace \texttt{012321212}\enspace $n{=}9$, $|H|{=}4$, len 4, br 2, $|\mathrm{Aut}|{=}2$, orb 7/5, hist-orb 3/2; M 7168, res 3136, rev 2248: ok}\\*
0 & -- & 0 & 4 & 1 & 198 & 2 & 312 & \texttt{-????????}\\*
1 & 0 & 1 & 2 & 1 & 192 & 64 & 256 & \texttt{--???----}\\
2 & 1 & 2 & 1 & 1 & 256 & 256 & 0 & \texttt{---+-----}\\
3 & 2 & 3 & 1 & 1 & 0 & 512 & 0 & \texttt{---------}\\
4 & 1 & 2 & 1 & 1 & 0 & 512 & 0 & \texttt{---------}\\
5 & 0 & 1 & 1 & 2 & 256 & 256 & 0 & \texttt{------+--}\\
6 & 5 & 2 & 1 & 2 & 0 & 512 & 0 & \texttt{---------}\\
7 & 0 & 1 & 1 & 2 & 256 & 256 & 0 & \texttt{--------+}\\
8 & 7 & 2 & 1 & 2 & 0 & 512 & 0 & \texttt{---------}\\
\multicolumn{9}{@{}l@{}}{\rule{0pt}{2.6ex}\textbf{T9.254}\enspace \texttt{012321211}\enspace $n{=}9$, $|H|{=}5$, len 4, br 2, $|\mathrm{Aut}|{=}2$, orb 8/7, hist-orb 4/3; M 6656, res 4192, rev 1348: ok}\\*
0 & -- & 0 & 5 & 1 & 66 & 2 & 444 & \texttt{-????????}\\*
1 & 0 & 1 & 2 & 1 & 192 & 64 & 256 & \texttt{--???----}\\
2 & 1 & 2 & 1 & 1 & 256 & 256 & 0 & \texttt{---+-----}\\
3 & 2 & 3 & 1 & 1 & 0 & 512 & 0 & \texttt{---------}\\
4 & 1 & 2 & 1 & 1 & 0 & 512 & 0 & \texttt{---------}\\
5 & 0 & 1 & 1 & 1 & 256 & 256 & 0 & \texttt{------+--}\\
6 & 5 & 2 & 1 & 1 & 0 & 512 & 0 & \texttt{---------}\\
7 & 0 & 1 & 1 & 2 & 0 & 512 & 0 & \texttt{---------}\\
8 & 0 & 1 & 1 & 2 & 0 & 512 & 0 & \texttt{---------}\\
\multicolumn{9}{@{}l@{}}{\rule{0pt}{2.6ex}\textbf{T9.255}\enspace \texttt{012321111}\enspace $n{=}9$, $|H|{=}6$, len 4, br 2, $|\mathrm{Aut}|{=}24$, orb 6/5, hist-orb 3/2; M 6144, res 4544, rev 848: ok}\\*
0 & -- & 0 & 6 & 1 & 22 & 2 & 488 & \texttt{-????????}\\*
1 & 0 & 1 & 2 & 1 & 192 & 64 & 256 & \texttt{--???----}\\
2 & 1 & 2 & 1 & 1 & 256 & 256 & 0 & \texttt{---+-----}\\
3 & 2 & 3 & 1 & 1 & 0 & 512 & 0 & \texttt{---------}\\
4 & 1 & 2 & 1 & 1 & 0 & 512 & 0 & \texttt{---------}\\
5 & 0 & 1 & 1 & 4 & 0 & 512 & 0 & \texttt{---------}\\
6 & 0 & 1 & 1 & 4 & 0 & 512 & 0 & \texttt{---------}\\
7 & 0 & 1 & 1 & 4 & 0 & 512 & 0 & \texttt{---------}\\
8 & 0 & 1 & 1 & 4 & 0 & 512 & 0 & \texttt{---------}\\
\multicolumn{9}{@{}l@{}}{\rule{0pt}{2.6ex}\textbf{T9.256}\enspace \texttt{012312312}\enspace $n{=}9$, $|H|{=}3$, len 4, br 1, $|\mathrm{Aut}|{=}2$, orb 6/3, hist-orb 2/1; M 7680, res 1728, rev 3108: ok}\\*
0 & -- & 0 & 3 & 1 & 294 & 2 & 216 & \texttt{-????????}\\*
1 & 0 & 1 & 1 & 2 & 384 & 128 & 0 & \texttt{--++-----}\\
2 & 1 & 2 & 1 & 2 & 256 & 256 & 0 & \texttt{---+-----}\\
3 & 2 & 3 & 1 & 2 & 0 & 512 & 0 & \texttt{---------}\\
4 & 0 & 1 & 1 & 2 & 384 & 128 & 0 & \texttt{-----++--}\\
5 & 4 & 2 & 1 & 2 & 256 & 256 & 0 & \texttt{------+--}\\
6 & 5 & 3 & 1 & 2 & 0 & 512 & 0 & \texttt{---------}\\
7 & 0 & 1 & 1 & 1 & 256 & 256 & 0 & \texttt{--------+}\\
8 & 7 & 2 & 1 & 1 & 0 & 512 & 0 & \texttt{---------}\\
\multicolumn{9}{@{}l@{}}{\rule{0pt}{2.6ex}\textbf{T9.257}\enspace \texttt{012312311}\enspace $n{=}9$, $|H|{=}4$, len 4, br 1, $|\mathrm{Aut}|{=}4$, orb 5/1, hist-orb 2/0; M 7168, res 3296, rev 2040: ok}\\*
0 & -- & 0 & 4 & 1 & 98 & 2 & 412 & \texttt{-????????}\\*
1 & 0 & 1 & 1 & 2 & 384 & 128 & 0 & \texttt{--++-----}\\
2 & 1 & 2 & 1 & 2 & 256 & 256 & 0 & \texttt{---+-----}\\
3 & 2 & 3 & 1 & 2 & 0 & 512 & 0 & \texttt{---------}\\
4 & 0 & 1 & 1 & 2 & 384 & 128 & 0 & \texttt{-----++--}\\
5 & 4 & 2 & 1 & 2 & 256 & 256 & 0 & \texttt{------+--}\\
6 & 5 & 3 & 1 & 2 & 0 & 512 & 0 & \texttt{---------}\\
7 & 0 & 1 & 1 & 2 & 0 & 512 & 0 & \texttt{---------}\\
8 & 0 & 1 & 1 & 2 & 0 & 512 & 0 & \texttt{---------}\\
\multicolumn{9}{@{}l@{}}{\rule{0pt}{2.6ex}\textbf{T9.258}\enspace \texttt{012312222}\enspace $n{=}9$, $|H|{=}5$, len 4, br 2, $|\mathrm{Aut}|{=}24$, orb 6/5, hist-orb 2/1; M 7680, res 3744, rev 2236: ok}\\*
0 & -- & 0 & 5 & 1 & 238 & 2 & 272 & \texttt{-????????}\\*
1 & 0 & 1 & 1 & 1 & 384 & 128 & 0 & \texttt{--++-----}\\
2 & 1 & 2 & 1 & 1 & 256 & 256 & 0 & \texttt{---+-----}\\
3 & 2 & 3 & 1 & 1 & 0 & 512 & 0 & \texttt{---------}\\
4 & 0 & 1 & 4 & 1 & 32 & 32 & 448 & \texttt{-----????}\\
5 & 4 & 2 & 1 & 4 & 0 & 512 & 0 & \texttt{---------}\\
6 & 4 & 2 & 1 & 4 & 0 & 512 & 0 & \texttt{---------}\\
7 & 4 & 2 & 1 & 4 & 0 & 512 & 0 & \texttt{---------}\\
8 & 4 & 2 & 1 & 4 & 0 & 512 & 0 & \texttt{---------}\\
\multicolumn{9}{@{}l@{}}{\rule{0pt}{2.6ex}\textbf{T9.259}\enspace \texttt{012312221}\enspace $n{=}9$, $|H|{=}5$, len 4, br 2, $|\mathrm{Aut}|{=}6$, orb 7/6, hist-orb 3/2; M 7168, res 4032, rev 1676: ok}\\*
0 & -- & 0 & 5 & 1 & 126 & 2 & 384 & \texttt{-????????}\\*
1 & 0 & 1 & 1 & 1 & 384 & 128 & 0 & \texttt{--++-----}\\
2 & 1 & 2 & 1 & 1 & 256 & 256 & 0 & \texttt{---+-----}\\
3 & 2 & 3 & 1 & 1 & 0 & 512 & 0 & \texttt{---------}\\
4 & 0 & 1 & 3 & 1 & 64 & 64 & 384 & \texttt{-----???-}\\
5 & 4 & 2 & 1 & 3 & 0 & 512 & 0 & \texttt{---------}\\
6 & 4 & 2 & 1 & 3 & 0 & 512 & 0 & \texttt{---------}\\
7 & 4 & 2 & 1 & 3 & 0 & 512 & 0 & \texttt{---------}\\
8 & 0 & 1 & 1 & 1 & 0 & 512 & 0 & \texttt{---------}\\
\multicolumn{9}{@{}l@{}}{\rule{0pt}{2.6ex}\textbf{T9.260}\enspace \texttt{012312212}\enspace $n{=}9$, $|H|{=}4$, len 4, br 2, $|\mathrm{Aut}|{=}2$, orb 8/7, hist-orb 3/2; M 7168, res 2784, rev 2352: ok}\\*
0 & -- & 0 & 4 & 1 & 210 & 2 & 300 & \texttt{-????????}\\*
1 & 0 & 1 & 1 & 1 & 384 & 128 & 0 & \texttt{--++-----}\\
2 & 1 & 2 & 1 & 1 & 256 & 256 & 0 & \texttt{---+-----}\\
3 & 2 & 3 & 1 & 1 & 0 & 512 & 0 & \texttt{---------}\\
4 & 0 & 1 & 2 & 1 & 128 & 128 & 256 & \texttt{-----??--}\\
5 & 4 & 2 & 1 & 2 & 0 & 512 & 0 & \texttt{---------}\\
6 & 4 & 2 & 1 & 2 & 0 & 512 & 0 & \texttt{---------}\\
7 & 0 & 1 & 1 & 1 & 256 & 256 & 0 & \texttt{--------+}\\
8 & 7 & 2 & 1 & 1 & 0 & 512 & 0 & \texttt{---------}\\
\multicolumn{9}{@{}l@{}}{\rule{0pt}{2.6ex}\textbf{T9.261}\enspace \texttt{012312211}\enspace $n{=}9$, $|H|{=}5$, len 4, br 2, $|\mathrm{Aut}|{=}4$, orb 7/5, hist-orb 3/1; M 6656, res 3904, rev 1428: ok}\\*
0 & -- & 0 & 5 & 1 & 70 & 2 & 440 & \texttt{-????????}\\*
1 & 0 & 1 & 1 & 1 & 384 & 128 & 0 & \texttt{--++-----}\\
2 & 1 & 2 & 1 & 1 & 256 & 256 & 0 & \texttt{---+-----}\\
3 & 2 & 3 & 1 & 1 & 0 & 512 & 0 & \texttt{---------}\\
4 & 0 & 1 & 2 & 1 & 128 & 128 & 256 & \texttt{-----??--}\\
5 & 4 & 2 & 1 & 2 & 0 & 512 & 0 & \texttt{---------}\\
6 & 4 & 2 & 1 & 2 & 0 & 512 & 0 & \texttt{---------}\\
7 & 0 & 1 & 1 & 2 & 0 & 512 & 0 & \texttt{---------}\\
8 & 0 & 1 & 1 & 2 & 0 & 512 & 0 & \texttt{---------}\\
\multicolumn{9}{@{}l@{}}{\rule{0pt}{2.6ex}\textbf{T9.262}\enspace \texttt{012312121}\enspace $n{=}9$, $|H|{=}4$, len 4, br 1, $|\mathrm{Aut}|{=}2$, orb 7/5, hist-orb 3/2; M 6656, res 3072, rev 1940: ok}\\*
0 & -- & 0 & 4 & 1 & 126 & 2 & 384 & \texttt{-????????}\\*
1 & 0 & 1 & 1 & 1 & 384 & 128 & 0 & \texttt{--++-----}\\
2 & 1 & 2 & 1 & 1 & 256 & 256 & 0 & \texttt{---+-----}\\
3 & 2 & 3 & 1 & 1 & 0 & 512 & 0 & \texttt{---------}\\
4 & 0 & 1 & 1 & 2 & 256 & 256 & 0 & \texttt{-----+---}\\
5 & 4 & 2 & 1 & 2 & 0 & 512 & 0 & \texttt{---------}\\
6 & 0 & 1 & 1 & 2 & 256 & 256 & 0 & \texttt{-------+-}\\
7 & 6 & 2 & 1 & 2 & 0 & 512 & 0 & \texttt{---------}\\
8 & 0 & 1 & 1 & 1 & 0 & 512 & 0 & \texttt{---------}\\
\multicolumn{9}{@{}l@{}}{\rule{0pt}{2.6ex}\textbf{T9.263}\enspace \texttt{012312111}\enspace $n{=}9$, $|H|{=}5$, len 4, br 1, $|\mathrm{Aut}|{=}6$, orb 7/6, hist-orb 3/2; M 6144, res 3744, rev 1272: ok}\\*
0 & -- & 0 & 5 & 1 & 42 & 2 & 468 & \texttt{-????????}\\*
1 & 0 & 1 & 1 & 1 & 384 & 128 & 0 & \texttt{--++-----}\\
2 & 1 & 2 & 1 & 1 & 256 & 256 & 0 & \texttt{---+-----}\\
3 & 2 & 3 & 1 & 1 & 0 & 512 & 0 & \texttt{---------}\\
4 & 0 & 1 & 1 & 1 & 256 & 256 & 0 & \texttt{-----+---}\\
5 & 4 & 2 & 1 & 1 & 0 & 512 & 0 & \texttt{---------}\\
6 & 0 & 1 & 1 & 3 & 0 & 512 & 0 & \texttt{---------}\\
7 & 0 & 1 & 1 & 3 & 0 & 512 & 0 & \texttt{---------}\\
8 & 0 & 1 & 1 & 3 & 0 & 512 & 0 & \texttt{---------}\\
\multicolumn{9}{@{}l@{}}{\rule{0pt}{2.6ex}\textbf{T9.264}\enspace \texttt{012311111}\enspace $n{=}9$, $|H|{=}6$, len 4, br 1, $|\mathrm{Aut}|{=}120$, orb 5/4, hist-orb 2/1; M 5632, res 3968, rev 860: ok}\\*
0 & -- & 0 & 6 & 1 & 14 & 2 & 496 & \texttt{-????????}\\*
1 & 0 & 1 & 1 & 1 & 384 & 128 & 0 & \texttt{--++-----}\\
2 & 1 & 2 & 1 & 1 & 256 & 256 & 0 & \texttt{---+-----}\\
3 & 2 & 3 & 1 & 1 & 0 & 512 & 0 & \texttt{---------}\\
4 & 0 & 1 & 1 & 5 & 0 & 512 & 0 & \texttt{---------}\\
5 & 0 & 1 & 1 & 5 & 0 & 512 & 0 & \texttt{---------}\\
6 & 0 & 1 & 1 & 5 & 0 & 512 & 0 & \texttt{---------}\\
7 & 0 & 1 & 1 & 5 & 0 & 512 & 0 & \texttt{---------}\\
8 & 0 & 1 & 1 & 5 & 0 & 512 & 0 & \texttt{---------}\\
\multicolumn{9}{@{}l@{}}{\rule{0pt}{2.6ex}\textbf{T9.265}\enspace \texttt{012222222}\enspace $n{=}9$, $|H|{=}7$, len 3, br 1, $|\mathrm{Aut}|{=}5040$, orb 3/2, hist-orb 1/0; M 7680, res 5292, rev 1836: ok}\\*
0 & -- & 0 & 7 & 1 & 258 & 2 & 252 & \texttt{-+???????}\\*
1 & 0 & 1 & 7 & 1 & 4 & 4 & 504 & \texttt{--???????}\\
2 & 1 & 2 & 1 & 7 & 0 & 512 & 0 & \texttt{---------}\\
3 & 1 & 2 & 1 & 7 & 0 & 512 & 0 & \texttt{---------}\\
4 & 1 & 2 & 1 & 7 & 0 & 512 & 0 & \texttt{---------}\\
5 & 1 & 2 & 1 & 7 & 0 & 512 & 0 & \texttt{---------}\\
6 & 1 & 2 & 1 & 7 & 0 & 512 & 0 & \texttt{---------}\\
7 & 1 & 2 & 1 & 7 & 0 & 512 & 0 & \texttt{---------}\\
8 & 1 & 2 & 1 & 7 & 0 & 512 & 0 & \texttt{---------}\\
\multicolumn{9}{@{}l@{}}{\rule{0pt}{2.6ex}\textbf{T9.266}\enspace \texttt{012222221}\enspace $n{=}9$, $|H|{=}7$, len 3, br 2, $|\mathrm{Aut}|{=}720$, orb 4/3, hist-orb 2/1; M 7168, res 5768, rev 960: ok}\\*
0 & -- & 0 & 7 & 1 & 130 & 2 & 380 & \texttt{-????????}\\*
1 & 0 & 1 & 6 & 1 & 8 & 8 & 496 & \texttt{--??????-}\\
2 & 1 & 2 & 1 & 6 & 0 & 512 & 0 & \texttt{---------}\\
3 & 1 & 2 & 1 & 6 & 0 & 512 & 0 & \texttt{---------}\\
4 & 1 & 2 & 1 & 6 & 0 & 512 & 0 & \texttt{---------}\\
5 & 1 & 2 & 1 & 6 & 0 & 512 & 0 & \texttt{---------}\\
6 & 1 & 2 & 1 & 6 & 0 & 512 & 0 & \texttt{---------}\\
7 & 1 & 2 & 1 & 6 & 0 & 512 & 0 & \texttt{---------}\\
8 & 0 & 1 & 1 & 1 & 0 & 512 & 0 & \texttt{---------}\\
\multicolumn{9}{@{}l@{}}{\rule{0pt}{2.6ex}\textbf{T9.267}\enspace \texttt{012222212}\enspace $n{=}9$, $|H|{=}6$, len 3, br 2, $|\mathrm{Aut}|{=}120$, orb 5/4, hist-orb 2/1; M 7168, res 4656, rev 1596: ok}\\*
0 & -- & 0 & 6 & 1 & 198 & 2 & 312 & \texttt{-????????}\\*
1 & 0 & 1 & 5 & 1 & 16 & 16 & 480 & \texttt{--?????--}\\
2 & 1 & 2 & 1 & 5 & 0 & 512 & 0 & \texttt{---------}\\
3 & 1 & 2 & 1 & 5 & 0 & 512 & 0 & \texttt{---------}\\
4 & 1 & 2 & 1 & 5 & 0 & 512 & 0 & \texttt{---------}\\
5 & 1 & 2 & 1 & 5 & 0 & 512 & 0 & \texttt{---------}\\
6 & 1 & 2 & 1 & 5 & 0 & 512 & 0 & \texttt{---------}\\
7 & 0 & 1 & 1 & 1 & 256 & 256 & 0 & \texttt{--------+}\\
8 & 7 & 2 & 1 & 1 & 0 & 512 & 0 & \texttt{---------}\\
\multicolumn{9}{@{}l@{}}{\rule{0pt}{2.6ex}\textbf{T9.268}\enspace \texttt{012222211}\enspace $n{=}9$, $|H|{=}7$, len 3, br 2, $|\mathrm{Aut}|{=}240$, orb 4/2, hist-orb 2/0; M 6656, res 5712, rev 544: ok}\\*
0 & -- & 0 & 7 & 1 & 66 & 2 & 444 & \texttt{-????????}\\*
1 & 0 & 1 & 5 & 1 & 16 & 16 & 480 & \texttt{--?????--}\\
2 & 1 & 2 & 1 & 5 & 0 & 512 & 0 & \texttt{---------}\\
3 & 1 & 2 & 1 & 5 & 0 & 512 & 0 & \texttt{---------}\\
4 & 1 & 2 & 1 & 5 & 0 & 512 & 0 & \texttt{---------}\\
5 & 1 & 2 & 1 & 5 & 0 & 512 & 0 & \texttt{---------}\\
6 & 1 & 2 & 1 & 5 & 0 & 512 & 0 & \texttt{---------}\\
7 & 0 & 1 & 1 & 2 & 0 & 512 & 0 & \texttt{---------}\\
8 & 0 & 1 & 1 & 2 & 0 & 512 & 0 & \texttt{---------}\\
\multicolumn{9}{@{}l@{}}{\rule{0pt}{2.6ex}\textbf{T9.269}\enspace \texttt{012222122}\enspace $n{=}9$, $|H|{=}6$, len 3, br 3, $|\mathrm{Aut}|{=}48$, orb 5/3, hist-orb 2/0; M 7168, res 4672, rev 1448: ok}\\*
0 & -- & 0 & 6 & 1 & 170 & 2 & 340 & \texttt{-????????}\\*
1 & 0 & 1 & 4 & 1 & 32 & 32 & 448 & \texttt{--????---}\\
2 & 1 & 2 & 1 & 4 & 0 & 512 & 0 & \texttt{---------}\\
3 & 1 & 2 & 1 & 4 & 0 & 512 & 0 & \texttt{---------}\\
4 & 1 & 2 & 1 & 4 & 0 & 512 & 0 & \texttt{---------}\\
5 & 1 & 2 & 1 & 4 & 0 & 512 & 0 & \texttt{---------}\\
6 & 0 & 1 & 2 & 1 & 128 & 128 & 256 & \texttt{-------??}\\
7 & 6 & 2 & 1 & 2 & 0 & 512 & 0 & \texttt{---------}\\
8 & 6 & 2 & 1 & 2 & 0 & 512 & 0 & \texttt{---------}\\
\multicolumn{9}{@{}l@{}}{\rule{0pt}{2.6ex}\textbf{T9.270}\enspace \texttt{012222121}\enspace $n{=}9$, $|H|{=}6$, len 3, br 2, $|\mathrm{Aut}|{=}24$, orb 6/5, hist-orb 3/2; M 6656, res 4832, rev 1036: ok}\\*
0 & -- & 0 & 6 & 1 & 102 & 2 & 408 & \texttt{-????????}\\*
1 & 0 & 1 & 4 & 1 & 32 & 32 & 448 & \texttt{--????---}\\
2 & 1 & 2 & 1 & 4 & 0 & 512 & 0 & \texttt{---------}\\
3 & 1 & 2 & 1 & 4 & 0 & 512 & 0 & \texttt{---------}\\
4 & 1 & 2 & 1 & 4 & 0 & 512 & 0 & \texttt{---------}\\
5 & 1 & 2 & 1 & 4 & 0 & 512 & 0 & \texttt{---------}\\
6 & 0 & 1 & 1 & 1 & 256 & 256 & 0 & \texttt{-------+-}\\
7 & 6 & 2 & 1 & 1 & 0 & 512 & 0 & \texttt{---------}\\
8 & 0 & 1 & 1 & 1 & 0 & 512 & 0 & \texttt{---------}\\
\multicolumn{9}{@{}l@{}}{\rule{0pt}{2.6ex}\textbf{T9.271}\enspace \texttt{012222111}\enspace $n{=}9$, $|H|{=}7$, len 3, br 2, $|\mathrm{Aut}|{=}144$, orb 4/2, hist-orb 2/0; M 6144, res 5376, rev 368: ok}\\*
0 & -- & 0 & 7 & 1 & 34 & 2 & 476 & \texttt{-????????}\\*
1 & 0 & 1 & 4 & 1 & 32 & 32 & 448 & \texttt{--????---}\\
2 & 1 & 2 & 1 & 4 & 0 & 512 & 0 & \texttt{---------}\\
3 & 1 & 2 & 1 & 4 & 0 & 512 & 0 & \texttt{---------}\\
4 & 1 & 2 & 1 & 4 & 0 & 512 & 0 & \texttt{---------}\\
5 & 1 & 2 & 1 & 4 & 0 & 512 & 0 & \texttt{---------}\\
6 & 0 & 1 & 1 & 3 & 0 & 512 & 0 & \texttt{---------}\\
7 & 0 & 1 & 1 & 3 & 0 & 512 & 0 & \texttt{---------}\\
8 & 0 & 1 & 1 & 3 & 0 & 512 & 0 & \texttt{---------}\\
\multicolumn{9}{@{}l@{}}{\rule{0pt}{2.6ex}\textbf{T9.272}\enspace \texttt{012221222}\enspace $n{=}9$, $|H|{=}6$, len 3, br 3, $|\mathrm{Aut}|{=}72$, orb 3/1, hist-orb 1/0; M 7168, res 4704, rev 1392: ok}\\*
0 & -- & 0 & 6 & 1 & 162 & 2 & 348 & \texttt{-????????}\\*
1 & 0 & 1 & 3 & 2 & 64 & 64 & 384 & \texttt{--???----}\\
2 & 1 & 2 & 1 & 6 & 0 & 512 & 0 & \texttt{---------}\\
3 & 1 & 2 & 1 & 6 & 0 & 512 & 0 & \texttt{---------}\\
4 & 1 & 2 & 1 & 6 & 0 & 512 & 0 & \texttt{---------}\\
5 & 0 & 1 & 3 & 2 & 64 & 64 & 384 & \texttt{------???}\\
6 & 5 & 2 & 1 & 6 & 0 & 512 & 0 & \texttt{---------}\\
7 & 5 & 2 & 1 & 6 & 0 & 512 & 0 & \texttt{---------}\\
8 & 5 & 2 & 1 & 6 & 0 & 512 & 0 & \texttt{---------}\\
\multicolumn{9}{@{}l@{}}{\rule{0pt}{2.6ex}\textbf{T9.273}\enspace \texttt{012221221}\enspace $n{=}9$, $|H|{=}6$, len 3, br 3, $|\mathrm{Aut}|{=}12$, orb 6/4, hist-orb 3/1; M 6656, res 4704, rev 1016: ok}\\*
0 & -- & 0 & 6 & 1 & 90 & 2 & 420 & \texttt{-????????}\\*
1 & 0 & 1 & 3 & 1 & 64 & 64 & 384 & \texttt{--???----}\\
2 & 1 & 2 & 1 & 3 & 0 & 512 & 0 & \texttt{---------}\\
3 & 1 & 2 & 1 & 3 & 0 & 512 & 0 & \texttt{---------}\\
4 & 1 & 2 & 1 & 3 & 0 & 512 & 0 & \texttt{---------}\\
5 & 0 & 1 & 2 & 1 & 128 & 128 & 256 & \texttt{------??-}\\
6 & 5 & 2 & 1 & 2 & 0 & 512 & 0 & \texttt{---------}\\
7 & 5 & 2 & 1 & 2 & 0 & 512 & 0 & \texttt{---------}\\
8 & 0 & 1 & 1 & 1 & 0 & 512 & 0 & \texttt{---------}\\
\multicolumn{9}{@{}l@{}}{\rule{0pt}{2.6ex}\textbf{T9.274}\enspace \texttt{012221212}\enspace $n{=}9$, $|H|{=}5$, len 3, br 2, $|\mathrm{Aut}|{=}12$, orb 5/2, hist-orb 2/0; M 6656, res 3744, rev 1640: ok}\\*
0 & -- & 0 & 5 & 1 & 162 & 2 & 348 & \texttt{-????????}\\*
1 & 0 & 1 & 3 & 1 & 64 & 64 & 384 & \texttt{--???----}\\
2 & 1 & 2 & 1 & 3 & 0 & 512 & 0 & \texttt{---------}\\
3 & 1 & 2 & 1 & 3 & 0 & 512 & 0 & \texttt{---------}\\
4 & 1 & 2 & 1 & 3 & 0 & 512 & 0 & \texttt{---------}\\
5 & 0 & 1 & 1 & 2 & 256 & 256 & 0 & \texttt{------+--}\\
6 & 5 & 2 & 1 & 2 & 0 & 512 & 0 & \texttt{---------}\\
7 & 0 & 1 & 1 & 2 & 256 & 256 & 0 & \texttt{--------+}\\
8 & 7 & 2 & 1 & 2 & 0 & 512 & 0 & \texttt{---------}\\
\multicolumn{9}{@{}l@{}}{\rule{0pt}{2.6ex}\textbf{T9.275}\enspace \texttt{012221211}\enspace $n{=}9$, $|H|{=}6$, len 3, br 2, $|\mathrm{Aut}|{=}12$, orb 6/4, hist-orb 3/1; M 6144, res 4608, rev 796: ok}\\*
0 & -- & 0 & 6 & 1 & 54 & 2 & 456 & \texttt{-????????}\\*
1 & 0 & 1 & 3 & 1 & 64 & 64 & 384 & \texttt{--???----}\\
2 & 1 & 2 & 1 & 3 & 0 & 512 & 0 & \texttt{---------}\\
3 & 1 & 2 & 1 & 3 & 0 & 512 & 0 & \texttt{---------}\\
4 & 1 & 2 & 1 & 3 & 0 & 512 & 0 & \texttt{---------}\\
5 & 0 & 1 & 1 & 1 & 256 & 256 & 0 & \texttt{------+--}\\
6 & 5 & 2 & 1 & 1 & 0 & 512 & 0 & \texttt{---------}\\
7 & 0 & 1 & 1 & 2 & 0 & 512 & 0 & \texttt{---------}\\
8 & 0 & 1 & 1 & 2 & 0 & 512 & 0 & \texttt{---------}\\
\multicolumn{9}{@{}l@{}}{\rule{0pt}{2.6ex}\textbf{T9.276}\enspace \texttt{012221111}\enspace $n{=}9$, $|H|{=}7$, len 3, br 2, $|\mathrm{Aut}|{=}144$, orb 4/2, hist-orb 2/0; M 5632, res 4896, rev 320: ok}\\*
0 & -- & 0 & 7 & 1 & 18 & 2 & 492 & \texttt{-????????}\\*
1 & 0 & 1 & 3 & 1 & 64 & 64 & 384 & \texttt{--???----}\\
2 & 1 & 2 & 1 & 3 & 0 & 512 & 0 & \texttt{---------}\\
3 & 1 & 2 & 1 & 3 & 0 & 512 & 0 & \texttt{---------}\\
4 & 1 & 2 & 1 & 3 & 0 & 512 & 0 & \texttt{---------}\\
5 & 0 & 1 & 1 & 4 & 0 & 512 & 0 & \texttt{---------}\\
6 & 0 & 1 & 1 & 4 & 0 & 512 & 0 & \texttt{---------}\\
7 & 0 & 1 & 1 & 4 & 0 & 512 & 0 & \texttt{---------}\\
8 & 0 & 1 & 1 & 4 & 0 & 512 & 0 & \texttt{---------}\\
\multicolumn{9}{@{}l@{}}{\rule{0pt}{2.6ex}\textbf{T9.277}\enspace \texttt{012212212}\enspace $n{=}9$, $|H|{=}5$, len 3, br 3, $|\mathrm{Aut}|{=}8$, orb 5/3, hist-orb 2/1; M 6656, res 3648, rev 1628: ok}\\*
0 & -- & 0 & 5 & 1 & 150 & 2 & 360 & \texttt{-????????}\\*
1 & 0 & 1 & 2 & 2 & 128 & 128 & 256 & \texttt{--??-----}\\
2 & 1 & 2 & 1 & 4 & 0 & 512 & 0 & \texttt{---------}\\
3 & 1 & 2 & 1 & 4 & 0 & 512 & 0 & \texttt{---------}\\
4 & 0 & 1 & 2 & 2 & 128 & 128 & 256 & \texttt{-----??--}\\
5 & 4 & 2 & 1 & 4 & 0 & 512 & 0 & \texttt{---------}\\
6 & 4 & 2 & 1 & 4 & 0 & 512 & 0 & \texttt{---------}\\
7 & 0 & 1 & 1 & 1 & 256 & 256 & 0 & \texttt{--------+}\\
8 & 7 & 2 & 1 & 1 & 0 & 512 & 0 & \texttt{---------}\\
\multicolumn{9}{@{}l@{}}{\rule{0pt}{2.6ex}\textbf{T9.278}\enspace \texttt{012212211}\enspace $n{=}9$, $|H|{=}6$, len 3, br 3, $|\mathrm{Aut}|{=}16$, orb 4/1, hist-orb 2/0; M 6144, res 4448, rev 832: ok}\\*
0 & -- & 0 & 6 & 1 & 50 & 2 & 460 & \texttt{-????????}\\*
1 & 0 & 1 & 2 & 2 & 128 & 128 & 256 & \texttt{--??-----}\\
2 & 1 & 2 & 1 & 4 & 0 & 512 & 0 & \texttt{---------}\\
3 & 1 & 2 & 1 & 4 & 0 & 512 & 0 & \texttt{---------}\\
4 & 0 & 1 & 2 & 2 & 128 & 128 & 256 & \texttt{-----??--}\\
5 & 4 & 2 & 1 & 4 & 0 & 512 & 0 & \texttt{---------}\\
6 & 4 & 2 & 1 & 4 & 0 & 512 & 0 & \texttt{---------}\\
7 & 0 & 1 & 1 & 2 & 0 & 512 & 0 & \texttt{---------}\\
8 & 0 & 1 & 1 & 2 & 0 & 512 & 0 & \texttt{---------}\\
\multicolumn{9}{@{}l@{}}{\rule{0pt}{2.6ex}\textbf{T9.279}\enspace \texttt{012212121}\enspace $n{=}9$, $|H|{=}5$, len 3, br 2, $|\mathrm{Aut}|{=}4$, orb 6/3, hist-orb 3/1; M 6144, res 3744, rev 1296: ok}\\*
0 & -- & 0 & 5 & 1 & 90 & 2 & 420 & \texttt{-????????}\\*
1 & 0 & 1 & 2 & 1 & 128 & 128 & 256 & \texttt{--??-----}\\
2 & 1 & 2 & 1 & 2 & 0 & 512 & 0 & \texttt{---------}\\
3 & 1 & 2 & 1 & 2 & 0 & 512 & 0 & \texttt{---------}\\
4 & 0 & 1 & 1 & 2 & 256 & 256 & 0 & \texttt{-----+---}\\
5 & 4 & 2 & 1 & 2 & 0 & 512 & 0 & \texttt{---------}\\
6 & 0 & 1 & 1 & 2 & 256 & 256 & 0 & \texttt{-------+-}\\
7 & 6 & 2 & 1 & 2 & 0 & 512 & 0 & \texttt{---------}\\
8 & 0 & 1 & 1 & 1 & 0 & 512 & 0 & \texttt{---------}\\
\multicolumn{9}{@{}l@{}}{\rule{0pt}{2.6ex}\textbf{T9.280}\enspace \texttt{012212111}\enspace $n{=}9$, $|H|{=}6$, len 3, br 2, $|\mathrm{Aut}|{=}12$, orb 6/4, hist-orb 3/1; M 5632, res 4224, rev 708: ok}\\*
0 & -- & 0 & 6 & 1 & 30 & 2 & 480 & \texttt{-????????}\\*
1 & 0 & 1 & 2 & 1 & 128 & 128 & 256 & \texttt{--??-----}\\
2 & 1 & 2 & 1 & 2 & 0 & 512 & 0 & \texttt{---------}\\
3 & 1 & 2 & 1 & 2 & 0 & 512 & 0 & \texttt{---------}\\
4 & 0 & 1 & 1 & 1 & 256 & 256 & 0 & \texttt{-----+---}\\
5 & 4 & 2 & 1 & 1 & 0 & 512 & 0 & \texttt{---------}\\
6 & 0 & 1 & 1 & 3 & 0 & 512 & 0 & \texttt{---------}\\
7 & 0 & 1 & 1 & 3 & 0 & 512 & 0 & \texttt{---------}\\
8 & 0 & 1 & 1 & 3 & 0 & 512 & 0 & \texttt{---------}\\
\multicolumn{9}{@{}l@{}}{\rule{0pt}{2.6ex}\textbf{T9.281}\enspace \texttt{012211111}\enspace $n{=}9$, $|H|{=}7$, len 3, br 2, $|\mathrm{Aut}|{=}240$, orb 4/2, hist-orb 2/0; M 5120, res 4384, rev 328: ok}\\*
0 & -- & 0 & 7 & 1 & 10 & 2 & 500 & \texttt{-????????}\\*
1 & 0 & 1 & 2 & 1 & 128 & 128 & 256 & \texttt{--??-----}\\
2 & 1 & 2 & 1 & 2 & 0 & 512 & 0 & \texttt{---------}\\
3 & 1 & 2 & 1 & 2 & 0 & 512 & 0 & \texttt{---------}\\
4 & 0 & 1 & 1 & 5 & 0 & 512 & 0 & \texttt{---------}\\
5 & 0 & 1 & 1 & 5 & 0 & 512 & 0 & \texttt{---------}\\
6 & 0 & 1 & 1 & 5 & 0 & 512 & 0 & \texttt{---------}\\
7 & 0 & 1 & 1 & 5 & 0 & 512 & 0 & \texttt{---------}\\
8 & 0 & 1 & 1 & 5 & 0 & 512 & 0 & \texttt{---------}\\
\multicolumn{9}{@{}l@{}}{\rule{0pt}{2.6ex}\textbf{T9.282}\enspace \texttt{012121212}\enspace $n{=}9$, $|H|{=}4$, len 3, br 1, $|\mathrm{Aut}|{=}24$, orb 3/1, hist-orb 1/0; M 6144, res 2784, rev 1888: ok}\\*
0 & -- & 0 & 4 & 1 & 162 & 2 & 348 & \texttt{-????????}\\*
1 & 0 & 1 & 1 & 4 & 256 & 256 & 0 & \texttt{--+------}\\
2 & 1 & 2 & 1 & 4 & 0 & 512 & 0 & \texttt{---------}\\
3 & 0 & 1 & 1 & 4 & 256 & 256 & 0 & \texttt{----+----}\\
4 & 3 & 2 & 1 & 4 & 0 & 512 & 0 & \texttt{---------}\\
5 & 0 & 1 & 1 & 4 & 256 & 256 & 0 & \texttt{------+--}\\
6 & 5 & 2 & 1 & 4 & 0 & 512 & 0 & \texttt{---------}\\
7 & 0 & 1 & 1 & 4 & 256 & 256 & 0 & \texttt{--------+}\\
8 & 7 & 2 & 1 & 4 & 0 & 512 & 0 & \texttt{---------}\\
\multicolumn{9}{@{}l@{}}{\rule{0pt}{2.6ex}\textbf{T9.283}\enspace \texttt{012121211}\enspace $n{=}9$, $|H|{=}5$, len 3, br 1, $|\mathrm{Aut}|{=}12$, orb 4/1, hist-orb 2/0; M 5632, res 3648, rev 1092: ok}\\*
0 & -- & 0 & 5 & 1 & 54 & 2 & 456 & \texttt{-????????}\\*
1 & 0 & 1 & 1 & 3 & 256 & 256 & 0 & \texttt{--+------}\\
2 & 1 & 2 & 1 & 3 & 0 & 512 & 0 & \texttt{---------}\\
3 & 0 & 1 & 1 & 3 & 256 & 256 & 0 & \texttt{----+----}\\
4 & 3 & 2 & 1 & 3 & 0 & 512 & 0 & \texttt{---------}\\
5 & 0 & 1 & 1 & 3 & 256 & 256 & 0 & \texttt{------+--}\\
6 & 5 & 2 & 1 & 3 & 0 & 512 & 0 & \texttt{---------}\\
7 & 0 & 1 & 1 & 2 & 0 & 512 & 0 & \texttt{---------}\\
8 & 0 & 1 & 1 & 2 & 0 & 512 & 0 & \texttt{---------}\\
\multicolumn{9}{@{}l@{}}{\rule{0pt}{2.6ex}\textbf{T9.284}\enspace \texttt{012121111}\enspace $n{=}9$, $|H|{=}6$, len 3, br 1, $|\mathrm{Aut}|{=}48$, orb 4/1, hist-orb 2/0; M 5120, res 3936, rev 632: ok}\\*
0 & -- & 0 & 6 & 1 & 18 & 2 & 492 & \texttt{-????????}\\*
1 & 0 & 1 & 1 & 2 & 256 & 256 & 0 & \texttt{--+------}\\
2 & 1 & 2 & 1 & 2 & 0 & 512 & 0 & \texttt{---------}\\
3 & 0 & 1 & 1 & 2 & 256 & 256 & 0 & \texttt{----+----}\\
4 & 3 & 2 & 1 & 2 & 0 & 512 & 0 & \texttt{---------}\\
5 & 0 & 1 & 1 & 4 & 0 & 512 & 0 & \texttt{---------}\\
6 & 0 & 1 & 1 & 4 & 0 & 512 & 0 & \texttt{---------}\\
7 & 0 & 1 & 1 & 4 & 0 & 512 & 0 & \texttt{---------}\\
8 & 0 & 1 & 1 & 4 & 0 & 512 & 0 & \texttt{---------}\\
\multicolumn{9}{@{}l@{}}{\rule{0pt}{2.6ex}\textbf{T9.285}\enspace \texttt{012111111}\enspace $n{=}9$, $|H|{=}7$, len 3, br 1, $|\mathrm{Aut}|{=}720$, orb 4/3, hist-orb 2/1; M 4608, res 4032, rev 300: ok}\\*
0 & -- & 0 & 7 & 1 & 6 & 2 & 504 & \texttt{-????????}\\*
1 & 0 & 1 & 1 & 1 & 256 & 256 & 0 & \texttt{--+------}\\
2 & 1 & 2 & 1 & 1 & 0 & 512 & 0 & \texttt{---------}\\
3 & 0 & 1 & 1 & 6 & 0 & 512 & 0 & \texttt{---------}\\
4 & 0 & 1 & 1 & 6 & 0 & 512 & 0 & \texttt{---------}\\
5 & 0 & 1 & 1 & 6 & 0 & 512 & 0 & \texttt{---------}\\
6 & 0 & 1 & 1 & 6 & 0 & 512 & 0 & \texttt{---------}\\
7 & 0 & 1 & 1 & 6 & 0 & 512 & 0 & \texttt{---------}\\
8 & 0 & 1 & 1 & 6 & 0 & 512 & 0 & \texttt{---------}\\
\multicolumn{9}{@{}l@{}}{\rule{0pt}{2.6ex}\textbf{T9.286}\enspace \texttt{011111111}\enspace $n{=}9$, $|H|{=}8$, len 2, br 1, $|\mathrm{Aut}|{=}40320$, orb 2/1, hist-orb 1/0; M 4096, res 4064, rev 16: ok}\\*
0 & -- & 0 & 8 & 1 & 2 & 2 & 508 & \texttt{-????????}\\*
1 & 0 & 1 & 1 & 8 & 0 & 512 & 0 & \texttt{---------}\\
2 & 0 & 1 & 1 & 8 & 0 & 512 & 0 & \texttt{---------}\\
3 & 0 & 1 & 1 & 8 & 0 & 512 & 0 & \texttt{---------}\\
4 & 0 & 1 & 1 & 8 & 0 & 512 & 0 & \texttt{---------}\\
5 & 0 & 1 & 1 & 8 & 0 & 512 & 0 & \texttt{---------}\\
6 & 0 & 1 & 1 & 8 & 0 & 512 & 0 & \texttt{---------}\\
7 & 0 & 1 & 1 & 8 & 0 & 512 & 0 & \texttt{---------}\\
8 & 0 & 1 & 1 & 8 & 0 & 512 & 0 & \texttt{---------}\\
\multicolumn{9}{@{}l@{}}{\rule{0pt}{3.4ex}\normalsize\bfseries Frames with 10 moments (719 classes)}\\*
\multicolumn{9}{@{}l@{}}{\rule{0pt}{2.6ex}\textbf{T10.001}\enspace \texttt{0123456789}\enspace $n{=}10$, $|H|{=}1$, len 10, br 0, $|\mathrm{Aut}|{=}1$, orb 10/10, hist-orb 1/1; M 46080, res 0, rev 14362: ok}\\*
0 & -- & 0 & 1 & 1 & 1022 & 2 & 0 & \texttt{-+++++++++}\\*
1 & 0 & 1 & 1 & 1 & 1020 & 4 & 0 & \texttt{--++++++++}\\
2 & 1 & 2 & 1 & 1 & 1016 & 8 & 0 & \texttt{---+++++++}\\
3 & 2 & 3 & 1 & 1 & 1008 & 16 & 0 & \texttt{----++++++}\\
4 & 3 & 4 & 1 & 1 & 992 & 32 & 0 & \texttt{-----+++++}\\
5 & 4 & 5 & 1 & 1 & 960 & 64 & 0 & \texttt{------++++}\\
6 & 5 & 6 & 1 & 1 & 896 & 128 & 0 & \texttt{-------+++}\\
7 & 6 & 7 & 1 & 1 & 768 & 256 & 0 & \texttt{--------++}\\
8 & 7 & 8 & 1 & 1 & 512 & 512 & 0 & \texttt{---------+}\\
9 & 8 & 9 & 1 & 1 & 0 & 1024 & 0 & \texttt{----------}\\
\multicolumn{9}{@{}l@{}}{\rule{0pt}{2.6ex}\textbf{T10.002}\enspace \texttt{0123456788}\enspace $n{=}10$, $|H|{=}2$, len 9, br 1, $|\mathrm{Aut}|{=}2$, orb 9/8, hist-orb 1/0; M 45056, res 2040, rev 15906: ok}\\*
0 & -- & 0 & 2 & 1 & 1018 & 2 & 4 & \texttt{-+++++++??}\\*
1 & 0 & 1 & 2 & 1 & 1012 & 4 & 8 & \texttt{--++++++??}\\
2 & 1 & 2 & 2 & 1 & 1000 & 8 & 16 & \texttt{---+++++??}\\
3 & 2 & 3 & 2 & 1 & 976 & 16 & 32 & \texttt{----++++??}\\
4 & 3 & 4 & 2 & 1 & 928 & 32 & 64 & \texttt{-----+++??}\\
5 & 4 & 5 & 2 & 1 & 832 & 64 & 128 & \texttt{------++??}\\
6 & 5 & 6 & 2 & 1 & 640 & 128 & 256 & \texttt{-------+??}\\
7 & 6 & 7 & 2 & 1 & 256 & 256 & 512 & \texttt{--------??}\\
8 & 7 & 8 & 1 & 2 & 0 & 1024 & 0 & \texttt{----------}\\
9 & 7 & 8 & 1 & 2 & 0 & 1024 & 0 & \texttt{----------}\\
\multicolumn{9}{@{}l@{}}{\rule{0pt}{2.6ex}\textbf{T10.003}\enspace \texttt{0123456787}\enspace $n{=}10$, $|H|{=}2$, len 9, br 1, $|\mathrm{Aut}|{=}1$, orb 10/10, hist-orb 2/2; M 44032, res 3048, rev 16174: ok}\\*
0 & -- & 0 & 2 & 1 & 1014 & 2 & 8 & \texttt{-++++++???}\\*
1 & 0 & 1 & 2 & 1 & 1004 & 4 & 16 & \texttt{--+++++???}\\
2 & 1 & 2 & 2 & 1 & 984 & 8 & 32 & \texttt{---++++???}\\
3 & 2 & 3 & 2 & 1 & 944 & 16 & 64 & \texttt{----+++???}\\
4 & 3 & 4 & 2 & 1 & 864 & 32 & 128 & \texttt{-----++???}\\
5 & 4 & 5 & 2 & 1 & 704 & 64 & 256 & \texttt{------+???}\\
6 & 5 & 6 & 2 & 1 & 384 & 128 & 512 & \texttt{-------???}\\
7 & 6 & 7 & 1 & 1 & 512 & 512 & 0 & \texttt{--------+-}\\
8 & 7 & 8 & 1 & 1 & 0 & 1024 & 0 & \texttt{----------}\\
9 & 6 & 7 & 1 & 1 & 0 & 1024 & 0 & \texttt{----------}\\
\multicolumn{9}{@{}l@{}}{\rule{0pt}{2.6ex}\textbf{T10.004}\enspace \texttt{0123456786}\enspace $n{=}10$, $|H|{=}2$, len 9, br 1, $|\mathrm{Aut}|{=}1$, orb 10/10, hist-orb 2/2; M 43008, res 4032, rev 15558: ok}\\*
0 & -- & 0 & 2 & 1 & 1006 & 2 & 16 & \texttt{-+++++????}\\*
1 & 0 & 1 & 2 & 1 & 988 & 4 & 32 & \texttt{--++++????}\\
2 & 1 & 2 & 2 & 1 & 952 & 8 & 64 & \texttt{---+++????}\\
3 & 2 & 3 & 2 & 1 & 880 & 16 & 128 & \texttt{----++????}\\
4 & 3 & 4 & 2 & 1 & 736 & 32 & 256 & \texttt{-----+????}\\
5 & 4 & 5 & 2 & 1 & 448 & 64 & 512 & \texttt{------????}\\
6 & 5 & 6 & 1 & 1 & 768 & 256 & 0 & \texttt{-------++-}\\
7 & 6 & 7 & 1 & 1 & 512 & 512 & 0 & \texttt{--------+-}\\
8 & 7 & 8 & 1 & 1 & 0 & 1024 & 0 & \texttt{----------}\\
9 & 5 & 6 & 1 & 1 & 0 & 1024 & 0 & \texttt{----------}\\
\multicolumn{9}{@{}l@{}}{\rule{0pt}{2.6ex}\textbf{T10.005}\enspace \texttt{0123456785}\enspace $n{=}10$, $|H|{=}2$, len 9, br 1, $|\mathrm{Aut}|{=}1$, orb 10/10, hist-orb 2/2; M 41984, res 4960, rev 14646: ok}\\*
0 & -- & 0 & 2 & 1 & 990 & 2 & 32 & \texttt{-++++?????}\\*
1 & 0 & 1 & 2 & 1 & 956 & 4 & 64 & \texttt{--+++?????}\\
2 & 1 & 2 & 2 & 1 & 888 & 8 & 128 & \texttt{---++?????}\\
3 & 2 & 3 & 2 & 1 & 752 & 16 & 256 & \texttt{----+?????}\\
4 & 3 & 4 & 2 & 1 & 480 & 32 & 512 & \texttt{-----?????}\\
5 & 4 & 5 & 1 & 1 & 896 & 128 & 0 & \texttt{------+++-}\\
6 & 5 & 6 & 1 & 1 & 768 & 256 & 0 & \texttt{-------++-}\\
7 & 6 & 7 & 1 & 1 & 512 & 512 & 0 & \texttt{--------+-}\\
8 & 7 & 8 & 1 & 1 & 0 & 1024 & 0 & \texttt{----------}\\
9 & 4 & 5 & 1 & 1 & 0 & 1024 & 0 & \texttt{----------}\\
\multicolumn{9}{@{}l@{}}{\rule{0pt}{2.6ex}\textbf{T10.006}\enspace \texttt{0123456784}\enspace $n{=}10$, $|H|{=}2$, len 9, br 1, $|\mathrm{Aut}|{=}1$, orb 10/10, hist-orb 2/2; M 40960, res 5760, rev 13686: ok}\\*
0 & -- & 0 & 2 & 1 & 958 & 2 & 64 & \texttt{-+++??????}\\*
1 & 0 & 1 & 2 & 1 & 892 & 4 & 128 & \texttt{--++??????}\\
2 & 1 & 2 & 2 & 1 & 760 & 8 & 256 & \texttt{---+??????}\\
3 & 2 & 3 & 2 & 1 & 496 & 16 & 512 & \texttt{----??????}\\
4 & 3 & 4 & 1 & 1 & 960 & 64 & 0 & \texttt{-----++++-}\\
5 & 4 & 5 & 1 & 1 & 896 & 128 & 0 & \texttt{------+++-}\\
6 & 5 & 6 & 1 & 1 & 768 & 256 & 0 & \texttt{-------++-}\\
7 & 6 & 7 & 1 & 1 & 512 & 512 & 0 & \texttt{--------+-}\\
8 & 7 & 8 & 1 & 1 & 0 & 1024 & 0 & \texttt{----------}\\
9 & 3 & 4 & 1 & 1 & 0 & 1024 & 0 & \texttt{----------}\\
\multicolumn{9}{@{}l@{}}{\rule{0pt}{2.6ex}\textbf{T10.007}\enspace \texttt{0123456783}\enspace $n{=}10$, $|H|{=}2$, len 9, br 1, $|\mathrm{Aut}|{=}1$, orb 10/10, hist-orb 2/2; M 39936, res 6272, rev 12806: ok}\\*
0 & -- & 0 & 2 & 1 & 894 & 2 & 128 & \texttt{-++???????}\\*
1 & 0 & 1 & 2 & 1 & 764 & 4 & 256 & \texttt{--+???????}\\
2 & 1 & 2 & 2 & 1 & 504 & 8 & 512 & \texttt{---???????}\\
3 & 2 & 3 & 1 & 1 & 992 & 32 & 0 & \texttt{----+++++-}\\
4 & 3 & 4 & 1 & 1 & 960 & 64 & 0 & \texttt{-----++++-}\\
5 & 4 & 5 & 1 & 1 & 896 & 128 & 0 & \texttt{------+++-}\\
6 & 5 & 6 & 1 & 1 & 768 & 256 & 0 & \texttt{-------++-}\\
7 & 6 & 7 & 1 & 1 & 512 & 512 & 0 & \texttt{--------+-}\\
8 & 7 & 8 & 1 & 1 & 0 & 1024 & 0 & \texttt{----------}\\
9 & 2 & 3 & 1 & 1 & 0 & 1024 & 0 & \texttt{----------}\\
\multicolumn{9}{@{}l@{}}{\rule{0pt}{2.6ex}\textbf{T10.008}\enspace \texttt{0123456782}\enspace $n{=}10$, $|H|{=}2$, len 9, br 1, $|\mathrm{Aut}|{=}1$, orb 10/10, hist-orb 2/2; M 38912, res 6144, rev 12126: ok}\\*
0 & -- & 0 & 2 & 1 & 766 & 2 & 256 & \texttt{-+????????}\\*
1 & 0 & 1 & 2 & 1 & 508 & 4 & 512 & \texttt{--????????}\\
2 & 1 & 2 & 1 & 1 & 1008 & 16 & 0 & \texttt{---++++++-}\\
3 & 2 & 3 & 1 & 1 & 992 & 32 & 0 & \texttt{----+++++-}\\
4 & 3 & 4 & 1 & 1 & 960 & 64 & 0 & \texttt{-----++++-}\\
5 & 4 & 5 & 1 & 1 & 896 & 128 & 0 & \texttt{------+++-}\\
6 & 5 & 6 & 1 & 1 & 768 & 256 & 0 & \texttt{-------++-}\\
7 & 6 & 7 & 1 & 1 & 512 & 512 & 0 & \texttt{--------+-}\\
8 & 7 & 8 & 1 & 1 & 0 & 1024 & 0 & \texttt{----------}\\
9 & 1 & 2 & 1 & 1 & 0 & 1024 & 0 & \texttt{----------}\\
\multicolumn{9}{@{}l@{}}{\rule{0pt}{2.6ex}\textbf{T10.009}\enspace \texttt{0123456781}\enspace $n{=}10$, $|H|{=}2$, len 9, br 1, $|\mathrm{Aut}|{=}1$, orb 10/10, hist-orb 2/2; M 37888, res 4608, rev 11842: ok}\\*
0 & -- & 0 & 2 & 1 & 510 & 2 & 512 & \texttt{-?????????}\\*
1 & 0 & 1 & 1 & 1 & 1016 & 8 & 0 & \texttt{--+++++++-}\\
2 & 1 & 2 & 1 & 1 & 1008 & 16 & 0 & \texttt{---++++++-}\\
3 & 2 & 3 & 1 & 1 & 992 & 32 & 0 & \texttt{----+++++-}\\
4 & 3 & 4 & 1 & 1 & 960 & 64 & 0 & \texttt{-----++++-}\\
5 & 4 & 5 & 1 & 1 & 896 & 128 & 0 & \texttt{------+++-}\\
6 & 5 & 6 & 1 & 1 & 768 & 256 & 0 & \texttt{-------++-}\\
7 & 6 & 7 & 1 & 1 & 512 & 512 & 0 & \texttt{--------+-}\\
8 & 7 & 8 & 1 & 1 & 0 & 1024 & 0 & \texttt{----------}\\
9 & 0 & 1 & 1 & 1 & 0 & 1024 & 0 & \texttt{----------}\\
\multicolumn{9}{@{}l@{}}{\rule{0pt}{2.6ex}\textbf{T10.010}\enspace \texttt{0123456777}\enspace $n{=}10$, $|H|{=}3$, len 8, br 1, $|\mathrm{Aut}|{=}6$, orb 8/7, hist-orb 1/0; M 43008, res 4572, rev 17214: ok}\\*
0 & -- & 0 & 3 & 1 & 1010 & 2 & 12 & \texttt{-++++++???}\\*
1 & 0 & 1 & 3 & 1 & 996 & 4 & 24 & \texttt{--+++++???}\\
2 & 1 & 2 & 3 & 1 & 968 & 8 & 48 & \texttt{---++++???}\\
3 & 2 & 3 & 3 & 1 & 912 & 16 & 96 & \texttt{----+++???}\\
4 & 3 & 4 & 3 & 1 & 800 & 32 & 192 & \texttt{-----++???}\\
5 & 4 & 5 & 3 & 1 & 576 & 64 & 384 & \texttt{------+???}\\
6 & 5 & 6 & 3 & 1 & 128 & 128 & 768 & \texttt{-------???}\\
7 & 6 & 7 & 1 & 3 & 0 & 1024 & 0 & \texttt{----------}\\
8 & 6 & 7 & 1 & 3 & 0 & 1024 & 0 & \texttt{----------}\\
9 & 6 & 7 & 1 & 3 & 0 & 1024 & 0 & \texttt{----------}\\
\multicolumn{9}{@{}l@{}}{\rule{0pt}{2.6ex}\textbf{T10.011}\enspace \texttt{0123456776}\enspace $n{=}10$, $|H|{=}3$, len 8, br 2, $|\mathrm{Aut}|{=}2$, orb 9/8, hist-orb 2/1; M 41984, res 5560, rev 16346: ok}\\*
0 & -- & 0 & 3 & 1 & 1002 & 2 & 20 & \texttt{-+++++????}\\*
1 & 0 & 1 & 3 & 1 & 980 & 4 & 40 & \texttt{--++++????}\\
2 & 1 & 2 & 3 & 1 & 936 & 8 & 80 & \texttt{---+++????}\\
3 & 2 & 3 & 3 & 1 & 848 & 16 & 160 & \texttt{----++????}\\
4 & 3 & 4 & 3 & 1 & 672 & 32 & 320 & \texttt{-----+????}\\
5 & 4 & 5 & 3 & 1 & 320 & 64 & 640 & \texttt{------????}\\
6 & 5 & 6 & 2 & 1 & 256 & 256 & 512 & \texttt{-------??-}\\
7 & 6 & 7 & 1 & 2 & 0 & 1024 & 0 & \texttt{----------}\\
8 & 6 & 7 & 1 & 2 & 0 & 1024 & 0 & \texttt{----------}\\
9 & 5 & 6 & 1 & 1 & 0 & 1024 & 0 & \texttt{----------}\\
\multicolumn{9}{@{}l@{}}{\rule{0pt}{2.6ex}\textbf{T10.012}\enspace \texttt{0123456775}\enspace $n{=}10$, $|H|{=}3$, len 8, br 2, $|\mathrm{Aut}|{=}2$, orb 9/8, hist-orb 2/1; M 40960, res 6620, rev 15314: ok}\\*
0 & -- & 0 & 3 & 1 & 986 & 2 & 36 & \texttt{-++++?????}\\*
1 & 0 & 1 & 3 & 1 & 948 & 4 & 72 & \texttt{--+++?????}\\
2 & 1 & 2 & 3 & 1 & 872 & 8 & 144 & \texttt{---++?????}\\
3 & 2 & 3 & 3 & 1 & 720 & 16 & 288 & \texttt{----+?????}\\
4 & 3 & 4 & 3 & 1 & 416 & 32 & 576 & \texttt{-----?????}\\
5 & 4 & 5 & 2 & 1 & 640 & 128 & 256 & \texttt{------+??-}\\
6 & 5 & 6 & 2 & 1 & 256 & 256 & 512 & \texttt{-------??-}\\
7 & 6 & 7 & 1 & 2 & 0 & 1024 & 0 & \texttt{----------}\\
8 & 6 & 7 & 1 & 2 & 0 & 1024 & 0 & \texttt{----------}\\
9 & 4 & 5 & 1 & 1 & 0 & 1024 & 0 & \texttt{----------}\\
\multicolumn{9}{@{}l@{}}{\rule{0pt}{2.6ex}\textbf{T10.013}\enspace \texttt{0123456774}\enspace $n{=}10$, $|H|{=}3$, len 8, br 2, $|\mathrm{Aut}|{=}2$, orb 9/8, hist-orb 2/1; M 39936, res 7552, rev 14306: ok}\\*
0 & -- & 0 & 3 & 1 & 954 & 2 & 68 & \texttt{-+++??????}\\*
1 & 0 & 1 & 3 & 1 & 884 & 4 & 136 & \texttt{--++??????}\\
2 & 1 & 2 & 3 & 1 & 744 & 8 & 272 & \texttt{---+??????}\\
3 & 2 & 3 & 3 & 1 & 464 & 16 & 544 & \texttt{----??????}\\
4 & 3 & 4 & 2 & 1 & 832 & 64 & 128 & \texttt{-----++??-}\\
5 & 4 & 5 & 2 & 1 & 640 & 128 & 256 & \texttt{------+??-}\\
6 & 5 & 6 & 2 & 1 & 256 & 256 & 512 & \texttt{-------??-}\\
7 & 6 & 7 & 1 & 2 & 0 & 1024 & 0 & \texttt{----------}\\
8 & 6 & 7 & 1 & 2 & 0 & 1024 & 0 & \texttt{----------}\\
9 & 3 & 4 & 1 & 1 & 0 & 1024 & 0 & \texttt{----------}\\
\multicolumn{9}{@{}l@{}}{\rule{0pt}{2.6ex}\textbf{T10.014}\enspace \texttt{0123456773}\enspace $n{=}10$, $|H|{=}3$, len 8, br 2, $|\mathrm{Aut}|{=}2$, orb 9/8, hist-orb 2/1; M 38912, res 8164, rev 13426: ok}\\*
0 & -- & 0 & 3 & 1 & 890 & 2 & 132 & \texttt{-++???????}\\*
1 & 0 & 1 & 3 & 1 & 756 & 4 & 264 & \texttt{--+???????}\\
2 & 1 & 2 & 3 & 1 & 488 & 8 & 528 & \texttt{---???????}\\
3 & 2 & 3 & 2 & 1 & 928 & 32 & 64 & \texttt{----+++??-}\\
4 & 3 & 4 & 2 & 1 & 832 & 64 & 128 & \texttt{-----++??-}\\
5 & 4 & 5 & 2 & 1 & 640 & 128 & 256 & \texttt{------+??-}\\
6 & 5 & 6 & 2 & 1 & 256 & 256 & 512 & \texttt{-------??-}\\
7 & 6 & 7 & 1 & 2 & 0 & 1024 & 0 & \texttt{----------}\\
8 & 6 & 7 & 1 & 2 & 0 & 1024 & 0 & \texttt{----------}\\
9 & 2 & 3 & 1 & 1 & 0 & 1024 & 0 & \texttt{----------}\\
\multicolumn{9}{@{}l@{}}{\rule{0pt}{2.6ex}\textbf{T10.015}\enspace \texttt{0123456772}\enspace $n{=}10$, $|H|{=}3$, len 8, br 2, $|\mathrm{Aut}|{=}2$, orb 9/8, hist-orb 2/1; M 37888, res 8104, rev 12794: ok}\\*
0 & -- & 0 & 3 & 1 & 762 & 2 & 260 & \texttt{-+????????}\\*
1 & 0 & 1 & 3 & 1 & 500 & 4 & 520 & \texttt{--????????}\\
2 & 1 & 2 & 2 & 1 & 976 & 16 & 32 & \texttt{---++++??-}\\
3 & 2 & 3 & 2 & 1 & 928 & 32 & 64 & \texttt{----+++??-}\\
4 & 3 & 4 & 2 & 1 & 832 & 64 & 128 & \texttt{-----++??-}\\
5 & 4 & 5 & 2 & 1 & 640 & 128 & 256 & \texttt{------+??-}\\
6 & 5 & 6 & 2 & 1 & 256 & 256 & 512 & \texttt{-------??-}\\
7 & 6 & 7 & 1 & 2 & 0 & 1024 & 0 & \texttt{----------}\\
8 & 6 & 7 & 1 & 2 & 0 & 1024 & 0 & \texttt{----------}\\
9 & 1 & 2 & 1 & 1 & 0 & 1024 & 0 & \texttt{----------}\\
\multicolumn{9}{@{}l@{}}{\rule{0pt}{2.6ex}\textbf{T10.016}\enspace \texttt{0123456771}\enspace $n{=}10$, $|H|{=}3$, len 8, br 2, $|\mathrm{Aut}|{=}2$, orb 9/8, hist-orb 2/1; M 36864, res 6612, rev 12630: ok}\\*
0 & -- & 0 & 3 & 1 & 506 & 2 & 516 & \texttt{-?????????}\\*
1 & 0 & 1 & 2 & 1 & 1000 & 8 & 16 & \texttt{--+++++??-}\\
2 & 1 & 2 & 2 & 1 & 976 & 16 & 32 & \texttt{---++++??-}\\
3 & 2 & 3 & 2 & 1 & 928 & 32 & 64 & \texttt{----+++??-}\\
4 & 3 & 4 & 2 & 1 & 832 & 64 & 128 & \texttt{-----++??-}\\
5 & 4 & 5 & 2 & 1 & 640 & 128 & 256 & \texttt{------+??-}\\
6 & 5 & 6 & 2 & 1 & 256 & 256 & 512 & \texttt{-------??-}\\
7 & 6 & 7 & 1 & 2 & 0 & 1024 & 0 & \texttt{----------}\\
8 & 6 & 7 & 1 & 2 & 0 & 1024 & 0 & \texttt{----------}\\
9 & 0 & 1 & 1 & 1 & 0 & 1024 & 0 & \texttt{----------}\\
\multicolumn{9}{@{}l@{}}{\rule{0pt}{2.6ex}\textbf{T10.017}\enspace \texttt{0123456767}\enspace $n{=}10$, $|H|{=}2$, len 8, br 1, $|\mathrm{Aut}|{=}2$, orb 8/6, hist-orb 1/0; M 41984, res 3024, rev 16834: ok}\\*
0 & -- & 0 & 2 & 1 & 1010 & 2 & 12 & \texttt{-+++++????}\\*
1 & 0 & 1 & 2 & 1 & 996 & 4 & 24 & \texttt{--++++????}\\
2 & 1 & 2 & 2 & 1 & 968 & 8 & 48 & \texttt{---+++????}\\
3 & 2 & 3 & 2 & 1 & 912 & 16 & 96 & \texttt{----++????}\\
4 & 3 & 4 & 2 & 1 & 800 & 32 & 192 & \texttt{-----+????}\\
5 & 4 & 5 & 2 & 1 & 576 & 64 & 384 & \texttt{------????}\\
6 & 5 & 6 & 1 & 2 & 512 & 512 & 0 & \texttt{-------+--}\\
7 & 6 & 7 & 1 & 2 & 0 & 1024 & 0 & \texttt{----------}\\
8 & 5 & 6 & 1 & 2 & 512 & 512 & 0 & \texttt{---------+}\\
9 & 8 & 7 & 1 & 2 & 0 & 1024 & 0 & \texttt{----------}\\
\multicolumn{9}{@{}l@{}}{\rule{0pt}{2.6ex}\textbf{T10.018}\enspace \texttt{0123456766}\enspace $n{=}10$, $|H|{=}3$, len 8, br 1, $|\mathrm{Aut}|{=}2$, orb 9/8, hist-orb 2/1; M 40960, res 6048, rev 16622: ok}\\*
0 & -- & 0 & 3 & 1 & 998 & 2 & 24 & \texttt{-+++++????}\\*
1 & 0 & 1 & 3 & 1 & 972 & 4 & 48 & \texttt{--++++????}\\
2 & 1 & 2 & 3 & 1 & 920 & 8 & 96 & \texttt{---+++????}\\
3 & 2 & 3 & 3 & 1 & 816 & 16 & 192 & \texttt{----++????}\\
4 & 3 & 4 & 3 & 1 & 608 & 32 & 384 & \texttt{-----+????}\\
5 & 4 & 5 & 3 & 1 & 192 & 64 & 768 & \texttt{------????}\\
6 & 5 & 6 & 1 & 1 & 512 & 512 & 0 & \texttt{-------+--}\\
7 & 6 & 7 & 1 & 1 & 0 & 1024 & 0 & \texttt{----------}\\
8 & 5 & 6 & 1 & 2 & 0 & 1024 & 0 & \texttt{----------}\\
9 & 5 & 6 & 1 & 2 & 0 & 1024 & 0 & \texttt{----------}\\
\multicolumn{9}{@{}l@{}}{\rule{0pt}{2.6ex}\textbf{T10.019}\enspace \texttt{0123456765}\enspace $n{=}10$, $|H|{=}3$, len 8, br 2, $|\mathrm{Aut}|{=}1$, orb 10/10, hist-orb 3/3; M 39936, res 7240, rev 15338: ok}\\*
0 & -- & 0 & 3 & 1 & 982 & 2 & 40 & \texttt{-++++?????}\\*
1 & 0 & 1 & 3 & 1 & 940 & 4 & 80 & \texttt{--+++?????}\\
2 & 1 & 2 & 3 & 1 & 856 & 8 & 160 & \texttt{---++?????}\\
3 & 2 & 3 & 3 & 1 & 688 & 16 & 320 & \texttt{----+?????}\\
4 & 3 & 4 & 3 & 1 & 352 & 32 & 640 & \texttt{-----?????}\\
5 & 4 & 5 & 2 & 1 & 384 & 128 & 512 & \texttt{------???-}\\
6 & 5 & 6 & 1 & 1 & 512 & 512 & 0 & \texttt{-------+--}\\
7 & 6 & 7 & 1 & 1 & 0 & 1024 & 0 & \texttt{----------}\\
8 & 5 & 6 & 1 & 1 & 0 & 1024 & 0 & \texttt{----------}\\
9 & 4 & 5 & 1 & 1 & 0 & 1024 & 0 & \texttt{----------}\\
\multicolumn{9}{@{}l@{}}{\rule{0pt}{2.6ex}\textbf{T10.020}\enspace \texttt{0123456764}\enspace $n{=}10$, $|H|{=}3$, len 8, br 2, $|\mathrm{Aut}|{=}1$, orb 10/10, hist-orb 3/3; M 38912, res 8304, rev 14210: ok}\\*
0 & -- & 0 & 3 & 1 & 950 & 2 & 72 & \texttt{-+++??????}\\*
1 & 0 & 1 & 3 & 1 & 876 & 4 & 144 & \texttt{--++??????}\\
2 & 1 & 2 & 3 & 1 & 728 & 8 & 288 & \texttt{---+??????}\\
3 & 2 & 3 & 3 & 1 & 432 & 16 & 576 & \texttt{----??????}\\
4 & 3 & 4 & 2 & 1 & 704 & 64 & 256 & \texttt{-----+???-}\\
5 & 4 & 5 & 2 & 1 & 384 & 128 & 512 & \texttt{------???-}\\
6 & 5 & 6 & 1 & 1 & 512 & 512 & 0 & \texttt{-------+--}\\
7 & 6 & 7 & 1 & 1 & 0 & 1024 & 0 & \texttt{----------}\\
8 & 5 & 6 & 1 & 1 & 0 & 1024 & 0 & \texttt{----------}\\
9 & 3 & 4 & 1 & 1 & 0 & 1024 & 0 & \texttt{----------}\\
\multicolumn{9}{@{}l@{}}{\rule{0pt}{2.6ex}\textbf{T10.021}\enspace \texttt{0123456763}\enspace $n{=}10$, $|H|{=}3$, len 8, br 2, $|\mathrm{Aut}|{=}1$, orb 10/10, hist-orb 3/3; M 37888, res 9016, rev 13282: ok}\\*
0 & -- & 0 & 3 & 1 & 886 & 2 & 136 & \texttt{-++???????}\\*
1 & 0 & 1 & 3 & 1 & 748 & 4 & 272 & \texttt{--+???????}\\
2 & 1 & 2 & 3 & 1 & 472 & 8 & 544 & \texttt{---???????}\\
3 & 2 & 3 & 2 & 1 & 864 & 32 & 128 & \texttt{----++???-}\\
4 & 3 & 4 & 2 & 1 & 704 & 64 & 256 & \texttt{-----+???-}\\
5 & 4 & 5 & 2 & 1 & 384 & 128 & 512 & \texttt{------???-}\\
6 & 5 & 6 & 1 & 1 & 512 & 512 & 0 & \texttt{-------+--}\\
7 & 6 & 7 & 1 & 1 & 0 & 1024 & 0 & \texttt{----------}\\
8 & 5 & 6 & 1 & 1 & 0 & 1024 & 0 & \texttt{----------}\\
9 & 2 & 3 & 1 & 1 & 0 & 1024 & 0 & \texttt{----------}\\
\multicolumn{9}{@{}l@{}}{\rule{0pt}{2.6ex}\textbf{T10.022}\enspace \texttt{0123456762}\enspace $n{=}10$, $|H|{=}3$, len 8, br 2, $|\mathrm{Aut}|{=}1$, orb 10/10, hist-orb 3/3; M 36864, res 9024, rev 12650: ok}\\*
0 & -- & 0 & 3 & 1 & 758 & 2 & 264 & \texttt{-+????????}\\*
1 & 0 & 1 & 3 & 1 & 492 & 4 & 528 & \texttt{--????????}\\
2 & 1 & 2 & 2 & 1 & 944 & 16 & 64 & \texttt{---+++???-}\\
3 & 2 & 3 & 2 & 1 & 864 & 32 & 128 & \texttt{----++???-}\\
4 & 3 & 4 & 2 & 1 & 704 & 64 & 256 & \texttt{-----+???-}\\
5 & 4 & 5 & 2 & 1 & 384 & 128 & 512 & \texttt{------???-}\\
6 & 5 & 6 & 1 & 1 & 512 & 512 & 0 & \texttt{-------+--}\\
7 & 6 & 7 & 1 & 1 & 0 & 1024 & 0 & \texttt{----------}\\
8 & 5 & 6 & 1 & 1 & 0 & 1024 & 0 & \texttt{----------}\\
9 & 1 & 2 & 1 & 1 & 0 & 1024 & 0 & \texttt{----------}\\
\multicolumn{9}{@{}l@{}}{\rule{0pt}{2.6ex}\textbf{T10.023}\enspace \texttt{0123456761}\enspace $n{=}10$, $|H|{=}3$, len 8, br 2, $|\mathrm{Aut}|{=}1$, orb 10/10, hist-orb 3/3; M 35840, res 7576, rev 12534: ok}\\*
0 & -- & 0 & 3 & 1 & 502 & 2 & 520 & \texttt{-?????????}\\*
1 & 0 & 1 & 2 & 1 & 984 & 8 & 32 & \texttt{--++++???-}\\
2 & 1 & 2 & 2 & 1 & 944 & 16 & 64 & \texttt{---+++???-}\\
3 & 2 & 3 & 2 & 1 & 864 & 32 & 128 & \texttt{----++???-}\\
4 & 3 & 4 & 2 & 1 & 704 & 64 & 256 & \texttt{-----+???-}\\
5 & 4 & 5 & 2 & 1 & 384 & 128 & 512 & \texttt{------???-}\\
6 & 5 & 6 & 1 & 1 & 512 & 512 & 0 & \texttt{-------+--}\\
7 & 6 & 7 & 1 & 1 & 0 & 1024 & 0 & \texttt{----------}\\
8 & 5 & 6 & 1 & 1 & 0 & 1024 & 0 & \texttt{----------}\\
9 & 0 & 1 & 1 & 1 & 0 & 1024 & 0 & \texttt{----------}\\
\multicolumn{9}{@{}l@{}}{\rule{0pt}{2.6ex}\textbf{T10.024}\enspace \texttt{0123456756}\enspace $n{=}10$, $|H|{=}2$, len 8, br 1, $|\mathrm{Aut}|{=}1$, orb 10/10, hist-orb 2/2; M 39936, res 3100, rev 16166: ok}\\*
0 & -- & 0 & 2 & 1 & 1002 & 2 & 20 & \texttt{-++++?????}\\*
1 & 0 & 1 & 2 & 1 & 980 & 4 & 40 & \texttt{--+++?????}\\
2 & 1 & 2 & 2 & 1 & 936 & 8 & 80 & \texttt{---++?????}\\
3 & 2 & 3 & 2 & 1 & 848 & 16 & 160 & \texttt{----+?????}\\
4 & 3 & 4 & 2 & 1 & 672 & 32 & 320 & \texttt{-----?????}\\
5 & 4 & 5 & 1 & 1 & 768 & 256 & 0 & \texttt{------++--}\\
6 & 5 & 6 & 1 & 1 & 512 & 512 & 0 & \texttt{-------+--}\\
7 & 6 & 7 & 1 & 1 & 0 & 1024 & 0 & \texttt{----------}\\
8 & 4 & 5 & 1 & 1 & 512 & 512 & 0 & \texttt{---------+}\\
9 & 8 & 6 & 1 & 1 & 0 & 1024 & 0 & \texttt{----------}\\
\multicolumn{9}{@{}l@{}}{\rule{0pt}{2.6ex}\textbf{T10.025}\enspace \texttt{0123456755}\enspace $n{=}10$, $|H|{=}3$, len 8, br 1, $|\mathrm{Aut}|{=}2$, orb 9/8, hist-orb 2/1; M 38912, res 7440, rev 14990: ok}\\*
0 & -- & 0 & 3 & 1 & 974 & 2 & 48 & \texttt{-++++?????}\\*
1 & 0 & 1 & 3 & 1 & 924 & 4 & 96 & \texttt{--+++?????}\\
2 & 1 & 2 & 3 & 1 & 824 & 8 & 192 & \texttt{---++?????}\\
3 & 2 & 3 & 3 & 1 & 624 & 16 & 384 & \texttt{----+?????}\\
4 & 3 & 4 & 3 & 1 & 224 & 32 & 768 & \texttt{-----?????}\\
5 & 4 & 5 & 1 & 1 & 768 & 256 & 0 & \texttt{------++--}\\
6 & 5 & 6 & 1 & 1 & 512 & 512 & 0 & \texttt{-------+--}\\
7 & 6 & 7 & 1 & 1 & 0 & 1024 & 0 & \texttt{----------}\\
8 & 4 & 5 & 1 & 2 & 0 & 1024 & 0 & \texttt{----------}\\
9 & 4 & 5 & 1 & 2 & 0 & 1024 & 0 & \texttt{----------}\\
\multicolumn{9}{@{}l@{}}{\rule{0pt}{2.6ex}\textbf{T10.026}\enspace \texttt{0123456754}\enspace $n{=}10$, $|H|{=}3$, len 8, br 2, $|\mathrm{Aut}|{=}1$, orb 10/10, hist-orb 3/3; M 37888, res 8768, rev 13610: ok}\\*
0 & -- & 0 & 3 & 1 & 942 & 2 & 80 & \texttt{-+++??????}\\*
1 & 0 & 1 & 3 & 1 & 860 & 4 & 160 & \texttt{--++??????}\\
2 & 1 & 2 & 3 & 1 & 696 & 8 & 320 & \texttt{---+??????}\\
3 & 2 & 3 & 3 & 1 & 368 & 16 & 640 & \texttt{----??????}\\
4 & 3 & 4 & 2 & 1 & 448 & 64 & 512 & \texttt{-----????-}\\
5 & 4 & 5 & 1 & 1 & 768 & 256 & 0 & \texttt{------++--}\\
6 & 5 & 6 & 1 & 1 & 512 & 512 & 0 & \texttt{-------+--}\\
7 & 6 & 7 & 1 & 1 & 0 & 1024 & 0 & \texttt{----------}\\
8 & 4 & 5 & 1 & 1 & 0 & 1024 & 0 & \texttt{----------}\\
9 & 3 & 4 & 1 & 1 & 0 & 1024 & 0 & \texttt{----------}\\
\multicolumn{9}{@{}l@{}}{\rule{0pt}{2.6ex}\textbf{T10.027}\enspace \texttt{0123456753}\enspace $n{=}10$, $|H|{=}3$, len 8, br 2, $|\mathrm{Aut}|{=}1$, orb 10/10, hist-orb 3/3; M 36864, res 9680, rev 12562: ok}\\*
0 & -- & 0 & 3 & 1 & 878 & 2 & 144 & \texttt{-++???????}\\*
1 & 0 & 1 & 3 & 1 & 732 & 4 & 288 & \texttt{--+???????}\\
2 & 1 & 2 & 3 & 1 & 440 & 8 & 576 & \texttt{---???????}\\
3 & 2 & 3 & 2 & 1 & 736 & 32 & 256 & \texttt{----+????-}\\
4 & 3 & 4 & 2 & 1 & 448 & 64 & 512 & \texttt{-----????-}\\
5 & 4 & 5 & 1 & 1 & 768 & 256 & 0 & \texttt{------++--}\\
6 & 5 & 6 & 1 & 1 & 512 & 512 & 0 & \texttt{-------+--}\\
7 & 6 & 7 & 1 & 1 & 0 & 1024 & 0 & \texttt{----------}\\
8 & 4 & 5 & 1 & 1 & 0 & 1024 & 0 & \texttt{----------}\\
9 & 2 & 3 & 1 & 1 & 0 & 1024 & 0 & \texttt{----------}\\
\multicolumn{9}{@{}l@{}}{\rule{0pt}{2.6ex}\textbf{T10.028}\enspace \texttt{0123456752}\enspace $n{=}10$, $|H|{=}3$, len 8, br 2, $|\mathrm{Aut}|{=}1$, orb 10/10, hist-orb 3/3; M 35840, res 9824, rev 11882: ok}\\*
0 & -- & 0 & 3 & 1 & 750 & 2 & 272 & \texttt{-+????????}\\*
1 & 0 & 1 & 3 & 1 & 476 & 4 & 544 & \texttt{--????????}\\
2 & 1 & 2 & 2 & 1 & 880 & 16 & 128 & \texttt{---++????-}\\
3 & 2 & 3 & 2 & 1 & 736 & 32 & 256 & \texttt{----+????-}\\
4 & 3 & 4 & 2 & 1 & 448 & 64 & 512 & \texttt{-----????-}\\
5 & 4 & 5 & 1 & 1 & 768 & 256 & 0 & \texttt{------++--}\\
6 & 5 & 6 & 1 & 1 & 512 & 512 & 0 & \texttt{-------+--}\\
7 & 6 & 7 & 1 & 1 & 0 & 1024 & 0 & \texttt{----------}\\
8 & 4 & 5 & 1 & 1 & 0 & 1024 & 0 & \texttt{----------}\\
9 & 1 & 2 & 1 & 1 & 0 & 1024 & 0 & \texttt{----------}\\
\multicolumn{9}{@{}l@{}}{\rule{0pt}{2.6ex}\textbf{T10.029}\enspace \texttt{0123456751}\enspace $n{=}10$, $|H|{=}3$, len 8, br 2, $|\mathrm{Aut}|{=}1$, orb 10/10, hist-orb 3/3; M 34816, res 8464, rev 11766: ok}\\*
0 & -- & 0 & 3 & 1 & 494 & 2 & 528 & \texttt{-?????????}\\*
1 & 0 & 1 & 2 & 1 & 952 & 8 & 64 & \texttt{--+++????-}\\
2 & 1 & 2 & 2 & 1 & 880 & 16 & 128 & \texttt{---++????-}\\
3 & 2 & 3 & 2 & 1 & 736 & 32 & 256 & \texttt{----+????-}\\
4 & 3 & 4 & 2 & 1 & 448 & 64 & 512 & \texttt{-----????-}\\
5 & 4 & 5 & 1 & 1 & 768 & 256 & 0 & \texttt{------++--}\\
6 & 5 & 6 & 1 & 1 & 512 & 512 & 0 & \texttt{-------+--}\\
7 & 6 & 7 & 1 & 1 & 0 & 1024 & 0 & \texttt{----------}\\
8 & 4 & 5 & 1 & 1 & 0 & 1024 & 0 & \texttt{----------}\\
9 & 0 & 1 & 1 & 1 & 0 & 1024 & 0 & \texttt{----------}\\
\multicolumn{9}{@{}l@{}}{\rule{0pt}{2.6ex}\textbf{T10.030}\enspace \texttt{0123456745}\enspace $n{=}10$, $|H|{=}2$, len 8, br 1, $|\mathrm{Aut}|{=}1$, orb 10/10, hist-orb 2/2; M 37888, res 3240, rev 14986: ok}\\*
0 & -- & 0 & 2 & 1 & 986 & 2 & 36 & \texttt{-+++??????}\\*
1 & 0 & 1 & 2 & 1 & 948 & 4 & 72 & \texttt{--++??????}\\
2 & 1 & 2 & 2 & 1 & 872 & 8 & 144 & \texttt{---+??????}\\
3 & 2 & 3 & 2 & 1 & 720 & 16 & 288 & \texttt{----??????}\\
4 & 3 & 4 & 1 & 1 & 896 & 128 & 0 & \texttt{-----+++--}\\
5 & 4 & 5 & 1 & 1 & 768 & 256 & 0 & \texttt{------++--}\\
6 & 5 & 6 & 1 & 1 & 512 & 512 & 0 & \texttt{-------+--}\\
7 & 6 & 7 & 1 & 1 & 0 & 1024 & 0 & \texttt{----------}\\
8 & 3 & 4 & 1 & 1 & 512 & 512 & 0 & \texttt{---------+}\\
9 & 8 & 5 & 1 & 1 & 0 & 1024 & 0 & \texttt{----------}\\
\multicolumn{9}{@{}l@{}}{\rule{0pt}{2.6ex}\textbf{T10.031}\enspace \texttt{0123456744}\enspace $n{=}10$, $|H|{=}3$, len 8, br 1, $|\mathrm{Aut}|{=}2$, orb 9/8, hist-orb 2/1; M 36864, res 8640, rev 13102: ok}\\*
0 & -- & 0 & 3 & 1 & 926 & 2 & 96 & \texttt{-+++??????}\\*
1 & 0 & 1 & 3 & 1 & 828 & 4 & 192 & \texttt{--++??????}\\
2 & 1 & 2 & 3 & 1 & 632 & 8 & 384 & \texttt{---+??????}\\
3 & 2 & 3 & 3 & 1 & 240 & 16 & 768 & \texttt{----??????}\\
4 & 3 & 4 & 1 & 1 & 896 & 128 & 0 & \texttt{-----+++--}\\
5 & 4 & 5 & 1 & 1 & 768 & 256 & 0 & \texttt{------++--}\\
6 & 5 & 6 & 1 & 1 & 512 & 512 & 0 & \texttt{-------+--}\\
7 & 6 & 7 & 1 & 1 & 0 & 1024 & 0 & \texttt{----------}\\
8 & 3 & 4 & 1 & 2 & 0 & 1024 & 0 & \texttt{----------}\\
9 & 3 & 4 & 1 & 2 & 0 & 1024 & 0 & \texttt{----------}\\
\multicolumn{9}{@{}l@{}}{\rule{0pt}{2.6ex}\textbf{T10.032}\enspace \texttt{0123456743}\enspace $n{=}10$, $|H|{=}3$, len 8, br 2, $|\mathrm{Aut}|{=}1$, orb 10/10, hist-orb 3/3; M 35840, res 9952, rev 11802: ok}\\*
0 & -- & 0 & 3 & 1 & 862 & 2 & 160 & \texttt{-++???????}\\*
1 & 0 & 1 & 3 & 1 & 700 & 4 & 320 & \texttt{--+???????}\\
2 & 1 & 2 & 3 & 1 & 376 & 8 & 640 & \texttt{---???????}\\
3 & 2 & 3 & 2 & 1 & 480 & 32 & 512 & \texttt{----?????-}\\
4 & 3 & 4 & 1 & 1 & 896 & 128 & 0 & \texttt{-----+++--}\\
5 & 4 & 5 & 1 & 1 & 768 & 256 & 0 & \texttt{------++--}\\
6 & 5 & 6 & 1 & 1 & 512 & 512 & 0 & \texttt{-------+--}\\
7 & 6 & 7 & 1 & 1 & 0 & 1024 & 0 & \texttt{----------}\\
8 & 3 & 4 & 1 & 1 & 0 & 1024 & 0 & \texttt{----------}\\
9 & 2 & 3 & 1 & 1 & 0 & 1024 & 0 & \texttt{----------}\\
\multicolumn{9}{@{}l@{}}{\rule{0pt}{2.6ex}\textbf{T10.033}\enspace \texttt{0123456742}\enspace $n{=}10$, $|H|{=}3$, len 8, br 2, $|\mathrm{Aut}|{=}1$, orb 10/10, hist-orb 3/3; M 34816, res 10368, rev 11002: ok}\\*
0 & -- & 0 & 3 & 1 & 734 & 2 & 288 & \texttt{-+????????}\\*
1 & 0 & 1 & 3 & 1 & 444 & 4 & 576 & \texttt{--????????}\\
2 & 1 & 2 & 2 & 1 & 752 & 16 & 256 & \texttt{---+?????-}\\
3 & 2 & 3 & 2 & 1 & 480 & 32 & 512 & \texttt{----?????-}\\
4 & 3 & 4 & 1 & 1 & 896 & 128 & 0 & \texttt{-----+++--}\\
5 & 4 & 5 & 1 & 1 & 768 & 256 & 0 & \texttt{------++--}\\
6 & 5 & 6 & 1 & 1 & 512 & 512 & 0 & \texttt{-------+--}\\
7 & 6 & 7 & 1 & 1 & 0 & 1024 & 0 & \texttt{----------}\\
8 & 3 & 4 & 1 & 1 & 0 & 1024 & 0 & \texttt{----------}\\
9 & 1 & 2 & 1 & 1 & 0 & 1024 & 0 & \texttt{----------}\\
\multicolumn{9}{@{}l@{}}{\rule{0pt}{2.6ex}\textbf{T10.034}\enspace \texttt{0123456741}\enspace $n{=}10$, $|H|{=}3$, len 8, br 2, $|\mathrm{Aut}|{=}1$, orb 10/10, hist-orb 3/3; M 33792, res 9184, rev 10838: ok}\\*
0 & -- & 0 & 3 & 1 & 478 & 2 & 544 & \texttt{-?????????}\\*
1 & 0 & 1 & 2 & 1 & 888 & 8 & 128 & \texttt{--++?????-}\\
2 & 1 & 2 & 2 & 1 & 752 & 16 & 256 & \texttt{---+?????-}\\
3 & 2 & 3 & 2 & 1 & 480 & 32 & 512 & \texttt{----?????-}\\
4 & 3 & 4 & 1 & 1 & 896 & 128 & 0 & \texttt{-----+++--}\\
5 & 4 & 5 & 1 & 1 & 768 & 256 & 0 & \texttt{------++--}\\
6 & 5 & 6 & 1 & 1 & 512 & 512 & 0 & \texttt{-------+--}\\
7 & 6 & 7 & 1 & 1 & 0 & 1024 & 0 & \texttt{----------}\\
8 & 3 & 4 & 1 & 1 & 0 & 1024 & 0 & \texttt{----------}\\
9 & 0 & 1 & 1 & 1 & 0 & 1024 & 0 & \texttt{----------}\\
\multicolumn{9}{@{}l@{}}{\rule{0pt}{2.6ex}\textbf{T10.035}\enspace \texttt{0123456734}\enspace $n{=}10$, $|H|{=}2$, len 8, br 1, $|\mathrm{Aut}|{=}1$, orb 10/10, hist-orb 2/2; M 35840, res 3332, rev 13662: ok}\\*
0 & -- & 0 & 2 & 1 & 954 & 2 & 68 & \texttt{-++???????}\\*
1 & 0 & 1 & 2 & 1 & 884 & 4 & 136 & \texttt{--+???????}\\
2 & 1 & 2 & 2 & 1 & 744 & 8 & 272 & \texttt{---???????}\\
3 & 2 & 3 & 1 & 1 & 960 & 64 & 0 & \texttt{----++++--}\\
4 & 3 & 4 & 1 & 1 & 896 & 128 & 0 & \texttt{-----+++--}\\
5 & 4 & 5 & 1 & 1 & 768 & 256 & 0 & \texttt{------++--}\\
6 & 5 & 6 & 1 & 1 & 512 & 512 & 0 & \texttt{-------+--}\\
7 & 6 & 7 & 1 & 1 & 0 & 1024 & 0 & \texttt{----------}\\
8 & 2 & 3 & 1 & 1 & 512 & 512 & 0 & \texttt{---------+}\\
9 & 8 & 4 & 1 & 1 & 0 & 1024 & 0 & \texttt{----------}\\
\multicolumn{9}{@{}l@{}}{\rule{0pt}{2.6ex}\textbf{T10.036}\enspace \texttt{0123456733}\enspace $n{=}10$, $|H|{=}3$, len 8, br 1, $|\mathrm{Aut}|{=}2$, orb 9/8, hist-orb 2/1; M 34816, res 9408, rev 11326: ok}\\*
0 & -- & 0 & 3 & 1 & 830 & 2 & 192 & \texttt{-++???????}\\*
1 & 0 & 1 & 3 & 1 & 636 & 4 & 384 & \texttt{--+???????}\\
2 & 1 & 2 & 3 & 1 & 248 & 8 & 768 & \texttt{---???????}\\
3 & 2 & 3 & 1 & 1 & 960 & 64 & 0 & \texttt{----++++--}\\
4 & 3 & 4 & 1 & 1 & 896 & 128 & 0 & \texttt{-----+++--}\\
5 & 4 & 5 & 1 & 1 & 768 & 256 & 0 & \texttt{------++--}\\
6 & 5 & 6 & 1 & 1 & 512 & 512 & 0 & \texttt{-------+--}\\
7 & 6 & 7 & 1 & 1 & 0 & 1024 & 0 & \texttt{----------}\\
8 & 2 & 3 & 1 & 2 & 0 & 1024 & 0 & \texttt{----------}\\
9 & 2 & 3 & 1 & 2 & 0 & 1024 & 0 & \texttt{----------}\\
\multicolumn{9}{@{}l@{}}{\rule{0pt}{2.6ex}\textbf{T10.037}\enspace \texttt{0123456732}\enspace $n{=}10$, $|H|{=}3$, len 8, br 2, $|\mathrm{Aut}|{=}1$, orb 10/10, hist-orb 3/3; M 33792, res 10368, rev 10274: ok}\\*
0 & -- & 0 & 3 & 1 & 702 & 2 & 320 & \texttt{-+????????}\\*
1 & 0 & 1 & 3 & 1 & 380 & 4 & 640 & \texttt{--????????}\\
2 & 1 & 2 & 2 & 1 & 496 & 16 & 512 & \texttt{---??????-}\\
3 & 2 & 3 & 1 & 1 & 960 & 64 & 0 & \texttt{----++++--}\\
4 & 3 & 4 & 1 & 1 & 896 & 128 & 0 & \texttt{-----+++--}\\
5 & 4 & 5 & 1 & 1 & 768 & 256 & 0 & \texttt{------++--}\\
6 & 5 & 6 & 1 & 1 & 512 & 512 & 0 & \texttt{-------+--}\\
7 & 6 & 7 & 1 & 1 & 0 & 1024 & 0 & \texttt{----------}\\
8 & 2 & 3 & 1 & 1 & 0 & 1024 & 0 & \texttt{----------}\\
9 & 1 & 2 & 1 & 1 & 0 & 1024 & 0 & \texttt{----------}\\
\multicolumn{9}{@{}l@{}}{\rule{0pt}{2.6ex}\textbf{T10.038}\enspace \texttt{0123456731}\enspace $n{=}10$, $|H|{=}3$, len 8, br 2, $|\mathrm{Aut}|{=}1$, orb 10/10, hist-orb 3/3; M 32768, res 9536, rev 9990: ok}\\*
0 & -- & 0 & 3 & 1 & 446 & 2 & 576 & \texttt{-?????????}\\*
1 & 0 & 1 & 2 & 1 & 760 & 8 & 256 & \texttt{--+??????-}\\
2 & 1 & 2 & 2 & 1 & 496 & 16 & 512 & \texttt{---??????-}\\
3 & 2 & 3 & 1 & 1 & 960 & 64 & 0 & \texttt{----++++--}\\
4 & 3 & 4 & 1 & 1 & 896 & 128 & 0 & \texttt{-----+++--}\\
5 & 4 & 5 & 1 & 1 & 768 & 256 & 0 & \texttt{------++--}\\
6 & 5 & 6 & 1 & 1 & 512 & 512 & 0 & \texttt{-------+--}\\
7 & 6 & 7 & 1 & 1 & 0 & 1024 & 0 & \texttt{----------}\\
8 & 2 & 3 & 1 & 1 & 0 & 1024 & 0 & \texttt{----------}\\
9 & 0 & 1 & 1 & 1 & 0 & 1024 & 0 & \texttt{----------}\\
\multicolumn{9}{@{}l@{}}{\rule{0pt}{2.6ex}\textbf{T10.039}\enspace \texttt{0123456723}\enspace $n{=}10$, $|H|{=}2$, len 8, br 1, $|\mathrm{Aut}|{=}1$, orb 10/10, hist-orb 2/2; M 33792, res 3168, rev 12394: ok}\\*
0 & -- & 0 & 2 & 1 & 890 & 2 & 132 & \texttt{-+????????}\\*
1 & 0 & 1 & 2 & 1 & 756 & 4 & 264 & \texttt{--????????}\\
2 & 1 & 2 & 1 & 1 & 992 & 32 & 0 & \texttt{---+++++--}\\
3 & 2 & 3 & 1 & 1 & 960 & 64 & 0 & \texttt{----++++--}\\
4 & 3 & 4 & 1 & 1 & 896 & 128 & 0 & \texttt{-----+++--}\\
5 & 4 & 5 & 1 & 1 & 768 & 256 & 0 & \texttt{------++--}\\
6 & 5 & 6 & 1 & 1 & 512 & 512 & 0 & \texttt{-------+--}\\
7 & 6 & 7 & 1 & 1 & 0 & 1024 & 0 & \texttt{----------}\\
8 & 1 & 2 & 1 & 1 & 512 & 512 & 0 & \texttt{---------+}\\
9 & 8 & 3 & 1 & 1 & 0 & 1024 & 0 & \texttt{----------}\\
\multicolumn{9}{@{}l@{}}{\rule{0pt}{2.6ex}\textbf{T10.040}\enspace \texttt{0123456722}\enspace $n{=}10$, $|H|{=}3$, len 8, br 1, $|\mathrm{Aut}|{=}2$, orb 9/8, hist-orb 2/1; M 32768, res 9216, rev 9942: ok}\\*
0 & -- & 0 & 3 & 1 & 638 & 2 & 384 & \texttt{-+????????}\\*
1 & 0 & 1 & 3 & 1 & 252 & 4 & 768 & \texttt{--????????}\\
2 & 1 & 2 & 1 & 1 & 992 & 32 & 0 & \texttt{---+++++--}\\
3 & 2 & 3 & 1 & 1 & 960 & 64 & 0 & \texttt{----++++--}\\
4 & 3 & 4 & 1 & 1 & 896 & 128 & 0 & \texttt{-----+++--}\\
5 & 4 & 5 & 1 & 1 & 768 & 256 & 0 & \texttt{------++--}\\
6 & 5 & 6 & 1 & 1 & 512 & 512 & 0 & \texttt{-------+--}\\
7 & 6 & 7 & 1 & 1 & 0 & 1024 & 0 & \texttt{----------}\\
8 & 1 & 2 & 1 & 2 & 0 & 1024 & 0 & \texttt{----------}\\
9 & 1 & 2 & 1 & 2 & 0 & 1024 & 0 & \texttt{----------}\\
\multicolumn{9}{@{}l@{}}{\rule{0pt}{2.6ex}\textbf{T10.041}\enspace \texttt{0123456721}\enspace $n{=}10$, $|H|{=}3$, len 8, br 2, $|\mathrm{Aut}|{=}1$, orb 10/10, hist-orb 3/3; M 31744, res 9088, rev 9406: ok}\\*
0 & -- & 0 & 3 & 1 & 382 & 2 & 640 & \texttt{-?????????}\\*
1 & 0 & 1 & 2 & 1 & 504 & 8 & 512 & \texttt{--???????-}\\
2 & 1 & 2 & 1 & 1 & 992 & 32 & 0 & \texttt{---+++++--}\\
3 & 2 & 3 & 1 & 1 & 960 & 64 & 0 & \texttt{----++++--}\\
4 & 3 & 4 & 1 & 1 & 896 & 128 & 0 & \texttt{-----+++--}\\
5 & 4 & 5 & 1 & 1 & 768 & 256 & 0 & \texttt{------++--}\\
6 & 5 & 6 & 1 & 1 & 512 & 512 & 0 & \texttt{-------+--}\\
7 & 6 & 7 & 1 & 1 & 0 & 1024 & 0 & \texttt{----------}\\
8 & 1 & 2 & 1 & 1 & 0 & 1024 & 0 & \texttt{----------}\\
9 & 0 & 1 & 1 & 1 & 0 & 1024 & 0 & \texttt{----------}\\
\multicolumn{9}{@{}l@{}}{\rule{0pt}{2.6ex}\textbf{T10.042}\enspace \texttt{0123456712}\enspace $n{=}10$, $|H|{=}2$, len 8, br 1, $|\mathrm{Aut}|{=}1$, orb 10/10, hist-orb 2/2; M 31744, res 2340, rev 11362: ok}\\*
0 & -- & 0 & 2 & 1 & 762 & 2 & 260 & \texttt{-?????????}\\*
1 & 0 & 1 & 1 & 1 & 1008 & 16 & 0 & \texttt{--++++++--}\\
2 & 1 & 2 & 1 & 1 & 992 & 32 & 0 & \texttt{---+++++--}\\
3 & 2 & 3 & 1 & 1 & 960 & 64 & 0 & \texttt{----++++--}\\
4 & 3 & 4 & 1 & 1 & 896 & 128 & 0 & \texttt{-----+++--}\\
5 & 4 & 5 & 1 & 1 & 768 & 256 & 0 & \texttt{------++--}\\
6 & 5 & 6 & 1 & 1 & 512 & 512 & 0 & \texttt{-------+--}\\
7 & 6 & 7 & 1 & 1 & 0 & 1024 & 0 & \texttt{----------}\\
8 & 0 & 1 & 1 & 1 & 512 & 512 & 0 & \texttt{---------+}\\
9 & 8 & 2 & 1 & 1 & 0 & 1024 & 0 & \texttt{----------}\\
\multicolumn{9}{@{}l@{}}{\rule{0pt}{2.6ex}\textbf{T10.043}\enspace \texttt{0123456711}\enspace $n{=}10$, $|H|{=}3$, len 8, br 1, $|\mathrm{Aut}|{=}2$, orb 9/8, hist-orb 2/1; M 30720, res 6912, rev 9354: ok}\\*
0 & -- & 0 & 3 & 1 & 254 & 2 & 768 & \texttt{-?????????}\\*
1 & 0 & 1 & 1 & 1 & 1008 & 16 & 0 & \texttt{--++++++--}\\
2 & 1 & 2 & 1 & 1 & 992 & 32 & 0 & \texttt{---+++++--}\\
3 & 2 & 3 & 1 & 1 & 960 & 64 & 0 & \texttt{----++++--}\\
4 & 3 & 4 & 1 & 1 & 896 & 128 & 0 & \texttt{-----+++--}\\
5 & 4 & 5 & 1 & 1 & 768 & 256 & 0 & \texttt{------++--}\\
6 & 5 & 6 & 1 & 1 & 512 & 512 & 0 & \texttt{-------+--}\\
7 & 6 & 7 & 1 & 1 & 0 & 1024 & 0 & \texttt{----------}\\
8 & 0 & 1 & 1 & 2 & 0 & 1024 & 0 & \texttt{----------}\\
9 & 0 & 1 & 1 & 2 & 0 & 1024 & 0 & \texttt{----------}\\
\multicolumn{9}{@{}l@{}}{\rule{0pt}{2.6ex}\textbf{T10.044}\enspace \texttt{0123456666}\enspace $n{=}10$, $|H|{=}4$, len 7, br 1, $|\mathrm{Aut}|{=}24$, orb 7/6, hist-orb 1/0; M 39936, res 7056, rev 17418: ok}\\*
0 & -- & 0 & 4 & 1 & 994 & 2 & 28 & \texttt{-+++++????}\\*
1 & 0 & 1 & 4 & 1 & 964 & 4 & 56 & \texttt{--++++????}\\
2 & 1 & 2 & 4 & 1 & 904 & 8 & 112 & \texttt{---+++????}\\
3 & 2 & 3 & 4 & 1 & 784 & 16 & 224 & \texttt{----++????}\\
4 & 3 & 4 & 4 & 1 & 544 & 32 & 448 & \texttt{-----+????}\\
5 & 4 & 5 & 4 & 1 & 64 & 64 & 896 & \texttt{------????}\\
6 & 5 & 6 & 1 & 4 & 0 & 1024 & 0 & \texttt{----------}\\
7 & 5 & 6 & 1 & 4 & 0 & 1024 & 0 & \texttt{----------}\\
8 & 5 & 6 & 1 & 4 & 0 & 1024 & 0 & \texttt{----------}\\
9 & 5 & 6 & 1 & 4 & 0 & 1024 & 0 & \texttt{----------}\\
\multicolumn{9}{@{}l@{}}{\rule{0pt}{2.6ex}\textbf{T10.045}\enspace \texttt{0123456665}\enspace $n{=}10$, $|H|{=}4$, len 7, br 2, $|\mathrm{Aut}|{=}6$, orb 8/7, hist-orb 2/1; M 38912, res 8380, rev 15758: ok}\\*
0 & -- & 0 & 4 & 1 & 978 & 2 & 44 & \texttt{-++++?????}\\*
1 & 0 & 1 & 4 & 1 & 932 & 4 & 88 & \texttt{--+++?????}\\
2 & 1 & 2 & 4 & 1 & 840 & 8 & 176 & \texttt{---++?????}\\
3 & 2 & 3 & 4 & 1 & 656 & 16 & 352 & \texttt{----+?????}\\
4 & 3 & 4 & 4 & 1 & 288 & 32 & 704 & \texttt{-----?????}\\
5 & 4 & 5 & 3 & 1 & 128 & 128 & 768 & \texttt{------???-}\\
6 & 5 & 6 & 1 & 3 & 0 & 1024 & 0 & \texttt{----------}\\
7 & 5 & 6 & 1 & 3 & 0 & 1024 & 0 & \texttt{----------}\\
8 & 5 & 6 & 1 & 3 & 0 & 1024 & 0 & \texttt{----------}\\
9 & 4 & 5 & 1 & 1 & 0 & 1024 & 0 & \texttt{----------}\\
\multicolumn{9}{@{}l@{}}{\rule{0pt}{2.6ex}\textbf{T10.046}\enspace \texttt{0123456664}\enspace $n{=}10$, $|H|{=}4$, len 7, br 2, $|\mathrm{Aut}|{=}6$, orb 8/7, hist-orb 2/1; M 37888, res 9576, rev 14454: ok}\\*
0 & -- & 0 & 4 & 1 & 946 & 2 & 76 & \texttt{-+++??????}\\*
1 & 0 & 1 & 4 & 1 & 868 & 4 & 152 & \texttt{--++??????}\\
2 & 1 & 2 & 4 & 1 & 712 & 8 & 304 & \texttt{---+??????}\\
3 & 2 & 3 & 4 & 1 & 400 & 16 & 608 & \texttt{----??????}\\
4 & 3 & 4 & 3 & 1 & 576 & 64 & 384 & \texttt{-----+???-}\\
5 & 4 & 5 & 3 & 1 & 128 & 128 & 768 & \texttt{------???-}\\
6 & 5 & 6 & 1 & 3 & 0 & 1024 & 0 & \texttt{----------}\\
7 & 5 & 6 & 1 & 3 & 0 & 1024 & 0 & \texttt{----------}\\
8 & 5 & 6 & 1 & 3 & 0 & 1024 & 0 & \texttt{----------}\\
9 & 3 & 4 & 1 & 1 & 0 & 1024 & 0 & \texttt{----------}\\
\multicolumn{9}{@{}l@{}}{\rule{0pt}{2.6ex}\textbf{T10.047}\enspace \texttt{0123456663}\enspace $n{=}10$, $|H|{=}4$, len 7, br 2, $|\mathrm{Aut}|{=}6$, orb 8/7, hist-orb 2/1; M 36864, res 10388, rev 13462: ok}\\*
0 & -- & 0 & 4 & 1 & 882 & 2 & 140 & \texttt{-++???????}\\*
1 & 0 & 1 & 4 & 1 & 740 & 4 & 280 & \texttt{--+???????}\\
2 & 1 & 2 & 4 & 1 & 456 & 8 & 560 & \texttt{---???????}\\
3 & 2 & 3 & 3 & 1 & 800 & 32 & 192 & \texttt{----++???-}\\
4 & 3 & 4 & 3 & 1 & 576 & 64 & 384 & \texttt{-----+???-}\\
5 & 4 & 5 & 3 & 1 & 128 & 128 & 768 & \texttt{------???-}\\
6 & 5 & 6 & 1 & 3 & 0 & 1024 & 0 & \texttt{----------}\\
7 & 5 & 6 & 1 & 3 & 0 & 1024 & 0 & \texttt{----------}\\
8 & 5 & 6 & 1 & 3 & 0 & 1024 & 0 & \texttt{----------}\\
9 & 2 & 3 & 1 & 1 & 0 & 1024 & 0 & \texttt{----------}\\
\multicolumn{9}{@{}l@{}}{\rule{0pt}{2.6ex}\textbf{T10.048}\enspace \texttt{0123456662}\enspace $n{=}10$, $|H|{=}4$, len 7, br 2, $|\mathrm{Aut}|{=}6$, orb 8/7, hist-orb 2/1; M 35840, res 10464, rev 12846: ok}\\*
0 & -- & 0 & 4 & 1 & 754 & 2 & 268 & \texttt{-+????????}\\*
1 & 0 & 1 & 4 & 1 & 484 & 4 & 536 & \texttt{--????????}\\
2 & 1 & 2 & 3 & 1 & 912 & 16 & 96 & \texttt{---+++???-}\\
3 & 2 & 3 & 3 & 1 & 800 & 32 & 192 & \texttt{----++???-}\\
4 & 3 & 4 & 3 & 1 & 576 & 64 & 384 & \texttt{-----+???-}\\
5 & 4 & 5 & 3 & 1 & 128 & 128 & 768 & \texttt{------???-}\\
6 & 5 & 6 & 1 & 3 & 0 & 1024 & 0 & \texttt{----------}\\
7 & 5 & 6 & 1 & 3 & 0 & 1024 & 0 & \texttt{----------}\\
8 & 5 & 6 & 1 & 3 & 0 & 1024 & 0 & \texttt{----------}\\
9 & 1 & 2 & 1 & 1 & 0 & 1024 & 0 & \texttt{----------}\\
\multicolumn{9}{@{}l@{}}{\rule{0pt}{2.6ex}\textbf{T10.049}\enspace \texttt{0123456661}\enspace $n{=}10$, $|H|{=}4$, len 7, br 2, $|\mathrm{Aut}|{=}6$, orb 8/7, hist-orb 2/1; M 34816, res 9060, rev 12834: ok}\\*
0 & -- & 0 & 4 & 1 & 498 & 2 & 524 & \texttt{-?????????}\\*
1 & 0 & 1 & 3 & 1 & 968 & 8 & 48 & \texttt{--++++???-}\\
2 & 1 & 2 & 3 & 1 & 912 & 16 & 96 & \texttt{---+++???-}\\
3 & 2 & 3 & 3 & 1 & 800 & 32 & 192 & \texttt{----++???-}\\
4 & 3 & 4 & 3 & 1 & 576 & 64 & 384 & \texttt{-----+???-}\\
5 & 4 & 5 & 3 & 1 & 128 & 128 & 768 & \texttt{------???-}\\
6 & 5 & 6 & 1 & 3 & 0 & 1024 & 0 & \texttt{----------}\\
7 & 5 & 6 & 1 & 3 & 0 & 1024 & 0 & \texttt{----------}\\
8 & 5 & 6 & 1 & 3 & 0 & 1024 & 0 & \texttt{----------}\\
9 & 0 & 1 & 1 & 1 & 0 & 1024 & 0 & \texttt{----------}\\
\multicolumn{9}{@{}l@{}}{\rule{0pt}{2.6ex}\textbf{T10.050}\enspace \texttt{0123456656}\enspace $n{=}10$, $|H|{=}3$, len 7, br 2, $|\mathrm{Aut}|{=}2$, orb 9/8, hist-orb 2/1; M 38912, res 5488, rev 16214: ok}\\*
0 & -- & 0 & 3 & 1 & 990 & 2 & 32 & \texttt{-++++?????}\\*
1 & 0 & 1 & 3 & 1 & 956 & 4 & 64 & \texttt{--+++?????}\\
2 & 1 & 2 & 3 & 1 & 888 & 8 & 128 & \texttt{---++?????}\\
3 & 2 & 3 & 3 & 1 & 752 & 16 & 256 & \texttt{----+?????}\\
4 & 3 & 4 & 3 & 1 & 480 & 32 & 512 & \texttt{-----?????}\\
5 & 4 & 5 & 2 & 1 & 256 & 256 & 512 & \texttt{------??--}\\
6 & 5 & 6 & 1 & 2 & 0 & 1024 & 0 & \texttt{----------}\\
7 & 5 & 6 & 1 & 2 & 0 & 1024 & 0 & \texttt{----------}\\
8 & 4 & 5 & 1 & 1 & 512 & 512 & 0 & \texttt{---------+}\\
9 & 8 & 6 & 1 & 1 & 0 & 1024 & 0 & \texttt{----------}\\
\multicolumn{9}{@{}l@{}}{\rule{0pt}{2.6ex}\textbf{T10.051}\enspace \texttt{0123456655}\enspace $n{=}10$, $|H|{=}4$, len 7, br 2, $|\mathrm{Aut}|{=}4$, orb 8/6, hist-orb 2/0; M 37888, res 8588, rev 15286: ok}\\*
0 & -- & 0 & 4 & 1 & 970 & 2 & 52 & \texttt{-++++?????}\\*
1 & 0 & 1 & 4 & 1 & 916 & 4 & 104 & \texttt{--+++?????}\\
2 & 1 & 2 & 4 & 1 & 808 & 8 & 208 & \texttt{---++?????}\\
3 & 2 & 3 & 4 & 1 & 592 & 16 & 416 & \texttt{----+?????}\\
4 & 3 & 4 & 4 & 1 & 160 & 32 & 832 & \texttt{-----?????}\\
5 & 4 & 5 & 2 & 1 & 256 & 256 & 512 & \texttt{------??--}\\
6 & 5 & 6 & 1 & 2 & 0 & 1024 & 0 & \texttt{----------}\\
7 & 5 & 6 & 1 & 2 & 0 & 1024 & 0 & \texttt{----------}\\
8 & 4 & 5 & 1 & 2 & 0 & 1024 & 0 & \texttt{----------}\\
9 & 4 & 5 & 1 & 2 & 0 & 1024 & 0 & \texttt{----------}\\
\multicolumn{9}{@{}l@{}}{\rule{0pt}{2.6ex}\textbf{T10.052}\enspace \texttt{0123456654}\enspace $n{=}10$, $|H|{=}4$, len 7, br 3, $|\mathrm{Aut}|{=}2$, orb 9/8, hist-orb 3/2; M 36864, res 10048, rev 13666: ok}\\*
0 & -- & 0 & 4 & 1 & 938 & 2 & 84 & \texttt{-+++??????}\\*
1 & 0 & 1 & 4 & 1 & 852 & 4 & 168 & \texttt{--++??????}\\
2 & 1 & 2 & 4 & 1 & 680 & 8 & 336 & \texttt{---+??????}\\
3 & 2 & 3 & 4 & 1 & 336 & 16 & 672 & \texttt{----??????}\\
4 & 3 & 4 & 3 & 1 & 320 & 64 & 640 & \texttt{-----????-}\\
5 & 4 & 5 & 2 & 1 & 256 & 256 & 512 & \texttt{------??--}\\
6 & 5 & 6 & 1 & 2 & 0 & 1024 & 0 & \texttt{----------}\\
7 & 5 & 6 & 1 & 2 & 0 & 1024 & 0 & \texttt{----------}\\
8 & 4 & 5 & 1 & 1 & 0 & 1024 & 0 & \texttt{----------}\\
9 & 3 & 4 & 1 & 1 & 0 & 1024 & 0 & \texttt{----------}\\
\multicolumn{9}{@{}l@{}}{\rule{0pt}{2.6ex}\textbf{T10.053}\enspace \texttt{0123456653}\enspace $n{=}10$, $|H|{=}4$, len 7, br 3, $|\mathrm{Aut}|{=}2$, orb 9/8, hist-orb 3/2; M 35840, res 11060, rev 12522: ok}\\*
0 & -- & 0 & 4 & 1 & 874 & 2 & 148 & \texttt{-++???????}\\*
1 & 0 & 1 & 4 & 1 & 724 & 4 & 296 & \texttt{--+???????}\\
2 & 1 & 2 & 4 & 1 & 424 & 8 & 592 & \texttt{---???????}\\
3 & 2 & 3 & 3 & 1 & 672 & 32 & 320 & \texttt{----+????-}\\
4 & 3 & 4 & 3 & 1 & 320 & 64 & 640 & \texttt{-----????-}\\
5 & 4 & 5 & 2 & 1 & 256 & 256 & 512 & \texttt{------??--}\\
6 & 5 & 6 & 1 & 2 & 0 & 1024 & 0 & \texttt{----------}\\
7 & 5 & 6 & 1 & 2 & 0 & 1024 & 0 & \texttt{----------}\\
8 & 4 & 5 & 1 & 1 & 0 & 1024 & 0 & \texttt{----------}\\
9 & 2 & 3 & 1 & 1 & 0 & 1024 & 0 & \texttt{----------}\\
\multicolumn{9}{@{}l@{}}{\rule{0pt}{2.6ex}\textbf{T10.054}\enspace \texttt{0123456652}\enspace $n{=}10$, $|H|{=}4$, len 7, br 3, $|\mathrm{Aut}|{=}2$, orb 9/8, hist-orb 3/2; M 34816, res 11272, rev 11842: ok}\\*
0 & -- & 0 & 4 & 1 & 746 & 2 & 276 & \texttt{-+????????}\\*
1 & 0 & 1 & 4 & 1 & 468 & 4 & 552 & \texttt{--????????}\\
2 & 1 & 2 & 3 & 1 & 848 & 16 & 160 & \texttt{---++????-}\\
3 & 2 & 3 & 3 & 1 & 672 & 32 & 320 & \texttt{----+????-}\\
4 & 3 & 4 & 3 & 1 & 320 & 64 & 640 & \texttt{-----????-}\\
5 & 4 & 5 & 2 & 1 & 256 & 256 & 512 & \texttt{------??--}\\
6 & 5 & 6 & 1 & 2 & 0 & 1024 & 0 & \texttt{----------}\\
7 & 5 & 6 & 1 & 2 & 0 & 1024 & 0 & \texttt{----------}\\
8 & 4 & 5 & 1 & 1 & 0 & 1024 & 0 & \texttt{----------}\\
9 & 1 & 2 & 1 & 1 & 0 & 1024 & 0 & \texttt{----------}\\
\multicolumn{9}{@{}l@{}}{\rule{0pt}{2.6ex}\textbf{T10.055}\enspace \texttt{0123456651}\enspace $n{=}10$, $|H|{=}4$, len 7, br 3, $|\mathrm{Aut}|{=}2$, orb 9/8, hist-orb 3/2; M 33792, res 9956, rev 11822: ok}\\*
0 & -- & 0 & 4 & 1 & 490 & 2 & 532 & \texttt{-?????????}\\*
1 & 0 & 1 & 3 & 1 & 936 & 8 & 80 & \texttt{--+++????-}\\
2 & 1 & 2 & 3 & 1 & 848 & 16 & 160 & \texttt{---++????-}\\
3 & 2 & 3 & 3 & 1 & 672 & 32 & 320 & \texttt{----+????-}\\
4 & 3 & 4 & 3 & 1 & 320 & 64 & 640 & \texttt{-----????-}\\
5 & 4 & 5 & 2 & 1 & 256 & 256 & 512 & \texttt{------??--}\\
6 & 5 & 6 & 1 & 2 & 0 & 1024 & 0 & \texttt{----------}\\
7 & 5 & 6 & 1 & 2 & 0 & 1024 & 0 & \texttt{----------}\\
8 & 4 & 5 & 1 & 1 & 0 & 1024 & 0 & \texttt{----------}\\
9 & 0 & 1 & 1 & 1 & 0 & 1024 & 0 & \texttt{----------}\\
\multicolumn{9}{@{}l@{}}{\rule{0pt}{2.6ex}\textbf{T10.056}\enspace \texttt{0123456645}\enspace $n{=}10$, $|H|{=}3$, len 7, br 2, $|\mathrm{Aut}|{=}2$, orb 9/8, hist-orb 2/1; M 36864, res 5376, rev 15170: ok}\\*
0 & -- & 0 & 3 & 1 & 974 & 2 & 48 & \texttt{-+++??????}\\*
1 & 0 & 1 & 3 & 1 & 924 & 4 & 96 & \texttt{--++??????}\\
2 & 1 & 2 & 3 & 1 & 824 & 8 & 192 & \texttt{---+??????}\\
3 & 2 & 3 & 3 & 1 & 624 & 16 & 384 & \texttt{----??????}\\
4 & 3 & 4 & 2 & 1 & 640 & 128 & 256 & \texttt{-----+??--}\\
5 & 4 & 5 & 2 & 1 & 256 & 256 & 512 & \texttt{------??--}\\
6 & 5 & 6 & 1 & 2 & 0 & 1024 & 0 & \texttt{----------}\\
7 & 5 & 6 & 1 & 2 & 0 & 1024 & 0 & \texttt{----------}\\
8 & 3 & 4 & 1 & 1 & 512 & 512 & 0 & \texttt{---------+}\\
9 & 8 & 5 & 1 & 1 & 0 & 1024 & 0 & \texttt{----------}\\
\multicolumn{9}{@{}l@{}}{\rule{0pt}{2.6ex}\textbf{T10.057}\enspace \texttt{0123456644}\enspace $n{=}10$, $|H|{=}4$, len 7, br 2, $|\mathrm{Aut}|{=}4$, orb 8/6, hist-orb 2/0; M 35840, res 10056, rev 13166: ok}\\*
0 & -- & 0 & 4 & 1 & 922 & 2 & 100 & \texttt{-+++??????}\\*
1 & 0 & 1 & 4 & 1 & 820 & 4 & 200 & \texttt{--++??????}\\
2 & 1 & 2 & 4 & 1 & 616 & 8 & 400 & \texttt{---+??????}\\
3 & 2 & 3 & 4 & 1 & 208 & 16 & 800 & \texttt{----??????}\\
4 & 3 & 4 & 2 & 1 & 640 & 128 & 256 & \texttt{-----+??--}\\
5 & 4 & 5 & 2 & 1 & 256 & 256 & 512 & \texttt{------??--}\\
6 & 5 & 6 & 1 & 2 & 0 & 1024 & 0 & \texttt{----------}\\
7 & 5 & 6 & 1 & 2 & 0 & 1024 & 0 & \texttt{----------}\\
8 & 3 & 4 & 1 & 2 & 0 & 1024 & 0 & \texttt{----------}\\
9 & 3 & 4 & 1 & 2 & 0 & 1024 & 0 & \texttt{----------}\\
\multicolumn{9}{@{}l@{}}{\rule{0pt}{2.6ex}\textbf{T10.058}\enspace \texttt{0123456643}\enspace $n{=}10$, $|H|{=}4$, len 7, br 3, $|\mathrm{Aut}|{=}2$, orb 9/8, hist-orb 3/2; M 34816, res 11468, rev 11706: ok}\\*
0 & -- & 0 & 4 & 1 & 858 & 2 & 164 & \texttt{-++???????}\\*
1 & 0 & 1 & 4 & 1 & 692 & 4 & 328 & \texttt{--+???????}\\
2 & 1 & 2 & 4 & 1 & 360 & 8 & 656 & \texttt{---???????}\\
3 & 2 & 3 & 3 & 1 & 416 & 32 & 576 & \texttt{----?????-}\\
4 & 3 & 4 & 2 & 1 & 640 & 128 & 256 & \texttt{-----+??--}\\
5 & 4 & 5 & 2 & 1 & 256 & 256 & 512 & \texttt{------??--}\\
6 & 5 & 6 & 1 & 2 & 0 & 1024 & 0 & \texttt{----------}\\
7 & 5 & 6 & 1 & 2 & 0 & 1024 & 0 & \texttt{----------}\\
8 & 3 & 4 & 1 & 1 & 0 & 1024 & 0 & \texttt{----------}\\
9 & 2 & 3 & 1 & 1 & 0 & 1024 & 0 & \texttt{----------}\\
\multicolumn{9}{@{}l@{}}{\rule{0pt}{2.6ex}\textbf{T10.059}\enspace \texttt{0123456642}\enspace $n{=}10$, $|H|{=}4$, len 7, br 3, $|\mathrm{Aut}|{=}2$, orb 9/8, hist-orb 3/2; M 33792, res 11952, rev 10874: ok}\\*
0 & -- & 0 & 4 & 1 & 730 & 2 & 292 & \texttt{-+????????}\\*
1 & 0 & 1 & 4 & 1 & 436 & 4 & 584 & \texttt{--????????}\\
2 & 1 & 2 & 3 & 1 & 720 & 16 & 288 & \texttt{---+?????-}\\
3 & 2 & 3 & 3 & 1 & 416 & 32 & 576 & \texttt{----?????-}\\
4 & 3 & 4 & 2 & 1 & 640 & 128 & 256 & \texttt{-----+??--}\\
5 & 4 & 5 & 2 & 1 & 256 & 256 & 512 & \texttt{------??--}\\
6 & 5 & 6 & 1 & 2 & 0 & 1024 & 0 & \texttt{----------}\\
7 & 5 & 6 & 1 & 2 & 0 & 1024 & 0 & \texttt{----------}\\
8 & 3 & 4 & 1 & 1 & 0 & 1024 & 0 & \texttt{----------}\\
9 & 1 & 2 & 1 & 1 & 0 & 1024 & 0 & \texttt{----------}\\
\multicolumn{9}{@{}l@{}}{\rule{0pt}{2.6ex}\textbf{T10.060}\enspace \texttt{0123456641}\enspace $n{=}10$, $|H|{=}4$, len 7, br 3, $|\mathrm{Aut}|{=}2$, orb 9/8, hist-orb 3/2; M 32768, res 10812, rev 10790: ok}\\*
0 & -- & 0 & 4 & 1 & 474 & 2 & 548 & \texttt{-?????????}\\*
1 & 0 & 1 & 3 & 1 & 872 & 8 & 144 & \texttt{--++?????-}\\
2 & 1 & 2 & 3 & 1 & 720 & 16 & 288 & \texttt{---+?????-}\\
3 & 2 & 3 & 3 & 1 & 416 & 32 & 576 & \texttt{----?????-}\\
4 & 3 & 4 & 2 & 1 & 640 & 128 & 256 & \texttt{-----+??--}\\
5 & 4 & 5 & 2 & 1 & 256 & 256 & 512 & \texttt{------??--}\\
6 & 5 & 6 & 1 & 2 & 0 & 1024 & 0 & \texttt{----------}\\
7 & 5 & 6 & 1 & 2 & 0 & 1024 & 0 & \texttt{----------}\\
8 & 3 & 4 & 1 & 1 & 0 & 1024 & 0 & \texttt{----------}\\
9 & 0 & 1 & 1 & 1 & 0 & 1024 & 0 & \texttt{----------}\\
\multicolumn{9}{@{}l@{}}{\rule{0pt}{2.6ex}\textbf{T10.061}\enspace \texttt{0123456634}\enspace $n{=}10$, $|H|{=}3$, len 7, br 2, $|\mathrm{Aut}|{=}2$, orb 9/8, hist-orb 2/1; M 34816, res 5376, rev 13926: ok}\\*
0 & -- & 0 & 3 & 1 & 942 & 2 & 80 & \texttt{-++???????}\\*
1 & 0 & 1 & 3 & 1 & 860 & 4 & 160 & \texttt{--+???????}\\
2 & 1 & 2 & 3 & 1 & 696 & 8 & 320 & \texttt{---???????}\\
3 & 2 & 3 & 2 & 1 & 832 & 64 & 128 & \texttt{----++??--}\\
4 & 3 & 4 & 2 & 1 & 640 & 128 & 256 & \texttt{-----+??--}\\
5 & 4 & 5 & 2 & 1 & 256 & 256 & 512 & \texttt{------??--}\\
6 & 5 & 6 & 1 & 2 & 0 & 1024 & 0 & \texttt{----------}\\
7 & 5 & 6 & 1 & 2 & 0 & 1024 & 0 & \texttt{----------}\\
8 & 2 & 3 & 1 & 1 & 512 & 512 & 0 & \texttt{---------+}\\
9 & 8 & 4 & 1 & 1 & 0 & 1024 & 0 & \texttt{----------}\\
\multicolumn{9}{@{}l@{}}{\rule{0pt}{2.6ex}\textbf{T10.062}\enspace \texttt{0123456633}\enspace $n{=}10$, $|H|{=}4$, len 7, br 2, $|\mathrm{Aut}|{=}4$, orb 8/6, hist-orb 2/0; M 33792, res 11060, rev 11310: ok}\\*
0 & -- & 0 & 4 & 1 & 826 & 2 & 196 & \texttt{-++???????}\\*
1 & 0 & 1 & 4 & 1 & 628 & 4 & 392 & \texttt{--+???????}\\
2 & 1 & 2 & 4 & 1 & 232 & 8 & 784 & \texttt{---???????}\\
3 & 2 & 3 & 2 & 1 & 832 & 64 & 128 & \texttt{----++??--}\\
4 & 3 & 4 & 2 & 1 & 640 & 128 & 256 & \texttt{-----+??--}\\
5 & 4 & 5 & 2 & 1 & 256 & 256 & 512 & \texttt{------??--}\\
6 & 5 & 6 & 1 & 2 & 0 & 1024 & 0 & \texttt{----------}\\
7 & 5 & 6 & 1 & 2 & 0 & 1024 & 0 & \texttt{----------}\\
8 & 2 & 3 & 1 & 2 & 0 & 1024 & 0 & \texttt{----------}\\
9 & 2 & 3 & 1 & 2 & 0 & 1024 & 0 & \texttt{----------}\\
\multicolumn{9}{@{}l@{}}{\rule{0pt}{2.6ex}\textbf{T10.063}\enspace \texttt{0123456632}\enspace $n{=}10$, $|H|{=}4$, len 7, br 3, $|\mathrm{Aut}|{=}2$, orb 9/8, hist-orb 3/2; M 32768, res 12088, rev 10162: ok}\\*
0 & -- & 0 & 4 & 1 & 698 & 2 & 324 & \texttt{-+????????}\\*
1 & 0 & 1 & 4 & 1 & 372 & 4 & 648 & \texttt{--????????}\\
2 & 1 & 2 & 3 & 1 & 464 & 16 & 544 & \texttt{---??????-}\\
3 & 2 & 3 & 2 & 1 & 832 & 64 & 128 & \texttt{----++??--}\\
4 & 3 & 4 & 2 & 1 & 640 & 128 & 256 & \texttt{-----+??--}\\
5 & 4 & 5 & 2 & 1 & 256 & 256 & 512 & \texttt{------??--}\\
6 & 5 & 6 & 1 & 2 & 0 & 1024 & 0 & \texttt{----------}\\
7 & 5 & 6 & 1 & 2 & 0 & 1024 & 0 & \texttt{----------}\\
8 & 2 & 3 & 1 & 1 & 0 & 1024 & 0 & \texttt{----------}\\
9 & 1 & 2 & 1 & 1 & 0 & 1024 & 0 & \texttt{----------}\\
\multicolumn{9}{@{}l@{}}{\rule{0pt}{2.6ex}\textbf{T10.064}\enspace \texttt{0123456631}\enspace $n{=}10$, $|H|{=}4$, len 7, br 3, $|\mathrm{Aut}|{=}2$, orb 9/8, hist-orb 3/2; M 31744, res 11300, rev 9926: ok}\\*
0 & -- & 0 & 4 & 1 & 442 & 2 & 580 & \texttt{-?????????}\\*
1 & 0 & 1 & 3 & 1 & 744 & 8 & 272 & \texttt{--+??????-}\\
2 & 1 & 2 & 3 & 1 & 464 & 16 & 544 & \texttt{---??????-}\\
3 & 2 & 3 & 2 & 1 & 832 & 64 & 128 & \texttt{----++??--}\\
4 & 3 & 4 & 2 & 1 & 640 & 128 & 256 & \texttt{-----+??--}\\
5 & 4 & 5 & 2 & 1 & 256 & 256 & 512 & \texttt{------??--}\\
6 & 5 & 6 & 1 & 2 & 0 & 1024 & 0 & \texttt{----------}\\
7 & 5 & 6 & 1 & 2 & 0 & 1024 & 0 & \texttt{----------}\\
8 & 2 & 3 & 1 & 1 & 0 & 1024 & 0 & \texttt{----------}\\
9 & 0 & 1 & 1 & 1 & 0 & 1024 & 0 & \texttt{----------}\\
\multicolumn{9}{@{}l@{}}{\rule{0pt}{2.6ex}\textbf{T10.065}\enspace \texttt{0123456623}\enspace $n{=}10$, $|H|{=}3$, len 7, br 2, $|\mathrm{Aut}|{=}2$, orb 9/8, hist-orb 2/1; M 32768, res 5184, rev 12722: ok}\\*
0 & -- & 0 & 3 & 1 & 878 & 2 & 144 & \texttt{-+????????}\\*
1 & 0 & 1 & 3 & 1 & 732 & 4 & 288 & \texttt{--????????}\\
2 & 1 & 2 & 2 & 1 & 928 & 32 & 64 & \texttt{---+++??--}\\
3 & 2 & 3 & 2 & 1 & 832 & 64 & 128 & \texttt{----++??--}\\
4 & 3 & 4 & 2 & 1 & 640 & 128 & 256 & \texttt{-----+??--}\\
5 & 4 & 5 & 2 & 1 & 256 & 256 & 512 & \texttt{------??--}\\
6 & 5 & 6 & 1 & 2 & 0 & 1024 & 0 & \texttt{----------}\\
7 & 5 & 6 & 1 & 2 & 0 & 1024 & 0 & \texttt{----------}\\
8 & 1 & 2 & 1 & 1 & 512 & 512 & 0 & \texttt{---------+}\\
9 & 8 & 3 & 1 & 1 & 0 & 1024 & 0 & \texttt{----------}\\
\multicolumn{9}{@{}l@{}}{\rule{0pt}{2.6ex}\textbf{T10.066}\enspace \texttt{0123456622}\enspace $n{=}10$, $|H|{=}4$, len 7, br 2, $|\mathrm{Aut}|{=}4$, orb 8/6, hist-orb 2/0; M 31744, res 11040, rev 9958: ok}\\*
0 & -- & 0 & 4 & 1 & 634 & 2 & 388 & \texttt{-+????????}\\*
1 & 0 & 1 & 4 & 1 & 244 & 4 & 776 & \texttt{--????????}\\
2 & 1 & 2 & 2 & 1 & 928 & 32 & 64 & \texttt{---+++??--}\\
3 & 2 & 3 & 2 & 1 & 832 & 64 & 128 & \texttt{----++??--}\\
4 & 3 & 4 & 2 & 1 & 640 & 128 & 256 & \texttt{-----+??--}\\
5 & 4 & 5 & 2 & 1 & 256 & 256 & 512 & \texttt{------??--}\\
6 & 5 & 6 & 1 & 2 & 0 & 1024 & 0 & \texttt{----------}\\
7 & 5 & 6 & 1 & 2 & 0 & 1024 & 0 & \texttt{----------}\\
8 & 1 & 2 & 1 & 2 & 0 & 1024 & 0 & \texttt{----------}\\
9 & 1 & 2 & 1 & 2 & 0 & 1024 & 0 & \texttt{----------}\\
\multicolumn{9}{@{}l@{}}{\rule{0pt}{2.6ex}\textbf{T10.067}\enspace \texttt{0123456621}\enspace $n{=}10$, $|H|{=}4$, len 7, br 3, $|\mathrm{Aut}|{=}2$, orb 9/8, hist-orb 3/2; M 30720, res 10956, rev 9406: ok}\\*
0 & -- & 0 & 4 & 1 & 378 & 2 & 644 & \texttt{-?????????}\\*
1 & 0 & 1 & 3 & 1 & 488 & 8 & 528 & \texttt{--???????-}\\
2 & 1 & 2 & 2 & 1 & 928 & 32 & 64 & \texttt{---+++??--}\\
3 & 2 & 3 & 2 & 1 & 832 & 64 & 128 & \texttt{----++??--}\\
4 & 3 & 4 & 2 & 1 & 640 & 128 & 256 & \texttt{-----+??--}\\
5 & 4 & 5 & 2 & 1 & 256 & 256 & 512 & \texttt{------??--}\\
6 & 5 & 6 & 1 & 2 & 0 & 1024 & 0 & \texttt{----------}\\
7 & 5 & 6 & 1 & 2 & 0 & 1024 & 0 & \texttt{----------}\\
8 & 1 & 2 & 1 & 1 & 0 & 1024 & 0 & \texttt{----------}\\
9 & 0 & 1 & 1 & 1 & 0 & 1024 & 0 & \texttt{----------}\\
\multicolumn{9}{@{}l@{}}{\rule{0pt}{2.6ex}\textbf{T10.068}\enspace \texttt{0123456612}\enspace $n{=}10$, $|H|{=}3$, len 7, br 2, $|\mathrm{Aut}|{=}2$, orb 9/8, hist-orb 2/1; M 30720, res 4352, rev 11770: ok}\\*
0 & -- & 0 & 3 & 1 & 750 & 2 & 272 & \texttt{-?????????}\\*
1 & 0 & 1 & 2 & 1 & 976 & 16 & 32 & \texttt{--++++??--}\\
2 & 1 & 2 & 2 & 1 & 928 & 32 & 64 & \texttt{---+++??--}\\
3 & 2 & 3 & 2 & 1 & 832 & 64 & 128 & \texttt{----++??--}\\
4 & 3 & 4 & 2 & 1 & 640 & 128 & 256 & \texttt{-----+??--}\\
5 & 4 & 5 & 2 & 1 & 256 & 256 & 512 & \texttt{------??--}\\
6 & 5 & 6 & 1 & 2 & 0 & 1024 & 0 & \texttt{----------}\\
7 & 5 & 6 & 1 & 2 & 0 & 1024 & 0 & \texttt{----------}\\
8 & 0 & 1 & 1 & 1 & 512 & 512 & 0 & \texttt{---------+}\\
9 & 8 & 2 & 1 & 1 & 0 & 1024 & 0 & \texttt{----------}\\
\multicolumn{9}{@{}l@{}}{\rule{0pt}{2.6ex}\textbf{T10.069}\enspace \texttt{0123456611}\enspace $n{=}10$, $|H|{=}4$, len 7, br 2, $|\mathrm{Aut}|{=}4$, orb 8/6, hist-orb 2/0; M 29696, res 8852, rev 9530: ok}\\*
0 & -- & 0 & 4 & 1 & 250 & 2 & 772 & \texttt{-?????????}\\*
1 & 0 & 1 & 2 & 1 & 976 & 16 & 32 & \texttt{--++++??--}\\
2 & 1 & 2 & 2 & 1 & 928 & 32 & 64 & \texttt{---+++??--}\\
3 & 2 & 3 & 2 & 1 & 832 & 64 & 128 & \texttt{----++??--}\\
4 & 3 & 4 & 2 & 1 & 640 & 128 & 256 & \texttt{-----+??--}\\
5 & 4 & 5 & 2 & 1 & 256 & 256 & 512 & \texttt{------??--}\\
6 & 5 & 6 & 1 & 2 & 0 & 1024 & 0 & \texttt{----------}\\
7 & 5 & 6 & 1 & 2 & 0 & 1024 & 0 & \texttt{----------}\\
8 & 0 & 1 & 1 & 2 & 0 & 1024 & 0 & \texttt{----------}\\
9 & 0 & 1 & 1 & 2 & 0 & 1024 & 0 & \texttt{----------}\\
\multicolumn{9}{@{}l@{}}{\rule{0pt}{2.6ex}\textbf{T10.070}\enspace \texttt{0123456565}\enspace $n{=}10$, $|H|{=}3$, len 7, br 1, $|\mathrm{Aut}|{=}2$, orb 8/6, hist-orb 2/1; M 37888, res 6820, rev 15758: ok}\\*
0 & -- & 0 & 3 & 1 & 978 & 2 & 44 & \texttt{-++++?????}\\*
1 & 0 & 1 & 3 & 1 & 932 & 4 & 88 & \texttt{--+++?????}\\
2 & 1 & 2 & 3 & 1 & 840 & 8 & 176 & \texttt{---++?????}\\
3 & 2 & 3 & 3 & 1 & 656 & 16 & 352 & \texttt{----+?????}\\
4 & 3 & 4 & 3 & 1 & 288 & 32 & 704 & \texttt{-----?????}\\
5 & 4 & 5 & 1 & 2 & 512 & 512 & 0 & \texttt{------+---}\\
6 & 5 & 6 & 1 & 2 & 0 & 1024 & 0 & \texttt{----------}\\
7 & 4 & 5 & 1 & 2 & 512 & 512 & 0 & \texttt{--------+-}\\
8 & 7 & 6 & 1 & 2 & 0 & 1024 & 0 & \texttt{----------}\\
9 & 4 & 5 & 1 & 1 & 0 & 1024 & 0 & \texttt{----------}\\
\multicolumn{9}{@{}l@{}}{\rule{0pt}{2.6ex}\textbf{T10.071}\enspace \texttt{0123456564}\enspace $n{=}10$, $|H|{=}3$, len 7, br 2, $|\mathrm{Aut}|{=}2$, orb 8/6, hist-orb 2/1; M 36864, res 8016, rev 14386: ok}\\*
0 & -- & 0 & 3 & 1 & 946 & 2 & 76 & \texttt{-+++??????}\\*
1 & 0 & 1 & 3 & 1 & 868 & 4 & 152 & \texttt{--++??????}\\
2 & 1 & 2 & 3 & 1 & 712 & 8 & 304 & \texttt{---+??????}\\
3 & 2 & 3 & 3 & 1 & 400 & 16 & 608 & \texttt{----??????}\\
4 & 3 & 4 & 2 & 1 & 576 & 64 & 384 & \texttt{-----????-}\\
5 & 4 & 5 & 1 & 2 & 512 & 512 & 0 & \texttt{------+---}\\
6 & 5 & 6 & 1 & 2 & 0 & 1024 & 0 & \texttt{----------}\\
7 & 4 & 5 & 1 & 2 & 512 & 512 & 0 & \texttt{--------+-}\\
8 & 7 & 6 & 1 & 2 & 0 & 1024 & 0 & \texttt{----------}\\
9 & 3 & 4 & 1 & 1 & 0 & 1024 & 0 & \texttt{----------}\\
\multicolumn{9}{@{}l@{}}{\rule{0pt}{2.6ex}\textbf{T10.072}\enspace \texttt{0123456563}\enspace $n{=}10$, $|H|{=}3$, len 7, br 2, $|\mathrm{Aut}|{=}2$, orb 8/6, hist-orb 2/1; M 35840, res 8828, rev 13354: ok}\\*
0 & -- & 0 & 3 & 1 & 882 & 2 & 140 & \texttt{-++???????}\\*
1 & 0 & 1 & 3 & 1 & 740 & 4 & 280 & \texttt{--+???????}\\
2 & 1 & 2 & 3 & 1 & 456 & 8 & 560 & \texttt{---???????}\\
3 & 2 & 3 & 2 & 1 & 800 & 32 & 192 & \texttt{----+????-}\\
4 & 3 & 4 & 2 & 1 & 576 & 64 & 384 & \texttt{-----????-}\\
5 & 4 & 5 & 1 & 2 & 512 & 512 & 0 & \texttt{------+---}\\
6 & 5 & 6 & 1 & 2 & 0 & 1024 & 0 & \texttt{----------}\\
7 & 4 & 5 & 1 & 2 & 512 & 512 & 0 & \texttt{--------+-}\\
8 & 7 & 6 & 1 & 2 & 0 & 1024 & 0 & \texttt{----------}\\
9 & 2 & 3 & 1 & 1 & 0 & 1024 & 0 & \texttt{----------}\\
\multicolumn{9}{@{}l@{}}{\rule{0pt}{2.6ex}\textbf{T10.073}\enspace \texttt{0123456562}\enspace $n{=}10$, $|H|{=}3$, len 7, br 2, $|\mathrm{Aut}|{=}2$, orb 8/6, hist-orb 2/1; M 34816, res 8904, rev 12706: ok}\\*
0 & -- & 0 & 3 & 1 & 754 & 2 & 268 & \texttt{-+????????}\\*
1 & 0 & 1 & 3 & 1 & 484 & 4 & 536 & \texttt{--????????}\\
2 & 1 & 2 & 2 & 1 & 912 & 16 & 96 & \texttt{---++????-}\\
3 & 2 & 3 & 2 & 1 & 800 & 32 & 192 & \texttt{----+????-}\\
4 & 3 & 4 & 2 & 1 & 576 & 64 & 384 & \texttt{-----????-}\\
5 & 4 & 5 & 1 & 2 & 512 & 512 & 0 & \texttt{------+---}\\
6 & 5 & 6 & 1 & 2 & 0 & 1024 & 0 & \texttt{----------}\\
7 & 4 & 5 & 1 & 2 & 512 & 512 & 0 & \texttt{--------+-}\\
8 & 7 & 6 & 1 & 2 & 0 & 1024 & 0 & \texttt{----------}\\
9 & 1 & 2 & 1 & 1 & 0 & 1024 & 0 & \texttt{----------}\\
\multicolumn{9}{@{}l@{}}{\rule{0pt}{2.6ex}\textbf{T10.074}\enspace \texttt{0123456561}\enspace $n{=}10$, $|H|{=}3$, len 7, br 2, $|\mathrm{Aut}|{=}2$, orb 8/6, hist-orb 2/1; M 33792, res 7500, rev 12654: ok}\\*
0 & -- & 0 & 3 & 1 & 498 & 2 & 524 & \texttt{-?????????}\\*
1 & 0 & 1 & 2 & 1 & 968 & 8 & 48 & \texttt{--+++????-}\\
2 & 1 & 2 & 2 & 1 & 912 & 16 & 96 & \texttt{---++????-}\\
3 & 2 & 3 & 2 & 1 & 800 & 32 & 192 & \texttt{----+????-}\\
4 & 3 & 4 & 2 & 1 & 576 & 64 & 384 & \texttt{-----????-}\\
5 & 4 & 5 & 1 & 2 & 512 & 512 & 0 & \texttt{------+---}\\
6 & 5 & 6 & 1 & 2 & 0 & 1024 & 0 & \texttt{----------}\\
7 & 4 & 5 & 1 & 2 & 512 & 512 & 0 & \texttt{--------+-}\\
8 & 7 & 6 & 1 & 2 & 0 & 1024 & 0 & \texttt{----------}\\
9 & 0 & 1 & 1 & 1 & 0 & 1024 & 0 & \texttt{----------}\\
\multicolumn{9}{@{}l@{}}{\rule{0pt}{2.6ex}\textbf{T10.075}\enspace \texttt{0123456555}\enspace $n{=}10$, $|H|{=}4$, len 7, br 1, $|\mathrm{Aut}|{=}6$, orb 8/7, hist-orb 2/1; M 36864, res 8680, rev 15574: ok}\\*
0 & -- & 0 & 4 & 1 & 966 & 2 & 56 & \texttt{-++++?????}\\*
1 & 0 & 1 & 4 & 1 & 908 & 4 & 112 & \texttt{--+++?????}\\
2 & 1 & 2 & 4 & 1 & 792 & 8 & 224 & \texttt{---++?????}\\
3 & 2 & 3 & 4 & 1 & 560 & 16 & 448 & \texttt{----+?????}\\
4 & 3 & 4 & 4 & 1 & 96 & 32 & 896 & \texttt{-----?????}\\
5 & 4 & 5 & 1 & 1 & 512 & 512 & 0 & \texttt{------+---}\\
6 & 5 & 6 & 1 & 1 & 0 & 1024 & 0 & \texttt{----------}\\
7 & 4 & 5 & 1 & 3 & 0 & 1024 & 0 & \texttt{----------}\\
8 & 4 & 5 & 1 & 3 & 0 & 1024 & 0 & \texttt{----------}\\
9 & 4 & 5 & 1 & 3 & 0 & 1024 & 0 & \texttt{----------}\\
\multicolumn{9}{@{}l@{}}{\rule{0pt}{2.6ex}\textbf{T10.076}\enspace \texttt{0123456554}\enspace $n{=}10$, $|H|{=}4$, len 7, br 2, $|\mathrm{Aut}|{=}2$, orb 9/8, hist-orb 3/2; M 35840, res 10272, rev 13706: ok}\\*
0 & -- & 0 & 4 & 1 & 934 & 2 & 88 & \texttt{-+++??????}\\*
1 & 0 & 1 & 4 & 1 & 844 & 4 & 176 & \texttt{--++??????}\\
2 & 1 & 2 & 4 & 1 & 664 & 8 & 352 & \texttt{---+??????}\\
3 & 2 & 3 & 4 & 1 & 304 & 16 & 704 & \texttt{----??????}\\
4 & 3 & 4 & 3 & 1 & 192 & 64 & 768 & \texttt{-----????-}\\
5 & 4 & 5 & 1 & 1 & 512 & 512 & 0 & \texttt{------+---}\\
6 & 5 & 6 & 1 & 1 & 0 & 1024 & 0 & \texttt{----------}\\
7 & 4 & 5 & 1 & 2 & 0 & 1024 & 0 & \texttt{----------}\\
8 & 4 & 5 & 1 & 2 & 0 & 1024 & 0 & \texttt{----------}\\
9 & 3 & 4 & 1 & 1 & 0 & 1024 & 0 & \texttt{----------}\\
\multicolumn{9}{@{}l@{}}{\rule{0pt}{2.6ex}\textbf{T10.077}\enspace \texttt{0123456553}\enspace $n{=}10$, $|H|{=}4$, len 7, br 2, $|\mathrm{Aut}|{=}2$, orb 9/8, hist-orb 3/2; M 34816, res 11384, rev 12450: ok}\\*
0 & -- & 0 & 4 & 1 & 870 & 2 & 152 & \texttt{-++???????}\\*
1 & 0 & 1 & 4 & 1 & 716 & 4 & 304 & \texttt{--+???????}\\
2 & 1 & 2 & 4 & 1 & 408 & 8 & 608 & \texttt{---???????}\\
3 & 2 & 3 & 3 & 1 & 608 & 32 & 384 & \texttt{----+????-}\\
4 & 3 & 4 & 3 & 1 & 192 & 64 & 768 & \texttt{-----????-}\\
5 & 4 & 5 & 1 & 1 & 512 & 512 & 0 & \texttt{------+---}\\
6 & 5 & 6 & 1 & 1 & 0 & 1024 & 0 & \texttt{----------}\\
7 & 4 & 5 & 1 & 2 & 0 & 1024 & 0 & \texttt{----------}\\
8 & 4 & 5 & 1 & 2 & 0 & 1024 & 0 & \texttt{----------}\\
9 & 2 & 3 & 1 & 1 & 0 & 1024 & 0 & \texttt{----------}\\
\multicolumn{9}{@{}l@{}}{\rule{0pt}{2.6ex}\textbf{T10.078}\enspace \texttt{0123456552}\enspace $n{=}10$, $|H|{=}4$, len 7, br 2, $|\mathrm{Aut}|{=}2$, orb 9/8, hist-orb 3/2; M 33792, res 11664, rev 11738: ok}\\*
0 & -- & 0 & 4 & 1 & 742 & 2 & 280 & \texttt{-+????????}\\*
1 & 0 & 1 & 4 & 1 & 460 & 4 & 560 & \texttt{--????????}\\
2 & 1 & 2 & 3 & 1 & 816 & 16 & 192 & \texttt{---++????-}\\
3 & 2 & 3 & 3 & 1 & 608 & 32 & 384 & \texttt{----+????-}\\
4 & 3 & 4 & 3 & 1 & 192 & 64 & 768 & \texttt{-----????-}\\
5 & 4 & 5 & 1 & 1 & 512 & 512 & 0 & \texttt{------+---}\\
6 & 5 & 6 & 1 & 1 & 0 & 1024 & 0 & \texttt{----------}\\
7 & 4 & 5 & 1 & 2 & 0 & 1024 & 0 & \texttt{----------}\\
8 & 4 & 5 & 1 & 2 & 0 & 1024 & 0 & \texttt{----------}\\
9 & 1 & 2 & 1 & 1 & 0 & 1024 & 0 & \texttt{----------}\\
\multicolumn{9}{@{}l@{}}{\rule{0pt}{2.6ex}\textbf{T10.079}\enspace \texttt{0123456551}\enspace $n{=}10$, $|H|{=}4$, len 7, br 2, $|\mathrm{Aut}|{=}2$, orb 9/8, hist-orb 3/2; M 32768, res 10392, rev 11750: ok}\\*
0 & -- & 0 & 4 & 1 & 486 & 2 & 536 & \texttt{-?????????}\\*
1 & 0 & 1 & 3 & 1 & 920 & 8 & 96 & \texttt{--+++????-}\\
2 & 1 & 2 & 3 & 1 & 816 & 16 & 192 & \texttt{---++????-}\\
3 & 2 & 3 & 3 & 1 & 608 & 32 & 384 & \texttt{----+????-}\\
4 & 3 & 4 & 3 & 1 & 192 & 64 & 768 & \texttt{-----????-}\\
5 & 4 & 5 & 1 & 1 & 512 & 512 & 0 & \texttt{------+---}\\
6 & 5 & 6 & 1 & 1 & 0 & 1024 & 0 & \texttt{----------}\\
7 & 4 & 5 & 1 & 2 & 0 & 1024 & 0 & \texttt{----------}\\
8 & 4 & 5 & 1 & 2 & 0 & 1024 & 0 & \texttt{----------}\\
9 & 0 & 1 & 1 & 1 & 0 & 1024 & 0 & \texttt{----------}\\
\multicolumn{9}{@{}l@{}}{\rule{0pt}{2.6ex}\textbf{T10.080}\enspace \texttt{0123456545}\enspace $n{=}10$, $|H|{=}3$, len 7, br 2, $|\mathrm{Aut}|{=}1$, orb 10/10, hist-orb 3/3; M 35840, res 6456, rev 14782: ok}\\*
0 & -- & 0 & 3 & 1 & 962 & 2 & 60 & \texttt{-+++??????}\\*
1 & 0 & 1 & 3 & 1 & 900 & 4 & 120 & \texttt{--++??????}\\
2 & 1 & 2 & 3 & 1 & 776 & 8 & 240 & \texttt{---+??????}\\
3 & 2 & 3 & 3 & 1 & 528 & 16 & 480 & \texttt{----??????}\\
4 & 3 & 4 & 2 & 1 & 384 & 128 & 512 & \texttt{-----???--}\\
5 & 4 & 5 & 1 & 1 & 512 & 512 & 0 & \texttt{------+---}\\
6 & 5 & 6 & 1 & 1 & 0 & 1024 & 0 & \texttt{----------}\\
7 & 4 & 5 & 1 & 1 & 0 & 1024 & 0 & \texttt{----------}\\
8 & 3 & 4 & 1 & 1 & 512 & 512 & 0 & \texttt{---------+}\\
9 & 8 & 5 & 1 & 1 & 0 & 1024 & 0 & \texttt{----------}\\
\multicolumn{9}{@{}l@{}}{\rule{0pt}{2.6ex}\textbf{T10.081}\enspace \texttt{0123456544}\enspace $n{=}10$, $|H|{=}4$, len 7, br 2, $|\mathrm{Aut}|{=}2$, orb 9/8, hist-orb 3/2; M 34816, res 10416, rev 13018: ok}\\*
0 & -- & 0 & 4 & 1 & 918 & 2 & 104 & \texttt{-+++??????}\\*
1 & 0 & 1 & 4 & 1 & 812 & 4 & 208 & \texttt{--++??????}\\
2 & 1 & 2 & 4 & 1 & 600 & 8 & 416 & \texttt{---+??????}\\
3 & 2 & 3 & 4 & 1 & 176 & 16 & 832 & \texttt{----??????}\\
4 & 3 & 4 & 2 & 1 & 384 & 128 & 512 & \texttt{-----???--}\\
5 & 4 & 5 & 1 & 1 & 512 & 512 & 0 & \texttt{------+---}\\
6 & 5 & 6 & 1 & 1 & 0 & 1024 & 0 & \texttt{----------}\\
7 & 4 & 5 & 1 & 1 & 0 & 1024 & 0 & \texttt{----------}\\
8 & 3 & 4 & 1 & 2 & 0 & 1024 & 0 & \texttt{----------}\\
9 & 3 & 4 & 1 & 2 & 0 & 1024 & 0 & \texttt{----------}\\
\multicolumn{9}{@{}l@{}}{\rule{0pt}{2.6ex}\textbf{T10.082}\enspace \texttt{0123456543}\enspace $n{=}10$, $|H|{=}4$, len 7, br 3, $|\mathrm{Aut}|{=}1$, orb 10/10, hist-orb 4/4; M 33792, res 11928, rev 11414: ok}\\*
0 & -- & 0 & 4 & 1 & 854 & 2 & 168 & \texttt{-++???????}\\*
1 & 0 & 1 & 4 & 1 & 684 & 4 & 336 & \texttt{--+???????}\\
2 & 1 & 2 & 4 & 1 & 344 & 8 & 672 & \texttt{---???????}\\
3 & 2 & 3 & 3 & 1 & 352 & 32 & 640 & \texttt{----?????-}\\
4 & 3 & 4 & 2 & 1 & 384 & 128 & 512 & \texttt{-----???--}\\
5 & 4 & 5 & 1 & 1 & 512 & 512 & 0 & \texttt{------+---}\\
6 & 5 & 6 & 1 & 1 & 0 & 1024 & 0 & \texttt{----------}\\
7 & 4 & 5 & 1 & 1 & 0 & 1024 & 0 & \texttt{----------}\\
8 & 3 & 4 & 1 & 1 & 0 & 1024 & 0 & \texttt{----------}\\
9 & 2 & 3 & 1 & 1 & 0 & 1024 & 0 & \texttt{----------}\\
\multicolumn{9}{@{}l@{}}{\rule{0pt}{2.6ex}\textbf{T10.083}\enspace \texttt{0123456542}\enspace $n{=}10$, $|H|{=}4$, len 7, br 3, $|\mathrm{Aut}|{=}1$, orb 10/10, hist-orb 4/4; M 32768, res 12480, rev 10534: ok}\\*
0 & -- & 0 & 4 & 1 & 726 & 2 & 296 & \texttt{-+????????}\\*
1 & 0 & 1 & 4 & 1 & 428 & 4 & 592 & \texttt{--????????}\\
2 & 1 & 2 & 3 & 1 & 688 & 16 & 320 & \texttt{---+?????-}\\
3 & 2 & 3 & 3 & 1 & 352 & 32 & 640 & \texttt{----?????-}\\
4 & 3 & 4 & 2 & 1 & 384 & 128 & 512 & \texttt{-----???--}\\
5 & 4 & 5 & 1 & 1 & 512 & 512 & 0 & \texttt{------+---}\\
6 & 5 & 6 & 1 & 1 & 0 & 1024 & 0 & \texttt{----------}\\
7 & 4 & 5 & 1 & 1 & 0 & 1024 & 0 & \texttt{----------}\\
8 & 3 & 4 & 1 & 1 & 0 & 1024 & 0 & \texttt{----------}\\
9 & 1 & 2 & 1 & 1 & 0 & 1024 & 0 & \texttt{----------}\\
\multicolumn{9}{@{}l@{}}{\rule{0pt}{2.6ex}\textbf{T10.084}\enspace \texttt{0123456541}\enspace $n{=}10$, $|H|{=}4$, len 7, br 3, $|\mathrm{Aut}|{=}1$, orb 10/10, hist-orb 4/4; M 31744, res 11384, rev 10474: ok}\\*
0 & -- & 0 & 4 & 1 & 470 & 2 & 552 & \texttt{-?????????}\\*
1 & 0 & 1 & 3 & 1 & 856 & 8 & 160 & \texttt{--++?????-}\\
2 & 1 & 2 & 3 & 1 & 688 & 16 & 320 & \texttt{---+?????-}\\
3 & 2 & 3 & 3 & 1 & 352 & 32 & 640 & \texttt{----?????-}\\
4 & 3 & 4 & 2 & 1 & 384 & 128 & 512 & \texttt{-----???--}\\
5 & 4 & 5 & 1 & 1 & 512 & 512 & 0 & \texttt{------+---}\\
6 & 5 & 6 & 1 & 1 & 0 & 1024 & 0 & \texttt{----------}\\
7 & 4 & 5 & 1 & 1 & 0 & 1024 & 0 & \texttt{----------}\\
8 & 3 & 4 & 1 & 1 & 0 & 1024 & 0 & \texttt{----------}\\
9 & 0 & 1 & 1 & 1 & 0 & 1024 & 0 & \texttt{----------}\\
\multicolumn{9}{@{}l@{}}{\rule{0pt}{2.6ex}\textbf{T10.085}\enspace \texttt{0123456534}\enspace $n{=}10$, $|H|{=}3$, len 7, br 2, $|\mathrm{Aut}|{=}1$, orb 10/10, hist-orb 3/3; M 33792, res 6364, rev 13578: ok}\\*
0 & -- & 0 & 3 & 1 & 930 & 2 & 92 & \texttt{-++???????}\\*
1 & 0 & 1 & 3 & 1 & 836 & 4 & 184 & \texttt{--+???????}\\
2 & 1 & 2 & 3 & 1 & 648 & 8 & 368 & \texttt{---???????}\\
3 & 2 & 3 & 2 & 1 & 704 & 64 & 256 & \texttt{----+???--}\\
4 & 3 & 4 & 2 & 1 & 384 & 128 & 512 & \texttt{-----???--}\\
5 & 4 & 5 & 1 & 1 & 512 & 512 & 0 & \texttt{------+---}\\
6 & 5 & 6 & 1 & 1 & 0 & 1024 & 0 & \texttt{----------}\\
7 & 4 & 5 & 1 & 1 & 0 & 1024 & 0 & \texttt{----------}\\
8 & 2 & 3 & 1 & 1 & 512 & 512 & 0 & \texttt{---------+}\\
9 & 8 & 4 & 1 & 1 & 0 & 1024 & 0 & \texttt{----------}\\
\multicolumn{9}{@{}l@{}}{\rule{0pt}{2.6ex}\textbf{T10.086}\enspace \texttt{0123456533}\enspace $n{=}10$, $|H|{=}4$, len 7, br 2, $|\mathrm{Aut}|{=}2$, orb 9/8, hist-orb 3/2; M 32768, res 11656, rev 10962: ok}\\*
0 & -- & 0 & 4 & 1 & 822 & 2 & 200 & \texttt{-++???????}\\*
1 & 0 & 1 & 4 & 1 & 620 & 4 & 400 & \texttt{--+???????}\\
2 & 1 & 2 & 4 & 1 & 216 & 8 & 800 & \texttt{---???????}\\
3 & 2 & 3 & 2 & 1 & 704 & 64 & 256 & \texttt{----+???--}\\
4 & 3 & 4 & 2 & 1 & 384 & 128 & 512 & \texttt{-----???--}\\
5 & 4 & 5 & 1 & 1 & 512 & 512 & 0 & \texttt{------+---}\\
6 & 5 & 6 & 1 & 1 & 0 & 1024 & 0 & \texttt{----------}\\
7 & 4 & 5 & 1 & 1 & 0 & 1024 & 0 & \texttt{----------}\\
8 & 2 & 3 & 1 & 2 & 0 & 1024 & 0 & \texttt{----------}\\
9 & 2 & 3 & 1 & 2 & 0 & 1024 & 0 & \texttt{----------}\\
\multicolumn{9}{@{}l@{}}{\rule{0pt}{2.6ex}\textbf{T10.087}\enspace \texttt{0123456532}\enspace $n{=}10$, $|H|{=}4$, len 7, br 3, $|\mathrm{Aut}|{=}1$, orb 10/10, hist-orb 4/4; M 31744, res 12752, rev 9734: ok}\\*
0 & -- & 0 & 4 & 1 & 694 & 2 & 328 & \texttt{-+????????}\\*
1 & 0 & 1 & 4 & 1 & 364 & 4 & 656 & \texttt{--????????}\\
2 & 1 & 2 & 3 & 1 & 432 & 16 & 576 & \texttt{---??????-}\\
3 & 2 & 3 & 2 & 1 & 704 & 64 & 256 & \texttt{----+???--}\\
4 & 3 & 4 & 2 & 1 & 384 & 128 & 512 & \texttt{-----???--}\\
5 & 4 & 5 & 1 & 1 & 512 & 512 & 0 & \texttt{------+---}\\
6 & 5 & 6 & 1 & 1 & 0 & 1024 & 0 & \texttt{----------}\\
7 & 4 & 5 & 1 & 1 & 0 & 1024 & 0 & \texttt{----------}\\
8 & 2 & 3 & 1 & 1 & 0 & 1024 & 0 & \texttt{----------}\\
9 & 1 & 2 & 1 & 1 & 0 & 1024 & 0 & \texttt{----------}\\
\multicolumn{9}{@{}l@{}}{\rule{0pt}{2.6ex}\textbf{T10.088}\enspace \texttt{0123456531}\enspace $n{=}10$, $|H|{=}4$, len 7, br 3, $|\mathrm{Aut}|{=}1$, orb 10/10, hist-orb 4/4; M 30720, res 12008, rev 9506: ok}\\*
0 & -- & 0 & 4 & 1 & 438 & 2 & 584 & \texttt{-?????????}\\*
1 & 0 & 1 & 3 & 1 & 728 & 8 & 288 & \texttt{--+??????-}\\
2 & 1 & 2 & 3 & 1 & 432 & 16 & 576 & \texttt{---??????-}\\
3 & 2 & 3 & 2 & 1 & 704 & 64 & 256 & \texttt{----+???--}\\
4 & 3 & 4 & 2 & 1 & 384 & 128 & 512 & \texttt{-----???--}\\
5 & 4 & 5 & 1 & 1 & 512 & 512 & 0 & \texttt{------+---}\\
6 & 5 & 6 & 1 & 1 & 0 & 1024 & 0 & \texttt{----------}\\
7 & 4 & 5 & 1 & 1 & 0 & 1024 & 0 & \texttt{----------}\\
8 & 2 & 3 & 1 & 1 & 0 & 1024 & 0 & \texttt{----------}\\
9 & 0 & 1 & 1 & 1 & 0 & 1024 & 0 & \texttt{----------}\\
\multicolumn{9}{@{}l@{}}{\rule{0pt}{2.6ex}\textbf{T10.089}\enspace \texttt{0123456523}\enspace $n{=}10$, $|H|{=}3$, len 7, br 2, $|\mathrm{Aut}|{=}1$, orb 10/10, hist-orb 3/3; M 31744, res 6144, rev 12406: ok}\\*
0 & -- & 0 & 3 & 1 & 866 & 2 & 156 & \texttt{-+????????}\\*
1 & 0 & 1 & 3 & 1 & 708 & 4 & 312 & \texttt{--????????}\\
2 & 1 & 2 & 2 & 1 & 864 & 32 & 128 & \texttt{---++???--}\\
3 & 2 & 3 & 2 & 1 & 704 & 64 & 256 & \texttt{----+???--}\\
4 & 3 & 4 & 2 & 1 & 384 & 128 & 512 & \texttt{-----???--}\\
5 & 4 & 5 & 1 & 1 & 512 & 512 & 0 & \texttt{------+---}\\
6 & 5 & 6 & 1 & 1 & 0 & 1024 & 0 & \texttt{----------}\\
7 & 4 & 5 & 1 & 1 & 0 & 1024 & 0 & \texttt{----------}\\
8 & 1 & 2 & 1 & 1 & 512 & 512 & 0 & \texttt{---------+}\\
9 & 8 & 3 & 1 & 1 & 0 & 1024 & 0 & \texttt{----------}\\
\multicolumn{9}{@{}l@{}}{\rule{0pt}{2.6ex}\textbf{T10.090}\enspace \texttt{0123456522}\enspace $n{=}10$, $|H|{=}4$, len 7, br 2, $|\mathrm{Aut}|{=}2$, orb 9/8, hist-orb 3/2; M 30720, res 11808, rev 9546: ok}\\*
0 & -- & 0 & 4 & 1 & 630 & 2 & 392 & \texttt{-+????????}\\*
1 & 0 & 1 & 4 & 1 & 236 & 4 & 784 & \texttt{--????????}\\
2 & 1 & 2 & 2 & 1 & 864 & 32 & 128 & \texttt{---++???--}\\
3 & 2 & 3 & 2 & 1 & 704 & 64 & 256 & \texttt{----+???--}\\
4 & 3 & 4 & 2 & 1 & 384 & 128 & 512 & \texttt{-----???--}\\
5 & 4 & 5 & 1 & 1 & 512 & 512 & 0 & \texttt{------+---}\\
6 & 5 & 6 & 1 & 1 & 0 & 1024 & 0 & \texttt{----------}\\
7 & 4 & 5 & 1 & 1 & 0 & 1024 & 0 & \texttt{----------}\\
8 & 1 & 2 & 1 & 2 & 0 & 1024 & 0 & \texttt{----------}\\
9 & 1 & 2 & 1 & 2 & 0 & 1024 & 0 & \texttt{----------}\\
\multicolumn{9}{@{}l@{}}{\rule{0pt}{2.6ex}\textbf{T10.091}\enspace \texttt{0123456521}\enspace $n{=}10$, $|H|{=}4$, len 7, br 3, $|\mathrm{Aut}|{=}1$, orb 10/10, hist-orb 4/4; M 29696, res 11768, rev 8970: ok}\\*
0 & -- & 0 & 4 & 1 & 374 & 2 & 648 & \texttt{-?????????}\\*
1 & 0 & 1 & 3 & 1 & 472 & 8 & 544 & \texttt{--???????-}\\
2 & 1 & 2 & 2 & 1 & 864 & 32 & 128 & \texttt{---++???--}\\
3 & 2 & 3 & 2 & 1 & 704 & 64 & 256 & \texttt{----+???--}\\
4 & 3 & 4 & 2 & 1 & 384 & 128 & 512 & \texttt{-----???--}\\
5 & 4 & 5 & 1 & 1 & 512 & 512 & 0 & \texttt{------+---}\\
6 & 5 & 6 & 1 & 1 & 0 & 1024 & 0 & \texttt{----------}\\
7 & 4 & 5 & 1 & 1 & 0 & 1024 & 0 & \texttt{----------}\\
8 & 1 & 2 & 1 & 1 & 0 & 1024 & 0 & \texttt{----------}\\
9 & 0 & 1 & 1 & 1 & 0 & 1024 & 0 & \texttt{----------}\\
\multicolumn{9}{@{}l@{}}{\rule{0pt}{2.6ex}\textbf{T10.092}\enspace \texttt{0123456512}\enspace $n{=}10$, $|H|{=}3$, len 7, br 2, $|\mathrm{Aut}|{=}1$, orb 10/10, hist-orb 3/3; M 29696, res 5308, rev 11494: ok}\\*
0 & -- & 0 & 3 & 1 & 738 & 2 & 284 & \texttt{-?????????}\\*
1 & 0 & 1 & 2 & 1 & 944 & 16 & 64 & \texttt{--+++???--}\\
2 & 1 & 2 & 2 & 1 & 864 & 32 & 128 & \texttt{---++???--}\\
3 & 2 & 3 & 2 & 1 & 704 & 64 & 256 & \texttt{----+???--}\\
4 & 3 & 4 & 2 & 1 & 384 & 128 & 512 & \texttt{-----???--}\\
5 & 4 & 5 & 1 & 1 & 512 & 512 & 0 & \texttt{------+---}\\
6 & 5 & 6 & 1 & 1 & 0 & 1024 & 0 & \texttt{----------}\\
7 & 4 & 5 & 1 & 1 & 0 & 1024 & 0 & \texttt{----------}\\
8 & 0 & 1 & 1 & 1 & 512 & 512 & 0 & \texttt{---------+}\\
9 & 8 & 2 & 1 & 1 & 0 & 1024 & 0 & \texttt{----------}\\
\multicolumn{9}{@{}l@{}}{\rule{0pt}{2.6ex}\textbf{T10.093}\enspace \texttt{0123456511}\enspace $n{=}10$, $|H|{=}4$, len 7, br 2, $|\mathrm{Aut}|{=}2$, orb 9/8, hist-orb 3/2; M 28672, res 9736, rev 9158: ok}\\*
0 & -- & 0 & 4 & 1 & 246 & 2 & 776 & \texttt{-?????????}\\*
1 & 0 & 1 & 2 & 1 & 944 & 16 & 64 & \texttt{--+++???--}\\
2 & 1 & 2 & 2 & 1 & 864 & 32 & 128 & \texttt{---++???--}\\
3 & 2 & 3 & 2 & 1 & 704 & 64 & 256 & \texttt{----+???--}\\
4 & 3 & 4 & 2 & 1 & 384 & 128 & 512 & \texttt{-----???--}\\
5 & 4 & 5 & 1 & 1 & 512 & 512 & 0 & \texttt{------+---}\\
6 & 5 & 6 & 1 & 1 & 0 & 1024 & 0 & \texttt{----------}\\
7 & 4 & 5 & 1 & 1 & 0 & 1024 & 0 & \texttt{----------}\\
8 & 0 & 1 & 1 & 2 & 0 & 1024 & 0 & \texttt{----------}\\
9 & 0 & 1 & 1 & 2 & 0 & 1024 & 0 & \texttt{----------}\\
\multicolumn{9}{@{}l@{}}{\rule{0pt}{2.6ex}\textbf{T10.094}\enspace \texttt{0123456456}\enspace $n{=}10$, $|H|{=}2$, len 7, br 1, $|\mathrm{Aut}|{=}2$, orb 7/4, hist-orb 1/0; M 36864, res 2520, rev 15354: ok}\\*
0 & -- & 0 & 2 & 1 & 994 & 2 & 28 & \texttt{-+++??????}\\*
1 & 0 & 1 & 2 & 1 & 964 & 4 & 56 & \texttt{--++??????}\\
2 & 1 & 2 & 2 & 1 & 904 & 8 & 112 & \texttt{---+??????}\\
3 & 2 & 3 & 2 & 1 & 784 & 16 & 224 & \texttt{----??????}\\
4 & 3 & 4 & 1 & 2 & 768 & 256 & 0 & \texttt{-----++---}\\
5 & 4 & 5 & 1 & 2 & 512 & 512 & 0 & \texttt{------+---}\\
6 & 5 & 6 & 1 & 2 & 0 & 1024 & 0 & \texttt{----------}\\
7 & 3 & 4 & 1 & 2 & 768 & 256 & 0 & \texttt{--------++}\\
8 & 7 & 5 & 1 & 2 & 512 & 512 & 0 & \texttt{---------+}\\
9 & 8 & 6 & 1 & 2 & 0 & 1024 & 0 & \texttt{----------}\\
\multicolumn{9}{@{}l@{}}{\rule{0pt}{2.6ex}\textbf{T10.095}\enspace \texttt{0123456455}\enspace $n{=}10$, $|H|{=}3$, len 7, br 2, $|\mathrm{Aut}|{=}2$, orb 9/8, hist-orb 2/1; M 35840, res 5584, rev 14758: ok}\\*
0 & -- & 0 & 3 & 1 & 966 & 2 & 56 & \texttt{-+++??????}\\*
1 & 0 & 1 & 3 & 1 & 908 & 4 & 112 & \texttt{--++??????}\\
2 & 1 & 2 & 3 & 1 & 792 & 8 & 224 & \texttt{---+??????}\\
3 & 2 & 3 & 3 & 1 & 560 & 16 & 448 & \texttt{----??????}\\
4 & 3 & 4 & 1 & 1 & 768 & 256 & 0 & \texttt{-----++---}\\
5 & 4 & 5 & 1 & 1 & 512 & 512 & 0 & \texttt{------+---}\\
6 & 5 & 6 & 1 & 1 & 0 & 1024 & 0 & \texttt{----------}\\
7 & 3 & 4 & 2 & 1 & 256 & 256 & 512 & \texttt{--------??}\\
8 & 7 & 5 & 1 & 2 & 0 & 1024 & 0 & \texttt{----------}\\
9 & 7 & 5 & 1 & 2 & 0 & 1024 & 0 & \texttt{----------}\\
\multicolumn{9}{@{}l@{}}{\rule{0pt}{2.6ex}\textbf{T10.096}\enspace \texttt{0123456454}\enspace $n{=}10$, $|H|{=}3$, len 7, br 1, $|\mathrm{Aut}|{=}1$, orb 10/10, hist-orb 3/3; M 34816, res 7560, rev 13802: ok}\\*
0 & -- & 0 & 3 & 1 & 938 & 2 & 84 & \texttt{-+++??????}\\*
1 & 0 & 1 & 3 & 1 & 852 & 4 & 168 & \texttt{--++??????}\\
2 & 1 & 2 & 3 & 1 & 680 & 8 & 336 & \texttt{---+??????}\\
3 & 2 & 3 & 3 & 1 & 336 & 16 & 672 & \texttt{----??????}\\
4 & 3 & 4 & 1 & 1 & 768 & 256 & 0 & \texttt{-----++---}\\
5 & 4 & 5 & 1 & 1 & 512 & 512 & 0 & \texttt{------+---}\\
6 & 5 & 6 & 1 & 1 & 0 & 1024 & 0 & \texttt{----------}\\
7 & 3 & 4 & 1 & 1 & 512 & 512 & 0 & \texttt{--------+-}\\
8 & 7 & 5 & 1 & 1 & 0 & 1024 & 0 & \texttt{----------}\\
9 & 3 & 4 & 1 & 1 & 0 & 1024 & 0 & \texttt{----------}\\
\multicolumn{9}{@{}l@{}}{\rule{0pt}{2.6ex}\textbf{T10.097}\enspace \texttt{0123456453}\enspace $n{=}10$, $|H|{=}3$, len 7, br 2, $|\mathrm{Aut}|{=}1$, orb 10/10, hist-orb 3/3; M 33792, res 8572, rev 12558: ok}\\*
0 & -- & 0 & 3 & 1 & 874 & 2 & 148 & \texttt{-++???????}\\*
1 & 0 & 1 & 3 & 1 & 724 & 4 & 296 & \texttt{--+???????}\\
2 & 1 & 2 & 3 & 1 & 424 & 8 & 592 & \texttt{---???????}\\
3 & 2 & 3 & 2 & 1 & 672 & 32 & 320 & \texttt{----?????-}\\
4 & 3 & 4 & 1 & 1 & 768 & 256 & 0 & \texttt{-----++---}\\
5 & 4 & 5 & 1 & 1 & 512 & 512 & 0 & \texttt{------+---}\\
6 & 5 & 6 & 1 & 1 & 0 & 1024 & 0 & \texttt{----------}\\
7 & 3 & 4 & 1 & 1 & 512 & 512 & 0 & \texttt{--------+-}\\
8 & 7 & 5 & 1 & 1 & 0 & 1024 & 0 & \texttt{----------}\\
9 & 2 & 3 & 1 & 1 & 0 & 1024 & 0 & \texttt{----------}\\
\multicolumn{9}{@{}l@{}}{\rule{0pt}{2.6ex}\textbf{T10.098}\enspace \texttt{0123456452}\enspace $n{=}10$, $|H|{=}3$, len 7, br 2, $|\mathrm{Aut}|{=}1$, orb 10/10, hist-orb 3/3; M 32768, res 8784, rev 11822: ok}\\*
0 & -- & 0 & 3 & 1 & 746 & 2 & 276 & \texttt{-+????????}\\*
1 & 0 & 1 & 3 & 1 & 468 & 4 & 552 & \texttt{--????????}\\
2 & 1 & 2 & 2 & 1 & 848 & 16 & 160 & \texttt{---+?????-}\\
3 & 2 & 3 & 2 & 1 & 672 & 32 & 320 & \texttt{----?????-}\\
4 & 3 & 4 & 1 & 1 & 768 & 256 & 0 & \texttt{-----++---}\\
5 & 4 & 5 & 1 & 1 & 512 & 512 & 0 & \texttt{------+---}\\
6 & 5 & 6 & 1 & 1 & 0 & 1024 & 0 & \texttt{----------}\\
7 & 3 & 4 & 1 & 1 & 512 & 512 & 0 & \texttt{--------+-}\\
8 & 7 & 5 & 1 & 1 & 0 & 1024 & 0 & \texttt{----------}\\
9 & 1 & 2 & 1 & 1 & 0 & 1024 & 0 & \texttt{----------}\\
\multicolumn{9}{@{}l@{}}{\rule{0pt}{2.6ex}\textbf{T10.099}\enspace \texttt{0123456451}\enspace $n{=}10$, $|H|{=}3$, len 7, br 2, $|\mathrm{Aut}|{=}1$, orb 10/10, hist-orb 3/3; M 31744, res 7468, rev 11762: ok}\\*
0 & -- & 0 & 3 & 1 & 490 & 2 & 532 & \texttt{-?????????}\\*
1 & 0 & 1 & 2 & 1 & 936 & 8 & 80 & \texttt{--++?????-}\\
2 & 1 & 2 & 2 & 1 & 848 & 16 & 160 & \texttt{---+?????-}\\
3 & 2 & 3 & 2 & 1 & 672 & 32 & 320 & \texttt{----?????-}\\
4 & 3 & 4 & 1 & 1 & 768 & 256 & 0 & \texttt{-----++---}\\
5 & 4 & 5 & 1 & 1 & 512 & 512 & 0 & \texttt{------+---}\\
6 & 5 & 6 & 1 & 1 & 0 & 1024 & 0 & \texttt{----------}\\
7 & 3 & 4 & 1 & 1 & 512 & 512 & 0 & \texttt{--------+-}\\
8 & 7 & 5 & 1 & 1 & 0 & 1024 & 0 & \texttt{----------}\\
9 & 0 & 1 & 1 & 1 & 0 & 1024 & 0 & \texttt{----------}\\
\multicolumn{9}{@{}l@{}}{\rule{0pt}{2.6ex}\textbf{T10.100}\enspace \texttt{0123456444}\enspace $n{=}10$, $|H|{=}4$, len 7, br 1, $|\mathrm{Aut}|{=}6$, orb 8/7, hist-orb 2/1; M 33792, res 10080, rev 12878: ok}\\*
0 & -- & 0 & 4 & 1 & 910 & 2 & 112 & \texttt{-+++??????}\\*
1 & 0 & 1 & 4 & 1 & 796 & 4 & 224 & \texttt{--++??????}\\
2 & 1 & 2 & 4 & 1 & 568 & 8 & 448 & \texttt{---+??????}\\
3 & 2 & 3 & 4 & 1 & 112 & 16 & 896 & \texttt{----??????}\\
4 & 3 & 4 & 1 & 1 & 768 & 256 & 0 & \texttt{-----++---}\\
5 & 4 & 5 & 1 & 1 & 512 & 512 & 0 & \texttt{------+---}\\
6 & 5 & 6 & 1 & 1 & 0 & 1024 & 0 & \texttt{----------}\\
7 & 3 & 4 & 1 & 3 & 0 & 1024 & 0 & \texttt{----------}\\
8 & 3 & 4 & 1 & 3 & 0 & 1024 & 0 & \texttt{----------}\\
9 & 3 & 4 & 1 & 3 & 0 & 1024 & 0 & \texttt{----------}\\
\multicolumn{9}{@{}l@{}}{\rule{0pt}{2.6ex}\textbf{T10.101}\enspace \texttt{0123456443}\enspace $n{=}10$, $|H|{=}4$, len 7, br 2, $|\mathrm{Aut}|{=}2$, orb 9/8, hist-orb 3/2; M 32768, res 11792, rev 11058: ok}\\*
0 & -- & 0 & 4 & 1 & 846 & 2 & 176 & \texttt{-++???????}\\*
1 & 0 & 1 & 4 & 1 & 668 & 4 & 352 & \texttt{--+???????}\\
2 & 1 & 2 & 4 & 1 & 312 & 8 & 704 & \texttt{---???????}\\
3 & 2 & 3 & 3 & 1 & 224 & 32 & 768 & \texttt{----?????-}\\
4 & 3 & 4 & 1 & 1 & 768 & 256 & 0 & \texttt{-----++---}\\
5 & 4 & 5 & 1 & 1 & 512 & 512 & 0 & \texttt{------+---}\\
6 & 5 & 6 & 1 & 1 & 0 & 1024 & 0 & \texttt{----------}\\
7 & 3 & 4 & 1 & 2 & 0 & 1024 & 0 & \texttt{----------}\\
8 & 3 & 4 & 1 & 2 & 0 & 1024 & 0 & \texttt{----------}\\
9 & 2 & 3 & 1 & 1 & 0 & 1024 & 0 & \texttt{----------}\\
\multicolumn{9}{@{}l@{}}{\rule{0pt}{2.6ex}\textbf{T10.102}\enspace \texttt{0123456442}\enspace $n{=}10$, $|H|{=}4$, len 7, br 2, $|\mathrm{Aut}|{=}2$, orb 9/8, hist-orb 3/2; M 31744, res 12480, rev 10082: ok}\\*
0 & -- & 0 & 4 & 1 & 718 & 2 & 304 & \texttt{-+????????}\\*
1 & 0 & 1 & 4 & 1 & 412 & 4 & 608 & \texttt{--????????}\\
2 & 1 & 2 & 3 & 1 & 624 & 16 & 384 & \texttt{---+?????-}\\
3 & 2 & 3 & 3 & 1 & 224 & 32 & 768 & \texttt{----?????-}\\
4 & 3 & 4 & 1 & 1 & 768 & 256 & 0 & \texttt{-----++---}\\
5 & 4 & 5 & 1 & 1 & 512 & 512 & 0 & \texttt{------+---}\\
6 & 5 & 6 & 1 & 1 & 0 & 1024 & 0 & \texttt{----------}\\
7 & 3 & 4 & 1 & 2 & 0 & 1024 & 0 & \texttt{----------}\\
8 & 3 & 4 & 1 & 2 & 0 & 1024 & 0 & \texttt{----------}\\
9 & 1 & 2 & 1 & 1 & 0 & 1024 & 0 & \texttt{----------}\\
\multicolumn{9}{@{}l@{}}{\rule{0pt}{2.6ex}\textbf{T10.103}\enspace \texttt{0123456441}\enspace $n{=}10$, $|H|{=}4$, len 7, br 2, $|\mathrm{Aut}|{=}2$, orb 9/8, hist-orb 3/2; M 30720, res 11472, rev 9998: ok}\\*
0 & -- & 0 & 4 & 1 & 462 & 2 & 560 & \texttt{-?????????}\\*
1 & 0 & 1 & 3 & 1 & 824 & 8 & 192 & \texttt{--++?????-}\\
2 & 1 & 2 & 3 & 1 & 624 & 16 & 384 & \texttt{---+?????-}\\
3 & 2 & 3 & 3 & 1 & 224 & 32 & 768 & \texttt{----?????-}\\
4 & 3 & 4 & 1 & 1 & 768 & 256 & 0 & \texttt{-----++---}\\
5 & 4 & 5 & 1 & 1 & 512 & 512 & 0 & \texttt{------+---}\\
6 & 5 & 6 & 1 & 1 & 0 & 1024 & 0 & \texttt{----------}\\
7 & 3 & 4 & 1 & 2 & 0 & 1024 & 0 & \texttt{----------}\\
8 & 3 & 4 & 1 & 2 & 0 & 1024 & 0 & \texttt{----------}\\
9 & 0 & 1 & 1 & 1 & 0 & 1024 & 0 & \texttt{----------}\\
\multicolumn{9}{@{}l@{}}{\rule{0pt}{2.6ex}\textbf{T10.104}\enspace \texttt{0123456434}\enspace $n{=}10$, $|H|{=}3$, len 7, br 2, $|\mathrm{Aut}|{=}1$, orb 10/10, hist-orb 3/3; M 32768, res 7284, rev 12666: ok}\\*
0 & -- & 0 & 3 & 1 & 906 & 2 & 116 & \texttt{-++???????}\\*
1 & 0 & 1 & 3 & 1 & 788 & 4 & 232 & \texttt{--+???????}\\
2 & 1 & 2 & 3 & 1 & 552 & 8 & 464 & \texttt{---???????}\\
3 & 2 & 3 & 2 & 1 & 448 & 64 & 512 & \texttt{----????--}\\
4 & 3 & 4 & 1 & 1 & 768 & 256 & 0 & \texttt{-----++---}\\
5 & 4 & 5 & 1 & 1 & 512 & 512 & 0 & \texttt{------+---}\\
6 & 5 & 6 & 1 & 1 & 0 & 1024 & 0 & \texttt{----------}\\
7 & 3 & 4 & 1 & 1 & 0 & 1024 & 0 & \texttt{----------}\\
8 & 2 & 3 & 1 & 1 & 512 & 512 & 0 & \texttt{---------+}\\
9 & 8 & 4 & 1 & 1 & 0 & 1024 & 0 & \texttt{----------}\\
\multicolumn{9}{@{}l@{}}{\rule{0pt}{2.6ex}\textbf{T10.105}\enspace \texttt{0123456433}\enspace $n{=}10$, $|H|{=}4$, len 7, br 2, $|\mathrm{Aut}|{=}2$, orb 9/8, hist-orb 3/2; M 31744, res 11792, rev 10386: ok}\\*
0 & -- & 0 & 4 & 1 & 814 & 2 & 208 & \texttt{-++???????}\\*
1 & 0 & 1 & 4 & 1 & 604 & 4 & 416 & \texttt{--+???????}\\
2 & 1 & 2 & 4 & 1 & 184 & 8 & 832 & \texttt{---???????}\\
3 & 2 & 3 & 2 & 1 & 448 & 64 & 512 & \texttt{----????--}\\
4 & 3 & 4 & 1 & 1 & 768 & 256 & 0 & \texttt{-----++---}\\
5 & 4 & 5 & 1 & 1 & 512 & 512 & 0 & \texttt{------+---}\\
6 & 5 & 6 & 1 & 1 & 0 & 1024 & 0 & \texttt{----------}\\
7 & 3 & 4 & 1 & 1 & 0 & 1024 & 0 & \texttt{----------}\\
8 & 2 & 3 & 1 & 2 & 0 & 1024 & 0 & \texttt{----------}\\
9 & 2 & 3 & 1 & 2 & 0 & 1024 & 0 & \texttt{----------}\\
\multicolumn{9}{@{}l@{}}{\rule{0pt}{2.6ex}\textbf{T10.106}\enspace \texttt{0123456432}\enspace $n{=}10$, $|H|{=}4$, len 7, br 3, $|\mathrm{Aut}|{=}1$, orb 10/10, hist-orb 4/4; M 30720, res 13024, rev 9046: ok}\\*
0 & -- & 0 & 4 & 1 & 686 & 2 & 336 & \texttt{-+????????}\\*
1 & 0 & 1 & 4 & 1 & 348 & 4 & 672 & \texttt{--????????}\\
2 & 1 & 2 & 3 & 1 & 368 & 16 & 640 & \texttt{---??????-}\\
3 & 2 & 3 & 2 & 1 & 448 & 64 & 512 & \texttt{----????--}\\
4 & 3 & 4 & 1 & 1 & 768 & 256 & 0 & \texttt{-----++---}\\
5 & 4 & 5 & 1 & 1 & 512 & 512 & 0 & \texttt{------+---}\\
6 & 5 & 6 & 1 & 1 & 0 & 1024 & 0 & \texttt{----------}\\
7 & 3 & 4 & 1 & 1 & 0 & 1024 & 0 & \texttt{----------}\\
8 & 2 & 3 & 1 & 1 & 0 & 1024 & 0 & \texttt{----------}\\
9 & 1 & 2 & 1 & 1 & 0 & 1024 & 0 & \texttt{----------}\\
\multicolumn{9}{@{}l@{}}{\rule{0pt}{2.6ex}\textbf{T10.107}\enspace \texttt{0123456431}\enspace $n{=}10$, $|H|{=}4$, len 7, br 3, $|\mathrm{Aut}|{=}1$, orb 10/10, hist-orb 4/4; M 29696, res 12368, rev 8786: ok}\\*
0 & -- & 0 & 4 & 1 & 430 & 2 & 592 & \texttt{-?????????}\\*
1 & 0 & 1 & 3 & 1 & 696 & 8 & 320 & \texttt{--+??????-}\\
2 & 1 & 2 & 3 & 1 & 368 & 16 & 640 & \texttt{---??????-}\\
3 & 2 & 3 & 2 & 1 & 448 & 64 & 512 & \texttt{----????--}\\
4 & 3 & 4 & 1 & 1 & 768 & 256 & 0 & \texttt{-----++---}\\
5 & 4 & 5 & 1 & 1 & 512 & 512 & 0 & \texttt{------+---}\\
6 & 5 & 6 & 1 & 1 & 0 & 1024 & 0 & \texttt{----------}\\
7 & 3 & 4 & 1 & 1 & 0 & 1024 & 0 & \texttt{----------}\\
8 & 2 & 3 & 1 & 1 & 0 & 1024 & 0 & \texttt{----------}\\
9 & 0 & 1 & 1 & 1 & 0 & 1024 & 0 & \texttt{----------}\\
\multicolumn{9}{@{}l@{}}{\rule{0pt}{2.6ex}\textbf{T10.108}\enspace \texttt{0123456423}\enspace $n{=}10$, $|H|{=}3$, len 7, br 2, $|\mathrm{Aut}|{=}1$, orb 10/10, hist-orb 3/3; M 30720, res 7008, rev 11534: ok}\\*
0 & -- & 0 & 3 & 1 & 842 & 2 & 180 & \texttt{-+????????}\\*
1 & 0 & 1 & 3 & 1 & 660 & 4 & 360 & \texttt{--????????}\\
2 & 1 & 2 & 2 & 1 & 736 & 32 & 256 & \texttt{---+????--}\\
3 & 2 & 3 & 2 & 1 & 448 & 64 & 512 & \texttt{----????--}\\
4 & 3 & 4 & 1 & 1 & 768 & 256 & 0 & \texttt{-----++---}\\
5 & 4 & 5 & 1 & 1 & 512 & 512 & 0 & \texttt{------+---}\\
6 & 5 & 6 & 1 & 1 & 0 & 1024 & 0 & \texttt{----------}\\
7 & 3 & 4 & 1 & 1 & 0 & 1024 & 0 & \texttt{----------}\\
8 & 1 & 2 & 1 & 1 & 512 & 512 & 0 & \texttt{---------+}\\
9 & 8 & 3 & 1 & 1 & 0 & 1024 & 0 & \texttt{----------}\\
\multicolumn{9}{@{}l@{}}{\rule{0pt}{2.6ex}\textbf{T10.109}\enspace \texttt{0123456422}\enspace $n{=}10$, $|H|{=}4$, len 7, br 2, $|\mathrm{Aut}|{=}2$, orb 9/8, hist-orb 3/2; M 29696, res 12288, rev 8770: ok}\\*
0 & -- & 0 & 4 & 1 & 622 & 2 & 400 & \texttt{-+????????}\\*
1 & 0 & 1 & 4 & 1 & 220 & 4 & 800 & \texttt{--????????}\\
2 & 1 & 2 & 2 & 1 & 736 & 32 & 256 & \texttt{---+????--}\\
3 & 2 & 3 & 2 & 1 & 448 & 64 & 512 & \texttt{----????--}\\
4 & 3 & 4 & 1 & 1 & 768 & 256 & 0 & \texttt{-----++---}\\
5 & 4 & 5 & 1 & 1 & 512 & 512 & 0 & \texttt{------+---}\\
6 & 5 & 6 & 1 & 1 & 0 & 1024 & 0 & \texttt{----------}\\
7 & 3 & 4 & 1 & 1 & 0 & 1024 & 0 & \texttt{----------}\\
8 & 1 & 2 & 1 & 2 & 0 & 1024 & 0 & \texttt{----------}\\
9 & 1 & 2 & 1 & 2 & 0 & 1024 & 0 & \texttt{----------}\\
\multicolumn{9}{@{}l@{}}{\rule{0pt}{2.6ex}\textbf{T10.110}\enspace \texttt{0123456421}\enspace $n{=}10$, $|H|{=}4$, len 7, br 3, $|\mathrm{Aut}|{=}1$, orb 10/10, hist-orb 4/4; M 28672, res 12336, rev 8146: ok}\\*
0 & -- & 0 & 4 & 1 & 366 & 2 & 656 & \texttt{-?????????}\\*
1 & 0 & 1 & 3 & 1 & 440 & 8 & 576 & \texttt{--???????-}\\
2 & 1 & 2 & 2 & 1 & 736 & 32 & 256 & \texttt{---+????--}\\
3 & 2 & 3 & 2 & 1 & 448 & 64 & 512 & \texttt{----????--}\\
4 & 3 & 4 & 1 & 1 & 768 & 256 & 0 & \texttt{-----++---}\\
5 & 4 & 5 & 1 & 1 & 512 & 512 & 0 & \texttt{------+---}\\
6 & 5 & 6 & 1 & 1 & 0 & 1024 & 0 & \texttt{----------}\\
7 & 3 & 4 & 1 & 1 & 0 & 1024 & 0 & \texttt{----------}\\
8 & 1 & 2 & 1 & 1 & 0 & 1024 & 0 & \texttt{----------}\\
9 & 0 & 1 & 1 & 1 & 0 & 1024 & 0 & \texttt{----------}\\
\multicolumn{9}{@{}l@{}}{\rule{0pt}{2.6ex}\textbf{T10.111}\enspace \texttt{0123456412}\enspace $n{=}10$, $|H|{=}3$, len 7, br 2, $|\mathrm{Aut}|{=}1$, orb 10/10, hist-orb 3/3; M 28672, res 6164, rev 10654: ok}\\*
0 & -- & 0 & 3 & 1 & 714 & 2 & 308 & \texttt{-?????????}\\*
1 & 0 & 1 & 2 & 1 & 880 & 16 & 128 & \texttt{--++????--}\\
2 & 1 & 2 & 2 & 1 & 736 & 32 & 256 & \texttt{---+????--}\\
3 & 2 & 3 & 2 & 1 & 448 & 64 & 512 & \texttt{----????--}\\
4 & 3 & 4 & 1 & 1 & 768 & 256 & 0 & \texttt{-----++---}\\
5 & 4 & 5 & 1 & 1 & 512 & 512 & 0 & \texttt{------+---}\\
6 & 5 & 6 & 1 & 1 & 0 & 1024 & 0 & \texttt{----------}\\
7 & 3 & 4 & 1 & 1 & 0 & 1024 & 0 & \texttt{----------}\\
8 & 0 & 1 & 1 & 1 & 512 & 512 & 0 & \texttt{---------+}\\
9 & 8 & 2 & 1 & 1 & 0 & 1024 & 0 & \texttt{----------}\\
\multicolumn{9}{@{}l@{}}{\rule{0pt}{2.6ex}\textbf{T10.112}\enspace \texttt{0123456411}\enspace $n{=}10$, $|H|{=}4$, len 7, br 2, $|\mathrm{Aut}|{=}2$, orb 9/8, hist-orb 3/2; M 27648, res 10448, rev 8318: ok}\\*
0 & -- & 0 & 4 & 1 & 238 & 2 & 784 & \texttt{-?????????}\\*
1 & 0 & 1 & 2 & 1 & 880 & 16 & 128 & \texttt{--++????--}\\
2 & 1 & 2 & 2 & 1 & 736 & 32 & 256 & \texttt{---+????--}\\
3 & 2 & 3 & 2 & 1 & 448 & 64 & 512 & \texttt{----????--}\\
4 & 3 & 4 & 1 & 1 & 768 & 256 & 0 & \texttt{-----++---}\\
5 & 4 & 5 & 1 & 1 & 512 & 512 & 0 & \texttt{------+---}\\
6 & 5 & 6 & 1 & 1 & 0 & 1024 & 0 & \texttt{----------}\\
7 & 3 & 4 & 1 & 1 & 0 & 1024 & 0 & \texttt{----------}\\
8 & 0 & 1 & 1 & 2 & 0 & 1024 & 0 & \texttt{----------}\\
9 & 0 & 1 & 1 & 2 & 0 & 1024 & 0 & \texttt{----------}\\
\multicolumn{9}{@{}l@{}}{\rule{0pt}{2.6ex}\textbf{T10.113}\enspace \texttt{0123456345}\enspace $n{=}10$, $|H|{=}2$, len 7, br 1, $|\mathrm{Aut}|{=}1$, orb 10/10, hist-orb 2/2; M 33792, res 2156, rev 13990: ok}\\*
0 & -- & 0 & 2 & 1 & 978 & 2 & 44 & \texttt{-++???????}\\*
1 & 0 & 1 & 2 & 1 & 932 & 4 & 88 & \texttt{--+???????}\\
2 & 1 & 2 & 2 & 1 & 840 & 8 & 176 & \texttt{---???????}\\
3 & 2 & 3 & 1 & 1 & 896 & 128 & 0 & \texttt{----+++---}\\
4 & 3 & 4 & 1 & 1 & 768 & 256 & 0 & \texttt{-----++---}\\
5 & 4 & 5 & 1 & 1 & 512 & 512 & 0 & \texttt{------+---}\\
6 & 5 & 6 & 1 & 1 & 0 & 1024 & 0 & \texttt{----------}\\
7 & 2 & 3 & 1 & 1 & 768 & 256 & 0 & \texttt{--------++}\\
8 & 7 & 4 & 1 & 1 & 512 & 512 & 0 & \texttt{---------+}\\
9 & 8 & 5 & 1 & 1 & 0 & 1024 & 0 & \texttt{----------}\\
\multicolumn{9}{@{}l@{}}{\rule{0pt}{2.6ex}\textbf{T10.114}\enspace \texttt{0123456344}\enspace $n{=}10$, $|H|{=}3$, len 7, br 2, $|\mathrm{Aut}|{=}2$, orb 9/8, hist-orb 2/1; M 32768, res 5672, rev 12878: ok}\\*
0 & -- & 0 & 3 & 1 & 918 & 2 & 104 & \texttt{-++???????}\\*
1 & 0 & 1 & 3 & 1 & 812 & 4 & 208 & \texttt{--+???????}\\
2 & 1 & 2 & 3 & 1 & 600 & 8 & 416 & \texttt{---???????}\\
3 & 2 & 3 & 1 & 1 & 896 & 128 & 0 & \texttt{----+++---}\\
4 & 3 & 4 & 1 & 1 & 768 & 256 & 0 & \texttt{-----++---}\\
5 & 4 & 5 & 1 & 1 & 512 & 512 & 0 & \texttt{------+---}\\
6 & 5 & 6 & 1 & 1 & 0 & 1024 & 0 & \texttt{----------}\\
7 & 2 & 3 & 2 & 1 & 256 & 256 & 512 & \texttt{--------??}\\
8 & 7 & 4 & 1 & 2 & 0 & 1024 & 0 & \texttt{----------}\\
9 & 7 & 4 & 1 & 2 & 0 & 1024 & 0 & \texttt{----------}\\
\multicolumn{9}{@{}l@{}}{\rule{0pt}{2.6ex}\textbf{T10.115}\enspace \texttt{0123456343}\enspace $n{=}10$, $|H|{=}3$, len 7, br 1, $|\mathrm{Aut}|{=}1$, orb 10/10, hist-orb 3/3; M 31744, res 8036, rev 11630: ok}\\*
0 & -- & 0 & 3 & 1 & 858 & 2 & 164 & \texttt{-++???????}\\*
1 & 0 & 1 & 3 & 1 & 692 & 4 & 328 & \texttt{--+???????}\\
2 & 1 & 2 & 3 & 1 & 360 & 8 & 656 & \texttt{---???????}\\
3 & 2 & 3 & 1 & 1 & 896 & 128 & 0 & \texttt{----+++---}\\
4 & 3 & 4 & 1 & 1 & 768 & 256 & 0 & \texttt{-----++---}\\
5 & 4 & 5 & 1 & 1 & 512 & 512 & 0 & \texttt{------+---}\\
6 & 5 & 6 & 1 & 1 & 0 & 1024 & 0 & \texttt{----------}\\
7 & 2 & 3 & 1 & 1 & 512 & 512 & 0 & \texttt{--------+-}\\
8 & 7 & 4 & 1 & 1 & 0 & 1024 & 0 & \texttt{----------}\\
9 & 2 & 3 & 1 & 1 & 0 & 1024 & 0 & \texttt{----------}\\
\multicolumn{9}{@{}l@{}}{\rule{0pt}{2.6ex}\textbf{T10.116}\enspace \texttt{0123456342}\enspace $n{=}10$, $|H|{=}3$, len 7, br 2, $|\mathrm{Aut}|{=}1$, orb 10/10, hist-orb 3/3; M 30720, res 8520, rev 10698: ok}\\*
0 & -- & 0 & 3 & 1 & 730 & 2 & 292 & \texttt{-+????????}\\*
1 & 0 & 1 & 3 & 1 & 436 & 4 & 584 & \texttt{--????????}\\
2 & 1 & 2 & 2 & 1 & 720 & 16 & 288 & \texttt{---??????-}\\
3 & 2 & 3 & 1 & 1 & 896 & 128 & 0 & \texttt{----+++---}\\
4 & 3 & 4 & 1 & 1 & 768 & 256 & 0 & \texttt{-----++---}\\
5 & 4 & 5 & 1 & 1 & 512 & 512 & 0 & \texttt{------+---}\\
6 & 5 & 6 & 1 & 1 & 0 & 1024 & 0 & \texttt{----------}\\
7 & 2 & 3 & 1 & 1 & 512 & 512 & 0 & \texttt{--------+-}\\
8 & 7 & 4 & 1 & 1 & 0 & 1024 & 0 & \texttt{----------}\\
9 & 1 & 2 & 1 & 1 & 0 & 1024 & 0 & \texttt{----------}\\
\multicolumn{9}{@{}l@{}}{\rule{0pt}{2.6ex}\textbf{T10.117}\enspace \texttt{0123456341}\enspace $n{=}10$, $|H|{=}3$, len 7, br 2, $|\mathrm{Aut}|{=}1$, orb 10/10, hist-orb 3/3; M 29696, res 7380, rev 10558: ok}\\*
0 & -- & 0 & 3 & 1 & 474 & 2 & 548 & \texttt{-?????????}\\*
1 & 0 & 1 & 2 & 1 & 872 & 8 & 144 & \texttt{--+??????-}\\
2 & 1 & 2 & 2 & 1 & 720 & 16 & 288 & \texttt{---??????-}\\
3 & 2 & 3 & 1 & 1 & 896 & 128 & 0 & \texttt{----+++---}\\
4 & 3 & 4 & 1 & 1 & 768 & 256 & 0 & \texttt{-----++---}\\
5 & 4 & 5 & 1 & 1 & 512 & 512 & 0 & \texttt{------+---}\\
6 & 5 & 6 & 1 & 1 & 0 & 1024 & 0 & \texttt{----------}\\
7 & 2 & 3 & 1 & 1 & 512 & 512 & 0 & \texttt{--------+-}\\
8 & 7 & 4 & 1 & 1 & 0 & 1024 & 0 & \texttt{----------}\\
9 & 0 & 1 & 1 & 1 & 0 & 1024 & 0 & \texttt{----------}\\
\multicolumn{9}{@{}l@{}}{\rule{0pt}{2.6ex}\textbf{T10.118}\enspace \texttt{0123456333}\enspace $n{=}10$, $|H|{=}4$, len 7, br 1, $|\mathrm{Aut}|{=}6$, orb 8/7, hist-orb 2/1; M 30720, res 10976, rev 10190: ok}\\*
0 & -- & 0 & 4 & 1 & 798 & 2 & 224 & \texttt{-++???????}\\*
1 & 0 & 1 & 4 & 1 & 572 & 4 & 448 & \texttt{--+???????}\\
2 & 1 & 2 & 4 & 1 & 120 & 8 & 896 & \texttt{---???????}\\
3 & 2 & 3 & 1 & 1 & 896 & 128 & 0 & \texttt{----+++---}\\
4 & 3 & 4 & 1 & 1 & 768 & 256 & 0 & \texttt{-----++---}\\
5 & 4 & 5 & 1 & 1 & 512 & 512 & 0 & \texttt{------+---}\\
6 & 5 & 6 & 1 & 1 & 0 & 1024 & 0 & \texttt{----------}\\
7 & 2 & 3 & 1 & 3 & 0 & 1024 & 0 & \texttt{----------}\\
8 & 2 & 3 & 1 & 3 & 0 & 1024 & 0 & \texttt{----------}\\
9 & 2 & 3 & 1 & 3 & 0 & 1024 & 0 & \texttt{----------}\\
\multicolumn{9}{@{}l@{}}{\rule{0pt}{2.6ex}\textbf{T10.119}\enspace \texttt{0123456332}\enspace $n{=}10$, $|H|{=}4$, len 7, br 2, $|\mathrm{Aut}|{=}2$, orb 9/8, hist-orb 3/2; M 29696, res 12480, rev 8650: ok}\\*
0 & -- & 0 & 4 & 1 & 670 & 2 & 352 & \texttt{-+????????}\\*
1 & 0 & 1 & 4 & 1 & 316 & 4 & 704 & \texttt{--????????}\\
2 & 1 & 2 & 3 & 1 & 240 & 16 & 768 & \texttt{---??????-}\\
3 & 2 & 3 & 1 & 1 & 896 & 128 & 0 & \texttt{----+++---}\\
4 & 3 & 4 & 1 & 1 & 768 & 256 & 0 & \texttt{-----++---}\\
5 & 4 & 5 & 1 & 1 & 512 & 512 & 0 & \texttt{------+---}\\
6 & 5 & 6 & 1 & 1 & 0 & 1024 & 0 & \texttt{----------}\\
7 & 2 & 3 & 1 & 2 & 0 & 1024 & 0 & \texttt{----------}\\
8 & 2 & 3 & 1 & 2 & 0 & 1024 & 0 & \texttt{----------}\\
9 & 1 & 2 & 1 & 1 & 0 & 1024 & 0 & \texttt{----------}\\
\multicolumn{9}{@{}l@{}}{\rule{0pt}{2.6ex}\textbf{T10.120}\enspace \texttt{0123456331}\enspace $n{=}10$, $|H|{=}4$, len 7, br 2, $|\mathrm{Aut}|{=}2$, orb 9/8, hist-orb 3/2; M 28672, res 12000, rev 8302: ok}\\*
0 & -- & 0 & 4 & 1 & 414 & 2 & 608 & \texttt{-?????????}\\*
1 & 0 & 1 & 3 & 1 & 632 & 8 & 384 & \texttt{--+??????-}\\
2 & 1 & 2 & 3 & 1 & 240 & 16 & 768 & \texttt{---??????-}\\
3 & 2 & 3 & 1 & 1 & 896 & 128 & 0 & \texttt{----+++---}\\
4 & 3 & 4 & 1 & 1 & 768 & 256 & 0 & \texttt{-----++---}\\
5 & 4 & 5 & 1 & 1 & 512 & 512 & 0 & \texttt{------+---}\\
6 & 5 & 6 & 1 & 1 & 0 & 1024 & 0 & \texttt{----------}\\
7 & 2 & 3 & 1 & 2 & 0 & 1024 & 0 & \texttt{----------}\\
8 & 2 & 3 & 1 & 2 & 0 & 1024 & 0 & \texttt{----------}\\
9 & 0 & 1 & 1 & 1 & 0 & 1024 & 0 & \texttt{----------}\\
\multicolumn{9}{@{}l@{}}{\rule{0pt}{2.6ex}\textbf{T10.121}\enspace \texttt{0123456323}\enspace $n{=}10$, $|H|{=}3$, len 7, br 2, $|\mathrm{Aut}|{=}1$, orb 10/10, hist-orb 3/3; M 29696, res 7648, rev 10566: ok}\\*
0 & -- & 0 & 3 & 1 & 794 & 2 & 228 & \texttt{-+????????}\\*
1 & 0 & 1 & 3 & 1 & 564 & 4 & 456 & \texttt{--????????}\\
2 & 1 & 2 & 2 & 1 & 480 & 32 & 512 & \texttt{---?????--}\\
3 & 2 & 3 & 1 & 1 & 896 & 128 & 0 & \texttt{----+++---}\\
4 & 3 & 4 & 1 & 1 & 768 & 256 & 0 & \texttt{-----++---}\\
5 & 4 & 5 & 1 & 1 & 512 & 512 & 0 & \texttt{------+---}\\
6 & 5 & 6 & 1 & 1 & 0 & 1024 & 0 & \texttt{----------}\\
7 & 2 & 3 & 1 & 1 & 0 & 1024 & 0 & \texttt{----------}\\
8 & 1 & 2 & 1 & 1 & 512 & 512 & 0 & \texttt{---------+}\\
9 & 8 & 3 & 1 & 1 & 0 & 1024 & 0 & \texttt{----------}\\
\multicolumn{9}{@{}l@{}}{\rule{0pt}{2.6ex}\textbf{T10.122}\enspace \texttt{0123456322}\enspace $n{=}10$, $|H|{=}4$, len 7, br 2, $|\mathrm{Aut}|{=}2$, orb 9/8, hist-orb 3/2; M 28672, res 12160, rev 8138: ok}\\*
0 & -- & 0 & 4 & 1 & 606 & 2 & 416 & \texttt{-+????????}\\*
1 & 0 & 1 & 4 & 1 & 188 & 4 & 832 & \texttt{--????????}\\
2 & 1 & 2 & 2 & 1 & 480 & 32 & 512 & \texttt{---?????--}\\
3 & 2 & 3 & 1 & 1 & 896 & 128 & 0 & \texttt{----+++---}\\
4 & 3 & 4 & 1 & 1 & 768 & 256 & 0 & \texttt{-----++---}\\
5 & 4 & 5 & 1 & 1 & 512 & 512 & 0 & \texttt{------+---}\\
6 & 5 & 6 & 1 & 1 & 0 & 1024 & 0 & \texttt{----------}\\
7 & 2 & 3 & 1 & 1 & 0 & 1024 & 0 & \texttt{----------}\\
8 & 1 & 2 & 1 & 2 & 0 & 1024 & 0 & \texttt{----------}\\
9 & 1 & 2 & 1 & 2 & 0 & 1024 & 0 & \texttt{----------}\\
\multicolumn{9}{@{}l@{}}{\rule{0pt}{2.6ex}\textbf{T10.123}\enspace \texttt{0123456321}\enspace $n{=}10$, $|H|{=}4$, len 7, br 3, $|\mathrm{Aut}|{=}1$, orb 10/10, hist-orb 4/4; M 27648, res 12384, rev 7418: ok}\\*
0 & -- & 0 & 4 & 1 & 350 & 2 & 672 & \texttt{-?????????}\\*
1 & 0 & 1 & 3 & 1 & 376 & 8 & 640 & \texttt{--???????-}\\
2 & 1 & 2 & 2 & 1 & 480 & 32 & 512 & \texttt{---?????--}\\
3 & 2 & 3 & 1 & 1 & 896 & 128 & 0 & \texttt{----+++---}\\
4 & 3 & 4 & 1 & 1 & 768 & 256 & 0 & \texttt{-----++---}\\
5 & 4 & 5 & 1 & 1 & 512 & 512 & 0 & \texttt{------+---}\\
6 & 5 & 6 & 1 & 1 & 0 & 1024 & 0 & \texttt{----------}\\
7 & 2 & 3 & 1 & 1 & 0 & 1024 & 0 & \texttt{----------}\\
8 & 1 & 2 & 1 & 1 & 0 & 1024 & 0 & \texttt{----------}\\
9 & 0 & 1 & 1 & 1 & 0 & 1024 & 0 & \texttt{----------}\\
\multicolumn{9}{@{}l@{}}{\rule{0pt}{2.6ex}\textbf{T10.124}\enspace \texttt{0123456312}\enspace $n{=}10$, $|H|{=}3$, len 7, br 2, $|\mathrm{Aut}|{=}1$, orb 10/10, hist-orb 3/3; M 27648, res 6788, rev 9726: ok}\\*
0 & -- & 0 & 3 & 1 & 666 & 2 & 356 & \texttt{-?????????}\\*
1 & 0 & 1 & 2 & 1 & 752 & 16 & 256 & \texttt{--+?????--}\\
2 & 1 & 2 & 2 & 1 & 480 & 32 & 512 & \texttt{---?????--}\\
3 & 2 & 3 & 1 & 1 & 896 & 128 & 0 & \texttt{----+++---}\\
4 & 3 & 4 & 1 & 1 & 768 & 256 & 0 & \texttt{-----++---}\\
5 & 4 & 5 & 1 & 1 & 512 & 512 & 0 & \texttt{------+---}\\
6 & 5 & 6 & 1 & 1 & 0 & 1024 & 0 & \texttt{----------}\\
7 & 2 & 3 & 1 & 1 & 0 & 1024 & 0 & \texttt{----------}\\
8 & 0 & 1 & 1 & 1 & 512 & 512 & 0 & \texttt{---------+}\\
9 & 8 & 2 & 1 & 1 & 0 & 1024 & 0 & \texttt{----------}\\
\multicolumn{9}{@{}l@{}}{\rule{0pt}{2.6ex}\textbf{T10.125}\enspace \texttt{0123456311}\enspace $n{=}10$, $|H|{=}4$, len 7, br 2, $|\mathrm{Aut}|{=}2$, orb 9/8, hist-orb 3/2; M 26624, res 10784, rev 7486: ok}\\*
0 & -- & 0 & 4 & 1 & 222 & 2 & 800 & \texttt{-?????????}\\*
1 & 0 & 1 & 2 & 1 & 752 & 16 & 256 & \texttt{--+?????--}\\
2 & 1 & 2 & 2 & 1 & 480 & 32 & 512 & \texttt{---?????--}\\
3 & 2 & 3 & 1 & 1 & 896 & 128 & 0 & \texttt{----+++---}\\
4 & 3 & 4 & 1 & 1 & 768 & 256 & 0 & \texttt{-----++---}\\
5 & 4 & 5 & 1 & 1 & 512 & 512 & 0 & \texttt{------+---}\\
6 & 5 & 6 & 1 & 1 & 0 & 1024 & 0 & \texttt{----------}\\
7 & 2 & 3 & 1 & 1 & 0 & 1024 & 0 & \texttt{----------}\\
8 & 0 & 1 & 1 & 2 & 0 & 1024 & 0 & \texttt{----------}\\
9 & 0 & 1 & 1 & 2 & 0 & 1024 & 0 & \texttt{----------}\\
\multicolumn{9}{@{}l@{}}{\rule{0pt}{2.6ex}\textbf{T10.126}\enspace \texttt{0123456234}\enspace $n{=}10$, $|H|{=}2$, len 7, br 1, $|\mathrm{Aut}|{=}1$, orb 10/10, hist-orb 2/2; M 30720, res 1824, rev 12474: ok}\\*
0 & -- & 0 & 2 & 1 & 946 & 2 & 76 & \texttt{-+????????}\\*
1 & 0 & 1 & 2 & 1 & 868 & 4 & 152 & \texttt{--????????}\\
2 & 1 & 2 & 1 & 1 & 960 & 64 & 0 & \texttt{---++++---}\\
3 & 2 & 3 & 1 & 1 & 896 & 128 & 0 & \texttt{----+++---}\\
4 & 3 & 4 & 1 & 1 & 768 & 256 & 0 & \texttt{-----++---}\\
5 & 4 & 5 & 1 & 1 & 512 & 512 & 0 & \texttt{------+---}\\
6 & 5 & 6 & 1 & 1 & 0 & 1024 & 0 & \texttt{----------}\\
7 & 1 & 2 & 1 & 1 & 768 & 256 & 0 & \texttt{--------++}\\
8 & 7 & 3 & 1 & 1 & 512 & 512 & 0 & \texttt{---------+}\\
9 & 8 & 4 & 1 & 1 & 0 & 1024 & 0 & \texttt{----------}\\
\multicolumn{9}{@{}l@{}}{\rule{0pt}{2.6ex}\textbf{T10.127}\enspace \texttt{0123456233}\enspace $n{=}10$, $|H|{=}3$, len 7, br 2, $|\mathrm{Aut}|{=}2$, orb 9/8, hist-orb 2/1; M 29696, res 5440, rev 11022: ok}\\*
0 & -- & 0 & 3 & 1 & 822 & 2 & 200 & \texttt{-+????????}\\*
1 & 0 & 1 & 3 & 1 & 620 & 4 & 400 & \texttt{--????????}\\
2 & 1 & 2 & 1 & 1 & 960 & 64 & 0 & \texttt{---++++---}\\
3 & 2 & 3 & 1 & 1 & 896 & 128 & 0 & \texttt{----+++---}\\
4 & 3 & 4 & 1 & 1 & 768 & 256 & 0 & \texttt{-----++---}\\
5 & 4 & 5 & 1 & 1 & 512 & 512 & 0 & \texttt{------+---}\\
6 & 5 & 6 & 1 & 1 & 0 & 1024 & 0 & \texttt{----------}\\
7 & 1 & 2 & 2 & 1 & 256 & 256 & 512 & \texttt{--------??}\\
8 & 7 & 3 & 1 & 2 & 0 & 1024 & 0 & \texttt{----------}\\
9 & 7 & 3 & 1 & 2 & 0 & 1024 & 0 & \texttt{----------}\\
\multicolumn{9}{@{}l@{}}{\rule{0pt}{2.6ex}\textbf{T10.128}\enspace \texttt{0123456232}\enspace $n{=}10$, $|H|{=}3$, len 7, br 1, $|\mathrm{Aut}|{=}1$, orb 10/10, hist-orb 3/3; M 28672, res 7776, rev 9722: ok}\\*
0 & -- & 0 & 3 & 1 & 698 & 2 & 324 & \texttt{-+????????}\\*
1 & 0 & 1 & 3 & 1 & 372 & 4 & 648 & \texttt{--????????}\\
2 & 1 & 2 & 1 & 1 & 960 & 64 & 0 & \texttt{---++++---}\\
3 & 2 & 3 & 1 & 1 & 896 & 128 & 0 & \texttt{----+++---}\\
4 & 3 & 4 & 1 & 1 & 768 & 256 & 0 & \texttt{-----++---}\\
5 & 4 & 5 & 1 & 1 & 512 & 512 & 0 & \texttt{------+---}\\
6 & 5 & 6 & 1 & 1 & 0 & 1024 & 0 & \texttt{----------}\\
7 & 1 & 2 & 1 & 1 & 512 & 512 & 0 & \texttt{--------+-}\\
8 & 7 & 3 & 1 & 1 & 0 & 1024 & 0 & \texttt{----------}\\
9 & 1 & 2 & 1 & 1 & 0 & 1024 & 0 & \texttt{----------}\\
\multicolumn{9}{@{}l@{}}{\rule{0pt}{2.6ex}\textbf{T10.129}\enspace \texttt{0123456231}\enspace $n{=}10$, $|H|{=}3$, len 7, br 2, $|\mathrm{Aut}|{=}1$, orb 10/10, hist-orb 3/3; M 27648, res 6988, rev 9394: ok}\\*
0 & -- & 0 & 3 & 1 & 442 & 2 & 580 & \texttt{-?????????}\\*
1 & 0 & 1 & 2 & 1 & 744 & 8 & 272 & \texttt{--???????-}\\
2 & 1 & 2 & 1 & 1 & 960 & 64 & 0 & \texttt{---++++---}\\
3 & 2 & 3 & 1 & 1 & 896 & 128 & 0 & \texttt{----+++---}\\
4 & 3 & 4 & 1 & 1 & 768 & 256 & 0 & \texttt{-----++---}\\
5 & 4 & 5 & 1 & 1 & 512 & 512 & 0 & \texttt{------+---}\\
6 & 5 & 6 & 1 & 1 & 0 & 1024 & 0 & \texttt{----------}\\
7 & 1 & 2 & 1 & 1 & 512 & 512 & 0 & \texttt{--------+-}\\
8 & 7 & 3 & 1 & 1 & 0 & 1024 & 0 & \texttt{----------}\\
9 & 0 & 1 & 1 & 1 & 0 & 1024 & 0 & \texttt{----------}\\
\multicolumn{9}{@{}l@{}}{\rule{0pt}{2.6ex}\textbf{T10.130}\enspace \texttt{0123456222}\enspace $n{=}10$, $|H|{=}4$, len 7, br 1, $|\mathrm{Aut}|{=}6$, orb 8/7, hist-orb 2/1; M 27648, res 10752, rev 8038: ok}\\*
0 & -- & 0 & 4 & 1 & 574 & 2 & 448 & \texttt{-+????????}\\*
1 & 0 & 1 & 4 & 1 & 124 & 4 & 896 & \texttt{--????????}\\
2 & 1 & 2 & 1 & 1 & 960 & 64 & 0 & \texttt{---++++---}\\
3 & 2 & 3 & 1 & 1 & 896 & 128 & 0 & \texttt{----+++---}\\
4 & 3 & 4 & 1 & 1 & 768 & 256 & 0 & \texttt{-----++---}\\
5 & 4 & 5 & 1 & 1 & 512 & 512 & 0 & \texttt{------+---}\\
6 & 5 & 6 & 1 & 1 & 0 & 1024 & 0 & \texttt{----------}\\
7 & 1 & 2 & 1 & 3 & 0 & 1024 & 0 & \texttt{----------}\\
8 & 1 & 2 & 1 & 3 & 0 & 1024 & 0 & \texttt{----------}\\
9 & 1 & 2 & 1 & 3 & 0 & 1024 & 0 & \texttt{----------}\\
\multicolumn{9}{@{}l@{}}{\rule{0pt}{2.6ex}\textbf{T10.131}\enspace \texttt{0123456221}\enspace $n{=}10$, $|H|{=}4$, len 7, br 2, $|\mathrm{Aut}|{=}2$, orb 9/8, hist-orb 3/2; M 26624, res 11328, rev 7126: ok}\\*
0 & -- & 0 & 4 & 1 & 318 & 2 & 704 & \texttt{-?????????}\\*
1 & 0 & 1 & 3 & 1 & 248 & 8 & 768 & \texttt{--???????-}\\
2 & 1 & 2 & 1 & 1 & 960 & 64 & 0 & \texttt{---++++---}\\
3 & 2 & 3 & 1 & 1 & 896 & 128 & 0 & \texttt{----+++---}\\
4 & 3 & 4 & 1 & 1 & 768 & 256 & 0 & \texttt{-----++---}\\
5 & 4 & 5 & 1 & 1 & 512 & 512 & 0 & \texttt{------+---}\\
6 & 5 & 6 & 1 & 1 & 0 & 1024 & 0 & \texttt{----------}\\
7 & 1 & 2 & 1 & 2 & 0 & 1024 & 0 & \texttt{----------}\\
8 & 1 & 2 & 1 & 2 & 0 & 1024 & 0 & \texttt{----------}\\
9 & 0 & 1 & 1 & 1 & 0 & 1024 & 0 & \texttt{----------}\\
\multicolumn{9}{@{}l@{}}{\rule{0pt}{2.6ex}\textbf{T10.132}\enspace \texttt{0123456212}\enspace $n{=}10$, $|H|{=}3$, len 7, br 2, $|\mathrm{Aut}|{=}1$, orb 10/10, hist-orb 3/3; M 26624, res 6884, rev 8950: ok}\\*
0 & -- & 0 & 3 & 1 & 570 & 2 & 452 & \texttt{-?????????}\\*
1 & 0 & 1 & 2 & 1 & 496 & 16 & 512 & \texttt{--??????--}\\
2 & 1 & 2 & 1 & 1 & 960 & 64 & 0 & \texttt{---++++---}\\
3 & 2 & 3 & 1 & 1 & 896 & 128 & 0 & \texttt{----+++---}\\
4 & 3 & 4 & 1 & 1 & 768 & 256 & 0 & \texttt{-----++---}\\
5 & 4 & 5 & 1 & 1 & 512 & 512 & 0 & \texttt{------+---}\\
6 & 5 & 6 & 1 & 1 & 0 & 1024 & 0 & \texttt{----------}\\
7 & 1 & 2 & 1 & 1 & 0 & 1024 & 0 & \texttt{----------}\\
8 & 0 & 1 & 1 & 1 & 512 & 512 & 0 & \texttt{---------+}\\
9 & 8 & 2 & 1 & 1 & 0 & 1024 & 0 & \texttt{----------}\\
\multicolumn{9}{@{}l@{}}{\rule{0pt}{2.6ex}\textbf{T10.133}\enspace \texttt{0123456211}\enspace $n{=}10$, $|H|{=}4$, len 7, br 2, $|\mathrm{Aut}|{=}2$, orb 9/8, hist-orb 3/2; M 25600, res 10304, rev 6950: ok}\\*
0 & -- & 0 & 4 & 1 & 190 & 2 & 832 & \texttt{-?????????}\\*
1 & 0 & 1 & 2 & 1 & 496 & 16 & 512 & \texttt{--??????--}\\
2 & 1 & 2 & 1 & 1 & 960 & 64 & 0 & \texttt{---++++---}\\
3 & 2 & 3 & 1 & 1 & 896 & 128 & 0 & \texttt{----+++---}\\
4 & 3 & 4 & 1 & 1 & 768 & 256 & 0 & \texttt{-----++---}\\
5 & 4 & 5 & 1 & 1 & 512 & 512 & 0 & \texttt{------+---}\\
6 & 5 & 6 & 1 & 1 & 0 & 1024 & 0 & \texttt{----------}\\
7 & 1 & 2 & 1 & 1 & 0 & 1024 & 0 & \texttt{----------}\\
8 & 0 & 1 & 1 & 2 & 0 & 1024 & 0 & \texttt{----------}\\
9 & 0 & 1 & 1 & 2 & 0 & 1024 & 0 & \texttt{----------}\\
\multicolumn{9}{@{}l@{}}{\rule{0pt}{2.6ex}\textbf{T10.134}\enspace \texttt{0123456123}\enspace $n{=}10$, $|H|{=}2$, len 7, br 1, $|\mathrm{Aut}|{=}1$, orb 10/10, hist-orb 2/2; M 27648, res 1260, rev 11034: ok}\\*
0 & -- & 0 & 2 & 1 & 882 & 2 & 140 & \texttt{-?????????}\\*
1 & 0 & 1 & 1 & 1 & 992 & 32 & 0 & \texttt{--+++++---}\\
2 & 1 & 2 & 1 & 1 & 960 & 64 & 0 & \texttt{---++++---}\\
3 & 2 & 3 & 1 & 1 & 896 & 128 & 0 & \texttt{----+++---}\\
4 & 3 & 4 & 1 & 1 & 768 & 256 & 0 & \texttt{-----++---}\\
5 & 4 & 5 & 1 & 1 & 512 & 512 & 0 & \texttt{------+---}\\
6 & 5 & 6 & 1 & 1 & 0 & 1024 & 0 & \texttt{----------}\\
7 & 0 & 1 & 1 & 1 & 768 & 256 & 0 & \texttt{--------++}\\
8 & 7 & 2 & 1 & 1 & 512 & 512 & 0 & \texttt{---------+}\\
9 & 8 & 3 & 1 & 1 & 0 & 1024 & 0 & \texttt{----------}\\
\multicolumn{9}{@{}l@{}}{\rule{0pt}{2.6ex}\textbf{T10.135}\enspace \texttt{0123456122}\enspace $n{=}10$, $|H|{=}3$, len 7, br 2, $|\mathrm{Aut}|{=}2$, orb 9/8, hist-orb 2/1; M 26624, res 4296, rev 9530: ok}\\*
0 & -- & 0 & 3 & 1 & 630 & 2 & 392 & \texttt{-?????????}\\*
1 & 0 & 1 & 1 & 1 & 992 & 32 & 0 & \texttt{--+++++---}\\
2 & 1 & 2 & 1 & 1 & 960 & 64 & 0 & \texttt{---++++---}\\
3 & 2 & 3 & 1 & 1 & 896 & 128 & 0 & \texttt{----+++---}\\
4 & 3 & 4 & 1 & 1 & 768 & 256 & 0 & \texttt{-----++---}\\
5 & 4 & 5 & 1 & 1 & 512 & 512 & 0 & \texttt{------+---}\\
6 & 5 & 6 & 1 & 1 & 0 & 1024 & 0 & \texttt{----------}\\
7 & 0 & 1 & 2 & 1 & 256 & 256 & 512 & \texttt{--------??}\\
8 & 7 & 2 & 1 & 2 & 0 & 1024 & 0 & \texttt{----------}\\
9 & 7 & 2 & 1 & 2 & 0 & 1024 & 0 & \texttt{----------}\\
\multicolumn{9}{@{}l@{}}{\rule{0pt}{2.6ex}\textbf{T10.136}\enspace \texttt{0123456121}\enspace $n{=}10$, $|H|{=}3$, len 7, br 1, $|\mathrm{Aut}|{=}1$, orb 10/10, hist-orb 3/3; M 25600, res 5796, rev 8546: ok}\\*
0 & -- & 0 & 3 & 1 & 378 & 2 & 644 & \texttt{-?????????}\\*
1 & 0 & 1 & 1 & 1 & 992 & 32 & 0 & \texttt{--+++++---}\\
2 & 1 & 2 & 1 & 1 & 960 & 64 & 0 & \texttt{---++++---}\\
3 & 2 & 3 & 1 & 1 & 896 & 128 & 0 & \texttt{----+++---}\\
4 & 3 & 4 & 1 & 1 & 768 & 256 & 0 & \texttt{-----++---}\\
5 & 4 & 5 & 1 & 1 & 512 & 512 & 0 & \texttt{------+---}\\
6 & 5 & 6 & 1 & 1 & 0 & 1024 & 0 & \texttt{----------}\\
7 & 0 & 1 & 1 & 1 & 512 & 512 & 0 & \texttt{--------+-}\\
8 & 7 & 2 & 1 & 1 & 0 & 1024 & 0 & \texttt{----------}\\
9 & 0 & 1 & 1 & 1 & 0 & 1024 & 0 & \texttt{----------}\\
\multicolumn{9}{@{}l@{}}{\rule{0pt}{2.6ex}\textbf{T10.137}\enspace \texttt{0123456111}\enspace $n{=}10$, $|H|{=}4$, len 7, br 1, $|\mathrm{Aut}|{=}6$, orb 8/7, hist-orb 2/1; M 24576, res 8064, rev 7050: ok}\\*
0 & -- & 0 & 4 & 1 & 126 & 2 & 896 & \texttt{-?????????}\\*
1 & 0 & 1 & 1 & 1 & 992 & 32 & 0 & \texttt{--+++++---}\\
2 & 1 & 2 & 1 & 1 & 960 & 64 & 0 & \texttt{---++++---}\\
3 & 2 & 3 & 1 & 1 & 896 & 128 & 0 & \texttt{----+++---}\\
4 & 3 & 4 & 1 & 1 & 768 & 256 & 0 & \texttt{-----++---}\\
5 & 4 & 5 & 1 & 1 & 512 & 512 & 0 & \texttt{------+---}\\
6 & 5 & 6 & 1 & 1 & 0 & 1024 & 0 & \texttt{----------}\\
7 & 0 & 1 & 1 & 3 & 0 & 1024 & 0 & \texttt{----------}\\
8 & 0 & 1 & 1 & 3 & 0 & 1024 & 0 & \texttt{----------}\\
9 & 0 & 1 & 1 & 3 & 0 & 1024 & 0 & \texttt{----------}\\
\multicolumn{9}{@{}l@{}}{\rule{0pt}{2.6ex}\textbf{T10.138}\enspace \texttt{0123455555}\enspace $n{=}10$, $|H|{=}5$, len 6, br 1, $|\mathrm{Aut}|{=}120$, orb 6/5, hist-orb 1/0; M 35840, res 9300, rev 16134: ok}\\*
0 & -- & 0 & 5 & 1 & 962 & 2 & 60 & \texttt{-++++?????}\\*
1 & 0 & 1 & 5 & 1 & 900 & 4 & 120 & \texttt{--+++?????}\\
2 & 1 & 2 & 5 & 1 & 776 & 8 & 240 & \texttt{---++?????}\\
3 & 2 & 3 & 5 & 1 & 528 & 16 & 480 & \texttt{----+?????}\\
4 & 3 & 4 & 5 & 1 & 32 & 32 & 960 & \texttt{-----?????}\\
5 & 4 & 5 & 1 & 5 & 0 & 1024 & 0 & \texttt{----------}\\
6 & 4 & 5 & 1 & 5 & 0 & 1024 & 0 & \texttt{----------}\\
7 & 4 & 5 & 1 & 5 & 0 & 1024 & 0 & \texttt{----------}\\
8 & 4 & 5 & 1 & 5 & 0 & 1024 & 0 & \texttt{----------}\\
9 & 4 & 5 & 1 & 5 & 0 & 1024 & 0 & \texttt{----------}\\
\multicolumn{9}{@{}l@{}}{\rule{0pt}{2.6ex}\textbf{T10.139}\enspace \texttt{0123455554}\enspace $n{=}10$, $|H|{=}5$, len 6, br 2, $|\mathrm{Aut}|{=}24$, orb 7/6, hist-orb 2/1; M 34816, res 11024, rev 13898: ok}\\*
0 & -- & 0 & 5 & 1 & 930 & 2 & 92 & \texttt{-+++??????}\\*
1 & 0 & 1 & 5 & 1 & 836 & 4 & 184 & \texttt{--++??????}\\
2 & 1 & 2 & 5 & 1 & 648 & 8 & 368 & \texttt{---+??????}\\
3 & 2 & 3 & 5 & 1 & 272 & 16 & 736 & \texttt{----??????}\\
4 & 3 & 4 & 4 & 1 & 64 & 64 & 896 & \texttt{-----????-}\\
5 & 4 & 5 & 1 & 4 & 0 & 1024 & 0 & \texttt{----------}\\
6 & 4 & 5 & 1 & 4 & 0 & 1024 & 0 & \texttt{----------}\\
7 & 4 & 5 & 1 & 4 & 0 & 1024 & 0 & \texttt{----------}\\
8 & 4 & 5 & 1 & 4 & 0 & 1024 & 0 & \texttt{----------}\\
9 & 3 & 4 & 1 & 1 & 0 & 1024 & 0 & \texttt{----------}\\
\multicolumn{9}{@{}l@{}}{\rule{0pt}{2.6ex}\textbf{T10.140}\enspace \texttt{0123455553}\enspace $n{=}10$, $|H|{=}5$, len 6, br 2, $|\mathrm{Aut}|{=}24$, orb 7/6, hist-orb 2/1; M 33792, res 12236, rev 12482: ok}\\*
0 & -- & 0 & 5 & 1 & 866 & 2 & 156 & \texttt{-++???????}\\*
1 & 0 & 1 & 5 & 1 & 708 & 4 & 312 & \texttt{--+???????}\\
2 & 1 & 2 & 5 & 1 & 392 & 8 & 624 & \texttt{---???????}\\
3 & 2 & 3 & 4 & 1 & 544 & 32 & 448 & \texttt{----+????-}\\
4 & 3 & 4 & 4 & 1 & 64 & 64 & 896 & \texttt{-----????-}\\
5 & 4 & 5 & 1 & 4 & 0 & 1024 & 0 & \texttt{----------}\\
6 & 4 & 5 & 1 & 4 & 0 & 1024 & 0 & \texttt{----------}\\
7 & 4 & 5 & 1 & 4 & 0 & 1024 & 0 & \texttt{----------}\\
8 & 4 & 5 & 1 & 4 & 0 & 1024 & 0 & \texttt{----------}\\
9 & 2 & 3 & 1 & 1 & 0 & 1024 & 0 & \texttt{----------}\\
\multicolumn{9}{@{}l@{}}{\rule{0pt}{2.6ex}\textbf{T10.141}\enspace \texttt{0123455552}\enspace $n{=}10$, $|H|{=}5$, len 6, br 2, $|\mathrm{Aut}|{=}24$, orb 7/6, hist-orb 2/1; M 32768, res 12584, rev 11738: ok}\\*
0 & -- & 0 & 5 & 1 & 738 & 2 & 284 & \texttt{-+????????}\\*
1 & 0 & 1 & 5 & 1 & 452 & 4 & 568 & \texttt{--????????}\\
2 & 1 & 2 & 4 & 1 & 784 & 16 & 224 & \texttt{---++????-}\\
3 & 2 & 3 & 4 & 1 & 544 & 32 & 448 & \texttt{----+????-}\\
4 & 3 & 4 & 4 & 1 & 64 & 64 & 896 & \texttt{-----????-}\\
5 & 4 & 5 & 1 & 4 & 0 & 1024 & 0 & \texttt{----------}\\
6 & 4 & 5 & 1 & 4 & 0 & 1024 & 0 & \texttt{----------}\\
7 & 4 & 5 & 1 & 4 & 0 & 1024 & 0 & \texttt{----------}\\
8 & 4 & 5 & 1 & 4 & 0 & 1024 & 0 & \texttt{----------}\\
9 & 1 & 2 & 1 & 1 & 0 & 1024 & 0 & \texttt{----------}\\
\multicolumn{9}{@{}l@{}}{\rule{0pt}{2.6ex}\textbf{T10.142}\enspace \texttt{0123455551}\enspace $n{=}10$, $|H|{=}5$, len 6, br 2, $|\mathrm{Aut}|{=}24$, orb 7/6, hist-orb 2/1; M 31744, res 11356, rev 11830: ok}\\*
0 & -- & 0 & 5 & 1 & 482 & 2 & 540 & \texttt{-?????????}\\*
1 & 0 & 1 & 4 & 1 & 904 & 8 & 112 & \texttt{--+++????-}\\
2 & 1 & 2 & 4 & 1 & 784 & 16 & 224 & \texttt{---++????-}\\
3 & 2 & 3 & 4 & 1 & 544 & 32 & 448 & \texttt{----+????-}\\
4 & 3 & 4 & 4 & 1 & 64 & 64 & 896 & \texttt{-----????-}\\
5 & 4 & 5 & 1 & 4 & 0 & 1024 & 0 & \texttt{----------}\\
6 & 4 & 5 & 1 & 4 & 0 & 1024 & 0 & \texttt{----------}\\
7 & 4 & 5 & 1 & 4 & 0 & 1024 & 0 & \texttt{----------}\\
8 & 4 & 5 & 1 & 4 & 0 & 1024 & 0 & \texttt{----------}\\
9 & 0 & 1 & 1 & 1 & 0 & 1024 & 0 & \texttt{----------}\\
\multicolumn{9}{@{}l@{}}{\rule{0pt}{2.6ex}\textbf{T10.143}\enspace \texttt{0123455545}\enspace $n{=}10$, $|H|{=}4$, len 6, br 2, $|\mathrm{Aut}|{=}6$, orb 8/7, hist-orb 2/1; M 34816, res 8064, rev 14606: ok}\\*
0 & -- & 0 & 4 & 1 & 950 & 2 & 72 & \texttt{-+++??????}\\*
1 & 0 & 1 & 4 & 1 & 876 & 4 & 144 & \texttt{--++??????}\\
2 & 1 & 2 & 4 & 1 & 728 & 8 & 288 & \texttt{---+??????}\\
3 & 2 & 3 & 4 & 1 & 432 & 16 & 576 & \texttt{----??????}\\
4 & 3 & 4 & 3 & 1 & 128 & 128 & 768 & \texttt{-----???--}\\
5 & 4 & 5 & 1 & 3 & 0 & 1024 & 0 & \texttt{----------}\\
6 & 4 & 5 & 1 & 3 & 0 & 1024 & 0 & \texttt{----------}\\
7 & 4 & 5 & 1 & 3 & 0 & 1024 & 0 & \texttt{----------}\\
8 & 3 & 4 & 1 & 1 & 512 & 512 & 0 & \texttt{---------+}\\
9 & 8 & 5 & 1 & 1 & 0 & 1024 & 0 & \texttt{----------}\\
\multicolumn{9}{@{}l@{}}{\rule{0pt}{2.6ex}\textbf{T10.144}\enspace \texttt{0123455544}\enspace $n{=}10$, $|H|{=}5$, len 6, br 2, $|\mathrm{Aut}|{=}12$, orb 7/5, hist-orb 2/0; M 33792, res 11304, rev 12962: ok}\\*
0 & -- & 0 & 5 & 1 & 914 & 2 & 108 & \texttt{-+++??????}\\*
1 & 0 & 1 & 5 & 1 & 804 & 4 & 216 & \texttt{--++??????}\\
2 & 1 & 2 & 5 & 1 & 584 & 8 & 432 & \texttt{---+??????}\\
3 & 2 & 3 & 5 & 1 & 144 & 16 & 864 & \texttt{----??????}\\
4 & 3 & 4 & 3 & 1 & 128 & 128 & 768 & \texttt{-----???--}\\
5 & 4 & 5 & 1 & 3 & 0 & 1024 & 0 & \texttt{----------}\\
6 & 4 & 5 & 1 & 3 & 0 & 1024 & 0 & \texttt{----------}\\
7 & 4 & 5 & 1 & 3 & 0 & 1024 & 0 & \texttt{----------}\\
8 & 3 & 4 & 1 & 2 & 0 & 1024 & 0 & \texttt{----------}\\
9 & 3 & 4 & 1 & 2 & 0 & 1024 & 0 & \texttt{----------}\\
\multicolumn{9}{@{}l@{}}{\rule{0pt}{2.6ex}\textbf{T10.145}\enspace \texttt{0123455543}\enspace $n{=}10$, $|H|{=}5$, len 6, br 3, $|\mathrm{Aut}|{=}6$, orb 8/7, hist-orb 3/2; M 32768, res 12916, rev 11134: ok}\\*
0 & -- & 0 & 5 & 1 & 850 & 2 & 172 & \texttt{-++???????}\\*
1 & 0 & 1 & 5 & 1 & 676 & 4 & 344 & \texttt{--+???????}\\
2 & 1 & 2 & 5 & 1 & 328 & 8 & 688 & \texttt{---???????}\\
3 & 2 & 3 & 4 & 1 & 288 & 32 & 704 & \texttt{----?????-}\\
4 & 3 & 4 & 3 & 1 & 128 & 128 & 768 & \texttt{-----???--}\\
5 & 4 & 5 & 1 & 3 & 0 & 1024 & 0 & \texttt{----------}\\
6 & 4 & 5 & 1 & 3 & 0 & 1024 & 0 & \texttt{----------}\\
7 & 4 & 5 & 1 & 3 & 0 & 1024 & 0 & \texttt{----------}\\
8 & 3 & 4 & 1 & 1 & 0 & 1024 & 0 & \texttt{----------}\\
9 & 2 & 3 & 1 & 1 & 0 & 1024 & 0 & \texttt{----------}\\
\multicolumn{9}{@{}l@{}}{\rule{0pt}{2.6ex}\textbf{T10.146}\enspace \texttt{0123455542}\enspace $n{=}10$, $|H|{=}5$, len 6, br 3, $|\mathrm{Aut}|{=}6$, orb 8/7, hist-orb 3/2; M 31744, res 13536, rev 10190: ok}\\*
0 & -- & 0 & 5 & 1 & 722 & 2 & 300 & \texttt{-+????????}\\*
1 & 0 & 1 & 5 & 1 & 420 & 4 & 600 & \texttt{--????????}\\
2 & 1 & 2 & 4 & 1 & 656 & 16 & 352 & \texttt{---+?????-}\\
3 & 2 & 3 & 4 & 1 & 288 & 32 & 704 & \texttt{----?????-}\\
4 & 3 & 4 & 3 & 1 & 128 & 128 & 768 & \texttt{-----???--}\\
5 & 4 & 5 & 1 & 3 & 0 & 1024 & 0 & \texttt{----------}\\
6 & 4 & 5 & 1 & 3 & 0 & 1024 & 0 & \texttt{----------}\\
7 & 4 & 5 & 1 & 3 & 0 & 1024 & 0 & \texttt{----------}\\
8 & 3 & 4 & 1 & 1 & 0 & 1024 & 0 & \texttt{----------}\\
9 & 1 & 2 & 1 & 1 & 0 & 1024 & 0 & \texttt{----------}\\
\multicolumn{9}{@{}l@{}}{\rule{0pt}{2.6ex}\textbf{T10.147}\enspace \texttt{0123455541}\enspace $n{=}10$, $|H|{=}5$, len 6, br 3, $|\mathrm{Aut}|{=}6$, orb 8/7, hist-orb 3/2; M 30720, res 12484, rev 10194: ok}\\*
0 & -- & 0 & 5 & 1 & 466 & 2 & 556 & \texttt{-?????????}\\*
1 & 0 & 1 & 4 & 1 & 840 & 8 & 176 & \texttt{--++?????-}\\
2 & 1 & 2 & 4 & 1 & 656 & 16 & 352 & \texttt{---+?????-}\\
3 & 2 & 3 & 4 & 1 & 288 & 32 & 704 & \texttt{----?????-}\\
4 & 3 & 4 & 3 & 1 & 128 & 128 & 768 & \texttt{-----???--}\\
5 & 4 & 5 & 1 & 3 & 0 & 1024 & 0 & \texttt{----------}\\
6 & 4 & 5 & 1 & 3 & 0 & 1024 & 0 & \texttt{----------}\\
7 & 4 & 5 & 1 & 3 & 0 & 1024 & 0 & \texttt{----------}\\
8 & 3 & 4 & 1 & 1 & 0 & 1024 & 0 & \texttt{----------}\\
9 & 0 & 1 & 1 & 1 & 0 & 1024 & 0 & \texttt{----------}\\
\multicolumn{9}{@{}l@{}}{\rule{0pt}{2.6ex}\textbf{T10.148}\enspace \texttt{0123455534}\enspace $n{=}10$, $|H|{=}4$, len 6, br 2, $|\mathrm{Aut}|{=}6$, orb 8/7, hist-orb 2/1; M 32768, res 7880, rev 13418: ok}\\*
0 & -- & 0 & 4 & 1 & 918 & 2 & 104 & \texttt{-++???????}\\*
1 & 0 & 1 & 4 & 1 & 812 & 4 & 208 & \texttt{--+???????}\\
2 & 1 & 2 & 4 & 1 & 600 & 8 & 416 & \texttt{---???????}\\
3 & 2 & 3 & 3 & 1 & 576 & 64 & 384 & \texttt{----+???--}\\
4 & 3 & 4 & 3 & 1 & 128 & 128 & 768 & \texttt{-----???--}\\
5 & 4 & 5 & 1 & 3 & 0 & 1024 & 0 & \texttt{----------}\\
6 & 4 & 5 & 1 & 3 & 0 & 1024 & 0 & \texttt{----------}\\
7 & 4 & 5 & 1 & 3 & 0 & 1024 & 0 & \texttt{----------}\\
8 & 2 & 3 & 1 & 1 & 512 & 512 & 0 & \texttt{---------+}\\
9 & 8 & 4 & 1 & 1 & 0 & 1024 & 0 & \texttt{----------}\\
\multicolumn{9}{@{}l@{}}{\rule{0pt}{2.6ex}\textbf{T10.149}\enspace \texttt{0123455533}\enspace $n{=}10$, $|H|{=}5$, len 6, br 2, $|\mathrm{Aut}|{=}12$, orb 7/5, hist-orb 2/0; M 31744, res 12780, rev 10634: ok}\\*
0 & -- & 0 & 5 & 1 & 818 & 2 & 204 & \texttt{-++???????}\\*
1 & 0 & 1 & 5 & 1 & 612 & 4 & 408 & \texttt{--+???????}\\
2 & 1 & 2 & 5 & 1 & 200 & 8 & 816 & \texttt{---???????}\\
3 & 2 & 3 & 3 & 1 & 576 & 64 & 384 & \texttt{----+???--}\\
4 & 3 & 4 & 3 & 1 & 128 & 128 & 768 & \texttt{-----???--}\\
5 & 4 & 5 & 1 & 3 & 0 & 1024 & 0 & \texttt{----------}\\
6 & 4 & 5 & 1 & 3 & 0 & 1024 & 0 & \texttt{----------}\\
7 & 4 & 5 & 1 & 3 & 0 & 1024 & 0 & \texttt{----------}\\
8 & 2 & 3 & 1 & 2 & 0 & 1024 & 0 & \texttt{----------}\\
9 & 2 & 3 & 1 & 2 & 0 & 1024 & 0 & \texttt{----------}\\
\multicolumn{9}{@{}l@{}}{\rule{0pt}{2.6ex}\textbf{T10.150}\enspace \texttt{0123455532}\enspace $n{=}10$, $|H|{=}5$, len 6, br 3, $|\mathrm{Aut}|{=}6$, orb 8/7, hist-orb 3/2; M 30720, res 13944, rev 9278: ok}\\*
0 & -- & 0 & 5 & 1 & 690 & 2 & 332 & \texttt{-+????????}\\*
1 & 0 & 1 & 5 & 1 & 356 & 4 & 664 & \texttt{--????????}\\
2 & 1 & 2 & 4 & 1 & 400 & 16 & 608 & \texttt{---??????-}\\
3 & 2 & 3 & 3 & 1 & 576 & 64 & 384 & \texttt{----+???--}\\
4 & 3 & 4 & 3 & 1 & 128 & 128 & 768 & \texttt{-----???--}\\
5 & 4 & 5 & 1 & 3 & 0 & 1024 & 0 & \texttt{----------}\\
6 & 4 & 5 & 1 & 3 & 0 & 1024 & 0 & \texttt{----------}\\
7 & 4 & 5 & 1 & 3 & 0 & 1024 & 0 & \texttt{----------}\\
8 & 2 & 3 & 1 & 1 & 0 & 1024 & 0 & \texttt{----------}\\
9 & 1 & 2 & 1 & 1 & 0 & 1024 & 0 & \texttt{----------}\\
\multicolumn{9}{@{}l@{}}{\rule{0pt}{2.6ex}\textbf{T10.151}\enspace \texttt{0123455531}\enspace $n{=}10$, $|H|{=}5$, len 6, br 3, $|\mathrm{Aut}|{=}6$, orb 8/7, hist-orb 3/2; M 29696, res 13244, rev 9082: ok}\\*
0 & -- & 0 & 5 & 1 & 434 & 2 & 588 & \texttt{-?????????}\\*
1 & 0 & 1 & 4 & 1 & 712 & 8 & 304 & \texttt{--+??????-}\\
2 & 1 & 2 & 4 & 1 & 400 & 16 & 608 & \texttt{---??????-}\\
3 & 2 & 3 & 3 & 1 & 576 & 64 & 384 & \texttt{----+???--}\\
4 & 3 & 4 & 3 & 1 & 128 & 128 & 768 & \texttt{-----???--}\\
5 & 4 & 5 & 1 & 3 & 0 & 1024 & 0 & \texttt{----------}\\
6 & 4 & 5 & 1 & 3 & 0 & 1024 & 0 & \texttt{----------}\\
7 & 4 & 5 & 1 & 3 & 0 & 1024 & 0 & \texttt{----------}\\
8 & 2 & 3 & 1 & 1 & 0 & 1024 & 0 & \texttt{----------}\\
9 & 0 & 1 & 1 & 1 & 0 & 1024 & 0 & \texttt{----------}\\
\multicolumn{9}{@{}l@{}}{\rule{0pt}{2.6ex}\textbf{T10.152}\enspace \texttt{0123455523}\enspace $n{=}10$, $|H|{=}4$, len 6, br 2, $|\mathrm{Aut}|{=}6$, orb 8/7, hist-orb 2/1; M 30720, res 7632, rev 12278: ok}\\*
0 & -- & 0 & 4 & 1 & 854 & 2 & 168 & \texttt{-+????????}\\*
1 & 0 & 1 & 4 & 1 & 684 & 4 & 336 & \texttt{--????????}\\
2 & 1 & 2 & 3 & 1 & 800 & 32 & 192 & \texttt{---++???--}\\
3 & 2 & 3 & 3 & 1 & 576 & 64 & 384 & \texttt{----+???--}\\
4 & 3 & 4 & 3 & 1 & 128 & 128 & 768 & \texttt{-----???--}\\
5 & 4 & 5 & 1 & 3 & 0 & 1024 & 0 & \texttt{----------}\\
6 & 4 & 5 & 1 & 3 & 0 & 1024 & 0 & \texttt{----------}\\
7 & 4 & 5 & 1 & 3 & 0 & 1024 & 0 & \texttt{----------}\\
8 & 1 & 2 & 1 & 1 & 512 & 512 & 0 & \texttt{---------+}\\
9 & 8 & 3 & 1 & 1 & 0 & 1024 & 0 & \texttt{----------}\\
\multicolumn{9}{@{}l@{}}{\rule{0pt}{2.6ex}\textbf{T10.153}\enspace \texttt{0123455522}\enspace $n{=}10$, $|H|{=}5$, len 6, br 2, $|\mathrm{Aut}|{=}12$, orb 7/5, hist-orb 2/0; M 29696, res 13104, rev 9154: ok}\\*
0 & -- & 0 & 5 & 1 & 626 & 2 & 396 & \texttt{-+????????}\\*
1 & 0 & 1 & 5 & 1 & 228 & 4 & 792 & \texttt{--????????}\\
2 & 1 & 2 & 3 & 1 & 800 & 32 & 192 & \texttt{---++???--}\\
3 & 2 & 3 & 3 & 1 & 576 & 64 & 384 & \texttt{----+???--}\\
4 & 3 & 4 & 3 & 1 & 128 & 128 & 768 & \texttt{-----???--}\\
5 & 4 & 5 & 1 & 3 & 0 & 1024 & 0 & \texttt{----------}\\
6 & 4 & 5 & 1 & 3 & 0 & 1024 & 0 & \texttt{----------}\\
7 & 4 & 5 & 1 & 3 & 0 & 1024 & 0 & \texttt{----------}\\
8 & 1 & 2 & 1 & 2 & 0 & 1024 & 0 & \texttt{----------}\\
9 & 1 & 2 & 1 & 2 & 0 & 1024 & 0 & \texttt{----------}\\
\multicolumn{9}{@{}l@{}}{\rule{0pt}{2.6ex}\textbf{T10.154}\enspace \texttt{0123455521}\enspace $n{=}10$, $|H|{=}5$, len 6, br 3, $|\mathrm{Aut}|{=}6$, orb 8/7, hist-orb 3/2; M 28672, res 13108, rev 8546: ok}\\*
0 & -- & 0 & 5 & 1 & 370 & 2 & 652 & \texttt{-?????????}\\*
1 & 0 & 1 & 4 & 1 & 456 & 8 & 560 & \texttt{--???????-}\\
2 & 1 & 2 & 3 & 1 & 800 & 32 & 192 & \texttt{---++???--}\\
3 & 2 & 3 & 3 & 1 & 576 & 64 & 384 & \texttt{----+???--}\\
4 & 3 & 4 & 3 & 1 & 128 & 128 & 768 & \texttt{-----???--}\\
5 & 4 & 5 & 1 & 3 & 0 & 1024 & 0 & \texttt{----------}\\
6 & 4 & 5 & 1 & 3 & 0 & 1024 & 0 & \texttt{----------}\\
7 & 4 & 5 & 1 & 3 & 0 & 1024 & 0 & \texttt{----------}\\
8 & 1 & 2 & 1 & 1 & 0 & 1024 & 0 & \texttt{----------}\\
9 & 0 & 1 & 1 & 1 & 0 & 1024 & 0 & \texttt{----------}\\
\multicolumn{9}{@{}l@{}}{\rule{0pt}{2.6ex}\textbf{T10.155}\enspace \texttt{0123455512}\enspace $n{=}10$, $|H|{=}4$, len 6, br 2, $|\mathrm{Aut}|{=}6$, orb 8/7, hist-orb 2/1; M 28672, res 6792, rev 11430: ok}\\*
0 & -- & 0 & 4 & 1 & 726 & 2 & 296 & \texttt{-?????????}\\*
1 & 0 & 1 & 3 & 1 & 912 & 16 & 96 & \texttt{--+++???--}\\
2 & 1 & 2 & 3 & 1 & 800 & 32 & 192 & \texttt{---++???--}\\
3 & 2 & 3 & 3 & 1 & 576 & 64 & 384 & \texttt{----+???--}\\
4 & 3 & 4 & 3 & 1 & 128 & 128 & 768 & \texttt{-----???--}\\
5 & 4 & 5 & 1 & 3 & 0 & 1024 & 0 & \texttt{----------}\\
6 & 4 & 5 & 1 & 3 & 0 & 1024 & 0 & \texttt{----------}\\
7 & 4 & 5 & 1 & 3 & 0 & 1024 & 0 & \texttt{----------}\\
8 & 0 & 1 & 1 & 1 & 512 & 512 & 0 & \texttt{---------+}\\
9 & 8 & 2 & 1 & 1 & 0 & 1024 & 0 & \texttt{----------}\\
\multicolumn{9}{@{}l@{}}{\rule{0pt}{2.6ex}\textbf{T10.156}\enspace \texttt{0123455511}\enspace $n{=}10$, $|H|{=}5$, len 6, br 2, $|\mathrm{Aut}|{=}12$, orb 7/5, hist-orb 2/0; M 27648, res 11148, rev 8878: ok}\\*
0 & -- & 0 & 5 & 1 & 242 & 2 & 780 & \texttt{-?????????}\\*
1 & 0 & 1 & 3 & 1 & 912 & 16 & 96 & \texttt{--+++???--}\\
2 & 1 & 2 & 3 & 1 & 800 & 32 & 192 & \texttt{---++???--}\\
3 & 2 & 3 & 3 & 1 & 576 & 64 & 384 & \texttt{----+???--}\\
4 & 3 & 4 & 3 & 1 & 128 & 128 & 768 & \texttt{-----???--}\\
5 & 4 & 5 & 1 & 3 & 0 & 1024 & 0 & \texttt{----------}\\
6 & 4 & 5 & 1 & 3 & 0 & 1024 & 0 & \texttt{----------}\\
7 & 4 & 5 & 1 & 3 & 0 & 1024 & 0 & \texttt{----------}\\
8 & 0 & 1 & 1 & 2 & 0 & 1024 & 0 & \texttt{----------}\\
9 & 0 & 1 & 1 & 2 & 0 & 1024 & 0 & \texttt{----------}\\
\multicolumn{9}{@{}l@{}}{\rule{0pt}{2.6ex}\textbf{T10.157}\enspace \texttt{0123455455}\enspace $n{=}10$, $|H|{=}4$, len 6, br 3, $|\mathrm{Aut}|{=}8$, orb 6/4, hist-orb 1/0; M 34816, res 7928, rev 14162: ok}\\*
0 & -- & 0 & 4 & 1 & 946 & 2 & 76 & \texttt{-+++??????}\\*
1 & 0 & 1 & 4 & 1 & 868 & 4 & 152 & \texttt{--++??????}\\
2 & 1 & 2 & 4 & 1 & 712 & 8 & 304 & \texttt{---+??????}\\
3 & 2 & 3 & 4 & 1 & 400 & 16 & 608 & \texttt{----??????}\\
4 & 3 & 4 & 2 & 2 & 256 & 256 & 512 & \texttt{-----??---}\\
5 & 4 & 5 & 1 & 4 & 0 & 1024 & 0 & \texttt{----------}\\
6 & 4 & 5 & 1 & 4 & 0 & 1024 & 0 & \texttt{----------}\\
7 & 3 & 4 & 2 & 2 & 256 & 256 & 512 & \texttt{--------??}\\
8 & 7 & 5 & 1 & 4 & 0 & 1024 & 0 & \texttt{----------}\\
9 & 7 & 5 & 1 & 4 & 0 & 1024 & 0 & \texttt{----------}\\
\multicolumn{9}{@{}l@{}}{\rule{0pt}{2.6ex}\textbf{T10.158}\enspace \texttt{0123455454}\enspace $n{=}10$, $|H|{=}4$, len 6, br 2, $|\mathrm{Aut}|{=}2$, orb 9/8, hist-orb 3/2; M 33792, res 9184, rev 13446: ok}\\*
0 & -- & 0 & 4 & 1 & 926 & 2 & 96 & \texttt{-+++??????}\\*
1 & 0 & 1 & 4 & 1 & 828 & 4 & 192 & \texttt{--++??????}\\
2 & 1 & 2 & 4 & 1 & 632 & 8 & 384 & \texttt{---+??????}\\
3 & 2 & 3 & 4 & 1 & 240 & 16 & 768 & \texttt{----??????}\\
4 & 3 & 4 & 2 & 1 & 256 & 256 & 512 & \texttt{-----??---}\\
5 & 4 & 5 & 1 & 2 & 0 & 1024 & 0 & \texttt{----------}\\
6 & 4 & 5 & 1 & 2 & 0 & 1024 & 0 & \texttt{----------}\\
7 & 3 & 4 & 1 & 1 & 512 & 512 & 0 & \texttt{--------+-}\\
8 & 7 & 5 & 1 & 1 & 0 & 1024 & 0 & \texttt{----------}\\
9 & 3 & 4 & 1 & 1 & 0 & 1024 & 0 & \texttt{----------}\\
\multicolumn{9}{@{}l@{}}{\rule{0pt}{2.6ex}\textbf{T10.159}\enspace \texttt{0123455453}\enspace $n{=}10$, $|H|{=}4$, len 6, br 3, $|\mathrm{Aut}|{=}2$, orb 9/8, hist-orb 3/2; M 32768, res 10496, rev 11946: ok}\\*
0 & -- & 0 & 4 & 1 & 862 & 2 & 160 & \texttt{-++???????}\\*
1 & 0 & 1 & 4 & 1 & 700 & 4 & 320 & \texttt{--+???????}\\
2 & 1 & 2 & 4 & 1 & 376 & 8 & 640 & \texttt{---???????}\\
3 & 2 & 3 & 3 & 1 & 480 & 32 & 512 & \texttt{----?????-}\\
4 & 3 & 4 & 2 & 1 & 256 & 256 & 512 & \texttt{-----??---}\\
5 & 4 & 5 & 1 & 2 & 0 & 1024 & 0 & \texttt{----------}\\
6 & 4 & 5 & 1 & 2 & 0 & 1024 & 0 & \texttt{----------}\\
7 & 3 & 4 & 1 & 1 & 512 & 512 & 0 & \texttt{--------+-}\\
8 & 7 & 5 & 1 & 1 & 0 & 1024 & 0 & \texttt{----------}\\
9 & 2 & 3 & 1 & 1 & 0 & 1024 & 0 & \texttt{----------}\\
\multicolumn{9}{@{}l@{}}{\rule{0pt}{2.6ex}\textbf{T10.160}\enspace \texttt{0123455452}\enspace $n{=}10$, $|H|{=}4$, len 6, br 3, $|\mathrm{Aut}|{=}2$, orb 9/8, hist-orb 3/2; M 31744, res 10912, rev 11130: ok}\\*
0 & -- & 0 & 4 & 1 & 734 & 2 & 288 & \texttt{-+????????}\\*
1 & 0 & 1 & 4 & 1 & 444 & 4 & 576 & \texttt{--????????}\\
2 & 1 & 2 & 3 & 1 & 752 & 16 & 256 & \texttt{---+?????-}\\
3 & 2 & 3 & 3 & 1 & 480 & 32 & 512 & \texttt{----?????-}\\
4 & 3 & 4 & 2 & 1 & 256 & 256 & 512 & \texttt{-----??---}\\
5 & 4 & 5 & 1 & 2 & 0 & 1024 & 0 & \texttt{----------}\\
6 & 4 & 5 & 1 & 2 & 0 & 1024 & 0 & \texttt{----------}\\
7 & 3 & 4 & 1 & 1 & 512 & 512 & 0 & \texttt{--------+-}\\
8 & 7 & 5 & 1 & 1 & 0 & 1024 & 0 & \texttt{----------}\\
9 & 1 & 2 & 1 & 1 & 0 & 1024 & 0 & \texttt{----------}\\
\multicolumn{9}{@{}l@{}}{\rule{0pt}{2.6ex}\textbf{T10.161}\enspace \texttt{0123455451}\enspace $n{=}10$, $|H|{=}4$, len 6, br 3, $|\mathrm{Aut}|{=}2$, orb 9/8, hist-orb 3/2; M 30720, res 9728, rev 11126: ok}\\*
0 & -- & 0 & 4 & 1 & 478 & 2 & 544 & \texttt{-?????????}\\*
1 & 0 & 1 & 3 & 1 & 888 & 8 & 128 & \texttt{--++?????-}\\
2 & 1 & 2 & 3 & 1 & 752 & 16 & 256 & \texttt{---+?????-}\\
3 & 2 & 3 & 3 & 1 & 480 & 32 & 512 & \texttt{----?????-}\\
4 & 3 & 4 & 2 & 1 & 256 & 256 & 512 & \texttt{-----??---}\\
5 & 4 & 5 & 1 & 2 & 0 & 1024 & 0 & \texttt{----------}\\
6 & 4 & 5 & 1 & 2 & 0 & 1024 & 0 & \texttt{----------}\\
7 & 3 & 4 & 1 & 1 & 512 & 512 & 0 & \texttt{--------+-}\\
8 & 7 & 5 & 1 & 1 & 0 & 1024 & 0 & \texttt{----------}\\
9 & 0 & 1 & 1 & 1 & 0 & 1024 & 0 & \texttt{----------}\\
\multicolumn{9}{@{}l@{}}{\rule{0pt}{2.6ex}\textbf{T10.162}\enspace \texttt{0123455444}\enspace $n{=}10$, $|H|{=}5$, len 6, br 2, $|\mathrm{Aut}|{=}12$, orb 7/5, hist-orb 2/0; M 32768, res 10984, rev 12762: ok}\\*
0 & -- & 0 & 5 & 1 & 906 & 2 & 116 & \texttt{-+++??????}\\*
1 & 0 & 1 & 5 & 1 & 788 & 4 & 232 & \texttt{--++??????}\\
2 & 1 & 2 & 5 & 1 & 552 & 8 & 464 & \texttt{---+??????}\\
3 & 2 & 3 & 5 & 1 & 80 & 16 & 928 & \texttt{----??????}\\
4 & 3 & 4 & 2 & 1 & 256 & 256 & 512 & \texttt{-----??---}\\
5 & 4 & 5 & 1 & 2 & 0 & 1024 & 0 & \texttt{----------}\\
6 & 4 & 5 & 1 & 2 & 0 & 1024 & 0 & \texttt{----------}\\
7 & 3 & 4 & 1 & 3 & 0 & 1024 & 0 & \texttt{----------}\\
8 & 3 & 4 & 1 & 3 & 0 & 1024 & 0 & \texttt{----------}\\
9 & 3 & 4 & 1 & 3 & 0 & 1024 & 0 & \texttt{----------}\\
\multicolumn{9}{@{}l@{}}{\rule{0pt}{2.6ex}\textbf{T10.163}\enspace \texttt{0123455443}\enspace $n{=}10$, $|H|{=}5$, len 6, br 3, $|\mathrm{Aut}|{=}4$, orb 8/6, hist-orb 3/1; M 31744, res 12796, rev 10686: ok}\\*
0 & -- & 0 & 5 & 1 & 842 & 2 & 180 & \texttt{-++???????}\\*
1 & 0 & 1 & 5 & 1 & 660 & 4 & 360 & \texttt{--+???????}\\
2 & 1 & 2 & 5 & 1 & 296 & 8 & 720 & \texttt{---???????}\\
3 & 2 & 3 & 4 & 1 & 160 & 32 & 832 & \texttt{----?????-}\\
4 & 3 & 4 & 2 & 1 & 256 & 256 & 512 & \texttt{-----??---}\\
5 & 4 & 5 & 1 & 2 & 0 & 1024 & 0 & \texttt{----------}\\
6 & 4 & 5 & 1 & 2 & 0 & 1024 & 0 & \texttt{----------}\\
7 & 3 & 4 & 1 & 2 & 0 & 1024 & 0 & \texttt{----------}\\
8 & 3 & 4 & 1 & 2 & 0 & 1024 & 0 & \texttt{----------}\\
9 & 2 & 3 & 1 & 1 & 0 & 1024 & 0 & \texttt{----------}\\
\multicolumn{9}{@{}l@{}}{\rule{0pt}{2.6ex}\textbf{T10.164}\enspace \texttt{0123455442}\enspace $n{=}10$, $|H|{=}5$, len 6, br 3, $|\mathrm{Aut}|{=}4$, orb 8/6, hist-orb 3/1; M 30720, res 13552, rev 9630: ok}\\*
0 & -- & 0 & 5 & 1 & 714 & 2 & 308 & \texttt{-+????????}\\*
1 & 0 & 1 & 5 & 1 & 404 & 4 & 616 & \texttt{--????????}\\
2 & 1 & 2 & 4 & 1 & 592 & 16 & 416 & \texttt{---+?????-}\\
3 & 2 & 3 & 4 & 1 & 160 & 32 & 832 & \texttt{----?????-}\\
4 & 3 & 4 & 2 & 1 & 256 & 256 & 512 & \texttt{-----??---}\\
5 & 4 & 5 & 1 & 2 & 0 & 1024 & 0 & \texttt{----------}\\
6 & 4 & 5 & 1 & 2 & 0 & 1024 & 0 & \texttt{----------}\\
7 & 3 & 4 & 1 & 2 & 0 & 1024 & 0 & \texttt{----------}\\
8 & 3 & 4 & 1 & 2 & 0 & 1024 & 0 & \texttt{----------}\\
9 & 1 & 2 & 1 & 1 & 0 & 1024 & 0 & \texttt{----------}\\
\multicolumn{9}{@{}l@{}}{\rule{0pt}{2.6ex}\textbf{T10.165}\enspace \texttt{0123455441}\enspace $n{=}10$, $|H|{=}5$, len 6, br 3, $|\mathrm{Aut}|{=}4$, orb 8/6, hist-orb 3/1; M 29696, res 12588, rev 9602: ok}\\*
0 & -- & 0 & 5 & 1 & 458 & 2 & 564 & \texttt{-?????????}\\*
1 & 0 & 1 & 4 & 1 & 808 & 8 & 208 & \texttt{--++?????-}\\
2 & 1 & 2 & 4 & 1 & 592 & 16 & 416 & \texttt{---+?????-}\\
3 & 2 & 3 & 4 & 1 & 160 & 32 & 832 & \texttt{----?????-}\\
4 & 3 & 4 & 2 & 1 & 256 & 256 & 512 & \texttt{-----??---}\\
5 & 4 & 5 & 1 & 2 & 0 & 1024 & 0 & \texttt{----------}\\
6 & 4 & 5 & 1 & 2 & 0 & 1024 & 0 & \texttt{----------}\\
7 & 3 & 4 & 1 & 2 & 0 & 1024 & 0 & \texttt{----------}\\
8 & 3 & 4 & 1 & 2 & 0 & 1024 & 0 & \texttt{----------}\\
9 & 0 & 1 & 1 & 1 & 0 & 1024 & 0 & \texttt{----------}\\
\multicolumn{9}{@{}l@{}}{\rule{0pt}{2.6ex}\textbf{T10.166}\enspace \texttt{0123455434}\enspace $n{=}10$, $|H|{=}4$, len 6, br 3, $|\mathrm{Aut}|{=}2$, orb 9/8, hist-orb 3/2; M 31744, res 8816, rev 12294: ok}\\*
0 & -- & 0 & 4 & 1 & 894 & 2 & 128 & \texttt{-++???????}\\*
1 & 0 & 1 & 4 & 1 & 764 & 4 & 256 & \texttt{--+???????}\\
2 & 1 & 2 & 4 & 1 & 504 & 8 & 512 & \texttt{---???????}\\
3 & 2 & 3 & 3 & 1 & 320 & 64 & 640 & \texttt{----????--}\\
4 & 3 & 4 & 2 & 1 & 256 & 256 & 512 & \texttt{-----??---}\\
5 & 4 & 5 & 1 & 2 & 0 & 1024 & 0 & \texttt{----------}\\
6 & 4 & 5 & 1 & 2 & 0 & 1024 & 0 & \texttt{----------}\\
7 & 3 & 4 & 1 & 1 & 0 & 1024 & 0 & \texttt{----------}\\
8 & 2 & 3 & 1 & 1 & 512 & 512 & 0 & \texttt{---------+}\\
9 & 8 & 4 & 1 & 1 & 0 & 1024 & 0 & \texttt{----------}\\
\multicolumn{9}{@{}l@{}}{\rule{0pt}{2.6ex}\textbf{T10.167}\enspace \texttt{0123455433}\enspace $n{=}10$, $|H|{=}5$, len 6, br 3, $|\mathrm{Aut}|{=}4$, orb 8/6, hist-orb 3/1; M 30720, res 12932, rev 9902: ok}\\*
0 & -- & 0 & 5 & 1 & 810 & 2 & 212 & \texttt{-++???????}\\*
1 & 0 & 1 & 5 & 1 & 596 & 4 & 424 & \texttt{--+???????}\\
2 & 1 & 2 & 5 & 1 & 168 & 8 & 848 & \texttt{---???????}\\
3 & 2 & 3 & 3 & 1 & 320 & 64 & 640 & \texttt{----????--}\\
4 & 3 & 4 & 2 & 1 & 256 & 256 & 512 & \texttt{-----??---}\\
5 & 4 & 5 & 1 & 2 & 0 & 1024 & 0 & \texttt{----------}\\
6 & 4 & 5 & 1 & 2 & 0 & 1024 & 0 & \texttt{----------}\\
7 & 3 & 4 & 1 & 1 & 0 & 1024 & 0 & \texttt{----------}\\
8 & 2 & 3 & 1 & 2 & 0 & 1024 & 0 & \texttt{----------}\\
9 & 2 & 3 & 1 & 2 & 0 & 1024 & 0 & \texttt{----------}\\
\multicolumn{9}{@{}l@{}}{\rule{0pt}{2.6ex}\textbf{T10.168}\enspace \texttt{0123455432}\enspace $n{=}10$, $|H|{=}5$, len 6, br 4, $|\mathrm{Aut}|{=}2$, orb 9/8, hist-orb 4/3; M 29696, res 14232, rev 8418: ok}\\*
0 & -- & 0 & 5 & 1 & 682 & 2 & 340 & \texttt{-+????????}\\*
1 & 0 & 1 & 5 & 1 & 340 & 4 & 680 & \texttt{--????????}\\
2 & 1 & 2 & 4 & 1 & 336 & 16 & 672 & \texttt{---??????-}\\
3 & 2 & 3 & 3 & 1 & 320 & 64 & 640 & \texttt{----????--}\\
4 & 3 & 4 & 2 & 1 & 256 & 256 & 512 & \texttt{-----??---}\\
5 & 4 & 5 & 1 & 2 & 0 & 1024 & 0 & \texttt{----------}\\
6 & 4 & 5 & 1 & 2 & 0 & 1024 & 0 & \texttt{----------}\\
7 & 3 & 4 & 1 & 1 & 0 & 1024 & 0 & \texttt{----------}\\
8 & 2 & 3 & 1 & 1 & 0 & 1024 & 0 & \texttt{----------}\\
9 & 1 & 2 & 1 & 1 & 0 & 1024 & 0 & \texttt{----------}\\
\multicolumn{9}{@{}l@{}}{\rule{0pt}{2.6ex}\textbf{T10.169}\enspace \texttt{0123455431}\enspace $n{=}10$, $|H|{=}5$, len 6, br 4, $|\mathrm{Aut}|{=}2$, orb 9/8, hist-orb 4/3; M 28672, res 13620, rev 8182: ok}\\*
0 & -- & 0 & 5 & 1 & 426 & 2 & 596 & \texttt{-?????????}\\*
1 & 0 & 1 & 4 & 1 & 680 & 8 & 336 & \texttt{--+??????-}\\
2 & 1 & 2 & 4 & 1 & 336 & 16 & 672 & \texttt{---??????-}\\
3 & 2 & 3 & 3 & 1 & 320 & 64 & 640 & \texttt{----????--}\\
4 & 3 & 4 & 2 & 1 & 256 & 256 & 512 & \texttt{-----??---}\\
5 & 4 & 5 & 1 & 2 & 0 & 1024 & 0 & \texttt{----------}\\
6 & 4 & 5 & 1 & 2 & 0 & 1024 & 0 & \texttt{----------}\\
7 & 3 & 4 & 1 & 1 & 0 & 1024 & 0 & \texttt{----------}\\
8 & 2 & 3 & 1 & 1 & 0 & 1024 & 0 & \texttt{----------}\\
9 & 0 & 1 & 1 & 1 & 0 & 1024 & 0 & \texttt{----------}\\
\multicolumn{9}{@{}l@{}}{\rule{0pt}{2.6ex}\textbf{T10.170}\enspace \texttt{0123455423}\enspace $n{=}10$, $|H|{=}4$, len 6, br 3, $|\mathrm{Aut}|{=}2$, orb 9/8, hist-orb 3/2; M 29696, res 8512, rev 11178: ok}\\*
0 & -- & 0 & 4 & 1 & 830 & 2 & 192 & \texttt{-+????????}\\*
1 & 0 & 1 & 4 & 1 & 636 & 4 & 384 & \texttt{--????????}\\
2 & 1 & 2 & 3 & 1 & 672 & 32 & 320 & \texttt{---+????--}\\
3 & 2 & 3 & 3 & 1 & 320 & 64 & 640 & \texttt{----????--}\\
4 & 3 & 4 & 2 & 1 & 256 & 256 & 512 & \texttt{-----??---}\\
5 & 4 & 5 & 1 & 2 & 0 & 1024 & 0 & \texttt{----------}\\
6 & 4 & 5 & 1 & 2 & 0 & 1024 & 0 & \texttt{----------}\\
7 & 3 & 4 & 1 & 1 & 0 & 1024 & 0 & \texttt{----------}\\
8 & 1 & 2 & 1 & 1 & 512 & 512 & 0 & \texttt{---------+}\\
9 & 8 & 3 & 1 & 1 & 0 & 1024 & 0 & \texttt{----------}\\
\multicolumn{9}{@{}l@{}}{\rule{0pt}{2.6ex}\textbf{T10.171}\enspace \texttt{0123455422}\enspace $n{=}10$, $|H|{=}5$, len 6, br 3, $|\mathrm{Aut}|{=}4$, orb 8/6, hist-orb 3/1; M 28672, res 13600, rev 8174: ok}\\*
0 & -- & 0 & 5 & 1 & 618 & 2 & 404 & \texttt{-+????????}\\*
1 & 0 & 1 & 5 & 1 & 212 & 4 & 808 & \texttt{--????????}\\
2 & 1 & 2 & 3 & 1 & 672 & 32 & 320 & \texttt{---+????--}\\
3 & 2 & 3 & 3 & 1 & 320 & 64 & 640 & \texttt{----????--}\\
4 & 3 & 4 & 2 & 1 & 256 & 256 & 512 & \texttt{-----??---}\\
5 & 4 & 5 & 1 & 2 & 0 & 1024 & 0 & \texttt{----------}\\
6 & 4 & 5 & 1 & 2 & 0 & 1024 & 0 & \texttt{----------}\\
7 & 3 & 4 & 1 & 1 & 0 & 1024 & 0 & \texttt{----------}\\
8 & 1 & 2 & 1 & 2 & 0 & 1024 & 0 & \texttt{----------}\\
9 & 1 & 2 & 1 & 2 & 0 & 1024 & 0 & \texttt{----------}\\
\multicolumn{9}{@{}l@{}}{\rule{0pt}{2.6ex}\textbf{T10.172}\enspace \texttt{0123455421}\enspace $n{=}10$, $|H|{=}5$, len 6, br 4, $|\mathrm{Aut}|{=}2$, orb 9/8, hist-orb 4/3; M 27648, res 13692, rev 7510: ok}\\*
0 & -- & 0 & 5 & 1 & 362 & 2 & 660 & \texttt{-?????????}\\*
1 & 0 & 1 & 4 & 1 & 424 & 8 & 592 & \texttt{--???????-}\\
2 & 1 & 2 & 3 & 1 & 672 & 32 & 320 & \texttt{---+????--}\\
3 & 2 & 3 & 3 & 1 & 320 & 64 & 640 & \texttt{----????--}\\
4 & 3 & 4 & 2 & 1 & 256 & 256 & 512 & \texttt{-----??---}\\
5 & 4 & 5 & 1 & 2 & 0 & 1024 & 0 & \texttt{----------}\\
6 & 4 & 5 & 1 & 2 & 0 & 1024 & 0 & \texttt{----------}\\
7 & 3 & 4 & 1 & 1 & 0 & 1024 & 0 & \texttt{----------}\\
8 & 1 & 2 & 1 & 1 & 0 & 1024 & 0 & \texttt{----------}\\
9 & 0 & 1 & 1 & 1 & 0 & 1024 & 0 & \texttt{----------}\\
\multicolumn{9}{@{}l@{}}{\rule{0pt}{2.6ex}\textbf{T10.173}\enspace \texttt{0123455412}\enspace $n{=}10$, $|H|{=}4$, len 6, br 3, $|\mathrm{Aut}|{=}2$, orb 9/8, hist-orb 3/2; M 27648, res 7664, rev 10354: ok}\\*
0 & -- & 0 & 4 & 1 & 702 & 2 & 320 & \texttt{-?????????}\\*
1 & 0 & 1 & 3 & 1 & 848 & 16 & 160 & \texttt{--++????--}\\
2 & 1 & 2 & 3 & 1 & 672 & 32 & 320 & \texttt{---+????--}\\
3 & 2 & 3 & 3 & 1 & 320 & 64 & 640 & \texttt{----????--}\\
4 & 3 & 4 & 2 & 1 & 256 & 256 & 512 & \texttt{-----??---}\\
5 & 4 & 5 & 1 & 2 & 0 & 1024 & 0 & \texttt{----------}\\
6 & 4 & 5 & 1 & 2 & 0 & 1024 & 0 & \texttt{----------}\\
7 & 3 & 4 & 1 & 1 & 0 & 1024 & 0 & \texttt{----------}\\
8 & 0 & 1 & 1 & 1 & 512 & 512 & 0 & \texttt{---------+}\\
9 & 8 & 2 & 1 & 1 & 0 & 1024 & 0 & \texttt{----------}\\
\multicolumn{9}{@{}l@{}}{\rule{0pt}{2.6ex}\textbf{T10.174}\enspace \texttt{0123455411}\enspace $n{=}10$, $|H|{=}5$, len 6, br 3, $|\mathrm{Aut}|{=}4$, orb 8/6, hist-orb 3/1; M 26624, res 11876, rev 7810: ok}\\*
0 & -- & 0 & 5 & 1 & 234 & 2 & 788 & \texttt{-?????????}\\*
1 & 0 & 1 & 3 & 1 & 848 & 16 & 160 & \texttt{--++????--}\\
2 & 1 & 2 & 3 & 1 & 672 & 32 & 320 & \texttt{---+????--}\\
3 & 2 & 3 & 3 & 1 & 320 & 64 & 640 & \texttt{----????--}\\
4 & 3 & 4 & 2 & 1 & 256 & 256 & 512 & \texttt{-----??---}\\
5 & 4 & 5 & 1 & 2 & 0 & 1024 & 0 & \texttt{----------}\\
6 & 4 & 5 & 1 & 2 & 0 & 1024 & 0 & \texttt{----------}\\
7 & 3 & 4 & 1 & 1 & 0 & 1024 & 0 & \texttt{----------}\\
8 & 0 & 1 & 1 & 2 & 0 & 1024 & 0 & \texttt{----------}\\
9 & 0 & 1 & 1 & 2 & 0 & 1024 & 0 & \texttt{----------}\\
\multicolumn{9}{@{}l@{}}{\rule{0pt}{2.6ex}\textbf{T10.175}\enspace \texttt{0123455345}\enspace $n{=}10$, $|H|{=}3$, len 6, br 2, $|\mathrm{Aut}|{=}2$, orb 9/8, hist-orb 2/1; M 32768, res 4616, rev 13690: ok}\\*
0 & -- & 0 & 3 & 1 & 950 & 2 & 72 & \texttt{-++???????}\\*
1 & 0 & 1 & 3 & 1 & 876 & 4 & 144 & \texttt{--+???????}\\
2 & 1 & 2 & 3 & 1 & 728 & 8 & 288 & \texttt{---???????}\\
3 & 2 & 3 & 2 & 1 & 640 & 128 & 256 & \texttt{----+??---}\\
4 & 3 & 4 & 2 & 1 & 256 & 256 & 512 & \texttt{-----??---}\\
5 & 4 & 5 & 1 & 2 & 0 & 1024 & 0 & \texttt{----------}\\
6 & 4 & 5 & 1 & 2 & 0 & 1024 & 0 & \texttt{----------}\\
7 & 2 & 3 & 1 & 1 & 768 & 256 & 0 & \texttt{--------++}\\
8 & 7 & 4 & 1 & 1 & 512 & 512 & 0 & \texttt{---------+}\\
9 & 8 & 5 & 1 & 1 & 0 & 1024 & 0 & \texttt{----------}\\
\multicolumn{9}{@{}l@{}}{\rule{0pt}{2.6ex}\textbf{T10.176}\enspace \texttt{0123455344}\enspace $n{=}10$, $|H|{=}4$, len 6, br 3, $|\mathrm{Aut}|{=}4$, orb 8/6, hist-orb 2/0; M 31744, res 7740, rev 12466: ok}\\*
0 & -- & 0 & 4 & 1 & 898 & 2 & 124 & \texttt{-++???????}\\*
1 & 0 & 1 & 4 & 1 & 772 & 4 & 248 & \texttt{--+???????}\\
2 & 1 & 2 & 4 & 1 & 520 & 8 & 496 & \texttt{---???????}\\
3 & 2 & 3 & 2 & 1 & 640 & 128 & 256 & \texttt{----+??---}\\
4 & 3 & 4 & 2 & 1 & 256 & 256 & 512 & \texttt{-----??---}\\
5 & 4 & 5 & 1 & 2 & 0 & 1024 & 0 & \texttt{----------}\\
6 & 4 & 5 & 1 & 2 & 0 & 1024 & 0 & \texttt{----------}\\
7 & 2 & 3 & 2 & 1 & 256 & 256 & 512 & \texttt{--------??}\\
8 & 7 & 4 & 1 & 2 & 0 & 1024 & 0 & \texttt{----------}\\
9 & 7 & 4 & 1 & 2 & 0 & 1024 & 0 & \texttt{----------}\\
\multicolumn{9}{@{}l@{}}{\rule{0pt}{2.6ex}\textbf{T10.177}\enspace \texttt{0123455343}\enspace $n{=}10$, $|H|{=}4$, len 6, br 2, $|\mathrm{Aut}|{=}2$, orb 9/8, hist-orb 3/2; M 30720, res 9712, rev 11218: ok}\\*
0 & -- & 0 & 4 & 1 & 846 & 2 & 176 & \texttt{-++???????}\\*
1 & 0 & 1 & 4 & 1 & 668 & 4 & 352 & \texttt{--+???????}\\
2 & 1 & 2 & 4 & 1 & 312 & 8 & 704 & \texttt{---???????}\\
3 & 2 & 3 & 2 & 1 & 640 & 128 & 256 & \texttt{----+??---}\\
4 & 3 & 4 & 2 & 1 & 256 & 256 & 512 & \texttt{-----??---}\\
5 & 4 & 5 & 1 & 2 & 0 & 1024 & 0 & \texttt{----------}\\
6 & 4 & 5 & 1 & 2 & 0 & 1024 & 0 & \texttt{----------}\\
7 & 2 & 3 & 1 & 1 & 512 & 512 & 0 & \texttt{--------+-}\\
8 & 7 & 4 & 1 & 1 & 0 & 1024 & 0 & \texttt{----------}\\
9 & 2 & 3 & 1 & 1 & 0 & 1024 & 0 & \texttt{----------}\\
\multicolumn{9}{@{}l@{}}{\rule{0pt}{2.6ex}\textbf{T10.178}\enspace \texttt{0123455342}\enspace $n{=}10$, $|H|{=}4$, len 6, br 3, $|\mathrm{Aut}|{=}2$, orb 9/8, hist-orb 3/2; M 29696, res 10400, rev 10174: ok}\\*
0 & -- & 0 & 4 & 1 & 718 & 2 & 304 & \texttt{-+????????}\\*
1 & 0 & 1 & 4 & 1 & 412 & 4 & 608 & \texttt{--????????}\\
2 & 1 & 2 & 3 & 1 & 624 & 16 & 384 & \texttt{---??????-}\\
3 & 2 & 3 & 2 & 1 & 640 & 128 & 256 & \texttt{----+??---}\\
4 & 3 & 4 & 2 & 1 & 256 & 256 & 512 & \texttt{-----??---}\\
5 & 4 & 5 & 1 & 2 & 0 & 1024 & 0 & \texttt{----------}\\
6 & 4 & 5 & 1 & 2 & 0 & 1024 & 0 & \texttt{----------}\\
7 & 2 & 3 & 1 & 1 & 512 & 512 & 0 & \texttt{--------+-}\\
8 & 7 & 4 & 1 & 1 & 0 & 1024 & 0 & \texttt{----------}\\
9 & 1 & 2 & 1 & 1 & 0 & 1024 & 0 & \texttt{----------}\\
\multicolumn{9}{@{}l@{}}{\rule{0pt}{2.6ex}\textbf{T10.179}\enspace \texttt{0123455341}\enspace $n{=}10$, $|H|{=}4$, len 6, br 3, $|\mathrm{Aut}|{=}2$, orb 9/8, hist-orb 3/2; M 28672, res 9392, rev 10074: ok}\\*
0 & -- & 0 & 4 & 1 & 462 & 2 & 560 & \texttt{-?????????}\\*
1 & 0 & 1 & 3 & 1 & 824 & 8 & 192 & \texttt{--+??????-}\\
2 & 1 & 2 & 3 & 1 & 624 & 16 & 384 & \texttt{---??????-}\\
3 & 2 & 3 & 2 & 1 & 640 & 128 & 256 & \texttt{----+??---}\\
4 & 3 & 4 & 2 & 1 & 256 & 256 & 512 & \texttt{-----??---}\\
5 & 4 & 5 & 1 & 2 & 0 & 1024 & 0 & \texttt{----------}\\
6 & 4 & 5 & 1 & 2 & 0 & 1024 & 0 & \texttt{----------}\\
7 & 2 & 3 & 1 & 1 & 512 & 512 & 0 & \texttt{--------+-}\\
8 & 7 & 4 & 1 & 1 & 0 & 1024 & 0 & \texttt{----------}\\
9 & 0 & 1 & 1 & 1 & 0 & 1024 & 0 & \texttt{----------}\\
\multicolumn{9}{@{}l@{}}{\rule{0pt}{2.6ex}\textbf{T10.180}\enspace \texttt{0123455333}\enspace $n{=}10$, $|H|{=}5$, len 6, br 2, $|\mathrm{Aut}|{=}12$, orb 7/5, hist-orb 2/0; M 29696, res 12260, rev 9778: ok}\\*
0 & -- & 0 & 5 & 1 & 794 & 2 & 228 & \texttt{-++???????}\\*
1 & 0 & 1 & 5 & 1 & 564 & 4 & 456 & \texttt{--+???????}\\
2 & 1 & 2 & 5 & 1 & 104 & 8 & 912 & \texttt{---???????}\\
3 & 2 & 3 & 2 & 1 & 640 & 128 & 256 & \texttt{----+??---}\\
4 & 3 & 4 & 2 & 1 & 256 & 256 & 512 & \texttt{-----??---}\\
5 & 4 & 5 & 1 & 2 & 0 & 1024 & 0 & \texttt{----------}\\
6 & 4 & 5 & 1 & 2 & 0 & 1024 & 0 & \texttt{----------}\\
7 & 2 & 3 & 1 & 3 & 0 & 1024 & 0 & \texttt{----------}\\
8 & 2 & 3 & 1 & 3 & 0 & 1024 & 0 & \texttt{----------}\\
9 & 2 & 3 & 1 & 3 & 0 & 1024 & 0 & \texttt{----------}\\
\multicolumn{9}{@{}l@{}}{\rule{0pt}{2.6ex}\textbf{T10.181}\enspace \texttt{0123455332}\enspace $n{=}10$, $|H|{=}5$, len 6, br 3, $|\mathrm{Aut}|{=}4$, orb 8/6, hist-orb 3/1; M 28672, res 13832, rev 8062: ok}\\*
0 & -- & 0 & 5 & 1 & 666 & 2 & 356 & \texttt{-+????????}\\*
1 & 0 & 1 & 5 & 1 & 308 & 4 & 712 & \texttt{--????????}\\
2 & 1 & 2 & 4 & 1 & 208 & 16 & 800 & \texttt{---??????-}\\
3 & 2 & 3 & 2 & 1 & 640 & 128 & 256 & \texttt{----+??---}\\
4 & 3 & 4 & 2 & 1 & 256 & 256 & 512 & \texttt{-----??---}\\
5 & 4 & 5 & 1 & 2 & 0 & 1024 & 0 & \texttt{----------}\\
6 & 4 & 5 & 1 & 2 & 0 & 1024 & 0 & \texttt{----------}\\
7 & 2 & 3 & 1 & 2 & 0 & 1024 & 0 & \texttt{----------}\\
8 & 2 & 3 & 1 & 2 & 0 & 1024 & 0 & \texttt{----------}\\
9 & 1 & 2 & 1 & 1 & 0 & 1024 & 0 & \texttt{----------}\\
\multicolumn{9}{@{}l@{}}{\rule{0pt}{2.6ex}\textbf{T10.182}\enspace \texttt{0123455331}\enspace $n{=}10$, $|H|{=}5$, len 6, br 3, $|\mathrm{Aut}|{=}4$, orb 8/6, hist-orb 3/1; M 27648, res 13396, rev 7722: ok}\\*
0 & -- & 0 & 5 & 1 & 410 & 2 & 612 & \texttt{-?????????}\\*
1 & 0 & 1 & 4 & 1 & 616 & 8 & 400 & \texttt{--+??????-}\\
2 & 1 & 2 & 4 & 1 & 208 & 16 & 800 & \texttt{---??????-}\\
3 & 2 & 3 & 2 & 1 & 640 & 128 & 256 & \texttt{----+??---}\\
4 & 3 & 4 & 2 & 1 & 256 & 256 & 512 & \texttt{-----??---}\\
5 & 4 & 5 & 1 & 2 & 0 & 1024 & 0 & \texttt{----------}\\
6 & 4 & 5 & 1 & 2 & 0 & 1024 & 0 & \texttt{----------}\\
7 & 2 & 3 & 1 & 2 & 0 & 1024 & 0 & \texttt{----------}\\
8 & 2 & 3 & 1 & 2 & 0 & 1024 & 0 & \texttt{----------}\\
9 & 0 & 1 & 1 & 1 & 0 & 1024 & 0 & \texttt{----------}\\
\multicolumn{9}{@{}l@{}}{\rule{0pt}{2.6ex}\textbf{T10.183}\enspace \texttt{0123455323}\enspace $n{=}10$, $|H|{=}4$, len 6, br 3, $|\mathrm{Aut}|{=}2$, orb 9/8, hist-orb 3/2; M 28672, res 9296, rev 10138: ok}\\*
0 & -- & 0 & 4 & 1 & 782 & 2 & 240 & \texttt{-+????????}\\*
1 & 0 & 1 & 4 & 1 & 540 & 4 & 480 & \texttt{--????????}\\
2 & 1 & 2 & 3 & 1 & 416 & 32 & 576 & \texttt{---?????--}\\
3 & 2 & 3 & 2 & 1 & 640 & 128 & 256 & \texttt{----+??---}\\
4 & 3 & 4 & 2 & 1 & 256 & 256 & 512 & \texttt{-----??---}\\
5 & 4 & 5 & 1 & 2 & 0 & 1024 & 0 & \texttt{----------}\\
6 & 4 & 5 & 1 & 2 & 0 & 1024 & 0 & \texttt{----------}\\
7 & 2 & 3 & 1 & 1 & 0 & 1024 & 0 & \texttt{----------}\\
8 & 1 & 2 & 1 & 1 & 512 & 512 & 0 & \texttt{---------+}\\
9 & 8 & 3 & 1 & 1 & 0 & 1024 & 0 & \texttt{----------}\\
\multicolumn{9}{@{}l@{}}{\rule{0pt}{2.6ex}\textbf{T10.184}\enspace \texttt{0123455322}\enspace $n{=}10$, $|H|{=}5$, len 6, br 3, $|\mathrm{Aut}|{=}4$, orb 8/6, hist-orb 3/1; M 27648, res 13616, rev 7518: ok}\\*
0 & -- & 0 & 5 & 1 & 602 & 2 & 420 & \texttt{-+????????}\\*
1 & 0 & 1 & 5 & 1 & 180 & 4 & 840 & \texttt{--????????}\\
2 & 1 & 2 & 3 & 1 & 416 & 32 & 576 & \texttt{---?????--}\\
3 & 2 & 3 & 2 & 1 & 640 & 128 & 256 & \texttt{----+??---}\\
4 & 3 & 4 & 2 & 1 & 256 & 256 & 512 & \texttt{-----??---}\\
5 & 4 & 5 & 1 & 2 & 0 & 1024 & 0 & \texttt{----------}\\
6 & 4 & 5 & 1 & 2 & 0 & 1024 & 0 & \texttt{----------}\\
7 & 2 & 3 & 1 & 1 & 0 & 1024 & 0 & \texttt{----------}\\
8 & 1 & 2 & 1 & 2 & 0 & 1024 & 0 & \texttt{----------}\\
9 & 1 & 2 & 1 & 2 & 0 & 1024 & 0 & \texttt{----------}\\
\multicolumn{9}{@{}l@{}}{\rule{0pt}{2.6ex}\textbf{T10.185}\enspace \texttt{0123455321}\enspace $n{=}10$, $|H|{=}5$, len 6, br 4, $|\mathrm{Aut}|{=}2$, orb 9/8, hist-orb 4/3; M 26624, res 13884, rev 6742: ok}\\*
0 & -- & 0 & 5 & 1 & 346 & 2 & 676 & \texttt{-?????????}\\*
1 & 0 & 1 & 4 & 1 & 360 & 8 & 656 & \texttt{--???????-}\\
2 & 1 & 2 & 3 & 1 & 416 & 32 & 576 & \texttt{---?????--}\\
3 & 2 & 3 & 2 & 1 & 640 & 128 & 256 & \texttt{----+??---}\\
4 & 3 & 4 & 2 & 1 & 256 & 256 & 512 & \texttt{-----??---}\\
5 & 4 & 5 & 1 & 2 & 0 & 1024 & 0 & \texttt{----------}\\
6 & 4 & 5 & 1 & 2 & 0 & 1024 & 0 & \texttt{----------}\\
7 & 2 & 3 & 1 & 1 & 0 & 1024 & 0 & \texttt{----------}\\
8 & 1 & 2 & 1 & 1 & 0 & 1024 & 0 & \texttt{----------}\\
9 & 0 & 1 & 1 & 1 & 0 & 1024 & 0 & \texttt{----------}\\
\multicolumn{9}{@{}l@{}}{\rule{0pt}{2.6ex}\textbf{T10.186}\enspace \texttt{0123455312}\enspace $n{=}10$, $|H|{=}4$, len 6, br 3, $|\mathrm{Aut}|{=}2$, orb 9/8, hist-orb 3/2; M 26624, res 8432, rev 9338: ok}\\*
0 & -- & 0 & 4 & 1 & 654 & 2 & 368 & \texttt{-?????????}\\*
1 & 0 & 1 & 3 & 1 & 720 & 16 & 288 & \texttt{--+?????--}\\
2 & 1 & 2 & 3 & 1 & 416 & 32 & 576 & \texttt{---?????--}\\
3 & 2 & 3 & 2 & 1 & 640 & 128 & 256 & \texttt{----+??---}\\
4 & 3 & 4 & 2 & 1 & 256 & 256 & 512 & \texttt{-----??---}\\
5 & 4 & 5 & 1 & 2 & 0 & 1024 & 0 & \texttt{----------}\\
6 & 4 & 5 & 1 & 2 & 0 & 1024 & 0 & \texttt{----------}\\
7 & 2 & 3 & 1 & 1 & 0 & 1024 & 0 & \texttt{----------}\\
8 & 0 & 1 & 1 & 1 & 512 & 512 & 0 & \texttt{---------+}\\
9 & 8 & 2 & 1 & 1 & 0 & 1024 & 0 & \texttt{----------}\\
\multicolumn{9}{@{}l@{}}{\rule{0pt}{2.6ex}\textbf{T10.187}\enspace \texttt{0123455311}\enspace $n{=}10$, $|H|{=}5$, len 6, br 3, $|\mathrm{Aut}|{=}4$, orb 8/6, hist-orb 3/1; M 25600, res 12356, rev 6906: ok}\\*
0 & -- & 0 & 5 & 1 & 218 & 2 & 804 & \texttt{-?????????}\\*
1 & 0 & 1 & 3 & 1 & 720 & 16 & 288 & \texttt{--+?????--}\\
2 & 1 & 2 & 3 & 1 & 416 & 32 & 576 & \texttt{---?????--}\\
3 & 2 & 3 & 2 & 1 & 640 & 128 & 256 & \texttt{----+??---}\\
4 & 3 & 4 & 2 & 1 & 256 & 256 & 512 & \texttt{-----??---}\\
5 & 4 & 5 & 1 & 2 & 0 & 1024 & 0 & \texttt{----------}\\
6 & 4 & 5 & 1 & 2 & 0 & 1024 & 0 & \texttt{----------}\\
7 & 2 & 3 & 1 & 1 & 0 & 1024 & 0 & \texttt{----------}\\
8 & 0 & 1 & 1 & 2 & 0 & 1024 & 0 & \texttt{----------}\\
9 & 0 & 1 & 1 & 2 & 0 & 1024 & 0 & \texttt{----------}\\
\multicolumn{9}{@{}l@{}}{\rule{0pt}{2.6ex}\textbf{T10.188}\enspace \texttt{0123455234}\enspace $n{=}10$, $|H|{=}3$, len 6, br 2, $|\mathrm{Aut}|{=}2$, orb 9/8, hist-orb 2/1; M 29696, res 4000, rev 12334: ok}\\*
0 & -- & 0 & 3 & 1 & 918 & 2 & 104 & \texttt{-+????????}\\*
1 & 0 & 1 & 3 & 1 & 812 & 4 & 208 & \texttt{--????????}\\
2 & 1 & 2 & 2 & 1 & 832 & 64 & 128 & \texttt{---++??---}\\
3 & 2 & 3 & 2 & 1 & 640 & 128 & 256 & \texttt{----+??---}\\
4 & 3 & 4 & 2 & 1 & 256 & 256 & 512 & \texttt{-----??---}\\
5 & 4 & 5 & 1 & 2 & 0 & 1024 & 0 & \texttt{----------}\\
6 & 4 & 5 & 1 & 2 & 0 & 1024 & 0 & \texttt{----------}\\
7 & 1 & 2 & 1 & 1 & 768 & 256 & 0 & \texttt{--------++}\\
8 & 7 & 3 & 1 & 1 & 512 & 512 & 0 & \texttt{---------+}\\
9 & 8 & 4 & 1 & 1 & 0 & 1024 & 0 & \texttt{----------}\\
\multicolumn{9}{@{}l@{}}{\rule{0pt}{2.6ex}\textbf{T10.189}\enspace \texttt{0123455233}\enspace $n{=}10$, $|H|{=}4$, len 6, br 3, $|\mathrm{Aut}|{=}4$, orb 8/6, hist-orb 2/0; M 28672, res 7424, rev 10738: ok}\\*
0 & -- & 0 & 4 & 1 & 802 & 2 & 220 & \texttt{-+????????}\\*
1 & 0 & 1 & 4 & 1 & 580 & 4 & 440 & \texttt{--????????}\\
2 & 1 & 2 & 2 & 1 & 832 & 64 & 128 & \texttt{---++??---}\\
3 & 2 & 3 & 2 & 1 & 640 & 128 & 256 & \texttt{----+??---}\\
4 & 3 & 4 & 2 & 1 & 256 & 256 & 512 & \texttt{-----??---}\\
5 & 4 & 5 & 1 & 2 & 0 & 1024 & 0 & \texttt{----------}\\
6 & 4 & 5 & 1 & 2 & 0 & 1024 & 0 & \texttt{----------}\\
7 & 1 & 2 & 2 & 1 & 256 & 256 & 512 & \texttt{--------??}\\
8 & 7 & 3 & 1 & 2 & 0 & 1024 & 0 & \texttt{----------}\\
9 & 7 & 3 & 1 & 2 & 0 & 1024 & 0 & \texttt{----------}\\
\multicolumn{9}{@{}l@{}}{\rule{0pt}{2.6ex}\textbf{T10.190}\enspace \texttt{0123455232}\enspace $n{=}10$, $|H|{=}4$, len 6, br 2, $|\mathrm{Aut}|{=}2$, orb 9/8, hist-orb 3/2; M 27648, res 9568, rev 9342: ok}\\*
0 & -- & 0 & 4 & 1 & 686 & 2 & 336 & \texttt{-+????????}\\*
1 & 0 & 1 & 4 & 1 & 348 & 4 & 672 & \texttt{--????????}\\
2 & 1 & 2 & 2 & 1 & 832 & 64 & 128 & \texttt{---++??---}\\
3 & 2 & 3 & 2 & 1 & 640 & 128 & 256 & \texttt{----+??---}\\
4 & 3 & 4 & 2 & 1 & 256 & 256 & 512 & \texttt{-----??---}\\
5 & 4 & 5 & 1 & 2 & 0 & 1024 & 0 & \texttt{----------}\\
6 & 4 & 5 & 1 & 2 & 0 & 1024 & 0 & \texttt{----------}\\
7 & 1 & 2 & 1 & 1 & 512 & 512 & 0 & \texttt{--------+-}\\
8 & 7 & 3 & 1 & 1 & 0 & 1024 & 0 & \texttt{----------}\\
9 & 1 & 2 & 1 & 1 & 0 & 1024 & 0 & \texttt{----------}\\
\multicolumn{9}{@{}l@{}}{\rule{0pt}{2.6ex}\textbf{T10.191}\enspace \texttt{0123455231}\enspace $n{=}10$, $|H|{=}4$, len 6, br 3, $|\mathrm{Aut}|{=}2$, orb 9/8, hist-orb 3/2; M 26624, res 8912, rev 9022: ok}\\*
0 & -- & 0 & 4 & 1 & 430 & 2 & 592 & \texttt{-?????????}\\*
1 & 0 & 1 & 3 & 1 & 696 & 8 & 320 & \texttt{--???????-}\\
2 & 1 & 2 & 2 & 1 & 832 & 64 & 128 & \texttt{---++??---}\\
3 & 2 & 3 & 2 & 1 & 640 & 128 & 256 & \texttt{----+??---}\\
4 & 3 & 4 & 2 & 1 & 256 & 256 & 512 & \texttt{-----??---}\\
5 & 4 & 5 & 1 & 2 & 0 & 1024 & 0 & \texttt{----------}\\
6 & 4 & 5 & 1 & 2 & 0 & 1024 & 0 & \texttt{----------}\\
7 & 1 & 2 & 1 & 1 & 512 & 512 & 0 & \texttt{--------+-}\\
8 & 7 & 3 & 1 & 1 & 0 & 1024 & 0 & \texttt{----------}\\
9 & 0 & 1 & 1 & 1 & 0 & 1024 & 0 & \texttt{----------}\\
\multicolumn{9}{@{}l@{}}{\rule{0pt}{2.6ex}\textbf{T10.192}\enspace \texttt{0123455222}\enspace $n{=}10$, $|H|{=}5$, len 6, br 2, $|\mathrm{Aut}|{=}12$, orb 7/5, hist-orb 2/0; M 26624, res 12352, rev 7562: ok}\\*
0 & -- & 0 & 5 & 1 & 570 & 2 & 452 & \texttt{-+????????}\\*
1 & 0 & 1 & 5 & 1 & 116 & 4 & 904 & \texttt{--????????}\\
2 & 1 & 2 & 2 & 1 & 832 & 64 & 128 & \texttt{---++??---}\\
3 & 2 & 3 & 2 & 1 & 640 & 128 & 256 & \texttt{----+??---}\\
4 & 3 & 4 & 2 & 1 & 256 & 256 & 512 & \texttt{-----??---}\\
5 & 4 & 5 & 1 & 2 & 0 & 1024 & 0 & \texttt{----------}\\
6 & 4 & 5 & 1 & 2 & 0 & 1024 & 0 & \texttt{----------}\\
7 & 1 & 2 & 1 & 3 & 0 & 1024 & 0 & \texttt{----------}\\
8 & 1 & 2 & 1 & 3 & 0 & 1024 & 0 & \texttt{----------}\\
9 & 1 & 2 & 1 & 3 & 0 & 1024 & 0 & \texttt{----------}\\
\multicolumn{9}{@{}l@{}}{\rule{0pt}{2.6ex}\textbf{T10.193}\enspace \texttt{0123455221}\enspace $n{=}10$, $|H|{=}5$, len 6, br 3, $|\mathrm{Aut}|{=}4$, orb 8/6, hist-orb 3/1; M 25600, res 12972, rev 6562: ok}\\*
0 & -- & 0 & 5 & 1 & 314 & 2 & 708 & \texttt{-?????????}\\*
1 & 0 & 1 & 4 & 1 & 232 & 8 & 784 & \texttt{--???????-}\\
2 & 1 & 2 & 2 & 1 & 832 & 64 & 128 & \texttt{---++??---}\\
3 & 2 & 3 & 2 & 1 & 640 & 128 & 256 & \texttt{----+??---}\\
4 & 3 & 4 & 2 & 1 & 256 & 256 & 512 & \texttt{-----??---}\\
5 & 4 & 5 & 1 & 2 & 0 & 1024 & 0 & \texttt{----------}\\
6 & 4 & 5 & 1 & 2 & 0 & 1024 & 0 & \texttt{----------}\\
7 & 1 & 2 & 1 & 2 & 0 & 1024 & 0 & \texttt{----------}\\
8 & 1 & 2 & 1 & 2 & 0 & 1024 & 0 & \texttt{----------}\\
9 & 0 & 1 & 1 & 1 & 0 & 1024 & 0 & \texttt{----------}\\
\multicolumn{9}{@{}l@{}}{\rule{0pt}{2.6ex}\textbf{T10.194}\enspace \texttt{0123455212}\enspace $n{=}10$, $|H|{=}4$, len 6, br 3, $|\mathrm{Aut}|{=}2$, orb 9/8, hist-orb 3/2; M 25600, res 8672, rev 8578: ok}\\*
0 & -- & 0 & 4 & 1 & 558 & 2 & 464 & \texttt{-?????????}\\*
1 & 0 & 1 & 3 & 1 & 464 & 16 & 544 & \texttt{--??????--}\\
2 & 1 & 2 & 2 & 1 & 832 & 64 & 128 & \texttt{---++??---}\\
3 & 2 & 3 & 2 & 1 & 640 & 128 & 256 & \texttt{----+??---}\\
4 & 3 & 4 & 2 & 1 & 256 & 256 & 512 & \texttt{-----??---}\\
5 & 4 & 5 & 1 & 2 & 0 & 1024 & 0 & \texttt{----------}\\
6 & 4 & 5 & 1 & 2 & 0 & 1024 & 0 & \texttt{----------}\\
7 & 1 & 2 & 1 & 1 & 0 & 1024 & 0 & \texttt{----------}\\
8 & 0 & 1 & 1 & 1 & 512 & 512 & 0 & \texttt{---------+}\\
9 & 8 & 2 & 1 & 1 & 0 & 1024 & 0 & \texttt{----------}\\
\multicolumn{9}{@{}l@{}}{\rule{0pt}{2.6ex}\textbf{T10.195}\enspace \texttt{0123455211}\enspace $n{=}10$, $|H|{=}5$, len 6, br 3, $|\mathrm{Aut}|{=}4$, orb 8/6, hist-orb 3/1; M 24576, res 12020, rev 6418: ok}\\*
0 & -- & 0 & 5 & 1 & 186 & 2 & 836 & \texttt{-?????????}\\*
1 & 0 & 1 & 3 & 1 & 464 & 16 & 544 & \texttt{--??????--}\\
2 & 1 & 2 & 2 & 1 & 832 & 64 & 128 & \texttt{---++??---}\\
3 & 2 & 3 & 2 & 1 & 640 & 128 & 256 & \texttt{----+??---}\\
4 & 3 & 4 & 2 & 1 & 256 & 256 & 512 & \texttt{-----??---}\\
5 & 4 & 5 & 1 & 2 & 0 & 1024 & 0 & \texttt{----------}\\
6 & 4 & 5 & 1 & 2 & 0 & 1024 & 0 & \texttt{----------}\\
7 & 1 & 2 & 1 & 1 & 0 & 1024 & 0 & \texttt{----------}\\
8 & 0 & 1 & 1 & 2 & 0 & 1024 & 0 & \texttt{----------}\\
9 & 0 & 1 & 1 & 2 & 0 & 1024 & 0 & \texttt{----------}\\
\multicolumn{9}{@{}l@{}}{\rule{0pt}{2.6ex}\textbf{T10.196}\enspace \texttt{0123455123}\enspace $n{=}10$, $|H|{=}3$, len 6, br 2, $|\mathrm{Aut}|{=}2$, orb 9/8, hist-orb 2/1; M 26624, res 3304, rev 10998: ok}\\*
0 & -- & 0 & 3 & 1 & 854 & 2 & 168 & \texttt{-?????????}\\*
1 & 0 & 1 & 2 & 1 & 928 & 32 & 64 & \texttt{--+++??---}\\
2 & 1 & 2 & 2 & 1 & 832 & 64 & 128 & \texttt{---++??---}\\
3 & 2 & 3 & 2 & 1 & 640 & 128 & 256 & \texttt{----+??---}\\
4 & 3 & 4 & 2 & 1 & 256 & 256 & 512 & \texttt{-----??---}\\
5 & 4 & 5 & 1 & 2 & 0 & 1024 & 0 & \texttt{----------}\\
6 & 4 & 5 & 1 & 2 & 0 & 1024 & 0 & \texttt{----------}\\
7 & 0 & 1 & 1 & 1 & 768 & 256 & 0 & \texttt{--------++}\\
8 & 7 & 2 & 1 & 1 & 512 & 512 & 0 & \texttt{---------+}\\
9 & 8 & 3 & 1 & 1 & 0 & 1024 & 0 & \texttt{----------}\\
\multicolumn{9}{@{}l@{}}{\rule{0pt}{2.6ex}\textbf{T10.197}\enspace \texttt{0123455122}\enspace $n{=}10$, $|H|{=}4$, len 6, br 3, $|\mathrm{Aut}|{=}4$, orb 8/6, hist-orb 2/0; M 25600, res 6268, rev 9382: ok}\\*
0 & -- & 0 & 4 & 1 & 610 & 2 & 412 & \texttt{-?????????}\\*
1 & 0 & 1 & 2 & 1 & 928 & 32 & 64 & \texttt{--+++??---}\\
2 & 1 & 2 & 2 & 1 & 832 & 64 & 128 & \texttt{---++??---}\\
3 & 2 & 3 & 2 & 1 & 640 & 128 & 256 & \texttt{----+??---}\\
4 & 3 & 4 & 2 & 1 & 256 & 256 & 512 & \texttt{-----??---}\\
5 & 4 & 5 & 1 & 2 & 0 & 1024 & 0 & \texttt{----------}\\
6 & 4 & 5 & 1 & 2 & 0 & 1024 & 0 & \texttt{----------}\\
7 & 0 & 1 & 2 & 1 & 256 & 256 & 512 & \texttt{--------??}\\
8 & 7 & 2 & 1 & 2 & 0 & 1024 & 0 & \texttt{----------}\\
9 & 7 & 2 & 1 & 2 & 0 & 1024 & 0 & \texttt{----------}\\
\multicolumn{9}{@{}l@{}}{\rule{0pt}{2.6ex}\textbf{T10.198}\enspace \texttt{0123455121}\enspace $n{=}10$, $|H|{=}4$, len 6, br 2, $|\mathrm{Aut}|{=}2$, orb 9/8, hist-orb 3/2; M 24576, res 7696, rev 8302: ok}\\*
0 & -- & 0 & 4 & 1 & 366 & 2 & 656 & \texttt{-?????????}\\*
1 & 0 & 1 & 2 & 1 & 928 & 32 & 64 & \texttt{--+++??---}\\
2 & 1 & 2 & 2 & 1 & 832 & 64 & 128 & \texttt{---++??---}\\
3 & 2 & 3 & 2 & 1 & 640 & 128 & 256 & \texttt{----+??---}\\
4 & 3 & 4 & 2 & 1 & 256 & 256 & 512 & \texttt{-----??---}\\
5 & 4 & 5 & 1 & 2 & 0 & 1024 & 0 & \texttt{----------}\\
6 & 4 & 5 & 1 & 2 & 0 & 1024 & 0 & \texttt{----------}\\
7 & 0 & 1 & 1 & 1 & 512 & 512 & 0 & \texttt{--------+-}\\
8 & 7 & 2 & 1 & 1 & 0 & 1024 & 0 & \texttt{----------}\\
9 & 0 & 1 & 1 & 1 & 0 & 1024 & 0 & \texttt{----------}\\
\multicolumn{9}{@{}l@{}}{\rule{0pt}{2.6ex}\textbf{T10.199}\enspace \texttt{0123455111}\enspace $n{=}10$, $|H|{=}5$, len 6, br 2, $|\mathrm{Aut}|{=}12$, orb 7/5, hist-orb 2/0; M 23552, res 9892, rev 6710: ok}\\*
0 & -- & 0 & 5 & 1 & 122 & 2 & 900 & \texttt{-?????????}\\*
1 & 0 & 1 & 2 & 1 & 928 & 32 & 64 & \texttt{--+++??---}\\
2 & 1 & 2 & 2 & 1 & 832 & 64 & 128 & \texttt{---++??---}\\
3 & 2 & 3 & 2 & 1 & 640 & 128 & 256 & \texttt{----+??---}\\
4 & 3 & 4 & 2 & 1 & 256 & 256 & 512 & \texttt{-----??---}\\
5 & 4 & 5 & 1 & 2 & 0 & 1024 & 0 & \texttt{----------}\\
6 & 4 & 5 & 1 & 2 & 0 & 1024 & 0 & \texttt{----------}\\
7 & 0 & 1 & 1 & 3 & 0 & 1024 & 0 & \texttt{----------}\\
8 & 0 & 1 & 1 & 3 & 0 & 1024 & 0 & \texttt{----------}\\
9 & 0 & 1 & 1 & 3 & 0 & 1024 & 0 & \texttt{----------}\\
\multicolumn{9}{@{}l@{}}{\rule{0pt}{2.6ex}\textbf{T10.200}\enspace \texttt{0123454545}\enspace $n{=}10$, $|H|{=}3$, len 6, br 1, $|\mathrm{Aut}|{=}6$, orb 6/4, hist-orb 1/0; M 33792, res 6480, rev 14614: ok}\\*
0 & -- & 0 & 3 & 1 & 950 & 2 & 72 & \texttt{-+++??????}\\*
1 & 0 & 1 & 3 & 1 & 876 & 4 & 144 & \texttt{--++??????}\\
2 & 1 & 2 & 3 & 1 & 728 & 8 & 288 & \texttt{---+??????}\\
3 & 2 & 3 & 3 & 1 & 432 & 16 & 576 & \texttt{----??????}\\
4 & 3 & 4 & 1 & 3 & 512 & 512 & 0 & \texttt{-----+----}\\
5 & 4 & 5 & 1 & 3 & 0 & 1024 & 0 & \texttt{----------}\\
6 & 3 & 4 & 1 & 3 & 512 & 512 & 0 & \texttt{-------+--}\\
7 & 6 & 5 & 1 & 3 & 0 & 1024 & 0 & \texttt{----------}\\
8 & 3 & 4 & 1 & 3 & 512 & 512 & 0 & \texttt{---------+}\\
9 & 8 & 5 & 1 & 3 & 0 & 1024 & 0 & \texttt{----------}\\
\multicolumn{9}{@{}l@{}}{\rule{0pt}{2.6ex}\textbf{T10.201}\enspace \texttt{0123454544}\enspace $n{=}10$, $|H|{=}4$, len 6, br 1, $|\mathrm{Aut}|{=}4$, orb 7/4, hist-orb 2/0; M 32768, res 9720, rev 13210: ok}\\*
0 & -- & 0 & 4 & 1 & 914 & 2 & 108 & \texttt{-+++??????}\\*
1 & 0 & 1 & 4 & 1 & 804 & 4 & 216 & \texttt{--++??????}\\
2 & 1 & 2 & 4 & 1 & 584 & 8 & 432 & \texttt{---+??????}\\
3 & 2 & 3 & 4 & 1 & 144 & 16 & 864 & \texttt{----??????}\\
4 & 3 & 4 & 1 & 2 & 512 & 512 & 0 & \texttt{-----+----}\\
5 & 4 & 5 & 1 & 2 & 0 & 1024 & 0 & \texttt{----------}\\
6 & 3 & 4 & 1 & 2 & 512 & 512 & 0 & \texttt{-------+--}\\
7 & 6 & 5 & 1 & 2 & 0 & 1024 & 0 & \texttt{----------}\\
8 & 3 & 4 & 1 & 2 & 0 & 1024 & 0 & \texttt{----------}\\
9 & 3 & 4 & 1 & 2 & 0 & 1024 & 0 & \texttt{----------}\\
\multicolumn{9}{@{}l@{}}{\rule{0pt}{2.6ex}\textbf{T10.202}\enspace \texttt{0123454543}\enspace $n{=}10$, $|H|{=}4$, len 6, br 2, $|\mathrm{Aut}|{=}2$, orb 8/6, hist-orb 3/2; M 31744, res 11332, rev 11374: ok}\\*
0 & -- & 0 & 4 & 1 & 850 & 2 & 172 & \texttt{-++???????}\\*
1 & 0 & 1 & 4 & 1 & 676 & 4 & 344 & \texttt{--+???????}\\
2 & 1 & 2 & 4 & 1 & 328 & 8 & 688 & \texttt{---???????}\\
3 & 2 & 3 & 3 & 1 & 288 & 32 & 704 & \texttt{----?????-}\\
4 & 3 & 4 & 1 & 2 & 512 & 512 & 0 & \texttt{-----+----}\\
5 & 4 & 5 & 1 & 2 & 0 & 1024 & 0 & \texttt{----------}\\
6 & 3 & 4 & 1 & 2 & 512 & 512 & 0 & \texttt{-------+--}\\
7 & 6 & 5 & 1 & 2 & 0 & 1024 & 0 & \texttt{----------}\\
8 & 3 & 4 & 1 & 1 & 0 & 1024 & 0 & \texttt{----------}\\
9 & 2 & 3 & 1 & 1 & 0 & 1024 & 0 & \texttt{----------}\\
\multicolumn{9}{@{}l@{}}{\rule{0pt}{2.6ex}\textbf{T10.203}\enspace \texttt{0123454542}\enspace $n{=}10$, $|H|{=}4$, len 6, br 2, $|\mathrm{Aut}|{=}2$, orb 8/6, hist-orb 3/2; M 30720, res 11952, rev 10414: ok}\\*
0 & -- & 0 & 4 & 1 & 722 & 2 & 300 & \texttt{-+????????}\\*
1 & 0 & 1 & 4 & 1 & 420 & 4 & 600 & \texttt{--????????}\\
2 & 1 & 2 & 3 & 1 & 656 & 16 & 352 & \texttt{---+?????-}\\
3 & 2 & 3 & 3 & 1 & 288 & 32 & 704 & \texttt{----?????-}\\
4 & 3 & 4 & 1 & 2 & 512 & 512 & 0 & \texttt{-----+----}\\
5 & 4 & 5 & 1 & 2 & 0 & 1024 & 0 & \texttt{----------}\\
6 & 3 & 4 & 1 & 2 & 512 & 512 & 0 & \texttt{-------+--}\\
7 & 6 & 5 & 1 & 2 & 0 & 1024 & 0 & \texttt{----------}\\
8 & 3 & 4 & 1 & 1 & 0 & 1024 & 0 & \texttt{----------}\\
9 & 1 & 2 & 1 & 1 & 0 & 1024 & 0 & \texttt{----------}\\
\multicolumn{9}{@{}l@{}}{\rule{0pt}{2.6ex}\textbf{T10.204}\enspace \texttt{0123454541}\enspace $n{=}10$, $|H|{=}4$, len 6, br 2, $|\mathrm{Aut}|{=}2$, orb 8/6, hist-orb 3/2; M 29696, res 10900, rev 10386: ok}\\*
0 & -- & 0 & 4 & 1 & 466 & 2 & 556 & \texttt{-?????????}\\*
1 & 0 & 1 & 3 & 1 & 840 & 8 & 176 & \texttt{--++?????-}\\
2 & 1 & 2 & 3 & 1 & 656 & 16 & 352 & \texttt{---+?????-}\\
3 & 2 & 3 & 3 & 1 & 288 & 32 & 704 & \texttt{----?????-}\\
4 & 3 & 4 & 1 & 2 & 512 & 512 & 0 & \texttt{-----+----}\\
5 & 4 & 5 & 1 & 2 & 0 & 1024 & 0 & \texttt{----------}\\
6 & 3 & 4 & 1 & 2 & 512 & 512 & 0 & \texttt{-------+--}\\
7 & 6 & 5 & 1 & 2 & 0 & 1024 & 0 & \texttt{----------}\\
8 & 3 & 4 & 1 & 1 & 0 & 1024 & 0 & \texttt{----------}\\
9 & 0 & 1 & 1 & 1 & 0 & 1024 & 0 & \texttt{----------}\\
\multicolumn{9}{@{}l@{}}{\rule{0pt}{2.6ex}\textbf{T10.205}\enspace \texttt{0123454534}\enspace $n{=}10$, $|H|{=}3$, len 6, br 2, $|\mathrm{Aut}|{=}2$, orb 8/6, hist-orb 2/1; M 31744, res 6296, rev 13390: ok}\\*
0 & -- & 0 & 3 & 1 & 918 & 2 & 104 & \texttt{-++???????}\\*
1 & 0 & 1 & 3 & 1 & 812 & 4 & 208 & \texttt{--+???????}\\
2 & 1 & 2 & 3 & 1 & 600 & 8 & 416 & \texttt{---???????}\\
3 & 2 & 3 & 2 & 1 & 576 & 64 & 384 & \texttt{----????--}\\
4 & 3 & 4 & 1 & 2 & 512 & 512 & 0 & \texttt{-----+----}\\
5 & 4 & 5 & 1 & 2 & 0 & 1024 & 0 & \texttt{----------}\\
6 & 3 & 4 & 1 & 2 & 512 & 512 & 0 & \texttt{-------+--}\\
7 & 6 & 5 & 1 & 2 & 0 & 1024 & 0 & \texttt{----------}\\
8 & 2 & 3 & 1 & 1 & 512 & 512 & 0 & \texttt{---------+}\\
9 & 8 & 4 & 1 & 1 & 0 & 1024 & 0 & \texttt{----------}\\
\multicolumn{9}{@{}l@{}}{\rule{0pt}{2.6ex}\textbf{T10.206}\enspace \texttt{0123454533}\enspace $n{=}10$, $|H|{=}4$, len 6, br 2, $|\mathrm{Aut}|{=}4$, orb 7/4, hist-orb 2/0; M 30720, res 11196, rev 10774: ok}\\*
0 & -- & 0 & 4 & 1 & 818 & 2 & 204 & \texttt{-++???????}\\*
1 & 0 & 1 & 4 & 1 & 612 & 4 & 408 & \texttt{--+???????}\\
2 & 1 & 2 & 4 & 1 & 200 & 8 & 816 & \texttt{---???????}\\
3 & 2 & 3 & 2 & 1 & 576 & 64 & 384 & \texttt{----????--}\\
4 & 3 & 4 & 1 & 2 & 512 & 512 & 0 & \texttt{-----+----}\\
5 & 4 & 5 & 1 & 2 & 0 & 1024 & 0 & \texttt{----------}\\
6 & 3 & 4 & 1 & 2 & 512 & 512 & 0 & \texttt{-------+--}\\
7 & 6 & 5 & 1 & 2 & 0 & 1024 & 0 & \texttt{----------}\\
8 & 2 & 3 & 1 & 2 & 0 & 1024 & 0 & \texttt{----------}\\
9 & 2 & 3 & 1 & 2 & 0 & 1024 & 0 & \texttt{----------}\\
\multicolumn{9}{@{}l@{}}{\rule{0pt}{2.6ex}\textbf{T10.207}\enspace \texttt{0123454532}\enspace $n{=}10$, $|H|{=}4$, len 6, br 3, $|\mathrm{Aut}|{=}2$, orb 8/6, hist-orb 3/2; M 29696, res 12360, rev 9418: ok}\\*
0 & -- & 0 & 4 & 1 & 690 & 2 & 332 & \texttt{-+????????}\\*
1 & 0 & 1 & 4 & 1 & 356 & 4 & 664 & \texttt{--????????}\\
2 & 1 & 2 & 3 & 1 & 400 & 16 & 608 & \texttt{---??????-}\\
3 & 2 & 3 & 2 & 1 & 576 & 64 & 384 & \texttt{----????--}\\
4 & 3 & 4 & 1 & 2 & 512 & 512 & 0 & \texttt{-----+----}\\
5 & 4 & 5 & 1 & 2 & 0 & 1024 & 0 & \texttt{----------}\\
6 & 3 & 4 & 1 & 2 & 512 & 512 & 0 & \texttt{-------+--}\\
7 & 6 & 5 & 1 & 2 & 0 & 1024 & 0 & \texttt{----------}\\
8 & 2 & 3 & 1 & 1 & 0 & 1024 & 0 & \texttt{----------}\\
9 & 1 & 2 & 1 & 1 & 0 & 1024 & 0 & \texttt{----------}\\
\multicolumn{9}{@{}l@{}}{\rule{0pt}{2.6ex}\textbf{T10.208}\enspace \texttt{0123454531}\enspace $n{=}10$, $|H|{=}4$, len 6, br 3, $|\mathrm{Aut}|{=}2$, orb 8/6, hist-orb 3/2; M 28672, res 11660, rev 9198: ok}\\*
0 & -- & 0 & 4 & 1 & 434 & 2 & 588 & \texttt{-?????????}\\*
1 & 0 & 1 & 3 & 1 & 712 & 8 & 304 & \texttt{--+??????-}\\
2 & 1 & 2 & 3 & 1 & 400 & 16 & 608 & \texttt{---??????-}\\
3 & 2 & 3 & 2 & 1 & 576 & 64 & 384 & \texttt{----????--}\\
4 & 3 & 4 & 1 & 2 & 512 & 512 & 0 & \texttt{-----+----}\\
5 & 4 & 5 & 1 & 2 & 0 & 1024 & 0 & \texttt{----------}\\
6 & 3 & 4 & 1 & 2 & 512 & 512 & 0 & \texttt{-------+--}\\
7 & 6 & 5 & 1 & 2 & 0 & 1024 & 0 & \texttt{----------}\\
8 & 2 & 3 & 1 & 1 & 0 & 1024 & 0 & \texttt{----------}\\
9 & 0 & 1 & 1 & 1 & 0 & 1024 & 0 & \texttt{----------}\\
\multicolumn{9}{@{}l@{}}{\rule{0pt}{2.6ex}\textbf{T10.209}\enspace \texttt{0123454523}\enspace $n{=}10$, $|H|{=}3$, len 6, br 2, $|\mathrm{Aut}|{=}2$, orb 8/6, hist-orb 2/1; M 29696, res 6048, rev 12226: ok}\\*
0 & -- & 0 & 3 & 1 & 854 & 2 & 168 & \texttt{-+????????}\\*
1 & 0 & 1 & 3 & 1 & 684 & 4 & 336 & \texttt{--????????}\\
2 & 1 & 2 & 2 & 1 & 800 & 32 & 192 & \texttt{---+????--}\\
3 & 2 & 3 & 2 & 1 & 576 & 64 & 384 & \texttt{----????--}\\
4 & 3 & 4 & 1 & 2 & 512 & 512 & 0 & \texttt{-----+----}\\
5 & 4 & 5 & 1 & 2 & 0 & 1024 & 0 & \texttt{----------}\\
6 & 3 & 4 & 1 & 2 & 512 & 512 & 0 & \texttt{-------+--}\\
7 & 6 & 5 & 1 & 2 & 0 & 1024 & 0 & \texttt{----------}\\
8 & 1 & 2 & 1 & 1 & 512 & 512 & 0 & \texttt{---------+}\\
9 & 8 & 3 & 1 & 1 & 0 & 1024 & 0 & \texttt{----------}\\
\multicolumn{9}{@{}l@{}}{\rule{0pt}{2.6ex}\textbf{T10.210}\enspace \texttt{0123454522}\enspace $n{=}10$, $|H|{=}4$, len 6, br 2, $|\mathrm{Aut}|{=}4$, orb 7/4, hist-orb 2/0; M 28672, res 11520, rev 9222: ok}\\*
0 & -- & 0 & 4 & 1 & 626 & 2 & 396 & \texttt{-+????????}\\*
1 & 0 & 1 & 4 & 1 & 228 & 4 & 792 & \texttt{--????????}\\
2 & 1 & 2 & 2 & 1 & 800 & 32 & 192 & \texttt{---+????--}\\
3 & 2 & 3 & 2 & 1 & 576 & 64 & 384 & \texttt{----????--}\\
4 & 3 & 4 & 1 & 2 & 512 & 512 & 0 & \texttt{-----+----}\\
5 & 4 & 5 & 1 & 2 & 0 & 1024 & 0 & \texttt{----------}\\
6 & 3 & 4 & 1 & 2 & 512 & 512 & 0 & \texttt{-------+--}\\
7 & 6 & 5 & 1 & 2 & 0 & 1024 & 0 & \texttt{----------}\\
8 & 1 & 2 & 1 & 2 & 0 & 1024 & 0 & \texttt{----------}\\
9 & 1 & 2 & 1 & 2 & 0 & 1024 & 0 & \texttt{----------}\\
\multicolumn{9}{@{}l@{}}{\rule{0pt}{2.6ex}\textbf{T10.211}\enspace \texttt{0123454521}\enspace $n{=}10$, $|H|{=}4$, len 6, br 3, $|\mathrm{Aut}|{=}2$, orb 8/6, hist-orb 3/2; M 27648, res 11524, rev 8606: ok}\\*
0 & -- & 0 & 4 & 1 & 370 & 2 & 652 & \texttt{-?????????}\\*
1 & 0 & 1 & 3 & 1 & 456 & 8 & 560 & \texttt{--???????-}\\
2 & 1 & 2 & 2 & 1 & 800 & 32 & 192 & \texttt{---+????--}\\
3 & 2 & 3 & 2 & 1 & 576 & 64 & 384 & \texttt{----????--}\\
4 & 3 & 4 & 1 & 2 & 512 & 512 & 0 & \texttt{-----+----}\\
5 & 4 & 5 & 1 & 2 & 0 & 1024 & 0 & \texttt{----------}\\
6 & 3 & 4 & 1 & 2 & 512 & 512 & 0 & \texttt{-------+--}\\
7 & 6 & 5 & 1 & 2 & 0 & 1024 & 0 & \texttt{----------}\\
8 & 1 & 2 & 1 & 1 & 0 & 1024 & 0 & \texttt{----------}\\
9 & 0 & 1 & 1 & 1 & 0 & 1024 & 0 & \texttt{----------}\\
\multicolumn{9}{@{}l@{}}{\rule{0pt}{2.6ex}\textbf{T10.212}\enspace \texttt{0123454512}\enspace $n{=}10$, $|H|{=}3$, len 6, br 2, $|\mathrm{Aut}|{=}2$, orb 8/6, hist-orb 2/1; M 27648, res 5208, rev 11354: ok}\\*
0 & -- & 0 & 3 & 1 & 726 & 2 & 296 & \texttt{-?????????}\\*
1 & 0 & 1 & 2 & 1 & 912 & 16 & 96 & \texttt{--++????--}\\
2 & 1 & 2 & 2 & 1 & 800 & 32 & 192 & \texttt{---+????--}\\
3 & 2 & 3 & 2 & 1 & 576 & 64 & 384 & \texttt{----????--}\\
4 & 3 & 4 & 1 & 2 & 512 & 512 & 0 & \texttt{-----+----}\\
5 & 4 & 5 & 1 & 2 & 0 & 1024 & 0 & \texttt{----------}\\
6 & 3 & 4 & 1 & 2 & 512 & 512 & 0 & \texttt{-------+--}\\
7 & 6 & 5 & 1 & 2 & 0 & 1024 & 0 & \texttt{----------}\\
8 & 0 & 1 & 1 & 1 & 512 & 512 & 0 & \texttt{---------+}\\
9 & 8 & 2 & 1 & 1 & 0 & 1024 & 0 & \texttt{----------}\\
\multicolumn{9}{@{}l@{}}{\rule{0pt}{2.6ex}\textbf{T10.213}\enspace \texttt{0123454511}\enspace $n{=}10$, $|H|{=}4$, len 6, br 2, $|\mathrm{Aut}|{=}4$, orb 7/4, hist-orb 2/0; M 26624, res 9564, rev 8874: ok}\\*
0 & -- & 0 & 4 & 1 & 242 & 2 & 780 & \texttt{-?????????}\\*
1 & 0 & 1 & 2 & 1 & 912 & 16 & 96 & \texttt{--++????--}\\
2 & 1 & 2 & 2 & 1 & 800 & 32 & 192 & \texttt{---+????--}\\
3 & 2 & 3 & 2 & 1 & 576 & 64 & 384 & \texttt{----????--}\\
4 & 3 & 4 & 1 & 2 & 512 & 512 & 0 & \texttt{-----+----}\\
5 & 4 & 5 & 1 & 2 & 0 & 1024 & 0 & \texttt{----------}\\
6 & 3 & 4 & 1 & 2 & 512 & 512 & 0 & \texttt{-------+--}\\
7 & 6 & 5 & 1 & 2 & 0 & 1024 & 0 & \texttt{----------}\\
8 & 0 & 1 & 1 & 2 & 0 & 1024 & 0 & \texttt{----------}\\
9 & 0 & 1 & 1 & 2 & 0 & 1024 & 0 & \texttt{----------}\\
\multicolumn{9}{@{}l@{}}{\rule{0pt}{2.6ex}\textbf{T10.214}\enspace \texttt{0123454444}\enspace $n{=}10$, $|H|{=}5$, len 6, br 1, $|\mathrm{Aut}|{=}24$, orb 7/6, hist-orb 2/1; M 31744, res 10800, rev 13006: ok}\\*
0 & -- & 0 & 5 & 1 & 902 & 2 & 120 & \texttt{-+++??????}\\*
1 & 0 & 1 & 5 & 1 & 780 & 4 & 240 & \texttt{--++??????}\\
2 & 1 & 2 & 5 & 1 & 536 & 8 & 480 & \texttt{---+??????}\\
3 & 2 & 3 & 5 & 1 & 48 & 16 & 960 & \texttt{----??????}\\
4 & 3 & 4 & 1 & 1 & 512 & 512 & 0 & \texttt{-----+----}\\
5 & 4 & 5 & 1 & 1 & 0 & 1024 & 0 & \texttt{----------}\\
6 & 3 & 4 & 1 & 4 & 0 & 1024 & 0 & \texttt{----------}\\
7 & 3 & 4 & 1 & 4 & 0 & 1024 & 0 & \texttt{----------}\\
8 & 3 & 4 & 1 & 4 & 0 & 1024 & 0 & \texttt{----------}\\
9 & 3 & 4 & 1 & 4 & 0 & 1024 & 0 & \texttt{----------}\\
\multicolumn{9}{@{}l@{}}{\rule{0pt}{2.6ex}\textbf{T10.215}\enspace \texttt{0123454443}\enspace $n{=}10$, $|H|{=}5$, len 6, br 2, $|\mathrm{Aut}|{=}6$, orb 8/7, hist-orb 3/2; M 30720, res 12712, rev 10722: ok}\\*
0 & -- & 0 & 5 & 1 & 838 & 2 & 184 & \texttt{-++???????}\\*
1 & 0 & 1 & 5 & 1 & 652 & 4 & 368 & \texttt{--+???????}\\
2 & 1 & 2 & 5 & 1 & 280 & 8 & 736 & \texttt{---???????}\\
3 & 2 & 3 & 4 & 1 & 96 & 32 & 896 & \texttt{----?????-}\\
4 & 3 & 4 & 1 & 1 & 512 & 512 & 0 & \texttt{-----+----}\\
5 & 4 & 5 & 1 & 1 & 0 & 1024 & 0 & \texttt{----------}\\
6 & 3 & 4 & 1 & 3 & 0 & 1024 & 0 & \texttt{----------}\\
7 & 3 & 4 & 1 & 3 & 0 & 1024 & 0 & \texttt{----------}\\
8 & 3 & 4 & 1 & 3 & 0 & 1024 & 0 & \texttt{----------}\\
9 & 2 & 3 & 1 & 1 & 0 & 1024 & 0 & \texttt{----------}\\
\multicolumn{9}{@{}l@{}}{\rule{0pt}{2.6ex}\textbf{T10.216}\enspace \texttt{0123454442}\enspace $n{=}10$, $|H|{=}5$, len 6, br 2, $|\mathrm{Aut}|{=}6$, orb 8/7, hist-orb 3/2; M 29696, res 13536, rev 9586: ok}\\*
0 & -- & 0 & 5 & 1 & 710 & 2 & 312 & \texttt{-+????????}\\*
1 & 0 & 1 & 5 & 1 & 396 & 4 & 624 & \texttt{--????????}\\
2 & 1 & 2 & 4 & 1 & 560 & 16 & 448 & \texttt{---+?????-}\\
3 & 2 & 3 & 4 & 1 & 96 & 32 & 896 & \texttt{----?????-}\\
4 & 3 & 4 & 1 & 1 & 512 & 512 & 0 & \texttt{-----+----}\\
5 & 4 & 5 & 1 & 1 & 0 & 1024 & 0 & \texttt{----------}\\
6 & 3 & 4 & 1 & 3 & 0 & 1024 & 0 & \texttt{----------}\\
7 & 3 & 4 & 1 & 3 & 0 & 1024 & 0 & \texttt{----------}\\
8 & 3 & 4 & 1 & 3 & 0 & 1024 & 0 & \texttt{----------}\\
9 & 1 & 2 & 1 & 1 & 0 & 1024 & 0 & \texttt{----------}\\
\multicolumn{9}{@{}l@{}}{\rule{0pt}{2.6ex}\textbf{T10.217}\enspace \texttt{0123454441}\enspace $n{=}10$, $|H|{=}5$, len 6, br 2, $|\mathrm{Aut}|{=}6$, orb 8/7, hist-orb 3/2; M 28672, res 12616, rev 9566: ok}\\*
0 & -- & 0 & 5 & 1 & 454 & 2 & 568 & \texttt{-?????????}\\*
1 & 0 & 1 & 4 & 1 & 792 & 8 & 224 & \texttt{--++?????-}\\
2 & 1 & 2 & 4 & 1 & 560 & 16 & 448 & \texttt{---+?????-}\\
3 & 2 & 3 & 4 & 1 & 96 & 32 & 896 & \texttt{----?????-}\\
4 & 3 & 4 & 1 & 1 & 512 & 512 & 0 & \texttt{-----+----}\\
5 & 4 & 5 & 1 & 1 & 0 & 1024 & 0 & \texttt{----------}\\
6 & 3 & 4 & 1 & 3 & 0 & 1024 & 0 & \texttt{----------}\\
7 & 3 & 4 & 1 & 3 & 0 & 1024 & 0 & \texttt{----------}\\
8 & 3 & 4 & 1 & 3 & 0 & 1024 & 0 & \texttt{----------}\\
9 & 0 & 1 & 1 & 1 & 0 & 1024 & 0 & \texttt{----------}\\
\multicolumn{9}{@{}l@{}}{\rule{0pt}{2.6ex}\textbf{T10.218}\enspace \texttt{0123454434}\enspace $n{=}10$, $|H|{=}4$, len 6, br 2, $|\mathrm{Aut}|{=}2$, orb 9/8, hist-orb 3/2; M 30720, res 9260, rev 12034: ok}\\*
0 & -- & 0 & 4 & 1 & 882 & 2 & 140 & \texttt{-++???????}\\*
1 & 0 & 1 & 4 & 1 & 740 & 4 & 280 & \texttt{--+???????}\\
2 & 1 & 2 & 4 & 1 & 456 & 8 & 560 & \texttt{---???????}\\
3 & 2 & 3 & 3 & 1 & 192 & 64 & 768 & \texttt{----????--}\\
4 & 3 & 4 & 1 & 1 & 512 & 512 & 0 & \texttt{-----+----}\\
5 & 4 & 5 & 1 & 1 & 0 & 1024 & 0 & \texttt{----------}\\
6 & 3 & 4 & 1 & 2 & 0 & 1024 & 0 & \texttt{----------}\\
7 & 3 & 4 & 1 & 2 & 0 & 1024 & 0 & \texttt{----------}\\
8 & 2 & 3 & 1 & 1 & 512 & 512 & 0 & \texttt{---------+}\\
9 & 8 & 4 & 1 & 1 & 0 & 1024 & 0 & \texttt{----------}\\
\multicolumn{9}{@{}l@{}}{\rule{0pt}{2.6ex}\textbf{T10.219}\enspace \texttt{0123454433}\enspace $n{=}10$, $|H|{=}5$, len 6, br 2, $|\mathrm{Aut}|{=}4$, orb 8/6, hist-orb 3/1; M 29696, res 12984, rev 9754: ok}\\*
0 & -- & 0 & 5 & 1 & 806 & 2 & 216 & \texttt{-++???????}\\*
1 & 0 & 1 & 5 & 1 & 588 & 4 & 432 & \texttt{--+???????}\\
2 & 1 & 2 & 5 & 1 & 152 & 8 & 864 & \texttt{---???????}\\
3 & 2 & 3 & 3 & 1 & 192 & 64 & 768 & \texttt{----????--}\\
4 & 3 & 4 & 1 & 1 & 512 & 512 & 0 & \texttt{-----+----}\\
5 & 4 & 5 & 1 & 1 & 0 & 1024 & 0 & \texttt{----------}\\
6 & 3 & 4 & 1 & 2 & 0 & 1024 & 0 & \texttt{----------}\\
7 & 3 & 4 & 1 & 2 & 0 & 1024 & 0 & \texttt{----------}\\
8 & 2 & 3 & 1 & 2 & 0 & 1024 & 0 & \texttt{----------}\\
9 & 2 & 3 & 1 & 2 & 0 & 1024 & 0 & \texttt{----------}\\
\multicolumn{9}{@{}l@{}}{\rule{0pt}{2.6ex}\textbf{T10.220}\enspace \texttt{0123454432}\enspace $n{=}10$, $|H|{=}5$, len 6, br 3, $|\mathrm{Aut}|{=}2$, orb 9/8, hist-orb 4/3; M 28672, res 14352, rev 8158: ok}\\*
0 & -- & 0 & 5 & 1 & 678 & 2 & 344 & \texttt{-+????????}\\*
1 & 0 & 1 & 5 & 1 & 332 & 4 & 688 & \texttt{--????????}\\
2 & 1 & 2 & 4 & 1 & 304 & 16 & 704 & \texttt{---??????-}\\
3 & 2 & 3 & 3 & 1 & 192 & 64 & 768 & \texttt{----????--}\\
4 & 3 & 4 & 1 & 1 & 512 & 512 & 0 & \texttt{-----+----}\\
5 & 4 & 5 & 1 & 1 & 0 & 1024 & 0 & \texttt{----------}\\
6 & 3 & 4 & 1 & 2 & 0 & 1024 & 0 & \texttt{----------}\\
7 & 3 & 4 & 1 & 2 & 0 & 1024 & 0 & \texttt{----------}\\
8 & 2 & 3 & 1 & 1 & 0 & 1024 & 0 & \texttt{----------}\\
9 & 1 & 2 & 1 & 1 & 0 & 1024 & 0 & \texttt{----------}\\
\multicolumn{9}{@{}l@{}}{\rule{0pt}{2.6ex}\textbf{T10.221}\enspace \texttt{0123454431}\enspace $n{=}10$, $|H|{=}5$, len 6, br 3, $|\mathrm{Aut}|{=}2$, orb 9/8, hist-orb 4/3; M 27648, res 13784, rev 7914: ok}\\*
0 & -- & 0 & 5 & 1 & 422 & 2 & 600 & \texttt{-?????????}\\*
1 & 0 & 1 & 4 & 1 & 664 & 8 & 352 & \texttt{--+??????-}\\
2 & 1 & 2 & 4 & 1 & 304 & 16 & 704 & \texttt{---??????-}\\
3 & 2 & 3 & 3 & 1 & 192 & 64 & 768 & \texttt{----????--}\\
4 & 3 & 4 & 1 & 1 & 512 & 512 & 0 & \texttt{-----+----}\\
5 & 4 & 5 & 1 & 1 & 0 & 1024 & 0 & \texttt{----------}\\
6 & 3 & 4 & 1 & 2 & 0 & 1024 & 0 & \texttt{----------}\\
7 & 3 & 4 & 1 & 2 & 0 & 1024 & 0 & \texttt{----------}\\
8 & 2 & 3 & 1 & 1 & 0 & 1024 & 0 & \texttt{----------}\\
9 & 0 & 1 & 1 & 1 & 0 & 1024 & 0 & \texttt{----------}\\
\multicolumn{9}{@{}l@{}}{\rule{0pt}{2.6ex}\textbf{T10.222}\enspace \texttt{0123454423}\enspace $n{=}10$, $|H|{=}4$, len 6, br 2, $|\mathrm{Aut}|{=}2$, orb 9/8, hist-orb 3/2; M 28672, res 8928, rev 10918: ok}\\*
0 & -- & 0 & 4 & 1 & 818 & 2 & 204 & \texttt{-+????????}\\*
1 & 0 & 1 & 4 & 1 & 612 & 4 & 408 & \texttt{--????????}\\
2 & 1 & 2 & 3 & 1 & 608 & 32 & 384 & \texttt{---+????--}\\
3 & 2 & 3 & 3 & 1 & 192 & 64 & 768 & \texttt{----????--}\\
4 & 3 & 4 & 1 & 1 & 512 & 512 & 0 & \texttt{-----+----}\\
5 & 4 & 5 & 1 & 1 & 0 & 1024 & 0 & \texttt{----------}\\
6 & 3 & 4 & 1 & 2 & 0 & 1024 & 0 & \texttt{----------}\\
7 & 3 & 4 & 1 & 2 & 0 & 1024 & 0 & \texttt{----------}\\
8 & 1 & 2 & 1 & 1 & 512 & 512 & 0 & \texttt{---------+}\\
9 & 8 & 3 & 1 & 1 & 0 & 1024 & 0 & \texttt{----------}\\
\multicolumn{9}{@{}l@{}}{\rule{0pt}{2.6ex}\textbf{T10.223}\enspace \texttt{0123454422}\enspace $n{=}10$, $|H|{=}5$, len 6, br 2, $|\mathrm{Aut}|{=}4$, orb 8/6, hist-orb 3/1; M 27648, res 13824, rev 7866: ok}\\*
0 & -- & 0 & 5 & 1 & 614 & 2 & 408 & \texttt{-+????????}\\*
1 & 0 & 1 & 5 & 1 & 204 & 4 & 816 & \texttt{--????????}\\
2 & 1 & 2 & 3 & 1 & 608 & 32 & 384 & \texttt{---+????--}\\
3 & 2 & 3 & 3 & 1 & 192 & 64 & 768 & \texttt{----????--}\\
4 & 3 & 4 & 1 & 1 & 512 & 512 & 0 & \texttt{-----+----}\\
5 & 4 & 5 & 1 & 1 & 0 & 1024 & 0 & \texttt{----------}\\
6 & 3 & 4 & 1 & 2 & 0 & 1024 & 0 & \texttt{----------}\\
7 & 3 & 4 & 1 & 2 & 0 & 1024 & 0 & \texttt{----------}\\
8 & 1 & 2 & 1 & 2 & 0 & 1024 & 0 & \texttt{----------}\\
9 & 1 & 2 & 1 & 2 & 0 & 1024 & 0 & \texttt{----------}\\
\multicolumn{9}{@{}l@{}}{\rule{0pt}{2.6ex}\textbf{T10.224}\enspace \texttt{0123454421}\enspace $n{=}10$, $|H|{=}5$, len 6, br 3, $|\mathrm{Aut}|{=}2$, orb 9/8, hist-orb 4/3; M 26624, res 13960, rev 7162: ok}\\*
0 & -- & 0 & 5 & 1 & 358 & 2 & 664 & \texttt{-?????????}\\*
1 & 0 & 1 & 4 & 1 & 408 & 8 & 608 & \texttt{--???????-}\\
2 & 1 & 2 & 3 & 1 & 608 & 32 & 384 & \texttt{---+????--}\\
3 & 2 & 3 & 3 & 1 & 192 & 64 & 768 & \texttt{----????--}\\
4 & 3 & 4 & 1 & 1 & 512 & 512 & 0 & \texttt{-----+----}\\
5 & 4 & 5 & 1 & 1 & 0 & 1024 & 0 & \texttt{----------}\\
6 & 3 & 4 & 1 & 2 & 0 & 1024 & 0 & \texttt{----------}\\
7 & 3 & 4 & 1 & 2 & 0 & 1024 & 0 & \texttt{----------}\\
8 & 1 & 2 & 1 & 1 & 0 & 1024 & 0 & \texttt{----------}\\
9 & 0 & 1 & 1 & 1 & 0 & 1024 & 0 & \texttt{----------}\\
\multicolumn{9}{@{}l@{}}{\rule{0pt}{2.6ex}\textbf{T10.225}\enspace \texttt{0123454412}\enspace $n{=}10$, $|H|{=}4$, len 6, br 2, $|\mathrm{Aut}|{=}2$, orb 9/8, hist-orb 3/2; M 26624, res 8076, rev 10118: ok}\\*
0 & -- & 0 & 4 & 1 & 690 & 2 & 332 & \texttt{-?????????}\\*
1 & 0 & 1 & 3 & 1 & 816 & 16 & 192 & \texttt{--++????--}\\
2 & 1 & 2 & 3 & 1 & 608 & 32 & 384 & \texttt{---+????--}\\
3 & 2 & 3 & 3 & 1 & 192 & 64 & 768 & \texttt{----????--}\\
4 & 3 & 4 & 1 & 1 & 512 & 512 & 0 & \texttt{-----+----}\\
5 & 4 & 5 & 1 & 1 & 0 & 1024 & 0 & \texttt{----------}\\
6 & 3 & 4 & 1 & 2 & 0 & 1024 & 0 & \texttt{----------}\\
7 & 3 & 4 & 1 & 2 & 0 & 1024 & 0 & \texttt{----------}\\
8 & 0 & 1 & 1 & 1 & 512 & 512 & 0 & \texttt{---------+}\\
9 & 8 & 2 & 1 & 1 & 0 & 1024 & 0 & \texttt{----------}\\
\multicolumn{9}{@{}l@{}}{\rule{0pt}{2.6ex}\textbf{T10.226}\enspace \texttt{0123454411}\enspace $n{=}10$, $|H|{=}5$, len 6, br 2, $|\mathrm{Aut}|{=}4$, orb 8/6, hist-orb 3/1; M 25600, res 12216, rev 7494: ok}\\*
0 & -- & 0 & 5 & 1 & 230 & 2 & 792 & \texttt{-?????????}\\*
1 & 0 & 1 & 3 & 1 & 816 & 16 & 192 & \texttt{--++????--}\\
2 & 1 & 2 & 3 & 1 & 608 & 32 & 384 & \texttt{---+????--}\\
3 & 2 & 3 & 3 & 1 & 192 & 64 & 768 & \texttt{----????--}\\
4 & 3 & 4 & 1 & 1 & 512 & 512 & 0 & \texttt{-----+----}\\
5 & 4 & 5 & 1 & 1 & 0 & 1024 & 0 & \texttt{----------}\\
6 & 3 & 4 & 1 & 2 & 0 & 1024 & 0 & \texttt{----------}\\
7 & 3 & 4 & 1 & 2 & 0 & 1024 & 0 & \texttt{----------}\\
8 & 0 & 1 & 1 & 2 & 0 & 1024 & 0 & \texttt{----------}\\
9 & 0 & 1 & 1 & 2 & 0 & 1024 & 0 & \texttt{----------}\\
\multicolumn{9}{@{}l@{}}{\rule{0pt}{2.6ex}\textbf{T10.227}\enspace \texttt{0123454345}\enspace $n{=}10$, $|H|{=}3$, len 6, br 2, $|\mathrm{Aut}|{=}1$, orb 10/10, hist-orb 3/3; M 31744, res 5988, rev 12994: ok}\\*
0 & -- & 0 & 3 & 1 & 922 & 2 & 100 & \texttt{-++???????}\\*
1 & 0 & 1 & 3 & 1 & 820 & 4 & 200 & \texttt{--+???????}\\
2 & 1 & 2 & 3 & 1 & 616 & 8 & 400 & \texttt{---???????}\\
3 & 2 & 3 & 2 & 1 & 384 & 128 & 512 & \texttt{----???---}\\
4 & 3 & 4 & 1 & 1 & 512 & 512 & 0 & \texttt{-----+----}\\
5 & 4 & 5 & 1 & 1 & 0 & 1024 & 0 & \texttt{----------}\\
6 & 3 & 4 & 1 & 1 & 0 & 1024 & 0 & \texttt{----------}\\
7 & 2 & 3 & 1 & 1 & 768 & 256 & 0 & \texttt{--------++}\\
8 & 7 & 4 & 1 & 1 & 512 & 512 & 0 & \texttt{---------+}\\
9 & 8 & 5 & 1 & 1 & 0 & 1024 & 0 & \texttt{----------}\\
\multicolumn{9}{@{}l@{}}{\rule{0pt}{2.6ex}\textbf{T10.228}\enspace \texttt{0123454344}\enspace $n{=}10$, $|H|{=}4$, len 6, br 3, $|\mathrm{Aut}|{=}2$, orb 9/8, hist-orb 3/2; M 30720, res 8720, rev 11826: ok}\\*
0 & -- & 0 & 4 & 1 & 878 & 2 & 144 & \texttt{-++???????}\\*
1 & 0 & 1 & 4 & 1 & 732 & 4 & 288 & \texttt{--+???????}\\
2 & 1 & 2 & 4 & 1 & 440 & 8 & 576 & \texttt{---???????}\\
3 & 2 & 3 & 2 & 1 & 384 & 128 & 512 & \texttt{----???---}\\
4 & 3 & 4 & 1 & 1 & 512 & 512 & 0 & \texttt{-----+----}\\
5 & 4 & 5 & 1 & 1 & 0 & 1024 & 0 & \texttt{----------}\\
6 & 3 & 4 & 1 & 1 & 0 & 1024 & 0 & \texttt{----------}\\
7 & 2 & 3 & 2 & 1 & 256 & 256 & 512 & \texttt{--------??}\\
8 & 7 & 4 & 1 & 2 & 0 & 1024 & 0 & \texttt{----------}\\
9 & 7 & 4 & 1 & 2 & 0 & 1024 & 0 & \texttt{----------}\\
\multicolumn{9}{@{}l@{}}{\rule{0pt}{2.6ex}\textbf{T10.229}\enspace \texttt{0123454343}\enspace $n{=}10$, $|H|{=}4$, len 6, br 2, $|\mathrm{Aut}|{=}1$, orb 10/10, hist-orb 4/4; M 29696, res 10300, rev 10746: ok}\\*
0 & -- & 0 & 4 & 1 & 834 & 2 & 188 & \texttt{-++???????}\\*
1 & 0 & 1 & 4 & 1 & 644 & 4 & 376 & \texttt{--+???????}\\
2 & 1 & 2 & 4 & 1 & 264 & 8 & 752 & \texttt{---???????}\\
3 & 2 & 3 & 2 & 1 & 384 & 128 & 512 & \texttt{----???---}\\
4 & 3 & 4 & 1 & 1 & 512 & 512 & 0 & \texttt{-----+----}\\
5 & 4 & 5 & 1 & 1 & 0 & 1024 & 0 & \texttt{----------}\\
6 & 3 & 4 & 1 & 1 & 0 & 1024 & 0 & \texttt{----------}\\
7 & 2 & 3 & 1 & 1 & 512 & 512 & 0 & \texttt{--------+-}\\
8 & 7 & 4 & 1 & 1 & 0 & 1024 & 0 & \texttt{----------}\\
9 & 2 & 3 & 1 & 1 & 0 & 1024 & 0 & \texttt{----------}\\
\multicolumn{9}{@{}l@{}}{\rule{0pt}{2.6ex}\textbf{T10.230}\enspace \texttt{0123454342}\enspace $n{=}10$, $|H|{=}4$, len 6, br 3, $|\mathrm{Aut}|{=}1$, orb 10/10, hist-orb 4/4; M 28672, res 11192, rev 9542: ok}\\*
0 & -- & 0 & 4 & 1 & 706 & 2 & 316 & \texttt{-+????????}\\*
1 & 0 & 1 & 4 & 1 & 388 & 4 & 632 & \texttt{--????????}\\
2 & 1 & 2 & 3 & 1 & 528 & 16 & 480 & \texttt{---??????-}\\
3 & 2 & 3 & 2 & 1 & 384 & 128 & 512 & \texttt{----???---}\\
4 & 3 & 4 & 1 & 1 & 512 & 512 & 0 & \texttt{-----+----}\\
5 & 4 & 5 & 1 & 1 & 0 & 1024 & 0 & \texttt{----------}\\
6 & 3 & 4 & 1 & 1 & 0 & 1024 & 0 & \texttt{----------}\\
7 & 2 & 3 & 1 & 1 & 512 & 512 & 0 & \texttt{--------+-}\\
8 & 7 & 4 & 1 & 1 & 0 & 1024 & 0 & \texttt{----------}\\
9 & 1 & 2 & 1 & 1 & 0 & 1024 & 0 & \texttt{----------}\\
\multicolumn{9}{@{}l@{}}{\rule{0pt}{2.6ex}\textbf{T10.231}\enspace \texttt{0123454341}\enspace $n{=}10$, $|H|{=}4$, len 6, br 3, $|\mathrm{Aut}|{=}1$, orb 10/10, hist-orb 4/4; M 27648, res 10316, rev 9410: ok}\\*
0 & -- & 0 & 4 & 1 & 450 & 2 & 572 & \texttt{-?????????}\\*
1 & 0 & 1 & 3 & 1 & 776 & 8 & 240 & \texttt{--+??????-}\\
2 & 1 & 2 & 3 & 1 & 528 & 16 & 480 & \texttt{---??????-}\\
3 & 2 & 3 & 2 & 1 & 384 & 128 & 512 & \texttt{----???---}\\
4 & 3 & 4 & 1 & 1 & 512 & 512 & 0 & \texttt{-----+----}\\
5 & 4 & 5 & 1 & 1 & 0 & 1024 & 0 & \texttt{----------}\\
6 & 3 & 4 & 1 & 1 & 0 & 1024 & 0 & \texttt{----------}\\
7 & 2 & 3 & 1 & 1 & 512 & 512 & 0 & \texttt{--------+-}\\
8 & 7 & 4 & 1 & 1 & 0 & 1024 & 0 & \texttt{----------}\\
9 & 0 & 1 & 1 & 1 & 0 & 1024 & 0 & \texttt{----------}\\
\multicolumn{9}{@{}l@{}}{\rule{0pt}{2.6ex}\textbf{T10.232}\enspace \texttt{0123454333}\enspace $n{=}10$, $|H|{=}5$, len 6, br 2, $|\mathrm{Aut}|{=}6$, orb 8/7, hist-orb 3/2; M 28672, res 12456, rev 9474: ok}\\*
0 & -- & 0 & 5 & 1 & 790 & 2 & 232 & \texttt{-++???????}\\*
1 & 0 & 1 & 5 & 1 & 556 & 4 & 464 & \texttt{--+???????}\\
2 & 1 & 2 & 5 & 1 & 88 & 8 & 928 & \texttt{---???????}\\
3 & 2 & 3 & 2 & 1 & 384 & 128 & 512 & \texttt{----???---}\\
4 & 3 & 4 & 1 & 1 & 512 & 512 & 0 & \texttt{-----+----}\\
5 & 4 & 5 & 1 & 1 & 0 & 1024 & 0 & \texttt{----------}\\
6 & 3 & 4 & 1 & 1 & 0 & 1024 & 0 & \texttt{----------}\\
7 & 2 & 3 & 1 & 3 & 0 & 1024 & 0 & \texttt{----------}\\
8 & 2 & 3 & 1 & 3 & 0 & 1024 & 0 & \texttt{----------}\\
9 & 2 & 3 & 1 & 3 & 0 & 1024 & 0 & \texttt{----------}\\
\multicolumn{9}{@{}l@{}}{\rule{0pt}{2.6ex}\textbf{T10.233}\enspace \texttt{0123454332}\enspace $n{=}10$, $|H|{=}5$, len 6, br 3, $|\mathrm{Aut}|{=}2$, orb 9/8, hist-orb 4/3; M 27648, res 14096, rev 7630: ok}\\*
0 & -- & 0 & 5 & 1 & 662 & 2 & 360 & \texttt{-+????????}\\*
1 & 0 & 1 & 5 & 1 & 300 & 4 & 720 & \texttt{--????????}\\
2 & 1 & 2 & 4 & 1 & 176 & 16 & 832 & \texttt{---??????-}\\
3 & 2 & 3 & 2 & 1 & 384 & 128 & 512 & \texttt{----???---}\\
4 & 3 & 4 & 1 & 1 & 512 & 512 & 0 & \texttt{-----+----}\\
5 & 4 & 5 & 1 & 1 & 0 & 1024 & 0 & \texttt{----------}\\
6 & 3 & 4 & 1 & 1 & 0 & 1024 & 0 & \texttt{----------}\\
7 & 2 & 3 & 1 & 2 & 0 & 1024 & 0 & \texttt{----------}\\
8 & 2 & 3 & 1 & 2 & 0 & 1024 & 0 & \texttt{----------}\\
9 & 1 & 2 & 1 & 1 & 0 & 1024 & 0 & \texttt{----------}\\
\multicolumn{9}{@{}l@{}}{\rule{0pt}{2.6ex}\textbf{T10.234}\enspace \texttt{0123454331}\enspace $n{=}10$, $|H|{=}5$, len 6, br 3, $|\mathrm{Aut}|{=}2$, orb 9/8, hist-orb 4/3; M 26624, res 13704, rev 7274: ok}\\*
0 & -- & 0 & 5 & 1 & 406 & 2 & 616 & \texttt{-?????????}\\*
1 & 0 & 1 & 4 & 1 & 600 & 8 & 416 & \texttt{--+??????-}\\
2 & 1 & 2 & 4 & 1 & 176 & 16 & 832 & \texttt{---??????-}\\
3 & 2 & 3 & 2 & 1 & 384 & 128 & 512 & \texttt{----???---}\\
4 & 3 & 4 & 1 & 1 & 512 & 512 & 0 & \texttt{-----+----}\\
5 & 4 & 5 & 1 & 1 & 0 & 1024 & 0 & \texttt{----------}\\
6 & 3 & 4 & 1 & 1 & 0 & 1024 & 0 & \texttt{----------}\\
7 & 2 & 3 & 1 & 2 & 0 & 1024 & 0 & \texttt{----------}\\
8 & 2 & 3 & 1 & 2 & 0 & 1024 & 0 & \texttt{----------}\\
9 & 0 & 1 & 1 & 1 & 0 & 1024 & 0 & \texttt{----------}\\
\multicolumn{9}{@{}l@{}}{\rule{0pt}{2.6ex}\textbf{T10.235}\enspace \texttt{0123454323}\enspace $n{=}10$, $|H|{=}4$, len 6, br 3, $|\mathrm{Aut}|{=}1$, orb 10/10, hist-orb 4/4; M 27648, res 9856, rev 9650: ok}\\*
0 & -- & 0 & 4 & 1 & 770 & 2 & 252 & \texttt{-+????????}\\*
1 & 0 & 1 & 4 & 1 & 516 & 4 & 504 & \texttt{--????????}\\
2 & 1 & 2 & 3 & 1 & 352 & 32 & 640 & \texttt{---?????--}\\
3 & 2 & 3 & 2 & 1 & 384 & 128 & 512 & \texttt{----???---}\\
4 & 3 & 4 & 1 & 1 & 512 & 512 & 0 & \texttt{-----+----}\\
5 & 4 & 5 & 1 & 1 & 0 & 1024 & 0 & \texttt{----------}\\
6 & 3 & 4 & 1 & 1 & 0 & 1024 & 0 & \texttt{----------}\\
7 & 2 & 3 & 1 & 1 & 0 & 1024 & 0 & \texttt{----------}\\
8 & 1 & 2 & 1 & 1 & 512 & 512 & 0 & \texttt{---------+}\\
9 & 8 & 3 & 1 & 1 & 0 & 1024 & 0 & \texttt{----------}\\
\multicolumn{9}{@{}l@{}}{\rule{0pt}{2.6ex}\textbf{T10.236}\enspace \texttt{0123454322}\enspace $n{=}10$, $|H|{=}5$, len 6, br 3, $|\mathrm{Aut}|{=}2$, orb 9/8, hist-orb 4/3; M 26624, res 13984, rev 7006: ok}\\*
0 & -- & 0 & 5 & 1 & 598 & 2 & 424 & \texttt{-+????????}\\*
1 & 0 & 1 & 5 & 1 & 172 & 4 & 848 & \texttt{--????????}\\
2 & 1 & 2 & 3 & 1 & 352 & 32 & 640 & \texttt{---?????--}\\
3 & 2 & 3 & 2 & 1 & 384 & 128 & 512 & \texttt{----???---}\\
4 & 3 & 4 & 1 & 1 & 512 & 512 & 0 & \texttt{-----+----}\\
5 & 4 & 5 & 1 & 1 & 0 & 1024 & 0 & \texttt{----------}\\
6 & 3 & 4 & 1 & 1 & 0 & 1024 & 0 & \texttt{----------}\\
7 & 2 & 3 & 1 & 1 & 0 & 1024 & 0 & \texttt{----------}\\
8 & 1 & 2 & 1 & 2 & 0 & 1024 & 0 & \texttt{----------}\\
9 & 1 & 2 & 1 & 2 & 0 & 1024 & 0 & \texttt{----------}\\
\multicolumn{9}{@{}l@{}}{\rule{0pt}{2.6ex}\textbf{T10.237}\enspace \texttt{0123454321}\enspace $n{=}10$, $|H|{=}5$, len 6, br 4, $|\mathrm{Aut}|{=}1$, orb 10/10, hist-orb 5/5; M 25600, res 14296, rev 6182: ok}\\*
0 & -- & 0 & 5 & 1 & 342 & 2 & 680 & \texttt{-?????????}\\*
1 & 0 & 1 & 4 & 1 & 344 & 8 & 672 & \texttt{--???????-}\\
2 & 1 & 2 & 3 & 1 & 352 & 32 & 640 & \texttt{---?????--}\\
3 & 2 & 3 & 2 & 1 & 384 & 128 & 512 & \texttt{----???---}\\
4 & 3 & 4 & 1 & 1 & 512 & 512 & 0 & \texttt{-----+----}\\
5 & 4 & 5 & 1 & 1 & 0 & 1024 & 0 & \texttt{----------}\\
6 & 3 & 4 & 1 & 1 & 0 & 1024 & 0 & \texttt{----------}\\
7 & 2 & 3 & 1 & 1 & 0 & 1024 & 0 & \texttt{----------}\\
8 & 1 & 2 & 1 & 1 & 0 & 1024 & 0 & \texttt{----------}\\
9 & 0 & 1 & 1 & 1 & 0 & 1024 & 0 & \texttt{----------}\\
\multicolumn{9}{@{}l@{}}{\rule{0pt}{2.6ex}\textbf{T10.238}\enspace \texttt{0123454312}\enspace $n{=}10$, $|H|{=}4$, len 6, br 3, $|\mathrm{Aut}|{=}1$, orb 10/10, hist-orb 4/4; M 25600, res 8988, rev 8866: ok}\\*
0 & -- & 0 & 4 & 1 & 642 & 2 & 380 & \texttt{-?????????}\\*
1 & 0 & 1 & 3 & 1 & 688 & 16 & 320 & \texttt{--+?????--}\\
2 & 1 & 2 & 3 & 1 & 352 & 32 & 640 & \texttt{---?????--}\\
3 & 2 & 3 & 2 & 1 & 384 & 128 & 512 & \texttt{----???---}\\
4 & 3 & 4 & 1 & 1 & 512 & 512 & 0 & \texttt{-----+----}\\
5 & 4 & 5 & 1 & 1 & 0 & 1024 & 0 & \texttt{----------}\\
6 & 3 & 4 & 1 & 1 & 0 & 1024 & 0 & \texttt{----------}\\
7 & 2 & 3 & 1 & 1 & 0 & 1024 & 0 & \texttt{----------}\\
8 & 0 & 1 & 1 & 1 & 512 & 512 & 0 & \texttt{---------+}\\
9 & 8 & 2 & 1 & 1 & 0 & 1024 & 0 & \texttt{----------}\\
\multicolumn{9}{@{}l@{}}{\rule{0pt}{2.6ex}\textbf{T10.239}\enspace \texttt{0123454311}\enspace $n{=}10$, $|H|{=}5$, len 6, br 3, $|\mathrm{Aut}|{=}2$, orb 9/8, hist-orb 4/3; M 24576, res 12840, rev 6362: ok}\\*
0 & -- & 0 & 5 & 1 & 214 & 2 & 808 & \texttt{-?????????}\\*
1 & 0 & 1 & 3 & 1 & 688 & 16 & 320 & \texttt{--+?????--}\\
2 & 1 & 2 & 3 & 1 & 352 & 32 & 640 & \texttt{---?????--}\\
3 & 2 & 3 & 2 & 1 & 384 & 128 & 512 & \texttt{----???---}\\
4 & 3 & 4 & 1 & 1 & 512 & 512 & 0 & \texttt{-----+----}\\
5 & 4 & 5 & 1 & 1 & 0 & 1024 & 0 & \texttt{----------}\\
6 & 3 & 4 & 1 & 1 & 0 & 1024 & 0 & \texttt{----------}\\
7 & 2 & 3 & 1 & 1 & 0 & 1024 & 0 & \texttt{----------}\\
8 & 0 & 1 & 1 & 2 & 0 & 1024 & 0 & \texttt{----------}\\
9 & 0 & 1 & 1 & 2 & 0 & 1024 & 0 & \texttt{----------}\\
\multicolumn{9}{@{}l@{}}{\rule{0pt}{2.6ex}\textbf{T10.240}\enspace \texttt{0123454234}\enspace $n{=}10$, $|H|{=}3$, len 6, br 2, $|\mathrm{Aut}|{=}1$, orb 10/10, hist-orb 3/3; M 28672, res 5088, rev 11774: ok}\\*
0 & -- & 0 & 3 & 1 & 890 & 2 & 132 & \texttt{-+????????}\\*
1 & 0 & 1 & 3 & 1 & 756 & 4 & 264 & \texttt{--????????}\\
2 & 1 & 2 & 2 & 1 & 704 & 64 & 256 & \texttt{---+???---}\\
3 & 2 & 3 & 2 & 1 & 384 & 128 & 512 & \texttt{----???---}\\
4 & 3 & 4 & 1 & 1 & 512 & 512 & 0 & \texttt{-----+----}\\
5 & 4 & 5 & 1 & 1 & 0 & 1024 & 0 & \texttt{----------}\\
6 & 3 & 4 & 1 & 1 & 0 & 1024 & 0 & \texttt{----------}\\
7 & 1 & 2 & 1 & 1 & 768 & 256 & 0 & \texttt{--------++}\\
8 & 7 & 3 & 1 & 1 & 512 & 512 & 0 & \texttt{---------+}\\
9 & 8 & 4 & 1 & 1 & 0 & 1024 & 0 & \texttt{----------}\\
\multicolumn{9}{@{}l@{}}{\rule{0pt}{2.6ex}\textbf{T10.241}\enspace \texttt{0123454233}\enspace $n{=}10$, $|H|{=}4$, len 6, br 3, $|\mathrm{Aut}|{=}2$, orb 9/8, hist-orb 3/2; M 27648, res 8320, rev 10154: ok}\\*
0 & -- & 0 & 4 & 1 & 782 & 2 & 240 & \texttt{-+????????}\\*
1 & 0 & 1 & 4 & 1 & 540 & 4 & 480 & \texttt{--????????}\\
2 & 1 & 2 & 2 & 1 & 704 & 64 & 256 & \texttt{---+???---}\\
3 & 2 & 3 & 2 & 1 & 384 & 128 & 512 & \texttt{----???---}\\
4 & 3 & 4 & 1 & 1 & 512 & 512 & 0 & \texttt{-----+----}\\
5 & 4 & 5 & 1 & 1 & 0 & 1024 & 0 & \texttt{----------}\\
6 & 3 & 4 & 1 & 1 & 0 & 1024 & 0 & \texttt{----------}\\
7 & 1 & 2 & 2 & 1 & 256 & 256 & 512 & \texttt{--------??}\\
8 & 7 & 3 & 1 & 2 & 0 & 1024 & 0 & \texttt{----------}\\
9 & 7 & 3 & 1 & 2 & 0 & 1024 & 0 & \texttt{----------}\\
\multicolumn{9}{@{}l@{}}{\rule{0pt}{2.6ex}\textbf{T10.242}\enspace \texttt{0123454232}\enspace $n{=}10$, $|H|{=}4$, len 6, br 2, $|\mathrm{Aut}|{=}1$, orb 10/10, hist-orb 4/4; M 26624, res 10272, rev 8782: ok}\\*
0 & -- & 0 & 4 & 1 & 674 & 2 & 348 & \texttt{-+????????}\\*
1 & 0 & 1 & 4 & 1 & 324 & 4 & 696 & \texttt{--????????}\\
2 & 1 & 2 & 2 & 1 & 704 & 64 & 256 & \texttt{---+???---}\\
3 & 2 & 3 & 2 & 1 & 384 & 128 & 512 & \texttt{----???---}\\
4 & 3 & 4 & 1 & 1 & 512 & 512 & 0 & \texttt{-----+----}\\
5 & 4 & 5 & 1 & 1 & 0 & 1024 & 0 & \texttt{----------}\\
6 & 3 & 4 & 1 & 1 & 0 & 1024 & 0 & \texttt{----------}\\
7 & 1 & 2 & 1 & 1 & 512 & 512 & 0 & \texttt{--------+-}\\
8 & 7 & 3 & 1 & 1 & 0 & 1024 & 0 & \texttt{----------}\\
9 & 1 & 2 & 1 & 1 & 0 & 1024 & 0 & \texttt{----------}\\
\multicolumn{9}{@{}l@{}}{\rule{0pt}{2.6ex}\textbf{T10.243}\enspace \texttt{0123454231}\enspace $n{=}10$, $|H|{=}4$, len 6, br 3, $|\mathrm{Aut}|{=}1$, orb 10/10, hist-orb 4/4; M 25600, res 9748, rev 8414: ok}\\*
0 & -- & 0 & 4 & 1 & 418 & 2 & 604 & \texttt{-?????????}\\*
1 & 0 & 1 & 3 & 1 & 648 & 8 & 368 & \texttt{--???????-}\\
2 & 1 & 2 & 2 & 1 & 704 & 64 & 256 & \texttt{---+???---}\\
3 & 2 & 3 & 2 & 1 & 384 & 128 & 512 & \texttt{----???---}\\
4 & 3 & 4 & 1 & 1 & 512 & 512 & 0 & \texttt{-----+----}\\
5 & 4 & 5 & 1 & 1 & 0 & 1024 & 0 & \texttt{----------}\\
6 & 3 & 4 & 1 & 1 & 0 & 1024 & 0 & \texttt{----------}\\
7 & 1 & 2 & 1 & 1 & 512 & 512 & 0 & \texttt{--------+-}\\
8 & 7 & 3 & 1 & 1 & 0 & 1024 & 0 & \texttt{----------}\\
9 & 0 & 1 & 1 & 1 & 0 & 1024 & 0 & \texttt{----------}\\
\multicolumn{9}{@{}l@{}}{\rule{0pt}{2.6ex}\textbf{T10.244}\enspace \texttt{0123454222}\enspace $n{=}10$, $|H|{=}5$, len 6, br 2, $|\mathrm{Aut}|{=}6$, orb 8/7, hist-orb 3/2; M 25600, res 12864, rev 7026: ok}\\*
0 & -- & 0 & 5 & 1 & 566 & 2 & 456 & \texttt{-+????????}\\*
1 & 0 & 1 & 5 & 1 & 108 & 4 & 912 & \texttt{--????????}\\
2 & 1 & 2 & 2 & 1 & 704 & 64 & 256 & \texttt{---+???---}\\
3 & 2 & 3 & 2 & 1 & 384 & 128 & 512 & \texttt{----???---}\\
4 & 3 & 4 & 1 & 1 & 512 & 512 & 0 & \texttt{-----+----}\\
5 & 4 & 5 & 1 & 1 & 0 & 1024 & 0 & \texttt{----------}\\
6 & 3 & 4 & 1 & 1 & 0 & 1024 & 0 & \texttt{----------}\\
7 & 1 & 2 & 1 & 3 & 0 & 1024 & 0 & \texttt{----------}\\
8 & 1 & 2 & 1 & 3 & 0 & 1024 & 0 & \texttt{----------}\\
9 & 1 & 2 & 1 & 3 & 0 & 1024 & 0 & \texttt{----------}\\
\multicolumn{9}{@{}l@{}}{\rule{0pt}{2.6ex}\textbf{T10.245}\enspace \texttt{0123454221}\enspace $n{=}10$, $|H|{=}5$, len 6, br 3, $|\mathrm{Aut}|{=}2$, orb 9/8, hist-orb 4/3; M 24576, res 13528, rev 5962: ok}\\*
0 & -- & 0 & 5 & 1 & 310 & 2 & 712 & \texttt{-?????????}\\*
1 & 0 & 1 & 4 & 1 & 216 & 8 & 800 & \texttt{--???????-}\\
2 & 1 & 2 & 2 & 1 & 704 & 64 & 256 & \texttt{---+???---}\\
3 & 2 & 3 & 2 & 1 & 384 & 128 & 512 & \texttt{----???---}\\
4 & 3 & 4 & 1 & 1 & 512 & 512 & 0 & \texttt{-----+----}\\
5 & 4 & 5 & 1 & 1 & 0 & 1024 & 0 & \texttt{----------}\\
6 & 3 & 4 & 1 & 1 & 0 & 1024 & 0 & \texttt{----------}\\
7 & 1 & 2 & 1 & 2 & 0 & 1024 & 0 & \texttt{----------}\\
8 & 1 & 2 & 1 & 2 & 0 & 1024 & 0 & \texttt{----------}\\
9 & 0 & 1 & 1 & 1 & 0 & 1024 & 0 & \texttt{----------}\\
\multicolumn{9}{@{}l@{}}{\rule{0pt}{2.6ex}\textbf{T10.246}\enspace \texttt{0123454212}\enspace $n{=}10$, $|H|{=}4$, len 6, br 3, $|\mathrm{Aut}|{=}1$, orb 10/10, hist-orb 4/4; M 24576, res 9372, rev 8018: ok}\\*
0 & -- & 0 & 4 & 1 & 546 & 2 & 476 & \texttt{-?????????}\\*
1 & 0 & 1 & 3 & 1 & 432 & 16 & 576 & \texttt{--??????--}\\
2 & 1 & 2 & 2 & 1 & 704 & 64 & 256 & \texttt{---+???---}\\
3 & 2 & 3 & 2 & 1 & 384 & 128 & 512 & \texttt{----???---}\\
4 & 3 & 4 & 1 & 1 & 512 & 512 & 0 & \texttt{-----+----}\\
5 & 4 & 5 & 1 & 1 & 0 & 1024 & 0 & \texttt{----------}\\
6 & 3 & 4 & 1 & 1 & 0 & 1024 & 0 & \texttt{----------}\\
7 & 1 & 2 & 1 & 1 & 0 & 1024 & 0 & \texttt{----------}\\
8 & 0 & 1 & 1 & 1 & 512 & 512 & 0 & \texttt{---------+}\\
9 & 8 & 2 & 1 & 1 & 0 & 1024 & 0 & \texttt{----------}\\
\multicolumn{9}{@{}l@{}}{\rule{0pt}{2.6ex}\textbf{T10.247}\enspace \texttt{0123454211}\enspace $n{=}10$, $|H|{=}5$, len 6, br 3, $|\mathrm{Aut}|{=}2$, orb 9/8, hist-orb 4/3; M 23552, res 12648, rev 5802: ok}\\*
0 & -- & 0 & 5 & 1 & 182 & 2 & 840 & \texttt{-?????????}\\*
1 & 0 & 1 & 3 & 1 & 432 & 16 & 576 & \texttt{--??????--}\\
2 & 1 & 2 & 2 & 1 & 704 & 64 & 256 & \texttt{---+???---}\\
3 & 2 & 3 & 2 & 1 & 384 & 128 & 512 & \texttt{----???---}\\
4 & 3 & 4 & 1 & 1 & 512 & 512 & 0 & \texttt{-----+----}\\
5 & 4 & 5 & 1 & 1 & 0 & 1024 & 0 & \texttt{----------}\\
6 & 3 & 4 & 1 & 1 & 0 & 1024 & 0 & \texttt{----------}\\
7 & 1 & 2 & 1 & 1 & 0 & 1024 & 0 & \texttt{----------}\\
8 & 0 & 1 & 1 & 2 & 0 & 1024 & 0 & \texttt{----------}\\
9 & 0 & 1 & 1 & 2 & 0 & 1024 & 0 & \texttt{----------}\\
\multicolumn{9}{@{}l@{}}{\rule{0pt}{2.6ex}\textbf{T10.248}\enspace \texttt{0123454123}\enspace $n{=}10$, $|H|{=}3$, len 6, br 2, $|\mathrm{Aut}|{=}1$, orb 10/10, hist-orb 3/3; M 25600, res 4260, rev 10518: ok}\\*
0 & -- & 0 & 3 & 1 & 826 & 2 & 196 & \texttt{-?????????}\\*
1 & 0 & 1 & 2 & 1 & 864 & 32 & 128 & \texttt{--++???---}\\
2 & 1 & 2 & 2 & 1 & 704 & 64 & 256 & \texttt{---+???---}\\
3 & 2 & 3 & 2 & 1 & 384 & 128 & 512 & \texttt{----???---}\\
4 & 3 & 4 & 1 & 1 & 512 & 512 & 0 & \texttt{-----+----}\\
5 & 4 & 5 & 1 & 1 & 0 & 1024 & 0 & \texttt{----------}\\
6 & 3 & 4 & 1 & 1 & 0 & 1024 & 0 & \texttt{----------}\\
7 & 0 & 1 & 1 & 1 & 768 & 256 & 0 & \texttt{--------++}\\
8 & 7 & 2 & 1 & 1 & 512 & 512 & 0 & \texttt{---------+}\\
9 & 8 & 3 & 1 & 1 & 0 & 1024 & 0 & \texttt{----------}\\
\multicolumn{9}{@{}l@{}}{\rule{0pt}{2.6ex}\textbf{T10.249}\enspace \texttt{0123454122}\enspace $n{=}10$, $|H|{=}4$, len 6, br 3, $|\mathrm{Aut}|{=}2$, orb 9/8, hist-orb 3/2; M 24576, res 7152, rev 8862: ok}\\*
0 & -- & 0 & 4 & 1 & 590 & 2 & 432 & \texttt{-?????????}\\*
1 & 0 & 1 & 2 & 1 & 864 & 32 & 128 & \texttt{--++???---}\\
2 & 1 & 2 & 2 & 1 & 704 & 64 & 256 & \texttt{---+???---}\\
3 & 2 & 3 & 2 & 1 & 384 & 128 & 512 & \texttt{----???---}\\
4 & 3 & 4 & 1 & 1 & 512 & 512 & 0 & \texttt{-----+----}\\
5 & 4 & 5 & 1 & 1 & 0 & 1024 & 0 & \texttt{----------}\\
6 & 3 & 4 & 1 & 1 & 0 & 1024 & 0 & \texttt{----------}\\
7 & 0 & 1 & 2 & 1 & 256 & 256 & 512 & \texttt{--------??}\\
8 & 7 & 2 & 1 & 2 & 0 & 1024 & 0 & \texttt{----------}\\
9 & 7 & 2 & 1 & 2 & 0 & 1024 & 0 & \texttt{----------}\\
\multicolumn{9}{@{}l@{}}{\rule{0pt}{2.6ex}\textbf{T10.250}\enspace \texttt{0123454121}\enspace $n{=}10$, $|H|{=}4$, len 6, br 2, $|\mathrm{Aut}|{=}1$, orb 10/10, hist-orb 4/4; M 23552, res 8508, rev 7758: ok}\\*
0 & -- & 0 & 4 & 1 & 354 & 2 & 668 & \texttt{-?????????}\\*
1 & 0 & 1 & 2 & 1 & 864 & 32 & 128 & \texttt{--++???---}\\
2 & 1 & 2 & 2 & 1 & 704 & 64 & 256 & \texttt{---+???---}\\
3 & 2 & 3 & 2 & 1 & 384 & 128 & 512 & \texttt{----???---}\\
4 & 3 & 4 & 1 & 1 & 512 & 512 & 0 & \texttt{-----+----}\\
5 & 4 & 5 & 1 & 1 & 0 & 1024 & 0 & \texttt{----------}\\
6 & 3 & 4 & 1 & 1 & 0 & 1024 & 0 & \texttt{----------}\\
7 & 0 & 1 & 1 & 1 & 512 & 512 & 0 & \texttt{--------+-}\\
8 & 7 & 2 & 1 & 1 & 0 & 1024 & 0 & \texttt{----------}\\
9 & 0 & 1 & 1 & 1 & 0 & 1024 & 0 & \texttt{----------}\\
\multicolumn{9}{@{}l@{}}{\rule{0pt}{2.6ex}\textbf{T10.251}\enspace \texttt{0123454111}\enspace $n{=}10$, $|H|{=}5$, len 6, br 2, $|\mathrm{Aut}|{=}6$, orb 8/7, hist-orb 3/2; M 22528, res 10632, rev 6142: ok}\\*
0 & -- & 0 & 5 & 1 & 118 & 2 & 904 & \texttt{-?????????}\\*
1 & 0 & 1 & 2 & 1 & 864 & 32 & 128 & \texttt{--++???---}\\
2 & 1 & 2 & 2 & 1 & 704 & 64 & 256 & \texttt{---+???---}\\
3 & 2 & 3 & 2 & 1 & 384 & 128 & 512 & \texttt{----???---}\\
4 & 3 & 4 & 1 & 1 & 512 & 512 & 0 & \texttt{-----+----}\\
5 & 4 & 5 & 1 & 1 & 0 & 1024 & 0 & \texttt{----------}\\
6 & 3 & 4 & 1 & 1 & 0 & 1024 & 0 & \texttt{----------}\\
7 & 0 & 1 & 1 & 3 & 0 & 1024 & 0 & \texttt{----------}\\
8 & 0 & 1 & 1 & 3 & 0 & 1024 & 0 & \texttt{----------}\\
9 & 0 & 1 & 1 & 3 & 0 & 1024 & 0 & \texttt{----------}\\
\multicolumn{9}{@{}l@{}}{\rule{0pt}{2.6ex}\textbf{T10.252}\enspace \texttt{0123453453}\enspace $n{=}10$, $|H|{=}3$, len 6, br 1, $|\mathrm{Aut}|{=}2$, orb 7/4, hist-orb 2/1; M 30720, res 7644, rev 11686: ok}\\*
0 & -- & 0 & 3 & 1 & 866 & 2 & 156 & \texttt{-++???????}\\*
1 & 0 & 1 & 3 & 1 & 708 & 4 & 312 & \texttt{--+???????}\\
2 & 1 & 2 & 3 & 1 & 392 & 8 & 624 & \texttt{---???????}\\
3 & 2 & 3 & 1 & 2 & 768 & 256 & 0 & \texttt{----++----}\\
4 & 3 & 4 & 1 & 2 & 512 & 512 & 0 & \texttt{-----+----}\\
5 & 4 & 5 & 1 & 2 & 0 & 1024 & 0 & \texttt{----------}\\
6 & 2 & 3 & 1 & 2 & 768 & 256 & 0 & \texttt{-------++-}\\
7 & 6 & 4 & 1 & 2 & 512 & 512 & 0 & \texttt{--------+-}\\
8 & 7 & 5 & 1 & 2 & 0 & 1024 & 0 & \texttt{----------}\\
9 & 2 & 3 & 1 & 1 & 0 & 1024 & 0 & \texttt{----------}\\
\multicolumn{9}{@{}l@{}}{\rule{0pt}{2.6ex}\textbf{T10.253}\enspace \texttt{0123453452}\enspace $n{=}10$, $|H|{=}3$, len 6, br 2, $|\mathrm{Aut}|{=}2$, orb 7/4, hist-orb 2/1; M 29696, res 7992, rev 10802: ok}\\*
0 & -- & 0 & 3 & 1 & 738 & 2 & 284 & \texttt{-+????????}\\*
1 & 0 & 1 & 3 & 1 & 452 & 4 & 568 & \texttt{--????????}\\
2 & 1 & 2 & 2 & 1 & 784 & 16 & 224 & \texttt{---??????-}\\
3 & 2 & 3 & 1 & 2 & 768 & 256 & 0 & \texttt{----++----}\\
4 & 3 & 4 & 1 & 2 & 512 & 512 & 0 & \texttt{-----+----}\\
5 & 4 & 5 & 1 & 2 & 0 & 1024 & 0 & \texttt{----------}\\
6 & 2 & 3 & 1 & 2 & 768 & 256 & 0 & \texttt{-------++-}\\
7 & 6 & 4 & 1 & 2 & 512 & 512 & 0 & \texttt{--------+-}\\
8 & 7 & 5 & 1 & 2 & 0 & 1024 & 0 & \texttt{----------}\\
9 & 1 & 2 & 1 & 1 & 0 & 1024 & 0 & \texttt{----------}\\
\multicolumn{9}{@{}l@{}}{\rule{0pt}{2.6ex}\textbf{T10.254}\enspace \texttt{0123453451}\enspace $n{=}10$, $|H|{=}3$, len 6, br 2, $|\mathrm{Aut}|{=}2$, orb 7/4, hist-orb 2/1; M 28672, res 6764, rev 10710: ok}\\*
0 & -- & 0 & 3 & 1 & 482 & 2 & 540 & \texttt{-?????????}\\*
1 & 0 & 1 & 2 & 1 & 904 & 8 & 112 & \texttt{--+??????-}\\
2 & 1 & 2 & 2 & 1 & 784 & 16 & 224 & \texttt{---??????-}\\
3 & 2 & 3 & 1 & 2 & 768 & 256 & 0 & \texttt{----++----}\\
4 & 3 & 4 & 1 & 2 & 512 & 512 & 0 & \texttt{-----+----}\\
5 & 4 & 5 & 1 & 2 & 0 & 1024 & 0 & \texttt{----------}\\
6 & 2 & 3 & 1 & 2 & 768 & 256 & 0 & \texttt{-------++-}\\
7 & 6 & 4 & 1 & 2 & 512 & 512 & 0 & \texttt{--------+-}\\
8 & 7 & 5 & 1 & 2 & 0 & 1024 & 0 & \texttt{----------}\\
9 & 0 & 1 & 1 & 1 & 0 & 1024 & 0 & \texttt{----------}\\
\multicolumn{9}{@{}l@{}}{\rule{0pt}{2.6ex}\textbf{T10.255}\enspace \texttt{0123453444}\enspace $n{=}10$, $|H|{=}4$, len 6, br 2, $|\mathrm{Aut}|{=}6$, orb 8/7, hist-orb 2/1; M 30720, res 7904, rev 12302: ok}\\*
0 & -- & 0 & 4 & 1 & 894 & 2 & 128 & \texttt{-++???????}\\*
1 & 0 & 1 & 4 & 1 & 764 & 4 & 256 & \texttt{--+???????}\\
2 & 1 & 2 & 4 & 1 & 504 & 8 & 512 & \texttt{---???????}\\
3 & 2 & 3 & 1 & 1 & 768 & 256 & 0 & \texttt{----++----}\\
4 & 3 & 4 & 1 & 1 & 512 & 512 & 0 & \texttt{-----+----}\\
5 & 4 & 5 & 1 & 1 & 0 & 1024 & 0 & \texttt{----------}\\
6 & 2 & 3 & 3 & 1 & 128 & 128 & 768 & \texttt{-------???}\\
7 & 6 & 4 & 1 & 3 & 0 & 1024 & 0 & \texttt{----------}\\
8 & 6 & 4 & 1 & 3 & 0 & 1024 & 0 & \texttt{----------}\\
9 & 6 & 4 & 1 & 3 & 0 & 1024 & 0 & \texttt{----------}\\
\multicolumn{9}{@{}l@{}}{\rule{0pt}{2.6ex}\textbf{T10.256}\enspace \texttt{0123453443}\enspace $n{=}10$, $|H|{=}4$, len 6, br 2, $|\mathrm{Aut}|{=}2$, orb 9/8, hist-orb 3/2; M 29696, res 9592, rev 10798: ok}\\*
0 & -- & 0 & 4 & 1 & 838 & 2 & 184 & \texttt{-++???????}\\*
1 & 0 & 1 & 4 & 1 & 652 & 4 & 368 & \texttt{--+???????}\\
2 & 1 & 2 & 4 & 1 & 280 & 8 & 736 & \texttt{---???????}\\
3 & 2 & 3 & 1 & 1 & 768 & 256 & 0 & \texttt{----++----}\\
4 & 3 & 4 & 1 & 1 & 512 & 512 & 0 & \texttt{-----+----}\\
5 & 4 & 5 & 1 & 1 & 0 & 1024 & 0 & \texttt{----------}\\
6 & 2 & 3 & 2 & 1 & 256 & 256 & 512 & \texttt{-------??-}\\
7 & 6 & 4 & 1 & 2 & 0 & 1024 & 0 & \texttt{----------}\\
8 & 6 & 4 & 1 & 2 & 0 & 1024 & 0 & \texttt{----------}\\
9 & 2 & 3 & 1 & 1 & 0 & 1024 & 0 & \texttt{----------}\\
\multicolumn{9}{@{}l@{}}{\rule{0pt}{2.6ex}\textbf{T10.257}\enspace \texttt{0123453442}\enspace $n{=}10$, $|H|{=}4$, len 6, br 3, $|\mathrm{Aut}|{=}2$, orb 9/8, hist-orb 3/2; M 28672, res 10416, rev 9658: ok}\\*
0 & -- & 0 & 4 & 1 & 710 & 2 & 312 & \texttt{-+????????}\\*
1 & 0 & 1 & 4 & 1 & 396 & 4 & 624 & \texttt{--????????}\\
2 & 1 & 2 & 3 & 1 & 560 & 16 & 448 & \texttt{---??????-}\\
3 & 2 & 3 & 1 & 1 & 768 & 256 & 0 & \texttt{----++----}\\
4 & 3 & 4 & 1 & 1 & 512 & 512 & 0 & \texttt{-----+----}\\
5 & 4 & 5 & 1 & 1 & 0 & 1024 & 0 & \texttt{----------}\\
6 & 2 & 3 & 2 & 1 & 256 & 256 & 512 & \texttt{-------??-}\\
7 & 6 & 4 & 1 & 2 & 0 & 1024 & 0 & \texttt{----------}\\
8 & 6 & 4 & 1 & 2 & 0 & 1024 & 0 & \texttt{----------}\\
9 & 1 & 2 & 1 & 1 & 0 & 1024 & 0 & \texttt{----------}\\
\multicolumn{9}{@{}l@{}}{\rule{0pt}{2.6ex}\textbf{T10.258}\enspace \texttt{0123453441}\enspace $n{=}10$, $|H|{=}4$, len 6, br 3, $|\mathrm{Aut}|{=}2$, orb 9/8, hist-orb 3/2; M 27648, res 9496, rev 9534: ok}\\*
0 & -- & 0 & 4 & 1 & 454 & 2 & 568 & \texttt{-?????????}\\*
1 & 0 & 1 & 3 & 1 & 792 & 8 & 224 & \texttt{--+??????-}\\
2 & 1 & 2 & 3 & 1 & 560 & 16 & 448 & \texttt{---??????-}\\
3 & 2 & 3 & 1 & 1 & 768 & 256 & 0 & \texttt{----++----}\\
4 & 3 & 4 & 1 & 1 & 512 & 512 & 0 & \texttt{-----+----}\\
5 & 4 & 5 & 1 & 1 & 0 & 1024 & 0 & \texttt{----------}\\
6 & 2 & 3 & 2 & 1 & 256 & 256 & 512 & \texttt{-------??-}\\
7 & 6 & 4 & 1 & 2 & 0 & 1024 & 0 & \texttt{----------}\\
8 & 6 & 4 & 1 & 2 & 0 & 1024 & 0 & \texttt{----------}\\
9 & 0 & 1 & 1 & 1 & 0 & 1024 & 0 & \texttt{----------}\\
\multicolumn{9}{@{}l@{}}{\rule{0pt}{2.6ex}\textbf{T10.259}\enspace \texttt{0123453434}\enspace $n{=}10$, $|H|{=}3$, len 6, br 1, $|\mathrm{Aut}|{=}2$, orb 8/6, hist-orb 2/1; M 29696, res 6272, rev 12390: ok}\\*
0 & -- & 0 & 3 & 1 & 894 & 2 & 128 & \texttt{-++???????}\\*
1 & 0 & 1 & 3 & 1 & 764 & 4 & 256 & \texttt{--+???????}\\
2 & 1 & 2 & 3 & 1 & 504 & 8 & 512 & \texttt{---???????}\\
3 & 2 & 3 & 1 & 1 & 768 & 256 & 0 & \texttt{----++----}\\
4 & 3 & 4 & 1 & 1 & 512 & 512 & 0 & \texttt{-----+----}\\
5 & 4 & 5 & 1 & 1 & 0 & 1024 & 0 & \texttt{----------}\\
6 & 2 & 3 & 1 & 2 & 512 & 512 & 0 & \texttt{-------+--}\\
7 & 6 & 4 & 1 & 2 & 0 & 1024 & 0 & \texttt{----------}\\
8 & 2 & 3 & 1 & 2 & 512 & 512 & 0 & \texttt{---------+}\\
9 & 8 & 4 & 1 & 2 & 0 & 1024 & 0 & \texttt{----------}\\
\multicolumn{9}{@{}l@{}}{\rule{0pt}{2.6ex}\textbf{T10.260}\enspace \texttt{0123453433}\enspace $n{=}10$, $|H|{=}4$, len 6, br 1, $|\mathrm{Aut}|{=}2$, orb 9/8, hist-orb 3/2; M 28672, res 10388, rev 10222: ok}\\*
0 & -- & 0 & 4 & 1 & 810 & 2 & 212 & \texttt{-++???????}\\*
1 & 0 & 1 & 4 & 1 & 596 & 4 & 424 & \texttt{--+???????}\\
2 & 1 & 2 & 4 & 1 & 168 & 8 & 848 & \texttt{---???????}\\
3 & 2 & 3 & 1 & 1 & 768 & 256 & 0 & \texttt{----++----}\\
4 & 3 & 4 & 1 & 1 & 512 & 512 & 0 & \texttt{-----+----}\\
5 & 4 & 5 & 1 & 1 & 0 & 1024 & 0 & \texttt{----------}\\
6 & 2 & 3 & 1 & 1 & 512 & 512 & 0 & \texttt{-------+--}\\
7 & 6 & 4 & 1 & 1 & 0 & 1024 & 0 & \texttt{----------}\\
8 & 2 & 3 & 1 & 2 & 0 & 1024 & 0 & \texttt{----------}\\
9 & 2 & 3 & 1 & 2 & 0 & 1024 & 0 & \texttt{----------}\\
\multicolumn{9}{@{}l@{}}{\rule{0pt}{2.6ex}\textbf{T10.261}\enspace \texttt{0123453432}\enspace $n{=}10$, $|H|{=}4$, len 6, br 2, $|\mathrm{Aut}|{=}1$, orb 10/10, hist-orb 4/4; M 27648, res 11688, rev 8698: ok}\\*
0 & -- & 0 & 4 & 1 & 682 & 2 & 340 & \texttt{-+????????}\\*
1 & 0 & 1 & 4 & 1 & 340 & 4 & 680 & \texttt{--????????}\\
2 & 1 & 2 & 3 & 1 & 336 & 16 & 672 & \texttt{---??????-}\\
3 & 2 & 3 & 1 & 1 & 768 & 256 & 0 & \texttt{----++----}\\
4 & 3 & 4 & 1 & 1 & 512 & 512 & 0 & \texttt{-----+----}\\
5 & 4 & 5 & 1 & 1 & 0 & 1024 & 0 & \texttt{----------}\\
6 & 2 & 3 & 1 & 1 & 512 & 512 & 0 & \texttt{-------+--}\\
7 & 6 & 4 & 1 & 1 & 0 & 1024 & 0 & \texttt{----------}\\
8 & 2 & 3 & 1 & 1 & 0 & 1024 & 0 & \texttt{----------}\\
9 & 1 & 2 & 1 & 1 & 0 & 1024 & 0 & \texttt{----------}\\
\multicolumn{9}{@{}l@{}}{\rule{0pt}{2.6ex}\textbf{T10.262}\enspace \texttt{0123453431}\enspace $n{=}10$, $|H|{=}4$, len 6, br 2, $|\mathrm{Aut}|{=}1$, orb 10/10, hist-orb 4/4; M 26624, res 11076, rev 8430: ok}\\*
0 & -- & 0 & 4 & 1 & 426 & 2 & 596 & \texttt{-?????????}\\*
1 & 0 & 1 & 3 & 1 & 680 & 8 & 336 & \texttt{--+??????-}\\
2 & 1 & 2 & 3 & 1 & 336 & 16 & 672 & \texttt{---??????-}\\
3 & 2 & 3 & 1 & 1 & 768 & 256 & 0 & \texttt{----++----}\\
4 & 3 & 4 & 1 & 1 & 512 & 512 & 0 & \texttt{-----+----}\\
5 & 4 & 5 & 1 & 1 & 0 & 1024 & 0 & \texttt{----------}\\
6 & 2 & 3 & 1 & 1 & 512 & 512 & 0 & \texttt{-------+--}\\
7 & 6 & 4 & 1 & 1 & 0 & 1024 & 0 & \texttt{----------}\\
8 & 2 & 3 & 1 & 1 & 0 & 1024 & 0 & \texttt{----------}\\
9 & 0 & 1 & 1 & 1 & 0 & 1024 & 0 & \texttt{----------}\\
\multicolumn{9}{@{}l@{}}{\rule{0pt}{2.6ex}\textbf{T10.263}\enspace \texttt{0123453423}\enspace $n{=}10$, $|H|{=}3$, len 6, br 2, $|\mathrm{Aut}|{=}1$, orb 10/10, hist-orb 3/3; M 27648, res 5968, rev 11222: ok}\\*
0 & -- & 0 & 3 & 1 & 830 & 2 & 192 & \texttt{-+????????}\\*
1 & 0 & 1 & 3 & 1 & 636 & 4 & 384 & \texttt{--????????}\\
2 & 1 & 2 & 2 & 1 & 672 & 32 & 320 & \texttt{---?????--}\\
3 & 2 & 3 & 1 & 1 & 768 & 256 & 0 & \texttt{----++----}\\
4 & 3 & 4 & 1 & 1 & 512 & 512 & 0 & \texttt{-----+----}\\
5 & 4 & 5 & 1 & 1 & 0 & 1024 & 0 & \texttt{----------}\\
6 & 2 & 3 & 1 & 1 & 512 & 512 & 0 & \texttt{-------+--}\\
7 & 6 & 4 & 1 & 1 & 0 & 1024 & 0 & \texttt{----------}\\
8 & 1 & 2 & 1 & 1 & 512 & 512 & 0 & \texttt{---------+}\\
9 & 8 & 3 & 1 & 1 & 0 & 1024 & 0 & \texttt{----------}\\
\multicolumn{9}{@{}l@{}}{\rule{0pt}{2.6ex}\textbf{T10.264}\enspace \texttt{0123453422}\enspace $n{=}10$, $|H|{=}4$, len 6, br 2, $|\mathrm{Aut}|{=}2$, orb 9/8, hist-orb 3/2; M 26624, res 11056, rev 8338: ok}\\*
0 & -- & 0 & 4 & 1 & 618 & 2 & 404 & \texttt{-+????????}\\*
1 & 0 & 1 & 4 & 1 & 212 & 4 & 808 & \texttt{--????????}\\
2 & 1 & 2 & 2 & 1 & 672 & 32 & 320 & \texttt{---?????--}\\
3 & 2 & 3 & 1 & 1 & 768 & 256 & 0 & \texttt{----++----}\\
4 & 3 & 4 & 1 & 1 & 512 & 512 & 0 & \texttt{-----+----}\\
5 & 4 & 5 & 1 & 1 & 0 & 1024 & 0 & \texttt{----------}\\
6 & 2 & 3 & 1 & 1 & 512 & 512 & 0 & \texttt{-------+--}\\
7 & 6 & 4 & 1 & 1 & 0 & 1024 & 0 & \texttt{----------}\\
8 & 1 & 2 & 1 & 2 & 0 & 1024 & 0 & \texttt{----------}\\
9 & 1 & 2 & 1 & 2 & 0 & 1024 & 0 & \texttt{----------}\\
\multicolumn{9}{@{}l@{}}{\rule{0pt}{2.6ex}\textbf{T10.265}\enspace \texttt{0123453421}\enspace $n{=}10$, $|H|{=}4$, len 6, br 3, $|\mathrm{Aut}|{=}1$, orb 10/10, hist-orb 4/4; M 25600, res 11148, rev 7650: ok}\\*
0 & -- & 0 & 4 & 1 & 362 & 2 & 660 & \texttt{-?????????}\\*
1 & 0 & 1 & 3 & 1 & 424 & 8 & 592 & \texttt{--???????-}\\
2 & 1 & 2 & 2 & 1 & 672 & 32 & 320 & \texttt{---?????--}\\
3 & 2 & 3 & 1 & 1 & 768 & 256 & 0 & \texttt{----++----}\\
4 & 3 & 4 & 1 & 1 & 512 & 512 & 0 & \texttt{-----+----}\\
5 & 4 & 5 & 1 & 1 & 0 & 1024 & 0 & \texttt{----------}\\
6 & 2 & 3 & 1 & 1 & 512 & 512 & 0 & \texttt{-------+--}\\
7 & 6 & 4 & 1 & 1 & 0 & 1024 & 0 & \texttt{----------}\\
8 & 1 & 2 & 1 & 1 & 0 & 1024 & 0 & \texttt{----------}\\
9 & 0 & 1 & 1 & 1 & 0 & 1024 & 0 & \texttt{----------}\\
\multicolumn{9}{@{}l@{}}{\rule{0pt}{2.6ex}\textbf{T10.266}\enspace \texttt{0123453412}\enspace $n{=}10$, $|H|{=}3$, len 6, br 2, $|\mathrm{Aut}|{=}1$, orb 10/10, hist-orb 3/3; M 25600, res 5120, rev 10366: ok}\\*
0 & -- & 0 & 3 & 1 & 702 & 2 & 320 & \texttt{-?????????}\\*
1 & 0 & 1 & 2 & 1 & 848 & 16 & 160 & \texttt{--+?????--}\\
2 & 1 & 2 & 2 & 1 & 672 & 32 & 320 & \texttt{---?????--}\\
3 & 2 & 3 & 1 & 1 & 768 & 256 & 0 & \texttt{----++----}\\
4 & 3 & 4 & 1 & 1 & 512 & 512 & 0 & \texttt{-----+----}\\
5 & 4 & 5 & 1 & 1 & 0 & 1024 & 0 & \texttt{----------}\\
6 & 2 & 3 & 1 & 1 & 512 & 512 & 0 & \texttt{-------+--}\\
7 & 6 & 4 & 1 & 1 & 0 & 1024 & 0 & \texttt{----------}\\
8 & 0 & 1 & 1 & 1 & 512 & 512 & 0 & \texttt{---------+}\\
9 & 8 & 2 & 1 & 1 & 0 & 1024 & 0 & \texttt{----------}\\
\multicolumn{9}{@{}l@{}}{\rule{0pt}{2.6ex}\textbf{T10.267}\enspace \texttt{0123453411}\enspace $n{=}10$, $|H|{=}4$, len 6, br 2, $|\mathrm{Aut}|{=}2$, orb 9/8, hist-orb 3/2; M 24576, res 9332, rev 7878: ok}\\*
0 & -- & 0 & 4 & 1 & 234 & 2 & 788 & \texttt{-?????????}\\*
1 & 0 & 1 & 2 & 1 & 848 & 16 & 160 & \texttt{--+?????--}\\
2 & 1 & 2 & 2 & 1 & 672 & 32 & 320 & \texttt{---?????--}\\
3 & 2 & 3 & 1 & 1 & 768 & 256 & 0 & \texttt{----++----}\\
4 & 3 & 4 & 1 & 1 & 512 & 512 & 0 & \texttt{-----+----}\\
5 & 4 & 5 & 1 & 1 & 0 & 1024 & 0 & \texttt{----------}\\
6 & 2 & 3 & 1 & 1 & 512 & 512 & 0 & \texttt{-------+--}\\
7 & 6 & 4 & 1 & 1 & 0 & 1024 & 0 & \texttt{----------}\\
8 & 0 & 1 & 1 & 2 & 0 & 1024 & 0 & \texttt{----------}\\
9 & 0 & 1 & 1 & 2 & 0 & 1024 & 0 & \texttt{----------}\\
\multicolumn{9}{@{}l@{}}{\rule{0pt}{2.6ex}\textbf{T10.268}\enspace \texttt{0123453333}\enspace $n{=}10$, $|H|{=}5$, len 6, br 1, $|\mathrm{Aut}|{=}24$, orb 7/6, hist-orb 2/1; M 27648, res 11760, rev 9454: ok}\\*
0 & -- & 0 & 5 & 1 & 782 & 2 & 240 & \texttt{-++???????}\\*
1 & 0 & 1 & 5 & 1 & 540 & 4 & 480 & \texttt{--+???????}\\
2 & 1 & 2 & 5 & 1 & 56 & 8 & 960 & \texttt{---???????}\\
3 & 2 & 3 & 1 & 1 & 768 & 256 & 0 & \texttt{----++----}\\
4 & 3 & 4 & 1 & 1 & 512 & 512 & 0 & \texttt{-----+----}\\
5 & 4 & 5 & 1 & 1 & 0 & 1024 & 0 & \texttt{----------}\\
6 & 2 & 3 & 1 & 4 & 0 & 1024 & 0 & \texttt{----------}\\
7 & 2 & 3 & 1 & 4 & 0 & 1024 & 0 & \texttt{----------}\\
8 & 2 & 3 & 1 & 4 & 0 & 1024 & 0 & \texttt{----------}\\
9 & 2 & 3 & 1 & 4 & 0 & 1024 & 0 & \texttt{----------}\\
\multicolumn{9}{@{}l@{}}{\rule{0pt}{2.6ex}\textbf{T10.269}\enspace \texttt{0123453332}\enspace $n{=}10$, $|H|{=}5$, len 6, br 2, $|\mathrm{Aut}|{=}6$, orb 8/7, hist-orb 3/2; M 26624, res 13536, rev 7450: ok}\\*
0 & -- & 0 & 5 & 1 & 654 & 2 & 368 & \texttt{-+????????}\\*
1 & 0 & 1 & 5 & 1 & 284 & 4 & 736 & \texttt{--????????}\\
2 & 1 & 2 & 4 & 1 & 112 & 16 & 896 & \texttt{---??????-}\\
3 & 2 & 3 & 1 & 1 & 768 & 256 & 0 & \texttt{----++----}\\
4 & 3 & 4 & 1 & 1 & 512 & 512 & 0 & \texttt{-----+----}\\
5 & 4 & 5 & 1 & 1 & 0 & 1024 & 0 & \texttt{----------}\\
6 & 2 & 3 & 1 & 3 & 0 & 1024 & 0 & \texttt{----------}\\
7 & 2 & 3 & 1 & 3 & 0 & 1024 & 0 & \texttt{----------}\\
8 & 2 & 3 & 1 & 3 & 0 & 1024 & 0 & \texttt{----------}\\
9 & 1 & 2 & 1 & 1 & 0 & 1024 & 0 & \texttt{----------}\\
\multicolumn{9}{@{}l@{}}{\rule{0pt}{2.6ex}\textbf{T10.270}\enspace \texttt{0123453331}\enspace $n{=}10$, $|H|{=}5$, len 6, br 2, $|\mathrm{Aut}|{=}6$, orb 8/7, hist-orb 3/2; M 25600, res 13232, rev 7038: ok}\\*
0 & -- & 0 & 5 & 1 & 398 & 2 & 624 & \texttt{-?????????}\\*
1 & 0 & 1 & 4 & 1 & 568 & 8 & 448 & \texttt{--+??????-}\\
2 & 1 & 2 & 4 & 1 & 112 & 16 & 896 & \texttt{---??????-}\\
3 & 2 & 3 & 1 & 1 & 768 & 256 & 0 & \texttt{----++----}\\
4 & 3 & 4 & 1 & 1 & 512 & 512 & 0 & \texttt{-----+----}\\
5 & 4 & 5 & 1 & 1 & 0 & 1024 & 0 & \texttt{----------}\\
6 & 2 & 3 & 1 & 3 & 0 & 1024 & 0 & \texttt{----------}\\
7 & 2 & 3 & 1 & 3 & 0 & 1024 & 0 & \texttt{----------}\\
8 & 2 & 3 & 1 & 3 & 0 & 1024 & 0 & \texttt{----------}\\
9 & 0 & 1 & 1 & 1 & 0 & 1024 & 0 & \texttt{----------}\\
\multicolumn{9}{@{}l@{}}{\rule{0pt}{2.6ex}\textbf{T10.271}\enspace \texttt{0123453323}\enspace $n{=}10$, $|H|{=}4$, len 6, br 2, $|\mathrm{Aut}|{=}2$, orb 9/8, hist-orb 3/2; M 26624, res 9888, rev 9118: ok}\\*
0 & -- & 0 & 4 & 1 & 746 & 2 & 276 & \texttt{-+????????}\\*
1 & 0 & 1 & 4 & 1 & 468 & 4 & 552 & \texttt{--????????}\\
2 & 1 & 2 & 3 & 1 & 224 & 32 & 768 & \texttt{---?????--}\\
3 & 2 & 3 & 1 & 1 & 768 & 256 & 0 & \texttt{----++----}\\
4 & 3 & 4 & 1 & 1 & 512 & 512 & 0 & \texttt{-----+----}\\
5 & 4 & 5 & 1 & 1 & 0 & 1024 & 0 & \texttt{----------}\\
6 & 2 & 3 & 1 & 2 & 0 & 1024 & 0 & \texttt{----------}\\
7 & 2 & 3 & 1 & 2 & 0 & 1024 & 0 & \texttt{----------}\\
8 & 1 & 2 & 1 & 1 & 512 & 512 & 0 & \texttt{---------+}\\
9 & 8 & 3 & 1 & 1 & 0 & 1024 & 0 & \texttt{----------}\\
\multicolumn{9}{@{}l@{}}{\rule{0pt}{2.6ex}\textbf{T10.272}\enspace \texttt{0123453322}\enspace $n{=}10$, $|H|{=}5$, len 6, br 2, $|\mathrm{Aut}|{=}4$, orb 8/6, hist-orb 3/1; M 25600, res 13632, rev 6642: ok}\\*
0 & -- & 0 & 5 & 1 & 590 & 2 & 432 & \texttt{-+????????}\\*
1 & 0 & 1 & 5 & 1 & 156 & 4 & 864 & \texttt{--????????}\\
2 & 1 & 2 & 3 & 1 & 224 & 32 & 768 & \texttt{---?????--}\\
3 & 2 & 3 & 1 & 1 & 768 & 256 & 0 & \texttt{----++----}\\
4 & 3 & 4 & 1 & 1 & 512 & 512 & 0 & \texttt{-----+----}\\
5 & 4 & 5 & 1 & 1 & 0 & 1024 & 0 & \texttt{----------}\\
6 & 2 & 3 & 1 & 2 & 0 & 1024 & 0 & \texttt{----------}\\
7 & 2 & 3 & 1 & 2 & 0 & 1024 & 0 & \texttt{----------}\\
8 & 1 & 2 & 1 & 2 & 0 & 1024 & 0 & \texttt{----------}\\
9 & 1 & 2 & 1 & 2 & 0 & 1024 & 0 & \texttt{----------}\\
\multicolumn{9}{@{}l@{}}{\rule{0pt}{2.6ex}\textbf{T10.273}\enspace \texttt{0123453321}\enspace $n{=}10$, $|H|{=}5$, len 6, br 3, $|\mathrm{Aut}|{=}2$, orb 9/8, hist-orb 4/3; M 24576, res 14032, rev 5746: ok}\\*
0 & -- & 0 & 5 & 1 & 334 & 2 & 688 & \texttt{-?????????}\\*
1 & 0 & 1 & 4 & 1 & 312 & 8 & 704 & \texttt{--???????-}\\
2 & 1 & 2 & 3 & 1 & 224 & 32 & 768 & \texttt{---?????--}\\
3 & 2 & 3 & 1 & 1 & 768 & 256 & 0 & \texttt{----++----}\\
4 & 3 & 4 & 1 & 1 & 512 & 512 & 0 & \texttt{-----+----}\\
5 & 4 & 5 & 1 & 1 & 0 & 1024 & 0 & \texttt{----------}\\
6 & 2 & 3 & 1 & 2 & 0 & 1024 & 0 & \texttt{----------}\\
7 & 2 & 3 & 1 & 2 & 0 & 1024 & 0 & \texttt{----------}\\
8 & 1 & 2 & 1 & 1 & 0 & 1024 & 0 & \texttt{----------}\\
9 & 0 & 1 & 1 & 1 & 0 & 1024 & 0 & \texttt{----------}\\
\multicolumn{9}{@{}l@{}}{\rule{0pt}{2.6ex}\textbf{T10.274}\enspace \texttt{0123453312}\enspace $n{=}10$, $|H|{=}4$, len 6, br 2, $|\mathrm{Aut}|{=}2$, orb 9/8, hist-orb 3/2; M 24576, res 9012, rev 8342: ok}\\*
0 & -- & 0 & 4 & 1 & 618 & 2 & 404 & \texttt{-?????????}\\*
1 & 0 & 1 & 3 & 1 & 624 & 16 & 384 & \texttt{--+?????--}\\
2 & 1 & 2 & 3 & 1 & 224 & 32 & 768 & \texttt{---?????--}\\
3 & 2 & 3 & 1 & 1 & 768 & 256 & 0 & \texttt{----++----}\\
4 & 3 & 4 & 1 & 1 & 512 & 512 & 0 & \texttt{-----+----}\\
5 & 4 & 5 & 1 & 1 & 0 & 1024 & 0 & \texttt{----------}\\
6 & 2 & 3 & 1 & 2 & 0 & 1024 & 0 & \texttt{----------}\\
7 & 2 & 3 & 1 & 2 & 0 & 1024 & 0 & \texttt{----------}\\
8 & 0 & 1 & 1 & 1 & 512 & 512 & 0 & \texttt{---------+}\\
9 & 8 & 2 & 1 & 1 & 0 & 1024 & 0 & \texttt{----------}\\
\multicolumn{9}{@{}l@{}}{\rule{0pt}{2.6ex}\textbf{T10.275}\enspace \texttt{0123453311}\enspace $n{=}10$, $|H|{=}5$, len 6, br 2, $|\mathrm{Aut}|{=}4$, orb 8/6, hist-orb 3/1; M 23552, res 12720, rev 5862: ok}\\*
0 & -- & 0 & 5 & 1 & 206 & 2 & 816 & \texttt{-?????????}\\*
1 & 0 & 1 & 3 & 1 & 624 & 16 & 384 & \texttt{--+?????--}\\
2 & 1 & 2 & 3 & 1 & 224 & 32 & 768 & \texttt{---?????--}\\
3 & 2 & 3 & 1 & 1 & 768 & 256 & 0 & \texttt{----++----}\\
4 & 3 & 4 & 1 & 1 & 512 & 512 & 0 & \texttt{-----+----}\\
5 & 4 & 5 & 1 & 1 & 0 & 1024 & 0 & \texttt{----------}\\
6 & 2 & 3 & 1 & 2 & 0 & 1024 & 0 & \texttt{----------}\\
7 & 2 & 3 & 1 & 2 & 0 & 1024 & 0 & \texttt{----------}\\
8 & 0 & 1 & 1 & 2 & 0 & 1024 & 0 & \texttt{----------}\\
9 & 0 & 1 & 1 & 2 & 0 & 1024 & 0 & \texttt{----------}\\
\multicolumn{9}{@{}l@{}}{\rule{0pt}{2.6ex}\textbf{T10.276}\enspace \texttt{0123453234}\enspace $n{=}10$, $|H|{=}3$, len 6, br 2, $|\mathrm{Aut}|{=}1$, orb 10/10, hist-orb 3/3; M 27648, res 6176, rev 10726: ok}\\*
0 & -- & 0 & 3 & 1 & 834 & 2 & 188 & \texttt{-+????????}\\*
1 & 0 & 1 & 3 & 1 & 644 & 4 & 376 & \texttt{--????????}\\
2 & 1 & 2 & 2 & 1 & 448 & 64 & 512 & \texttt{---????---}\\
3 & 2 & 3 & 1 & 1 & 768 & 256 & 0 & \texttt{----++----}\\
4 & 3 & 4 & 1 & 1 & 512 & 512 & 0 & \texttt{-----+----}\\
5 & 4 & 5 & 1 & 1 & 0 & 1024 & 0 & \texttt{----------}\\
6 & 2 & 3 & 1 & 1 & 0 & 1024 & 0 & \texttt{----------}\\
7 & 1 & 2 & 1 & 1 & 768 & 256 & 0 & \texttt{--------++}\\
8 & 7 & 3 & 1 & 1 & 512 & 512 & 0 & \texttt{---------+}\\
9 & 8 & 4 & 1 & 1 & 0 & 1024 & 0 & \texttt{----------}\\
\multicolumn{9}{@{}l@{}}{\rule{0pt}{2.6ex}\textbf{T10.277}\enspace \texttt{0123453233}\enspace $n{=}10$, $|H|{=}4$, len 6, br 3, $|\mathrm{Aut}|{=}2$, orb 9/8, hist-orb 3/2; M 26624, res 9024, rev 9202: ok}\\*
0 & -- & 0 & 4 & 1 & 742 & 2 & 280 & \texttt{-+????????}\\*
1 & 0 & 1 & 4 & 1 & 460 & 4 & 560 & \texttt{--????????}\\
2 & 1 & 2 & 2 & 1 & 448 & 64 & 512 & \texttt{---????---}\\
3 & 2 & 3 & 1 & 1 & 768 & 256 & 0 & \texttt{----++----}\\
4 & 3 & 4 & 1 & 1 & 512 & 512 & 0 & \texttt{-----+----}\\
5 & 4 & 5 & 1 & 1 & 0 & 1024 & 0 & \texttt{----------}\\
6 & 2 & 3 & 1 & 1 & 0 & 1024 & 0 & \texttt{----------}\\
7 & 1 & 2 & 2 & 1 & 256 & 256 & 512 & \texttt{--------??}\\
8 & 7 & 3 & 1 & 2 & 0 & 1024 & 0 & \texttt{----------}\\
9 & 7 & 3 & 1 & 2 & 0 & 1024 & 0 & \texttt{----------}\\
\multicolumn{9}{@{}l@{}}{\rule{0pt}{2.6ex}\textbf{T10.278}\enspace \texttt{0123453232}\enspace $n{=}10$, $|H|{=}4$, len 6, br 2, $|\mathrm{Aut}|{=}1$, orb 10/10, hist-orb 4/4; M 25600, res 10592, rev 8022: ok}\\*
0 & -- & 0 & 4 & 1 & 650 & 2 & 372 & \texttt{-+????????}\\*
1 & 0 & 1 & 4 & 1 & 276 & 4 & 744 & \texttt{--????????}\\
2 & 1 & 2 & 2 & 1 & 448 & 64 & 512 & \texttt{---????---}\\
3 & 2 & 3 & 1 & 1 & 768 & 256 & 0 & \texttt{----++----}\\
4 & 3 & 4 & 1 & 1 & 512 & 512 & 0 & \texttt{-----+----}\\
5 & 4 & 5 & 1 & 1 & 0 & 1024 & 0 & \texttt{----------}\\
6 & 2 & 3 & 1 & 1 & 0 & 1024 & 0 & \texttt{----------}\\
7 & 1 & 2 & 1 & 1 & 512 & 512 & 0 & \texttt{--------+-}\\
8 & 7 & 3 & 1 & 1 & 0 & 1024 & 0 & \texttt{----------}\\
9 & 1 & 2 & 1 & 1 & 0 & 1024 & 0 & \texttt{----------}\\
\multicolumn{9}{@{}l@{}}{\rule{0pt}{2.6ex}\textbf{T10.279}\enspace \texttt{0123453231}\enspace $n{=}10$, $|H|{=}4$, len 6, br 3, $|\mathrm{Aut}|{=}1$, orb 10/10, hist-orb 4/4; M 24576, res 10332, rev 7510: ok}\\*
0 & -- & 0 & 4 & 1 & 394 & 2 & 628 & \texttt{-?????????}\\*
1 & 0 & 1 & 3 & 1 & 552 & 8 & 464 & \texttt{--???????-}\\
2 & 1 & 2 & 2 & 1 & 448 & 64 & 512 & \texttt{---????---}\\
3 & 2 & 3 & 1 & 1 & 768 & 256 & 0 & \texttt{----++----}\\
4 & 3 & 4 & 1 & 1 & 512 & 512 & 0 & \texttt{-----+----}\\
5 & 4 & 5 & 1 & 1 & 0 & 1024 & 0 & \texttt{----------}\\
6 & 2 & 3 & 1 & 1 & 0 & 1024 & 0 & \texttt{----------}\\
7 & 1 & 2 & 1 & 1 & 512 & 512 & 0 & \texttt{--------+-}\\
8 & 7 & 3 & 1 & 1 & 0 & 1024 & 0 & \texttt{----------}\\
9 & 0 & 1 & 1 & 1 & 0 & 1024 & 0 & \texttt{----------}\\
\multicolumn{9}{@{}l@{}}{\rule{0pt}{2.6ex}\textbf{T10.280}\enspace \texttt{0123453222}\enspace $n{=}10$, $|H|{=}5$, len 6, br 2, $|\mathrm{Aut}|{=}6$, orb 8/7, hist-orb 3/2; M 24576, res 12800, rev 6458: ok}\\*
0 & -- & 0 & 5 & 1 & 558 & 2 & 464 & \texttt{-+????????}\\*
1 & 0 & 1 & 5 & 1 & 92 & 4 & 928 & \texttt{--????????}\\
2 & 1 & 2 & 2 & 1 & 448 & 64 & 512 & \texttt{---????---}\\
3 & 2 & 3 & 1 & 1 & 768 & 256 & 0 & \texttt{----++----}\\
4 & 3 & 4 & 1 & 1 & 512 & 512 & 0 & \texttt{-----+----}\\
5 & 4 & 5 & 1 & 1 & 0 & 1024 & 0 & \texttt{----------}\\
6 & 2 & 3 & 1 & 1 & 0 & 1024 & 0 & \texttt{----------}\\
7 & 1 & 2 & 1 & 3 & 0 & 1024 & 0 & \texttt{----------}\\
8 & 1 & 2 & 1 & 3 & 0 & 1024 & 0 & \texttt{----------}\\
9 & 1 & 2 & 1 & 3 & 0 & 1024 & 0 & \texttt{----------}\\
\multicolumn{9}{@{}l@{}}{\rule{0pt}{2.6ex}\textbf{T10.281}\enspace \texttt{0123453221}\enspace $n{=}10$, $|H|{=}5$, len 6, br 3, $|\mathrm{Aut}|{=}2$, orb 9/8, hist-orb 4/3; M 23552, res 13552, rev 5314: ok}\\*
0 & -- & 0 & 5 & 1 & 302 & 2 & 720 & \texttt{-?????????}\\*
1 & 0 & 1 & 4 & 1 & 184 & 8 & 832 & \texttt{--???????-}\\
2 & 1 & 2 & 2 & 1 & 448 & 64 & 512 & \texttt{---????---}\\
3 & 2 & 3 & 1 & 1 & 768 & 256 & 0 & \texttt{----++----}\\
4 & 3 & 4 & 1 & 1 & 512 & 512 & 0 & \texttt{-----+----}\\
5 & 4 & 5 & 1 & 1 & 0 & 1024 & 0 & \texttt{----------}\\
6 & 2 & 3 & 1 & 1 & 0 & 1024 & 0 & \texttt{----------}\\
7 & 1 & 2 & 1 & 2 & 0 & 1024 & 0 & \texttt{----------}\\
8 & 1 & 2 & 1 & 2 & 0 & 1024 & 0 & \texttt{----------}\\
9 & 0 & 1 & 1 & 1 & 0 & 1024 & 0 & \texttt{----------}\\
\multicolumn{9}{@{}l@{}}{\rule{0pt}{2.6ex}\textbf{T10.282}\enspace \texttt{0123453212}\enspace $n{=}10$, $|H|{=}4$, len 6, br 3, $|\mathrm{Aut}|{=}1$, orb 10/10, hist-orb 4/4; M 23552, res 9684, rev 7258: ok}\\*
0 & -- & 0 & 4 & 1 & 522 & 2 & 500 & \texttt{-?????????}\\*
1 & 0 & 1 & 3 & 1 & 368 & 16 & 640 & \texttt{--??????--}\\
2 & 1 & 2 & 2 & 1 & 448 & 64 & 512 & \texttt{---????---}\\
3 & 2 & 3 & 1 & 1 & 768 & 256 & 0 & \texttt{----++----}\\
4 & 3 & 4 & 1 & 1 & 512 & 512 & 0 & \texttt{-----+----}\\
5 & 4 & 5 & 1 & 1 & 0 & 1024 & 0 & \texttt{----------}\\
6 & 2 & 3 & 1 & 1 & 0 & 1024 & 0 & \texttt{----------}\\
7 & 1 & 2 & 1 & 1 & 0 & 1024 & 0 & \texttt{----------}\\
8 & 0 & 1 & 1 & 1 & 512 & 512 & 0 & \texttt{---------+}\\
9 & 8 & 2 & 1 & 1 & 0 & 1024 & 0 & \texttt{----------}\\
\multicolumn{9}{@{}l@{}}{\rule{0pt}{2.6ex}\textbf{T10.283}\enspace \texttt{0123453211}\enspace $n{=}10$, $|H|{=}5$, len 6, br 3, $|\mathrm{Aut}|{=}2$, orb 9/8, hist-orb 4/3; M 22528, res 12816, rev 5074: ok}\\*
0 & -- & 0 & 5 & 1 & 174 & 2 & 848 & \texttt{-?????????}\\*
1 & 0 & 1 & 3 & 1 & 368 & 16 & 640 & \texttt{--??????--}\\
2 & 1 & 2 & 2 & 1 & 448 & 64 & 512 & \texttt{---????---}\\
3 & 2 & 3 & 1 & 1 & 768 & 256 & 0 & \texttt{----++----}\\
4 & 3 & 4 & 1 & 1 & 512 & 512 & 0 & \texttt{-----+----}\\
5 & 4 & 5 & 1 & 1 & 0 & 1024 & 0 & \texttt{----------}\\
6 & 2 & 3 & 1 & 1 & 0 & 1024 & 0 & \texttt{----------}\\
7 & 1 & 2 & 1 & 1 & 0 & 1024 & 0 & \texttt{----------}\\
8 & 0 & 1 & 1 & 2 & 0 & 1024 & 0 & \texttt{----------}\\
9 & 0 & 1 & 1 & 2 & 0 & 1024 & 0 & \texttt{----------}\\
\multicolumn{9}{@{}l@{}}{\rule{0pt}{2.6ex}\textbf{T10.284}\enspace \texttt{0123453123}\enspace $n{=}10$, $|H|{=}3$, len 6, br 2, $|\mathrm{Aut}|{=}1$, orb 10/10, hist-orb 3/3; M 24576, res 5084, rev 9606: ok}\\*
0 & -- & 0 & 3 & 1 & 770 & 2 & 252 & \texttt{-?????????}\\*
1 & 0 & 1 & 2 & 1 & 736 & 32 & 256 & \texttt{--+????---}\\
2 & 1 & 2 & 2 & 1 & 448 & 64 & 512 & \texttt{---????---}\\
3 & 2 & 3 & 1 & 1 & 768 & 256 & 0 & \texttt{----++----}\\
4 & 3 & 4 & 1 & 1 & 512 & 512 & 0 & \texttt{-----+----}\\
5 & 4 & 5 & 1 & 1 & 0 & 1024 & 0 & \texttt{----------}\\
6 & 2 & 3 & 1 & 1 & 0 & 1024 & 0 & \texttt{----------}\\
7 & 0 & 1 & 1 & 1 & 768 & 256 & 0 & \texttt{--------++}\\
8 & 7 & 2 & 1 & 1 & 512 & 512 & 0 & \texttt{---------+}\\
9 & 8 & 3 & 1 & 1 & 0 & 1024 & 0 & \texttt{----------}\\
\multicolumn{9}{@{}l@{}}{\rule{0pt}{2.6ex}\textbf{T10.285}\enspace \texttt{0123453122}\enspace $n{=}10$, $|H|{=}4$, len 6, br 3, $|\mathrm{Aut}|{=}2$, orb 9/8, hist-orb 3/2; M 23552, res 7832, rev 7966: ok}\\*
0 & -- & 0 & 4 & 1 & 550 & 2 & 472 & \texttt{-?????????}\\*
1 & 0 & 1 & 2 & 1 & 736 & 32 & 256 & \texttt{--+????---}\\
2 & 1 & 2 & 2 & 1 & 448 & 64 & 512 & \texttt{---????---}\\
3 & 2 & 3 & 1 & 1 & 768 & 256 & 0 & \texttt{----++----}\\
4 & 3 & 4 & 1 & 1 & 512 & 512 & 0 & \texttt{-----+----}\\
5 & 4 & 5 & 1 & 1 & 0 & 1024 & 0 & \texttt{----------}\\
6 & 2 & 3 & 1 & 1 & 0 & 1024 & 0 & \texttt{----------}\\
7 & 0 & 1 & 2 & 1 & 256 & 256 & 512 & \texttt{--------??}\\
8 & 7 & 2 & 1 & 2 & 0 & 1024 & 0 & \texttt{----------}\\
9 & 7 & 2 & 1 & 2 & 0 & 1024 & 0 & \texttt{----------}\\
\multicolumn{9}{@{}l@{}}{\rule{0pt}{2.6ex}\textbf{T10.286}\enspace \texttt{0123453121}\enspace $n{=}10$, $|H|{=}4$, len 6, br 2, $|\mathrm{Aut}|{=}1$, orb 10/10, hist-orb 4/4; M 22528, res 9044, rev 6910: ok}\\*
0 & -- & 0 & 4 & 1 & 330 & 2 & 692 & \texttt{-?????????}\\*
1 & 0 & 1 & 2 & 1 & 736 & 32 & 256 & \texttt{--+????---}\\
2 & 1 & 2 & 2 & 1 & 448 & 64 & 512 & \texttt{---????---}\\
3 & 2 & 3 & 1 & 1 & 768 & 256 & 0 & \texttt{----++----}\\
4 & 3 & 4 & 1 & 1 & 512 & 512 & 0 & \texttt{-----+----}\\
5 & 4 & 5 & 1 & 1 & 0 & 1024 & 0 & \texttt{----------}\\
6 & 2 & 3 & 1 & 1 & 0 & 1024 & 0 & \texttt{----------}\\
7 & 0 & 1 & 1 & 1 & 512 & 512 & 0 & \texttt{--------+-}\\
8 & 7 & 2 & 1 & 1 & 0 & 1024 & 0 & \texttt{----------}\\
9 & 0 & 1 & 1 & 1 & 0 & 1024 & 0 & \texttt{----------}\\
\multicolumn{9}{@{}l@{}}{\rule{0pt}{2.6ex}\textbf{T10.287}\enspace \texttt{0123453111}\enspace $n{=}10$, $|H|{=}5$, len 6, br 2, $|\mathrm{Aut}|{=}6$, orb 8/7, hist-orb 3/2; M 21504, res 11024, rev 5342: ok}\\*
0 & -- & 0 & 5 & 1 & 110 & 2 & 912 & \texttt{-?????????}\\*
1 & 0 & 1 & 2 & 1 & 736 & 32 & 256 & \texttt{--+????---}\\
2 & 1 & 2 & 2 & 1 & 448 & 64 & 512 & \texttt{---????---}\\
3 & 2 & 3 & 1 & 1 & 768 & 256 & 0 & \texttt{----++----}\\
4 & 3 & 4 & 1 & 1 & 512 & 512 & 0 & \texttt{-----+----}\\
5 & 4 & 5 & 1 & 1 & 0 & 1024 & 0 & \texttt{----------}\\
6 & 2 & 3 & 1 & 1 & 0 & 1024 & 0 & \texttt{----------}\\
7 & 0 & 1 & 1 & 3 & 0 & 1024 & 0 & \texttt{----------}\\
8 & 0 & 1 & 1 & 3 & 0 & 1024 & 0 & \texttt{----------}\\
9 & 0 & 1 & 1 & 3 & 0 & 1024 & 0 & \texttt{----------}\\
\multicolumn{9}{@{}l@{}}{\rule{0pt}{2.6ex}\textbf{T10.288}\enspace \texttt{0123452345}\enspace $n{=}10$, $|H|{=}2$, len 6, br 1, $|\mathrm{Aut}|{=}2$, orb 6/2, hist-orb 1/0; M 29696, res 1440, rev 12490: ok}\\*
0 & -- & 0 & 2 & 1 & 962 & 2 & 60 & \texttt{-+????????}\\*
1 & 0 & 1 & 2 & 1 & 900 & 4 & 120 & \texttt{--????????}\\
2 & 1 & 2 & 1 & 2 & 896 & 128 & 0 & \texttt{---+++----}\\
3 & 2 & 3 & 1 & 2 & 768 & 256 & 0 & \texttt{----++----}\\
4 & 3 & 4 & 1 & 2 & 512 & 512 & 0 & \texttt{-----+----}\\
5 & 4 & 5 & 1 & 2 & 0 & 1024 & 0 & \texttt{----------}\\
6 & 1 & 2 & 1 & 2 & 896 & 128 & 0 & \texttt{-------+++}\\
7 & 6 & 3 & 1 & 2 & 768 & 256 & 0 & \texttt{--------++}\\
8 & 7 & 4 & 1 & 2 & 512 & 512 & 0 & \texttt{---------+}\\
9 & 8 & 5 & 1 & 2 & 0 & 1024 & 0 & \texttt{----------}\\
\multicolumn{9}{@{}l@{}}{\rule{0pt}{2.6ex}\textbf{T10.289}\enspace \texttt{0123452344}\enspace $n{=}10$, $|H|{=}3$, len 6, br 2, $|\mathrm{Aut}|{=}2$, orb 9/8, hist-orb 2/1; M 28672, res 4032, rev 11770: ok}\\*
0 & -- & 0 & 3 & 1 & 902 & 2 & 120 & \texttt{-+????????}\\*
1 & 0 & 1 & 3 & 1 & 780 & 4 & 240 & \texttt{--????????}\\
2 & 1 & 2 & 1 & 1 & 896 & 128 & 0 & \texttt{---+++----}\\
3 & 2 & 3 & 1 & 1 & 768 & 256 & 0 & \texttt{----++----}\\
4 & 3 & 4 & 1 & 1 & 512 & 512 & 0 & \texttt{-----+----}\\
5 & 4 & 5 & 1 & 1 & 0 & 1024 & 0 & \texttt{----------}\\
6 & 1 & 2 & 2 & 1 & 640 & 128 & 256 & \texttt{-------+??}\\
7 & 6 & 3 & 2 & 1 & 256 & 256 & 512 & \texttt{--------??}\\
8 & 7 & 4 & 1 & 2 & 0 & 1024 & 0 & \texttt{----------}\\
9 & 7 & 4 & 1 & 2 & 0 & 1024 & 0 & \texttt{----------}\\
\multicolumn{9}{@{}l@{}}{\rule{0pt}{2.6ex}\textbf{T10.290}\enspace \texttt{0123452343}\enspace $n{=}10$, $|H|{=}3$, len 6, br 2, $|\mathrm{Aut}|{=}1$, orb 10/10, hist-orb 3/3; M 27648, res 5472, rev 10894: ok}\\*
0 & -- & 0 & 3 & 1 & 842 & 2 & 180 & \texttt{-+????????}\\*
1 & 0 & 1 & 3 & 1 & 660 & 4 & 360 & \texttt{--????????}\\
2 & 1 & 2 & 1 & 1 & 896 & 128 & 0 & \texttt{---+++----}\\
3 & 2 & 3 & 1 & 1 & 768 & 256 & 0 & \texttt{----++----}\\
4 & 3 & 4 & 1 & 1 & 512 & 512 & 0 & \texttt{-----+----}\\
5 & 4 & 5 & 1 & 1 & 0 & 1024 & 0 & \texttt{----------}\\
6 & 1 & 2 & 2 & 1 & 384 & 128 & 512 & \texttt{-------???}\\
7 & 6 & 3 & 1 & 1 & 512 & 512 & 0 & \texttt{--------+-}\\
8 & 7 & 4 & 1 & 1 & 0 & 1024 & 0 & \texttt{----------}\\
9 & 6 & 3 & 1 & 1 & 0 & 1024 & 0 & \texttt{----------}\\
\multicolumn{9}{@{}l@{}}{\rule{0pt}{2.6ex}\textbf{T10.291}\enspace \texttt{0123452342}\enspace $n{=}10$, $|H|{=}3$, len 6, br 1, $|\mathrm{Aut}|{=}1$, orb 10/10, hist-orb 3/3; M 26624, res 7200, rev 9562: ok}\\*
0 & -- & 0 & 3 & 1 & 722 & 2 & 300 & \texttt{-+????????}\\*
1 & 0 & 1 & 3 & 1 & 420 & 4 & 600 & \texttt{--????????}\\
2 & 1 & 2 & 1 & 1 & 896 & 128 & 0 & \texttt{---+++----}\\
3 & 2 & 3 & 1 & 1 & 768 & 256 & 0 & \texttt{----++----}\\
4 & 3 & 4 & 1 & 1 & 512 & 512 & 0 & \texttt{-----+----}\\
5 & 4 & 5 & 1 & 1 & 0 & 1024 & 0 & \texttt{----------}\\
6 & 1 & 2 & 1 & 1 & 768 & 256 & 0 & \texttt{-------++-}\\
7 & 6 & 3 & 1 & 1 & 512 & 512 & 0 & \texttt{--------+-}\\
8 & 7 & 4 & 1 & 1 & 0 & 1024 & 0 & \texttt{----------}\\
9 & 1 & 2 & 1 & 1 & 0 & 1024 & 0 & \texttt{----------}\\
\multicolumn{9}{@{}l@{}}{\rule{0pt}{2.6ex}\textbf{T10.292}\enspace \texttt{0123452341}\enspace $n{=}10$, $|H|{=}3$, len 6, br 2, $|\mathrm{Aut}|{=}1$, orb 10/10, hist-orb 3/3; M 25600, res 6148, rev 9346: ok}\\*
0 & -- & 0 & 3 & 1 & 466 & 2 & 556 & \texttt{-?????????}\\*
1 & 0 & 1 & 2 & 1 & 840 & 8 & 176 & \texttt{--???????-}\\
2 & 1 & 2 & 1 & 1 & 896 & 128 & 0 & \texttt{---+++----}\\
3 & 2 & 3 & 1 & 1 & 768 & 256 & 0 & \texttt{----++----}\\
4 & 3 & 4 & 1 & 1 & 512 & 512 & 0 & \texttt{-----+----}\\
5 & 4 & 5 & 1 & 1 & 0 & 1024 & 0 & \texttt{----------}\\
6 & 1 & 2 & 1 & 1 & 768 & 256 & 0 & \texttt{-------++-}\\
7 & 6 & 3 & 1 & 1 & 512 & 512 & 0 & \texttt{--------+-}\\
8 & 7 & 4 & 1 & 1 & 0 & 1024 & 0 & \texttt{----------}\\
9 & 0 & 1 & 1 & 1 & 0 & 1024 & 0 & \texttt{----------}\\
\multicolumn{9}{@{}l@{}}{\rule{0pt}{2.6ex}\textbf{T10.293}\enspace \texttt{0123452333}\enspace $n{=}10$, $|H|{=}4$, len 6, br 2, $|\mathrm{Aut}|{=}6$, orb 8/7, hist-orb 2/1; M 26624, res 7488, rev 9814: ok}\\*
0 & -- & 0 & 4 & 1 & 782 & 2 & 240 & \texttt{-+????????}\\*
1 & 0 & 1 & 4 & 1 & 540 & 4 & 480 & \texttt{--????????}\\
2 & 1 & 2 & 1 & 1 & 896 & 128 & 0 & \texttt{---+++----}\\
3 & 2 & 3 & 1 & 1 & 768 & 256 & 0 & \texttt{----++----}\\
4 & 3 & 4 & 1 & 1 & 512 & 512 & 0 & \texttt{-----+----}\\
5 & 4 & 5 & 1 & 1 & 0 & 1024 & 0 & \texttt{----------}\\
6 & 1 & 2 & 3 & 1 & 128 & 128 & 768 & \texttt{-------???}\\
7 & 6 & 3 & 1 & 3 & 0 & 1024 & 0 & \texttt{----------}\\
8 & 6 & 3 & 1 & 3 & 0 & 1024 & 0 & \texttt{----------}\\
9 & 6 & 3 & 1 & 3 & 0 & 1024 & 0 & \texttt{----------}\\
\multicolumn{9}{@{}l@{}}{\rule{0pt}{2.6ex}\textbf{T10.294}\enspace \texttt{0123452332}\enspace $n{=}10$, $|H|{=}4$, len 6, br 2, $|\mathrm{Aut}|{=}2$, orb 9/8, hist-orb 3/2; M 25600, res 9280, rev 8302: ok}\\*
0 & -- & 0 & 4 & 1 & 662 & 2 & 360 & \texttt{-+????????}\\*
1 & 0 & 1 & 4 & 1 & 300 & 4 & 720 & \texttt{--????????}\\
2 & 1 & 2 & 1 & 1 & 896 & 128 & 0 & \texttt{---+++----}\\
3 & 2 & 3 & 1 & 1 & 768 & 256 & 0 & \texttt{----++----}\\
4 & 3 & 4 & 1 & 1 & 512 & 512 & 0 & \texttt{-----+----}\\
5 & 4 & 5 & 1 & 1 & 0 & 1024 & 0 & \texttt{----------}\\
6 & 1 & 2 & 2 & 1 & 256 & 256 & 512 & \texttt{-------??-}\\
7 & 6 & 3 & 1 & 2 & 0 & 1024 & 0 & \texttt{----------}\\
8 & 6 & 3 & 1 & 2 & 0 & 1024 & 0 & \texttt{----------}\\
9 & 1 & 2 & 1 & 1 & 0 & 1024 & 0 & \texttt{----------}\\
\multicolumn{9}{@{}l@{}}{\rule{0pt}{2.6ex}\textbf{T10.295}\enspace \texttt{0123452331}\enspace $n{=}10$, $|H|{=}4$, len 6, br 3, $|\mathrm{Aut}|{=}2$, orb 9/8, hist-orb 3/2; M 24576, res 8888, rev 7870: ok}\\*
0 & -- & 0 & 4 & 1 & 406 & 2 & 616 & \texttt{-?????????}\\*
1 & 0 & 1 & 3 & 1 & 600 & 8 & 416 & \texttt{--???????-}\\
2 & 1 & 2 & 1 & 1 & 896 & 128 & 0 & \texttt{---+++----}\\
3 & 2 & 3 & 1 & 1 & 768 & 256 & 0 & \texttt{----++----}\\
4 & 3 & 4 & 1 & 1 & 512 & 512 & 0 & \texttt{-----+----}\\
5 & 4 & 5 & 1 & 1 & 0 & 1024 & 0 & \texttt{----------}\\
6 & 1 & 2 & 2 & 1 & 256 & 256 & 512 & \texttt{-------??-}\\
7 & 6 & 3 & 1 & 2 & 0 & 1024 & 0 & \texttt{----------}\\
8 & 6 & 3 & 1 & 2 & 0 & 1024 & 0 & \texttt{----------}\\
9 & 0 & 1 & 1 & 1 & 0 & 1024 & 0 & \texttt{----------}\\
\multicolumn{9}{@{}l@{}}{\rule{0pt}{2.6ex}\textbf{T10.296}\enspace \texttt{0123452323}\enspace $n{=}10$, $|H|{=}3$, len 6, br 1, $|\mathrm{Aut}|{=}2$, orb 8/6, hist-orb 2/1; M 25600, res 5760, rev 10030: ok}\\*
0 & -- & 0 & 3 & 1 & 782 & 2 & 240 & \texttt{-+????????}\\*
1 & 0 & 1 & 3 & 1 & 540 & 4 & 480 & \texttt{--????????}\\
2 & 1 & 2 & 1 & 1 & 896 & 128 & 0 & \texttt{---+++----}\\
3 & 2 & 3 & 1 & 1 & 768 & 256 & 0 & \texttt{----++----}\\
4 & 3 & 4 & 1 & 1 & 512 & 512 & 0 & \texttt{-----+----}\\
5 & 4 & 5 & 1 & 1 & 0 & 1024 & 0 & \texttt{----------}\\
6 & 1 & 2 & 1 & 2 & 512 & 512 & 0 & \texttt{-------+--}\\
7 & 6 & 3 & 1 & 2 & 0 & 1024 & 0 & \texttt{----------}\\
8 & 1 & 2 & 1 & 2 & 512 & 512 & 0 & \texttt{---------+}\\
9 & 8 & 3 & 1 & 2 & 0 & 1024 & 0 & \texttt{----------}\\
\multicolumn{9}{@{}l@{}}{\rule{0pt}{2.6ex}\textbf{T10.297}\enspace \texttt{0123452322}\enspace $n{=}10$, $|H|{=}4$, len 6, br 1, $|\mathrm{Aut}|{=}2$, orb 9/8, hist-orb 3/2; M 24576, res 10080, rev 7578: ok}\\*
0 & -- & 0 & 4 & 1 & 602 & 2 & 420 & \texttt{-+????????}\\*
1 & 0 & 1 & 4 & 1 & 180 & 4 & 840 & \texttt{--????????}\\
2 & 1 & 2 & 1 & 1 & 896 & 128 & 0 & \texttt{---+++----}\\
3 & 2 & 3 & 1 & 1 & 768 & 256 & 0 & \texttt{----++----}\\
4 & 3 & 4 & 1 & 1 & 512 & 512 & 0 & \texttt{-----+----}\\
5 & 4 & 5 & 1 & 1 & 0 & 1024 & 0 & \texttt{----------}\\
6 & 1 & 2 & 1 & 1 & 512 & 512 & 0 & \texttt{-------+--}\\
7 & 6 & 3 & 1 & 1 & 0 & 1024 & 0 & \texttt{----------}\\
8 & 1 & 2 & 1 & 2 & 0 & 1024 & 0 & \texttt{----------}\\
9 & 1 & 2 & 1 & 2 & 0 & 1024 & 0 & \texttt{----------}\\
\multicolumn{9}{@{}l@{}}{\rule{0pt}{2.6ex}\textbf{T10.298}\enspace \texttt{0123452321}\enspace $n{=}10$, $|H|{=}4$, len 6, br 2, $|\mathrm{Aut}|{=}1$, orb 10/10, hist-orb 4/4; M 23552, res 10348, rev 6754: ok}\\*
0 & -- & 0 & 4 & 1 & 346 & 2 & 676 & \texttt{-?????????}\\*
1 & 0 & 1 & 3 & 1 & 360 & 8 & 656 & \texttt{--???????-}\\
2 & 1 & 2 & 1 & 1 & 896 & 128 & 0 & \texttt{---+++----}\\
3 & 2 & 3 & 1 & 1 & 768 & 256 & 0 & \texttt{----++----}\\
4 & 3 & 4 & 1 & 1 & 512 & 512 & 0 & \texttt{-----+----}\\
5 & 4 & 5 & 1 & 1 & 0 & 1024 & 0 & \texttt{----------}\\
6 & 1 & 2 & 1 & 1 & 512 & 512 & 0 & \texttt{-------+--}\\
7 & 6 & 3 & 1 & 1 & 0 & 1024 & 0 & \texttt{----------}\\
8 & 1 & 2 & 1 & 1 & 0 & 1024 & 0 & \texttt{----------}\\
9 & 0 & 1 & 1 & 1 & 0 & 1024 & 0 & \texttt{----------}\\
\multicolumn{9}{@{}l@{}}{\rule{0pt}{2.6ex}\textbf{T10.299}\enspace \texttt{0123452312}\enspace $n{=}10$, $|H|{=}3$, len 6, br 2, $|\mathrm{Aut}|{=}1$, orb 10/10, hist-orb 3/3; M 23552, res 4896, rev 9178: ok}\\*
0 & -- & 0 & 3 & 1 & 654 & 2 & 368 & \texttt{-?????????}\\*
1 & 0 & 1 & 2 & 1 & 720 & 16 & 288 & \texttt{--??????--}\\
2 & 1 & 2 & 1 & 1 & 896 & 128 & 0 & \texttt{---+++----}\\
3 & 2 & 3 & 1 & 1 & 768 & 256 & 0 & \texttt{----++----}\\
4 & 3 & 4 & 1 & 1 & 512 & 512 & 0 & \texttt{-----+----}\\
5 & 4 & 5 & 1 & 1 & 0 & 1024 & 0 & \texttt{----------}\\
6 & 1 & 2 & 1 & 1 & 512 & 512 & 0 & \texttt{-------+--}\\
7 & 6 & 3 & 1 & 1 & 0 & 1024 & 0 & \texttt{----------}\\
8 & 0 & 1 & 1 & 1 & 512 & 512 & 0 & \texttt{---------+}\\
9 & 8 & 2 & 1 & 1 & 0 & 1024 & 0 & \texttt{----------}\\
\multicolumn{9}{@{}l@{}}{\rule{0pt}{2.6ex}\textbf{T10.300}\enspace \texttt{0123452311}\enspace $n{=}10$, $|H|{=}4$, len 6, br 2, $|\mathrm{Aut}|{=}2$, orb 9/8, hist-orb 3/2; M 22528, res 8820, rev 6810: ok}\\*
0 & -- & 0 & 4 & 1 & 218 & 2 & 804 & \texttt{-?????????}\\*
1 & 0 & 1 & 2 & 1 & 720 & 16 & 288 & \texttt{--??????--}\\
2 & 1 & 2 & 1 & 1 & 896 & 128 & 0 & \texttt{---+++----}\\
3 & 2 & 3 & 1 & 1 & 768 & 256 & 0 & \texttt{----++----}\\
4 & 3 & 4 & 1 & 1 & 512 & 512 & 0 & \texttt{-----+----}\\
5 & 4 & 5 & 1 & 1 & 0 & 1024 & 0 & \texttt{----------}\\
6 & 1 & 2 & 1 & 1 & 512 & 512 & 0 & \texttt{-------+--}\\
7 & 6 & 3 & 1 & 1 & 0 & 1024 & 0 & \texttt{----------}\\
8 & 0 & 1 & 1 & 2 & 0 & 1024 & 0 & \texttt{----------}\\
9 & 0 & 1 & 1 & 2 & 0 & 1024 & 0 & \texttt{----------}\\
\multicolumn{9}{@{}l@{}}{\rule{0pt}{2.6ex}\textbf{T10.301}\enspace \texttt{0123452222}\enspace $n{=}10$, $|H|{=}5$, len 6, br 1, $|\mathrm{Aut}|{=}24$, orb 7/6, hist-orb 2/1; M 23552, res 11520, rev 6470: ok}\\*
0 & -- & 0 & 5 & 1 & 542 & 2 & 480 & \texttt{-+????????}\\*
1 & 0 & 1 & 5 & 1 & 60 & 4 & 960 & \texttt{--????????}\\
2 & 1 & 2 & 1 & 1 & 896 & 128 & 0 & \texttt{---+++----}\\
3 & 2 & 3 & 1 & 1 & 768 & 256 & 0 & \texttt{----++----}\\
4 & 3 & 4 & 1 & 1 & 512 & 512 & 0 & \texttt{-----+----}\\
5 & 4 & 5 & 1 & 1 & 0 & 1024 & 0 & \texttt{----------}\\
6 & 1 & 2 & 1 & 4 & 0 & 1024 & 0 & \texttt{----------}\\
7 & 1 & 2 & 1 & 4 & 0 & 1024 & 0 & \texttt{----------}\\
8 & 1 & 2 & 1 & 4 & 0 & 1024 & 0 & \texttt{----------}\\
9 & 1 & 2 & 1 & 4 & 0 & 1024 & 0 & \texttt{----------}\\
\multicolumn{9}{@{}l@{}}{\rule{0pt}{2.6ex}\textbf{T10.302}\enspace \texttt{0123452221}\enspace $n{=}10$, $|H|{=}5$, len 6, br 2, $|\mathrm{Aut}|{=}6$, orb 8/7, hist-orb 3/2; M 22528, res 12448, rev 5190: ok}\\*
0 & -- & 0 & 5 & 1 & 286 & 2 & 736 & \texttt{-?????????}\\*
1 & 0 & 1 & 4 & 1 & 120 & 8 & 896 & \texttt{--???????-}\\
2 & 1 & 2 & 1 & 1 & 896 & 128 & 0 & \texttt{---+++----}\\
3 & 2 & 3 & 1 & 1 & 768 & 256 & 0 & \texttt{----++----}\\
4 & 3 & 4 & 1 & 1 & 512 & 512 & 0 & \texttt{-----+----}\\
5 & 4 & 5 & 1 & 1 & 0 & 1024 & 0 & \texttt{----------}\\
6 & 1 & 2 & 1 & 3 & 0 & 1024 & 0 & \texttt{----------}\\
7 & 1 & 2 & 1 & 3 & 0 & 1024 & 0 & \texttt{----------}\\
8 & 1 & 2 & 1 & 3 & 0 & 1024 & 0 & \texttt{----------}\\
9 & 0 & 1 & 1 & 1 & 0 & 1024 & 0 & \texttt{----------}\\
\multicolumn{9}{@{}l@{}}{\rule{0pt}{2.6ex}\textbf{T10.303}\enspace \texttt{0123452212}\enspace $n{=}10$, $|H|{=}4$, len 6, br 2, $|\mathrm{Aut}|{=}2$, orb 9/8, hist-orb 3/2; M 22528, res 9156, rev 6814: ok}\\*
0 & -- & 0 & 4 & 1 & 474 & 2 & 548 & \texttt{-?????????}\\*
1 & 0 & 1 & 3 & 1 & 240 & 16 & 768 & \texttt{--??????--}\\
2 & 1 & 2 & 1 & 1 & 896 & 128 & 0 & \texttt{---+++----}\\
3 & 2 & 3 & 1 & 1 & 768 & 256 & 0 & \texttt{----++----}\\
4 & 3 & 4 & 1 & 1 & 512 & 512 & 0 & \texttt{-----+----}\\
5 & 4 & 5 & 1 & 1 & 0 & 1024 & 0 & \texttt{----------}\\
6 & 1 & 2 & 1 & 2 & 0 & 1024 & 0 & \texttt{----------}\\
7 & 1 & 2 & 1 & 2 & 0 & 1024 & 0 & \texttt{----------}\\
8 & 0 & 1 & 1 & 1 & 512 & 512 & 0 & \texttt{---------+}\\
9 & 8 & 2 & 1 & 1 & 0 & 1024 & 0 & \texttt{----------}\\
\multicolumn{9}{@{}l@{}}{\rule{0pt}{2.6ex}\textbf{T10.304}\enspace \texttt{0123452211}\enspace $n{=}10$, $|H|{=}5$, len 6, br 2, $|\mathrm{Aut}|{=}4$, orb 8/6, hist-orb 3/1; M 21504, res 12000, rev 4766: ok}\\*
0 & -- & 0 & 5 & 1 & 158 & 2 & 864 & \texttt{-?????????}\\*
1 & 0 & 1 & 3 & 1 & 240 & 16 & 768 & \texttt{--??????--}\\
2 & 1 & 2 & 1 & 1 & 896 & 128 & 0 & \texttt{---+++----}\\
3 & 2 & 3 & 1 & 1 & 768 & 256 & 0 & \texttt{----++----}\\
4 & 3 & 4 & 1 & 1 & 512 & 512 & 0 & \texttt{-----+----}\\
5 & 4 & 5 & 1 & 1 & 0 & 1024 & 0 & \texttt{----------}\\
6 & 1 & 2 & 1 & 2 & 0 & 1024 & 0 & \texttt{----------}\\
7 & 1 & 2 & 1 & 2 & 0 & 1024 & 0 & \texttt{----------}\\
8 & 0 & 1 & 1 & 2 & 0 & 1024 & 0 & \texttt{----------}\\
9 & 0 & 1 & 1 & 2 & 0 & 1024 & 0 & \texttt{----------}\\
\multicolumn{9}{@{}l@{}}{\rule{0pt}{2.6ex}\textbf{T10.305}\enspace \texttt{0123452123}\enspace $n{=}10$, $|H|{=}3$, len 6, br 2, $|\mathrm{Aut}|{=}1$, orb 10/10, hist-orb 3/3; M 23552, res 5580, rev 8702: ok}\\*
0 & -- & 0 & 3 & 1 & 658 & 2 & 364 & \texttt{-?????????}\\*
1 & 0 & 1 & 2 & 1 & 480 & 32 & 512 & \texttt{--?????---}\\
2 & 1 & 2 & 1 & 1 & 896 & 128 & 0 & \texttt{---+++----}\\
3 & 2 & 3 & 1 & 1 & 768 & 256 & 0 & \texttt{----++----}\\
4 & 3 & 4 & 1 & 1 & 512 & 512 & 0 & \texttt{-----+----}\\
5 & 4 & 5 & 1 & 1 & 0 & 1024 & 0 & \texttt{----------}\\
6 & 1 & 2 & 1 & 1 & 0 & 1024 & 0 & \texttt{----------}\\
7 & 0 & 1 & 1 & 1 & 768 & 256 & 0 & \texttt{--------++}\\
8 & 7 & 2 & 1 & 1 & 512 & 512 & 0 & \texttt{---------+}\\
9 & 8 & 3 & 1 & 1 & 0 & 1024 & 0 & \texttt{----------}\\
\multicolumn{9}{@{}l@{}}{\rule{0pt}{2.6ex}\textbf{T10.306}\enspace \texttt{0123452122}\enspace $n{=}10$, $|H|{=}4$, len 6, br 3, $|\mathrm{Aut}|{=}2$, orb 9/8, hist-orb 3/2; M 22528, res 8040, rev 7142: ok}\\*
0 & -- & 0 & 4 & 1 & 470 & 2 & 552 & \texttt{-?????????}\\*
1 & 0 & 1 & 2 & 1 & 480 & 32 & 512 & \texttt{--?????---}\\
2 & 1 & 2 & 1 & 1 & 896 & 128 & 0 & \texttt{---+++----}\\
3 & 2 & 3 & 1 & 1 & 768 & 256 & 0 & \texttt{----++----}\\
4 & 3 & 4 & 1 & 1 & 512 & 512 & 0 & \texttt{-----+----}\\
5 & 4 & 5 & 1 & 1 & 0 & 1024 & 0 & \texttt{----------}\\
6 & 1 & 2 & 1 & 1 & 0 & 1024 & 0 & \texttt{----------}\\
7 & 0 & 1 & 2 & 1 & 256 & 256 & 512 & \texttt{--------??}\\
8 & 7 & 2 & 1 & 2 & 0 & 1024 & 0 & \texttt{----------}\\
9 & 7 & 2 & 1 & 2 & 0 & 1024 & 0 & \texttt{----------}\\
\multicolumn{9}{@{}l@{}}{\rule{0pt}{2.6ex}\textbf{T10.307}\enspace \texttt{0123452121}\enspace $n{=}10$, $|H|{=}4$, len 6, br 2, $|\mathrm{Aut}|{=}1$, orb 10/10, hist-orb 4/4; M 21504, res 8964, rev 6230: ok}\\*
0 & -- & 0 & 4 & 1 & 282 & 2 & 740 & \texttt{-?????????}\\*
1 & 0 & 1 & 2 & 1 & 480 & 32 & 512 & \texttt{--?????---}\\
2 & 1 & 2 & 1 & 1 & 896 & 128 & 0 & \texttt{---+++----}\\
3 & 2 & 3 & 1 & 1 & 768 & 256 & 0 & \texttt{----++----}\\
4 & 3 & 4 & 1 & 1 & 512 & 512 & 0 & \texttt{-----+----}\\
5 & 4 & 5 & 1 & 1 & 0 & 1024 & 0 & \texttt{----------}\\
6 & 1 & 2 & 1 & 1 & 0 & 1024 & 0 & \texttt{----------}\\
7 & 0 & 1 & 1 & 1 & 512 & 512 & 0 & \texttt{--------+-}\\
8 & 7 & 2 & 1 & 1 & 0 & 1024 & 0 & \texttt{----------}\\
9 & 0 & 1 & 1 & 1 & 0 & 1024 & 0 & \texttt{----------}\\
\multicolumn{9}{@{}l@{}}{\rule{0pt}{2.6ex}\textbf{T10.308}\enspace \texttt{0123452111}\enspace $n{=}10$, $|H|{=}5$, len 6, br 2, $|\mathrm{Aut}|{=}6$, orb 8/7, hist-orb 3/2; M 20480, res 10656, rev 4806: ok}\\*
0 & -- & 0 & 5 & 1 & 94 & 2 & 928 & \texttt{-?????????}\\*
1 & 0 & 1 & 2 & 1 & 480 & 32 & 512 & \texttt{--?????---}\\
2 & 1 & 2 & 1 & 1 & 896 & 128 & 0 & \texttt{---+++----}\\
3 & 2 & 3 & 1 & 1 & 768 & 256 & 0 & \texttt{----++----}\\
4 & 3 & 4 & 1 & 1 & 512 & 512 & 0 & \texttt{-----+----}\\
5 & 4 & 5 & 1 & 1 & 0 & 1024 & 0 & \texttt{----------}\\
6 & 1 & 2 & 1 & 1 & 0 & 1024 & 0 & \texttt{----------}\\
7 & 0 & 1 & 1 & 3 & 0 & 1024 & 0 & \texttt{----------}\\
8 & 0 & 1 & 1 & 3 & 0 & 1024 & 0 & \texttt{----------}\\
9 & 0 & 1 & 1 & 3 & 0 & 1024 & 0 & \texttt{----------}\\
\multicolumn{9}{@{}l@{}}{\rule{0pt}{2.6ex}\textbf{T10.309}\enspace \texttt{0123451234}\enspace $n{=}10$, $|H|{=}2$, len 6, br 1, $|\mathrm{Aut}|{=}1$, orb 10/10, hist-orb 2/2; M 25600, res 828, rev 10874: ok}\\*
0 & -- & 0 & 2 & 1 & 930 & 2 & 92 & \texttt{-?????????}\\*
1 & 0 & 1 & 1 & 1 & 960 & 64 & 0 & \texttt{--++++----}\\
2 & 1 & 2 & 1 & 1 & 896 & 128 & 0 & \texttt{---+++----}\\
3 & 2 & 3 & 1 & 1 & 768 & 256 & 0 & \texttt{----++----}\\
4 & 3 & 4 & 1 & 1 & 512 & 512 & 0 & \texttt{-----+----}\\
5 & 4 & 5 & 1 & 1 & 0 & 1024 & 0 & \texttt{----------}\\
6 & 0 & 1 & 1 & 1 & 896 & 128 & 0 & \texttt{-------+++}\\
7 & 6 & 2 & 1 & 1 & 768 & 256 & 0 & \texttt{--------++}\\
8 & 7 & 3 & 1 & 1 & 512 & 512 & 0 & \texttt{---------+}\\
9 & 8 & 4 & 1 & 1 & 0 & 1024 & 0 & \texttt{----------}\\
\multicolumn{9}{@{}l@{}}{\rule{0pt}{2.6ex}\textbf{T10.310}\enspace \texttt{0123451233}\enspace $n{=}10$, $|H|{=}3$, len 6, br 2, $|\mathrm{Aut}|{=}2$, orb 9/8, hist-orb 2/1; M 24576, res 3224, rev 9870: ok}\\*
0 & -- & 0 & 3 & 1 & 806 & 2 & 216 & \texttt{-?????????}\\*
1 & 0 & 1 & 1 & 1 & 960 & 64 & 0 & \texttt{--++++----}\\
2 & 1 & 2 & 1 & 1 & 896 & 128 & 0 & \texttt{---+++----}\\
3 & 2 & 3 & 1 & 1 & 768 & 256 & 0 & \texttt{----++----}\\
4 & 3 & 4 & 1 & 1 & 512 & 512 & 0 & \texttt{-----+----}\\
5 & 4 & 5 & 1 & 1 & 0 & 1024 & 0 & \texttt{----------}\\
6 & 0 & 1 & 2 & 1 & 640 & 128 & 256 & \texttt{-------+??}\\
7 & 6 & 2 & 2 & 1 & 256 & 256 & 512 & \texttt{--------??}\\
8 & 7 & 3 & 1 & 2 & 0 & 1024 & 0 & \texttt{----------}\\
9 & 7 & 3 & 1 & 2 & 0 & 1024 & 0 & \texttt{----------}\\
\multicolumn{9}{@{}l@{}}{\rule{0pt}{2.6ex}\textbf{T10.311}\enspace \texttt{0123451232}\enspace $n{=}10$, $|H|{=}3$, len 6, br 2, $|\mathrm{Aut}|{=}1$, orb 10/10, hist-orb 3/3; M 23552, res 4340, rev 9006: ok}\\*
0 & -- & 0 & 3 & 1 & 682 & 2 & 340 & \texttt{-?????????}\\*
1 & 0 & 1 & 1 & 1 & 960 & 64 & 0 & \texttt{--++++----}\\
2 & 1 & 2 & 1 & 1 & 896 & 128 & 0 & \texttt{---+++----}\\
3 & 2 & 3 & 1 & 1 & 768 & 256 & 0 & \texttt{----++----}\\
4 & 3 & 4 & 1 & 1 & 512 & 512 & 0 & \texttt{-----+----}\\
5 & 4 & 5 & 1 & 1 & 0 & 1024 & 0 & \texttt{----------}\\
6 & 0 & 1 & 2 & 1 & 384 & 128 & 512 & \texttt{-------???}\\
7 & 6 & 2 & 1 & 1 & 512 & 512 & 0 & \texttt{--------+-}\\
8 & 7 & 3 & 1 & 1 & 0 & 1024 & 0 & \texttt{----------}\\
9 & 6 & 2 & 1 & 1 & 0 & 1024 & 0 & \texttt{----------}\\
\multicolumn{9}{@{}l@{}}{\rule{0pt}{2.6ex}\textbf{T10.312}\enspace \texttt{0123451231}\enspace $n{=}10$, $|H|{=}3$, len 6, br 1, $|\mathrm{Aut}|{=}1$, orb 10/10, hist-orb 3/3; M 22528, res 5292, rev 8058: ok}\\*
0 & -- & 0 & 3 & 1 & 434 & 2 & 588 & \texttt{-?????????}\\*
1 & 0 & 1 & 1 & 1 & 960 & 64 & 0 & \texttt{--++++----}\\
2 & 1 & 2 & 1 & 1 & 896 & 128 & 0 & \texttt{---+++----}\\
3 & 2 & 3 & 1 & 1 & 768 & 256 & 0 & \texttt{----++----}\\
4 & 3 & 4 & 1 & 1 & 512 & 512 & 0 & \texttt{-----+----}\\
5 & 4 & 5 & 1 & 1 & 0 & 1024 & 0 & \texttt{----------}\\
6 & 0 & 1 & 1 & 1 & 768 & 256 & 0 & \texttt{-------++-}\\
7 & 6 & 2 & 1 & 1 & 512 & 512 & 0 & \texttt{--------+-}\\
8 & 7 & 3 & 1 & 1 & 0 & 1024 & 0 & \texttt{----------}\\
9 & 0 & 1 & 1 & 1 & 0 & 1024 & 0 & \texttt{----------}\\
\multicolumn{9}{@{}l@{}}{\rule{0pt}{2.6ex}\textbf{T10.313}\enspace \texttt{0123451222}\enspace $n{=}10$, $|H|{=}4$, len 6, br 2, $|\mathrm{Aut}|{=}6$, orb 8/7, hist-orb 2/1; M 22528, res 6096, rev 7754: ok}\\*
0 & -- & 0 & 4 & 1 & 558 & 2 & 464 & \texttt{-?????????}\\*
1 & 0 & 1 & 1 & 1 & 960 & 64 & 0 & \texttt{--++++----}\\
2 & 1 & 2 & 1 & 1 & 896 & 128 & 0 & \texttt{---+++----}\\
3 & 2 & 3 & 1 & 1 & 768 & 256 & 0 & \texttt{----++----}\\
4 & 3 & 4 & 1 & 1 & 512 & 512 & 0 & \texttt{-----+----}\\
5 & 4 & 5 & 1 & 1 & 0 & 1024 & 0 & \texttt{----------}\\
6 & 0 & 1 & 3 & 1 & 128 & 128 & 768 & \texttt{-------???}\\
7 & 6 & 2 & 1 & 3 & 0 & 1024 & 0 & \texttt{----------}\\
8 & 6 & 2 & 1 & 3 & 0 & 1024 & 0 & \texttt{----------}\\
9 & 6 & 2 & 1 & 3 & 0 & 1024 & 0 & \texttt{----------}\\
\multicolumn{9}{@{}l@{}}{\rule{0pt}{2.6ex}\textbf{T10.314}\enspace \texttt{0123451221}\enspace $n{=}10$, $|H|{=}4$, len 6, br 2, $|\mathrm{Aut}|{=}2$, orb 9/8, hist-orb 3/2; M 21504, res 7176, rev 6682: ok}\\*
0 & -- & 0 & 4 & 1 & 310 & 2 & 712 & \texttt{-?????????}\\*
1 & 0 & 1 & 1 & 1 & 960 & 64 & 0 & \texttt{--++++----}\\
2 & 1 & 2 & 1 & 1 & 896 & 128 & 0 & \texttt{---+++----}\\
3 & 2 & 3 & 1 & 1 & 768 & 256 & 0 & \texttt{----++----}\\
4 & 3 & 4 & 1 & 1 & 512 & 512 & 0 & \texttt{-----+----}\\
5 & 4 & 5 & 1 & 1 & 0 & 1024 & 0 & \texttt{----------}\\
6 & 0 & 1 & 2 & 1 & 256 & 256 & 512 & \texttt{-------??-}\\
7 & 6 & 2 & 1 & 2 & 0 & 1024 & 0 & \texttt{----------}\\
8 & 6 & 2 & 1 & 2 & 0 & 1024 & 0 & \texttt{----------}\\
9 & 0 & 1 & 1 & 1 & 0 & 1024 & 0 & \texttt{----------}\\
\multicolumn{9}{@{}l@{}}{\rule{0pt}{2.6ex}\textbf{T10.315}\enspace \texttt{0123451212}\enspace $n{=}10$, $|H|{=}3$, len 6, br 1, $|\mathrm{Aut}|{=}2$, orb 8/6, hist-orb 2/1; M 21504, res 4176, rev 8146: ok}\\*
0 & -- & 0 & 3 & 1 & 558 & 2 & 464 & \texttt{-?????????}\\*
1 & 0 & 1 & 1 & 1 & 960 & 64 & 0 & \texttt{--++++----}\\
2 & 1 & 2 & 1 & 1 & 896 & 128 & 0 & \texttt{---+++----}\\
3 & 2 & 3 & 1 & 1 & 768 & 256 & 0 & \texttt{----++----}\\
4 & 3 & 4 & 1 & 1 & 512 & 512 & 0 & \texttt{-----+----}\\
5 & 4 & 5 & 1 & 1 & 0 & 1024 & 0 & \texttt{----------}\\
6 & 0 & 1 & 1 & 2 & 512 & 512 & 0 & \texttt{-------+--}\\
7 & 6 & 2 & 1 & 2 & 0 & 1024 & 0 & \texttt{----------}\\
8 & 0 & 1 & 1 & 2 & 512 & 512 & 0 & \texttt{---------+}\\
9 & 8 & 2 & 1 & 2 & 0 & 1024 & 0 & \texttt{----------}\\
\multicolumn{9}{@{}l@{}}{\rule{0pt}{2.6ex}\textbf{T10.316}\enspace \texttt{0123451211}\enspace $n{=}10$, $|H|{=}4$, len 6, br 1, $|\mathrm{Aut}|{=}2$, orb 9/8, hist-orb 3/2; M 20480, res 7524, rev 6082: ok}\\*
0 & -- & 0 & 4 & 1 & 186 & 2 & 836 & \texttt{-?????????}\\*
1 & 0 & 1 & 1 & 1 & 960 & 64 & 0 & \texttt{--++++----}\\
2 & 1 & 2 & 1 & 1 & 896 & 128 & 0 & \texttt{---+++----}\\
3 & 2 & 3 & 1 & 1 & 768 & 256 & 0 & \texttt{----++----}\\
4 & 3 & 4 & 1 & 1 & 512 & 512 & 0 & \texttt{-----+----}\\
5 & 4 & 5 & 1 & 1 & 0 & 1024 & 0 & \texttt{----------}\\
6 & 0 & 1 & 1 & 1 & 512 & 512 & 0 & \texttt{-------+--}\\
7 & 6 & 2 & 1 & 1 & 0 & 1024 & 0 & \texttt{----------}\\
8 & 0 & 1 & 1 & 2 & 0 & 1024 & 0 & \texttt{----------}\\
9 & 0 & 1 & 1 & 2 & 0 & 1024 & 0 & \texttt{----------}\\
\multicolumn{9}{@{}l@{}}{\rule{0pt}{2.6ex}\textbf{T10.317}\enspace \texttt{0123451111}\enspace $n{=}10$, $|H|{=}5$, len 6, br 1, $|\mathrm{Aut}|{=}24$, orb 7/6, hist-orb 2/1; M 19456, res 8640, rev 4970: ok}\\*
0 & -- & 0 & 5 & 1 & 62 & 2 & 960 & \texttt{-?????????}\\*
1 & 0 & 1 & 1 & 1 & 960 & 64 & 0 & \texttt{--++++----}\\
2 & 1 & 2 & 1 & 1 & 896 & 128 & 0 & \texttt{---+++----}\\
3 & 2 & 3 & 1 & 1 & 768 & 256 & 0 & \texttt{----++----}\\
4 & 3 & 4 & 1 & 1 & 512 & 512 & 0 & \texttt{-----+----}\\
5 & 4 & 5 & 1 & 1 & 0 & 1024 & 0 & \texttt{----------}\\
6 & 0 & 1 & 1 & 4 & 0 & 1024 & 0 & \texttt{----------}\\
7 & 0 & 1 & 1 & 4 & 0 & 1024 & 0 & \texttt{----------}\\
8 & 0 & 1 & 1 & 4 & 0 & 1024 & 0 & \texttt{----------}\\
9 & 0 & 1 & 1 & 4 & 0 & 1024 & 0 & \texttt{----------}\\
\multicolumn{9}{@{}l@{}}{\rule{0pt}{2.6ex}\textbf{T10.318}\enspace \texttt{0123444444}\enspace $n{=}10$, $|H|{=}6$, len 5, br 1, $|\mathrm{Aut}|{=}720$, orb 5/4, hist-orb 1/0; M 30720, res 11160, rev 13282: ok}\\*
0 & -- & 0 & 6 & 1 & 898 & 2 & 124 & \texttt{-+++??????}\\*
1 & 0 & 1 & 6 & 1 & 772 & 4 & 248 & \texttt{--++??????}\\
2 & 1 & 2 & 6 & 1 & 520 & 8 & 496 & \texttt{---+??????}\\
3 & 2 & 3 & 6 & 1 & 16 & 16 & 992 & \texttt{----??????}\\
4 & 3 & 4 & 1 & 6 & 0 & 1024 & 0 & \texttt{----------}\\
5 & 3 & 4 & 1 & 6 & 0 & 1024 & 0 & \texttt{----------}\\
6 & 3 & 4 & 1 & 6 & 0 & 1024 & 0 & \texttt{----------}\\
7 & 3 & 4 & 1 & 6 & 0 & 1024 & 0 & \texttt{----------}\\
8 & 3 & 4 & 1 & 6 & 0 & 1024 & 0 & \texttt{----------}\\
9 & 3 & 4 & 1 & 6 & 0 & 1024 & 0 & \texttt{----------}\\
\multicolumn{9}{@{}l@{}}{\rule{0pt}{2.6ex}\textbf{T10.319}\enspace \texttt{0123444443}\enspace $n{=}10$, $|H|{=}6$, len 5, br 2, $|\mathrm{Aut}|{=}120$, orb 6/5, hist-orb 2/1; M 29696, res 13172, rev 10678: ok}\\*
0 & -- & 0 & 6 & 1 & 834 & 2 & 188 & \texttt{-++???????}\\*
1 & 0 & 1 & 6 & 1 & 644 & 4 & 376 & \texttt{--+???????}\\
2 & 1 & 2 & 6 & 1 & 264 & 8 & 752 & \texttt{---???????}\\
3 & 2 & 3 & 5 & 1 & 32 & 32 & 960 & \texttt{----?????-}\\
4 & 3 & 4 & 1 & 5 & 0 & 1024 & 0 & \texttt{----------}\\
5 & 3 & 4 & 1 & 5 & 0 & 1024 & 0 & \texttt{----------}\\
6 & 3 & 4 & 1 & 5 & 0 & 1024 & 0 & \texttt{----------}\\
7 & 3 & 4 & 1 & 5 & 0 & 1024 & 0 & \texttt{----------}\\
8 & 3 & 4 & 1 & 5 & 0 & 1024 & 0 & \texttt{----------}\\
9 & 2 & 3 & 1 & 1 & 0 & 1024 & 0 & \texttt{----------}\\
\multicolumn{9}{@{}l@{}}{\rule{0pt}{2.6ex}\textbf{T10.320}\enspace \texttt{0123444442}\enspace $n{=}10$, $|H|{=}6$, len 5, br 2, $|\mathrm{Aut}|{=}120$, orb 6/5, hist-orb 2/1; M 28672, res 14064, rev 9430: ok}\\*
0 & -- & 0 & 6 & 1 & 706 & 2 & 316 & \texttt{-+????????}\\*
1 & 0 & 1 & 6 & 1 & 388 & 4 & 632 & \texttt{--????????}\\
2 & 1 & 2 & 5 & 1 & 528 & 16 & 480 & \texttt{---+?????-}\\
3 & 2 & 3 & 5 & 1 & 32 & 32 & 960 & \texttt{----?????-}\\
4 & 3 & 4 & 1 & 5 & 0 & 1024 & 0 & \texttt{----------}\\
5 & 3 & 4 & 1 & 5 & 0 & 1024 & 0 & \texttt{----------}\\
6 & 3 & 4 & 1 & 5 & 0 & 1024 & 0 & \texttt{----------}\\
7 & 3 & 4 & 1 & 5 & 0 & 1024 & 0 & \texttt{----------}\\
8 & 3 & 4 & 1 & 5 & 0 & 1024 & 0 & \texttt{----------}\\
9 & 1 & 2 & 1 & 1 & 0 & 1024 & 0 & \texttt{----------}\\
\multicolumn{9}{@{}l@{}}{\rule{0pt}{2.6ex}\textbf{T10.321}\enspace \texttt{0123444441}\enspace $n{=}10$, $|H|{=}6$, len 5, br 2, $|\mathrm{Aut}|{=}120$, orb 6/5, hist-orb 2/1; M 27648, res 13188, rev 9450: ok}\\*
0 & -- & 0 & 6 & 1 & 450 & 2 & 572 & \texttt{-?????????}\\*
1 & 0 & 1 & 5 & 1 & 776 & 8 & 240 & \texttt{--++?????-}\\
2 & 1 & 2 & 5 & 1 & 528 & 16 & 480 & \texttt{---+?????-}\\
3 & 2 & 3 & 5 & 1 & 32 & 32 & 960 & \texttt{----?????-}\\
4 & 3 & 4 & 1 & 5 & 0 & 1024 & 0 & \texttt{----------}\\
5 & 3 & 4 & 1 & 5 & 0 & 1024 & 0 & \texttt{----------}\\
6 & 3 & 4 & 1 & 5 & 0 & 1024 & 0 & \texttt{----------}\\
7 & 3 & 4 & 1 & 5 & 0 & 1024 & 0 & \texttt{----------}\\
8 & 3 & 4 & 1 & 5 & 0 & 1024 & 0 & \texttt{----------}\\
9 & 0 & 1 & 1 & 1 & 0 & 1024 & 0 & \texttt{----------}\\
\multicolumn{9}{@{}l@{}}{\rule{0pt}{2.6ex}\textbf{T10.322}\enspace \texttt{0123444434}\enspace $n{=}10$, $|H|{=}5$, len 5, br 2, $|\mathrm{Aut}|{=}24$, orb 7/6, hist-orb 2/1; M 29696, res 10248, rev 11750: ok}\\*
0 & -- & 0 & 5 & 1 & 870 & 2 & 152 & \texttt{-++???????}\\*
1 & 0 & 1 & 5 & 1 & 716 & 4 & 304 & \texttt{--+???????}\\
2 & 1 & 2 & 5 & 1 & 408 & 8 & 608 & \texttt{---???????}\\
3 & 2 & 3 & 4 & 1 & 64 & 64 & 896 & \texttt{----????--}\\
4 & 3 & 4 & 1 & 4 & 0 & 1024 & 0 & \texttt{----------}\\
5 & 3 & 4 & 1 & 4 & 0 & 1024 & 0 & \texttt{----------}\\
6 & 3 & 4 & 1 & 4 & 0 & 1024 & 0 & \texttt{----------}\\
7 & 3 & 4 & 1 & 4 & 0 & 1024 & 0 & \texttt{----------}\\
8 & 2 & 3 & 1 & 1 & 512 & 512 & 0 & \texttt{---------+}\\
9 & 8 & 4 & 1 & 1 & 0 & 1024 & 0 & \texttt{----------}\\
\multicolumn{9}{@{}l@{}}{\rule{0pt}{2.6ex}\textbf{T10.323}\enspace \texttt{0123444433}\enspace $n{=}10$, $|H|{=}6$, len 5, br 2, $|\mathrm{Aut}|{=}48$, orb 6/4, hist-orb 2/0; M 28672, res 13580, rev 9470: ok}\\*
0 & -- & 0 & 6 & 1 & 802 & 2 & 220 & \texttt{-++???????}\\*
1 & 0 & 1 & 6 & 1 & 580 & 4 & 440 & \texttt{--+???????}\\
2 & 1 & 2 & 6 & 1 & 136 & 8 & 880 & \texttt{---???????}\\
3 & 2 & 3 & 4 & 1 & 64 & 64 & 896 & \texttt{----????--}\\
4 & 3 & 4 & 1 & 4 & 0 & 1024 & 0 & \texttt{----------}\\
5 & 3 & 4 & 1 & 4 & 0 & 1024 & 0 & \texttt{----------}\\
6 & 3 & 4 & 1 & 4 & 0 & 1024 & 0 & \texttt{----------}\\
7 & 3 & 4 & 1 & 4 & 0 & 1024 & 0 & \texttt{----------}\\
8 & 2 & 3 & 1 & 2 & 0 & 1024 & 0 & \texttt{----------}\\
9 & 2 & 3 & 1 & 2 & 0 & 1024 & 0 & \texttt{----------}\\
\multicolumn{9}{@{}l@{}}{\rule{0pt}{2.6ex}\textbf{T10.324}\enspace \texttt{0123444432}\enspace $n{=}10$, $|H|{=}6$, len 5, br 3, $|\mathrm{Aut}|{=}24$, orb 7/6, hist-orb 3/2; M 27648, res 15016, rev 7698: ok}\\*
0 & -- & 0 & 6 & 1 & 674 & 2 & 348 & \texttt{-+????????}\\*
1 & 0 & 1 & 6 & 1 & 324 & 4 & 696 & \texttt{--????????}\\
2 & 1 & 2 & 5 & 1 & 272 & 16 & 736 & \texttt{---??????-}\\
3 & 2 & 3 & 4 & 1 & 64 & 64 & 896 & \texttt{----????--}\\
4 & 3 & 4 & 1 & 4 & 0 & 1024 & 0 & \texttt{----------}\\
5 & 3 & 4 & 1 & 4 & 0 & 1024 & 0 & \texttt{----------}\\
6 & 3 & 4 & 1 & 4 & 0 & 1024 & 0 & \texttt{----------}\\
7 & 3 & 4 & 1 & 4 & 0 & 1024 & 0 & \texttt{----------}\\
8 & 2 & 3 & 1 & 1 & 0 & 1024 & 0 & \texttt{----------}\\
9 & 1 & 2 & 1 & 1 & 0 & 1024 & 0 & \texttt{----------}\\
\multicolumn{9}{@{}l@{}}{\rule{0pt}{2.6ex}\textbf{T10.325}\enspace \texttt{0123444431}\enspace $n{=}10$, $|H|{=}6$, len 5, br 3, $|\mathrm{Aut}|{=}24$, orb 7/6, hist-orb 3/2; M 26624, res 14492, rev 7462: ok}\\*
0 & -- & 0 & 6 & 1 & 418 & 2 & 604 & \texttt{-?????????}\\*
1 & 0 & 1 & 5 & 1 & 648 & 8 & 368 & \texttt{--+??????-}\\
2 & 1 & 2 & 5 & 1 & 272 & 16 & 736 & \texttt{---??????-}\\
3 & 2 & 3 & 4 & 1 & 64 & 64 & 896 & \texttt{----????--}\\
4 & 3 & 4 & 1 & 4 & 0 & 1024 & 0 & \texttt{----------}\\
5 & 3 & 4 & 1 & 4 & 0 & 1024 & 0 & \texttt{----------}\\
6 & 3 & 4 & 1 & 4 & 0 & 1024 & 0 & \texttt{----------}\\
7 & 3 & 4 & 1 & 4 & 0 & 1024 & 0 & \texttt{----------}\\
8 & 2 & 3 & 1 & 1 & 0 & 1024 & 0 & \texttt{----------}\\
9 & 0 & 1 & 1 & 1 & 0 & 1024 & 0 & \texttt{----------}\\
\multicolumn{9}{@{}l@{}}{\rule{0pt}{2.6ex}\textbf{T10.326}\enspace \texttt{0123444423}\enspace $n{=}10$, $|H|{=}5$, len 5, br 2, $|\mathrm{Aut}|{=}24$, orb 7/6, hist-orb 2/1; M 27648, res 9888, rev 10618: ok}\\*
0 & -- & 0 & 5 & 1 & 806 & 2 & 216 & \texttt{-+????????}\\*
1 & 0 & 1 & 5 & 1 & 588 & 4 & 432 & \texttt{--????????}\\
2 & 1 & 2 & 4 & 1 & 544 & 32 & 448 & \texttt{---+????--}\\
3 & 2 & 3 & 4 & 1 & 64 & 64 & 896 & \texttt{----????--}\\
4 & 3 & 4 & 1 & 4 & 0 & 1024 & 0 & \texttt{----------}\\
5 & 3 & 4 & 1 & 4 & 0 & 1024 & 0 & \texttt{----------}\\
6 & 3 & 4 & 1 & 4 & 0 & 1024 & 0 & \texttt{----------}\\
7 & 3 & 4 & 1 & 4 & 0 & 1024 & 0 & \texttt{----------}\\
8 & 1 & 2 & 1 & 1 & 512 & 512 & 0 & \texttt{---------+}\\
9 & 8 & 3 & 1 & 1 & 0 & 1024 & 0 & \texttt{----------}\\
\multicolumn{9}{@{}l@{}}{\rule{0pt}{2.6ex}\textbf{T10.327}\enspace \texttt{0123444422}\enspace $n{=}10$, $|H|{=}6$, len 5, br 2, $|\mathrm{Aut}|{=}48$, orb 6/4, hist-orb 2/0; M 26624, res 14592, rev 7374: ok}\\*
0 & -- & 0 & 6 & 1 & 610 & 2 & 412 & \texttt{-+????????}\\*
1 & 0 & 1 & 6 & 1 & 196 & 4 & 824 & \texttt{--????????}\\
2 & 1 & 2 & 4 & 1 & 544 & 32 & 448 & \texttt{---+????--}\\
3 & 2 & 3 & 4 & 1 & 64 & 64 & 896 & \texttt{----????--}\\
4 & 3 & 4 & 1 & 4 & 0 & 1024 & 0 & \texttt{----------}\\
5 & 3 & 4 & 1 & 4 & 0 & 1024 & 0 & \texttt{----------}\\
6 & 3 & 4 & 1 & 4 & 0 & 1024 & 0 & \texttt{----------}\\
7 & 3 & 4 & 1 & 4 & 0 & 1024 & 0 & \texttt{----------}\\
8 & 1 & 2 & 1 & 2 & 0 & 1024 & 0 & \texttt{----------}\\
9 & 1 & 2 & 1 & 2 & 0 & 1024 & 0 & \texttt{----------}\\
\multicolumn{9}{@{}l@{}}{\rule{0pt}{2.6ex}\textbf{T10.328}\enspace \texttt{0123444421}\enspace $n{=}10$, $|H|{=}6$, len 5, br 3, $|\mathrm{Aut}|{=}24$, orb 7/6, hist-orb 3/2; M 25600, res 14772, rev 6614: ok}\\*
0 & -- & 0 & 6 & 1 & 354 & 2 & 668 & \texttt{-?????????}\\*
1 & 0 & 1 & 5 & 1 & 392 & 8 & 624 & \texttt{--???????-}\\
2 & 1 & 2 & 4 & 1 & 544 & 32 & 448 & \texttt{---+????--}\\
3 & 2 & 3 & 4 & 1 & 64 & 64 & 896 & \texttt{----????--}\\
4 & 3 & 4 & 1 & 4 & 0 & 1024 & 0 & \texttt{----------}\\
5 & 3 & 4 & 1 & 4 & 0 & 1024 & 0 & \texttt{----------}\\
6 & 3 & 4 & 1 & 4 & 0 & 1024 & 0 & \texttt{----------}\\
7 & 3 & 4 & 1 & 4 & 0 & 1024 & 0 & \texttt{----------}\\
8 & 1 & 2 & 1 & 1 & 0 & 1024 & 0 & \texttt{----------}\\
9 & 0 & 1 & 1 & 1 & 0 & 1024 & 0 & \texttt{----------}\\
\multicolumn{9}{@{}l@{}}{\rule{0pt}{2.6ex}\textbf{T10.329}\enspace \texttt{0123444412}\enspace $n{=}10$, $|H|{=}5$, len 5, br 2, $|\mathrm{Aut}|{=}24$, orb 7/6, hist-orb 2/1; M 25600, res 9032, rev 9858: ok}\\*
0 & -- & 0 & 5 & 1 & 678 & 2 & 344 & \texttt{-?????????}\\*
1 & 0 & 1 & 4 & 1 & 784 & 16 & 224 & \texttt{--++????--}\\
2 & 1 & 2 & 4 & 1 & 544 & 32 & 448 & \texttt{---+????--}\\
3 & 2 & 3 & 4 & 1 & 64 & 64 & 896 & \texttt{----????--}\\
4 & 3 & 4 & 1 & 4 & 0 & 1024 & 0 & \texttt{----------}\\
5 & 3 & 4 & 1 & 4 & 0 & 1024 & 0 & \texttt{----------}\\
6 & 3 & 4 & 1 & 4 & 0 & 1024 & 0 & \texttt{----------}\\
7 & 3 & 4 & 1 & 4 & 0 & 1024 & 0 & \texttt{----------}\\
8 & 0 & 1 & 1 & 1 & 512 & 512 & 0 & \texttt{---------+}\\
9 & 8 & 2 & 1 & 1 & 0 & 1024 & 0 & \texttt{----------}\\
\multicolumn{9}{@{}l@{}}{\rule{0pt}{2.6ex}\textbf{T10.330}\enspace \texttt{0123444411}\enspace $n{=}10$, $|H|{=}6$, len 5, br 2, $|\mathrm{Aut}|{=}48$, orb 6/4, hist-orb 2/0; M 24576, res 13100, rev 7042: ok}\\*
0 & -- & 0 & 6 & 1 & 226 & 2 & 796 & \texttt{-?????????}\\*
1 & 0 & 1 & 4 & 1 & 784 & 16 & 224 & \texttt{--++????--}\\
2 & 1 & 2 & 4 & 1 & 544 & 32 & 448 & \texttt{---+????--}\\
3 & 2 & 3 & 4 & 1 & 64 & 64 & 896 & \texttt{----????--}\\
4 & 3 & 4 & 1 & 4 & 0 & 1024 & 0 & \texttt{----------}\\
5 & 3 & 4 & 1 & 4 & 0 & 1024 & 0 & \texttt{----------}\\
6 & 3 & 4 & 1 & 4 & 0 & 1024 & 0 & \texttt{----------}\\
7 & 3 & 4 & 1 & 4 & 0 & 1024 & 0 & \texttt{----------}\\
8 & 0 & 1 & 1 & 2 & 0 & 1024 & 0 & \texttt{----------}\\
9 & 0 & 1 & 1 & 2 & 0 & 1024 & 0 & \texttt{----------}\\
\multicolumn{9}{@{}l@{}}{\rule{0pt}{2.6ex}\textbf{T10.331}\enspace \texttt{0123444344}\enspace $n{=}10$, $|H|{=}5$, len 5, br 3, $|\mathrm{Aut}|{=}12$, orb 7/5, hist-orb 2/0; M 29696, res 10244, rev 11134: ok}\\*
0 & -- & 0 & 5 & 1 & 858 & 2 & 164 & \texttt{-++???????}\\*
1 & 0 & 1 & 5 & 1 & 692 & 4 & 328 & \texttt{--+???????}\\
2 & 1 & 2 & 5 & 1 & 360 & 8 & 656 & \texttt{---???????}\\
3 & 2 & 3 & 3 & 1 & 128 & 128 & 768 & \texttt{----???---}\\
4 & 3 & 4 & 1 & 3 & 0 & 1024 & 0 & \texttt{----------}\\
5 & 3 & 4 & 1 & 3 & 0 & 1024 & 0 & \texttt{----------}\\
6 & 3 & 4 & 1 & 3 & 0 & 1024 & 0 & \texttt{----------}\\
7 & 2 & 3 & 2 & 1 & 256 & 256 & 512 & \texttt{--------??}\\
8 & 7 & 4 & 1 & 2 & 0 & 1024 & 0 & \texttt{----------}\\
9 & 7 & 4 & 1 & 2 & 0 & 1024 & 0 & \texttt{----------}\\
\multicolumn{9}{@{}l@{}}{\rule{0pt}{2.6ex}\textbf{T10.332}\enspace \texttt{0123444343}\enspace $n{=}10$, $|H|{=}5$, len 5, br 2, $|\mathrm{Aut}|{=}6$, orb 8/7, hist-orb 3/2; M 28672, res 11432, rev 10166: ok}\\*
0 & -- & 0 & 5 & 1 & 822 & 2 & 200 & \texttt{-++???????}\\*
1 & 0 & 1 & 5 & 1 & 620 & 4 & 400 & \texttt{--+???????}\\
2 & 1 & 2 & 5 & 1 & 216 & 8 & 800 & \texttt{---???????}\\
3 & 2 & 3 & 3 & 1 & 128 & 128 & 768 & \texttt{----???---}\\
4 & 3 & 4 & 1 & 3 & 0 & 1024 & 0 & \texttt{----------}\\
5 & 3 & 4 & 1 & 3 & 0 & 1024 & 0 & \texttt{----------}\\
6 & 3 & 4 & 1 & 3 & 0 & 1024 & 0 & \texttt{----------}\\
7 & 2 & 3 & 1 & 1 & 512 & 512 & 0 & \texttt{--------+-}\\
8 & 7 & 4 & 1 & 1 & 0 & 1024 & 0 & \texttt{----------}\\
9 & 2 & 3 & 1 & 1 & 0 & 1024 & 0 & \texttt{----------}\\
\multicolumn{9}{@{}l@{}}{\rule{0pt}{2.6ex}\textbf{T10.333}\enspace \texttt{0123444342}\enspace $n{=}10$, $|H|{=}5$, len 5, br 3, $|\mathrm{Aut}|{=}6$, orb 8/7, hist-orb 3/2; M 27648, res 12528, rev 8754: ok}\\*
0 & -- & 0 & 5 & 1 & 694 & 2 & 328 & \texttt{-+????????}\\*
1 & 0 & 1 & 5 & 1 & 364 & 4 & 656 & \texttt{--????????}\\
2 & 1 & 2 & 4 & 1 & 432 & 16 & 576 & \texttt{---??????-}\\
3 & 2 & 3 & 3 & 1 & 128 & 128 & 768 & \texttt{----???---}\\
4 & 3 & 4 & 1 & 3 & 0 & 1024 & 0 & \texttt{----------}\\
5 & 3 & 4 & 1 & 3 & 0 & 1024 & 0 & \texttt{----------}\\
6 & 3 & 4 & 1 & 3 & 0 & 1024 & 0 & \texttt{----------}\\
7 & 2 & 3 & 1 & 1 & 512 & 512 & 0 & \texttt{--------+-}\\
8 & 7 & 4 & 1 & 1 & 0 & 1024 & 0 & \texttt{----------}\\
9 & 1 & 2 & 1 & 1 & 0 & 1024 & 0 & \texttt{----------}\\
\multicolumn{9}{@{}l@{}}{\rule{0pt}{2.6ex}\textbf{T10.334}\enspace \texttt{0123444341}\enspace $n{=}10$, $|H|{=}5$, len 5, br 3, $|\mathrm{Aut}|{=}6$, orb 8/7, hist-orb 3/2; M 26624, res 11784, rev 8614: ok}\\*
0 & -- & 0 & 5 & 1 & 438 & 2 & 584 & \texttt{-?????????}\\*
1 & 0 & 1 & 4 & 1 & 728 & 8 & 288 & \texttt{--+??????-}\\
2 & 1 & 2 & 4 & 1 & 432 & 16 & 576 & \texttt{---??????-}\\
3 & 2 & 3 & 3 & 1 & 128 & 128 & 768 & \texttt{----???---}\\
4 & 3 & 4 & 1 & 3 & 0 & 1024 & 0 & \texttt{----------}\\
5 & 3 & 4 & 1 & 3 & 0 & 1024 & 0 & \texttt{----------}\\
6 & 3 & 4 & 1 & 3 & 0 & 1024 & 0 & \texttt{----------}\\
7 & 2 & 3 & 1 & 1 & 512 & 512 & 0 & \texttt{--------+-}\\
8 & 7 & 4 & 1 & 1 & 0 & 1024 & 0 & \texttt{----------}\\
9 & 0 & 1 & 1 & 1 & 0 & 1024 & 0 & \texttt{----------}\\
\multicolumn{9}{@{}l@{}}{\rule{0pt}{2.6ex}\textbf{T10.335}\enspace \texttt{0123444333}\enspace $n{=}10$, $|H|{=}6$, len 5, br 2, $|\mathrm{Aut}|{=}36$, orb 6/4, hist-orb 2/0; M 27648, res 13196, rev 9006: ok}\\*
0 & -- & 0 & 6 & 1 & 786 & 2 & 236 & \texttt{-++???????}\\*
1 & 0 & 1 & 6 & 1 & 548 & 4 & 472 & \texttt{--+???????}\\
2 & 1 & 2 & 6 & 1 & 72 & 8 & 944 & \texttt{---???????}\\
3 & 2 & 3 & 3 & 1 & 128 & 128 & 768 & \texttt{----???---}\\
4 & 3 & 4 & 1 & 3 & 0 & 1024 & 0 & \texttt{----------}\\
5 & 3 & 4 & 1 & 3 & 0 & 1024 & 0 & \texttt{----------}\\
6 & 3 & 4 & 1 & 3 & 0 & 1024 & 0 & \texttt{----------}\\
7 & 2 & 3 & 1 & 3 & 0 & 1024 & 0 & \texttt{----------}\\
8 & 2 & 3 & 1 & 3 & 0 & 1024 & 0 & \texttt{----------}\\
9 & 2 & 3 & 1 & 3 & 0 & 1024 & 0 & \texttt{----------}\\
\multicolumn{9}{@{}l@{}}{\rule{0pt}{2.6ex}\textbf{T10.336}\enspace \texttt{0123444332}\enspace $n{=}10$, $|H|{=}6$, len 5, br 3, $|\mathrm{Aut}|{=}12$, orb 7/5, hist-orb 3/1; M 26624, res 14904, rev 6954: ok}\\*
0 & -- & 0 & 6 & 1 & 658 & 2 & 364 & \texttt{-+????????}\\*
1 & 0 & 1 & 6 & 1 & 292 & 4 & 728 & \texttt{--????????}\\
2 & 1 & 2 & 5 & 1 & 144 & 16 & 864 & \texttt{---??????-}\\
3 & 2 & 3 & 3 & 1 & 128 & 128 & 768 & \texttt{----???---}\\
4 & 3 & 4 & 1 & 3 & 0 & 1024 & 0 & \texttt{----------}\\
5 & 3 & 4 & 1 & 3 & 0 & 1024 & 0 & \texttt{----------}\\
6 & 3 & 4 & 1 & 3 & 0 & 1024 & 0 & \texttt{----------}\\
7 & 2 & 3 & 1 & 2 & 0 & 1024 & 0 & \texttt{----------}\\
8 & 2 & 3 & 1 & 2 & 0 & 1024 & 0 & \texttt{----------}\\
9 & 1 & 2 & 1 & 1 & 0 & 1024 & 0 & \texttt{----------}\\
\multicolumn{9}{@{}l@{}}{\rule{0pt}{2.6ex}\textbf{T10.337}\enspace \texttt{0123444331}\enspace $n{=}10$, $|H|{=}6$, len 5, br 3, $|\mathrm{Aut}|{=}12$, orb 7/5, hist-orb 3/1; M 25600, res 14556, rev 6590: ok}\\*
0 & -- & 0 & 6 & 1 & 402 & 2 & 620 & \texttt{-?????????}\\*
1 & 0 & 1 & 5 & 1 & 584 & 8 & 432 & \texttt{--+??????-}\\
2 & 1 & 2 & 5 & 1 & 144 & 16 & 864 & \texttt{---??????-}\\
3 & 2 & 3 & 3 & 1 & 128 & 128 & 768 & \texttt{----???---}\\
4 & 3 & 4 & 1 & 3 & 0 & 1024 & 0 & \texttt{----------}\\
5 & 3 & 4 & 1 & 3 & 0 & 1024 & 0 & \texttt{----------}\\
6 & 3 & 4 & 1 & 3 & 0 & 1024 & 0 & \texttt{----------}\\
7 & 2 & 3 & 1 & 2 & 0 & 1024 & 0 & \texttt{----------}\\
8 & 2 & 3 & 1 & 2 & 0 & 1024 & 0 & \texttt{----------}\\
9 & 0 & 1 & 1 & 1 & 0 & 1024 & 0 & \texttt{----------}\\
\multicolumn{9}{@{}l@{}}{\rule{0pt}{2.6ex}\textbf{T10.338}\enspace \texttt{0123444323}\enspace $n{=}10$, $|H|{=}5$, len 5, br 3, $|\mathrm{Aut}|{=}6$, orb 8/7, hist-orb 3/2; M 26624, res 10960, rev 9022: ok}\\*
0 & -- & 0 & 5 & 1 & 758 & 2 & 264 & \texttt{-+????????}\\*
1 & 0 & 1 & 5 & 1 & 492 & 4 & 528 & \texttt{--????????}\\
2 & 1 & 2 & 4 & 1 & 288 & 32 & 704 & \texttt{---?????--}\\
3 & 2 & 3 & 3 & 1 & 128 & 128 & 768 & \texttt{----???---}\\
4 & 3 & 4 & 1 & 3 & 0 & 1024 & 0 & \texttt{----------}\\
5 & 3 & 4 & 1 & 3 & 0 & 1024 & 0 & \texttt{----------}\\
6 & 3 & 4 & 1 & 3 & 0 & 1024 & 0 & \texttt{----------}\\
7 & 2 & 3 & 1 & 1 & 0 & 1024 & 0 & \texttt{----------}\\
8 & 1 & 2 & 1 & 1 & 512 & 512 & 0 & \texttt{---------+}\\
9 & 8 & 3 & 1 & 1 & 0 & 1024 & 0 & \texttt{----------}\\
\multicolumn{9}{@{}l@{}}{\rule{0pt}{2.6ex}\textbf{T10.339}\enspace \texttt{0123444322}\enspace $n{=}10$, $|H|{=}6$, len 5, br 3, $|\mathrm{Aut}|{=}12$, orb 7/5, hist-orb 3/1; M 25600, res 14896, rev 6234: ok}\\*
0 & -- & 0 & 6 & 1 & 594 & 2 & 428 & \texttt{-+????????}\\*
1 & 0 & 1 & 6 & 1 & 164 & 4 & 856 & \texttt{--????????}\\
2 & 1 & 2 & 4 & 1 & 288 & 32 & 704 & \texttt{---?????--}\\
3 & 2 & 3 & 3 & 1 & 128 & 128 & 768 & \texttt{----???---}\\
4 & 3 & 4 & 1 & 3 & 0 & 1024 & 0 & \texttt{----------}\\
5 & 3 & 4 & 1 & 3 & 0 & 1024 & 0 & \texttt{----------}\\
6 & 3 & 4 & 1 & 3 & 0 & 1024 & 0 & \texttt{----------}\\
7 & 2 & 3 & 1 & 1 & 0 & 1024 & 0 & \texttt{----------}\\
8 & 1 & 2 & 1 & 2 & 0 & 1024 & 0 & \texttt{----------}\\
9 & 1 & 2 & 1 & 2 & 0 & 1024 & 0 & \texttt{----------}\\
\multicolumn{9}{@{}l@{}}{\rule{0pt}{2.6ex}\textbf{T10.340}\enspace \texttt{0123444321}\enspace $n{=}10$, $|H|{=}6$, len 5, br 4, $|\mathrm{Aut}|{=}6$, orb 8/7, hist-orb 4/3; M 24576, res 15252, rev 5338: ok}\\*
0 & -- & 0 & 6 & 1 & 338 & 2 & 684 & \texttt{-?????????}\\*
1 & 0 & 1 & 5 & 1 & 328 & 8 & 688 & \texttt{--???????-}\\
2 & 1 & 2 & 4 & 1 & 288 & 32 & 704 & \texttt{---?????--}\\
3 & 2 & 3 & 3 & 1 & 128 & 128 & 768 & \texttt{----???---}\\
4 & 3 & 4 & 1 & 3 & 0 & 1024 & 0 & \texttt{----------}\\
5 & 3 & 4 & 1 & 3 & 0 & 1024 & 0 & \texttt{----------}\\
6 & 3 & 4 & 1 & 3 & 0 & 1024 & 0 & \texttt{----------}\\
7 & 2 & 3 & 1 & 1 & 0 & 1024 & 0 & \texttt{----------}\\
8 & 1 & 2 & 1 & 1 & 0 & 1024 & 0 & \texttt{----------}\\
9 & 0 & 1 & 1 & 1 & 0 & 1024 & 0 & \texttt{----------}\\
\multicolumn{9}{@{}l@{}}{\rule{0pt}{2.6ex}\textbf{T10.341}\enspace \texttt{0123444312}\enspace $n{=}10$, $|H|{=}5$, len 5, br 3, $|\mathrm{Aut}|{=}6$, orb 8/7, hist-orb 3/2; M 24576, res 10088, rev 8262: ok}\\*
0 & -- & 0 & 5 & 1 & 630 & 2 & 392 & \texttt{-?????????}\\*
1 & 0 & 1 & 4 & 1 & 656 & 16 & 352 & \texttt{--+?????--}\\
2 & 1 & 2 & 4 & 1 & 288 & 32 & 704 & \texttt{---?????--}\\
3 & 2 & 3 & 3 & 1 & 128 & 128 & 768 & \texttt{----???---}\\
4 & 3 & 4 & 1 & 3 & 0 & 1024 & 0 & \texttt{----------}\\
5 & 3 & 4 & 1 & 3 & 0 & 1024 & 0 & \texttt{----------}\\
6 & 3 & 4 & 1 & 3 & 0 & 1024 & 0 & \texttt{----------}\\
7 & 2 & 3 & 1 & 1 & 0 & 1024 & 0 & \texttt{----------}\\
8 & 0 & 1 & 1 & 1 & 512 & 512 & 0 & \texttt{---------+}\\
9 & 8 & 2 & 1 & 1 & 0 & 1024 & 0 & \texttt{----------}\\
\multicolumn{9}{@{}l@{}}{\rule{0pt}{2.6ex}\textbf{T10.342}\enspace \texttt{0123444311}\enspace $n{=}10$, $|H|{=}6$, len 5, br 3, $|\mathrm{Aut}|{=}12$, orb 7/5, hist-orb 3/1; M 23552, res 13868, rev 5582: ok}\\*
0 & -- & 0 & 6 & 1 & 210 & 2 & 812 & \texttt{-?????????}\\*
1 & 0 & 1 & 4 & 1 & 656 & 16 & 352 & \texttt{--+?????--}\\
2 & 1 & 2 & 4 & 1 & 288 & 32 & 704 & \texttt{---?????--}\\
3 & 2 & 3 & 3 & 1 & 128 & 128 & 768 & \texttt{----???---}\\
4 & 3 & 4 & 1 & 3 & 0 & 1024 & 0 & \texttt{----------}\\
5 & 3 & 4 & 1 & 3 & 0 & 1024 & 0 & \texttt{----------}\\
6 & 3 & 4 & 1 & 3 & 0 & 1024 & 0 & \texttt{----------}\\
7 & 2 & 3 & 1 & 1 & 0 & 1024 & 0 & \texttt{----------}\\
8 & 0 & 1 & 1 & 2 & 0 & 1024 & 0 & \texttt{----------}\\
9 & 0 & 1 & 1 & 2 & 0 & 1024 & 0 & \texttt{----------}\\
\multicolumn{9}{@{}l@{}}{\rule{0pt}{2.6ex}\textbf{T10.343}\enspace \texttt{0123444234}\enspace $n{=}10$, $|H|{=}4$, len 5, br 2, $|\mathrm{Aut}|{=}6$, orb 8/7, hist-orb 2/1; M 27648, res 6720, rev 11210: ok}\\*
0 & -- & 0 & 4 & 1 & 862 & 2 & 160 & \texttt{-+????????}\\*
1 & 0 & 1 & 4 & 1 & 700 & 4 & 320 & \texttt{--????????}\\
2 & 1 & 2 & 3 & 1 & 576 & 64 & 384 & \texttt{---+???---}\\
3 & 2 & 3 & 3 & 1 & 128 & 128 & 768 & \texttt{----???---}\\
4 & 3 & 4 & 1 & 3 & 0 & 1024 & 0 & \texttt{----------}\\
5 & 3 & 4 & 1 & 3 & 0 & 1024 & 0 & \texttt{----------}\\
6 & 3 & 4 & 1 & 3 & 0 & 1024 & 0 & \texttt{----------}\\
7 & 1 & 2 & 1 & 1 & 768 & 256 & 0 & \texttt{--------++}\\
8 & 7 & 3 & 1 & 1 & 512 & 512 & 0 & \texttt{---------+}\\
9 & 8 & 4 & 1 & 1 & 0 & 1024 & 0 & \texttt{----------}\\
\multicolumn{9}{@{}l@{}}{\rule{0pt}{2.6ex}\textbf{T10.344}\enspace \texttt{0123444233}\enspace $n{=}10$, $|H|{=}5$, len 5, br 3, $|\mathrm{Aut}|{=}12$, orb 7/5, hist-orb 2/0; M 26624, res 9760, rev 9494: ok}\\*
0 & -- & 0 & 5 & 1 & 762 & 2 & 260 & \texttt{-+????????}\\*
1 & 0 & 1 & 5 & 1 & 500 & 4 & 520 & \texttt{--????????}\\
2 & 1 & 2 & 3 & 1 & 576 & 64 & 384 & \texttt{---+???---}\\
3 & 2 & 3 & 3 & 1 & 128 & 128 & 768 & \texttt{----???---}\\
4 & 3 & 4 & 1 & 3 & 0 & 1024 & 0 & \texttt{----------}\\
5 & 3 & 4 & 1 & 3 & 0 & 1024 & 0 & \texttt{----------}\\
6 & 3 & 4 & 1 & 3 & 0 & 1024 & 0 & \texttt{----------}\\
7 & 1 & 2 & 2 & 1 & 256 & 256 & 512 & \texttt{--------??}\\
8 & 7 & 3 & 1 & 2 & 0 & 1024 & 0 & \texttt{----------}\\
9 & 7 & 3 & 1 & 2 & 0 & 1024 & 0 & \texttt{----------}\\
\multicolumn{9}{@{}l@{}}{\rule{0pt}{2.6ex}\textbf{T10.345}\enspace \texttt{0123444232}\enspace $n{=}10$, $|H|{=}5$, len 5, br 2, $|\mathrm{Aut}|{=}6$, orb 8/7, hist-orb 3/2; M 25600, res 11520, rev 8074: ok}\\*
0 & -- & 0 & 5 & 1 & 662 & 2 & 360 & \texttt{-+????????}\\*
1 & 0 & 1 & 5 & 1 & 300 & 4 & 720 & \texttt{--????????}\\
2 & 1 & 2 & 3 & 1 & 576 & 64 & 384 & \texttt{---+???---}\\
3 & 2 & 3 & 3 & 1 & 128 & 128 & 768 & \texttt{----???---}\\
4 & 3 & 4 & 1 & 3 & 0 & 1024 & 0 & \texttt{----------}\\
5 & 3 & 4 & 1 & 3 & 0 & 1024 & 0 & \texttt{----------}\\
6 & 3 & 4 & 1 & 3 & 0 & 1024 & 0 & \texttt{----------}\\
7 & 1 & 2 & 1 & 1 & 512 & 512 & 0 & \texttt{--------+-}\\
8 & 7 & 3 & 1 & 1 & 0 & 1024 & 0 & \texttt{----------}\\
9 & 1 & 2 & 1 & 1 & 0 & 1024 & 0 & \texttt{----------}\\
\multicolumn{9}{@{}l@{}}{\rule{0pt}{2.6ex}\textbf{T10.346}\enspace \texttt{0123444231}\enspace $n{=}10$, $|H|{=}5$, len 5, br 3, $|\mathrm{Aut}|{=}6$, orb 8/7, hist-orb 3/2; M 24576, res 11128, rev 7666: ok}\\*
0 & -- & 0 & 5 & 1 & 406 & 2 & 616 & \texttt{-?????????}\\*
1 & 0 & 1 & 4 & 1 & 600 & 8 & 416 & \texttt{--???????-}\\
2 & 1 & 2 & 3 & 1 & 576 & 64 & 384 & \texttt{---+???---}\\
3 & 2 & 3 & 3 & 1 & 128 & 128 & 768 & \texttt{----???---}\\
4 & 3 & 4 & 1 & 3 & 0 & 1024 & 0 & \texttt{----------}\\
5 & 3 & 4 & 1 & 3 & 0 & 1024 & 0 & \texttt{----------}\\
6 & 3 & 4 & 1 & 3 & 0 & 1024 & 0 & \texttt{----------}\\
7 & 1 & 2 & 1 & 1 & 512 & 512 & 0 & \texttt{--------+-}\\
8 & 7 & 3 & 1 & 1 & 0 & 1024 & 0 & \texttt{----------}\\
9 & 0 & 1 & 1 & 1 & 0 & 1024 & 0 & \texttt{----------}\\
\multicolumn{9}{@{}l@{}}{\rule{0pt}{2.6ex}\textbf{T10.347}\enspace \texttt{0123444222}\enspace $n{=}10$, $|H|{=}6$, len 5, br 2, $|\mathrm{Aut}|{=}36$, orb 6/4, hist-orb 2/0; M 24576, res 13920, rev 6270: ok}\\*
0 & -- & 0 & 6 & 1 & 562 & 2 & 460 & \texttt{-+????????}\\*
1 & 0 & 1 & 6 & 1 & 100 & 4 & 920 & \texttt{--????????}\\
2 & 1 & 2 & 3 & 1 & 576 & 64 & 384 & \texttt{---+???---}\\
3 & 2 & 3 & 3 & 1 & 128 & 128 & 768 & \texttt{----???---}\\
4 & 3 & 4 & 1 & 3 & 0 & 1024 & 0 & \texttt{----------}\\
5 & 3 & 4 & 1 & 3 & 0 & 1024 & 0 & \texttt{----------}\\
6 & 3 & 4 & 1 & 3 & 0 & 1024 & 0 & \texttt{----------}\\
7 & 1 & 2 & 1 & 3 & 0 & 1024 & 0 & \texttt{----------}\\
8 & 1 & 2 & 1 & 3 & 0 & 1024 & 0 & \texttt{----------}\\
9 & 1 & 2 & 1 & 3 & 0 & 1024 & 0 & \texttt{----------}\\
\multicolumn{9}{@{}l@{}}{\rule{0pt}{2.6ex}\textbf{T10.348}\enspace \texttt{0123444221}\enspace $n{=}10$, $|H|{=}6$, len 5, br 3, $|\mathrm{Aut}|{=}12$, orb 7/5, hist-orb 3/1; M 23552, res 14628, rev 5102: ok}\\*
0 & -- & 0 & 6 & 1 & 306 & 2 & 716 & \texttt{-?????????}\\*
1 & 0 & 1 & 5 & 1 & 200 & 8 & 816 & \texttt{--???????-}\\
2 & 1 & 2 & 3 & 1 & 576 & 64 & 384 & \texttt{---+???---}\\
3 & 2 & 3 & 3 & 1 & 128 & 128 & 768 & \texttt{----???---}\\
4 & 3 & 4 & 1 & 3 & 0 & 1024 & 0 & \texttt{----------}\\
5 & 3 & 4 & 1 & 3 & 0 & 1024 & 0 & \texttt{----------}\\
6 & 3 & 4 & 1 & 3 & 0 & 1024 & 0 & \texttt{----------}\\
7 & 1 & 2 & 1 & 2 & 0 & 1024 & 0 & \texttt{----------}\\
8 & 1 & 2 & 1 & 2 & 0 & 1024 & 0 & \texttt{----------}\\
9 & 0 & 1 & 1 & 1 & 0 & 1024 & 0 & \texttt{----------}\\
\multicolumn{9}{@{}l@{}}{\rule{0pt}{2.6ex}\textbf{T10.349}\enspace \texttt{0123444212}\enspace $n{=}10$, $|H|{=}5$, len 5, br 3, $|\mathrm{Aut}|{=}6$, orb 8/7, hist-orb 3/2; M 23552, res 10616, rev 7302: ok}\\*
0 & -- & 0 & 5 & 1 & 534 & 2 & 488 & \texttt{-?????????}\\*
1 & 0 & 1 & 4 & 1 & 400 & 16 & 608 & \texttt{--??????--}\\
2 & 1 & 2 & 3 & 1 & 576 & 64 & 384 & \texttt{---+???---}\\
3 & 2 & 3 & 3 & 1 & 128 & 128 & 768 & \texttt{----???---}\\
4 & 3 & 4 & 1 & 3 & 0 & 1024 & 0 & \texttt{----------}\\
5 & 3 & 4 & 1 & 3 & 0 & 1024 & 0 & \texttt{----------}\\
6 & 3 & 4 & 1 & 3 & 0 & 1024 & 0 & \texttt{----------}\\
7 & 1 & 2 & 1 & 1 & 0 & 1024 & 0 & \texttt{----------}\\
8 & 0 & 1 & 1 & 1 & 512 & 512 & 0 & \texttt{---------+}\\
9 & 8 & 2 & 1 & 1 & 0 & 1024 & 0 & \texttt{----------}\\
\multicolumn{9}{@{}l@{}}{\rule{0pt}{2.6ex}\textbf{T10.350}\enspace \texttt{0123444211}\enspace $n{=}10$, $|H|{=}6$, len 5, br 3, $|\mathrm{Aut}|{=}12$, orb 7/5, hist-orb 3/1; M 22528, res 13820, rev 4942: ok}\\*
0 & -- & 0 & 6 & 1 & 178 & 2 & 844 & \texttt{-?????????}\\*
1 & 0 & 1 & 4 & 1 & 400 & 16 & 608 & \texttt{--??????--}\\
2 & 1 & 2 & 3 & 1 & 576 & 64 & 384 & \texttt{---+???---}\\
3 & 2 & 3 & 3 & 1 & 128 & 128 & 768 & \texttt{----???---}\\
4 & 3 & 4 & 1 & 3 & 0 & 1024 & 0 & \texttt{----------}\\
5 & 3 & 4 & 1 & 3 & 0 & 1024 & 0 & \texttt{----------}\\
6 & 3 & 4 & 1 & 3 & 0 & 1024 & 0 & \texttt{----------}\\
7 & 1 & 2 & 1 & 1 & 0 & 1024 & 0 & \texttt{----------}\\
8 & 0 & 1 & 1 & 2 & 0 & 1024 & 0 & \texttt{----------}\\
9 & 0 & 1 & 1 & 2 & 0 & 1024 & 0 & \texttt{----------}\\
\multicolumn{9}{@{}l@{}}{\rule{0pt}{2.6ex}\textbf{T10.351}\enspace \texttt{0123444123}\enspace $n{=}10$, $|H|{=}4$, len 5, br 2, $|\mathrm{Aut}|{=}6$, orb 8/7, hist-orb 2/1; M 24576, res 5760, rev 10042: ok}\\*
0 & -- & 0 & 4 & 1 & 798 & 2 & 224 & \texttt{-?????????}\\*
1 & 0 & 1 & 3 & 1 & 800 & 32 & 192 & \texttt{--++???---}\\
2 & 1 & 2 & 3 & 1 & 576 & 64 & 384 & \texttt{---+???---}\\
3 & 2 & 3 & 3 & 1 & 128 & 128 & 768 & \texttt{----???---}\\
4 & 3 & 4 & 1 & 3 & 0 & 1024 & 0 & \texttt{----------}\\
5 & 3 & 4 & 1 & 3 & 0 & 1024 & 0 & \texttt{----------}\\
6 & 3 & 4 & 1 & 3 & 0 & 1024 & 0 & \texttt{----------}\\
7 & 0 & 1 & 1 & 1 & 768 & 256 & 0 & \texttt{--------++}\\
8 & 7 & 2 & 1 & 1 & 512 & 512 & 0 & \texttt{---------+}\\
9 & 8 & 3 & 1 & 1 & 0 & 1024 & 0 & \texttt{----------}\\
\multicolumn{9}{@{}l@{}}{\rule{0pt}{2.6ex}\textbf{T10.352}\enspace \texttt{0123444122}\enspace $n{=}10$, $|H|{=}5$, len 5, br 3, $|\mathrm{Aut}|{=}12$, orb 7/5, hist-orb 2/0; M 23552, res 8580, rev 8290: ok}\\*
0 & -- & 0 & 5 & 1 & 570 & 2 & 452 & \texttt{-?????????}\\*
1 & 0 & 1 & 3 & 1 & 800 & 32 & 192 & \texttt{--++???---}\\
2 & 1 & 2 & 3 & 1 & 576 & 64 & 384 & \texttt{---+???---}\\
3 & 2 & 3 & 3 & 1 & 128 & 128 & 768 & \texttt{----???---}\\
4 & 3 & 4 & 1 & 3 & 0 & 1024 & 0 & \texttt{----------}\\
5 & 3 & 4 & 1 & 3 & 0 & 1024 & 0 & \texttt{----------}\\
6 & 3 & 4 & 1 & 3 & 0 & 1024 & 0 & \texttt{----------}\\
7 & 0 & 1 & 2 & 1 & 256 & 256 & 512 & \texttt{--------??}\\
8 & 7 & 2 & 1 & 2 & 0 & 1024 & 0 & \texttt{----------}\\
9 & 7 & 2 & 1 & 2 & 0 & 1024 & 0 & \texttt{----------}\\
\multicolumn{9}{@{}l@{}}{\rule{0pt}{2.6ex}\textbf{T10.353}\enspace \texttt{0123444121}\enspace $n{=}10$, $|H|{=}5$, len 5, br 2, $|\mathrm{Aut}|{=}6$, orb 8/7, hist-orb 3/2; M 22528, res 9864, rev 7106: ok}\\*
0 & -- & 0 & 5 & 1 & 342 & 2 & 680 & \texttt{-?????????}\\*
1 & 0 & 1 & 3 & 1 & 800 & 32 & 192 & \texttt{--++???---}\\
2 & 1 & 2 & 3 & 1 & 576 & 64 & 384 & \texttt{---+???---}\\
3 & 2 & 3 & 3 & 1 & 128 & 128 & 768 & \texttt{----???---}\\
4 & 3 & 4 & 1 & 3 & 0 & 1024 & 0 & \texttt{----------}\\
5 & 3 & 4 & 1 & 3 & 0 & 1024 & 0 & \texttt{----------}\\
6 & 3 & 4 & 1 & 3 & 0 & 1024 & 0 & \texttt{----------}\\
7 & 0 & 1 & 1 & 1 & 512 & 512 & 0 & \texttt{--------+-}\\
8 & 7 & 2 & 1 & 1 & 0 & 1024 & 0 & \texttt{----------}\\
9 & 0 & 1 & 1 & 1 & 0 & 1024 & 0 & \texttt{----------}\\
\multicolumn{9}{@{}l@{}}{\rule{0pt}{2.6ex}\textbf{T10.354}\enspace \texttt{0123444111}\enspace $n{=}10$, $|H|{=}6$, len 5, br 2, $|\mathrm{Aut}|{=}36$, orb 6/4, hist-orb 2/0; M 21504, res 11916, rev 5410: ok}\\*
0 & -- & 0 & 6 & 1 & 114 & 2 & 908 & \texttt{-?????????}\\*
1 & 0 & 1 & 3 & 1 & 800 & 32 & 192 & \texttt{--++???---}\\
2 & 1 & 2 & 3 & 1 & 576 & 64 & 384 & \texttt{---+???---}\\
3 & 2 & 3 & 3 & 1 & 128 & 128 & 768 & \texttt{----???---}\\
4 & 3 & 4 & 1 & 3 & 0 & 1024 & 0 & \texttt{----------}\\
5 & 3 & 4 & 1 & 3 & 0 & 1024 & 0 & \texttt{----------}\\
6 & 3 & 4 & 1 & 3 & 0 & 1024 & 0 & \texttt{----------}\\
7 & 0 & 1 & 1 & 3 & 0 & 1024 & 0 & \texttt{----------}\\
8 & 0 & 1 & 1 & 3 & 0 & 1024 & 0 & \texttt{----------}\\
9 & 0 & 1 & 1 & 3 & 0 & 1024 & 0 & \texttt{----------}\\
\multicolumn{9}{@{}l@{}}{\rule{0pt}{2.6ex}\textbf{T10.355}\enspace \texttt{0123443443}\enspace $n{=}10$, $|H|{=}5$, len 5, br 3, $|\mathrm{Aut}|{=}8$, orb 6/4, hist-orb 2/1; M 28672, res 11148, rev 9966: ok}\\*
0 & -- & 0 & 5 & 1 & 818 & 2 & 204 & \texttt{-++???????}\\*
1 & 0 & 1 & 5 & 1 & 612 & 4 & 408 & \texttt{--+???????}\\
2 & 1 & 2 & 5 & 1 & 200 & 8 & 816 & \texttt{---???????}\\
3 & 2 & 3 & 2 & 2 & 256 & 256 & 512 & \texttt{----??----}\\
4 & 3 & 4 & 1 & 4 & 0 & 1024 & 0 & \texttt{----------}\\
5 & 3 & 4 & 1 & 4 & 0 & 1024 & 0 & \texttt{----------}\\
6 & 2 & 3 & 2 & 2 & 256 & 256 & 512 & \texttt{-------??-}\\
7 & 6 & 4 & 1 & 4 & 0 & 1024 & 0 & \texttt{----------}\\
8 & 6 & 4 & 1 & 4 & 0 & 1024 & 0 & \texttt{----------}\\
9 & 2 & 3 & 1 & 1 & 0 & 1024 & 0 & \texttt{----------}\\
\multicolumn{9}{@{}l@{}}{\rule{0pt}{2.6ex}\textbf{T10.356}\enspace \texttt{0123443442}\enspace $n{=}10$, $|H|{=}5$, len 5, br 4, $|\mathrm{Aut}|{=}8$, orb 6/4, hist-orb 2/1; M 27648, res 12312, rev 8538: ok}\\*
0 & -- & 0 & 5 & 1 & 690 & 2 & 332 & \texttt{-+????????}\\*
1 & 0 & 1 & 5 & 1 & 356 & 4 & 664 & \texttt{--????????}\\
2 & 1 & 2 & 4 & 1 & 400 & 16 & 608 & \texttt{---??????-}\\
3 & 2 & 3 & 2 & 2 & 256 & 256 & 512 & \texttt{----??----}\\
4 & 3 & 4 & 1 & 4 & 0 & 1024 & 0 & \texttt{----------}\\
5 & 3 & 4 & 1 & 4 & 0 & 1024 & 0 & \texttt{----------}\\
6 & 2 & 3 & 2 & 2 & 256 & 256 & 512 & \texttt{-------??-}\\
7 & 6 & 4 & 1 & 4 & 0 & 1024 & 0 & \texttt{----------}\\
8 & 6 & 4 & 1 & 4 & 0 & 1024 & 0 & \texttt{----------}\\
9 & 1 & 2 & 1 & 1 & 0 & 1024 & 0 & \texttt{----------}\\
\multicolumn{9}{@{}l@{}}{\rule{0pt}{2.6ex}\textbf{T10.357}\enspace \texttt{0123443441}\enspace $n{=}10$, $|H|{=}5$, len 5, br 4, $|\mathrm{Aut}|{=}8$, orb 6/4, hist-orb 2/1; M 26624, res 11612, rev 8366: ok}\\*
0 & -- & 0 & 5 & 1 & 434 & 2 & 588 & \texttt{-?????????}\\*
1 & 0 & 1 & 4 & 1 & 712 & 8 & 304 & \texttt{--+??????-}\\
2 & 1 & 2 & 4 & 1 & 400 & 16 & 608 & \texttt{---??????-}\\
3 & 2 & 3 & 2 & 2 & 256 & 256 & 512 & \texttt{----??----}\\
4 & 3 & 4 & 1 & 4 & 0 & 1024 & 0 & \texttt{----------}\\
5 & 3 & 4 & 1 & 4 & 0 & 1024 & 0 & \texttt{----------}\\
6 & 2 & 3 & 2 & 2 & 256 & 256 & 512 & \texttt{-------??-}\\
7 & 6 & 4 & 1 & 4 & 0 & 1024 & 0 & \texttt{----------}\\
8 & 6 & 4 & 1 & 4 & 0 & 1024 & 0 & \texttt{----------}\\
9 & 0 & 1 & 1 & 1 & 0 & 1024 & 0 & \texttt{----------}\\
\multicolumn{9}{@{}l@{}}{\rule{0pt}{2.6ex}\textbf{T10.358}\enspace \texttt{0123443434}\enspace $n{=}10$, $|H|{=}4$, len 5, br 2, $|\mathrm{Aut}|{=}4$, orb 7/4, hist-orb 2/0; M 28672, res 8612, rev 11334: ok}\\*
0 & -- & 0 & 4 & 1 & 858 & 2 & 164 & \texttt{-++???????}\\*
1 & 0 & 1 & 4 & 1 & 692 & 4 & 328 & \texttt{--+???????}\\
2 & 1 & 2 & 4 & 1 & 360 & 8 & 656 & \texttt{---???????}\\
3 & 2 & 3 & 2 & 1 & 256 & 256 & 512 & \texttt{----??----}\\
4 & 3 & 4 & 1 & 2 & 0 & 1024 & 0 & \texttt{----------}\\
5 & 3 & 4 & 1 & 2 & 0 & 1024 & 0 & \texttt{----------}\\
6 & 2 & 3 & 1 & 2 & 512 & 512 & 0 & \texttt{-------+--}\\
7 & 6 & 4 & 1 & 2 & 0 & 1024 & 0 & \texttt{----------}\\
8 & 2 & 3 & 1 & 2 & 512 & 512 & 0 & \texttt{---------+}\\
9 & 8 & 4 & 1 & 2 & 0 & 1024 & 0 & \texttt{----------}\\
\multicolumn{9}{@{}l@{}}{\rule{0pt}{2.6ex}\textbf{T10.359}\enspace \texttt{0123443433}\enspace $n{=}10$, $|H|{=}5$, len 5, br 2, $|\mathrm{Aut}|{=}4$, orb 8/6, hist-orb 3/1; M 27648, res 11552, rev 9558: ok}\\*
0 & -- & 0 & 5 & 1 & 798 & 2 & 224 & \texttt{-++???????}\\*
1 & 0 & 1 & 5 & 1 & 572 & 4 & 448 & \texttt{--+???????}\\
2 & 1 & 2 & 5 & 1 & 120 & 8 & 896 & \texttt{---???????}\\
3 & 2 & 3 & 2 & 1 & 256 & 256 & 512 & \texttt{----??----}\\
4 & 3 & 4 & 1 & 2 & 0 & 1024 & 0 & \texttt{----------}\\
5 & 3 & 4 & 1 & 2 & 0 & 1024 & 0 & \texttt{----------}\\
6 & 2 & 3 & 1 & 1 & 512 & 512 & 0 & \texttt{-------+--}\\
7 & 6 & 4 & 1 & 1 & 0 & 1024 & 0 & \texttt{----------}\\
8 & 2 & 3 & 1 & 2 & 0 & 1024 & 0 & \texttt{----------}\\
9 & 2 & 3 & 1 & 2 & 0 & 1024 & 0 & \texttt{----------}\\
\multicolumn{9}{@{}l@{}}{\rule{0pt}{2.6ex}\textbf{T10.360}\enspace \texttt{0123443432}\enspace $n{=}10$, $|H|{=}5$, len 5, br 3, $|\mathrm{Aut}|{=}2$, orb 9/8, hist-orb 4/3; M 26624, res 13056, rev 7778: ok}\\*
0 & -- & 0 & 5 & 1 & 670 & 2 & 352 & \texttt{-+????????}\\*
1 & 0 & 1 & 5 & 1 & 316 & 4 & 704 & \texttt{--????????}\\
2 & 1 & 2 & 4 & 1 & 240 & 16 & 768 & \texttt{---??????-}\\
3 & 2 & 3 & 2 & 1 & 256 & 256 & 512 & \texttt{----??----}\\
4 & 3 & 4 & 1 & 2 & 0 & 1024 & 0 & \texttt{----------}\\
5 & 3 & 4 & 1 & 2 & 0 & 1024 & 0 & \texttt{----------}\\
6 & 2 & 3 & 1 & 1 & 512 & 512 & 0 & \texttt{-------+--}\\
7 & 6 & 4 & 1 & 1 & 0 & 1024 & 0 & \texttt{----------}\\
8 & 2 & 3 & 1 & 1 & 0 & 1024 & 0 & \texttt{----------}\\
9 & 1 & 2 & 1 & 1 & 0 & 1024 & 0 & \texttt{----------}\\
\multicolumn{9}{@{}l@{}}{\rule{0pt}{2.6ex}\textbf{T10.361}\enspace \texttt{0123443431}\enspace $n{=}10$, $|H|{=}5$, len 5, br 3, $|\mathrm{Aut}|{=}2$, orb 9/8, hist-orb 4/3; M 25600, res 12576, rev 7478: ok}\\*
0 & -- & 0 & 5 & 1 & 414 & 2 & 608 & \texttt{-?????????}\\*
1 & 0 & 1 & 4 & 1 & 632 & 8 & 384 & \texttt{--+??????-}\\
2 & 1 & 2 & 4 & 1 & 240 & 16 & 768 & \texttt{---??????-}\\
3 & 2 & 3 & 2 & 1 & 256 & 256 & 512 & \texttt{----??----}\\
4 & 3 & 4 & 1 & 2 & 0 & 1024 & 0 & \texttt{----------}\\
5 & 3 & 4 & 1 & 2 & 0 & 1024 & 0 & \texttt{----------}\\
6 & 2 & 3 & 1 & 1 & 512 & 512 & 0 & \texttt{-------+--}\\
7 & 6 & 4 & 1 & 1 & 0 & 1024 & 0 & \texttt{----------}\\
8 & 2 & 3 & 1 & 1 & 0 & 1024 & 0 & \texttt{----------}\\
9 & 0 & 1 & 1 & 1 & 0 & 1024 & 0 & \texttt{----------}\\
\multicolumn{9}{@{}l@{}}{\rule{0pt}{2.6ex}\textbf{T10.362}\enspace \texttt{0123443423}\enspace $n{=}10$, $|H|{=}4$, len 5, br 3, $|\mathrm{Aut}|{=}2$, orb 9/8, hist-orb 3/2; M 26624, res 8224, rev 10198: ok}\\*
0 & -- & 0 & 4 & 1 & 794 & 2 & 228 & \texttt{-+????????}\\*
1 & 0 & 1 & 4 & 1 & 564 & 4 & 456 & \texttt{--????????}\\
2 & 1 & 2 & 3 & 1 & 480 & 32 & 512 & \texttt{---?????--}\\
3 & 2 & 3 & 2 & 1 & 256 & 256 & 512 & \texttt{----??----}\\
4 & 3 & 4 & 1 & 2 & 0 & 1024 & 0 & \texttt{----------}\\
5 & 3 & 4 & 1 & 2 & 0 & 1024 & 0 & \texttt{----------}\\
6 & 2 & 3 & 1 & 1 & 512 & 512 & 0 & \texttt{-------+--}\\
7 & 6 & 4 & 1 & 1 & 0 & 1024 & 0 & \texttt{----------}\\
8 & 1 & 2 & 1 & 1 & 512 & 512 & 0 & \texttt{---------+}\\
9 & 8 & 3 & 1 & 1 & 0 & 1024 & 0 & \texttt{----------}\\
\multicolumn{9}{@{}l@{}}{\rule{0pt}{2.6ex}\textbf{T10.363}\enspace \texttt{0123443422}\enspace $n{=}10$, $|H|{=}5$, len 5, br 3, $|\mathrm{Aut}|{=}4$, orb 8/6, hist-orb 3/1; M 25600, res 12736, rev 7290: ok}\\*
0 & -- & 0 & 5 & 1 & 606 & 2 & 416 & \texttt{-+????????}\\*
1 & 0 & 1 & 5 & 1 & 188 & 4 & 832 & \texttt{--????????}\\
2 & 1 & 2 & 3 & 1 & 480 & 32 & 512 & \texttt{---?????--}\\
3 & 2 & 3 & 2 & 1 & 256 & 256 & 512 & \texttt{----??----}\\
4 & 3 & 4 & 1 & 2 & 0 & 1024 & 0 & \texttt{----------}\\
5 & 3 & 4 & 1 & 2 & 0 & 1024 & 0 & \texttt{----------}\\
6 & 2 & 3 & 1 & 1 & 512 & 512 & 0 & \texttt{-------+--}\\
7 & 6 & 4 & 1 & 1 & 0 & 1024 & 0 & \texttt{----------}\\
8 & 1 & 2 & 1 & 2 & 0 & 1024 & 0 & \texttt{----------}\\
9 & 1 & 2 & 1 & 2 & 0 & 1024 & 0 & \texttt{----------}\\
\multicolumn{9}{@{}l@{}}{\rule{0pt}{2.6ex}\textbf{T10.364}\enspace \texttt{0123443421}\enspace $n{=}10$, $|H|{=}5$, len 5, br 4, $|\mathrm{Aut}|{=}2$, orb 9/8, hist-orb 4/3; M 24576, res 12960, rev 6506: ok}\\*
0 & -- & 0 & 5 & 1 & 350 & 2 & 672 & \texttt{-?????????}\\*
1 & 0 & 1 & 4 & 1 & 376 & 8 & 640 & \texttt{--???????-}\\
2 & 1 & 2 & 3 & 1 & 480 & 32 & 512 & \texttt{---?????--}\\
3 & 2 & 3 & 2 & 1 & 256 & 256 & 512 & \texttt{----??----}\\
4 & 3 & 4 & 1 & 2 & 0 & 1024 & 0 & \texttt{----------}\\
5 & 3 & 4 & 1 & 2 & 0 & 1024 & 0 & \texttt{----------}\\
6 & 2 & 3 & 1 & 1 & 512 & 512 & 0 & \texttt{-------+--}\\
7 & 6 & 4 & 1 & 1 & 0 & 1024 & 0 & \texttt{----------}\\
8 & 1 & 2 & 1 & 1 & 0 & 1024 & 0 & \texttt{----------}\\
9 & 0 & 1 & 1 & 1 & 0 & 1024 & 0 & \texttt{----------}\\
\multicolumn{9}{@{}l@{}}{\rule{0pt}{2.6ex}\textbf{T10.365}\enspace \texttt{0123443412}\enspace $n{=}10$, $|H|{=}4$, len 5, br 3, $|\mathrm{Aut}|{=}2$, orb 9/8, hist-orb 3/2; M 24576, res 7364, rev 9406: ok}\\*
0 & -- & 0 & 4 & 1 & 666 & 2 & 356 & \texttt{-?????????}\\*
1 & 0 & 1 & 3 & 1 & 752 & 16 & 256 & \texttt{--+?????--}\\
2 & 1 & 2 & 3 & 1 & 480 & 32 & 512 & \texttt{---?????--}\\
3 & 2 & 3 & 2 & 1 & 256 & 256 & 512 & \texttt{----??----}\\
4 & 3 & 4 & 1 & 2 & 0 & 1024 & 0 & \texttt{----------}\\
5 & 3 & 4 & 1 & 2 & 0 & 1024 & 0 & \texttt{----------}\\
6 & 2 & 3 & 1 & 1 & 512 & 512 & 0 & \texttt{-------+--}\\
7 & 6 & 4 & 1 & 1 & 0 & 1024 & 0 & \texttt{----------}\\
8 & 0 & 1 & 1 & 1 & 512 & 512 & 0 & \texttt{---------+}\\
9 & 8 & 2 & 1 & 1 & 0 & 1024 & 0 & \texttt{----------}\\
\multicolumn{9}{@{}l@{}}{\rule{0pt}{2.6ex}\textbf{T10.366}\enspace \texttt{0123443411}\enspace $n{=}10$, $|H|{=}5$, len 5, br 3, $|\mathrm{Aut}|{=}4$, orb 8/6, hist-orb 3/1; M 23552, res 11360, rev 6782: ok}\\*
0 & -- & 0 & 5 & 1 & 222 & 2 & 800 & \texttt{-?????????}\\*
1 & 0 & 1 & 3 & 1 & 752 & 16 & 256 & \texttt{--+?????--}\\
2 & 1 & 2 & 3 & 1 & 480 & 32 & 512 & \texttt{---?????--}\\
3 & 2 & 3 & 2 & 1 & 256 & 256 & 512 & \texttt{----??----}\\
4 & 3 & 4 & 1 & 2 & 0 & 1024 & 0 & \texttt{----------}\\
5 & 3 & 4 & 1 & 2 & 0 & 1024 & 0 & \texttt{----------}\\
6 & 2 & 3 & 1 & 1 & 512 & 512 & 0 & \texttt{-------+--}\\
7 & 6 & 4 & 1 & 1 & 0 & 1024 & 0 & \texttt{----------}\\
8 & 0 & 1 & 1 & 2 & 0 & 1024 & 0 & \texttt{----------}\\
9 & 0 & 1 & 1 & 2 & 0 & 1024 & 0 & \texttt{----------}\\
\multicolumn{9}{@{}l@{}}{\rule{0pt}{2.6ex}\textbf{T10.367}\enspace \texttt{0123443333}\enspace $n{=}10$, $|H|{=}6$, len 5, br 2, $|\mathrm{Aut}|{=}48$, orb 6/4, hist-orb 2/0; M 26624, res 12532, rev 8958: ok}\\*
0 & -- & 0 & 6 & 1 & 778 & 2 & 244 & \texttt{-++???????}\\*
1 & 0 & 1 & 6 & 1 & 532 & 4 & 488 & \texttt{--+???????}\\
2 & 1 & 2 & 6 & 1 & 40 & 8 & 976 & \texttt{---???????}\\
3 & 2 & 3 & 2 & 1 & 256 & 256 & 512 & \texttt{----??----}\\
4 & 3 & 4 & 1 & 2 & 0 & 1024 & 0 & \texttt{----------}\\
5 & 3 & 4 & 1 & 2 & 0 & 1024 & 0 & \texttt{----------}\\
6 & 2 & 3 & 1 & 4 & 0 & 1024 & 0 & \texttt{----------}\\
7 & 2 & 3 & 1 & 4 & 0 & 1024 & 0 & \texttt{----------}\\
8 & 2 & 3 & 1 & 4 & 0 & 1024 & 0 & \texttt{----------}\\
9 & 2 & 3 & 1 & 4 & 0 & 1024 & 0 & \texttt{----------}\\
\multicolumn{9}{@{}l@{}}{\rule{0pt}{2.6ex}\textbf{T10.368}\enspace \texttt{0123443332}\enspace $n{=}10$, $|H|{=}6$, len 5, br 3, $|\mathrm{Aut}|{=}12$, orb 7/5, hist-orb 3/1; M 25600, res 14376, rev 6730: ok}\\*
0 & -- & 0 & 6 & 1 & 650 & 2 & 372 & \texttt{-+????????}\\*
1 & 0 & 1 & 6 & 1 & 276 & 4 & 744 & \texttt{--????????}\\
2 & 1 & 2 & 5 & 1 & 80 & 16 & 928 & \texttt{---??????-}\\
3 & 2 & 3 & 2 & 1 & 256 & 256 & 512 & \texttt{----??----}\\
4 & 3 & 4 & 1 & 2 & 0 & 1024 & 0 & \texttt{----------}\\
5 & 3 & 4 & 1 & 2 & 0 & 1024 & 0 & \texttt{----------}\\
6 & 2 & 3 & 1 & 3 & 0 & 1024 & 0 & \texttt{----------}\\
7 & 2 & 3 & 1 & 3 & 0 & 1024 & 0 & \texttt{----------}\\
8 & 2 & 3 & 1 & 3 & 0 & 1024 & 0 & \texttt{----------}\\
9 & 1 & 2 & 1 & 1 & 0 & 1024 & 0 & \texttt{----------}\\
\multicolumn{9}{@{}l@{}}{\rule{0pt}{2.6ex}\textbf{T10.369}\enspace \texttt{0123443331}\enspace $n{=}10$, $|H|{=}6$, len 5, br 3, $|\mathrm{Aut}|{=}12$, orb 7/5, hist-orb 3/1; M 24576, res 14116, rev 6302: ok}\\*
0 & -- & 0 & 6 & 1 & 394 & 2 & 628 & \texttt{-?????????}\\*
1 & 0 & 1 & 5 & 1 & 552 & 8 & 464 & \texttt{--+??????-}\\
2 & 1 & 2 & 5 & 1 & 80 & 16 & 928 & \texttt{---??????-}\\
3 & 2 & 3 & 2 & 1 & 256 & 256 & 512 & \texttt{----??----}\\
4 & 3 & 4 & 1 & 2 & 0 & 1024 & 0 & \texttt{----------}\\
5 & 3 & 4 & 1 & 2 & 0 & 1024 & 0 & \texttt{----------}\\
6 & 2 & 3 & 1 & 3 & 0 & 1024 & 0 & \texttt{----------}\\
7 & 2 & 3 & 1 & 3 & 0 & 1024 & 0 & \texttt{----------}\\
8 & 2 & 3 & 1 & 3 & 0 & 1024 & 0 & \texttt{----------}\\
9 & 0 & 1 & 1 & 1 & 0 & 1024 & 0 & \texttt{----------}\\
\multicolumn{9}{@{}l@{}}{\rule{0pt}{2.6ex}\textbf{T10.370}\enspace \texttt{0123443323}\enspace $n{=}10$, $|H|{=}5$, len 5, br 3, $|\mathrm{Aut}|{=}4$, orb 8/6, hist-orb 3/1; M 25600, res 11024, rev 8390: ok}\\*
0 & -- & 0 & 5 & 1 & 734 & 2 & 288 & \texttt{-+????????}\\*
1 & 0 & 1 & 5 & 1 & 444 & 4 & 576 & \texttt{--????????}\\
2 & 1 & 2 & 4 & 1 & 160 & 32 & 832 & \texttt{---?????--}\\
3 & 2 & 3 & 2 & 1 & 256 & 256 & 512 & \texttt{----??----}\\
4 & 3 & 4 & 1 & 2 & 0 & 1024 & 0 & \texttt{----------}\\
5 & 3 & 4 & 1 & 2 & 0 & 1024 & 0 & \texttt{----------}\\
6 & 2 & 3 & 1 & 2 & 0 & 1024 & 0 & \texttt{----------}\\
7 & 2 & 3 & 1 & 2 & 0 & 1024 & 0 & \texttt{----------}\\
8 & 1 & 2 & 1 & 1 & 512 & 512 & 0 & \texttt{---------+}\\
9 & 8 & 3 & 1 & 1 & 0 & 1024 & 0 & \texttt{----------}\\
\multicolumn{9}{@{}l@{}}{\rule{0pt}{2.6ex}\textbf{T10.371}\enspace \texttt{0123443322}\enspace $n{=}10$, $|H|{=}6$, len 5, br 3, $|\mathrm{Aut}|{=}8$, orb 7/4, hist-orb 3/0; M 24576, res 14576, rev 5794: ok}\\*
0 & -- & 0 & 6 & 1 & 586 & 2 & 436 & \texttt{-+????????}\\*
1 & 0 & 1 & 6 & 1 & 148 & 4 & 872 & \texttt{--????????}\\
2 & 1 & 2 & 4 & 1 & 160 & 32 & 832 & \texttt{---?????--}\\
3 & 2 & 3 & 2 & 1 & 256 & 256 & 512 & \texttt{----??----}\\
4 & 3 & 4 & 1 & 2 & 0 & 1024 & 0 & \texttt{----------}\\
5 & 3 & 4 & 1 & 2 & 0 & 1024 & 0 & \texttt{----------}\\
6 & 2 & 3 & 1 & 2 & 0 & 1024 & 0 & \texttt{----------}\\
7 & 2 & 3 & 1 & 2 & 0 & 1024 & 0 & \texttt{----------}\\
8 & 1 & 2 & 1 & 2 & 0 & 1024 & 0 & \texttt{----------}\\
9 & 1 & 2 & 1 & 2 & 0 & 1024 & 0 & \texttt{----------}\\
\multicolumn{9}{@{}l@{}}{\rule{0pt}{2.6ex}\textbf{T10.372}\enspace \texttt{0123443321}\enspace $n{=}10$, $|H|{=}6$, len 5, br 4, $|\mathrm{Aut}|{=}4$, orb 8/6, hist-orb 4/2; M 23552, res 15020, rev 4818: ok}\\*
0 & -- & 0 & 6 & 1 & 330 & 2 & 692 & \texttt{-?????????}\\*
1 & 0 & 1 & 5 & 1 & 296 & 8 & 720 & \texttt{--???????-}\\
2 & 1 & 2 & 4 & 1 & 160 & 32 & 832 & \texttt{---?????--}\\
3 & 2 & 3 & 2 & 1 & 256 & 256 & 512 & \texttt{----??----}\\
4 & 3 & 4 & 1 & 2 & 0 & 1024 & 0 & \texttt{----------}\\
5 & 3 & 4 & 1 & 2 & 0 & 1024 & 0 & \texttt{----------}\\
6 & 2 & 3 & 1 & 2 & 0 & 1024 & 0 & \texttt{----------}\\
7 & 2 & 3 & 1 & 2 & 0 & 1024 & 0 & \texttt{----------}\\
8 & 1 & 2 & 1 & 1 & 0 & 1024 & 0 & \texttt{----------}\\
9 & 0 & 1 & 1 & 1 & 0 & 1024 & 0 & \texttt{----------}\\
\multicolumn{9}{@{}l@{}}{\rule{0pt}{2.6ex}\textbf{T10.373}\enspace \texttt{0123443312}\enspace $n{=}10$, $|H|{=}5$, len 5, br 3, $|\mathrm{Aut}|{=}4$, orb 8/6, hist-orb 3/1; M 23552, res 10144, rev 7630: ok}\\*
0 & -- & 0 & 5 & 1 & 606 & 2 & 416 & \texttt{-?????????}\\*
1 & 0 & 1 & 4 & 1 & 592 & 16 & 416 & \texttt{--+?????--}\\
2 & 1 & 2 & 4 & 1 & 160 & 32 & 832 & \texttt{---?????--}\\
3 & 2 & 3 & 2 & 1 & 256 & 256 & 512 & \texttt{----??----}\\
4 & 3 & 4 & 1 & 2 & 0 & 1024 & 0 & \texttt{----------}\\
5 & 3 & 4 & 1 & 2 & 0 & 1024 & 0 & \texttt{----------}\\
6 & 2 & 3 & 1 & 2 & 0 & 1024 & 0 & \texttt{----------}\\
7 & 2 & 3 & 1 & 2 & 0 & 1024 & 0 & \texttt{----------}\\
8 & 0 & 1 & 1 & 1 & 512 & 512 & 0 & \texttt{---------+}\\
9 & 8 & 2 & 1 & 1 & 0 & 1024 & 0 & \texttt{----------}\\
\multicolumn{9}{@{}l@{}}{\rule{0pt}{2.6ex}\textbf{T10.374}\enspace \texttt{0123443311}\enspace $n{=}10$, $|H|{=}6$, len 5, br 3, $|\mathrm{Aut}|{=}8$, orb 7/4, hist-orb 3/0; M 22528, res 13780, rev 4982: ok}\\*
0 & -- & 0 & 6 & 1 & 202 & 2 & 820 & \texttt{-?????????}\\*
1 & 0 & 1 & 4 & 1 & 592 & 16 & 416 & \texttt{--+?????--}\\
2 & 1 & 2 & 4 & 1 & 160 & 32 & 832 & \texttt{---?????--}\\
3 & 2 & 3 & 2 & 1 & 256 & 256 & 512 & \texttt{----??----}\\
4 & 3 & 4 & 1 & 2 & 0 & 1024 & 0 & \texttt{----------}\\
5 & 3 & 4 & 1 & 2 & 0 & 1024 & 0 & \texttt{----------}\\
6 & 2 & 3 & 1 & 2 & 0 & 1024 & 0 & \texttt{----------}\\
7 & 2 & 3 & 1 & 2 & 0 & 1024 & 0 & \texttt{----------}\\
8 & 0 & 1 & 1 & 2 & 0 & 1024 & 0 & \texttt{----------}\\
9 & 0 & 1 & 1 & 2 & 0 & 1024 & 0 & \texttt{----------}\\
\multicolumn{9}{@{}l@{}}{\rule{0pt}{2.6ex}\textbf{T10.375}\enspace \texttt{0123443234}\enspace $n{=}10$, $|H|{=}4$, len 5, br 3, $|\mathrm{Aut}|{=}2$, orb 9/8, hist-orb 3/2; M 26624, res 7840, rev 9950: ok}\\*
0 & -- & 0 & 4 & 1 & 806 & 2 & 216 & \texttt{-+????????}\\*
1 & 0 & 1 & 4 & 1 & 588 & 4 & 432 & \texttt{--????????}\\
2 & 1 & 2 & 3 & 1 & 320 & 64 & 640 & \texttt{---????---}\\
3 & 2 & 3 & 2 & 1 & 256 & 256 & 512 & \texttt{----??----}\\
4 & 3 & 4 & 1 & 2 & 0 & 1024 & 0 & \texttt{----------}\\
5 & 3 & 4 & 1 & 2 & 0 & 1024 & 0 & \texttt{----------}\\
6 & 2 & 3 & 1 & 1 & 0 & 1024 & 0 & \texttt{----------}\\
7 & 1 & 2 & 1 & 1 & 768 & 256 & 0 & \texttt{--------++}\\
8 & 7 & 3 & 1 & 1 & 512 & 512 & 0 & \texttt{---------+}\\
9 & 8 & 4 & 1 & 1 & 0 & 1024 & 0 & \texttt{----------}\\
\multicolumn{9}{@{}l@{}}{\rule{0pt}{2.6ex}\textbf{T10.376}\enspace \texttt{0123443233}\enspace $n{=}10$, $|H|{=}5$, len 5, br 4, $|\mathrm{Aut}|{=}4$, orb 8/6, hist-orb 3/1; M 25600, res 10496, rev 8354: ok}\\*
0 & -- & 0 & 5 & 1 & 722 & 2 & 300 & \texttt{-+????????}\\*
1 & 0 & 1 & 5 & 1 & 420 & 4 & 600 & \texttt{--????????}\\
2 & 1 & 2 & 3 & 1 & 320 & 64 & 640 & \texttt{---????---}\\
3 & 2 & 3 & 2 & 1 & 256 & 256 & 512 & \texttt{----??----}\\
4 & 3 & 4 & 1 & 2 & 0 & 1024 & 0 & \texttt{----------}\\
5 & 3 & 4 & 1 & 2 & 0 & 1024 & 0 & \texttt{----------}\\
6 & 2 & 3 & 1 & 1 & 0 & 1024 & 0 & \texttt{----------}\\
7 & 1 & 2 & 2 & 1 & 256 & 256 & 512 & \texttt{--------??}\\
8 & 7 & 3 & 1 & 2 & 0 & 1024 & 0 & \texttt{----------}\\
9 & 7 & 3 & 1 & 2 & 0 & 1024 & 0 & \texttt{----------}\\
\multicolumn{9}{@{}l@{}}{\rule{0pt}{2.6ex}\textbf{T10.377}\enspace \texttt{0123443232}\enspace $n{=}10$, $|H|{=}5$, len 5, br 3, $|\mathrm{Aut}|{=}2$, orb 9/8, hist-orb 4/3; M 24576, res 11872, rev 7150: ok}\\*
0 & -- & 0 & 5 & 1 & 638 & 2 & 384 & \texttt{-+????????}\\*
1 & 0 & 1 & 5 & 1 & 252 & 4 & 768 & \texttt{--????????}\\
2 & 1 & 2 & 3 & 1 & 320 & 64 & 640 & \texttt{---????---}\\
3 & 2 & 3 & 2 & 1 & 256 & 256 & 512 & \texttt{----??----}\\
4 & 3 & 4 & 1 & 2 & 0 & 1024 & 0 & \texttt{----------}\\
5 & 3 & 4 & 1 & 2 & 0 & 1024 & 0 & \texttt{----------}\\
6 & 2 & 3 & 1 & 1 & 0 & 1024 & 0 & \texttt{----------}\\
7 & 1 & 2 & 1 & 1 & 512 & 512 & 0 & \texttt{--------+-}\\
8 & 7 & 3 & 1 & 1 & 0 & 1024 & 0 & \texttt{----------}\\
9 & 1 & 2 & 1 & 1 & 0 & 1024 & 0 & \texttt{----------}\\
\multicolumn{9}{@{}l@{}}{\rule{0pt}{2.6ex}\textbf{T10.378}\enspace \texttt{0123443231}\enspace $n{=}10$, $|H|{=}5$, len 5, br 4, $|\mathrm{Aut}|{=}2$, orb 9/8, hist-orb 4/3; M 23552, res 11744, rev 6574: ok}\\*
0 & -- & 0 & 5 & 1 & 382 & 2 & 640 & \texttt{-?????????}\\*
1 & 0 & 1 & 4 & 1 & 504 & 8 & 512 & \texttt{--???????-}\\
2 & 1 & 2 & 3 & 1 & 320 & 64 & 640 & \texttt{---????---}\\
3 & 2 & 3 & 2 & 1 & 256 & 256 & 512 & \texttt{----??----}\\
4 & 3 & 4 & 1 & 2 & 0 & 1024 & 0 & \texttt{----------}\\
5 & 3 & 4 & 1 & 2 & 0 & 1024 & 0 & \texttt{----------}\\
6 & 2 & 3 & 1 & 1 & 0 & 1024 & 0 & \texttt{----------}\\
7 & 1 & 2 & 1 & 1 & 512 & 512 & 0 & \texttt{--------+-}\\
8 & 7 & 3 & 1 & 1 & 0 & 1024 & 0 & \texttt{----------}\\
9 & 0 & 1 & 1 & 1 & 0 & 1024 & 0 & \texttt{----------}\\
\multicolumn{9}{@{}l@{}}{\rule{0pt}{2.6ex}\textbf{T10.379}\enspace \texttt{0123443222}\enspace $n{=}10$, $|H|{=}6$, len 5, br 3, $|\mathrm{Aut}|{=}12$, orb 7/5, hist-orb 3/1; M 23552, res 13888, rev 5562: ok}\\*
0 & -- & 0 & 6 & 1 & 554 & 2 & 468 & \texttt{-+????????}\\*
1 & 0 & 1 & 6 & 1 & 84 & 4 & 936 & \texttt{--????????}\\
2 & 1 & 2 & 3 & 1 & 320 & 64 & 640 & \texttt{---????---}\\
3 & 2 & 3 & 2 & 1 & 256 & 256 & 512 & \texttt{----??----}\\
4 & 3 & 4 & 1 & 2 & 0 & 1024 & 0 & \texttt{----------}\\
5 & 3 & 4 & 1 & 2 & 0 & 1024 & 0 & \texttt{----------}\\
6 & 2 & 3 & 1 & 1 & 0 & 1024 & 0 & \texttt{----------}\\
7 & 1 & 2 & 1 & 3 & 0 & 1024 & 0 & \texttt{----------}\\
8 & 1 & 2 & 1 & 3 & 0 & 1024 & 0 & \texttt{----------}\\
9 & 1 & 2 & 1 & 3 & 0 & 1024 & 0 & \texttt{----------}\\
\multicolumn{9}{@{}l@{}}{\rule{0pt}{2.6ex}\textbf{T10.380}\enspace \texttt{0123443221}\enspace $n{=}10$, $|H|{=}6$, len 5, br 4, $|\mathrm{Aut}|{=}4$, orb 8/6, hist-orb 4/2; M 22528, res 14684, rev 4306: ok}\\*
0 & -- & 0 & 6 & 1 & 298 & 2 & 724 & \texttt{-?????????}\\*
1 & 0 & 1 & 5 & 1 & 168 & 8 & 848 & \texttt{--???????-}\\
2 & 1 & 2 & 3 & 1 & 320 & 64 & 640 & \texttt{---????---}\\
3 & 2 & 3 & 2 & 1 & 256 & 256 & 512 & \texttt{----??----}\\
4 & 3 & 4 & 1 & 2 & 0 & 1024 & 0 & \texttt{----------}\\
5 & 3 & 4 & 1 & 2 & 0 & 1024 & 0 & \texttt{----------}\\
6 & 2 & 3 & 1 & 1 & 0 & 1024 & 0 & \texttt{----------}\\
7 & 1 & 2 & 1 & 2 & 0 & 1024 & 0 & \texttt{----------}\\
8 & 1 & 2 & 1 & 2 & 0 & 1024 & 0 & \texttt{----------}\\
9 & 0 & 1 & 1 & 1 & 0 & 1024 & 0 & \texttt{----------}\\
\multicolumn{9}{@{}l@{}}{\rule{0pt}{2.6ex}\textbf{T10.381}\enspace \texttt{0123443212}\enspace $n{=}10$, $|H|{=}5$, len 5, br 4, $|\mathrm{Aut}|{=}2$, orb 9/8, hist-orb 4/3; M 22528, res 10960, rev 6370: ok}\\*
0 & -- & 0 & 5 & 1 & 510 & 2 & 512 & \texttt{-?????????}\\*
1 & 0 & 1 & 4 & 1 & 336 & 16 & 672 & \texttt{--??????--}\\
2 & 1 & 2 & 3 & 1 & 320 & 64 & 640 & \texttt{---????---}\\
3 & 2 & 3 & 2 & 1 & 256 & 256 & 512 & \texttt{----??----}\\
4 & 3 & 4 & 1 & 2 & 0 & 1024 & 0 & \texttt{----------}\\
5 & 3 & 4 & 1 & 2 & 0 & 1024 & 0 & \texttt{----------}\\
6 & 2 & 3 & 1 & 1 & 0 & 1024 & 0 & \texttt{----------}\\
7 & 1 & 2 & 1 & 1 & 0 & 1024 & 0 & \texttt{----------}\\
8 & 0 & 1 & 1 & 1 & 512 & 512 & 0 & \texttt{---------+}\\
9 & 8 & 2 & 1 & 1 & 0 & 1024 & 0 & \texttt{----------}\\
\multicolumn{9}{@{}l@{}}{\rule{0pt}{2.6ex}\textbf{T10.382}\enspace \texttt{0123443211}\enspace $n{=}10$, $|H|{=}6$, len 5, br 4, $|\mathrm{Aut}|{=}4$, orb 8/6, hist-orb 4/2; M 21504, res 14020, rev 4050: ok}\\*
0 & -- & 0 & 6 & 1 & 170 & 2 & 852 & \texttt{-?????????}\\*
1 & 0 & 1 & 4 & 1 & 336 & 16 & 672 & \texttt{--??????--}\\
2 & 1 & 2 & 3 & 1 & 320 & 64 & 640 & \texttt{---????---}\\
3 & 2 & 3 & 2 & 1 & 256 & 256 & 512 & \texttt{----??----}\\
4 & 3 & 4 & 1 & 2 & 0 & 1024 & 0 & \texttt{----------}\\
5 & 3 & 4 & 1 & 2 & 0 & 1024 & 0 & \texttt{----------}\\
6 & 2 & 3 & 1 & 1 & 0 & 1024 & 0 & \texttt{----------}\\
7 & 1 & 2 & 1 & 1 & 0 & 1024 & 0 & \texttt{----------}\\
8 & 0 & 1 & 1 & 2 & 0 & 1024 & 0 & \texttt{----------}\\
9 & 0 & 1 & 1 & 2 & 0 & 1024 & 0 & \texttt{----------}\\
\multicolumn{9}{@{}l@{}}{\rule{0pt}{2.6ex}\textbf{T10.383}\enspace \texttt{0123443123}\enspace $n{=}10$, $|H|{=}4$, len 5, br 3, $|\mathrm{Aut}|{=}2$, orb 9/8, hist-orb 3/2; M 23552, res 6616, rev 8910: ok}\\*
0 & -- & 0 & 4 & 1 & 742 & 2 & 280 & \texttt{-?????????}\\*
1 & 0 & 1 & 3 & 1 & 672 & 32 & 320 & \texttt{--+????---}\\
2 & 1 & 2 & 3 & 1 & 320 & 64 & 640 & \texttt{---????---}\\
3 & 2 & 3 & 2 & 1 & 256 & 256 & 512 & \texttt{----??----}\\
4 & 3 & 4 & 1 & 2 & 0 & 1024 & 0 & \texttt{----------}\\
5 & 3 & 4 & 1 & 2 & 0 & 1024 & 0 & \texttt{----------}\\
6 & 2 & 3 & 1 & 1 & 0 & 1024 & 0 & \texttt{----------}\\
7 & 0 & 1 & 1 & 1 & 768 & 256 & 0 & \texttt{--------++}\\
8 & 7 & 2 & 1 & 1 & 512 & 512 & 0 & \texttt{---------+}\\
9 & 8 & 3 & 1 & 1 & 0 & 1024 & 0 & \texttt{----------}\\
\multicolumn{9}{@{}l@{}}{\rule{0pt}{2.6ex}\textbf{T10.384}\enspace \texttt{0123443122}\enspace $n{=}10$, $|H|{=}5$, len 5, br 4, $|\mathrm{Aut}|{=}4$, orb 8/6, hist-orb 3/1; M 22528, res 9292, rev 7182: ok}\\*
0 & -- & 0 & 5 & 1 & 530 & 2 & 492 & \texttt{-?????????}\\*
1 & 0 & 1 & 3 & 1 & 672 & 32 & 320 & \texttt{--+????---}\\
2 & 1 & 2 & 3 & 1 & 320 & 64 & 640 & \texttt{---????---}\\
3 & 2 & 3 & 2 & 1 & 256 & 256 & 512 & \texttt{----??----}\\
4 & 3 & 4 & 1 & 2 & 0 & 1024 & 0 & \texttt{----------}\\
5 & 3 & 4 & 1 & 2 & 0 & 1024 & 0 & \texttt{----------}\\
6 & 2 & 3 & 1 & 1 & 0 & 1024 & 0 & \texttt{----------}\\
7 & 0 & 1 & 2 & 1 & 256 & 256 & 512 & \texttt{--------??}\\
8 & 7 & 2 & 1 & 2 & 0 & 1024 & 0 & \texttt{----------}\\
9 & 7 & 2 & 1 & 2 & 0 & 1024 & 0 & \texttt{----------}\\
\multicolumn{9}{@{}l@{}}{\rule{0pt}{2.6ex}\textbf{T10.385}\enspace \texttt{0123443121}\enspace $n{=}10$, $|H|{=}5$, len 5, br 3, $|\mathrm{Aut}|{=}2$, orb 9/8, hist-orb 4/3; M 21504, res 10432, rev 6054: ok}\\*
0 & -- & 0 & 5 & 1 & 318 & 2 & 704 & \texttt{-?????????}\\*
1 & 0 & 1 & 3 & 1 & 672 & 32 & 320 & \texttt{--+????---}\\
2 & 1 & 2 & 3 & 1 & 320 & 64 & 640 & \texttt{---????---}\\
3 & 2 & 3 & 2 & 1 & 256 & 256 & 512 & \texttt{----??----}\\
4 & 3 & 4 & 1 & 2 & 0 & 1024 & 0 & \texttt{----------}\\
5 & 3 & 4 & 1 & 2 & 0 & 1024 & 0 & \texttt{----------}\\
6 & 2 & 3 & 1 & 1 & 0 & 1024 & 0 & \texttt{----------}\\
7 & 0 & 1 & 1 & 1 & 512 & 512 & 0 & \texttt{--------+-}\\
8 & 7 & 2 & 1 & 1 & 0 & 1024 & 0 & \texttt{----------}\\
9 & 0 & 1 & 1 & 1 & 0 & 1024 & 0 & \texttt{----------}\\
\multicolumn{9}{@{}l@{}}{\rule{0pt}{2.6ex}\textbf{T10.386}\enspace \texttt{0123443111}\enspace $n{=}10$, $|H|{=}6$, len 5, br 3, $|\mathrm{Aut}|{=}12$, orb 7/5, hist-orb 3/1; M 20480, res 12340, rev 4414: ok}\\*
0 & -- & 0 & 6 & 1 & 106 & 2 & 916 & \texttt{-?????????}\\*
1 & 0 & 1 & 3 & 1 & 672 & 32 & 320 & \texttt{--+????---}\\
2 & 1 & 2 & 3 & 1 & 320 & 64 & 640 & \texttt{---????---}\\
3 & 2 & 3 & 2 & 1 & 256 & 256 & 512 & \texttt{----??----}\\
4 & 3 & 4 & 1 & 2 & 0 & 1024 & 0 & \texttt{----------}\\
5 & 3 & 4 & 1 & 2 & 0 & 1024 & 0 & \texttt{----------}\\
6 & 2 & 3 & 1 & 1 & 0 & 1024 & 0 & \texttt{----------}\\
7 & 0 & 1 & 1 & 3 & 0 & 1024 & 0 & \texttt{----------}\\
8 & 0 & 1 & 1 & 3 & 0 & 1024 & 0 & \texttt{----------}\\
9 & 0 & 1 & 1 & 3 & 0 & 1024 & 0 & \texttt{----------}\\
\multicolumn{9}{@{}l@{}}{\rule{0pt}{2.6ex}\textbf{T10.387}\enspace \texttt{0123442344}\enspace $n{=}10$, $|H|{=}4$, len 5, br 3, $|\mathrm{Aut}|{=}8$, orb 5/2, hist-orb 1/0; M 27648, res 6432, rev 10954: ok}\\*
0 & -- & 0 & 4 & 1 & 850 & 2 & 172 & \texttt{-+????????}\\*
1 & 0 & 1 & 4 & 1 & 676 & 4 & 344 & \texttt{--????????}\\
2 & 1 & 2 & 2 & 2 & 640 & 128 & 256 & \texttt{---+??----}\\
3 & 2 & 3 & 2 & 2 & 256 & 256 & 512 & \texttt{----??----}\\
4 & 3 & 4 & 1 & 4 & 0 & 1024 & 0 & \texttt{----------}\\
5 & 3 & 4 & 1 & 4 & 0 & 1024 & 0 & \texttt{----------}\\
6 & 1 & 2 & 2 & 2 & 640 & 128 & 256 & \texttt{-------+??}\\
7 & 6 & 3 & 2 & 2 & 256 & 256 & 512 & \texttt{--------??}\\
8 & 7 & 4 & 1 & 4 & 0 & 1024 & 0 & \texttt{----------}\\
9 & 7 & 4 & 1 & 4 & 0 & 1024 & 0 & \texttt{----------}\\
\multicolumn{9}{@{}l@{}}{\rule{0pt}{2.6ex}\textbf{T10.388}\enspace \texttt{0123442343}\enspace $n{=}10$, $|H|{=}4$, len 5, br 3, $|\mathrm{Aut}|{=}2$, orb 9/8, hist-orb 3/2; M 26624, res 7680, rev 10054: ok}\\*
0 & -- & 0 & 4 & 1 & 798 & 2 & 224 & \texttt{-+????????}\\*
1 & 0 & 1 & 4 & 1 & 572 & 4 & 448 & \texttt{--????????}\\
2 & 1 & 2 & 2 & 1 & 640 & 128 & 256 & \texttt{---+??----}\\
3 & 2 & 3 & 2 & 1 & 256 & 256 & 512 & \texttt{----??----}\\
4 & 3 & 4 & 1 & 2 & 0 & 1024 & 0 & \texttt{----------}\\
5 & 3 & 4 & 1 & 2 & 0 & 1024 & 0 & \texttt{----------}\\
6 & 1 & 2 & 2 & 1 & 384 & 128 & 512 & \texttt{-------???}\\
7 & 6 & 3 & 1 & 1 & 512 & 512 & 0 & \texttt{--------+-}\\
8 & 7 & 4 & 1 & 1 & 0 & 1024 & 0 & \texttt{----------}\\
9 & 6 & 3 & 1 & 1 & 0 & 1024 & 0 & \texttt{----------}\\
\multicolumn{9}{@{}l@{}}{\rule{0pt}{2.6ex}\textbf{T10.389}\enspace \texttt{0123442342}\enspace $n{=}10$, $|H|{=}4$, len 5, br 2, $|\mathrm{Aut}|{=}2$, orb 9/8, hist-orb 3/2; M 25600, res 9024, rev 8746: ok}\\*
0 & -- & 0 & 4 & 1 & 694 & 2 & 328 & \texttt{-+????????}\\*
1 & 0 & 1 & 4 & 1 & 364 & 4 & 656 & \texttt{--????????}\\
2 & 1 & 2 & 2 & 1 & 640 & 128 & 256 & \texttt{---+??----}\\
3 & 2 & 3 & 2 & 1 & 256 & 256 & 512 & \texttt{----??----}\\
4 & 3 & 4 & 1 & 2 & 0 & 1024 & 0 & \texttt{----------}\\
5 & 3 & 4 & 1 & 2 & 0 & 1024 & 0 & \texttt{----------}\\
6 & 1 & 2 & 1 & 1 & 768 & 256 & 0 & \texttt{-------++-}\\
7 & 6 & 3 & 1 & 1 & 512 & 512 & 0 & \texttt{--------+-}\\
8 & 7 & 4 & 1 & 1 & 0 & 1024 & 0 & \texttt{----------}\\
9 & 1 & 2 & 1 & 1 & 0 & 1024 & 0 & \texttt{----------}\\
\multicolumn{9}{@{}l@{}}{\rule{0pt}{2.6ex}\textbf{T10.390}\enspace \texttt{0123442341}\enspace $n{=}10$, $|H|{=}4$, len 5, br 3, $|\mathrm{Aut}|{=}2$, orb 9/8, hist-orb 3/2; M 24576, res 8280, rev 8482: ok}\\*
0 & -- & 0 & 4 & 1 & 438 & 2 & 584 & \texttt{-?????????}\\*
1 & 0 & 1 & 3 & 1 & 728 & 8 & 288 & \texttt{--???????-}\\
2 & 1 & 2 & 2 & 1 & 640 & 128 & 256 & \texttt{---+??----}\\
3 & 2 & 3 & 2 & 1 & 256 & 256 & 512 & \texttt{----??----}\\
4 & 3 & 4 & 1 & 2 & 0 & 1024 & 0 & \texttt{----------}\\
5 & 3 & 4 & 1 & 2 & 0 & 1024 & 0 & \texttt{----------}\\
6 & 1 & 2 & 1 & 1 & 768 & 256 & 0 & \texttt{-------++-}\\
7 & 6 & 3 & 1 & 1 & 512 & 512 & 0 & \texttt{--------+-}\\
8 & 7 & 4 & 1 & 1 & 0 & 1024 & 0 & \texttt{----------}\\
9 & 0 & 1 & 1 & 1 & 0 & 1024 & 0 & \texttt{----------}\\
\multicolumn{9}{@{}l@{}}{\rule{0pt}{2.6ex}\textbf{T10.391}\enspace \texttt{0123442333}\enspace $n{=}10$, $|H|{=}5$, len 5, br 3, $|\mathrm{Aut}|{=}12$, orb 7/5, hist-orb 2/0; M 25600, res 9504, rev 8926: ok}\\*
0 & -- & 0 & 5 & 1 & 746 & 2 & 276 & \texttt{-+????????}\\*
1 & 0 & 1 & 5 & 1 & 468 & 4 & 552 & \texttt{--????????}\\
2 & 1 & 2 & 2 & 1 & 640 & 128 & 256 & \texttt{---+??----}\\
3 & 2 & 3 & 2 & 1 & 256 & 256 & 512 & \texttt{----??----}\\
4 & 3 & 4 & 1 & 2 & 0 & 1024 & 0 & \texttt{----------}\\
5 & 3 & 4 & 1 & 2 & 0 & 1024 & 0 & \texttt{----------}\\
6 & 1 & 2 & 3 & 1 & 128 & 128 & 768 & \texttt{-------???}\\
7 & 6 & 3 & 1 & 3 & 0 & 1024 & 0 & \texttt{----------}\\
8 & 6 & 3 & 1 & 3 & 0 & 1024 & 0 & \texttt{----------}\\
9 & 6 & 3 & 1 & 3 & 0 & 1024 & 0 & \texttt{----------}\\
\multicolumn{9}{@{}l@{}}{\rule{0pt}{2.6ex}\textbf{T10.392}\enspace \texttt{0123442332}\enspace $n{=}10$, $|H|{=}5$, len 5, br 3, $|\mathrm{Aut}|{=}4$, orb 8/6, hist-orb 3/1; M 24576, res 10912, rev 7462: ok}\\*
0 & -- & 0 & 5 & 1 & 642 & 2 & 380 & \texttt{-+????????}\\*
1 & 0 & 1 & 5 & 1 & 260 & 4 & 760 & \texttt{--????????}\\
2 & 1 & 2 & 2 & 1 & 640 & 128 & 256 & \texttt{---+??----}\\
3 & 2 & 3 & 2 & 1 & 256 & 256 & 512 & \texttt{----??----}\\
4 & 3 & 4 & 1 & 2 & 0 & 1024 & 0 & \texttt{----------}\\
5 & 3 & 4 & 1 & 2 & 0 & 1024 & 0 & \texttt{----------}\\
6 & 1 & 2 & 2 & 1 & 256 & 256 & 512 & \texttt{-------??-}\\
7 & 6 & 3 & 1 & 2 & 0 & 1024 & 0 & \texttt{----------}\\
8 & 6 & 3 & 1 & 2 & 0 & 1024 & 0 & \texttt{----------}\\
9 & 1 & 2 & 1 & 1 & 0 & 1024 & 0 & \texttt{----------}\\
\multicolumn{9}{@{}l@{}}{\rule{0pt}{2.6ex}\textbf{T10.393}\enspace \texttt{0123442331}\enspace $n{=}10$, $|H|{=}5$, len 5, br 4, $|\mathrm{Aut}|{=}4$, orb 8/6, hist-orb 3/1; M 23552, res 10740, rev 6934: ok}\\*
0 & -- & 0 & 5 & 1 & 386 & 2 & 636 & \texttt{-?????????}\\*
1 & 0 & 1 & 4 & 1 & 520 & 8 & 496 & \texttt{--???????-}\\
2 & 1 & 2 & 2 & 1 & 640 & 128 & 256 & \texttt{---+??----}\\
3 & 2 & 3 & 2 & 1 & 256 & 256 & 512 & \texttt{----??----}\\
4 & 3 & 4 & 1 & 2 & 0 & 1024 & 0 & \texttt{----------}\\
5 & 3 & 4 & 1 & 2 & 0 & 1024 & 0 & \texttt{----------}\\
6 & 1 & 2 & 2 & 1 & 256 & 256 & 512 & \texttt{-------??-}\\
7 & 6 & 3 & 1 & 2 & 0 & 1024 & 0 & \texttt{----------}\\
8 & 6 & 3 & 1 & 2 & 0 & 1024 & 0 & \texttt{----------}\\
9 & 0 & 1 & 1 & 1 & 0 & 1024 & 0 & \texttt{----------}\\
\multicolumn{9}{@{}l@{}}{\rule{0pt}{2.6ex}\textbf{T10.394}\enspace \texttt{0123442323}\enspace $n{=}10$, $|H|{=}4$, len 5, br 2, $|\mathrm{Aut}|{=}4$, orb 7/4, hist-orb 2/0; M 24576, res 7776, rev 9190: ok}\\*
0 & -- & 0 & 4 & 1 & 746 & 2 & 276 & \texttt{-+????????}\\*
1 & 0 & 1 & 4 & 1 & 468 & 4 & 552 & \texttt{--????????}\\
2 & 1 & 2 & 2 & 1 & 640 & 128 & 256 & \texttt{---+??----}\\
3 & 2 & 3 & 2 & 1 & 256 & 256 & 512 & \texttt{----??----}\\
4 & 3 & 4 & 1 & 2 & 0 & 1024 & 0 & \texttt{----------}\\
5 & 3 & 4 & 1 & 2 & 0 & 1024 & 0 & \texttt{----------}\\
6 & 1 & 2 & 1 & 2 & 512 & 512 & 0 & \texttt{-------+--}\\
7 & 6 & 3 & 1 & 2 & 0 & 1024 & 0 & \texttt{----------}\\
8 & 1 & 2 & 1 & 2 & 512 & 512 & 0 & \texttt{---------+}\\
9 & 8 & 3 & 1 & 2 & 0 & 1024 & 0 & \texttt{----------}\\
\multicolumn{9}{@{}l@{}}{\rule{0pt}{2.6ex}\textbf{T10.395}\enspace \texttt{0123442322}\enspace $n{=}10$, $|H|{=}5$, len 5, br 2, $|\mathrm{Aut}|{=}4$, orb 8/6, hist-orb 3/1; M 23552, res 11520, rev 6762: ok}\\*
0 & -- & 0 & 5 & 1 & 590 & 2 & 432 & \texttt{-+????????}\\*
1 & 0 & 1 & 5 & 1 & 156 & 4 & 864 & \texttt{--????????}\\
2 & 1 & 2 & 2 & 1 & 640 & 128 & 256 & \texttt{---+??----}\\
3 & 2 & 3 & 2 & 1 & 256 & 256 & 512 & \texttt{----??----}\\
4 & 3 & 4 & 1 & 2 & 0 & 1024 & 0 & \texttt{----------}\\
5 & 3 & 4 & 1 & 2 & 0 & 1024 & 0 & \texttt{----------}\\
6 & 1 & 2 & 1 & 1 & 512 & 512 & 0 & \texttt{-------+--}\\
7 & 6 & 3 & 1 & 1 & 0 & 1024 & 0 & \texttt{----------}\\
8 & 1 & 2 & 1 & 2 & 0 & 1024 & 0 & \texttt{----------}\\
9 & 1 & 2 & 1 & 2 & 0 & 1024 & 0 & \texttt{----------}\\
\multicolumn{9}{@{}l@{}}{\rule{0pt}{2.6ex}\textbf{T10.396}\enspace \texttt{0123442321}\enspace $n{=}10$, $|H|{=}5$, len 5, br 3, $|\mathrm{Aut}|{=}2$, orb 9/8, hist-orb 4/3; M 22528, res 11920, rev 5826: ok}\\*
0 & -- & 0 & 5 & 1 & 334 & 2 & 688 & \texttt{-?????????}\\*
1 & 0 & 1 & 4 & 1 & 312 & 8 & 704 & \texttt{--???????-}\\
2 & 1 & 2 & 2 & 1 & 640 & 128 & 256 & \texttt{---+??----}\\
3 & 2 & 3 & 2 & 1 & 256 & 256 & 512 & \texttt{----??----}\\
4 & 3 & 4 & 1 & 2 & 0 & 1024 & 0 & \texttt{----------}\\
5 & 3 & 4 & 1 & 2 & 0 & 1024 & 0 & \texttt{----------}\\
6 & 1 & 2 & 1 & 1 & 512 & 512 & 0 & \texttt{-------+--}\\
7 & 6 & 3 & 1 & 1 & 0 & 1024 & 0 & \texttt{----------}\\
8 & 1 & 2 & 1 & 1 & 0 & 1024 & 0 & \texttt{----------}\\
9 & 0 & 1 & 1 & 1 & 0 & 1024 & 0 & \texttt{----------}\\
\multicolumn{9}{@{}l@{}}{\rule{0pt}{2.6ex}\textbf{T10.397}\enspace \texttt{0123442312}\enspace $n{=}10$, $|H|{=}4$, len 5, br 3, $|\mathrm{Aut}|{=}2$, orb 9/8, hist-orb 3/2; M 22528, res 6900, rev 8386: ok}\\*
0 & -- & 0 & 4 & 1 & 618 & 2 & 404 & \texttt{-?????????}\\*
1 & 0 & 1 & 3 & 1 & 624 & 16 & 384 & \texttt{--??????--}\\
2 & 1 & 2 & 2 & 1 & 640 & 128 & 256 & \texttt{---+??----}\\
3 & 2 & 3 & 2 & 1 & 256 & 256 & 512 & \texttt{----??----}\\
4 & 3 & 4 & 1 & 2 & 0 & 1024 & 0 & \texttt{----------}\\
5 & 3 & 4 & 1 & 2 & 0 & 1024 & 0 & \texttt{----------}\\
6 & 1 & 2 & 1 & 1 & 512 & 512 & 0 & \texttt{-------+--}\\
7 & 6 & 3 & 1 & 1 & 0 & 1024 & 0 & \texttt{----------}\\
8 & 0 & 1 & 1 & 1 & 512 & 512 & 0 & \texttt{---------+}\\
9 & 8 & 2 & 1 & 1 & 0 & 1024 & 0 & \texttt{----------}\\
\multicolumn{9}{@{}l@{}}{\rule{0pt}{2.6ex}\textbf{T10.398}\enspace \texttt{0123442311}\enspace $n{=}10$, $|H|{=}5$, len 5, br 3, $|\mathrm{Aut}|{=}4$, orb 8/6, hist-orb 3/1; M 21504, res 10608, rev 5898: ok}\\*
0 & -- & 0 & 5 & 1 & 206 & 2 & 816 & \texttt{-?????????}\\*
1 & 0 & 1 & 3 & 1 & 624 & 16 & 384 & \texttt{--??????--}\\
2 & 1 & 2 & 2 & 1 & 640 & 128 & 256 & \texttt{---+??----}\\
3 & 2 & 3 & 2 & 1 & 256 & 256 & 512 & \texttt{----??----}\\
4 & 3 & 4 & 1 & 2 & 0 & 1024 & 0 & \texttt{----------}\\
5 & 3 & 4 & 1 & 2 & 0 & 1024 & 0 & \texttt{----------}\\
6 & 1 & 2 & 1 & 1 & 512 & 512 & 0 & \texttt{-------+--}\\
7 & 6 & 3 & 1 & 1 & 0 & 1024 & 0 & \texttt{----------}\\
8 & 0 & 1 & 1 & 2 & 0 & 1024 & 0 & \texttt{----------}\\
9 & 0 & 1 & 1 & 2 & 0 & 1024 & 0 & \texttt{----------}\\
\multicolumn{9}{@{}l@{}}{\rule{0pt}{2.6ex}\textbf{T10.399}\enspace \texttt{0123442222}\enspace $n{=}10$, $|H|{=}6$, len 5, br 2, $|\mathrm{Aut}|{=}48$, orb 6/4, hist-orb 2/0; M 22528, res 12768, rev 5678: ok}\\*
0 & -- & 0 & 6 & 1 & 538 & 2 & 484 & \texttt{-+????????}\\*
1 & 0 & 1 & 6 & 1 & 52 & 4 & 968 & \texttt{--????????}\\
2 & 1 & 2 & 2 & 1 & 640 & 128 & 256 & \texttt{---+??----}\\
3 & 2 & 3 & 2 & 1 & 256 & 256 & 512 & \texttt{----??----}\\
4 & 3 & 4 & 1 & 2 & 0 & 1024 & 0 & \texttt{----------}\\
5 & 3 & 4 & 1 & 2 & 0 & 1024 & 0 & \texttt{----------}\\
6 & 1 & 2 & 1 & 4 & 0 & 1024 & 0 & \texttt{----------}\\
7 & 1 & 2 & 1 & 4 & 0 & 1024 & 0 & \texttt{----------}\\
8 & 1 & 2 & 1 & 4 & 0 & 1024 & 0 & \texttt{----------}\\
9 & 1 & 2 & 1 & 4 & 0 & 1024 & 0 & \texttt{----------}\\
\multicolumn{9}{@{}l@{}}{\rule{0pt}{2.6ex}\textbf{T10.400}\enspace \texttt{0123442221}\enspace $n{=}10$, $|H|{=}6$, len 5, br 3, $|\mathrm{Aut}|{=}12$, orb 7/5, hist-orb 3/1; M 21504, res 13740, rev 4270: ok}\\*
0 & -- & 0 & 6 & 1 & 282 & 2 & 740 & \texttt{-?????????}\\*
1 & 0 & 1 & 5 & 1 & 104 & 8 & 912 & \texttt{--???????-}\\
2 & 1 & 2 & 2 & 1 & 640 & 128 & 256 & \texttt{---+??----}\\
3 & 2 & 3 & 2 & 1 & 256 & 256 & 512 & \texttt{----??----}\\
4 & 3 & 4 & 1 & 2 & 0 & 1024 & 0 & \texttt{----------}\\
5 & 3 & 4 & 1 & 2 & 0 & 1024 & 0 & \texttt{----------}\\
6 & 1 & 2 & 1 & 3 & 0 & 1024 & 0 & \texttt{----------}\\
7 & 1 & 2 & 1 & 3 & 0 & 1024 & 0 & \texttt{----------}\\
8 & 1 & 2 & 1 & 3 & 0 & 1024 & 0 & \texttt{----------}\\
9 & 0 & 1 & 1 & 1 & 0 & 1024 & 0 & \texttt{----------}\\
\multicolumn{9}{@{}l@{}}{\rule{0pt}{2.6ex}\textbf{T10.401}\enspace \texttt{0123442212}\enspace $n{=}10$, $|H|{=}5$, len 5, br 3, $|\mathrm{Aut}|{=}4$, orb 8/6, hist-orb 3/1; M 21504, res 10592, rev 5966: ok}\\*
0 & -- & 0 & 5 & 1 & 462 & 2 & 560 & \texttt{-?????????}\\*
1 & 0 & 1 & 4 & 1 & 208 & 16 & 800 & \texttt{--??????--}\\
2 & 1 & 2 & 2 & 1 & 640 & 128 & 256 & \texttt{---+??----}\\
3 & 2 & 3 & 2 & 1 & 256 & 256 & 512 & \texttt{----??----}\\
4 & 3 & 4 & 1 & 2 & 0 & 1024 & 0 & \texttt{----------}\\
5 & 3 & 4 & 1 & 2 & 0 & 1024 & 0 & \texttt{----------}\\
6 & 1 & 2 & 1 & 2 & 0 & 1024 & 0 & \texttt{----------}\\
7 & 1 & 2 & 1 & 2 & 0 & 1024 & 0 & \texttt{----------}\\
8 & 0 & 1 & 1 & 1 & 512 & 512 & 0 & \texttt{---------+}\\
9 & 8 & 2 & 1 & 1 & 0 & 1024 & 0 & \texttt{----------}\\
\multicolumn{9}{@{}l@{}}{\rule{0pt}{2.6ex}\textbf{T10.402}\enspace \texttt{0123442211}\enspace $n{=}10$, $|H|{=}6$, len 5, br 3, $|\mathrm{Aut}|{=}8$, orb 7/4, hist-orb 3/0; M 20480, res 13364, rev 3798: ok}\\*
0 & -- & 0 & 6 & 1 & 154 & 2 & 868 & \texttt{-?????????}\\*
1 & 0 & 1 & 4 & 1 & 208 & 16 & 800 & \texttt{--??????--}\\
2 & 1 & 2 & 2 & 1 & 640 & 128 & 256 & \texttt{---+??----}\\
3 & 2 & 3 & 2 & 1 & 256 & 256 & 512 & \texttt{----??----}\\
4 & 3 & 4 & 1 & 2 & 0 & 1024 & 0 & \texttt{----------}\\
5 & 3 & 4 & 1 & 2 & 0 & 1024 & 0 & \texttt{----------}\\
6 & 1 & 2 & 1 & 2 & 0 & 1024 & 0 & \texttt{----------}\\
7 & 1 & 2 & 1 & 2 & 0 & 1024 & 0 & \texttt{----------}\\
8 & 0 & 1 & 1 & 2 & 0 & 1024 & 0 & \texttt{----------}\\
9 & 0 & 1 & 1 & 2 & 0 & 1024 & 0 & \texttt{----------}\\
\multicolumn{9}{@{}l@{}}{\rule{0pt}{2.6ex}\textbf{T10.403}\enspace \texttt{0123442123}\enspace $n{=}10$, $|H|{=}4$, len 5, br 3, $|\mathrm{Aut}|{=}2$, orb 9/8, hist-orb 3/2; M 22528, res 7272, rev 7950: ok}\\*
0 & -- & 0 & 4 & 1 & 630 & 2 & 392 & \texttt{-?????????}\\*
1 & 0 & 1 & 3 & 1 & 416 & 32 & 576 & \texttt{--?????---}\\
2 & 1 & 2 & 2 & 1 & 640 & 128 & 256 & \texttt{---+??----}\\
3 & 2 & 3 & 2 & 1 & 256 & 256 & 512 & \texttt{----??----}\\
4 & 3 & 4 & 1 & 2 & 0 & 1024 & 0 & \texttt{----------}\\
5 & 3 & 4 & 1 & 2 & 0 & 1024 & 0 & \texttt{----------}\\
6 & 1 & 2 & 1 & 1 & 0 & 1024 & 0 & \texttt{----------}\\
7 & 0 & 1 & 1 & 1 & 768 & 256 & 0 & \texttt{--------++}\\
8 & 7 & 2 & 1 & 1 & 512 & 512 & 0 & \texttt{---------+}\\
9 & 8 & 3 & 1 & 1 & 0 & 1024 & 0 & \texttt{----------}\\
\multicolumn{9}{@{}l@{}}{\rule{0pt}{2.6ex}\textbf{T10.404}\enspace \texttt{0123442122}\enspace $n{=}10$, $|H|{=}5$, len 5, br 4, $|\mathrm{Aut}|{=}4$, orb 8/6, hist-orb 3/1; M 21504, res 9660, rev 6318: ok}\\*
0 & -- & 0 & 5 & 1 & 450 & 2 & 572 & \texttt{-?????????}\\*
1 & 0 & 1 & 3 & 1 & 416 & 32 & 576 & \texttt{--?????---}\\
2 & 1 & 2 & 2 & 1 & 640 & 128 & 256 & \texttt{---+??----}\\
3 & 2 & 3 & 2 & 1 & 256 & 256 & 512 & \texttt{----??----}\\
4 & 3 & 4 & 1 & 2 & 0 & 1024 & 0 & \texttt{----------}\\
5 & 3 & 4 & 1 & 2 & 0 & 1024 & 0 & \texttt{----------}\\
6 & 1 & 2 & 1 & 1 & 0 & 1024 & 0 & \texttt{----------}\\
7 & 0 & 1 & 2 & 1 & 256 & 256 & 512 & \texttt{--------??}\\
8 & 7 & 2 & 1 & 2 & 0 & 1024 & 0 & \texttt{----------}\\
9 & 7 & 2 & 1 & 2 & 0 & 1024 & 0 & \texttt{----------}\\
\multicolumn{9}{@{}l@{}}{\rule{0pt}{2.6ex}\textbf{T10.405}\enspace \texttt{0123442121}\enspace $n{=}10$, $|H|{=}5$, len 5, br 3, $|\mathrm{Aut}|{=}2$, orb 9/8, hist-orb 4/3; M 20480, res 10512, rev 5350: ok}\\*
0 & -- & 0 & 5 & 1 & 270 & 2 & 752 & \texttt{-?????????}\\*
1 & 0 & 1 & 3 & 1 & 416 & 32 & 576 & \texttt{--?????---}\\
2 & 1 & 2 & 2 & 1 & 640 & 128 & 256 & \texttt{---+??----}\\
3 & 2 & 3 & 2 & 1 & 256 & 256 & 512 & \texttt{----??----}\\
4 & 3 & 4 & 1 & 2 & 0 & 1024 & 0 & \texttt{----------}\\
5 & 3 & 4 & 1 & 2 & 0 & 1024 & 0 & \texttt{----------}\\
6 & 1 & 2 & 1 & 1 & 0 & 1024 & 0 & \texttt{----------}\\
7 & 0 & 1 & 1 & 1 & 512 & 512 & 0 & \texttt{--------+-}\\
8 & 7 & 2 & 1 & 1 & 0 & 1024 & 0 & \texttt{----------}\\
9 & 0 & 1 & 1 & 1 & 0 & 1024 & 0 & \texttt{----------}\\
\multicolumn{9}{@{}l@{}}{\rule{0pt}{2.6ex}\textbf{T10.406}\enspace \texttt{0123442111}\enspace $n{=}10$, $|H|{=}6$, len 5, br 3, $|\mathrm{Aut}|{=}12$, orb 7/5, hist-orb 3/1; M 19456, res 12132, rev 3870: ok}\\*
0 & -- & 0 & 6 & 1 & 90 & 2 & 932 & \texttt{-?????????}\\*
1 & 0 & 1 & 3 & 1 & 416 & 32 & 576 & \texttt{--?????---}\\
2 & 1 & 2 & 2 & 1 & 640 & 128 & 256 & \texttt{---+??----}\\
3 & 2 & 3 & 2 & 1 & 256 & 256 & 512 & \texttt{----??----}\\
4 & 3 & 4 & 1 & 2 & 0 & 1024 & 0 & \texttt{----------}\\
5 & 3 & 4 & 1 & 2 & 0 & 1024 & 0 & \texttt{----------}\\
6 & 1 & 2 & 1 & 1 & 0 & 1024 & 0 & \texttt{----------}\\
7 & 0 & 1 & 1 & 3 & 0 & 1024 & 0 & \texttt{----------}\\
8 & 0 & 1 & 1 & 3 & 0 & 1024 & 0 & \texttt{----------}\\
9 & 0 & 1 & 1 & 3 & 0 & 1024 & 0 & \texttt{----------}\\
\multicolumn{9}{@{}l@{}}{\rule{0pt}{2.6ex}\textbf{T10.407}\enspace \texttt{0123441234}\enspace $n{=}10$, $|H|{=}3$, len 5, br 2, $|\mathrm{Aut}|{=}2$, orb 9/8, hist-orb 2/1; M 24576, res 2968, rev 10362: ok}\\*
0 & -- & 0 & 3 & 1 & 870 & 2 & 152 & \texttt{-?????????}\\*
1 & 0 & 1 & 2 & 1 & 832 & 64 & 128 & \texttt{--++??----}\\
2 & 1 & 2 & 2 & 1 & 640 & 128 & 256 & \texttt{---+??----}\\
3 & 2 & 3 & 2 & 1 & 256 & 256 & 512 & \texttt{----??----}\\
4 & 3 & 4 & 1 & 2 & 0 & 1024 & 0 & \texttt{----------}\\
5 & 3 & 4 & 1 & 2 & 0 & 1024 & 0 & \texttt{----------}\\
6 & 0 & 1 & 1 & 1 & 896 & 128 & 0 & \texttt{-------+++}\\
7 & 6 & 2 & 1 & 1 & 768 & 256 & 0 & \texttt{--------++}\\
8 & 7 & 3 & 1 & 1 & 512 & 512 & 0 & \texttt{---------+}\\
9 & 8 & 4 & 1 & 1 & 0 & 1024 & 0 & \texttt{----------}\\
\multicolumn{9}{@{}l@{}}{\rule{0pt}{2.6ex}\textbf{T10.408}\enspace \texttt{0123441233}\enspace $n{=}10$, $|H|{=}4$, len 5, br 3, $|\mathrm{Aut}|{=}4$, orb 8/6, hist-orb 2/0; M 23552, res 5292, rev 9294: ok}\\*
0 & -- & 0 & 4 & 1 & 754 & 2 & 268 & \texttt{-?????????}\\*
1 & 0 & 1 & 2 & 1 & 832 & 64 & 128 & \texttt{--++??----}\\
2 & 1 & 2 & 2 & 1 & 640 & 128 & 256 & \texttt{---+??----}\\
3 & 2 & 3 & 2 & 1 & 256 & 256 & 512 & \texttt{----??----}\\
4 & 3 & 4 & 1 & 2 & 0 & 1024 & 0 & \texttt{----------}\\
5 & 3 & 4 & 1 & 2 & 0 & 1024 & 0 & \texttt{----------}\\
6 & 0 & 1 & 2 & 1 & 640 & 128 & 256 & \texttt{-------+??}\\
7 & 6 & 2 & 2 & 1 & 256 & 256 & 512 & \texttt{--------??}\\
8 & 7 & 3 & 1 & 2 & 0 & 1024 & 0 & \texttt{----------}\\
9 & 7 & 3 & 1 & 2 & 0 & 1024 & 0 & \texttt{----------}\\
\multicolumn{9}{@{}l@{}}{\rule{0pt}{2.6ex}\textbf{T10.409}\enspace \texttt{0123441232}\enspace $n{=}10$, $|H|{=}4$, len 5, br 3, $|\mathrm{Aut}|{=}2$, orb 9/8, hist-orb 3/2; M 22528, res 6336, rev 8390: ok}\\*
0 & -- & 0 & 4 & 1 & 638 & 2 & 384 & \texttt{-?????????}\\*
1 & 0 & 1 & 2 & 1 & 832 & 64 & 128 & \texttt{--++??----}\\
2 & 1 & 2 & 2 & 1 & 640 & 128 & 256 & \texttt{---+??----}\\
3 & 2 & 3 & 2 & 1 & 256 & 256 & 512 & \texttt{----??----}\\
4 & 3 & 4 & 1 & 2 & 0 & 1024 & 0 & \texttt{----------}\\
5 & 3 & 4 & 1 & 2 & 0 & 1024 & 0 & \texttt{----------}\\
6 & 0 & 1 & 2 & 1 & 384 & 128 & 512 & \texttt{-------???}\\
7 & 6 & 2 & 1 & 1 & 512 & 512 & 0 & \texttt{--------+-}\\
8 & 7 & 3 & 1 & 1 & 0 & 1024 & 0 & \texttt{----------}\\
9 & 6 & 2 & 1 & 1 & 0 & 1024 & 0 & \texttt{----------}\\
\multicolumn{9}{@{}l@{}}{\rule{0pt}{2.6ex}\textbf{T10.410}\enspace \texttt{0123441231}\enspace $n{=}10$, $|H|{=}4$, len 5, br 2, $|\mathrm{Aut}|{=}2$, orb 9/8, hist-orb 3/2; M 21504, res 7144, rev 7386: ok}\\*
0 & -- & 0 & 4 & 1 & 406 & 2 & 616 & \texttt{-?????????}\\*
1 & 0 & 1 & 2 & 1 & 832 & 64 & 128 & \texttt{--++??----}\\
2 & 1 & 2 & 2 & 1 & 640 & 128 & 256 & \texttt{---+??----}\\
3 & 2 & 3 & 2 & 1 & 256 & 256 & 512 & \texttt{----??----}\\
4 & 3 & 4 & 1 & 2 & 0 & 1024 & 0 & \texttt{----------}\\
5 & 3 & 4 & 1 & 2 & 0 & 1024 & 0 & \texttt{----------}\\
6 & 0 & 1 & 1 & 1 & 768 & 256 & 0 & \texttt{-------++-}\\
7 & 6 & 2 & 1 & 1 & 512 & 512 & 0 & \texttt{--------+-}\\
8 & 7 & 3 & 1 & 1 & 0 & 1024 & 0 & \texttt{----------}\\
9 & 0 & 1 & 1 & 1 & 0 & 1024 & 0 & \texttt{----------}\\
\multicolumn{9}{@{}l@{}}{\rule{0pt}{2.6ex}\textbf{T10.411}\enspace \texttt{0123441222}\enspace $n{=}10$, $|H|{=}5$, len 5, br 3, $|\mathrm{Aut}|{=}12$, orb 7/5, hist-orb 2/0; M 21504, res 8020, rev 7090: ok}\\*
0 & -- & 0 & 5 & 1 & 522 & 2 & 500 & \texttt{-?????????}\\*
1 & 0 & 1 & 2 & 1 & 832 & 64 & 128 & \texttt{--++??----}\\
2 & 1 & 2 & 2 & 1 & 640 & 128 & 256 & \texttt{---+??----}\\
3 & 2 & 3 & 2 & 1 & 256 & 256 & 512 & \texttt{----??----}\\
4 & 3 & 4 & 1 & 2 & 0 & 1024 & 0 & \texttt{----------}\\
5 & 3 & 4 & 1 & 2 & 0 & 1024 & 0 & \texttt{----------}\\
6 & 0 & 1 & 3 & 1 & 128 & 128 & 768 & \texttt{-------???}\\
7 & 6 & 2 & 1 & 3 & 0 & 1024 & 0 & \texttt{----------}\\
8 & 6 & 2 & 1 & 3 & 0 & 1024 & 0 & \texttt{----------}\\
9 & 6 & 2 & 1 & 3 & 0 & 1024 & 0 & \texttt{----------}\\
\multicolumn{9}{@{}l@{}}{\rule{0pt}{2.6ex}\textbf{T10.412}\enspace \texttt{0123441221}\enspace $n{=}10$, $|H|{=}5$, len 5, br 3, $|\mathrm{Aut}|{=}4$, orb 8/6, hist-orb 3/1; M 20480, res 8956, rev 5970: ok}\\*
0 & -- & 0 & 5 & 1 & 290 & 2 & 732 & \texttt{-?????????}\\*
1 & 0 & 1 & 2 & 1 & 832 & 64 & 128 & \texttt{--++??----}\\
2 & 1 & 2 & 2 & 1 & 640 & 128 & 256 & \texttt{---+??----}\\
3 & 2 & 3 & 2 & 1 & 256 & 256 & 512 & \texttt{----??----}\\
4 & 3 & 4 & 1 & 2 & 0 & 1024 & 0 & \texttt{----------}\\
5 & 3 & 4 & 1 & 2 & 0 & 1024 & 0 & \texttt{----------}\\
6 & 0 & 1 & 2 & 1 & 256 & 256 & 512 & \texttt{-------??-}\\
7 & 6 & 2 & 1 & 2 & 0 & 1024 & 0 & \texttt{----------}\\
8 & 6 & 2 & 1 & 2 & 0 & 1024 & 0 & \texttt{----------}\\
9 & 0 & 1 & 1 & 1 & 0 & 1024 & 0 & \texttt{----------}\\
\multicolumn{9}{@{}l@{}}{\rule{0pt}{2.6ex}\textbf{T10.413}\enspace \texttt{0123441212}\enspace $n{=}10$, $|H|{=}4$, len 5, br 2, $|\mathrm{Aut}|{=}4$, orb 7/4, hist-orb 2/0; M 20480, res 6100, rev 7498: ok}\\*
0 & -- & 0 & 4 & 1 & 522 & 2 & 500 & \texttt{-?????????}\\*
1 & 0 & 1 & 2 & 1 & 832 & 64 & 128 & \texttt{--++??----}\\
2 & 1 & 2 & 2 & 1 & 640 & 128 & 256 & \texttt{---+??----}\\
3 & 2 & 3 & 2 & 1 & 256 & 256 & 512 & \texttt{----??----}\\
4 & 3 & 4 & 1 & 2 & 0 & 1024 & 0 & \texttt{----------}\\
5 & 3 & 4 & 1 & 2 & 0 & 1024 & 0 & \texttt{----------}\\
6 & 0 & 1 & 1 & 2 & 512 & 512 & 0 & \texttt{-------+--}\\
7 & 6 & 2 & 1 & 2 & 0 & 1024 & 0 & \texttt{----------}\\
8 & 0 & 1 & 1 & 2 & 512 & 512 & 0 & \texttt{---------+}\\
9 & 8 & 2 & 1 & 2 & 0 & 1024 & 0 & \texttt{----------}\\
\multicolumn{9}{@{}l@{}}{\rule{0pt}{2.6ex}\textbf{T10.414}\enspace \texttt{0123441211}\enspace $n{=}10$, $|H|{=}5$, len 5, br 2, $|\mathrm{Aut}|{=}4$, orb 8/6, hist-orb 3/1; M 19456, res 9232, rev 5346: ok}\\*
0 & -- & 0 & 5 & 1 & 174 & 2 & 848 & \texttt{-?????????}\\*
1 & 0 & 1 & 2 & 1 & 832 & 64 & 128 & \texttt{--++??----}\\
2 & 1 & 2 & 2 & 1 & 640 & 128 & 256 & \texttt{---+??----}\\
3 & 2 & 3 & 2 & 1 & 256 & 256 & 512 & \texttt{----??----}\\
4 & 3 & 4 & 1 & 2 & 0 & 1024 & 0 & \texttt{----------}\\
5 & 3 & 4 & 1 & 2 & 0 & 1024 & 0 & \texttt{----------}\\
6 & 0 & 1 & 1 & 1 & 512 & 512 & 0 & \texttt{-------+--}\\
7 & 6 & 2 & 1 & 1 & 0 & 1024 & 0 & \texttt{----------}\\
8 & 0 & 1 & 1 & 2 & 0 & 1024 & 0 & \texttt{----------}\\
9 & 0 & 1 & 1 & 2 & 0 & 1024 & 0 & \texttt{----------}\\
\multicolumn{9}{@{}l@{}}{\rule{0pt}{2.6ex}\textbf{T10.415}\enspace \texttt{0123441111}\enspace $n{=}10$, $|H|{=}6$, len 5, br 2, $|\mathrm{Aut}|{=}48$, orb 6/4, hist-orb 2/0; M 18432, res 10276, rev 4210: ok}\\*
0 & -- & 0 & 6 & 1 & 58 & 2 & 964 & \texttt{-?????????}\\*
1 & 0 & 1 & 2 & 1 & 832 & 64 & 128 & \texttt{--++??----}\\
2 & 1 & 2 & 2 & 1 & 640 & 128 & 256 & \texttt{---+??----}\\
3 & 2 & 3 & 2 & 1 & 256 & 256 & 512 & \texttt{----??----}\\
4 & 3 & 4 & 1 & 2 & 0 & 1024 & 0 & \texttt{----------}\\
5 & 3 & 4 & 1 & 2 & 0 & 1024 & 0 & \texttt{----------}\\
6 & 0 & 1 & 1 & 4 & 0 & 1024 & 0 & \texttt{----------}\\
7 & 0 & 1 & 1 & 4 & 0 & 1024 & 0 & \texttt{----------}\\
8 & 0 & 1 & 1 & 4 & 0 & 1024 & 0 & \texttt{----------}\\
9 & 0 & 1 & 1 & 4 & 0 & 1024 & 0 & \texttt{----------}\\
\multicolumn{9}{@{}l@{}}{\rule{0pt}{2.6ex}\textbf{T10.416}\enspace \texttt{0123434343}\enspace $n{=}10$, $|H|{=}4$, len 5, br 1, $|\mathrm{Aut}|{=}6$, orb 6/4, hist-orb 2/1; M 27648, res 9800, rev 10478: ok}\\*
0 & -- & 0 & 4 & 1 & 822 & 2 & 200 & \texttt{-++???????}\\*
1 & 0 & 1 & 4 & 1 & 620 & 4 & 400 & \texttt{--+???????}\\
2 & 1 & 2 & 4 & 1 & 216 & 8 & 800 & \texttt{---???????}\\
3 & 2 & 3 & 1 & 3 & 512 & 512 & 0 & \texttt{----+-----}\\
4 & 3 & 4 & 1 & 3 & 0 & 1024 & 0 & \texttt{----------}\\
5 & 2 & 3 & 1 & 3 & 512 & 512 & 0 & \texttt{------+---}\\
6 & 5 & 4 & 1 & 3 & 0 & 1024 & 0 & \texttt{----------}\\
7 & 2 & 3 & 1 & 3 & 512 & 512 & 0 & \texttt{--------+-}\\
8 & 7 & 4 & 1 & 3 & 0 & 1024 & 0 & \texttt{----------}\\
9 & 2 & 3 & 1 & 1 & 0 & 1024 & 0 & \texttt{----------}\\
\multicolumn{9}{@{}l@{}}{\rule{0pt}{2.6ex}\textbf{T10.417}\enspace \texttt{0123434342}\enspace $n{=}10$, $|H|{=}4$, len 5, br 2, $|\mathrm{Aut}|{=}6$, orb 6/4, hist-orb 2/1; M 26624, res 10896, rev 9018: ok}\\*
0 & -- & 0 & 4 & 1 & 694 & 2 & 328 & \texttt{-+????????}\\*
1 & 0 & 1 & 4 & 1 & 364 & 4 & 656 & \texttt{--????????}\\
2 & 1 & 2 & 3 & 1 & 432 & 16 & 576 & \texttt{---??????-}\\
3 & 2 & 3 & 1 & 3 & 512 & 512 & 0 & \texttt{----+-----}\\
4 & 3 & 4 & 1 & 3 & 0 & 1024 & 0 & \texttt{----------}\\
5 & 2 & 3 & 1 & 3 & 512 & 512 & 0 & \texttt{------+---}\\
6 & 5 & 4 & 1 & 3 & 0 & 1024 & 0 & \texttt{----------}\\
7 & 2 & 3 & 1 & 3 & 512 & 512 & 0 & \texttt{--------+-}\\
8 & 7 & 4 & 1 & 3 & 0 & 1024 & 0 & \texttt{----------}\\
9 & 1 & 2 & 1 & 1 & 0 & 1024 & 0 & \texttt{----------}\\
\multicolumn{9}{@{}l@{}}{\rule{0pt}{2.6ex}\textbf{T10.418}\enspace \texttt{0123434341}\enspace $n{=}10$, $|H|{=}4$, len 5, br 2, $|\mathrm{Aut}|{=}6$, orb 6/4, hist-orb 2/1; M 25600, res 10152, rev 8830: ok}\\*
0 & -- & 0 & 4 & 1 & 438 & 2 & 584 & \texttt{-?????????}\\*
1 & 0 & 1 & 3 & 1 & 728 & 8 & 288 & \texttt{--+??????-}\\
2 & 1 & 2 & 3 & 1 & 432 & 16 & 576 & \texttt{---??????-}\\
3 & 2 & 3 & 1 & 3 & 512 & 512 & 0 & \texttt{----+-----}\\
4 & 3 & 4 & 1 & 3 & 0 & 1024 & 0 & \texttt{----------}\\
5 & 2 & 3 & 1 & 3 & 512 & 512 & 0 & \texttt{------+---}\\
6 & 5 & 4 & 1 & 3 & 0 & 1024 & 0 & \texttt{----------}\\
7 & 2 & 3 & 1 & 3 & 512 & 512 & 0 & \texttt{--------+-}\\
8 & 7 & 4 & 1 & 3 & 0 & 1024 & 0 & \texttt{----------}\\
9 & 0 & 1 & 1 & 1 & 0 & 1024 & 0 & \texttt{----------}\\
\multicolumn{9}{@{}l@{}}{\rule{0pt}{2.6ex}\textbf{T10.419}\enspace \texttt{0123434333}\enspace $n{=}10$, $|H|{=}5$, len 5, br 1, $|\mathrm{Aut}|{=}12$, orb 6/3, hist-orb 2/0; M 26624, res 11564, rev 9430: ok}\\*
0 & -- & 0 & 5 & 1 & 786 & 2 & 236 & \texttt{-++???????}\\*
1 & 0 & 1 & 5 & 1 & 548 & 4 & 472 & \texttt{--+???????}\\
2 & 1 & 2 & 5 & 1 & 72 & 8 & 944 & \texttt{---???????}\\
3 & 2 & 3 & 1 & 2 & 512 & 512 & 0 & \texttt{----+-----}\\
4 & 3 & 4 & 1 & 2 & 0 & 1024 & 0 & \texttt{----------}\\
5 & 2 & 3 & 1 & 2 & 512 & 512 & 0 & \texttt{------+---}\\
6 & 5 & 4 & 1 & 2 & 0 & 1024 & 0 & \texttt{----------}\\
7 & 2 & 3 & 1 & 3 & 0 & 1024 & 0 & \texttt{----------}\\
8 & 2 & 3 & 1 & 3 & 0 & 1024 & 0 & \texttt{----------}\\
9 & 2 & 3 & 1 & 3 & 0 & 1024 & 0 & \texttt{----------}\\
\multicolumn{9}{@{}l@{}}{\rule{0pt}{2.6ex}\textbf{T10.420}\enspace \texttt{0123434332}\enspace $n{=}10$, $|H|{=}5$, len 5, br 2, $|\mathrm{Aut}|{=}4$, orb 7/4, hist-orb 3/1; M 25600, res 13272, rev 7394: ok}\\*
0 & -- & 0 & 5 & 1 & 658 & 2 & 364 & \texttt{-+????????}\\*
1 & 0 & 1 & 5 & 1 & 292 & 4 & 728 & \texttt{--????????}\\
2 & 1 & 2 & 4 & 1 & 144 & 16 & 864 & \texttt{---??????-}\\
3 & 2 & 3 & 1 & 2 & 512 & 512 & 0 & \texttt{----+-----}\\
4 & 3 & 4 & 1 & 2 & 0 & 1024 & 0 & \texttt{----------}\\
5 & 2 & 3 & 1 & 2 & 512 & 512 & 0 & \texttt{------+---}\\
6 & 5 & 4 & 1 & 2 & 0 & 1024 & 0 & \texttt{----------}\\
7 & 2 & 3 & 1 & 2 & 0 & 1024 & 0 & \texttt{----------}\\
8 & 2 & 3 & 1 & 2 & 0 & 1024 & 0 & \texttt{----------}\\
9 & 1 & 2 & 1 & 1 & 0 & 1024 & 0 & \texttt{----------}\\
\multicolumn{9}{@{}l@{}}{\rule{0pt}{2.6ex}\textbf{T10.421}\enspace \texttt{0123434331}\enspace $n{=}10$, $|H|{=}5$, len 5, br 2, $|\mathrm{Aut}|{=}4$, orb 7/4, hist-orb 3/1; M 24576, res 12924, rev 7014: ok}\\*
0 & -- & 0 & 5 & 1 & 402 & 2 & 620 & \texttt{-?????????}\\*
1 & 0 & 1 & 4 & 1 & 584 & 8 & 432 & \texttt{--+??????-}\\
2 & 1 & 2 & 4 & 1 & 144 & 16 & 864 & \texttt{---??????-}\\
3 & 2 & 3 & 1 & 2 & 512 & 512 & 0 & \texttt{----+-----}\\
4 & 3 & 4 & 1 & 2 & 0 & 1024 & 0 & \texttt{----------}\\
5 & 2 & 3 & 1 & 2 & 512 & 512 & 0 & \texttt{------+---}\\
6 & 5 & 4 & 1 & 2 & 0 & 1024 & 0 & \texttt{----------}\\
7 & 2 & 3 & 1 & 2 & 0 & 1024 & 0 & \texttt{----------}\\
8 & 2 & 3 & 1 & 2 & 0 & 1024 & 0 & \texttt{----------}\\
9 & 0 & 1 & 1 & 1 & 0 & 1024 & 0 & \texttt{----------}\\
\multicolumn{9}{@{}l@{}}{\rule{0pt}{2.6ex}\textbf{T10.422}\enspace \texttt{0123434323}\enspace $n{=}10$, $|H|{=}4$, len 5, br 2, $|\mathrm{Aut}|{=}2$, orb 8/6, hist-orb 3/2; M 25600, res 9328, rev 9326: ok}\\*
0 & -- & 0 & 4 & 1 & 758 & 2 & 264 & \texttt{-+????????}\\*
1 & 0 & 1 & 4 & 1 & 492 & 4 & 528 & \texttt{--????????}\\
2 & 1 & 2 & 3 & 1 & 288 & 32 & 704 & \texttt{---?????--}\\
3 & 2 & 3 & 1 & 2 & 512 & 512 & 0 & \texttt{----+-----}\\
4 & 3 & 4 & 1 & 2 & 0 & 1024 & 0 & \texttt{----------}\\
5 & 2 & 3 & 1 & 2 & 512 & 512 & 0 & \texttt{------+---}\\
6 & 5 & 4 & 1 & 2 & 0 & 1024 & 0 & \texttt{----------}\\
7 & 2 & 3 & 1 & 1 & 0 & 1024 & 0 & \texttt{----------}\\
8 & 1 & 2 & 1 & 1 & 512 & 512 & 0 & \texttt{---------+}\\
9 & 8 & 3 & 1 & 1 & 0 & 1024 & 0 & \texttt{----------}\\
\multicolumn{9}{@{}l@{}}{\rule{0pt}{2.6ex}\textbf{T10.423}\enspace \texttt{0123434322}\enspace $n{=}10$, $|H|{=}5$, len 5, br 2, $|\mathrm{Aut}|{=}4$, orb 7/4, hist-orb 3/1; M 24576, res 13264, rev 6634: ok}\\*
0 & -- & 0 & 5 & 1 & 594 & 2 & 428 & \texttt{-+????????}\\*
1 & 0 & 1 & 5 & 1 & 164 & 4 & 856 & \texttt{--????????}\\
2 & 1 & 2 & 3 & 1 & 288 & 32 & 704 & \texttt{---?????--}\\
3 & 2 & 3 & 1 & 2 & 512 & 512 & 0 & \texttt{----+-----}\\
4 & 3 & 4 & 1 & 2 & 0 & 1024 & 0 & \texttt{----------}\\
5 & 2 & 3 & 1 & 2 & 512 & 512 & 0 & \texttt{------+---}\\
6 & 5 & 4 & 1 & 2 & 0 & 1024 & 0 & \texttt{----------}\\
7 & 2 & 3 & 1 & 1 & 0 & 1024 & 0 & \texttt{----------}\\
8 & 1 & 2 & 1 & 2 & 0 & 1024 & 0 & \texttt{----------}\\
9 & 1 & 2 & 1 & 2 & 0 & 1024 & 0 & \texttt{----------}\\
\multicolumn{9}{@{}l@{}}{\rule{0pt}{2.6ex}\textbf{T10.424}\enspace \texttt{0123434321}\enspace $n{=}10$, $|H|{=}5$, len 5, br 3, $|\mathrm{Aut}|{=}2$, orb 8/6, hist-orb 4/3; M 23552, res 13620, rev 5738: ok}\\*
0 & -- & 0 & 5 & 1 & 338 & 2 & 684 & \texttt{-?????????}\\*
1 & 0 & 1 & 4 & 1 & 328 & 8 & 688 & \texttt{--???????-}\\
2 & 1 & 2 & 3 & 1 & 288 & 32 & 704 & \texttt{---?????--}\\
3 & 2 & 3 & 1 & 2 & 512 & 512 & 0 & \texttt{----+-----}\\
4 & 3 & 4 & 1 & 2 & 0 & 1024 & 0 & \texttt{----------}\\
5 & 2 & 3 & 1 & 2 & 512 & 512 & 0 & \texttt{------+---}\\
6 & 5 & 4 & 1 & 2 & 0 & 1024 & 0 & \texttt{----------}\\
7 & 2 & 3 & 1 & 1 & 0 & 1024 & 0 & \texttt{----------}\\
8 & 1 & 2 & 1 & 1 & 0 & 1024 & 0 & \texttt{----------}\\
9 & 0 & 1 & 1 & 1 & 0 & 1024 & 0 & \texttt{----------}\\
\multicolumn{9}{@{}l@{}}{\rule{0pt}{2.6ex}\textbf{T10.425}\enspace \texttt{0123434312}\enspace $n{=}10$, $|H|{=}4$, len 5, br 2, $|\mathrm{Aut}|{=}2$, orb 8/6, hist-orb 3/2; M 23552, res 8456, rev 8550: ok}\\*
0 & -- & 0 & 4 & 1 & 630 & 2 & 392 & \texttt{-?????????}\\*
1 & 0 & 1 & 3 & 1 & 656 & 16 & 352 & \texttt{--+?????--}\\
2 & 1 & 2 & 3 & 1 & 288 & 32 & 704 & \texttt{---?????--}\\
3 & 2 & 3 & 1 & 2 & 512 & 512 & 0 & \texttt{----+-----}\\
4 & 3 & 4 & 1 & 2 & 0 & 1024 & 0 & \texttt{----------}\\
5 & 2 & 3 & 1 & 2 & 512 & 512 & 0 & \texttt{------+---}\\
6 & 5 & 4 & 1 & 2 & 0 & 1024 & 0 & \texttt{----------}\\
7 & 2 & 3 & 1 & 1 & 0 & 1024 & 0 & \texttt{----------}\\
8 & 0 & 1 & 1 & 1 & 512 & 512 & 0 & \texttt{---------+}\\
9 & 8 & 2 & 1 & 1 & 0 & 1024 & 0 & \texttt{----------}\\
\multicolumn{9}{@{}l@{}}{\rule{0pt}{2.6ex}\textbf{T10.426}\enspace \texttt{0123434311}\enspace $n{=}10$, $|H|{=}5$, len 5, br 2, $|\mathrm{Aut}|{=}4$, orb 7/4, hist-orb 3/1; M 22528, res 12236, rev 5934: ok}\\*
0 & -- & 0 & 5 & 1 & 210 & 2 & 812 & \texttt{-?????????}\\*
1 & 0 & 1 & 3 & 1 & 656 & 16 & 352 & \texttt{--+?????--}\\
2 & 1 & 2 & 3 & 1 & 288 & 32 & 704 & \texttt{---?????--}\\
3 & 2 & 3 & 1 & 2 & 512 & 512 & 0 & \texttt{----+-----}\\
4 & 3 & 4 & 1 & 2 & 0 & 1024 & 0 & \texttt{----------}\\
5 & 2 & 3 & 1 & 2 & 512 & 512 & 0 & \texttt{------+---}\\
6 & 5 & 4 & 1 & 2 & 0 & 1024 & 0 & \texttt{----------}\\
7 & 2 & 3 & 1 & 1 & 0 & 1024 & 0 & \texttt{----------}\\
8 & 0 & 1 & 1 & 2 & 0 & 1024 & 0 & \texttt{----------}\\
9 & 0 & 1 & 1 & 2 & 0 & 1024 & 0 & \texttt{----------}\\
\multicolumn{9}{@{}l@{}}{\rule{0pt}{2.6ex}\textbf{T10.427}\enspace \texttt{0123434234}\enspace $n{=}10$, $|H|{=}3$, len 5, br 2, $|\mathrm{Aut}|{=}2$, orb 8/6, hist-orb 2/1; M 26624, res 5088, rev 11278: ok}\\*
0 & -- & 0 & 3 & 1 & 862 & 2 & 160 & \texttt{-+????????}\\*
1 & 0 & 1 & 3 & 1 & 700 & 4 & 320 & \texttt{--????????}\\
2 & 1 & 2 & 2 & 1 & 576 & 64 & 384 & \texttt{---????---}\\
3 & 2 & 3 & 1 & 2 & 512 & 512 & 0 & \texttt{----+-----}\\
4 & 3 & 4 & 1 & 2 & 0 & 1024 & 0 & \texttt{----------}\\
5 & 2 & 3 & 1 & 2 & 512 & 512 & 0 & \texttt{------+---}\\
6 & 5 & 4 & 1 & 2 & 0 & 1024 & 0 & \texttt{----------}\\
7 & 1 & 2 & 1 & 1 & 768 & 256 & 0 & \texttt{--------++}\\
8 & 7 & 3 & 1 & 1 & 512 & 512 & 0 & \texttt{---------+}\\
9 & 8 & 4 & 1 & 1 & 0 & 1024 & 0 & \texttt{----------}\\
\multicolumn{9}{@{}l@{}}{\rule{0pt}{2.6ex}\textbf{T10.428}\enspace \texttt{0123434233}\enspace $n{=}10$, $|H|{=}4$, len 5, br 3, $|\mathrm{Aut}|{=}4$, orb 7/4, hist-orb 2/0; M 25600, res 8128, rev 9634: ok}\\*
0 & -- & 0 & 4 & 1 & 762 & 2 & 260 & \texttt{-+????????}\\*
1 & 0 & 1 & 4 & 1 & 500 & 4 & 520 & \texttt{--????????}\\
2 & 1 & 2 & 2 & 1 & 576 & 64 & 384 & \texttt{---????---}\\
3 & 2 & 3 & 1 & 2 & 512 & 512 & 0 & \texttt{----+-----}\\
4 & 3 & 4 & 1 & 2 & 0 & 1024 & 0 & \texttt{----------}\\
5 & 2 & 3 & 1 & 2 & 512 & 512 & 0 & \texttt{------+---}\\
6 & 5 & 4 & 1 & 2 & 0 & 1024 & 0 & \texttt{----------}\\
7 & 1 & 2 & 2 & 1 & 256 & 256 & 512 & \texttt{--------??}\\
8 & 7 & 3 & 1 & 2 & 0 & 1024 & 0 & \texttt{----------}\\
9 & 7 & 3 & 1 & 2 & 0 & 1024 & 0 & \texttt{----------}\\
\multicolumn{9}{@{}l@{}}{\rule{0pt}{2.6ex}\textbf{T10.429}\enspace \texttt{0123434232}\enspace $n{=}10$, $|H|{=}4$, len 5, br 2, $|\mathrm{Aut}|{=}2$, orb 8/6, hist-orb 3/2; M 24576, res 9888, rev 8286: ok}\\*
0 & -- & 0 & 4 & 1 & 662 & 2 & 360 & \texttt{-+????????}\\*
1 & 0 & 1 & 4 & 1 & 300 & 4 & 720 & \texttt{--????????}\\
2 & 1 & 2 & 2 & 1 & 576 & 64 & 384 & \texttt{---????---}\\
3 & 2 & 3 & 1 & 2 & 512 & 512 & 0 & \texttt{----+-----}\\
4 & 3 & 4 & 1 & 2 & 0 & 1024 & 0 & \texttt{----------}\\
5 & 2 & 3 & 1 & 2 & 512 & 512 & 0 & \texttt{------+---}\\
6 & 5 & 4 & 1 & 2 & 0 & 1024 & 0 & \texttt{----------}\\
7 & 1 & 2 & 1 & 1 & 512 & 512 & 0 & \texttt{--------+-}\\
8 & 7 & 3 & 1 & 1 & 0 & 1024 & 0 & \texttt{----------}\\
9 & 1 & 2 & 1 & 1 & 0 & 1024 & 0 & \texttt{----------}\\
\multicolumn{9}{@{}l@{}}{\rule{0pt}{2.6ex}\textbf{T10.430}\enspace \texttt{0123434231}\enspace $n{=}10$, $|H|{=}4$, len 5, br 3, $|\mathrm{Aut}|{=}2$, orb 8/6, hist-orb 3/2; M 23552, res 9496, rev 7838: ok}\\*
0 & -- & 0 & 4 & 1 & 406 & 2 & 616 & \texttt{-?????????}\\*
1 & 0 & 1 & 3 & 1 & 600 & 8 & 416 & \texttt{--???????-}\\
2 & 1 & 2 & 2 & 1 & 576 & 64 & 384 & \texttt{---????---}\\
3 & 2 & 3 & 1 & 2 & 512 & 512 & 0 & \texttt{----+-----}\\
4 & 3 & 4 & 1 & 2 & 0 & 1024 & 0 & \texttt{----------}\\
5 & 2 & 3 & 1 & 2 & 512 & 512 & 0 & \texttt{------+---}\\
6 & 5 & 4 & 1 & 2 & 0 & 1024 & 0 & \texttt{----------}\\
7 & 1 & 2 & 1 & 1 & 512 & 512 & 0 & \texttt{--------+-}\\
8 & 7 & 3 & 1 & 1 & 0 & 1024 & 0 & \texttt{----------}\\
9 & 0 & 1 & 1 & 1 & 0 & 1024 & 0 & \texttt{----------}\\
\multicolumn{9}{@{}l@{}}{\rule{0pt}{2.6ex}\textbf{T10.431}\enspace \texttt{0123434222}\enspace $n{=}10$, $|H|{=}5$, len 5, br 2, $|\mathrm{Aut}|{=}12$, orb 6/3, hist-orb 2/0; M 23552, res 12288, rev 6554: ok}\\*
0 & -- & 0 & 5 & 1 & 562 & 2 & 460 & \texttt{-+????????}\\*
1 & 0 & 1 & 5 & 1 & 100 & 4 & 920 & \texttt{--????????}\\
2 & 1 & 2 & 2 & 1 & 576 & 64 & 384 & \texttt{---????---}\\
3 & 2 & 3 & 1 & 2 & 512 & 512 & 0 & \texttt{----+-----}\\
4 & 3 & 4 & 1 & 2 & 0 & 1024 & 0 & \texttt{----------}\\
5 & 2 & 3 & 1 & 2 & 512 & 512 & 0 & \texttt{------+---}\\
6 & 5 & 4 & 1 & 2 & 0 & 1024 & 0 & \texttt{----------}\\
7 & 1 & 2 & 1 & 3 & 0 & 1024 & 0 & \texttt{----------}\\
8 & 1 & 2 & 1 & 3 & 0 & 1024 & 0 & \texttt{----------}\\
9 & 1 & 2 & 1 & 3 & 0 & 1024 & 0 & \texttt{----------}\\
\multicolumn{9}{@{}l@{}}{\rule{0pt}{2.6ex}\textbf{T10.432}\enspace \texttt{0123434221}\enspace $n{=}10$, $|H|{=}5$, len 5, br 3, $|\mathrm{Aut}|{=}4$, orb 7/4, hist-orb 3/1; M 22528, res 12996, rev 5394: ok}\\*
0 & -- & 0 & 5 & 1 & 306 & 2 & 716 & \texttt{-?????????}\\*
1 & 0 & 1 & 4 & 1 & 200 & 8 & 816 & \texttt{--???????-}\\
2 & 1 & 2 & 2 & 1 & 576 & 64 & 384 & \texttt{---????---}\\
3 & 2 & 3 & 1 & 2 & 512 & 512 & 0 & \texttt{----+-----}\\
4 & 3 & 4 & 1 & 2 & 0 & 1024 & 0 & \texttt{----------}\\
5 & 2 & 3 & 1 & 2 & 512 & 512 & 0 & \texttt{------+---}\\
6 & 5 & 4 & 1 & 2 & 0 & 1024 & 0 & \texttt{----------}\\
7 & 1 & 2 & 1 & 2 & 0 & 1024 & 0 & \texttt{----------}\\
8 & 1 & 2 & 1 & 2 & 0 & 1024 & 0 & \texttt{----------}\\
9 & 0 & 1 & 1 & 1 & 0 & 1024 & 0 & \texttt{----------}\\
\multicolumn{9}{@{}l@{}}{\rule{0pt}{2.6ex}\textbf{T10.433}\enspace \texttt{0123434212}\enspace $n{=}10$, $|H|{=}4$, len 5, br 3, $|\mathrm{Aut}|{=}2$, orb 8/6, hist-orb 3/2; M 22528, res 8984, rev 7506: ok}\\*
0 & -- & 0 & 4 & 1 & 534 & 2 & 488 & \texttt{-?????????}\\*
1 & 0 & 1 & 3 & 1 & 400 & 16 & 608 & \texttt{--??????--}\\
2 & 1 & 2 & 2 & 1 & 576 & 64 & 384 & \texttt{---????---}\\
3 & 2 & 3 & 1 & 2 & 512 & 512 & 0 & \texttt{----+-----}\\
4 & 3 & 4 & 1 & 2 & 0 & 1024 & 0 & \texttt{----------}\\
5 & 2 & 3 & 1 & 2 & 512 & 512 & 0 & \texttt{------+---}\\
6 & 5 & 4 & 1 & 2 & 0 & 1024 & 0 & \texttt{----------}\\
7 & 1 & 2 & 1 & 1 & 0 & 1024 & 0 & \texttt{----------}\\
8 & 0 & 1 & 1 & 1 & 512 & 512 & 0 & \texttt{---------+}\\
9 & 8 & 2 & 1 & 1 & 0 & 1024 & 0 & \texttt{----------}\\
\multicolumn{9}{@{}l@{}}{\rule{0pt}{2.6ex}\textbf{T10.434}\enspace \texttt{0123434211}\enspace $n{=}10$, $|H|{=}5$, len 5, br 3, $|\mathrm{Aut}|{=}4$, orb 7/4, hist-orb 3/1; M 21504, res 12188, rev 5202: ok}\\*
0 & -- & 0 & 5 & 1 & 178 & 2 & 844 & \texttt{-?????????}\\*
1 & 0 & 1 & 3 & 1 & 400 & 16 & 608 & \texttt{--??????--}\\
2 & 1 & 2 & 2 & 1 & 576 & 64 & 384 & \texttt{---????---}\\
3 & 2 & 3 & 1 & 2 & 512 & 512 & 0 & \texttt{----+-----}\\
4 & 3 & 4 & 1 & 2 & 0 & 1024 & 0 & \texttt{----------}\\
5 & 2 & 3 & 1 & 2 & 512 & 512 & 0 & \texttt{------+---}\\
6 & 5 & 4 & 1 & 2 & 0 & 1024 & 0 & \texttt{----------}\\
7 & 1 & 2 & 1 & 1 & 0 & 1024 & 0 & \texttt{----------}\\
8 & 0 & 1 & 1 & 2 & 0 & 1024 & 0 & \texttt{----------}\\
9 & 0 & 1 & 1 & 2 & 0 & 1024 & 0 & \texttt{----------}\\
\multicolumn{9}{@{}l@{}}{\rule{0pt}{2.6ex}\textbf{T10.435}\enspace \texttt{0123434123}\enspace $n{=}10$, $|H|{=}3$, len 5, br 2, $|\mathrm{Aut}|{=}2$, orb 8/6, hist-orb 2/1; M 23552, res 4128, rev 10094: ok}\\*
0 & -- & 0 & 3 & 1 & 798 & 2 & 224 & \texttt{-?????????}\\*
1 & 0 & 1 & 2 & 1 & 800 & 32 & 192 & \texttt{--+????---}\\
2 & 1 & 2 & 2 & 1 & 576 & 64 & 384 & \texttt{---????---}\\
3 & 2 & 3 & 1 & 2 & 512 & 512 & 0 & \texttt{----+-----}\\
4 & 3 & 4 & 1 & 2 & 0 & 1024 & 0 & \texttt{----------}\\
5 & 2 & 3 & 1 & 2 & 512 & 512 & 0 & \texttt{------+---}\\
6 & 5 & 4 & 1 & 2 & 0 & 1024 & 0 & \texttt{----------}\\
7 & 0 & 1 & 1 & 1 & 768 & 256 & 0 & \texttt{--------++}\\
8 & 7 & 2 & 1 & 1 & 512 & 512 & 0 & \texttt{---------+}\\
9 & 8 & 3 & 1 & 1 & 0 & 1024 & 0 & \texttt{----------}\\
\multicolumn{9}{@{}l@{}}{\rule{0pt}{2.6ex}\textbf{T10.436}\enspace \texttt{0123434122}\enspace $n{=}10$, $|H|{=}4$, len 5, br 3, $|\mathrm{Aut}|{=}4$, orb 7/4, hist-orb 2/0; M 22528, res 6948, rev 8382: ok}\\*
0 & -- & 0 & 4 & 1 & 570 & 2 & 452 & \texttt{-?????????}\\*
1 & 0 & 1 & 2 & 1 & 800 & 32 & 192 & \texttt{--+????---}\\
2 & 1 & 2 & 2 & 1 & 576 & 64 & 384 & \texttt{---????---}\\
3 & 2 & 3 & 1 & 2 & 512 & 512 & 0 & \texttt{----+-----}\\
4 & 3 & 4 & 1 & 2 & 0 & 1024 & 0 & \texttt{----------}\\
5 & 2 & 3 & 1 & 2 & 512 & 512 & 0 & \texttt{------+---}\\
6 & 5 & 4 & 1 & 2 & 0 & 1024 & 0 & \texttt{----------}\\
7 & 0 & 1 & 2 & 1 & 256 & 256 & 512 & \texttt{--------??}\\
8 & 7 & 2 & 1 & 2 & 0 & 1024 & 0 & \texttt{----------}\\
9 & 7 & 2 & 1 & 2 & 0 & 1024 & 0 & \texttt{----------}\\
\multicolumn{9}{@{}l@{}}{\rule{0pt}{2.6ex}\textbf{T10.437}\enspace \texttt{0123434121}\enspace $n{=}10$, $|H|{=}4$, len 5, br 2, $|\mathrm{Aut}|{=}2$, orb 8/6, hist-orb 3/2; M 21504, res 8232, rev 7238: ok}\\*
0 & -- & 0 & 4 & 1 & 342 & 2 & 680 & \texttt{-?????????}\\*
1 & 0 & 1 & 2 & 1 & 800 & 32 & 192 & \texttt{--+????---}\\
2 & 1 & 2 & 2 & 1 & 576 & 64 & 384 & \texttt{---????---}\\
3 & 2 & 3 & 1 & 2 & 512 & 512 & 0 & \texttt{----+-----}\\
4 & 3 & 4 & 1 & 2 & 0 & 1024 & 0 & \texttt{----------}\\
5 & 2 & 3 & 1 & 2 & 512 & 512 & 0 & \texttt{------+---}\\
6 & 5 & 4 & 1 & 2 & 0 & 1024 & 0 & \texttt{----------}\\
7 & 0 & 1 & 1 & 1 & 512 & 512 & 0 & \texttt{--------+-}\\
8 & 7 & 2 & 1 & 1 & 0 & 1024 & 0 & \texttt{----------}\\
9 & 0 & 1 & 1 & 1 & 0 & 1024 & 0 & \texttt{----------}\\
\multicolumn{9}{@{}l@{}}{\rule{0pt}{2.6ex}\textbf{T10.438}\enspace \texttt{0123434111}\enspace $n{=}10$, $|H|{=}5$, len 5, br 2, $|\mathrm{Aut}|{=}12$, orb 6/3, hist-orb 2/0; M 20480, res 10284, rev 5582: ok}\\*
0 & -- & 0 & 5 & 1 & 114 & 2 & 908 & \texttt{-?????????}\\*
1 & 0 & 1 & 2 & 1 & 800 & 32 & 192 & \texttt{--+????---}\\
2 & 1 & 2 & 2 & 1 & 576 & 64 & 384 & \texttt{---????---}\\
3 & 2 & 3 & 1 & 2 & 512 & 512 & 0 & \texttt{----+-----}\\
4 & 3 & 4 & 1 & 2 & 0 & 1024 & 0 & \texttt{----------}\\
5 & 2 & 3 & 1 & 2 & 512 & 512 & 0 & \texttt{------+---}\\
6 & 5 & 4 & 1 & 2 & 0 & 1024 & 0 & \texttt{----------}\\
7 & 0 & 1 & 1 & 3 & 0 & 1024 & 0 & \texttt{----------}\\
8 & 0 & 1 & 1 & 3 & 0 & 1024 & 0 & \texttt{----------}\\
9 & 0 & 1 & 1 & 3 & 0 & 1024 & 0 & \texttt{----------}\\
\multicolumn{9}{@{}l@{}}{\rule{0pt}{2.6ex}\textbf{T10.439}\enspace \texttt{0123433333}\enspace $n{=}10$, $|H|{=}6$, len 5, br 1, $|\mathrm{Aut}|{=}120$, orb 6/5, hist-orb 2/1; M 25600, res 12152, rev 9110: ok}\\*
0 & -- & 0 & 6 & 1 & 774 & 2 & 248 & \texttt{-++???????}\\*
1 & 0 & 1 & 6 & 1 & 524 & 4 & 496 & \texttt{--+???????}\\
2 & 1 & 2 & 6 & 1 & 24 & 8 & 992 & \texttt{---???????}\\
3 & 2 & 3 & 1 & 1 & 512 & 512 & 0 & \texttt{----+-----}\\
4 & 3 & 4 & 1 & 1 & 0 & 1024 & 0 & \texttt{----------}\\
5 & 2 & 3 & 1 & 5 & 0 & 1024 & 0 & \texttt{----------}\\
6 & 2 & 3 & 1 & 5 & 0 & 1024 & 0 & \texttt{----------}\\
7 & 2 & 3 & 1 & 5 & 0 & 1024 & 0 & \texttt{----------}\\
8 & 2 & 3 & 1 & 5 & 0 & 1024 & 0 & \texttt{----------}\\
9 & 2 & 3 & 1 & 5 & 0 & 1024 & 0 & \texttt{----------}\\
\multicolumn{9}{@{}l@{}}{\rule{0pt}{2.6ex}\textbf{T10.440}\enspace \texttt{0123433332}\enspace $n{=}10$, $|H|{=}6$, len 5, br 2, $|\mathrm{Aut}|{=}24$, orb 7/6, hist-orb 3/2; M 24576, res 14064, rev 6722: ok}\\*
0 & -- & 0 & 6 & 1 & 646 & 2 & 376 & \texttt{-+????????}\\*
1 & 0 & 1 & 6 & 1 & 268 & 4 & 752 & \texttt{--????????}\\
2 & 1 & 2 & 5 & 1 & 48 & 16 & 960 & \texttt{---??????-}\\
3 & 2 & 3 & 1 & 1 & 512 & 512 & 0 & \texttt{----+-----}\\
4 & 3 & 4 & 1 & 1 & 0 & 1024 & 0 & \texttt{----------}\\
5 & 2 & 3 & 1 & 4 & 0 & 1024 & 0 & \texttt{----------}\\
6 & 2 & 3 & 1 & 4 & 0 & 1024 & 0 & \texttt{----------}\\
7 & 2 & 3 & 1 & 4 & 0 & 1024 & 0 & \texttt{----------}\\
8 & 2 & 3 & 1 & 4 & 0 & 1024 & 0 & \texttt{----------}\\
9 & 1 & 2 & 1 & 1 & 0 & 1024 & 0 & \texttt{----------}\\
\multicolumn{9}{@{}l@{}}{\rule{0pt}{2.6ex}\textbf{T10.441}\enspace \texttt{0123433331}\enspace $n{=}10$, $|H|{=}6$, len 5, br 2, $|\mathrm{Aut}|{=}24$, orb 7/6, hist-orb 3/2; M 23552, res 13848, rev 6262: ok}\\*
0 & -- & 0 & 6 & 1 & 390 & 2 & 632 & \texttt{-?????????}\\*
1 & 0 & 1 & 5 & 1 & 536 & 8 & 480 & \texttt{--+??????-}\\
2 & 1 & 2 & 5 & 1 & 48 & 16 & 960 & \texttt{---??????-}\\
3 & 2 & 3 & 1 & 1 & 512 & 512 & 0 & \texttt{----+-----}\\
4 & 3 & 4 & 1 & 1 & 0 & 1024 & 0 & \texttt{----------}\\
5 & 2 & 3 & 1 & 4 & 0 & 1024 & 0 & \texttt{----------}\\
6 & 2 & 3 & 1 & 4 & 0 & 1024 & 0 & \texttt{----------}\\
7 & 2 & 3 & 1 & 4 & 0 & 1024 & 0 & \texttt{----------}\\
8 & 2 & 3 & 1 & 4 & 0 & 1024 & 0 & \texttt{----------}\\
9 & 0 & 1 & 1 & 1 & 0 & 1024 & 0 & \texttt{----------}\\
\multicolumn{9}{@{}l@{}}{\rule{0pt}{2.6ex}\textbf{T10.442}\enspace \texttt{0123433323}\enspace $n{=}10$, $|H|{=}5$, len 5, br 2, $|\mathrm{Aut}|{=}6$, orb 8/7, hist-orb 3/2; M 24576, res 11008, rev 8214: ok}\\*
0 & -- & 0 & 5 & 1 & 722 & 2 & 300 & \texttt{-+????????}\\*
1 & 0 & 1 & 5 & 1 & 420 & 4 & 600 & \texttt{--????????}\\
2 & 1 & 2 & 4 & 1 & 96 & 32 & 896 & \texttt{---?????--}\\
3 & 2 & 3 & 1 & 1 & 512 & 512 & 0 & \texttt{----+-----}\\
4 & 3 & 4 & 1 & 1 & 0 & 1024 & 0 & \texttt{----------}\\
5 & 2 & 3 & 1 & 3 & 0 & 1024 & 0 & \texttt{----------}\\
6 & 2 & 3 & 1 & 3 & 0 & 1024 & 0 & \texttt{----------}\\
7 & 2 & 3 & 1 & 3 & 0 & 1024 & 0 & \texttt{----------}\\
8 & 1 & 2 & 1 & 1 & 512 & 512 & 0 & \texttt{---------+}\\
9 & 8 & 3 & 1 & 1 & 0 & 1024 & 0 & \texttt{----------}\\
\multicolumn{9}{@{}l@{}}{\rule{0pt}{2.6ex}\textbf{T10.443}\enspace \texttt{0123433322}\enspace $n{=}10$, $|H|{=}6$, len 5, br 2, $|\mathrm{Aut}|{=}12$, orb 7/5, hist-orb 3/1; M 23552, res 14368, rev 5642: ok}\\*
0 & -- & 0 & 6 & 1 & 582 & 2 & 440 & \texttt{-+????????}\\*
1 & 0 & 1 & 6 & 1 & 140 & 4 & 880 & \texttt{--????????}\\
2 & 1 & 2 & 4 & 1 & 96 & 32 & 896 & \texttt{---?????--}\\
3 & 2 & 3 & 1 & 1 & 512 & 512 & 0 & \texttt{----+-----}\\
4 & 3 & 4 & 1 & 1 & 0 & 1024 & 0 & \texttt{----------}\\
5 & 2 & 3 & 1 & 3 & 0 & 1024 & 0 & \texttt{----------}\\
6 & 2 & 3 & 1 & 3 & 0 & 1024 & 0 & \texttt{----------}\\
7 & 2 & 3 & 1 & 3 & 0 & 1024 & 0 & \texttt{----------}\\
8 & 1 & 2 & 1 & 2 & 0 & 1024 & 0 & \texttt{----------}\\
9 & 1 & 2 & 1 & 2 & 0 & 1024 & 0 & \texttt{----------}\\
\multicolumn{9}{@{}l@{}}{\rule{0pt}{2.6ex}\textbf{T10.444}\enspace \texttt{0123433321}\enspace $n{=}10$, $|H|{=}6$, len 5, br 3, $|\mathrm{Aut}|{=}6$, orb 8/7, hist-orb 4/3; M 22528, res 14856, rev 4602: ok}\\*
0 & -- & 0 & 6 & 1 & 326 & 2 & 696 & \texttt{-?????????}\\*
1 & 0 & 1 & 5 & 1 & 280 & 8 & 736 & \texttt{--???????-}\\
2 & 1 & 2 & 4 & 1 & 96 & 32 & 896 & \texttt{---?????--}\\
3 & 2 & 3 & 1 & 1 & 512 & 512 & 0 & \texttt{----+-----}\\
4 & 3 & 4 & 1 & 1 & 0 & 1024 & 0 & \texttt{----------}\\
5 & 2 & 3 & 1 & 3 & 0 & 1024 & 0 & \texttt{----------}\\
6 & 2 & 3 & 1 & 3 & 0 & 1024 & 0 & \texttt{----------}\\
7 & 2 & 3 & 1 & 3 & 0 & 1024 & 0 & \texttt{----------}\\
8 & 1 & 2 & 1 & 1 & 0 & 1024 & 0 & \texttt{----------}\\
9 & 0 & 1 & 1 & 1 & 0 & 1024 & 0 & \texttt{----------}\\
\multicolumn{9}{@{}l@{}}{\rule{0pt}{2.6ex}\textbf{T10.445}\enspace \texttt{0123433312}\enspace $n{=}10$, $|H|{=}5$, len 5, br 2, $|\mathrm{Aut}|{=}6$, orb 8/7, hist-orb 3/2; M 22528, res 10124, rev 7454: ok}\\*
0 & -- & 0 & 5 & 1 & 594 & 2 & 428 & \texttt{-?????????}\\*
1 & 0 & 1 & 4 & 1 & 560 & 16 & 448 & \texttt{--+?????--}\\
2 & 1 & 2 & 4 & 1 & 96 & 32 & 896 & \texttt{---?????--}\\
3 & 2 & 3 & 1 & 1 & 512 & 512 & 0 & \texttt{----+-----}\\
4 & 3 & 4 & 1 & 1 & 0 & 1024 & 0 & \texttt{----------}\\
5 & 2 & 3 & 1 & 3 & 0 & 1024 & 0 & \texttt{----------}\\
6 & 2 & 3 & 1 & 3 & 0 & 1024 & 0 & \texttt{----------}\\
7 & 2 & 3 & 1 & 3 & 0 & 1024 & 0 & \texttt{----------}\\
8 & 0 & 1 & 1 & 1 & 512 & 512 & 0 & \texttt{---------+}\\
9 & 8 & 2 & 1 & 1 & 0 & 1024 & 0 & \texttt{----------}\\
\multicolumn{9}{@{}l@{}}{\rule{0pt}{2.6ex}\textbf{T10.446}\enspace \texttt{0123433311}\enspace $n{=}10$, $|H|{=}6$, len 5, br 2, $|\mathrm{Aut}|{=}12$, orb 7/5, hist-orb 3/1; M 21504, res 13688, rev 4750: ok}\\*
0 & -- & 0 & 6 & 1 & 198 & 2 & 824 & \texttt{-?????????}\\*
1 & 0 & 1 & 4 & 1 & 560 & 16 & 448 & \texttt{--+?????--}\\
2 & 1 & 2 & 4 & 1 & 96 & 32 & 896 & \texttt{---?????--}\\
3 & 2 & 3 & 1 & 1 & 512 & 512 & 0 & \texttt{----+-----}\\
4 & 3 & 4 & 1 & 1 & 0 & 1024 & 0 & \texttt{----------}\\
5 & 2 & 3 & 1 & 3 & 0 & 1024 & 0 & \texttt{----------}\\
6 & 2 & 3 & 1 & 3 & 0 & 1024 & 0 & \texttt{----------}\\
7 & 2 & 3 & 1 & 3 & 0 & 1024 & 0 & \texttt{----------}\\
8 & 0 & 1 & 1 & 2 & 0 & 1024 & 0 & \texttt{----------}\\
9 & 0 & 1 & 1 & 2 & 0 & 1024 & 0 & \texttt{----------}\\
\multicolumn{9}{@{}l@{}}{\rule{0pt}{2.6ex}\textbf{T10.447}\enspace \texttt{0123433234}\enspace $n{=}10$, $|H|{=}4$, len 5, br 2, $|\mathrm{Aut}|{=}2$, orb 9/8, hist-orb 3/2; M 25600, res 8352, rev 9478: ok}\\*
0 & -- & 0 & 4 & 1 & 778 & 2 & 244 & \texttt{-+????????}\\*
1 & 0 & 1 & 4 & 1 & 532 & 4 & 488 & \texttt{--????????}\\
2 & 1 & 2 & 3 & 1 & 192 & 64 & 768 & \texttt{---????---}\\
3 & 2 & 3 & 1 & 1 & 512 & 512 & 0 & \texttt{----+-----}\\
4 & 3 & 4 & 1 & 1 & 0 & 1024 & 0 & \texttt{----------}\\
5 & 2 & 3 & 1 & 2 & 0 & 1024 & 0 & \texttt{----------}\\
6 & 2 & 3 & 1 & 2 & 0 & 1024 & 0 & \texttt{----------}\\
7 & 1 & 2 & 1 & 1 & 768 & 256 & 0 & \texttt{--------++}\\
8 & 7 & 3 & 1 & 1 & 512 & 512 & 0 & \texttt{---------+}\\
9 & 8 & 4 & 1 & 1 & 0 & 1024 & 0 & \texttt{----------}\\
\multicolumn{9}{@{}l@{}}{\rule{0pt}{2.6ex}\textbf{T10.448}\enspace \texttt{0123433233}\enspace $n{=}10$, $|H|{=}5$, len 5, br 3, $|\mathrm{Aut}|{=}4$, orb 8/6, hist-orb 3/1; M 24576, res 10816, rev 7906: ok}\\*
0 & -- & 0 & 5 & 1 & 702 & 2 & 320 & \texttt{-+????????}\\*
1 & 0 & 1 & 5 & 1 & 380 & 4 & 640 & \texttt{--????????}\\
2 & 1 & 2 & 3 & 1 & 192 & 64 & 768 & \texttt{---????---}\\
3 & 2 & 3 & 1 & 1 & 512 & 512 & 0 & \texttt{----+-----}\\
4 & 3 & 4 & 1 & 1 & 0 & 1024 & 0 & \texttt{----------}\\
5 & 2 & 3 & 1 & 2 & 0 & 1024 & 0 & \texttt{----------}\\
6 & 2 & 3 & 1 & 2 & 0 & 1024 & 0 & \texttt{----------}\\
7 & 1 & 2 & 2 & 1 & 256 & 256 & 512 & \texttt{--------??}\\
8 & 7 & 3 & 1 & 2 & 0 & 1024 & 0 & \texttt{----------}\\
9 & 7 & 3 & 1 & 2 & 0 & 1024 & 0 & \texttt{----------}\\
\multicolumn{9}{@{}l@{}}{\rule{0pt}{2.6ex}\textbf{T10.449}\enspace \texttt{0123433232}\enspace $n{=}10$, $|H|{=}5$, len 5, br 2, $|\mathrm{Aut}|{=}2$, orb 9/8, hist-orb 4/3; M 23552, res 12000, rev 6774: ok}\\*
0 & -- & 0 & 5 & 1 & 626 & 2 & 396 & \texttt{-+????????}\\*
1 & 0 & 1 & 5 & 1 & 228 & 4 & 792 & \texttt{--????????}\\
2 & 1 & 2 & 3 & 1 & 192 & 64 & 768 & \texttt{---????---}\\
3 & 2 & 3 & 1 & 1 & 512 & 512 & 0 & \texttt{----+-----}\\
4 & 3 & 4 & 1 & 1 & 0 & 1024 & 0 & \texttt{----------}\\
5 & 2 & 3 & 1 & 2 & 0 & 1024 & 0 & \texttt{----------}\\
6 & 2 & 3 & 1 & 2 & 0 & 1024 & 0 & \texttt{----------}\\
7 & 1 & 2 & 1 & 1 & 512 & 512 & 0 & \texttt{--------+-}\\
8 & 7 & 3 & 1 & 1 & 0 & 1024 & 0 & \texttt{----------}\\
9 & 1 & 2 & 1 & 1 & 0 & 1024 & 0 & \texttt{----------}\\
\multicolumn{9}{@{}l@{}}{\rule{0pt}{2.6ex}\textbf{T10.450}\enspace \texttt{0123433231}\enspace $n{=}10$, $|H|{=}5$, len 5, br 3, $|\mathrm{Aut}|{=}2$, orb 9/8, hist-orb 4/3; M 22528, res 12004, rev 6102: ok}\\*
0 & -- & 0 & 5 & 1 & 370 & 2 & 652 & \texttt{-?????????}\\*
1 & 0 & 1 & 4 & 1 & 456 & 8 & 560 & \texttt{--???????-}\\
2 & 1 & 2 & 3 & 1 & 192 & 64 & 768 & \texttt{---????---}\\
3 & 2 & 3 & 1 & 1 & 512 & 512 & 0 & \texttt{----+-----}\\
4 & 3 & 4 & 1 & 1 & 0 & 1024 & 0 & \texttt{----------}\\
5 & 2 & 3 & 1 & 2 & 0 & 1024 & 0 & \texttt{----------}\\
6 & 2 & 3 & 1 & 2 & 0 & 1024 & 0 & \texttt{----------}\\
7 & 1 & 2 & 1 & 1 & 512 & 512 & 0 & \texttt{--------+-}\\
8 & 7 & 3 & 1 & 1 & 0 & 1024 & 0 & \texttt{----------}\\
9 & 0 & 1 & 1 & 1 & 0 & 1024 & 0 & \texttt{----------}\\
\multicolumn{9}{@{}l@{}}{\rule{0pt}{2.6ex}\textbf{T10.451}\enspace \texttt{0123433222}\enspace $n{=}10$, $|H|{=}6$, len 5, br 2, $|\mathrm{Aut}|{=}12$, orb 7/5, hist-orb 3/1; M 22528, res 13824, rev 5258: ok}\\*
0 & -- & 0 & 6 & 1 & 550 & 2 & 472 & \texttt{-+????????}\\*
1 & 0 & 1 & 6 & 1 & 76 & 4 & 944 & \texttt{--????????}\\
2 & 1 & 2 & 3 & 1 & 192 & 64 & 768 & \texttt{---????---}\\
3 & 2 & 3 & 1 & 1 & 512 & 512 & 0 & \texttt{----+-----}\\
4 & 3 & 4 & 1 & 1 & 0 & 1024 & 0 & \texttt{----------}\\
5 & 2 & 3 & 1 & 2 & 0 & 1024 & 0 & \texttt{----------}\\
6 & 2 & 3 & 1 & 2 & 0 & 1024 & 0 & \texttt{----------}\\
7 & 1 & 2 & 1 & 3 & 0 & 1024 & 0 & \texttt{----------}\\
8 & 1 & 2 & 1 & 3 & 0 & 1024 & 0 & \texttt{----------}\\
9 & 1 & 2 & 1 & 3 & 0 & 1024 & 0 & \texttt{----------}\\
\multicolumn{9}{@{}l@{}}{\rule{0pt}{2.6ex}\textbf{T10.452}\enspace \texttt{0123433221}\enspace $n{=}10$, $|H|{=}6$, len 5, br 3, $|\mathrm{Aut}|{=}4$, orb 8/6, hist-orb 4/2; M 21504, res 14664, rev 3922: ok}\\*
0 & -- & 0 & 6 & 1 & 294 & 2 & 728 & \texttt{-?????????}\\*
1 & 0 & 1 & 5 & 1 & 152 & 8 & 864 & \texttt{--???????-}\\
2 & 1 & 2 & 3 & 1 & 192 & 64 & 768 & \texttt{---????---}\\
3 & 2 & 3 & 1 & 1 & 512 & 512 & 0 & \texttt{----+-----}\\
4 & 3 & 4 & 1 & 1 & 0 & 1024 & 0 & \texttt{----------}\\
5 & 2 & 3 & 1 & 2 & 0 & 1024 & 0 & \texttt{----------}\\
6 & 2 & 3 & 1 & 2 & 0 & 1024 & 0 & \texttt{----------}\\
7 & 1 & 2 & 1 & 2 & 0 & 1024 & 0 & \texttt{----------}\\
8 & 1 & 2 & 1 & 2 & 0 & 1024 & 0 & \texttt{----------}\\
9 & 0 & 1 & 1 & 1 & 0 & 1024 & 0 & \texttt{----------}\\
\multicolumn{9}{@{}l@{}}{\rule{0pt}{2.6ex}\textbf{T10.453}\enspace \texttt{0123433212}\enspace $n{=}10$, $|H|{=}5$, len 5, br 3, $|\mathrm{Aut}|{=}2$, orb 9/8, hist-orb 4/3; M 21504, res 11084, rev 5978: ok}\\*
0 & -- & 0 & 5 & 1 & 498 & 2 & 524 & \texttt{-?????????}\\*
1 & 0 & 1 & 4 & 1 & 304 & 16 & 704 & \texttt{--??????--}\\
2 & 1 & 2 & 3 & 1 & 192 & 64 & 768 & \texttt{---????---}\\
3 & 2 & 3 & 1 & 1 & 512 & 512 & 0 & \texttt{----+-----}\\
4 & 3 & 4 & 1 & 1 & 0 & 1024 & 0 & \texttt{----------}\\
5 & 2 & 3 & 1 & 2 & 0 & 1024 & 0 & \texttt{----------}\\
6 & 2 & 3 & 1 & 2 & 0 & 1024 & 0 & \texttt{----------}\\
7 & 1 & 2 & 1 & 1 & 0 & 1024 & 0 & \texttt{----------}\\
8 & 0 & 1 & 1 & 1 & 512 & 512 & 0 & \texttt{---------+}\\
9 & 8 & 2 & 1 & 1 & 0 & 1024 & 0 & \texttt{----------}\\
\multicolumn{9}{@{}l@{}}{\rule{0pt}{2.6ex}\textbf{T10.454}\enspace \texttt{0123433211}\enspace $n{=}10$, $|H|{=}6$, len 5, br 3, $|\mathrm{Aut}|{=}4$, orb 8/6, hist-orb 4/2; M 20480, res 14072, rev 3618: ok}\\*
0 & -- & 0 & 6 & 1 & 166 & 2 & 856 & \texttt{-?????????}\\*
1 & 0 & 1 & 4 & 1 & 304 & 16 & 704 & \texttt{--??????--}\\
2 & 1 & 2 & 3 & 1 & 192 & 64 & 768 & \texttt{---????---}\\
3 & 2 & 3 & 1 & 1 & 512 & 512 & 0 & \texttt{----+-----}\\
4 & 3 & 4 & 1 & 1 & 0 & 1024 & 0 & \texttt{----------}\\
5 & 2 & 3 & 1 & 2 & 0 & 1024 & 0 & \texttt{----------}\\
6 & 2 & 3 & 1 & 2 & 0 & 1024 & 0 & \texttt{----------}\\
7 & 1 & 2 & 1 & 1 & 0 & 1024 & 0 & \texttt{----------}\\
8 & 0 & 1 & 1 & 2 & 0 & 1024 & 0 & \texttt{----------}\\
9 & 0 & 1 & 1 & 2 & 0 & 1024 & 0 & \texttt{----------}\\
\multicolumn{9}{@{}l@{}}{\rule{0pt}{2.6ex}\textbf{T10.455}\enspace \texttt{0123433123}\enspace $n{=}10$, $|H|{=}4$, len 5, br 2, $|\mathrm{Aut}|{=}2$, orb 9/8, hist-orb 3/2; M 22528, res 6996, rev 8502: ok}\\*
0 & -- & 0 & 4 & 1 & 714 & 2 & 308 & \texttt{-?????????}\\*
1 & 0 & 1 & 3 & 1 & 608 & 32 & 384 & \texttt{--+????---}\\
2 & 1 & 2 & 3 & 1 & 192 & 64 & 768 & \texttt{---????---}\\
3 & 2 & 3 & 1 & 1 & 512 & 512 & 0 & \texttt{----+-----}\\
4 & 3 & 4 & 1 & 1 & 0 & 1024 & 0 & \texttt{----------}\\
5 & 2 & 3 & 1 & 2 & 0 & 1024 & 0 & \texttt{----------}\\
6 & 2 & 3 & 1 & 2 & 0 & 1024 & 0 & \texttt{----------}\\
7 & 0 & 1 & 1 & 1 & 768 & 256 & 0 & \texttt{--------++}\\
8 & 7 & 2 & 1 & 1 & 512 & 512 & 0 & \texttt{---------+}\\
9 & 8 & 3 & 1 & 1 & 0 & 1024 & 0 & \texttt{----------}\\
\multicolumn{9}{@{}l@{}}{\rule{0pt}{2.6ex}\textbf{T10.456}\enspace \texttt{0123433122}\enspace $n{=}10$, $|H|{=}5$, len 5, br 3, $|\mathrm{Aut}|{=}4$, orb 8/6, hist-orb 3/1; M 21504, res 9600, rev 6750: ok}\\*
0 & -- & 0 & 5 & 1 & 510 & 2 & 512 & \texttt{-?????????}\\*
1 & 0 & 1 & 3 & 1 & 608 & 32 & 384 & \texttt{--+????---}\\
2 & 1 & 2 & 3 & 1 & 192 & 64 & 768 & \texttt{---????---}\\
3 & 2 & 3 & 1 & 1 & 512 & 512 & 0 & \texttt{----+-----}\\
4 & 3 & 4 & 1 & 1 & 0 & 1024 & 0 & \texttt{----------}\\
5 & 2 & 3 & 1 & 2 & 0 & 1024 & 0 & \texttt{----------}\\
6 & 2 & 3 & 1 & 2 & 0 & 1024 & 0 & \texttt{----------}\\
7 & 0 & 1 & 2 & 1 & 256 & 256 & 512 & \texttt{--------??}\\
8 & 7 & 2 & 1 & 2 & 0 & 1024 & 0 & \texttt{----------}\\
9 & 7 & 2 & 1 & 2 & 0 & 1024 & 0 & \texttt{----------}\\
\multicolumn{9}{@{}l@{}}{\rule{0pt}{2.6ex}\textbf{T10.457}\enspace \texttt{0123433121}\enspace $n{=}10$, $|H|{=}5$, len 5, br 2, $|\mathrm{Aut}|{=}2$, orb 9/8, hist-orb 4/3; M 20480, res 10668, rev 5614: ok}\\*
0 & -- & 0 & 5 & 1 & 306 & 2 & 716 & \texttt{-?????????}\\*
1 & 0 & 1 & 3 & 1 & 608 & 32 & 384 & \texttt{--+????---}\\
2 & 1 & 2 & 3 & 1 & 192 & 64 & 768 & \texttt{---????---}\\
3 & 2 & 3 & 1 & 1 & 512 & 512 & 0 & \texttt{----+-----}\\
4 & 3 & 4 & 1 & 1 & 0 & 1024 & 0 & \texttt{----------}\\
5 & 2 & 3 & 1 & 2 & 0 & 1024 & 0 & \texttt{----------}\\
6 & 2 & 3 & 1 & 2 & 0 & 1024 & 0 & \texttt{----------}\\
7 & 0 & 1 & 1 & 1 & 512 & 512 & 0 & \texttt{--------+-}\\
8 & 7 & 2 & 1 & 1 & 0 & 1024 & 0 & \texttt{----------}\\
9 & 0 & 1 & 1 & 1 & 0 & 1024 & 0 & \texttt{----------}\\
\multicolumn{9}{@{}l@{}}{\rule{0pt}{2.6ex}\textbf{T10.458}\enspace \texttt{0123433111}\enspace $n{=}10$, $|H|{=}6$, len 5, br 2, $|\mathrm{Aut}|{=}12$, orb 7/5, hist-orb 3/1; M 19456, res 12504, rev 3966: ok}\\*
0 & -- & 0 & 6 & 1 & 102 & 2 & 920 & \texttt{-?????????}\\*
1 & 0 & 1 & 3 & 1 & 608 & 32 & 384 & \texttt{--+????---}\\
2 & 1 & 2 & 3 & 1 & 192 & 64 & 768 & \texttt{---????---}\\
3 & 2 & 3 & 1 & 1 & 512 & 512 & 0 & \texttt{----+-----}\\
4 & 3 & 4 & 1 & 1 & 0 & 1024 & 0 & \texttt{----------}\\
5 & 2 & 3 & 1 & 2 & 0 & 1024 & 0 & \texttt{----------}\\
6 & 2 & 3 & 1 & 2 & 0 & 1024 & 0 & \texttt{----------}\\
7 & 0 & 1 & 1 & 3 & 0 & 1024 & 0 & \texttt{----------}\\
8 & 0 & 1 & 1 & 3 & 0 & 1024 & 0 & \texttt{----------}\\
9 & 0 & 1 & 1 & 3 & 0 & 1024 & 0 & \texttt{----------}\\
\multicolumn{9}{@{}l@{}}{\rule{0pt}{2.6ex}\textbf{T10.459}\enspace \texttt{0123432343}\enspace $n{=}10$, $|H|{=}4$, len 5, br 3, $|\mathrm{Aut}|{=}2$, orb 6/2, hist-orb 2/0; M 25600, res 8736, rev 9202: ok}\\*
0 & -- & 0 & 4 & 1 & 754 & 2 & 268 & \texttt{-+????????}\\*
1 & 0 & 1 & 4 & 1 & 484 & 4 & 536 & \texttt{--????????}\\
2 & 1 & 2 & 2 & 2 & 384 & 128 & 512 & \texttt{---???----}\\
3 & 2 & 3 & 1 & 2 & 512 & 512 & 0 & \texttt{----+-----}\\
4 & 3 & 4 & 1 & 2 & 0 & 1024 & 0 & \texttt{----------}\\
5 & 2 & 3 & 1 & 2 & 0 & 1024 & 0 & \texttt{----------}\\
6 & 1 & 2 & 2 & 2 & 384 & 128 & 512 & \texttt{-------???}\\
7 & 6 & 3 & 1 & 2 & 512 & 512 & 0 & \texttt{--------+-}\\
8 & 7 & 4 & 1 & 2 & 0 & 1024 & 0 & \texttt{----------}\\
9 & 6 & 3 & 1 & 2 & 0 & 1024 & 0 & \texttt{----------}\\
\multicolumn{9}{@{}l@{}}{\rule{0pt}{2.6ex}\textbf{T10.460}\enspace \texttt{0123432342}\enspace $n{=}10$, $|H|{=}4$, len 5, br 2, $|\mathrm{Aut}|{=}1$, orb 10/10, hist-orb 4/4; M 24576, res 9696, rev 8062: ok}\\*
0 & -- & 0 & 4 & 1 & 666 & 2 & 356 & \texttt{-+????????}\\*
1 & 0 & 1 & 4 & 1 & 308 & 4 & 712 & \texttt{--????????}\\
2 & 1 & 2 & 2 & 1 & 384 & 128 & 512 & \texttt{---???----}\\
3 & 2 & 3 & 1 & 1 & 512 & 512 & 0 & \texttt{----+-----}\\
4 & 3 & 4 & 1 & 1 & 0 & 1024 & 0 & \texttt{----------}\\
5 & 2 & 3 & 1 & 1 & 0 & 1024 & 0 & \texttt{----------}\\
6 & 1 & 2 & 1 & 1 & 768 & 256 & 0 & \texttt{-------++-}\\
7 & 6 & 3 & 1 & 1 & 512 & 512 & 0 & \texttt{--------+-}\\
8 & 7 & 4 & 1 & 1 & 0 & 1024 & 0 & \texttt{----------}\\
9 & 1 & 2 & 1 & 1 & 0 & 1024 & 0 & \texttt{----------}\\
\multicolumn{9}{@{}l@{}}{\rule{0pt}{2.6ex}\textbf{T10.461}\enspace \texttt{0123432341}\enspace $n{=}10$, $|H|{=}4$, len 5, br 3, $|\mathrm{Aut}|{=}1$, orb 10/10, hist-orb 4/4; M 23552, res 9260, rev 7646: ok}\\*
0 & -- & 0 & 4 & 1 & 410 & 2 & 612 & \texttt{-?????????}\\*
1 & 0 & 1 & 3 & 1 & 616 & 8 & 400 & \texttt{--???????-}\\
2 & 1 & 2 & 2 & 1 & 384 & 128 & 512 & \texttt{---???----}\\
3 & 2 & 3 & 1 & 1 & 512 & 512 & 0 & \texttt{----+-----}\\
4 & 3 & 4 & 1 & 1 & 0 & 1024 & 0 & \texttt{----------}\\
5 & 2 & 3 & 1 & 1 & 0 & 1024 & 0 & \texttt{----------}\\
6 & 1 & 2 & 1 & 1 & 768 & 256 & 0 & \texttt{-------++-}\\
7 & 6 & 3 & 1 & 1 & 512 & 512 & 0 & \texttt{--------+-}\\
8 & 7 & 4 & 1 & 1 & 0 & 1024 & 0 & \texttt{----------}\\
9 & 0 & 1 & 1 & 1 & 0 & 1024 & 0 & \texttt{----------}\\
\multicolumn{9}{@{}l@{}}{\rule{0pt}{2.6ex}\textbf{T10.462}\enspace \texttt{0123432333}\enspace $n{=}10$, $|H|{=}5$, len 5, br 3, $|\mathrm{Aut}|{=}6$, orb 8/7, hist-orb 3/2; M 24576, res 10368, rev 8098: ok}\\*
0 & -- & 0 & 5 & 1 & 710 & 2 & 312 & \texttt{-+????????}\\*
1 & 0 & 1 & 5 & 1 & 396 & 4 & 624 & \texttt{--????????}\\
2 & 1 & 2 & 2 & 1 & 384 & 128 & 512 & \texttt{---???----}\\
3 & 2 & 3 & 1 & 1 & 512 & 512 & 0 & \texttt{----+-----}\\
4 & 3 & 4 & 1 & 1 & 0 & 1024 & 0 & \texttt{----------}\\
5 & 2 & 3 & 1 & 1 & 0 & 1024 & 0 & \texttt{----------}\\
6 & 1 & 2 & 3 & 1 & 128 & 128 & 768 & \texttt{-------???}\\
7 & 6 & 3 & 1 & 3 & 0 & 1024 & 0 & \texttt{----------}\\
8 & 6 & 3 & 1 & 3 & 0 & 1024 & 0 & \texttt{----------}\\
9 & 6 & 3 & 1 & 3 & 0 & 1024 & 0 & \texttt{----------}\\
\multicolumn{9}{@{}l@{}}{\rule{0pt}{2.6ex}\textbf{T10.463}\enspace \texttt{0123432332}\enspace $n{=}10$, $|H|{=}5$, len 5, br 3, $|\mathrm{Aut}|{=}2$, orb 9/8, hist-orb 4/3; M 23552, res 11392, rev 6826: ok}\\*
0 & -- & 0 & 5 & 1 & 622 & 2 & 400 & \texttt{-+????????}\\*
1 & 0 & 1 & 5 & 1 & 220 & 4 & 800 & \texttt{--????????}\\
2 & 1 & 2 & 2 & 1 & 384 & 128 & 512 & \texttt{---???----}\\
3 & 2 & 3 & 1 & 1 & 512 & 512 & 0 & \texttt{----+-----}\\
4 & 3 & 4 & 1 & 1 & 0 & 1024 & 0 & \texttt{----------}\\
5 & 2 & 3 & 1 & 1 & 0 & 1024 & 0 & \texttt{----------}\\
6 & 1 & 2 & 2 & 1 & 256 & 256 & 512 & \texttt{-------??-}\\
7 & 6 & 3 & 1 & 2 & 0 & 1024 & 0 & \texttt{----------}\\
8 & 6 & 3 & 1 & 2 & 0 & 1024 & 0 & \texttt{----------}\\
9 & 1 & 2 & 1 & 1 & 0 & 1024 & 0 & \texttt{----------}\\
\multicolumn{9}{@{}l@{}}{\rule{0pt}{2.6ex}\textbf{T10.464}\enspace \texttt{0123432331}\enspace $n{=}10$, $|H|{=}5$, len 5, br 4, $|\mathrm{Aut}|{=}2$, orb 9/8, hist-orb 4/3; M 22528, res 11440, rev 6146: ok}\\*
0 & -- & 0 & 5 & 1 & 366 & 2 & 656 & \texttt{-?????????}\\*
1 & 0 & 1 & 4 & 1 & 440 & 8 & 576 & \texttt{--???????-}\\
2 & 1 & 2 & 2 & 1 & 384 & 128 & 512 & \texttt{---???----}\\
3 & 2 & 3 & 1 & 1 & 512 & 512 & 0 & \texttt{----+-----}\\
4 & 3 & 4 & 1 & 1 & 0 & 1024 & 0 & \texttt{----------}\\
5 & 2 & 3 & 1 & 1 & 0 & 1024 & 0 & \texttt{----------}\\
6 & 1 & 2 & 2 & 1 & 256 & 256 & 512 & \texttt{-------??-}\\
7 & 6 & 3 & 1 & 2 & 0 & 1024 & 0 & \texttt{----------}\\
8 & 6 & 3 & 1 & 2 & 0 & 1024 & 0 & \texttt{----------}\\
9 & 0 & 1 & 1 & 1 & 0 & 1024 & 0 & \texttt{----------}\\
\multicolumn{9}{@{}l@{}}{\rule{0pt}{2.6ex}\textbf{T10.465}\enspace \texttt{0123432323}\enspace $n{=}10$, $|H|{=}4$, len 5, br 2, $|\mathrm{Aut}|{=}2$, orb 8/6, hist-orb 3/2; M 23552, res 8640, rev 8410: ok}\\*
0 & -- & 0 & 4 & 1 & 710 & 2 & 312 & \texttt{-+????????}\\*
1 & 0 & 1 & 4 & 1 & 396 & 4 & 624 & \texttt{--????????}\\
2 & 1 & 2 & 2 & 1 & 384 & 128 & 512 & \texttt{---???----}\\
3 & 2 & 3 & 1 & 1 & 512 & 512 & 0 & \texttt{----+-----}\\
4 & 3 & 4 & 1 & 1 & 0 & 1024 & 0 & \texttt{----------}\\
5 & 2 & 3 & 1 & 1 & 0 & 1024 & 0 & \texttt{----------}\\
6 & 1 & 2 & 1 & 2 & 512 & 512 & 0 & \texttt{-------+--}\\
7 & 6 & 3 & 1 & 2 & 0 & 1024 & 0 & \texttt{----------}\\
8 & 1 & 2 & 1 & 2 & 512 & 512 & 0 & \texttt{---------+}\\
9 & 8 & 3 & 1 & 2 & 0 & 1024 & 0 & \texttt{----------}\\
\multicolumn{9}{@{}l@{}}{\rule{0pt}{2.6ex}\textbf{T10.466}\enspace \texttt{0123432322}\enspace $n{=}10$, $|H|{=}5$, len 5, br 2, $|\mathrm{Aut}|{=}2$, orb 9/8, hist-orb 4/3; M 22528, res 11808, rev 6222: ok}\\*
0 & -- & 0 & 5 & 1 & 578 & 2 & 444 & \texttt{-+????????}\\*
1 & 0 & 1 & 5 & 1 & 132 & 4 & 888 & \texttt{--????????}\\
2 & 1 & 2 & 2 & 1 & 384 & 128 & 512 & \texttt{---???----}\\
3 & 2 & 3 & 1 & 1 & 512 & 512 & 0 & \texttt{----+-----}\\
4 & 3 & 4 & 1 & 1 & 0 & 1024 & 0 & \texttt{----------}\\
5 & 2 & 3 & 1 & 1 & 0 & 1024 & 0 & \texttt{----------}\\
6 & 1 & 2 & 1 & 1 & 512 & 512 & 0 & \texttt{-------+--}\\
7 & 6 & 3 & 1 & 1 & 0 & 1024 & 0 & \texttt{----------}\\
8 & 1 & 2 & 1 & 2 & 0 & 1024 & 0 & \texttt{----------}\\
9 & 1 & 2 & 1 & 2 & 0 & 1024 & 0 & \texttt{----------}\\
\multicolumn{9}{@{}l@{}}{\rule{0pt}{2.6ex}\textbf{T10.467}\enspace \texttt{0123432321}\enspace $n{=}10$, $|H|{=}5$, len 5, br 3, $|\mathrm{Aut}|{=}1$, orb 10/10, hist-orb 5/5; M 21504, res 12340, rev 5166: ok}\\*
0 & -- & 0 & 5 & 1 & 322 & 2 & 700 & \texttt{-?????????}\\*
1 & 0 & 1 & 4 & 1 & 264 & 8 & 752 & \texttt{--???????-}\\
2 & 1 & 2 & 2 & 1 & 384 & 128 & 512 & \texttt{---???----}\\
3 & 2 & 3 & 1 & 1 & 512 & 512 & 0 & \texttt{----+-----}\\
4 & 3 & 4 & 1 & 1 & 0 & 1024 & 0 & \texttt{----------}\\
5 & 2 & 3 & 1 & 1 & 0 & 1024 & 0 & \texttt{----------}\\
6 & 1 & 2 & 1 & 1 & 512 & 512 & 0 & \texttt{-------+--}\\
7 & 6 & 3 & 1 & 1 & 0 & 1024 & 0 & \texttt{----------}\\
8 & 1 & 2 & 1 & 1 & 0 & 1024 & 0 & \texttt{----------}\\
9 & 0 & 1 & 1 & 1 & 0 & 1024 & 0 & \texttt{----------}\\
\multicolumn{9}{@{}l@{}}{\rule{0pt}{2.6ex}\textbf{T10.468}\enspace \texttt{0123432312}\enspace $n{=}10$, $|H|{=}4$, len 5, br 3, $|\mathrm{Aut}|{=}1$, orb 10/10, hist-orb 4/4; M 21504, res 7752, rev 7614: ok}\\*
0 & -- & 0 & 4 & 1 & 582 & 2 & 440 & \texttt{-?????????}\\*
1 & 0 & 1 & 3 & 1 & 528 & 16 & 480 & \texttt{--??????--}\\
2 & 1 & 2 & 2 & 1 & 384 & 128 & 512 & \texttt{---???----}\\
3 & 2 & 3 & 1 & 1 & 512 & 512 & 0 & \texttt{----+-----}\\
4 & 3 & 4 & 1 & 1 & 0 & 1024 & 0 & \texttt{----------}\\
5 & 2 & 3 & 1 & 1 & 0 & 1024 & 0 & \texttt{----------}\\
6 & 1 & 2 & 1 & 1 & 512 & 512 & 0 & \texttt{-------+--}\\
7 & 6 & 3 & 1 & 1 & 0 & 1024 & 0 & \texttt{----------}\\
8 & 0 & 1 & 1 & 1 & 512 & 512 & 0 & \texttt{---------+}\\
9 & 8 & 2 & 1 & 1 & 0 & 1024 & 0 & \texttt{----------}\\
\multicolumn{9}{@{}l@{}}{\rule{0pt}{2.6ex}\textbf{T10.469}\enspace \texttt{0123432311}\enspace $n{=}10$, $|H|{=}5$, len 5, br 3, $|\mathrm{Aut}|{=}2$, orb 9/8, hist-orb 4/3; M 20480, res 11244, rev 5142: ok}\\*
0 & -- & 0 & 5 & 1 & 194 & 2 & 828 & \texttt{-?????????}\\*
1 & 0 & 1 & 3 & 1 & 528 & 16 & 480 & \texttt{--??????--}\\
2 & 1 & 2 & 2 & 1 & 384 & 128 & 512 & \texttt{---???----}\\
3 & 2 & 3 & 1 & 1 & 512 & 512 & 0 & \texttt{----+-----}\\
4 & 3 & 4 & 1 & 1 & 0 & 1024 & 0 & \texttt{----------}\\
5 & 2 & 3 & 1 & 1 & 0 & 1024 & 0 & \texttt{----------}\\
6 & 1 & 2 & 1 & 1 & 512 & 512 & 0 & \texttt{-------+--}\\
7 & 6 & 3 & 1 & 1 & 0 & 1024 & 0 & \texttt{----------}\\
8 & 0 & 1 & 1 & 2 & 0 & 1024 & 0 & \texttt{----------}\\
9 & 0 & 1 & 1 & 2 & 0 & 1024 & 0 & \texttt{----------}\\
\multicolumn{9}{@{}l@{}}{\rule{0pt}{2.6ex}\textbf{T10.470}\enspace \texttt{0123432222}\enspace $n{=}10$, $|H|{=}6$, len 5, br 2, $|\mathrm{Aut}|{=}24$, orb 7/6, hist-orb 3/2; M 21504, res 12864, rev 5234: ok}\\*
0 & -- & 0 & 6 & 1 & 534 & 2 & 488 & \texttt{-+????????}\\*
1 & 0 & 1 & 6 & 1 & 44 & 4 & 976 & \texttt{--????????}\\
2 & 1 & 2 & 2 & 1 & 384 & 128 & 512 & \texttt{---???----}\\
3 & 2 & 3 & 1 & 1 & 512 & 512 & 0 & \texttt{----+-----}\\
4 & 3 & 4 & 1 & 1 & 0 & 1024 & 0 & \texttt{----------}\\
5 & 2 & 3 & 1 & 1 & 0 & 1024 & 0 & \texttt{----------}\\
6 & 1 & 2 & 1 & 4 & 0 & 1024 & 0 & \texttt{----------}\\
7 & 1 & 2 & 1 & 4 & 0 & 1024 & 0 & \texttt{----------}\\
8 & 1 & 2 & 1 & 4 & 0 & 1024 & 0 & \texttt{----------}\\
9 & 1 & 2 & 1 & 4 & 0 & 1024 & 0 & \texttt{----------}\\
\multicolumn{9}{@{}l@{}}{\rule{0pt}{2.6ex}\textbf{T10.471}\enspace \texttt{0123432221}\enspace $n{=}10$, $|H|{=}6$, len 5, br 3, $|\mathrm{Aut}|{=}6$, orb 8/7, hist-orb 4/3; M 20480, res 13880, rev 3738: ok}\\*
0 & -- & 0 & 6 & 1 & 278 & 2 & 744 & \texttt{-?????????}\\*
1 & 0 & 1 & 5 & 1 & 88 & 8 & 928 & \texttt{--???????-}\\
2 & 1 & 2 & 2 & 1 & 384 & 128 & 512 & \texttt{---???----}\\
3 & 2 & 3 & 1 & 1 & 512 & 512 & 0 & \texttt{----+-----}\\
4 & 3 & 4 & 1 & 1 & 0 & 1024 & 0 & \texttt{----------}\\
5 & 2 & 3 & 1 & 1 & 0 & 1024 & 0 & \texttt{----------}\\
6 & 1 & 2 & 1 & 3 & 0 & 1024 & 0 & \texttt{----------}\\
7 & 1 & 2 & 1 & 3 & 0 & 1024 & 0 & \texttt{----------}\\
8 & 1 & 2 & 1 & 3 & 0 & 1024 & 0 & \texttt{----------}\\
9 & 0 & 1 & 1 & 1 & 0 & 1024 & 0 & \texttt{----------}\\
\multicolumn{9}{@{}l@{}}{\rule{0pt}{2.6ex}\textbf{T10.472}\enspace \texttt{0123432212}\enspace $n{=}10$, $|H|{=}5$, len 5, br 3, $|\mathrm{Aut}|{=}2$, orb 9/8, hist-orb 4/3; M 20480, res 10876, rev 5402: ok}\\*
0 & -- & 0 & 5 & 1 & 450 & 2 & 572 & \texttt{-?????????}\\*
1 & 0 & 1 & 4 & 1 & 176 & 16 & 832 & \texttt{--??????--}\\
2 & 1 & 2 & 2 & 1 & 384 & 128 & 512 & \texttt{---???----}\\
3 & 2 & 3 & 1 & 1 & 512 & 512 & 0 & \texttt{----+-----}\\
4 & 3 & 4 & 1 & 1 & 0 & 1024 & 0 & \texttt{----------}\\
5 & 2 & 3 & 1 & 1 & 0 & 1024 & 0 & \texttt{----------}\\
6 & 1 & 2 & 1 & 2 & 0 & 1024 & 0 & \texttt{----------}\\
7 & 1 & 2 & 1 & 2 & 0 & 1024 & 0 & \texttt{----------}\\
8 & 0 & 1 & 1 & 1 & 512 & 512 & 0 & \texttt{---------+}\\
9 & 8 & 2 & 1 & 1 & 0 & 1024 & 0 & \texttt{----------}\\
\multicolumn{9}{@{}l@{}}{\rule{0pt}{2.6ex}\textbf{T10.473}\enspace \texttt{0123432211}\enspace $n{=}10$, $|H|{=}6$, len 5, br 3, $|\mathrm{Aut}|{=}4$, orb 8/6, hist-orb 4/2; M 19456, res 13576, rev 3202: ok}\\*
0 & -- & 0 & 6 & 1 & 150 & 2 & 872 & \texttt{-?????????}\\*
1 & 0 & 1 & 4 & 1 & 176 & 16 & 832 & \texttt{--??????--}\\
2 & 1 & 2 & 2 & 1 & 384 & 128 & 512 & \texttt{---???----}\\
3 & 2 & 3 & 1 & 1 & 512 & 512 & 0 & \texttt{----+-----}\\
4 & 3 & 4 & 1 & 1 & 0 & 1024 & 0 & \texttt{----------}\\
5 & 2 & 3 & 1 & 1 & 0 & 1024 & 0 & \texttt{----------}\\
6 & 1 & 2 & 1 & 2 & 0 & 1024 & 0 & \texttt{----------}\\
7 & 1 & 2 & 1 & 2 & 0 & 1024 & 0 & \texttt{----------}\\
8 & 0 & 1 & 1 & 2 & 0 & 1024 & 0 & \texttt{----------}\\
9 & 0 & 1 & 1 & 2 & 0 & 1024 & 0 & \texttt{----------}\\
\multicolumn{9}{@{}l@{}}{\rule{0pt}{2.6ex}\textbf{T10.474}\enspace \texttt{0123432123}\enspace $n{=}10$, $|H|{=}4$, len 5, br 3, $|\mathrm{Aut}|{=}1$, orb 10/10, hist-orb 4/4; M 21504, res 7812, rev 7322: ok}\\*
0 & -- & 0 & 4 & 1 & 602 & 2 & 420 & \texttt{-?????????}\\*
1 & 0 & 1 & 3 & 1 & 352 & 32 & 640 & \texttt{--?????---}\\
2 & 1 & 2 & 2 & 1 & 384 & 128 & 512 & \texttt{---???----}\\
3 & 2 & 3 & 1 & 1 & 512 & 512 & 0 & \texttt{----+-----}\\
4 & 3 & 4 & 1 & 1 & 0 & 1024 & 0 & \texttt{----------}\\
5 & 2 & 3 & 1 & 1 & 0 & 1024 & 0 & \texttt{----------}\\
6 & 1 & 2 & 1 & 1 & 0 & 1024 & 0 & \texttt{----------}\\
7 & 0 & 1 & 1 & 1 & 768 & 256 & 0 & \texttt{--------++}\\
8 & 7 & 2 & 1 & 1 & 512 & 512 & 0 & \texttt{---------+}\\
9 & 8 & 3 & 1 & 1 & 0 & 1024 & 0 & \texttt{----------}\\
\multicolumn{9}{@{}l@{}}{\rule{0pt}{2.6ex}\textbf{T10.475}\enspace \texttt{0123432122}\enspace $n{=}10$, $|H|{=}5$, len 5, br 4, $|\mathrm{Aut}|{=}2$, orb 9/8, hist-orb 4/3; M 20480, res 10128, rev 5674: ok}\\*
0 & -- & 0 & 5 & 1 & 430 & 2 & 592 & \texttt{-?????????}\\*
1 & 0 & 1 & 3 & 1 & 352 & 32 & 640 & \texttt{--?????---}\\
2 & 1 & 2 & 2 & 1 & 384 & 128 & 512 & \texttt{---???----}\\
3 & 2 & 3 & 1 & 1 & 512 & 512 & 0 & \texttt{----+-----}\\
4 & 3 & 4 & 1 & 1 & 0 & 1024 & 0 & \texttt{----------}\\
5 & 2 & 3 & 1 & 1 & 0 & 1024 & 0 & \texttt{----------}\\
6 & 1 & 2 & 1 & 1 & 0 & 1024 & 0 & \texttt{----------}\\
7 & 0 & 1 & 2 & 1 & 256 & 256 & 512 & \texttt{--------??}\\
8 & 7 & 2 & 1 & 2 & 0 & 1024 & 0 & \texttt{----------}\\
9 & 7 & 2 & 1 & 2 & 0 & 1024 & 0 & \texttt{----------}\\
\multicolumn{9}{@{}l@{}}{\rule{0pt}{2.6ex}\textbf{T10.476}\enspace \texttt{0123432121}\enspace $n{=}10$, $|H|{=}5$, len 5, br 3, $|\mathrm{Aut}|{=}1$, orb 10/10, hist-orb 5/5; M 19456, res 10908, rev 4706: ok}\\*
0 & -- & 0 & 5 & 1 & 258 & 2 & 764 & \texttt{-?????????}\\*
1 & 0 & 1 & 3 & 1 & 352 & 32 & 640 & \texttt{--?????---}\\
2 & 1 & 2 & 2 & 1 & 384 & 128 & 512 & \texttt{---???----}\\
3 & 2 & 3 & 1 & 1 & 512 & 512 & 0 & \texttt{----+-----}\\
4 & 3 & 4 & 1 & 1 & 0 & 1024 & 0 & \texttt{----------}\\
5 & 2 & 3 & 1 & 1 & 0 & 1024 & 0 & \texttt{----------}\\
6 & 1 & 2 & 1 & 1 & 0 & 1024 & 0 & \texttt{----------}\\
7 & 0 & 1 & 1 & 1 & 512 & 512 & 0 & \texttt{--------+-}\\
8 & 7 & 2 & 1 & 1 & 0 & 1024 & 0 & \texttt{----------}\\
9 & 0 & 1 & 1 & 1 & 0 & 1024 & 0 & \texttt{----------}\\
\multicolumn{9}{@{}l@{}}{\rule{0pt}{2.6ex}\textbf{T10.477}\enspace \texttt{0123432111}\enspace $n{=}10$, $|H|{=}6$, len 5, br 3, $|\mathrm{Aut}|{=}6$, orb 8/7, hist-orb 4/3; M 18432, res 12456, rev 3226: ok}\\*
0 & -- & 0 & 6 & 1 & 86 & 2 & 936 & \texttt{-?????????}\\*
1 & 0 & 1 & 3 & 1 & 352 & 32 & 640 & \texttt{--?????---}\\
2 & 1 & 2 & 2 & 1 & 384 & 128 & 512 & \texttt{---???----}\\
3 & 2 & 3 & 1 & 1 & 512 & 512 & 0 & \texttt{----+-----}\\
4 & 3 & 4 & 1 & 1 & 0 & 1024 & 0 & \texttt{----------}\\
5 & 2 & 3 & 1 & 1 & 0 & 1024 & 0 & \texttt{----------}\\
6 & 1 & 2 & 1 & 1 & 0 & 1024 & 0 & \texttt{----------}\\
7 & 0 & 1 & 1 & 3 & 0 & 1024 & 0 & \texttt{----------}\\
8 & 0 & 1 & 1 & 3 & 0 & 1024 & 0 & \texttt{----------}\\
9 & 0 & 1 & 1 & 3 & 0 & 1024 & 0 & \texttt{----------}\\
\multicolumn{9}{@{}l@{}}{\rule{0pt}{2.6ex}\textbf{T10.478}\enspace \texttt{0123431234}\enspace $n{=}10$, $|H|{=}3$, len 5, br 2, $|\mathrm{Aut}|{=}1$, orb 10/10, hist-orb 3/3; M 23552, res 3956, rev 9678: ok}\\*
0 & -- & 0 & 3 & 1 & 810 & 2 & 212 & \texttt{-?????????}\\*
1 & 0 & 1 & 2 & 1 & 704 & 64 & 256 & \texttt{--+???----}\\
2 & 1 & 2 & 2 & 1 & 384 & 128 & 512 & \texttt{---???----}\\
3 & 2 & 3 & 1 & 1 & 512 & 512 & 0 & \texttt{----+-----}\\
4 & 3 & 4 & 1 & 1 & 0 & 1024 & 0 & \texttt{----------}\\
5 & 2 & 3 & 1 & 1 & 0 & 1024 & 0 & \texttt{----------}\\
6 & 0 & 1 & 1 & 1 & 896 & 128 & 0 & \texttt{-------+++}\\
7 & 6 & 2 & 1 & 1 & 768 & 256 & 0 & \texttt{--------++}\\
8 & 7 & 3 & 1 & 1 & 512 & 512 & 0 & \texttt{---------+}\\
9 & 8 & 4 & 1 & 1 & 0 & 1024 & 0 & \texttt{----------}\\
\multicolumn{9}{@{}l@{}}{\rule{0pt}{2.6ex}\textbf{T10.479}\enspace \texttt{0123431233}\enspace $n{=}10$, $|H|{=}4$, len 5, br 3, $|\mathrm{Aut}|{=}2$, orb 9/8, hist-orb 3/2; M 22528, res 6208, rev 8586: ok}\\*
0 & -- & 0 & 4 & 1 & 702 & 2 & 320 & \texttt{-?????????}\\*
1 & 0 & 1 & 2 & 1 & 704 & 64 & 256 & \texttt{--+???----}\\
2 & 1 & 2 & 2 & 1 & 384 & 128 & 512 & \texttt{---???----}\\
3 & 2 & 3 & 1 & 1 & 512 & 512 & 0 & \texttt{----+-----}\\
4 & 3 & 4 & 1 & 1 & 0 & 1024 & 0 & \texttt{----------}\\
5 & 2 & 3 & 1 & 1 & 0 & 1024 & 0 & \texttt{----------}\\
6 & 0 & 1 & 2 & 1 & 640 & 128 & 256 & \texttt{-------+??}\\
7 & 6 & 2 & 2 & 1 & 256 & 256 & 512 & \texttt{--------??}\\
8 & 7 & 3 & 1 & 2 & 0 & 1024 & 0 & \texttt{----------}\\
9 & 7 & 3 & 1 & 2 & 0 & 1024 & 0 & \texttt{----------}\\
\multicolumn{9}{@{}l@{}}{\rule{0pt}{2.6ex}\textbf{T10.480}\enspace \texttt{0123431232}\enspace $n{=}10$, $|H|{=}4$, len 5, br 3, $|\mathrm{Aut}|{=}1$, orb 10/10, hist-orb 4/4; M 21504, res 7180, rev 7682: ok}\\*
0 & -- & 0 & 4 & 1 & 594 & 2 & 428 & \texttt{-?????????}\\*
1 & 0 & 1 & 2 & 1 & 704 & 64 & 256 & \texttt{--+???----}\\
2 & 1 & 2 & 2 & 1 & 384 & 128 & 512 & \texttt{---???----}\\
3 & 2 & 3 & 1 & 1 & 512 & 512 & 0 & \texttt{----+-----}\\
4 & 3 & 4 & 1 & 1 & 0 & 1024 & 0 & \texttt{----------}\\
5 & 2 & 3 & 1 & 1 & 0 & 1024 & 0 & \texttt{----------}\\
6 & 0 & 1 & 2 & 1 & 384 & 128 & 512 & \texttt{-------???}\\
7 & 6 & 2 & 1 & 1 & 512 & 512 & 0 & \texttt{--------+-}\\
8 & 7 & 3 & 1 & 1 & 0 & 1024 & 0 & \texttt{----------}\\
9 & 6 & 2 & 1 & 1 & 0 & 1024 & 0 & \texttt{----------}\\
\multicolumn{9}{@{}l@{}}{\rule{0pt}{2.6ex}\textbf{T10.481}\enspace \texttt{0123431231}\enspace $n{=}10$, $|H|{=}4$, len 5, br 2, $|\mathrm{Aut}|{=}1$, orb 10/10, hist-orb 4/4; M 20480, res 7844, rev 6702: ok}\\*
0 & -- & 0 & 4 & 1 & 378 & 2 & 644 & \texttt{-?????????}\\*
1 & 0 & 1 & 2 & 1 & 704 & 64 & 256 & \texttt{--+???----}\\
2 & 1 & 2 & 2 & 1 & 384 & 128 & 512 & \texttt{---???----}\\
3 & 2 & 3 & 1 & 1 & 512 & 512 & 0 & \texttt{----+-----}\\
4 & 3 & 4 & 1 & 1 & 0 & 1024 & 0 & \texttt{----------}\\
5 & 2 & 3 & 1 & 1 & 0 & 1024 & 0 & \texttt{----------}\\
6 & 0 & 1 & 1 & 1 & 768 & 256 & 0 & \texttt{-------++-}\\
7 & 6 & 2 & 1 & 1 & 512 & 512 & 0 & \texttt{--------+-}\\
8 & 7 & 3 & 1 & 1 & 0 & 1024 & 0 & \texttt{----------}\\
9 & 0 & 1 & 1 & 1 & 0 & 1024 & 0 & \texttt{----------}\\
\multicolumn{9}{@{}l@{}}{\rule{0pt}{2.6ex}\textbf{T10.482}\enspace \texttt{0123431222}\enspace $n{=}10$, $|H|{=}5$, len 5, br 3, $|\mathrm{Aut}|{=}6$, orb 8/7, hist-orb 3/2; M 20480, res 8792, rev 6374: ok}\\*
0 & -- & 0 & 5 & 1 & 486 & 2 & 536 & \texttt{-?????????}\\*
1 & 0 & 1 & 2 & 1 & 704 & 64 & 256 & \texttt{--+???----}\\
2 & 1 & 2 & 2 & 1 & 384 & 128 & 512 & \texttt{---???----}\\
3 & 2 & 3 & 1 & 1 & 512 & 512 & 0 & \texttt{----+-----}\\
4 & 3 & 4 & 1 & 1 & 0 & 1024 & 0 & \texttt{----------}\\
5 & 2 & 3 & 1 & 1 & 0 & 1024 & 0 & \texttt{----------}\\
6 & 0 & 1 & 3 & 1 & 128 & 128 & 768 & \texttt{-------???}\\
7 & 6 & 2 & 1 & 3 & 0 & 1024 & 0 & \texttt{----------}\\
8 & 6 & 2 & 1 & 3 & 0 & 1024 & 0 & \texttt{----------}\\
9 & 6 & 2 & 1 & 3 & 0 & 1024 & 0 & \texttt{----------}\\
\multicolumn{9}{@{}l@{}}{\rule{0pt}{2.6ex}\textbf{T10.483}\enspace \texttt{0123431221}\enspace $n{=}10$, $|H|{=}5$, len 5, br 3, $|\mathrm{Aut}|{=}2$, orb 9/8, hist-orb 4/3; M 19456, res 9584, rev 5286: ok}\\*
0 & -- & 0 & 5 & 1 & 270 & 2 & 752 & \texttt{-?????????}\\*
1 & 0 & 1 & 2 & 1 & 704 & 64 & 256 & \texttt{--+???----}\\
2 & 1 & 2 & 2 & 1 & 384 & 128 & 512 & \texttt{---???----}\\
3 & 2 & 3 & 1 & 1 & 512 & 512 & 0 & \texttt{----+-----}\\
4 & 3 & 4 & 1 & 1 & 0 & 1024 & 0 & \texttt{----------}\\
5 & 2 & 3 & 1 & 1 & 0 & 1024 & 0 & \texttt{----------}\\
6 & 0 & 1 & 2 & 1 & 256 & 256 & 512 & \texttt{-------??-}\\
7 & 6 & 2 & 1 & 2 & 0 & 1024 & 0 & \texttt{----------}\\
8 & 6 & 2 & 1 & 2 & 0 & 1024 & 0 & \texttt{----------}\\
9 & 0 & 1 & 1 & 1 & 0 & 1024 & 0 & \texttt{----------}\\
\multicolumn{9}{@{}l@{}}{\rule{0pt}{2.6ex}\textbf{T10.484}\enspace \texttt{0123431212}\enspace $n{=}10$, $|H|{=}4$, len 5, br 2, $|\mathrm{Aut}|{=}2$, orb 8/6, hist-orb 3/2; M 19456, res 6872, rev 6798: ok}\\*
0 & -- & 0 & 4 & 1 & 486 & 2 & 536 & \texttt{-?????????}\\*
1 & 0 & 1 & 2 & 1 & 704 & 64 & 256 & \texttt{--+???----}\\
2 & 1 & 2 & 2 & 1 & 384 & 128 & 512 & \texttt{---???----}\\
3 & 2 & 3 & 1 & 1 & 512 & 512 & 0 & \texttt{----+-----}\\
4 & 3 & 4 & 1 & 1 & 0 & 1024 & 0 & \texttt{----------}\\
5 & 2 & 3 & 1 & 1 & 0 & 1024 & 0 & \texttt{----------}\\
6 & 0 & 1 & 1 & 2 & 512 & 512 & 0 & \texttt{-------+--}\\
7 & 6 & 2 & 1 & 2 & 0 & 1024 & 0 & \texttt{----------}\\
8 & 0 & 1 & 1 & 2 & 512 & 512 & 0 & \texttt{---------+}\\
9 & 8 & 2 & 1 & 2 & 0 & 1024 & 0 & \texttt{----------}\\
\multicolumn{9}{@{}l@{}}{\rule{0pt}{2.6ex}\textbf{T10.485}\enspace \texttt{0123431211}\enspace $n{=}10$, $|H|{=}5$, len 5, br 2, $|\mathrm{Aut}|{=}2$, orb 9/8, hist-orb 4/3; M 18432, res 9788, rev 4678: ok}\\*
0 & -- & 0 & 5 & 1 & 162 & 2 & 860 & \texttt{-?????????}\\*
1 & 0 & 1 & 2 & 1 & 704 & 64 & 256 & \texttt{--+???----}\\
2 & 1 & 2 & 2 & 1 & 384 & 128 & 512 & \texttt{---???----}\\
3 & 2 & 3 & 1 & 1 & 512 & 512 & 0 & \texttt{----+-----}\\
4 & 3 & 4 & 1 & 1 & 0 & 1024 & 0 & \texttt{----------}\\
5 & 2 & 3 & 1 & 1 & 0 & 1024 & 0 & \texttt{----------}\\
6 & 0 & 1 & 1 & 1 & 512 & 512 & 0 & \texttt{-------+--}\\
7 & 6 & 2 & 1 & 1 & 0 & 1024 & 0 & \texttt{----------}\\
8 & 0 & 1 & 1 & 2 & 0 & 1024 & 0 & \texttt{----------}\\
9 & 0 & 1 & 1 & 2 & 0 & 1024 & 0 & \texttt{----------}\\
\multicolumn{9}{@{}l@{}}{\rule{0pt}{2.6ex}\textbf{T10.486}\enspace \texttt{0123431111}\enspace $n{=}10$, $|H|{=}6$, len 5, br 2, $|\mathrm{Aut}|{=}24$, orb 7/6, hist-orb 3/2; M 17408, res 10760, rev 3558: ok}\\*
0 & -- & 0 & 6 & 1 & 54 & 2 & 968 & \texttt{-?????????}\\*
1 & 0 & 1 & 2 & 1 & 704 & 64 & 256 & \texttt{--+???----}\\
2 & 1 & 2 & 2 & 1 & 384 & 128 & 512 & \texttt{---???----}\\
3 & 2 & 3 & 1 & 1 & 512 & 512 & 0 & \texttt{----+-----}\\
4 & 3 & 4 & 1 & 1 & 0 & 1024 & 0 & \texttt{----------}\\
5 & 2 & 3 & 1 & 1 & 0 & 1024 & 0 & \texttt{----------}\\
6 & 0 & 1 & 1 & 4 & 0 & 1024 & 0 & \texttt{----------}\\
7 & 0 & 1 & 1 & 4 & 0 & 1024 & 0 & \texttt{----------}\\
8 & 0 & 1 & 1 & 4 & 0 & 1024 & 0 & \texttt{----------}\\
9 & 0 & 1 & 1 & 4 & 0 & 1024 & 0 & \texttt{----------}\\
\multicolumn{9}{@{}l@{}}{\rule{0pt}{2.6ex}\textbf{T10.487}\enspace \texttt{0123423423}\enspace $n{=}10$, $|H|{=}3$, len 5, br 1, $|\mathrm{Aut}|{=}2$, orb 7/4, hist-orb 2/1; M 24576, res 5184, rev 10094: ok}\\*
0 & -- & 0 & 3 & 1 & 806 & 2 & 216 & \texttt{-+????????}\\*
1 & 0 & 1 & 3 & 1 & 588 & 4 & 432 & \texttt{--????????}\\
2 & 1 & 2 & 1 & 2 & 768 & 256 & 0 & \texttt{---++-----}\\
3 & 2 & 3 & 1 & 2 & 512 & 512 & 0 & \texttt{----+-----}\\
4 & 3 & 4 & 1 & 2 & 0 & 1024 & 0 & \texttt{----------}\\
5 & 1 & 2 & 1 & 2 & 768 & 256 & 0 & \texttt{------++--}\\
6 & 5 & 3 & 1 & 2 & 512 & 512 & 0 & \texttt{-------+--}\\
7 & 6 & 4 & 1 & 2 & 0 & 1024 & 0 & \texttt{----------}\\
8 & 1 & 2 & 1 & 1 & 512 & 512 & 0 & \texttt{---------+}\\
9 & 8 & 3 & 1 & 1 & 0 & 1024 & 0 & \texttt{----------}\\
\multicolumn{9}{@{}l@{}}{\rule{0pt}{2.6ex}\textbf{T10.488}\enspace \texttt{0123423422}\enspace $n{=}10$, $|H|{=}4$, len 5, br 1, $|\mathrm{Aut}|{=}4$, orb 6/2, hist-orb 2/0; M 23552, res 9888, rev 7402: ok}\\*
0 & -- & 0 & 4 & 1 & 610 & 2 & 412 & \texttt{-+????????}\\*
1 & 0 & 1 & 4 & 1 & 196 & 4 & 824 & \texttt{--????????}\\
2 & 1 & 2 & 1 & 2 & 768 & 256 & 0 & \texttt{---++-----}\\
3 & 2 & 3 & 1 & 2 & 512 & 512 & 0 & \texttt{----+-----}\\
4 & 3 & 4 & 1 & 2 & 0 & 1024 & 0 & \texttt{----------}\\
5 & 1 & 2 & 1 & 2 & 768 & 256 & 0 & \texttt{------++--}\\
6 & 5 & 3 & 1 & 2 & 512 & 512 & 0 & \texttt{-------+--}\\
7 & 6 & 4 & 1 & 2 & 0 & 1024 & 0 & \texttt{----------}\\
8 & 1 & 2 & 1 & 2 & 0 & 1024 & 0 & \texttt{----------}\\
9 & 1 & 2 & 1 & 2 & 0 & 1024 & 0 & \texttt{----------}\\
\multicolumn{9}{@{}l@{}}{\rule{0pt}{2.6ex}\textbf{T10.489}\enspace \texttt{0123423421}\enspace $n{=}10$, $|H|{=}4$, len 5, br 2, $|\mathrm{Aut}|{=}2$, orb 7/4, hist-orb 3/2; M 22528, res 10068, rev 6610: ok}\\*
0 & -- & 0 & 4 & 1 & 354 & 2 & 668 & \texttt{-?????????}\\*
1 & 0 & 1 & 3 & 1 & 392 & 8 & 624 & \texttt{--???????-}\\
2 & 1 & 2 & 1 & 2 & 768 & 256 & 0 & \texttt{---++-----}\\
3 & 2 & 3 & 1 & 2 & 512 & 512 & 0 & \texttt{----+-----}\\
4 & 3 & 4 & 1 & 2 & 0 & 1024 & 0 & \texttt{----------}\\
5 & 1 & 2 & 1 & 2 & 768 & 256 & 0 & \texttt{------++--}\\
6 & 5 & 3 & 1 & 2 & 512 & 512 & 0 & \texttt{-------+--}\\
7 & 6 & 4 & 1 & 2 & 0 & 1024 & 0 & \texttt{----------}\\
8 & 1 & 2 & 1 & 1 & 0 & 1024 & 0 & \texttt{----------}\\
9 & 0 & 1 & 1 & 1 & 0 & 1024 & 0 & \texttt{----------}\\
\multicolumn{9}{@{}l@{}}{\rule{0pt}{2.6ex}\textbf{T10.490}\enspace \texttt{0123423412}\enspace $n{=}10$, $|H|{=}3$, len 5, br 2, $|\mathrm{Aut}|{=}2$, orb 7/4, hist-orb 2/1; M 22528, res 4328, rev 9226: ok}\\*
0 & -- & 0 & 3 & 1 & 678 & 2 & 344 & \texttt{-?????????}\\*
1 & 0 & 1 & 2 & 1 & 784 & 16 & 224 & \texttt{--??????--}\\
2 & 1 & 2 & 1 & 2 & 768 & 256 & 0 & \texttt{---++-----}\\
3 & 2 & 3 & 1 & 2 & 512 & 512 & 0 & \texttt{----+-----}\\
4 & 3 & 4 & 1 & 2 & 0 & 1024 & 0 & \texttt{----------}\\
5 & 1 & 2 & 1 & 2 & 768 & 256 & 0 & \texttt{------++--}\\
6 & 5 & 3 & 1 & 2 & 512 & 512 & 0 & \texttt{-------+--}\\
7 & 6 & 4 & 1 & 2 & 0 & 1024 & 0 & \texttt{----------}\\
8 & 0 & 1 & 1 & 1 & 512 & 512 & 0 & \texttt{---------+}\\
9 & 8 & 2 & 1 & 1 & 0 & 1024 & 0 & \texttt{----------}\\
\multicolumn{9}{@{}l@{}}{\rule{0pt}{2.6ex}\textbf{T10.491}\enspace \texttt{0123423411}\enspace $n{=}10$, $|H|{=}4$, len 5, br 2, $|\mathrm{Aut}|{=}4$, orb 6/2, hist-orb 2/0; M 21504, res 8396, rev 6746: ok}\\*
0 & -- & 0 & 4 & 1 & 226 & 2 & 796 & \texttt{-?????????}\\*
1 & 0 & 1 & 2 & 1 & 784 & 16 & 224 & \texttt{--??????--}\\
2 & 1 & 2 & 1 & 2 & 768 & 256 & 0 & \texttt{---++-----}\\
3 & 2 & 3 & 1 & 2 & 512 & 512 & 0 & \texttt{----+-----}\\
4 & 3 & 4 & 1 & 2 & 0 & 1024 & 0 & \texttt{----------}\\
5 & 1 & 2 & 1 & 2 & 768 & 256 & 0 & \texttt{------++--}\\
6 & 5 & 3 & 1 & 2 & 512 & 512 & 0 & \texttt{-------+--}\\
7 & 6 & 4 & 1 & 2 & 0 & 1024 & 0 & \texttt{----------}\\
8 & 0 & 1 & 1 & 2 & 0 & 1024 & 0 & \texttt{----------}\\
9 & 0 & 1 & 1 & 2 & 0 & 1024 & 0 & \texttt{----------}\\
\multicolumn{9}{@{}l@{}}{\rule{0pt}{2.6ex}\textbf{T10.492}\enspace \texttt{0123423333}\enspace $n{=}10$, $|H|{=}5$, len 5, br 2, $|\mathrm{Aut}|{=}24$, orb 7/6, hist-orb 2/1; M 24576, res 9440, rev 8790: ok}\\*
0 & -- & 0 & 5 & 1 & 750 & 2 & 272 & \texttt{-+????????}\\*
1 & 0 & 1 & 5 & 1 & 476 & 4 & 544 & \texttt{--????????}\\
2 & 1 & 2 & 1 & 1 & 768 & 256 & 0 & \texttt{---++-----}\\
3 & 2 & 3 & 1 & 1 & 512 & 512 & 0 & \texttt{----+-----}\\
4 & 3 & 4 & 1 & 1 & 0 & 1024 & 0 & \texttt{----------}\\
5 & 1 & 2 & 4 & 1 & 64 & 64 & 896 & \texttt{------????}\\
6 & 5 & 3 & 1 & 4 & 0 & 1024 & 0 & \texttt{----------}\\
7 & 5 & 3 & 1 & 4 & 0 & 1024 & 0 & \texttt{----------}\\
8 & 5 & 3 & 1 & 4 & 0 & 1024 & 0 & \texttt{----------}\\
9 & 5 & 3 & 1 & 4 & 0 & 1024 & 0 & \texttt{----------}\\
\multicolumn{9}{@{}l@{}}{\rule{0pt}{2.6ex}\textbf{T10.493}\enspace \texttt{0123423332}\enspace $n{=}10$, $|H|{=}5$, len 5, br 2, $|\mathrm{Aut}|{=}6$, orb 8/7, hist-orb 3/2; M 23552, res 10944, rev 7078: ok}\\*
0 & -- & 0 & 5 & 1 & 638 & 2 & 384 & \texttt{-+????????}\\*
1 & 0 & 1 & 5 & 1 & 252 & 4 & 768 & \texttt{--????????}\\
2 & 1 & 2 & 1 & 1 & 768 & 256 & 0 & \texttt{---++-----}\\
3 & 2 & 3 & 1 & 1 & 512 & 512 & 0 & \texttt{----+-----}\\
4 & 3 & 4 & 1 & 1 & 0 & 1024 & 0 & \texttt{----------}\\
5 & 1 & 2 & 3 & 1 & 128 & 128 & 768 & \texttt{------???-}\\
6 & 5 & 3 & 1 & 3 & 0 & 1024 & 0 & \texttt{----------}\\
7 & 5 & 3 & 1 & 3 & 0 & 1024 & 0 & \texttt{----------}\\
8 & 5 & 3 & 1 & 3 & 0 & 1024 & 0 & \texttt{----------}\\
9 & 1 & 2 & 1 & 1 & 0 & 1024 & 0 & \texttt{----------}\\
\multicolumn{9}{@{}l@{}}{\rule{0pt}{2.6ex}\textbf{T10.494}\enspace \texttt{0123423331}\enspace $n{=}10$, $|H|{=}5$, len 5, br 3, $|\mathrm{Aut}|{=}6$, orb 8/7, hist-orb 3/2; M 22528, res 10816, rev 6502: ok}\\*
0 & -- & 0 & 5 & 1 & 382 & 2 & 640 & \texttt{-?????????}\\*
1 & 0 & 1 & 4 & 1 & 504 & 8 & 512 & \texttt{--???????-}\\
2 & 1 & 2 & 1 & 1 & 768 & 256 & 0 & \texttt{---++-----}\\
3 & 2 & 3 & 1 & 1 & 512 & 512 & 0 & \texttt{----+-----}\\
4 & 3 & 4 & 1 & 1 & 0 & 1024 & 0 & \texttt{----------}\\
5 & 1 & 2 & 3 & 1 & 128 & 128 & 768 & \texttt{------???-}\\
6 & 5 & 3 & 1 & 3 & 0 & 1024 & 0 & \texttt{----------}\\
7 & 5 & 3 & 1 & 3 & 0 & 1024 & 0 & \texttt{----------}\\
8 & 5 & 3 & 1 & 3 & 0 & 1024 & 0 & \texttt{----------}\\
9 & 0 & 1 & 1 & 1 & 0 & 1024 & 0 & \texttt{----------}\\
\multicolumn{9}{@{}l@{}}{\rule{0pt}{2.6ex}\textbf{T10.495}\enspace \texttt{0123423323}\enspace $n{=}10$, $|H|{=}4$, len 5, br 2, $|\mathrm{Aut}|{=}2$, orb 9/8, hist-orb 3/2; M 23552, res 7840, rev 8594: ok}\\*
0 & -- & 0 & 4 & 1 & 722 & 2 & 300 & \texttt{-+????????}\\*
1 & 0 & 1 & 4 & 1 & 420 & 4 & 600 & \texttt{--????????}\\
2 & 1 & 2 & 1 & 1 & 768 & 256 & 0 & \texttt{---++-----}\\
3 & 2 & 3 & 1 & 1 & 512 & 512 & 0 & \texttt{----+-----}\\
4 & 3 & 4 & 1 & 1 & 0 & 1024 & 0 & \texttt{----------}\\
5 & 1 & 2 & 2 & 1 & 256 & 256 & 512 & \texttt{------??--}\\
6 & 5 & 3 & 1 & 2 & 0 & 1024 & 0 & \texttt{----------}\\
7 & 5 & 3 & 1 & 2 & 0 & 1024 & 0 & \texttt{----------}\\
8 & 1 & 2 & 1 & 1 & 512 & 512 & 0 & \texttt{---------+}\\
9 & 8 & 3 & 1 & 1 & 0 & 1024 & 0 & \texttt{----------}\\
\multicolumn{9}{@{}l@{}}{\rule{0pt}{2.6ex}\textbf{T10.496}\enspace \texttt{0123423322}\enspace $n{=}10$, $|H|{=}5$, len 5, br 2, $|\mathrm{Aut}|{=}4$, orb 8/6, hist-orb 3/1; M 22528, res 11200, rev 6334: ok}\\*
0 & -- & 0 & 5 & 1 & 582 & 2 & 440 & \texttt{-+????????}\\*
1 & 0 & 1 & 5 & 1 & 140 & 4 & 880 & \texttt{--????????}\\
2 & 1 & 2 & 1 & 1 & 768 & 256 & 0 & \texttt{---++-----}\\
3 & 2 & 3 & 1 & 1 & 512 & 512 & 0 & \texttt{----+-----}\\
4 & 3 & 4 & 1 & 1 & 0 & 1024 & 0 & \texttt{----------}\\
5 & 1 & 2 & 2 & 1 & 256 & 256 & 512 & \texttt{------??--}\\
6 & 5 & 3 & 1 & 2 & 0 & 1024 & 0 & \texttt{----------}\\
7 & 5 & 3 & 1 & 2 & 0 & 1024 & 0 & \texttt{----------}\\
8 & 1 & 2 & 1 & 2 & 0 & 1024 & 0 & \texttt{----------}\\
9 & 1 & 2 & 1 & 2 & 0 & 1024 & 0 & \texttt{----------}\\
\multicolumn{9}{@{}l@{}}{\rule{0pt}{2.6ex}\textbf{T10.497}\enspace \texttt{0123423321}\enspace $n{=}10$, $|H|{=}5$, len 5, br 3, $|\mathrm{Aut}|{=}2$, orb 9/8, hist-orb 4/3; M 21504, res 11688, rev 5326: ok}\\*
0 & -- & 0 & 5 & 1 & 326 & 2 & 696 & \texttt{-?????????}\\*
1 & 0 & 1 & 4 & 1 & 280 & 8 & 736 & \texttt{--???????-}\\
2 & 1 & 2 & 1 & 1 & 768 & 256 & 0 & \texttt{---++-----}\\
3 & 2 & 3 & 1 & 1 & 512 & 512 & 0 & \texttt{----+-----}\\
4 & 3 & 4 & 1 & 1 & 0 & 1024 & 0 & \texttt{----------}\\
5 & 1 & 2 & 2 & 1 & 256 & 256 & 512 & \texttt{------??--}\\
6 & 5 & 3 & 1 & 2 & 0 & 1024 & 0 & \texttt{----------}\\
7 & 5 & 3 & 1 & 2 & 0 & 1024 & 0 & \texttt{----------}\\
8 & 1 & 2 & 1 & 1 & 0 & 1024 & 0 & \texttt{----------}\\
9 & 0 & 1 & 1 & 1 & 0 & 1024 & 0 & \texttt{----------}\\
\multicolumn{9}{@{}l@{}}{\rule{0pt}{2.6ex}\textbf{T10.498}\enspace \texttt{0123423312}\enspace $n{=}10$, $|H|{=}4$, len 5, br 3, $|\mathrm{Aut}|{=}2$, orb 9/8, hist-orb 3/2; M 21504, res 6956, rev 7798: ok}\\*
0 & -- & 0 & 4 & 1 & 594 & 2 & 428 & \texttt{-?????????}\\*
1 & 0 & 1 & 3 & 1 & 560 & 16 & 448 & \texttt{--??????--}\\
2 & 1 & 2 & 1 & 1 & 768 & 256 & 0 & \texttt{---++-----}\\
3 & 2 & 3 & 1 & 1 & 512 & 512 & 0 & \texttt{----+-----}\\
4 & 3 & 4 & 1 & 1 & 0 & 1024 & 0 & \texttt{----------}\\
5 & 1 & 2 & 2 & 1 & 256 & 256 & 512 & \texttt{------??--}\\
6 & 5 & 3 & 1 & 2 & 0 & 1024 & 0 & \texttt{----------}\\
7 & 5 & 3 & 1 & 2 & 0 & 1024 & 0 & \texttt{----------}\\
8 & 0 & 1 & 1 & 1 & 512 & 512 & 0 & \texttt{---------+}\\
9 & 8 & 2 & 1 & 1 & 0 & 1024 & 0 & \texttt{----------}\\
\multicolumn{9}{@{}l@{}}{\rule{0pt}{2.6ex}\textbf{T10.499}\enspace \texttt{0123423311}\enspace $n{=}10$, $|H|{=}5$, len 5, br 3, $|\mathrm{Aut}|{=}4$, orb 8/6, hist-orb 3/1; M 20480, res 10520, rev 5334: ok}\\*
0 & -- & 0 & 5 & 1 & 198 & 2 & 824 & \texttt{-?????????}\\*
1 & 0 & 1 & 3 & 1 & 560 & 16 & 448 & \texttt{--??????--}\\
2 & 1 & 2 & 1 & 1 & 768 & 256 & 0 & \texttt{---++-----}\\
3 & 2 & 3 & 1 & 1 & 512 & 512 & 0 & \texttt{----+-----}\\
4 & 3 & 4 & 1 & 1 & 0 & 1024 & 0 & \texttt{----------}\\
5 & 1 & 2 & 2 & 1 & 256 & 256 & 512 & \texttt{------??--}\\
6 & 5 & 3 & 1 & 2 & 0 & 1024 & 0 & \texttt{----------}\\
7 & 5 & 3 & 1 & 2 & 0 & 1024 & 0 & \texttt{----------}\\
8 & 0 & 1 & 1 & 2 & 0 & 1024 & 0 & \texttt{----------}\\
9 & 0 & 1 & 1 & 2 & 0 & 1024 & 0 & \texttt{----------}\\
\multicolumn{9}{@{}l@{}}{\rule{0pt}{2.6ex}\textbf{T10.500}\enspace \texttt{0123423232}\enspace $n{=}10$, $|H|{=}4$, len 5, br 1, $|\mathrm{Aut}|{=}2$, orb 8/6, hist-orb 3/2; M 22528, res 9216, rev 7486: ok}\\*
0 & -- & 0 & 4 & 1 & 638 & 2 & 384 & \texttt{-+????????}\\*
1 & 0 & 1 & 4 & 1 & 252 & 4 & 768 & \texttt{--????????}\\
2 & 1 & 2 & 1 & 1 & 768 & 256 & 0 & \texttt{---++-----}\\
3 & 2 & 3 & 1 & 1 & 512 & 512 & 0 & \texttt{----+-----}\\
4 & 3 & 4 & 1 & 1 & 0 & 1024 & 0 & \texttt{----------}\\
5 & 1 & 2 & 1 & 2 & 512 & 512 & 0 & \texttt{------+---}\\
6 & 5 & 3 & 1 & 2 & 0 & 1024 & 0 & \texttt{----------}\\
7 & 1 & 2 & 1 & 2 & 512 & 512 & 0 & \texttt{--------+-}\\
8 & 7 & 3 & 1 & 2 & 0 & 1024 & 0 & \texttt{----------}\\
9 & 1 & 2 & 1 & 1 & 0 & 1024 & 0 & \texttt{----------}\\
\multicolumn{9}{@{}l@{}}{\rule{0pt}{2.6ex}\textbf{T10.501}\enspace \texttt{0123423231}\enspace $n{=}10$, $|H|{=}4$, len 5, br 2, $|\mathrm{Aut}|{=}2$, orb 8/6, hist-orb 3/2; M 21504, res 9088, rev 6830: ok}\\*
0 & -- & 0 & 4 & 1 & 382 & 2 & 640 & \texttt{-?????????}\\*
1 & 0 & 1 & 3 & 1 & 504 & 8 & 512 & \texttt{--???????-}\\
2 & 1 & 2 & 1 & 1 & 768 & 256 & 0 & \texttt{---++-----}\\
3 & 2 & 3 & 1 & 1 & 512 & 512 & 0 & \texttt{----+-----}\\
4 & 3 & 4 & 1 & 1 & 0 & 1024 & 0 & \texttt{----------}\\
5 & 1 & 2 & 1 & 2 & 512 & 512 & 0 & \texttt{------+---}\\
6 & 5 & 3 & 1 & 2 & 0 & 1024 & 0 & \texttt{----------}\\
7 & 1 & 2 & 1 & 2 & 512 & 512 & 0 & \texttt{--------+-}\\
8 & 7 & 3 & 1 & 2 & 0 & 1024 & 0 & \texttt{----------}\\
9 & 0 & 1 & 1 & 1 & 0 & 1024 & 0 & \texttt{----------}\\
\multicolumn{9}{@{}l@{}}{\rule{0pt}{2.6ex}\textbf{T10.502}\enspace \texttt{0123423222}\enspace $n{=}10$, $|H|{=}5$, len 5, br 1, $|\mathrm{Aut}|{=}6$, orb 8/7, hist-orb 3/2; M 21504, res 11232, rev 5994: ok}\\*
0 & -- & 0 & 5 & 1 & 554 & 2 & 468 & \texttt{-+????????}\\*
1 & 0 & 1 & 5 & 1 & 84 & 4 & 936 & \texttt{--????????}\\
2 & 1 & 2 & 1 & 1 & 768 & 256 & 0 & \texttt{---++-----}\\
3 & 2 & 3 & 1 & 1 & 512 & 512 & 0 & \texttt{----+-----}\\
4 & 3 & 4 & 1 & 1 & 0 & 1024 & 0 & \texttt{----------}\\
5 & 1 & 2 & 1 & 1 & 512 & 512 & 0 & \texttt{------+---}\\
6 & 5 & 3 & 1 & 1 & 0 & 1024 & 0 & \texttt{----------}\\
7 & 1 & 2 & 1 & 3 & 0 & 1024 & 0 & \texttt{----------}\\
8 & 1 & 2 & 1 & 3 & 0 & 1024 & 0 & \texttt{----------}\\
9 & 1 & 2 & 1 & 3 & 0 & 1024 & 0 & \texttt{----------}\\
\multicolumn{9}{@{}l@{}}{\rule{0pt}{2.6ex}\textbf{T10.503}\enspace \texttt{0123423221}\enspace $n{=}10$, $|H|{=}5$, len 5, br 2, $|\mathrm{Aut}|{=}2$, orb 9/8, hist-orb 4/3; M 20480, res 12028, rev 4722: ok}\\*
0 & -- & 0 & 5 & 1 & 298 & 2 & 724 & \texttt{-?????????}\\*
1 & 0 & 1 & 4 & 1 & 168 & 8 & 848 & \texttt{--???????-}\\
2 & 1 & 2 & 1 & 1 & 768 & 256 & 0 & \texttt{---++-----}\\
3 & 2 & 3 & 1 & 1 & 512 & 512 & 0 & \texttt{----+-----}\\
4 & 3 & 4 & 1 & 1 & 0 & 1024 & 0 & \texttt{----------}\\
5 & 1 & 2 & 1 & 1 & 512 & 512 & 0 & \texttt{------+---}\\
6 & 5 & 3 & 1 & 1 & 0 & 1024 & 0 & \texttt{----------}\\
7 & 1 & 2 & 1 & 2 & 0 & 1024 & 0 & \texttt{----------}\\
8 & 1 & 2 & 1 & 2 & 0 & 1024 & 0 & \texttt{----------}\\
9 & 0 & 1 & 1 & 1 & 0 & 1024 & 0 & \texttt{----------}\\
\multicolumn{9}{@{}l@{}}{\rule{0pt}{2.6ex}\textbf{T10.504}\enspace \texttt{0123423212}\enspace $n{=}10$, $|H|{=}4$, len 5, br 2, $|\mathrm{Aut}|{=}1$, orb 10/10, hist-orb 4/4; M 20480, res 8304, rev 6682: ok}\\*
0 & -- & 0 & 4 & 1 & 510 & 2 & 512 & \texttt{-?????????}\\*
1 & 0 & 1 & 3 & 1 & 336 & 16 & 672 & \texttt{--??????--}\\
2 & 1 & 2 & 1 & 1 & 768 & 256 & 0 & \texttt{---++-----}\\
3 & 2 & 3 & 1 & 1 & 512 & 512 & 0 & \texttt{----+-----}\\
4 & 3 & 4 & 1 & 1 & 0 & 1024 & 0 & \texttt{----------}\\
5 & 1 & 2 & 1 & 1 & 512 & 512 & 0 & \texttt{------+---}\\
6 & 5 & 3 & 1 & 1 & 0 & 1024 & 0 & \texttt{----------}\\
7 & 1 & 2 & 1 & 1 & 0 & 1024 & 0 & \texttt{----------}\\
8 & 0 & 1 & 1 & 1 & 512 & 512 & 0 & \texttt{---------+}\\
9 & 8 & 2 & 1 & 1 & 0 & 1024 & 0 & \texttt{----------}\\
\multicolumn{9}{@{}l@{}}{\rule{0pt}{2.6ex}\textbf{T10.505}\enspace \texttt{0123423211}\enspace $n{=}10$, $|H|{=}5$, len 5, br 2, $|\mathrm{Aut}|{=}2$, orb 9/8, hist-orb 4/3; M 19456, res 11364, rev 4410: ok}\\*
0 & -- & 0 & 5 & 1 & 170 & 2 & 852 & \texttt{-?????????}\\*
1 & 0 & 1 & 3 & 1 & 336 & 16 & 672 & \texttt{--??????--}\\
2 & 1 & 2 & 1 & 1 & 768 & 256 & 0 & \texttt{---++-----}\\
3 & 2 & 3 & 1 & 1 & 512 & 512 & 0 & \texttt{----+-----}\\
4 & 3 & 4 & 1 & 1 & 0 & 1024 & 0 & \texttt{----------}\\
5 & 1 & 2 & 1 & 1 & 512 & 512 & 0 & \texttt{------+---}\\
6 & 5 & 3 & 1 & 1 & 0 & 1024 & 0 & \texttt{----------}\\
7 & 1 & 2 & 1 & 1 & 0 & 1024 & 0 & \texttt{----------}\\
8 & 0 & 1 & 1 & 2 & 0 & 1024 & 0 & \texttt{----------}\\
9 & 0 & 1 & 1 & 2 & 0 & 1024 & 0 & \texttt{----------}\\
\multicolumn{9}{@{}l@{}}{\rule{0pt}{2.6ex}\textbf{T10.506}\enspace \texttt{0123423123}\enspace $n{=}10$, $|H|{=}3$, len 5, br 2, $|\mathrm{Aut}|{=}1$, orb 10/10, hist-orb 3/3; M 21504, res 3960, rev 9026: ok}\\*
0 & -- & 0 & 3 & 1 & 742 & 2 & 280 & \texttt{-?????????}\\*
1 & 0 & 1 & 2 & 1 & 672 & 32 & 320 & \texttt{--?????---}\\
2 & 1 & 2 & 1 & 1 & 768 & 256 & 0 & \texttt{---++-----}\\
3 & 2 & 3 & 1 & 1 & 512 & 512 & 0 & \texttt{----+-----}\\
4 & 3 & 4 & 1 & 1 & 0 & 1024 & 0 & \texttt{----------}\\
5 & 1 & 2 & 1 & 1 & 512 & 512 & 0 & \texttt{------+---}\\
6 & 5 & 3 & 1 & 1 & 0 & 1024 & 0 & \texttt{----------}\\
7 & 0 & 1 & 1 & 1 & 768 & 256 & 0 & \texttt{--------++}\\
8 & 7 & 2 & 1 & 1 & 512 & 512 & 0 & \texttt{---------+}\\
9 & 8 & 3 & 1 & 1 & 0 & 1024 & 0 & \texttt{----------}\\
\multicolumn{9}{@{}l@{}}{\rule{0pt}{2.6ex}\textbf{T10.507}\enspace \texttt{0123423122}\enspace $n{=}10$, $|H|{=}4$, len 5, br 3, $|\mathrm{Aut}|{=}2$, orb 9/8, hist-orb 3/2; M 20480, res 6636, rev 7338: ok}\\*
0 & -- & 0 & 4 & 1 & 530 & 2 & 492 & \texttt{-?????????}\\*
1 & 0 & 1 & 2 & 1 & 672 & 32 & 320 & \texttt{--?????---}\\
2 & 1 & 2 & 1 & 1 & 768 & 256 & 0 & \texttt{---++-----}\\
3 & 2 & 3 & 1 & 1 & 512 & 512 & 0 & \texttt{----+-----}\\
4 & 3 & 4 & 1 & 1 & 0 & 1024 & 0 & \texttt{----------}\\
5 & 1 & 2 & 1 & 1 & 512 & 512 & 0 & \texttt{------+---}\\
6 & 5 & 3 & 1 & 1 & 0 & 1024 & 0 & \texttt{----------}\\
7 & 0 & 1 & 2 & 1 & 256 & 256 & 512 & \texttt{--------??}\\
8 & 7 & 2 & 1 & 2 & 0 & 1024 & 0 & \texttt{----------}\\
9 & 7 & 2 & 1 & 2 & 0 & 1024 & 0 & \texttt{----------}\\
\multicolumn{9}{@{}l@{}}{\rule{0pt}{2.6ex}\textbf{T10.508}\enspace \texttt{0123423121}\enspace $n{=}10$, $|H|{=}4$, len 5, br 2, $|\mathrm{Aut}|{=}1$, orb 10/10, hist-orb 4/4; M 19456, res 7776, rev 6250: ok}\\*
0 & -- & 0 & 4 & 1 & 318 & 2 & 704 & \texttt{-?????????}\\*
1 & 0 & 1 & 2 & 1 & 672 & 32 & 320 & \texttt{--?????---}\\
2 & 1 & 2 & 1 & 1 & 768 & 256 & 0 & \texttt{---++-----}\\
3 & 2 & 3 & 1 & 1 & 512 & 512 & 0 & \texttt{----+-----}\\
4 & 3 & 4 & 1 & 1 & 0 & 1024 & 0 & \texttt{----------}\\
5 & 1 & 2 & 1 & 1 & 512 & 512 & 0 & \texttt{------+---}\\
6 & 5 & 3 & 1 & 1 & 0 & 1024 & 0 & \texttt{----------}\\
7 & 0 & 1 & 1 & 1 & 512 & 512 & 0 & \texttt{--------+-}\\
8 & 7 & 2 & 1 & 1 & 0 & 1024 & 0 & \texttt{----------}\\
9 & 0 & 1 & 1 & 1 & 0 & 1024 & 0 & \texttt{----------}\\
\multicolumn{9}{@{}l@{}}{\rule{0pt}{2.6ex}\textbf{T10.509}\enspace \texttt{0123423111}\enspace $n{=}10$, $|H|{=}5$, len 5, br 2, $|\mathrm{Aut}|{=}6$, orb 8/7, hist-orb 3/2; M 18432, res 9684, rev 4650: ok}\\*
0 & -- & 0 & 5 & 1 & 106 & 2 & 916 & \texttt{-?????????}\\*
1 & 0 & 1 & 2 & 1 & 672 & 32 & 320 & \texttt{--?????---}\\
2 & 1 & 2 & 1 & 1 & 768 & 256 & 0 & \texttt{---++-----}\\
3 & 2 & 3 & 1 & 1 & 512 & 512 & 0 & \texttt{----+-----}\\
4 & 3 & 4 & 1 & 1 & 0 & 1024 & 0 & \texttt{----------}\\
5 & 1 & 2 & 1 & 1 & 512 & 512 & 0 & \texttt{------+---}\\
6 & 5 & 3 & 1 & 1 & 0 & 1024 & 0 & \texttt{----------}\\
7 & 0 & 1 & 1 & 3 & 0 & 1024 & 0 & \texttt{----------}\\
8 & 0 & 1 & 1 & 3 & 0 & 1024 & 0 & \texttt{----------}\\
9 & 0 & 1 & 1 & 3 & 0 & 1024 & 0 & \texttt{----------}\\
\multicolumn{9}{@{}l@{}}{\rule{0pt}{2.6ex}\textbf{T10.510}\enspace \texttt{0123422222}\enspace $n{=}10$, $|H|{=}6$, len 5, br 1, $|\mathrm{Aut}|{=}120$, orb 6/5, hist-orb 2/1; M 20480, res 11904, rev 5270: ok}\\*
0 & -- & 0 & 6 & 1 & 526 & 2 & 496 & \texttt{-+????????}\\*
1 & 0 & 1 & 6 & 1 & 28 & 4 & 992 & \texttt{--????????}\\
2 & 1 & 2 & 1 & 1 & 768 & 256 & 0 & \texttt{---++-----}\\
3 & 2 & 3 & 1 & 1 & 512 & 512 & 0 & \texttt{----+-----}\\
4 & 3 & 4 & 1 & 1 & 0 & 1024 & 0 & \texttt{----------}\\
5 & 1 & 2 & 1 & 5 & 0 & 1024 & 0 & \texttt{----------}\\
6 & 1 & 2 & 1 & 5 & 0 & 1024 & 0 & \texttt{----------}\\
7 & 1 & 2 & 1 & 5 & 0 & 1024 & 0 & \texttt{----------}\\
8 & 1 & 2 & 1 & 5 & 0 & 1024 & 0 & \texttt{----------}\\
9 & 1 & 2 & 1 & 5 & 0 & 1024 & 0 & \texttt{----------}\\
\multicolumn{9}{@{}l@{}}{\rule{0pt}{2.6ex}\textbf{T10.511}\enspace \texttt{0123422221}\enspace $n{=}10$, $|H|{=}6$, len 5, br 2, $|\mathrm{Aut}|{=}24$, orb 7/6, hist-orb 3/2; M 19456, res 13008, rev 3670: ok}\\*
0 & -- & 0 & 6 & 1 & 270 & 2 & 752 & \texttt{-?????????}\\*
1 & 0 & 1 & 5 & 1 & 56 & 8 & 960 & \texttt{--???????-}\\
2 & 1 & 2 & 1 & 1 & 768 & 256 & 0 & \texttt{---++-----}\\
3 & 2 & 3 & 1 & 1 & 512 & 512 & 0 & \texttt{----+-----}\\
4 & 3 & 4 & 1 & 1 & 0 & 1024 & 0 & \texttt{----------}\\
5 & 1 & 2 & 1 & 4 & 0 & 1024 & 0 & \texttt{----------}\\
6 & 1 & 2 & 1 & 4 & 0 & 1024 & 0 & \texttt{----------}\\
7 & 1 & 2 & 1 & 4 & 0 & 1024 & 0 & \texttt{----------}\\
8 & 1 & 2 & 1 & 4 & 0 & 1024 & 0 & \texttt{----------}\\
9 & 0 & 1 & 1 & 1 & 0 & 1024 & 0 & \texttt{----------}\\
\multicolumn{9}{@{}l@{}}{\rule{0pt}{2.6ex}\textbf{T10.512}\enspace \texttt{0123422212}\enspace $n{=}10$, $|H|{=}5$, len 5, br 2, $|\mathrm{Aut}|{=}6$, orb 8/7, hist-orb 3/2; M 19456, res 10292, rev 5150: ok}\\*
0 & -- & 0 & 5 & 1 & 426 & 2 & 596 & \texttt{-?????????}\\*
1 & 0 & 1 & 4 & 1 & 112 & 16 & 896 & \texttt{--??????--}\\
2 & 1 & 2 & 1 & 1 & 768 & 256 & 0 & \texttt{---++-----}\\
3 & 2 & 3 & 1 & 1 & 512 & 512 & 0 & \texttt{----+-----}\\
4 & 3 & 4 & 1 & 1 & 0 & 1024 & 0 & \texttt{----------}\\
5 & 1 & 2 & 1 & 3 & 0 & 1024 & 0 & \texttt{----------}\\
6 & 1 & 2 & 1 & 3 & 0 & 1024 & 0 & \texttt{----------}\\
7 & 1 & 2 & 1 & 3 & 0 & 1024 & 0 & \texttt{----------}\\
8 & 0 & 1 & 1 & 1 & 512 & 512 & 0 & \texttt{---------+}\\
9 & 8 & 2 & 1 & 1 & 0 & 1024 & 0 & \texttt{----------}\\
\multicolumn{9}{@{}l@{}}{\rule{0pt}{2.6ex}\textbf{T10.513}\enspace \texttt{0123422211}\enspace $n{=}10$, $|H|{=}6$, len 5, br 2, $|\mathrm{Aut}|{=}12$, orb 7/5, hist-orb 3/1; M 18432, res 12848, rev 3006: ok}\\*
0 & -- & 0 & 6 & 1 & 142 & 2 & 880 & \texttt{-?????????}\\*
1 & 0 & 1 & 4 & 1 & 112 & 16 & 896 & \texttt{--??????--}\\
2 & 1 & 2 & 1 & 1 & 768 & 256 & 0 & \texttt{---++-----}\\
3 & 2 & 3 & 1 & 1 & 512 & 512 & 0 & \texttt{----+-----}\\
4 & 3 & 4 & 1 & 1 & 0 & 1024 & 0 & \texttt{----------}\\
5 & 1 & 2 & 1 & 3 & 0 & 1024 & 0 & \texttt{----------}\\
6 & 1 & 2 & 1 & 3 & 0 & 1024 & 0 & \texttt{----------}\\
7 & 1 & 2 & 1 & 3 & 0 & 1024 & 0 & \texttt{----------}\\
8 & 0 & 1 & 1 & 2 & 0 & 1024 & 0 & \texttt{----------}\\
9 & 0 & 1 & 1 & 2 & 0 & 1024 & 0 & \texttt{----------}\\
\multicolumn{9}{@{}l@{}}{\rule{0pt}{2.6ex}\textbf{T10.514}\enspace \texttt{0123422123}\enspace $n{=}10$, $|H|{=}4$, len 5, br 2, $|\mathrm{Aut}|{=}2$, orb 9/8, hist-orb 3/2; M 20480, res 7740, rev 6774: ok}\\*
0 & -- & 0 & 4 & 1 & 546 & 2 & 476 & \texttt{-?????????}\\*
1 & 0 & 1 & 3 & 1 & 224 & 32 & 768 & \texttt{--?????---}\\
2 & 1 & 2 & 1 & 1 & 768 & 256 & 0 & \texttt{---++-----}\\
3 & 2 & 3 & 1 & 1 & 512 & 512 & 0 & \texttt{----+-----}\\
4 & 3 & 4 & 1 & 1 & 0 & 1024 & 0 & \texttt{----------}\\
5 & 1 & 2 & 1 & 2 & 0 & 1024 & 0 & \texttt{----------}\\
6 & 1 & 2 & 1 & 2 & 0 & 1024 & 0 & \texttt{----------}\\
7 & 0 & 1 & 1 & 1 & 768 & 256 & 0 & \texttt{--------++}\\
8 & 7 & 2 & 1 & 1 & 512 & 512 & 0 & \texttt{---------+}\\
9 & 8 & 3 & 1 & 1 & 0 & 1024 & 0 & \texttt{----------}\\
\multicolumn{9}{@{}l@{}}{\rule{0pt}{2.6ex}\textbf{T10.515}\enspace \texttt{0123422122}\enspace $n{=}10$, $|H|{=}5$, len 5, br 3, $|\mathrm{Aut}|{=}4$, orb 8/6, hist-orb 3/1; M 19456, res 9912, rev 5166: ok}\\*
0 & -- & 0 & 5 & 1 & 390 & 2 & 632 & \texttt{-?????????}\\*
1 & 0 & 1 & 3 & 1 & 224 & 32 & 768 & \texttt{--?????---}\\
2 & 1 & 2 & 1 & 1 & 768 & 256 & 0 & \texttt{---++-----}\\
3 & 2 & 3 & 1 & 1 & 512 & 512 & 0 & \texttt{----+-----}\\
4 & 3 & 4 & 1 & 1 & 0 & 1024 & 0 & \texttt{----------}\\
5 & 1 & 2 & 1 & 2 & 0 & 1024 & 0 & \texttt{----------}\\
6 & 1 & 2 & 1 & 2 & 0 & 1024 & 0 & \texttt{----------}\\
7 & 0 & 1 & 2 & 1 & 256 & 256 & 512 & \texttt{--------??}\\
8 & 7 & 2 & 1 & 2 & 0 & 1024 & 0 & \texttt{----------}\\
9 & 7 & 2 & 1 & 2 & 0 & 1024 & 0 & \texttt{----------}\\
\multicolumn{9}{@{}l@{}}{\rule{0pt}{2.6ex}\textbf{T10.516}\enspace \texttt{0123422121}\enspace $n{=}10$, $|H|{=}5$, len 5, br 2, $|\mathrm{Aut}|{=}2$, orb 9/8, hist-orb 4/3; M 18432, res 10548, rev 4270: ok}\\*
0 & -- & 0 & 5 & 1 & 234 & 2 & 788 & \texttt{-?????????}\\*
1 & 0 & 1 & 3 & 1 & 224 & 32 & 768 & \texttt{--?????---}\\
2 & 1 & 2 & 1 & 1 & 768 & 256 & 0 & \texttt{---++-----}\\
3 & 2 & 3 & 1 & 1 & 512 & 512 & 0 & \texttt{----+-----}\\
4 & 3 & 4 & 1 & 1 & 0 & 1024 & 0 & \texttt{----------}\\
5 & 1 & 2 & 1 & 2 & 0 & 1024 & 0 & \texttt{----------}\\
6 & 1 & 2 & 1 & 2 & 0 & 1024 & 0 & \texttt{----------}\\
7 & 0 & 1 & 1 & 1 & 512 & 512 & 0 & \texttt{--------+-}\\
8 & 7 & 2 & 1 & 1 & 0 & 1024 & 0 & \texttt{----------}\\
9 & 0 & 1 & 1 & 1 & 0 & 1024 & 0 & \texttt{----------}\\
\multicolumn{9}{@{}l@{}}{\rule{0pt}{2.6ex}\textbf{T10.517}\enspace \texttt{0123422111}\enspace $n{=}10$, $|H|{=}6$, len 5, br 2, $|\mathrm{Aut}|{=}12$, orb 7/5, hist-orb 3/1; M 17408, res 11952, rev 2862: ok}\\*
0 & -- & 0 & 6 & 1 & 78 & 2 & 944 & \texttt{-?????????}\\*
1 & 0 & 1 & 3 & 1 & 224 & 32 & 768 & \texttt{--?????---}\\
2 & 1 & 2 & 1 & 1 & 768 & 256 & 0 & \texttt{---++-----}\\
3 & 2 & 3 & 1 & 1 & 512 & 512 & 0 & \texttt{----+-----}\\
4 & 3 & 4 & 1 & 1 & 0 & 1024 & 0 & \texttt{----------}\\
5 & 1 & 2 & 1 & 2 & 0 & 1024 & 0 & \texttt{----------}\\
6 & 1 & 2 & 1 & 2 & 0 & 1024 & 0 & \texttt{----------}\\
7 & 0 & 1 & 1 & 3 & 0 & 1024 & 0 & \texttt{----------}\\
8 & 0 & 1 & 1 & 3 & 0 & 1024 & 0 & \texttt{----------}\\
9 & 0 & 1 & 1 & 3 & 0 & 1024 & 0 & \texttt{----------}\\
\multicolumn{9}{@{}l@{}}{\rule{0pt}{2.6ex}\textbf{T10.518}\enspace \texttt{0123421234}\enspace $n{=}10$, $|H|{=}3$, len 5, br 2, $|\mathrm{Aut}|{=}1$, orb 10/10, hist-orb 3/3; M 22528, res 4780, rev 8718: ok}\\*
0 & -- & 0 & 3 & 1 & 690 & 2 & 332 & \texttt{-?????????}\\*
1 & 0 & 1 & 2 & 1 & 448 & 64 & 512 & \texttt{--????----}\\
2 & 1 & 2 & 1 & 1 & 768 & 256 & 0 & \texttt{---++-----}\\
3 & 2 & 3 & 1 & 1 & 512 & 512 & 0 & \texttt{----+-----}\\
4 & 3 & 4 & 1 & 1 & 0 & 1024 & 0 & \texttt{----------}\\
5 & 1 & 2 & 1 & 1 & 0 & 1024 & 0 & \texttt{----------}\\
6 & 0 & 1 & 1 & 1 & 896 & 128 & 0 & \texttt{-------+++}\\
7 & 6 & 2 & 1 & 1 & 768 & 256 & 0 & \texttt{--------++}\\
8 & 7 & 3 & 1 & 1 & 512 & 512 & 0 & \texttt{---------+}\\
9 & 8 & 4 & 1 & 1 & 0 & 1024 & 0 & \texttt{----------}\\
\multicolumn{9}{@{}l@{}}{\rule{0pt}{2.6ex}\textbf{T10.519}\enspace \texttt{0123421233}\enspace $n{=}10$, $|H|{=}4$, len 5, br 3, $|\mathrm{Aut}|{=}2$, orb 9/8, hist-orb 3/2; M 21504, res 6888, rev 7626: ok}\\*
0 & -- & 0 & 4 & 1 & 598 & 2 & 424 & \texttt{-?????????}\\*
1 & 0 & 1 & 2 & 1 & 448 & 64 & 512 & \texttt{--????----}\\
2 & 1 & 2 & 1 & 1 & 768 & 256 & 0 & \texttt{---++-----}\\
3 & 2 & 3 & 1 & 1 & 512 & 512 & 0 & \texttt{----+-----}\\
4 & 3 & 4 & 1 & 1 & 0 & 1024 & 0 & \texttt{----------}\\
5 & 1 & 2 & 1 & 1 & 0 & 1024 & 0 & \texttt{----------}\\
6 & 0 & 1 & 2 & 1 & 640 & 128 & 256 & \texttt{-------+??}\\
7 & 6 & 2 & 2 & 1 & 256 & 256 & 512 & \texttt{--------??}\\
8 & 7 & 3 & 1 & 2 & 0 & 1024 & 0 & \texttt{----------}\\
9 & 7 & 3 & 1 & 2 & 0 & 1024 & 0 & \texttt{----------}\\
\multicolumn{9}{@{}l@{}}{\rule{0pt}{2.6ex}\textbf{T10.520}\enspace \texttt{0123421232}\enspace $n{=}10$, $|H|{=}4$, len 5, br 3, $|\mathrm{Aut}|{=}1$, orb 10/10, hist-orb 4/4; M 20480, res 7716, rev 6770: ok}\\*
0 & -- & 0 & 4 & 1 & 506 & 2 & 516 & \texttt{-?????????}\\*
1 & 0 & 1 & 2 & 1 & 448 & 64 & 512 & \texttt{--????----}\\
2 & 1 & 2 & 1 & 1 & 768 & 256 & 0 & \texttt{---++-----}\\
3 & 2 & 3 & 1 & 1 & 512 & 512 & 0 & \texttt{----+-----}\\
4 & 3 & 4 & 1 & 1 & 0 & 1024 & 0 & \texttt{----------}\\
5 & 1 & 2 & 1 & 1 & 0 & 1024 & 0 & \texttt{----------}\\
6 & 0 & 1 & 2 & 1 & 384 & 128 & 512 & \texttt{-------???}\\
7 & 6 & 2 & 1 & 1 & 512 & 512 & 0 & \texttt{--------+-}\\
8 & 7 & 3 & 1 & 1 & 0 & 1024 & 0 & \texttt{----------}\\
9 & 6 & 2 & 1 & 1 & 0 & 1024 & 0 & \texttt{----------}\\
\multicolumn{9}{@{}l@{}}{\rule{0pt}{2.6ex}\textbf{T10.521}\enspace \texttt{0123421231}\enspace $n{=}10$, $|H|{=}4$, len 5, br 2, $|\mathrm{Aut}|{=}1$, orb 10/10, hist-orb 4/4; M 19456, res 8092, rev 5934: ok}\\*
0 & -- & 0 & 4 & 1 & 322 & 2 & 700 & \texttt{-?????????}\\*
1 & 0 & 1 & 2 & 1 & 448 & 64 & 512 & \texttt{--????----}\\
2 & 1 & 2 & 1 & 1 & 768 & 256 & 0 & \texttt{---++-----}\\
3 & 2 & 3 & 1 & 1 & 512 & 512 & 0 & \texttt{----+-----}\\
4 & 3 & 4 & 1 & 1 & 0 & 1024 & 0 & \texttt{----------}\\
5 & 1 & 2 & 1 & 1 & 0 & 1024 & 0 & \texttt{----------}\\
6 & 0 & 1 & 1 & 1 & 768 & 256 & 0 & \texttt{-------++-}\\
7 & 6 & 2 & 1 & 1 & 512 & 512 & 0 & \texttt{--------+-}\\
8 & 7 & 3 & 1 & 1 & 0 & 1024 & 0 & \texttt{----------}\\
9 & 0 & 1 & 1 & 1 & 0 & 1024 & 0 & \texttt{----------}\\
\multicolumn{9}{@{}l@{}}{\rule{0pt}{2.6ex}\textbf{T10.522}\enspace \texttt{0123421222}\enspace $n{=}10$, $|H|{=}5$, len 5, br 3, $|\mathrm{Aut}|{=}6$, orb 8/7, hist-orb 3/2; M 19456, res 9184, rev 5494: ok}\\*
0 & -- & 0 & 5 & 1 & 414 & 2 & 608 & \texttt{-?????????}\\*
1 & 0 & 1 & 2 & 1 & 448 & 64 & 512 & \texttt{--????----}\\
2 & 1 & 2 & 1 & 1 & 768 & 256 & 0 & \texttt{---++-----}\\
3 & 2 & 3 & 1 & 1 & 512 & 512 & 0 & \texttt{----+-----}\\
4 & 3 & 4 & 1 & 1 & 0 & 1024 & 0 & \texttt{----------}\\
5 & 1 & 2 & 1 & 1 & 0 & 1024 & 0 & \texttt{----------}\\
6 & 0 & 1 & 3 & 1 & 128 & 128 & 768 & \texttt{-------???}\\
7 & 6 & 2 & 1 & 3 & 0 & 1024 & 0 & \texttt{----------}\\
8 & 6 & 2 & 1 & 3 & 0 & 1024 & 0 & \texttt{----------}\\
9 & 6 & 2 & 1 & 3 & 0 & 1024 & 0 & \texttt{----------}\\
\multicolumn{9}{@{}l@{}}{\rule{0pt}{2.6ex}\textbf{T10.523}\enspace \texttt{0123421221}\enspace $n{=}10$, $|H|{=}5$, len 5, br 3, $|\mathrm{Aut}|{=}2$, orb 9/8, hist-orb 4/3; M 18432, res 9688, rev 4566: ok}\\*
0 & -- & 0 & 5 & 1 & 230 & 2 & 792 & \texttt{-?????????}\\*
1 & 0 & 1 & 2 & 1 & 448 & 64 & 512 & \texttt{--????----}\\
2 & 1 & 2 & 1 & 1 & 768 & 256 & 0 & \texttt{---++-----}\\
3 & 2 & 3 & 1 & 1 & 512 & 512 & 0 & \texttt{----+-----}\\
4 & 3 & 4 & 1 & 1 & 0 & 1024 & 0 & \texttt{----------}\\
5 & 1 & 2 & 1 & 1 & 0 & 1024 & 0 & \texttt{----------}\\
6 & 0 & 1 & 2 & 1 & 256 & 256 & 512 & \texttt{-------??-}\\
7 & 6 & 2 & 1 & 2 & 0 & 1024 & 0 & \texttt{----------}\\
8 & 6 & 2 & 1 & 2 & 0 & 1024 & 0 & \texttt{----------}\\
9 & 0 & 1 & 1 & 1 & 0 & 1024 & 0 & \texttt{----------}\\
\multicolumn{9}{@{}l@{}}{\rule{0pt}{2.6ex}\textbf{T10.524}\enspace \texttt{0123421212}\enspace $n{=}10$, $|H|{=}4$, len 5, br 2, $|\mathrm{Aut}|{=}2$, orb 8/6, hist-orb 3/2; M 18432, res 7264, rev 5950: ok}\\*
0 & -- & 0 & 4 & 1 & 414 & 2 & 608 & \texttt{-?????????}\\*
1 & 0 & 1 & 2 & 1 & 448 & 64 & 512 & \texttt{--????----}\\
2 & 1 & 2 & 1 & 1 & 768 & 256 & 0 & \texttt{---++-----}\\
3 & 2 & 3 & 1 & 1 & 512 & 512 & 0 & \texttt{----+-----}\\
4 & 3 & 4 & 1 & 1 & 0 & 1024 & 0 & \texttt{----------}\\
5 & 1 & 2 & 1 & 1 & 0 & 1024 & 0 & \texttt{----------}\\
6 & 0 & 1 & 1 & 2 & 512 & 512 & 0 & \texttt{-------+--}\\
7 & 6 & 2 & 1 & 2 & 0 & 1024 & 0 & \texttt{----------}\\
8 & 0 & 1 & 1 & 2 & 512 & 512 & 0 & \texttt{---------+}\\
9 & 8 & 2 & 1 & 2 & 0 & 1024 & 0 & \texttt{----------}\\
\multicolumn{9}{@{}l@{}}{\rule{0pt}{2.6ex}\textbf{T10.525}\enspace \texttt{0123421211}\enspace $n{=}10$, $|H|{=}5$, len 5, br 2, $|\mathrm{Aut}|{=}2$, orb 9/8, hist-orb 4/3; M 17408, res 9748, rev 4038: ok}\\*
0 & -- & 0 & 5 & 1 & 138 & 2 & 884 & \texttt{-?????????}\\*
1 & 0 & 1 & 2 & 1 & 448 & 64 & 512 & \texttt{--????----}\\
2 & 1 & 2 & 1 & 1 & 768 & 256 & 0 & \texttt{---++-----}\\
3 & 2 & 3 & 1 & 1 & 512 & 512 & 0 & \texttt{----+-----}\\
4 & 3 & 4 & 1 & 1 & 0 & 1024 & 0 & \texttt{----------}\\
5 & 1 & 2 & 1 & 1 & 0 & 1024 & 0 & \texttt{----------}\\
6 & 0 & 1 & 1 & 1 & 512 & 512 & 0 & \texttt{-------+--}\\
7 & 6 & 2 & 1 & 1 & 0 & 1024 & 0 & \texttt{----------}\\
8 & 0 & 1 & 1 & 2 & 0 & 1024 & 0 & \texttt{----------}\\
9 & 0 & 1 & 1 & 2 & 0 & 1024 & 0 & \texttt{----------}\\
\multicolumn{9}{@{}l@{}}{\rule{0pt}{2.6ex}\textbf{T10.526}\enspace \texttt{0123421111}\enspace $n{=}10$, $|H|{=}6$, len 5, br 2, $|\mathrm{Aut}|{=}24$, orb 7/6, hist-orb 3/2; M 16384, res 10576, rev 2998: ok}\\*
0 & -- & 0 & 6 & 1 & 46 & 2 & 976 & \texttt{-?????????}\\*
1 & 0 & 1 & 2 & 1 & 448 & 64 & 512 & \texttt{--????----}\\
2 & 1 & 2 & 1 & 1 & 768 & 256 & 0 & \texttt{---++-----}\\
3 & 2 & 3 & 1 & 1 & 512 & 512 & 0 & \texttt{----+-----}\\
4 & 3 & 4 & 1 & 1 & 0 & 1024 & 0 & \texttt{----------}\\
5 & 1 & 2 & 1 & 1 & 0 & 1024 & 0 & \texttt{----------}\\
6 & 0 & 1 & 1 & 4 & 0 & 1024 & 0 & \texttt{----------}\\
7 & 0 & 1 & 1 & 4 & 0 & 1024 & 0 & \texttt{----------}\\
8 & 0 & 1 & 1 & 4 & 0 & 1024 & 0 & \texttt{----------}\\
9 & 0 & 1 & 1 & 4 & 0 & 1024 & 0 & \texttt{----------}\\
\multicolumn{9}{@{}l@{}}{\rule{0pt}{2.6ex}\textbf{T10.527}\enspace \texttt{0123412341}\enspace $n{=}10$, $|H|{=}3$, len 5, br 1, $|\mathrm{Aut}|{=}2$, orb 6/2, hist-orb 2/1; M 21504, res 5148, rev 7898: ok}\\*
0 & -- & 0 & 3 & 1 & 450 & 2 & 572 & \texttt{-?????????}\\*
1 & 0 & 1 & 1 & 2 & 896 & 128 & 0 & \texttt{--+++-----}\\
2 & 1 & 2 & 1 & 2 & 768 & 256 & 0 & \texttt{---++-----}\\
3 & 2 & 3 & 1 & 2 & 512 & 512 & 0 & \texttt{----+-----}\\
4 & 3 & 4 & 1 & 2 & 0 & 1024 & 0 & \texttt{----------}\\
5 & 0 & 1 & 1 & 2 & 896 & 128 & 0 & \texttt{------+++-}\\
6 & 5 & 2 & 1 & 2 & 768 & 256 & 0 & \texttt{-------++-}\\
7 & 6 & 3 & 1 & 2 & 512 & 512 & 0 & \texttt{--------+-}\\
8 & 7 & 4 & 1 & 2 & 0 & 1024 & 0 & \texttt{----------}\\
9 & 0 & 1 & 1 & 1 & 0 & 1024 & 0 & \texttt{----------}\\
\multicolumn{9}{@{}l@{}}{\rule{0pt}{2.6ex}\textbf{T10.528}\enspace \texttt{0123412333}\enspace $n{=}10$, $|H|{=}4$, len 5, br 2, $|\mathrm{Aut}|{=}6$, orb 8/7, hist-orb 2/1; M 22528, res 5520, rev 8790: ok}\\*
0 & -- & 0 & 4 & 1 & 750 & 2 & 272 & \texttt{-?????????}\\*
1 & 0 & 1 & 1 & 1 & 896 & 128 & 0 & \texttt{--+++-----}\\
2 & 1 & 2 & 1 & 1 & 768 & 256 & 0 & \texttt{---++-----}\\
3 & 2 & 3 & 1 & 1 & 512 & 512 & 0 & \texttt{----+-----}\\
4 & 3 & 4 & 1 & 1 & 0 & 1024 & 0 & \texttt{----------}\\
5 & 0 & 1 & 3 & 1 & 576 & 64 & 384 & \texttt{------+???}\\
6 & 5 & 2 & 3 & 1 & 128 & 128 & 768 & \texttt{-------???}\\
7 & 6 & 3 & 1 & 3 & 0 & 1024 & 0 & \texttt{----------}\\
8 & 6 & 3 & 1 & 3 & 0 & 1024 & 0 & \texttt{----------}\\
9 & 6 & 3 & 1 & 3 & 0 & 1024 & 0 & \texttt{----------}\\
\multicolumn{9}{@{}l@{}}{\rule{0pt}{2.6ex}\textbf{T10.529}\enspace \texttt{0123412332}\enspace $n{=}10$, $|H|{=}4$, len 5, br 3, $|\mathrm{Aut}|{=}2$, orb 9/8, hist-orb 3/2; M 21504, res 6408, rev 7642: ok}\\*
0 & -- & 0 & 4 & 1 & 630 & 2 & 392 & \texttt{-?????????}\\*
1 & 0 & 1 & 1 & 1 & 896 & 128 & 0 & \texttt{--+++-----}\\
2 & 1 & 2 & 1 & 1 & 768 & 256 & 0 & \texttt{---++-----}\\
3 & 2 & 3 & 1 & 1 & 512 & 512 & 0 & \texttt{----+-----}\\
4 & 3 & 4 & 1 & 1 & 0 & 1024 & 0 & \texttt{----------}\\
5 & 0 & 1 & 3 & 1 & 320 & 64 & 640 & \texttt{------????}\\
6 & 5 & 2 & 2 & 1 & 256 & 256 & 512 & \texttt{-------??-}\\
7 & 6 & 3 & 1 & 2 & 0 & 1024 & 0 & \texttt{----------}\\
8 & 6 & 3 & 1 & 2 & 0 & 1024 & 0 & \texttt{----------}\\
9 & 5 & 2 & 1 & 1 & 0 & 1024 & 0 & \texttt{----------}\\
\multicolumn{9}{@{}l@{}}{\rule{0pt}{2.6ex}\textbf{T10.530}\enspace \texttt{0123412331}\enspace $n{=}10$, $|H|{=}4$, len 5, br 2, $|\mathrm{Aut}|{=}2$, orb 9/8, hist-orb 3/2; M 20480, res 6968, rev 6830: ok}\\*
0 & -- & 0 & 4 & 1 & 390 & 2 & 632 & \texttt{-?????????}\\*
1 & 0 & 1 & 1 & 1 & 896 & 128 & 0 & \texttt{--+++-----}\\
2 & 1 & 2 & 1 & 1 & 768 & 256 & 0 & \texttt{---++-----}\\
3 & 2 & 3 & 1 & 1 & 512 & 512 & 0 & \texttt{----+-----}\\
4 & 3 & 4 & 1 & 1 & 0 & 1024 & 0 & \texttt{----------}\\
5 & 0 & 1 & 2 & 1 & 640 & 128 & 256 & \texttt{------+??-}\\
6 & 5 & 2 & 2 & 1 & 256 & 256 & 512 & \texttt{-------??-}\\
7 & 6 & 3 & 1 & 2 & 0 & 1024 & 0 & \texttt{----------}\\
8 & 6 & 3 & 1 & 2 & 0 & 1024 & 0 & \texttt{----------}\\
9 & 0 & 1 & 1 & 1 & 0 & 1024 & 0 & \texttt{----------}\\
\multicolumn{9}{@{}l@{}}{\rule{0pt}{2.6ex}\textbf{T10.531}\enspace \texttt{0123412323}\enspace $n{=}10$, $|H|{=}3$, len 5, br 2, $|\mathrm{Aut}|{=}2$, orb 8/6, hist-orb 2/1; M 21504, res 3792, rev 8994: ok}\\*
0 & -- & 0 & 3 & 1 & 750 & 2 & 272 & \texttt{-?????????}\\*
1 & 0 & 1 & 1 & 1 & 896 & 128 & 0 & \texttt{--+++-----}\\
2 & 1 & 2 & 1 & 1 & 768 & 256 & 0 & \texttt{---++-----}\\
3 & 2 & 3 & 1 & 1 & 512 & 512 & 0 & \texttt{----+-----}\\
4 & 3 & 4 & 1 & 1 & 0 & 1024 & 0 & \texttt{----------}\\
5 & 0 & 1 & 2 & 1 & 576 & 64 & 384 & \texttt{------????}\\
6 & 5 & 2 & 1 & 2 & 512 & 512 & 0 & \texttt{-------+--}\\
7 & 6 & 3 & 1 & 2 & 0 & 1024 & 0 & \texttt{----------}\\
8 & 5 & 2 & 1 & 2 & 512 & 512 & 0 & \texttt{---------+}\\
9 & 8 & 3 & 1 & 2 & 0 & 1024 & 0 & \texttt{----------}\\
\multicolumn{9}{@{}l@{}}{\rule{0pt}{2.6ex}\textbf{T10.532}\enspace \texttt{0123412322}\enspace $n{=}10$, $|H|{=}4$, len 5, br 2, $|\mathrm{Aut}|{=}2$, orb 9/8, hist-orb 3/2; M 20480, res 6756, rev 7094: ok}\\*
0 & -- & 0 & 4 & 1 & 570 & 2 & 452 & \texttt{-?????????}\\*
1 & 0 & 1 & 1 & 1 & 896 & 128 & 0 & \texttt{--+++-----}\\
2 & 1 & 2 & 1 & 1 & 768 & 256 & 0 & \texttt{---++-----}\\
3 & 2 & 3 & 1 & 1 & 512 & 512 & 0 & \texttt{----+-----}\\
4 & 3 & 4 & 1 & 1 & 0 & 1024 & 0 & \texttt{----------}\\
5 & 0 & 1 & 3 & 1 & 192 & 64 & 768 & \texttt{------????}\\
6 & 5 & 2 & 1 & 1 & 512 & 512 & 0 & \texttt{-------+--}\\
7 & 6 & 3 & 1 & 1 & 0 & 1024 & 0 & \texttt{----------}\\
8 & 5 & 2 & 1 & 2 & 0 & 1024 & 0 & \texttt{----------}\\
9 & 5 & 2 & 1 & 2 & 0 & 1024 & 0 & \texttt{----------}\\
\multicolumn{9}{@{}l@{}}{\rule{0pt}{2.6ex}\textbf{T10.533}\enspace \texttt{0123412321}\enspace $n{=}10$, $|H|{=}4$, len 5, br 2, $|\mathrm{Aut}|{=}1$, orb 10/10, hist-orb 4/4; M 19456, res 7508, rev 6094: ok}\\*
0 & -- & 0 & 4 & 1 & 330 & 2 & 692 & \texttt{-?????????}\\*
1 & 0 & 1 & 1 & 1 & 896 & 128 & 0 & \texttt{--+++-----}\\
2 & 1 & 2 & 1 & 1 & 768 & 256 & 0 & \texttt{---++-----}\\
3 & 2 & 3 & 1 & 1 & 512 & 512 & 0 & \texttt{----+-----}\\
4 & 3 & 4 & 1 & 1 & 0 & 1024 & 0 & \texttt{----------}\\
5 & 0 & 1 & 2 & 1 & 384 & 128 & 512 & \texttt{------???-}\\
6 & 5 & 2 & 1 & 1 & 512 & 512 & 0 & \texttt{-------+--}\\
7 & 6 & 3 & 1 & 1 & 0 & 1024 & 0 & \texttt{----------}\\
8 & 5 & 2 & 1 & 1 & 0 & 1024 & 0 & \texttt{----------}\\
9 & 0 & 1 & 1 & 1 & 0 & 1024 & 0 & \texttt{----------}\\
\multicolumn{9}{@{}l@{}}{\rule{0pt}{2.6ex}\textbf{T10.534}\enspace \texttt{0123412312}\enspace $n{=}10$, $|H|{=}3$, len 5, br 1, $|\mathrm{Aut}|{=}1$, orb 10/10, hist-orb 3/3; M 19456, res 3528, rev 7914: ok}\\*
0 & -- & 0 & 3 & 1 & 630 & 2 & 392 & \texttt{-?????????}\\*
1 & 0 & 1 & 1 & 1 & 896 & 128 & 0 & \texttt{--+++-----}\\
2 & 1 & 2 & 1 & 1 & 768 & 256 & 0 & \texttt{---++-----}\\
3 & 2 & 3 & 1 & 1 & 512 & 512 & 0 & \texttt{----+-----}\\
4 & 3 & 4 & 1 & 1 & 0 & 1024 & 0 & \texttt{----------}\\
5 & 0 & 1 & 1 & 1 & 768 & 256 & 0 & \texttt{------++--}\\
6 & 5 & 2 & 1 & 1 & 512 & 512 & 0 & \texttt{-------+--}\\
7 & 6 & 3 & 1 & 1 & 0 & 1024 & 0 & \texttt{----------}\\
8 & 0 & 1 & 1 & 1 & 512 & 512 & 0 & \texttt{---------+}\\
9 & 8 & 2 & 1 & 1 & 0 & 1024 & 0 & \texttt{----------}\\
\multicolumn{9}{@{}l@{}}{\rule{0pt}{2.6ex}\textbf{T10.535}\enspace \texttt{0123412311}\enspace $n{=}10$, $|H|{=}4$, len 5, br 1, $|\mathrm{Aut}|{=}2$, orb 9/8, hist-orb 3/2; M 18432, res 7308, rev 5594: ok}\\*
0 & -- & 0 & 4 & 1 & 210 & 2 & 812 & \texttt{-?????????}\\*
1 & 0 & 1 & 1 & 1 & 896 & 128 & 0 & \texttt{--+++-----}\\
2 & 1 & 2 & 1 & 1 & 768 & 256 & 0 & \texttt{---++-----}\\
3 & 2 & 3 & 1 & 1 & 512 & 512 & 0 & \texttt{----+-----}\\
4 & 3 & 4 & 1 & 1 & 0 & 1024 & 0 & \texttt{----------}\\
5 & 0 & 1 & 1 & 1 & 768 & 256 & 0 & \texttt{------++--}\\
6 & 5 & 2 & 1 & 1 & 512 & 512 & 0 & \texttt{-------+--}\\
7 & 6 & 3 & 1 & 1 & 0 & 1024 & 0 & \texttt{----------}\\
8 & 0 & 1 & 1 & 2 & 0 & 1024 & 0 & \texttt{----------}\\
9 & 0 & 1 & 1 & 2 & 0 & 1024 & 0 & \texttt{----------}\\
\multicolumn{9}{@{}l@{}}{\rule{0pt}{2.6ex}\textbf{T10.536}\enspace \texttt{0123412222}\enspace $n{=}10$, $|H|{=}5$, len 5, br 2, $|\mathrm{Aut}|{=}24$, orb 7/6, hist-orb 2/1; M 19456, res 7744, rev 6154: ok}\\*
0 & -- & 0 & 5 & 1 & 510 & 2 & 512 & \texttt{-?????????}\\*
1 & 0 & 1 & 1 & 1 & 896 & 128 & 0 & \texttt{--+++-----}\\
2 & 1 & 2 & 1 & 1 & 768 & 256 & 0 & \texttt{---++-----}\\
3 & 2 & 3 & 1 & 1 & 512 & 512 & 0 & \texttt{----+-----}\\
4 & 3 & 4 & 1 & 1 & 0 & 1024 & 0 & \texttt{----------}\\
5 & 0 & 1 & 4 & 1 & 64 & 64 & 896 & \texttt{------????}\\
6 & 5 & 2 & 1 & 4 & 0 & 1024 & 0 & \texttt{----------}\\
7 & 5 & 2 & 1 & 4 & 0 & 1024 & 0 & \texttt{----------}\\
8 & 5 & 2 & 1 & 4 & 0 & 1024 & 0 & \texttt{----------}\\
9 & 5 & 2 & 1 & 4 & 0 & 1024 & 0 & \texttt{----------}\\
\multicolumn{9}{@{}l@{}}{\rule{0pt}{2.6ex}\textbf{T10.537}\enspace \texttt{0123412221}\enspace $n{=}10$, $|H|{=}5$, len 5, br 2, $|\mathrm{Aut}|{=}6$, orb 8/7, hist-orb 3/2; M 18432, res 8688, rev 4906: ok}\\*
0 & -- & 0 & 5 & 1 & 270 & 2 & 752 & \texttt{-?????????}\\*
1 & 0 & 1 & 1 & 1 & 896 & 128 & 0 & \texttt{--+++-----}\\
2 & 1 & 2 & 1 & 1 & 768 & 256 & 0 & \texttt{---++-----}\\
3 & 2 & 3 & 1 & 1 & 512 & 512 & 0 & \texttt{----+-----}\\
4 & 3 & 4 & 1 & 1 & 0 & 1024 & 0 & \texttt{----------}\\
5 & 0 & 1 & 3 & 1 & 128 & 128 & 768 & \texttt{------???-}\\
6 & 5 & 2 & 1 & 3 & 0 & 1024 & 0 & \texttt{----------}\\
7 & 5 & 2 & 1 & 3 & 0 & 1024 & 0 & \texttt{----------}\\
8 & 5 & 2 & 1 & 3 & 0 & 1024 & 0 & \texttt{----------}\\
9 & 0 & 1 & 1 & 1 & 0 & 1024 & 0 & \texttt{----------}\\
\multicolumn{9}{@{}l@{}}{\rule{0pt}{2.6ex}\textbf{T10.538}\enspace \texttt{0123412212}\enspace $n{=}10$, $|H|{=}4$, len 5, br 2, $|\mathrm{Aut}|{=}2$, orb 9/8, hist-orb 3/2; M 18432, res 5916, rev 6338: ok}\\*
0 & -- & 0 & 4 & 1 & 450 & 2 & 572 & \texttt{-?????????}\\*
1 & 0 & 1 & 1 & 1 & 896 & 128 & 0 & \texttt{--+++-----}\\
2 & 1 & 2 & 1 & 1 & 768 & 256 & 0 & \texttt{---++-----}\\
3 & 2 & 3 & 1 & 1 & 512 & 512 & 0 & \texttt{----+-----}\\
4 & 3 & 4 & 1 & 1 & 0 & 1024 & 0 & \texttt{----------}\\
5 & 0 & 1 & 2 & 1 & 256 & 256 & 512 & \texttt{------??--}\\
6 & 5 & 2 & 1 & 2 & 0 & 1024 & 0 & \texttt{----------}\\
7 & 5 & 2 & 1 & 2 & 0 & 1024 & 0 & \texttt{----------}\\
8 & 0 & 1 & 1 & 1 & 512 & 512 & 0 & \texttt{---------+}\\
9 & 8 & 2 & 1 & 1 & 0 & 1024 & 0 & \texttt{----------}\\
\multicolumn{9}{@{}l@{}}{\rule{0pt}{2.6ex}\textbf{T10.539}\enspace \texttt{0123412211}\enspace $n{=}10$, $|H|{=}5$, len 5, br 2, $|\mathrm{Aut}|{=}4$, orb 8/6, hist-orb 3/1; M 17408, res 8616, rev 4346: ok}\\*
0 & -- & 0 & 5 & 1 & 150 & 2 & 872 & \texttt{-?????????}\\*
1 & 0 & 1 & 1 & 1 & 896 & 128 & 0 & \texttt{--+++-----}\\
2 & 1 & 2 & 1 & 1 & 768 & 256 & 0 & \texttt{---++-----}\\
3 & 2 & 3 & 1 & 1 & 512 & 512 & 0 & \texttt{----+-----}\\
4 & 3 & 4 & 1 & 1 & 0 & 1024 & 0 & \texttt{----------}\\
5 & 0 & 1 & 2 & 1 & 256 & 256 & 512 & \texttt{------??--}\\
6 & 5 & 2 & 1 & 2 & 0 & 1024 & 0 & \texttt{----------}\\
7 & 5 & 2 & 1 & 2 & 0 & 1024 & 0 & \texttt{----------}\\
8 & 0 & 1 & 1 & 2 & 0 & 1024 & 0 & \texttt{----------}\\
9 & 0 & 1 & 1 & 2 & 0 & 1024 & 0 & \texttt{----------}\\
\multicolumn{9}{@{}l@{}}{\rule{0pt}{2.6ex}\textbf{T10.540}\enspace \texttt{0123412121}\enspace $n{=}10$, $|H|{=}4$, len 5, br 1, $|\mathrm{Aut}|{=}2$, orb 8/6, hist-orb 3/2; M 17408, res 6768, rev 5426: ok}\\*
0 & -- & 0 & 4 & 1 & 270 & 2 & 752 & \texttt{-?????????}\\*
1 & 0 & 1 & 1 & 1 & 896 & 128 & 0 & \texttt{--+++-----}\\
2 & 1 & 2 & 1 & 1 & 768 & 256 & 0 & \texttt{---++-----}\\
3 & 2 & 3 & 1 & 1 & 512 & 512 & 0 & \texttt{----+-----}\\
4 & 3 & 4 & 1 & 1 & 0 & 1024 & 0 & \texttt{----------}\\
5 & 0 & 1 & 1 & 2 & 512 & 512 & 0 & \texttt{------+---}\\
6 & 5 & 2 & 1 & 2 & 0 & 1024 & 0 & \texttt{----------}\\
7 & 0 & 1 & 1 & 2 & 512 & 512 & 0 & \texttt{--------+-}\\
8 & 7 & 2 & 1 & 2 & 0 & 1024 & 0 & \texttt{----------}\\
9 & 0 & 1 & 1 & 1 & 0 & 1024 & 0 & \texttt{----------}\\
\multicolumn{9}{@{}l@{}}{\rule{0pt}{2.6ex}\textbf{T10.541}\enspace \texttt{0123412111}\enspace $n{=}10$, $|H|{=}5$, len 5, br 1, $|\mathrm{Aut}|{=}6$, orb 8/7, hist-orb 3/2; M 16384, res 8388, rev 4002: ok}\\*
0 & -- & 0 & 5 & 1 & 90 & 2 & 932 & \texttt{-?????????}\\*
1 & 0 & 1 & 1 & 1 & 896 & 128 & 0 & \texttt{--+++-----}\\
2 & 1 & 2 & 1 & 1 & 768 & 256 & 0 & \texttt{---++-----}\\
3 & 2 & 3 & 1 & 1 & 512 & 512 & 0 & \texttt{----+-----}\\
4 & 3 & 4 & 1 & 1 & 0 & 1024 & 0 & \texttt{----------}\\
5 & 0 & 1 & 1 & 1 & 512 & 512 & 0 & \texttt{------+---}\\
6 & 5 & 2 & 1 & 1 & 0 & 1024 & 0 & \texttt{----------}\\
7 & 0 & 1 & 1 & 3 & 0 & 1024 & 0 & \texttt{----------}\\
8 & 0 & 1 & 1 & 3 & 0 & 1024 & 0 & \texttt{----------}\\
9 & 0 & 1 & 1 & 3 & 0 & 1024 & 0 & \texttt{----------}\\
\multicolumn{9}{@{}l@{}}{\rule{0pt}{2.6ex}\textbf{T10.542}\enspace \texttt{0123411111}\enspace $n{=}10$, $|H|{=}6$, len 5, br 1, $|\mathrm{Aut}|{=}120$, orb 6/5, hist-orb 2/1; M 15360, res 8928, rev 3146: ok}\\*
0 & -- & 0 & 6 & 1 & 30 & 2 & 992 & \texttt{-?????????}\\*
1 & 0 & 1 & 1 & 1 & 896 & 128 & 0 & \texttt{--+++-----}\\
2 & 1 & 2 & 1 & 1 & 768 & 256 & 0 & \texttt{---++-----}\\
3 & 2 & 3 & 1 & 1 & 512 & 512 & 0 & \texttt{----+-----}\\
4 & 3 & 4 & 1 & 1 & 0 & 1024 & 0 & \texttt{----------}\\
5 & 0 & 1 & 1 & 5 & 0 & 1024 & 0 & \texttt{----------}\\
6 & 0 & 1 & 1 & 5 & 0 & 1024 & 0 & \texttt{----------}\\
7 & 0 & 1 & 1 & 5 & 0 & 1024 & 0 & \texttt{----------}\\
8 & 0 & 1 & 1 & 5 & 0 & 1024 & 0 & \texttt{----------}\\
9 & 0 & 1 & 1 & 5 & 0 & 1024 & 0 & \texttt{----------}\\
\multicolumn{9}{@{}l@{}}{\rule{0pt}{2.6ex}\textbf{T10.543}\enspace \texttt{0123333333}\enspace $n{=}10$, $|H|{=}7$, len 4, br 1, $|\mathrm{Aut}|{=}5040$, orb 4/3, hist-orb 1/0; M 24576, res 12348, rev 9070: ok}\\*
0 & -- & 0 & 7 & 1 & 770 & 2 & 252 & \texttt{-++???????}\\*
1 & 0 & 1 & 7 & 1 & 516 & 4 & 504 & \texttt{--+???????}\\
2 & 1 & 2 & 7 & 1 & 8 & 8 & 1008 & \texttt{---???????}\\
3 & 2 & 3 & 1 & 7 & 0 & 1024 & 0 & \texttt{----------}\\
4 & 2 & 3 & 1 & 7 & 0 & 1024 & 0 & \texttt{----------}\\
5 & 2 & 3 & 1 & 7 & 0 & 1024 & 0 & \texttt{----------}\\
6 & 2 & 3 & 1 & 7 & 0 & 1024 & 0 & \texttt{----------}\\
7 & 2 & 3 & 1 & 7 & 0 & 1024 & 0 & \texttt{----------}\\
8 & 2 & 3 & 1 & 7 & 0 & 1024 & 0 & \texttt{----------}\\
9 & 2 & 3 & 1 & 7 & 0 & 1024 & 0 & \texttt{----------}\\
\multicolumn{9}{@{}l@{}}{\rule{0pt}{2.6ex}\textbf{T10.544}\enspace \texttt{0123333332}\enspace $n{=}10$, $|H|{=}7$, len 4, br 2, $|\mathrm{Aut}|{=}720$, orb 5/4, hist-orb 2/1; M 23552, res 14328, rev 6426: ok}\\*
0 & -- & 0 & 7 & 1 & 642 & 2 & 380 & \texttt{-+????????}\\*
1 & 0 & 1 & 7 & 1 & 260 & 4 & 760 & \texttt{--????????}\\
2 & 1 & 2 & 6 & 1 & 16 & 16 & 992 & \texttt{---??????-}\\
3 & 2 & 3 & 1 & 6 & 0 & 1024 & 0 & \texttt{----------}\\
4 & 2 & 3 & 1 & 6 & 0 & 1024 & 0 & \texttt{----------}\\
5 & 2 & 3 & 1 & 6 & 0 & 1024 & 0 & \texttt{----------}\\
6 & 2 & 3 & 1 & 6 & 0 & 1024 & 0 & \texttt{----------}\\
7 & 2 & 3 & 1 & 6 & 0 & 1024 & 0 & \texttt{----------}\\
8 & 2 & 3 & 1 & 6 & 0 & 1024 & 0 & \texttt{----------}\\
9 & 1 & 2 & 1 & 1 & 0 & 1024 & 0 & \texttt{----------}\\
\multicolumn{9}{@{}l@{}}{\rule{0pt}{2.6ex}\textbf{T10.545}\enspace \texttt{0123333331}\enspace $n{=}10$, $|H|{=}7$, len 4, br 2, $|\mathrm{Aut}|{=}720$, orb 5/4, hist-orb 2/1; M 22528, res 14156, rev 5934: ok}\\*
0 & -- & 0 & 7 & 1 & 386 & 2 & 636 & \texttt{-?????????}\\*
1 & 0 & 1 & 6 & 1 & 520 & 8 & 496 & \texttt{--+??????-}\\
2 & 1 & 2 & 6 & 1 & 16 & 16 & 992 & \texttt{---??????-}\\
3 & 2 & 3 & 1 & 6 & 0 & 1024 & 0 & \texttt{----------}\\
4 & 2 & 3 & 1 & 6 & 0 & 1024 & 0 & \texttt{----------}\\
5 & 2 & 3 & 1 & 6 & 0 & 1024 & 0 & \texttt{----------}\\
6 & 2 & 3 & 1 & 6 & 0 & 1024 & 0 & \texttt{----------}\\
7 & 2 & 3 & 1 & 6 & 0 & 1024 & 0 & \texttt{----------}\\
8 & 2 & 3 & 1 & 6 & 0 & 1024 & 0 & \texttt{----------}\\
9 & 0 & 1 & 1 & 1 & 0 & 1024 & 0 & \texttt{----------}\\
\multicolumn{9}{@{}l@{}}{\rule{0pt}{2.6ex}\textbf{T10.546}\enspace \texttt{0123333323}\enspace $n{=}10$, $|H|{=}6$, len 4, br 2, $|\mathrm{Aut}|{=}120$, orb 6/5, hist-orb 2/1; M 23552, res 11568, rev 7798: ok}\\*
0 & -- & 0 & 6 & 1 & 710 & 2 & 312 & \texttt{-+????????}\\*
1 & 0 & 1 & 6 & 1 & 396 & 4 & 624 & \texttt{--????????}\\
2 & 1 & 2 & 5 & 1 & 32 & 32 & 960 & \texttt{---?????--}\\
3 & 2 & 3 & 1 & 5 & 0 & 1024 & 0 & \texttt{----------}\\
4 & 2 & 3 & 1 & 5 & 0 & 1024 & 0 & \texttt{----------}\\
5 & 2 & 3 & 1 & 5 & 0 & 1024 & 0 & \texttt{----------}\\
6 & 2 & 3 & 1 & 5 & 0 & 1024 & 0 & \texttt{----------}\\
7 & 2 & 3 & 1 & 5 & 0 & 1024 & 0 & \texttt{----------}\\
8 & 1 & 2 & 1 & 1 & 512 & 512 & 0 & \texttt{---------+}\\
9 & 8 & 3 & 1 & 1 & 0 & 1024 & 0 & \texttt{----------}\\
\multicolumn{9}{@{}l@{}}{\rule{0pt}{2.6ex}\textbf{T10.547}\enspace \texttt{0123333322}\enspace $n{=}10$, $|H|{=}7$, len 4, br 2, $|\mathrm{Aut}|{=}240$, orb 5/3, hist-orb 2/0; M 22528, res 14736, rev 5154: ok}\\*
0 & -- & 0 & 7 & 1 & 578 & 2 & 444 & \texttt{-+????????}\\*
1 & 0 & 1 & 7 & 1 & 132 & 4 & 888 & \texttt{--????????}\\
2 & 1 & 2 & 5 & 1 & 32 & 32 & 960 & \texttt{---?????--}\\
3 & 2 & 3 & 1 & 5 & 0 & 1024 & 0 & \texttt{----------}\\
4 & 2 & 3 & 1 & 5 & 0 & 1024 & 0 & \texttt{----------}\\
5 & 2 & 3 & 1 & 5 & 0 & 1024 & 0 & \texttt{----------}\\
6 & 2 & 3 & 1 & 5 & 0 & 1024 & 0 & \texttt{----------}\\
7 & 2 & 3 & 1 & 5 & 0 & 1024 & 0 & \texttt{----------}\\
8 & 1 & 2 & 1 & 2 & 0 & 1024 & 0 & \texttt{----------}\\
9 & 1 & 2 & 1 & 2 & 0 & 1024 & 0 & \texttt{----------}\\
\multicolumn{9}{@{}l@{}}{\rule{0pt}{2.6ex}\textbf{T10.548}\enspace \texttt{0123333321}\enspace $n{=}10$, $|H|{=}7$, len 4, br 3, $|\mathrm{Aut}|{=}120$, orb 6/5, hist-orb 3/2; M 21504, res 15268, rev 4018: ok}\\*
0 & -- & 0 & 7 & 1 & 322 & 2 & 700 & \texttt{-?????????}\\*
1 & 0 & 1 & 6 & 1 & 264 & 8 & 752 & \texttt{--???????-}\\
2 & 1 & 2 & 5 & 1 & 32 & 32 & 960 & \texttt{---?????--}\\
3 & 2 & 3 & 1 & 5 & 0 & 1024 & 0 & \texttt{----------}\\
4 & 2 & 3 & 1 & 5 & 0 & 1024 & 0 & \texttt{----------}\\
5 & 2 & 3 & 1 & 5 & 0 & 1024 & 0 & \texttt{----------}\\
6 & 2 & 3 & 1 & 5 & 0 & 1024 & 0 & \texttt{----------}\\
7 & 2 & 3 & 1 & 5 & 0 & 1024 & 0 & \texttt{----------}\\
8 & 1 & 2 & 1 & 1 & 0 & 1024 & 0 & \texttt{----------}\\
9 & 0 & 1 & 1 & 1 & 0 & 1024 & 0 & \texttt{----------}\\
\multicolumn{9}{@{}l@{}}{\rule{0pt}{2.6ex}\textbf{T10.549}\enspace \texttt{0123333312}\enspace $n{=}10$, $|H|{=}6$, len 4, br 2, $|\mathrm{Aut}|{=}120$, orb 6/5, hist-orb 2/1; M 21504, res 10680, rev 7038: ok}\\*
0 & -- & 0 & 6 & 1 & 582 & 2 & 440 & \texttt{-?????????}\\*
1 & 0 & 1 & 5 & 1 & 528 & 16 & 480 & \texttt{--+?????--}\\
2 & 1 & 2 & 5 & 1 & 32 & 32 & 960 & \texttt{---?????--}\\
3 & 2 & 3 & 1 & 5 & 0 & 1024 & 0 & \texttt{----------}\\
4 & 2 & 3 & 1 & 5 & 0 & 1024 & 0 & \texttt{----------}\\
5 & 2 & 3 & 1 & 5 & 0 & 1024 & 0 & \texttt{----------}\\
6 & 2 & 3 & 1 & 5 & 0 & 1024 & 0 & \texttt{----------}\\
7 & 2 & 3 & 1 & 5 & 0 & 1024 & 0 & \texttt{----------}\\
8 & 0 & 1 & 1 & 1 & 512 & 512 & 0 & \texttt{---------+}\\
9 & 8 & 2 & 1 & 1 & 0 & 1024 & 0 & \texttt{----------}\\
\multicolumn{9}{@{}l@{}}{\rule{0pt}{2.6ex}\textbf{T10.550}\enspace \texttt{0123333311}\enspace $n{=}10$, $|H|{=}7$, len 4, br 2, $|\mathrm{Aut}|{=}240$, orb 5/3, hist-orb 2/0; M 20480, res 14172, rev 4182: ok}\\*
0 & -- & 0 & 7 & 1 & 194 & 2 & 828 & \texttt{-?????????}\\*
1 & 0 & 1 & 5 & 1 & 528 & 16 & 480 & \texttt{--+?????--}\\
2 & 1 & 2 & 5 & 1 & 32 & 32 & 960 & \texttt{---?????--}\\
3 & 2 & 3 & 1 & 5 & 0 & 1024 & 0 & \texttt{----------}\\
4 & 2 & 3 & 1 & 5 & 0 & 1024 & 0 & \texttt{----------}\\
5 & 2 & 3 & 1 & 5 & 0 & 1024 & 0 & \texttt{----------}\\
6 & 2 & 3 & 1 & 5 & 0 & 1024 & 0 & \texttt{----------}\\
7 & 2 & 3 & 1 & 5 & 0 & 1024 & 0 & \texttt{----------}\\
8 & 0 & 1 & 1 & 2 & 0 & 1024 & 0 & \texttt{----------}\\
9 & 0 & 1 & 1 & 2 & 0 & 1024 & 0 & \texttt{----------}\\
\multicolumn{9}{@{}l@{}}{\rule{0pt}{2.6ex}\textbf{T10.551}\enspace \texttt{0123333233}\enspace $n{=}10$, $|H|{=}6$, len 4, br 3, $|\mathrm{Aut}|{=}48$, orb 6/4, hist-orb 2/0; M 23552, res 11712, rev 7194: ok}\\*
0 & -- & 0 & 6 & 1 & 682 & 2 & 340 & \texttt{-+????????}\\*
1 & 0 & 1 & 6 & 1 & 340 & 4 & 680 & \texttt{--????????}\\
2 & 1 & 2 & 4 & 1 & 64 & 64 & 896 & \texttt{---????---}\\
3 & 2 & 3 & 1 & 4 & 0 & 1024 & 0 & \texttt{----------}\\
4 & 2 & 3 & 1 & 4 & 0 & 1024 & 0 & \texttt{----------}\\
5 & 2 & 3 & 1 & 4 & 0 & 1024 & 0 & \texttt{----------}\\
6 & 2 & 3 & 1 & 4 & 0 & 1024 & 0 & \texttt{----------}\\
7 & 1 & 2 & 2 & 1 & 256 & 256 & 512 & \texttt{--------??}\\
8 & 7 & 3 & 1 & 2 & 0 & 1024 & 0 & \texttt{----------}\\
9 & 7 & 3 & 1 & 2 & 0 & 1024 & 0 & \texttt{----------}\\
\multicolumn{9}{@{}l@{}}{\rule{0pt}{2.6ex}\textbf{T10.552}\enspace \texttt{0123333232}\enspace $n{=}10$, $|H|{=}6$, len 4, br 2, $|\mathrm{Aut}|{=}24$, orb 7/6, hist-orb 3/2; M 22528, res 12704, rev 6086: ok}\\*
0 & -- & 0 & 6 & 1 & 614 & 2 & 408 & \texttt{-+????????}\\*
1 & 0 & 1 & 6 & 1 & 204 & 4 & 816 & \texttt{--????????}\\
2 & 1 & 2 & 4 & 1 & 64 & 64 & 896 & \texttt{---????---}\\
3 & 2 & 3 & 1 & 4 & 0 & 1024 & 0 & \texttt{----------}\\
4 & 2 & 3 & 1 & 4 & 0 & 1024 & 0 & \texttt{----------}\\
5 & 2 & 3 & 1 & 4 & 0 & 1024 & 0 & \texttt{----------}\\
6 & 2 & 3 & 1 & 4 & 0 & 1024 & 0 & \texttt{----------}\\
7 & 1 & 2 & 1 & 1 & 512 & 512 & 0 & \texttt{--------+-}\\
8 & 7 & 3 & 1 & 1 & 0 & 1024 & 0 & \texttt{----------}\\
9 & 1 & 2 & 1 & 1 & 0 & 1024 & 0 & \texttt{----------}\\
\multicolumn{9}{@{}l@{}}{\rule{0pt}{2.6ex}\textbf{T10.553}\enspace \texttt{0123333231}\enspace $n{=}10$, $|H|{=}6$, len 4, br 3, $|\mathrm{Aut}|{=}24$, orb 7/6, hist-orb 3/2; M 21504, res 12840, rev 5302: ok}\\*
0 & -- & 0 & 6 & 1 & 358 & 2 & 664 & \texttt{-?????????}\\*
1 & 0 & 1 & 5 & 1 & 408 & 8 & 608 & \texttt{--???????-}\\
2 & 1 & 2 & 4 & 1 & 64 & 64 & 896 & \texttt{---????---}\\
3 & 2 & 3 & 1 & 4 & 0 & 1024 & 0 & \texttt{----------}\\
4 & 2 & 3 & 1 & 4 & 0 & 1024 & 0 & \texttt{----------}\\
5 & 2 & 3 & 1 & 4 & 0 & 1024 & 0 & \texttt{----------}\\
6 & 2 & 3 & 1 & 4 & 0 & 1024 & 0 & \texttt{----------}\\
7 & 1 & 2 & 1 & 1 & 512 & 512 & 0 & \texttt{--------+-}\\
8 & 7 & 3 & 1 & 1 & 0 & 1024 & 0 & \texttt{----------}\\
9 & 0 & 1 & 1 & 1 & 0 & 1024 & 0 & \texttt{----------}\\
\multicolumn{9}{@{}l@{}}{\rule{0pt}{2.6ex}\textbf{T10.554}\enspace \texttt{0123333222}\enspace $n{=}10$, $|H|{=}7$, len 4, br 2, $|\mathrm{Aut}|{=}144$, orb 5/3, hist-orb 2/0; M 21504, res 14336, rev 4594: ok}\\*
0 & -- & 0 & 7 & 1 & 546 & 2 & 476 & \texttt{-+????????}\\*
1 & 0 & 1 & 7 & 1 & 68 & 4 & 952 & \texttt{--????????}\\
2 & 1 & 2 & 4 & 1 & 64 & 64 & 896 & \texttt{---????---}\\
3 & 2 & 3 & 1 & 4 & 0 & 1024 & 0 & \texttt{----------}\\
4 & 2 & 3 & 1 & 4 & 0 & 1024 & 0 & \texttt{----------}\\
5 & 2 & 3 & 1 & 4 & 0 & 1024 & 0 & \texttt{----------}\\
6 & 2 & 3 & 1 & 4 & 0 & 1024 & 0 & \texttt{----------}\\
7 & 1 & 2 & 1 & 3 & 0 & 1024 & 0 & \texttt{----------}\\
8 & 1 & 2 & 1 & 3 & 0 & 1024 & 0 & \texttt{----------}\\
9 & 1 & 2 & 1 & 3 & 0 & 1024 & 0 & \texttt{----------}\\
\multicolumn{9}{@{}l@{}}{\rule{0pt}{2.6ex}\textbf{T10.555}\enspace \texttt{0123333221}\enspace $n{=}10$, $|H|{=}7$, len 4, br 3, $|\mathrm{Aut}|{=}48$, orb 6/4, hist-orb 3/1; M 20480, res 15220, rev 3130: ok}\\*
0 & -- & 0 & 7 & 1 & 290 & 2 & 732 & \texttt{-?????????}\\*
1 & 0 & 1 & 6 & 1 & 136 & 8 & 880 & \texttt{--???????-}\\
2 & 1 & 2 & 4 & 1 & 64 & 64 & 896 & \texttt{---????---}\\
3 & 2 & 3 & 1 & 4 & 0 & 1024 & 0 & \texttt{----------}\\
4 & 2 & 3 & 1 & 4 & 0 & 1024 & 0 & \texttt{----------}\\
5 & 2 & 3 & 1 & 4 & 0 & 1024 & 0 & \texttt{----------}\\
6 & 2 & 3 & 1 & 4 & 0 & 1024 & 0 & \texttt{----------}\\
7 & 1 & 2 & 1 & 2 & 0 & 1024 & 0 & \texttt{----------}\\
8 & 1 & 2 & 1 & 2 & 0 & 1024 & 0 & \texttt{----------}\\
9 & 0 & 1 & 1 & 1 & 0 & 1024 & 0 & \texttt{----------}\\
\multicolumn{9}{@{}l@{}}{\rule{0pt}{2.6ex}\textbf{T10.556}\enspace \texttt{0123333212}\enspace $n{=}10$, $|H|{=}6$, len 4, br 3, $|\mathrm{Aut}|{=}24$, orb 7/6, hist-orb 3/2; M 20480, res 11784, rev 5258: ok}\\*
0 & -- & 0 & 6 & 1 & 486 & 2 & 536 & \texttt{-?????????}\\*
1 & 0 & 1 & 5 & 1 & 272 & 16 & 736 & \texttt{--??????--}\\
2 & 1 & 2 & 4 & 1 & 64 & 64 & 896 & \texttt{---????---}\\
3 & 2 & 3 & 1 & 4 & 0 & 1024 & 0 & \texttt{----------}\\
4 & 2 & 3 & 1 & 4 & 0 & 1024 & 0 & \texttt{----------}\\
5 & 2 & 3 & 1 & 4 & 0 & 1024 & 0 & \texttt{----------}\\
6 & 2 & 3 & 1 & 4 & 0 & 1024 & 0 & \texttt{----------}\\
7 & 1 & 2 & 1 & 1 & 0 & 1024 & 0 & \texttt{----------}\\
8 & 0 & 1 & 1 & 1 & 512 & 512 & 0 & \texttt{---------+}\\
9 & 8 & 2 & 1 & 1 & 0 & 1024 & 0 & \texttt{----------}\\
\multicolumn{9}{@{}l@{}}{\rule{0pt}{2.6ex}\textbf{T10.557}\enspace \texttt{0123333211}\enspace $n{=}10$, $|H|{=}7$, len 4, br 3, $|\mathrm{Aut}|{=}48$, orb 6/4, hist-orb 3/1; M 19456, res 14700, rev 2778: ok}\\*
0 & -- & 0 & 7 & 1 & 162 & 2 & 860 & \texttt{-?????????}\\*
1 & 0 & 1 & 5 & 1 & 272 & 16 & 736 & \texttt{--??????--}\\
2 & 1 & 2 & 4 & 1 & 64 & 64 & 896 & \texttt{---????---}\\
3 & 2 & 3 & 1 & 4 & 0 & 1024 & 0 & \texttt{----------}\\
4 & 2 & 3 & 1 & 4 & 0 & 1024 & 0 & \texttt{----------}\\
5 & 2 & 3 & 1 & 4 & 0 & 1024 & 0 & \texttt{----------}\\
6 & 2 & 3 & 1 & 4 & 0 & 1024 & 0 & \texttt{----------}\\
7 & 1 & 2 & 1 & 1 & 0 & 1024 & 0 & \texttt{----------}\\
8 & 0 & 1 & 1 & 2 & 0 & 1024 & 0 & \texttt{----------}\\
9 & 0 & 1 & 1 & 2 & 0 & 1024 & 0 & \texttt{----------}\\
\multicolumn{9}{@{}l@{}}{\rule{0pt}{2.6ex}\textbf{T10.558}\enspace \texttt{0123333123}\enspace $n{=}10$, $|H|{=}5$, len 4, br 2, $|\mathrm{Aut}|{=}24$, orb 7/6, hist-orb 2/1; M 21504, res 7952, rev 7878: ok}\\*
0 & -- & 0 & 5 & 1 & 686 & 2 & 336 & \texttt{-?????????}\\*
1 & 0 & 1 & 4 & 1 & 544 & 32 & 448 & \texttt{--+????---}\\
2 & 1 & 2 & 4 & 1 & 64 & 64 & 896 & \texttt{---????---}\\
3 & 2 & 3 & 1 & 4 & 0 & 1024 & 0 & \texttt{----------}\\
4 & 2 & 3 & 1 & 4 & 0 & 1024 & 0 & \texttt{----------}\\
5 & 2 & 3 & 1 & 4 & 0 & 1024 & 0 & \texttt{----------}\\
6 & 2 & 3 & 1 & 4 & 0 & 1024 & 0 & \texttt{----------}\\
7 & 0 & 1 & 1 & 1 & 768 & 256 & 0 & \texttt{--------++}\\
8 & 7 & 2 & 1 & 1 & 512 & 512 & 0 & \texttt{---------+}\\
9 & 8 & 3 & 1 & 1 & 0 & 1024 & 0 & \texttt{----------}\\
\multicolumn{9}{@{}l@{}}{\rule{0pt}{2.6ex}\textbf{T10.559}\enspace \texttt{0123333122}\enspace $n{=}10$, $|H|{=}6$, len 4, br 3, $|\mathrm{Aut}|{=}48$, orb 6/4, hist-orb 2/0; M 20480, res 10484, rev 6054: ok}\\*
0 & -- & 0 & 6 & 1 & 490 & 2 & 532 & \texttt{-?????????}\\*
1 & 0 & 1 & 4 & 1 & 544 & 32 & 448 & \texttt{--+????---}\\
2 & 1 & 2 & 4 & 1 & 64 & 64 & 896 & \texttt{---????---}\\
3 & 2 & 3 & 1 & 4 & 0 & 1024 & 0 & \texttt{----------}\\
4 & 2 & 3 & 1 & 4 & 0 & 1024 & 0 & \texttt{----------}\\
5 & 2 & 3 & 1 & 4 & 0 & 1024 & 0 & \texttt{----------}\\
6 & 2 & 3 & 1 & 4 & 0 & 1024 & 0 & \texttt{----------}\\
7 & 0 & 1 & 2 & 1 & 256 & 256 & 512 & \texttt{--------??}\\
8 & 7 & 2 & 1 & 2 & 0 & 1024 & 0 & \texttt{----------}\\
9 & 7 & 2 & 1 & 2 & 0 & 1024 & 0 & \texttt{----------}\\
\multicolumn{9}{@{}l@{}}{\rule{0pt}{2.6ex}\textbf{T10.560}\enspace \texttt{0123333121}\enspace $n{=}10$, $|H|{=}6$, len 4, br 2, $|\mathrm{Aut}|{=}24$, orb 7/6, hist-orb 3/2; M 19456, res 11480, rev 4862: ok}\\*
0 & -- & 0 & 6 & 1 & 294 & 2 & 728 & \texttt{-?????????}\\*
1 & 0 & 1 & 4 & 1 & 544 & 32 & 448 & \texttt{--+????---}\\
2 & 1 & 2 & 4 & 1 & 64 & 64 & 896 & \texttt{---????---}\\
3 & 2 & 3 & 1 & 4 & 0 & 1024 & 0 & \texttt{----------}\\
4 & 2 & 3 & 1 & 4 & 0 & 1024 & 0 & \texttt{----------}\\
5 & 2 & 3 & 1 & 4 & 0 & 1024 & 0 & \texttt{----------}\\
6 & 2 & 3 & 1 & 4 & 0 & 1024 & 0 & \texttt{----------}\\
7 & 0 & 1 & 1 & 1 & 512 & 512 & 0 & \texttt{--------+-}\\
8 & 7 & 2 & 1 & 1 & 0 & 1024 & 0 & \texttt{----------}\\
9 & 0 & 1 & 1 & 1 & 0 & 1024 & 0 & \texttt{----------}\\
\multicolumn{9}{@{}l@{}}{\rule{0pt}{2.6ex}\textbf{T10.561}\enspace \texttt{0123333111}\enspace $n{=}10$, $|H|{=}7$, len 4, br 2, $|\mathrm{Aut}|{=}144$, orb 5/3, hist-orb 2/0; M 18432, res 13244, rev 3158: ok}\\*
0 & -- & 0 & 7 & 1 & 98 & 2 & 924 & \texttt{-?????????}\\*
1 & 0 & 1 & 4 & 1 & 544 & 32 & 448 & \texttt{--+????---}\\
2 & 1 & 2 & 4 & 1 & 64 & 64 & 896 & \texttt{---????---}\\
3 & 2 & 3 & 1 & 4 & 0 & 1024 & 0 & \texttt{----------}\\
4 & 2 & 3 & 1 & 4 & 0 & 1024 & 0 & \texttt{----------}\\
5 & 2 & 3 & 1 & 4 & 0 & 1024 & 0 & \texttt{----------}\\
6 & 2 & 3 & 1 & 4 & 0 & 1024 & 0 & \texttt{----------}\\
7 & 0 & 1 & 1 & 3 & 0 & 1024 & 0 & \texttt{----------}\\
8 & 0 & 1 & 1 & 3 & 0 & 1024 & 0 & \texttt{----------}\\
9 & 0 & 1 & 1 & 3 & 0 & 1024 & 0 & \texttt{----------}\\
\multicolumn{9}{@{}l@{}}{\rule{0pt}{2.6ex}\textbf{T10.562}\enspace \texttt{0123332333}\enspace $n{=}10$, $|H|{=}6$, len 4, br 3, $|\mathrm{Aut}|{=}72$, orb 4/2, hist-orb 1/0; M 23552, res 11808, rev 6994: ok}\\*
0 & -- & 0 & 6 & 1 & 674 & 2 & 348 & \texttt{-+????????}\\*
1 & 0 & 1 & 6 & 1 & 324 & 4 & 696 & \texttt{--????????}\\
2 & 1 & 2 & 3 & 2 & 128 & 128 & 768 & \texttt{---???----}\\
3 & 2 & 3 & 1 & 6 & 0 & 1024 & 0 & \texttt{----------}\\
4 & 2 & 3 & 1 & 6 & 0 & 1024 & 0 & \texttt{----------}\\
5 & 2 & 3 & 1 & 6 & 0 & 1024 & 0 & \texttt{----------}\\
6 & 1 & 2 & 3 & 2 & 128 & 128 & 768 & \texttt{-------???}\\
7 & 6 & 3 & 1 & 6 & 0 & 1024 & 0 & \texttt{----------}\\
8 & 6 & 3 & 1 & 6 & 0 & 1024 & 0 & \texttt{----------}\\
9 & 6 & 3 & 1 & 6 & 0 & 1024 & 0 & \texttt{----------}\\
\multicolumn{9}{@{}l@{}}{\rule{0pt}{2.6ex}\textbf{T10.563}\enspace \texttt{0123332332}\enspace $n{=}10$, $|H|{=}6$, len 4, br 3, $|\mathrm{Aut}|{=}12$, orb 7/5, hist-orb 3/1; M 22528, res 12448, rev 5866: ok}\\*
0 & -- & 0 & 6 & 1 & 602 & 2 & 420 & \texttt{-+????????}\\*
1 & 0 & 1 & 6 & 1 & 180 & 4 & 840 & \texttt{--????????}\\
2 & 1 & 2 & 3 & 1 & 128 & 128 & 768 & \texttt{---???----}\\
3 & 2 & 3 & 1 & 3 & 0 & 1024 & 0 & \texttt{----------}\\
4 & 2 & 3 & 1 & 3 & 0 & 1024 & 0 & \texttt{----------}\\
5 & 2 & 3 & 1 & 3 & 0 & 1024 & 0 & \texttt{----------}\\
6 & 1 & 2 & 2 & 1 & 256 & 256 & 512 & \texttt{-------??-}\\
7 & 6 & 3 & 1 & 2 & 0 & 1024 & 0 & \texttt{----------}\\
8 & 6 & 3 & 1 & 2 & 0 & 1024 & 0 & \texttt{----------}\\
9 & 1 & 2 & 1 & 1 & 0 & 1024 & 0 & \texttt{----------}\\
\multicolumn{9}{@{}l@{}}{\rule{0pt}{2.6ex}\textbf{T10.564}\enspace \texttt{0123332331}\enspace $n{=}10$, $|H|{=}6$, len 4, br 4, $|\mathrm{Aut}|{=}12$, orb 7/5, hist-orb 3/1; M 21504, res 12716, rev 5010: ok}\\*
0 & -- & 0 & 6 & 1 & 346 & 2 & 676 & \texttt{-?????????}\\*
1 & 0 & 1 & 5 & 1 & 360 & 8 & 656 & \texttt{--???????-}\\
2 & 1 & 2 & 3 & 1 & 128 & 128 & 768 & \texttt{---???----}\\
3 & 2 & 3 & 1 & 3 & 0 & 1024 & 0 & \texttt{----------}\\
4 & 2 & 3 & 1 & 3 & 0 & 1024 & 0 & \texttt{----------}\\
5 & 2 & 3 & 1 & 3 & 0 & 1024 & 0 & \texttt{----------}\\
6 & 1 & 2 & 2 & 1 & 256 & 256 & 512 & \texttt{-------??-}\\
7 & 6 & 3 & 1 & 2 & 0 & 1024 & 0 & \texttt{----------}\\
8 & 6 & 3 & 1 & 2 & 0 & 1024 & 0 & \texttt{----------}\\
9 & 0 & 1 & 1 & 1 & 0 & 1024 & 0 & \texttt{----------}\\
\multicolumn{9}{@{}l@{}}{\rule{0pt}{2.6ex}\textbf{T10.565}\enspace \texttt{0123332323}\enspace $n{=}10$, $|H|{=}5$, len 4, br 2, $|\mathrm{Aut}|{=}12$, orb 6/3, hist-orb 2/0; M 22528, res 10080, rev 7354: ok}\\*
0 & -- & 0 & 5 & 1 & 674 & 2 & 348 & \texttt{-+????????}\\*
1 & 0 & 1 & 5 & 1 & 324 & 4 & 696 & \texttt{--????????}\\
2 & 1 & 2 & 3 & 1 & 128 & 128 & 768 & \texttt{---???----}\\
3 & 2 & 3 & 1 & 3 & 0 & 1024 & 0 & \texttt{----------}\\
4 & 2 & 3 & 1 & 3 & 0 & 1024 & 0 & \texttt{----------}\\
5 & 2 & 3 & 1 & 3 & 0 & 1024 & 0 & \texttt{----------}\\
6 & 1 & 2 & 1 & 2 & 512 & 512 & 0 & \texttt{-------+--}\\
7 & 6 & 3 & 1 & 2 & 0 & 1024 & 0 & \texttt{----------}\\
8 & 1 & 2 & 1 & 2 & 512 & 512 & 0 & \texttt{---------+}\\
9 & 8 & 3 & 1 & 2 & 0 & 1024 & 0 & \texttt{----------}\\
\multicolumn{9}{@{}l@{}}{\rule{0pt}{2.6ex}\textbf{T10.566}\enspace \texttt{0123332322}\enspace $n{=}10$, $|H|{=}6$, len 4, br 2, $|\mathrm{Aut}|{=}12$, orb 7/5, hist-orb 3/1; M 21504, res 12672, rev 5334: ok}\\*
0 & -- & 0 & 6 & 1 & 566 & 2 & 456 & \texttt{-+????????}\\*
1 & 0 & 1 & 6 & 1 & 108 & 4 & 912 & \texttt{--????????}\\
2 & 1 & 2 & 3 & 1 & 128 & 128 & 768 & \texttt{---???----}\\
3 & 2 & 3 & 1 & 3 & 0 & 1024 & 0 & \texttt{----------}\\
4 & 2 & 3 & 1 & 3 & 0 & 1024 & 0 & \texttt{----------}\\
5 & 2 & 3 & 1 & 3 & 0 & 1024 & 0 & \texttt{----------}\\
6 & 1 & 2 & 1 & 1 & 512 & 512 & 0 & \texttt{-------+--}\\
7 & 6 & 3 & 1 & 1 & 0 & 1024 & 0 & \texttt{----------}\\
8 & 1 & 2 & 1 & 2 & 0 & 1024 & 0 & \texttt{----------}\\
9 & 1 & 2 & 1 & 2 & 0 & 1024 & 0 & \texttt{----------}\\
\multicolumn{9}{@{}l@{}}{\rule{0pt}{2.6ex}\textbf{T10.567}\enspace \texttt{0123332321}\enspace $n{=}10$, $|H|{=}6$, len 4, br 3, $|\mathrm{Aut}|{=}6$, orb 8/7, hist-orb 4/3; M 20480, res 13336, rev 4118: ok}\\*
0 & -- & 0 & 6 & 1 & 310 & 2 & 712 & \texttt{-?????????}\\*
1 & 0 & 1 & 5 & 1 & 216 & 8 & 800 & \texttt{--???????-}\\
2 & 1 & 2 & 3 & 1 & 128 & 128 & 768 & \texttt{---???----}\\
3 & 2 & 3 & 1 & 3 & 0 & 1024 & 0 & \texttt{----------}\\
4 & 2 & 3 & 1 & 3 & 0 & 1024 & 0 & \texttt{----------}\\
5 & 2 & 3 & 1 & 3 & 0 & 1024 & 0 & \texttt{----------}\\
6 & 1 & 2 & 1 & 1 & 512 & 512 & 0 & \texttt{-------+--}\\
7 & 6 & 3 & 1 & 1 & 0 & 1024 & 0 & \texttt{----------}\\
8 & 1 & 2 & 1 & 1 & 0 & 1024 & 0 & \texttt{----------}\\
9 & 0 & 1 & 1 & 1 & 0 & 1024 & 0 & \texttt{----------}\\
\multicolumn{9}{@{}l@{}}{\rule{0pt}{2.6ex}\textbf{T10.568}\enspace \texttt{0123332312}\enspace $n{=}10$, $|H|{=}5$, len 4, br 3, $|\mathrm{Aut}|{=}6$, orb 8/7, hist-orb 3/2; M 20480, res 9180, rev 6558: ok}\\*
0 & -- & 0 & 5 & 1 & 546 & 2 & 476 & \texttt{-?????????}\\*
1 & 0 & 1 & 4 & 1 & 432 & 16 & 576 & \texttt{--??????--}\\
2 & 1 & 2 & 3 & 1 & 128 & 128 & 768 & \texttt{---???----}\\
3 & 2 & 3 & 1 & 3 & 0 & 1024 & 0 & \texttt{----------}\\
4 & 2 & 3 & 1 & 3 & 0 & 1024 & 0 & \texttt{----------}\\
5 & 2 & 3 & 1 & 3 & 0 & 1024 & 0 & \texttt{----------}\\
6 & 1 & 2 & 1 & 1 & 512 & 512 & 0 & \texttt{-------+--}\\
7 & 6 & 3 & 1 & 1 & 0 & 1024 & 0 & \texttt{----------}\\
8 & 0 & 1 & 1 & 1 & 512 & 512 & 0 & \texttt{---------+}\\
9 & 8 & 2 & 1 & 1 & 0 & 1024 & 0 & \texttt{----------}\\
\multicolumn{9}{@{}l@{}}{\rule{0pt}{2.6ex}\textbf{T10.569}\enspace \texttt{0123332311}\enspace $n{=}10$, $|H|{=}6$, len 4, br 3, $|\mathrm{Aut}|{=}12$, orb 7/5, hist-orb 3/1; M 19456, res 12456, rev 4014: ok}\\*
0 & -- & 0 & 6 & 1 & 182 & 2 & 840 & \texttt{-?????????}\\*
1 & 0 & 1 & 4 & 1 & 432 & 16 & 576 & \texttt{--??????--}\\
2 & 1 & 2 & 3 & 1 & 128 & 128 & 768 & \texttt{---???----}\\
3 & 2 & 3 & 1 & 3 & 0 & 1024 & 0 & \texttt{----------}\\
4 & 2 & 3 & 1 & 3 & 0 & 1024 & 0 & \texttt{----------}\\
5 & 2 & 3 & 1 & 3 & 0 & 1024 & 0 & \texttt{----------}\\
6 & 1 & 2 & 1 & 1 & 512 & 512 & 0 & \texttt{-------+--}\\
7 & 6 & 3 & 1 & 1 & 0 & 1024 & 0 & \texttt{----------}\\
8 & 0 & 1 & 1 & 2 & 0 & 1024 & 0 & \texttt{----------}\\
9 & 0 & 1 & 1 & 2 & 0 & 1024 & 0 & \texttt{----------}\\
\multicolumn{9}{@{}l@{}}{\rule{0pt}{2.6ex}\textbf{T10.570}\enspace \texttt{0123332222}\enspace $n{=}10$, $|H|{=}7$, len 4, br 2, $|\mathrm{Aut}|{=}144$, orb 5/3, hist-orb 2/0; M 20480, res 13536, rev 4418: ok}\\*
0 & -- & 0 & 7 & 1 & 530 & 2 & 492 & \texttt{-+????????}\\*
1 & 0 & 1 & 7 & 1 & 36 & 4 & 984 & \texttt{--????????}\\
2 & 1 & 2 & 3 & 1 & 128 & 128 & 768 & \texttt{---???----}\\
3 & 2 & 3 & 1 & 3 & 0 & 1024 & 0 & \texttt{----------}\\
4 & 2 & 3 & 1 & 3 & 0 & 1024 & 0 & \texttt{----------}\\
5 & 2 & 3 & 1 & 3 & 0 & 1024 & 0 & \texttt{----------}\\
6 & 1 & 2 & 1 & 4 & 0 & 1024 & 0 & \texttt{----------}\\
7 & 1 & 2 & 1 & 4 & 0 & 1024 & 0 & \texttt{----------}\\
8 & 1 & 2 & 1 & 4 & 0 & 1024 & 0 & \texttt{----------}\\
9 & 1 & 2 & 1 & 4 & 0 & 1024 & 0 & \texttt{----------}\\
\multicolumn{9}{@{}l@{}}{\rule{0pt}{2.6ex}\textbf{T10.571}\enspace \texttt{0123332221}\enspace $n{=}10$, $|H|{=}7$, len 4, br 3, $|\mathrm{Aut}|{=}36$, orb 6/4, hist-orb 3/1; M 19456, res 14596, rev 2778: ok}\\*
0 & -- & 0 & 7 & 1 & 274 & 2 & 748 & \texttt{-?????????}\\*
1 & 0 & 1 & 6 & 1 & 72 & 8 & 944 & \texttt{--???????-}\\
2 & 1 & 2 & 3 & 1 & 128 & 128 & 768 & \texttt{---???----}\\
3 & 2 & 3 & 1 & 3 & 0 & 1024 & 0 & \texttt{----------}\\
4 & 2 & 3 & 1 & 3 & 0 & 1024 & 0 & \texttt{----------}\\
5 & 2 & 3 & 1 & 3 & 0 & 1024 & 0 & \texttt{----------}\\
6 & 1 & 2 & 1 & 3 & 0 & 1024 & 0 & \texttt{----------}\\
7 & 1 & 2 & 1 & 3 & 0 & 1024 & 0 & \texttt{----------}\\
8 & 1 & 2 & 1 & 3 & 0 & 1024 & 0 & \texttt{----------}\\
9 & 0 & 1 & 1 & 1 & 0 & 1024 & 0 & \texttt{----------}\\
\multicolumn{9}{@{}l@{}}{\rule{0pt}{2.6ex}\textbf{T10.572}\enspace \texttt{0123332212}\enspace $n{=}10$, $|H|{=}6$, len 4, br 3, $|\mathrm{Aut}|{=}12$, orb 7/5, hist-orb 3/1; M 19456, res 11736, rev 4466: ok}\\*
0 & -- & 0 & 6 & 1 & 438 & 2 & 584 & \texttt{-?????????}\\*
1 & 0 & 1 & 5 & 1 & 144 & 16 & 864 & \texttt{--??????--}\\
2 & 1 & 2 & 3 & 1 & 128 & 128 & 768 & \texttt{---???----}\\
3 & 2 & 3 & 1 & 3 & 0 & 1024 & 0 & \texttt{----------}\\
4 & 2 & 3 & 1 & 3 & 0 & 1024 & 0 & \texttt{----------}\\
5 & 2 & 3 & 1 & 3 & 0 & 1024 & 0 & \texttt{----------}\\
6 & 1 & 2 & 1 & 2 & 0 & 1024 & 0 & \texttt{----------}\\
7 & 1 & 2 & 1 & 2 & 0 & 1024 & 0 & \texttt{----------}\\
8 & 0 & 1 & 1 & 1 & 512 & 512 & 0 & \texttt{---------+}\\
9 & 8 & 2 & 1 & 1 & 0 & 1024 & 0 & \texttt{----------}\\
\multicolumn{9}{@{}l@{}}{\rule{0pt}{2.6ex}\textbf{T10.573}\enspace \texttt{0123332211}\enspace $n{=}10$, $|H|{=}7$, len 4, br 3, $|\mathrm{Aut}|{=}24$, orb 6/3, hist-orb 3/0; M 18432, res 14364, rev 2162: ok}\\*
0 & -- & 0 & 7 & 1 & 146 & 2 & 876 & \texttt{-?????????}\\*
1 & 0 & 1 & 5 & 1 & 144 & 16 & 864 & \texttt{--??????--}\\
2 & 1 & 2 & 3 & 1 & 128 & 128 & 768 & \texttt{---???----}\\
3 & 2 & 3 & 1 & 3 & 0 & 1024 & 0 & \texttt{----------}\\
4 & 2 & 3 & 1 & 3 & 0 & 1024 & 0 & \texttt{----------}\\
5 & 2 & 3 & 1 & 3 & 0 & 1024 & 0 & \texttt{----------}\\
6 & 1 & 2 & 1 & 2 & 0 & 1024 & 0 & \texttt{----------}\\
7 & 1 & 2 & 1 & 2 & 0 & 1024 & 0 & \texttt{----------}\\
8 & 0 & 1 & 1 & 2 & 0 & 1024 & 0 & \texttt{----------}\\
9 & 0 & 1 & 1 & 2 & 0 & 1024 & 0 & \texttt{----------}\\
\multicolumn{9}{@{}l@{}}{\rule{0pt}{2.6ex}\textbf{T10.574}\enspace \texttt{0123332123}\enspace $n{=}10$, $|H|{=}5$, len 4, br 3, $|\mathrm{Aut}|{=}6$, orb 8/7, hist-orb 3/2; M 20480, res 8928, rev 6386: ok}\\*
0 & -- & 0 & 5 & 1 & 574 & 2 & 448 & \texttt{-?????????}\\*
1 & 0 & 1 & 4 & 1 & 288 & 32 & 704 & \texttt{--?????---}\\
2 & 1 & 2 & 3 & 1 & 128 & 128 & 768 & \texttt{---???----}\\
3 & 2 & 3 & 1 & 3 & 0 & 1024 & 0 & \texttt{----------}\\
4 & 2 & 3 & 1 & 3 & 0 & 1024 & 0 & \texttt{----------}\\
5 & 2 & 3 & 1 & 3 & 0 & 1024 & 0 & \texttt{----------}\\
6 & 1 & 2 & 1 & 1 & 0 & 1024 & 0 & \texttt{----------}\\
7 & 0 & 1 & 1 & 1 & 768 & 256 & 0 & \texttt{--------++}\\
8 & 7 & 2 & 1 & 1 & 512 & 512 & 0 & \texttt{---------+}\\
9 & 8 & 3 & 1 & 1 & 0 & 1024 & 0 & \texttt{----------}\\
\multicolumn{9}{@{}l@{}}{\rule{0pt}{2.6ex}\textbf{T10.575}\enspace \texttt{0123332122}\enspace $n{=}10$, $|H|{=}6$, len 4, br 4, $|\mathrm{Aut}|{=}12$, orb 7/5, hist-orb 3/1; M 19456, res 11172, rev 4682: ok}\\*
0 & -- & 0 & 6 & 1 & 410 & 2 & 612 & \texttt{-?????????}\\*
1 & 0 & 1 & 4 & 1 & 288 & 32 & 704 & \texttt{--?????---}\\
2 & 1 & 2 & 3 & 1 & 128 & 128 & 768 & \texttt{---???----}\\
3 & 2 & 3 & 1 & 3 & 0 & 1024 & 0 & \texttt{----------}\\
4 & 2 & 3 & 1 & 3 & 0 & 1024 & 0 & \texttt{----------}\\
5 & 2 & 3 & 1 & 3 & 0 & 1024 & 0 & \texttt{----------}\\
6 & 1 & 2 & 1 & 1 & 0 & 1024 & 0 & \texttt{----------}\\
7 & 0 & 1 & 2 & 1 & 256 & 256 & 512 & \texttt{--------??}\\
8 & 7 & 2 & 1 & 2 & 0 & 1024 & 0 & \texttt{----------}\\
9 & 7 & 2 & 1 & 2 & 0 & 1024 & 0 & \texttt{----------}\\
\multicolumn{9}{@{}l@{}}{\rule{0pt}{2.6ex}\textbf{T10.576}\enspace \texttt{0123332121}\enspace $n{=}10$, $|H|{=}6$, len 4, br 3, $|\mathrm{Aut}|{=}6$, orb 8/7, hist-orb 4/3; M 18432, res 11880, rev 3674: ok}\\*
0 & -- & 0 & 6 & 1 & 246 & 2 & 776 & \texttt{-?????????}\\*
1 & 0 & 1 & 4 & 1 & 288 & 32 & 704 & \texttt{--?????---}\\
2 & 1 & 2 & 3 & 1 & 128 & 128 & 768 & \texttt{---???----}\\
3 & 2 & 3 & 1 & 3 & 0 & 1024 & 0 & \texttt{----------}\\
4 & 2 & 3 & 1 & 3 & 0 & 1024 & 0 & \texttt{----------}\\
5 & 2 & 3 & 1 & 3 & 0 & 1024 & 0 & \texttt{----------}\\
6 & 1 & 2 & 1 & 1 & 0 & 1024 & 0 & \texttt{----------}\\
7 & 0 & 1 & 1 & 1 & 512 & 512 & 0 & \texttt{--------+-}\\
8 & 7 & 2 & 1 & 1 & 0 & 1024 & 0 & \texttt{----------}\\
9 & 0 & 1 & 1 & 1 & 0 & 1024 & 0 & \texttt{----------}\\
\multicolumn{9}{@{}l@{}}{\rule{0pt}{2.6ex}\textbf{T10.577}\enspace \texttt{0123332111}\enspace $n{=}10$, $|H|{=}7$, len 4, br 3, $|\mathrm{Aut}|{=}36$, orb 6/4, hist-orb 3/1; M 17408, res 13356, rev 2154: ok}\\*
0 & -- & 0 & 7 & 1 & 82 & 2 & 940 & \texttt{-?????????}\\*
1 & 0 & 1 & 4 & 1 & 288 & 32 & 704 & \texttt{--?????---}\\
2 & 1 & 2 & 3 & 1 & 128 & 128 & 768 & \texttt{---???----}\\
3 & 2 & 3 & 1 & 3 & 0 & 1024 & 0 & \texttt{----------}\\
4 & 2 & 3 & 1 & 3 & 0 & 1024 & 0 & \texttt{----------}\\
5 & 2 & 3 & 1 & 3 & 0 & 1024 & 0 & \texttt{----------}\\
6 & 1 & 2 & 1 & 1 & 0 & 1024 & 0 & \texttt{----------}\\
7 & 0 & 1 & 1 & 3 & 0 & 1024 & 0 & \texttt{----------}\\
8 & 0 & 1 & 1 & 3 & 0 & 1024 & 0 & \texttt{----------}\\
9 & 0 & 1 & 1 & 3 & 0 & 1024 & 0 & \texttt{----------}\\
\multicolumn{9}{@{}l@{}}{\rule{0pt}{2.6ex}\textbf{T10.578}\enspace \texttt{0123331233}\enspace $n{=}10$, $|H|{=}5$, len 4, br 3, $|\mathrm{Aut}|{=}12$, orb 7/5, hist-orb 2/0; M 21504, res 7700, rev 7650: ok}\\*
0 & -- & 0 & 5 & 1 & 650 & 2 & 372 & \texttt{-?????????}\\*
1 & 0 & 1 & 3 & 1 & 576 & 64 & 384 & \texttt{--+???----}\\
2 & 1 & 2 & 3 & 1 & 128 & 128 & 768 & \texttt{---???----}\\
3 & 2 & 3 & 1 & 3 & 0 & 1024 & 0 & \texttt{----------}\\
4 & 2 & 3 & 1 & 3 & 0 & 1024 & 0 & \texttt{----------}\\
5 & 2 & 3 & 1 & 3 & 0 & 1024 & 0 & \texttt{----------}\\
6 & 0 & 1 & 2 & 1 & 640 & 128 & 256 & \texttt{-------+??}\\
7 & 6 & 2 & 2 & 1 & 256 & 256 & 512 & \texttt{--------??}\\
8 & 7 & 3 & 1 & 2 & 0 & 1024 & 0 & \texttt{----------}\\
9 & 7 & 3 & 1 & 2 & 0 & 1024 & 0 & \texttt{----------}\\
\multicolumn{9}{@{}l@{}}{\rule{0pt}{2.6ex}\textbf{T10.579}\enspace \texttt{0123331232}\enspace $n{=}10$, $|H|{=}5$, len 4, br 3, $|\mathrm{Aut}|{=}6$, orb 8/7, hist-orb 3/2; M 20480, res 8600, rev 6722: ok}\\*
0 & -- & 0 & 5 & 1 & 550 & 2 & 472 & \texttt{-?????????}\\*
1 & 0 & 1 & 3 & 1 & 576 & 64 & 384 & \texttt{--+???----}\\
2 & 1 & 2 & 3 & 1 & 128 & 128 & 768 & \texttt{---???----}\\
3 & 2 & 3 & 1 & 3 & 0 & 1024 & 0 & \texttt{----------}\\
4 & 2 & 3 & 1 & 3 & 0 & 1024 & 0 & \texttt{----------}\\
5 & 2 & 3 & 1 & 3 & 0 & 1024 & 0 & \texttt{----------}\\
6 & 0 & 1 & 2 & 1 & 384 & 128 & 512 & \texttt{-------???}\\
7 & 6 & 2 & 1 & 1 & 512 & 512 & 0 & \texttt{--------+-}\\
8 & 7 & 3 & 1 & 1 & 0 & 1024 & 0 & \texttt{----------}\\
9 & 6 & 2 & 1 & 1 & 0 & 1024 & 0 & \texttt{----------}\\
\multicolumn{9}{@{}l@{}}{\rule{0pt}{2.6ex}\textbf{T10.580}\enspace \texttt{0123331231}\enspace $n{=}10$, $|H|{=}5$, len 4, br 2, $|\mathrm{Aut}|{=}6$, orb 8/7, hist-orb 3/2; M 19456, res 9120, rev 5718: ok}\\*
0 & -- & 0 & 5 & 1 & 350 & 2 & 672 & \texttt{-?????????}\\*
1 & 0 & 1 & 3 & 1 & 576 & 64 & 384 & \texttt{--+???----}\\
2 & 1 & 2 & 3 & 1 & 128 & 128 & 768 & \texttt{---???----}\\
3 & 2 & 3 & 1 & 3 & 0 & 1024 & 0 & \texttt{----------}\\
4 & 2 & 3 & 1 & 3 & 0 & 1024 & 0 & \texttt{----------}\\
5 & 2 & 3 & 1 & 3 & 0 & 1024 & 0 & \texttt{----------}\\
6 & 0 & 1 & 1 & 1 & 768 & 256 & 0 & \texttt{-------++-}\\
7 & 6 & 2 & 1 & 1 & 512 & 512 & 0 & \texttt{--------+-}\\
8 & 7 & 3 & 1 & 1 & 0 & 1024 & 0 & \texttt{----------}\\
9 & 0 & 1 & 1 & 1 & 0 & 1024 & 0 & \texttt{----------}\\
\multicolumn{9}{@{}l@{}}{\rule{0pt}{2.6ex}\textbf{T10.581}\enspace \texttt{0123331222}\enspace $n{=}10$, $|H|{=}6$, len 4, br 3, $|\mathrm{Aut}|{=}36$, orb 6/4, hist-orb 2/0; M 19456, res 10140, rev 5382: ok}\\*
0 & -- & 0 & 6 & 1 & 450 & 2 & 572 & \texttt{-?????????}\\*
1 & 0 & 1 & 3 & 1 & 576 & 64 & 384 & \texttt{--+???----}\\
2 & 1 & 2 & 3 & 1 & 128 & 128 & 768 & \texttt{---???----}\\
3 & 2 & 3 & 1 & 3 & 0 & 1024 & 0 & \texttt{----------}\\
4 & 2 & 3 & 1 & 3 & 0 & 1024 & 0 & \texttt{----------}\\
5 & 2 & 3 & 1 & 3 & 0 & 1024 & 0 & \texttt{----------}\\
6 & 0 & 1 & 3 & 1 & 128 & 128 & 768 & \texttt{-------???}\\
7 & 6 & 2 & 1 & 3 & 0 & 1024 & 0 & \texttt{----------}\\
8 & 6 & 2 & 1 & 3 & 0 & 1024 & 0 & \texttt{----------}\\
9 & 6 & 2 & 1 & 3 & 0 & 1024 & 0 & \texttt{----------}\\
\multicolumn{9}{@{}l@{}}{\rule{0pt}{2.6ex}\textbf{T10.582}\enspace \texttt{0123331221}\enspace $n{=}10$, $|H|{=}6$, len 4, br 3, $|\mathrm{Aut}|{=}12$, orb 7/5, hist-orb 3/1; M 18432, res 10788, rev 4278: ok}\\*
0 & -- & 0 & 6 & 1 & 250 & 2 & 772 & \texttt{-?????????}\\*
1 & 0 & 1 & 3 & 1 & 576 & 64 & 384 & \texttt{--+???----}\\
2 & 1 & 2 & 3 & 1 & 128 & 128 & 768 & \texttt{---???----}\\
3 & 2 & 3 & 1 & 3 & 0 & 1024 & 0 & \texttt{----------}\\
4 & 2 & 3 & 1 & 3 & 0 & 1024 & 0 & \texttt{----------}\\
5 & 2 & 3 & 1 & 3 & 0 & 1024 & 0 & \texttt{----------}\\
6 & 0 & 1 & 2 & 1 & 256 & 256 & 512 & \texttt{-------??-}\\
7 & 6 & 2 & 1 & 2 & 0 & 1024 & 0 & \texttt{----------}\\
8 & 6 & 2 & 1 & 2 & 0 & 1024 & 0 & \texttt{----------}\\
9 & 0 & 1 & 1 & 1 & 0 & 1024 & 0 & \texttt{----------}\\
\multicolumn{9}{@{}l@{}}{\rule{0pt}{2.6ex}\textbf{T10.583}\enspace \texttt{0123331212}\enspace $n{=}10$, $|H|{=}5$, len 4, br 2, $|\mathrm{Aut}|{=}12$, orb 6/3, hist-orb 2/0; M 18432, res 8220, rev 5822: ok}\\*
0 & -- & 0 & 5 & 1 & 450 & 2 & 572 & \texttt{-?????????}\\*
1 & 0 & 1 & 3 & 1 & 576 & 64 & 384 & \texttt{--+???----}\\
2 & 1 & 2 & 3 & 1 & 128 & 128 & 768 & \texttt{---???----}\\
3 & 2 & 3 & 1 & 3 & 0 & 1024 & 0 & \texttt{----------}\\
4 & 2 & 3 & 1 & 3 & 0 & 1024 & 0 & \texttt{----------}\\
5 & 2 & 3 & 1 & 3 & 0 & 1024 & 0 & \texttt{----------}\\
6 & 0 & 1 & 1 & 2 & 512 & 512 & 0 & \texttt{-------+--}\\
7 & 6 & 2 & 1 & 2 & 0 & 1024 & 0 & \texttt{----------}\\
8 & 0 & 1 & 1 & 2 & 512 & 512 & 0 & \texttt{---------+}\\
9 & 8 & 2 & 1 & 2 & 0 & 1024 & 0 & \texttt{----------}\\
\multicolumn{9}{@{}l@{}}{\rule{0pt}{2.6ex}\textbf{T10.584}\enspace \texttt{0123331211}\enspace $n{=}10$, $|H|{=}6$, len 4, br 2, $|\mathrm{Aut}|{=}12$, orb 7/5, hist-orb 3/1; M 17408, res 10920, rev 3662: ok}\\*
0 & -- & 0 & 6 & 1 & 150 & 2 & 872 & \texttt{-?????????}\\*
1 & 0 & 1 & 3 & 1 & 576 & 64 & 384 & \texttt{--+???----}\\
2 & 1 & 2 & 3 & 1 & 128 & 128 & 768 & \texttt{---???----}\\
3 & 2 & 3 & 1 & 3 & 0 & 1024 & 0 & \texttt{----------}\\
4 & 2 & 3 & 1 & 3 & 0 & 1024 & 0 & \texttt{----------}\\
5 & 2 & 3 & 1 & 3 & 0 & 1024 & 0 & \texttt{----------}\\
6 & 0 & 1 & 1 & 1 & 512 & 512 & 0 & \texttt{-------+--}\\
7 & 6 & 2 & 1 & 1 & 0 & 1024 & 0 & \texttt{----------}\\
8 & 0 & 1 & 1 & 2 & 0 & 1024 & 0 & \texttt{----------}\\
9 & 0 & 1 & 1 & 2 & 0 & 1024 & 0 & \texttt{----------}\\
\multicolumn{9}{@{}l@{}}{\rule{0pt}{2.6ex}\textbf{T10.585}\enspace \texttt{0123331111}\enspace $n{=}10$, $|H|{=}7$, len 4, br 2, $|\mathrm{Aut}|{=}144$, orb 5/3, hist-orb 2/0; M 16384, res 11820, rev 2534: ok}\\*
0 & -- & 0 & 7 & 1 & 50 & 2 & 972 & \texttt{-?????????}\\*
1 & 0 & 1 & 3 & 1 & 576 & 64 & 384 & \texttt{--+???----}\\
2 & 1 & 2 & 3 & 1 & 128 & 128 & 768 & \texttt{---???----}\\
3 & 2 & 3 & 1 & 3 & 0 & 1024 & 0 & \texttt{----------}\\
4 & 2 & 3 & 1 & 3 & 0 & 1024 & 0 & \texttt{----------}\\
5 & 2 & 3 & 1 & 3 & 0 & 1024 & 0 & \texttt{----------}\\
6 & 0 & 1 & 1 & 4 & 0 & 1024 & 0 & \texttt{----------}\\
7 & 0 & 1 & 1 & 4 & 0 & 1024 & 0 & \texttt{----------}\\
8 & 0 & 1 & 1 & 4 & 0 & 1024 & 0 & \texttt{----------}\\
9 & 0 & 1 & 1 & 4 & 0 & 1024 & 0 & \texttt{----------}\\
\multicolumn{9}{@{}l@{}}{\rule{0pt}{2.6ex}\textbf{T10.586}\enspace \texttt{0123323323}\enspace $n{=}10$, $|H|{=}5$, len 4, br 3, $|\mathrm{Aut}|{=}8$, orb 6/4, hist-orb 2/1; M 22528, res 9920, rev 7190: ok}\\*
0 & -- & 0 & 5 & 1 & 662 & 2 & 360 & \texttt{-+????????}\\*
1 & 0 & 1 & 5 & 1 & 300 & 4 & 720 & \texttt{--????????}\\
2 & 1 & 2 & 2 & 2 & 256 & 256 & 512 & \texttt{---??-----}\\
3 & 2 & 3 & 1 & 4 & 0 & 1024 & 0 & \texttt{----------}\\
4 & 2 & 3 & 1 & 4 & 0 & 1024 & 0 & \texttt{----------}\\
5 & 1 & 2 & 2 & 2 & 256 & 256 & 512 & \texttt{------??--}\\
6 & 5 & 3 & 1 & 4 & 0 & 1024 & 0 & \texttt{----------}\\
7 & 5 & 3 & 1 & 4 & 0 & 1024 & 0 & \texttt{----------}\\
8 & 1 & 2 & 1 & 1 & 512 & 512 & 0 & \texttt{---------+}\\
9 & 8 & 3 & 1 & 1 & 0 & 1024 & 0 & \texttt{----------}\\
\multicolumn{9}{@{}l@{}}{\rule{0pt}{2.6ex}\textbf{T10.587}\enspace \texttt{0123323322}\enspace $n{=}10$, $|H|{=}6$, len 4, br 3, $|\mathrm{Aut}|{=}16$, orb 5/2, hist-orb 2/0; M 21504, res 12320, rev 5314: ok}\\*
0 & -- & 0 & 6 & 1 & 562 & 2 & 460 & \texttt{-+????????}\\*
1 & 0 & 1 & 6 & 1 & 100 & 4 & 920 & \texttt{--????????}\\
2 & 1 & 2 & 2 & 2 & 256 & 256 & 512 & \texttt{---??-----}\\
3 & 2 & 3 & 1 & 4 & 0 & 1024 & 0 & \texttt{----------}\\
4 & 2 & 3 & 1 & 4 & 0 & 1024 & 0 & \texttt{----------}\\
5 & 1 & 2 & 2 & 2 & 256 & 256 & 512 & \texttt{------??--}\\
6 & 5 & 3 & 1 & 4 & 0 & 1024 & 0 & \texttt{----------}\\
7 & 5 & 3 & 1 & 4 & 0 & 1024 & 0 & \texttt{----------}\\
8 & 1 & 2 & 1 & 2 & 0 & 1024 & 0 & \texttt{----------}\\
9 & 1 & 2 & 1 & 2 & 0 & 1024 & 0 & \texttt{----------}\\
\multicolumn{9}{@{}l@{}}{\rule{0pt}{2.6ex}\textbf{T10.588}\enspace \texttt{0123323321}\enspace $n{=}10$, $|H|{=}6$, len 4, br 4, $|\mathrm{Aut}|{=}8$, orb 6/4, hist-orb 3/2; M 20480, res 13028, rev 4090: ok}\\*
0 & -- & 0 & 6 & 1 & 306 & 2 & 716 & \texttt{-?????????}\\*
1 & 0 & 1 & 5 & 1 & 200 & 8 & 816 & \texttt{--???????-}\\
2 & 1 & 2 & 2 & 2 & 256 & 256 & 512 & \texttt{---??-----}\\
3 & 2 & 3 & 1 & 4 & 0 & 1024 & 0 & \texttt{----------}\\
4 & 2 & 3 & 1 & 4 & 0 & 1024 & 0 & \texttt{----------}\\
5 & 1 & 2 & 2 & 2 & 256 & 256 & 512 & \texttt{------??--}\\
6 & 5 & 3 & 1 & 4 & 0 & 1024 & 0 & \texttt{----------}\\
7 & 5 & 3 & 1 & 4 & 0 & 1024 & 0 & \texttt{----------}\\
8 & 1 & 2 & 1 & 1 & 0 & 1024 & 0 & \texttt{----------}\\
9 & 0 & 1 & 1 & 1 & 0 & 1024 & 0 & \texttt{----------}\\
\multicolumn{9}{@{}l@{}}{\rule{0pt}{2.6ex}\textbf{T10.589}\enspace \texttt{0123323312}\enspace $n{=}10$, $|H|{=}5$, len 4, br 4, $|\mathrm{Aut}|{=}8$, orb 6/4, hist-orb 2/1; M 20480, res 9016, rev 6402: ok}\\*
0 & -- & 0 & 5 & 1 & 534 & 2 & 488 & \texttt{-?????????}\\*
1 & 0 & 1 & 4 & 1 & 400 & 16 & 608 & \texttt{--??????--}\\
2 & 1 & 2 & 2 & 2 & 256 & 256 & 512 & \texttt{---??-----}\\
3 & 2 & 3 & 1 & 4 & 0 & 1024 & 0 & \texttt{----------}\\
4 & 2 & 3 & 1 & 4 & 0 & 1024 & 0 & \texttt{----------}\\
5 & 1 & 2 & 2 & 2 & 256 & 256 & 512 & \texttt{------??--}\\
6 & 5 & 3 & 1 & 4 & 0 & 1024 & 0 & \texttt{----------}\\
7 & 5 & 3 & 1 & 4 & 0 & 1024 & 0 & \texttt{----------}\\
8 & 0 & 1 & 1 & 1 & 512 & 512 & 0 & \texttt{---------+}\\
9 & 8 & 2 & 1 & 1 & 0 & 1024 & 0 & \texttt{----------}\\
\multicolumn{9}{@{}l@{}}{\rule{0pt}{2.6ex}\textbf{T10.590}\enspace \texttt{0123323311}\enspace $n{=}10$, $|H|{=}6$, len 4, br 4, $|\mathrm{Aut}|{=}16$, orb 5/2, hist-orb 2/0; M 19456, res 12220, rev 3938: ok}\\*
0 & -- & 0 & 6 & 1 & 178 & 2 & 844 & \texttt{-?????????}\\*
1 & 0 & 1 & 4 & 1 & 400 & 16 & 608 & \texttt{--??????--}\\
2 & 1 & 2 & 2 & 2 & 256 & 256 & 512 & \texttt{---??-----}\\
3 & 2 & 3 & 1 & 4 & 0 & 1024 & 0 & \texttt{----------}\\
4 & 2 & 3 & 1 & 4 & 0 & 1024 & 0 & \texttt{----------}\\
5 & 1 & 2 & 2 & 2 & 256 & 256 & 512 & \texttt{------??--}\\
6 & 5 & 3 & 1 & 4 & 0 & 1024 & 0 & \texttt{----------}\\
7 & 5 & 3 & 1 & 4 & 0 & 1024 & 0 & \texttt{----------}\\
8 & 0 & 1 & 1 & 2 & 0 & 1024 & 0 & \texttt{----------}\\
9 & 0 & 1 & 1 & 2 & 0 & 1024 & 0 & \texttt{----------}\\
\multicolumn{9}{@{}l@{}}{\rule{0pt}{2.6ex}\textbf{T10.591}\enspace \texttt{0123323232}\enspace $n{=}10$, $|H|{=}5$, len 4, br 2, $|\mathrm{Aut}|{=}4$, orb 7/4, hist-orb 3/1; M 21504, res 10720, rev 6322: ok}\\*
0 & -- & 0 & 5 & 1 & 602 & 2 & 420 & \texttt{-+????????}\\*
1 & 0 & 1 & 5 & 1 & 180 & 4 & 840 & \texttt{--????????}\\
2 & 1 & 2 & 2 & 1 & 256 & 256 & 512 & \texttt{---??-----}\\
3 & 2 & 3 & 1 & 2 & 0 & 1024 & 0 & \texttt{----------}\\
4 & 2 & 3 & 1 & 2 & 0 & 1024 & 0 & \texttt{----------}\\
5 & 1 & 2 & 1 & 2 & 512 & 512 & 0 & \texttt{------+---}\\
6 & 5 & 3 & 1 & 2 & 0 & 1024 & 0 & \texttt{----------}\\
7 & 1 & 2 & 1 & 2 & 512 & 512 & 0 & \texttt{--------+-}\\
8 & 7 & 3 & 1 & 2 & 0 & 1024 & 0 & \texttt{----------}\\
9 & 1 & 2 & 1 & 1 & 0 & 1024 & 0 & \texttt{----------}\\
\multicolumn{9}{@{}l@{}}{\rule{0pt}{2.6ex}\textbf{T10.592}\enspace \texttt{0123323231}\enspace $n{=}10$, $|H|{=}5$, len 4, br 3, $|\mathrm{Aut}|{=}4$, orb 7/4, hist-orb 3/1; M 20480, res 10988, rev 5418: ok}\\*
0 & -- & 0 & 5 & 1 & 346 & 2 & 676 & \texttt{-?????????}\\*
1 & 0 & 1 & 4 & 1 & 360 & 8 & 656 & \texttt{--???????-}\\
2 & 1 & 2 & 2 & 1 & 256 & 256 & 512 & \texttt{---??-----}\\
3 & 2 & 3 & 1 & 2 & 0 & 1024 & 0 & \texttt{----------}\\
4 & 2 & 3 & 1 & 2 & 0 & 1024 & 0 & \texttt{----------}\\
5 & 1 & 2 & 1 & 2 & 512 & 512 & 0 & \texttt{------+---}\\
6 & 5 & 3 & 1 & 2 & 0 & 1024 & 0 & \texttt{----------}\\
7 & 1 & 2 & 1 & 2 & 512 & 512 & 0 & \texttt{--------+-}\\
8 & 7 & 3 & 1 & 2 & 0 & 1024 & 0 & \texttt{----------}\\
9 & 0 & 1 & 1 & 1 & 0 & 1024 & 0 & \texttt{----------}\\
\multicolumn{9}{@{}l@{}}{\rule{0pt}{2.6ex}\textbf{T10.593}\enspace \texttt{0123323222}\enspace $n{=}10$, $|H|{=}6$, len 4, br 2, $|\mathrm{Aut}|{=}12$, orb 7/5, hist-orb 3/1; M 20480, res 12160, rev 5070: ok}\\*
0 & -- & 0 & 6 & 1 & 542 & 2 & 480 & \texttt{-+????????}\\*
1 & 0 & 1 & 6 & 1 & 60 & 4 & 960 & \texttt{--????????}\\
2 & 1 & 2 & 2 & 1 & 256 & 256 & 512 & \texttt{---??-----}\\
3 & 2 & 3 & 1 & 2 & 0 & 1024 & 0 & \texttt{----------}\\
4 & 2 & 3 & 1 & 2 & 0 & 1024 & 0 & \texttt{----------}\\
5 & 1 & 2 & 1 & 1 & 512 & 512 & 0 & \texttt{------+---}\\
6 & 5 & 3 & 1 & 1 & 0 & 1024 & 0 & \texttt{----------}\\
7 & 1 & 2 & 1 & 3 & 0 & 1024 & 0 & \texttt{----------}\\
8 & 1 & 2 & 1 & 3 & 0 & 1024 & 0 & \texttt{----------}\\
9 & 1 & 2 & 1 & 3 & 0 & 1024 & 0 & \texttt{----------}\\
\multicolumn{9}{@{}l@{}}{\rule{0pt}{2.6ex}\textbf{T10.594}\enspace \texttt{0123323221}\enspace $n{=}10$, $|H|{=}6$, len 4, br 3, $|\mathrm{Aut}|{=}4$, orb 8/6, hist-orb 4/2; M 19456, res 13088, rev 3614: ok}\\*
0 & -- & 0 & 6 & 1 & 286 & 2 & 736 & \texttt{-?????????}\\*
1 & 0 & 1 & 5 & 1 & 120 & 8 & 896 & \texttt{--???????-}\\
2 & 1 & 2 & 2 & 1 & 256 & 256 & 512 & \texttt{---??-----}\\
3 & 2 & 3 & 1 & 2 & 0 & 1024 & 0 & \texttt{----------}\\
4 & 2 & 3 & 1 & 2 & 0 & 1024 & 0 & \texttt{----------}\\
5 & 1 & 2 & 1 & 1 & 512 & 512 & 0 & \texttt{------+---}\\
6 & 5 & 3 & 1 & 1 & 0 & 1024 & 0 & \texttt{----------}\\
7 & 1 & 2 & 1 & 2 & 0 & 1024 & 0 & \texttt{----------}\\
8 & 1 & 2 & 1 & 2 & 0 & 1024 & 0 & \texttt{----------}\\
9 & 0 & 1 & 1 & 1 & 0 & 1024 & 0 & \texttt{----------}\\
\multicolumn{9}{@{}l@{}}{\rule{0pt}{2.6ex}\textbf{T10.595}\enspace \texttt{0123323212}\enspace $n{=}10$, $|H|{=}5$, len 4, br 3, $|\mathrm{Aut}|{=}2$, orb 9/8, hist-orb 4/3; M 19456, res 9796, rev 5494: ok}\\*
0 & -- & 0 & 5 & 1 & 474 & 2 & 548 & \texttt{-?????????}\\*
1 & 0 & 1 & 4 & 1 & 240 & 16 & 768 & \texttt{--??????--}\\
2 & 1 & 2 & 2 & 1 & 256 & 256 & 512 & \texttt{---??-----}\\
3 & 2 & 3 & 1 & 2 & 0 & 1024 & 0 & \texttt{----------}\\
4 & 2 & 3 & 1 & 2 & 0 & 1024 & 0 & \texttt{----------}\\
5 & 1 & 2 & 1 & 1 & 512 & 512 & 0 & \texttt{------+---}\\
6 & 5 & 3 & 1 & 1 & 0 & 1024 & 0 & \texttt{----------}\\
7 & 1 & 2 & 1 & 1 & 0 & 1024 & 0 & \texttt{----------}\\
8 & 0 & 1 & 1 & 1 & 512 & 512 & 0 & \texttt{---------+}\\
9 & 8 & 2 & 1 & 1 & 0 & 1024 & 0 & \texttt{----------}\\
\multicolumn{9}{@{}l@{}}{\rule{0pt}{2.6ex}\textbf{T10.596}\enspace \texttt{0123323211}\enspace $n{=}10$, $|H|{=}6$, len 4, br 3, $|\mathrm{Aut}|{=}4$, orb 8/6, hist-orb 4/2; M 18432, res 12640, rev 3174: ok}\\*
0 & -- & 0 & 6 & 1 & 158 & 2 & 864 & \texttt{-?????????}\\*
1 & 0 & 1 & 4 & 1 & 240 & 16 & 768 & \texttt{--??????--}\\
2 & 1 & 2 & 2 & 1 & 256 & 256 & 512 & \texttt{---??-----}\\
3 & 2 & 3 & 1 & 2 & 0 & 1024 & 0 & \texttt{----------}\\
4 & 2 & 3 & 1 & 2 & 0 & 1024 & 0 & \texttt{----------}\\
5 & 1 & 2 & 1 & 1 & 512 & 512 & 0 & \texttt{------+---}\\
6 & 5 & 3 & 1 & 1 & 0 & 1024 & 0 & \texttt{----------}\\
7 & 1 & 2 & 1 & 1 & 0 & 1024 & 0 & \texttt{----------}\\
8 & 0 & 1 & 1 & 2 & 0 & 1024 & 0 & \texttt{----------}\\
9 & 0 & 1 & 1 & 2 & 0 & 1024 & 0 & \texttt{----------}\\
\multicolumn{9}{@{}l@{}}{\rule{0pt}{2.6ex}\textbf{T10.597}\enspace \texttt{0123323123}\enspace $n{=}10$, $|H|{=}4$, len 4, br 3, $|\mathrm{Aut}|{=}2$, orb 9/8, hist-orb 3/2; M 20480, res 6220, rev 7726: ok}\\*
0 & -- & 0 & 4 & 1 & 658 & 2 & 364 & \texttt{-?????????}\\*
1 & 0 & 1 & 3 & 1 & 480 & 32 & 512 & \texttt{--?????---}\\
2 & 1 & 2 & 2 & 1 & 256 & 256 & 512 & \texttt{---??-----}\\
3 & 2 & 3 & 1 & 2 & 0 & 1024 & 0 & \texttt{----------}\\
4 & 2 & 3 & 1 & 2 & 0 & 1024 & 0 & \texttt{----------}\\
5 & 1 & 2 & 1 & 1 & 512 & 512 & 0 & \texttt{------+---}\\
6 & 5 & 3 & 1 & 1 & 0 & 1024 & 0 & \texttt{----------}\\
7 & 0 & 1 & 1 & 1 & 768 & 256 & 0 & \texttt{--------++}\\
8 & 7 & 2 & 1 & 1 & 512 & 512 & 0 & \texttt{---------+}\\
9 & 8 & 3 & 1 & 1 & 0 & 1024 & 0 & \texttt{----------}\\
\multicolumn{9}{@{}l@{}}{\rule{0pt}{2.6ex}\textbf{T10.598}\enspace \texttt{0123323122}\enspace $n{=}10$, $|H|{=}5$, len 4, br 4, $|\mathrm{Aut}|{=}4$, orb 8/6, hist-orb 3/1; M 19456, res 8680, rev 6006: ok}\\*
0 & -- & 0 & 5 & 1 & 470 & 2 & 552 & \texttt{-?????????}\\*
1 & 0 & 1 & 3 & 1 & 480 & 32 & 512 & \texttt{--?????---}\\
2 & 1 & 2 & 2 & 1 & 256 & 256 & 512 & \texttt{---??-----}\\
3 & 2 & 3 & 1 & 2 & 0 & 1024 & 0 & \texttt{----------}\\
4 & 2 & 3 & 1 & 2 & 0 & 1024 & 0 & \texttt{----------}\\
5 & 1 & 2 & 1 & 1 & 512 & 512 & 0 & \texttt{------+---}\\
6 & 5 & 3 & 1 & 1 & 0 & 1024 & 0 & \texttt{----------}\\
7 & 0 & 1 & 2 & 1 & 256 & 256 & 512 & \texttt{--------??}\\
8 & 7 & 2 & 1 & 2 & 0 & 1024 & 0 & \texttt{----------}\\
9 & 7 & 2 & 1 & 2 & 0 & 1024 & 0 & \texttt{----------}\\
\multicolumn{9}{@{}l@{}}{\rule{0pt}{2.6ex}\textbf{T10.599}\enspace \texttt{0123323121}\enspace $n{=}10$, $|H|{=}5$, len 4, br 3, $|\mathrm{Aut}|{=}2$, orb 9/8, hist-orb 4/3; M 18432, res 9604, rev 4934: ok}\\*
0 & -- & 0 & 5 & 1 & 282 & 2 & 740 & \texttt{-?????????}\\*
1 & 0 & 1 & 3 & 1 & 480 & 32 & 512 & \texttt{--?????---}\\
2 & 1 & 2 & 2 & 1 & 256 & 256 & 512 & \texttt{---??-----}\\
3 & 2 & 3 & 1 & 2 & 0 & 1024 & 0 & \texttt{----------}\\
4 & 2 & 3 & 1 & 2 & 0 & 1024 & 0 & \texttt{----------}\\
5 & 1 & 2 & 1 & 1 & 512 & 512 & 0 & \texttt{------+---}\\
6 & 5 & 3 & 1 & 1 & 0 & 1024 & 0 & \texttt{----------}\\
7 & 0 & 1 & 1 & 1 & 512 & 512 & 0 & \texttt{--------+-}\\
8 & 7 & 2 & 1 & 1 & 0 & 1024 & 0 & \texttt{----------}\\
9 & 0 & 1 & 1 & 1 & 0 & 1024 & 0 & \texttt{----------}\\
\multicolumn{9}{@{}l@{}}{\rule{0pt}{2.6ex}\textbf{T10.600}\enspace \texttt{0123323111}\enspace $n{=}10$, $|H|{=}6$, len 4, br 3, $|\mathrm{Aut}|{=}12$, orb 7/5, hist-orb 3/1; M 17408, res 11296, rev 3350: ok}\\*
0 & -- & 0 & 6 & 1 & 94 & 2 & 928 & \texttt{-?????????}\\*
1 & 0 & 1 & 3 & 1 & 480 & 32 & 512 & \texttt{--?????---}\\
2 & 1 & 2 & 2 & 1 & 256 & 256 & 512 & \texttt{---??-----}\\
3 & 2 & 3 & 1 & 2 & 0 & 1024 & 0 & \texttt{----------}\\
4 & 2 & 3 & 1 & 2 & 0 & 1024 & 0 & \texttt{----------}\\
5 & 1 & 2 & 1 & 1 & 512 & 512 & 0 & \texttt{------+---}\\
6 & 5 & 3 & 1 & 1 & 0 & 1024 & 0 & \texttt{----------}\\
7 & 0 & 1 & 1 & 3 & 0 & 1024 & 0 & \texttt{----------}\\
8 & 0 & 1 & 1 & 3 & 0 & 1024 & 0 & \texttt{----------}\\
9 & 0 & 1 & 1 & 3 & 0 & 1024 & 0 & \texttt{----------}\\
\multicolumn{9}{@{}l@{}}{\rule{0pt}{2.6ex}\textbf{T10.601}\enspace \texttt{0123322222}\enspace $n{=}10$, $|H|{=}7$, len 4, br 2, $|\mathrm{Aut}|{=}240$, orb 5/3, hist-orb 2/0; M 19456, res 12640, rev 4442: ok}\\*
0 & -- & 0 & 7 & 1 & 522 & 2 & 500 & \texttt{-+????????}\\*
1 & 0 & 1 & 7 & 1 & 20 & 4 & 1000 & \texttt{--????????}\\
2 & 1 & 2 & 2 & 1 & 256 & 256 & 512 & \texttt{---??-----}\\
3 & 2 & 3 & 1 & 2 & 0 & 1024 & 0 & \texttt{----------}\\
4 & 2 & 3 & 1 & 2 & 0 & 1024 & 0 & \texttt{----------}\\
5 & 1 & 2 & 1 & 5 & 0 & 1024 & 0 & \texttt{----------}\\
6 & 1 & 2 & 1 & 5 & 0 & 1024 & 0 & \texttt{----------}\\
7 & 1 & 2 & 1 & 5 & 0 & 1024 & 0 & \texttt{----------}\\
8 & 1 & 2 & 1 & 5 & 0 & 1024 & 0 & \texttt{----------}\\
9 & 1 & 2 & 1 & 5 & 0 & 1024 & 0 & \texttt{----------}\\
\multicolumn{9}{@{}l@{}}{\rule{0pt}{2.6ex}\textbf{T10.602}\enspace \texttt{0123322221}\enspace $n{=}10$, $|H|{=}7$, len 4, br 3, $|\mathrm{Aut}|{=}48$, orb 6/4, hist-orb 3/1; M 18432, res 13788, rev 2690: ok}\\*
0 & -- & 0 & 7 & 1 & 266 & 2 & 756 & \texttt{-?????????}\\*
1 & 0 & 1 & 6 & 1 & 40 & 8 & 976 & \texttt{--???????-}\\
2 & 1 & 2 & 2 & 1 & 256 & 256 & 512 & \texttt{---??-----}\\
3 & 2 & 3 & 1 & 2 & 0 & 1024 & 0 & \texttt{----------}\\
4 & 2 & 3 & 1 & 2 & 0 & 1024 & 0 & \texttt{----------}\\
5 & 1 & 2 & 1 & 4 & 0 & 1024 & 0 & \texttt{----------}\\
6 & 1 & 2 & 1 & 4 & 0 & 1024 & 0 & \texttt{----------}\\
7 & 1 & 2 & 1 & 4 & 0 & 1024 & 0 & \texttt{----------}\\
8 & 1 & 2 & 1 & 4 & 0 & 1024 & 0 & \texttt{----------}\\
9 & 0 & 1 & 1 & 1 & 0 & 1024 & 0 & \texttt{----------}\\
\multicolumn{9}{@{}l@{}}{\rule{0pt}{2.6ex}\textbf{T10.603}\enspace \texttt{0123322212}\enspace $n{=}10$, $|H|{=}6$, len 4, br 3, $|\mathrm{Aut}|{=}12$, orb 7/5, hist-orb 3/1; M 18432, res 11216, rev 4170: ok}\\*
0 & -- & 0 & 6 & 1 & 414 & 2 & 608 & \texttt{-?????????}\\*
1 & 0 & 1 & 5 & 1 & 80 & 16 & 928 & \texttt{--??????--}\\
2 & 1 & 2 & 2 & 1 & 256 & 256 & 512 & \texttt{---??-----}\\
3 & 2 & 3 & 1 & 2 & 0 & 1024 & 0 & \texttt{----------}\\
4 & 2 & 3 & 1 & 2 & 0 & 1024 & 0 & \texttt{----------}\\
5 & 1 & 2 & 1 & 3 & 0 & 1024 & 0 & \texttt{----------}\\
6 & 1 & 2 & 1 & 3 & 0 & 1024 & 0 & \texttt{----------}\\
7 & 1 & 2 & 1 & 3 & 0 & 1024 & 0 & \texttt{----------}\\
8 & 0 & 1 & 1 & 1 & 512 & 512 & 0 & \texttt{---------+}\\
9 & 8 & 2 & 1 & 1 & 0 & 1024 & 0 & \texttt{----------}\\
\multicolumn{9}{@{}l@{}}{\rule{0pt}{2.6ex}\textbf{T10.604}\enspace \texttt{0123322211}\enspace $n{=}10$, $|H|{=}7$, len 4, br 3, $|\mathrm{Aut}|{=}24$, orb 6/3, hist-orb 3/0; M 17408, res 13700, rev 1930: ok}\\*
0 & -- & 0 & 7 & 1 & 138 & 2 & 884 & \texttt{-?????????}\\*
1 & 0 & 1 & 5 & 1 & 80 & 16 & 928 & \texttt{--??????--}\\
2 & 1 & 2 & 2 & 1 & 256 & 256 & 512 & \texttt{---??-----}\\
3 & 2 & 3 & 1 & 2 & 0 & 1024 & 0 & \texttt{----------}\\
4 & 2 & 3 & 1 & 2 & 0 & 1024 & 0 & \texttt{----------}\\
5 & 1 & 2 & 1 & 3 & 0 & 1024 & 0 & \texttt{----------}\\
6 & 1 & 2 & 1 & 3 & 0 & 1024 & 0 & \texttt{----------}\\
7 & 1 & 2 & 1 & 3 & 0 & 1024 & 0 & \texttt{----------}\\
8 & 0 & 1 & 1 & 2 & 0 & 1024 & 0 & \texttt{----------}\\
9 & 0 & 1 & 1 & 2 & 0 & 1024 & 0 & \texttt{----------}\\
\multicolumn{9}{@{}l@{}}{\rule{0pt}{2.6ex}\textbf{T10.605}\enspace \texttt{0123322123}\enspace $n{=}10$, $|H|{=}5$, len 4, br 3, $|\mathrm{Aut}|{=}4$, orb 8/6, hist-orb 3/1; M 19456, res 8920, rev 5746: ok}\\*
0 & -- & 0 & 5 & 1 & 518 & 2 & 504 & \texttt{-?????????}\\*
1 & 0 & 1 & 4 & 1 & 160 & 32 & 832 & \texttt{--?????---}\\
2 & 1 & 2 & 2 & 1 & 256 & 256 & 512 & \texttt{---??-----}\\
3 & 2 & 3 & 1 & 2 & 0 & 1024 & 0 & \texttt{----------}\\
4 & 2 & 3 & 1 & 2 & 0 & 1024 & 0 & \texttt{----------}\\
5 & 1 & 2 & 1 & 2 & 0 & 1024 & 0 & \texttt{----------}\\
6 & 1 & 2 & 1 & 2 & 0 & 1024 & 0 & \texttt{----------}\\
7 & 0 & 1 & 1 & 1 & 768 & 256 & 0 & \texttt{--------++}\\
8 & 7 & 2 & 1 & 1 & 512 & 512 & 0 & \texttt{---------+}\\
9 & 8 & 3 & 1 & 1 & 0 & 1024 & 0 & \texttt{----------}\\
\multicolumn{9}{@{}l@{}}{\rule{0pt}{2.6ex}\textbf{T10.606}\enspace \texttt{0123322122}\enspace $n{=}10$, $|H|{=}6$, len 4, br 4, $|\mathrm{Aut}|{=}8$, orb 7/4, hist-orb 3/0; M 18432, res 11020, rev 4090: ok}\\*
0 & -- & 0 & 6 & 1 & 370 & 2 & 652 & \texttt{-?????????}\\*
1 & 0 & 1 & 4 & 1 & 160 & 32 & 832 & \texttt{--?????---}\\
2 & 1 & 2 & 2 & 1 & 256 & 256 & 512 & \texttt{---??-----}\\
3 & 2 & 3 & 1 & 2 & 0 & 1024 & 0 & \texttt{----------}\\
4 & 2 & 3 & 1 & 2 & 0 & 1024 & 0 & \texttt{----------}\\
5 & 1 & 2 & 1 & 2 & 0 & 1024 & 0 & \texttt{----------}\\
6 & 1 & 2 & 1 & 2 & 0 & 1024 & 0 & \texttt{----------}\\
7 & 0 & 1 & 2 & 1 & 256 & 256 & 512 & \texttt{--------??}\\
8 & 7 & 2 & 1 & 2 & 0 & 1024 & 0 & \texttt{----------}\\
9 & 7 & 2 & 1 & 2 & 0 & 1024 & 0 & \texttt{----------}\\
\multicolumn{9}{@{}l@{}}{\rule{0pt}{2.6ex}\textbf{T10.607}\enspace \texttt{0123322121}\enspace $n{=}10$, $|H|{=}6$, len 4, br 3, $|\mathrm{Aut}|{=}4$, orb 8/6, hist-orb 4/2; M 17408, res 11584, rev 3162: ok}\\*
0 & -- & 0 & 6 & 1 & 222 & 2 & 800 & \texttt{-?????????}\\*
1 & 0 & 1 & 4 & 1 & 160 & 32 & 832 & \texttt{--?????---}\\
2 & 1 & 2 & 2 & 1 & 256 & 256 & 512 & \texttt{---??-----}\\
3 & 2 & 3 & 1 & 2 & 0 & 1024 & 0 & \texttt{----------}\\
4 & 2 & 3 & 1 & 2 & 0 & 1024 & 0 & \texttt{----------}\\
5 & 1 & 2 & 1 & 2 & 0 & 1024 & 0 & \texttt{----------}\\
6 & 1 & 2 & 1 & 2 & 0 & 1024 & 0 & \texttt{----------}\\
7 & 0 & 1 & 1 & 1 & 512 & 512 & 0 & \texttt{--------+-}\\
8 & 7 & 2 & 1 & 1 & 0 & 1024 & 0 & \texttt{----------}\\
9 & 0 & 1 & 1 & 1 & 0 & 1024 & 0 & \texttt{----------}\\
\multicolumn{9}{@{}l@{}}{\rule{0pt}{2.6ex}\textbf{T10.608}\enspace \texttt{0123322111}\enspace $n{=}10$, $|H|{=}7$, len 4, br 3, $|\mathrm{Aut}|{=}24$, orb 6/3, hist-orb 3/0; M 16384, res 12916, rev 1722: ok}\\*
0 & -- & 0 & 7 & 1 & 74 & 2 & 948 & \texttt{-?????????}\\*
1 & 0 & 1 & 4 & 1 & 160 & 32 & 832 & \texttt{--?????---}\\
2 & 1 & 2 & 2 & 1 & 256 & 256 & 512 & \texttt{---??-----}\\
3 & 2 & 3 & 1 & 2 & 0 & 1024 & 0 & \texttt{----------}\\
4 & 2 & 3 & 1 & 2 & 0 & 1024 & 0 & \texttt{----------}\\
5 & 1 & 2 & 1 & 2 & 0 & 1024 & 0 & \texttt{----------}\\
6 & 1 & 2 & 1 & 2 & 0 & 1024 & 0 & \texttt{----------}\\
7 & 0 & 1 & 1 & 3 & 0 & 1024 & 0 & \texttt{----------}\\
8 & 0 & 1 & 1 & 3 & 0 & 1024 & 0 & \texttt{----------}\\
9 & 0 & 1 & 1 & 3 & 0 & 1024 & 0 & \texttt{----------}\\
\multicolumn{9}{@{}l@{}}{\rule{0pt}{2.6ex}\textbf{T10.609}\enspace \texttt{0123321233}\enspace $n{=}10$, $|H|{=}5$, len 4, br 4, $|\mathrm{Aut}|{=}4$, orb 8/6, hist-orb 3/1; M 20480, res 8444, rev 6510: ok}\\*
0 & -- & 0 & 5 & 1 & 546 & 2 & 476 & \texttt{-?????????}\\*
1 & 0 & 1 & 3 & 1 & 320 & 64 & 640 & \texttt{--????----}\\
2 & 1 & 2 & 2 & 1 & 256 & 256 & 512 & \texttt{---??-----}\\
3 & 2 & 3 & 1 & 2 & 0 & 1024 & 0 & \texttt{----------}\\
4 & 2 & 3 & 1 & 2 & 0 & 1024 & 0 & \texttt{----------}\\
5 & 1 & 2 & 1 & 1 & 0 & 1024 & 0 & \texttt{----------}\\
6 & 0 & 1 & 2 & 1 & 640 & 128 & 256 & \texttt{-------+??}\\
7 & 6 & 2 & 2 & 1 & 256 & 256 & 512 & \texttt{--------??}\\
8 & 7 & 3 & 1 & 2 & 0 & 1024 & 0 & \texttt{----------}\\
9 & 7 & 3 & 1 & 2 & 0 & 1024 & 0 & \texttt{----------}\\
\multicolumn{9}{@{}l@{}}{\rule{0pt}{2.6ex}\textbf{T10.610}\enspace \texttt{0123321232}\enspace $n{=}10$, $|H|{=}5$, len 4, br 4, $|\mathrm{Aut}|{=}2$, orb 9/8, hist-orb 4/3; M 19456, res 9200, rev 5638: ok}\\*
0 & -- & 0 & 5 & 1 & 462 & 2 & 560 & \texttt{-?????????}\\*
1 & 0 & 1 & 3 & 1 & 320 & 64 & 640 & \texttt{--????----}\\
2 & 1 & 2 & 2 & 1 & 256 & 256 & 512 & \texttt{---??-----}\\
3 & 2 & 3 & 1 & 2 & 0 & 1024 & 0 & \texttt{----------}\\
4 & 2 & 3 & 1 & 2 & 0 & 1024 & 0 & \texttt{----------}\\
5 & 1 & 2 & 1 & 1 & 0 & 1024 & 0 & \texttt{----------}\\
6 & 0 & 1 & 2 & 1 & 384 & 128 & 512 & \texttt{-------???}\\
7 & 6 & 2 & 1 & 1 & 512 & 512 & 0 & \texttt{--------+-}\\
8 & 7 & 3 & 1 & 1 & 0 & 1024 & 0 & \texttt{----------}\\
9 & 6 & 2 & 1 & 1 & 0 & 1024 & 0 & \texttt{----------}\\
\multicolumn{9}{@{}l@{}}{\rule{0pt}{2.6ex}\textbf{T10.611}\enspace \texttt{0123321231}\enspace $n{=}10$, $|H|{=}5$, len 4, br 3, $|\mathrm{Aut}|{=}2$, orb 9/8, hist-orb 4/3; M 18432, res 9432, rev 4794: ok}\\*
0 & -- & 0 & 5 & 1 & 294 & 2 & 728 & \texttt{-?????????}\\*
1 & 0 & 1 & 3 & 1 & 320 & 64 & 640 & \texttt{--????----}\\
2 & 1 & 2 & 2 & 1 & 256 & 256 & 512 & \texttt{---??-----}\\
3 & 2 & 3 & 1 & 2 & 0 & 1024 & 0 & \texttt{----------}\\
4 & 2 & 3 & 1 & 2 & 0 & 1024 & 0 & \texttt{----------}\\
5 & 1 & 2 & 1 & 1 & 0 & 1024 & 0 & \texttt{----------}\\
6 & 0 & 1 & 1 & 1 & 768 & 256 & 0 & \texttt{-------++-}\\
7 & 6 & 2 & 1 & 1 & 512 & 512 & 0 & \texttt{--------+-}\\
8 & 7 & 3 & 1 & 1 & 0 & 1024 & 0 & \texttt{----------}\\
9 & 0 & 1 & 1 & 1 & 0 & 1024 & 0 & \texttt{----------}\\
\multicolumn{9}{@{}l@{}}{\rule{0pt}{2.6ex}\textbf{T10.612}\enspace \texttt{0123321222}\enspace $n{=}10$, $|H|{=}6$, len 4, br 4, $|\mathrm{Aut}|{=}12$, orb 7/5, hist-orb 3/1; M 18432, res 10596, rev 4338: ok}\\*
0 & -- & 0 & 6 & 1 & 378 & 2 & 644 & \texttt{-?????????}\\*
1 & 0 & 1 & 3 & 1 & 320 & 64 & 640 & \texttt{--????----}\\
2 & 1 & 2 & 2 & 1 & 256 & 256 & 512 & \texttt{---??-----}\\
3 & 2 & 3 & 1 & 2 & 0 & 1024 & 0 & \texttt{----------}\\
4 & 2 & 3 & 1 & 2 & 0 & 1024 & 0 & \texttt{----------}\\
5 & 1 & 2 & 1 & 1 & 0 & 1024 & 0 & \texttt{----------}\\
6 & 0 & 1 & 3 & 1 & 128 & 128 & 768 & \texttt{-------???}\\
7 & 6 & 2 & 1 & 3 & 0 & 1024 & 0 & \texttt{----------}\\
8 & 6 & 2 & 1 & 3 & 0 & 1024 & 0 & \texttt{----------}\\
9 & 6 & 2 & 1 & 3 & 0 & 1024 & 0 & \texttt{----------}\\
\multicolumn{9}{@{}l@{}}{\rule{0pt}{2.6ex}\textbf{T10.613}\enspace \texttt{0123321221}\enspace $n{=}10$, $|H|{=}6$, len 4, br 4, $|\mathrm{Aut}|{=}4$, orb 8/6, hist-orb 4/2; M 17408, res 10956, rev 3410: ok}\\*
0 & -- & 0 & 6 & 1 & 210 & 2 & 812 & \texttt{-?????????}\\*
1 & 0 & 1 & 3 & 1 & 320 & 64 & 640 & \texttt{--????----}\\
2 & 1 & 2 & 2 & 1 & 256 & 256 & 512 & \texttt{---??-----}\\
3 & 2 & 3 & 1 & 2 & 0 & 1024 & 0 & \texttt{----------}\\
4 & 2 & 3 & 1 & 2 & 0 & 1024 & 0 & \texttt{----------}\\
5 & 1 & 2 & 1 & 1 & 0 & 1024 & 0 & \texttt{----------}\\
6 & 0 & 1 & 2 & 1 & 256 & 256 & 512 & \texttt{-------??-}\\
7 & 6 & 2 & 1 & 2 & 0 & 1024 & 0 & \texttt{----------}\\
8 & 6 & 2 & 1 & 2 & 0 & 1024 & 0 & \texttt{----------}\\
9 & 0 & 1 & 1 & 1 & 0 & 1024 & 0 & \texttt{----------}\\
\multicolumn{9}{@{}l@{}}{\rule{0pt}{2.6ex}\textbf{T10.614}\enspace \texttt{0123321212}\enspace $n{=}10$, $|H|{=}5$, len 4, br 3, $|\mathrm{Aut}|{=}4$, orb 7/4, hist-orb 3/1; M 17408, res 8676, rev 4810: ok}\\*
0 & -- & 0 & 5 & 1 & 378 & 2 & 644 & \texttt{-?????????}\\*
1 & 0 & 1 & 3 & 1 & 320 & 64 & 640 & \texttt{--????----}\\
2 & 1 & 2 & 2 & 1 & 256 & 256 & 512 & \texttt{---??-----}\\
3 & 2 & 3 & 1 & 2 & 0 & 1024 & 0 & \texttt{----------}\\
4 & 2 & 3 & 1 & 2 & 0 & 1024 & 0 & \texttt{----------}\\
5 & 1 & 2 & 1 & 1 & 0 & 1024 & 0 & \texttt{----------}\\
6 & 0 & 1 & 1 & 2 & 512 & 512 & 0 & \texttt{-------+--}\\
7 & 6 & 2 & 1 & 2 & 0 & 1024 & 0 & \texttt{----------}\\
8 & 0 & 1 & 1 & 2 & 512 & 512 & 0 & \texttt{---------+}\\
9 & 8 & 2 & 1 & 2 & 0 & 1024 & 0 & \texttt{----------}\\
\multicolumn{9}{@{}l@{}}{\rule{0pt}{2.6ex}\textbf{T10.615}\enspace \texttt{0123321211}\enspace $n{=}10$, $|H|{=}6$, len 4, br 3, $|\mathrm{Aut}|{=}4$, orb 8/6, hist-orb 4/2; M 16384, res 10944, rev 2882: ok}\\*
0 & -- & 0 & 6 & 1 & 126 & 2 & 896 & \texttt{-?????????}\\*
1 & 0 & 1 & 3 & 1 & 320 & 64 & 640 & \texttt{--????----}\\
2 & 1 & 2 & 2 & 1 & 256 & 256 & 512 & \texttt{---??-----}\\
3 & 2 & 3 & 1 & 2 & 0 & 1024 & 0 & \texttt{----------}\\
4 & 2 & 3 & 1 & 2 & 0 & 1024 & 0 & \texttt{----------}\\
5 & 1 & 2 & 1 & 1 & 0 & 1024 & 0 & \texttt{----------}\\
6 & 0 & 1 & 1 & 1 & 512 & 512 & 0 & \texttt{-------+--}\\
7 & 6 & 2 & 1 & 1 & 0 & 1024 & 0 & \texttt{----------}\\
8 & 0 & 1 & 1 & 2 & 0 & 1024 & 0 & \texttt{----------}\\
9 & 0 & 1 & 1 & 2 & 0 & 1024 & 0 & \texttt{----------}\\
\multicolumn{9}{@{}l@{}}{\rule{0pt}{2.6ex}\textbf{T10.616}\enspace \texttt{0123321111}\enspace $n{=}10$, $|H|{=}7$, len 4, br 3, $|\mathrm{Aut}|{=}48$, orb 6/4, hist-orb 3/1; M 15360, res 11700, rev 1842: ok}\\*
0 & -- & 0 & 7 & 1 & 42 & 2 & 980 & \texttt{-?????????}\\*
1 & 0 & 1 & 3 & 1 & 320 & 64 & 640 & \texttt{--????----}\\
2 & 1 & 2 & 2 & 1 & 256 & 256 & 512 & \texttt{---??-----}\\
3 & 2 & 3 & 1 & 2 & 0 & 1024 & 0 & \texttt{----------}\\
4 & 2 & 3 & 1 & 2 & 0 & 1024 & 0 & \texttt{----------}\\
5 & 1 & 2 & 1 & 1 & 0 & 1024 & 0 & \texttt{----------}\\
6 & 0 & 1 & 1 & 4 & 0 & 1024 & 0 & \texttt{----------}\\
7 & 0 & 1 & 1 & 4 & 0 & 1024 & 0 & \texttt{----------}\\
8 & 0 & 1 & 1 & 4 & 0 & 1024 & 0 & \texttt{----------}\\
9 & 0 & 1 & 1 & 4 & 0 & 1024 & 0 & \texttt{----------}\\
\multicolumn{9}{@{}l@{}}{\rule{0pt}{2.6ex}\textbf{T10.617}\enspace \texttt{0123312331}\enspace $n{=}10$, $|H|{=}5$, len 4, br 3, $|\mathrm{Aut}|{=}8$, orb 5/2, hist-orb 2/1; M 19456, res 8716, rev 5738: ok}\\*
0 & -- & 0 & 5 & 1 & 338 & 2 & 684 & \texttt{-?????????}\\*
1 & 0 & 1 & 2 & 2 & 640 & 128 & 256 & \texttt{--+??-----}\\
2 & 1 & 2 & 2 & 2 & 256 & 256 & 512 & \texttt{---??-----}\\
3 & 2 & 3 & 1 & 4 & 0 & 1024 & 0 & \texttt{----------}\\
4 & 2 & 3 & 1 & 4 & 0 & 1024 & 0 & \texttt{----------}\\
5 & 0 & 1 & 2 & 2 & 640 & 128 & 256 & \texttt{------+??-}\\
6 & 5 & 2 & 2 & 2 & 256 & 256 & 512 & \texttt{-------??-}\\
7 & 6 & 3 & 1 & 4 & 0 & 1024 & 0 & \texttt{----------}\\
8 & 6 & 3 & 1 & 4 & 0 & 1024 & 0 & \texttt{----------}\\
9 & 0 & 1 & 1 & 1 & 0 & 1024 & 0 & \texttt{----------}\\
\multicolumn{9}{@{}l@{}}{\rule{0pt}{2.6ex}\textbf{T10.618}\enspace \texttt{0123312323}\enspace $n{=}10$, $|H|{=}4$, len 4, br 3, $|\mathrm{Aut}|{=}4$, orb 7/4, hist-orb 2/0; M 20480, res 5972, rev 7878: ok}\\*
0 & -- & 0 & 4 & 1 & 650 & 2 & 372 & \texttt{-?????????}\\*
1 & 0 & 1 & 2 & 1 & 640 & 128 & 256 & \texttt{--+??-----}\\
2 & 1 & 2 & 2 & 1 & 256 & 256 & 512 & \texttt{---??-----}\\
3 & 2 & 3 & 1 & 2 & 0 & 1024 & 0 & \texttt{----------}\\
4 & 2 & 3 & 1 & 2 & 0 & 1024 & 0 & \texttt{----------}\\
5 & 0 & 1 & 2 & 1 & 576 & 64 & 384 & \texttt{------????}\\
6 & 5 & 2 & 1 & 2 & 512 & 512 & 0 & \texttt{-------+--}\\
7 & 6 & 3 & 1 & 2 & 0 & 1024 & 0 & \texttt{----------}\\
8 & 5 & 2 & 1 & 2 & 512 & 512 & 0 & \texttt{---------+}\\
9 & 8 & 3 & 1 & 2 & 0 & 1024 & 0 & \texttt{----------}\\
\multicolumn{9}{@{}l@{}}{\rule{0pt}{2.6ex}\textbf{T10.619}\enspace \texttt{0123312322}\enspace $n{=}10$, $|H|{=}5$, len 4, br 3, $|\mathrm{Aut}|{=}4$, orb 8/6, hist-orb 3/1; M 19456, res 8720, rev 5954: ok}\\*
0 & -- & 0 & 5 & 1 & 494 & 2 & 528 & \texttt{-?????????}\\*
1 & 0 & 1 & 2 & 1 & 640 & 128 & 256 & \texttt{--+??-----}\\
2 & 1 & 2 & 2 & 1 & 256 & 256 & 512 & \texttt{---??-----}\\
3 & 2 & 3 & 1 & 2 & 0 & 1024 & 0 & \texttt{----------}\\
4 & 2 & 3 & 1 & 2 & 0 & 1024 & 0 & \texttt{----------}\\
5 & 0 & 1 & 3 & 1 & 192 & 64 & 768 & \texttt{------????}\\
6 & 5 & 2 & 1 & 1 & 512 & 512 & 0 & \texttt{-------+--}\\
7 & 6 & 3 & 1 & 1 & 0 & 1024 & 0 & \texttt{----------}\\
8 & 5 & 2 & 1 & 2 & 0 & 1024 & 0 & \texttt{----------}\\
9 & 5 & 2 & 1 & 2 & 0 & 1024 & 0 & \texttt{----------}\\
\multicolumn{9}{@{}l@{}}{\rule{0pt}{2.6ex}\textbf{T10.620}\enspace \texttt{0123312321}\enspace $n{=}10$, $|H|{=}5$, len 4, br 3, $|\mathrm{Aut}|{=}2$, orb 9/8, hist-orb 4/3; M 18432, res 9184, rev 5002: ok}\\*
0 & -- & 0 & 5 & 1 & 286 & 2 & 736 & \texttt{-?????????}\\*
1 & 0 & 1 & 2 & 1 & 640 & 128 & 256 & \texttt{--+??-----}\\
2 & 1 & 2 & 2 & 1 & 256 & 256 & 512 & \texttt{---??-----}\\
3 & 2 & 3 & 1 & 2 & 0 & 1024 & 0 & \texttt{----------}\\
4 & 2 & 3 & 1 & 2 & 0 & 1024 & 0 & \texttt{----------}\\
5 & 0 & 1 & 2 & 1 & 384 & 128 & 512 & \texttt{------???-}\\
6 & 5 & 2 & 1 & 1 & 512 & 512 & 0 & \texttt{-------+--}\\
7 & 6 & 3 & 1 & 1 & 0 & 1024 & 0 & \texttt{----------}\\
8 & 5 & 2 & 1 & 1 & 0 & 1024 & 0 & \texttt{----------}\\
9 & 0 & 1 & 1 & 1 & 0 & 1024 & 0 & \texttt{----------}\\
\multicolumn{9}{@{}l@{}}{\rule{0pt}{2.6ex}\textbf{T10.621}\enspace \texttt{0123312312}\enspace $n{=}10$, $|H|{=}4$, len 4, br 2, $|\mathrm{Aut}|{=}2$, orb 9/8, hist-orb 3/2; M 18432, res 5564, rev 6814: ok}\\*
0 & -- & 0 & 4 & 1 & 546 & 2 & 476 & \texttt{-?????????}\\*
1 & 0 & 1 & 2 & 1 & 640 & 128 & 256 & \texttt{--+??-----}\\
2 & 1 & 2 & 2 & 1 & 256 & 256 & 512 & \texttt{---??-----}\\
3 & 2 & 3 & 1 & 2 & 0 & 1024 & 0 & \texttt{----------}\\
4 & 2 & 3 & 1 & 2 & 0 & 1024 & 0 & \texttt{----------}\\
5 & 0 & 1 & 1 & 1 & 768 & 256 & 0 & \texttt{------++--}\\
6 & 5 & 2 & 1 & 1 & 512 & 512 & 0 & \texttt{-------+--}\\
7 & 6 & 3 & 1 & 1 & 0 & 1024 & 0 & \texttt{----------}\\
8 & 0 & 1 & 1 & 1 & 512 & 512 & 0 & \texttt{---------+}\\
9 & 8 & 2 & 1 & 1 & 0 & 1024 & 0 & \texttt{----------}\\
\multicolumn{9}{@{}l@{}}{\rule{0pt}{2.6ex}\textbf{T10.622}\enspace \texttt{0123312311}\enspace $n{=}10$, $|H|{=}5$, len 4, br 2, $|\mathrm{Aut}|{=}4$, orb 8/6, hist-orb 3/1; M 17408, res 8840, rev 4526: ok}\\*
0 & -- & 0 & 5 & 1 & 182 & 2 & 840 & \texttt{-?????????}\\*
1 & 0 & 1 & 2 & 1 & 640 & 128 & 256 & \texttt{--+??-----}\\
2 & 1 & 2 & 2 & 1 & 256 & 256 & 512 & \texttt{---??-----}\\
3 & 2 & 3 & 1 & 2 & 0 & 1024 & 0 & \texttt{----------}\\
4 & 2 & 3 & 1 & 2 & 0 & 1024 & 0 & \texttt{----------}\\
5 & 0 & 1 & 1 & 1 & 768 & 256 & 0 & \texttt{------++--}\\
6 & 5 & 2 & 1 & 1 & 512 & 512 & 0 & \texttt{-------+--}\\
7 & 6 & 3 & 1 & 1 & 0 & 1024 & 0 & \texttt{----------}\\
8 & 0 & 1 & 1 & 2 & 0 & 1024 & 0 & \texttt{----------}\\
9 & 0 & 1 & 1 & 2 & 0 & 1024 & 0 & \texttt{----------}\\
\multicolumn{9}{@{}l@{}}{\rule{0pt}{2.6ex}\textbf{T10.623}\enspace \texttt{0123312222}\enspace $n{=}10$, $|H|{=}6$, len 4, br 3, $|\mathrm{Aut}|{=}48$, orb 6/4, hist-orb 2/0; M 18432, res 9636, rev 4990: ok}\\*
0 & -- & 0 & 6 & 1 & 442 & 2 & 580 & \texttt{-?????????}\\*
1 & 0 & 1 & 2 & 1 & 640 & 128 & 256 & \texttt{--+??-----}\\
2 & 1 & 2 & 2 & 1 & 256 & 256 & 512 & \texttt{---??-----}\\
3 & 2 & 3 & 1 & 2 & 0 & 1024 & 0 & \texttt{----------}\\
4 & 2 & 3 & 1 & 2 & 0 & 1024 & 0 & \texttt{----------}\\
5 & 0 & 1 & 4 & 1 & 64 & 64 & 896 & \texttt{------????}\\
6 & 5 & 2 & 1 & 4 & 0 & 1024 & 0 & \texttt{----------}\\
7 & 5 & 2 & 1 & 4 & 0 & 1024 & 0 & \texttt{----------}\\
8 & 5 & 2 & 1 & 4 & 0 & 1024 & 0 & \texttt{----------}\\
9 & 5 & 2 & 1 & 4 & 0 & 1024 & 0 & \texttt{----------}\\
\multicolumn{9}{@{}l@{}}{\rule{0pt}{2.6ex}\textbf{T10.624}\enspace \texttt{0123312221}\enspace $n{=}10$, $|H|{=}6$, len 4, br 3, $|\mathrm{Aut}|{=}12$, orb 7/5, hist-orb 3/1; M 17408, res 10292, rev 3806: ok}\\*
0 & -- & 0 & 6 & 1 & 234 & 2 & 788 & \texttt{-?????????}\\*
1 & 0 & 1 & 2 & 1 & 640 & 128 & 256 & \texttt{--+??-----}\\
2 & 1 & 2 & 2 & 1 & 256 & 256 & 512 & \texttt{---??-----}\\
3 & 2 & 3 & 1 & 2 & 0 & 1024 & 0 & \texttt{----------}\\
4 & 2 & 3 & 1 & 2 & 0 & 1024 & 0 & \texttt{----------}\\
5 & 0 & 1 & 3 & 1 & 128 & 128 & 768 & \texttt{------???-}\\
6 & 5 & 2 & 1 & 3 & 0 & 1024 & 0 & \texttt{----------}\\
7 & 5 & 2 & 1 & 3 & 0 & 1024 & 0 & \texttt{----------}\\
8 & 5 & 2 & 1 & 3 & 0 & 1024 & 0 & \texttt{----------}\\
9 & 0 & 1 & 1 & 1 & 0 & 1024 & 0 & \texttt{----------}\\
\multicolumn{9}{@{}l@{}}{\rule{0pt}{2.6ex}\textbf{T10.625}\enspace \texttt{0123312212}\enspace $n{=}10$, $|H|{=}5$, len 4, br 3, $|\mathrm{Aut}|{=}4$, orb 8/6, hist-orb 3/1; M 17408, res 7736, rev 5222: ok}\\*
0 & -- & 0 & 5 & 1 & 390 & 2 & 632 & \texttt{-?????????}\\*
1 & 0 & 1 & 2 & 1 & 640 & 128 & 256 & \texttt{--+??-----}\\
2 & 1 & 2 & 2 & 1 & 256 & 256 & 512 & \texttt{---??-----}\\
3 & 2 & 3 & 1 & 2 & 0 & 1024 & 0 & \texttt{----------}\\
4 & 2 & 3 & 1 & 2 & 0 & 1024 & 0 & \texttt{----------}\\
5 & 0 & 1 & 2 & 1 & 256 & 256 & 512 & \texttt{------??--}\\
6 & 5 & 2 & 1 & 2 & 0 & 1024 & 0 & \texttt{----------}\\
7 & 5 & 2 & 1 & 2 & 0 & 1024 & 0 & \texttt{----------}\\
8 & 0 & 1 & 1 & 1 & 512 & 512 & 0 & \texttt{---------+}\\
9 & 8 & 2 & 1 & 1 & 0 & 1024 & 0 & \texttt{----------}\\
\multicolumn{9}{@{}l@{}}{\rule{0pt}{2.6ex}\textbf{T10.626}\enspace \texttt{0123312211}\enspace $n{=}10$, $|H|{=}6$, len 4, br 3, $|\mathrm{Aut}|{=}8$, orb 7/4, hist-orb 3/0; M 16384, res 10076, rev 3278: ok}\\*
0 & -- & 0 & 6 & 1 & 130 & 2 & 892 & \texttt{-?????????}\\*
1 & 0 & 1 & 2 & 1 & 640 & 128 & 256 & \texttt{--+??-----}\\
2 & 1 & 2 & 2 & 1 & 256 & 256 & 512 & \texttt{---??-----}\\
3 & 2 & 3 & 1 & 2 & 0 & 1024 & 0 & \texttt{----------}\\
4 & 2 & 3 & 1 & 2 & 0 & 1024 & 0 & \texttt{----------}\\
5 & 0 & 1 & 2 & 1 & 256 & 256 & 512 & \texttt{------??--}\\
6 & 5 & 2 & 1 & 2 & 0 & 1024 & 0 & \texttt{----------}\\
7 & 5 & 2 & 1 & 2 & 0 & 1024 & 0 & \texttt{----------}\\
8 & 0 & 1 & 1 & 2 & 0 & 1024 & 0 & \texttt{----------}\\
9 & 0 & 1 & 1 & 2 & 0 & 1024 & 0 & \texttt{----------}\\
\multicolumn{9}{@{}l@{}}{\rule{0pt}{2.6ex}\textbf{T10.627}\enspace \texttt{0123312121}\enspace $n{=}10$, $|H|{=}5$, len 4, br 2, $|\mathrm{Aut}|{=}4$, orb 7/4, hist-orb 3/1; M 16384, res 8372, rev 4342: ok}\\*
0 & -- & 0 & 5 & 1 & 234 & 2 & 788 & \texttt{-?????????}\\*
1 & 0 & 1 & 2 & 1 & 640 & 128 & 256 & \texttt{--+??-----}\\
2 & 1 & 2 & 2 & 1 & 256 & 256 & 512 & \texttt{---??-----}\\
3 & 2 & 3 & 1 & 2 & 0 & 1024 & 0 & \texttt{----------}\\
4 & 2 & 3 & 1 & 2 & 0 & 1024 & 0 & \texttt{----------}\\
5 & 0 & 1 & 1 & 2 & 512 & 512 & 0 & \texttt{------+---}\\
6 & 5 & 2 & 1 & 2 & 0 & 1024 & 0 & \texttt{----------}\\
7 & 0 & 1 & 1 & 2 & 512 & 512 & 0 & \texttt{--------+-}\\
8 & 7 & 2 & 1 & 2 & 0 & 1024 & 0 & \texttt{----------}\\
9 & 0 & 1 & 1 & 1 & 0 & 1024 & 0 & \texttt{----------}\\
\multicolumn{9}{@{}l@{}}{\rule{0pt}{2.6ex}\textbf{T10.628}\enspace \texttt{0123312111}\enspace $n{=}10$, $|H|{=}6$, len 4, br 2, $|\mathrm{Aut}|{=}12$, orb 7/5, hist-orb 3/1; M 15360, res 9776, rev 2950: ok}\\*
0 & -- & 0 & 6 & 1 & 78 & 2 & 944 & \texttt{-?????????}\\*
1 & 0 & 1 & 2 & 1 & 640 & 128 & 256 & \texttt{--+??-----}\\
2 & 1 & 2 & 2 & 1 & 256 & 256 & 512 & \texttt{---??-----}\\
3 & 2 & 3 & 1 & 2 & 0 & 1024 & 0 & \texttt{----------}\\
4 & 2 & 3 & 1 & 2 & 0 & 1024 & 0 & \texttt{----------}\\
5 & 0 & 1 & 1 & 1 & 512 & 512 & 0 & \texttt{------+---}\\
6 & 5 & 2 & 1 & 1 & 0 & 1024 & 0 & \texttt{----------}\\
7 & 0 & 1 & 1 & 3 & 0 & 1024 & 0 & \texttt{----------}\\
8 & 0 & 1 & 1 & 3 & 0 & 1024 & 0 & \texttt{----------}\\
9 & 0 & 1 & 1 & 3 & 0 & 1024 & 0 & \texttt{----------}\\
\multicolumn{9}{@{}l@{}}{\rule{0pt}{2.6ex}\textbf{T10.629}\enspace \texttt{0123311111}\enspace $n{=}10$, $|H|{=}7$, len 4, br 2, $|\mathrm{Aut}|{=}240$, orb 5/3, hist-orb 2/0; M 14336, res 10244, rev 2110: ok}\\*
0 & -- & 0 & 7 & 1 & 26 & 2 & 996 & \texttt{-?????????}\\*
1 & 0 & 1 & 2 & 1 & 640 & 128 & 256 & \texttt{--+??-----}\\
2 & 1 & 2 & 2 & 1 & 256 & 256 & 512 & \texttt{---??-----}\\
3 & 2 & 3 & 1 & 2 & 0 & 1024 & 0 & \texttt{----------}\\
4 & 2 & 3 & 1 & 2 & 0 & 1024 & 0 & \texttt{----------}\\
5 & 0 & 1 & 1 & 5 & 0 & 1024 & 0 & \texttt{----------}\\
6 & 0 & 1 & 1 & 5 & 0 & 1024 & 0 & \texttt{----------}\\
7 & 0 & 1 & 1 & 5 & 0 & 1024 & 0 & \texttt{----------}\\
8 & 0 & 1 & 1 & 5 & 0 & 1024 & 0 & \texttt{----------}\\
9 & 0 & 1 & 1 & 5 & 0 & 1024 & 0 & \texttt{----------}\\
\multicolumn{9}{@{}l@{}}{\rule{0pt}{2.6ex}\textbf{T10.630}\enspace \texttt{0123232323}\enspace $n{=}10$, $|H|{=}4$, len 4, br 1, $|\mathrm{Aut}|{=}24$, orb 4/2, hist-orb 1/0; M 21504, res 8352, rev 7714: ok}\\*
0 & -- & 0 & 4 & 1 & 674 & 2 & 348 & \texttt{-+????????}\\*
1 & 0 & 1 & 4 & 1 & 324 & 4 & 696 & \texttt{--????????}\\
2 & 1 & 2 & 1 & 4 & 512 & 512 & 0 & \texttt{---+------}\\
3 & 2 & 3 & 1 & 4 & 0 & 1024 & 0 & \texttt{----------}\\
4 & 1 & 2 & 1 & 4 & 512 & 512 & 0 & \texttt{-----+----}\\
5 & 4 & 3 & 1 & 4 & 0 & 1024 & 0 & \texttt{----------}\\
6 & 1 & 2 & 1 & 4 & 512 & 512 & 0 & \texttt{-------+--}\\
7 & 6 & 3 & 1 & 4 & 0 & 1024 & 0 & \texttt{----------}\\
8 & 1 & 2 & 1 & 4 & 512 & 512 & 0 & \texttt{---------+}\\
9 & 8 & 3 & 1 & 4 & 0 & 1024 & 0 & \texttt{----------}\\
\multicolumn{9}{@{}l@{}}{\rule{0pt}{2.6ex}\textbf{T10.631}\enspace \texttt{0123232322}\enspace $n{=}10$, $|H|{=}5$, len 4, br 1, $|\mathrm{Aut}|{=}12$, orb 5/2, hist-orb 2/0; M 20480, res 10944, rev 5838: ok}\\*
0 & -- & 0 & 5 & 1 & 566 & 2 & 456 & \texttt{-+????????}\\*
1 & 0 & 1 & 5 & 1 & 108 & 4 & 912 & \texttt{--????????}\\
2 & 1 & 2 & 1 & 3 & 512 & 512 & 0 & \texttt{---+------}\\
3 & 2 & 3 & 1 & 3 & 0 & 1024 & 0 & \texttt{----------}\\
4 & 1 & 2 & 1 & 3 & 512 & 512 & 0 & \texttt{-----+----}\\
5 & 4 & 3 & 1 & 3 & 0 & 1024 & 0 & \texttt{----------}\\
6 & 1 & 2 & 1 & 3 & 512 & 512 & 0 & \texttt{-------+--}\\
7 & 6 & 3 & 1 & 3 & 0 & 1024 & 0 & \texttt{----------}\\
8 & 1 & 2 & 1 & 2 & 0 & 1024 & 0 & \texttt{----------}\\
9 & 1 & 2 & 1 & 2 & 0 & 1024 & 0 & \texttt{----------}\\
\multicolumn{9}{@{}l@{}}{\rule{0pt}{2.6ex}\textbf{T10.632}\enspace \texttt{0123232321}\enspace $n{=}10$, $|H|{=}5$, len 4, br 2, $|\mathrm{Aut}|{=}6$, orb 6/4, hist-orb 3/2; M 19456, res 11608, rev 4606: ok}\\*
0 & -- & 0 & 5 & 1 & 310 & 2 & 712 & \texttt{-?????????}\\*
1 & 0 & 1 & 4 & 1 & 216 & 8 & 800 & \texttt{--???????-}\\
2 & 1 & 2 & 1 & 3 & 512 & 512 & 0 & \texttt{---+------}\\
3 & 2 & 3 & 1 & 3 & 0 & 1024 & 0 & \texttt{----------}\\
4 & 1 & 2 & 1 & 3 & 512 & 512 & 0 & \texttt{-----+----}\\
5 & 4 & 3 & 1 & 3 & 0 & 1024 & 0 & \texttt{----------}\\
6 & 1 & 2 & 1 & 3 & 512 & 512 & 0 & \texttt{-------+--}\\
7 & 6 & 3 & 1 & 3 & 0 & 1024 & 0 & \texttt{----------}\\
8 & 1 & 2 & 1 & 1 & 0 & 1024 & 0 & \texttt{----------}\\
9 & 0 & 1 & 1 & 1 & 0 & 1024 & 0 & \texttt{----------}\\
\multicolumn{9}{@{}l@{}}{\rule{0pt}{2.6ex}\textbf{T10.633}\enspace \texttt{0123232312}\enspace $n{=}10$, $|H|{=}4$, len 4, br 2, $|\mathrm{Aut}|{=}6$, orb 6/4, hist-orb 2/1; M 19456, res 7452, rev 6886: ok}\\*
0 & -- & 0 & 4 & 1 & 546 & 2 & 476 & \texttt{-?????????}\\*
1 & 0 & 1 & 3 & 1 & 432 & 16 & 576 & \texttt{--??????--}\\
2 & 1 & 2 & 1 & 3 & 512 & 512 & 0 & \texttt{---+------}\\
3 & 2 & 3 & 1 & 3 & 0 & 1024 & 0 & \texttt{----------}\\
4 & 1 & 2 & 1 & 3 & 512 & 512 & 0 & \texttt{-----+----}\\
5 & 4 & 3 & 1 & 3 & 0 & 1024 & 0 & \texttt{----------}\\
6 & 1 & 2 & 1 & 3 & 512 & 512 & 0 & \texttt{-------+--}\\
7 & 6 & 3 & 1 & 3 & 0 & 1024 & 0 & \texttt{----------}\\
8 & 0 & 1 & 1 & 1 & 512 & 512 & 0 & \texttt{---------+}\\
9 & 8 & 2 & 1 & 1 & 0 & 1024 & 0 & \texttt{----------}\\
\multicolumn{9}{@{}l@{}}{\rule{0pt}{2.6ex}\textbf{T10.634}\enspace \texttt{0123232311}\enspace $n{=}10$, $|H|{=}5$, len 4, br 2, $|\mathrm{Aut}|{=}12$, orb 5/2, hist-orb 2/0; M 18432, res 10728, rev 4422: ok}\\*
0 & -- & 0 & 5 & 1 & 182 & 2 & 840 & \texttt{-?????????}\\*
1 & 0 & 1 & 3 & 1 & 432 & 16 & 576 & \texttt{--??????--}\\
2 & 1 & 2 & 1 & 3 & 512 & 512 & 0 & \texttt{---+------}\\
3 & 2 & 3 & 1 & 3 & 0 & 1024 & 0 & \texttt{----------}\\
4 & 1 & 2 & 1 & 3 & 512 & 512 & 0 & \texttt{-----+----}\\
5 & 4 & 3 & 1 & 3 & 0 & 1024 & 0 & \texttt{----------}\\
6 & 1 & 2 & 1 & 3 & 512 & 512 & 0 & \texttt{-------+--}\\
7 & 6 & 3 & 1 & 3 & 0 & 1024 & 0 & \texttt{----------}\\
8 & 0 & 1 & 1 & 2 & 0 & 1024 & 0 & \texttt{----------}\\
9 & 0 & 1 & 1 & 2 & 0 & 1024 & 0 & \texttt{----------}\\
\multicolumn{9}{@{}l@{}}{\rule{0pt}{2.6ex}\textbf{T10.635}\enspace \texttt{0123232222}\enspace $n{=}10$, $|H|{=}6$, len 4, br 1, $|\mathrm{Aut}|{=}48$, orb 5/2, hist-orb 2/0; M 19456, res 11808, rev 4970: ok}\\*
0 & -- & 0 & 6 & 1 & 530 & 2 & 492 & \texttt{-+????????}\\*
1 & 0 & 1 & 6 & 1 & 36 & 4 & 984 & \texttt{--????????}\\
2 & 1 & 2 & 1 & 2 & 512 & 512 & 0 & \texttt{---+------}\\
3 & 2 & 3 & 1 & 2 & 0 & 1024 & 0 & \texttt{----------}\\
4 & 1 & 2 & 1 & 2 & 512 & 512 & 0 & \texttt{-----+----}\\
5 & 4 & 3 & 1 & 2 & 0 & 1024 & 0 & \texttt{----------}\\
6 & 1 & 2 & 1 & 4 & 0 & 1024 & 0 & \texttt{----------}\\
7 & 1 & 2 & 1 & 4 & 0 & 1024 & 0 & \texttt{----------}\\
8 & 1 & 2 & 1 & 4 & 0 & 1024 & 0 & \texttt{----------}\\
9 & 1 & 2 & 1 & 4 & 0 & 1024 & 0 & \texttt{----------}\\
\multicolumn{9}{@{}l@{}}{\rule{0pt}{2.6ex}\textbf{T10.636}\enspace \texttt{0123232221}\enspace $n{=}10$, $|H|{=}6$, len 4, br 2, $|\mathrm{Aut}|{=}12$, orb 6/3, hist-orb 3/1; M 18432, res 12868, rev 3346: ok}\\*
0 & -- & 0 & 6 & 1 & 274 & 2 & 748 & \texttt{-?????????}\\*
1 & 0 & 1 & 5 & 1 & 72 & 8 & 944 & \texttt{--???????-}\\
2 & 1 & 2 & 1 & 2 & 512 & 512 & 0 & \texttt{---+------}\\
3 & 2 & 3 & 1 & 2 & 0 & 1024 & 0 & \texttt{----------}\\
4 & 1 & 2 & 1 & 2 & 512 & 512 & 0 & \texttt{-----+----}\\
5 & 4 & 3 & 1 & 2 & 0 & 1024 & 0 & \texttt{----------}\\
6 & 1 & 2 & 1 & 3 & 0 & 1024 & 0 & \texttt{----------}\\
7 & 1 & 2 & 1 & 3 & 0 & 1024 & 0 & \texttt{----------}\\
8 & 1 & 2 & 1 & 3 & 0 & 1024 & 0 & \texttt{----------}\\
9 & 0 & 1 & 1 & 1 & 0 & 1024 & 0 & \texttt{----------}\\
\multicolumn{9}{@{}l@{}}{\rule{0pt}{2.6ex}\textbf{T10.637}\enspace \texttt{0123232212}\enspace $n{=}10$, $|H|{=}5$, len 4, br 2, $|\mathrm{Aut}|{=}4$, orb 7/4, hist-orb 3/1; M 18432, res 10008, rev 4970: ok}\\*
0 & -- & 0 & 5 & 1 & 438 & 2 & 584 & \texttt{-?????????}\\*
1 & 0 & 1 & 4 & 1 & 144 & 16 & 864 & \texttt{--??????--}\\
2 & 1 & 2 & 1 & 2 & 512 & 512 & 0 & \texttt{---+------}\\
3 & 2 & 3 & 1 & 2 & 0 & 1024 & 0 & \texttt{----------}\\
4 & 1 & 2 & 1 & 2 & 512 & 512 & 0 & \texttt{-----+----}\\
5 & 4 & 3 & 1 & 2 & 0 & 1024 & 0 & \texttt{----------}\\
6 & 1 & 2 & 1 & 2 & 0 & 1024 & 0 & \texttt{----------}\\
7 & 1 & 2 & 1 & 2 & 0 & 1024 & 0 & \texttt{----------}\\
8 & 0 & 1 & 1 & 1 & 512 & 512 & 0 & \texttt{---------+}\\
9 & 8 & 2 & 1 & 1 & 0 & 1024 & 0 & \texttt{----------}\\
\multicolumn{9}{@{}l@{}}{\rule{0pt}{2.6ex}\textbf{T10.638}\enspace \texttt{0123232211}\enspace $n{=}10$, $|H|{=}6$, len 4, br 2, $|\mathrm{Aut}|{=}8$, orb 6/2, hist-orb 3/0; M 17408, res 12636, rev 2714: ok}\\*
0 & -- & 0 & 6 & 1 & 146 & 2 & 876 & \texttt{-?????????}\\*
1 & 0 & 1 & 4 & 1 & 144 & 16 & 864 & \texttt{--??????--}\\
2 & 1 & 2 & 1 & 2 & 512 & 512 & 0 & \texttt{---+------}\\
3 & 2 & 3 & 1 & 2 & 0 & 1024 & 0 & \texttt{----------}\\
4 & 1 & 2 & 1 & 2 & 512 & 512 & 0 & \texttt{-----+----}\\
5 & 4 & 3 & 1 & 2 & 0 & 1024 & 0 & \texttt{----------}\\
6 & 1 & 2 & 1 & 2 & 0 & 1024 & 0 & \texttt{----------}\\
7 & 1 & 2 & 1 & 2 & 0 & 1024 & 0 & \texttt{----------}\\
8 & 0 & 1 & 1 & 2 & 0 & 1024 & 0 & \texttt{----------}\\
9 & 0 & 1 & 1 & 2 & 0 & 1024 & 0 & \texttt{----------}\\
\multicolumn{9}{@{}l@{}}{\rule{0pt}{2.6ex}\textbf{T10.639}\enspace \texttt{0123232123}\enspace $n{=}10$, $|H|{=}4$, len 4, br 2, $|\mathrm{Aut}|{=}2$, orb 8/6, hist-orb 3/2; M 19456, res 7200, rev 6786: ok}\\*
0 & -- & 0 & 4 & 1 & 574 & 2 & 448 & \texttt{-?????????}\\*
1 & 0 & 1 & 3 & 1 & 288 & 32 & 704 & \texttt{--?????---}\\
2 & 1 & 2 & 1 & 2 & 512 & 512 & 0 & \texttt{---+------}\\
3 & 2 & 3 & 1 & 2 & 0 & 1024 & 0 & \texttt{----------}\\
4 & 1 & 2 & 1 & 2 & 512 & 512 & 0 & \texttt{-----+----}\\
5 & 4 & 3 & 1 & 2 & 0 & 1024 & 0 & \texttt{----------}\\
6 & 1 & 2 & 1 & 1 & 0 & 1024 & 0 & \texttt{----------}\\
7 & 0 & 1 & 1 & 1 & 768 & 256 & 0 & \texttt{--------++}\\
8 & 7 & 2 & 1 & 1 & 512 & 512 & 0 & \texttt{---------+}\\
9 & 8 & 3 & 1 & 1 & 0 & 1024 & 0 & \texttt{----------}\\
\multicolumn{9}{@{}l@{}}{\rule{0pt}{2.6ex}\textbf{T10.640}\enspace \texttt{0123232122}\enspace $n{=}10$, $|H|{=}5$, len 4, br 3, $|\mathrm{Aut}|{=}4$, orb 7/4, hist-orb 3/1; M 18432, res 9444, rev 5114: ok}\\*
0 & -- & 0 & 5 & 1 & 410 & 2 & 612 & \texttt{-?????????}\\*
1 & 0 & 1 & 3 & 1 & 288 & 32 & 704 & \texttt{--?????---}\\
2 & 1 & 2 & 1 & 2 & 512 & 512 & 0 & \texttt{---+------}\\
3 & 2 & 3 & 1 & 2 & 0 & 1024 & 0 & \texttt{----------}\\
4 & 1 & 2 & 1 & 2 & 512 & 512 & 0 & \texttt{-----+----}\\
5 & 4 & 3 & 1 & 2 & 0 & 1024 & 0 & \texttt{----------}\\
6 & 1 & 2 & 1 & 1 & 0 & 1024 & 0 & \texttt{----------}\\
7 & 0 & 1 & 2 & 1 & 256 & 256 & 512 & \texttt{--------??}\\
8 & 7 & 2 & 1 & 2 & 0 & 1024 & 0 & \texttt{----------}\\
9 & 7 & 2 & 1 & 2 & 0 & 1024 & 0 & \texttt{----------}\\
\multicolumn{9}{@{}l@{}}{\rule{0pt}{2.6ex}\textbf{T10.641}\enspace \texttt{0123232121}\enspace $n{=}10$, $|H|{=}5$, len 4, br 2, $|\mathrm{Aut}|{=}2$, orb 8/6, hist-orb 4/3; M 17408, res 10152, rev 4138: ok}\\*
0 & -- & 0 & 5 & 1 & 246 & 2 & 776 & \texttt{-?????????}\\*
1 & 0 & 1 & 3 & 1 & 288 & 32 & 704 & \texttt{--?????---}\\
2 & 1 & 2 & 1 & 2 & 512 & 512 & 0 & \texttt{---+------}\\
3 & 2 & 3 & 1 & 2 & 0 & 1024 & 0 & \texttt{----------}\\
4 & 1 & 2 & 1 & 2 & 512 & 512 & 0 & \texttt{-----+----}\\
5 & 4 & 3 & 1 & 2 & 0 & 1024 & 0 & \texttt{----------}\\
6 & 1 & 2 & 1 & 1 & 0 & 1024 & 0 & \texttt{----------}\\
7 & 0 & 1 & 1 & 1 & 512 & 512 & 0 & \texttt{--------+-}\\
8 & 7 & 2 & 1 & 1 & 0 & 1024 & 0 & \texttt{----------}\\
9 & 0 & 1 & 1 & 1 & 0 & 1024 & 0 & \texttt{----------}\\
\multicolumn{9}{@{}l@{}}{\rule{0pt}{2.6ex}\textbf{T10.642}\enspace \texttt{0123232111}\enspace $n{=}10$, $|H|{=}6$, len 4, br 2, $|\mathrm{Aut}|{=}12$, orb 6/3, hist-orb 3/1; M 16384, res 11628, rev 2650: ok}\\*
0 & -- & 0 & 6 & 1 & 82 & 2 & 940 & \texttt{-?????????}\\*
1 & 0 & 1 & 3 & 1 & 288 & 32 & 704 & \texttt{--?????---}\\
2 & 1 & 2 & 1 & 2 & 512 & 512 & 0 & \texttt{---+------}\\
3 & 2 & 3 & 1 & 2 & 0 & 1024 & 0 & \texttt{----------}\\
4 & 1 & 2 & 1 & 2 & 512 & 512 & 0 & \texttt{-----+----}\\
5 & 4 & 3 & 1 & 2 & 0 & 1024 & 0 & \texttt{----------}\\
6 & 1 & 2 & 1 & 1 & 0 & 1024 & 0 & \texttt{----------}\\
7 & 0 & 1 & 1 & 3 & 0 & 1024 & 0 & \texttt{----------}\\
8 & 0 & 1 & 1 & 3 & 0 & 1024 & 0 & \texttt{----------}\\
9 & 0 & 1 & 1 & 3 & 0 & 1024 & 0 & \texttt{----------}\\
\multicolumn{9}{@{}l@{}}{\rule{0pt}{2.6ex}\textbf{T10.643}\enspace \texttt{0123231232}\enspace $n{=}10$, $|H|{=}4$, len 4, br 3, $|\mathrm{Aut}|{=}2$, orb 8/6, hist-orb 3/2; M 19456, res 6872, rev 6974: ok}\\*
0 & -- & 0 & 4 & 1 & 550 & 2 & 472 & \texttt{-?????????}\\*
1 & 0 & 1 & 2 & 1 & 576 & 64 & 384 & \texttt{--????----}\\
2 & 1 & 2 & 1 & 2 & 512 & 512 & 0 & \texttt{---+------}\\
3 & 2 & 3 & 1 & 2 & 0 & 1024 & 0 & \texttt{----------}\\
4 & 1 & 2 & 1 & 2 & 512 & 512 & 0 & \texttt{-----+----}\\
5 & 4 & 3 & 1 & 2 & 0 & 1024 & 0 & \texttt{----------}\\
6 & 0 & 1 & 2 & 1 & 384 & 128 & 512 & \texttt{-------???}\\
7 & 6 & 2 & 1 & 1 & 512 & 512 & 0 & \texttt{--------+-}\\
8 & 7 & 3 & 1 & 1 & 0 & 1024 & 0 & \texttt{----------}\\
9 & 6 & 2 & 1 & 1 & 0 & 1024 & 0 & \texttt{----------}\\
\multicolumn{9}{@{}l@{}}{\rule{0pt}{2.6ex}\textbf{T10.644}\enspace \texttt{0123231231}\enspace $n{=}10$, $|H|{=}4$, len 4, br 2, $|\mathrm{Aut}|{=}2$, orb 8/6, hist-orb 3/2; M 18432, res 7392, rev 6018: ok}\\*
0 & -- & 0 & 4 & 1 & 350 & 2 & 672 & \texttt{-?????????}\\*
1 & 0 & 1 & 2 & 1 & 576 & 64 & 384 & \texttt{--????----}\\
2 & 1 & 2 & 1 & 2 & 512 & 512 & 0 & \texttt{---+------}\\
3 & 2 & 3 & 1 & 2 & 0 & 1024 & 0 & \texttt{----------}\\
4 & 1 & 2 & 1 & 2 & 512 & 512 & 0 & \texttt{-----+----}\\
5 & 4 & 3 & 1 & 2 & 0 & 1024 & 0 & \texttt{----------}\\
6 & 0 & 1 & 1 & 1 & 768 & 256 & 0 & \texttt{-------++-}\\
7 & 6 & 2 & 1 & 1 & 512 & 512 & 0 & \texttt{--------+-}\\
8 & 7 & 3 & 1 & 1 & 0 & 1024 & 0 & \texttt{----------}\\
9 & 0 & 1 & 1 & 1 & 0 & 1024 & 0 & \texttt{----------}\\
\multicolumn{9}{@{}l@{}}{\rule{0pt}{2.6ex}\textbf{T10.645}\enspace \texttt{0123231222}\enspace $n{=}10$, $|H|{=}5$, len 4, br 3, $|\mathrm{Aut}|{=}12$, orb 6/3, hist-orb 2/0; M 18432, res 8412, rev 5658: ok}\\*
0 & -- & 0 & 5 & 1 & 450 & 2 & 572 & \texttt{-?????????}\\*
1 & 0 & 1 & 2 & 1 & 576 & 64 & 384 & \texttt{--????----}\\
2 & 1 & 2 & 1 & 2 & 512 & 512 & 0 & \texttt{---+------}\\
3 & 2 & 3 & 1 & 2 & 0 & 1024 & 0 & \texttt{----------}\\
4 & 1 & 2 & 1 & 2 & 512 & 512 & 0 & \texttt{-----+----}\\
5 & 4 & 3 & 1 & 2 & 0 & 1024 & 0 & \texttt{----------}\\
6 & 0 & 1 & 3 & 1 & 128 & 128 & 768 & \texttt{-------???}\\
7 & 6 & 2 & 1 & 3 & 0 & 1024 & 0 & \texttt{----------}\\
8 & 6 & 2 & 1 & 3 & 0 & 1024 & 0 & \texttt{----------}\\
9 & 6 & 2 & 1 & 3 & 0 & 1024 & 0 & \texttt{----------}\\
\multicolumn{9}{@{}l@{}}{\rule{0pt}{2.6ex}\textbf{T10.646}\enspace \texttt{0123231221}\enspace $n{=}10$, $|H|{=}5$, len 4, br 3, $|\mathrm{Aut}|{=}4$, orb 7/4, hist-orb 3/1; M 17408, res 9060, rev 4602: ok}\\*
0 & -- & 0 & 5 & 1 & 250 & 2 & 772 & \texttt{-?????????}\\*
1 & 0 & 1 & 2 & 1 & 576 & 64 & 384 & \texttt{--????----}\\
2 & 1 & 2 & 1 & 2 & 512 & 512 & 0 & \texttt{---+------}\\
3 & 2 & 3 & 1 & 2 & 0 & 1024 & 0 & \texttt{----------}\\
4 & 1 & 2 & 1 & 2 & 512 & 512 & 0 & \texttt{-----+----}\\
5 & 4 & 3 & 1 & 2 & 0 & 1024 & 0 & \texttt{----------}\\
6 & 0 & 1 & 2 & 1 & 256 & 256 & 512 & \texttt{-------??-}\\
7 & 6 & 2 & 1 & 2 & 0 & 1024 & 0 & \texttt{----------}\\
8 & 6 & 2 & 1 & 2 & 0 & 1024 & 0 & \texttt{----------}\\
9 & 0 & 1 & 1 & 1 & 0 & 1024 & 0 & \texttt{----------}\\
\multicolumn{9}{@{}l@{}}{\rule{0pt}{2.6ex}\textbf{T10.647}\enspace \texttt{0123231212}\enspace $n{=}10$, $|H|{=}4$, len 4, br 2, $|\mathrm{Aut}|{=}4$, orb 6/2, hist-orb 2/0; M 17408, res 6492, rev 6098: ok}\\*
0 & -- & 0 & 4 & 1 & 450 & 2 & 572 & \texttt{-?????????}\\*
1 & 0 & 1 & 2 & 1 & 576 & 64 & 384 & \texttt{--????----}\\
2 & 1 & 2 & 1 & 2 & 512 & 512 & 0 & \texttt{---+------}\\
3 & 2 & 3 & 1 & 2 & 0 & 1024 & 0 & \texttt{----------}\\
4 & 1 & 2 & 1 & 2 & 512 & 512 & 0 & \texttt{-----+----}\\
5 & 4 & 3 & 1 & 2 & 0 & 1024 & 0 & \texttt{----------}\\
6 & 0 & 1 & 1 & 2 & 512 & 512 & 0 & \texttt{-------+--}\\
7 & 6 & 2 & 1 & 2 & 0 & 1024 & 0 & \texttt{----------}\\
8 & 0 & 1 & 1 & 2 & 512 & 512 & 0 & \texttt{---------+}\\
9 & 8 & 2 & 1 & 2 & 0 & 1024 & 0 & \texttt{----------}\\
\multicolumn{9}{@{}l@{}}{\rule{0pt}{2.6ex}\textbf{T10.648}\enspace \texttt{0123231211}\enspace $n{=}10$, $|H|{=}5$, len 4, br 2, $|\mathrm{Aut}|{=}4$, orb 7/4, hist-orb 3/1; M 16384, res 9192, rev 4010: ok}\\*
0 & -- & 0 & 5 & 1 & 150 & 2 & 872 & \texttt{-?????????}\\*
1 & 0 & 1 & 2 & 1 & 576 & 64 & 384 & \texttt{--????----}\\
2 & 1 & 2 & 1 & 2 & 512 & 512 & 0 & \texttt{---+------}\\
3 & 2 & 3 & 1 & 2 & 0 & 1024 & 0 & \texttt{----------}\\
4 & 1 & 2 & 1 & 2 & 512 & 512 & 0 & \texttt{-----+----}\\
5 & 4 & 3 & 1 & 2 & 0 & 1024 & 0 & \texttt{----------}\\
6 & 0 & 1 & 1 & 1 & 512 & 512 & 0 & \texttt{-------+--}\\
7 & 6 & 2 & 1 & 1 & 0 & 1024 & 0 & \texttt{----------}\\
8 & 0 & 1 & 1 & 2 & 0 & 1024 & 0 & \texttt{----------}\\
9 & 0 & 1 & 1 & 2 & 0 & 1024 & 0 & \texttt{----------}\\
\multicolumn{9}{@{}l@{}}{\rule{0pt}{2.6ex}\textbf{T10.649}\enspace \texttt{0123231111}\enspace $n{=}10$, $|H|{=}6$, len 4, br 2, $|\mathrm{Aut}|{=}48$, orb 5/2, hist-orb 2/0; M 15360, res 10092, rev 2906: ok}\\*
0 & -- & 0 & 6 & 1 & 50 & 2 & 972 & \texttt{-?????????}\\*
1 & 0 & 1 & 2 & 1 & 576 & 64 & 384 & \texttt{--????----}\\
2 & 1 & 2 & 1 & 2 & 512 & 512 & 0 & \texttt{---+------}\\
3 & 2 & 3 & 1 & 2 & 0 & 1024 & 0 & \texttt{----------}\\
4 & 1 & 2 & 1 & 2 & 512 & 512 & 0 & \texttt{-----+----}\\
5 & 4 & 3 & 1 & 2 & 0 & 1024 & 0 & \texttt{----------}\\
6 & 0 & 1 & 1 & 4 & 0 & 1024 & 0 & \texttt{----------}\\
7 & 0 & 1 & 1 & 4 & 0 & 1024 & 0 & \texttt{----------}\\
8 & 0 & 1 & 1 & 4 & 0 & 1024 & 0 & \texttt{----------}\\
9 & 0 & 1 & 1 & 4 & 0 & 1024 & 0 & \texttt{----------}\\
\multicolumn{9}{@{}l@{}}{\rule{0pt}{2.6ex}\textbf{T10.650}\enspace \texttt{0123222222}\enspace $n{=}10$, $|H|{=}7$, len 4, br 1, $|\mathrm{Aut}|{=}720$, orb 5/4, hist-orb 2/1; M 18432, res 12096, rev 4486: ok}\\*
0 & -- & 0 & 7 & 1 & 518 & 2 & 504 & \texttt{-+????????}\\*
1 & 0 & 1 & 7 & 1 & 12 & 4 & 1008 & \texttt{--????????}\\
2 & 1 & 2 & 1 & 1 & 512 & 512 & 0 & \texttt{---+------}\\
3 & 2 & 3 & 1 & 1 & 0 & 1024 & 0 & \texttt{----------}\\
4 & 1 & 2 & 1 & 6 & 0 & 1024 & 0 & \texttt{----------}\\
5 & 1 & 2 & 1 & 6 & 0 & 1024 & 0 & \texttt{----------}\\
6 & 1 & 2 & 1 & 6 & 0 & 1024 & 0 & \texttt{----------}\\
7 & 1 & 2 & 1 & 6 & 0 & 1024 & 0 & \texttt{----------}\\
8 & 1 & 2 & 1 & 6 & 0 & 1024 & 0 & \texttt{----------}\\
9 & 1 & 2 & 1 & 6 & 0 & 1024 & 0 & \texttt{----------}\\
\multicolumn{9}{@{}l@{}}{\rule{0pt}{2.6ex}\textbf{T10.651}\enspace \texttt{0123222221}\enspace $n{=}10$, $|H|{=}7$, len 4, br 2, $|\mathrm{Aut}|{=}120$, orb 6/5, hist-orb 3/2; M 17408, res 13288, rev 2630: ok}\\*
0 & -- & 0 & 7 & 1 & 262 & 2 & 760 & \texttt{-?????????}\\*
1 & 0 & 1 & 6 & 1 & 24 & 8 & 992 & \texttt{--???????-}\\
2 & 1 & 2 & 1 & 1 & 512 & 512 & 0 & \texttt{---+------}\\
3 & 2 & 3 & 1 & 1 & 0 & 1024 & 0 & \texttt{----------}\\
4 & 1 & 2 & 1 & 5 & 0 & 1024 & 0 & \texttt{----------}\\
5 & 1 & 2 & 1 & 5 & 0 & 1024 & 0 & \texttt{----------}\\
6 & 1 & 2 & 1 & 5 & 0 & 1024 & 0 & \texttt{----------}\\
7 & 1 & 2 & 1 & 5 & 0 & 1024 & 0 & \texttt{----------}\\
8 & 1 & 2 & 1 & 5 & 0 & 1024 & 0 & \texttt{----------}\\
9 & 0 & 1 & 1 & 1 & 0 & 1024 & 0 & \texttt{----------}\\
\multicolumn{9}{@{}l@{}}{\rule{0pt}{2.6ex}\textbf{T10.652}\enspace \texttt{0123222212}\enspace $n{=}10$, $|H|{=}6$, len 4, br 2, $|\mathrm{Aut}|{=}24$, orb 7/6, hist-orb 3/2; M 17408, res 10860, rev 4030: ok}\\*
0 & -- & 0 & 6 & 1 & 402 & 2 & 620 & \texttt{-?????????}\\*
1 & 0 & 1 & 5 & 1 & 48 & 16 & 960 & \texttt{--??????--}\\
2 & 1 & 2 & 1 & 1 & 512 & 512 & 0 & \texttt{---+------}\\
3 & 2 & 3 & 1 & 1 & 0 & 1024 & 0 & \texttt{----------}\\
4 & 1 & 2 & 1 & 4 & 0 & 1024 & 0 & \texttt{----------}\\
5 & 1 & 2 & 1 & 4 & 0 & 1024 & 0 & \texttt{----------}\\
6 & 1 & 2 & 1 & 4 & 0 & 1024 & 0 & \texttt{----------}\\
7 & 1 & 2 & 1 & 4 & 0 & 1024 & 0 & \texttt{----------}\\
8 & 0 & 1 & 1 & 1 & 512 & 512 & 0 & \texttt{---------+}\\
9 & 8 & 2 & 1 & 1 & 0 & 1024 & 0 & \texttt{----------}\\
\multicolumn{9}{@{}l@{}}{\rule{0pt}{2.6ex}\textbf{T10.653}\enspace \texttt{0123222211}\enspace $n{=}10$, $|H|{=}7$, len 4, br 2, $|\mathrm{Aut}|{=}48$, orb 6/4, hist-orb 3/1; M 16384, res 13272, rev 1774: ok}\\*
0 & -- & 0 & 7 & 1 & 134 & 2 & 888 & \texttt{-?????????}\\*
1 & 0 & 1 & 5 & 1 & 48 & 16 & 960 & \texttt{--??????--}\\
2 & 1 & 2 & 1 & 1 & 512 & 512 & 0 & \texttt{---+------}\\
3 & 2 & 3 & 1 & 1 & 0 & 1024 & 0 & \texttt{----------}\\
4 & 1 & 2 & 1 & 4 & 0 & 1024 & 0 & \texttt{----------}\\
5 & 1 & 2 & 1 & 4 & 0 & 1024 & 0 & \texttt{----------}\\
6 & 1 & 2 & 1 & 4 & 0 & 1024 & 0 & \texttt{----------}\\
7 & 1 & 2 & 1 & 4 & 0 & 1024 & 0 & \texttt{----------}\\
8 & 0 & 1 & 1 & 2 & 0 & 1024 & 0 & \texttt{----------}\\
9 & 0 & 1 & 1 & 2 & 0 & 1024 & 0 & \texttt{----------}\\
\multicolumn{9}{@{}l@{}}{\rule{0pt}{2.6ex}\textbf{T10.654}\enspace \texttt{0123222123}\enspace $n{=}10$, $|H|{=}5$, len 4, br 2, $|\mathrm{Aut}|{=}6$, orb 8/7, hist-orb 3/2; M 18432, res 8820, rev 5446: ok}\\*
0 & -- & 0 & 5 & 1 & 490 & 2 & 532 & \texttt{-?????????}\\*
1 & 0 & 1 & 4 & 1 & 96 & 32 & 896 & \texttt{--?????---}\\
2 & 1 & 2 & 1 & 1 & 512 & 512 & 0 & \texttt{---+------}\\
3 & 2 & 3 & 1 & 1 & 0 & 1024 & 0 & \texttt{----------}\\
4 & 1 & 2 & 1 & 3 & 0 & 1024 & 0 & \texttt{----------}\\
5 & 1 & 2 & 1 & 3 & 0 & 1024 & 0 & \texttt{----------}\\
6 & 1 & 2 & 1 & 3 & 0 & 1024 & 0 & \texttt{----------}\\
7 & 0 & 1 & 1 & 1 & 768 & 256 & 0 & \texttt{--------++}\\
8 & 7 & 2 & 1 & 1 & 512 & 512 & 0 & \texttt{---------+}\\
9 & 8 & 3 & 1 & 1 & 0 & 1024 & 0 & \texttt{----------}\\
\multicolumn{9}{@{}l@{}}{\rule{0pt}{2.6ex}\textbf{T10.655}\enspace \texttt{0123222122}\enspace $n{=}10$, $|H|{=}6$, len 4, br 3, $|\mathrm{Aut}|{=}12$, orb 7/5, hist-orb 3/1; M 17408, res 10848, rev 3790: ok}\\*
0 & -- & 0 & 6 & 1 & 350 & 2 & 672 & \texttt{-?????????}\\*
1 & 0 & 1 & 4 & 1 & 96 & 32 & 896 & \texttt{--?????---}\\
2 & 1 & 2 & 1 & 1 & 512 & 512 & 0 & \texttt{---+------}\\
3 & 2 & 3 & 1 & 1 & 0 & 1024 & 0 & \texttt{----------}\\
4 & 1 & 2 & 1 & 3 & 0 & 1024 & 0 & \texttt{----------}\\
5 & 1 & 2 & 1 & 3 & 0 & 1024 & 0 & \texttt{----------}\\
6 & 1 & 2 & 1 & 3 & 0 & 1024 & 0 & \texttt{----------}\\
7 & 0 & 1 & 2 & 1 & 256 & 256 & 512 & \texttt{--------??}\\
8 & 7 & 2 & 1 & 2 & 0 & 1024 & 0 & \texttt{----------}\\
9 & 7 & 2 & 1 & 2 & 0 & 1024 & 0 & \texttt{----------}\\
\multicolumn{9}{@{}l@{}}{\rule{0pt}{2.6ex}\textbf{T10.656}\enspace \texttt{0123222121}\enspace $n{=}10$, $|H|{=}6$, len 4, br 2, $|\mathrm{Aut}|{=}6$, orb 8/7, hist-orb 4/3; M 16384, res 11340, rev 2878: ok}\\*
0 & -- & 0 & 6 & 1 & 210 & 2 & 812 & \texttt{-?????????}\\*
1 & 0 & 1 & 4 & 1 & 96 & 32 & 896 & \texttt{--?????---}\\
2 & 1 & 2 & 1 & 1 & 512 & 512 & 0 & \texttt{---+------}\\
3 & 2 & 3 & 1 & 1 & 0 & 1024 & 0 & \texttt{----------}\\
4 & 1 & 2 & 1 & 3 & 0 & 1024 & 0 & \texttt{----------}\\
5 & 1 & 2 & 1 & 3 & 0 & 1024 & 0 & \texttt{----------}\\
6 & 1 & 2 & 1 & 3 & 0 & 1024 & 0 & \texttt{----------}\\
7 & 0 & 1 & 1 & 1 & 512 & 512 & 0 & \texttt{--------+-}\\
8 & 7 & 2 & 1 & 1 & 0 & 1024 & 0 & \texttt{----------}\\
9 & 0 & 1 & 1 & 1 & 0 & 1024 & 0 & \texttt{----------}\\
\multicolumn{9}{@{}l@{}}{\rule{0pt}{2.6ex}\textbf{T10.657}\enspace \texttt{0123222111}\enspace $n{=}10$, $|H|{=}7$, len 4, br 2, $|\mathrm{Aut}|{=}36$, orb 6/4, hist-orb 3/1; M 15360, res 12600, rev 1454: ok}\\*
0 & -- & 0 & 7 & 1 & 70 & 2 & 952 & \texttt{-?????????}\\*
1 & 0 & 1 & 4 & 1 & 96 & 32 & 896 & \texttt{--?????---}\\
2 & 1 & 2 & 1 & 1 & 512 & 512 & 0 & \texttt{---+------}\\
3 & 2 & 3 & 1 & 1 & 0 & 1024 & 0 & \texttt{----------}\\
4 & 1 & 2 & 1 & 3 & 0 & 1024 & 0 & \texttt{----------}\\
5 & 1 & 2 & 1 & 3 & 0 & 1024 & 0 & \texttt{----------}\\
6 & 1 & 2 & 1 & 3 & 0 & 1024 & 0 & \texttt{----------}\\
7 & 0 & 1 & 1 & 3 & 0 & 1024 & 0 & \texttt{----------}\\
8 & 0 & 1 & 1 & 3 & 0 & 1024 & 0 & \texttt{----------}\\
9 & 0 & 1 & 1 & 3 & 0 & 1024 & 0 & \texttt{----------}\\
\multicolumn{9}{@{}l@{}}{\rule{0pt}{2.6ex}\textbf{T10.658}\enspace \texttt{0123221232}\enspace $n{=}10$, $|H|{=}5$, len 4, br 3, $|\mathrm{Aut}|{=}2$, orb 9/8, hist-orb 4/3; M 18432, res 9404, rev 5098: ok}\\*
0 & -- & 0 & 5 & 1 & 418 & 2 & 604 & \texttt{-?????????}\\*
1 & 0 & 1 & 3 & 1 & 192 & 64 & 768 & \texttt{--????----}\\
2 & 1 & 2 & 1 & 1 & 512 & 512 & 0 & \texttt{---+------}\\
3 & 2 & 3 & 1 & 1 & 0 & 1024 & 0 & \texttt{----------}\\
4 & 1 & 2 & 1 & 2 & 0 & 1024 & 0 & \texttt{----------}\\
5 & 1 & 2 & 1 & 2 & 0 & 1024 & 0 & \texttt{----------}\\
6 & 0 & 1 & 2 & 1 & 384 & 128 & 512 & \texttt{-------???}\\
7 & 6 & 2 & 1 & 1 & 512 & 512 & 0 & \texttt{--------+-}\\
8 & 7 & 3 & 1 & 1 & 0 & 1024 & 0 & \texttt{----------}\\
9 & 6 & 2 & 1 & 1 & 0 & 1024 & 0 & \texttt{----------}\\
\multicolumn{9}{@{}l@{}}{\rule{0pt}{2.6ex}\textbf{T10.659}\enspace \texttt{0123221231}\enspace $n{=}10$, $|H|{=}5$, len 4, br 2, $|\mathrm{Aut}|{=}2$, orb 9/8, hist-orb 4/3; M 17408, res 9492, rev 4310: ok}\\*
0 & -- & 0 & 5 & 1 & 266 & 2 & 756 & \texttt{-?????????}\\*
1 & 0 & 1 & 3 & 1 & 192 & 64 & 768 & \texttt{--????----}\\
2 & 1 & 2 & 1 & 1 & 512 & 512 & 0 & \texttt{---+------}\\
3 & 2 & 3 & 1 & 1 & 0 & 1024 & 0 & \texttt{----------}\\
4 & 1 & 2 & 1 & 2 & 0 & 1024 & 0 & \texttt{----------}\\
5 & 1 & 2 & 1 & 2 & 0 & 1024 & 0 & \texttt{----------}\\
6 & 0 & 1 & 1 & 1 & 768 & 256 & 0 & \texttt{-------++-}\\
7 & 6 & 2 & 1 & 1 & 512 & 512 & 0 & \texttt{--------+-}\\
8 & 7 & 3 & 1 & 1 & 0 & 1024 & 0 & \texttt{----------}\\
9 & 0 & 1 & 1 & 1 & 0 & 1024 & 0 & \texttt{----------}\\
\multicolumn{9}{@{}l@{}}{\rule{0pt}{2.6ex}\textbf{T10.660}\enspace \texttt{0123221222}\enspace $n{=}10$, $|H|{=}6$, len 4, br 3, $|\mathrm{Aut}|{=}12$, orb 7/5, hist-orb 3/1; M 17408, res 10728, rev 3806: ok}\\*
0 & -- & 0 & 6 & 1 & 342 & 2 & 680 & \texttt{-?????????}\\*
1 & 0 & 1 & 3 & 1 & 192 & 64 & 768 & \texttt{--????----}\\
2 & 1 & 2 & 1 & 1 & 512 & 512 & 0 & \texttt{---+------}\\
3 & 2 & 3 & 1 & 1 & 0 & 1024 & 0 & \texttt{----------}\\
4 & 1 & 2 & 1 & 2 & 0 & 1024 & 0 & \texttt{----------}\\
5 & 1 & 2 & 1 & 2 & 0 & 1024 & 0 & \texttt{----------}\\
6 & 0 & 1 & 3 & 1 & 128 & 128 & 768 & \texttt{-------???}\\
7 & 6 & 2 & 1 & 3 & 0 & 1024 & 0 & \texttt{----------}\\
8 & 6 & 2 & 1 & 3 & 0 & 1024 & 0 & \texttt{----------}\\
9 & 6 & 2 & 1 & 3 & 0 & 1024 & 0 & \texttt{----------}\\
\multicolumn{9}{@{}l@{}}{\rule{0pt}{2.6ex}\textbf{T10.661}\enspace \texttt{0123221221}\enspace $n{=}10$, $|H|{=}6$, len 4, br 3, $|\mathrm{Aut}|{=}4$, orb 8/6, hist-orb 4/2; M 16384, res 10944, rev 2942: ok}\\*
0 & -- & 0 & 6 & 1 & 190 & 2 & 832 & \texttt{-?????????}\\*
1 & 0 & 1 & 3 & 1 & 192 & 64 & 768 & \texttt{--????----}\\
2 & 1 & 2 & 1 & 1 & 512 & 512 & 0 & \texttt{---+------}\\
3 & 2 & 3 & 1 & 1 & 0 & 1024 & 0 & \texttt{----------}\\
4 & 1 & 2 & 1 & 2 & 0 & 1024 & 0 & \texttt{----------}\\
5 & 1 & 2 & 1 & 2 & 0 & 1024 & 0 & \texttt{----------}\\
6 & 0 & 1 & 2 & 1 & 256 & 256 & 512 & \texttt{-------??-}\\
7 & 6 & 2 & 1 & 2 & 0 & 1024 & 0 & \texttt{----------}\\
8 & 6 & 2 & 1 & 2 & 0 & 1024 & 0 & \texttt{----------}\\
9 & 0 & 1 & 1 & 1 & 0 & 1024 & 0 & \texttt{----------}\\
\multicolumn{9}{@{}l@{}}{\rule{0pt}{2.6ex}\textbf{T10.662}\enspace \texttt{0123221212}\enspace $n{=}10$, $|H|{=}5$, len 4, br 2, $|\mathrm{Aut}|{=}4$, orb 7/4, hist-orb 3/1; M 16384, res 8808, rev 4294: ok}\\*
0 & -- & 0 & 5 & 1 & 342 & 2 & 680 & \texttt{-?????????}\\*
1 & 0 & 1 & 3 & 1 & 192 & 64 & 768 & \texttt{--????----}\\
2 & 1 & 2 & 1 & 1 & 512 & 512 & 0 & \texttt{---+------}\\
3 & 2 & 3 & 1 & 1 & 0 & 1024 & 0 & \texttt{----------}\\
4 & 1 & 2 & 1 & 2 & 0 & 1024 & 0 & \texttt{----------}\\
5 & 1 & 2 & 1 & 2 & 0 & 1024 & 0 & \texttt{----------}\\
6 & 0 & 1 & 1 & 2 & 512 & 512 & 0 & \texttt{-------+--}\\
7 & 6 & 2 & 1 & 2 & 0 & 1024 & 0 & \texttt{----------}\\
8 & 0 & 1 & 1 & 2 & 512 & 512 & 0 & \texttt{---------+}\\
9 & 8 & 2 & 1 & 2 & 0 & 1024 & 0 & \texttt{----------}\\
\multicolumn{9}{@{}l@{}}{\rule{0pt}{2.6ex}\textbf{T10.663}\enspace \texttt{0123221211}\enspace $n{=}10$, $|H|{=}6$, len 4, br 2, $|\mathrm{Aut}|{=}4$, orb 8/6, hist-orb 4/2; M 15360, res 10860, rev 2446: ok}\\*
0 & -- & 0 & 6 & 1 & 114 & 2 & 908 & \texttt{-?????????}\\*
1 & 0 & 1 & 3 & 1 & 192 & 64 & 768 & \texttt{--????----}\\
2 & 1 & 2 & 1 & 1 & 512 & 512 & 0 & \texttt{---+------}\\
3 & 2 & 3 & 1 & 1 & 0 & 1024 & 0 & \texttt{----------}\\
4 & 1 & 2 & 1 & 2 & 0 & 1024 & 0 & \texttt{----------}\\
5 & 1 & 2 & 1 & 2 & 0 & 1024 & 0 & \texttt{----------}\\
6 & 0 & 1 & 1 & 1 & 512 & 512 & 0 & \texttt{-------+--}\\
7 & 6 & 2 & 1 & 1 & 0 & 1024 & 0 & \texttt{----------}\\
8 & 0 & 1 & 1 & 2 & 0 & 1024 & 0 & \texttt{----------}\\
9 & 0 & 1 & 1 & 2 & 0 & 1024 & 0 & \texttt{----------}\\
\multicolumn{9}{@{}l@{}}{\rule{0pt}{2.6ex}\textbf{T10.664}\enspace \texttt{0123221111}\enspace $n{=}10$, $|H|{=}7$, len 4, br 2, $|\mathrm{Aut}|{=}48$, orb 6/4, hist-orb 3/1; M 14336, res 11544, rev 1438: ok}\\*
0 & -- & 0 & 7 & 1 & 38 & 2 & 984 & \texttt{-?????????}\\*
1 & 0 & 1 & 3 & 1 & 192 & 64 & 768 & \texttt{--????----}\\
2 & 1 & 2 & 1 & 1 & 512 & 512 & 0 & \texttt{---+------}\\
3 & 2 & 3 & 1 & 1 & 0 & 1024 & 0 & \texttt{----------}\\
4 & 1 & 2 & 1 & 2 & 0 & 1024 & 0 & \texttt{----------}\\
5 & 1 & 2 & 1 & 2 & 0 & 1024 & 0 & \texttt{----------}\\
6 & 0 & 1 & 1 & 4 & 0 & 1024 & 0 & \texttt{----------}\\
7 & 0 & 1 & 1 & 4 & 0 & 1024 & 0 & \texttt{----------}\\
8 & 0 & 1 & 1 & 4 & 0 & 1024 & 0 & \texttt{----------}\\
9 & 0 & 1 & 1 & 4 & 0 & 1024 & 0 & \texttt{----------}\\
\multicolumn{9}{@{}l@{}}{\rule{0pt}{2.6ex}\textbf{T10.665}\enspace \texttt{0123212321}\enspace $n{=}10$, $|H|{=}5$, len 4, br 3, $|\mathrm{Aut}|{=}2$, orb 6/2, hist-orb 3/1; M 17408, res 9580, rev 4290: ok}\\*
0 & -- & 0 & 5 & 1 & 242 & 2 & 780 & \texttt{-?????????}\\*
1 & 0 & 1 & 2 & 2 & 384 & 128 & 512 & \texttt{--???-----}\\
2 & 1 & 2 & 1 & 2 & 512 & 512 & 0 & \texttt{---+------}\\
3 & 2 & 3 & 1 & 2 & 0 & 1024 & 0 & \texttt{----------}\\
4 & 1 & 2 & 1 & 2 & 0 & 1024 & 0 & \texttt{----------}\\
5 & 0 & 1 & 2 & 2 & 384 & 128 & 512 & \texttt{------???-}\\
6 & 5 & 2 & 1 & 2 & 512 & 512 & 0 & \texttt{-------+--}\\
7 & 6 & 3 & 1 & 2 & 0 & 1024 & 0 & \texttt{----------}\\
8 & 5 & 2 & 1 & 2 & 0 & 1024 & 0 & \texttt{----------}\\
9 & 0 & 1 & 1 & 1 & 0 & 1024 & 0 & \texttt{----------}\\
\multicolumn{9}{@{}l@{}}{\rule{0pt}{2.6ex}\textbf{T10.666}\enspace \texttt{0123212312}\enspace $n{=}10$, $|H|{=}4$, len 4, br 2, $|\mathrm{Aut}|{=}1$, orb 10/10, hist-orb 4/4; M 17408, res 6320, rev 5974: ok}\\*
0 & -- & 0 & 4 & 1 & 462 & 2 & 560 & \texttt{-?????????}\\*
1 & 0 & 1 & 2 & 1 & 384 & 128 & 512 & \texttt{--???-----}\\
2 & 1 & 2 & 1 & 1 & 512 & 512 & 0 & \texttt{---+------}\\
3 & 2 & 3 & 1 & 1 & 0 & 1024 & 0 & \texttt{----------}\\
4 & 1 & 2 & 1 & 1 & 0 & 1024 & 0 & \texttt{----------}\\
5 & 0 & 1 & 1 & 1 & 768 & 256 & 0 & \texttt{------++--}\\
6 & 5 & 2 & 1 & 1 & 512 & 512 & 0 & \texttt{-------+--}\\
7 & 6 & 3 & 1 & 1 & 0 & 1024 & 0 & \texttt{----------}\\
8 & 0 & 1 & 1 & 1 & 512 & 512 & 0 & \texttt{---------+}\\
9 & 8 & 2 & 1 & 1 & 0 & 1024 & 0 & \texttt{----------}\\
\multicolumn{9}{@{}l@{}}{\rule{0pt}{2.6ex}\textbf{T10.667}\enspace \texttt{0123212311}\enspace $n{=}10$, $|H|{=}5$, len 4, br 2, $|\mathrm{Aut}|{=}2$, orb 9/8, hist-orb 4/3; M 16384, res 9092, rev 3886: ok}\\*
0 & -- & 0 & 5 & 1 & 154 & 2 & 868 & \texttt{-?????????}\\*
1 & 0 & 1 & 2 & 1 & 384 & 128 & 512 & \texttt{--???-----}\\
2 & 1 & 2 & 1 & 1 & 512 & 512 & 0 & \texttt{---+------}\\
3 & 2 & 3 & 1 & 1 & 0 & 1024 & 0 & \texttt{----------}\\
4 & 1 & 2 & 1 & 1 & 0 & 1024 & 0 & \texttt{----------}\\
5 & 0 & 1 & 1 & 1 & 768 & 256 & 0 & \texttt{------++--}\\
6 & 5 & 2 & 1 & 1 & 512 & 512 & 0 & \texttt{-------+--}\\
7 & 6 & 3 & 1 & 1 & 0 & 1024 & 0 & \texttt{----------}\\
8 & 0 & 1 & 1 & 2 & 0 & 1024 & 0 & \texttt{----------}\\
9 & 0 & 1 & 1 & 2 & 0 & 1024 & 0 & \texttt{----------}\\
\multicolumn{9}{@{}l@{}}{\rule{0pt}{2.6ex}\textbf{T10.668}\enspace \texttt{0123212222}\enspace $n{=}10$, $|H|{=}6$, len 4, br 3, $|\mathrm{Aut}|{=}24$, orb 7/6, hist-orb 3/2; M 17408, res 10248, rev 4134: ok}\\*
0 & -- & 0 & 6 & 1 & 374 & 2 & 648 & \texttt{-?????????}\\*
1 & 0 & 1 & 2 & 1 & 384 & 128 & 512 & \texttt{--???-----}\\
2 & 1 & 2 & 1 & 1 & 512 & 512 & 0 & \texttt{---+------}\\
3 & 2 & 3 & 1 & 1 & 0 & 1024 & 0 & \texttt{----------}\\
4 & 1 & 2 & 1 & 1 & 0 & 1024 & 0 & \texttt{----------}\\
5 & 0 & 1 & 4 & 1 & 64 & 64 & 896 & \texttt{------????}\\
6 & 5 & 2 & 1 & 4 & 0 & 1024 & 0 & \texttt{----------}\\
7 & 5 & 2 & 1 & 4 & 0 & 1024 & 0 & \texttt{----------}\\
8 & 5 & 2 & 1 & 4 & 0 & 1024 & 0 & \texttt{----------}\\
9 & 5 & 2 & 1 & 4 & 0 & 1024 & 0 & \texttt{----------}\\
\multicolumn{9}{@{}l@{}}{\rule{0pt}{2.6ex}\textbf{T10.669}\enspace \texttt{0123212221}\enspace $n{=}10$, $|H|{=}6$, len 4, br 3, $|\mathrm{Aut}|{=}6$, orb 8/7, hist-orb 4/3; M 16384, res 10616, rev 3110: ok}\\*
0 & -- & 0 & 6 & 1 & 198 & 2 & 824 & \texttt{-?????????}\\*
1 & 0 & 1 & 2 & 1 & 384 & 128 & 512 & \texttt{--???-----}\\
2 & 1 & 2 & 1 & 1 & 512 & 512 & 0 & \texttt{---+------}\\
3 & 2 & 3 & 1 & 1 & 0 & 1024 & 0 & \texttt{----------}\\
4 & 1 & 2 & 1 & 1 & 0 & 1024 & 0 & \texttt{----------}\\
5 & 0 & 1 & 3 & 1 & 128 & 128 & 768 & \texttt{------???-}\\
6 & 5 & 2 & 1 & 3 & 0 & 1024 & 0 & \texttt{----------}\\
7 & 5 & 2 & 1 & 3 & 0 & 1024 & 0 & \texttt{----------}\\
8 & 5 & 2 & 1 & 3 & 0 & 1024 & 0 & \texttt{----------}\\
9 & 0 & 1 & 1 & 1 & 0 & 1024 & 0 & \texttt{----------}\\
\multicolumn{9}{@{}l@{}}{\rule{0pt}{2.6ex}\textbf{T10.670}\enspace \texttt{0123212212}\enspace $n{=}10$, $|H|{=}5$, len 4, br 3, $|\mathrm{Aut}|{=}2$, orb 9/8, hist-orb 4/3; M 16384, res 8276, rev 4438: ok}\\*
0 & -- & 0 & 5 & 1 & 330 & 2 & 692 & \texttt{-?????????}\\*
1 & 0 & 1 & 2 & 1 & 384 & 128 & 512 & \texttt{--???-----}\\
2 & 1 & 2 & 1 & 1 & 512 & 512 & 0 & \texttt{---+------}\\
3 & 2 & 3 & 1 & 1 & 0 & 1024 & 0 & \texttt{----------}\\
4 & 1 & 2 & 1 & 1 & 0 & 1024 & 0 & \texttt{----------}\\
5 & 0 & 1 & 2 & 1 & 256 & 256 & 512 & \texttt{------??--}\\
6 & 5 & 2 & 1 & 2 & 0 & 1024 & 0 & \texttt{----------}\\
7 & 5 & 2 & 1 & 2 & 0 & 1024 & 0 & \texttt{----------}\\
8 & 0 & 1 & 1 & 1 & 512 & 512 & 0 & \texttt{---------+}\\
9 & 8 & 2 & 1 & 1 & 0 & 1024 & 0 & \texttt{----------}\\
\multicolumn{9}{@{}l@{}}{\rule{0pt}{2.6ex}\textbf{T10.671}\enspace \texttt{0123212211}\enspace $n{=}10$, $|H|{=}6$, len 4, br 3, $|\mathrm{Aut}|{=}4$, orb 8/6, hist-orb 4/2; M 15360, res 10256, rev 2662: ok}\\*
0 & -- & 0 & 6 & 1 & 110 & 2 & 912 & \texttt{-?????????}\\*
1 & 0 & 1 & 2 & 1 & 384 & 128 & 512 & \texttt{--???-----}\\
2 & 1 & 2 & 1 & 1 & 512 & 512 & 0 & \texttt{---+------}\\
3 & 2 & 3 & 1 & 1 & 0 & 1024 & 0 & \texttt{----------}\\
4 & 1 & 2 & 1 & 1 & 0 & 1024 & 0 & \texttt{----------}\\
5 & 0 & 1 & 2 & 1 & 256 & 256 & 512 & \texttt{------??--}\\
6 & 5 & 2 & 1 & 2 & 0 & 1024 & 0 & \texttt{----------}\\
7 & 5 & 2 & 1 & 2 & 0 & 1024 & 0 & \texttt{----------}\\
8 & 0 & 1 & 1 & 2 & 0 & 1024 & 0 & \texttt{----------}\\
9 & 0 & 1 & 1 & 2 & 0 & 1024 & 0 & \texttt{----------}\\
\multicolumn{9}{@{}l@{}}{\rule{0pt}{2.6ex}\textbf{T10.672}\enspace \texttt{0123212121}\enspace $n{=}10$, $|H|{=}5$, len 4, br 2, $|\mathrm{Aut}|{=}2$, orb 8/6, hist-orb 4/3; M 15360, res 8696, rev 3662: ok}\\*
0 & -- & 0 & 5 & 1 & 198 & 2 & 824 & \texttt{-?????????}\\*
1 & 0 & 1 & 2 & 1 & 384 & 128 & 512 & \texttt{--???-----}\\
2 & 1 & 2 & 1 & 1 & 512 & 512 & 0 & \texttt{---+------}\\
3 & 2 & 3 & 1 & 1 & 0 & 1024 & 0 & \texttt{----------}\\
4 & 1 & 2 & 1 & 1 & 0 & 1024 & 0 & \texttt{----------}\\
5 & 0 & 1 & 1 & 2 & 512 & 512 & 0 & \texttt{------+---}\\
6 & 5 & 2 & 1 & 2 & 0 & 1024 & 0 & \texttt{----------}\\
7 & 0 & 1 & 1 & 2 & 512 & 512 & 0 & \texttt{--------+-}\\
8 & 7 & 2 & 1 & 2 & 0 & 1024 & 0 & \texttt{----------}\\
9 & 0 & 1 & 1 & 1 & 0 & 1024 & 0 & \texttt{----------}\\
\multicolumn{9}{@{}l@{}}{\rule{0pt}{2.6ex}\textbf{T10.673}\enspace \texttt{0123212111}\enspace $n{=}10$, $|H|{=}6$, len 4, br 2, $|\mathrm{Aut}|{=}6$, orb 8/7, hist-orb 4/3; M 14336, res 9884, rev 2374: ok}\\*
0 & -- & 0 & 6 & 1 & 66 & 2 & 956 & \texttt{-?????????}\\*
1 & 0 & 1 & 2 & 1 & 384 & 128 & 512 & \texttt{--???-----}\\
2 & 1 & 2 & 1 & 1 & 512 & 512 & 0 & \texttt{---+------}\\
3 & 2 & 3 & 1 & 1 & 0 & 1024 & 0 & \texttt{----------}\\
4 & 1 & 2 & 1 & 1 & 0 & 1024 & 0 & \texttt{----------}\\
5 & 0 & 1 & 1 & 1 & 512 & 512 & 0 & \texttt{------+---}\\
6 & 5 & 2 & 1 & 1 & 0 & 1024 & 0 & \texttt{----------}\\
7 & 0 & 1 & 1 & 3 & 0 & 1024 & 0 & \texttt{----------}\\
8 & 0 & 1 & 1 & 3 & 0 & 1024 & 0 & \texttt{----------}\\
9 & 0 & 1 & 1 & 3 & 0 & 1024 & 0 & \texttt{----------}\\
\multicolumn{9}{@{}l@{}}{\rule{0pt}{2.6ex}\textbf{T10.674}\enspace \texttt{0123211111}\enspace $n{=}10$, $|H|{=}7$, len 4, br 2, $|\mathrm{Aut}|{=}120$, orb 6/5, hist-orb 3/2; M 13312, res 10280, rev 1574: ok}\\*
0 & -- & 0 & 7 & 1 & 22 & 2 & 1000 & \texttt{-?????????}\\*
1 & 0 & 1 & 2 & 1 & 384 & 128 & 512 & \texttt{--???-----}\\
2 & 1 & 2 & 1 & 1 & 512 & 512 & 0 & \texttt{---+------}\\
3 & 2 & 3 & 1 & 1 & 0 & 1024 & 0 & \texttt{----------}\\
4 & 1 & 2 & 1 & 1 & 0 & 1024 & 0 & \texttt{----------}\\
5 & 0 & 1 & 1 & 5 & 0 & 1024 & 0 & \texttt{----------}\\
6 & 0 & 1 & 1 & 5 & 0 & 1024 & 0 & \texttt{----------}\\
7 & 0 & 1 & 1 & 5 & 0 & 1024 & 0 & \texttt{----------}\\
8 & 0 & 1 & 1 & 5 & 0 & 1024 & 0 & \texttt{----------}\\
9 & 0 & 1 & 1 & 5 & 0 & 1024 & 0 & \texttt{----------}\\
\multicolumn{9}{@{}l@{}}{\rule{0pt}{2.6ex}\textbf{T10.675}\enspace \texttt{0123123123}\enspace $n{=}10$, $|H|{=}3$, len 4, br 1, $|\mathrm{Aut}|{=}6$, orb 4/1, hist-orb 1/0; M 18432, res 3024, rev 7842: ok}\\*
0 & -- & 0 & 3 & 1 & 686 & 2 & 336 & \texttt{-?????????}\\*
1 & 0 & 1 & 1 & 3 & 768 & 256 & 0 & \texttt{--++------}\\
2 & 1 & 2 & 1 & 3 & 512 & 512 & 0 & \texttt{---+------}\\
3 & 2 & 3 & 1 & 3 & 0 & 1024 & 0 & \texttt{----------}\\
4 & 0 & 1 & 1 & 3 & 768 & 256 & 0 & \texttt{-----++---}\\
5 & 4 & 2 & 1 & 3 & 512 & 512 & 0 & \texttt{------+---}\\
6 & 5 & 3 & 1 & 3 & 0 & 1024 & 0 & \texttt{----------}\\
7 & 0 & 1 & 1 & 3 & 768 & 256 & 0 & \texttt{--------++}\\
8 & 7 & 2 & 1 & 3 & 512 & 512 & 0 & \texttt{---------+}\\
9 & 8 & 3 & 1 & 3 & 0 & 1024 & 0 & \texttt{----------}\\
\multicolumn{9}{@{}l@{}}{\rule{0pt}{2.6ex}\textbf{T10.676}\enspace \texttt{0123123122}\enspace $n{=}10$, $|H|{=}4$, len 4, br 2, $|\mathrm{Aut}|{=}4$, orb 6/2, hist-orb 2/0; M 17408, res 5556, rev 6202: ok}\\*
0 & -- & 0 & 4 & 1 & 490 & 2 & 532 & \texttt{-?????????}\\*
1 & 0 & 1 & 1 & 2 & 768 & 256 & 0 & \texttt{--++------}\\
2 & 1 & 2 & 1 & 2 & 512 & 512 & 0 & \texttt{---+------}\\
3 & 2 & 3 & 1 & 2 & 0 & 1024 & 0 & \texttt{----------}\\
4 & 0 & 1 & 1 & 2 & 768 & 256 & 0 & \texttt{-----++---}\\
5 & 4 & 2 & 1 & 2 & 512 & 512 & 0 & \texttt{------+---}\\
6 & 5 & 3 & 1 & 2 & 0 & 1024 & 0 & \texttt{----------}\\
7 & 0 & 1 & 2 & 1 & 256 & 256 & 512 & \texttt{--------??}\\
8 & 7 & 2 & 1 & 2 & 0 & 1024 & 0 & \texttt{----------}\\
9 & 7 & 2 & 1 & 2 & 0 & 1024 & 0 & \texttt{----------}\\
\multicolumn{9}{@{}l@{}}{\rule{0pt}{2.6ex}\textbf{T10.677}\enspace \texttt{0123123121}\enspace $n{=}10$, $|H|{=}4$, len 4, br 1, $|\mathrm{Aut}|{=}2$, orb 7/4, hist-orb 3/2; M 16384, res 6552, rev 5194: ok}\\*
0 & -- & 0 & 4 & 1 & 294 & 2 & 728 & \texttt{-?????????}\\*
1 & 0 & 1 & 1 & 2 & 768 & 256 & 0 & \texttt{--++------}\\
2 & 1 & 2 & 1 & 2 & 512 & 512 & 0 & \texttt{---+------}\\
3 & 2 & 3 & 1 & 2 & 0 & 1024 & 0 & \texttt{----------}\\
4 & 0 & 1 & 1 & 2 & 768 & 256 & 0 & \texttt{-----++---}\\
5 & 4 & 2 & 1 & 2 & 512 & 512 & 0 & \texttt{------+---}\\
6 & 5 & 3 & 1 & 2 & 0 & 1024 & 0 & \texttt{----------}\\
7 & 0 & 1 & 1 & 1 & 512 & 512 & 0 & \texttt{--------+-}\\
8 & 7 & 2 & 1 & 1 & 0 & 1024 & 0 & \texttt{----------}\\
9 & 0 & 1 & 1 & 1 & 0 & 1024 & 0 & \texttt{----------}\\
\multicolumn{9}{@{}l@{}}{\rule{0pt}{2.6ex}\textbf{T10.678}\enspace \texttt{0123123111}\enspace $n{=}10$, $|H|{=}5$, len 4, br 1, $|\mathrm{Aut}|{=}12$, orb 5/1, hist-orb 2/0; M 15360, res 8316, rev 3674: ok}\\*
0 & -- & 0 & 5 & 1 & 98 & 2 & 924 & \texttt{-?????????}\\*
1 & 0 & 1 & 1 & 2 & 768 & 256 & 0 & \texttt{--++------}\\
2 & 1 & 2 & 1 & 2 & 512 & 512 & 0 & \texttt{---+------}\\
3 & 2 & 3 & 1 & 2 & 0 & 1024 & 0 & \texttt{----------}\\
4 & 0 & 1 & 1 & 2 & 768 & 256 & 0 & \texttt{-----++---}\\
5 & 4 & 2 & 1 & 2 & 512 & 512 & 0 & \texttt{------+---}\\
6 & 5 & 3 & 1 & 2 & 0 & 1024 & 0 & \texttt{----------}\\
7 & 0 & 1 & 1 & 3 & 0 & 1024 & 0 & \texttt{----------}\\
8 & 0 & 1 & 1 & 3 & 0 & 1024 & 0 & \texttt{----------}\\
9 & 0 & 1 & 1 & 3 & 0 & 1024 & 0 & \texttt{----------}\\
\multicolumn{9}{@{}l@{}}{\rule{0pt}{2.6ex}\textbf{T10.679}\enspace \texttt{0123122222}\enspace $n{=}10$, $|H|{=}6$, len 4, br 2, $|\mathrm{Aut}|{=}120$, orb 6/5, hist-orb 2/1; M 17408, res 9360, rev 4746: ok}\\*
0 & -- & 0 & 6 & 1 & 462 & 2 & 560 & \texttt{-?????????}\\*
1 & 0 & 1 & 1 & 1 & 768 & 256 & 0 & \texttt{--++------}\\
2 & 1 & 2 & 1 & 1 & 512 & 512 & 0 & \texttt{---+------}\\
3 & 2 & 3 & 1 & 1 & 0 & 1024 & 0 & \texttt{----------}\\
4 & 0 & 1 & 5 & 1 & 32 & 32 & 960 & \texttt{-----?????}\\
5 & 4 & 2 & 1 & 5 & 0 & 1024 & 0 & \texttt{----------}\\
6 & 4 & 2 & 1 & 5 & 0 & 1024 & 0 & \texttt{----------}\\
7 & 4 & 2 & 1 & 5 & 0 & 1024 & 0 & \texttt{----------}\\
8 & 4 & 2 & 1 & 5 & 0 & 1024 & 0 & \texttt{----------}\\
9 & 4 & 2 & 1 & 5 & 0 & 1024 & 0 & \texttt{----------}\\
\multicolumn{9}{@{}l@{}}{\rule{0pt}{2.6ex}\textbf{T10.680}\enspace \texttt{0123122221}\enspace $n{=}10$, $|H|{=}6$, len 4, br 2, $|\mathrm{Aut}|{=}24$, orb 7/6, hist-orb 3/2; M 16384, res 10192, rev 3370: ok}\\*
0 & -- & 0 & 6 & 1 & 238 & 2 & 784 & \texttt{-?????????}\\*
1 & 0 & 1 & 1 & 1 & 768 & 256 & 0 & \texttt{--++------}\\
2 & 1 & 2 & 1 & 1 & 512 & 512 & 0 & \texttt{---+------}\\
3 & 2 & 3 & 1 & 1 & 0 & 1024 & 0 & \texttt{----------}\\
4 & 0 & 1 & 4 & 1 & 64 & 64 & 896 & \texttt{-----????-}\\
5 & 4 & 2 & 1 & 4 & 0 & 1024 & 0 & \texttt{----------}\\
6 & 4 & 2 & 1 & 4 & 0 & 1024 & 0 & \texttt{----------}\\
7 & 4 & 2 & 1 & 4 & 0 & 1024 & 0 & \texttt{----------}\\
8 & 4 & 2 & 1 & 4 & 0 & 1024 & 0 & \texttt{----------}\\
9 & 0 & 1 & 1 & 1 & 0 & 1024 & 0 & \texttt{----------}\\
\multicolumn{9}{@{}l@{}}{\rule{0pt}{2.6ex}\textbf{T10.681}\enspace \texttt{0123122212}\enspace $n{=}10$, $|H|{=}5$, len 4, br 2, $|\mathrm{Aut}|{=}6$, orb 8/7, hist-orb 3/2; M 16384, res 7716, rev 4706: ok}\\*
0 & -- & 0 & 5 & 1 & 378 & 2 & 644 & \texttt{-?????????}\\*
1 & 0 & 1 & 1 & 1 & 768 & 256 & 0 & \texttt{--++------}\\
2 & 1 & 2 & 1 & 1 & 512 & 512 & 0 & \texttt{---+------}\\
3 & 2 & 3 & 1 & 1 & 0 & 1024 & 0 & \texttt{----------}\\
4 & 0 & 1 & 3 & 1 & 128 & 128 & 768 & \texttt{-----???--}\\
5 & 4 & 2 & 1 & 3 & 0 & 1024 & 0 & \texttt{----------}\\
6 & 4 & 2 & 1 & 3 & 0 & 1024 & 0 & \texttt{----------}\\
7 & 4 & 2 & 1 & 3 & 0 & 1024 & 0 & \texttt{----------}\\
8 & 0 & 1 & 1 & 1 & 512 & 512 & 0 & \texttt{---------+}\\
9 & 8 & 2 & 1 & 1 & 0 & 1024 & 0 & \texttt{----------}\\
\multicolumn{9}{@{}l@{}}{\rule{0pt}{2.6ex}\textbf{T10.682}\enspace \texttt{0123122211}\enspace $n{=}10$, $|H|{=}6$, len 4, br 2, $|\mathrm{Aut}|{=}12$, orb 7/5, hist-orb 3/1; M 15360, res 9984, rev 2762: ok}\\*
0 & -- & 0 & 6 & 1 & 126 & 2 & 896 & \texttt{-?????????}\\*
1 & 0 & 1 & 1 & 1 & 768 & 256 & 0 & \texttt{--++------}\\
2 & 1 & 2 & 1 & 1 & 512 & 512 & 0 & \texttt{---+------}\\
3 & 2 & 3 & 1 & 1 & 0 & 1024 & 0 & \texttt{----------}\\
4 & 0 & 1 & 3 & 1 & 128 & 128 & 768 & \texttt{-----???--}\\
5 & 4 & 2 & 1 & 3 & 0 & 1024 & 0 & \texttt{----------}\\
6 & 4 & 2 & 1 & 3 & 0 & 1024 & 0 & \texttt{----------}\\
7 & 4 & 2 & 1 & 3 & 0 & 1024 & 0 & \texttt{----------}\\
8 & 0 & 1 & 1 & 2 & 0 & 1024 & 0 & \texttt{----------}\\
9 & 0 & 1 & 1 & 2 & 0 & 1024 & 0 & \texttt{----------}\\
\multicolumn{9}{@{}l@{}}{\rule{0pt}{2.6ex}\textbf{T10.683}\enspace \texttt{0123122122}\enspace $n{=}10$, $|H|{=}5$, len 4, br 3, $|\mathrm{Aut}|{=}8$, orb 6/4, hist-orb 2/1; M 16384, res 7584, rev 4650: ok}\\*
0 & -- & 0 & 5 & 1 & 350 & 2 & 672 & \texttt{-?????????}\\*
1 & 0 & 1 & 1 & 1 & 768 & 256 & 0 & \texttt{--++------}\\
2 & 1 & 2 & 1 & 1 & 512 & 512 & 0 & \texttt{---+------}\\
3 & 2 & 3 & 1 & 1 & 0 & 1024 & 0 & \texttt{----------}\\
4 & 0 & 1 & 2 & 2 & 256 & 256 & 512 & \texttt{-----??---}\\
5 & 4 & 2 & 1 & 4 & 0 & 1024 & 0 & \texttt{----------}\\
6 & 4 & 2 & 1 & 4 & 0 & 1024 & 0 & \texttt{----------}\\
7 & 0 & 1 & 2 & 2 & 256 & 256 & 512 & \texttt{--------??}\\
8 & 7 & 2 & 1 & 4 & 0 & 1024 & 0 & \texttt{----------}\\
9 & 7 & 2 & 1 & 4 & 0 & 1024 & 0 & \texttt{----------}\\
\multicolumn{9}{@{}l@{}}{\rule{0pt}{2.6ex}\textbf{T10.684}\enspace \texttt{0123122121}\enspace $n{=}10$, $|H|{=}5$, len 4, br 2, $|\mathrm{Aut}|{=}2$, orb 9/8, hist-orb 4/3; M 15360, res 8076, rev 3842: ok}\\*
0 & -- & 0 & 5 & 1 & 210 & 2 & 812 & \texttt{-?????????}\\*
1 & 0 & 1 & 1 & 1 & 768 & 256 & 0 & \texttt{--++------}\\
2 & 1 & 2 & 1 & 1 & 512 & 512 & 0 & \texttt{---+------}\\
3 & 2 & 3 & 1 & 1 & 0 & 1024 & 0 & \texttt{----------}\\
4 & 0 & 1 & 2 & 1 & 256 & 256 & 512 & \texttt{-----??---}\\
5 & 4 & 2 & 1 & 2 & 0 & 1024 & 0 & \texttt{----------}\\
6 & 4 & 2 & 1 & 2 & 0 & 1024 & 0 & \texttt{----------}\\
7 & 0 & 1 & 1 & 1 & 512 & 512 & 0 & \texttt{--------+-}\\
8 & 7 & 2 & 1 & 1 & 0 & 1024 & 0 & \texttt{----------}\\
9 & 0 & 1 & 1 & 1 & 0 & 1024 & 0 & \texttt{----------}\\
\multicolumn{9}{@{}l@{}}{\rule{0pt}{2.6ex}\textbf{T10.685}\enspace \texttt{0123122111}\enspace $n{=}10$, $|H|{=}6$, len 4, br 2, $|\mathrm{Aut}|{=}12$, orb 7/5, hist-orb 3/1; M 14336, res 9336, rev 2522: ok}\\*
0 & -- & 0 & 6 & 1 & 70 & 2 & 952 & \texttt{-?????????}\\*
1 & 0 & 1 & 1 & 1 & 768 & 256 & 0 & \texttt{--++------}\\
2 & 1 & 2 & 1 & 1 & 512 & 512 & 0 & \texttt{---+------}\\
3 & 2 & 3 & 1 & 1 & 0 & 1024 & 0 & \texttt{----------}\\
4 & 0 & 1 & 2 & 1 & 256 & 256 & 512 & \texttt{-----??---}\\
5 & 4 & 2 & 1 & 2 & 0 & 1024 & 0 & \texttt{----------}\\
6 & 4 & 2 & 1 & 2 & 0 & 1024 & 0 & \texttt{----------}\\
7 & 0 & 1 & 1 & 3 & 0 & 1024 & 0 & \texttt{----------}\\
8 & 0 & 1 & 1 & 3 & 0 & 1024 & 0 & \texttt{----------}\\
9 & 0 & 1 & 1 & 3 & 0 & 1024 & 0 & \texttt{----------}\\
\multicolumn{9}{@{}l@{}}{\rule{0pt}{2.6ex}\textbf{T10.686}\enspace \texttt{0123121212}\enspace $n{=}10$, $|H|{=}4$, len 4, br 1, $|\mathrm{Aut}|{=}6$, orb 6/4, hist-orb 2/1; M 15360, res 5796, rev 5178: ok}\\*
0 & -- & 0 & 4 & 1 & 378 & 2 & 644 & \texttt{-?????????}\\*
1 & 0 & 1 & 1 & 1 & 768 & 256 & 0 & \texttt{--++------}\\
2 & 1 & 2 & 1 & 1 & 512 & 512 & 0 & \texttt{---+------}\\
3 & 2 & 3 & 1 & 1 & 0 & 1024 & 0 & \texttt{----------}\\
4 & 0 & 1 & 1 & 3 & 512 & 512 & 0 & \texttt{-----+----}\\
5 & 4 & 2 & 1 & 3 & 0 & 1024 & 0 & \texttt{----------}\\
6 & 0 & 1 & 1 & 3 & 512 & 512 & 0 & \texttt{-------+--}\\
7 & 6 & 2 & 1 & 3 & 0 & 1024 & 0 & \texttt{----------}\\
8 & 0 & 1 & 1 & 3 & 512 & 512 & 0 & \texttt{---------+}\\
9 & 8 & 2 & 1 & 3 & 0 & 1024 & 0 & \texttt{----------}\\
\multicolumn{9}{@{}l@{}}{\rule{0pt}{2.6ex}\textbf{T10.687}\enspace \texttt{0123121211}\enspace $n{=}10$, $|H|{=}5$, len 4, br 1, $|\mathrm{Aut}|{=}4$, orb 7/4, hist-orb 3/1; M 14336, res 8064, rev 3346: ok}\\*
0 & -- & 0 & 5 & 1 & 126 & 2 & 896 & \texttt{-?????????}\\*
1 & 0 & 1 & 1 & 1 & 768 & 256 & 0 & \texttt{--++------}\\
2 & 1 & 2 & 1 & 1 & 512 & 512 & 0 & \texttt{---+------}\\
3 & 2 & 3 & 1 & 1 & 0 & 1024 & 0 & \texttt{----------}\\
4 & 0 & 1 & 1 & 2 & 512 & 512 & 0 & \texttt{-----+----}\\
5 & 4 & 2 & 1 & 2 & 0 & 1024 & 0 & \texttt{----------}\\
6 & 0 & 1 & 1 & 2 & 512 & 512 & 0 & \texttt{-------+--}\\
7 & 6 & 2 & 1 & 2 & 0 & 1024 & 0 & \texttt{----------}\\
8 & 0 & 1 & 1 & 2 & 0 & 1024 & 0 & \texttt{----------}\\
9 & 0 & 1 & 1 & 2 & 0 & 1024 & 0 & \texttt{----------}\\
\multicolumn{9}{@{}l@{}}{\rule{0pt}{2.6ex}\textbf{T10.688}\enspace \texttt{0123121111}\enspace $n{=}10$, $|H|{=}6$, len 4, br 1, $|\mathrm{Aut}|{=}24$, orb 7/6, hist-orb 3/2; M 13312, res 8820, rev 2338: ok}\\*
0 & -- & 0 & 6 & 1 & 42 & 2 & 980 & \texttt{-?????????}\\*
1 & 0 & 1 & 1 & 1 & 768 & 256 & 0 & \texttt{--++------}\\
2 & 1 & 2 & 1 & 1 & 512 & 512 & 0 & \texttt{---+------}\\
3 & 2 & 3 & 1 & 1 & 0 & 1024 & 0 & \texttt{----------}\\
4 & 0 & 1 & 1 & 1 & 512 & 512 & 0 & \texttt{-----+----}\\
5 & 4 & 2 & 1 & 1 & 0 & 1024 & 0 & \texttt{----------}\\
6 & 0 & 1 & 1 & 4 & 0 & 1024 & 0 & \texttt{----------}\\
7 & 0 & 1 & 1 & 4 & 0 & 1024 & 0 & \texttt{----------}\\
8 & 0 & 1 & 1 & 4 & 0 & 1024 & 0 & \texttt{----------}\\
9 & 0 & 1 & 1 & 4 & 0 & 1024 & 0 & \texttt{----------}\\
\multicolumn{9}{@{}l@{}}{\rule{0pt}{2.6ex}\textbf{T10.689}\enspace \texttt{0123111111}\enspace $n{=}10$, $|H|{=}7$, len 4, br 1, $|\mathrm{Aut}|{=}720$, orb 5/4, hist-orb 2/1; M 12288, res 9072, rev 1642: ok}\\*
0 & -- & 0 & 7 & 1 & 14 & 2 & 1008 & \texttt{-?????????}\\*
1 & 0 & 1 & 1 & 1 & 768 & 256 & 0 & \texttt{--++------}\\
2 & 1 & 2 & 1 & 1 & 512 & 512 & 0 & \texttt{---+------}\\
3 & 2 & 3 & 1 & 1 & 0 & 1024 & 0 & \texttt{----------}\\
4 & 0 & 1 & 1 & 6 & 0 & 1024 & 0 & \texttt{----------}\\
5 & 0 & 1 & 1 & 6 & 0 & 1024 & 0 & \texttt{----------}\\
6 & 0 & 1 & 1 & 6 & 0 & 1024 & 0 & \texttt{----------}\\
7 & 0 & 1 & 1 & 6 & 0 & 1024 & 0 & \texttt{----------}\\
8 & 0 & 1 & 1 & 6 & 0 & 1024 & 0 & \texttt{----------}\\
9 & 0 & 1 & 1 & 6 & 0 & 1024 & 0 & \texttt{----------}\\
\multicolumn{9}{@{}l@{}}{\rule{0pt}{2.6ex}\textbf{T10.690}\enspace \texttt{0122222222}\enspace $n{=}10$, $|H|{=}8$, len 3, br 1, $|\mathrm{Aut}|{=}40320$, orb 3/2, hist-orb 1/0; M 17408, res 12192, rev 4146: ok}\\*
0 & -- & 0 & 8 & 1 & 514 & 2 & 508 & \texttt{-+????????}\\*
1 & 0 & 1 & 8 & 1 & 4 & 4 & 1016 & \texttt{--????????}\\
2 & 1 & 2 & 1 & 8 & 0 & 1024 & 0 & \texttt{----------}\\
3 & 1 & 2 & 1 & 8 & 0 & 1024 & 0 & \texttt{----------}\\
4 & 1 & 2 & 1 & 8 & 0 & 1024 & 0 & \texttt{----------}\\
5 & 1 & 2 & 1 & 8 & 0 & 1024 & 0 & \texttt{----------}\\
6 & 1 & 2 & 1 & 8 & 0 & 1024 & 0 & \texttt{----------}\\
7 & 1 & 2 & 1 & 8 & 0 & 1024 & 0 & \texttt{----------}\\
8 & 1 & 2 & 1 & 8 & 0 & 1024 & 0 & \texttt{----------}\\
9 & 1 & 2 & 1 & 8 & 0 & 1024 & 0 & \texttt{----------}\\
\multicolumn{9}{@{}l@{}}{\rule{0pt}{2.6ex}\textbf{T10.691}\enspace \texttt{0122222221}\enspace $n{=}10$, $|H|{=}8$, len 3, br 2, $|\mathrm{Aut}|{=}5040$, orb 4/3, hist-orb 2/1; M 16384, res 13428, rev 2122: ok}\\*
0 & -- & 0 & 8 & 1 & 258 & 2 & 764 & \texttt{-?????????}\\*
1 & 0 & 1 & 7 & 1 & 8 & 8 & 1008 & \texttt{--???????-}\\
2 & 1 & 2 & 1 & 7 & 0 & 1024 & 0 & \texttt{----------}\\
3 & 1 & 2 & 1 & 7 & 0 & 1024 & 0 & \texttt{----------}\\
4 & 1 & 2 & 1 & 7 & 0 & 1024 & 0 & \texttt{----------}\\
5 & 1 & 2 & 1 & 7 & 0 & 1024 & 0 & \texttt{----------}\\
6 & 1 & 2 & 1 & 7 & 0 & 1024 & 0 & \texttt{----------}\\
7 & 1 & 2 & 1 & 7 & 0 & 1024 & 0 & \texttt{----------}\\
8 & 1 & 2 & 1 & 7 & 0 & 1024 & 0 & \texttt{----------}\\
9 & 0 & 1 & 1 & 1 & 0 & 1024 & 0 & \texttt{----------}\\
\multicolumn{9}{@{}l@{}}{\rule{0pt}{2.6ex}\textbf{T10.692}\enspace \texttt{0122222212}\enspace $n{=}10$, $|H|{=}7$, len 3, br 2, $|\mathrm{Aut}|{=}720$, orb 5/4, hist-orb 2/1; M 16384, res 11144, rev 3474: ok}\\*
0 & -- & 0 & 7 & 1 & 390 & 2 & 632 & \texttt{-?????????}\\*
1 & 0 & 1 & 6 & 1 & 16 & 16 & 992 & \texttt{--??????--}\\
2 & 1 & 2 & 1 & 6 & 0 & 1024 & 0 & \texttt{----------}\\
3 & 1 & 2 & 1 & 6 & 0 & 1024 & 0 & \texttt{----------}\\
4 & 1 & 2 & 1 & 6 & 0 & 1024 & 0 & \texttt{----------}\\
5 & 1 & 2 & 1 & 6 & 0 & 1024 & 0 & \texttt{----------}\\
6 & 1 & 2 & 1 & 6 & 0 & 1024 & 0 & \texttt{----------}\\
7 & 1 & 2 & 1 & 6 & 0 & 1024 & 0 & \texttt{----------}\\
8 & 0 & 1 & 1 & 1 & 512 & 512 & 0 & \texttt{---------+}\\
9 & 8 & 2 & 1 & 1 & 0 & 1024 & 0 & \texttt{----------}\\
\multicolumn{9}{@{}l@{}}{\rule{0pt}{2.6ex}\textbf{T10.693}\enspace \texttt{0122222211}\enspace $n{=}10$, $|H|{=}8$, len 3, br 2, $|\mathrm{Aut}|{=}1440$, orb 4/2, hist-orb 2/0; M 15360, res 13484, rev 1138: ok}\\*
0 & -- & 0 & 8 & 1 & 130 & 2 & 892 & \texttt{-?????????}\\*
1 & 0 & 1 & 6 & 1 & 16 & 16 & 992 & \texttt{--??????--}\\
2 & 1 & 2 & 1 & 6 & 0 & 1024 & 0 & \texttt{----------}\\
3 & 1 & 2 & 1 & 6 & 0 & 1024 & 0 & \texttt{----------}\\
4 & 1 & 2 & 1 & 6 & 0 & 1024 & 0 & \texttt{----------}\\
5 & 1 & 2 & 1 & 6 & 0 & 1024 & 0 & \texttt{----------}\\
6 & 1 & 2 & 1 & 6 & 0 & 1024 & 0 & \texttt{----------}\\
7 & 1 & 2 & 1 & 6 & 0 & 1024 & 0 & \texttt{----------}\\
8 & 0 & 1 & 1 & 2 & 0 & 1024 & 0 & \texttt{----------}\\
9 & 0 & 1 & 1 & 2 & 0 & 1024 & 0 & \texttt{----------}\\
\multicolumn{9}{@{}l@{}}{\rule{0pt}{2.6ex}\textbf{T10.694}\enspace \texttt{0122222122}\enspace $n{=}10$, $|H|{=}7$, len 3, br 3, $|\mathrm{Aut}|{=}240$, orb 5/3, hist-orb 2/0; M 16384, res 11316, rev 3058: ok}\\*
0 & -- & 0 & 7 & 1 & 330 & 2 & 692 & \texttt{-?????????}\\*
1 & 0 & 1 & 5 & 1 & 32 & 32 & 960 & \texttt{--?????---}\\
2 & 1 & 2 & 1 & 5 & 0 & 1024 & 0 & \texttt{----------}\\
3 & 1 & 2 & 1 & 5 & 0 & 1024 & 0 & \texttt{----------}\\
4 & 1 & 2 & 1 & 5 & 0 & 1024 & 0 & \texttt{----------}\\
5 & 1 & 2 & 1 & 5 & 0 & 1024 & 0 & \texttt{----------}\\
6 & 1 & 2 & 1 & 5 & 0 & 1024 & 0 & \texttt{----------}\\
7 & 0 & 1 & 2 & 1 & 256 & 256 & 512 & \texttt{--------??}\\
8 & 7 & 2 & 1 & 2 & 0 & 1024 & 0 & \texttt{----------}\\
9 & 7 & 2 & 1 & 2 & 0 & 1024 & 0 & \texttt{----------}\\
\multicolumn{9}{@{}l@{}}{\rule{0pt}{2.6ex}\textbf{T10.695}\enspace \texttt{0122222121}\enspace $n{=}10$, $|H|{=}7$, len 3, br 2, $|\mathrm{Aut}|{=}120$, orb 6/5, hist-orb 3/2; M 15360, res 11736, rev 2130: ok}\\*
0 & -- & 0 & 7 & 1 & 198 & 2 & 824 & \texttt{-?????????}\\*
1 & 0 & 1 & 5 & 1 & 32 & 32 & 960 & \texttt{--?????---}\\
2 & 1 & 2 & 1 & 5 & 0 & 1024 & 0 & \texttt{----------}\\
3 & 1 & 2 & 1 & 5 & 0 & 1024 & 0 & \texttt{----------}\\
4 & 1 & 2 & 1 & 5 & 0 & 1024 & 0 & \texttt{----------}\\
5 & 1 & 2 & 1 & 5 & 0 & 1024 & 0 & \texttt{----------}\\
6 & 1 & 2 & 1 & 5 & 0 & 1024 & 0 & \texttt{----------}\\
7 & 0 & 1 & 1 & 1 & 512 & 512 & 0 & \texttt{--------+-}\\
8 & 7 & 2 & 1 & 1 & 0 & 1024 & 0 & \texttt{----------}\\
9 & 0 & 1 & 1 & 1 & 0 & 1024 & 0 & \texttt{----------}\\
\multicolumn{9}{@{}l@{}}{\rule{0pt}{2.6ex}\textbf{T10.696}\enspace \texttt{0122222111}\enspace $n{=}10$, $|H|{=}8$, len 3, br 2, $|\mathrm{Aut}|{=}720$, orb 4/2, hist-orb 2/0; M 14336, res 12924, rev 690: ok}\\*
0 & -- & 0 & 8 & 1 & 66 & 2 & 956 & \texttt{-?????????}\\*
1 & 0 & 1 & 5 & 1 & 32 & 32 & 960 & \texttt{--?????---}\\
2 & 1 & 2 & 1 & 5 & 0 & 1024 & 0 & \texttt{----------}\\
3 & 1 & 2 & 1 & 5 & 0 & 1024 & 0 & \texttt{----------}\\
4 & 1 & 2 & 1 & 5 & 0 & 1024 & 0 & \texttt{----------}\\
5 & 1 & 2 & 1 & 5 & 0 & 1024 & 0 & \texttt{----------}\\
6 & 1 & 2 & 1 & 5 & 0 & 1024 & 0 & \texttt{----------}\\
7 & 0 & 1 & 1 & 3 & 0 & 1024 & 0 & \texttt{----------}\\
8 & 0 & 1 & 1 & 3 & 0 & 1024 & 0 & \texttt{----------}\\
9 & 0 & 1 & 1 & 3 & 0 & 1024 & 0 & \texttt{----------}\\
\multicolumn{9}{@{}l@{}}{\rule{0pt}{2.6ex}\textbf{T10.697}\enspace \texttt{0122221222}\enspace $n{=}10$, $|H|{=}7$, len 3, br 3, $|\mathrm{Aut}|{=}144$, orb 5/3, hist-orb 2/0; M 16384, res 11500, rev 2834: ok}\\*
0 & -- & 0 & 7 & 1 & 306 & 2 & 716 & \texttt{-?????????}\\*
1 & 0 & 1 & 4 & 1 & 64 & 64 & 896 & \texttt{--????----}\\
2 & 1 & 2 & 1 & 4 & 0 & 1024 & 0 & \texttt{----------}\\
3 & 1 & 2 & 1 & 4 & 0 & 1024 & 0 & \texttt{----------}\\
4 & 1 & 2 & 1 & 4 & 0 & 1024 & 0 & \texttt{----------}\\
5 & 1 & 2 & 1 & 4 & 0 & 1024 & 0 & \texttt{----------}\\
6 & 0 & 1 & 3 & 1 & 128 & 128 & 768 & \texttt{-------???}\\
7 & 6 & 2 & 1 & 3 & 0 & 1024 & 0 & \texttt{----------}\\
8 & 6 & 2 & 1 & 3 & 0 & 1024 & 0 & \texttt{----------}\\
9 & 6 & 2 & 1 & 3 & 0 & 1024 & 0 & \texttt{----------}\\
\multicolumn{9}{@{}l@{}}{\rule{0pt}{2.6ex}\textbf{T10.698}\enspace \texttt{0122221221}\enspace $n{=}10$, $|H|{=}7$, len 3, br 3, $|\mathrm{Aut}|{=}48$, orb 6/4, hist-orb 3/1; M 15360, res 11572, rev 2002: ok}\\*
0 & -- & 0 & 7 & 1 & 170 & 2 & 852 & \texttt{-?????????}\\*
1 & 0 & 1 & 4 & 1 & 64 & 64 & 896 & \texttt{--????----}\\
2 & 1 & 2 & 1 & 4 & 0 & 1024 & 0 & \texttt{----------}\\
3 & 1 & 2 & 1 & 4 & 0 & 1024 & 0 & \texttt{----------}\\
4 & 1 & 2 & 1 & 4 & 0 & 1024 & 0 & \texttt{----------}\\
5 & 1 & 2 & 1 & 4 & 0 & 1024 & 0 & \texttt{----------}\\
6 & 0 & 1 & 2 & 1 & 256 & 256 & 512 & \texttt{-------??-}\\
7 & 6 & 2 & 1 & 2 & 0 & 1024 & 0 & \texttt{----------}\\
8 & 6 & 2 & 1 & 2 & 0 & 1024 & 0 & \texttt{----------}\\
9 & 0 & 1 & 1 & 1 & 0 & 1024 & 0 & \texttt{----------}\\
\multicolumn{9}{@{}l@{}}{\rule{0pt}{2.6ex}\textbf{T10.699}\enspace \texttt{0122221212}\enspace $n{=}10$, $|H|{=}6$, len 3, br 2, $|\mathrm{Aut}|{=}48$, orb 5/2, hist-orb 2/0; M 15360, res 9580, rev 3338: ok}\\*
0 & -- & 0 & 6 & 1 & 306 & 2 & 716 & \texttt{-?????????}\\*
1 & 0 & 1 & 4 & 1 & 64 & 64 & 896 & \texttt{--????----}\\
2 & 1 & 2 & 1 & 4 & 0 & 1024 & 0 & \texttt{----------}\\
3 & 1 & 2 & 1 & 4 & 0 & 1024 & 0 & \texttt{----------}\\
4 & 1 & 2 & 1 & 4 & 0 & 1024 & 0 & \texttt{----------}\\
5 & 1 & 2 & 1 & 4 & 0 & 1024 & 0 & \texttt{----------}\\
6 & 0 & 1 & 1 & 2 & 512 & 512 & 0 & \texttt{-------+--}\\
7 & 6 & 2 & 1 & 2 & 0 & 1024 & 0 & \texttt{----------}\\
8 & 0 & 1 & 1 & 2 & 512 & 512 & 0 & \texttt{---------+}\\
9 & 8 & 2 & 1 & 2 & 0 & 1024 & 0 & \texttt{----------}\\
\multicolumn{9}{@{}l@{}}{\rule{0pt}{2.6ex}\textbf{T10.700}\enspace \texttt{0122221211}\enspace $n{=}10$, $|H|{=}7$, len 3, br 2, $|\mathrm{Aut}|{=}48$, orb 6/4, hist-orb 3/1; M 14336, res 11416, rev 1522: ok}\\*
0 & -- & 0 & 7 & 1 & 102 & 2 & 920 & \texttt{-?????????}\\*
1 & 0 & 1 & 4 & 1 & 64 & 64 & 896 & \texttt{--????----}\\
2 & 1 & 2 & 1 & 4 & 0 & 1024 & 0 & \texttt{----------}\\
3 & 1 & 2 & 1 & 4 & 0 & 1024 & 0 & \texttt{----------}\\
4 & 1 & 2 & 1 & 4 & 0 & 1024 & 0 & \texttt{----------}\\
5 & 1 & 2 & 1 & 4 & 0 & 1024 & 0 & \texttt{----------}\\
6 & 0 & 1 & 1 & 1 & 512 & 512 & 0 & \texttt{-------+--}\\
7 & 6 & 2 & 1 & 1 & 0 & 1024 & 0 & \texttt{----------}\\
8 & 0 & 1 & 1 & 2 & 0 & 1024 & 0 & \texttt{----------}\\
9 & 0 & 1 & 1 & 2 & 0 & 1024 & 0 & \texttt{----------}\\
\multicolumn{9}{@{}l@{}}{\rule{0pt}{2.6ex}\textbf{T10.701}\enspace \texttt{0122221111}\enspace $n{=}10$, $|H|{=}8$, len 3, br 2, $|\mathrm{Aut}|{=}576$, orb 4/2, hist-orb 2/0; M 13312, res 12028, rev 530: ok}\\*
0 & -- & 0 & 8 & 1 & 34 & 2 & 988 & \texttt{-?????????}\\*
1 & 0 & 1 & 4 & 1 & 64 & 64 & 896 & \texttt{--????----}\\
2 & 1 & 2 & 1 & 4 & 0 & 1024 & 0 & \texttt{----------}\\
3 & 1 & 2 & 1 & 4 & 0 & 1024 & 0 & \texttt{----------}\\
4 & 1 & 2 & 1 & 4 & 0 & 1024 & 0 & \texttt{----------}\\
5 & 1 & 2 & 1 & 4 & 0 & 1024 & 0 & \texttt{----------}\\
6 & 0 & 1 & 1 & 4 & 0 & 1024 & 0 & \texttt{----------}\\
7 & 0 & 1 & 1 & 4 & 0 & 1024 & 0 & \texttt{----------}\\
8 & 0 & 1 & 1 & 4 & 0 & 1024 & 0 & \texttt{----------}\\
9 & 0 & 1 & 1 & 4 & 0 & 1024 & 0 & \texttt{----------}\\
\multicolumn{9}{@{}l@{}}{\rule{0pt}{2.6ex}\textbf{T10.702}\enspace \texttt{0122212221}\enspace $n{=}10$, $|H|{=}7$, len 3, br 3, $|\mathrm{Aut}|{=}72$, orb 4/2, hist-orb 2/1; M 15360, res 11580, rev 1938: ok}\\*
0 & -- & 0 & 7 & 1 & 162 & 2 & 860 & \texttt{-?????????}\\*
1 & 0 & 1 & 3 & 2 & 128 & 128 & 768 & \texttt{--???-----}\\
2 & 1 & 2 & 1 & 6 & 0 & 1024 & 0 & \texttt{----------}\\
3 & 1 & 2 & 1 & 6 & 0 & 1024 & 0 & \texttt{----------}\\
4 & 1 & 2 & 1 & 6 & 0 & 1024 & 0 & \texttt{----------}\\
5 & 0 & 1 & 3 & 2 & 128 & 128 & 768 & \texttt{------???-}\\
6 & 5 & 2 & 1 & 6 & 0 & 1024 & 0 & \texttt{----------}\\
7 & 5 & 2 & 1 & 6 & 0 & 1024 & 0 & \texttt{----------}\\
8 & 5 & 2 & 1 & 6 & 0 & 1024 & 0 & \texttt{----------}\\
9 & 0 & 1 & 1 & 1 & 0 & 1024 & 0 & \texttt{----------}\\
\multicolumn{9}{@{}l@{}}{\rule{0pt}{2.6ex}\textbf{T10.703}\enspace \texttt{0122212212}\enspace $n{=}10$, $|H|{=}6$, len 3, br 3, $|\mathrm{Aut}|{=}12$, orb 7/5, hist-orb 3/1; M 15360, res 9456, rev 3202: ok}\\*
0 & -- & 0 & 6 & 1 & 270 & 2 & 752 & \texttt{-?????????}\\*
1 & 0 & 1 & 3 & 1 & 128 & 128 & 768 & \texttt{--???-----}\\
2 & 1 & 2 & 1 & 3 & 0 & 1024 & 0 & \texttt{----------}\\
3 & 1 & 2 & 1 & 3 & 0 & 1024 & 0 & \texttt{----------}\\
4 & 1 & 2 & 1 & 3 & 0 & 1024 & 0 & \texttt{----------}\\
5 & 0 & 1 & 2 & 1 & 256 & 256 & 512 & \texttt{------??--}\\
6 & 5 & 2 & 1 & 2 & 0 & 1024 & 0 & \texttt{----------}\\
7 & 5 & 2 & 1 & 2 & 0 & 1024 & 0 & \texttt{----------}\\
8 & 0 & 1 & 1 & 1 & 512 & 512 & 0 & \texttt{---------+}\\
9 & 8 & 2 & 1 & 1 & 0 & 1024 & 0 & \texttt{----------}\\
\multicolumn{9}{@{}l@{}}{\rule{0pt}{2.6ex}\textbf{T10.704}\enspace \texttt{0122212211}\enspace $n{=}10$, $|H|{=}7$, len 3, br 3, $|\mathrm{Aut}|{=}24$, orb 6/3, hist-orb 3/0; M 14336, res 11076, rev 1554: ok}\\*
0 & -- & 0 & 7 & 1 & 90 & 2 & 932 & \texttt{-?????????}\\*
1 & 0 & 1 & 3 & 1 & 128 & 128 & 768 & \texttt{--???-----}\\
2 & 1 & 2 & 1 & 3 & 0 & 1024 & 0 & \texttt{----------}\\
3 & 1 & 2 & 1 & 3 & 0 & 1024 & 0 & \texttt{----------}\\
4 & 1 & 2 & 1 & 3 & 0 & 1024 & 0 & \texttt{----------}\\
5 & 0 & 1 & 2 & 1 & 256 & 256 & 512 & \texttt{------??--}\\
6 & 5 & 2 & 1 & 2 & 0 & 1024 & 0 & \texttt{----------}\\
7 & 5 & 2 & 1 & 2 & 0 & 1024 & 0 & \texttt{----------}\\
8 & 0 & 1 & 1 & 2 & 0 & 1024 & 0 & \texttt{----------}\\
9 & 0 & 1 & 1 & 2 & 0 & 1024 & 0 & \texttt{----------}\\
\multicolumn{9}{@{}l@{}}{\rule{0pt}{2.6ex}\textbf{T10.705}\enspace \texttt{0122212121}\enspace $n{=}10$, $|H|{=}6$, len 3, br 2, $|\mathrm{Aut}|{=}12$, orb 6/3, hist-orb 3/1; M 14336, res 9660, rev 2506: ok}\\*
0 & -- & 0 & 6 & 1 & 162 & 2 & 860 & \texttt{-?????????}\\*
1 & 0 & 1 & 3 & 1 & 128 & 128 & 768 & \texttt{--???-----}\\
2 & 1 & 2 & 1 & 3 & 0 & 1024 & 0 & \texttt{----------}\\
3 & 1 & 2 & 1 & 3 & 0 & 1024 & 0 & \texttt{----------}\\
4 & 1 & 2 & 1 & 3 & 0 & 1024 & 0 & \texttt{----------}\\
5 & 0 & 1 & 1 & 2 & 512 & 512 & 0 & \texttt{------+---}\\
6 & 5 & 2 & 1 & 2 & 0 & 1024 & 0 & \texttt{----------}\\
7 & 0 & 1 & 1 & 2 & 512 & 512 & 0 & \texttt{--------+-}\\
8 & 7 & 2 & 1 & 2 & 0 & 1024 & 0 & \texttt{----------}\\
9 & 0 & 1 & 1 & 1 & 0 & 1024 & 0 & \texttt{----------}\\
\multicolumn{9}{@{}l@{}}{\rule{0pt}{2.6ex}\textbf{T10.706}\enspace \texttt{0122212111}\enspace $n{=}10$, $|H|{=}7$, len 3, br 2, $|\mathrm{Aut}|{=}36$, orb 6/4, hist-orb 3/1; M 13312, res 10632, rev 1298: ok}\\*
0 & -- & 0 & 7 & 1 & 54 & 2 & 968 & \texttt{-?????????}\\*
1 & 0 & 1 & 3 & 1 & 128 & 128 & 768 & \texttt{--???-----}\\
2 & 1 & 2 & 1 & 3 & 0 & 1024 & 0 & \texttt{----------}\\
3 & 1 & 2 & 1 & 3 & 0 & 1024 & 0 & \texttt{----------}\\
4 & 1 & 2 & 1 & 3 & 0 & 1024 & 0 & \texttt{----------}\\
5 & 0 & 1 & 1 & 1 & 512 & 512 & 0 & \texttt{------+---}\\
6 & 5 & 2 & 1 & 1 & 0 & 1024 & 0 & \texttt{----------}\\
7 & 0 & 1 & 1 & 3 & 0 & 1024 & 0 & \texttt{----------}\\
8 & 0 & 1 & 1 & 3 & 0 & 1024 & 0 & \texttt{----------}\\
9 & 0 & 1 & 1 & 3 & 0 & 1024 & 0 & \texttt{----------}\\
\multicolumn{9}{@{}l@{}}{\rule{0pt}{2.6ex}\textbf{T10.707}\enspace \texttt{0122211111}\enspace $n{=}10$, $|H|{=}8$, len 3, br 2, $|\mathrm{Aut}|{=}720$, orb 4/2, hist-orb 2/0; M 12288, res 10956, rev 530: ok}\\*
0 & -- & 0 & 8 & 1 & 18 & 2 & 1004 & \texttt{-?????????}\\*
1 & 0 & 1 & 3 & 1 & 128 & 128 & 768 & \texttt{--???-----}\\
2 & 1 & 2 & 1 & 3 & 0 & 1024 & 0 & \texttt{----------}\\
3 & 1 & 2 & 1 & 3 & 0 & 1024 & 0 & \texttt{----------}\\
4 & 1 & 2 & 1 & 3 & 0 & 1024 & 0 & \texttt{----------}\\
5 & 0 & 1 & 1 & 5 & 0 & 1024 & 0 & \texttt{----------}\\
6 & 0 & 1 & 1 & 5 & 0 & 1024 & 0 & \texttt{----------}\\
7 & 0 & 1 & 1 & 5 & 0 & 1024 & 0 & \texttt{----------}\\
8 & 0 & 1 & 1 & 5 & 0 & 1024 & 0 & \texttt{----------}\\
9 & 0 & 1 & 1 & 5 & 0 & 1024 & 0 & \texttt{----------}\\
\multicolumn{9}{@{}l@{}}{\rule{0pt}{2.6ex}\textbf{T10.708}\enspace \texttt{0122122122}\enspace $n{=}10$, $|H|{=}6$, len 3, br 4, $|\mathrm{Aut}|{=}48$, orb 3/1, hist-orb 1/0; M 15360, res 9252, rev 3186: ok}\\*
0 & -- & 0 & 6 & 1 & 250 & 2 & 772 & \texttt{-?????????}\\*
1 & 0 & 1 & 2 & 3 & 256 & 256 & 512 & \texttt{--??------}\\
2 & 1 & 2 & 1 & 6 & 0 & 1024 & 0 & \texttt{----------}\\
3 & 1 & 2 & 1 & 6 & 0 & 1024 & 0 & \texttt{----------}\\
4 & 0 & 1 & 2 & 3 & 256 & 256 & 512 & \texttt{-----??---}\\
5 & 4 & 2 & 1 & 6 & 0 & 1024 & 0 & \texttt{----------}\\
6 & 4 & 2 & 1 & 6 & 0 & 1024 & 0 & \texttt{----------}\\
7 & 0 & 1 & 2 & 3 & 256 & 256 & 512 & \texttt{--------??}\\
8 & 7 & 2 & 1 & 6 & 0 & 1024 & 0 & \texttt{----------}\\
9 & 7 & 2 & 1 & 6 & 0 & 1024 & 0 & \texttt{----------}\\
\multicolumn{9}{@{}l@{}}{\rule{0pt}{2.6ex}\textbf{T10.709}\enspace \texttt{0122122121}\enspace $n{=}10$, $|H|{=}6$, len 3, br 3, $|\mathrm{Aut}|{=}8$, orb 6/4, hist-orb 3/2; M 14336, res 9384, rev 2546: ok}\\*
0 & -- & 0 & 6 & 1 & 150 & 2 & 872 & \texttt{-?????????}\\*
1 & 0 & 1 & 2 & 2 & 256 & 256 & 512 & \texttt{--??------}\\
2 & 1 & 2 & 1 & 4 & 0 & 1024 & 0 & \texttt{----------}\\
3 & 1 & 2 & 1 & 4 & 0 & 1024 & 0 & \texttt{----------}\\
4 & 0 & 1 & 2 & 2 & 256 & 256 & 512 & \texttt{-----??---}\\
5 & 4 & 2 & 1 & 4 & 0 & 1024 & 0 & \texttt{----------}\\
6 & 4 & 2 & 1 & 4 & 0 & 1024 & 0 & \texttt{----------}\\
7 & 0 & 1 & 1 & 1 & 512 & 512 & 0 & \texttt{--------+-}\\
8 & 7 & 2 & 1 & 1 & 0 & 1024 & 0 & \texttt{----------}\\
9 & 0 & 1 & 1 & 1 & 0 & 1024 & 0 & \texttt{----------}\\
\multicolumn{9}{@{}l@{}}{\rule{0pt}{2.6ex}\textbf{T10.710}\enspace \texttt{0122122111}\enspace $n{=}10$, $|H|{=}7$, len 3, br 3, $|\mathrm{Aut}|{=}48$, orb 4/1, hist-orb 2/0; M 13312, res 10284, rev 1394: ok}\\*
0 & -- & 0 & 7 & 1 & 50 & 2 & 972 & \texttt{-?????????}\\*
1 & 0 & 1 & 2 & 2 & 256 & 256 & 512 & \texttt{--??------}\\
2 & 1 & 2 & 1 & 4 & 0 & 1024 & 0 & \texttt{----------}\\
3 & 1 & 2 & 1 & 4 & 0 & 1024 & 0 & \texttt{----------}\\
4 & 0 & 1 & 2 & 2 & 256 & 256 & 512 & \texttt{-----??---}\\
5 & 4 & 2 & 1 & 4 & 0 & 1024 & 0 & \texttt{----------}\\
6 & 4 & 2 & 1 & 4 & 0 & 1024 & 0 & \texttt{----------}\\
7 & 0 & 1 & 1 & 3 & 0 & 1024 & 0 & \texttt{----------}\\
8 & 0 & 1 & 1 & 3 & 0 & 1024 & 0 & \texttt{----------}\\
9 & 0 & 1 & 1 & 3 & 0 & 1024 & 0 & \texttt{----------}\\
\multicolumn{9}{@{}l@{}}{\rule{0pt}{2.6ex}\textbf{T10.711}\enspace \texttt{0122121212}\enspace $n{=}10$, $|H|{=}5$, len 3, br 2, $|\mathrm{Aut}|{=}12$, orb 5/2, hist-orb 2/0; M 14336, res 7536, rev 3722: ok}\\*
0 & -- & 0 & 5 & 1 & 270 & 2 & 752 & \texttt{-?????????}\\*
1 & 0 & 1 & 2 & 1 & 256 & 256 & 512 & \texttt{--??------}\\
2 & 1 & 2 & 1 & 2 & 0 & 1024 & 0 & \texttt{----------}\\
3 & 1 & 2 & 1 & 2 & 0 & 1024 & 0 & \texttt{----------}\\
4 & 0 & 1 & 1 & 3 & 512 & 512 & 0 & \texttt{-----+----}\\
5 & 4 & 2 & 1 & 3 & 0 & 1024 & 0 & \texttt{----------}\\
6 & 0 & 1 & 1 & 3 & 512 & 512 & 0 & \texttt{-------+--}\\
7 & 6 & 2 & 1 & 3 & 0 & 1024 & 0 & \texttt{----------}\\
8 & 0 & 1 & 1 & 3 & 512 & 512 & 0 & \texttt{---------+}\\
9 & 8 & 2 & 1 & 3 & 0 & 1024 & 0 & \texttt{----------}\\
\multicolumn{9}{@{}l@{}}{\rule{0pt}{2.6ex}\textbf{T10.712}\enspace \texttt{0122121211}\enspace $n{=}10$, $|H|{=}6$, len 3, br 2, $|\mathrm{Aut}|{=}8$, orb 6/2, hist-orb 3/0; M 13312, res 9156, rev 2154: ok}\\*
0 & -- & 0 & 6 & 1 & 90 & 2 & 932 & \texttt{-?????????}\\*
1 & 0 & 1 & 2 & 1 & 256 & 256 & 512 & \texttt{--??------}\\
2 & 1 & 2 & 1 & 2 & 0 & 1024 & 0 & \texttt{----------}\\
3 & 1 & 2 & 1 & 2 & 0 & 1024 & 0 & \texttt{----------}\\
4 & 0 & 1 & 1 & 2 & 512 & 512 & 0 & \texttt{-----+----}\\
5 & 4 & 2 & 1 & 2 & 0 & 1024 & 0 & \texttt{----------}\\
6 & 0 & 1 & 1 & 2 & 512 & 512 & 0 & \texttt{-------+--}\\
7 & 6 & 2 & 1 & 2 & 0 & 1024 & 0 & \texttt{----------}\\
8 & 0 & 1 & 1 & 2 & 0 & 1024 & 0 & \texttt{----------}\\
9 & 0 & 1 & 1 & 2 & 0 & 1024 & 0 & \texttt{----------}\\
\multicolumn{9}{@{}l@{}}{\rule{0pt}{2.6ex}\textbf{T10.713}\enspace \texttt{0122121111}\enspace $n{=}10$, $|H|{=}7$, len 3, br 2, $|\mathrm{Aut}|{=}48$, orb 6/4, hist-orb 3/1; M 12288, res 9696, rev 1250: ok}\\*
0 & -- & 0 & 7 & 1 & 30 & 2 & 992 & \texttt{-?????????}\\*
1 & 0 & 1 & 2 & 1 & 256 & 256 & 512 & \texttt{--??------}\\
2 & 1 & 2 & 1 & 2 & 0 & 1024 & 0 & \texttt{----------}\\
3 & 1 & 2 & 1 & 2 & 0 & 1024 & 0 & \texttt{----------}\\
4 & 0 & 1 & 1 & 1 & 512 & 512 & 0 & \texttt{-----+----}\\
5 & 4 & 2 & 1 & 1 & 0 & 1024 & 0 & \texttt{----------}\\
6 & 0 & 1 & 1 & 4 & 0 & 1024 & 0 & \texttt{----------}\\
7 & 0 & 1 & 1 & 4 & 0 & 1024 & 0 & \texttt{----------}\\
8 & 0 & 1 & 1 & 4 & 0 & 1024 & 0 & \texttt{----------}\\
9 & 0 & 1 & 1 & 4 & 0 & 1024 & 0 & \texttt{----------}\\
\multicolumn{9}{@{}l@{}}{\rule{0pt}{2.6ex}\textbf{T10.714}\enspace \texttt{0122111111}\enspace $n{=}10$, $|H|{=}8$, len 3, br 2, $|\mathrm{Aut}|{=}1440$, orb 4/2, hist-orb 2/0; M 11264, res 9876, rev 594: ok}\\*
0 & -- & 0 & 8 & 1 & 10 & 2 & 1012 & \texttt{-?????????}\\*
1 & 0 & 1 & 2 & 1 & 256 & 256 & 512 & \texttt{--??------}\\
2 & 1 & 2 & 1 & 2 & 0 & 1024 & 0 & \texttt{----------}\\
3 & 1 & 2 & 1 & 2 & 0 & 1024 & 0 & \texttt{----------}\\
4 & 0 & 1 & 1 & 6 & 0 & 1024 & 0 & \texttt{----------}\\
5 & 0 & 1 & 1 & 6 & 0 & 1024 & 0 & \texttt{----------}\\
6 & 0 & 1 & 1 & 6 & 0 & 1024 & 0 & \texttt{----------}\\
7 & 0 & 1 & 1 & 6 & 0 & 1024 & 0 & \texttt{----------}\\
8 & 0 & 1 & 1 & 6 & 0 & 1024 & 0 & \texttt{----------}\\
9 & 0 & 1 & 1 & 6 & 0 & 1024 & 0 & \texttt{----------}\\
\multicolumn{9}{@{}l@{}}{\rule{0pt}{2.6ex}\textbf{T10.715}\enspace \texttt{0121212121}\enspace $n{=}10$, $|H|{=}5$, len 3, br 1, $|\mathrm{Aut}|{=}24$, orb 4/2, hist-orb 2/1; M 13312, res 7740, rev 3074: ok}\\*
0 & -- & 0 & 5 & 1 & 162 & 2 & 860 & \texttt{-?????????}\\*
1 & 0 & 1 & 1 & 4 & 512 & 512 & 0 & \texttt{--+-------}\\
2 & 1 & 2 & 1 & 4 & 0 & 1024 & 0 & \texttt{----------}\\
3 & 0 & 1 & 1 & 4 & 512 & 512 & 0 & \texttt{----+-----}\\
4 & 3 & 2 & 1 & 4 & 0 & 1024 & 0 & \texttt{----------}\\
5 & 0 & 1 & 1 & 4 & 512 & 512 & 0 & \texttt{------+---}\\
6 & 5 & 2 & 1 & 4 & 0 & 1024 & 0 & \texttt{----------}\\
7 & 0 & 1 & 1 & 4 & 512 & 512 & 0 & \texttt{--------+-}\\
8 & 7 & 2 & 1 & 4 & 0 & 1024 & 0 & \texttt{----------}\\
9 & 0 & 1 & 1 & 1 & 0 & 1024 & 0 & \texttt{----------}\\
\multicolumn{9}{@{}l@{}}{\rule{0pt}{2.6ex}\textbf{T10.716}\enspace \texttt{0121212111}\enspace $n{=}10$, $|H|{=}6$, len 3, br 1, $|\mathrm{Aut}|{=}36$, orb 4/1, hist-orb 2/0; M 12288, res 8712, rev 1914: ok}\\*
0 & -- & 0 & 6 & 1 & 54 & 2 & 968 & \texttt{-?????????}\\*
1 & 0 & 1 & 1 & 3 & 512 & 512 & 0 & \texttt{--+-------}\\
2 & 1 & 2 & 1 & 3 & 0 & 1024 & 0 & \texttt{----------}\\
3 & 0 & 1 & 1 & 3 & 512 & 512 & 0 & \texttt{----+-----}\\
4 & 3 & 2 & 1 & 3 & 0 & 1024 & 0 & \texttt{----------}\\
5 & 0 & 1 & 1 & 3 & 512 & 512 & 0 & \texttt{------+---}\\
6 & 5 & 2 & 1 & 3 & 0 & 1024 & 0 & \texttt{----------}\\
7 & 0 & 1 & 1 & 3 & 0 & 1024 & 0 & \texttt{----------}\\
8 & 0 & 1 & 1 & 3 & 0 & 1024 & 0 & \texttt{----------}\\
9 & 0 & 1 & 1 & 3 & 0 & 1024 & 0 & \texttt{----------}\\
\multicolumn{9}{@{}l@{}}{\rule{0pt}{2.6ex}\textbf{T10.717}\enspace \texttt{0121211111}\enspace $n{=}10$, $|H|{=}7$, len 3, br 1, $|\mathrm{Aut}|{=}240$, orb 4/1, hist-orb 2/0; M 11264, res 9036, rev 1162: ok}\\*
0 & -- & 0 & 7 & 1 & 18 & 2 & 1004 & \texttt{-?????????}\\*
1 & 0 & 1 & 1 & 2 & 512 & 512 & 0 & \texttt{--+-------}\\
2 & 1 & 2 & 1 & 2 & 0 & 1024 & 0 & \texttt{----------}\\
3 & 0 & 1 & 1 & 2 & 512 & 512 & 0 & \texttt{----+-----}\\
4 & 3 & 2 & 1 & 2 & 0 & 1024 & 0 & \texttt{----------}\\
5 & 0 & 1 & 1 & 5 & 0 & 1024 & 0 & \texttt{----------}\\
6 & 0 & 1 & 1 & 5 & 0 & 1024 & 0 & \texttt{----------}\\
7 & 0 & 1 & 1 & 5 & 0 & 1024 & 0 & \texttt{----------}\\
8 & 0 & 1 & 1 & 5 & 0 & 1024 & 0 & \texttt{----------}\\
9 & 0 & 1 & 1 & 5 & 0 & 1024 & 0 & \texttt{----------}\\
\multicolumn{9}{@{}l@{}}{\rule{0pt}{2.6ex}\textbf{T10.718}\enspace \texttt{0121111111}\enspace $n{=}10$, $|H|{=}8$, len 3, br 1, $|\mathrm{Aut}|{=}5040$, orb 4/3, hist-orb 2/1; M 10240, res 9144, rev 562: ok}\\*
0 & -- & 0 & 8 & 1 & 6 & 2 & 1016 & \texttt{-?????????}\\*
1 & 0 & 1 & 1 & 1 & 512 & 512 & 0 & \texttt{--+-------}\\
2 & 1 & 2 & 1 & 1 & 0 & 1024 & 0 & \texttt{----------}\\
3 & 0 & 1 & 1 & 7 & 0 & 1024 & 0 & \texttt{----------}\\
4 & 0 & 1 & 1 & 7 & 0 & 1024 & 0 & \texttt{----------}\\
5 & 0 & 1 & 1 & 7 & 0 & 1024 & 0 & \texttt{----------}\\
6 & 0 & 1 & 1 & 7 & 0 & 1024 & 0 & \texttt{----------}\\
7 & 0 & 1 & 1 & 7 & 0 & 1024 & 0 & \texttt{----------}\\
8 & 0 & 1 & 1 & 7 & 0 & 1024 & 0 & \texttt{----------}\\
9 & 0 & 1 & 1 & 7 & 0 & 1024 & 0 & \texttt{----------}\\
\multicolumn{9}{@{}l@{}}{\rule{0pt}{2.6ex}\textbf{T10.719}\enspace \texttt{0111111111}\enspace $n{=}10$, $|H|{=}9$, len 2, br 1, $|\mathrm{Aut}|{=}362880$, orb 2/1, hist-orb 1/0; M 9216, res 9180, rev 18: ok}\\*
0 & -- & 0 & 9 & 1 & 2 & 2 & 1020 & \texttt{-?????????}\\*
1 & 0 & 1 & 1 & 9 & 0 & 1024 & 0 & \texttt{----------}\\
2 & 0 & 1 & 1 & 9 & 0 & 1024 & 0 & \texttt{----------}\\
3 & 0 & 1 & 1 & 9 & 0 & 1024 & 0 & \texttt{----------}\\
4 & 0 & 1 & 1 & 9 & 0 & 1024 & 0 & \texttt{----------}\\
5 & 0 & 1 & 1 & 9 & 0 & 1024 & 0 & \texttt{----------}\\
6 & 0 & 1 & 1 & 9 & 0 & 1024 & 0 & \texttt{----------}\\
7 & 0 & 1 & 1 & 9 & 0 & 1024 & 0 & \texttt{----------}\\
8 & 0 & 1 & 1 & 9 & 0 & 1024 & 0 & \texttt{----------}\\
9 & 0 & 1 & 1 & 9 & 0 & 1024 & 0 & \texttt{----------}\\
\bottomrule
\end{longtable}}


\section{Volume of histories: profiles and settling moments}\label{app:histories}

\paragraph{What is listed.} For every frame with at most nine moments (\NHistRows{} histories in all), this appendix lists each history: its number $h$ (histories are numbered in the order of their maximal moments); the moments it contains, from the least moment to its maximal moment; its \emph{profile}, the values of $h(m)$ at those moments in the same order; its \emph{settling moment}, the first moment on it through which no other history passes; and the size of its orbit under the automorphism group. A history with orbit size 1 is singled out by the shape of the frame; a frame none of whose histories has orbit size 1 is one in which no thin red line can be specified by the B-structure (Proposition~\ref{p:trl}).

\paragraph{What is checked.} The program verifies that every profile is non-increasing, begins with the number of histories of the frame and ends with 1 (Corollary~\ref{c:profile}). From the settling moment on, every tenseless proposition is settled relative to the history, by Proposition~\ref{p:bivalence}(c). In a chain the settling moment is the least moment; in a frame with more than one history, the settling moment of a history is the immediate successor, on that history, of the last branching moment on it. The profiles display Theorem~\ref{t:monotone} in its most compact form: the number of possible continuations never increases along a history.

% generated by gen/generate.py -- do not edit
{\scriptsize\setlength{\tabcolsep}{4pt}
\begin{longtable}{@{}llrlllr@{}}
\toprule
frame & level sequence & $h$ & moments of $h$ & profile $h(m)$ along $h$ & settles at & orb\\
\midrule
\endfirsthead
\toprule
frame & level sequence & $h$ & moments of $h$ & profile $h(m)$ along $h$ & settles at & orb\\
\midrule
\endhead
T1.1 & \texttt{0} & 0 & \texttt{0} & \texttt{1} & 0 & 1\\
T2.1 & \texttt{01} & 0 & \texttt{0.1} & \texttt{1.1} & 0 & 1\\
T3.1 & \texttt{012} & 0 & \texttt{0.1.2} & \texttt{1.1.1} & 0 & 1\\
T3.2 & \texttt{011} & 0 & \texttt{0.1} & \texttt{2.1} & 1 & 2\\*
 &  & 1 & \texttt{0.2} & \texttt{2.1} & 2 & 2\\
T4.1 & \texttt{0123} & 0 & \texttt{0.1.2.3} & \texttt{1.1.1.1} & 0 & 1\\
T4.2 & \texttt{0122} & 0 & \texttt{0.1.2} & \texttt{2.2.1} & 2 & 2\\*
 &  & 1 & \texttt{0.1.3} & \texttt{2.2.1} & 3 & 2\\
T4.3 & \texttt{0121} & 0 & \texttt{0.1.2} & \texttt{2.1.1} & 1 & 1\\*
 &  & 1 & \texttt{0.3} & \texttt{2.1} & 3 & 1\\
T4.4 & \texttt{0111} & 0 & \texttt{0.1} & \texttt{3.1} & 1 & 3\\*
 &  & 1 & \texttt{0.2} & \texttt{3.1} & 2 & 3\\
 &  & 2 & \texttt{0.3} & \texttt{3.1} & 3 & 3\\
T5.1 & \texttt{01234} & 0 & \texttt{0.1.2.3.4} & \texttt{1.1.1.1.1} & 0 & 1\\
T5.2 & \texttt{01233} & 0 & \texttt{0.1.2.3} & \texttt{2.2.2.1} & 3 & 2\\*
 &  & 1 & \texttt{0.1.2.4} & \texttt{2.2.2.1} & 4 & 2\\
T5.3 & \texttt{01232} & 0 & \texttt{0.1.2.3} & \texttt{2.2.1.1} & 2 & 1\\*
 &  & 1 & \texttt{0.1.4} & \texttt{2.2.1} & 4 & 1\\
T5.4 & \texttt{01231} & 0 & \texttt{0.1.2.3} & \texttt{2.1.1.1} & 1 & 1\\*
 &  & 1 & \texttt{0.4} & \texttt{2.1} & 4 & 1\\
T5.5 & \texttt{01222} & 0 & \texttt{0.1.2} & \texttt{3.3.1} & 2 & 3\\*
 &  & 1 & \texttt{0.1.3} & \texttt{3.3.1} & 3 & 3\\
 &  & 2 & \texttt{0.1.4} & \texttt{3.3.1} & 4 & 3\\
T5.6 & \texttt{01221} & 0 & \texttt{0.1.2} & \texttt{3.2.1} & 2 & 2\\*
 &  & 1 & \texttt{0.1.3} & \texttt{3.2.1} & 3 & 2\\
 &  & 2 & \texttt{0.4} & \texttt{3.1} & 4 & 1\\
T5.7 & \texttt{01212} & 0 & \texttt{0.1.2} & \texttt{2.1.1} & 1 & 2\\*
 &  & 1 & \texttt{0.3.4} & \texttt{2.1.1} & 3 & 2\\
T5.8 & \texttt{01211} & 0 & \texttt{0.1.2} & \texttt{3.1.1} & 1 & 1\\*
 &  & 1 & \texttt{0.3} & \texttt{3.1} & 3 & 2\\
 &  & 2 & \texttt{0.4} & \texttt{3.1} & 4 & 2\\
T5.9 & \texttt{01111} & 0 & \texttt{0.1} & \texttt{4.1} & 1 & 4\\*
 &  & 1 & \texttt{0.2} & \texttt{4.1} & 2 & 4\\
 &  & 2 & \texttt{0.3} & \texttt{4.1} & 3 & 4\\
 &  & 3 & \texttt{0.4} & \texttt{4.1} & 4 & 4\\
T6.01 & \texttt{012345} & 0 & \texttt{0.1.2.3.4.5} & \texttt{1.1.1.1.1.1} & 0 & 1\\
T6.02 & \texttt{012344} & 0 & \texttt{0.1.2.3.4} & \texttt{2.2.2.2.1} & 4 & 2\\*
 &  & 1 & \texttt{0.1.2.3.5} & \texttt{2.2.2.2.1} & 5 & 2\\
T6.03 & \texttt{012343} & 0 & \texttt{0.1.2.3.4} & \texttt{2.2.2.1.1} & 3 & 1\\*
 &  & 1 & \texttt{0.1.2.5} & \texttt{2.2.2.1} & 5 & 1\\
T6.04 & \texttt{012342} & 0 & \texttt{0.1.2.3.4} & \texttt{2.2.1.1.1} & 2 & 1\\*
 &  & 1 & \texttt{0.1.5} & \texttt{2.2.1} & 5 & 1\\
T6.05 & \texttt{012341} & 0 & \texttt{0.1.2.3.4} & \texttt{2.1.1.1.1} & 1 & 1\\*
 &  & 1 & \texttt{0.5} & \texttt{2.1} & 5 & 1\\
T6.06 & \texttt{012333} & 0 & \texttt{0.1.2.3} & \texttt{3.3.3.1} & 3 & 3\\*
 &  & 1 & \texttt{0.1.2.4} & \texttt{3.3.3.1} & 4 & 3\\
 &  & 2 & \texttt{0.1.2.5} & \texttt{3.3.3.1} & 5 & 3\\
T6.07 & \texttt{012332} & 0 & \texttt{0.1.2.3} & \texttt{3.3.2.1} & 3 & 2\\*
 &  & 1 & \texttt{0.1.2.4} & \texttt{3.3.2.1} & 4 & 2\\
 &  & 2 & \texttt{0.1.5} & \texttt{3.3.1} & 5 & 1\\
T6.08 & \texttt{012331} & 0 & \texttt{0.1.2.3} & \texttt{3.2.2.1} & 3 & 2\\*
 &  & 1 & \texttt{0.1.2.4} & \texttt{3.2.2.1} & 4 & 2\\
 &  & 2 & \texttt{0.5} & \texttt{3.1} & 5 & 1\\
T6.09 & \texttt{012323} & 0 & \texttt{0.1.2.3} & \texttt{2.2.1.1} & 2 & 2\\*
 &  & 1 & \texttt{0.1.4.5} & \texttt{2.2.1.1} & 4 & 2\\
T6.10 & \texttt{012322} & 0 & \texttt{0.1.2.3} & \texttt{3.3.1.1} & 2 & 1\\*
 &  & 1 & \texttt{0.1.4} & \texttt{3.3.1} & 4 & 2\\
 &  & 2 & \texttt{0.1.5} & \texttt{3.3.1} & 5 & 2\\
T6.11 & \texttt{012321} & 0 & \texttt{0.1.2.3} & \texttt{3.2.1.1} & 2 & 1\\*
 &  & 1 & \texttt{0.1.4} & \texttt{3.2.1} & 4 & 1\\
 &  & 2 & \texttt{0.5} & \texttt{3.1} & 5 & 1\\
T6.12 & \texttt{012312} & 0 & \texttt{0.1.2.3} & \texttt{2.1.1.1} & 1 & 1\\*
 &  & 1 & \texttt{0.4.5} & \texttt{2.1.1} & 4 & 1\\
T6.13 & \texttt{012311} & 0 & \texttt{0.1.2.3} & \texttt{3.1.1.1} & 1 & 1\\*
 &  & 1 & \texttt{0.4} & \texttt{3.1} & 4 & 2\\
 &  & 2 & \texttt{0.5} & \texttt{3.1} & 5 & 2\\
T6.14 & \texttt{012222} & 0 & \texttt{0.1.2} & \texttt{4.4.1} & 2 & 4\\*
 &  & 1 & \texttt{0.1.3} & \texttt{4.4.1} & 3 & 4\\
 &  & 2 & \texttt{0.1.4} & \texttt{4.4.1} & 4 & 4\\
 &  & 3 & \texttt{0.1.5} & \texttt{4.4.1} & 5 & 4\\
T6.15 & \texttt{012221} & 0 & \texttt{0.1.2} & \texttt{4.3.1} & 2 & 3\\*
 &  & 1 & \texttt{0.1.3} & \texttt{4.3.1} & 3 & 3\\
 &  & 2 & \texttt{0.1.4} & \texttt{4.3.1} & 4 & 3\\
 &  & 3 & \texttt{0.5} & \texttt{4.1} & 5 & 1\\
T6.16 & \texttt{012212} & 0 & \texttt{0.1.2} & \texttt{3.2.1} & 2 & 2\\*
 &  & 1 & \texttt{0.1.3} & \texttt{3.2.1} & 3 & 2\\
 &  & 2 & \texttt{0.4.5} & \texttt{3.1.1} & 4 & 1\\
T6.17 & \texttt{012211} & 0 & \texttt{0.1.2} & \texttt{4.2.1} & 2 & 2\\*
 &  & 1 & \texttt{0.1.3} & \texttt{4.2.1} & 3 & 2\\
 &  & 2 & \texttt{0.4} & \texttt{4.1} & 4 & 2\\
 &  & 3 & \texttt{0.5} & \texttt{4.1} & 5 & 2\\
T6.18 & \texttt{012121} & 0 & \texttt{0.1.2} & \texttt{3.1.1} & 1 & 2\\*
 &  & 1 & \texttt{0.3.4} & \texttt{3.1.1} & 3 & 2\\
 &  & 2 & \texttt{0.5} & \texttt{3.1} & 5 & 1\\
T6.19 & \texttt{012111} & 0 & \texttt{0.1.2} & \texttt{4.1.1} & 1 & 1\\*
 &  & 1 & \texttt{0.3} & \texttt{4.1} & 3 & 3\\
 &  & 2 & \texttt{0.4} & \texttt{4.1} & 4 & 3\\
 &  & 3 & \texttt{0.5} & \texttt{4.1} & 5 & 3\\
T6.20 & \texttt{011111} & 0 & \texttt{0.1} & \texttt{5.1} & 1 & 5\\*
 &  & 1 & \texttt{0.2} & \texttt{5.1} & 2 & 5\\
 &  & 2 & \texttt{0.3} & \texttt{5.1} & 3 & 5\\
 &  & 3 & \texttt{0.4} & \texttt{5.1} & 4 & 5\\
 &  & 4 & \texttt{0.5} & \texttt{5.1} & 5 & 5\\
T7.01 & \texttt{0123456} & 0 & \texttt{0.1.2.3.4.5.6} & \texttt{1.1.1.1.1.1.1} & 0 & 1\\
T7.02 & \texttt{0123455} & 0 & \texttt{0.1.2.3.4.5} & \texttt{2.2.2.2.2.1} & 5 & 2\\*
 &  & 1 & \texttt{0.1.2.3.4.6} & \texttt{2.2.2.2.2.1} & 6 & 2\\
T7.03 & \texttt{0123454} & 0 & \texttt{0.1.2.3.4.5} & \texttt{2.2.2.2.1.1} & 4 & 1\\*
 &  & 1 & \texttt{0.1.2.3.6} & \texttt{2.2.2.2.1} & 6 & 1\\
T7.04 & \texttt{0123453} & 0 & \texttt{0.1.2.3.4.5} & \texttt{2.2.2.1.1.1} & 3 & 1\\*
 &  & 1 & \texttt{0.1.2.6} & \texttt{2.2.2.1} & 6 & 1\\
T7.05 & \texttt{0123452} & 0 & \texttt{0.1.2.3.4.5} & \texttt{2.2.1.1.1.1} & 2 & 1\\*
 &  & 1 & \texttt{0.1.6} & \texttt{2.2.1} & 6 & 1\\
T7.06 & \texttt{0123451} & 0 & \texttt{0.1.2.3.4.5} & \texttt{2.1.1.1.1.1} & 1 & 1\\*
 &  & 1 & \texttt{0.6} & \texttt{2.1} & 6 & 1\\
T7.07 & \texttt{0123444} & 0 & \texttt{0.1.2.3.4} & \texttt{3.3.3.3.1} & 4 & 3\\*
 &  & 1 & \texttt{0.1.2.3.5} & \texttt{3.3.3.3.1} & 5 & 3\\
 &  & 2 & \texttt{0.1.2.3.6} & \texttt{3.3.3.3.1} & 6 & 3\\
T7.08 & \texttt{0123443} & 0 & \texttt{0.1.2.3.4} & \texttt{3.3.3.2.1} & 4 & 2\\*
 &  & 1 & \texttt{0.1.2.3.5} & \texttt{3.3.3.2.1} & 5 & 2\\
 &  & 2 & \texttt{0.1.2.6} & \texttt{3.3.3.1} & 6 & 1\\
T7.09 & \texttt{0123442} & 0 & \texttt{0.1.2.3.4} & \texttt{3.3.2.2.1} & 4 & 2\\*
 &  & 1 & \texttt{0.1.2.3.5} & \texttt{3.3.2.2.1} & 5 & 2\\
 &  & 2 & \texttt{0.1.6} & \texttt{3.3.1} & 6 & 1\\
T7.10 & \texttt{0123441} & 0 & \texttt{0.1.2.3.4} & \texttt{3.2.2.2.1} & 4 & 2\\*
 &  & 1 & \texttt{0.1.2.3.5} & \texttt{3.2.2.2.1} & 5 & 2\\
 &  & 2 & \texttt{0.6} & \texttt{3.1} & 6 & 1\\
T7.11 & \texttt{0123434} & 0 & \texttt{0.1.2.3.4} & \texttt{2.2.2.1.1} & 3 & 2\\*
 &  & 1 & \texttt{0.1.2.5.6} & \texttt{2.2.2.1.1} & 5 & 2\\
T7.12 & \texttt{0123433} & 0 & \texttt{0.1.2.3.4} & \texttt{3.3.3.1.1} & 3 & 1\\*
 &  & 1 & \texttt{0.1.2.5} & \texttt{3.3.3.1} & 5 & 2\\
 &  & 2 & \texttt{0.1.2.6} & \texttt{3.3.3.1} & 6 & 2\\
T7.13 & \texttt{0123432} & 0 & \texttt{0.1.2.3.4} & \texttt{3.3.2.1.1} & 3 & 1\\*
 &  & 1 & \texttt{0.1.2.5} & \texttt{3.3.2.1} & 5 & 1\\
 &  & 2 & \texttt{0.1.6} & \texttt{3.3.1} & 6 & 1\\
T7.14 & \texttt{0123431} & 0 & \texttt{0.1.2.3.4} & \texttt{3.2.2.1.1} & 3 & 1\\*
 &  & 1 & \texttt{0.1.2.5} & \texttt{3.2.2.1} & 5 & 1\\
 &  & 2 & \texttt{0.6} & \texttt{3.1} & 6 & 1\\
T7.15 & \texttt{0123423} & 0 & \texttt{0.1.2.3.4} & \texttt{2.2.1.1.1} & 2 & 1\\*
 &  & 1 & \texttt{0.1.5.6} & \texttt{2.2.1.1} & 5 & 1\\
T7.16 & \texttt{0123422} & 0 & \texttt{0.1.2.3.4} & \texttt{3.3.1.1.1} & 2 & 1\\*
 &  & 1 & \texttt{0.1.5} & \texttt{3.3.1} & 5 & 2\\
 &  & 2 & \texttt{0.1.6} & \texttt{3.3.1} & 6 & 2\\
T7.17 & \texttt{0123421} & 0 & \texttt{0.1.2.3.4} & \texttt{3.2.1.1.1} & 2 & 1\\*
 &  & 1 & \texttt{0.1.5} & \texttt{3.2.1} & 5 & 1\\
 &  & 2 & \texttt{0.6} & \texttt{3.1} & 6 & 1\\
T7.18 & \texttt{0123412} & 0 & \texttt{0.1.2.3.4} & \texttt{2.1.1.1.1} & 1 & 1\\*
 &  & 1 & \texttt{0.5.6} & \texttt{2.1.1} & 5 & 1\\
T7.19 & \texttt{0123411} & 0 & \texttt{0.1.2.3.4} & \texttt{3.1.1.1.1} & 1 & 1\\*
 &  & 1 & \texttt{0.5} & \texttt{3.1} & 5 & 2\\
 &  & 2 & \texttt{0.6} & \texttt{3.1} & 6 & 2\\
T7.20 & \texttt{0123333} & 0 & \texttt{0.1.2.3} & \texttt{4.4.4.1} & 3 & 4\\*
 &  & 1 & \texttt{0.1.2.4} & \texttt{4.4.4.1} & 4 & 4\\
 &  & 2 & \texttt{0.1.2.5} & \texttt{4.4.4.1} & 5 & 4\\
 &  & 3 & \texttt{0.1.2.6} & \texttt{4.4.4.1} & 6 & 4\\
T7.21 & \texttt{0123332} & 0 & \texttt{0.1.2.3} & \texttt{4.4.3.1} & 3 & 3\\*
 &  & 1 & \texttt{0.1.2.4} & \texttt{4.4.3.1} & 4 & 3\\
 &  & 2 & \texttt{0.1.2.5} & \texttt{4.4.3.1} & 5 & 3\\
 &  & 3 & \texttt{0.1.6} & \texttt{4.4.1} & 6 & 1\\
T7.22 & \texttt{0123331} & 0 & \texttt{0.1.2.3} & \texttt{4.3.3.1} & 3 & 3\\*
 &  & 1 & \texttt{0.1.2.4} & \texttt{4.3.3.1} & 4 & 3\\
 &  & 2 & \texttt{0.1.2.5} & \texttt{4.3.3.1} & 5 & 3\\
 &  & 3 & \texttt{0.6} & \texttt{4.1} & 6 & 1\\
T7.23 & \texttt{0123323} & 0 & \texttt{0.1.2.3} & \texttt{3.3.2.1} & 3 & 2\\*
 &  & 1 & \texttt{0.1.2.4} & \texttt{3.3.2.1} & 4 & 2\\
 &  & 2 & \texttt{0.1.5.6} & \texttt{3.3.1.1} & 5 & 1\\
T7.24 & \texttt{0123322} & 0 & \texttt{0.1.2.3} & \texttt{4.4.2.1} & 3 & 2\\*
 &  & 1 & \texttt{0.1.2.4} & \texttt{4.4.2.1} & 4 & 2\\
 &  & 2 & \texttt{0.1.5} & \texttt{4.4.1} & 5 & 2\\
 &  & 3 & \texttt{0.1.6} & \texttt{4.4.1} & 6 & 2\\
T7.25 & \texttt{0123321} & 0 & \texttt{0.1.2.3} & \texttt{4.3.2.1} & 3 & 2\\*
 &  & 1 & \texttt{0.1.2.4} & \texttt{4.3.2.1} & 4 & 2\\
 &  & 2 & \texttt{0.1.5} & \texttt{4.3.1} & 5 & 1\\
 &  & 3 & \texttt{0.6} & \texttt{4.1} & 6 & 1\\
T7.26 & \texttt{0123312} & 0 & \texttt{0.1.2.3} & \texttt{3.2.2.1} & 3 & 2\\*
 &  & 1 & \texttt{0.1.2.4} & \texttt{3.2.2.1} & 4 & 2\\
 &  & 2 & \texttt{0.5.6} & \texttt{3.1.1} & 5 & 1\\
T7.27 & \texttt{0123311} & 0 & \texttt{0.1.2.3} & \texttt{4.2.2.1} & 3 & 2\\*
 &  & 1 & \texttt{0.1.2.4} & \texttt{4.2.2.1} & 4 & 2\\
 &  & 2 & \texttt{0.5} & \texttt{4.1} & 5 & 2\\
 &  & 3 & \texttt{0.6} & \texttt{4.1} & 6 & 2\\
T7.28 & \texttt{0123232} & 0 & \texttt{0.1.2.3} & \texttt{3.3.1.1} & 2 & 2\\*
 &  & 1 & \texttt{0.1.4.5} & \texttt{3.3.1.1} & 4 & 2\\
 &  & 2 & \texttt{0.1.6} & \texttt{3.3.1} & 6 & 1\\
T7.29 & \texttt{0123231} & 0 & \texttt{0.1.2.3} & \texttt{3.2.1.1} & 2 & 2\\*
 &  & 1 & \texttt{0.1.4.5} & \texttt{3.2.1.1} & 4 & 2\\
 &  & 2 & \texttt{0.6} & \texttt{3.1} & 6 & 1\\
T7.30 & \texttt{0123222} & 0 & \texttt{0.1.2.3} & \texttt{4.4.1.1} & 2 & 1\\*
 &  & 1 & \texttt{0.1.4} & \texttt{4.4.1} & 4 & 3\\
 &  & 2 & \texttt{0.1.5} & \texttt{4.4.1} & 5 & 3\\
 &  & 3 & \texttt{0.1.6} & \texttt{4.4.1} & 6 & 3\\
T7.31 & \texttt{0123221} & 0 & \texttt{0.1.2.3} & \texttt{4.3.1.1} & 2 & 1\\*
 &  & 1 & \texttt{0.1.4} & \texttt{4.3.1} & 4 & 2\\
 &  & 2 & \texttt{0.1.5} & \texttt{4.3.1} & 5 & 2\\
 &  & 3 & \texttt{0.6} & \texttt{4.1} & 6 & 1\\
T7.32 & \texttt{0123212} & 0 & \texttt{0.1.2.3} & \texttt{3.2.1.1} & 2 & 1\\*
 &  & 1 & \texttt{0.1.4} & \texttt{3.2.1} & 4 & 1\\
 &  & 2 & \texttt{0.5.6} & \texttt{3.1.1} & 5 & 1\\
T7.33 & \texttt{0123211} & 0 & \texttt{0.1.2.3} & \texttt{4.2.1.1} & 2 & 1\\*
 &  & 1 & \texttt{0.1.4} & \texttt{4.2.1} & 4 & 1\\
 &  & 2 & \texttt{0.5} & \texttt{4.1} & 5 & 2\\
 &  & 3 & \texttt{0.6} & \texttt{4.1} & 6 & 2\\
T7.34 & \texttt{0123123} & 0 & \texttt{0.1.2.3} & \texttt{2.1.1.1} & 1 & 2\\*
 &  & 1 & \texttt{0.4.5.6} & \texttt{2.1.1.1} & 4 & 2\\
T7.35 & \texttt{0123122} & 0 & \texttt{0.1.2.3} & \texttt{3.1.1.1} & 1 & 1\\*
 &  & 1 & \texttt{0.4.5} & \texttt{3.2.1} & 5 & 2\\
 &  & 2 & \texttt{0.4.6} & \texttt{3.2.1} & 6 & 2\\
T7.36 & \texttt{0123121} & 0 & \texttt{0.1.2.3} & \texttt{3.1.1.1} & 1 & 1\\*
 &  & 1 & \texttt{0.4.5} & \texttt{3.1.1} & 4 & 1\\
 &  & 2 & \texttt{0.6} & \texttt{3.1} & 6 & 1\\
T7.37 & \texttt{0123111} & 0 & \texttt{0.1.2.3} & \texttt{4.1.1.1} & 1 & 1\\*
 &  & 1 & \texttt{0.4} & \texttt{4.1} & 4 & 3\\
 &  & 2 & \texttt{0.5} & \texttt{4.1} & 5 & 3\\
 &  & 3 & \texttt{0.6} & \texttt{4.1} & 6 & 3\\
T7.38 & \texttt{0122222} & 0 & \texttt{0.1.2} & \texttt{5.5.1} & 2 & 5\\*
 &  & 1 & \texttt{0.1.3} & \texttt{5.5.1} & 3 & 5\\
 &  & 2 & \texttt{0.1.4} & \texttt{5.5.1} & 4 & 5\\
 &  & 3 & \texttt{0.1.5} & \texttt{5.5.1} & 5 & 5\\
 &  & 4 & \texttt{0.1.6} & \texttt{5.5.1} & 6 & 5\\
T7.39 & \texttt{0122221} & 0 & \texttt{0.1.2} & \texttt{5.4.1} & 2 & 4\\*
 &  & 1 & \texttt{0.1.3} & \texttt{5.4.1} & 3 & 4\\
 &  & 2 & \texttt{0.1.4} & \texttt{5.4.1} & 4 & 4\\
 &  & 3 & \texttt{0.1.5} & \texttt{5.4.1} & 5 & 4\\
 &  & 4 & \texttt{0.6} & \texttt{5.1} & 6 & 1\\
T7.40 & \texttt{0122212} & 0 & \texttt{0.1.2} & \texttt{4.3.1} & 2 & 3\\*
 &  & 1 & \texttt{0.1.3} & \texttt{4.3.1} & 3 & 3\\
 &  & 2 & \texttt{0.1.4} & \texttt{4.3.1} & 4 & 3\\
 &  & 3 & \texttt{0.5.6} & \texttt{4.1.1} & 5 & 1\\
T7.41 & \texttt{0122211} & 0 & \texttt{0.1.2} & \texttt{5.3.1} & 2 & 3\\*
 &  & 1 & \texttt{0.1.3} & \texttt{5.3.1} & 3 & 3\\
 &  & 2 & \texttt{0.1.4} & \texttt{5.3.1} & 4 & 3\\
 &  & 3 & \texttt{0.5} & \texttt{5.1} & 5 & 2\\
 &  & 4 & \texttt{0.6} & \texttt{5.1} & 6 & 2\\
T7.42 & \texttt{0122122} & 0 & \texttt{0.1.2} & \texttt{4.2.1} & 2 & 4\\*
 &  & 1 & \texttt{0.1.3} & \texttt{4.2.1} & 3 & 4\\
 &  & 2 & \texttt{0.4.5} & \texttt{4.2.1} & 5 & 4\\
 &  & 3 & \texttt{0.4.6} & \texttt{4.2.1} & 6 & 4\\
T7.43 & \texttt{0122121} & 0 & \texttt{0.1.2} & \texttt{4.2.1} & 2 & 2\\*
 &  & 1 & \texttt{0.1.3} & \texttt{4.2.1} & 3 & 2\\
 &  & 2 & \texttt{0.4.5} & \texttt{4.1.1} & 4 & 1\\
 &  & 3 & \texttt{0.6} & \texttt{4.1} & 6 & 1\\
T7.44 & \texttt{0122111} & 0 & \texttt{0.1.2} & \texttt{5.2.1} & 2 & 2\\*
 &  & 1 & \texttt{0.1.3} & \texttt{5.2.1} & 3 & 2\\
 &  & 2 & \texttt{0.4} & \texttt{5.1} & 4 & 3\\
 &  & 3 & \texttt{0.5} & \texttt{5.1} & 5 & 3\\
 &  & 4 & \texttt{0.6} & \texttt{5.1} & 6 & 3\\
T7.45 & \texttt{0121212} & 0 & \texttt{0.1.2} & \texttt{3.1.1} & 1 & 3\\*
 &  & 1 & \texttt{0.3.4} & \texttt{3.1.1} & 3 & 3\\
 &  & 2 & \texttt{0.5.6} & \texttt{3.1.1} & 5 & 3\\
T7.46 & \texttt{0121211} & 0 & \texttt{0.1.2} & \texttt{4.1.1} & 1 & 2\\*
 &  & 1 & \texttt{0.3.4} & \texttt{4.1.1} & 3 & 2\\
 &  & 2 & \texttt{0.5} & \texttt{4.1} & 5 & 2\\
 &  & 3 & \texttt{0.6} & \texttt{4.1} & 6 & 2\\
T7.47 & \texttt{0121111} & 0 & \texttt{0.1.2} & \texttt{5.1.1} & 1 & 1\\*
 &  & 1 & \texttt{0.3} & \texttt{5.1} & 3 & 4\\
 &  & 2 & \texttt{0.4} & \texttt{5.1} & 4 & 4\\
 &  & 3 & \texttt{0.5} & \texttt{5.1} & 5 & 4\\
 &  & 4 & \texttt{0.6} & \texttt{5.1} & 6 & 4\\
T7.48 & \texttt{0111111} & 0 & \texttt{0.1} & \texttt{6.1} & 1 & 6\\*
 &  & 1 & \texttt{0.2} & \texttt{6.1} & 2 & 6\\
 &  & 2 & \texttt{0.3} & \texttt{6.1} & 3 & 6\\
 &  & 3 & \texttt{0.4} & \texttt{6.1} & 4 & 6\\
 &  & 4 & \texttt{0.5} & \texttt{6.1} & 5 & 6\\
 &  & 5 & \texttt{0.6} & \texttt{6.1} & 6 & 6\\
T8.001 & \texttt{01234567} & 0 & \texttt{0.1.2.3.4.5.6.7} & \texttt{1.1.1.1.1.1.1.1} & 0 & 1\\
T8.002 & \texttt{01234566} & 0 & \texttt{0.1.2.3.4.5.6} & \texttt{2.2.2.2.2.2.1} & 6 & 2\\*
 &  & 1 & \texttt{0.1.2.3.4.5.7} & \texttt{2.2.2.2.2.2.1} & 7 & 2\\
T8.003 & \texttt{01234565} & 0 & \texttt{0.1.2.3.4.5.6} & \texttt{2.2.2.2.2.1.1} & 5 & 1\\*
 &  & 1 & \texttt{0.1.2.3.4.7} & \texttt{2.2.2.2.2.1} & 7 & 1\\
T8.004 & \texttt{01234564} & 0 & \texttt{0.1.2.3.4.5.6} & \texttt{2.2.2.2.1.1.1} & 4 & 1\\*
 &  & 1 & \texttt{0.1.2.3.7} & \texttt{2.2.2.2.1} & 7 & 1\\
T8.005 & \texttt{01234563} & 0 & \texttt{0.1.2.3.4.5.6} & \texttt{2.2.2.1.1.1.1} & 3 & 1\\*
 &  & 1 & \texttt{0.1.2.7} & \texttt{2.2.2.1} & 7 & 1\\
T8.006 & \texttt{01234562} & 0 & \texttt{0.1.2.3.4.5.6} & \texttt{2.2.1.1.1.1.1} & 2 & 1\\*
 &  & 1 & \texttt{0.1.7} & \texttt{2.2.1} & 7 & 1\\
T8.007 & \texttt{01234561} & 0 & \texttt{0.1.2.3.4.5.6} & \texttt{2.1.1.1.1.1.1} & 1 & 1\\*
 &  & 1 & \texttt{0.7} & \texttt{2.1} & 7 & 1\\
T8.008 & \texttt{01234555} & 0 & \texttt{0.1.2.3.4.5} & \texttt{3.3.3.3.3.1} & 5 & 3\\*
 &  & 1 & \texttt{0.1.2.3.4.6} & \texttt{3.3.3.3.3.1} & 6 & 3\\
 &  & 2 & \texttt{0.1.2.3.4.7} & \texttt{3.3.3.3.3.1} & 7 & 3\\
T8.009 & \texttt{01234554} & 0 & \texttt{0.1.2.3.4.5} & \texttt{3.3.3.3.2.1} & 5 & 2\\*
 &  & 1 & \texttt{0.1.2.3.4.6} & \texttt{3.3.3.3.2.1} & 6 & 2\\
 &  & 2 & \texttt{0.1.2.3.7} & \texttt{3.3.3.3.1} & 7 & 1\\
T8.010 & \texttt{01234553} & 0 & \texttt{0.1.2.3.4.5} & \texttt{3.3.3.2.2.1} & 5 & 2\\*
 &  & 1 & \texttt{0.1.2.3.4.6} & \texttt{3.3.3.2.2.1} & 6 & 2\\
 &  & 2 & \texttt{0.1.2.7} & \texttt{3.3.3.1} & 7 & 1\\
T8.011 & \texttt{01234552} & 0 & \texttt{0.1.2.3.4.5} & \texttt{3.3.2.2.2.1} & 5 & 2\\*
 &  & 1 & \texttt{0.1.2.3.4.6} & \texttt{3.3.2.2.2.1} & 6 & 2\\
 &  & 2 & \texttt{0.1.7} & \texttt{3.3.1} & 7 & 1\\
T8.012 & \texttt{01234551} & 0 & \texttt{0.1.2.3.4.5} & \texttt{3.2.2.2.2.1} & 5 & 2\\*
 &  & 1 & \texttt{0.1.2.3.4.6} & \texttt{3.2.2.2.2.1} & 6 & 2\\
 &  & 2 & \texttt{0.7} & \texttt{3.1} & 7 & 1\\
T8.013 & \texttt{01234545} & 0 & \texttt{0.1.2.3.4.5} & \texttt{2.2.2.2.1.1} & 4 & 2\\*
 &  & 1 & \texttt{0.1.2.3.6.7} & \texttt{2.2.2.2.1.1} & 6 & 2\\
T8.014 & \texttt{01234544} & 0 & \texttt{0.1.2.3.4.5} & \texttt{3.3.3.3.1.1} & 4 & 1\\*
 &  & 1 & \texttt{0.1.2.3.6} & \texttt{3.3.3.3.1} & 6 & 2\\
 &  & 2 & \texttt{0.1.2.3.7} & \texttt{3.3.3.3.1} & 7 & 2\\
T8.015 & \texttt{01234543} & 0 & \texttt{0.1.2.3.4.5} & \texttt{3.3.3.2.1.1} & 4 & 1\\*
 &  & 1 & \texttt{0.1.2.3.6} & \texttt{3.3.3.2.1} & 6 & 1\\
 &  & 2 & \texttt{0.1.2.7} & \texttt{3.3.3.1} & 7 & 1\\
T8.016 & \texttt{01234542} & 0 & \texttt{0.1.2.3.4.5} & \texttt{3.3.2.2.1.1} & 4 & 1\\*
 &  & 1 & \texttt{0.1.2.3.6} & \texttt{3.3.2.2.1} & 6 & 1\\
 &  & 2 & \texttt{0.1.7} & \texttt{3.3.1} & 7 & 1\\
T8.017 & \texttt{01234541} & 0 & \texttt{0.1.2.3.4.5} & \texttt{3.2.2.2.1.1} & 4 & 1\\*
 &  & 1 & \texttt{0.1.2.3.6} & \texttt{3.2.2.2.1} & 6 & 1\\
 &  & 2 & \texttt{0.7} & \texttt{3.1} & 7 & 1\\
T8.018 & \texttt{01234534} & 0 & \texttt{0.1.2.3.4.5} & \texttt{2.2.2.1.1.1} & 3 & 1\\*
 &  & 1 & \texttt{0.1.2.6.7} & \texttt{2.2.2.1.1} & 6 & 1\\
T8.019 & \texttt{01234533} & 0 & \texttt{0.1.2.3.4.5} & \texttt{3.3.3.1.1.1} & 3 & 1\\*
 &  & 1 & \texttt{0.1.2.6} & \texttt{3.3.3.1} & 6 & 2\\
 &  & 2 & \texttt{0.1.2.7} & \texttt{3.3.3.1} & 7 & 2\\
T8.020 & \texttt{01234532} & 0 & \texttt{0.1.2.3.4.5} & \texttt{3.3.2.1.1.1} & 3 & 1\\*
 &  & 1 & \texttt{0.1.2.6} & \texttt{3.3.2.1} & 6 & 1\\
 &  & 2 & \texttt{0.1.7} & \texttt{3.3.1} & 7 & 1\\
T8.021 & \texttt{01234531} & 0 & \texttt{0.1.2.3.4.5} & \texttt{3.2.2.1.1.1} & 3 & 1\\*
 &  & 1 & \texttt{0.1.2.6} & \texttt{3.2.2.1} & 6 & 1\\
 &  & 2 & \texttt{0.7} & \texttt{3.1} & 7 & 1\\
T8.022 & \texttt{01234523} & 0 & \texttt{0.1.2.3.4.5} & \texttt{2.2.1.1.1.1} & 2 & 1\\*
 &  & 1 & \texttt{0.1.6.7} & \texttt{2.2.1.1} & 6 & 1\\
T8.023 & \texttt{01234522} & 0 & \texttt{0.1.2.3.4.5} & \texttt{3.3.1.1.1.1} & 2 & 1\\*
 &  & 1 & \texttt{0.1.6} & \texttt{3.3.1} & 6 & 2\\
 &  & 2 & \texttt{0.1.7} & \texttt{3.3.1} & 7 & 2\\
T8.024 & \texttt{01234521} & 0 & \texttt{0.1.2.3.4.5} & \texttt{3.2.1.1.1.1} & 2 & 1\\*
 &  & 1 & \texttt{0.1.6} & \texttt{3.2.1} & 6 & 1\\
 &  & 2 & \texttt{0.7} & \texttt{3.1} & 7 & 1\\
T8.025 & \texttt{01234512} & 0 & \texttt{0.1.2.3.4.5} & \texttt{2.1.1.1.1.1} & 1 & 1\\*
 &  & 1 & \texttt{0.6.7} & \texttt{2.1.1} & 6 & 1\\
T8.026 & \texttt{01234511} & 0 & \texttt{0.1.2.3.4.5} & \texttt{3.1.1.1.1.1} & 1 & 1\\*
 &  & 1 & \texttt{0.6} & \texttt{3.1} & 6 & 2\\
 &  & 2 & \texttt{0.7} & \texttt{3.1} & 7 & 2\\
T8.027 & \texttt{01234444} & 0 & \texttt{0.1.2.3.4} & \texttt{4.4.4.4.1} & 4 & 4\\*
 &  & 1 & \texttt{0.1.2.3.5} & \texttt{4.4.4.4.1} & 5 & 4\\
 &  & 2 & \texttt{0.1.2.3.6} & \texttt{4.4.4.4.1} & 6 & 4\\
 &  & 3 & \texttt{0.1.2.3.7} & \texttt{4.4.4.4.1} & 7 & 4\\
T8.028 & \texttt{01234443} & 0 & \texttt{0.1.2.3.4} & \texttt{4.4.4.3.1} & 4 & 3\\*
 &  & 1 & \texttt{0.1.2.3.5} & \texttt{4.4.4.3.1} & 5 & 3\\
 &  & 2 & \texttt{0.1.2.3.6} & \texttt{4.4.4.3.1} & 6 & 3\\
 &  & 3 & \texttt{0.1.2.7} & \texttt{4.4.4.1} & 7 & 1\\
T8.029 & \texttt{01234442} & 0 & \texttt{0.1.2.3.4} & \texttt{4.4.3.3.1} & 4 & 3\\*
 &  & 1 & \texttt{0.1.2.3.5} & \texttt{4.4.3.3.1} & 5 & 3\\
 &  & 2 & \texttt{0.1.2.3.6} & \texttt{4.4.3.3.1} & 6 & 3\\
 &  & 3 & \texttt{0.1.7} & \texttt{4.4.1} & 7 & 1\\
T8.030 & \texttt{01234441} & 0 & \texttt{0.1.2.3.4} & \texttt{4.3.3.3.1} & 4 & 3\\*
 &  & 1 & \texttt{0.1.2.3.5} & \texttt{4.3.3.3.1} & 5 & 3\\
 &  & 2 & \texttt{0.1.2.3.6} & \texttt{4.3.3.3.1} & 6 & 3\\
 &  & 3 & \texttt{0.7} & \texttt{4.1} & 7 & 1\\
T8.031 & \texttt{01234434} & 0 & \texttt{0.1.2.3.4} & \texttt{3.3.3.2.1} & 4 & 2\\*
 &  & 1 & \texttt{0.1.2.3.5} & \texttt{3.3.3.2.1} & 5 & 2\\
 &  & 2 & \texttt{0.1.2.6.7} & \texttt{3.3.3.1.1} & 6 & 1\\
T8.032 & \texttt{01234433} & 0 & \texttt{0.1.2.3.4} & \texttt{4.4.4.2.1} & 4 & 2\\*
 &  & 1 & \texttt{0.1.2.3.5} & \texttt{4.4.4.2.1} & 5 & 2\\
 &  & 2 & \texttt{0.1.2.6} & \texttt{4.4.4.1} & 6 & 2\\
 &  & 3 & \texttt{0.1.2.7} & \texttt{4.4.4.1} & 7 & 2\\
T8.033 & \texttt{01234432} & 0 & \texttt{0.1.2.3.4} & \texttt{4.4.3.2.1} & 4 & 2\\*
 &  & 1 & \texttt{0.1.2.3.5} & \texttt{4.4.3.2.1} & 5 & 2\\
 &  & 2 & \texttt{0.1.2.6} & \texttt{4.4.3.1} & 6 & 1\\
 &  & 3 & \texttt{0.1.7} & \texttt{4.4.1} & 7 & 1\\
T8.034 & \texttt{01234431} & 0 & \texttt{0.1.2.3.4} & \texttt{4.3.3.2.1} & 4 & 2\\*
 &  & 1 & \texttt{0.1.2.3.5} & \texttt{4.3.3.2.1} & 5 & 2\\
 &  & 2 & \texttt{0.1.2.6} & \texttt{4.3.3.1} & 6 & 1\\
 &  & 3 & \texttt{0.7} & \texttt{4.1} & 7 & 1\\
T8.035 & \texttt{01234423} & 0 & \texttt{0.1.2.3.4} & \texttt{3.3.2.2.1} & 4 & 2\\*
 &  & 1 & \texttt{0.1.2.3.5} & \texttt{3.3.2.2.1} & 5 & 2\\
 &  & 2 & \texttt{0.1.6.7} & \texttt{3.3.1.1} & 6 & 1\\
T8.036 & \texttt{01234422} & 0 & \texttt{0.1.2.3.4} & \texttt{4.4.2.2.1} & 4 & 2\\*
 &  & 1 & \texttt{0.1.2.3.5} & \texttt{4.4.2.2.1} & 5 & 2\\
 &  & 2 & \texttt{0.1.6} & \texttt{4.4.1} & 6 & 2\\
 &  & 3 & \texttt{0.1.7} & \texttt{4.4.1} & 7 & 2\\
T8.037 & \texttt{01234421} & 0 & \texttt{0.1.2.3.4} & \texttt{4.3.2.2.1} & 4 & 2\\*
 &  & 1 & \texttt{0.1.2.3.5} & \texttt{4.3.2.2.1} & 5 & 2\\
 &  & 2 & \texttt{0.1.6} & \texttt{4.3.1} & 6 & 1\\
 &  & 3 & \texttt{0.7} & \texttt{4.1} & 7 & 1\\
T8.038 & \texttt{01234412} & 0 & \texttt{0.1.2.3.4} & \texttt{3.2.2.2.1} & 4 & 2\\*
 &  & 1 & \texttt{0.1.2.3.5} & \texttt{3.2.2.2.1} & 5 & 2\\
 &  & 2 & \texttt{0.6.7} & \texttt{3.1.1} & 6 & 1\\
T8.039 & \texttt{01234411} & 0 & \texttt{0.1.2.3.4} & \texttt{4.2.2.2.1} & 4 & 2\\*
 &  & 1 & \texttt{0.1.2.3.5} & \texttt{4.2.2.2.1} & 5 & 2\\
 &  & 2 & \texttt{0.6} & \texttt{4.1} & 6 & 2\\
 &  & 3 & \texttt{0.7} & \texttt{4.1} & 7 & 2\\
T8.040 & \texttt{01234343} & 0 & \texttt{0.1.2.3.4} & \texttt{3.3.3.1.1} & 3 & 2\\*
 &  & 1 & \texttt{0.1.2.5.6} & \texttt{3.3.3.1.1} & 5 & 2\\
 &  & 2 & \texttt{0.1.2.7} & \texttt{3.3.3.1} & 7 & 1\\
T8.041 & \texttt{01234342} & 0 & \texttt{0.1.2.3.4} & \texttt{3.3.2.1.1} & 3 & 2\\*
 &  & 1 & \texttt{0.1.2.5.6} & \texttt{3.3.2.1.1} & 5 & 2\\
 &  & 2 & \texttt{0.1.7} & \texttt{3.3.1} & 7 & 1\\
T8.042 & \texttt{01234341} & 0 & \texttt{0.1.2.3.4} & \texttt{3.2.2.1.1} & 3 & 2\\*
 &  & 1 & \texttt{0.1.2.5.6} & \texttt{3.2.2.1.1} & 5 & 2\\
 &  & 2 & \texttt{0.7} & \texttt{3.1} & 7 & 1\\
T8.043 & \texttt{01234333} & 0 & \texttt{0.1.2.3.4} & \texttt{4.4.4.1.1} & 3 & 1\\*
 &  & 1 & \texttt{0.1.2.5} & \texttt{4.4.4.1} & 5 & 3\\
 &  & 2 & \texttt{0.1.2.6} & \texttt{4.4.4.1} & 6 & 3\\
 &  & 3 & \texttt{0.1.2.7} & \texttt{4.4.4.1} & 7 & 3\\
T8.044 & \texttt{01234332} & 0 & \texttt{0.1.2.3.4} & \texttt{4.4.3.1.1} & 3 & 1\\*
 &  & 1 & \texttt{0.1.2.5} & \texttt{4.4.3.1} & 5 & 2\\
 &  & 2 & \texttt{0.1.2.6} & \texttt{4.4.3.1} & 6 & 2\\
 &  & 3 & \texttt{0.1.7} & \texttt{4.4.1} & 7 & 1\\
T8.045 & \texttt{01234331} & 0 & \texttt{0.1.2.3.4} & \texttt{4.3.3.1.1} & 3 & 1\\*
 &  & 1 & \texttt{0.1.2.5} & \texttt{4.3.3.1} & 5 & 2\\
 &  & 2 & \texttt{0.1.2.6} & \texttt{4.3.3.1} & 6 & 2\\
 &  & 3 & \texttt{0.7} & \texttt{4.1} & 7 & 1\\
T8.046 & \texttt{01234323} & 0 & \texttt{0.1.2.3.4} & \texttt{3.3.2.1.1} & 3 & 1\\*
 &  & 1 & \texttt{0.1.2.5} & \texttt{3.3.2.1} & 5 & 1\\
 &  & 2 & \texttt{0.1.6.7} & \texttt{3.3.1.1} & 6 & 1\\
T8.047 & \texttt{01234322} & 0 & \texttt{0.1.2.3.4} & \texttt{4.4.2.1.1} & 3 & 1\\*
 &  & 1 & \texttt{0.1.2.5} & \texttt{4.4.2.1} & 5 & 1\\
 &  & 2 & \texttt{0.1.6} & \texttt{4.4.1} & 6 & 2\\
 &  & 3 & \texttt{0.1.7} & \texttt{4.4.1} & 7 & 2\\
T8.048 & \texttt{01234321} & 0 & \texttt{0.1.2.3.4} & \texttt{4.3.2.1.1} & 3 & 1\\*
 &  & 1 & \texttt{0.1.2.5} & \texttt{4.3.2.1} & 5 & 1\\
 &  & 2 & \texttt{0.1.6} & \texttt{4.3.1} & 6 & 1\\
 &  & 3 & \texttt{0.7} & \texttt{4.1} & 7 & 1\\
T8.049 & \texttt{01234312} & 0 & \texttt{0.1.2.3.4} & \texttt{3.2.2.1.1} & 3 & 1\\*
 &  & 1 & \texttt{0.1.2.5} & \texttt{3.2.2.1} & 5 & 1\\
 &  & 2 & \texttt{0.6.7} & \texttt{3.1.1} & 6 & 1\\
T8.050 & \texttt{01234311} & 0 & \texttt{0.1.2.3.4} & \texttt{4.2.2.1.1} & 3 & 1\\*
 &  & 1 & \texttt{0.1.2.5} & \texttt{4.2.2.1} & 5 & 1\\
 &  & 2 & \texttt{0.6} & \texttt{4.1} & 6 & 2\\
 &  & 3 & \texttt{0.7} & \texttt{4.1} & 7 & 2\\
T8.051 & \texttt{01234234} & 0 & \texttt{0.1.2.3.4} & \texttt{2.2.1.1.1} & 2 & 2\\*
 &  & 1 & \texttt{0.1.5.6.7} & \texttt{2.2.1.1.1} & 5 & 2\\
T8.052 & \texttt{01234233} & 0 & \texttt{0.1.2.3.4} & \texttt{3.3.1.1.1} & 2 & 1\\*
 &  & 1 & \texttt{0.1.5.6} & \texttt{3.3.2.1} & 6 & 2\\
 &  & 2 & \texttt{0.1.5.7} & \texttt{3.3.2.1} & 7 & 2\\
T8.053 & \texttt{01234232} & 0 & \texttt{0.1.2.3.4} & \texttt{3.3.1.1.1} & 2 & 1\\*
 &  & 1 & \texttt{0.1.5.6} & \texttt{3.3.1.1} & 5 & 1\\
 &  & 2 & \texttt{0.1.7} & \texttt{3.3.1} & 7 & 1\\
T8.054 & \texttt{01234231} & 0 & \texttt{0.1.2.3.4} & \texttt{3.2.1.1.1} & 2 & 1\\*
 &  & 1 & \texttt{0.1.5.6} & \texttt{3.2.1.1} & 5 & 1\\
 &  & 2 & \texttt{0.7} & \texttt{3.1} & 7 & 1\\
T8.055 & \texttt{01234222} & 0 & \texttt{0.1.2.3.4} & \texttt{4.4.1.1.1} & 2 & 1\\*
 &  & 1 & \texttt{0.1.5} & \texttt{4.4.1} & 5 & 3\\
 &  & 2 & \texttt{0.1.6} & \texttt{4.4.1} & 6 & 3\\
 &  & 3 & \texttt{0.1.7} & \texttt{4.4.1} & 7 & 3\\
T8.056 & \texttt{01234221} & 0 & \texttt{0.1.2.3.4} & \texttt{4.3.1.1.1} & 2 & 1\\*
 &  & 1 & \texttt{0.1.5} & \texttt{4.3.1} & 5 & 2\\
 &  & 2 & \texttt{0.1.6} & \texttt{4.3.1} & 6 & 2\\
 &  & 3 & \texttt{0.7} & \texttt{4.1} & 7 & 1\\
T8.057 & \texttt{01234212} & 0 & \texttt{0.1.2.3.4} & \texttt{3.2.1.1.1} & 2 & 1\\*
 &  & 1 & \texttt{0.1.5} & \texttt{3.2.1} & 5 & 1\\
 &  & 2 & \texttt{0.6.7} & \texttt{3.1.1} & 6 & 1\\
T8.058 & \texttt{01234211} & 0 & \texttt{0.1.2.3.4} & \texttt{4.2.1.1.1} & 2 & 1\\*
 &  & 1 & \texttt{0.1.5} & \texttt{4.2.1} & 5 & 1\\
 &  & 2 & \texttt{0.6} & \texttt{4.1} & 6 & 2\\
 &  & 3 & \texttt{0.7} & \texttt{4.1} & 7 & 2\\
T8.059 & \texttt{01234123} & 0 & \texttt{0.1.2.3.4} & \texttt{2.1.1.1.1} & 1 & 1\\*
 &  & 1 & \texttt{0.5.6.7} & \texttt{2.1.1.1} & 5 & 1\\
T8.060 & \texttt{01234122} & 0 & \texttt{0.1.2.3.4} & \texttt{3.1.1.1.1} & 1 & 1\\*
 &  & 1 & \texttt{0.5.6} & \texttt{3.2.1} & 6 & 2\\
 &  & 2 & \texttt{0.5.7} & \texttt{3.2.1} & 7 & 2\\
T8.061 & \texttt{01234121} & 0 & \texttt{0.1.2.3.4} & \texttt{3.1.1.1.1} & 1 & 1\\*
 &  & 1 & \texttt{0.5.6} & \texttt{3.1.1} & 5 & 1\\
 &  & 2 & \texttt{0.7} & \texttt{3.1} & 7 & 1\\
T8.062 & \texttt{01234111} & 0 & \texttt{0.1.2.3.4} & \texttt{4.1.1.1.1} & 1 & 1\\*
 &  & 1 & \texttt{0.5} & \texttt{4.1} & 5 & 3\\
 &  & 2 & \texttt{0.6} & \texttt{4.1} & 6 & 3\\
 &  & 3 & \texttt{0.7} & \texttt{4.1} & 7 & 3\\
T8.063 & \texttt{01233333} & 0 & \texttt{0.1.2.3} & \texttt{5.5.5.1} & 3 & 5\\*
 &  & 1 & \texttt{0.1.2.4} & \texttt{5.5.5.1} & 4 & 5\\
 &  & 2 & \texttt{0.1.2.5} & \texttt{5.5.5.1} & 5 & 5\\
 &  & 3 & \texttt{0.1.2.6} & \texttt{5.5.5.1} & 6 & 5\\
 &  & 4 & \texttt{0.1.2.7} & \texttt{5.5.5.1} & 7 & 5\\
T8.064 & \texttt{01233332} & 0 & \texttt{0.1.2.3} & \texttt{5.5.4.1} & 3 & 4\\*
 &  & 1 & \texttt{0.1.2.4} & \texttt{5.5.4.1} & 4 & 4\\
 &  & 2 & \texttt{0.1.2.5} & \texttt{5.5.4.1} & 5 & 4\\
 &  & 3 & \texttt{0.1.2.6} & \texttt{5.5.4.1} & 6 & 4\\
 &  & 4 & \texttt{0.1.7} & \texttt{5.5.1} & 7 & 1\\
T8.065 & \texttt{01233331} & 0 & \texttt{0.1.2.3} & \texttt{5.4.4.1} & 3 & 4\\*
 &  & 1 & \texttt{0.1.2.4} & \texttt{5.4.4.1} & 4 & 4\\
 &  & 2 & \texttt{0.1.2.5} & \texttt{5.4.4.1} & 5 & 4\\
 &  & 3 & \texttt{0.1.2.6} & \texttt{5.4.4.1} & 6 & 4\\
 &  & 4 & \texttt{0.7} & \texttt{5.1} & 7 & 1\\
T8.066 & \texttt{01233323} & 0 & \texttt{0.1.2.3} & \texttt{4.4.3.1} & 3 & 3\\*
 &  & 1 & \texttt{0.1.2.4} & \texttt{4.4.3.1} & 4 & 3\\
 &  & 2 & \texttt{0.1.2.5} & \texttt{4.4.3.1} & 5 & 3\\
 &  & 3 & \texttt{0.1.6.7} & \texttt{4.4.1.1} & 6 & 1\\
T8.067 & \texttt{01233322} & 0 & \texttt{0.1.2.3} & \texttt{5.5.3.1} & 3 & 3\\*
 &  & 1 & \texttt{0.1.2.4} & \texttt{5.5.3.1} & 4 & 3\\
 &  & 2 & \texttt{0.1.2.5} & \texttt{5.5.3.1} & 5 & 3\\
 &  & 3 & \texttt{0.1.6} & \texttt{5.5.1} & 6 & 2\\
 &  & 4 & \texttt{0.1.7} & \texttt{5.5.1} & 7 & 2\\
T8.068 & \texttt{01233321} & 0 & \texttt{0.1.2.3} & \texttt{5.4.3.1} & 3 & 3\\*
 &  & 1 & \texttt{0.1.2.4} & \texttt{5.4.3.1} & 4 & 3\\
 &  & 2 & \texttt{0.1.2.5} & \texttt{5.4.3.1} & 5 & 3\\
 &  & 3 & \texttt{0.1.6} & \texttt{5.4.1} & 6 & 1\\
 &  & 4 & \texttt{0.7} & \texttt{5.1} & 7 & 1\\
T8.069 & \texttt{01233312} & 0 & \texttt{0.1.2.3} & \texttt{4.3.3.1} & 3 & 3\\*
 &  & 1 & \texttt{0.1.2.4} & \texttt{4.3.3.1} & 4 & 3\\
 &  & 2 & \texttt{0.1.2.5} & \texttt{4.3.3.1} & 5 & 3\\
 &  & 3 & \texttt{0.6.7} & \texttt{4.1.1} & 6 & 1\\
T8.070 & \texttt{01233311} & 0 & \texttt{0.1.2.3} & \texttt{5.3.3.1} & 3 & 3\\*
 &  & 1 & \texttt{0.1.2.4} & \texttt{5.3.3.1} & 4 & 3\\
 &  & 2 & \texttt{0.1.2.5} & \texttt{5.3.3.1} & 5 & 3\\
 &  & 3 & \texttt{0.6} & \texttt{5.1} & 6 & 2\\
 &  & 4 & \texttt{0.7} & \texttt{5.1} & 7 & 2\\
T8.071 & \texttt{01233233} & 0 & \texttt{0.1.2.3} & \texttt{4.4.2.1} & 3 & 4\\*
 &  & 1 & \texttt{0.1.2.4} & \texttt{4.4.2.1} & 4 & 4\\
 &  & 2 & \texttt{0.1.5.6} & \texttt{4.4.2.1} & 6 & 4\\
 &  & 3 & \texttt{0.1.5.7} & \texttt{4.4.2.1} & 7 & 4\\
T8.072 & \texttt{01233232} & 0 & \texttt{0.1.2.3} & \texttt{4.4.2.1} & 3 & 2\\*
 &  & 1 & \texttt{0.1.2.4} & \texttt{4.4.2.1} & 4 & 2\\
 &  & 2 & \texttt{0.1.5.6} & \texttt{4.4.1.1} & 5 & 1\\
 &  & 3 & \texttt{0.1.7} & \texttt{4.4.1} & 7 & 1\\
T8.073 & \texttt{01233231} & 0 & \texttt{0.1.2.3} & \texttt{4.3.2.1} & 3 & 2\\*
 &  & 1 & \texttt{0.1.2.4} & \texttt{4.3.2.1} & 4 & 2\\
 &  & 2 & \texttt{0.1.5.6} & \texttt{4.3.1.1} & 5 & 1\\
 &  & 3 & \texttt{0.7} & \texttt{4.1} & 7 & 1\\
T8.074 & \texttt{01233222} & 0 & \texttt{0.1.2.3} & \texttt{5.5.2.1} & 3 & 2\\*
 &  & 1 & \texttt{0.1.2.4} & \texttt{5.5.2.1} & 4 & 2\\
 &  & 2 & \texttt{0.1.5} & \texttt{5.5.1} & 5 & 3\\
 &  & 3 & \texttt{0.1.6} & \texttt{5.5.1} & 6 & 3\\
 &  & 4 & \texttt{0.1.7} & \texttt{5.5.1} & 7 & 3\\
T8.075 & \texttt{01233221} & 0 & \texttt{0.1.2.3} & \texttt{5.4.2.1} & 3 & 2\\*
 &  & 1 & \texttt{0.1.2.4} & \texttt{5.4.2.1} & 4 & 2\\
 &  & 2 & \texttt{0.1.5} & \texttt{5.4.1} & 5 & 2\\
 &  & 3 & \texttt{0.1.6} & \texttt{5.4.1} & 6 & 2\\
 &  & 4 & \texttt{0.7} & \texttt{5.1} & 7 & 1\\
T8.076 & \texttt{01233212} & 0 & \texttt{0.1.2.3} & \texttt{4.3.2.1} & 3 & 2\\*
 &  & 1 & \texttt{0.1.2.4} & \texttt{4.3.2.1} & 4 & 2\\
 &  & 2 & \texttt{0.1.5} & \texttt{4.3.1} & 5 & 1\\
 &  & 3 & \texttt{0.6.7} & \texttt{4.1.1} & 6 & 1\\
T8.077 & \texttt{01233211} & 0 & \texttt{0.1.2.3} & \texttt{5.3.2.1} & 3 & 2\\*
 &  & 1 & \texttt{0.1.2.4} & \texttt{5.3.2.1} & 4 & 2\\
 &  & 2 & \texttt{0.1.5} & \texttt{5.3.1} & 5 & 1\\
 &  & 3 & \texttt{0.6} & \texttt{5.1} & 6 & 2\\
 &  & 4 & \texttt{0.7} & \texttt{5.1} & 7 & 2\\
T8.078 & \texttt{01233123} & 0 & \texttt{0.1.2.3} & \texttt{3.2.2.1} & 3 & 2\\*
 &  & 1 & \texttt{0.1.2.4} & \texttt{3.2.2.1} & 4 & 2\\
 &  & 2 & \texttt{0.5.6.7} & \texttt{3.1.1.1} & 5 & 1\\
T8.079 & \texttt{01233122} & 0 & \texttt{0.1.2.3} & \texttt{4.2.2.1} & 3 & 2\\*
 &  & 1 & \texttt{0.1.2.4} & \texttt{4.2.2.1} & 4 & 2\\
 &  & 2 & \texttt{0.5.6} & \texttt{4.2.1} & 6 & 2\\
 &  & 3 & \texttt{0.5.7} & \texttt{4.2.1} & 7 & 2\\
T8.080 & \texttt{01233121} & 0 & \texttt{0.1.2.3} & \texttt{4.2.2.1} & 3 & 2\\*
 &  & 1 & \texttt{0.1.2.4} & \texttt{4.2.2.1} & 4 & 2\\
 &  & 2 & \texttt{0.5.6} & \texttt{4.1.1} & 5 & 1\\
 &  & 3 & \texttt{0.7} & \texttt{4.1} & 7 & 1\\
T8.081 & \texttt{01233111} & 0 & \texttt{0.1.2.3} & \texttt{5.2.2.1} & 3 & 2\\*
 &  & 1 & \texttt{0.1.2.4} & \texttt{5.2.2.1} & 4 & 2\\
 &  & 2 & \texttt{0.5} & \texttt{5.1} & 5 & 3\\
 &  & 3 & \texttt{0.6} & \texttt{5.1} & 6 & 3\\
 &  & 4 & \texttt{0.7} & \texttt{5.1} & 7 & 3\\
T8.082 & \texttt{01232323} & 0 & \texttt{0.1.2.3} & \texttt{3.3.1.1} & 2 & 3\\*
 &  & 1 & \texttt{0.1.4.5} & \texttt{3.3.1.1} & 4 & 3\\
 &  & 2 & \texttt{0.1.6.7} & \texttt{3.3.1.1} & 6 & 3\\
T8.083 & \texttt{01232322} & 0 & \texttt{0.1.2.3} & \texttt{4.4.1.1} & 2 & 2\\*
 &  & 1 & \texttt{0.1.4.5} & \texttt{4.4.1.1} & 4 & 2\\
 &  & 2 & \texttt{0.1.6} & \texttt{4.4.1} & 6 & 2\\
 &  & 3 & \texttt{0.1.7} & \texttt{4.4.1} & 7 & 2\\
T8.084 & \texttt{01232321} & 0 & \texttt{0.1.2.3} & \texttt{4.3.1.1} & 2 & 2\\*
 &  & 1 & \texttt{0.1.4.5} & \texttt{4.3.1.1} & 4 & 2\\
 &  & 2 & \texttt{0.1.6} & \texttt{4.3.1} & 6 & 1\\
 &  & 3 & \texttt{0.7} & \texttt{4.1} & 7 & 1\\
T8.085 & \texttt{01232312} & 0 & \texttt{0.1.2.3} & \texttt{3.2.1.1} & 2 & 2\\*
 &  & 1 & \texttt{0.1.4.5} & \texttt{3.2.1.1} & 4 & 2\\
 &  & 2 & \texttt{0.6.7} & \texttt{3.1.1} & 6 & 1\\
T8.086 & \texttt{01232311} & 0 & \texttt{0.1.2.3} & \texttt{4.2.1.1} & 2 & 2\\*
 &  & 1 & \texttt{0.1.4.5} & \texttt{4.2.1.1} & 4 & 2\\
 &  & 2 & \texttt{0.6} & \texttt{4.1} & 6 & 2\\
 &  & 3 & \texttt{0.7} & \texttt{4.1} & 7 & 2\\
T8.087 & \texttt{01232222} & 0 & \texttt{0.1.2.3} & \texttt{5.5.1.1} & 2 & 1\\*
 &  & 1 & \texttt{0.1.4} & \texttt{5.5.1} & 4 & 4\\
 &  & 2 & \texttt{0.1.5} & \texttt{5.5.1} & 5 & 4\\
 &  & 3 & \texttt{0.1.6} & \texttt{5.5.1} & 6 & 4\\
 &  & 4 & \texttt{0.1.7} & \texttt{5.5.1} & 7 & 4\\
T8.088 & \texttt{01232221} & 0 & \texttt{0.1.2.3} & \texttt{5.4.1.1} & 2 & 1\\*
 &  & 1 & \texttt{0.1.4} & \texttt{5.4.1} & 4 & 3\\
 &  & 2 & \texttt{0.1.5} & \texttt{5.4.1} & 5 & 3\\
 &  & 3 & \texttt{0.1.6} & \texttt{5.4.1} & 6 & 3\\
 &  & 4 & \texttt{0.7} & \texttt{5.1} & 7 & 1\\
T8.089 & \texttt{01232212} & 0 & \texttt{0.1.2.3} & \texttt{4.3.1.1} & 2 & 1\\*
 &  & 1 & \texttt{0.1.4} & \texttt{4.3.1} & 4 & 2\\
 &  & 2 & \texttt{0.1.5} & \texttt{4.3.1} & 5 & 2\\
 &  & 3 & \texttt{0.6.7} & \texttt{4.1.1} & 6 & 1\\
T8.090 & \texttt{01232211} & 0 & \texttt{0.1.2.3} & \texttt{5.3.1.1} & 2 & 1\\*
 &  & 1 & \texttt{0.1.4} & \texttt{5.3.1} & 4 & 2\\
 &  & 2 & \texttt{0.1.5} & \texttt{5.3.1} & 5 & 2\\
 &  & 3 & \texttt{0.6} & \texttt{5.1} & 6 & 2\\
 &  & 4 & \texttt{0.7} & \texttt{5.1} & 7 & 2\\
T8.091 & \texttt{01232123} & 0 & \texttt{0.1.2.3} & \texttt{3.2.1.1} & 2 & 1\\*
 &  & 1 & \texttt{0.1.4} & \texttt{3.2.1} & 4 & 1\\
 &  & 2 & \texttt{0.5.6.7} & \texttt{3.1.1.1} & 5 & 1\\
T8.092 & \texttt{01232122} & 0 & \texttt{0.1.2.3} & \texttt{4.2.1.1} & 2 & 1\\*
 &  & 1 & \texttt{0.1.4} & \texttt{4.2.1} & 4 & 1\\
 &  & 2 & \texttt{0.5.6} & \texttt{4.2.1} & 6 & 2\\
 &  & 3 & \texttt{0.5.7} & \texttt{4.2.1} & 7 & 2\\
T8.093 & \texttt{01232121} & 0 & \texttt{0.1.2.3} & \texttt{4.2.1.1} & 2 & 1\\*
 &  & 1 & \texttt{0.1.4} & \texttt{4.2.1} & 4 & 1\\
 &  & 2 & \texttt{0.5.6} & \texttt{4.1.1} & 5 & 1\\
 &  & 3 & \texttt{0.7} & \texttt{4.1} & 7 & 1\\
T8.094 & \texttt{01232111} & 0 & \texttt{0.1.2.3} & \texttt{5.2.1.1} & 2 & 1\\*
 &  & 1 & \texttt{0.1.4} & \texttt{5.2.1} & 4 & 1\\
 &  & 2 & \texttt{0.5} & \texttt{5.1} & 5 & 3\\
 &  & 3 & \texttt{0.6} & \texttt{5.1} & 6 & 3\\
 &  & 4 & \texttt{0.7} & \texttt{5.1} & 7 & 3\\
T8.095 & \texttt{01231231} & 0 & \texttt{0.1.2.3} & \texttt{3.1.1.1} & 1 & 2\\*
 &  & 1 & \texttt{0.4.5.6} & \texttt{3.1.1.1} & 4 & 2\\
 &  & 2 & \texttt{0.7} & \texttt{3.1} & 7 & 1\\
T8.096 & \texttt{01231222} & 0 & \texttt{0.1.2.3} & \texttt{4.1.1.1} & 1 & 1\\*
 &  & 1 & \texttt{0.4.5} & \texttt{4.3.1} & 5 & 3\\
 &  & 2 & \texttt{0.4.6} & \texttt{4.3.1} & 6 & 3\\
 &  & 3 & \texttt{0.4.7} & \texttt{4.3.1} & 7 & 3\\
T8.097 & \texttt{01231221} & 0 & \texttt{0.1.2.3} & \texttt{4.1.1.1} & 1 & 1\\*
 &  & 1 & \texttt{0.4.5} & \texttt{4.2.1} & 5 & 2\\
 &  & 2 & \texttt{0.4.6} & \texttt{4.2.1} & 6 & 2\\
 &  & 3 & \texttt{0.7} & \texttt{4.1} & 7 & 1\\
T8.098 & \texttt{01231212} & 0 & \texttt{0.1.2.3} & \texttt{3.1.1.1} & 1 & 1\\*
 &  & 1 & \texttt{0.4.5} & \texttt{3.1.1} & 4 & 2\\
 &  & 2 & \texttt{0.6.7} & \texttt{3.1.1} & 6 & 2\\
T8.099 & \texttt{01231211} & 0 & \texttt{0.1.2.3} & \texttt{4.1.1.1} & 1 & 1\\*
 &  & 1 & \texttt{0.4.5} & \texttt{4.1.1} & 4 & 1\\
 &  & 2 & \texttt{0.6} & \texttt{4.1} & 6 & 2\\
 &  & 3 & \texttt{0.7} & \texttt{4.1} & 7 & 2\\
T8.100 & \texttt{01231111} & 0 & \texttt{0.1.2.3} & \texttt{5.1.1.1} & 1 & 1\\*
 &  & 1 & \texttt{0.4} & \texttt{5.1} & 4 & 4\\
 &  & 2 & \texttt{0.5} & \texttt{5.1} & 5 & 4\\
 &  & 3 & \texttt{0.6} & \texttt{5.1} & 6 & 4\\
 &  & 4 & \texttt{0.7} & \texttt{5.1} & 7 & 4\\
T8.101 & \texttt{01222222} & 0 & \texttt{0.1.2} & \texttt{6.6.1} & 2 & 6\\*
 &  & 1 & \texttt{0.1.3} & \texttt{6.6.1} & 3 & 6\\
 &  & 2 & \texttt{0.1.4} & \texttt{6.6.1} & 4 & 6\\
 &  & 3 & \texttt{0.1.5} & \texttt{6.6.1} & 5 & 6\\
 &  & 4 & \texttt{0.1.6} & \texttt{6.6.1} & 6 & 6\\
 &  & 5 & \texttt{0.1.7} & \texttt{6.6.1} & 7 & 6\\
T8.102 & \texttt{01222221} & 0 & \texttt{0.1.2} & \texttt{6.5.1} & 2 & 5\\*
 &  & 1 & \texttt{0.1.3} & \texttt{6.5.1} & 3 & 5\\
 &  & 2 & \texttt{0.1.4} & \texttt{6.5.1} & 4 & 5\\
 &  & 3 & \texttt{0.1.5} & \texttt{6.5.1} & 5 & 5\\
 &  & 4 & \texttt{0.1.6} & \texttt{6.5.1} & 6 & 5\\
 &  & 5 & \texttt{0.7} & \texttt{6.1} & 7 & 1\\
T8.103 & \texttt{01222212} & 0 & \texttt{0.1.2} & \texttt{5.4.1} & 2 & 4\\*
 &  & 1 & \texttt{0.1.3} & \texttt{5.4.1} & 3 & 4\\
 &  & 2 & \texttt{0.1.4} & \texttt{5.4.1} & 4 & 4\\
 &  & 3 & \texttt{0.1.5} & \texttt{5.4.1} & 5 & 4\\
 &  & 4 & \texttt{0.6.7} & \texttt{5.1.1} & 6 & 1\\
T8.104 & \texttt{01222211} & 0 & \texttt{0.1.2} & \texttt{6.4.1} & 2 & 4\\*
 &  & 1 & \texttt{0.1.3} & \texttt{6.4.1} & 3 & 4\\
 &  & 2 & \texttt{0.1.4} & \texttt{6.4.1} & 4 & 4\\
 &  & 3 & \texttt{0.1.5} & \texttt{6.4.1} & 5 & 4\\
 &  & 4 & \texttt{0.6} & \texttt{6.1} & 6 & 2\\
 &  & 5 & \texttt{0.7} & \texttt{6.1} & 7 & 2\\
T8.105 & \texttt{01222122} & 0 & \texttt{0.1.2} & \texttt{5.3.1} & 2 & 3\\*
 &  & 1 & \texttt{0.1.3} & \texttt{5.3.1} & 3 & 3\\
 &  & 2 & \texttt{0.1.4} & \texttt{5.3.1} & 4 & 3\\
 &  & 3 & \texttt{0.5.6} & \texttt{5.2.1} & 6 & 2\\
 &  & 4 & \texttt{0.5.7} & \texttt{5.2.1} & 7 & 2\\
T8.106 & \texttt{01222121} & 0 & \texttt{0.1.2} & \texttt{5.3.1} & 2 & 3\\*
 &  & 1 & \texttt{0.1.3} & \texttt{5.3.1} & 3 & 3\\
 &  & 2 & \texttt{0.1.4} & \texttt{5.3.1} & 4 & 3\\
 &  & 3 & \texttt{0.5.6} & \texttt{5.1.1} & 5 & 1\\
 &  & 4 & \texttt{0.7} & \texttt{5.1} & 7 & 1\\
T8.107 & \texttt{01222111} & 0 & \texttt{0.1.2} & \texttt{6.3.1} & 2 & 3\\*
 &  & 1 & \texttt{0.1.3} & \texttt{6.3.1} & 3 & 3\\
 &  & 2 & \texttt{0.1.4} & \texttt{6.3.1} & 4 & 3\\
 &  & 3 & \texttt{0.5} & \texttt{6.1} & 5 & 3\\
 &  & 4 & \texttt{0.6} & \texttt{6.1} & 6 & 3\\
 &  & 5 & \texttt{0.7} & \texttt{6.1} & 7 & 3\\
T8.108 & \texttt{01221221} & 0 & \texttt{0.1.2} & \texttt{5.2.1} & 2 & 4\\*
 &  & 1 & \texttt{0.1.3} & \texttt{5.2.1} & 3 & 4\\
 &  & 2 & \texttt{0.4.5} & \texttt{5.2.1} & 5 & 4\\
 &  & 3 & \texttt{0.4.6} & \texttt{5.2.1} & 6 & 4\\
 &  & 4 & \texttt{0.7} & \texttt{5.1} & 7 & 1\\
T8.109 & \texttt{01221212} & 0 & \texttt{0.1.2} & \texttt{4.2.1} & 2 & 2\\*
 &  & 1 & \texttt{0.1.3} & \texttt{4.2.1} & 3 & 2\\
 &  & 2 & \texttt{0.4.5} & \texttt{4.1.1} & 4 & 2\\
 &  & 3 & \texttt{0.6.7} & \texttt{4.1.1} & 6 & 2\\
T8.110 & \texttt{01221211} & 0 & \texttt{0.1.2} & \texttt{5.2.1} & 2 & 2\\*
 &  & 1 & \texttt{0.1.3} & \texttt{5.2.1} & 3 & 2\\
 &  & 2 & \texttt{0.4.5} & \texttt{5.1.1} & 4 & 1\\
 &  & 3 & \texttt{0.6} & \texttt{5.1} & 6 & 2\\
 &  & 4 & \texttt{0.7} & \texttt{5.1} & 7 & 2\\
T8.111 & \texttt{01221111} & 0 & \texttt{0.1.2} & \texttt{6.2.1} & 2 & 2\\*
 &  & 1 & \texttt{0.1.3} & \texttt{6.2.1} & 3 & 2\\
 &  & 2 & \texttt{0.4} & \texttt{6.1} & 4 & 4\\
 &  & 3 & \texttt{0.5} & \texttt{6.1} & 5 & 4\\
 &  & 4 & \texttt{0.6} & \texttt{6.1} & 6 & 4\\
 &  & 5 & \texttt{0.7} & \texttt{6.1} & 7 & 4\\
T8.112 & \texttt{01212121} & 0 & \texttt{0.1.2} & \texttt{4.1.1} & 1 & 3\\*
 &  & 1 & \texttt{0.3.4} & \texttt{4.1.1} & 3 & 3\\
 &  & 2 & \texttt{0.5.6} & \texttt{4.1.1} & 5 & 3\\
 &  & 3 & \texttt{0.7} & \texttt{4.1} & 7 & 1\\
T8.113 & \texttt{01212111} & 0 & \texttt{0.1.2} & \texttt{5.1.1} & 1 & 2\\*
 &  & 1 & \texttt{0.3.4} & \texttt{5.1.1} & 3 & 2\\
 &  & 2 & \texttt{0.5} & \texttt{5.1} & 5 & 3\\
 &  & 3 & \texttt{0.6} & \texttt{5.1} & 6 & 3\\
 &  & 4 & \texttt{0.7} & \texttt{5.1} & 7 & 3\\
T8.114 & \texttt{01211111} & 0 & \texttt{0.1.2} & \texttt{6.1.1} & 1 & 1\\*
 &  & 1 & \texttt{0.3} & \texttt{6.1} & 3 & 5\\
 &  & 2 & \texttt{0.4} & \texttt{6.1} & 4 & 5\\
 &  & 3 & \texttt{0.5} & \texttt{6.1} & 5 & 5\\
 &  & 4 & \texttt{0.6} & \texttt{6.1} & 6 & 5\\
 &  & 5 & \texttt{0.7} & \texttt{6.1} & 7 & 5\\
T8.115 & \texttt{01111111} & 0 & \texttt{0.1} & \texttt{7.1} & 1 & 7\\*
 &  & 1 & \texttt{0.2} & \texttt{7.1} & 2 & 7\\
 &  & 2 & \texttt{0.3} & \texttt{7.1} & 3 & 7\\
 &  & 3 & \texttt{0.4} & \texttt{7.1} & 4 & 7\\
 &  & 4 & \texttt{0.5} & \texttt{7.1} & 5 & 7\\
 &  & 5 & \texttt{0.6} & \texttt{7.1} & 6 & 7\\
 &  & 6 & \texttt{0.7} & \texttt{7.1} & 7 & 7\\
T9.001 & \texttt{012345678} & 0 & \texttt{0.1.2.3.4.5.6.7.8} & \texttt{1.1.1.1.1.1.1.1.1} & 0 & 1\\
T9.002 & \texttt{012345677} & 0 & \texttt{0.1.2.3.4.5.6.7} & \texttt{2.2.2.2.2.2.2.1} & 7 & 2\\*
 &  & 1 & \texttt{0.1.2.3.4.5.6.8} & \texttt{2.2.2.2.2.2.2.1} & 8 & 2\\
T9.003 & \texttt{012345676} & 0 & \texttt{0.1.2.3.4.5.6.7} & \texttt{2.2.2.2.2.2.1.1} & 6 & 1\\*
 &  & 1 & \texttt{0.1.2.3.4.5.8} & \texttt{2.2.2.2.2.2.1} & 8 & 1\\
T9.004 & \texttt{012345675} & 0 & \texttt{0.1.2.3.4.5.6.7} & \texttt{2.2.2.2.2.1.1.1} & 5 & 1\\*
 &  & 1 & \texttt{0.1.2.3.4.8} & \texttt{2.2.2.2.2.1} & 8 & 1\\
T9.005 & \texttt{012345674} & 0 & \texttt{0.1.2.3.4.5.6.7} & \texttt{2.2.2.2.1.1.1.1} & 4 & 1\\*
 &  & 1 & \texttt{0.1.2.3.8} & \texttt{2.2.2.2.1} & 8 & 1\\
T9.006 & \texttt{012345673} & 0 & \texttt{0.1.2.3.4.5.6.7} & \texttt{2.2.2.1.1.1.1.1} & 3 & 1\\*
 &  & 1 & \texttt{0.1.2.8} & \texttt{2.2.2.1} & 8 & 1\\
T9.007 & \texttt{012345672} & 0 & \texttt{0.1.2.3.4.5.6.7} & \texttt{2.2.1.1.1.1.1.1} & 2 & 1\\*
 &  & 1 & \texttt{0.1.8} & \texttt{2.2.1} & 8 & 1\\
T9.008 & \texttt{012345671} & 0 & \texttt{0.1.2.3.4.5.6.7} & \texttt{2.1.1.1.1.1.1.1} & 1 & 1\\*
 &  & 1 & \texttt{0.8} & \texttt{2.1} & 8 & 1\\
T9.009 & \texttt{012345666} & 0 & \texttt{0.1.2.3.4.5.6} & \texttt{3.3.3.3.3.3.1} & 6 & 3\\*
 &  & 1 & \texttt{0.1.2.3.4.5.7} & \texttt{3.3.3.3.3.3.1} & 7 & 3\\
 &  & 2 & \texttt{0.1.2.3.4.5.8} & \texttt{3.3.3.3.3.3.1} & 8 & 3\\
T9.010 & \texttt{012345665} & 0 & \texttt{0.1.2.3.4.5.6} & \texttt{3.3.3.3.3.2.1} & 6 & 2\\*
 &  & 1 & \texttt{0.1.2.3.4.5.7} & \texttt{3.3.3.3.3.2.1} & 7 & 2\\
 &  & 2 & \texttt{0.1.2.3.4.8} & \texttt{3.3.3.3.3.1} & 8 & 1\\
T9.011 & \texttt{012345664} & 0 & \texttt{0.1.2.3.4.5.6} & \texttt{3.3.3.3.2.2.1} & 6 & 2\\*
 &  & 1 & \texttt{0.1.2.3.4.5.7} & \texttt{3.3.3.3.2.2.1} & 7 & 2\\
 &  & 2 & \texttt{0.1.2.3.8} & \texttt{3.3.3.3.1} & 8 & 1\\
T9.012 & \texttt{012345663} & 0 & \texttt{0.1.2.3.4.5.6} & \texttt{3.3.3.2.2.2.1} & 6 & 2\\*
 &  & 1 & \texttt{0.1.2.3.4.5.7} & \texttt{3.3.3.2.2.2.1} & 7 & 2\\
 &  & 2 & \texttt{0.1.2.8} & \texttt{3.3.3.1} & 8 & 1\\
T9.013 & \texttt{012345662} & 0 & \texttt{0.1.2.3.4.5.6} & \texttt{3.3.2.2.2.2.1} & 6 & 2\\*
 &  & 1 & \texttt{0.1.2.3.4.5.7} & \texttt{3.3.2.2.2.2.1} & 7 & 2\\
 &  & 2 & \texttt{0.1.8} & \texttt{3.3.1} & 8 & 1\\
T9.014 & \texttt{012345661} & 0 & \texttt{0.1.2.3.4.5.6} & \texttt{3.2.2.2.2.2.1} & 6 & 2\\*
 &  & 1 & \texttt{0.1.2.3.4.5.7} & \texttt{3.2.2.2.2.2.1} & 7 & 2\\
 &  & 2 & \texttt{0.8} & \texttt{3.1} & 8 & 1\\
T9.015 & \texttt{012345656} & 0 & \texttt{0.1.2.3.4.5.6} & \texttt{2.2.2.2.2.1.1} & 5 & 2\\*
 &  & 1 & \texttt{0.1.2.3.4.7.8} & \texttt{2.2.2.2.2.1.1} & 7 & 2\\
T9.016 & \texttt{012345655} & 0 & \texttt{0.1.2.3.4.5.6} & \texttt{3.3.3.3.3.1.1} & 5 & 1\\*
 &  & 1 & \texttt{0.1.2.3.4.7} & \texttt{3.3.3.3.3.1} & 7 & 2\\
 &  & 2 & \texttt{0.1.2.3.4.8} & \texttt{3.3.3.3.3.1} & 8 & 2\\
T9.017 & \texttt{012345654} & 0 & \texttt{0.1.2.3.4.5.6} & \texttt{3.3.3.3.2.1.1} & 5 & 1\\*
 &  & 1 & \texttt{0.1.2.3.4.7} & \texttt{3.3.3.3.2.1} & 7 & 1\\
 &  & 2 & \texttt{0.1.2.3.8} & \texttt{3.3.3.3.1} & 8 & 1\\
T9.018 & \texttt{012345653} & 0 & \texttt{0.1.2.3.4.5.6} & \texttt{3.3.3.2.2.1.1} & 5 & 1\\*
 &  & 1 & \texttt{0.1.2.3.4.7} & \texttt{3.3.3.2.2.1} & 7 & 1\\
 &  & 2 & \texttt{0.1.2.8} & \texttt{3.3.3.1} & 8 & 1\\
T9.019 & \texttt{012345652} & 0 & \texttt{0.1.2.3.4.5.6} & \texttt{3.3.2.2.2.1.1} & 5 & 1\\*
 &  & 1 & \texttt{0.1.2.3.4.7} & \texttt{3.3.2.2.2.1} & 7 & 1\\
 &  & 2 & \texttt{0.1.8} & \texttt{3.3.1} & 8 & 1\\
T9.020 & \texttt{012345651} & 0 & \texttt{0.1.2.3.4.5.6} & \texttt{3.2.2.2.2.1.1} & 5 & 1\\*
 &  & 1 & \texttt{0.1.2.3.4.7} & \texttt{3.2.2.2.2.1} & 7 & 1\\
 &  & 2 & \texttt{0.8} & \texttt{3.1} & 8 & 1\\
T9.021 & \texttt{012345645} & 0 & \texttt{0.1.2.3.4.5.6} & \texttt{2.2.2.2.1.1.1} & 4 & 1\\*
 &  & 1 & \texttt{0.1.2.3.7.8} & \texttt{2.2.2.2.1.1} & 7 & 1\\
T9.022 & \texttt{012345644} & 0 & \texttt{0.1.2.3.4.5.6} & \texttt{3.3.3.3.1.1.1} & 4 & 1\\*
 &  & 1 & \texttt{0.1.2.3.7} & \texttt{3.3.3.3.1} & 7 & 2\\
 &  & 2 & \texttt{0.1.2.3.8} & \texttt{3.3.3.3.1} & 8 & 2\\
T9.023 & \texttt{012345643} & 0 & \texttt{0.1.2.3.4.5.6} & \texttt{3.3.3.2.1.1.1} & 4 & 1\\*
 &  & 1 & \texttt{0.1.2.3.7} & \texttt{3.3.3.2.1} & 7 & 1\\
 &  & 2 & \texttt{0.1.2.8} & \texttt{3.3.3.1} & 8 & 1\\
T9.024 & \texttt{012345642} & 0 & \texttt{0.1.2.3.4.5.6} & \texttt{3.3.2.2.1.1.1} & 4 & 1\\*
 &  & 1 & \texttt{0.1.2.3.7} & \texttt{3.3.2.2.1} & 7 & 1\\
 &  & 2 & \texttt{0.1.8} & \texttt{3.3.1} & 8 & 1\\
T9.025 & \texttt{012345641} & 0 & \texttt{0.1.2.3.4.5.6} & \texttt{3.2.2.2.1.1.1} & 4 & 1\\*
 &  & 1 & \texttt{0.1.2.3.7} & \texttt{3.2.2.2.1} & 7 & 1\\
 &  & 2 & \texttt{0.8} & \texttt{3.1} & 8 & 1\\
T9.026 & \texttt{012345634} & 0 & \texttt{0.1.2.3.4.5.6} & \texttt{2.2.2.1.1.1.1} & 3 & 1\\*
 &  & 1 & \texttt{0.1.2.7.8} & \texttt{2.2.2.1.1} & 7 & 1\\
T9.027 & \texttt{012345633} & 0 & \texttt{0.1.2.3.4.5.6} & \texttt{3.3.3.1.1.1.1} & 3 & 1\\*
 &  & 1 & \texttt{0.1.2.7} & \texttt{3.3.3.1} & 7 & 2\\
 &  & 2 & \texttt{0.1.2.8} & \texttt{3.3.3.1} & 8 & 2\\
T9.028 & \texttt{012345632} & 0 & \texttt{0.1.2.3.4.5.6} & \texttt{3.3.2.1.1.1.1} & 3 & 1\\*
 &  & 1 & \texttt{0.1.2.7} & \texttt{3.3.2.1} & 7 & 1\\
 &  & 2 & \texttt{0.1.8} & \texttt{3.3.1} & 8 & 1\\
T9.029 & \texttt{012345631} & 0 & \texttt{0.1.2.3.4.5.6} & \texttt{3.2.2.1.1.1.1} & 3 & 1\\*
 &  & 1 & \texttt{0.1.2.7} & \texttt{3.2.2.1} & 7 & 1\\
 &  & 2 & \texttt{0.8} & \texttt{3.1} & 8 & 1\\
T9.030 & \texttt{012345623} & 0 & \texttt{0.1.2.3.4.5.6} & \texttt{2.2.1.1.1.1.1} & 2 & 1\\*
 &  & 1 & \texttt{0.1.7.8} & \texttt{2.2.1.1} & 7 & 1\\
T9.031 & \texttt{012345622} & 0 & \texttt{0.1.2.3.4.5.6} & \texttt{3.3.1.1.1.1.1} & 2 & 1\\*
 &  & 1 & \texttt{0.1.7} & \texttt{3.3.1} & 7 & 2\\
 &  & 2 & \texttt{0.1.8} & \texttt{3.3.1} & 8 & 2\\
T9.032 & \texttt{012345621} & 0 & \texttt{0.1.2.3.4.5.6} & \texttt{3.2.1.1.1.1.1} & 2 & 1\\*
 &  & 1 & \texttt{0.1.7} & \texttt{3.2.1} & 7 & 1\\
 &  & 2 & \texttt{0.8} & \texttt{3.1} & 8 & 1\\
T9.033 & \texttt{012345612} & 0 & \texttt{0.1.2.3.4.5.6} & \texttt{2.1.1.1.1.1.1} & 1 & 1\\*
 &  & 1 & \texttt{0.7.8} & \texttt{2.1.1} & 7 & 1\\
T9.034 & \texttt{012345611} & 0 & \texttt{0.1.2.3.4.5.6} & \texttt{3.1.1.1.1.1.1} & 1 & 1\\*
 &  & 1 & \texttt{0.7} & \texttt{3.1} & 7 & 2\\
 &  & 2 & \texttt{0.8} & \texttt{3.1} & 8 & 2\\
T9.035 & \texttt{012345555} & 0 & \texttt{0.1.2.3.4.5} & \texttt{4.4.4.4.4.1} & 5 & 4\\*
 &  & 1 & \texttt{0.1.2.3.4.6} & \texttt{4.4.4.4.4.1} & 6 & 4\\
 &  & 2 & \texttt{0.1.2.3.4.7} & \texttt{4.4.4.4.4.1} & 7 & 4\\
 &  & 3 & \texttt{0.1.2.3.4.8} & \texttt{4.4.4.4.4.1} & 8 & 4\\
T9.036 & \texttt{012345554} & 0 & \texttt{0.1.2.3.4.5} & \texttt{4.4.4.4.3.1} & 5 & 3\\*
 &  & 1 & \texttt{0.1.2.3.4.6} & \texttt{4.4.4.4.3.1} & 6 & 3\\
 &  & 2 & \texttt{0.1.2.3.4.7} & \texttt{4.4.4.4.3.1} & 7 & 3\\
 &  & 3 & \texttt{0.1.2.3.8} & \texttt{4.4.4.4.1} & 8 & 1\\
T9.037 & \texttt{012345553} & 0 & \texttt{0.1.2.3.4.5} & \texttt{4.4.4.3.3.1} & 5 & 3\\*
 &  & 1 & \texttt{0.1.2.3.4.6} & \texttt{4.4.4.3.3.1} & 6 & 3\\
 &  & 2 & \texttt{0.1.2.3.4.7} & \texttt{4.4.4.3.3.1} & 7 & 3\\
 &  & 3 & \texttt{0.1.2.8} & \texttt{4.4.4.1} & 8 & 1\\
T9.038 & \texttt{012345552} & 0 & \texttt{0.1.2.3.4.5} & \texttt{4.4.3.3.3.1} & 5 & 3\\*
 &  & 1 & \texttt{0.1.2.3.4.6} & \texttt{4.4.3.3.3.1} & 6 & 3\\
 &  & 2 & \texttt{0.1.2.3.4.7} & \texttt{4.4.3.3.3.1} & 7 & 3\\
 &  & 3 & \texttt{0.1.8} & \texttt{4.4.1} & 8 & 1\\
T9.039 & \texttt{012345551} & 0 & \texttt{0.1.2.3.4.5} & \texttt{4.3.3.3.3.1} & 5 & 3\\*
 &  & 1 & \texttt{0.1.2.3.4.6} & \texttt{4.3.3.3.3.1} & 6 & 3\\
 &  & 2 & \texttt{0.1.2.3.4.7} & \texttt{4.3.3.3.3.1} & 7 & 3\\
 &  & 3 & \texttt{0.8} & \texttt{4.1} & 8 & 1\\
T9.040 & \texttt{012345545} & 0 & \texttt{0.1.2.3.4.5} & \texttt{3.3.3.3.2.1} & 5 & 2\\*
 &  & 1 & \texttt{0.1.2.3.4.6} & \texttt{3.3.3.3.2.1} & 6 & 2\\
 &  & 2 & \texttt{0.1.2.3.7.8} & \texttt{3.3.3.3.1.1} & 7 & 1\\
T9.041 & \texttt{012345544} & 0 & \texttt{0.1.2.3.4.5} & \texttt{4.4.4.4.2.1} & 5 & 2\\*
 &  & 1 & \texttt{0.1.2.3.4.6} & \texttt{4.4.4.4.2.1} & 6 & 2\\
 &  & 2 & \texttt{0.1.2.3.7} & \texttt{4.4.4.4.1} & 7 & 2\\
 &  & 3 & \texttt{0.1.2.3.8} & \texttt{4.4.4.4.1} & 8 & 2\\
T9.042 & \texttt{012345543} & 0 & \texttt{0.1.2.3.4.5} & \texttt{4.4.4.3.2.1} & 5 & 2\\*
 &  & 1 & \texttt{0.1.2.3.4.6} & \texttt{4.4.4.3.2.1} & 6 & 2\\
 &  & 2 & \texttt{0.1.2.3.7} & \texttt{4.4.4.3.1} & 7 & 1\\
 &  & 3 & \texttt{0.1.2.8} & \texttt{4.4.4.1} & 8 & 1\\
T9.043 & \texttt{012345542} & 0 & \texttt{0.1.2.3.4.5} & \texttt{4.4.3.3.2.1} & 5 & 2\\*
 &  & 1 & \texttt{0.1.2.3.4.6} & \texttt{4.4.3.3.2.1} & 6 & 2\\
 &  & 2 & \texttt{0.1.2.3.7} & \texttt{4.4.3.3.1} & 7 & 1\\
 &  & 3 & \texttt{0.1.8} & \texttt{4.4.1} & 8 & 1\\
T9.044 & \texttt{012345541} & 0 & \texttt{0.1.2.3.4.5} & \texttt{4.3.3.3.2.1} & 5 & 2\\*
 &  & 1 & \texttt{0.1.2.3.4.6} & \texttt{4.3.3.3.2.1} & 6 & 2\\
 &  & 2 & \texttt{0.1.2.3.7} & \texttt{4.3.3.3.1} & 7 & 1\\
 &  & 3 & \texttt{0.8} & \texttt{4.1} & 8 & 1\\
T9.045 & \texttt{012345534} & 0 & \texttt{0.1.2.3.4.5} & \texttt{3.3.3.2.2.1} & 5 & 2\\*
 &  & 1 & \texttt{0.1.2.3.4.6} & \texttt{3.3.3.2.2.1} & 6 & 2\\
 &  & 2 & \texttt{0.1.2.7.8} & \texttt{3.3.3.1.1} & 7 & 1\\
T9.046 & \texttt{012345533} & 0 & \texttt{0.1.2.3.4.5} & \texttt{4.4.4.2.2.1} & 5 & 2\\*
 &  & 1 & \texttt{0.1.2.3.4.6} & \texttt{4.4.4.2.2.1} & 6 & 2\\
 &  & 2 & \texttt{0.1.2.7} & \texttt{4.4.4.1} & 7 & 2\\
 &  & 3 & \texttt{0.1.2.8} & \texttt{4.4.4.1} & 8 & 2\\
T9.047 & \texttt{012345532} & 0 & \texttt{0.1.2.3.4.5} & \texttt{4.4.3.2.2.1} & 5 & 2\\*
 &  & 1 & \texttt{0.1.2.3.4.6} & \texttt{4.4.3.2.2.1} & 6 & 2\\
 &  & 2 & \texttt{0.1.2.7} & \texttt{4.4.3.1} & 7 & 1\\
 &  & 3 & \texttt{0.1.8} & \texttt{4.4.1} & 8 & 1\\
T9.048 & \texttt{012345531} & 0 & \texttt{0.1.2.3.4.5} & \texttt{4.3.3.2.2.1} & 5 & 2\\*
 &  & 1 & \texttt{0.1.2.3.4.6} & \texttt{4.3.3.2.2.1} & 6 & 2\\
 &  & 2 & \texttt{0.1.2.7} & \texttt{4.3.3.1} & 7 & 1\\
 &  & 3 & \texttt{0.8} & \texttt{4.1} & 8 & 1\\
T9.049 & \texttt{012345523} & 0 & \texttt{0.1.2.3.4.5} & \texttt{3.3.2.2.2.1} & 5 & 2\\*
 &  & 1 & \texttt{0.1.2.3.4.6} & \texttt{3.3.2.2.2.1} & 6 & 2\\
 &  & 2 & \texttt{0.1.7.8} & \texttt{3.3.1.1} & 7 & 1\\
T9.050 & \texttt{012345522} & 0 & \texttt{0.1.2.3.4.5} & \texttt{4.4.2.2.2.1} & 5 & 2\\*
 &  & 1 & \texttt{0.1.2.3.4.6} & \texttt{4.4.2.2.2.1} & 6 & 2\\
 &  & 2 & \texttt{0.1.7} & \texttt{4.4.1} & 7 & 2\\
 &  & 3 & \texttt{0.1.8} & \texttt{4.4.1} & 8 & 2\\
T9.051 & \texttt{012345521} & 0 & \texttt{0.1.2.3.4.5} & \texttt{4.3.2.2.2.1} & 5 & 2\\*
 &  & 1 & \texttt{0.1.2.3.4.6} & \texttt{4.3.2.2.2.1} & 6 & 2\\
 &  & 2 & \texttt{0.1.7} & \texttt{4.3.1} & 7 & 1\\
 &  & 3 & \texttt{0.8} & \texttt{4.1} & 8 & 1\\
T9.052 & \texttt{012345512} & 0 & \texttt{0.1.2.3.4.5} & \texttt{3.2.2.2.2.1} & 5 & 2\\*
 &  & 1 & \texttt{0.1.2.3.4.6} & \texttt{3.2.2.2.2.1} & 6 & 2\\
 &  & 2 & \texttt{0.7.8} & \texttt{3.1.1} & 7 & 1\\
T9.053 & \texttt{012345511} & 0 & \texttt{0.1.2.3.4.5} & \texttt{4.2.2.2.2.1} & 5 & 2\\*
 &  & 1 & \texttt{0.1.2.3.4.6} & \texttt{4.2.2.2.2.1} & 6 & 2\\
 &  & 2 & \texttt{0.7} & \texttt{4.1} & 7 & 2\\
 &  & 3 & \texttt{0.8} & \texttt{4.1} & 8 & 2\\
T9.054 & \texttt{012345454} & 0 & \texttt{0.1.2.3.4.5} & \texttt{3.3.3.3.1.1} & 4 & 2\\*
 &  & 1 & \texttt{0.1.2.3.6.7} & \texttt{3.3.3.3.1.1} & 6 & 2\\
 &  & 2 & \texttt{0.1.2.3.8} & \texttt{3.3.3.3.1} & 8 & 1\\
T9.055 & \texttt{012345453} & 0 & \texttt{0.1.2.3.4.5} & \texttt{3.3.3.2.1.1} & 4 & 2\\*
 &  & 1 & \texttt{0.1.2.3.6.7} & \texttt{3.3.3.2.1.1} & 6 & 2\\
 &  & 2 & \texttt{0.1.2.8} & \texttt{3.3.3.1} & 8 & 1\\
T9.056 & \texttt{012345452} & 0 & \texttt{0.1.2.3.4.5} & \texttt{3.3.2.2.1.1} & 4 & 2\\*
 &  & 1 & \texttt{0.1.2.3.6.7} & \texttt{3.3.2.2.1.1} & 6 & 2\\
 &  & 2 & \texttt{0.1.8} & \texttt{3.3.1} & 8 & 1\\
T9.057 & \texttt{012345451} & 0 & \texttt{0.1.2.3.4.5} & \texttt{3.2.2.2.1.1} & 4 & 2\\*
 &  & 1 & \texttt{0.1.2.3.6.7} & \texttt{3.2.2.2.1.1} & 6 & 2\\
 &  & 2 & \texttt{0.8} & \texttt{3.1} & 8 & 1\\
T9.058 & \texttt{012345444} & 0 & \texttt{0.1.2.3.4.5} & \texttt{4.4.4.4.1.1} & 4 & 1\\*
 &  & 1 & \texttt{0.1.2.3.6} & \texttt{4.4.4.4.1} & 6 & 3\\
 &  & 2 & \texttt{0.1.2.3.7} & \texttt{4.4.4.4.1} & 7 & 3\\
 &  & 3 & \texttt{0.1.2.3.8} & \texttt{4.4.4.4.1} & 8 & 3\\
T9.059 & \texttt{012345443} & 0 & \texttt{0.1.2.3.4.5} & \texttt{4.4.4.3.1.1} & 4 & 1\\*
 &  & 1 & \texttt{0.1.2.3.6} & \texttt{4.4.4.3.1} & 6 & 2\\
 &  & 2 & \texttt{0.1.2.3.7} & \texttt{4.4.4.3.1} & 7 & 2\\
 &  & 3 & \texttt{0.1.2.8} & \texttt{4.4.4.1} & 8 & 1\\
T9.060 & \texttt{012345442} & 0 & \texttt{0.1.2.3.4.5} & \texttt{4.4.3.3.1.1} & 4 & 1\\*
 &  & 1 & \texttt{0.1.2.3.6} & \texttt{4.4.3.3.1} & 6 & 2\\
 &  & 2 & \texttt{0.1.2.3.7} & \texttt{4.4.3.3.1} & 7 & 2\\
 &  & 3 & \texttt{0.1.8} & \texttt{4.4.1} & 8 & 1\\
T9.061 & \texttt{012345441} & 0 & \texttt{0.1.2.3.4.5} & \texttt{4.3.3.3.1.1} & 4 & 1\\*
 &  & 1 & \texttt{0.1.2.3.6} & \texttt{4.3.3.3.1} & 6 & 2\\
 &  & 2 & \texttt{0.1.2.3.7} & \texttt{4.3.3.3.1} & 7 & 2\\
 &  & 3 & \texttt{0.8} & \texttt{4.1} & 8 & 1\\
T9.062 & \texttt{012345434} & 0 & \texttt{0.1.2.3.4.5} & \texttt{3.3.3.2.1.1} & 4 & 1\\*
 &  & 1 & \texttt{0.1.2.3.6} & \texttt{3.3.3.2.1} & 6 & 1\\
 &  & 2 & \texttt{0.1.2.7.8} & \texttt{3.3.3.1.1} & 7 & 1\\
T9.063 & \texttt{012345433} & 0 & \texttt{0.1.2.3.4.5} & \texttt{4.4.4.2.1.1} & 4 & 1\\*
 &  & 1 & \texttt{0.1.2.3.6} & \texttt{4.4.4.2.1} & 6 & 1\\
 &  & 2 & \texttt{0.1.2.7} & \texttt{4.4.4.1} & 7 & 2\\
 &  & 3 & \texttt{0.1.2.8} & \texttt{4.4.4.1} & 8 & 2\\
T9.064 & \texttt{012345432} & 0 & \texttt{0.1.2.3.4.5} & \texttt{4.4.3.2.1.1} & 4 & 1\\*
 &  & 1 & \texttt{0.1.2.3.6} & \texttt{4.4.3.2.1} & 6 & 1\\
 &  & 2 & \texttt{0.1.2.7} & \texttt{4.4.3.1} & 7 & 1\\
 &  & 3 & \texttt{0.1.8} & \texttt{4.4.1} & 8 & 1\\
T9.065 & \texttt{012345431} & 0 & \texttt{0.1.2.3.4.5} & \texttt{4.3.3.2.1.1} & 4 & 1\\*
 &  & 1 & \texttt{0.1.2.3.6} & \texttt{4.3.3.2.1} & 6 & 1\\
 &  & 2 & \texttt{0.1.2.7} & \texttt{4.3.3.1} & 7 & 1\\
 &  & 3 & \texttt{0.8} & \texttt{4.1} & 8 & 1\\
T9.066 & \texttt{012345423} & 0 & \texttt{0.1.2.3.4.5} & \texttt{3.3.2.2.1.1} & 4 & 1\\*
 &  & 1 & \texttt{0.1.2.3.6} & \texttt{3.3.2.2.1} & 6 & 1\\
 &  & 2 & \texttt{0.1.7.8} & \texttt{3.3.1.1} & 7 & 1\\
T9.067 & \texttt{012345422} & 0 & \texttt{0.1.2.3.4.5} & \texttt{4.4.2.2.1.1} & 4 & 1\\*
 &  & 1 & \texttt{0.1.2.3.6} & \texttt{4.4.2.2.1} & 6 & 1\\
 &  & 2 & \texttt{0.1.7} & \texttt{4.4.1} & 7 & 2\\
 &  & 3 & \texttt{0.1.8} & \texttt{4.4.1} & 8 & 2\\
T9.068 & \texttt{012345421} & 0 & \texttt{0.1.2.3.4.5} & \texttt{4.3.2.2.1.1} & 4 & 1\\*
 &  & 1 & \texttt{0.1.2.3.6} & \texttt{4.3.2.2.1} & 6 & 1\\
 &  & 2 & \texttt{0.1.7} & \texttt{4.3.1} & 7 & 1\\
 &  & 3 & \texttt{0.8} & \texttt{4.1} & 8 & 1\\
T9.069 & \texttt{012345412} & 0 & \texttt{0.1.2.3.4.5} & \texttt{3.2.2.2.1.1} & 4 & 1\\*
 &  & 1 & \texttt{0.1.2.3.6} & \texttt{3.2.2.2.1} & 6 & 1\\
 &  & 2 & \texttt{0.7.8} & \texttt{3.1.1} & 7 & 1\\
T9.070 & \texttt{012345411} & 0 & \texttt{0.1.2.3.4.5} & \texttt{4.2.2.2.1.1} & 4 & 1\\*
 &  & 1 & \texttt{0.1.2.3.6} & \texttt{4.2.2.2.1} & 6 & 1\\
 &  & 2 & \texttt{0.7} & \texttt{4.1} & 7 & 2\\
 &  & 3 & \texttt{0.8} & \texttt{4.1} & 8 & 2\\
T9.071 & \texttt{012345345} & 0 & \texttt{0.1.2.3.4.5} & \texttt{2.2.2.1.1.1} & 3 & 2\\*
 &  & 1 & \texttt{0.1.2.6.7.8} & \texttt{2.2.2.1.1.1} & 6 & 2\\
T9.072 & \texttt{012345344} & 0 & \texttt{0.1.2.3.4.5} & \texttt{3.3.3.1.1.1} & 3 & 1\\*
 &  & 1 & \texttt{0.1.2.6.7} & \texttt{3.3.3.2.1} & 7 & 2\\
 &  & 2 & \texttt{0.1.2.6.8} & \texttt{3.3.3.2.1} & 8 & 2\\
T9.073 & \texttt{012345343} & 0 & \texttt{0.1.2.3.4.5} & \texttt{3.3.3.1.1.1} & 3 & 1\\*
 &  & 1 & \texttt{0.1.2.6.7} & \texttt{3.3.3.1.1} & 6 & 1\\
 &  & 2 & \texttt{0.1.2.8} & \texttt{3.3.3.1} & 8 & 1\\
T9.074 & \texttt{012345342} & 0 & \texttt{0.1.2.3.4.5} & \texttt{3.3.2.1.1.1} & 3 & 1\\*
 &  & 1 & \texttt{0.1.2.6.7} & \texttt{3.3.2.1.1} & 6 & 1\\
 &  & 2 & \texttt{0.1.8} & \texttt{3.3.1} & 8 & 1\\
T9.075 & \texttt{012345341} & 0 & \texttt{0.1.2.3.4.5} & \texttt{3.2.2.1.1.1} & 3 & 1\\*
 &  & 1 & \texttt{0.1.2.6.7} & \texttt{3.2.2.1.1} & 6 & 1\\
 &  & 2 & \texttt{0.8} & \texttt{3.1} & 8 & 1\\
T9.076 & \texttt{012345333} & 0 & \texttt{0.1.2.3.4.5} & \texttt{4.4.4.1.1.1} & 3 & 1\\*
 &  & 1 & \texttt{0.1.2.6} & \texttt{4.4.4.1} & 6 & 3\\
 &  & 2 & \texttt{0.1.2.7} & \texttt{4.4.4.1} & 7 & 3\\
 &  & 3 & \texttt{0.1.2.8} & \texttt{4.4.4.1} & 8 & 3\\
T9.077 & \texttt{012345332} & 0 & \texttt{0.1.2.3.4.5} & \texttt{4.4.3.1.1.1} & 3 & 1\\*
 &  & 1 & \texttt{0.1.2.6} & \texttt{4.4.3.1} & 6 & 2\\
 &  & 2 & \texttt{0.1.2.7} & \texttt{4.4.3.1} & 7 & 2\\
 &  & 3 & \texttt{0.1.8} & \texttt{4.4.1} & 8 & 1\\
T9.078 & \texttt{012345331} & 0 & \texttt{0.1.2.3.4.5} & \texttt{4.3.3.1.1.1} & 3 & 1\\*
 &  & 1 & \texttt{0.1.2.6} & \texttt{4.3.3.1} & 6 & 2\\
 &  & 2 & \texttt{0.1.2.7} & \texttt{4.3.3.1} & 7 & 2\\
 &  & 3 & \texttt{0.8} & \texttt{4.1} & 8 & 1\\
T9.079 & \texttt{012345323} & 0 & \texttt{0.1.2.3.4.5} & \texttt{3.3.2.1.1.1} & 3 & 1\\*
 &  & 1 & \texttt{0.1.2.6} & \texttt{3.3.2.1} & 6 & 1\\
 &  & 2 & \texttt{0.1.7.8} & \texttt{3.3.1.1} & 7 & 1\\
T9.080 & \texttt{012345322} & 0 & \texttt{0.1.2.3.4.5} & \texttt{4.4.2.1.1.1} & 3 & 1\\*
 &  & 1 & \texttt{0.1.2.6} & \texttt{4.4.2.1} & 6 & 1\\
 &  & 2 & \texttt{0.1.7} & \texttt{4.4.1} & 7 & 2\\
 &  & 3 & \texttt{0.1.8} & \texttt{4.4.1} & 8 & 2\\
T9.081 & \texttt{012345321} & 0 & \texttt{0.1.2.3.4.5} & \texttt{4.3.2.1.1.1} & 3 & 1\\*
 &  & 1 & \texttt{0.1.2.6} & \texttt{4.3.2.1} & 6 & 1\\
 &  & 2 & \texttt{0.1.7} & \texttt{4.3.1} & 7 & 1\\
 &  & 3 & \texttt{0.8} & \texttt{4.1} & 8 & 1\\
T9.082 & \texttt{012345312} & 0 & \texttt{0.1.2.3.4.5} & \texttt{3.2.2.1.1.1} & 3 & 1\\*
 &  & 1 & \texttt{0.1.2.6} & \texttt{3.2.2.1} & 6 & 1\\
 &  & 2 & \texttt{0.7.8} & \texttt{3.1.1} & 7 & 1\\
T9.083 & \texttt{012345311} & 0 & \texttt{0.1.2.3.4.5} & \texttt{4.2.2.1.1.1} & 3 & 1\\*
 &  & 1 & \texttt{0.1.2.6} & \texttt{4.2.2.1} & 6 & 1\\
 &  & 2 & \texttt{0.7} & \texttt{4.1} & 7 & 2\\
 &  & 3 & \texttt{0.8} & \texttt{4.1} & 8 & 2\\
T9.084 & \texttt{012345234} & 0 & \texttt{0.1.2.3.4.5} & \texttt{2.2.1.1.1.1} & 2 & 1\\*
 &  & 1 & \texttt{0.1.6.7.8} & \texttt{2.2.1.1.1} & 6 & 1\\
T9.085 & \texttt{012345233} & 0 & \texttt{0.1.2.3.4.5} & \texttt{3.3.1.1.1.1} & 2 & 1\\*
 &  & 1 & \texttt{0.1.6.7} & \texttt{3.3.2.1} & 7 & 2\\
 &  & 2 & \texttt{0.1.6.8} & \texttt{3.3.2.1} & 8 & 2\\
T9.086 & \texttt{012345232} & 0 & \texttt{0.1.2.3.4.5} & \texttt{3.3.1.1.1.1} & 2 & 1\\*
 &  & 1 & \texttt{0.1.6.7} & \texttt{3.3.1.1} & 6 & 1\\
 &  & 2 & \texttt{0.1.8} & \texttt{3.3.1} & 8 & 1\\
T9.087 & \texttt{012345231} & 0 & \texttt{0.1.2.3.4.5} & \texttt{3.2.1.1.1.1} & 2 & 1\\*
 &  & 1 & \texttt{0.1.6.7} & \texttt{3.2.1.1} & 6 & 1\\
 &  & 2 & \texttt{0.8} & \texttt{3.1} & 8 & 1\\
T9.088 & \texttt{012345222} & 0 & \texttt{0.1.2.3.4.5} & \texttt{4.4.1.1.1.1} & 2 & 1\\*
 &  & 1 & \texttt{0.1.6} & \texttt{4.4.1} & 6 & 3\\
 &  & 2 & \texttt{0.1.7} & \texttt{4.4.1} & 7 & 3\\
 &  & 3 & \texttt{0.1.8} & \texttt{4.4.1} & 8 & 3\\
T9.089 & \texttt{012345221} & 0 & \texttt{0.1.2.3.4.5} & \texttt{4.3.1.1.1.1} & 2 & 1\\*
 &  & 1 & \texttt{0.1.6} & \texttt{4.3.1} & 6 & 2\\
 &  & 2 & \texttt{0.1.7} & \texttt{4.3.1} & 7 & 2\\
 &  & 3 & \texttt{0.8} & \texttt{4.1} & 8 & 1\\
T9.090 & \texttt{012345212} & 0 & \texttt{0.1.2.3.4.5} & \texttt{3.2.1.1.1.1} & 2 & 1\\*
 &  & 1 & \texttt{0.1.6} & \texttt{3.2.1} & 6 & 1\\
 &  & 2 & \texttt{0.7.8} & \texttt{3.1.1} & 7 & 1\\
T9.091 & \texttt{012345211} & 0 & \texttt{0.1.2.3.4.5} & \texttt{4.2.1.1.1.1} & 2 & 1\\*
 &  & 1 & \texttt{0.1.6} & \texttt{4.2.1} & 6 & 1\\
 &  & 2 & \texttt{0.7} & \texttt{4.1} & 7 & 2\\
 &  & 3 & \texttt{0.8} & \texttt{4.1} & 8 & 2\\
T9.092 & \texttt{012345123} & 0 & \texttt{0.1.2.3.4.5} & \texttt{2.1.1.1.1.1} & 1 & 1\\*
 &  & 1 & \texttt{0.6.7.8} & \texttt{2.1.1.1} & 6 & 1\\
T9.093 & \texttt{012345122} & 0 & \texttt{0.1.2.3.4.5} & \texttt{3.1.1.1.1.1} & 1 & 1\\*
 &  & 1 & \texttt{0.6.7} & \texttt{3.2.1} & 7 & 2\\
 &  & 2 & \texttt{0.6.8} & \texttt{3.2.1} & 8 & 2\\
T9.094 & \texttt{012345121} & 0 & \texttt{0.1.2.3.4.5} & \texttt{3.1.1.1.1.1} & 1 & 1\\*
 &  & 1 & \texttt{0.6.7} & \texttt{3.1.1} & 6 & 1\\
 &  & 2 & \texttt{0.8} & \texttt{3.1} & 8 & 1\\
T9.095 & \texttt{012345111} & 0 & \texttt{0.1.2.3.4.5} & \texttt{4.1.1.1.1.1} & 1 & 1\\*
 &  & 1 & \texttt{0.6} & \texttt{4.1} & 6 & 3\\
 &  & 2 & \texttt{0.7} & \texttt{4.1} & 7 & 3\\
 &  & 3 & \texttt{0.8} & \texttt{4.1} & 8 & 3\\
T9.096 & \texttt{012344444} & 0 & \texttt{0.1.2.3.4} & \texttt{5.5.5.5.1} & 4 & 5\\*
 &  & 1 & \texttt{0.1.2.3.5} & \texttt{5.5.5.5.1} & 5 & 5\\
 &  & 2 & \texttt{0.1.2.3.6} & \texttt{5.5.5.5.1} & 6 & 5\\
 &  & 3 & \texttt{0.1.2.3.7} & \texttt{5.5.5.5.1} & 7 & 5\\
 &  & 4 & \texttt{0.1.2.3.8} & \texttt{5.5.5.5.1} & 8 & 5\\
T9.097 & \texttt{012344443} & 0 & \texttt{0.1.2.3.4} & \texttt{5.5.5.4.1} & 4 & 4\\*
 &  & 1 & \texttt{0.1.2.3.5} & \texttt{5.5.5.4.1} & 5 & 4\\
 &  & 2 & \texttt{0.1.2.3.6} & \texttt{5.5.5.4.1} & 6 & 4\\
 &  & 3 & \texttt{0.1.2.3.7} & \texttt{5.5.5.4.1} & 7 & 4\\
 &  & 4 & \texttt{0.1.2.8} & \texttt{5.5.5.1} & 8 & 1\\
T9.098 & \texttt{012344442} & 0 & \texttt{0.1.2.3.4} & \texttt{5.5.4.4.1} & 4 & 4\\*
 &  & 1 & \texttt{0.1.2.3.5} & \texttt{5.5.4.4.1} & 5 & 4\\
 &  & 2 & \texttt{0.1.2.3.6} & \texttt{5.5.4.4.1} & 6 & 4\\
 &  & 3 & \texttt{0.1.2.3.7} & \texttt{5.5.4.4.1} & 7 & 4\\
 &  & 4 & \texttt{0.1.8} & \texttt{5.5.1} & 8 & 1\\
T9.099 & \texttt{012344441} & 0 & \texttt{0.1.2.3.4} & \texttt{5.4.4.4.1} & 4 & 4\\*
 &  & 1 & \texttt{0.1.2.3.5} & \texttt{5.4.4.4.1} & 5 & 4\\
 &  & 2 & \texttt{0.1.2.3.6} & \texttt{5.4.4.4.1} & 6 & 4\\
 &  & 3 & \texttt{0.1.2.3.7} & \texttt{5.4.4.4.1} & 7 & 4\\
 &  & 4 & \texttt{0.8} & \texttt{5.1} & 8 & 1\\
T9.100 & \texttt{012344434} & 0 & \texttt{0.1.2.3.4} & \texttt{4.4.4.3.1} & 4 & 3\\*
 &  & 1 & \texttt{0.1.2.3.5} & \texttt{4.4.4.3.1} & 5 & 3\\
 &  & 2 & \texttt{0.1.2.3.6} & \texttt{4.4.4.3.1} & 6 & 3\\
 &  & 3 & \texttt{0.1.2.7.8} & \texttt{4.4.4.1.1} & 7 & 1\\
T9.101 & \texttt{012344433} & 0 & \texttt{0.1.2.3.4} & \texttt{5.5.5.3.1} & 4 & 3\\*
 &  & 1 & \texttt{0.1.2.3.5} & \texttt{5.5.5.3.1} & 5 & 3\\
 &  & 2 & \texttt{0.1.2.3.6} & \texttt{5.5.5.3.1} & 6 & 3\\
 &  & 3 & \texttt{0.1.2.7} & \texttt{5.5.5.1} & 7 & 2\\
 &  & 4 & \texttt{0.1.2.8} & \texttt{5.5.5.1} & 8 & 2\\
T9.102 & \texttt{012344432} & 0 & \texttt{0.1.2.3.4} & \texttt{5.5.4.3.1} & 4 & 3\\*
 &  & 1 & \texttt{0.1.2.3.5} & \texttt{5.5.4.3.1} & 5 & 3\\
 &  & 2 & \texttt{0.1.2.3.6} & \texttt{5.5.4.3.1} & 6 & 3\\
 &  & 3 & \texttt{0.1.2.7} & \texttt{5.5.4.1} & 7 & 1\\
 &  & 4 & \texttt{0.1.8} & \texttt{5.5.1} & 8 & 1\\
T9.103 & \texttt{012344431} & 0 & \texttt{0.1.2.3.4} & \texttt{5.4.4.3.1} & 4 & 3\\*
 &  & 1 & \texttt{0.1.2.3.5} & \texttt{5.4.4.3.1} & 5 & 3\\
 &  & 2 & \texttt{0.1.2.3.6} & \texttt{5.4.4.3.1} & 6 & 3\\
 &  & 3 & \texttt{0.1.2.7} & \texttt{5.4.4.1} & 7 & 1\\
 &  & 4 & \texttt{0.8} & \texttt{5.1} & 8 & 1\\
T9.104 & \texttt{012344423} & 0 & \texttt{0.1.2.3.4} & \texttt{4.4.3.3.1} & 4 & 3\\*
 &  & 1 & \texttt{0.1.2.3.5} & \texttt{4.4.3.3.1} & 5 & 3\\
 &  & 2 & \texttt{0.1.2.3.6} & \texttt{4.4.3.3.1} & 6 & 3\\
 &  & 3 & \texttt{0.1.7.8} & \texttt{4.4.1.1} & 7 & 1\\
T9.105 & \texttt{012344422} & 0 & \texttt{0.1.2.3.4} & \texttt{5.5.3.3.1} & 4 & 3\\*
 &  & 1 & \texttt{0.1.2.3.5} & \texttt{5.5.3.3.1} & 5 & 3\\
 &  & 2 & \texttt{0.1.2.3.6} & \texttt{5.5.3.3.1} & 6 & 3\\
 &  & 3 & \texttt{0.1.7} & \texttt{5.5.1} & 7 & 2\\
 &  & 4 & \texttt{0.1.8} & \texttt{5.5.1} & 8 & 2\\
T9.106 & \texttt{012344421} & 0 & \texttt{0.1.2.3.4} & \texttt{5.4.3.3.1} & 4 & 3\\*
 &  & 1 & \texttt{0.1.2.3.5} & \texttt{5.4.3.3.1} & 5 & 3\\
 &  & 2 & \texttt{0.1.2.3.6} & \texttt{5.4.3.3.1} & 6 & 3\\
 &  & 3 & \texttt{0.1.7} & \texttt{5.4.1} & 7 & 1\\
 &  & 4 & \texttt{0.8} & \texttt{5.1} & 8 & 1\\
T9.107 & \texttt{012344412} & 0 & \texttt{0.1.2.3.4} & \texttt{4.3.3.3.1} & 4 & 3\\*
 &  & 1 & \texttt{0.1.2.3.5} & \texttt{4.3.3.3.1} & 5 & 3\\
 &  & 2 & \texttt{0.1.2.3.6} & \texttt{4.3.3.3.1} & 6 & 3\\
 &  & 3 & \texttt{0.7.8} & \texttt{4.1.1} & 7 & 1\\
T9.108 & \texttt{012344411} & 0 & \texttt{0.1.2.3.4} & \texttt{5.3.3.3.1} & 4 & 3\\*
 &  & 1 & \texttt{0.1.2.3.5} & \texttt{5.3.3.3.1} & 5 & 3\\
 &  & 2 & \texttt{0.1.2.3.6} & \texttt{5.3.3.3.1} & 6 & 3\\
 &  & 3 & \texttt{0.7} & \texttt{5.1} & 7 & 2\\
 &  & 4 & \texttt{0.8} & \texttt{5.1} & 8 & 2\\
T9.109 & \texttt{012344344} & 0 & \texttt{0.1.2.3.4} & \texttt{4.4.4.2.1} & 4 & 4\\*
 &  & 1 & \texttt{0.1.2.3.5} & \texttt{4.4.4.2.1} & 5 & 4\\
 &  & 2 & \texttt{0.1.2.6.7} & \texttt{4.4.4.2.1} & 7 & 4\\
 &  & 3 & \texttt{0.1.2.6.8} & \texttt{4.4.4.2.1} & 8 & 4\\
T9.110 & \texttt{012344343} & 0 & \texttt{0.1.2.3.4} & \texttt{4.4.4.2.1} & 4 & 2\\*
 &  & 1 & \texttt{0.1.2.3.5} & \texttt{4.4.4.2.1} & 5 & 2\\
 &  & 2 & \texttt{0.1.2.6.7} & \texttt{4.4.4.1.1} & 6 & 1\\
 &  & 3 & \texttt{0.1.2.8} & \texttt{4.4.4.1} & 8 & 1\\
T9.111 & \texttt{012344342} & 0 & \texttt{0.1.2.3.4} & \texttt{4.4.3.2.1} & 4 & 2\\*
 &  & 1 & \texttt{0.1.2.3.5} & \texttt{4.4.3.2.1} & 5 & 2\\
 &  & 2 & \texttt{0.1.2.6.7} & \texttt{4.4.3.1.1} & 6 & 1\\
 &  & 3 & \texttt{0.1.8} & \texttt{4.4.1} & 8 & 1\\
T9.112 & \texttt{012344341} & 0 & \texttt{0.1.2.3.4} & \texttt{4.3.3.2.1} & 4 & 2\\*
 &  & 1 & \texttt{0.1.2.3.5} & \texttt{4.3.3.2.1} & 5 & 2\\
 &  & 2 & \texttt{0.1.2.6.7} & \texttt{4.3.3.1.1} & 6 & 1\\
 &  & 3 & \texttt{0.8} & \texttt{4.1} & 8 & 1\\
T9.113 & \texttt{012344333} & 0 & \texttt{0.1.2.3.4} & \texttt{5.5.5.2.1} & 4 & 2\\*
 &  & 1 & \texttt{0.1.2.3.5} & \texttt{5.5.5.2.1} & 5 & 2\\
 &  & 2 & \texttt{0.1.2.6} & \texttt{5.5.5.1} & 6 & 3\\
 &  & 3 & \texttt{0.1.2.7} & \texttt{5.5.5.1} & 7 & 3\\
 &  & 4 & \texttt{0.1.2.8} & \texttt{5.5.5.1} & 8 & 3\\
T9.114 & \texttt{012344332} & 0 & \texttt{0.1.2.3.4} & \texttt{5.5.4.2.1} & 4 & 2\\*
 &  & 1 & \texttt{0.1.2.3.5} & \texttt{5.5.4.2.1} & 5 & 2\\
 &  & 2 & \texttt{0.1.2.6} & \texttt{5.5.4.1} & 6 & 2\\
 &  & 3 & \texttt{0.1.2.7} & \texttt{5.5.4.1} & 7 & 2\\
 &  & 4 & \texttt{0.1.8} & \texttt{5.5.1} & 8 & 1\\
T9.115 & \texttt{012344331} & 0 & \texttt{0.1.2.3.4} & \texttt{5.4.4.2.1} & 4 & 2\\*
 &  & 1 & \texttt{0.1.2.3.5} & \texttt{5.4.4.2.1} & 5 & 2\\
 &  & 2 & \texttt{0.1.2.6} & \texttt{5.4.4.1} & 6 & 2\\
 &  & 3 & \texttt{0.1.2.7} & \texttt{5.4.4.1} & 7 & 2\\
 &  & 4 & \texttt{0.8} & \texttt{5.1} & 8 & 1\\
T9.116 & \texttt{012344323} & 0 & \texttt{0.1.2.3.4} & \texttt{4.4.3.2.1} & 4 & 2\\*
 &  & 1 & \texttt{0.1.2.3.5} & \texttt{4.4.3.2.1} & 5 & 2\\
 &  & 2 & \texttt{0.1.2.6} & \texttt{4.4.3.1} & 6 & 1\\
 &  & 3 & \texttt{0.1.7.8} & \texttt{4.4.1.1} & 7 & 1\\
T9.117 & \texttt{012344322} & 0 & \texttt{0.1.2.3.4} & \texttt{5.5.3.2.1} & 4 & 2\\*
 &  & 1 & \texttt{0.1.2.3.5} & \texttt{5.5.3.2.1} & 5 & 2\\
 &  & 2 & \texttt{0.1.2.6} & \texttt{5.5.3.1} & 6 & 1\\
 &  & 3 & \texttt{0.1.7} & \texttt{5.5.1} & 7 & 2\\
 &  & 4 & \texttt{0.1.8} & \texttt{5.5.1} & 8 & 2\\
T9.118 & \texttt{012344321} & 0 & \texttt{0.1.2.3.4} & \texttt{5.4.3.2.1} & 4 & 2\\*
 &  & 1 & \texttt{0.1.2.3.5} & \texttt{5.4.3.2.1} & 5 & 2\\
 &  & 2 & \texttt{0.1.2.6} & \texttt{5.4.3.1} & 6 & 1\\
 &  & 3 & \texttt{0.1.7} & \texttt{5.4.1} & 7 & 1\\
 &  & 4 & \texttt{0.8} & \texttt{5.1} & 8 & 1\\
T9.119 & \texttt{012344312} & 0 & \texttt{0.1.2.3.4} & \texttt{4.3.3.2.1} & 4 & 2\\*
 &  & 1 & \texttt{0.1.2.3.5} & \texttt{4.3.3.2.1} & 5 & 2\\
 &  & 2 & \texttt{0.1.2.6} & \texttt{4.3.3.1} & 6 & 1\\
 &  & 3 & \texttt{0.7.8} & \texttt{4.1.1} & 7 & 1\\
T9.120 & \texttt{012344311} & 0 & \texttt{0.1.2.3.4} & \texttt{5.3.3.2.1} & 4 & 2\\*
 &  & 1 & \texttt{0.1.2.3.5} & \texttt{5.3.3.2.1} & 5 & 2\\
 &  & 2 & \texttt{0.1.2.6} & \texttt{5.3.3.1} & 6 & 1\\
 &  & 3 & \texttt{0.7} & \texttt{5.1} & 7 & 2\\
 &  & 4 & \texttt{0.8} & \texttt{5.1} & 8 & 2\\
T9.121 & \texttt{012344234} & 0 & \texttt{0.1.2.3.4} & \texttt{3.3.2.2.1} & 4 & 2\\*
 &  & 1 & \texttt{0.1.2.3.5} & \texttt{3.3.2.2.1} & 5 & 2\\
 &  & 2 & \texttt{0.1.6.7.8} & \texttt{3.3.1.1.1} & 6 & 1\\
T9.122 & \texttt{012344233} & 0 & \texttt{0.1.2.3.4} & \texttt{4.4.2.2.1} & 4 & 2\\*
 &  & 1 & \texttt{0.1.2.3.5} & \texttt{4.4.2.2.1} & 5 & 2\\
 &  & 2 & \texttt{0.1.6.7} & \texttt{4.4.2.1} & 7 & 2\\
 &  & 3 & \texttt{0.1.6.8} & \texttt{4.4.2.1} & 8 & 2\\
T9.123 & \texttt{012344232} & 0 & \texttt{0.1.2.3.4} & \texttt{4.4.2.2.1} & 4 & 2\\*
 &  & 1 & \texttt{0.1.2.3.5} & \texttt{4.4.2.2.1} & 5 & 2\\
 &  & 2 & \texttt{0.1.6.7} & \texttt{4.4.1.1} & 6 & 1\\
 &  & 3 & \texttt{0.1.8} & \texttt{4.4.1} & 8 & 1\\
T9.124 & \texttt{012344231} & 0 & \texttt{0.1.2.3.4} & \texttt{4.3.2.2.1} & 4 & 2\\*
 &  & 1 & \texttt{0.1.2.3.5} & \texttt{4.3.2.2.1} & 5 & 2\\
 &  & 2 & \texttt{0.1.6.7} & \texttt{4.3.1.1} & 6 & 1\\
 &  & 3 & \texttt{0.8} & \texttt{4.1} & 8 & 1\\
T9.125 & \texttt{012344222} & 0 & \texttt{0.1.2.3.4} & \texttt{5.5.2.2.1} & 4 & 2\\*
 &  & 1 & \texttt{0.1.2.3.5} & \texttt{5.5.2.2.1} & 5 & 2\\
 &  & 2 & \texttt{0.1.6} & \texttt{5.5.1} & 6 & 3\\
 &  & 3 & \texttt{0.1.7} & \texttt{5.5.1} & 7 & 3\\
 &  & 4 & \texttt{0.1.8} & \texttt{5.5.1} & 8 & 3\\
T9.126 & \texttt{012344221} & 0 & \texttt{0.1.2.3.4} & \texttt{5.4.2.2.1} & 4 & 2\\*
 &  & 1 & \texttt{0.1.2.3.5} & \texttt{5.4.2.2.1} & 5 & 2\\
 &  & 2 & \texttt{0.1.6} & \texttt{5.4.1} & 6 & 2\\
 &  & 3 & \texttt{0.1.7} & \texttt{5.4.1} & 7 & 2\\
 &  & 4 & \texttt{0.8} & \texttt{5.1} & 8 & 1\\
T9.127 & \texttt{012344212} & 0 & \texttt{0.1.2.3.4} & \texttt{4.3.2.2.1} & 4 & 2\\*
 &  & 1 & \texttt{0.1.2.3.5} & \texttt{4.3.2.2.1} & 5 & 2\\
 &  & 2 & \texttt{0.1.6} & \texttt{4.3.1} & 6 & 1\\
 &  & 3 & \texttt{0.7.8} & \texttt{4.1.1} & 7 & 1\\
T9.128 & \texttt{012344211} & 0 & \texttt{0.1.2.3.4} & \texttt{5.3.2.2.1} & 4 & 2\\*
 &  & 1 & \texttt{0.1.2.3.5} & \texttt{5.3.2.2.1} & 5 & 2\\
 &  & 2 & \texttt{0.1.6} & \texttt{5.3.1} & 6 & 1\\
 &  & 3 & \texttt{0.7} & \texttt{5.1} & 7 & 2\\
 &  & 4 & \texttt{0.8} & \texttt{5.1} & 8 & 2\\
T9.129 & \texttt{012344123} & 0 & \texttt{0.1.2.3.4} & \texttt{3.2.2.2.1} & 4 & 2\\*
 &  & 1 & \texttt{0.1.2.3.5} & \texttt{3.2.2.2.1} & 5 & 2\\
 &  & 2 & \texttt{0.6.7.8} & \texttt{3.1.1.1} & 6 & 1\\
T9.130 & \texttt{012344122} & 0 & \texttt{0.1.2.3.4} & \texttt{4.2.2.2.1} & 4 & 2\\*
 &  & 1 & \texttt{0.1.2.3.5} & \texttt{4.2.2.2.1} & 5 & 2\\
 &  & 2 & \texttt{0.6.7} & \texttt{4.2.1} & 7 & 2\\
 &  & 3 & \texttt{0.6.8} & \texttt{4.2.1} & 8 & 2\\
T9.131 & \texttt{012344121} & 0 & \texttt{0.1.2.3.4} & \texttt{4.2.2.2.1} & 4 & 2\\*
 &  & 1 & \texttt{0.1.2.3.5} & \texttt{4.2.2.2.1} & 5 & 2\\
 &  & 2 & \texttt{0.6.7} & \texttt{4.1.1} & 6 & 1\\
 &  & 3 & \texttt{0.8} & \texttt{4.1} & 8 & 1\\
T9.132 & \texttt{012344111} & 0 & \texttt{0.1.2.3.4} & \texttt{5.2.2.2.1} & 4 & 2\\*
 &  & 1 & \texttt{0.1.2.3.5} & \texttt{5.2.2.2.1} & 5 & 2\\
 &  & 2 & \texttt{0.6} & \texttt{5.1} & 6 & 3\\
 &  & 3 & \texttt{0.7} & \texttt{5.1} & 7 & 3\\
 &  & 4 & \texttt{0.8} & \texttt{5.1} & 8 & 3\\
T9.133 & \texttt{012343434} & 0 & \texttt{0.1.2.3.4} & \texttt{3.3.3.1.1} & 3 & 3\\*
 &  & 1 & \texttt{0.1.2.5.6} & \texttt{3.3.3.1.1} & 5 & 3\\
 &  & 2 & \texttt{0.1.2.7.8} & \texttt{3.3.3.1.1} & 7 & 3\\
T9.134 & \texttt{012343433} & 0 & \texttt{0.1.2.3.4} & \texttt{4.4.4.1.1} & 3 & 2\\*
 &  & 1 & \texttt{0.1.2.5.6} & \texttt{4.4.4.1.1} & 5 & 2\\
 &  & 2 & \texttt{0.1.2.7} & \texttt{4.4.4.1} & 7 & 2\\
 &  & 3 & \texttt{0.1.2.8} & \texttt{4.4.4.1} & 8 & 2\\
T9.135 & \texttt{012343432} & 0 & \texttt{0.1.2.3.4} & \texttt{4.4.3.1.1} & 3 & 2\\*
 &  & 1 & \texttt{0.1.2.5.6} & \texttt{4.4.3.1.1} & 5 & 2\\
 &  & 2 & \texttt{0.1.2.7} & \texttt{4.4.3.1} & 7 & 1\\
 &  & 3 & \texttt{0.1.8} & \texttt{4.4.1} & 8 & 1\\
T9.136 & \texttt{012343431} & 0 & \texttt{0.1.2.3.4} & \texttt{4.3.3.1.1} & 3 & 2\\*
 &  & 1 & \texttt{0.1.2.5.6} & \texttt{4.3.3.1.1} & 5 & 2\\
 &  & 2 & \texttt{0.1.2.7} & \texttt{4.3.3.1} & 7 & 1\\
 &  & 3 & \texttt{0.8} & \texttt{4.1} & 8 & 1\\
T9.137 & \texttt{012343423} & 0 & \texttt{0.1.2.3.4} & \texttt{3.3.2.1.1} & 3 & 2\\*
 &  & 1 & \texttt{0.1.2.5.6} & \texttt{3.3.2.1.1} & 5 & 2\\
 &  & 2 & \texttt{0.1.7.8} & \texttt{3.3.1.1} & 7 & 1\\
T9.138 & \texttt{012343422} & 0 & \texttt{0.1.2.3.4} & \texttt{4.4.2.1.1} & 3 & 2\\*
 &  & 1 & \texttt{0.1.2.5.6} & \texttt{4.4.2.1.1} & 5 & 2\\
 &  & 2 & \texttt{0.1.7} & \texttt{4.4.1} & 7 & 2\\
 &  & 3 & \texttt{0.1.8} & \texttt{4.4.1} & 8 & 2\\
T9.139 & \texttt{012343421} & 0 & \texttt{0.1.2.3.4} & \texttt{4.3.2.1.1} & 3 & 2\\*
 &  & 1 & \texttt{0.1.2.5.6} & \texttt{4.3.2.1.1} & 5 & 2\\
 &  & 2 & \texttt{0.1.7} & \texttt{4.3.1} & 7 & 1\\
 &  & 3 & \texttt{0.8} & \texttt{4.1} & 8 & 1\\
T9.140 & \texttt{012343412} & 0 & \texttt{0.1.2.3.4} & \texttt{3.2.2.1.1} & 3 & 2\\*
 &  & 1 & \texttt{0.1.2.5.6} & \texttt{3.2.2.1.1} & 5 & 2\\
 &  & 2 & \texttt{0.7.8} & \texttt{3.1.1} & 7 & 1\\
T9.141 & \texttt{012343411} & 0 & \texttt{0.1.2.3.4} & \texttt{4.2.2.1.1} & 3 & 2\\*
 &  & 1 & \texttt{0.1.2.5.6} & \texttt{4.2.2.1.1} & 5 & 2\\
 &  & 2 & \texttt{0.7} & \texttt{4.1} & 7 & 2\\
 &  & 3 & \texttt{0.8} & \texttt{4.1} & 8 & 2\\
T9.142 & \texttt{012343333} & 0 & \texttt{0.1.2.3.4} & \texttt{5.5.5.1.1} & 3 & 1\\*
 &  & 1 & \texttt{0.1.2.5} & \texttt{5.5.5.1} & 5 & 4\\
 &  & 2 & \texttt{0.1.2.6} & \texttt{5.5.5.1} & 6 & 4\\
 &  & 3 & \texttt{0.1.2.7} & \texttt{5.5.5.1} & 7 & 4\\
 &  & 4 & \texttt{0.1.2.8} & \texttt{5.5.5.1} & 8 & 4\\
T9.143 & \texttt{012343332} & 0 & \texttt{0.1.2.3.4} & \texttt{5.5.4.1.1} & 3 & 1\\*
 &  & 1 & \texttt{0.1.2.5} & \texttt{5.5.4.1} & 5 & 3\\
 &  & 2 & \texttt{0.1.2.6} & \texttt{5.5.4.1} & 6 & 3\\
 &  & 3 & \texttt{0.1.2.7} & \texttt{5.5.4.1} & 7 & 3\\
 &  & 4 & \texttt{0.1.8} & \texttt{5.5.1} & 8 & 1\\
T9.144 & \texttt{012343331} & 0 & \texttt{0.1.2.3.4} & \texttt{5.4.4.1.1} & 3 & 1\\*
 &  & 1 & \texttt{0.1.2.5} & \texttt{5.4.4.1} & 5 & 3\\
 &  & 2 & \texttt{0.1.2.6} & \texttt{5.4.4.1} & 6 & 3\\
 &  & 3 & \texttt{0.1.2.7} & \texttt{5.4.4.1} & 7 & 3\\
 &  & 4 & \texttt{0.8} & \texttt{5.1} & 8 & 1\\
T9.145 & \texttt{012343323} & 0 & \texttt{0.1.2.3.4} & \texttt{4.4.3.1.1} & 3 & 1\\*
 &  & 1 & \texttt{0.1.2.5} & \texttt{4.4.3.1} & 5 & 2\\
 &  & 2 & \texttt{0.1.2.6} & \texttt{4.4.3.1} & 6 & 2\\
 &  & 3 & \texttt{0.1.7.8} & \texttt{4.4.1.1} & 7 & 1\\
T9.146 & \texttt{012343322} & 0 & \texttt{0.1.2.3.4} & \texttt{5.5.3.1.1} & 3 & 1\\*
 &  & 1 & \texttt{0.1.2.5} & \texttt{5.5.3.1} & 5 & 2\\
 &  & 2 & \texttt{0.1.2.6} & \texttt{5.5.3.1} & 6 & 2\\
 &  & 3 & \texttt{0.1.7} & \texttt{5.5.1} & 7 & 2\\
 &  & 4 & \texttt{0.1.8} & \texttt{5.5.1} & 8 & 2\\
T9.147 & \texttt{012343321} & 0 & \texttt{0.1.2.3.4} & \texttt{5.4.3.1.1} & 3 & 1\\*
 &  & 1 & \texttt{0.1.2.5} & \texttt{5.4.3.1} & 5 & 2\\
 &  & 2 & \texttt{0.1.2.6} & \texttt{5.4.3.1} & 6 & 2\\
 &  & 3 & \texttt{0.1.7} & \texttt{5.4.1} & 7 & 1\\
 &  & 4 & \texttt{0.8} & \texttt{5.1} & 8 & 1\\
T9.148 & \texttt{012343312} & 0 & \texttt{0.1.2.3.4} & \texttt{4.3.3.1.1} & 3 & 1\\*
 &  & 1 & \texttt{0.1.2.5} & \texttt{4.3.3.1} & 5 & 2\\
 &  & 2 & \texttt{0.1.2.6} & \texttt{4.3.3.1} & 6 & 2\\
 &  & 3 & \texttt{0.7.8} & \texttt{4.1.1} & 7 & 1\\
T9.149 & \texttt{012343311} & 0 & \texttt{0.1.2.3.4} & \texttt{5.3.3.1.1} & 3 & 1\\*
 &  & 1 & \texttt{0.1.2.5} & \texttt{5.3.3.1} & 5 & 2\\
 &  & 2 & \texttt{0.1.2.6} & \texttt{5.3.3.1} & 6 & 2\\
 &  & 3 & \texttt{0.7} & \texttt{5.1} & 7 & 2\\
 &  & 4 & \texttt{0.8} & \texttt{5.1} & 8 & 2\\
T9.150 & \texttt{012343234} & 0 & \texttt{0.1.2.3.4} & \texttt{3.3.2.1.1} & 3 & 1\\*
 &  & 1 & \texttt{0.1.2.5} & \texttt{3.3.2.1} & 5 & 1\\
 &  & 2 & \texttt{0.1.6.7.8} & \texttt{3.3.1.1.1} & 6 & 1\\
T9.151 & \texttt{012343233} & 0 & \texttt{0.1.2.3.4} & \texttt{4.4.2.1.1} & 3 & 1\\*
 &  & 1 & \texttt{0.1.2.5} & \texttt{4.4.2.1} & 5 & 1\\
 &  & 2 & \texttt{0.1.6.7} & \texttt{4.4.2.1} & 7 & 2\\
 &  & 3 & \texttt{0.1.6.8} & \texttt{4.4.2.1} & 8 & 2\\
T9.152 & \texttt{012343232} & 0 & \texttt{0.1.2.3.4} & \texttt{4.4.2.1.1} & 3 & 1\\*
 &  & 1 & \texttt{0.1.2.5} & \texttt{4.4.2.1} & 5 & 1\\
 &  & 2 & \texttt{0.1.6.7} & \texttt{4.4.1.1} & 6 & 1\\
 &  & 3 & \texttt{0.1.8} & \texttt{4.4.1} & 8 & 1\\
T9.153 & \texttt{012343231} & 0 & \texttt{0.1.2.3.4} & \texttt{4.3.2.1.1} & 3 & 1\\*
 &  & 1 & \texttt{0.1.2.5} & \texttt{4.3.2.1} & 5 & 1\\
 &  & 2 & \texttt{0.1.6.7} & \texttt{4.3.1.1} & 6 & 1\\
 &  & 3 & \texttt{0.8} & \texttt{4.1} & 8 & 1\\
T9.154 & \texttt{012343222} & 0 & \texttt{0.1.2.3.4} & \texttt{5.5.2.1.1} & 3 & 1\\*
 &  & 1 & \texttt{0.1.2.5} & \texttt{5.5.2.1} & 5 & 1\\
 &  & 2 & \texttt{0.1.6} & \texttt{5.5.1} & 6 & 3\\
 &  & 3 & \texttt{0.1.7} & \texttt{5.5.1} & 7 & 3\\
 &  & 4 & \texttt{0.1.8} & \texttt{5.5.1} & 8 & 3\\
T9.155 & \texttt{012343221} & 0 & \texttt{0.1.2.3.4} & \texttt{5.4.2.1.1} & 3 & 1\\*
 &  & 1 & \texttt{0.1.2.5} & \texttt{5.4.2.1} & 5 & 1\\
 &  & 2 & \texttt{0.1.6} & \texttt{5.4.1} & 6 & 2\\
 &  & 3 & \texttt{0.1.7} & \texttt{5.4.1} & 7 & 2\\
 &  & 4 & \texttt{0.8} & \texttt{5.1} & 8 & 1\\
T9.156 & \texttt{012343212} & 0 & \texttt{0.1.2.3.4} & \texttt{4.3.2.1.1} & 3 & 1\\*
 &  & 1 & \texttt{0.1.2.5} & \texttt{4.3.2.1} & 5 & 1\\
 &  & 2 & \texttt{0.1.6} & \texttt{4.3.1} & 6 & 1\\
 &  & 3 & \texttt{0.7.8} & \texttt{4.1.1} & 7 & 1\\
T9.157 & \texttt{012343211} & 0 & \texttt{0.1.2.3.4} & \texttt{5.3.2.1.1} & 3 & 1\\*
 &  & 1 & \texttt{0.1.2.5} & \texttt{5.3.2.1} & 5 & 1\\
 &  & 2 & \texttt{0.1.6} & \texttt{5.3.1} & 6 & 1\\
 &  & 3 & \texttt{0.7} & \texttt{5.1} & 7 & 2\\
 &  & 4 & \texttt{0.8} & \texttt{5.1} & 8 & 2\\
T9.158 & \texttt{012343123} & 0 & \texttt{0.1.2.3.4} & \texttt{3.2.2.1.1} & 3 & 1\\*
 &  & 1 & \texttt{0.1.2.5} & \texttt{3.2.2.1} & 5 & 1\\
 &  & 2 & \texttt{0.6.7.8} & \texttt{3.1.1.1} & 6 & 1\\
T9.159 & \texttt{012343122} & 0 & \texttt{0.1.2.3.4} & \texttt{4.2.2.1.1} & 3 & 1\\*
 &  & 1 & \texttt{0.1.2.5} & \texttt{4.2.2.1} & 5 & 1\\
 &  & 2 & \texttt{0.6.7} & \texttt{4.2.1} & 7 & 2\\
 &  & 3 & \texttt{0.6.8} & \texttt{4.2.1} & 8 & 2\\
T9.160 & \texttt{012343121} & 0 & \texttt{0.1.2.3.4} & \texttt{4.2.2.1.1} & 3 & 1\\*
 &  & 1 & \texttt{0.1.2.5} & \texttt{4.2.2.1} & 5 & 1\\
 &  & 2 & \texttt{0.6.7} & \texttt{4.1.1} & 6 & 1\\
 &  & 3 & \texttt{0.8} & \texttt{4.1} & 8 & 1\\
T9.161 & \texttt{012343111} & 0 & \texttt{0.1.2.3.4} & \texttt{5.2.2.1.1} & 3 & 1\\*
 &  & 1 & \texttt{0.1.2.5} & \texttt{5.2.2.1} & 5 & 1\\
 &  & 2 & \texttt{0.6} & \texttt{5.1} & 6 & 3\\
 &  & 3 & \texttt{0.7} & \texttt{5.1} & 7 & 3\\
 &  & 4 & \texttt{0.8} & \texttt{5.1} & 8 & 3\\
T9.162 & \texttt{012342342} & 0 & \texttt{0.1.2.3.4} & \texttt{3.3.1.1.1} & 2 & 2\\*
 &  & 1 & \texttt{0.1.5.6.7} & \texttt{3.3.1.1.1} & 5 & 2\\
 &  & 2 & \texttt{0.1.8} & \texttt{3.3.1} & 8 & 1\\
T9.163 & \texttt{012342341} & 0 & \texttt{0.1.2.3.4} & \texttt{3.2.1.1.1} & 2 & 2\\*
 &  & 1 & \texttt{0.1.5.6.7} & \texttt{3.2.1.1.1} & 5 & 2\\
 &  & 2 & \texttt{0.8} & \texttt{3.1} & 8 & 1\\
T9.164 & \texttt{012342333} & 0 & \texttt{0.1.2.3.4} & \texttt{4.4.1.1.1} & 2 & 1\\*
 &  & 1 & \texttt{0.1.5.6} & \texttt{4.4.3.1} & 6 & 3\\
 &  & 2 & \texttt{0.1.5.7} & \texttt{4.4.3.1} & 7 & 3\\
 &  & 3 & \texttt{0.1.5.8} & \texttt{4.4.3.1} & 8 & 3\\
T9.165 & \texttt{012342332} & 0 & \texttt{0.1.2.3.4} & \texttt{4.4.1.1.1} & 2 & 1\\*
 &  & 1 & \texttt{0.1.5.6} & \texttt{4.4.2.1} & 6 & 2\\
 &  & 2 & \texttt{0.1.5.7} & \texttt{4.4.2.1} & 7 & 2\\
 &  & 3 & \texttt{0.1.8} & \texttt{4.4.1} & 8 & 1\\
T9.166 & \texttt{012342331} & 0 & \texttt{0.1.2.3.4} & \texttt{4.3.1.1.1} & 2 & 1\\*
 &  & 1 & \texttt{0.1.5.6} & \texttt{4.3.2.1} & 6 & 2\\
 &  & 2 & \texttt{0.1.5.7} & \texttt{4.3.2.1} & 7 & 2\\
 &  & 3 & \texttt{0.8} & \texttt{4.1} & 8 & 1\\
T9.167 & \texttt{012342323} & 0 & \texttt{0.1.2.3.4} & \texttt{3.3.1.1.1} & 2 & 1\\*
 &  & 1 & \texttt{0.1.5.6} & \texttt{3.3.1.1} & 5 & 2\\
 &  & 2 & \texttt{0.1.7.8} & \texttt{3.3.1.1} & 7 & 2\\
T9.168 & \texttt{012342322} & 0 & \texttt{0.1.2.3.4} & \texttt{4.4.1.1.1} & 2 & 1\\*
 &  & 1 & \texttt{0.1.5.6} & \texttt{4.4.1.1} & 5 & 1\\
 &  & 2 & \texttt{0.1.7} & \texttt{4.4.1} & 7 & 2\\
 &  & 3 & \texttt{0.1.8} & \texttt{4.4.1} & 8 & 2\\
T9.169 & \texttt{012342321} & 0 & \texttt{0.1.2.3.4} & \texttt{4.3.1.1.1} & 2 & 1\\*
 &  & 1 & \texttt{0.1.5.6} & \texttt{4.3.1.1} & 5 & 1\\
 &  & 2 & \texttt{0.1.7} & \texttt{4.3.1} & 7 & 1\\
 &  & 3 & \texttt{0.8} & \texttt{4.1} & 8 & 1\\
T9.170 & \texttt{012342312} & 0 & \texttt{0.1.2.3.4} & \texttt{3.2.1.1.1} & 2 & 1\\*
 &  & 1 & \texttt{0.1.5.6} & \texttt{3.2.1.1} & 5 & 1\\
 &  & 2 & \texttt{0.7.8} & \texttt{3.1.1} & 7 & 1\\
T9.171 & \texttt{012342311} & 0 & \texttt{0.1.2.3.4} & \texttt{4.2.1.1.1} & 2 & 1\\*
 &  & 1 & \texttt{0.1.5.6} & \texttt{4.2.1.1} & 5 & 1\\
 &  & 2 & \texttt{0.7} & \texttt{4.1} & 7 & 2\\
 &  & 3 & \texttt{0.8} & \texttt{4.1} & 8 & 2\\
T9.172 & \texttt{012342222} & 0 & \texttt{0.1.2.3.4} & \texttt{5.5.1.1.1} & 2 & 1\\*
 &  & 1 & \texttt{0.1.5} & \texttt{5.5.1} & 5 & 4\\
 &  & 2 & \texttt{0.1.6} & \texttt{5.5.1} & 6 & 4\\
 &  & 3 & \texttt{0.1.7} & \texttt{5.5.1} & 7 & 4\\
 &  & 4 & \texttt{0.1.8} & \texttt{5.5.1} & 8 & 4\\
T9.173 & \texttt{012342221} & 0 & \texttt{0.1.2.3.4} & \texttt{5.4.1.1.1} & 2 & 1\\*
 &  & 1 & \texttt{0.1.5} & \texttt{5.4.1} & 5 & 3\\
 &  & 2 & \texttt{0.1.6} & \texttt{5.4.1} & 6 & 3\\
 &  & 3 & \texttt{0.1.7} & \texttt{5.4.1} & 7 & 3\\
 &  & 4 & \texttt{0.8} & \texttt{5.1} & 8 & 1\\
T9.174 & \texttt{012342212} & 0 & \texttt{0.1.2.3.4} & \texttt{4.3.1.1.1} & 2 & 1\\*
 &  & 1 & \texttt{0.1.5} & \texttt{4.3.1} & 5 & 2\\
 &  & 2 & \texttt{0.1.6} & \texttt{4.3.1} & 6 & 2\\
 &  & 3 & \texttt{0.7.8} & \texttt{4.1.1} & 7 & 1\\
T9.175 & \texttt{012342211} & 0 & \texttt{0.1.2.3.4} & \texttt{5.3.1.1.1} & 2 & 1\\*
 &  & 1 & \texttt{0.1.5} & \texttt{5.3.1} & 5 & 2\\
 &  & 2 & \texttt{0.1.6} & \texttt{5.3.1} & 6 & 2\\
 &  & 3 & \texttt{0.7} & \texttt{5.1} & 7 & 2\\
 &  & 4 & \texttt{0.8} & \texttt{5.1} & 8 & 2\\
T9.176 & \texttt{012342123} & 0 & \texttt{0.1.2.3.4} & \texttt{3.2.1.1.1} & 2 & 1\\*
 &  & 1 & \texttt{0.1.5} & \texttt{3.2.1} & 5 & 1\\
 &  & 2 & \texttt{0.6.7.8} & \texttt{3.1.1.1} & 6 & 1\\
T9.177 & \texttt{012342122} & 0 & \texttt{0.1.2.3.4} & \texttt{4.2.1.1.1} & 2 & 1\\*
 &  & 1 & \texttt{0.1.5} & \texttt{4.2.1} & 5 & 1\\
 &  & 2 & \texttt{0.6.7} & \texttt{4.2.1} & 7 & 2\\
 &  & 3 & \texttt{0.6.8} & \texttt{4.2.1} & 8 & 2\\
T9.178 & \texttt{012342121} & 0 & \texttt{0.1.2.3.4} & \texttt{4.2.1.1.1} & 2 & 1\\*
 &  & 1 & \texttt{0.1.5} & \texttt{4.2.1} & 5 & 1\\
 &  & 2 & \texttt{0.6.7} & \texttt{4.1.1} & 6 & 1\\
 &  & 3 & \texttt{0.8} & \texttt{4.1} & 8 & 1\\
T9.179 & \texttt{012342111} & 0 & \texttt{0.1.2.3.4} & \texttt{5.2.1.1.1} & 2 & 1\\*
 &  & 1 & \texttt{0.1.5} & \texttt{5.2.1} & 5 & 1\\
 &  & 2 & \texttt{0.6} & \texttt{5.1} & 6 & 3\\
 &  & 3 & \texttt{0.7} & \texttt{5.1} & 7 & 3\\
 &  & 4 & \texttt{0.8} & \texttt{5.1} & 8 & 3\\
T9.180 & \texttt{012341234} & 0 & \texttt{0.1.2.3.4} & \texttt{2.1.1.1.1} & 1 & 2\\*
 &  & 1 & \texttt{0.5.6.7.8} & \texttt{2.1.1.1.1} & 5 & 2\\
T9.181 & \texttt{012341233} & 0 & \texttt{0.1.2.3.4} & \texttt{3.1.1.1.1} & 1 & 1\\*
 &  & 1 & \texttt{0.5.6.7} & \texttt{3.2.2.1} & 7 & 2\\
 &  & 2 & \texttt{0.5.6.8} & \texttt{3.2.2.1} & 8 & 2\\
T9.182 & \texttt{012341232} & 0 & \texttt{0.1.2.3.4} & \texttt{3.1.1.1.1} & 1 & 1\\*
 &  & 1 & \texttt{0.5.6.7} & \texttt{3.2.1.1} & 6 & 1\\
 &  & 2 & \texttt{0.5.8} & \texttt{3.2.1} & 8 & 1\\
T9.183 & \texttt{012341231} & 0 & \texttt{0.1.2.3.4} & \texttt{3.1.1.1.1} & 1 & 1\\*
 &  & 1 & \texttt{0.5.6.7} & \texttt{3.1.1.1} & 5 & 1\\
 &  & 2 & \texttt{0.8} & \texttt{3.1} & 8 & 1\\
T9.184 & \texttt{012341222} & 0 & \texttt{0.1.2.3.4} & \texttt{4.1.1.1.1} & 1 & 1\\*
 &  & 1 & \texttt{0.5.6} & \texttt{4.3.1} & 6 & 3\\
 &  & 2 & \texttt{0.5.7} & \texttt{4.3.1} & 7 & 3\\
 &  & 3 & \texttt{0.5.8} & \texttt{4.3.1} & 8 & 3\\
T9.185 & \texttt{012341221} & 0 & \texttt{0.1.2.3.4} & \texttt{4.1.1.1.1} & 1 & 1\\*
 &  & 1 & \texttt{0.5.6} & \texttt{4.2.1} & 6 & 2\\
 &  & 2 & \texttt{0.5.7} & \texttt{4.2.1} & 7 & 2\\
 &  & 3 & \texttt{0.8} & \texttt{4.1} & 8 & 1\\
T9.186 & \texttt{012341212} & 0 & \texttt{0.1.2.3.4} & \texttt{3.1.1.1.1} & 1 & 1\\*
 &  & 1 & \texttt{0.5.6} & \texttt{3.1.1} & 5 & 2\\
 &  & 2 & \texttt{0.7.8} & \texttt{3.1.1} & 7 & 2\\
T9.187 & \texttt{012341211} & 0 & \texttt{0.1.2.3.4} & \texttt{4.1.1.1.1} & 1 & 1\\*
 &  & 1 & \texttt{0.5.6} & \texttt{4.1.1} & 5 & 1\\
 &  & 2 & \texttt{0.7} & \texttt{4.1} & 7 & 2\\
 &  & 3 & \texttt{0.8} & \texttt{4.1} & 8 & 2\\
T9.188 & \texttt{012341111} & 0 & \texttt{0.1.2.3.4} & \texttt{5.1.1.1.1} & 1 & 1\\*
 &  & 1 & \texttt{0.5} & \texttt{5.1} & 5 & 4\\
 &  & 2 & \texttt{0.6} & \texttt{5.1} & 6 & 4\\
 &  & 3 & \texttt{0.7} & \texttt{5.1} & 7 & 4\\
 &  & 4 & \texttt{0.8} & \texttt{5.1} & 8 & 4\\
T9.189 & \texttt{012333333} & 0 & \texttt{0.1.2.3} & \texttt{6.6.6.1} & 3 & 6\\*
 &  & 1 & \texttt{0.1.2.4} & \texttt{6.6.6.1} & 4 & 6\\
 &  & 2 & \texttt{0.1.2.5} & \texttt{6.6.6.1} & 5 & 6\\
 &  & 3 & \texttt{0.1.2.6} & \texttt{6.6.6.1} & 6 & 6\\
 &  & 4 & \texttt{0.1.2.7} & \texttt{6.6.6.1} & 7 & 6\\
 &  & 5 & \texttt{0.1.2.8} & \texttt{6.6.6.1} & 8 & 6\\
T9.190 & \texttt{012333332} & 0 & \texttt{0.1.2.3} & \texttt{6.6.5.1} & 3 & 5\\*
 &  & 1 & \texttt{0.1.2.4} & \texttt{6.6.5.1} & 4 & 5\\
 &  & 2 & \texttt{0.1.2.5} & \texttt{6.6.5.1} & 5 & 5\\
 &  & 3 & \texttt{0.1.2.6} & \texttt{6.6.5.1} & 6 & 5\\
 &  & 4 & \texttt{0.1.2.7} & \texttt{6.6.5.1} & 7 & 5\\
 &  & 5 & \texttt{0.1.8} & \texttt{6.6.1} & 8 & 1\\
T9.191 & \texttt{012333331} & 0 & \texttt{0.1.2.3} & \texttt{6.5.5.1} & 3 & 5\\*
 &  & 1 & \texttt{0.1.2.4} & \texttt{6.5.5.1} & 4 & 5\\
 &  & 2 & \texttt{0.1.2.5} & \texttt{6.5.5.1} & 5 & 5\\
 &  & 3 & \texttt{0.1.2.6} & \texttt{6.5.5.1} & 6 & 5\\
 &  & 4 & \texttt{0.1.2.7} & \texttt{6.5.5.1} & 7 & 5\\
 &  & 5 & \texttt{0.8} & \texttt{6.1} & 8 & 1\\
T9.192 & \texttt{012333323} & 0 & \texttt{0.1.2.3} & \texttt{5.5.4.1} & 3 & 4\\*
 &  & 1 & \texttt{0.1.2.4} & \texttt{5.5.4.1} & 4 & 4\\
 &  & 2 & \texttt{0.1.2.5} & \texttt{5.5.4.1} & 5 & 4\\
 &  & 3 & \texttt{0.1.2.6} & \texttt{5.5.4.1} & 6 & 4\\
 &  & 4 & \texttt{0.1.7.8} & \texttt{5.5.1.1} & 7 & 1\\
T9.193 & \texttt{012333322} & 0 & \texttt{0.1.2.3} & \texttt{6.6.4.1} & 3 & 4\\*
 &  & 1 & \texttt{0.1.2.4} & \texttt{6.6.4.1} & 4 & 4\\
 &  & 2 & \texttt{0.1.2.5} & \texttt{6.6.4.1} & 5 & 4\\
 &  & 3 & \texttt{0.1.2.6} & \texttt{6.6.4.1} & 6 & 4\\
 &  & 4 & \texttt{0.1.7} & \texttt{6.6.1} & 7 & 2\\
 &  & 5 & \texttt{0.1.8} & \texttt{6.6.1} & 8 & 2\\
T9.194 & \texttt{012333321} & 0 & \texttt{0.1.2.3} & \texttt{6.5.4.1} & 3 & 4\\*
 &  & 1 & \texttt{0.1.2.4} & \texttt{6.5.4.1} & 4 & 4\\
 &  & 2 & \texttt{0.1.2.5} & \texttt{6.5.4.1} & 5 & 4\\
 &  & 3 & \texttt{0.1.2.6} & \texttt{6.5.4.1} & 6 & 4\\
 &  & 4 & \texttt{0.1.7} & \texttt{6.5.1} & 7 & 1\\
 &  & 5 & \texttt{0.8} & \texttt{6.1} & 8 & 1\\
T9.195 & \texttt{012333312} & 0 & \texttt{0.1.2.3} & \texttt{5.4.4.1} & 3 & 4\\*
 &  & 1 & \texttt{0.1.2.4} & \texttt{5.4.4.1} & 4 & 4\\
 &  & 2 & \texttt{0.1.2.5} & \texttt{5.4.4.1} & 5 & 4\\
 &  & 3 & \texttt{0.1.2.6} & \texttt{5.4.4.1} & 6 & 4\\
 &  & 4 & \texttt{0.7.8} & \texttt{5.1.1} & 7 & 1\\
T9.196 & \texttt{012333311} & 0 & \texttt{0.1.2.3} & \texttt{6.4.4.1} & 3 & 4\\*
 &  & 1 & \texttt{0.1.2.4} & \texttt{6.4.4.1} & 4 & 4\\
 &  & 2 & \texttt{0.1.2.5} & \texttt{6.4.4.1} & 5 & 4\\
 &  & 3 & \texttt{0.1.2.6} & \texttt{6.4.4.1} & 6 & 4\\
 &  & 4 & \texttt{0.7} & \texttt{6.1} & 7 & 2\\
 &  & 5 & \texttt{0.8} & \texttt{6.1} & 8 & 2\\
T9.197 & \texttt{012333233} & 0 & \texttt{0.1.2.3} & \texttt{5.5.3.1} & 3 & 3\\*
 &  & 1 & \texttt{0.1.2.4} & \texttt{5.5.3.1} & 4 & 3\\
 &  & 2 & \texttt{0.1.2.5} & \texttt{5.5.3.1} & 5 & 3\\
 &  & 3 & \texttt{0.1.6.7} & \texttt{5.5.2.1} & 7 & 2\\
 &  & 4 & \texttt{0.1.6.8} & \texttt{5.5.2.1} & 8 & 2\\
T9.198 & \texttt{012333232} & 0 & \texttt{0.1.2.3} & \texttt{5.5.3.1} & 3 & 3\\*
 &  & 1 & \texttt{0.1.2.4} & \texttt{5.5.3.1} & 4 & 3\\
 &  & 2 & \texttt{0.1.2.5} & \texttt{5.5.3.1} & 5 & 3\\
 &  & 3 & \texttt{0.1.6.7} & \texttt{5.5.1.1} & 6 & 1\\
 &  & 4 & \texttt{0.1.8} & \texttt{5.5.1} & 8 & 1\\
T9.199 & \texttt{012333231} & 0 & \texttt{0.1.2.3} & \texttt{5.4.3.1} & 3 & 3\\*
 &  & 1 & \texttt{0.1.2.4} & \texttt{5.4.3.1} & 4 & 3\\
 &  & 2 & \texttt{0.1.2.5} & \texttt{5.4.3.1} & 5 & 3\\
 &  & 3 & \texttt{0.1.6.7} & \texttt{5.4.1.1} & 6 & 1\\
 &  & 4 & \texttt{0.8} & \texttt{5.1} & 8 & 1\\
T9.200 & \texttt{012333222} & 0 & \texttt{0.1.2.3} & \texttt{6.6.3.1} & 3 & 3\\*
 &  & 1 & \texttt{0.1.2.4} & \texttt{6.6.3.1} & 4 & 3\\
 &  & 2 & \texttt{0.1.2.5} & \texttt{6.6.3.1} & 5 & 3\\
 &  & 3 & \texttt{0.1.6} & \texttt{6.6.1} & 6 & 3\\
 &  & 4 & \texttt{0.1.7} & \texttt{6.6.1} & 7 & 3\\
 &  & 5 & \texttt{0.1.8} & \texttt{6.6.1} & 8 & 3\\
T9.201 & \texttt{012333221} & 0 & \texttt{0.1.2.3} & \texttt{6.5.3.1} & 3 & 3\\*
 &  & 1 & \texttt{0.1.2.4} & \texttt{6.5.3.1} & 4 & 3\\
 &  & 2 & \texttt{0.1.2.5} & \texttt{6.5.3.1} & 5 & 3\\
 &  & 3 & \texttt{0.1.6} & \texttt{6.5.1} & 6 & 2\\
 &  & 4 & \texttt{0.1.7} & \texttt{6.5.1} & 7 & 2\\
 &  & 5 & \texttt{0.8} & \texttt{6.1} & 8 & 1\\
T9.202 & \texttt{012333212} & 0 & \texttt{0.1.2.3} & \texttt{5.4.3.1} & 3 & 3\\*
 &  & 1 & \texttt{0.1.2.4} & \texttt{5.4.3.1} & 4 & 3\\
 &  & 2 & \texttt{0.1.2.5} & \texttt{5.4.3.1} & 5 & 3\\
 &  & 3 & \texttt{0.1.6} & \texttt{5.4.1} & 6 & 1\\
 &  & 4 & \texttt{0.7.8} & \texttt{5.1.1} & 7 & 1\\
T9.203 & \texttt{012333211} & 0 & \texttt{0.1.2.3} & \texttt{6.4.3.1} & 3 & 3\\*
 &  & 1 & \texttt{0.1.2.4} & \texttt{6.4.3.1} & 4 & 3\\
 &  & 2 & \texttt{0.1.2.5} & \texttt{6.4.3.1} & 5 & 3\\
 &  & 3 & \texttt{0.1.6} & \texttt{6.4.1} & 6 & 1\\
 &  & 4 & \texttt{0.7} & \texttt{6.1} & 7 & 2\\
 &  & 5 & \texttt{0.8} & \texttt{6.1} & 8 & 2\\
T9.204 & \texttt{012333123} & 0 & \texttt{0.1.2.3} & \texttt{4.3.3.1} & 3 & 3\\*
 &  & 1 & \texttt{0.1.2.4} & \texttt{4.3.3.1} & 4 & 3\\
 &  & 2 & \texttt{0.1.2.5} & \texttt{4.3.3.1} & 5 & 3\\
 &  & 3 & \texttt{0.6.7.8} & \texttt{4.1.1.1} & 6 & 1\\
T9.205 & \texttt{012333122} & 0 & \texttt{0.1.2.3} & \texttt{5.3.3.1} & 3 & 3\\*
 &  & 1 & \texttt{0.1.2.4} & \texttt{5.3.3.1} & 4 & 3\\
 &  & 2 & \texttt{0.1.2.5} & \texttt{5.3.3.1} & 5 & 3\\
 &  & 3 & \texttt{0.6.7} & \texttt{5.2.1} & 7 & 2\\
 &  & 4 & \texttt{0.6.8} & \texttt{5.2.1} & 8 & 2\\
T9.206 & \texttt{012333121} & 0 & \texttt{0.1.2.3} & \texttt{5.3.3.1} & 3 & 3\\*
 &  & 1 & \texttt{0.1.2.4} & \texttt{5.3.3.1} & 4 & 3\\
 &  & 2 & \texttt{0.1.2.5} & \texttt{5.3.3.1} & 5 & 3\\
 &  & 3 & \texttt{0.6.7} & \texttt{5.1.1} & 6 & 1\\
 &  & 4 & \texttt{0.8} & \texttt{5.1} & 8 & 1\\
T9.207 & \texttt{012333111} & 0 & \texttt{0.1.2.3} & \texttt{6.3.3.1} & 3 & 3\\*
 &  & 1 & \texttt{0.1.2.4} & \texttt{6.3.3.1} & 4 & 3\\
 &  & 2 & \texttt{0.1.2.5} & \texttt{6.3.3.1} & 5 & 3\\
 &  & 3 & \texttt{0.6} & \texttt{6.1} & 6 & 3\\
 &  & 4 & \texttt{0.7} & \texttt{6.1} & 7 & 3\\
 &  & 5 & \texttt{0.8} & \texttt{6.1} & 8 & 3\\
T9.208 & \texttt{012332332} & 0 & \texttt{0.1.2.3} & \texttt{5.5.2.1} & 3 & 4\\*
 &  & 1 & \texttt{0.1.2.4} & \texttt{5.5.2.1} & 4 & 4\\
 &  & 2 & \texttt{0.1.5.6} & \texttt{5.5.2.1} & 6 & 4\\
 &  & 3 & \texttt{0.1.5.7} & \texttt{5.5.2.1} & 7 & 4\\
 &  & 4 & \texttt{0.1.8} & \texttt{5.5.1} & 8 & 1\\
T9.209 & \texttt{012332331} & 0 & \texttt{0.1.2.3} & \texttt{5.4.2.1} & 3 & 4\\*
 &  & 1 & \texttt{0.1.2.4} & \texttt{5.4.2.1} & 4 & 4\\
 &  & 2 & \texttt{0.1.5.6} & \texttt{5.4.2.1} & 6 & 4\\
 &  & 3 & \texttt{0.1.5.7} & \texttt{5.4.2.1} & 7 & 4\\
 &  & 4 & \texttt{0.8} & \texttt{5.1} & 8 & 1\\
T9.210 & \texttt{012332323} & 0 & \texttt{0.1.2.3} & \texttt{4.4.2.1} & 3 & 2\\*
 &  & 1 & \texttt{0.1.2.4} & \texttt{4.4.2.1} & 4 & 2\\
 &  & 2 & \texttt{0.1.5.6} & \texttt{4.4.1.1} & 5 & 2\\
 &  & 3 & \texttt{0.1.7.8} & \texttt{4.4.1.1} & 7 & 2\\
T9.211 & \texttt{012332322} & 0 & \texttt{0.1.2.3} & \texttt{5.5.2.1} & 3 & 2\\*
 &  & 1 & \texttt{0.1.2.4} & \texttt{5.5.2.1} & 4 & 2\\
 &  & 2 & \texttt{0.1.5.6} & \texttt{5.5.1.1} & 5 & 1\\
 &  & 3 & \texttt{0.1.7} & \texttt{5.5.1} & 7 & 2\\
 &  & 4 & \texttt{0.1.8} & \texttt{5.5.1} & 8 & 2\\
T9.212 & \texttt{012332321} & 0 & \texttt{0.1.2.3} & \texttt{5.4.2.1} & 3 & 2\\*
 &  & 1 & \texttt{0.1.2.4} & \texttt{5.4.2.1} & 4 & 2\\
 &  & 2 & \texttt{0.1.5.6} & \texttt{5.4.1.1} & 5 & 1\\
 &  & 3 & \texttt{0.1.7} & \texttt{5.4.1} & 7 & 1\\
 &  & 4 & \texttt{0.8} & \texttt{5.1} & 8 & 1\\
T9.213 & \texttt{012332312} & 0 & \texttt{0.1.2.3} & \texttt{4.3.2.1} & 3 & 2\\*
 &  & 1 & \texttt{0.1.2.4} & \texttt{4.3.2.1} & 4 & 2\\
 &  & 2 & \texttt{0.1.5.6} & \texttt{4.3.1.1} & 5 & 1\\
 &  & 3 & \texttt{0.7.8} & \texttt{4.1.1} & 7 & 1\\
T9.214 & \texttt{012332311} & 0 & \texttt{0.1.2.3} & \texttt{5.3.2.1} & 3 & 2\\*
 &  & 1 & \texttt{0.1.2.4} & \texttt{5.3.2.1} & 4 & 2\\
 &  & 2 & \texttt{0.1.5.6} & \texttt{5.3.1.1} & 5 & 1\\
 &  & 3 & \texttt{0.7} & \texttt{5.1} & 7 & 2\\
 &  & 4 & \texttt{0.8} & \texttt{5.1} & 8 & 2\\
T9.215 & \texttt{012332222} & 0 & \texttt{0.1.2.3} & \texttt{6.6.2.1} & 3 & 2\\*
 &  & 1 & \texttt{0.1.2.4} & \texttt{6.6.2.1} & 4 & 2\\
 &  & 2 & \texttt{0.1.5} & \texttt{6.6.1} & 5 & 4\\
 &  & 3 & \texttt{0.1.6} & \texttt{6.6.1} & 6 & 4\\
 &  & 4 & \texttt{0.1.7} & \texttt{6.6.1} & 7 & 4\\
 &  & 5 & \texttt{0.1.8} & \texttt{6.6.1} & 8 & 4\\
T9.216 & \texttt{012332221} & 0 & \texttt{0.1.2.3} & \texttt{6.5.2.1} & 3 & 2\\*
 &  & 1 & \texttt{0.1.2.4} & \texttt{6.5.2.1} & 4 & 2\\
 &  & 2 & \texttt{0.1.5} & \texttt{6.5.1} & 5 & 3\\
 &  & 3 & \texttt{0.1.6} & \texttt{6.5.1} & 6 & 3\\
 &  & 4 & \texttt{0.1.7} & \texttt{6.5.1} & 7 & 3\\
 &  & 5 & \texttt{0.8} & \texttt{6.1} & 8 & 1\\
T9.217 & \texttt{012332212} & 0 & \texttt{0.1.2.3} & \texttt{5.4.2.1} & 3 & 2\\*
 &  & 1 & \texttt{0.1.2.4} & \texttt{5.4.2.1} & 4 & 2\\
 &  & 2 & \texttt{0.1.5} & \texttt{5.4.1} & 5 & 2\\
 &  & 3 & \texttt{0.1.6} & \texttt{5.4.1} & 6 & 2\\
 &  & 4 & \texttt{0.7.8} & \texttt{5.1.1} & 7 & 1\\
T9.218 & \texttt{012332211} & 0 & \texttt{0.1.2.3} & \texttt{6.4.2.1} & 3 & 2\\*
 &  & 1 & \texttt{0.1.2.4} & \texttt{6.4.2.1} & 4 & 2\\
 &  & 2 & \texttt{0.1.5} & \texttt{6.4.1} & 5 & 2\\
 &  & 3 & \texttt{0.1.6} & \texttt{6.4.1} & 6 & 2\\
 &  & 4 & \texttt{0.7} & \texttt{6.1} & 7 & 2\\
 &  & 5 & \texttt{0.8} & \texttt{6.1} & 8 & 2\\
T9.219 & \texttt{012332123} & 0 & \texttt{0.1.2.3} & \texttt{4.3.2.1} & 3 & 2\\*
 &  & 1 & \texttt{0.1.2.4} & \texttt{4.3.2.1} & 4 & 2\\
 &  & 2 & \texttt{0.1.5} & \texttt{4.3.1} & 5 & 1\\
 &  & 3 & \texttt{0.6.7.8} & \texttt{4.1.1.1} & 6 & 1\\
T9.220 & \texttt{012332122} & 0 & \texttt{0.1.2.3} & \texttt{5.3.2.1} & 3 & 2\\*
 &  & 1 & \texttt{0.1.2.4} & \texttt{5.3.2.1} & 4 & 2\\
 &  & 2 & \texttt{0.1.5} & \texttt{5.3.1} & 5 & 1\\
 &  & 3 & \texttt{0.6.7} & \texttt{5.2.1} & 7 & 2\\
 &  & 4 & \texttt{0.6.8} & \texttt{5.2.1} & 8 & 2\\
T9.221 & \texttt{012332121} & 0 & \texttt{0.1.2.3} & \texttt{5.3.2.1} & 3 & 2\\*
 &  & 1 & \texttt{0.1.2.4} & \texttt{5.3.2.1} & 4 & 2\\
 &  & 2 & \texttt{0.1.5} & \texttt{5.3.1} & 5 & 1\\
 &  & 3 & \texttt{0.6.7} & \texttt{5.1.1} & 6 & 1\\
 &  & 4 & \texttt{0.8} & \texttt{5.1} & 8 & 1\\
T9.222 & \texttt{012332111} & 0 & \texttt{0.1.2.3} & \texttt{6.3.2.1} & 3 & 2\\*
 &  & 1 & \texttt{0.1.2.4} & \texttt{6.3.2.1} & 4 & 2\\
 &  & 2 & \texttt{0.1.5} & \texttt{6.3.1} & 5 & 1\\
 &  & 3 & \texttt{0.6} & \texttt{6.1} & 6 & 3\\
 &  & 4 & \texttt{0.7} & \texttt{6.1} & 7 & 3\\
 &  & 5 & \texttt{0.8} & \texttt{6.1} & 8 & 3\\
T9.223 & \texttt{012331233} & 0 & \texttt{0.1.2.3} & \texttt{4.2.2.1} & 3 & 4\\*
 &  & 1 & \texttt{0.1.2.4} & \texttt{4.2.2.1} & 4 & 4\\
 &  & 2 & \texttt{0.5.6.7} & \texttt{4.2.2.1} & 7 & 4\\
 &  & 3 & \texttt{0.5.6.8} & \texttt{4.2.2.1} & 8 & 4\\
T9.224 & \texttt{012331232} & 0 & \texttt{0.1.2.3} & \texttt{4.2.2.1} & 3 & 2\\*
 &  & 1 & \texttt{0.1.2.4} & \texttt{4.2.2.1} & 4 & 2\\
 &  & 2 & \texttt{0.5.6.7} & \texttt{4.2.1.1} & 6 & 1\\
 &  & 3 & \texttt{0.5.8} & \texttt{4.2.1} & 8 & 1\\
T9.225 & \texttt{012331231} & 0 & \texttt{0.1.2.3} & \texttt{4.2.2.1} & 3 & 2\\*
 &  & 1 & \texttt{0.1.2.4} & \texttt{4.2.2.1} & 4 & 2\\
 &  & 2 & \texttt{0.5.6.7} & \texttt{4.1.1.1} & 5 & 1\\
 &  & 3 & \texttt{0.8} & \texttt{4.1} & 8 & 1\\
T9.226 & \texttt{012331222} & 0 & \texttt{0.1.2.3} & \texttt{5.2.2.1} & 3 & 2\\*
 &  & 1 & \texttt{0.1.2.4} & \texttt{5.2.2.1} & 4 & 2\\
 &  & 2 & \texttt{0.5.6} & \texttt{5.3.1} & 6 & 3\\
 &  & 3 & \texttt{0.5.7} & \texttt{5.3.1} & 7 & 3\\
 &  & 4 & \texttt{0.5.8} & \texttt{5.3.1} & 8 & 3\\
T9.227 & \texttt{012331221} & 0 & \texttt{0.1.2.3} & \texttt{5.2.2.1} & 3 & 2\\*
 &  & 1 & \texttt{0.1.2.4} & \texttt{5.2.2.1} & 4 & 2\\
 &  & 2 & \texttt{0.5.6} & \texttt{5.2.1} & 6 & 2\\
 &  & 3 & \texttt{0.5.7} & \texttt{5.2.1} & 7 & 2\\
 &  & 4 & \texttt{0.8} & \texttt{5.1} & 8 & 1\\
T9.228 & \texttt{012331212} & 0 & \texttt{0.1.2.3} & \texttt{4.2.2.1} & 3 & 2\\*
 &  & 1 & \texttt{0.1.2.4} & \texttt{4.2.2.1} & 4 & 2\\
 &  & 2 & \texttt{0.5.6} & \texttt{4.1.1} & 5 & 2\\
 &  & 3 & \texttt{0.7.8} & \texttt{4.1.1} & 7 & 2\\
T9.229 & \texttt{012331211} & 0 & \texttt{0.1.2.3} & \texttt{5.2.2.1} & 3 & 2\\*
 &  & 1 & \texttt{0.1.2.4} & \texttt{5.2.2.1} & 4 & 2\\
 &  & 2 & \texttt{0.5.6} & \texttt{5.1.1} & 5 & 1\\
 &  & 3 & \texttt{0.7} & \texttt{5.1} & 7 & 2\\
 &  & 4 & \texttt{0.8} & \texttt{5.1} & 8 & 2\\
T9.230 & \texttt{012331111} & 0 & \texttt{0.1.2.3} & \texttt{6.2.2.1} & 3 & 2\\*
 &  & 1 & \texttt{0.1.2.4} & \texttt{6.2.2.1} & 4 & 2\\
 &  & 2 & \texttt{0.5} & \texttt{6.1} & 5 & 4\\
 &  & 3 & \texttt{0.6} & \texttt{6.1} & 6 & 4\\
 &  & 4 & \texttt{0.7} & \texttt{6.1} & 7 & 4\\
 &  & 5 & \texttt{0.8} & \texttt{6.1} & 8 & 4\\
T9.231 & \texttt{012323232} & 0 & \texttt{0.1.2.3} & \texttt{4.4.1.1} & 2 & 3\\*
 &  & 1 & \texttt{0.1.4.5} & \texttt{4.4.1.1} & 4 & 3\\
 &  & 2 & \texttt{0.1.6.7} & \texttt{4.4.1.1} & 6 & 3\\
 &  & 3 & \texttt{0.1.8} & \texttt{4.4.1} & 8 & 1\\
T9.232 & \texttt{012323231} & 0 & \texttt{0.1.2.3} & \texttt{4.3.1.1} & 2 & 3\\*
 &  & 1 & \texttt{0.1.4.5} & \texttt{4.3.1.1} & 4 & 3\\
 &  & 2 & \texttt{0.1.6.7} & \texttt{4.3.1.1} & 6 & 3\\
 &  & 3 & \texttt{0.8} & \texttt{4.1} & 8 & 1\\
T9.233 & \texttt{012323222} & 0 & \texttt{0.1.2.3} & \texttt{5.5.1.1} & 2 & 2\\*
 &  & 1 & \texttt{0.1.4.5} & \texttt{5.5.1.1} & 4 & 2\\
 &  & 2 & \texttt{0.1.6} & \texttt{5.5.1} & 6 & 3\\
 &  & 3 & \texttt{0.1.7} & \texttt{5.5.1} & 7 & 3\\
 &  & 4 & \texttt{0.1.8} & \texttt{5.5.1} & 8 & 3\\
T9.234 & \texttt{012323221} & 0 & \texttt{0.1.2.3} & \texttt{5.4.1.1} & 2 & 2\\*
 &  & 1 & \texttt{0.1.4.5} & \texttt{5.4.1.1} & 4 & 2\\
 &  & 2 & \texttt{0.1.6} & \texttt{5.4.1} & 6 & 2\\
 &  & 3 & \texttt{0.1.7} & \texttt{5.4.1} & 7 & 2\\
 &  & 4 & \texttt{0.8} & \texttt{5.1} & 8 & 1\\
T9.235 & \texttt{012323212} & 0 & \texttt{0.1.2.3} & \texttt{4.3.1.1} & 2 & 2\\*
 &  & 1 & \texttt{0.1.4.5} & \texttt{4.3.1.1} & 4 & 2\\
 &  & 2 & \texttt{0.1.6} & \texttt{4.3.1} & 6 & 1\\
 &  & 3 & \texttt{0.7.8} & \texttt{4.1.1} & 7 & 1\\
T9.236 & \texttt{012323211} & 0 & \texttt{0.1.2.3} & \texttt{5.3.1.1} & 2 & 2\\*
 &  & 1 & \texttt{0.1.4.5} & \texttt{5.3.1.1} & 4 & 2\\
 &  & 2 & \texttt{0.1.6} & \texttt{5.3.1} & 6 & 1\\
 &  & 3 & \texttt{0.7} & \texttt{5.1} & 7 & 2\\
 &  & 4 & \texttt{0.8} & \texttt{5.1} & 8 & 2\\
T9.237 & \texttt{012323123} & 0 & \texttt{0.1.2.3} & \texttt{3.2.1.1} & 2 & 2\\*
 &  & 1 & \texttt{0.1.4.5} & \texttt{3.2.1.1} & 4 & 2\\
 &  & 2 & \texttt{0.6.7.8} & \texttt{3.1.1.1} & 6 & 1\\
T9.238 & \texttt{012323122} & 0 & \texttt{0.1.2.3} & \texttt{4.2.1.1} & 2 & 2\\*
 &  & 1 & \texttt{0.1.4.5} & \texttt{4.2.1.1} & 4 & 2\\
 &  & 2 & \texttt{0.6.7} & \texttt{4.2.1} & 7 & 2\\
 &  & 3 & \texttt{0.6.8} & \texttt{4.2.1} & 8 & 2\\
T9.239 & \texttt{012323121} & 0 & \texttt{0.1.2.3} & \texttt{4.2.1.1} & 2 & 2\\*
 &  & 1 & \texttt{0.1.4.5} & \texttt{4.2.1.1} & 4 & 2\\
 &  & 2 & \texttt{0.6.7} & \texttt{4.1.1} & 6 & 1\\
 &  & 3 & \texttt{0.8} & \texttt{4.1} & 8 & 1\\
T9.240 & \texttt{012323111} & 0 & \texttt{0.1.2.3} & \texttt{5.2.1.1} & 2 & 2\\*
 &  & 1 & \texttt{0.1.4.5} & \texttt{5.2.1.1} & 4 & 2\\
 &  & 2 & \texttt{0.6} & \texttt{5.1} & 6 & 3\\
 &  & 3 & \texttt{0.7} & \texttt{5.1} & 7 & 3\\
 &  & 4 & \texttt{0.8} & \texttt{5.1} & 8 & 3\\
T9.241 & \texttt{012322222} & 0 & \texttt{0.1.2.3} & \texttt{6.6.1.1} & 2 & 1\\*
 &  & 1 & \texttt{0.1.4} & \texttt{6.6.1} & 4 & 5\\
 &  & 2 & \texttt{0.1.5} & \texttt{6.6.1} & 5 & 5\\
 &  & 3 & \texttt{0.1.6} & \texttt{6.6.1} & 6 & 5\\
 &  & 4 & \texttt{0.1.7} & \texttt{6.6.1} & 7 & 5\\
 &  & 5 & \texttt{0.1.8} & \texttt{6.6.1} & 8 & 5\\
T9.242 & \texttt{012322221} & 0 & \texttt{0.1.2.3} & \texttt{6.5.1.1} & 2 & 1\\*
 &  & 1 & \texttt{0.1.4} & \texttt{6.5.1} & 4 & 4\\
 &  & 2 & \texttt{0.1.5} & \texttt{6.5.1} & 5 & 4\\
 &  & 3 & \texttt{0.1.6} & \texttt{6.5.1} & 6 & 4\\
 &  & 4 & \texttt{0.1.7} & \texttt{6.5.1} & 7 & 4\\
 &  & 5 & \texttt{0.8} & \texttt{6.1} & 8 & 1\\
T9.243 & \texttt{012322212} & 0 & \texttt{0.1.2.3} & \texttt{5.4.1.1} & 2 & 1\\*
 &  & 1 & \texttt{0.1.4} & \texttt{5.4.1} & 4 & 3\\
 &  & 2 & \texttt{0.1.5} & \texttt{5.4.1} & 5 & 3\\
 &  & 3 & \texttt{0.1.6} & \texttt{5.4.1} & 6 & 3\\
 &  & 4 & \texttt{0.7.8} & \texttt{5.1.1} & 7 & 1\\
T9.244 & \texttt{012322211} & 0 & \texttt{0.1.2.3} & \texttt{6.4.1.1} & 2 & 1\\*
 &  & 1 & \texttt{0.1.4} & \texttt{6.4.1} & 4 & 3\\
 &  & 2 & \texttt{0.1.5} & \texttt{6.4.1} & 5 & 3\\
 &  & 3 & \texttt{0.1.6} & \texttt{6.4.1} & 6 & 3\\
 &  & 4 & \texttt{0.7} & \texttt{6.1} & 7 & 2\\
 &  & 5 & \texttt{0.8} & \texttt{6.1} & 8 & 2\\
T9.245 & \texttt{012322123} & 0 & \texttt{0.1.2.3} & \texttt{4.3.1.1} & 2 & 1\\*
 &  & 1 & \texttt{0.1.4} & \texttt{4.3.1} & 4 & 2\\
 &  & 2 & \texttt{0.1.5} & \texttt{4.3.1} & 5 & 2\\
 &  & 3 & \texttt{0.6.7.8} & \texttt{4.1.1.1} & 6 & 1\\
T9.246 & \texttt{012322122} & 0 & \texttt{0.1.2.3} & \texttt{5.3.1.1} & 2 & 1\\*
 &  & 1 & \texttt{0.1.4} & \texttt{5.3.1} & 4 & 2\\
 &  & 2 & \texttt{0.1.5} & \texttt{5.3.1} & 5 & 2\\
 &  & 3 & \texttt{0.6.7} & \texttt{5.2.1} & 7 & 2\\
 &  & 4 & \texttt{0.6.8} & \texttt{5.2.1} & 8 & 2\\
T9.247 & \texttt{012322121} & 0 & \texttt{0.1.2.3} & \texttt{5.3.1.1} & 2 & 1\\*
 &  & 1 & \texttt{0.1.4} & \texttt{5.3.1} & 4 & 2\\
 &  & 2 & \texttt{0.1.5} & \texttt{5.3.1} & 5 & 2\\
 &  & 3 & \texttt{0.6.7} & \texttt{5.1.1} & 6 & 1\\
 &  & 4 & \texttt{0.8} & \texttt{5.1} & 8 & 1\\
T9.248 & \texttt{012322111} & 0 & \texttt{0.1.2.3} & \texttt{6.3.1.1} & 2 & 1\\*
 &  & 1 & \texttt{0.1.4} & \texttt{6.3.1} & 4 & 2\\
 &  & 2 & \texttt{0.1.5} & \texttt{6.3.1} & 5 & 2\\
 &  & 3 & \texttt{0.6} & \texttt{6.1} & 6 & 3\\
 &  & 4 & \texttt{0.7} & \texttt{6.1} & 7 & 3\\
 &  & 5 & \texttt{0.8} & \texttt{6.1} & 8 & 3\\
T9.249 & \texttt{012321232} & 0 & \texttt{0.1.2.3} & \texttt{4.2.1.1} & 2 & 2\\*
 &  & 1 & \texttt{0.1.4} & \texttt{4.2.1} & 4 & 2\\
 &  & 2 & \texttt{0.5.6.7} & \texttt{4.2.1.1} & 6 & 2\\
 &  & 3 & \texttt{0.5.8} & \texttt{4.2.1} & 8 & 2\\
T9.250 & \texttt{012321231} & 0 & \texttt{0.1.2.3} & \texttt{4.2.1.1} & 2 & 1\\*
 &  & 1 & \texttt{0.1.4} & \texttt{4.2.1} & 4 & 1\\
 &  & 2 & \texttt{0.5.6.7} & \texttt{4.1.1.1} & 5 & 1\\
 &  & 3 & \texttt{0.8} & \texttt{4.1} & 8 & 1\\
T9.251 & \texttt{012321222} & 0 & \texttt{0.1.2.3} & \texttt{5.2.1.1} & 2 & 1\\*
 &  & 1 & \texttt{0.1.4} & \texttt{5.2.1} & 4 & 1\\
 &  & 2 & \texttt{0.5.6} & \texttt{5.3.1} & 6 & 3\\
 &  & 3 & \texttt{0.5.7} & \texttt{5.3.1} & 7 & 3\\
 &  & 4 & \texttt{0.5.8} & \texttt{5.3.1} & 8 & 3\\
T9.252 & \texttt{012321221} & 0 & \texttt{0.1.2.3} & \texttt{5.2.1.1} & 2 & 1\\*
 &  & 1 & \texttt{0.1.4} & \texttt{5.2.1} & 4 & 1\\
 &  & 2 & \texttt{0.5.6} & \texttt{5.2.1} & 6 & 2\\
 &  & 3 & \texttt{0.5.7} & \texttt{5.2.1} & 7 & 2\\
 &  & 4 & \texttt{0.8} & \texttt{5.1} & 8 & 1\\
T9.253 & \texttt{012321212} & 0 & \texttt{0.1.2.3} & \texttt{4.2.1.1} & 2 & 1\\*
 &  & 1 & \texttt{0.1.4} & \texttt{4.2.1} & 4 & 1\\
 &  & 2 & \texttt{0.5.6} & \texttt{4.1.1} & 5 & 2\\
 &  & 3 & \texttt{0.7.8} & \texttt{4.1.1} & 7 & 2\\
T9.254 & \texttt{012321211} & 0 & \texttt{0.1.2.3} & \texttt{5.2.1.1} & 2 & 1\\*
 &  & 1 & \texttt{0.1.4} & \texttt{5.2.1} & 4 & 1\\
 &  & 2 & \texttt{0.5.6} & \texttt{5.1.1} & 5 & 1\\
 &  & 3 & \texttt{0.7} & \texttt{5.1} & 7 & 2\\
 &  & 4 & \texttt{0.8} & \texttt{5.1} & 8 & 2\\
T9.255 & \texttt{012321111} & 0 & \texttt{0.1.2.3} & \texttt{6.2.1.1} & 2 & 1\\*
 &  & 1 & \texttt{0.1.4} & \texttt{6.2.1} & 4 & 1\\
 &  & 2 & \texttt{0.5} & \texttt{6.1} & 5 & 4\\
 &  & 3 & \texttt{0.6} & \texttt{6.1} & 6 & 4\\
 &  & 4 & \texttt{0.7} & \texttt{6.1} & 7 & 4\\
 &  & 5 & \texttt{0.8} & \texttt{6.1} & 8 & 4\\
T9.256 & \texttt{012312312} & 0 & \texttt{0.1.2.3} & \texttt{3.1.1.1} & 1 & 2\\*
 &  & 1 & \texttt{0.4.5.6} & \texttt{3.1.1.1} & 4 & 2\\
 &  & 2 & \texttt{0.7.8} & \texttt{3.1.1} & 7 & 1\\
T9.257 & \texttt{012312311} & 0 & \texttt{0.1.2.3} & \texttt{4.1.1.1} & 1 & 2\\*
 &  & 1 & \texttt{0.4.5.6} & \texttt{4.1.1.1} & 4 & 2\\
 &  & 2 & \texttt{0.7} & \texttt{4.1} & 7 & 2\\
 &  & 3 & \texttt{0.8} & \texttt{4.1} & 8 & 2\\
T9.258 & \texttt{012312222} & 0 & \texttt{0.1.2.3} & \texttt{5.1.1.1} & 1 & 1\\*
 &  & 1 & \texttt{0.4.5} & \texttt{5.4.1} & 5 & 4\\
 &  & 2 & \texttt{0.4.6} & \texttt{5.4.1} & 6 & 4\\
 &  & 3 & \texttt{0.4.7} & \texttt{5.4.1} & 7 & 4\\
 &  & 4 & \texttt{0.4.8} & \texttt{5.4.1} & 8 & 4\\
T9.259 & \texttt{012312221} & 0 & \texttt{0.1.2.3} & \texttt{5.1.1.1} & 1 & 1\\*
 &  & 1 & \texttt{0.4.5} & \texttt{5.3.1} & 5 & 3\\
 &  & 2 & \texttt{0.4.6} & \texttt{5.3.1} & 6 & 3\\
 &  & 3 & \texttt{0.4.7} & \texttt{5.3.1} & 7 & 3\\
 &  & 4 & \texttt{0.8} & \texttt{5.1} & 8 & 1\\
T9.260 & \texttt{012312212} & 0 & \texttt{0.1.2.3} & \texttt{4.1.1.1} & 1 & 1\\*
 &  & 1 & \texttt{0.4.5} & \texttt{4.2.1} & 5 & 2\\
 &  & 2 & \texttt{0.4.6} & \texttt{4.2.1} & 6 & 2\\
 &  & 3 & \texttt{0.7.8} & \texttt{4.1.1} & 7 & 1\\
T9.261 & \texttt{012312211} & 0 & \texttt{0.1.2.3} & \texttt{5.1.1.1} & 1 & 1\\*
 &  & 1 & \texttt{0.4.5} & \texttt{5.2.1} & 5 & 2\\
 &  & 2 & \texttt{0.4.6} & \texttt{5.2.1} & 6 & 2\\
 &  & 3 & \texttt{0.7} & \texttt{5.1} & 7 & 2\\
 &  & 4 & \texttt{0.8} & \texttt{5.1} & 8 & 2\\
T9.262 & \texttt{012312121} & 0 & \texttt{0.1.2.3} & \texttt{4.1.1.1} & 1 & 1\\*
 &  & 1 & \texttt{0.4.5} & \texttt{4.1.1} & 4 & 2\\
 &  & 2 & \texttt{0.6.7} & \texttt{4.1.1} & 6 & 2\\
 &  & 3 & \texttt{0.8} & \texttt{4.1} & 8 & 1\\
T9.263 & \texttt{012312111} & 0 & \texttt{0.1.2.3} & \texttt{5.1.1.1} & 1 & 1\\*
 &  & 1 & \texttt{0.4.5} & \texttt{5.1.1} & 4 & 1\\
 &  & 2 & \texttt{0.6} & \texttt{5.1} & 6 & 3\\
 &  & 3 & \texttt{0.7} & \texttt{5.1} & 7 & 3\\
 &  & 4 & \texttt{0.8} & \texttt{5.1} & 8 & 3\\
T9.264 & \texttt{012311111} & 0 & \texttt{0.1.2.3} & \texttt{6.1.1.1} & 1 & 1\\*
 &  & 1 & \texttt{0.4} & \texttt{6.1} & 4 & 5\\
 &  & 2 & \texttt{0.5} & \texttt{6.1} & 5 & 5\\
 &  & 3 & \texttt{0.6} & \texttt{6.1} & 6 & 5\\
 &  & 4 & \texttt{0.7} & \texttt{6.1} & 7 & 5\\
 &  & 5 & \texttt{0.8} & \texttt{6.1} & 8 & 5\\
T9.265 & \texttt{012222222} & 0 & \texttt{0.1.2} & \texttt{7.7.1} & 2 & 7\\*
 &  & 1 & \texttt{0.1.3} & \texttt{7.7.1} & 3 & 7\\
 &  & 2 & \texttt{0.1.4} & \texttt{7.7.1} & 4 & 7\\
 &  & 3 & \texttt{0.1.5} & \texttt{7.7.1} & 5 & 7\\
 &  & 4 & \texttt{0.1.6} & \texttt{7.7.1} & 6 & 7\\
 &  & 5 & \texttt{0.1.7} & \texttt{7.7.1} & 7 & 7\\
 &  & 6 & \texttt{0.1.8} & \texttt{7.7.1} & 8 & 7\\
T9.266 & \texttt{012222221} & 0 & \texttt{0.1.2} & \texttt{7.6.1} & 2 & 6\\*
 &  & 1 & \texttt{0.1.3} & \texttt{7.6.1} & 3 & 6\\
 &  & 2 & \texttt{0.1.4} & \texttt{7.6.1} & 4 & 6\\
 &  & 3 & \texttt{0.1.5} & \texttt{7.6.1} & 5 & 6\\
 &  & 4 & \texttt{0.1.6} & \texttt{7.6.1} & 6 & 6\\
 &  & 5 & \texttt{0.1.7} & \texttt{7.6.1} & 7 & 6\\
 &  & 6 & \texttt{0.8} & \texttt{7.1} & 8 & 1\\
T9.267 & \texttt{012222212} & 0 & \texttt{0.1.2} & \texttt{6.5.1} & 2 & 5\\*
 &  & 1 & \texttt{0.1.3} & \texttt{6.5.1} & 3 & 5\\
 &  & 2 & \texttt{0.1.4} & \texttt{6.5.1} & 4 & 5\\
 &  & 3 & \texttt{0.1.5} & \texttt{6.5.1} & 5 & 5\\
 &  & 4 & \texttt{0.1.6} & \texttt{6.5.1} & 6 & 5\\
 &  & 5 & \texttt{0.7.8} & \texttt{6.1.1} & 7 & 1\\
T9.268 & \texttt{012222211} & 0 & \texttt{0.1.2} & \texttt{7.5.1} & 2 & 5\\*
 &  & 1 & \texttt{0.1.3} & \texttt{7.5.1} & 3 & 5\\
 &  & 2 & \texttt{0.1.4} & \texttt{7.5.1} & 4 & 5\\
 &  & 3 & \texttt{0.1.5} & \texttt{7.5.1} & 5 & 5\\
 &  & 4 & \texttt{0.1.6} & \texttt{7.5.1} & 6 & 5\\
 &  & 5 & \texttt{0.7} & \texttt{7.1} & 7 & 2\\
 &  & 6 & \texttt{0.8} & \texttt{7.1} & 8 & 2\\
T9.269 & \texttt{012222122} & 0 & \texttt{0.1.2} & \texttt{6.4.1} & 2 & 4\\*
 &  & 1 & \texttt{0.1.3} & \texttt{6.4.1} & 3 & 4\\
 &  & 2 & \texttt{0.1.4} & \texttt{6.4.1} & 4 & 4\\
 &  & 3 & \texttt{0.1.5} & \texttt{6.4.1} & 5 & 4\\
 &  & 4 & \texttt{0.6.7} & \texttt{6.2.1} & 7 & 2\\
 &  & 5 & \texttt{0.6.8} & \texttt{6.2.1} & 8 & 2\\
T9.270 & \texttt{012222121} & 0 & \texttt{0.1.2} & \texttt{6.4.1} & 2 & 4\\*
 &  & 1 & \texttt{0.1.3} & \texttt{6.4.1} & 3 & 4\\
 &  & 2 & \texttt{0.1.4} & \texttt{6.4.1} & 4 & 4\\
 &  & 3 & \texttt{0.1.5} & \texttt{6.4.1} & 5 & 4\\
 &  & 4 & \texttt{0.6.7} & \texttt{6.1.1} & 6 & 1\\
 &  & 5 & \texttt{0.8} & \texttt{6.1} & 8 & 1\\
T9.271 & \texttt{012222111} & 0 & \texttt{0.1.2} & \texttt{7.4.1} & 2 & 4\\*
 &  & 1 & \texttt{0.1.3} & \texttt{7.4.1} & 3 & 4\\
 &  & 2 & \texttt{0.1.4} & \texttt{7.4.1} & 4 & 4\\
 &  & 3 & \texttt{0.1.5} & \texttt{7.4.1} & 5 & 4\\
 &  & 4 & \texttt{0.6} & \texttt{7.1} & 6 & 3\\
 &  & 5 & \texttt{0.7} & \texttt{7.1} & 7 & 3\\
 &  & 6 & \texttt{0.8} & \texttt{7.1} & 8 & 3\\
T9.272 & \texttt{012221222} & 0 & \texttt{0.1.2} & \texttt{6.3.1} & 2 & 6\\*
 &  & 1 & \texttt{0.1.3} & \texttt{6.3.1} & 3 & 6\\
 &  & 2 & \texttt{0.1.4} & \texttt{6.3.1} & 4 & 6\\
 &  & 3 & \texttt{0.5.6} & \texttt{6.3.1} & 6 & 6\\
 &  & 4 & \texttt{0.5.7} & \texttt{6.3.1} & 7 & 6\\
 &  & 5 & \texttt{0.5.8} & \texttt{6.3.1} & 8 & 6\\
T9.273 & \texttt{012221221} & 0 & \texttt{0.1.2} & \texttt{6.3.1} & 2 & 3\\*
 &  & 1 & \texttt{0.1.3} & \texttt{6.3.1} & 3 & 3\\
 &  & 2 & \texttt{0.1.4} & \texttt{6.3.1} & 4 & 3\\
 &  & 3 & \texttt{0.5.6} & \texttt{6.2.1} & 6 & 2\\
 &  & 4 & \texttt{0.5.7} & \texttt{6.2.1} & 7 & 2\\
 &  & 5 & \texttt{0.8} & \texttt{6.1} & 8 & 1\\
T9.274 & \texttt{012221212} & 0 & \texttt{0.1.2} & \texttt{5.3.1} & 2 & 3\\*
 &  & 1 & \texttt{0.1.3} & \texttt{5.3.1} & 3 & 3\\
 &  & 2 & \texttt{0.1.4} & \texttt{5.3.1} & 4 & 3\\
 &  & 3 & \texttt{0.5.6} & \texttt{5.1.1} & 5 & 2\\
 &  & 4 & \texttt{0.7.8} & \texttt{5.1.1} & 7 & 2\\
T9.275 & \texttt{012221211} & 0 & \texttt{0.1.2} & \texttt{6.3.1} & 2 & 3\\*
 &  & 1 & \texttt{0.1.3} & \texttt{6.3.1} & 3 & 3\\
 &  & 2 & \texttt{0.1.4} & \texttt{6.3.1} & 4 & 3\\
 &  & 3 & \texttt{0.5.6} & \texttt{6.1.1} & 5 & 1\\
 &  & 4 & \texttt{0.7} & \texttt{6.1} & 7 & 2\\
 &  & 5 & \texttt{0.8} & \texttt{6.1} & 8 & 2\\
T9.276 & \texttt{012221111} & 0 & \texttt{0.1.2} & \texttt{7.3.1} & 2 & 3\\*
 &  & 1 & \texttt{0.1.3} & \texttt{7.3.1} & 3 & 3\\
 &  & 2 & \texttt{0.1.4} & \texttt{7.3.1} & 4 & 3\\
 &  & 3 & \texttt{0.5} & \texttt{7.1} & 5 & 4\\
 &  & 4 & \texttt{0.6} & \texttt{7.1} & 6 & 4\\
 &  & 5 & \texttt{0.7} & \texttt{7.1} & 7 & 4\\
 &  & 6 & \texttt{0.8} & \texttt{7.1} & 8 & 4\\
T9.277 & \texttt{012212212} & 0 & \texttt{0.1.2} & \texttt{5.2.1} & 2 & 4\\*
 &  & 1 & \texttt{0.1.3} & \texttt{5.2.1} & 3 & 4\\
 &  & 2 & \texttt{0.4.5} & \texttt{5.2.1} & 5 & 4\\
 &  & 3 & \texttt{0.4.6} & \texttt{5.2.1} & 6 & 4\\
 &  & 4 & \texttt{0.7.8} & \texttt{5.1.1} & 7 & 1\\
T9.278 & \texttt{012212211} & 0 & \texttt{0.1.2} & \texttt{6.2.1} & 2 & 4\\*
 &  & 1 & \texttt{0.1.3} & \texttt{6.2.1} & 3 & 4\\
 &  & 2 & \texttt{0.4.5} & \texttt{6.2.1} & 5 & 4\\
 &  & 3 & \texttt{0.4.6} & \texttt{6.2.1} & 6 & 4\\
 &  & 4 & \texttt{0.7} & \texttt{6.1} & 7 & 2\\
 &  & 5 & \texttt{0.8} & \texttt{6.1} & 8 & 2\\
T9.279 & \texttt{012212121} & 0 & \texttt{0.1.2} & \texttt{5.2.1} & 2 & 2\\*
 &  & 1 & \texttt{0.1.3} & \texttt{5.2.1} & 3 & 2\\
 &  & 2 & \texttt{0.4.5} & \texttt{5.1.1} & 4 & 2\\
 &  & 3 & \texttt{0.6.7} & \texttt{5.1.1} & 6 & 2\\
 &  & 4 & \texttt{0.8} & \texttt{5.1} & 8 & 1\\
T9.280 & \texttt{012212111} & 0 & \texttt{0.1.2} & \texttt{6.2.1} & 2 & 2\\*
 &  & 1 & \texttt{0.1.3} & \texttt{6.2.1} & 3 & 2\\
 &  & 2 & \texttt{0.4.5} & \texttt{6.1.1} & 4 & 1\\
 &  & 3 & \texttt{0.6} & \texttt{6.1} & 6 & 3\\
 &  & 4 & \texttt{0.7} & \texttt{6.1} & 7 & 3\\
 &  & 5 & \texttt{0.8} & \texttt{6.1} & 8 & 3\\
T9.281 & \texttt{012211111} & 0 & \texttt{0.1.2} & \texttt{7.2.1} & 2 & 2\\*
 &  & 1 & \texttt{0.1.3} & \texttt{7.2.1} & 3 & 2\\
 &  & 2 & \texttt{0.4} & \texttt{7.1} & 4 & 5\\
 &  & 3 & \texttt{0.5} & \texttt{7.1} & 5 & 5\\
 &  & 4 & \texttt{0.6} & \texttt{7.1} & 6 & 5\\
 &  & 5 & \texttt{0.7} & \texttt{7.1} & 7 & 5\\
 &  & 6 & \texttt{0.8} & \texttt{7.1} & 8 & 5\\
T9.282 & \texttt{012121212} & 0 & \texttt{0.1.2} & \texttt{4.1.1} & 1 & 4\\*
 &  & 1 & \texttt{0.3.4} & \texttt{4.1.1} & 3 & 4\\
 &  & 2 & \texttt{0.5.6} & \texttt{4.1.1} & 5 & 4\\
 &  & 3 & \texttt{0.7.8} & \texttt{4.1.1} & 7 & 4\\
T9.283 & \texttt{012121211} & 0 & \texttt{0.1.2} & \texttt{5.1.1} & 1 & 3\\*
 &  & 1 & \texttt{0.3.4} & \texttt{5.1.1} & 3 & 3\\
 &  & 2 & \texttt{0.5.6} & \texttt{5.1.1} & 5 & 3\\
 &  & 3 & \texttt{0.7} & \texttt{5.1} & 7 & 2\\
 &  & 4 & \texttt{0.8} & \texttt{5.1} & 8 & 2\\
T9.284 & \texttt{012121111} & 0 & \texttt{0.1.2} & \texttt{6.1.1} & 1 & 2\\*
 &  & 1 & \texttt{0.3.4} & \texttt{6.1.1} & 3 & 2\\
 &  & 2 & \texttt{0.5} & \texttt{6.1} & 5 & 4\\
 &  & 3 & \texttt{0.6} & \texttt{6.1} & 6 & 4\\
 &  & 4 & \texttt{0.7} & \texttt{6.1} & 7 & 4\\
 &  & 5 & \texttt{0.8} & \texttt{6.1} & 8 & 4\\
T9.285 & \texttt{012111111} & 0 & \texttt{0.1.2} & \texttt{7.1.1} & 1 & 1\\*
 &  & 1 & \texttt{0.3} & \texttt{7.1} & 3 & 6\\
 &  & 2 & \texttt{0.4} & \texttt{7.1} & 4 & 6\\
 &  & 3 & \texttt{0.5} & \texttt{7.1} & 5 & 6\\
 &  & 4 & \texttt{0.6} & \texttt{7.1} & 6 & 6\\
 &  & 5 & \texttt{0.7} & \texttt{7.1} & 7 & 6\\
 &  & 6 & \texttt{0.8} & \texttt{7.1} & 8 & 6\\
T9.286 & \texttt{011111111} & 0 & \texttt{0.1} & \texttt{8.1} & 1 & 8\\*
 &  & 1 & \texttt{0.2} & \texttt{8.1} & 2 & 8\\
 &  & 2 & \texttt{0.3} & \texttt{8.1} & 3 & 8\\
 &  & 3 & \texttt{0.4} & \texttt{8.1} & 4 & 8\\
 &  & 4 & \texttt{0.5} & \texttt{8.1} & 5 & 8\\
 &  & 5 & \texttt{0.6} & \texttt{8.1} & 6 & 8\\
 &  & 6 & \texttt{0.7} & \texttt{8.1} & 7 & 8\\
 &  & 7 & \texttt{0.8} & \texttt{8.1} & 8 & 8\\
\bottomrule
\end{longtable}}


\section{Frames with eleven moments}\label{app:eleven}

\paragraph{What is listed.} One line for each of the \NFramesEleven{} frames with eleven moments, with the same frame-level data as the headers of Appendix~\ref{app:detail}: identifier, canonical level sequence, number of histories, length of a longest history, number of branching moments, order of the automorphism group, orbits on moments and fixed moments, orbits on histories and fixed histories, and the counts M, res and rev. The per-moment values are computed and checked exactly as for the smaller frames but are not printed. All checks hold in every case. The entries include the chain T11.0001, which has 32,796 reversals and no resolutions, and the star T11.1842, whose automorphism group has order $10!$ and in which, of the 20,480 instances of monotonicity checked, 20,440 are resolutions of an open future contingent at the root by a later moment.

% generated by gen/generate.py -- do not edit
{\scriptsize\setlength{\tabcolsep}{3.2pt}
\begin{longtable}{@{}llrrrrrrrrrrr@{}}
\toprule
id & level sequence & $|H|$ & len & br & $|\mathrm{Aut}|$ & orb & fix & h-orb & h-fix & M & res & rev\\
\midrule
\endfirsthead
\toprule
id & level sequence & $|H|$ & len & br & $|\mathrm{Aut}|$ & orb & fix & h-orb & h-fix & M & res & rev\\
\midrule
\endhead
T11.0001 & \texttt{0123456789a} & 1 & 11 & 0 & 1 & 11 & 11 & 1 & 1 & 112640 & 0 & 32796\\
T11.0002 & \texttt{01234567899} & 2 & 10 & 1 & 2 & 10 & 9 & 1 & 0 & 110592 & 4088 & 36900\\
T11.0003 & \texttt{01234567898} & 2 & 10 & 1 & 1 & 11 & 11 & 2 & 2 & 108544 & 6120 & 37936\\
T11.0004 & \texttt{01234567897} & 2 & 10 & 1 & 1 & 11 & 11 & 2 & 2 & 106496 & 8128 & 36936\\
T11.0005 & \texttt{01234567896} & 2 & 10 & 1 & 1 & 11 & 11 & 2 & 2 & 104448 & 10080 & 35192\\
T11.0006 & \texttt{01234567895} & 2 & 10 & 1 & 1 & 11 & 11 & 2 & 2 & 102400 & 11904 & 33240\\
T11.0007 & \texttt{01234567894} & 2 & 10 & 1 & 1 & 11 & 11 & 2 & 2 & 100352 & 13440 & 31320\\
T11.0008 & \texttt{01234567893} & 2 & 10 & 1 & 1 & 11 & 11 & 2 & 2 & 98304 & 14336 & 29592\\
T11.0009 & \texttt{01234567892} & 2 & 10 & 1 & 1 & 11 & 11 & 2 & 2 & 96256 & 13824 & 28264\\
T11.0010 & \texttt{01234567891} & 2 & 10 & 1 & 1 & 11 & 11 & 2 & 2 & 94208 & 10240 & 27720\\
T11.0011 & \texttt{01234567888} & 3 & 9 & 1 & 6 & 9 & 8 & 1 & 0 & 106496 & 9180 & 41024\\
T11.0012 & \texttt{01234567887} & 3 & 9 & 2 & 2 & 10 & 9 & 2 & 1 & 104448 & 11192 & 39516\\
T11.0013 & \texttt{01234567886} & 3 & 9 & 2 & 2 & 10 & 9 & 2 & 1 & 102400 & 13404 & 37524\\
T11.0014 & \texttt{01234567885} & 3 & 9 & 2 & 2 & 10 & 9 & 2 & 1 & 100352 & 15488 & 35460\\
T11.0015 & \texttt{01234567884} & 3 & 9 & 2 & 2 & 10 & 9 & 2 & 1 & 98304 & 17220 & 33508\\
T11.0016 & \texttt{01234567883} & 3 & 9 & 2 & 2 & 10 & 9 & 2 & 1 & 96256 & 18248 & 31812\\
T11.0017 & \texttt{01234567882} & 3 & 9 & 2 & 2 & 10 & 9 & 2 & 1 & 94208 & 17820 & 30596\\
T11.0018 & \texttt{01234567881} & 3 & 9 & 2 & 2 & 10 & 9 & 2 & 1 & 92160 & 14288 & 30300\\
T11.0019 & \texttt{01234567878} & 2 & 9 & 1 & 2 & 9 & 7 & 1 & 0 & 104448 & 6096 & 40004\\
T11.0020 & \texttt{01234567877} & 3 & 9 & 1 & 2 & 10 & 9 & 2 & 1 & 102400 & 12192 & 40560\\
T11.0021 & \texttt{01234567876} & 3 & 9 & 2 & 1 & 11 & 11 & 3 & 3 & 100352 & 14664 & 38060\\
T11.0022 & \texttt{01234567875} & 3 & 9 & 2 & 1 & 11 & 11 & 3 & 3 & 98304 & 17008 & 35748\\
T11.0023 & \texttt{01234567874} & 3 & 9 & 2 & 1 & 11 & 11 & 3 & 3 & 96256 & 18936 & 33684\\
T11.0024 & \texttt{01234567873} & 3 & 9 & 2 & 1 & 11 & 11 & 3 & 3 & 94208 & 20096 & 31956\\
T11.0025 & \texttt{01234567872} & 3 & 9 & 2 & 1 & 11 & 11 & 3 & 3 & 92160 & 19752 & 30772\\
T11.0026 & \texttt{01234567871} & 3 & 9 & 2 & 1 & 11 & 11 & 3 & 3 & 90112 & 16272 & 30588\\
T11.0027 & \texttt{01234567867} & 2 & 9 & 1 & 1 & 11 & 11 & 2 & 2 & 100352 & 6300 & 39016\\
T11.0028 & \texttt{01234567866} & 3 & 9 & 1 & 2 & 10 & 9 & 2 & 1 & 98304 & 15120 & 37584\\
T11.0029 & \texttt{01234567865} & 3 & 9 & 2 & 1 & 11 & 11 & 3 & 3 & 96256 & 17984 & 34764\\
T11.0030 & \texttt{01234567864} & 3 & 9 & 2 & 1 & 11 & 11 & 3 & 3 & 94208 & 20304 & 32452\\
T11.0031 & \texttt{01234567863} & 3 & 9 & 2 & 1 & 11 & 11 & 3 & 3 & 92160 & 21728 & 30612\\
T11.0032 & \texttt{01234567862} & 3 & 9 & 2 & 1 & 11 & 11 & 3 & 3 & 90112 & 21552 & 29396\\
T11.0033 & \texttt{01234567861} & 3 & 9 & 2 & 1 & 11 & 11 & 3 & 3 & 88064 & 18176 & 29244\\
T11.0034 & \texttt{01234567856} & 2 & 9 & 1 & 1 & 11 & 11 & 2 & 2 & 96256 & 6696 & 36780\\
T11.0035 & \texttt{01234567855} & 3 & 9 & 1 & 2 & 10 & 9 & 2 & 1 & 94208 & 17856 & 33808\\
T11.0036 & \texttt{01234567854} & 3 & 9 & 2 & 1 & 11 & 11 & 3 & 3 & 92160 & 20960 & 30988\\
T11.0037 & \texttt{01234567853} & 3 & 9 & 2 & 1 & 11 & 11 & 3 & 3 & 90112 & 22912 & 28900\\
T11.0038 & \texttt{01234567852} & 3 & 9 & 2 & 1 & 11 & 11 & 3 & 3 & 88064 & 23072 & 27572\\
T11.0039 & \texttt{01234567851} & 3 & 9 & 2 & 1 & 11 & 11 & 3 & 3 & 86016 & 19904 & 27388\\
T11.0040 & \texttt{01234567845} & 2 & 9 & 1 & 1 & 11 & 11 & 2 & 2 & 92160 & 7140 & 34096\\
T11.0041 & \texttt{01234567844} & 3 & 9 & 1 & 2 & 10 & 9 & 2 & 1 & 90112 & 20160 & 29968\\
T11.0042 & \texttt{01234567843} & 3 & 9 & 2 & 1 & 11 & 11 & 3 & 3 & 88064 & 23168 & 27372\\
T11.0043 & \texttt{01234567842} & 3 & 9 & 2 & 1 & 11 & 11 & 3 & 3 & 86016 & 24000 & 25796\\
T11.0044 & \texttt{01234567841} & 3 & 9 & 2 & 1 & 11 & 11 & 3 & 3 & 83968 & 21248 & 25500\\
T11.0045 & \texttt{01234567834} & 2 & 9 & 1 & 1 & 11 & 11 & 2 & 2 & 88064 & 7392 & 31348\\
T11.0046 & \texttt{01234567833} & 3 & 9 & 1 & 2 & 10 & 9 & 2 & 1 & 86016 & 21504 & 26480\\
T11.0047 & \texttt{01234567832} & 3 & 9 & 2 & 1 & 11 & 11 & 3 & 3 & 83968 & 23680 & 24396\\
T11.0048 & \texttt{01234567831} & 3 & 9 & 2 & 1 & 11 & 11 & 3 & 3 & 81920 & 21760 & 23852\\
T11.0049 & \texttt{01234567823} & 2 & 9 & 1 & 1 & 11 & 11 & 2 & 2 & 83968 & 7020 & 28792\\
T11.0050 & \texttt{01234567822} & 3 & 9 & 1 & 2 & 10 & 9 & 2 & 1 & 81920 & 20736 & 23792\\
T11.0051 & \texttt{01234567821} & 3 & 9 & 2 & 1 & 11 & 11 & 3 & 3 & 79872 & 20480 & 22740\\
T11.0052 & \texttt{01234567812} & 2 & 9 & 1 & 1 & 11 & 11 & 2 & 2 & 79872 & 5160 & 26732\\
T11.0053 & \texttt{01234567811} & 3 & 9 & 1 & 2 & 10 & 9 & 2 & 1 & 77824 & 15360 & 22680\\
T11.0054 & \texttt{01234567777} & 4 & 8 & 1 & 24 & 8 & 7 & 1 & 0 & 100352 & 14224 & 43148\\
T11.0055 & \texttt{01234567776} & 4 & 8 & 2 & 6 & 9 & 8 & 2 & 1 & 98304 & 16956 & 39888\\
T11.0056 & \texttt{01234567775} & 4 & 8 & 2 & 6 & 9 & 8 & 2 & 1 & 96256 & 19560 & 37208\\
T11.0057 & \texttt{01234567774} & 4 & 8 & 2 & 6 & 9 & 8 & 2 & 1 & 94208 & 21684 & 34984\\
T11.0058 & \texttt{01234567773} & 4 & 8 & 2 & 6 & 9 & 8 & 2 & 1 & 92160 & 22976 & 33224\\
T11.0059 & \texttt{01234567772} & 4 & 8 & 2 & 6 & 9 & 8 & 2 & 1 & 90112 & 22716 & 32120\\
T11.0060 & \texttt{01234567771} & 4 & 8 & 2 & 6 & 9 & 8 & 2 & 1 & 88064 & 19288 & 32168\\
T11.0061 & \texttt{01234567767} & 3 & 8 & 2 & 2 & 10 & 9 & 2 & 1 & 98304 & 11120 & 40088\\
T11.0062 & \texttt{01234567766} & 4 & 8 & 2 & 4 & 9 & 7 & 2 & 0 & 96256 & 17420 & 39160\\
T11.0063 & \texttt{01234567765} & 4 & 8 & 3 & 2 & 10 & 9 & 3 & 2 & 94208 & 20544 & 35844\\
T11.0064 & \texttt{01234567764} & 4 & 8 & 3 & 2 & 10 & 9 & 3 & 2 & 92160 & 23060 & 33308\\
T11.0065 & \texttt{01234567763} & 4 & 8 & 3 & 2 & 10 & 9 & 3 & 2 & 90112 & 24616 & 31404\\
T11.0066 & \texttt{01234567762} & 4 & 8 & 3 & 2 & 10 & 9 & 3 & 2 & 88064 & 24524 & 30252\\
T11.0067 & \texttt{01234567761} & 4 & 8 & 3 & 2 & 10 & 9 & 3 & 2 & 86016 & 21200 & 30324\\
T11.0068 & \texttt{01234567756} & 3 & 8 & 2 & 2 & 10 & 9 & 2 & 1 & 94208 & 11008 & 38116\\
T11.0069 & \texttt{01234567755} & 4 & 8 & 2 & 4 & 9 & 7 & 2 & 0 & 92160 & 20680 & 34896\\
T11.0070 & \texttt{01234567754} & 4 & 8 & 3 & 2 & 10 & 9 & 3 & 2 & 90112 & 23980 & 31724\\
T11.0071 & \texttt{01234567753} & 4 & 8 & 3 & 2 & 10 & 9 & 3 & 2 & 88064 & 26064 & 29508\\
T11.0072 & \texttt{01234567752} & 4 & 8 & 3 & 2 & 10 & 9 & 3 & 2 & 86016 & 26308 & 28212\\
T11.0073 & \texttt{01234567751} & 4 & 8 & 3 & 2 & 10 & 9 & 3 & 2 & 83968 & 23192 & 28236\\
T11.0074 & \texttt{01234567745} & 3 & 8 & 2 & 2 & 10 & 9 & 2 & 1 & 90112 & 11264 & 35576\\
T11.0075 & \texttt{01234567744} & 4 & 8 & 2 & 4 & 9 & 7 & 2 & 0 & 88064 & 23444 & 30848\\
T11.0076 & \texttt{01234567743} & 4 & 8 & 3 & 2 & 10 & 9 & 3 & 2 & 86016 & 26584 & 27996\\
T11.0077 & \texttt{01234567742} & 4 & 8 & 3 & 2 & 10 & 9 & 3 & 2 & 83968 & 27500 & 26388\\
T11.0078 & \texttt{01234567741} & 4 & 8 & 3 & 2 & 10 & 9 & 3 & 2 & 81920 & 24800 & 26268\\
T11.0079 & \texttt{01234567734} & 3 & 8 & 2 & 2 & 10 & 9 & 2 & 1 & 86016 & 11456 & 32924\\
T11.0080 & \texttt{01234567733} & 4 & 8 & 2 & 4 & 9 & 7 & 2 & 0 & 83968 & 25120 & 27328\\
T11.0081 & \texttt{01234567732} & 4 & 8 & 3 & 2 & 10 & 9 & 3 & 2 & 81920 & 27380 & 25084\\
T11.0082 & \texttt{01234567731} & 4 & 8 & 3 & 2 & 10 & 9 & 3 & 2 & 79872 & 25512 & 24652\\
T11.0083 & \texttt{01234567723} & 3 & 8 & 2 & 2 & 10 & 9 & 2 & 1 & 81920 & 11072 & 30464\\
T11.0084 & \texttt{01234567722} & 4 & 8 & 2 & 4 & 9 & 7 & 2 & 0 & 79872 & 24572 & 24768\\
T11.0085 & \texttt{01234567721} & 4 & 8 & 3 & 2 & 10 & 9 & 3 & 2 & 77824 & 24368 & 23700\\
T11.0086 & \texttt{01234567712} & 3 & 8 & 2 & 2 & 10 & 9 & 2 & 1 & 77824 & 9216 & 28548\\
T11.0087 & \texttt{01234567711} & 4 & 8 & 2 & 4 & 9 & 7 & 2 & 0 & 75776 & 19336 & 24008\\
T11.0088 & \texttt{01234567676} & 3 & 8 & 1 & 2 & 9 & 7 & 2 & 1 & 96256 & 13860 & 39632\\
T11.0089 & \texttt{01234567675} & 3 & 8 & 2 & 2 & 9 & 7 & 2 & 1 & 94208 & 16464 & 36820\\
T11.0090 & \texttt{01234567674} & 3 & 8 & 2 & 2 & 9 & 7 & 2 & 1 & 92160 & 18588 & 34524\\
T11.0091 & \texttt{01234567673} & 3 & 8 & 2 & 2 & 9 & 7 & 2 & 1 & 90112 & 19880 & 32716\\
T11.0092 & \texttt{01234567672} & 3 & 8 & 2 & 2 & 9 & 7 & 2 & 1 & 88064 & 19620 & 31564\\
T11.0093 & \texttt{01234567671} & 3 & 8 & 2 & 2 & 9 & 7 & 2 & 1 & 86016 & 16192 & 31540\\
T11.0094 & \texttt{01234567666} & 4 & 8 & 1 & 6 & 9 & 8 & 2 & 1 & 94208 & 17640 & 40216\\
T11.0095 & \texttt{01234567665} & 4 & 8 & 2 & 2 & 10 & 9 & 3 & 2 & 92160 & 21024 & 36396\\
T11.0096 & \texttt{01234567664} & 4 & 8 & 2 & 2 & 10 & 9 & 3 & 2 & 90112 & 23736 & 33620\\
T11.0097 & \texttt{01234567663} & 4 & 8 & 2 & 2 & 10 & 9 & 3 & 2 & 88064 & 25424 & 31620\\
T11.0098 & \texttt{01234567662} & 4 & 8 & 2 & 2 & 10 & 9 & 3 & 2 & 86016 & 25416 & 30468\\
T11.0099 & \texttt{01234567661} & 4 & 8 & 2 & 2 & 10 & 9 & 3 & 2 & 83968 & 22144 & 30636\\
T11.0100 & \texttt{01234567656} & 3 & 8 & 2 & 1 & 11 & 11 & 3 & 3 & 92160 & 13240 & 37792\\
T11.0101 & \texttt{01234567655} & 4 & 8 & 2 & 2 & 10 & 9 & 3 & 2 & 90112 & 21424 & 35068\\
T11.0102 & \texttt{01234567654} & 4 & 8 & 3 & 1 & 11 & 11 & 4 & 4 & 88064 & 24920 & 31592\\
T11.0103 & \texttt{01234567653} & 4 & 8 & 3 & 1 & 11 & 11 & 4 & 4 & 86016 & 27136 & 29248\\
T11.0104 & \texttt{01234567652} & 4 & 8 & 3 & 1 & 11 & 11 & 4 & 4 & 83968 & 27464 & 27936\\
T11.0105 & \texttt{01234567651} & 4 & 8 & 3 & 1 & 11 & 11 & 4 & 4 & 81920 & 24400 & 28048\\
T11.0106 & \texttt{01234567645} & 3 & 8 & 2 & 1 & 11 & 11 & 3 & 3 & 88064 & 13308 & 35324\\
T11.0107 & \texttt{01234567644} & 4 & 8 & 2 & 2 & 10 & 9 & 3 & 2 & 86016 & 24648 & 30596\\
T11.0108 & \texttt{01234567643} & 4 & 8 & 3 & 1 & 11 & 11 & 4 & 4 & 83968 & 27920 & 27552\\
T11.0109 & \texttt{01234567642} & 4 & 8 & 3 & 1 & 11 & 11 & 4 & 4 & 81920 & 28920 & 25896\\
T11.0110 & \texttt{01234567641} & 4 & 8 & 3 & 1 & 11 & 11 & 4 & 4 & 79872 & 26272 & 25848\\
T11.0111 & \texttt{01234567634} & 3 & 8 & 2 & 1 & 11 & 11 & 3 & 3 & 83968 & 13440 & 32720\\
T11.0112 & \texttt{01234567633} & 4 & 8 & 2 & 2 & 10 & 9 & 3 & 2 & 81920 & 26656 & 26900\\
T11.0113 & \texttt{01234567632} & 4 & 8 & 3 & 1 & 11 & 11 & 4 & 4 & 79872 & 29000 & 24544\\
T11.0114 & \texttt{01234567631} & 4 & 8 & 3 & 1 & 11 & 11 & 4 & 4 & 77824 & 27184 & 24152\\
T11.0115 & \texttt{01234567623} & 3 & 8 & 2 & 1 & 11 & 11 & 3 & 3 & 79872 & 13044 & 30308\\
T11.0116 & \texttt{01234567622} & 4 & 8 & 2 & 2 & 10 & 9 & 3 & 2 & 77824 & 26328 & 24324\\
T11.0117 & \texttt{01234567621} & 4 & 8 & 3 & 1 & 11 & 11 & 4 & 4 & 75776 & 26176 & 23232\\
T11.0118 & \texttt{01234567612} & 3 & 8 & 2 & 1 & 11 & 11 & 3 & 3 & 75776 & 11192 & 28464\\
T11.0119 & \texttt{01234567611} & 4 & 8 & 2 & 2 & 10 & 9 & 3 & 2 & 73728 & 21232 & 23700\\
T11.0120 & \texttt{01234567567} & 2 & 8 & 1 & 2 & 8 & 5 & 1 & 0 & 94208 & 5208 & 37788\\
T11.0121 & \texttt{01234567566} & 3 & 8 & 2 & 2 & 10 & 9 & 2 & 1 & 92160 & 11472 & 37512\\
T11.0122 & \texttt{01234567565} & 3 & 8 & 1 & 1 & 11 & 11 & 3 & 3 & 90112 & 15624 & 35980\\
T11.0123 & \texttt{01234567564} & 3 & 8 & 2 & 1 & 11 & 11 & 3 & 3 & 88064 & 18140 & 33248\\
T11.0124 & \texttt{01234567563} & 3 & 8 & 2 & 1 & 11 & 11 & 3 & 3 & 86016 & 19696 & 31240\\
T11.0125 & \texttt{01234567562} & 3 & 8 & 2 & 1 & 11 & 11 & 3 & 3 & 83968 & 19604 & 30024\\
T11.0126 & \texttt{01234567561} & 3 & 8 & 2 & 1 & 11 & 11 & 3 & 3 & 81920 & 16280 & 30040\\
T11.0127 & \texttt{01234567555} & 4 & 8 & 1 & 6 & 9 & 8 & 2 & 1 & 88064 & 20832 & 34992\\
T11.0128 & \texttt{01234567554} & 4 & 8 & 2 & 2 & 10 & 9 & 3 & 2 & 86016 & 24720 & 31076\\
T11.0129 & \texttt{01234567553} & 4 & 8 & 2 & 2 & 10 & 9 & 3 & 2 & 83968 & 27200 & 28524\\
T11.0130 & \texttt{01234567552} & 4 & 8 & 2 & 2 & 10 & 9 & 3 & 2 & 81920 & 27696 & 27132\\
T11.0131 & \texttt{01234567551} & 4 & 8 & 2 & 2 & 10 & 9 & 3 & 2 & 79872 & 24736 & 27252\\
T11.0132 & \texttt{01234567545} & 3 & 8 & 2 & 1 & 11 & 11 & 3 & 3 & 86016 & 15316 & 33644\\
T11.0133 & \texttt{01234567544} & 4 & 8 & 2 & 2 & 10 & 9 & 3 & 2 & 83968 & 24976 & 29636\\
T11.0134 & \texttt{01234567543} & 4 & 8 & 3 & 1 & 11 & 11 & 4 & 4 & 81920 & 28512 & 26352\\
T11.0135 & \texttt{01234567542} & 4 & 8 & 3 & 1 & 11 & 11 & 4 & 4 & 79872 & 29680 & 24600\\
T11.0136 & \texttt{01234567541} & 4 & 8 & 3 & 1 & 11 & 11 & 4 & 4 & 77824 & 27136 & 24552\\
T11.0137 & \texttt{01234567534} & 3 & 8 & 2 & 1 & 11 & 11 & 3 & 3 & 81920 & 15328 & 31112\\
T11.0138 & \texttt{01234567533} & 4 & 8 & 2 & 2 & 10 & 9 & 3 & 2 & 79872 & 27648 & 25516\\
T11.0139 & \texttt{01234567532} & 4 & 8 & 3 & 1 & 11 & 11 & 4 & 4 & 77824 & 30160 & 23032\\
T11.0140 & \texttt{01234567531} & 4 & 8 & 3 & 1 & 11 & 11 & 4 & 4 & 75776 & 28448 & 22624\\
T11.0141 & \texttt{01234567523} & 3 & 8 & 2 & 1 & 11 & 11 & 3 & 3 & 77824 & 14908 & 28748\\
T11.0142 & \texttt{01234567522} & 4 & 8 & 2 & 2 & 10 & 9 & 3 & 2 & 75776 & 27760 & 22764\\
T11.0143 & \texttt{01234567521} & 4 & 8 & 3 & 1 & 11 & 11 & 4 & 4 & 73728 & 27712 & 21624\\
T11.0144 & \texttt{01234567512} & 3 & 8 & 2 & 1 & 11 & 11 & 3 & 3 & 73728 & 13064 & 26952\\
T11.0145 & \texttt{01234567511} & 4 & 8 & 2 & 2 & 10 & 9 & 3 & 2 & 71680 & 22944 & 22124\\
T11.0146 & \texttt{01234567456} & 2 & 8 & 1 & 1 & 11 & 11 & 2 & 2 & 88064 & 4620 & 35192\\
T11.0147 & \texttt{01234567455} & 3 & 8 & 2 & 2 & 10 & 9 & 2 & 1 & 86016 & 12008 & 33760\\
T11.0148 & \texttt{01234567454} & 3 & 8 & 1 & 1 & 11 & 11 & 3 & 3 & 83968 & 17220 & 31488\\
T11.0149 & \texttt{01234567453} & 3 & 8 & 2 & 1 & 11 & 11 & 3 & 3 & 81920 & 19304 & 29076\\
T11.0150 & \texttt{01234567452} & 3 & 8 & 2 & 1 & 11 & 11 & 3 & 3 & 79872 & 19548 & 27676\\
T11.0151 & \texttt{01234567451} & 3 & 8 & 2 & 1 & 11 & 11 & 3 & 3 & 77824 & 16432 & 27636\\
T11.0152 & \texttt{01234567444} & 4 & 8 & 1 & 6 & 9 & 8 & 2 & 1 & 81920 & 23520 & 29280\\
T11.0153 & \texttt{01234567443} & 4 & 8 & 2 & 2 & 10 & 9 & 3 & 2 & 79872 & 27584 & 25588\\
T11.0154 & \texttt{01234567442} & 4 & 8 & 2 & 2 & 10 & 9 & 3 & 2 & 77824 & 29088 & 23644\\
T11.0155 & \texttt{01234567441} & 4 & 8 & 2 & 2 & 10 & 9 & 3 & 2 & 75776 & 26752 & 23524\\
T11.0156 & \texttt{01234567434} & 3 & 8 & 2 & 1 & 11 & 11 & 3 & 3 & 79872 & 16992 & 29088\\
T11.0157 & \texttt{01234567433} & 4 & 8 & 2 & 2 & 10 & 9 & 3 & 2 & 77824 & 27520 & 24276\\
T11.0158 & \texttt{01234567432} & 4 & 8 & 3 & 1 & 11 & 11 & 4 & 4 & 75776 & 30368 & 21584\\
T11.0159 & \texttt{01234567431} & 4 & 8 & 3 & 1 & 11 & 11 & 4 & 4 & 73728 & 28864 & 21096\\
T11.0160 & \texttt{01234567423} & 3 & 8 & 2 & 1 & 11 & 11 & 3 & 3 & 75776 & 16524 & 26796\\
T11.0161 & \texttt{01234567422} & 4 & 8 & 2 & 2 & 10 & 9 & 3 & 2 & 73728 & 28512 & 21100\\
T11.0162 & \texttt{01234567421} & 4 & 8 & 3 & 1 & 11 & 11 & 4 & 4 & 71680 & 28672 & 19864\\
T11.0163 & \texttt{01234567412} & 3 & 8 & 2 & 1 & 11 & 11 & 3 & 3 & 71680 & 14696 & 25048\\
T11.0164 & \texttt{01234567411} & 4 & 8 & 2 & 2 & 10 & 9 & 3 & 2 & 69632 & 24256 & 20284\\
T11.0165 & \texttt{01234567345} & 2 & 8 & 1 & 1 & 11 & 11 & 2 & 2 & 81920 & 4256 & 32100\\
T11.0166 & \texttt{01234567344} & 3 & 8 & 2 & 2 & 10 & 9 & 2 & 1 & 79872 & 12352 & 29736\\
T11.0167 & \texttt{01234567343} & 3 & 8 & 1 & 1 & 11 & 11 & 3 & 3 & 77824 & 18144 & 27044\\
T11.0168 & \texttt{01234567342} & 3 & 8 & 2 & 1 & 11 & 11 & 3 & 3 & 75776 & 19060 & 25256\\
T11.0169 & \texttt{01234567341} & 3 & 8 & 2 & 1 & 11 & 11 & 3 & 3 & 73728 & 16360 & 25040\\
T11.0170 & \texttt{01234567333} & 4 & 8 & 1 & 6 & 9 & 8 & 2 & 1 & 75776 & 25088 & 23968\\
T11.0171 & \texttt{01234567332} & 4 & 8 & 2 & 2 & 10 & 9 & 3 & 2 & 73728 & 28608 & 20884\\
T11.0172 & \texttt{01234567331} & 4 & 8 & 2 & 2 & 10 & 9 & 3 & 2 & 71680 & 27520 & 20212\\
T11.0173 & \texttt{01234567323} & 3 & 8 & 2 & 1 & 11 & 11 & 3 & 3 & 73728 & 17580 & 24884\\
T11.0174 & \texttt{01234567322} & 4 & 8 & 2 & 2 & 10 & 9 & 3 & 2 & 71680 & 27840 & 19908\\
T11.0175 & \texttt{01234567321} & 4 & 8 & 3 & 1 & 11 & 11 & 4 & 4 & 69632 & 28416 & 18480\\
T11.0176 & \texttt{01234567312} & 3 & 8 & 2 & 1 & 11 & 11 & 3 & 3 & 69632 & 15784 & 23208\\
T11.0177 & \texttt{01234567311} & 4 & 8 & 2 & 2 & 10 & 9 & 3 & 2 & 67584 & 24704 & 18668\\
T11.0178 & \texttt{01234567234} & 2 & 8 & 1 & 1 & 11 & 11 & 2 & 2 & 75776 & 3780 & 28944\\
T11.0179 & \texttt{01234567233} & 3 & 8 & 2 & 2 & 10 & 9 & 2 & 1 & 73728 & 11864 & 25968\\
T11.0180 & \texttt{01234567232} & 3 & 8 & 1 & 1 & 11 & 11 & 3 & 3 & 71680 & 17388 & 23272\\
T11.0181 & \texttt{01234567231} & 3 & 8 & 2 & 1 & 11 & 11 & 3 & 3 & 69632 & 15520 & 22676\\
T11.0182 & \texttt{01234567222} & 4 & 8 & 1 & 6 & 9 & 8 & 2 & 1 & 69632 & 24192 & 19808\\
T11.0183 & \texttt{01234567221} & 4 & 8 & 2 & 2 & 10 & 9 & 3 & 2 & 67584 & 25600 & 17996\\
T11.0184 & \texttt{01234567212} & 3 & 8 & 2 & 1 & 11 & 11 & 3 & 3 & 67584 & 15656 & 21728\\
T11.0185 & \texttt{01234567211} & 4 & 8 & 2 & 2 & 10 & 9 & 3 & 2 & 65536 & 23296 & 17684\\
T11.0186 & \texttt{01234567123} & 2 & 8 & 1 & 1 & 11 & 11 & 2 & 2 & 69632 & 2680 & 26028\\
T11.0187 & \texttt{01234567122} & 3 & 8 & 2 & 2 & 10 & 9 & 2 & 1 & 67584 & 9296 & 22992\\
T11.0188 & \texttt{01234567121} & 3 & 8 & 1 & 1 & 11 & 11 & 3 & 3 & 65536 & 12840 & 20988\\
T11.0189 & \texttt{01234567111} & 4 & 8 & 1 & 6 & 9 & 8 & 2 & 1 & 63488 & 17920 & 17960\\
T11.0190 & \texttt{01234566666} & 5 & 7 & 1 & 120 & 7 & 6 & 1 & 0 & 92160 & 18900 & 42312\\
T11.0191 & \texttt{01234566665} & 5 & 7 & 2 & 24 & 8 & 7 & 2 & 1 & 90112 & 22544 & 37740\\
T11.0192 & \texttt{01234566664} & 5 & 7 & 2 & 24 & 8 & 7 & 2 & 1 & 88064 & 25452 & 34612\\
T11.0193 & \texttt{01234566663} & 5 & 7 & 2 & 24 & 8 & 7 & 2 & 1 & 86016 & 27272 & 32484\\
T11.0194 & \texttt{01234566662} & 5 & 7 & 2 & 24 & 8 & 7 & 2 & 1 & 83968 & 27348 & 31364\\
T11.0195 & \texttt{01234566661} & 5 & 7 & 2 & 24 & 8 & 7 & 2 & 1 & 81920 & 24128 & 31740\\
T11.0196 & \texttt{01234566656} & 4 & 7 & 2 & 6 & 9 & 8 & 2 & 1 & 90112 & 16512 & 38384\\
T11.0197 & \texttt{01234566655} & 5 & 7 & 2 & 12 & 8 & 6 & 2 & 0 & 88064 & 23208 & 35908\\
T11.0198 & \texttt{01234566654} & 5 & 7 & 3 & 6 & 9 & 8 & 3 & 2 & 86016 & 26900 & 31952\\
T11.0199 & \texttt{01234566653} & 5 & 7 & 3 & 6 & 9 & 8 & 3 & 2 & 83968 & 29248 & 29416\\
T11.0200 & \texttt{01234566652} & 5 & 7 & 3 & 6 & 9 & 8 & 3 & 2 & 81920 & 29660 & 28104\\
T11.0201 & \texttt{01234566651} & 5 & 7 & 3 & 6 & 9 & 8 & 3 & 2 & 79872 & 26648 & 28408\\
T11.0202 & \texttt{01234566645} & 4 & 7 & 2 & 6 & 9 & 8 & 2 & 1 & 86016 & 16392 & 35932\\
T11.0203 & \texttt{01234566644} & 5 & 7 & 2 & 12 & 8 & 6 & 2 & 0 & 83968 & 26892 & 30844\\
T11.0204 & \texttt{01234566643} & 5 & 7 & 3 & 6 & 9 & 8 & 3 & 2 & 81920 & 30296 & 27480\\
T11.0205 & \texttt{01234566642} & 5 & 7 & 3 & 6 & 9 & 8 & 3 & 2 & 79872 & 31380 & 25760\\
T11.0206 & \texttt{01234566641} & 5 & 7 & 3 & 6 & 9 & 8 & 3 & 2 & 77824 & 28784 & 25872\\
T11.0207 & \texttt{01234566634} & 4 & 7 & 2 & 6 & 9 & 8 & 2 & 1 & 81920 & 16464 & 33360\\
T11.0208 & \texttt{01234566633} & 5 & 7 & 2 & 12 & 8 & 6 & 2 & 0 & 79872 & 29232 & 26924\\
T11.0209 & \texttt{01234566632} & 5 & 7 & 3 & 6 & 9 & 8 & 3 & 2 & 77824 & 31660 & 24376\\
T11.0210 & \texttt{01234566631} & 5 & 7 & 3 & 6 & 9 & 8 & 3 & 2 & 75776 & 29896 & 24080\\
T11.0211 & \texttt{01234566623} & 4 & 7 & 2 & 6 & 9 & 8 & 2 & 1 & 77824 & 16056 & 31012\\
T11.0212 & \texttt{01234566622} & 5 & 7 & 2 & 12 & 8 & 6 & 2 & 0 & 75776 & 29124 & 24380\\
T11.0213 & \texttt{01234566621} & 5 & 7 & 3 & 6 & 9 & 8 & 3 & 2 & 73728 & 29024 & 23256\\
T11.0214 & \texttt{01234566612} & 4 & 7 & 2 & 6 & 9 & 8 & 2 & 1 & 73728 & 14208 & 29296\\
T11.0215 & \texttt{01234566611} & 5 & 7 & 2 & 12 & 8 & 6 & 2 & 0 & 71680 & 24168 & 24060\\
T11.0216 & \texttt{01234566566} & 4 & 7 & 3 & 8 & 7 & 5 & 1 & 0 & 90112 & 16248 & 37236\\
T11.0217 & \texttt{01234566565} & 4 & 7 & 2 & 2 & 10 & 9 & 3 & 2 & 88064 & 18912 & 36200\\
T11.0218 & \texttt{01234566564} & 4 & 7 & 3 & 2 & 10 & 9 & 3 & 2 & 86016 & 22016 & 32924\\
T11.0219 & \texttt{01234566563} & 4 & 7 & 3 & 2 & 10 & 9 & 3 & 2 & 83968 & 23968 & 30692\\
T11.0220 & \texttt{01234566562} & 4 & 7 & 3 & 2 & 10 & 9 & 3 & 2 & 81920 & 24128 & 29460\\
T11.0221 & \texttt{01234566561} & 4 & 7 & 3 & 2 & 10 & 9 & 3 & 2 & 79872 & 20960 & 29660\\
T11.0222 & \texttt{01234566555} & 5 & 7 & 2 & 12 & 8 & 6 & 2 & 0 & 86016 & 22632 & 35708\\
T11.0223 & \texttt{01234566554} & 5 & 7 & 3 & 4 & 9 & 7 & 3 & 1 & 83968 & 26716 & 31248\\
T11.0224 & \texttt{01234566553} & 5 & 7 & 3 & 4 & 9 & 7 & 3 & 1 & 81920 & 29328 & 28472\\
T11.0225 & \texttt{01234566552} & 5 & 7 & 3 & 4 & 9 & 7 & 3 & 1 & 79872 & 29908 & 27064\\
T11.0226 & \texttt{01234566551} & 5 & 7 & 3 & 4 & 9 & 7 & 3 & 1 & 77824 & 27000 & 27368\\
T11.0227 & \texttt{01234566545} & 4 & 7 & 3 & 2 & 10 & 9 & 3 & 2 & 83968 & 18416 & 33816\\
T11.0228 & \texttt{01234566544} & 5 & 7 & 3 & 4 & 9 & 7 & 3 & 1 & 81920 & 27236 & 29568\\
T11.0229 & \texttt{01234566543} & 5 & 7 & 4 & 2 & 10 & 9 & 4 & 3 & 79872 & 30904 & 25932\\
T11.0230 & \texttt{01234566542} & 5 & 7 & 4 & 2 & 10 & 9 & 4 & 3 & 77824 & 32156 & 24100\\
T11.0231 & \texttt{01234566541} & 5 & 7 & 4 & 2 & 10 & 9 & 4 & 3 & 75776 & 29664 & 24204\\
T11.0232 & \texttt{01234566534} & 4 & 7 & 3 & 2 & 10 & 9 & 3 & 2 & 79872 & 18368 & 31284\\
T11.0233 & \texttt{01234566533} & 5 & 7 & 3 & 4 & 9 & 7 & 3 & 1 & 77824 & 30240 & 25128\\
T11.0234 & \texttt{01234566532} & 5 & 7 & 4 & 2 & 10 & 9 & 4 & 3 & 75776 & 32836 & 22436\\
T11.0235 & \texttt{01234566531} & 5 & 7 & 4 & 2 & 10 & 9 & 4 & 3 & 73728 & 31176 & 22116\\
T11.0236 & \texttt{01234566523} & 4 & 7 & 3 & 2 & 10 & 9 & 3 & 2 & 75776 & 17936 & 28968\\
T11.0237 & \texttt{01234566522} & 5 & 7 & 3 & 4 & 9 & 7 & 3 & 1 & 73728 & 30572 & 22360\\
T11.0238 & \texttt{01234566521} & 5 & 7 & 4 & 2 & 10 & 9 & 4 & 3 & 71680 & 30576 & 21180\\
T11.0239 & \texttt{01234566512} & 4 & 7 & 3 & 2 & 10 & 9 & 3 & 2 & 71680 & 16096 & 27292\\
T11.0240 & \texttt{01234566511} & 5 & 7 & 3 & 4 & 9 & 7 & 3 & 1 & 69632 & 25896 & 22000\\
T11.0241 & \texttt{01234566456} & 3 & 7 & 2 & 2 & 10 & 9 & 2 & 1 & 86016 & 9672 & 35500\\
T11.0242 & \texttt{01234566455} & 4 & 7 & 3 & 4 & 9 & 7 & 2 & 0 & 83968 & 16220 & 33828\\
T11.0243 & \texttt{01234566454} & 4 & 7 & 2 & 2 & 10 & 9 & 3 & 2 & 81920 & 20592 & 31556\\
T11.0244 & \texttt{01234566453} & 4 & 7 & 3 & 2 & 10 & 9 & 3 & 2 & 79872 & 23072 & 28856\\
T11.0245 & \texttt{01234566452} & 4 & 7 & 3 & 2 & 10 & 9 & 3 & 2 & 77824 & 23568 & 27408\\
T11.0246 & \texttt{01234566451} & 4 & 7 & 3 & 2 & 10 & 9 & 3 & 2 & 75776 & 20608 & 27536\\
T11.0247 & \texttt{01234566444} & 5 & 7 & 2 & 12 & 8 & 6 & 2 & 0 & 79872 & 26052 & 29348\\
T11.0248 & \texttt{01234566443} & 5 & 7 & 3 & 4 & 9 & 7 & 3 & 1 & 77824 & 30248 & 25240\\
T11.0249 & \texttt{01234566442} & 5 & 7 & 3 & 4 & 9 & 7 & 3 & 1 & 75776 & 31836 & 23184\\
T11.0250 & \texttt{01234566441} & 5 & 7 & 3 & 4 & 9 & 7 & 3 & 1 & 73728 & 29552 & 23200\\
T11.0251 & \texttt{01234566434} & 4 & 7 & 3 & 2 & 10 & 9 & 3 & 2 & 77824 & 20304 & 29092\\
T11.0252 & \texttt{01234566433} & 5 & 7 & 3 & 4 & 9 & 7 & 3 & 1 & 75776 & 30384 & 23832\\
T11.0253 & \texttt{01234566432} & 5 & 7 & 4 & 2 & 10 & 9 & 4 & 3 & 73728 & 33316 & 20900\\
T11.0254 & \texttt{01234566431} & 5 & 7 & 4 & 2 & 10 & 9 & 4 & 3 & 71680 & 31864 & 20484\\
T11.0255 & \texttt{01234566423} & 4 & 7 & 3 & 2 & 10 & 9 & 3 & 2 & 73728 & 19824 & 26816\\
T11.0256 & \texttt{01234566422} & 5 & 7 & 3 & 4 & 9 & 7 & 3 & 1 & 71680 & 31596 & 20544\\
T11.0257 & \texttt{01234566421} & 5 & 7 & 4 & 2 & 10 & 9 & 4 & 3 & 69632 & 31808 & 19252\\
T11.0258 & \texttt{01234566412} & 4 & 7 & 3 & 2 & 10 & 9 & 3 & 2 & 69632 & 18000 & 25172\\
T11.0259 & \texttt{01234566411} & 5 & 7 & 3 & 4 & 9 & 7 & 3 & 1 & 67584 & 27480 & 19960\\
T11.0260 & \texttt{01234566345} & 3 & 7 & 2 & 2 & 10 & 9 & 2 & 1 & 79872 & 8736 & 32712\\
T11.0261 & \texttt{01234566344} & 4 & 7 & 3 & 4 & 9 & 7 & 2 & 0 & 77824 & 16384 & 30012\\
T11.0262 & \texttt{01234566343} & 4 & 7 & 2 & 2 & 10 & 9 & 3 & 2 & 75776 & 21728 & 27096\\
T11.0263 & \texttt{01234566342} & 4 & 7 & 3 & 2 & 10 & 9 & 3 & 2 & 73728 & 22896 & 25196\\
T11.0264 & \texttt{01234566341} & 4 & 7 & 3 & 2 & 10 & 9 & 3 & 2 & 71680 & 20352 & 25116\\
T11.0265 & \texttt{01234566333} & 5 & 7 & 2 & 12 & 8 & 6 & 2 & 0 & 73728 & 28224 & 23796\\
T11.0266 & \texttt{01234566332} & 5 & 7 & 3 & 4 & 9 & 7 & 3 & 1 & 71680 & 31828 & 20408\\
T11.0267 & \texttt{01234566331} & 5 & 7 & 3 & 4 & 9 & 7 & 3 & 1 & 69632 & 30792 & 19776\\
T11.0268 & \texttt{01234566323} & 4 & 7 & 3 & 2 & 10 & 9 & 3 & 2 & 71680 & 21152 & 24888\\
T11.0269 & \texttt{01234566322} & 5 & 7 & 3 & 4 & 9 & 7 & 3 & 1 & 69632 & 31196 & 19432\\
T11.0270 & \texttt{01234566321} & 5 & 7 & 4 & 2 & 10 & 9 & 4 & 3 & 67584 & 31824 & 17916\\
T11.0271 & \texttt{01234566312} & 4 & 7 & 3 & 2 & 10 & 9 & 3 & 2 & 67584 & 19360 & 23284\\
T11.0272 & \texttt{01234566311} & 5 & 7 & 3 & 4 & 9 & 7 & 3 & 1 & 65536 & 28200 & 18328\\
T11.0273 & \texttt{01234566234} & 3 & 7 & 2 & 2 & 10 & 9 & 2 & 1 & 73728 & 7992 & 29732\\
T11.0274 & \texttt{01234566233} & 4 & 7 & 3 & 4 & 9 & 7 & 2 & 0 & 71680 & 15860 & 26420\\
T11.0275 & \texttt{01234566232} & 4 & 7 & 2 & 2 & 10 & 9 & 3 & 2 & 69632 & 21168 & 23436\\
T11.0276 & \texttt{01234566231} & 4 & 7 & 3 & 2 & 10 & 9 & 3 & 2 & 67584 & 19456 & 22912\\
T11.0277 & \texttt{01234566222} & 5 & 7 & 2 & 12 & 8 & 6 & 2 & 0 & 67584 & 27756 & 19684\\
T11.0278 & \texttt{01234566221} & 5 & 7 & 3 & 4 & 9 & 7 & 3 & 1 & 65536 & 29216 & 17720\\
T11.0279 & \texttt{01234566212} & 4 & 7 & 3 & 2 & 10 & 9 & 3 & 2 & 65536 & 19440 & 21900\\
T11.0280 & \texttt{01234566211} & 5 & 7 & 3 & 4 & 9 & 7 & 3 & 1 & 63488 & 27000 & 17504\\
T11.0281 & \texttt{01234566123} & 3 & 7 & 2 & 2 & 10 & 9 & 2 & 1 & 67584 & 6768 & 26952\\
T11.0282 & \texttt{01234566122} & 4 & 7 & 3 & 4 & 9 & 7 & 2 & 0 & 65536 & 13304 & 23676\\
T11.0283 & \texttt{01234566121} & 4 & 7 & 2 & 2 & 10 & 9 & 3 & 2 & 63488 & 16768 & 21448\\
T11.0284 & \texttt{01234566111} & 5 & 7 & 2 & 12 & 8 & 6 & 2 & 0 & 61440 & 21768 & 18196\\
T11.0285 & \texttt{01234565656} & 3 & 7 & 1 & 6 & 7 & 5 & 1 & 0 & 88064 & 13392 & 38136\\
T11.0286 & \texttt{01234565655} & 4 & 7 & 1 & 4 & 8 & 5 & 2 & 0 & 86016 & 20088 & 36156\\
T11.0287 & \texttt{01234565654} & 4 & 7 & 2 & 2 & 9 & 7 & 3 & 2 & 83968 & 23780 & 32192\\
T11.0288 & \texttt{01234565653} & 4 & 7 & 2 & 2 & 9 & 7 & 3 & 2 & 81920 & 26128 & 29640\\
T11.0289 & \texttt{01234565652} & 4 & 7 & 2 & 2 & 9 & 7 & 3 & 2 & 79872 & 26540 & 28296\\
T11.0290 & \texttt{01234565651} & 4 & 7 & 2 & 2 & 9 & 7 & 3 & 2 & 77824 & 23528 & 28536\\
T11.0291 & \texttt{01234565645} & 3 & 7 & 2 & 2 & 9 & 7 & 2 & 1 & 83968 & 13272 & 35616\\
T11.0292 & \texttt{01234565644} & 4 & 7 & 2 & 4 & 8 & 5 & 2 & 0 & 81920 & 23772 & 30888\\
T11.0293 & \texttt{01234565643} & 4 & 7 & 3 & 2 & 9 & 7 & 3 & 2 & 79872 & 27176 & 27540\\
T11.0294 & \texttt{01234565642} & 4 & 7 & 3 & 2 & 9 & 7 & 3 & 2 & 77824 & 28260 & 25804\\
T11.0295 & \texttt{01234565641} & 4 & 7 & 3 & 2 & 9 & 7 & 3 & 2 & 75776 & 25664 & 25860\\
T11.0296 & \texttt{01234565634} & 3 & 7 & 2 & 2 & 9 & 7 & 2 & 1 & 79872 & 13344 & 33004\\
T11.0297 & \texttt{01234565633} & 4 & 7 & 2 & 4 & 8 & 5 & 2 & 0 & 77824 & 26112 & 26848\\
T11.0298 & \texttt{01234565632} & 4 & 7 & 3 & 2 & 9 & 7 & 3 & 2 & 75776 & 28540 & 24316\\
T11.0299 & \texttt{01234565631} & 4 & 7 & 3 & 2 & 9 & 7 & 3 & 2 & 73728 & 26776 & 23980\\
T11.0300 & \texttt{01234565623} & 3 & 7 & 2 & 2 & 9 & 7 & 2 & 1 & 75776 & 12936 & 30624\\
T11.0301 & \texttt{01234565622} & 4 & 7 & 2 & 4 & 8 & 5 & 2 & 0 & 73728 & 26004 & 24208\\
T11.0302 & \texttt{01234565621} & 4 & 7 & 3 & 2 & 9 & 7 & 3 & 2 & 71680 & 25904 & 23076\\
T11.0303 & \texttt{01234565612} & 3 & 7 & 2 & 2 & 9 & 7 & 2 & 1 & 71680 & 11088 & 28868\\
T11.0304 & \texttt{01234565611} & 4 & 7 & 2 & 4 & 8 & 5 & 2 & 0 & 69632 & 21048 & 23768\\
T11.0305 & \texttt{01234565555} & 5 & 7 & 1 & 24 & 8 & 7 & 2 & 1 & 83968 & 22320 & 36656\\
T11.0306 & \texttt{01234565554} & 5 & 7 & 2 & 6 & 9 & 8 & 3 & 2 & 81920 & 26600 & 31764\\
T11.0307 & \texttt{01234565553} & 5 & 7 & 2 & 6 & 9 & 8 & 3 & 2 & 79872 & 29344 & 28796\\
T11.0308 & \texttt{01234565552} & 5 & 7 & 2 & 6 & 9 & 8 & 3 & 2 & 77824 & 30008 & 27340\\
T11.0309 & \texttt{01234565551} & 5 & 7 & 2 & 6 & 9 & 8 & 3 & 2 & 75776 & 27152 & 27716\\
T11.0310 & \texttt{01234565545} & 4 & 7 & 2 & 2 & 10 & 9 & 3 & 2 & 81920 & 19404 & 33716\\
T11.0311 & \texttt{01234565544} & 5 & 7 & 2 & 4 & 9 & 7 & 3 & 1 & 79872 & 27384 & 29708\\
T11.0312 & \texttt{01234565543} & 5 & 7 & 3 & 2 & 10 & 9 & 4 & 3 & 77824 & 31184 & 25816\\
T11.0313 & \texttt{01234565542} & 5 & 7 & 3 & 2 & 10 & 9 & 4 & 3 & 75776 & 32520 & 23904\\
T11.0314 & \texttt{01234565541} & 5 & 7 & 3 & 2 & 10 & 9 & 4 & 3 & 73728 & 30080 & 24064\\
T11.0315 & \texttt{01234565534} & 4 & 7 & 2 & 2 & 10 & 9 & 3 & 2 & 77824 & 19296 & 31168\\
T11.0316 & \texttt{01234565533} & 5 & 7 & 2 & 4 & 9 & 7 & 3 & 1 & 75776 & 30720 & 24900\\
T11.0317 & \texttt{01234565532} & 5 & 7 & 3 & 2 & 10 & 9 & 4 & 3 & 73728 & 33400 & 22064\\
T11.0318 & \texttt{01234565531} & 5 & 7 & 3 & 2 & 10 & 9 & 4 & 3 & 71680 & 31792 & 21768\\
T11.0319 & \texttt{01234565523} & 4 & 7 & 2 & 2 & 10 & 9 & 3 & 2 & 73728 & 18852 & 28868\\
T11.0320 & \texttt{01234565522} & 5 & 7 & 2 & 4 & 9 & 7 & 3 & 1 & 71680 & 31272 & 22020\\
T11.0321 & \texttt{01234565521} & 5 & 7 & 3 & 2 & 10 & 9 & 4 & 3 & 69632 & 31328 & 20800\\
T11.0322 & \texttt{01234565512} & 4 & 7 & 2 & 2 & 10 & 9 & 3 & 2 & 69632 & 17016 & 27248\\
T11.0323 & \texttt{01234565511} & 5 & 7 & 2 & 4 & 9 & 7 & 3 & 1 & 67584 & 26736 & 21748\\
T11.0324 & \texttt{01234565456} & 3 & 7 & 2 & 1 & 11 & 11 & 3 & 3 & 83968 & 12612 & 34484\\
T11.0325 & \texttt{01234565455} & 4 & 7 & 3 & 2 & 10 & 9 & 3 & 2 & 81920 & 18320 & 32932\\
T11.0326 & \texttt{01234565454} & 4 & 7 & 2 & 1 & 11 & 11 & 4 & 4 & 79872 & 21852 & 31020\\
T11.0327 & \texttt{01234565453} & 4 & 7 & 3 & 1 & 11 & 11 & 4 & 4 & 77824 & 24728 & 27968\\
T11.0328 & \texttt{01234565452} & 4 & 7 & 3 & 1 & 11 & 11 & 4 & 4 & 75776 & 25476 & 26392\\
T11.0329 & \texttt{01234565451} & 4 & 7 & 3 & 1 & 11 & 11 & 4 & 4 & 73728 & 22672 & 26552\\
T11.0330 & \texttt{01234565444} & 5 & 7 & 2 & 6 & 9 & 8 & 3 & 2 & 77824 & 26472 & 29172\\
T11.0331 & \texttt{01234565443} & 5 & 7 & 3 & 2 & 10 & 9 & 4 & 3 & 75776 & 30800 & 24776\\
T11.0332 & \texttt{01234565442} & 5 & 7 & 3 & 2 & 10 & 9 & 4 & 3 & 73728 & 32472 & 22624\\
T11.0333 & \texttt{01234565441} & 5 & 7 & 3 & 2 & 10 & 9 & 4 & 3 & 71680 & 30240 & 22688\\
T11.0334 & \texttt{01234565434} & 4 & 7 & 3 & 1 & 11 & 11 & 4 & 4 & 75776 & 21504 & 28508\\
T11.0335 & \texttt{01234565433} & 5 & 7 & 3 & 2 & 10 & 9 & 4 & 3 & 73728 & 31136 & 23192\\
T11.0336 & \texttt{01234565432} & 5 & 7 & 4 & 1 & 11 & 11 & 5 & 5 & 71680 & 34152 & 20100\\
T11.0337 & \texttt{01234565431} & 5 & 7 & 4 & 1 & 11 & 11 & 5 & 5 & 69632 & 32752 & 19700\\
T11.0338 & \texttt{01234565423} & 4 & 7 & 3 & 1 & 11 & 11 & 4 & 4 & 71680 & 21012 & 26232\\
T11.0339 & \texttt{01234565422} & 5 & 7 & 3 & 2 & 10 & 9 & 4 & 3 & 69632 & 32568 & 19744\\
T11.0340 & \texttt{01234565421} & 5 & 7 & 4 & 1 & 11 & 11 & 5 & 5 & 67584 & 32832 & 18404\\
T11.0341 & \texttt{01234565412} & 4 & 7 & 3 & 1 & 11 & 11 & 4 & 4 & 67584 & 19192 & 24636\\
T11.0342 & \texttt{01234565411} & 5 & 7 & 3 & 2 & 10 & 9 & 4 & 3 & 65536 & 28592 & 19224\\
T11.0343 & \texttt{01234565345} & 3 & 7 & 2 & 1 & 11 & 11 & 3 & 3 & 77824 & 11104 & 31960\\
T11.0344 & \texttt{01234565344} & 4 & 7 & 3 & 2 & 10 & 9 & 3 & 2 & 75776 & 18304 & 29204\\
T11.0345 & \texttt{01234565343} & 4 & 7 & 2 & 1 & 11 & 11 & 4 & 4 & 73728 & 23200 & 26344\\
T11.0346 & \texttt{01234565342} & 4 & 7 & 3 & 1 & 11 & 11 & 4 & 4 & 71680 & 24620 & 24284\\
T11.0347 & \texttt{01234565341} & 4 & 7 & 3 & 1 & 11 & 11 & 4 & 4 & 69632 & 22232 & 24220\\
T11.0348 & \texttt{01234565333} & 5 & 7 & 2 & 6 & 9 & 8 & 3 & 2 & 71680 & 29248 & 23100\\
T11.0349 & \texttt{01234565332} & 5 & 7 & 3 & 2 & 10 & 9 & 4 & 3 & 69632 & 32936 & 19520\\
T11.0350 & \texttt{01234565331} & 5 & 7 & 3 & 2 & 10 & 9 & 4 & 3 & 67584 & 31952 & 18888\\
T11.0351 & \texttt{01234565323} & 4 & 7 & 3 & 1 & 11 & 11 & 4 & 4 & 69632 & 22612 & 24104\\
T11.0352 & \texttt{01234565322} & 5 & 7 & 3 & 2 & 10 & 9 & 4 & 3 & 67584 & 32440 & 18480\\
T11.0353 & \texttt{01234565321} & 5 & 7 & 4 & 1 & 11 & 11 & 5 & 5 & 65536 & 33120 & 16900\\
T11.0354 & \texttt{01234565312} & 4 & 7 & 3 & 1 & 11 & 11 & 4 & 4 & 65536 & 20824 & 22532\\
T11.0355 & \texttt{01234565311} & 5 & 7 & 3 & 2 & 10 & 9 & 4 & 3 & 63488 & 29584 & 17392\\
T11.0356 & \texttt{01234565234} & 3 & 7 & 2 & 1 & 11 & 11 & 3 & 3 & 71680 & 10092 & 29124\\
T11.0357 & \texttt{01234565233} & 4 & 7 & 3 & 2 & 10 & 9 & 3 & 2 & 69632 & 17744 & 25692\\
T11.0358 & \texttt{01234565232} & 4 & 7 & 2 & 1 & 11 & 11 & 4 & 4 & 67584 & 22836 & 22636\\
T11.0359 & \texttt{01234565231} & 4 & 7 & 3 & 1 & 11 & 11 & 4 & 4 & 65536 & 21280 & 22096\\
T11.0360 & \texttt{01234565222} & 5 & 7 & 2 & 6 & 9 & 8 & 3 & 2 & 65536 & 29208 & 18812\\
T11.0361 & \texttt{01234565221} & 5 & 7 & 3 & 2 & 10 & 9 & 4 & 3 & 63488 & 30720 & 16752\\
T11.0362 & \texttt{01234565212} & 4 & 7 & 3 & 1 & 11 & 11 & 4 & 4 & 63488 & 21112 & 21100\\
T11.0363 & \texttt{01234565211} & 5 & 7 & 3 & 2 & 10 & 9 & 4 & 3 & 61440 & 28592 & 16552\\
T11.0364 & \texttt{01234565123} & 3 & 7 & 2 & 1 & 11 & 11 & 3 & 3 & 65536 & 8744 & 26440\\
T11.0365 & \texttt{01234565122} & 4 & 7 & 3 & 2 & 10 & 9 & 3 & 2 & 63488 & 15200 & 23060\\
T11.0366 & \texttt{01234565121} & 4 & 7 & 2 & 1 & 11 & 11 & 4 & 4 & 61440 & 18584 & 20744\\
T11.0367 & \texttt{01234565111} & 5 & 7 & 2 & 6 & 9 & 8 & 3 & 2 & 59392 & 23504 & 17404\\
T11.0368 & \texttt{01234564564} & 3 & 7 & 1 & 2 & 8 & 5 & 2 & 1 & 81920 & 16380 & 31864\\
T11.0369 & \texttt{01234564563} & 3 & 7 & 2 & 2 & 8 & 5 & 2 & 1 & 79872 & 18200 & 29532\\
T11.0370 & \texttt{01234564562} & 3 & 7 & 2 & 2 & 8 & 5 & 2 & 1 & 77824 & 18276 & 28196\\
T11.0371 & \texttt{01234564561} & 3 & 7 & 2 & 2 & 8 & 5 & 2 & 1 & 75776 & 15056 & 28236\\
T11.0372 & \texttt{01234564555} & 4 & 7 & 2 & 6 & 9 & 8 & 2 & 1 & 81920 & 16608 & 33952\\
T11.0373 & \texttt{01234564554} & 4 & 7 & 2 & 2 & 10 & 9 & 3 & 2 & 79872 & 20408 & 30912\\
T11.0374 & \texttt{01234564553} & 4 & 7 & 3 & 2 & 10 & 9 & 3 & 2 & 77824 & 23152 & 28004\\
T11.0375 & \texttt{01234564552} & 4 & 7 & 3 & 2 & 10 & 9 & 3 & 2 & 75776 & 23816 & 26476\\
T11.0376 & \texttt{01234564551} & 4 & 7 & 3 & 2 & 10 & 9 & 3 & 2 & 73728 & 20960 & 26612\\
T11.0377 & \texttt{01234564545} & 3 & 7 & 1 & 2 & 9 & 7 & 2 & 1 & 79872 & 13440 & 33848\\
T11.0378 & \texttt{01234564544} & 4 & 7 & 1 & 2 & 10 & 9 & 3 & 2 & 77824 & 22260 & 30080\\
T11.0379 & \texttt{01234564543} & 4 & 7 & 2 & 1 & 11 & 11 & 4 & 4 & 75776 & 25928 & 26372\\
T11.0380 & \texttt{01234564542} & 4 & 7 & 2 & 1 & 11 & 11 & 4 & 4 & 73728 & 27180 & 24492\\
T11.0381 & \texttt{01234564541} & 4 & 7 & 2 & 1 & 11 & 11 & 4 & 4 & 71680 & 24688 & 24548\\
T11.0382 & \texttt{01234564534} & 3 & 7 & 2 & 1 & 11 & 11 & 3 & 3 & 75776 & 13392 & 31216\\
T11.0383 & \texttt{01234564533} & 4 & 7 & 2 & 2 & 10 & 9 & 3 & 2 & 73728 & 25264 & 25340\\
T11.0384 & \texttt{01234564532} & 4 & 7 & 3 & 1 & 11 & 11 & 4 & 4 & 71680 & 27860 & 22616\\
T11.0385 & \texttt{01234564531} & 4 & 7 & 3 & 1 & 11 & 11 & 4 & 4 & 69632 & 26200 & 22256\\
T11.0386 & \texttt{01234564523} & 3 & 7 & 2 & 1 & 11 & 11 & 3 & 3 & 71680 & 12960 & 28844\\
T11.0387 & \texttt{01234564522} & 4 & 7 & 2 & 2 & 10 & 9 & 3 & 2 & 69632 & 25596 & 22404\\
T11.0388 & \texttt{01234564521} & 4 & 7 & 3 & 1 & 11 & 11 & 4 & 4 & 67584 & 25600 & 21200\\
T11.0389 & \texttt{01234564512} & 3 & 7 & 2 & 1 & 11 & 11 & 3 & 3 & 67584 & 11120 & 27128\\
T11.0390 & \texttt{01234564511} & 4 & 7 & 2 & 2 & 10 & 9 & 3 & 2 & 65536 & 20920 & 21924\\
T11.0391 & \texttt{01234564444} & 5 & 7 & 1 & 24 & 8 & 7 & 2 & 1 & 75776 & 25200 & 29312\\
T11.0392 & \texttt{01234564443} & 5 & 7 & 2 & 6 & 9 & 8 & 3 & 2 & 73728 & 29792 & 24580\\
T11.0393 & \texttt{01234564442} & 5 & 7 & 2 & 6 & 9 & 8 & 3 & 2 & 71680 & 31632 & 22284\\
T11.0394 & \texttt{01234564441} & 5 & 7 & 2 & 6 & 9 & 8 & 3 & 2 & 69632 & 29504 & 22324\\
T11.0395 & \texttt{01234564434} & 4 & 7 & 2 & 2 & 10 & 9 & 3 & 2 & 73728 & 21792 & 27544\\
T11.0396 & \texttt{01234564433} & 5 & 7 & 2 & 4 & 9 & 7 & 3 & 1 & 71680 & 30528 & 22620\\
T11.0397 & \texttt{01234564432} & 5 & 7 & 3 & 2 & 10 & 9 & 4 & 3 & 69632 & 33712 & 19352\\
T11.0398 & \texttt{01234564431} & 5 & 7 & 3 & 2 & 10 & 9 & 4 & 3 & 67584 & 32416 & 18912\\
T11.0399 & \texttt{01234564423} & 4 & 7 & 2 & 2 & 10 & 9 & 3 & 2 & 69632 & 21276 & 25268\\
T11.0400 & \texttt{01234564422} & 5 & 7 & 2 & 4 & 9 & 7 & 3 & 1 & 67584 & 32400 & 18852\\
T11.0401 & \texttt{01234564421} & 5 & 7 & 3 & 2 & 10 & 9 & 4 & 3 & 65536 & 32768 & 17440\\
T11.0402 & \texttt{01234564412} & 4 & 7 & 2 & 2 & 10 & 9 & 3 & 2 & 65536 & 19464 & 23696\\
T11.0403 & \texttt{01234564411} & 5 & 7 & 2 & 4 & 9 & 7 & 3 & 1 & 63488 & 28704 & 18244\\
T11.0404 & \texttt{01234564345} & 3 & 7 & 2 & 1 & 11 & 11 & 3 & 3 & 75776 & 13728 & 29856\\
T11.0405 & \texttt{01234564344} & 4 & 7 & 3 & 2 & 10 & 9 & 3 & 2 & 73728 & 20032 & 27324\\
T11.0406 & \texttt{01234564343} & 4 & 7 & 2 & 1 & 11 & 11 & 4 & 4 & 71680 & 24032 & 24912\\
T11.0407 & \texttt{01234564342} & 4 & 7 & 3 & 1 & 11 & 11 & 4 & 4 & 69632 & 25956 & 22532\\
T11.0408 & \texttt{01234564341} & 4 & 7 & 3 & 1 & 11 & 11 & 4 & 4 & 67584 & 23880 & 22356\\
T11.0409 & \texttt{01234564333} & 5 & 7 & 2 & 6 & 9 & 8 & 3 & 2 & 69632 & 29184 & 22116\\
T11.0410 & \texttt{01234564332} & 5 & 7 & 3 & 2 & 10 & 9 & 4 & 3 & 67584 & 33040 & 18344\\
T11.0411 & \texttt{01234564331} & 5 & 7 & 3 & 2 & 10 & 9 & 4 & 3 & 65536 & 32160 & 17664\\
T11.0412 & \texttt{01234564323} & 4 & 7 & 3 & 1 & 11 & 11 & 4 & 4 & 67584 & 23420 & 22656\\
T11.0413 & \texttt{01234564322} & 5 & 7 & 3 & 2 & 10 & 9 & 4 & 3 & 65536 & 32816 & 17128\\
T11.0414 & \texttt{01234564321} & 5 & 7 & 4 & 1 & 11 & 11 & 5 & 5 & 63488 & 33600 & 15468\\
T11.0415 & \texttt{01234564312} & 4 & 7 & 3 & 1 & 11 & 11 & 4 & 4 & 63488 & 21640 & 21100\\
T11.0416 & \texttt{01234564311} & 5 & 7 & 3 & 2 & 10 & 9 & 4 & 3 & 61440 & 30240 & 15928\\
T11.0417 & \texttt{01234564234} & 3 & 7 & 2 & 1 & 11 & 11 & 3 & 3 & 69632 & 12180 & 27284\\
T11.0418 & \texttt{01234564233} & 4 & 7 & 3 & 2 & 10 & 9 & 3 & 2 & 67584 & 19400 & 23900\\
T11.0419 & \texttt{01234564232} & 4 & 7 & 2 & 1 & 11 & 11 & 4 & 4 & 65536 & 24060 & 20988\\
T11.0420 & \texttt{01234564231} & 4 & 7 & 3 & 1 & 11 & 11 & 4 & 4 & 63488 & 22816 & 20320\\
T11.0421 & \texttt{01234564222} & 5 & 7 & 2 & 6 & 9 & 8 & 3 & 2 & 63488 & 30000 & 17308\\
T11.0422 & \texttt{01234564221} & 5 & 7 & 3 & 2 & 10 & 9 & 4 & 3 & 61440 & 31616 & 15152\\
T11.0423 & \texttt{01234564212} & 4 & 7 & 3 & 1 & 11 & 11 & 4 & 4 & 61440 & 22344 & 19452\\
T11.0424 & \texttt{01234564211} & 5 & 7 & 3 & 2 & 10 & 9 & 4 & 3 & 59392 & 29664 & 14888\\
T11.0425 & \texttt{01234564123} & 3 & 7 & 2 & 1 & 11 & 11 & 3 & 3 & 63488 & 10584 & 24744\\
T11.0426 & \texttt{01234564122} & 4 & 7 & 3 & 2 & 10 & 9 & 3 & 2 & 61440 & 16880 & 21348\\
T11.0427 & \texttt{01234564121} & 4 & 7 & 2 & 1 & 11 & 11 & 4 & 4 & 59392 & 20104 & 19048\\
T11.0428 & \texttt{01234564111} & 5 & 7 & 2 & 6 & 9 & 8 & 3 & 2 & 57344 & 24864 & 15724\\
T11.0429 & \texttt{01234563456} & 2 & 7 & 1 & 2 & 7 & 3 & 1 & 0 & 79872 & 3360 & 32308\\
T11.0430 & \texttt{01234563455} & 3 & 7 & 2 & 2 & 10 & 9 & 2 & 1 & 77824 & 8896 & 31652\\
T11.0431 & \texttt{01234563454} & 3 & 7 & 2 & 1 & 11 & 11 & 3 & 3 & 75776 & 12256 & 30120\\
T11.0432 & \texttt{01234563453} & 3 & 7 & 1 & 1 & 11 & 11 & 3 & 3 & 73728 & 16800 & 27140\\
T11.0433 & \texttt{01234563452} & 3 & 7 & 2 & 1 & 11 & 11 & 3 & 3 & 71680 & 17212 & 25528\\
T11.0434 & \texttt{01234563451} & 3 & 7 & 2 & 1 & 11 & 11 & 3 & 3 & 69632 & 14200 & 25472\\
T11.0435 & \texttt{01234563444} & 4 & 7 & 2 & 6 & 9 & 8 & 2 & 1 & 73728 & 16704 & 28624\\
T11.0436 & \texttt{01234563443} & 4 & 7 & 2 & 2 & 10 & 9 & 3 & 2 & 71680 & 21312 & 25224\\
T11.0437 & \texttt{01234563442} & 4 & 7 & 3 & 2 & 10 & 9 & 3 & 2 & 69632 & 22984 & 23052\\
T11.0438 & \texttt{01234563441} & 4 & 7 & 3 & 2 & 10 & 9 & 3 & 2 & 67584 & 20752 & 22908\\
T11.0439 & \texttt{01234563434} & 3 & 7 & 1 & 2 & 9 & 7 & 2 & 1 & 71680 & 13440 & 28744\\
T11.0440 & \texttt{01234563433} & 4 & 7 & 1 & 2 & 10 & 9 & 3 & 2 & 69632 & 23520 & 23876\\
T11.0441 & \texttt{01234563432} & 4 & 7 & 2 & 1 & 11 & 11 & 4 & 4 & 67584 & 26452 & 20856\\
T11.0442 & \texttt{01234563431} & 4 & 7 & 2 & 1 & 11 & 11 & 4 & 4 & 65536 & 25000 & 20384\\
T11.0443 & \texttt{01234563423} & 3 & 7 & 2 & 1 & 11 & 11 & 3 & 3 & 67584 & 12960 & 26368\\
T11.0444 & \texttt{01234563422} & 4 & 7 & 2 & 2 & 10 & 9 & 3 & 2 & 65536 & 24732 & 20288\\
T11.0445 & \texttt{01234563421} & 4 & 7 & 3 & 1 & 11 & 11 & 4 & 4 & 63488 & 24944 & 18948\\
T11.0446 & \texttt{01234563412} & 3 & 7 & 2 & 1 & 11 & 11 & 3 & 3 & 63488 & 11136 & 24668\\
T11.0447 & \texttt{01234563411} & 4 & 7 & 2 & 2 & 10 & 9 & 3 & 2 & 61440 & 20616 & 19536\\
T11.0448 & \texttt{01234563333} & 5 & 7 & 1 & 24 & 8 & 7 & 2 & 1 & 67584 & 26880 & 22144\\
T11.0449 & \texttt{01234563332} & 5 & 7 & 2 & 6 & 9 & 8 & 3 & 2 & 65536 & 31072 & 18084\\
T11.0450 & \texttt{01234563331} & 5 & 7 & 2 & 6 & 9 & 8 & 3 & 2 & 63488 & 30400 & 17284\\
T11.0451 & \texttt{01234563323} & 4 & 7 & 2 & 2 & 10 & 9 & 3 & 2 & 65536 & 22860 & 21612\\
T11.0452 & \texttt{01234563322} & 5 & 7 & 2 & 4 & 9 & 7 & 3 & 1 & 63488 & 31392 & 16492\\
T11.0453 & \texttt{01234563321} & 5 & 7 & 3 & 2 & 10 & 9 & 4 & 3 & 61440 & 32384 & 14696\\
T11.0454 & \texttt{01234563312} & 4 & 7 & 2 & 2 & 10 & 9 & 3 & 2 & 61440 & 21096 & 20064\\
T11.0455 & \texttt{01234563311} & 5 & 7 & 2 & 4 & 9 & 7 & 3 & 1 & 59392 & 29376 & 14996\\
T11.0456 & \texttt{01234563234} & 3 & 7 & 2 & 1 & 11 & 11 & 3 & 3 & 67584 & 14180 & 25084\\
T11.0457 & \texttt{01234563233} & 4 & 7 & 3 & 2 & 10 & 9 & 3 & 2 & 65536 & 20536 & 21940\\
T11.0458 & \texttt{01234563232} & 4 & 7 & 2 & 1 & 11 & 11 & 4 & 4 & 63488 & 24332 & 19460\\
T11.0459 & \texttt{01234563231} & 4 & 7 & 3 & 1 & 11 & 11 & 4 & 4 & 61440 & 23712 & 18488\\
T11.0460 & \texttt{01234563222} & 5 & 7 & 2 & 6 & 9 & 8 & 3 & 2 & 61440 & 29408 & 16212\\
T11.0461 & \texttt{01234563221} & 5 & 7 & 3 & 2 & 10 & 9 & 4 & 3 & 59392 & 31232 & 13912\\
T11.0462 & \texttt{01234563212} & 4 & 7 & 3 & 1 & 11 & 11 & 4 & 4 & 59392 & 22632 & 17924\\
T11.0463 & \texttt{01234563211} & 5 & 7 & 3 & 2 & 10 & 9 & 4 & 3 & 57344 & 29632 & 13472\\
T11.0464 & \texttt{01234563123} & 3 & 7 & 2 & 1 & 11 & 11 & 3 & 3 & 61440 & 12088 & 22808\\
T11.0465 & \texttt{01234563122} & 4 & 7 & 3 & 2 & 10 & 9 & 3 & 2 & 59392 & 18064 & 19476\\
T11.0466 & \texttt{01234563121} & 4 & 7 & 2 & 1 & 11 & 11 & 4 & 4 & 57344 & 20968 & 17304\\
T11.0467 & \texttt{01234563111} & 5 & 7 & 2 & 6 & 9 & 8 & 3 & 2 & 55296 & 25408 & 14108\\
T11.0468 & \texttt{01234562345} & 2 & 7 & 1 & 1 & 11 & 11 & 2 & 2 & 71680 & 2484 & 28976\\
T11.0469 & \texttt{01234562344} & 3 & 7 & 2 & 2 & 10 & 9 & 2 & 1 & 69632 & 8136 & 27500\\
T11.0470 & \texttt{01234562343} & 3 & 7 & 2 & 1 & 11 & 11 & 3 & 3 & 67584 & 11484 & 25676\\
T11.0471 & \texttt{01234562342} & 3 & 7 & 1 & 1 & 11 & 11 & 3 & 3 & 65536 & 15876 & 22832\\
T11.0472 & \texttt{01234562341} & 3 & 7 & 2 & 1 & 11 & 11 & 3 & 3 & 63488 & 13280 & 22524\\
T11.0473 & \texttt{01234562333} & 4 & 7 & 2 & 6 & 9 & 8 & 2 & 1 & 65536 & 15984 & 23456\\
T11.0474 & \texttt{01234562332} & 4 & 7 & 2 & 2 & 10 & 9 & 3 & 2 & 63488 & 20504 & 20240\\
T11.0475 & \texttt{01234562331} & 4 & 7 & 3 & 2 & 10 & 9 & 3 & 2 & 61440 & 19520 & 19468\\
T11.0476 & \texttt{01234562323} & 3 & 7 & 1 & 2 & 9 & 7 & 2 & 1 & 63488 & 12528 & 23864\\
T11.0477 & \texttt{01234562322} & 4 & 7 & 1 & 2 & 10 & 9 & 3 & 2 & 61440 & 22572 & 18696\\
T11.0478 & \texttt{01234562321} & 4 & 7 & 2 & 1 & 11 & 11 & 4 & 4 & 59392 & 23200 & 17092\\
T11.0479 & \texttt{01234562312} & 3 & 7 & 2 & 1 & 11 & 11 & 3 & 3 & 59392 & 10736 & 22168\\
T11.0480 & \texttt{01234562311} & 4 & 7 & 2 & 2 & 10 & 9 & 3 & 2 & 57344 & 19576 & 17316\\
T11.0481 & \texttt{01234562222} & 5 & 7 & 1 & 24 & 8 & 7 & 2 & 1 & 59392 & 25920 & 16384\\
T11.0482 & \texttt{01234562221} & 5 & 7 & 2 & 6 & 9 & 8 & 3 & 2 & 57344 & 28160 & 13820\\
T11.0483 & \texttt{01234562212} & 4 & 7 & 2 & 2 & 10 & 9 & 3 & 2 & 57344 & 20904 & 17160\\
T11.0484 & \texttt{01234562211} & 5 & 7 & 2 & 4 & 9 & 7 & 3 & 1 & 55296 & 27264 & 13004\\
T11.0485 & \texttt{01234562123} & 3 & 7 & 2 & 1 & 11 & 11 & 3 & 3 & 59392 & 12792 & 21072\\
T11.0486 & \texttt{01234562122} & 4 & 7 & 3 & 2 & 10 & 9 & 3 & 2 & 57344 & 18128 & 17916\\
T11.0487 & \texttt{01234562121} & 4 & 7 & 2 & 1 & 11 & 11 & 4 & 4 & 55296 & 20392 & 16048\\
T11.0488 & \texttt{01234562111} & 5 & 7 & 2 & 6 & 9 & 8 & 3 & 2 & 53248 & 24192 & 13156\\
T11.0489 & \texttt{01234561234} & 2 & 7 & 1 & 1 & 11 & 11 & 2 & 2 & 63488 & 1560 & 25628\\
T11.0490 & \texttt{01234561233} & 3 & 7 & 2 & 2 & 10 & 9 & 2 & 1 & 61440 & 6640 & 23604\\
T11.0491 & \texttt{01234561232} & 3 & 7 & 2 & 1 & 11 & 11 & 3 & 3 & 59392 & 9160 & 21848\\
T11.0492 & \texttt{01234561231} & 3 & 7 & 1 & 1 & 11 & 11 & 3 & 3 & 57344 & 11640 & 19884\\
T11.0493 & \texttt{01234561222} & 4 & 7 & 2 & 6 & 9 & 8 & 2 & 1 & 57344 & 12960 & 19320\\
T11.0494 & \texttt{01234561221} & 4 & 7 & 2 & 2 & 10 & 9 & 3 & 2 & 55296 & 15696 & 17104\\
T11.0495 & \texttt{01234561212} & 3 & 7 & 1 & 2 & 9 & 7 & 2 & 1 & 55296 & 9120 & 20096\\
T11.0496 & \texttt{01234561211} & 4 & 7 & 1 & 2 & 10 & 9 & 3 & 2 & 53248 & 16680 & 15868\\
T11.0497 & \texttt{01234561111} & 5 & 7 & 1 & 24 & 8 & 7 & 2 & 1 & 51200 & 19200 & 13608\\
T11.0498 & \texttt{01234555555} & 6 & 6 & 1 & 720 & 6 & 5 & 1 & 0 & 81920 & 23064 & 38148\\
T11.0499 & \texttt{01234555554} & 6 & 6 & 2 & 120 & 7 & 6 & 2 & 1 & 79872 & 27540 & 32584\\
T11.0500 & \texttt{01234555553} & 6 & 6 & 2 & 120 & 7 & 6 & 2 & 1 & 77824 & 30416 & 29328\\
T11.0501 & \texttt{01234555552} & 6 & 6 & 2 & 120 & 7 & 6 & 2 & 1 & 75776 & 31164 & 27824\\
T11.0502 & \texttt{01234555551} & 6 & 6 & 2 & 120 & 7 & 6 & 2 & 1 & 73728 & 28360 & 28368\\
T11.0503 & \texttt{01234555545} & 5 & 6 & 2 & 24 & 8 & 7 & 2 & 1 & 79872 & 21448 & 34040\\
T11.0504 & \texttt{01234555544} & 6 & 6 & 2 & 48 & 7 & 5 & 2 & 0 & 77824 & 28588 & 30032\\
T11.0505 & \texttt{01234555543} & 6 & 6 & 3 & 24 & 8 & 7 & 3 & 2 & 75776 & 32520 & 25724\\
T11.0506 & \texttt{01234555542} & 6 & 6 & 3 & 24 & 8 & 7 & 3 & 2 & 73728 & 33940 & 23700\\
T11.0507 & \texttt{01234555541} & 6 & 6 & 3 & 24 & 8 & 7 & 3 & 2 & 71680 & 31552 & 23996\\
T11.0508 & \texttt{01234555534} & 5 & 6 & 2 & 24 & 8 & 7 & 2 & 1 & 75776 & 21280 & 31428\\
T11.0509 & \texttt{01234555533} & 6 & 6 & 2 & 48 & 7 & 5 & 2 & 0 & 73728 & 32256 & 24712\\
T11.0510 & \texttt{01234555532} & 6 & 6 & 3 & 24 & 8 & 7 & 3 & 2 & 71680 & 35020 & 21636\\
T11.0511 & \texttt{01234555531} & 6 & 6 & 3 & 24 & 8 & 7 & 3 & 2 & 69632 & 33464 & 21412\\
T11.0512 & \texttt{01234555523} & 5 & 6 & 2 & 24 & 8 & 7 & 2 & 1 & 71680 & 20824 & 29144\\
T11.0513 & \texttt{01234555522} & 6 & 6 & 2 & 48 & 7 & 5 & 2 & 0 & 69632 & 33028 & 21720\\
T11.0514 & \texttt{01234555521} & 6 & 6 & 3 & 24 & 8 & 7 & 3 & 2 & 67584 & 33136 & 20444\\
T11.0515 & \texttt{01234555512} & 5 & 6 & 2 & 24 & 8 & 7 & 2 & 1 & 67584 & 18992 & 27628\\
T11.0516 & \texttt{01234555511} & 6 & 6 & 2 & 48 & 7 & 5 & 2 & 0 & 65536 & 28632 & 21680\\
T11.0517 & \texttt{01234555455} & 5 & 6 & 3 & 12 & 8 & 6 & 2 & 0 & 79872 & 21476 & 32400\\
T11.0518 & \texttt{01234555454} & 5 & 6 & 2 & 6 & 9 & 8 & 3 & 2 & 77824 & 24168 & 30728\\
T11.0519 & \texttt{01234555453} & 5 & 6 & 3 & 6 & 9 & 8 & 3 & 2 & 75776 & 27440 & 27196\\
T11.0520 & \texttt{01234555452} & 5 & 6 & 3 & 6 & 9 & 8 & 3 & 2 & 73728 & 28440 & 25476\\
T11.0521 & \texttt{01234555451} & 5 & 6 & 3 & 6 & 9 & 8 & 3 & 2 & 71680 & 25792 & 25756\\
T11.0522 & \texttt{01234555444} & 6 & 6 & 2 & 36 & 7 & 5 & 2 & 0 & 75776 & 27948 & 29120\\
T11.0523 & \texttt{01234555443} & 6 & 6 & 3 & 12 & 8 & 6 & 3 & 1 & 73728 & 32408 & 24244\\
T11.0524 & \texttt{01234555442} & 6 & 6 & 3 & 12 & 8 & 6 & 3 & 1 & 71680 & 34164 & 21948\\
T11.0525 & \texttt{01234555441} & 6 & 6 & 3 & 12 & 8 & 6 & 3 & 1 & 69632 & 31984 & 22132\\
T11.0526 & \texttt{01234555434} & 5 & 6 & 3 & 6 & 9 & 8 & 3 & 2 & 73728 & 23760 & 28088\\
T11.0527 & \texttt{01234555433} & 6 & 6 & 3 & 12 & 8 & 6 & 3 & 1 & 71680 & 32944 & 22436\\
T11.0528 & \texttt{01234555432} & 6 & 6 & 4 & 6 & 9 & 8 & 4 & 3 & 69632 & 36044 & 19072\\
T11.0529 & \texttt{01234555431} & 6 & 6 & 4 & 6 & 9 & 8 & 4 & 3 & 67584 & 34696 & 18728\\
T11.0530 & \texttt{01234555423} & 5 & 6 & 3 & 6 & 9 & 8 & 3 & 2 & 69632 & 23256 & 25796\\
T11.0531 & \texttt{01234555422} & 6 & 6 & 3 & 12 & 8 & 6 & 3 & 1 & 67584 & 34596 & 18780\\
T11.0532 & \texttt{01234555421} & 6 & 6 & 4 & 6 & 9 & 8 & 4 & 3 & 65536 & 34912 & 17368\\
T11.0533 & \texttt{01234555412} & 5 & 6 & 3 & 6 & 9 & 8 & 3 & 2 & 65536 & 21440 & 24288\\
T11.0534 & \texttt{01234555411} & 6 & 6 & 3 & 12 & 8 & 6 & 3 & 1 & 63488 & 30760 & 18444\\
T11.0535 & \texttt{01234555345} & 4 & 6 & 2 & 6 & 9 & 8 & 2 & 1 & 75776 & 14528 & 31668\\
T11.0536 & \texttt{01234555344} & 5 & 6 & 3 & 12 & 8 & 6 & 2 & 0 & 73728 & 21280 & 28688\\
T11.0537 & \texttt{01234555343} & 5 & 6 & 2 & 6 & 9 & 8 & 3 & 2 & 71680 & 25728 & 25716\\
T11.0538 & \texttt{01234555342} & 5 & 6 & 3 & 6 & 9 & 8 & 3 & 2 & 69632 & 27400 & 23448\\
T11.0539 & \texttt{01234555341} & 5 & 6 & 3 & 6 & 9 & 8 & 3 & 2 & 67584 & 25168 & 23472\\
T11.0540 & \texttt{01234555333} & 6 & 6 & 2 & 36 & 7 & 5 & 2 & 0 & 69632 & 31328 & 22360\\
T11.0541 & \texttt{01234555332} & 6 & 6 & 3 & 12 & 8 & 6 & 3 & 1 & 67584 & 35100 & 18444\\
T11.0542 & \texttt{01234555331} & 6 & 6 & 3 & 12 & 8 & 6 & 3 & 1 & 65536 & 34168 & 17836\\
T11.0543 & \texttt{01234555323} & 5 & 6 & 3 & 6 & 9 & 8 & 3 & 2 & 67584 & 25128 & 23396\\
T11.0544 & \texttt{01234555322} & 6 & 6 & 3 & 12 & 8 & 6 & 3 & 1 & 65536 & 34740 & 17340\\
T11.0545 & \texttt{01234555321} & 6 & 6 & 4 & 6 & 9 & 8 & 4 & 3 & 63488 & 35472 & 15656\\
T11.0546 & \texttt{01234555312} & 5 & 6 & 3 & 6 & 9 & 8 & 3 & 2 & 63488 & 23344 & 21880\\
T11.0547 & \texttt{01234555311} & 6 & 6 & 3 & 12 & 8 & 6 & 3 & 1 & 61440 & 32024 & 16340\\
T11.0548 & \texttt{01234555234} & 4 & 6 & 2 & 6 & 9 & 8 & 2 & 1 & 69632 & 13248 & 28976\\
T11.0549 & \texttt{01234555233} & 5 & 6 & 3 & 12 & 8 & 6 & 2 & 0 & 67584 & 20684 & 25256\\
T11.0550 & \texttt{01234555232} & 5 & 6 & 2 & 6 & 9 & 8 & 3 & 2 & 65536 & 25560 & 21960\\
T11.0551 & \texttt{01234555231} & 5 & 6 & 3 & 6 & 9 & 8 & 3 & 2 & 63488 & 24160 & 21444\\
T11.0552 & \texttt{01234555222} & 6 & 6 & 2 & 36 & 7 & 5 & 2 & 0 & 63488 & 31716 & 17896\\
T11.0553 & \texttt{01234555221} & 6 & 6 & 3 & 12 & 8 & 6 & 3 & 1 & 61440 & 33280 & 15668\\
T11.0554 & \texttt{01234555212} & 5 & 6 & 3 & 6 & 9 & 8 & 3 & 2 & 61440 & 23840 & 20416\\
T11.0555 & \texttt{01234555211} & 6 & 6 & 3 & 12 & 8 & 6 & 3 & 1 & 59392 & 31240 & 15532\\
T11.0556 & \texttt{01234555123} & 4 & 6 & 2 & 6 & 9 & 8 & 2 & 1 & 63488 & 11776 & 26412\\
T11.0557 & \texttt{01234555122} & 5 & 6 & 3 & 12 & 8 & 6 & 2 & 0 & 61440 & 18152 & 22808\\
T11.0558 & \texttt{01234555121} & 5 & 6 & 2 & 6 & 9 & 8 & 3 & 2 & 59392 & 21456 & 20284\\
T11.0559 & \texttt{01234555111} & 6 & 6 & 2 & 36 & 7 & 5 & 2 & 0 & 57344 & 26296 & 16736\\
T11.0560 & \texttt{01234554554} & 5 & 6 & 3 & 8 & 7 & 5 & 2 & 1 & 77824 & 23596 & 30080\\
T11.0561 & \texttt{01234554553} & 5 & 6 & 4 & 8 & 7 & 5 & 2 & 1 & 75776 & 27000 & 26532\\
T11.0562 & \texttt{01234554552} & 5 & 6 & 4 & 8 & 7 & 5 & 2 & 1 & 73728 & 28084 & 24780\\
T11.0563 & \texttt{01234554551} & 5 & 6 & 4 & 8 & 7 & 5 & 2 & 1 & 71680 & 25488 & 24996\\
T11.0564 & \texttt{01234554545} & 4 & 6 & 2 & 4 & 8 & 5 & 2 & 0 & 77824 & 18308 & 32536\\
T11.0565 & \texttt{01234554544} & 5 & 6 & 2 & 4 & 9 & 7 & 3 & 1 & 75776 & 24608 & 29608\\
T11.0566 & \texttt{01234554543} & 5 & 6 & 3 & 2 & 10 & 9 & 4 & 3 & 73728 & 28672 & 25324\\
T11.0567 & \texttt{01234554542} & 5 & 6 & 3 & 2 & 10 & 9 & 4 & 3 & 71680 & 30176 & 23252\\
T11.0568 & \texttt{01234554541} & 5 & 6 & 3 & 2 & 10 & 9 & 4 & 3 & 69632 & 27840 & 23404\\
T11.0569 & \texttt{01234554534} & 4 & 6 & 3 & 2 & 10 & 9 & 3 & 2 & 73728 & 18080 & 29936\\
T11.0570 & \texttt{01234554533} & 5 & 6 & 3 & 4 & 9 & 7 & 3 & 1 & 71680 & 28608 & 24004\\
T11.0571 & \texttt{01234554532} & 5 & 6 & 4 & 2 & 10 & 9 & 4 & 3 & 69632 & 31456 & 20960\\
T11.0572 & \texttt{01234554531} & 5 & 6 & 4 & 2 & 10 & 9 & 4 & 3 & 67584 & 29952 & 20632\\
T11.0573 & \texttt{01234554523} & 4 & 6 & 3 & 2 & 10 & 9 & 3 & 2 & 69632 & 17612 & 27628\\
T11.0574 & \texttt{01234554522} & 5 & 6 & 3 & 4 & 9 & 7 & 3 & 1 & 67584 & 29600 & 20780\\
T11.0575 & \texttt{01234554521} & 5 & 6 & 4 & 2 & 10 & 9 & 4 & 3 & 65536 & 29760 & 19480\\
T11.0576 & \texttt{01234554512} & 4 & 6 & 3 & 2 & 10 & 9 & 3 & 2 & 65536 & 15784 & 26040\\
T11.0577 & \texttt{01234554511} & 5 & 6 & 3 & 4 & 9 & 7 & 3 & 1 & 63488 & 25344 & 20444\\
T11.0578 & \texttt{01234554444} & 6 & 6 & 2 & 48 & 7 & 5 & 2 & 0 & 73728 & 26708 & 29200\\
T11.0579 & \texttt{01234554443} & 6 & 6 & 3 & 12 & 8 & 6 & 3 & 1 & 71680 & 31432 & 23956\\
T11.0580 & \texttt{01234554442} & 6 & 6 & 3 & 12 & 8 & 6 & 3 & 1 & 69632 & 33356 & 21500\\
T11.0581 & \texttt{01234554441} & 6 & 6 & 3 & 12 & 8 & 6 & 3 & 1 & 67584 & 31280 & 21652\\
T11.0582 & \texttt{01234554434} & 5 & 6 & 3 & 4 & 9 & 7 & 3 & 1 & 71680 & 24080 & 26912\\
T11.0583 & \texttt{01234554433} & 6 & 6 & 3 & 8 & 8 & 5 & 3 & 0 & 69632 & 32368 & 21708\\
T11.0584 & \texttt{01234554432} & 6 & 6 & 4 & 4 & 9 & 7 & 4 & 2 & 67584 & 35636 & 18152\\
T11.0585 & \texttt{01234554431} & 6 & 6 & 4 & 4 & 9 & 7 & 4 & 2 & 65536 & 34392 & 17760\\
T11.0586 & \texttt{01234554423} & 5 & 6 & 3 & 4 & 9 & 7 & 3 & 1 & 67584 & 23552 & 24604\\
T11.0587 & \texttt{01234554422} & 6 & 6 & 3 & 8 & 8 & 5 & 3 & 0 & 65536 & 34460 & 17684\\
T11.0588 & \texttt{01234554421} & 6 & 6 & 4 & 4 & 9 & 7 & 4 & 2 & 63488 & 34880 & 16192\\
T11.0589 & \texttt{01234554412} & 5 & 6 & 3 & 4 & 9 & 7 & 3 & 1 & 63488 & 21744 & 23112\\
T11.0590 & \texttt{01234554411} & 6 & 6 & 3 & 8 & 8 & 5 & 3 & 0 & 61440 & 30904 & 17236\\
T11.0591 & \texttt{01234554345} & 4 & 6 & 3 & 2 & 10 & 9 & 3 & 2 & 73728 & 17184 & 29112\\
T11.0592 & \texttt{01234554344} & 5 & 6 & 4 & 4 & 9 & 7 & 3 & 1 & 71680 & 23040 & 26412\\
T11.0593 & \texttt{01234554343} & 5 & 6 & 3 & 2 & 10 & 9 & 4 & 3 & 69632 & 26592 & 23944\\
T11.0594 & \texttt{01234554342} & 5 & 6 & 4 & 2 & 10 & 9 & 4 & 3 & 67584 & 28768 & 21308\\
T11.0595 & \texttt{01234554341} & 5 & 6 & 4 & 2 & 10 & 9 & 4 & 3 & 65536 & 26848 & 21196\\
T11.0596 & \texttt{01234554333} & 6 & 6 & 3 & 12 & 8 & 6 & 3 & 1 & 67584 & 31296 & 21092\\
T11.0597 & \texttt{01234554332} & 6 & 6 & 4 & 4 & 9 & 7 & 4 & 2 & 65536 & 35236 & 16968\\
T11.0598 & \texttt{01234554331} & 6 & 6 & 4 & 4 & 9 & 7 & 4 & 2 & 63488 & 34408 & 16304\\
T11.0599 & \texttt{01234554323} & 5 & 6 & 4 & 2 & 10 & 9 & 4 & 3 & 65536 & 25968 & 21592\\
T11.0600 & \texttt{01234554322} & 6 & 6 & 4 & 4 & 9 & 7 & 4 & 2 & 63488 & 35148 & 15656\\
T11.0601 & \texttt{01234554321} & 6 & 6 & 5 & 2 & 10 & 9 & 5 & 4 & 61440 & 35984 & 13884\\
T11.0602 & \texttt{01234554312} & 5 & 6 & 4 & 2 & 10 & 9 & 4 & 3 & 61440 & 24192 & 20084\\
T11.0603 & \texttt{01234554311} & 6 & 6 & 4 & 4 & 9 & 7 & 4 & 2 & 59392 & 32712 & 14520\\
T11.0604 & \texttt{01234554234} & 4 & 6 & 3 & 2 & 10 & 9 & 3 & 2 & 67584 & 15368 & 26668\\
T11.0605 & \texttt{01234554233} & 5 & 6 & 4 & 4 & 9 & 7 & 3 & 1 & 65536 & 22372 & 23020\\
T11.0606 & \texttt{01234554232} & 5 & 6 & 3 & 2 & 10 & 9 & 4 & 3 & 63488 & 26816 & 19892\\
T11.0607 & \texttt{01234554231} & 5 & 6 & 4 & 2 & 10 & 9 & 4 & 3 & 61440 & 25728 & 19224\\
T11.0608 & \texttt{01234554222} & 6 & 6 & 3 & 12 & 8 & 6 & 3 & 1 & 61440 & 32540 & 15996\\
T11.0609 & \texttt{01234554221} & 6 & 6 & 4 & 4 & 9 & 7 & 4 & 2 & 59392 & 34208 & 13664\\
T11.0610 & \texttt{01234554212} & 5 & 6 & 4 & 2 & 10 & 9 & 4 & 3 & 59392 & 25104 & 18340\\
T11.0611 & \texttt{01234554211} & 6 & 6 & 4 & 4 & 9 & 7 & 4 & 2 & 57344 & 32344 & 13448\\
T11.0612 & \texttt{01234554123} & 4 & 6 & 3 & 2 & 10 & 9 & 3 & 2 & 61440 & 13648 & 24240\\
T11.0613 & \texttt{01234554122} & 5 & 6 & 4 & 4 & 9 & 7 & 3 & 1 & 59392 & 19864 & 20628\\
T11.0614 & \texttt{01234554121} & 5 & 6 & 3 & 2 & 10 & 9 & 4 & 3 & 57344 & 23008 & 18128\\
T11.0615 & \texttt{01234554111} & 6 & 6 & 3 & 12 & 8 & 6 & 3 & 1 & 55296 & 27688 & 14604\\
T11.0616 & \texttt{01234553455} & 4 & 6 & 3 & 8 & 6 & 3 & 1 & 0 & 75776 & 13984 & 30772\\
T11.0617 & \texttt{01234553454} & 4 & 6 & 3 & 2 & 10 & 9 & 3 & 2 & 73728 & 16896 & 29184\\
T11.0618 & \texttt{01234553453} & 4 & 6 & 2 & 2 & 10 & 9 & 3 & 2 & 71680 & 20544 & 26260\\
T11.0619 & \texttt{01234553452} & 4 & 6 & 3 & 2 & 10 & 9 & 3 & 2 & 69632 & 21544 & 24424\\
T11.0620 & \texttt{01234553451} & 4 & 6 & 3 & 2 & 10 & 9 & 3 & 2 & 67584 & 18896 & 24448\\
T11.0621 & \texttt{01234553444} & 5 & 6 & 3 & 12 & 8 & 6 & 2 & 0 & 71680 & 20896 & 27576\\
T11.0622 & \texttt{01234553443} & 5 & 6 & 3 & 4 & 9 & 7 & 3 & 1 & 69632 & 24608 & 24288\\
T11.0623 & \texttt{01234553442} & 5 & 6 & 4 & 4 & 9 & 7 & 3 & 1 & 67584 & 26700 & 21796\\
T11.0624 & \texttt{01234553441} & 5 & 6 & 4 & 4 & 9 & 7 & 3 & 1 & 65536 & 24728 & 21684\\
T11.0625 & \texttt{01234553434} & 4 & 6 & 2 & 4 & 8 & 5 & 2 & 0 & 69632 & 17632 & 27808\\
T11.0626 & \texttt{01234553433} & 5 & 6 & 2 & 4 & 9 & 7 & 3 & 1 & 67584 & 26368 & 22996\\
T11.0627 & \texttt{01234553432} & 5 & 6 & 3 & 2 & 10 & 9 & 4 & 3 & 65536 & 29552 & 19624\\
T11.0628 & \texttt{01234553431} & 5 & 6 & 3 & 2 & 10 & 9 & 4 & 3 & 63488 & 28256 & 19168\\
T11.0629 & \texttt{01234553423} & 4 & 6 & 3 & 2 & 10 & 9 & 3 & 2 & 65536 & 17116 & 25464\\
T11.0630 & \texttt{01234553422} & 5 & 6 & 3 & 4 & 9 & 7 & 3 & 1 & 63488 & 28240 & 19024\\
T11.0631 & \texttt{01234553421} & 5 & 6 & 4 & 2 & 10 & 9 & 4 & 3 & 61440 & 28608 & 17572\\
T11.0632 & \texttt{01234553412} & 4 & 6 & 3 & 2 & 10 & 9 & 3 & 2 & 61440 & 15304 & 23876\\
T11.0633 & \texttt{01234553411} & 5 & 6 & 3 & 4 & 9 & 7 & 3 & 1 & 59392 & 24544 & 18368\\
T11.0634 & \texttt{01234553333} & 6 & 6 & 2 & 48 & 7 & 5 & 2 & 0 & 65536 & 29280 & 21320\\
T11.0635 & \texttt{01234553332} & 6 & 6 & 3 & 12 & 8 & 6 & 3 & 1 & 63488 & 33556 & 16876\\
T11.0636 & \texttt{01234553331} & 6 & 6 & 3 & 12 & 8 & 6 & 3 & 1 & 61440 & 32936 & 16076\\
T11.0637 & \texttt{01234553323} & 5 & 6 & 3 & 4 & 9 & 7 & 3 & 1 & 63488 & 25696 & 20604\\
T11.0638 & \texttt{01234553322} & 6 & 6 & 3 & 8 & 8 & 5 & 3 & 0 & 61440 & 34012 & 15124\\
T11.0639 & \texttt{01234553321} & 6 & 6 & 4 & 4 & 9 & 7 & 4 & 2 & 59392 & 35056 & 13200\\
T11.0640 & \texttt{01234553312} & 5 & 6 & 3 & 4 & 9 & 7 & 3 & 1 & 59392 & 23936 & 19088\\
T11.0641 & \texttt{01234553311} & 6 & 6 & 3 & 8 & 8 & 5 & 3 & 0 & 57344 & 32136 & 13644\\
T11.0642 & \texttt{01234553234} & 4 & 6 & 3 & 2 & 10 & 9 & 3 & 2 & 65536 & 17656 & 24300\\
T11.0643 & \texttt{01234553233} & 5 & 6 & 4 & 4 & 9 & 7 & 3 & 1 & 63488 & 23796 & 20940\\
T11.0644 & \texttt{01234553232} & 5 & 6 & 3 & 2 & 10 & 9 & 4 & 3 & 61440 & 27376 & 18292\\
T11.0645 & \texttt{01234553231} & 5 & 6 & 4 & 2 & 10 & 9 & 4 & 3 & 59392 & 26912 & 17272\\
T11.0646 & \texttt{01234553222} & 6 & 6 & 3 & 12 & 8 & 6 & 3 & 1 & 59392 & 32236 & 14876\\
T11.0647 & \texttt{01234553221} & 6 & 6 & 4 & 4 & 9 & 7 & 4 & 2 & 57344 & 34112 & 12384\\
T11.0648 & \texttt{01234553212} & 5 & 6 & 4 & 2 & 10 & 9 & 4 & 3 & 57344 & 25680 & 16724\\
T11.0649 & \texttt{01234553211} & 6 & 6 & 4 & 4 & 9 & 7 & 4 & 2 & 55296 & 32600 & 11960\\
T11.0650 & \texttt{01234553123} & 4 & 6 & 3 & 2 & 10 & 9 & 3 & 2 & 59392 & 15440 & 22120\\
T11.0651 & \texttt{01234553122} & 5 & 6 & 4 & 4 & 9 & 7 & 3 & 1 & 57344 & 21336 & 18588\\
T11.0652 & \texttt{01234553121} & 5 & 6 & 3 & 2 & 10 & 9 & 4 & 3 & 55296 & 24160 & 16232\\
T11.0653 & \texttt{01234553111} & 6 & 6 & 3 & 12 & 8 & 6 & 3 & 1 & 53248 & 28520 & 12852\\
T11.0654 & \texttt{01234552345} & 3 & 6 & 2 & 2 & 10 & 9 & 2 & 1 & 69632 & 7112 & 28720\\
T11.0655 & \texttt{01234552344} & 4 & 6 & 3 & 4 & 9 & 7 & 2 & 0 & 67584 & 12548 & 27052\\
T11.0656 & \texttt{01234552343} & 4 & 6 & 3 & 2 & 10 & 9 & 3 & 2 & 65536 & 15680 & 25108\\
T11.0657 & \texttt{01234552342} & 4 & 6 & 2 & 2 & 10 & 9 & 3 & 2 & 63488 & 19640 & 22096\\
T11.0658 & \texttt{01234552341} & 4 & 6 & 3 & 2 & 10 & 9 & 3 & 2 & 61440 & 17408 & 21836\\
T11.0659 & \texttt{01234552333} & 5 & 6 & 3 & 12 & 8 & 6 & 2 & 0 & 63488 & 19964 & 22744\\
T11.0660 & \texttt{01234552332} & 5 & 6 & 3 & 4 & 9 & 7 & 3 & 1 & 61440 & 24052 & 19384\\
T11.0661 & \texttt{01234552331} & 5 & 6 & 4 & 4 & 9 & 7 & 3 & 1 & 59392 & 23328 & 18548\\
T11.0662 & \texttt{01234552323} & 4 & 6 & 2 & 4 & 8 & 5 & 2 & 0 & 61440 & 16508 & 23200\\
T11.0663 & \texttt{01234552322} & 5 & 6 & 2 & 4 & 9 & 7 & 3 & 1 & 59392 & 25904 & 17768\\
T11.0664 & \texttt{01234552321} & 5 & 6 & 3 & 2 & 10 & 9 & 4 & 3 & 57344 & 26688 & 16020\\
T11.0665 & \texttt{01234552312} & 4 & 6 & 3 & 2 & 10 & 9 & 3 & 2 & 57344 & 14728 & 21584\\
T11.0666 & \texttt{01234552311} & 5 & 6 & 3 & 4 & 9 & 7 & 3 & 1 & 55296 & 23328 & 16388\\
T11.0667 & \texttt{01234552222} & 6 & 6 & 2 & 48 & 7 & 5 & 2 & 0 & 57344 & 29036 & 15384\\
T11.0668 & \texttt{01234552221} & 6 & 6 & 3 & 12 & 8 & 6 & 3 & 1 & 55296 & 31328 & 12596\\
T11.0669 & \texttt{01234552212} & 5 & 6 & 3 & 4 & 9 & 7 & 3 & 1 & 55296 & 24240 & 16168\\
T11.0670 & \texttt{01234552211} & 6 & 6 & 3 & 8 & 8 & 5 & 3 & 0 & 53248 & 30520 & 11732\\
T11.0671 & \texttt{01234552123} & 4 & 6 & 3 & 2 & 10 & 9 & 3 & 2 & 57344 & 16432 & 20400\\
T11.0672 & \texttt{01234552122} & 5 & 6 & 4 & 4 & 9 & 7 & 3 & 1 & 55296 & 21688 & 17076\\
T11.0673 & \texttt{01234552121} & 5 & 6 & 3 & 2 & 10 & 9 & 4 & 3 & 53248 & 23872 & 15056\\
T11.0674 & \texttt{01234552111} & 6 & 6 & 3 & 12 & 8 & 6 & 3 & 1 & 51200 & 27592 & 12012\\
T11.0675 & \texttt{01234551234} & 3 & 6 & 2 & 2 & 10 & 9 & 2 & 1 & 61440 & 5744 & 25596\\
T11.0676 & \texttt{01234551233} & 4 & 6 & 3 & 4 & 9 & 7 & 2 & 0 & 59392 & 10744 & 23444\\
T11.0677 & \texttt{01234551232} & 4 & 6 & 3 & 2 & 10 & 9 & 3 & 2 & 57344 & 13184 & 21584\\
T11.0678 & \texttt{01234551231} & 4 & 6 & 2 & 2 & 10 & 9 & 3 & 2 & 55296 & 15504 & 19436\\
T11.0679 & \texttt{01234551222} & 5 & 6 & 3 & 12 & 8 & 6 & 2 & 0 & 55296 & 16904 & 18944\\
T11.0680 & \texttt{01234551221} & 5 & 6 & 3 & 4 & 9 & 7 & 3 & 1 & 53248 & 19480 & 16552\\
T11.0681 & \texttt{01234551212} & 4 & 6 & 2 & 4 & 8 & 5 & 2 & 0 & 53248 & 13064 & 19736\\
T11.0682 & \texttt{01234551211} & 5 & 6 & 2 & 4 & 9 & 7 & 3 & 1 & 51200 & 20384 & 15228\\
T11.0683 & \texttt{01234551111} & 6 & 6 & 2 & 48 & 7 & 5 & 2 & 0 & 49152 & 22824 & 12880\\
T11.0684 & \texttt{01234545454} & 4 & 6 & 1 & 6 & 7 & 5 & 2 & 1 & 75776 & 21000 & 31104\\
T11.0685 & \texttt{01234545453} & 4 & 6 & 2 & 6 & 7 & 5 & 2 & 1 & 73728 & 24272 & 27492\\
T11.0686 & \texttt{01234545452} & 4 & 6 & 2 & 6 & 7 & 5 & 2 & 1 & 71680 & 25272 & 25708\\
T11.0687 & \texttt{01234545451} & 4 & 6 & 2 & 6 & 7 & 5 & 2 & 1 & 69632 & 22624 & 25908\\
T11.0688 & \texttt{01234545444} & 5 & 6 & 1 & 12 & 7 & 4 & 2 & 0 & 73728 & 24780 & 29736\\
T11.0689 & \texttt{01234545443} & 5 & 6 & 2 & 4 & 8 & 5 & 3 & 1 & 71680 & 29240 & 24908\\
T11.0690 & \texttt{01234545442} & 5 & 6 & 2 & 4 & 8 & 5 & 3 & 1 & 69632 & 30996 & 22612\\
T11.0691 & \texttt{01234545441} & 5 & 6 & 2 & 4 & 8 & 5 & 3 & 1 & 67584 & 28816 & 22748\\
T11.0692 & \texttt{01234545434} & 4 & 6 & 2 & 2 & 9 & 7 & 3 & 2 & 71680 & 20592 & 28456\\
T11.0693 & \texttt{01234545433} & 5 & 6 & 2 & 4 & 8 & 5 & 3 & 1 & 69632 & 29776 & 23028\\
T11.0694 & \texttt{01234545432} & 5 & 6 & 3 & 2 & 9 & 7 & 4 & 3 & 67584 & 32876 & 19696\\
T11.0695 & \texttt{01234545431} & 5 & 6 & 3 & 2 & 9 & 7 & 4 & 3 & 65536 & 31528 & 19320\\
T11.0696 & \texttt{01234545423} & 4 & 6 & 2 & 2 & 9 & 7 & 3 & 2 & 67584 & 20088 & 26148\\
T11.0697 & \texttt{01234545422} & 5 & 6 & 2 & 4 & 8 & 5 & 3 & 1 & 65536 & 31428 & 19324\\
T11.0698 & \texttt{01234545421} & 5 & 6 & 3 & 2 & 9 & 7 & 4 & 3 & 63488 & 31744 & 17912\\
T11.0699 & \texttt{01234545412} & 4 & 6 & 2 & 2 & 9 & 7 & 3 & 2 & 63488 & 18272 & 24608\\
T11.0700 & \texttt{01234545411} & 5 & 6 & 2 & 4 & 8 & 5 & 3 & 1 & 61440 & 27592 & 18892\\
T11.0701 & \texttt{01234545345} & 3 & 6 & 2 & 2 & 9 & 7 & 2 & 1 & 73728 & 11360 & 31528\\
T11.0702 & \texttt{01234545344} & 4 & 6 & 3 & 4 & 8 & 5 & 2 & 0 & 71680 & 18112 & 28716\\
T11.0703 & \texttt{01234545343} & 4 & 6 & 2 & 2 & 9 & 7 & 3 & 2 & 69632 & 22560 & 25912\\
T11.0704 & \texttt{01234545342} & 4 & 6 & 3 & 2 & 9 & 7 & 3 & 2 & 67584 & 24232 & 23596\\
T11.0705 & \texttt{01234545341} & 4 & 6 & 3 & 2 & 9 & 7 & 3 & 2 & 65536 & 22000 & 23548\\
T11.0706 & \texttt{01234545333} & 5 & 6 & 2 & 12 & 7 & 4 & 2 & 0 & 67584 & 28160 & 22724\\
T11.0707 & \texttt{01234545332} & 5 & 6 & 3 & 4 & 8 & 5 & 3 & 1 & 65536 & 31932 & 18856\\
T11.0708 & \texttt{01234545331} & 5 & 6 & 3 & 4 & 8 & 5 & 3 & 1 & 63488 & 31000 & 18224\\
T11.0709 & \texttt{01234545323} & 4 & 6 & 3 & 2 & 9 & 7 & 3 & 2 & 65536 & 21960 & 23592\\
T11.0710 & \texttt{01234545322} & 5 & 6 & 3 & 4 & 8 & 5 & 3 & 1 & 63488 & 31572 & 17704\\
T11.0711 & \texttt{01234545321} & 5 & 6 & 4 & 2 & 9 & 7 & 4 & 3 & 61440 & 32304 & 16028\\
T11.0712 & \texttt{01234545312} & 4 & 6 & 3 & 2 & 9 & 7 & 3 & 2 & 61440 & 20176 & 22052\\
T11.0713 & \texttt{01234545311} & 5 & 6 & 3 & 4 & 8 & 5 & 3 & 1 & 59392 & 28856 & 16632\\
T11.0714 & \texttt{01234545234} & 3 & 6 & 2 & 2 & 9 & 7 & 2 & 1 & 67584 & 10080 & 28812\\
T11.0715 & \texttt{01234545233} & 4 & 6 & 3 & 4 & 8 & 5 & 2 & 0 & 65536 & 17516 & 25212\\
T11.0716 & \texttt{01234545232} & 4 & 6 & 2 & 2 & 9 & 7 & 3 & 2 & 63488 & 22392 & 22036\\
T11.0717 & \texttt{01234545231} & 4 & 6 & 3 & 2 & 9 & 7 & 3 & 2 & 61440 & 20992 & 21464\\
T11.0718 & \texttt{01234545222} & 5 & 6 & 2 & 12 & 7 & 4 & 2 & 0 & 61440 & 28548 & 18092\\
T11.0719 & \texttt{01234545221} & 5 & 6 & 3 & 4 & 8 & 5 & 3 & 1 & 59392 & 30112 & 15888\\
T11.0720 & \texttt{01234545212} & 4 & 6 & 3 & 2 & 9 & 7 & 3 & 2 & 59392 & 20672 & 20484\\
T11.0721 & \texttt{01234545211} & 5 & 6 & 3 & 4 & 8 & 5 & 3 & 1 & 57344 & 28072 & 15704\\
T11.0722 & \texttt{01234545123} & 3 & 6 & 2 & 2 & 9 & 7 & 2 & 1 & 61440 & 8608 & 26224\\
T11.0723 & \texttt{01234545122} & 4 & 6 & 3 & 4 & 8 & 5 & 2 & 0 & 59392 & 14984 & 22692\\
T11.0724 & \texttt{01234545121} & 4 & 6 & 2 & 2 & 9 & 7 & 3 & 2 & 57344 & 18288 & 20240\\
T11.0725 & \texttt{01234545111} & 5 & 6 & 2 & 12 & 7 & 4 & 2 & 0 & 55296 & 23128 & 16764\\
T11.0726 & \texttt{01234544444} & 6 & 6 & 1 & 120 & 7 & 6 & 2 & 1 & 71680 & 26040 & 29928\\
T11.0727 & \texttt{01234544443} & 6 & 6 & 2 & 24 & 8 & 7 & 3 & 2 & 69632 & 30896 & 24332\\
T11.0728 & \texttt{01234544442} & 6 & 6 & 2 & 24 & 8 & 7 & 3 & 2 & 67584 & 32904 & 21748\\
T11.0729 & \texttt{01234544441} & 6 & 6 & 2 & 24 & 8 & 7 & 3 & 2 & 65536 & 30880 & 21932\\
T11.0730 & \texttt{01234544434} & 5 & 6 & 2 & 6 & 9 & 8 & 3 & 2 & 69632 & 24192 & 26928\\
T11.0731 & \texttt{01234544433} & 6 & 6 & 2 & 12 & 8 & 6 & 3 & 1 & 67584 & 32032 & 21780\\
T11.0732 & \texttt{01234544432} & 6 & 6 & 3 & 6 & 9 & 8 & 4 & 3 & 65536 & 35384 & 18032\\
T11.0733 & \texttt{01234544431} & 6 & 6 & 3 & 6 & 9 & 8 & 4 & 3 & 63488 & 34192 & 17640\\
T11.0734 & \texttt{01234544423} & 5 & 6 & 2 & 6 & 9 & 8 & 3 & 2 & 65536 & 23652 & 24588\\
T11.0735 & \texttt{01234544422} & 6 & 6 & 2 & 12 & 8 & 6 & 3 & 1 & 63488 & 34344 & 17500\\
T11.0736 & \texttt{01234544421} & 6 & 6 & 3 & 6 & 9 & 8 & 4 & 3 & 61440 & 34816 & 15944\\
T11.0737 & \texttt{01234544412} & 5 & 6 & 2 & 6 & 9 & 8 & 3 & 2 & 61440 & 21848 & 23128\\
T11.0738 & \texttt{01234544411} & 6 & 6 & 2 & 12 & 8 & 6 & 3 & 1 & 59392 & 30928 & 17068\\
T11.0739 & \texttt{01234544345} & 4 & 6 & 2 & 2 & 10 & 9 & 3 & 2 & 71680 & 18464 & 28480\\
T11.0740 & \texttt{01234544344} & 5 & 6 & 3 & 4 & 9 & 7 & 3 & 1 & 69632 & 23872 & 25836\\
T11.0741 & \texttt{01234544343} & 5 & 6 & 2 & 2 & 10 & 9 & 4 & 3 & 67584 & 26976 & 23536\\
T11.0742 & \texttt{01234544342} & 5 & 6 & 3 & 2 & 10 & 9 & 4 & 3 & 65536 & 29404 & 20644\\
T11.0743 & \texttt{01234544341} & 5 & 6 & 3 & 2 & 10 & 9 & 4 & 3 & 63488 & 27640 & 20500\\
T11.0744 & \texttt{01234544333} & 6 & 6 & 2 & 12 & 8 & 6 & 3 & 1 & 65536 & 31232 & 20852\\
T11.0745 & \texttt{01234544332} & 6 & 6 & 3 & 4 & 9 & 7 & 4 & 2 & 63488 & 35256 & 16504\\
T11.0746 & \texttt{01234544331} & 6 & 6 & 3 & 4 & 9 & 7 & 4 & 2 & 61440 & 34480 & 15824\\
T11.0747 & \texttt{01234544323} & 5 & 6 & 3 & 2 & 10 & 9 & 4 & 3 & 63488 & 26340 & 21120\\
T11.0748 & \texttt{01234544322} & 6 & 6 & 3 & 4 & 9 & 7 & 4 & 2 & 61440 & 35304 & 15064\\
T11.0749 & \texttt{01234544321} & 6 & 6 & 4 & 2 & 10 & 9 & 5 & 4 & 59392 & 36192 & 13212\\
T11.0750 & \texttt{01234544312} & 5 & 6 & 3 & 2 & 10 & 9 & 4 & 3 & 59392 & 24568 & 19628\\
T11.0751 & \texttt{01234544311} & 6 & 6 & 3 & 4 & 9 & 7 & 4 & 2 & 57344 & 33008 & 13896\\
T11.0752 & \texttt{01234544234} & 4 & 6 & 2 & 2 & 10 & 9 & 3 & 2 & 65536 & 16380 & 26148\\
T11.0753 & \texttt{01234544233} & 5 & 6 & 3 & 4 & 9 & 7 & 3 & 1 & 63488 & 23168 & 22428\\
T11.0754 & \texttt{01234544232} & 5 & 6 & 2 & 2 & 10 & 9 & 4 & 3 & 61440 & 27396 & 19276\\
T11.0755 & \texttt{01234544231} & 5 & 6 & 3 & 2 & 10 & 9 & 4 & 3 & 59392 & 26464 & 18544\\
T11.0756 & \texttt{01234544222} & 6 & 6 & 2 & 12 & 8 & 6 & 3 & 1 & 59392 & 32904 & 15356\\
T11.0757 & \texttt{01234544221} & 6 & 6 & 3 & 4 & 9 & 7 & 4 & 2 & 57344 & 34624 & 12912\\
T11.0758 & \texttt{01234544212} & 5 & 6 & 3 & 2 & 10 & 9 & 4 & 3 & 57344 & 25688 & 17708\\
T11.0759 & \texttt{01234544211} & 6 & 6 & 3 & 4 & 9 & 7 & 4 & 2 & 55296 & 32848 & 12680\\
T11.0760 & \texttt{01234544123} & 4 & 6 & 2 & 2 & 10 & 9 & 3 & 2 & 59392 & 14536 & 23800\\
T11.0761 & \texttt{01234544122} & 5 & 6 & 3 & 4 & 9 & 7 & 3 & 1 & 57344 & 20672 & 20100\\
T11.0762 & \texttt{01234544121} & 5 & 6 & 2 & 2 & 10 & 9 & 4 & 3 & 55296 & 23736 & 17528\\
T11.0763 & \texttt{01234544111} & 6 & 6 & 2 & 12 & 8 & 6 & 3 & 1 & 53248 & 28336 & 13932\\
T11.0764 & \texttt{01234543454} & 4 & 6 & 3 & 2 & 7 & 3 & 2 & 0 & 71680 & 19360 & 27708\\
T11.0765 & \texttt{01234543453} & 4 & 6 & 2 & 1 & 11 & 11 & 4 & 4 & 69632 & 22112 & 25176\\
T11.0766 & \texttt{01234543452} & 4 & 6 & 3 & 1 & 11 & 11 & 4 & 4 & 67584 & 23700 & 22972\\
T11.0767 & \texttt{01234543451} & 4 & 6 & 3 & 1 & 11 & 11 & 4 & 4 & 65536 & 21416 & 22908\\
T11.0768 & \texttt{01234543444} & 5 & 6 & 3 & 6 & 9 & 8 & 3 & 2 & 69632 & 22912 & 26156\\
T11.0769 & \texttt{01234543443} & 5 & 6 & 3 & 2 & 10 & 9 & 4 & 3 & 67584 & 25728 & 23316\\
T11.0770 & \texttt{01234543442} & 5 & 6 & 4 & 2 & 10 & 9 & 4 & 3 & 65536 & 28240 & 20456\\
T11.0771 & \texttt{01234543441} & 5 & 6 & 4 & 2 & 10 & 9 & 4 & 3 & 63488 & 26528 & 20256\\
T11.0772 & \texttt{01234543434} & 4 & 6 & 2 & 2 & 9 & 7 & 3 & 2 & 67584 & 19648 & 26500\\
T11.0773 & \texttt{01234543433} & 5 & 6 & 2 & 2 & 10 & 9 & 4 & 3 & 65536 & 27040 & 22248\\
T11.0774 & \texttt{01234543432} & 5 & 6 & 3 & 1 & 11 & 11 & 5 & 5 & 63488 & 30476 & 18572\\
T11.0775 & \texttt{01234543431} & 5 & 6 & 3 & 1 & 11 & 11 & 5 & 5 & 61440 & 29336 & 18060\\
T11.0776 & \texttt{01234543423} & 4 & 6 & 3 & 1 & 11 & 11 & 4 & 4 & 63488 & 19096 & 24140\\
T11.0777 & \texttt{01234543422} & 5 & 6 & 3 & 2 & 10 & 9 & 4 & 3 & 61440 & 29572 & 17748\\
T11.0778 & \texttt{01234543421} & 5 & 6 & 4 & 1 & 11 & 11 & 5 & 5 & 59392 & 30096 & 16176\\
T11.0779 & \texttt{01234543412} & 4 & 6 & 3 & 1 & 11 & 11 & 4 & 4 & 59392 & 17296 & 22592\\
T11.0780 & \texttt{01234543411} & 5 & 6 & 3 & 2 & 10 & 9 & 4 & 3 & 57344 & 26296 & 16972\\
T11.0781 & \texttt{01234543333} & 6 & 6 & 2 & 24 & 8 & 7 & 3 & 2 & 63488 & 29504 & 20796\\
T11.0782 & \texttt{01234543332} & 6 & 6 & 3 & 6 & 9 & 8 & 4 & 3 & 61440 & 33864 & 16112\\
T11.0783 & \texttt{01234543331} & 6 & 6 & 3 & 6 & 9 & 8 & 4 & 3 & 59392 & 33296 & 15288\\
T11.0784 & \texttt{01234543323} & 5 & 6 & 3 & 2 & 10 & 9 & 4 & 3 & 61440 & 26356 & 19776\\
T11.0785 & \texttt{01234543322} & 6 & 6 & 3 & 4 & 9 & 7 & 4 & 2 & 59392 & 34456 & 14200\\
T11.0786 & \texttt{01234543321} & 6 & 6 & 4 & 2 & 10 & 9 & 5 & 4 & 57344 & 35552 & 12188\\
T11.0787 & \texttt{01234543312} & 5 & 6 & 3 & 2 & 10 & 9 & 4 & 3 & 57344 & 24600 & 18268\\
T11.0788 & \texttt{01234543311} & 6 & 6 & 3 & 4 & 9 & 7 & 4 & 2 & 55296 & 32720 & 12664\\
T11.0789 & \texttt{01234543234} & 4 & 6 & 3 & 1 & 11 & 11 & 4 & 4 & 63488 & 18956 & 23312\\
T11.0790 & \texttt{01234543233} & 5 & 6 & 4 & 2 & 10 & 9 & 4 & 3 & 61440 & 24880 & 19904\\
T11.0791 & \texttt{01234543232} & 5 & 6 & 3 & 1 & 11 & 11 & 5 & 5 & 59392 & 28244 & 17256\\
T11.0792 & \texttt{01234543231} & 5 & 6 & 4 & 1 & 11 & 11 & 5 & 5 & 57344 & 27936 & 16148\\
T11.0793 & \texttt{01234543222} & 6 & 6 & 3 & 6 & 9 & 8 & 4 & 3 & 57344 & 32888 & 13840\\
T11.0794 & \texttt{01234543221} & 6 & 6 & 4 & 2 & 10 & 9 & 5 & 4 & 55296 & 34816 & 11228\\
T11.0795 & \texttt{01234543212} & 5 & 6 & 4 & 1 & 11 & 11 & 5 & 5 & 55296 & 26552 & 15664\\
T11.0796 & \texttt{01234543211} & 6 & 6 & 4 & 2 & 10 & 9 & 5 & 4 & 53248 & 33392 & 10772\\
T11.0797 & \texttt{01234543123} & 4 & 6 & 3 & 1 & 11 & 11 & 4 & 4 & 57344 & 16616 & 21204\\
T11.0798 & \texttt{01234543122} & 5 & 6 & 4 & 2 & 10 & 9 & 4 & 3 & 55296 & 22432 & 17592\\
T11.0799 & \texttt{01234543121} & 5 & 6 & 3 & 1 & 11 & 11 & 5 & 5 & 53248 & 25176 & 15172\\
T11.0800 & \texttt{01234543111} & 6 & 6 & 3 & 6 & 9 & 8 & 4 & 3 & 51200 & 29456 & 11728\\
T11.0801 & \texttt{01234542345} & 3 & 6 & 2 & 1 & 11 & 11 & 3 & 3 & 67584 & 9564 & 27564\\
T11.0802 & \texttt{01234542344} & 4 & 6 & 3 & 2 & 10 & 9 & 3 & 2 & 65536 & 14784 & 25824\\
T11.0803 & \texttt{01234542343} & 4 & 6 & 3 & 1 & 11 & 11 & 4 & 4 & 63488 & 17700 & 23880\\
T11.0804 & \texttt{01234542342} & 4 & 6 & 2 & 1 & 11 & 11 & 4 & 4 & 61440 & 21228 & 20940\\
T11.0805 & \texttt{01234542341} & 4 & 6 & 3 & 1 & 11 & 11 & 4 & 4 & 59392 & 19360 & 20576\\
T11.0806 & \texttt{01234542333} & 5 & 6 & 3 & 6 & 9 & 8 & 3 & 2 & 61440 & 21768 & 21492\\
T11.0807 & \texttt{01234542332} & 5 & 6 & 3 & 2 & 10 & 9 & 4 & 3 & 59392 & 25424 & 18228\\
T11.0808 & \texttt{01234542331} & 5 & 6 & 4 & 2 & 10 & 9 & 4 & 3 & 57344 & 24960 & 17256\\
T11.0809 & \texttt{01234542323} & 4 & 6 & 2 & 2 & 9 & 7 & 3 & 2 & 59392 & 18312 & 21996\\
T11.0810 & \texttt{01234542322} & 5 & 6 & 2 & 2 & 10 & 9 & 4 & 3 & 57344 & 27060 & 16660\\
T11.0811 & \texttt{01234542321} & 5 & 6 & 3 & 1 & 11 & 11 & 5 & 5 & 55296 & 28000 & 14776\\
T11.0812 & \texttt{01234542312} & 4 & 6 & 3 & 1 & 11 & 11 & 4 & 4 & 55296 & 16544 & 20404\\
T11.0813 & \texttt{01234542311} & 5 & 6 & 3 & 2 & 10 & 9 & 4 & 3 & 53248 & 24904 & 15112\\
T11.0814 & \texttt{01234542222} & 6 & 6 & 2 & 24 & 8 & 7 & 3 & 2 & 55296 & 29976 & 14324\\
T11.0815 & \texttt{01234542221} & 6 & 6 & 3 & 6 & 9 & 8 & 4 & 3 & 53248 & 32320 & 11400\\
T11.0816 & \texttt{01234542212} & 5 & 6 & 3 & 2 & 10 & 9 & 4 & 3 & 53248 & 25400 & 15020\\
T11.0817 & \texttt{01234542211} & 6 & 6 & 3 & 4 & 9 & 7 & 4 & 2 & 51200 & 31600 & 10472\\
T11.0818 & \texttt{01234542123} & 4 & 6 & 3 & 1 & 11 & 11 & 4 & 4 & 55296 & 17896 & 19300\\
T11.0819 & \texttt{01234542122} & 5 & 6 & 4 & 2 & 10 & 9 & 4 & 3 & 53248 & 23072 & 15912\\
T11.0820 & \texttt{01234542121} & 5 & 6 & 3 & 1 & 11 & 11 & 5 & 5 & 51200 & 25176 & 13844\\
T11.0821 & \texttt{01234542111} & 6 & 6 & 3 & 6 & 9 & 8 & 4 & 3 & 49152 & 28816 & 10752\\
T11.0822 & \texttt{01234541234} & 3 & 6 & 2 & 1 & 11 & 11 & 3 & 3 & 59392 & 7752 & 24640\\
T11.0823 & \texttt{01234541233} & 4 & 6 & 3 & 2 & 10 & 9 & 3 & 2 & 57344 & 12672 & 22432\\
T11.0824 & \texttt{01234541232} & 4 & 6 & 3 & 1 & 11 & 11 & 4 & 4 & 55296 & 15032 & 20540\\
T11.0825 & \texttt{01234541231} & 4 & 6 & 2 & 1 & 11 & 11 & 4 & 4 & 53248 & 17192 & 18352\\
T11.0826 & \texttt{01234541222} & 5 & 6 & 3 & 6 & 9 & 8 & 3 & 2 & 53248 & 18672 & 17860\\
T11.0827 & \texttt{01234541221} & 5 & 6 & 3 & 2 & 10 & 9 & 4 & 3 & 51200 & 21088 & 15436\\
T11.0828 & \texttt{01234541212} & 4 & 6 & 2 & 2 & 9 & 7 & 3 & 2 & 51200 & 14832 & 18668\\
T11.0829 & \texttt{01234541211} & 5 & 6 & 2 & 2 & 10 & 9 & 4 & 3 & 49152 & 21912 & 14096\\
T11.0830 & \texttt{01234541111} & 6 & 6 & 2 & 24 & 8 & 7 & 3 & 2 & 47104 & 24272 & 11732\\
T11.0831 & \texttt{01234534534} & 3 & 6 & 1 & 2 & 8 & 5 & 2 & 1 & 69632 & 12096 & 29192\\
T11.0832 & \texttt{01234534533} & 4 & 6 & 1 & 4 & 7 & 3 & 2 & 0 & 67584 & 23072 & 23764\\
T11.0833 & \texttt{01234534532} & 4 & 6 & 2 & 2 & 8 & 5 & 3 & 2 & 65536 & 25836 & 20776\\
T11.0834 & \texttt{01234534531} & 4 & 6 & 2 & 2 & 8 & 5 & 3 & 2 & 63488 & 24280 & 20368\\
T11.0835 & \texttt{01234534523} & 3 & 6 & 2 & 2 & 8 & 5 & 2 & 1 & 65536 & 11640 & 26768\\
T11.0836 & \texttt{01234534522} & 4 & 6 & 2 & 4 & 7 & 3 & 2 & 0 & 63488 & 23844 & 20352\\
T11.0837 & \texttt{01234534521} & 4 & 6 & 3 & 2 & 8 & 5 & 3 & 2 & 61440 & 23952 & 19044\\
T11.0838 & \texttt{01234534512} & 3 & 6 & 2 & 2 & 8 & 5 & 2 & 1 & 61440 & 9808 & 25068\\
T11.0839 & \texttt{01234534511} & 4 & 6 & 2 & 4 & 7 & 3 & 2 & 0 & 59392 & 19448 & 19760\\
T11.0840 & \texttt{01234534444} & 5 & 6 & 2 & 24 & 8 & 7 & 2 & 1 & 69632 & 20832 & 27856\\
T11.0841 & \texttt{01234534443} & 5 & 6 & 2 & 6 & 9 & 8 & 3 & 2 & 67584 & 24768 & 23904\\
T11.0842 & \texttt{01234534442} & 5 & 6 & 3 & 6 & 9 & 8 & 3 & 2 & 65536 & 26944 & 21252\\
T11.0843 & \texttt{01234534441} & 5 & 6 & 3 & 6 & 9 & 8 & 3 & 2 & 63488 & 25024 & 21156\\
T11.0844 & \texttt{01234534434} & 4 & 6 & 2 & 2 & 10 & 9 & 3 & 2 & 67584 & 17952 & 26716\\
T11.0845 & \texttt{01234534433} & 5 & 6 & 2 & 4 & 9 & 7 & 3 & 1 & 65536 & 25792 & 22296\\
T11.0846 & \texttt{01234534432} & 5 & 6 & 3 & 2 & 10 & 9 & 4 & 3 & 63488 & 29144 & 18748\\
T11.0847 & \texttt{01234534431} & 5 & 6 & 3 & 2 & 10 & 9 & 4 & 3 & 61440 & 27952 & 18252\\
T11.0848 & \texttt{01234534423} & 4 & 6 & 3 & 2 & 10 & 9 & 3 & 2 & 63488 & 17412 & 24372\\
T11.0849 & \texttt{01234534422} & 5 & 6 & 3 & 4 & 9 & 7 & 3 & 1 & 61440 & 28104 & 18004\\
T11.0850 & \texttt{01234534421} & 5 & 6 & 4 & 2 & 10 & 9 & 4 & 3 & 59392 & 28576 & 16480\\
T11.0851 & \texttt{01234534412} & 4 & 6 & 3 & 2 & 10 & 9 & 3 & 2 & 59392 & 15608 & 22808\\
T11.0852 & \texttt{01234534411} & 5 & 6 & 3 & 4 & 9 & 7 & 3 & 1 & 57344 & 24688 & 17260\\
T11.0853 & \texttt{01234534343} & 4 & 6 & 1 & 2 & 9 & 7 & 3 & 2 & 65536 & 21504 & 24472\\
T11.0854 & \texttt{01234534342} & 4 & 6 & 2 & 2 & 9 & 7 & 3 & 2 & 63488 & 23680 & 21692\\
T11.0855 & \texttt{01234534341} & 4 & 6 & 2 & 2 & 9 & 7 & 3 & 2 & 61440 & 21760 & 21484\\
T11.0856 & \texttt{01234534333} & 5 & 6 & 1 & 6 & 9 & 8 & 3 & 2 & 63488 & 26208 & 21844\\
T11.0857 & \texttt{01234534332} & 5 & 6 & 2 & 2 & 10 & 9 & 4 & 3 & 61440 & 30148 & 17704\\
T11.0858 & \texttt{01234534331} & 5 & 6 & 2 & 2 & 10 & 9 & 4 & 3 & 59392 & 29320 & 17008\\
T11.0859 & \texttt{01234534323} & 4 & 6 & 2 & 1 & 11 & 11 & 4 & 4 & 61440 & 20880 & 22080\\
T11.0860 & \texttt{01234534322} & 5 & 6 & 2 & 2 & 10 & 9 & 4 & 3 & 59392 & 30060 & 16288\\
T11.0861 & \texttt{01234534321} & 5 & 6 & 3 & 1 & 11 & 11 & 5 & 5 & 57344 & 30896 & 14500\\
T11.0862 & \texttt{01234534312} & 4 & 6 & 2 & 1 & 11 & 11 & 4 & 4 & 57344 & 19104 & 20540\\
T11.0863 & \texttt{01234534311} & 5 & 6 & 2 & 2 & 10 & 9 & 4 & 3 & 55296 & 27624 & 15056\\
T11.0864 & \texttt{01234534234} & 3 & 6 & 2 & 1 & 11 & 11 & 3 & 3 & 63488 & 10280 & 26696\\
T11.0865 & \texttt{01234534233} & 4 & 6 & 3 & 2 & 10 & 9 & 3 & 2 & 61440 & 17284 & 23168\\
T11.0866 & \texttt{01234534232} & 4 & 6 & 2 & 1 & 11 & 11 & 4 & 4 & 59392 & 21728 & 20160\\
T11.0867 & \texttt{01234534231} & 4 & 6 & 3 & 1 & 11 & 11 & 4 & 4 & 57344 & 20640 & 19404\\
T11.0868 & \texttt{01234534222} & 5 & 6 & 2 & 6 & 9 & 8 & 3 & 2 & 57344 & 27452 & 16384\\
T11.0869 & \texttt{01234534221} & 5 & 6 & 3 & 2 & 10 & 9 & 4 & 3 & 55296 & 29120 & 14044\\
T11.0870 & \texttt{01234534212} & 4 & 6 & 3 & 1 & 11 & 11 & 4 & 4 & 55296 & 20016 & 18584\\
T11.0871 & \texttt{01234534211} & 5 & 6 & 3 & 2 & 10 & 9 & 4 & 3 & 53248 & 27256 & 13764\\
T11.0872 & \texttt{01234534123} & 3 & 6 & 2 & 1 & 11 & 11 & 3 & 3 & 57344 & 8560 & 24236\\
T11.0873 & \texttt{01234534122} & 4 & 6 & 3 & 2 & 10 & 9 & 3 & 2 & 55296 & 14776 & 20680\\
T11.0874 & \texttt{01234534121} & 4 & 6 & 2 & 1 & 11 & 11 & 4 & 4 & 53248 & 17920 & 18236\\
T11.0875 & \texttt{01234534111} & 5 & 6 & 2 & 6 & 9 & 8 & 3 & 2 & 51200 & 22600 & 14768\\
T11.0876 & \texttt{01234533333} & 6 & 6 & 1 & 120 & 7 & 6 & 2 & 1 & 61440 & 27776 & 21008\\
T11.0877 & \texttt{01234533332} & 6 & 6 & 2 & 24 & 8 & 7 & 3 & 2 & 59392 & 32304 & 16084\\
T11.0878 & \texttt{01234533331} & 6 & 6 & 2 & 24 & 8 & 7 & 3 & 2 & 57344 & 31840 & 15188\\
T11.0879 & \texttt{01234533323} & 5 & 6 & 2 & 6 & 9 & 8 & 3 & 2 & 59392 & 25500 & 19308\\
T11.0880 & \texttt{01234533322} & 6 & 6 & 2 & 12 & 8 & 6 & 3 & 1 & 57344 & 33168 & 13900\\
T11.0881 & \texttt{01234533321} & 6 & 6 & 3 & 6 & 9 & 8 & 4 & 3 & 55296 & 34368 & 11784\\
T11.0882 & \texttt{01234533312} & 5 & 6 & 2 & 6 & 9 & 8 & 3 & 2 & 55296 & 23752 & 17792\\
T11.0883 & \texttt{01234533311} & 6 & 6 & 2 & 12 & 8 & 6 & 3 & 1 & 53248 & 31712 & 12180\\
T11.0884 & \texttt{01234533234} & 4 & 6 & 2 & 2 & 10 & 9 & 3 & 2 & 61440 & 19380 & 22116\\
T11.0885 & \texttt{01234533233} & 5 & 6 & 3 & 4 & 9 & 7 & 3 & 1 & 59392 & 24872 & 18828\\
T11.0886 & \texttt{01234533232} & 5 & 6 & 2 & 2 & 10 & 9 & 4 & 3 & 57344 & 27804 & 16396\\
T11.0887 & \texttt{01234533231} & 5 & 6 & 3 & 2 & 10 & 9 & 4 & 3 & 55296 & 27808 & 15088\\
T11.0888 & \texttt{01234533222} & 6 & 6 & 2 & 12 & 8 & 6 & 3 & 1 & 55296 & 32016 & 13196\\
T11.0889 & \texttt{01234533221} & 6 & 6 & 3 & 4 & 9 & 7 & 4 & 2 & 53248 & 34048 & 10464\\
T11.0890 & \texttt{01234533212} & 5 & 6 & 3 & 2 & 10 & 9 & 4 & 3 & 53248 & 26120 & 14780\\
T11.0891 & \texttt{01234533211} & 6 & 6 & 3 & 4 & 9 & 7 & 4 & 2 & 51200 & 32800 & 9896\\
T11.0892 & \texttt{01234533123} & 4 & 6 & 2 & 2 & 10 & 9 & 3 & 2 & 55296 & 16792 & 20128\\
T11.0893 & \texttt{01234533122} & 5 & 6 & 3 & 4 & 9 & 7 & 3 & 1 & 53248 & 22448 & 16524\\
T11.0894 & \texttt{01234533121} & 5 & 6 & 2 & 2 & 10 & 9 & 4 & 3 & 51200 & 25032 & 14144\\
T11.0895 & \texttt{01234533111} & 6 & 6 & 2 & 12 & 8 & 6 & 3 & 1 & 49152 & 29152 & 10740\\
T11.0896 & \texttt{01234532345} & 3 & 6 & 2 & 1 & 11 & 11 & 3 & 3 & 65536 & 12292 & 25324\\
T11.0897 & \texttt{01234532344} & 4 & 6 & 3 & 2 & 10 & 9 & 3 & 2 & 63488 & 17080 & 23584\\
T11.0898 & \texttt{01234532343} & 4 & 6 & 3 & 1 & 11 & 11 & 4 & 4 & 61440 & 19564 & 21784\\
T11.0899 & \texttt{01234532342} & 4 & 6 & 2 & 1 & 11 & 11 & 4 & 4 & 59392 & 22228 & 19276\\
T11.0900 & \texttt{01234532341} & 4 & 6 & 3 & 1 & 11 & 11 & 4 & 4 & 57344 & 21088 & 18560\\
T11.0901 & \texttt{01234532333} & 5 & 6 & 3 & 6 & 9 & 8 & 3 & 2 & 59392 & 23200 & 19492\\
T11.0902 & \texttt{01234532332} & 5 & 6 & 3 & 2 & 10 & 9 & 4 & 3 & 57344 & 25992 & 16708\\
T11.0903 & \texttt{01234532331} & 5 & 6 & 4 & 2 & 10 & 9 & 4 & 3 & 55296 & 26048 & 15416\\
T11.0904 & \texttt{01234532323} & 4 & 6 & 2 & 2 & 9 & 7 & 3 & 2 & 57344 & 19744 & 20092\\
T11.0905 & \texttt{01234532322} & 5 & 6 & 2 & 2 & 10 & 9 & 4 & 3 & 55296 & 27196 & 15380\\
T11.0906 & \texttt{01234532321} & 5 & 6 & 3 & 1 & 11 & 11 & 5 & 5 & 53248 & 28448 & 13272\\
T11.0907 & \texttt{01234532312} & 4 & 6 & 3 & 1 & 11 & 11 & 4 & 4 & 53248 & 18000 & 18500\\
T11.0908 & \texttt{01234532311} & 5 & 6 & 3 & 2 & 10 & 9 & 4 & 3 & 51200 & 25880 & 13352\\
T11.0909 & \texttt{01234532222} & 6 & 6 & 2 & 24 & 8 & 7 & 3 & 2 & 53248 & 29680 & 13284\\
T11.0910 & \texttt{01234532221} & 6 & 6 & 3 & 6 & 9 & 8 & 4 & 3 & 51200 & 32128 & 10232\\
T11.0911 & \texttt{01234532212} & 5 & 6 & 3 & 2 & 10 & 9 & 4 & 3 & 51200 & 25544 & 13708\\
T11.0912 & \texttt{01234532211} & 6 & 6 & 3 & 4 & 9 & 7 & 4 & 2 & 49152 & 31584 & 9176\\
T11.0913 & \texttt{01234532123} & 4 & 6 & 3 & 1 & 11 & 11 & 4 & 4 & 53248 & 18648 & 17748\\
T11.0914 & \texttt{01234532122} & 5 & 6 & 4 & 2 & 10 & 9 & 4 & 3 & 51200 & 23664 & 14376\\
T11.0915 & \texttt{01234532121} & 5 & 6 & 3 & 1 & 11 & 11 & 5 & 5 & 49152 & 25608 & 12356\\
T11.0916 & \texttt{01234532111} & 6 & 6 & 3 & 6 & 9 & 8 & 4 & 3 & 47104 & 29088 & 9312\\
T11.0917 & \texttt{01234531234} & 3 & 6 & 2 & 1 & 11 & 11 & 3 & 3 & 57344 & 9592 & 22776\\
T11.0918 & \texttt{01234531233} & 4 & 6 & 3 & 2 & 10 & 9 & 3 & 2 & 55296 & 14352 & 20552\\
T11.0919 & \texttt{01234531232} & 4 & 6 & 3 & 1 & 11 & 11 & 4 & 4 & 53248 & 16552 & 18692\\
T11.0920 & \texttt{01234531231} & 4 & 6 & 2 & 1 & 11 & 11 & 4 & 4 & 51200 & 18392 & 16616\\
T11.0921 & \texttt{01234531222} & 5 & 6 & 3 & 6 & 9 & 8 & 3 & 2 & 51200 & 20032 & 16028\\
T11.0922 & \texttt{01234531221} & 5 & 6 & 3 & 2 & 10 & 9 & 4 & 3 & 49152 & 22128 & 13732\\
T11.0923 & \texttt{01234531212} & 4 & 6 & 2 & 2 & 9 & 7 & 3 & 2 & 49152 & 16192 & 16868\\
T11.0924 & \texttt{01234531211} & 5 & 6 & 2 & 2 & 10 & 9 & 4 & 3 & 47104 & 22792 & 12456\\
T11.0925 & \texttt{01234531111} & 6 & 6 & 2 & 24 & 8 & 7 & 3 & 2 & 45056 & 24992 & 10156\\
T11.0926 & \texttt{01234523452} & 3 & 6 & 1 & 2 & 7 & 3 & 2 & 1 & 63488 & 15444 & 22672\\
T11.0927 & \texttt{01234523451} & 3 & 6 & 2 & 2 & 7 & 3 & 2 & 1 & 61440 & 12640 & 22460\\
T11.0928 & \texttt{01234523444} & 4 & 6 & 2 & 6 & 9 & 8 & 2 & 1 & 65536 & 13104 & 26436\\
T11.0929 & \texttt{01234523443} & 4 & 6 & 3 & 2 & 10 & 9 & 3 & 2 & 63488 & 15896 & 23760\\
T11.0930 & \texttt{01234523442} & 4 & 6 & 2 & 2 & 10 & 9 & 3 & 2 & 61440 & 19368 & 21004\\
T11.0931 & \texttt{01234523441} & 4 & 6 & 3 & 2 & 10 & 9 & 3 & 2 & 59392 & 17344 & 20688\\
T11.0932 & \texttt{01234523434} & 3 & 6 & 2 & 2 & 9 & 7 & 2 & 1 & 63488 & 9840 & 26536\\
T11.0933 & \texttt{01234523433} & 4 & 6 & 2 & 2 & 10 & 9 & 3 & 2 & 61440 & 17196 & 22756\\
T11.0934 & \texttt{01234523432} & 4 & 6 & 2 & 1 & 11 & 11 & 4 & 4 & 59392 & 20988 & 19564\\
T11.0935 & \texttt{01234523431} & 4 & 6 & 3 & 1 & 11 & 11 & 4 & 4 & 57344 & 19744 & 18912\\
T11.0936 & \texttt{01234523423} & 3 & 6 & 1 & 1 & 11 & 11 & 3 & 3 & 59392 & 10584 & 24000\\
T11.0937 & \texttt{01234523422} & 4 & 6 & 1 & 2 & 10 & 9 & 3 & 2 & 57344 & 21924 & 18064\\
T11.0938 & \texttt{01234523421} & 4 & 6 & 2 & 1 & 11 & 11 & 4 & 4 & 55296 & 22240 & 16572\\
T11.0939 & \texttt{01234523412} & 3 & 6 & 2 & 1 & 11 & 11 & 3 & 3 & 55296 & 8768 & 22272\\
T11.0940 & \texttt{01234523411} & 4 & 6 & 2 & 2 & 10 & 9 & 3 & 2 & 53248 & 18088 & 17068\\
T11.0941 & \texttt{01234523333} & 5 & 6 & 2 & 24 & 8 & 7 & 2 & 1 & 59392 & 19648 & 21344\\
T11.0942 & \texttt{01234523332} & 5 & 6 & 2 & 6 & 9 & 8 & 3 & 2 & 57344 & 23760 & 17536\\
T11.0943 & \texttt{01234523331} & 5 & 6 & 3 & 6 & 9 & 8 & 3 & 2 & 55296 & 23296 & 16540\\
T11.0944 & \texttt{01234523323} & 4 & 6 & 2 & 2 & 10 & 9 & 3 & 2 & 57344 & 16724 & 20808\\
T11.0945 & \texttt{01234523322} & 5 & 6 & 2 & 4 & 9 & 7 & 3 & 1 & 55296 & 24824 & 15856\\
T11.0946 & \texttt{01234523321} & 5 & 6 & 3 & 2 & 10 & 9 & 4 & 3 & 53248 & 25920 & 13900\\
T11.0947 & \texttt{01234523312} & 4 & 6 & 3 & 2 & 10 & 9 & 3 & 2 & 53248 & 14968 & 19224\\
T11.0948 & \texttt{01234523311} & 5 & 6 & 3 & 4 & 9 & 7 & 3 & 1 & 51200 & 23088 & 14092\\
T11.0949 & \texttt{01234523232} & 4 & 6 & 1 & 2 & 9 & 7 & 3 & 2 & 55296 & 20304 & 18328\\
T11.0950 & \texttt{01234523231} & 4 & 6 & 2 & 2 & 9 & 7 & 3 & 2 & 53248 & 19840 & 17156\\
T11.0951 & \texttt{01234523222} & 5 & 6 & 1 & 6 & 9 & 8 & 3 & 2 & 53248 & 25164 & 15080\\
T11.0952 & \texttt{01234523221} & 5 & 6 & 2 & 2 & 10 & 9 & 4 & 3 & 51200 & 27040 & 12548\\
T11.0953 & \texttt{01234523212} & 4 & 6 & 2 & 1 & 11 & 11 & 4 & 4 & 51200 & 18608 & 16712\\
T11.0954 & \texttt{01234523211} & 5 & 6 & 2 & 2 & 10 & 9 & 4 & 3 & 49152 & 25528 & 12020\\
T11.0955 & \texttt{01234523123} & 3 & 6 & 2 & 1 & 11 & 11 & 3 & 3 & 53248 & 8368 & 21768\\
T11.0956 & \texttt{01234523122} & 4 & 6 & 3 & 2 & 10 & 9 & 3 & 2 & 51200 & 14264 & 18300\\
T11.0957 & \texttt{01234523121} & 4 & 6 & 2 & 1 & 11 & 11 & 4 & 4 & 49152 & 17088 & 16008\\
T11.0958 & \texttt{01234523111} & 5 & 6 & 2 & 6 & 9 & 8 & 3 & 2 & 47104 & 21448 & 12692\\
T11.0959 & \texttt{01234522222} & 6 & 6 & 1 & 120 & 7 & 6 & 2 & 1 & 51200 & 26784 & 13536\\
T11.0960 & \texttt{01234522221} & 6 & 6 & 2 & 24 & 8 & 7 & 3 & 2 & 49152 & 29440 & 10300\\
T11.0961 & \texttt{01234522212} & 5 & 6 & 2 & 6 & 9 & 8 & 3 & 2 & 49152 & 23528 & 13368\\
T11.0962 & \texttt{01234522211} & 6 & 6 & 2 & 12 & 8 & 6 & 3 & 1 & 47104 & 29248 & 8988\\
T11.0963 & \texttt{01234522123} & 4 & 6 & 2 & 2 & 10 & 9 & 3 & 2 & 51200 & 17848 & 16760\\
T11.0964 & \texttt{01234522122} & 5 & 6 & 3 & 4 & 9 & 7 & 3 & 1 & 49152 & 22544 & 13492\\
T11.0965 & \texttt{01234522121} & 5 & 6 & 2 & 2 & 10 & 9 & 4 & 3 & 47104 & 24168 & 11640\\
T11.0966 & \texttt{01234522111} & 6 & 6 & 2 & 12 & 8 & 6 & 3 & 1 & 45056 & 27328 & 8764\\
T11.0967 & \texttt{01234521234} & 3 & 6 & 2 & 1 & 11 & 11 & 3 & 3 & 55296 & 10968 & 20864\\
T11.0968 & \texttt{01234521233} & 4 & 6 & 3 & 2 & 10 & 9 & 3 & 2 & 53248 & 15408 & 18656\\
T11.0969 & \texttt{01234521232} & 4 & 6 & 3 & 1 & 11 & 11 & 4 & 4 & 51200 & 17288 & 16908\\
T11.0970 & \texttt{01234521231} & 4 & 6 & 2 & 1 & 11 & 11 & 4 & 4 & 49152 & 18488 & 15152\\
T11.0971 & \texttt{01234521222} & 5 & 6 & 3 & 6 & 9 & 8 & 3 & 2 & 49152 & 20448 & 14324\\
T11.0972 & \texttt{01234521221} & 5 & 6 & 3 & 2 & 10 & 9 & 4 & 3 & 47104 & 21904 & 12380\\
T11.0973 & \texttt{01234521212} & 4 & 6 & 2 & 2 & 9 & 7 & 3 & 2 & 47104 & 16608 & 15228\\
T11.0974 & \texttt{01234521211} & 5 & 6 & 2 & 2 & 10 & 9 & 4 & 3 & 45056 & 22248 & 11280\\
T11.0975 & \texttt{01234521111} & 6 & 6 & 2 & 24 & 8 & 7 & 3 & 2 & 43008 & 24128 & 9156\\
T11.0976 & \texttt{01234512345} & 2 & 6 & 1 & 2 & 6 & 1 & 1 & 0 & 61440 & 1240 & 25500\\
T11.0977 & \texttt{01234512344} & 3 & 6 & 2 & 2 & 10 & 9 & 2 & 1 & 59392 & 5680 & 24480\\
T11.0978 & \texttt{01234512343} & 3 & 6 & 2 & 1 & 11 & 11 & 3 & 3 & 57344 & 7816 & 23096\\
T11.0979 & \texttt{01234512342} & 3 & 6 & 2 & 1 & 11 & 11 & 3 & 3 & 55296 & 9784 & 21120\\
T11.0980 & \texttt{01234512341} & 3 & 6 & 1 & 1 & 11 & 11 & 3 & 3 & 53248 & 11160 & 19356\\
T11.0981 & \texttt{01234512333} & 4 & 6 & 2 & 6 & 9 & 8 & 2 & 1 & 55296 & 11104 & 21316\\
T11.0982 & \texttt{01234512332} & 4 & 6 & 3 & 2 & 10 & 9 & 3 & 2 & 53248 & 13200 & 18960\\
T11.0983 & \texttt{01234512331} & 4 & 6 & 2 & 2 & 10 & 9 & 3 & 2 & 51200 & 14960 & 17204\\
T11.0984 & \texttt{01234512323} & 3 & 6 & 2 & 2 & 9 & 7 & 2 & 1 & 53248 & 7648 & 21712\\
T11.0985 & \texttt{01234512322} & 4 & 6 & 2 & 2 & 10 & 9 & 3 & 2 & 51200 & 14056 & 17840\\
T11.0986 & \texttt{01234512321} & 4 & 6 & 2 & 1 & 11 & 11 & 4 & 4 & 49152 & 16200 & 15704\\
T11.0987 & \texttt{01234512312} & 3 & 6 & 1 & 1 & 11 & 11 & 3 & 3 & 49152 & 7440 & 19488\\
T11.0988 & \texttt{01234512311} & 4 & 6 & 1 & 2 & 10 & 9 & 3 & 2 & 47104 & 16120 & 14636\\
T11.0989 & \texttt{01234512222} & 5 & 6 & 2 & 24 & 8 & 7 & 2 & 1 & 49152 & 16192 & 15944\\
T11.0990 & \texttt{01234512221} & 5 & 6 & 2 & 6 & 9 & 8 & 3 & 2 & 47104 & 18720 & 13304\\
T11.0991 & \texttt{01234512212} & 4 & 6 & 2 & 2 & 10 & 9 & 3 & 2 & 47104 & 12696 & 16260\\
T11.0992 & \texttt{01234512211} & 5 & 6 & 2 & 4 & 9 & 7 & 3 & 1 & 45056 & 18896 & 12112\\
T11.0993 & \texttt{01234512121} & 4 & 6 & 1 & 2 & 9 & 7 & 3 & 2 & 45056 & 14880 & 14336\\
T11.0994 & \texttt{01234512111} & 5 & 6 & 1 & 6 & 9 & 8 & 3 & 2 & 43008 & 18600 & 11388\\
T11.0995 & \texttt{01234511111} & 6 & 6 & 1 & 120 & 7 & 6 & 2 & 1 & 40960 & 19840 & 9640\\
T11.0996 & \texttt{01234444444} & 7 & 5 & 1 & 5040 & 5 & 4 & 1 & 0 & 69632 & 26460 & 30720\\
T11.0997 & \texttt{01234444443} & 7 & 5 & 2 & 720 & 6 & 5 & 2 & 1 & 67584 & 31448 & 24548\\
T11.0998 & \texttt{01234444442} & 7 & 5 & 2 & 720 & 6 & 5 & 2 & 1 & 65536 & 33540 & 21772\\
T11.0999 & \texttt{01234444441} & 7 & 5 & 2 & 720 & 6 & 5 & 2 & 1 & 63488 & 31568 & 22052\\
T11.1000 & \texttt{01234444434} & 6 & 5 & 2 & 120 & 7 & 6 & 2 & 1 & 67584 & 25392 & 26896\\
T11.1001 & \texttt{01234444433} & 7 & 5 & 2 & 240 & 6 & 4 & 2 & 0 & 65536 & 32784 & 21580\\
T11.1002 & \texttt{01234444432} & 7 & 5 & 3 & 120 & 7 & 6 & 3 & 2 & 63488 & 36220 & 17512\\
T11.1003 & \texttt{01234444431} & 7 & 5 & 3 & 120 & 7 & 6 & 3 & 2 & 61440 & 35080 & 17152\\
T11.1004 & \texttt{01234444423} & 6 & 5 & 2 & 120 & 7 & 6 & 2 & 1 & 63488 & 24840 & 24492\\
T11.1005 & \texttt{01234444422} & 7 & 5 & 2 & 240 & 6 & 4 & 2 & 0 & 61440 & 35316 & 16948\\
T11.1006 & \texttt{01234444421} & 7 & 5 & 3 & 120 & 7 & 6 & 3 & 2 & 59392 & 35840 & 15296\\
T11.1007 & \texttt{01234444412} & 6 & 5 & 2 & 120 & 7 & 6 & 2 & 1 & 59392 & 23040 & 23096\\
T11.1008 & \texttt{01234444411} & 7 & 5 & 2 & 240 & 6 & 4 & 2 & 0 & 57344 & 32040 & 16628\\
T11.1009 & \texttt{01234444344} & 6 & 5 & 3 & 48 & 7 & 5 & 2 & 0 & 67584 & 25792 & 25156\\
T11.1010 & \texttt{01234444343} & 6 & 5 & 2 & 24 & 8 & 7 & 3 & 2 & 65536 & 28448 & 22912\\
T11.1011 & \texttt{01234444342} & 6 & 5 & 3 & 24 & 8 & 7 & 3 & 2 & 63488 & 31128 & 19668\\
T11.1012 & \texttt{01234444341} & 6 & 5 & 3 & 24 & 8 & 7 & 3 & 2 & 61440 & 29520 & 19540\\
T11.1013 & \texttt{01234444333} & 7 & 5 & 2 & 144 & 6 & 4 & 2 & 0 & 63488 & 32256 & 20284\\
T11.1014 & \texttt{01234444332} & 7 & 5 & 3 & 48 & 7 & 5 & 3 & 1 & 61440 & 36364 & 15552\\
T11.1015 & \texttt{01234444331} & 7 & 5 & 3 & 48 & 7 & 5 & 3 & 1 & 59392 & 35640 & 14872\\
T11.1016 & \texttt{01234444323} & 6 & 5 & 3 & 24 & 8 & 7 & 3 & 2 & 61440 & 27800 & 20368\\
T11.1017 & \texttt{01234444322} & 7 & 5 & 3 & 48 & 7 & 5 & 3 & 1 & 59392 & 36548 & 13952\\
T11.1018 & \texttt{01234444321} & 7 & 5 & 4 & 24 & 8 & 7 & 4 & 3 & 57344 & 37488 & 11972\\
T11.1019 & \texttt{01234444312} & 6 & 5 & 3 & 24 & 8 & 7 & 3 & 2 & 57344 & 26032 & 18908\\
T11.1020 & \texttt{01234444311} & 7 & 5 & 3 & 48 & 7 & 5 & 3 & 1 & 55296 & 34392 & 12800\\
T11.1021 & \texttt{01234444234} & 5 & 5 & 2 & 24 & 8 & 7 & 2 & 1 & 63488 & 18480 & 25620\\
T11.1022 & \texttt{01234444233} & 6 & 5 & 3 & 48 & 7 & 5 & 2 & 0 & 61440 & 25052 & 21684\\
T11.1023 & \texttt{01234444232} & 6 & 5 & 2 & 24 & 8 & 7 & 3 & 2 & 59392 & 29064 & 18364\\
T11.1024 & \texttt{01234444231} & 6 & 5 & 3 & 24 & 8 & 7 & 3 & 2 & 57344 & 28288 & 17584\\
T11.1025 & \texttt{01234444222} & 7 & 5 & 2 & 144 & 6 & 4 & 2 & 0 & 57344 & 34356 & 14276\\
T11.1026 & \texttt{01234444221} & 7 & 5 & 3 & 48 & 7 & 5 & 3 & 1 & 55296 & 36128 & 11640\\
T11.1027 & \texttt{01234444212} & 6 & 5 & 3 & 24 & 8 & 7 & 3 & 2 & 55296 & 27360 & 16764\\
T11.1028 & \texttt{01234444211} & 7 & 5 & 3 & 48 & 7 & 5 & 3 & 1 & 53248 & 34440 & 11424\\
T11.1029 & \texttt{01234444123} & 5 & 5 & 2 & 24 & 8 & 7 & 2 & 1 & 57344 & 16512 & 23368\\
T11.1030 & \texttt{01234444122} & 6 & 5 & 3 & 48 & 7 & 5 & 2 & 0 & 55296 & 22568 & 19468\\
T11.1031 & \texttt{01234444121} & 6 & 5 & 2 & 24 & 8 & 7 & 3 & 2 & 53248 & 25552 & 16712\\
T11.1032 & \texttt{01234444111} & 7 & 5 & 2 & 144 & 6 & 4 & 2 & 0 & 51200 & 30072 & 12932\\
T11.1033 & \texttt{01234443444} & 6 & 5 & 3 & 72 & 5 & 3 & 1 & 0 & 67584 & 26016 & 24604\\
T11.1034 & \texttt{01234443443} & 6 & 5 & 3 & 12 & 8 & 6 & 3 & 1 & 65536 & 27936 & 22100\\
T11.1035 & \texttt{01234443442} & 6 & 5 & 4 & 12 & 8 & 6 & 3 & 1 & 63488 & 30868 & 18760\\
T11.1036 & \texttt{01234443441} & 6 & 5 & 4 & 12 & 8 & 6 & 3 & 1 & 61440 & 29416 & 18512\\
T11.1037 & \texttt{01234443434} & 5 & 5 & 2 & 12 & 7 & 4 & 2 & 0 & 65536 & 22752 & 25060\\
T11.1038 & \texttt{01234443433} & 6 & 5 & 2 & 12 & 8 & 6 & 3 & 1 & 63488 & 28800 & 21200\\
T11.1039 & \texttt{01234443432} & 6 & 5 & 3 & 6 & 9 & 8 & 4 & 3 & 61440 & 32488 & 17076\\
T11.1040 & \texttt{01234443431} & 6 & 5 & 3 & 6 & 9 & 8 & 4 & 3 & 59392 & 31504 & 16532\\
T11.1041 & \texttt{01234443423} & 5 & 5 & 3 & 6 & 9 & 8 & 3 & 2 & 61440 & 22164 & 22636\\
T11.1042 & \texttt{01234443422} & 6 & 5 & 3 & 12 & 8 & 6 & 3 & 1 & 59392 & 31992 & 16028\\
T11.1043 & \texttt{01234443421} & 6 & 5 & 4 & 6 & 9 & 8 & 4 & 3 & 57344 & 32672 & 14296\\
T11.1044 & \texttt{01234443412} & 5 & 5 & 3 & 6 & 9 & 8 & 3 & 2 & 57344 & 20376 & 21152\\
T11.1045 & \texttt{01234443411} & 6 & 5 & 3 & 12 & 8 & 6 & 3 & 1 & 55296 & 29136 & 15204\\
T11.1046 & \texttt{01234443333} & 7 & 5 & 2 & 144 & 6 & 4 & 2 & 0 & 61440 & 30816 & 19916\\
T11.1047 & \texttt{01234443332} & 7 & 5 & 3 & 36 & 7 & 5 & 3 & 1 & 59392 & 35260 & 14816\\
T11.1048 & \texttt{01234443331} & 7 & 5 & 3 & 36 & 7 & 5 & 3 & 1 & 57344 & 34744 & 13976\\
T11.1049 & \texttt{01234443323} & 6 & 5 & 3 & 12 & 8 & 6 & 3 & 1 & 59392 & 28104 & 18568\\
T11.1050 & \texttt{01234443322} & 7 & 5 & 3 & 24 & 7 & 4 & 3 & 0 & 57344 & 35988 & 12680\\
T11.1051 & \texttt{01234443321} & 7 & 5 & 4 & 12 & 8 & 6 & 4 & 2 & 55296 & 37136 & 10524\\
T11.1052 & \texttt{01234443312} & 6 & 5 & 3 & 12 & 8 & 6 & 3 & 1 & 55296 & 26352 & 17076\\
T11.1053 & \texttt{01234443311} & 7 & 5 & 3 & 24 & 7 & 4 & 3 & 0 & 53248 & 34392 & 11112\\
T11.1054 & \texttt{01234443234} & 5 & 5 & 3 & 6 & 9 & 8 & 3 & 2 & 61440 & 21344 & 22104\\
T11.1055 & \texttt{01234443233} & 6 & 5 & 4 & 12 & 8 & 6 & 3 & 1 & 59392 & 27052 & 18528\\
T11.1056 & \texttt{01234443232} & 6 & 5 & 3 & 6 & 9 & 8 & 4 & 3 & 57344 & 30200 & 15760\\
T11.1057 & \texttt{01234443231} & 6 & 5 & 4 & 6 & 9 & 8 & 4 & 3 & 55296 & 30048 & 14556\\
T11.1058 & \texttt{01234443222} & 7 & 5 & 3 & 36 & 7 & 5 & 3 & 1 & 55296 & 34628 & 12224\\
T11.1059 & \texttt{01234443221} & 7 & 5 & 4 & 12 & 8 & 6 & 4 & 2 & 53248 & 36608 & 9404\\
T11.1060 & \texttt{01234443212} & 6 & 5 & 4 & 6 & 9 & 8 & 4 & 3 & 53248 & 28512 & 14120\\
T11.1061 & \texttt{01234443211} & 7 & 5 & 4 & 12 & 8 & 6 & 4 & 2 & 51200 & 35272 & 8932\\
T11.1062 & \texttt{01234443123} & 5 & 5 & 3 & 6 & 9 & 8 & 3 & 2 & 55296 & 18880 & 20076\\
T11.1063 & \texttt{01234443122} & 6 & 5 & 4 & 12 & 8 & 6 & 3 & 1 & 53248 & 24616 & 16280\\
T11.1064 & \texttt{01234443121} & 6 & 5 & 3 & 6 & 9 & 8 & 4 & 3 & 51200 & 27280 & 13692\\
T11.1065 & \texttt{01234443111} & 7 & 5 & 3 & 36 & 7 & 5 & 3 & 1 & 49152 & 31480 & 10080\\
T11.1066 & \texttt{01234442344} & 5 & 5 & 3 & 12 & 8 & 6 & 2 & 0 & 63488 & 18108 & 24552\\
T11.1067 & \texttt{01234442343} & 5 & 5 & 3 & 6 & 9 & 8 & 3 & 2 & 61440 & 20808 & 22536\\
T11.1068 & \texttt{01234442342} & 5 & 5 & 2 & 6 & 9 & 8 & 3 & 2 & 59392 & 23904 & 19524\\
T11.1069 & \texttt{01234442341} & 5 & 5 & 3 & 6 & 9 & 8 & 3 & 2 & 57344 & 22400 & 19096\\
T11.1070 & \texttt{01234442333} & 6 & 5 & 3 & 36 & 7 & 5 & 2 & 0 & 59392 & 24660 & 20052\\
T11.1071 & \texttt{01234442332} & 6 & 5 & 3 & 12 & 8 & 6 & 3 & 1 & 57344 & 27884 & 16740\\
T11.1072 & \texttt{01234442331} & 6 & 5 & 4 & 12 & 8 & 6 & 3 & 1 & 55296 & 27680 & 15624\\
T11.1073 & \texttt{01234442323} & 5 & 5 & 2 & 12 & 7 & 4 & 2 & 0 & 57344 & 21204 & 20604\\
T11.1074 & \texttt{01234442322} & 6 & 5 & 2 & 12 & 8 & 6 & 3 & 1 & 55296 & 29304 & 15148\\
T11.1075 & \texttt{01234442321} & 6 & 5 & 3 & 6 & 9 & 8 & 4 & 3 & 53248 & 30400 & 13072\\
T11.1076 & \texttt{01234442312} & 5 & 5 & 3 & 6 & 9 & 8 & 3 & 2 & 53248 & 19448 & 19044\\
T11.1077 & \texttt{01234442311} & 6 & 5 & 3 & 12 & 8 & 6 & 3 & 1 & 51200 & 27568 & 13456\\
T11.1078 & \texttt{01234442222} & 7 & 5 & 2 & 144 & 6 & 4 & 2 & 0 & 53248 & 32004 & 12788\\
T11.1079 & \texttt{01234442221} & 7 & 5 & 3 & 36 & 7 & 5 & 3 & 1 & 51200 & 34400 & 9624\\
T11.1080 & \texttt{01234442212} & 6 & 5 & 3 & 12 & 8 & 6 & 3 & 1 & 51200 & 27648 & 13428\\
T11.1081 & \texttt{01234442211} & 7 & 5 & 3 & 24 & 7 & 4 & 3 & 0 & 49152 & 33768 & 8616\\
T11.1082 & \texttt{01234442123} & 5 & 5 & 3 & 6 & 9 & 8 & 3 & 2 & 53248 & 20448 & 17932\\
T11.1083 & \texttt{01234442122} & 6 & 5 & 4 & 12 & 8 & 6 & 3 & 1 & 51200 & 25544 & 14392\\
T11.1084 & \texttt{01234442121} & 6 & 5 & 3 & 6 & 9 & 8 & 4 & 3 & 49152 & 27568 & 12188\\
T11.1085 & \texttt{01234442111} & 7 & 5 & 3 & 36 & 7 & 5 & 3 & 1 & 47104 & 31128 & 8960\\
T11.1086 & \texttt{01234441234} & 4 & 5 & 2 & 6 & 9 & 8 & 2 & 1 & 57344 & 10848 & 23720\\
T11.1087 & \texttt{01234441233} & 5 & 5 & 3 & 12 & 8 & 6 & 2 & 0 & 55296 & 15688 & 21400\\
T11.1088 & \texttt{01234441232} & 5 & 5 & 3 & 6 & 9 & 8 & 3 & 2 & 53248 & 17968 & 19420\\
T11.1089 & \texttt{01234441231} & 5 & 5 & 2 & 6 & 9 & 8 & 3 & 2 & 51200 & 19968 & 17080\\
T11.1090 & \texttt{01234441222} & 6 & 5 & 3 & 36 & 7 & 5 & 2 & 0 & 51200 & 21528 & 16644\\
T11.1091 & \texttt{01234441221} & 6 & 5 & 3 & 12 & 8 & 6 & 3 & 1 & 49152 & 23784 & 14076\\
T11.1092 & \texttt{01234441212} & 5 & 5 & 2 & 12 & 7 & 4 & 2 & 0 & 49152 & 17688 & 17468\\
T11.1093 & \texttt{01234441211} & 6 & 5 & 2 & 12 & 8 & 6 & 3 & 1 & 47104 & 24528 & 12664\\
T11.1094 & \texttt{01234441111} & 7 & 5 & 2 & 144 & 6 & 4 & 2 & 0 & 45056 & 26808 & 10228\\
T11.1095 & \texttt{01234434434} & 5 & 5 & 3 & 8 & 7 & 5 & 2 & 1 & 65536 & 22464 & 24464\\
T11.1096 & \texttt{01234434433} & 6 & 5 & 3 & 16 & 6 & 3 & 2 & 0 & 63488 & 28064 & 20940\\
T11.1097 & \texttt{01234434432} & 6 & 5 & 4 & 8 & 7 & 5 & 3 & 2 & 61440 & 31836 & 16832\\
T11.1098 & \texttt{01234434431} & 6 & 5 & 4 & 8 & 7 & 5 & 3 & 2 & 59392 & 30904 & 16248\\
T11.1099 & \texttt{01234434423} & 5 & 5 & 4 & 8 & 7 & 5 & 2 & 1 & 61440 & 21864 & 22072\\
T11.1100 & \texttt{01234434422} & 6 & 5 & 4 & 16 & 6 & 3 & 2 & 0 & 59392 & 31476 & 15704\\
T11.1101 & \texttt{01234434421} & 6 & 5 & 5 & 8 & 7 & 5 & 3 & 2 & 57344 & 32208 & 13964\\
T11.1102 & \texttt{01234434412} & 5 & 5 & 4 & 8 & 7 & 5 & 2 & 1 & 57344 & 20080 & 20580\\
T11.1103 & \texttt{01234434411} & 6 & 5 & 4 & 16 & 6 & 3 & 2 & 0 & 55296 & 28760 & 14776\\
T11.1104 & \texttt{01234434343} & 5 & 5 & 2 & 4 & 8 & 5 & 3 & 1 & 63488 & 24672 & 22780\\
T11.1105 & \texttt{01234434342} & 5 & 5 & 3 & 4 & 8 & 5 & 3 & 1 & 61440 & 27604 & 19376\\
T11.1106 & \texttt{01234434341} & 5 & 5 & 3 & 4 & 8 & 5 & 3 & 1 & 59392 & 26152 & 19048\\
T11.1107 & \texttt{01234434333} & 6 & 5 & 2 & 12 & 8 & 6 & 3 & 1 & 61440 & 28032 & 20712\\
T11.1108 & \texttt{01234434332} & 6 & 5 & 3 & 4 & 9 & 7 & 4 & 2 & 59392 & 32224 & 16076\\
T11.1109 & \texttt{01234434331} & 6 & 5 & 3 & 4 & 9 & 7 & 4 & 2 & 57344 & 31552 & 15324\\
T11.1110 & \texttt{01234434323} & 5 & 5 & 3 & 2 & 10 & 9 & 4 & 3 & 59392 & 24012 & 20276\\
T11.1111 & \texttt{01234434322} & 6 & 5 & 3 & 4 & 9 & 7 & 4 & 2 & 57344 & 32544 & 14340\\
T11.1112 & \texttt{01234434321} & 6 & 5 & 4 & 2 & 10 & 9 & 5 & 4 & 55296 & 33536 & 12368\\
T11.1113 & \texttt{01234434312} & 5 & 5 & 3 & 2 & 10 & 9 & 4 & 3 & 55296 & 22248 & 18776\\
T11.1114 & \texttt{01234434311} & 6 & 5 & 3 & 4 & 9 & 7 & 4 & 2 & 53248 & 30528 & 12988\\
T11.1115 & \texttt{01234434234} & 4 & 5 & 3 & 2 & 10 & 9 & 3 & 2 & 61440 & 15332 & 24588\\
T11.1116 & \texttt{01234434233} & 5 & 5 & 4 & 4 & 9 & 7 & 3 & 1 & 59392 & 21688 & 20964\\
T11.1117 & \texttt{01234434232} & 5 & 5 & 3 & 2 & 10 & 9 & 4 & 3 & 57344 & 25484 & 18004\\
T11.1118 & \texttt{01234434231} & 5 & 5 & 4 & 2 & 10 & 9 & 4 & 3 & 55296 & 24864 & 17064\\
T11.1119 & \texttt{01234434222} & 6 & 5 & 3 & 12 & 8 & 6 & 3 & 1 & 55296 & 30560 & 14276\\
T11.1120 & \texttt{01234434221} & 6 & 5 & 4 & 4 & 9 & 7 & 4 & 2 & 53248 & 32384 & 11688\\
T11.1121 & \texttt{01234434212} & 5 & 5 & 4 & 2 & 10 & 9 & 4 & 3 & 53248 & 23784 & 16404\\
T11.1122 & \texttt{01234434211} & 6 & 5 & 4 & 4 & 9 & 7 & 4 & 2 & 51200 & 30784 & 11344\\
T11.1123 & \texttt{01234434123} & 4 & 5 & 3 & 2 & 10 & 9 & 3 & 2 & 55296 & 13240 & 22360\\
T11.1124 & \texttt{01234434122} & 5 & 5 & 4 & 4 & 9 & 7 & 3 & 1 & 53248 & 19216 & 18644\\
T11.1125 & \texttt{01234434121} & 5 & 5 & 3 & 2 & 10 & 9 & 4 & 3 & 51200 & 22120 & 16088\\
T11.1126 & \texttt{01234434111} & 6 & 5 & 3 & 12 & 8 & 6 & 3 & 1 & 49152 & 26560 & 12508\\
T11.1127 & \texttt{01234433333} & 7 & 5 & 2 & 240 & 6 & 4 & 2 & 0 & 59392 & 29152 & 20100\\
T11.1128 & \texttt{01234433332} & 7 & 5 & 3 & 48 & 7 & 5 & 3 & 1 & 57344 & 33764 & 14744\\
T11.1129 & \texttt{01234433331} & 7 & 5 & 3 & 48 & 7 & 5 & 3 & 1 & 55296 & 33352 & 13824\\
T11.1130 & \texttt{01234433323} & 6 & 5 & 3 & 12 & 8 & 6 & 3 & 1 & 57344 & 27312 & 18000\\
T11.1131 & \texttt{01234433322} & 7 & 5 & 3 & 24 & 7 & 4 & 3 & 0 & 55296 & 34764 & 12304\\
T11.1132 & \texttt{01234433321} & 7 & 5 & 4 & 12 & 8 & 6 & 4 & 2 & 53248 & 36016 & 10036\\
T11.1133 & \texttt{01234433312} & 6 & 5 & 3 & 12 & 8 & 6 & 3 & 1 & 53248 & 25568 & 16492\\
T11.1134 & \texttt{01234433311} & 7 & 5 & 3 & 24 & 7 & 4 & 3 & 0 & 51200 & 33448 & 10528\\
T11.1135 & \texttt{01234433234} & 5 & 5 & 3 & 4 & 9 & 7 & 3 & 1 & 59392 & 21832 & 20696\\
T11.1136 & \texttt{01234433233} & 6 & 5 & 4 & 8 & 8 & 5 & 3 & 0 & 57344 & 27108 & 17264\\
T11.1137 & \texttt{01234433232} & 6 & 5 & 3 & 4 & 9 & 7 & 4 & 2 & 55296 & 29824 & 14736\\
T11.1138 & \texttt{01234433231} & 6 & 5 & 4 & 4 & 9 & 7 & 4 & 2 & 53248 & 29984 & 13308\\
T11.1139 & \texttt{01234433222} & 7 & 5 & 3 & 24 & 7 & 4 & 3 & 0 & 53248 & 33820 & 11440\\
T11.1140 & \texttt{01234433221} & 7 & 5 & 4 & 8 & 8 & 5 & 4 & 1 & 51200 & 35904 & 8492\\
T11.1141 & \texttt{01234433212} & 6 & 5 & 4 & 4 & 9 & 7 & 4 & 2 & 51200 & 28144 & 13064\\
T11.1142 & \texttt{01234433211} & 7 & 5 & 4 & 8 & 8 & 5 & 4 & 1 & 49152 & 34744 & 7892\\
T11.1143 & \texttt{01234433123} & 5 & 5 & 3 & 4 & 9 & 7 & 3 & 1 & 53248 & 19120 & 18780\\
T11.1144 & \texttt{01234433122} & 6 & 5 & 4 & 8 & 8 & 5 & 3 & 0 & 51200 & 24696 & 15000\\
T11.1145 & \texttt{01234433121} & 6 & 5 & 3 & 4 & 9 & 7 & 4 & 2 & 49152 & 27200 & 12460\\
T11.1146 & \texttt{01234433111} & 7 & 5 & 3 & 24 & 7 & 4 & 3 & 0 & 47104 & 31240 & 8896\\
T11.1147 & \texttt{01234432344} & 5 & 5 & 4 & 4 & 9 & 7 & 3 & 1 & 61440 & 20468 & 21900\\
T11.1148 & \texttt{01234432343} & 5 & 5 & 4 & 2 & 10 & 9 & 4 & 3 & 59392 & 22736 & 20052\\
T11.1149 & \texttt{01234432342} & 5 & 5 & 3 & 2 & 10 & 9 & 4 & 3 & 57344 & 24968 & 17520\\
T11.1150 & \texttt{01234432341} & 5 & 5 & 4 & 2 & 10 & 9 & 4 & 3 & 55296 & 24192 & 16684\\
T11.1151 & \texttt{01234432333} & 6 & 5 & 4 & 12 & 8 & 6 & 3 & 1 & 57344 & 26156 & 17688\\
T11.1152 & \texttt{01234432332} & 6 & 5 & 4 & 4 & 9 & 7 & 4 & 2 & 55296 & 28516 & 14904\\
T11.1153 & \texttt{01234432331} & 6 & 5 & 5 & 4 & 9 & 7 & 4 & 2 & 53248 & 28832 & 13428\\
T11.1154 & \texttt{01234432323} & 5 & 5 & 3 & 4 & 8 & 5 & 3 & 1 & 55296 & 22700 & 18336\\
T11.1155 & \texttt{01234432322} & 6 & 5 & 3 & 4 & 9 & 7 & 4 & 2 & 53248 & 29504 & 13576\\
T11.1156 & \texttt{01234432321} & 6 & 5 & 4 & 2 & 10 & 9 & 5 & 4 & 51200 & 30912 & 11252\\
T11.1157 & \texttt{01234432312} & 5 & 5 & 4 & 2 & 10 & 9 & 4 & 3 & 51200 & 20968 & 16752\\
T11.1158 & \texttt{01234432311} & 6 & 5 & 4 & 4 & 9 & 7 & 4 & 2 & 49152 & 28608 & 11332\\
T11.1159 & \texttt{01234432222} & 7 & 5 & 3 & 48 & 7 & 5 & 3 & 1 & 51200 & 31772 & 11480\\
T11.1160 & \texttt{01234432221} & 7 & 5 & 4 & 12 & 8 & 6 & 4 & 2 & 49152 & 34272 & 8180\\
T11.1161 & \texttt{01234432212} & 6 & 5 & 4 & 4 & 9 & 7 & 4 & 2 & 49152 & 27856 & 11816\\
T11.1162 & \texttt{01234432211} & 7 & 5 & 4 & 8 & 8 & 5 & 4 & 1 & 47104 & 33816 & 7028\\
T11.1163 & \texttt{01234432123} & 5 & 5 & 4 & 2 & 10 & 9 & 4 & 3 & 51200 & 21264 & 16032\\
T11.1164 & \texttt{01234432122} & 6 & 5 & 5 & 4 & 9 & 7 & 4 & 2 & 49152 & 26200 & 12516\\
T11.1165 & \texttt{01234432121} & 6 & 5 & 4 & 2 & 10 & 9 & 5 & 4 & 47104 & 28064 & 10368\\
T11.1166 & \texttt{01234432111} & 7 & 5 & 4 & 12 & 8 & 6 & 4 & 2 & 45056 & 31464 & 7196\\
T11.1167 & \texttt{01234431234} & 4 & 5 & 3 & 2 & 10 & 9 & 3 & 2 & 55296 & 12752 & 21412\\
T11.1168 & \texttt{01234431233} & 5 & 5 & 4 & 4 & 9 & 7 & 3 & 1 & 53248 & 17432 & 19084\\
T11.1169 & \texttt{01234431232} & 5 & 5 & 4 & 2 & 10 & 9 & 4 & 3 & 51200 & 19552 & 17144\\
T11.1170 & \texttt{01234431231} & 5 & 5 & 3 & 2 & 10 & 9 & 4 & 3 & 49152 & 21232 & 14932\\
T11.1171 & \texttt{01234431222} & 6 & 5 & 4 & 12 & 8 & 6 & 3 & 1 & 49152 & 22952 & 14392\\
T11.1172 & \texttt{01234431221} & 6 & 5 & 4 & 4 & 9 & 7 & 4 & 2 & 47104 & 24888 & 11968\\
T11.1173 & \texttt{01234431212} & 5 & 5 & 3 & 4 & 8 & 5 & 3 & 1 & 47104 & 19112 & 15248\\
T11.1174 & \texttt{01234431211} & 6 & 5 & 3 & 4 & 9 & 7 & 4 & 2 & 45056 & 25472 & 10628\\
T11.1175 & \texttt{01234431111} & 7 & 5 & 3 & 48 & 7 & 5 & 3 & 1 & 43008 & 27592 & 8264\\
T11.1176 & \texttt{01234423442} & 5 & 5 & 3 & 8 & 6 & 3 & 2 & 1 & 59392 & 23076 & 19264\\
T11.1177 & \texttt{01234423441} & 5 & 5 & 4 & 8 & 6 & 3 & 2 & 1 & 57344 & 21728 & 18764\\
T11.1178 & \texttt{01234423434} & 4 & 5 & 3 & 4 & 8 & 5 & 2 & 0 & 61440 & 14844 & 24724\\
T11.1179 & \texttt{01234423433} & 5 & 5 & 3 & 4 & 9 & 7 & 3 & 1 & 59392 & 21552 & 20872\\
T11.1180 & \texttt{01234423432} & 5 & 5 & 3 & 2 & 10 & 9 & 4 & 3 & 57344 & 24480 & 17824\\
T11.1181 & \texttt{01234423431} & 5 & 5 & 4 & 2 & 10 & 9 & 4 & 3 & 55296 & 23808 & 16956\\
T11.1182 & \texttt{01234423423} & 4 & 5 & 2 & 2 & 10 & 9 & 3 & 2 & 57344 & 15156 & 22236\\
T11.1183 & \texttt{01234423422} & 5 & 5 & 2 & 4 & 9 & 7 & 3 & 1 & 55296 & 24984 & 16396\\
T11.1184 & \texttt{01234423421} & 5 & 5 & 3 & 2 & 10 & 9 & 4 & 3 & 53248 & 25664 & 14672\\
T11.1185 & \texttt{01234423412} & 4 & 5 & 3 & 2 & 10 & 9 & 3 & 2 & 53248 & 13368 & 20628\\
T11.1186 & \texttt{01234423411} & 5 & 5 & 3 & 4 & 9 & 7 & 3 & 1 & 51200 & 22128 & 15200\\
T11.1187 & \texttt{01234423333} & 6 & 5 & 3 & 48 & 7 & 5 & 2 & 0 & 57344 & 23788 & 19388\\
T11.1188 & \texttt{01234423332} & 6 & 5 & 3 & 12 & 8 & 6 & 3 & 1 & 55296 & 27036 & 15772\\
T11.1189 & \texttt{01234423331} & 6 & 5 & 4 & 12 & 8 & 6 & 3 & 1 & 53248 & 27040 & 14512\\
T11.1190 & \texttt{01234423323} & 5 & 5 & 3 & 4 & 9 & 7 & 3 & 1 & 55296 & 20648 & 18996\\
T11.1191 & \texttt{01234423322} & 6 & 5 & 3 & 8 & 8 & 5 & 3 & 0 & 53248 & 27668 & 14188\\
T11.1192 & \texttt{01234423321} & 6 & 5 & 4 & 4 & 9 & 7 & 4 & 2 & 51200 & 29024 & 11968\\
T11.1193 & \texttt{01234423312} & 5 & 5 & 4 & 4 & 9 & 7 & 3 & 1 & 51200 & 18912 & 17436\\
T11.1194 & \texttt{01234423311} & 6 & 5 & 4 & 8 & 8 & 5 & 3 & 0 & 49152 & 26632 & 12096\\
T11.1195 & \texttt{01234423232} & 5 & 5 & 2 & 4 & 8 & 5 & 3 & 1 & 53248 & 23580 & 16612\\
T11.1196 & \texttt{01234423231} & 5 & 5 & 3 & 4 & 8 & 5 & 3 & 1 & 51200 & 23584 & 15208\\
T11.1197 & \texttt{01234423222} & 6 & 5 & 2 & 12 & 8 & 6 & 3 & 1 & 51200 & 27792 & 13460\\
T11.1198 & \texttt{01234423221} & 6 & 5 & 3 & 4 & 9 & 7 & 4 & 2 & 49152 & 29824 & 10664\\
T11.1199 & \texttt{01234423212} & 5 & 5 & 3 & 2 & 10 & 9 & 4 & 3 & 49152 & 21896 & 14956\\
T11.1200 & \texttt{01234423211} & 6 & 5 & 3 & 4 & 9 & 7 & 4 & 2 & 47104 & 28576 & 10040\\
T11.1201 & \texttt{01234423123} & 4 & 5 & 3 & 2 & 10 & 9 & 3 & 2 & 51200 & 12568 & 20220\\
T11.1202 & \texttt{01234423122} & 5 & 5 & 4 & 4 & 9 & 7 & 3 & 1 & 49152 & 18224 & 16608\\
T11.1203 & \texttt{01234423121} & 5 & 5 & 3 & 2 & 10 & 9 & 4 & 3 & 47104 & 20808 & 14220\\
T11.1204 & \texttt{01234423111} & 6 & 5 & 3 & 12 & 8 & 6 & 3 & 1 & 45056 & 24928 & 10808\\
T11.1205 & \texttt{01234422222} & 7 & 5 & 2 & 240 & 6 & 4 & 2 & 0 & 49152 & 29196 & 11964\\
T11.1206 & \texttt{01234422221} & 7 & 5 & 3 & 48 & 7 & 5 & 3 & 1 & 47104 & 31904 & 8464\\
T11.1207 & \texttt{01234422212} & 6 & 5 & 3 & 12 & 8 & 6 & 3 & 1 & 47104 & 26160 & 11644\\
T11.1208 & \texttt{01234422211} & 7 & 5 & 3 & 24 & 7 & 4 & 3 & 0 & 45056 & 31800 & 7024\\
T11.1209 & \texttt{01234422123} & 5 & 5 & 3 & 4 & 9 & 7 & 3 & 1 & 49152 & 20784 & 15116\\
T11.1210 & \texttt{01234422122} & 6 & 5 & 4 & 8 & 8 & 5 & 3 & 0 & 47104 & 25400 & 11720\\
T11.1211 & \texttt{01234422121} & 6 & 5 & 3 & 4 & 9 & 7 & 4 & 2 & 45056 & 26944 & 9756\\
T11.1212 & \texttt{01234422111} & 7 & 5 & 3 & 24 & 7 & 4 & 3 & 0 & 43008 & 30024 & 6768\\
T11.1213 & \texttt{01234421234} & 4 & 5 & 3 & 2 & 10 & 9 & 3 & 2 & 53248 & 14448 & 19380\\
T11.1214 & \texttt{01234421233} & 5 & 5 & 4 & 4 & 9 & 7 & 3 & 1 & 51200 & 18808 & 17084\\
T11.1215 & \texttt{01234421232} & 5 & 5 & 4 & 2 & 10 & 9 & 4 & 3 & 49152 & 20608 & 15272\\
T11.1216 & \texttt{01234421231} & 5 & 5 & 3 & 2 & 10 & 9 & 4 & 3 & 47104 & 21648 & 13412\\
T11.1217 & \texttt{01234421222} & 6 & 5 & 4 & 12 & 8 & 6 & 3 & 1 & 47104 & 23688 & 12616\\
T11.1218 & \texttt{01234421221} & 6 & 5 & 4 & 4 & 9 & 7 & 4 & 2 & 45056 & 24984 & 10576\\
T11.1219 & \texttt{01234421212} & 5 & 5 & 3 & 4 & 8 & 5 & 3 & 1 & 45056 & 19848 & 13536\\
T11.1220 & \texttt{01234421211} & 6 & 5 & 3 & 4 & 9 & 7 & 4 & 2 & 43008 & 25248 & 9428\\
T11.1221 & \texttt{01234421111} & 7 & 5 & 3 & 48 & 7 & 5 & 3 & 1 & 40960 & 27048 & 7256\\
T11.1222 & \texttt{01234412344} & 4 & 5 & 3 & 8 & 5 & 1 & 1 & 0 & 57344 & 10040 & 23364\\
T11.1223 & \texttt{01234412343} & 4 & 5 & 3 & 2 & 10 & 9 & 3 & 2 & 55296 & 12096 & 21924\\
T11.1224 & \texttt{01234412342} & 4 & 5 & 3 & 2 & 10 & 9 & 3 & 2 & 53248 & 13904 & 19884\\
T11.1225 & \texttt{01234412341} & 4 & 5 & 2 & 2 & 10 & 9 & 3 & 2 & 51200 & 14960 & 18016\\
T11.1226 & \texttt{01234412333} & 5 & 5 & 3 & 12 & 8 & 6 & 2 & 0 & 53248 & 15304 & 20064\\
T11.1227 & \texttt{01234412332} & 5 & 5 & 4 & 4 & 9 & 7 & 3 & 1 & 51200 & 17240 & 17652\\
T11.1228 & \texttt{01234412331} & 5 & 5 & 3 & 4 & 9 & 7 & 3 & 1 & 49152 & 18680 & 15808\\
T11.1229 & \texttt{01234412323} & 4 & 5 & 3 & 4 & 8 & 5 & 2 & 0 & 51200 & 11848 & 20484\\
T11.1230 & \texttt{01234412322} & 5 & 5 & 3 & 4 & 9 & 7 & 3 & 1 & 49152 & 18016 & 16492\\
T11.1231 & \texttt{01234412321} & 5 & 5 & 3 & 2 & 10 & 9 & 4 & 3 & 47104 & 19840 & 14276\\
T11.1232 & \texttt{01234412312} & 4 & 5 & 2 & 2 & 10 & 9 & 3 & 2 & 47104 & 11480 & 18212\\
T11.1233 & \texttt{01234412311} & 5 & 5 & 2 & 4 & 9 & 7 & 3 & 1 & 45056 & 19600 & 13168\\
T11.1234 & \texttt{01234412222} & 6 & 5 & 3 & 48 & 7 & 5 & 2 & 0 & 47104 & 20072 & 14540\\
T11.1235 & \texttt{01234412221} & 6 & 5 & 3 & 12 & 8 & 6 & 3 & 1 & 45056 & 22280 & 11836\\
T11.1236 & \texttt{01234412212} & 5 & 5 & 3 & 4 & 9 & 7 & 3 & 1 & 45056 & 16496 & 14872\\
T11.1237 & \texttt{01234412211} & 6 & 5 & 3 & 8 & 8 & 5 & 3 & 0 & 43008 & 22296 & 10612\\
T11.1238 & \texttt{01234412121} & 5 & 5 & 2 & 4 & 8 & 5 & 3 & 1 & 43008 & 18440 & 12884\\
T11.1239 & \texttt{01234412111} & 6 & 5 & 2 & 12 & 8 & 6 & 3 & 1 & 40960 & 21920 & 9872\\
T11.1240 & \texttt{01234411111} & 7 & 5 & 2 & 240 & 6 & 4 & 2 & 0 & 38912 & 23080 & 8108\\
T11.1241 & \texttt{01234343434} & 4 & 5 & 1 & 24 & 5 & 3 & 1 & 0 & 63488 & 19488 & 25516\\
T11.1242 & \texttt{01234343433} & 5 & 5 & 1 & 12 & 6 & 3 & 2 & 0 & 61440 & 25536 & 21992\\
T11.1243 & \texttt{01234343432} & 5 & 5 & 2 & 6 & 7 & 5 & 3 & 2 & 59392 & 29224 & 17868\\
T11.1244 & \texttt{01234343431} & 5 & 5 & 2 & 6 & 7 & 5 & 3 & 2 & 57344 & 28240 & 17276\\
T11.1245 & \texttt{01234343423} & 4 & 5 & 2 & 6 & 7 & 5 & 2 & 1 & 59392 & 18900 & 23044\\
T11.1246 & \texttt{01234343422} & 5 & 5 & 2 & 12 & 6 & 3 & 2 & 0 & 57344 & 28728 & 16676\\
T11.1247 & \texttt{01234343421} & 5 & 5 & 3 & 6 & 7 & 5 & 3 & 2 & 55296 & 29408 & 14928\\
T11.1248 & \texttt{01234343412} & 4 & 5 & 2 & 6 & 7 & 5 & 2 & 1 & 55296 & 17112 & 21512\\
T11.1249 & \texttt{01234343411} & 5 & 5 & 2 & 12 & 6 & 3 & 2 & 0 & 53248 & 25872 & 15708\\
T11.1250 & \texttt{01234343333} & 6 & 5 & 1 & 48 & 6 & 3 & 2 & 0 & 59392 & 27552 & 20820\\
T11.1251 & \texttt{01234343332} & 6 & 5 & 2 & 12 & 7 & 4 & 3 & 1 & 57344 & 31996 & 15784\\
T11.1252 & \texttt{01234343331} & 6 & 5 & 2 & 12 & 7 & 4 & 3 & 1 & 55296 & 31480 & 14928\\
T11.1253 & \texttt{01234343323} & 5 & 5 & 2 & 4 & 8 & 5 & 3 & 1 & 57344 & 24840 & 19376\\
T11.1254 & \texttt{01234343322} & 6 & 5 & 2 & 8 & 7 & 3 & 3 & 0 & 55296 & 32724 & 13632\\
T11.1255 & \texttt{01234343321} & 6 & 5 & 3 & 4 & 8 & 5 & 4 & 2 & 53248 & 33872 & 11492\\
T11.1256 & \texttt{01234343312} & 5 & 5 & 2 & 4 & 8 & 5 & 3 & 1 & 53248 & 23088 & 17868\\
T11.1257 & \texttt{01234343311} & 6 & 5 & 2 & 8 & 7 & 3 & 3 & 0 & 51200 & 31128 & 12016\\
T11.1258 & \texttt{01234343234} & 4 & 5 & 2 & 2 & 9 & 7 & 3 & 2 & 59392 & 18080 & 22664\\
T11.1259 & \texttt{01234343233} & 5 & 5 & 3 & 4 & 8 & 5 & 3 & 1 & 57344 & 23788 & 19184\\
T11.1260 & \texttt{01234343232} & 5 & 5 & 2 & 2 & 9 & 7 & 4 & 3 & 55296 & 26936 & 16512\\
T11.1261 & \texttt{01234343231} & 5 & 5 & 3 & 2 & 9 & 7 & 4 & 3 & 53248 & 26784 & 15276\\
T11.1262 & \texttt{01234343222} & 6 & 5 & 2 & 12 & 7 & 4 & 3 & 1 & 53248 & 31364 & 13072\\
T11.1263 & \texttt{01234343221} & 6 & 5 & 3 & 4 & 8 & 5 & 4 & 2 & 51200 & 33344 & 10284\\
T11.1264 & \texttt{01234343212} & 5 & 5 & 3 & 2 & 9 & 7 & 4 & 3 & 51200 & 25248 & 14872\\
T11.1265 & \texttt{01234343211} & 6 & 5 & 3 & 4 & 8 & 5 & 4 & 2 & 49152 & 32008 & 9780\\
T11.1266 & \texttt{01234343123} & 4 & 5 & 2 & 2 & 9 & 7 & 3 & 2 & 53248 & 15616 & 20620\\
T11.1267 & \texttt{01234343122} & 5 & 5 & 3 & 4 & 8 & 5 & 3 & 1 & 51200 & 21352 & 16888\\
T11.1268 & \texttt{01234343121} & 5 & 5 & 2 & 2 & 9 & 7 & 4 & 3 & 49152 & 24016 & 14364\\
T11.1269 & \texttt{01234343111} & 6 & 5 & 2 & 12 & 7 & 4 & 3 & 1 & 47104 & 28216 & 10816\\
T11.1270 & \texttt{01234342343} & 4 & 5 & 3 & 2 & 9 & 7 & 3 & 2 & 59392 & 17544 & 22780\\
T11.1271 & \texttt{01234342342} & 4 & 5 & 2 & 2 & 9 & 7 & 3 & 2 & 57344 & 20640 & 19912\\
T11.1272 & \texttt{01234342341} & 4 & 5 & 3 & 2 & 9 & 7 & 3 & 2 & 55296 & 19136 & 19380\\
T11.1273 & \texttt{01234342333} & 5 & 5 & 3 & 12 & 7 & 4 & 2 & 0 & 57344 & 21396 & 20368\\
T11.1274 & \texttt{01234342332} & 5 & 5 & 3 & 4 & 8 & 5 & 3 & 1 & 55296 & 24620 & 17200\\
T11.1275 & \texttt{01234342331} & 5 & 5 & 4 & 4 & 8 & 5 & 3 & 1 & 53248 & 24416 & 16028\\
T11.1276 & \texttt{01234342323} & 4 & 5 & 2 & 4 & 7 & 3 & 2 & 0 & 55296 & 17940 & 20920\\
T11.1277 & \texttt{01234342322} & 5 & 5 & 2 & 4 & 8 & 5 & 3 & 1 & 53248 & 26040 & 15680\\
T11.1278 & \texttt{01234342321} & 5 & 5 & 3 & 2 & 9 & 7 & 4 & 3 & 51200 & 27136 & 13596\\
T11.1279 & \texttt{01234342312} & 4 & 5 & 3 & 2 & 9 & 7 & 3 & 2 & 51200 & 16184 & 19320\\
T11.1280 & \texttt{01234342311} & 5 & 5 & 3 & 4 & 8 & 5 & 3 & 1 & 49152 & 24304 & 13868\\
T11.1281 & \texttt{01234342222} & 6 & 5 & 2 & 48 & 6 & 3 & 2 & 0 & 51200 & 28740 & 13392\\
T11.1282 & \texttt{01234342221} & 6 & 5 & 3 & 12 & 7 & 4 & 3 & 1 & 49152 & 31136 & 10268\\
T11.1283 & \texttt{01234342212} & 5 & 5 & 3 & 4 & 8 & 5 & 3 & 1 & 49152 & 24384 & 13968\\
T11.1284 & \texttt{01234342211} & 6 & 5 & 3 & 8 & 7 & 3 & 3 & 0 & 47104 & 30504 & 9244\\
T11.1285 & \texttt{01234342123} & 4 & 5 & 3 & 2 & 9 & 7 & 3 & 2 & 51200 & 17184 & 18312\\
T11.1286 & \texttt{01234342122} & 5 & 5 & 4 & 4 & 8 & 5 & 3 & 1 & 49152 & 22280 & 14828\\
T11.1287 & \texttt{01234342121} & 5 & 5 & 3 & 2 & 9 & 7 & 4 & 3 & 47104 & 24304 & 12680\\
T11.1288 & \texttt{01234342111} & 6 & 5 & 3 & 12 & 7 & 4 & 3 & 1 & 45056 & 27864 & 9508\\
T11.1289 & \texttt{01234341234} & 3 & 5 & 2 & 2 & 9 & 7 & 2 & 1 & 55296 & 7584 & 23804\\
T11.1290 & \texttt{01234341233} & 4 & 5 & 3 & 4 & 8 & 5 & 2 & 0 & 53248 & 12424 & 21524\\
T11.1291 & \texttt{01234341232} & 4 & 5 & 3 & 2 & 9 & 7 & 3 & 2 & 51200 & 14704 & 19584\\
T11.1292 & \texttt{01234341231} & 4 & 5 & 2 & 2 & 9 & 7 & 3 & 2 & 49152 & 16704 & 17324\\
T11.1293 & \texttt{01234341222} & 5 & 5 & 3 & 12 & 7 & 4 & 2 & 0 & 49152 & 18264 & 16848\\
T11.1294 & \texttt{01234341221} & 5 & 5 & 3 & 4 & 8 & 5 & 3 & 1 & 47104 & 20520 & 14360\\
T11.1295 & \texttt{01234341212} & 4 & 5 & 2 & 4 & 7 & 3 & 2 & 0 & 47104 & 14424 & 17672\\
T11.1296 & \texttt{01234341211} & 5 & 5 & 2 & 4 & 8 & 5 & 3 & 1 & 45056 & 21264 & 12988\\
T11.1297 & \texttt{01234341111} & 6 & 5 & 2 & 48 & 6 & 3 & 2 & 0 & 43008 & 23544 & 10592\\
T11.1298 & \texttt{01234333333} & 7 & 5 & 1 & 720 & 6 & 5 & 2 & 1 & 57344 & 28224 & 20544\\
T11.1299 & \texttt{01234333332} & 7 & 5 & 2 & 120 & 7 & 6 & 3 & 2 & 55296 & 32920 & 14916\\
T11.1300 & \texttt{01234333331} & 7 & 5 & 2 & 120 & 7 & 6 & 3 & 2 & 53248 & 32560 & 13956\\
T11.1301 & \texttt{01234333323} & 6 & 5 & 2 & 24 & 8 & 7 & 3 & 2 & 55296 & 26820 & 17996\\
T11.1302 & \texttt{01234333322} & 7 & 5 & 2 & 48 & 7 & 5 & 3 & 1 & 53248 & 34056 & 12252\\
T11.1303 & \texttt{01234333321} & 7 & 5 & 3 & 24 & 8 & 7 & 4 & 3 & 51200 & 35360 & 9880\\
T11.1304 & \texttt{01234333312} & 6 & 5 & 2 & 24 & 8 & 7 & 3 & 2 & 51200 & 25080 & 16480\\
T11.1305 & \texttt{01234333311} & 7 & 5 & 2 & 48 & 7 & 5 & 3 & 1 & 49152 & 32880 & 10372\\
T11.1306 & \texttt{01234333234} & 5 & 5 & 2 & 6 & 9 & 8 & 3 & 2 & 57344 & 21980 & 20308\\
T11.1307 & \texttt{01234333233} & 6 & 5 & 3 & 12 & 8 & 6 & 3 & 1 & 55296 & 27040 & 16876\\
T11.1308 & \texttt{01234333232} & 6 & 5 & 2 & 6 & 9 & 8 & 4 & 3 & 53248 & 29540 & 14396\\
T11.1309 & \texttt{01234333231} & 6 & 5 & 3 & 6 & 9 & 8 & 4 & 3 & 51200 & 29856 & 12832\\
T11.1310 & \texttt{01234333222} & 7 & 5 & 2 & 36 & 7 & 5 & 3 & 1 & 51200 & 33320 & 11148\\
T11.1311 & \texttt{01234333221} & 7 & 5 & 3 & 12 & 8 & 6 & 4 & 2 & 49152 & 35456 & 8064\\
T11.1312 & \texttt{01234333212} & 6 & 5 & 3 & 6 & 9 & 8 & 4 & 3 & 49152 & 27864 & 12684\\
T11.1313 & \texttt{01234333211} & 7 & 5 & 3 & 12 & 8 & 6 & 4 & 2 & 47104 & 34384 & 7400\\
T11.1314 & \texttt{01234333123} & 5 & 5 & 2 & 6 & 9 & 8 & 3 & 2 & 51200 & 19144 & 18448\\
T11.1315 & \texttt{01234333122} & 6 & 5 & 3 & 12 & 8 & 6 & 3 & 1 & 49152 & 24640 & 14604\\
T11.1316 & \texttt{01234333121} & 6 & 5 & 2 & 6 & 9 & 8 & 4 & 3 & 47104 & 27064 & 12016\\
T11.1317 & \texttt{01234333111} & 7 & 5 & 2 & 36 & 7 & 5 & 3 & 1 & 45056 & 31024 & 8404\\
T11.1318 & \texttt{01234332343} & 5 & 5 & 3 & 2 & 10 & 9 & 4 & 3 & 57344 & 23604 & 19072\\
T11.1319 & \texttt{01234332342} & 5 & 5 & 2 & 2 & 10 & 9 & 4 & 3 & 55296 & 25404 & 16708\\
T11.1320 & \texttt{01234332341} & 5 & 5 & 3 & 2 & 10 & 9 & 4 & 3 & 53248 & 24992 & 15656\\
T11.1321 & \texttt{01234332333} & 6 & 5 & 3 & 12 & 8 & 6 & 3 & 1 & 55296 & 26808 & 16732\\
T11.1322 & \texttt{01234332332} & 6 & 5 & 3 & 4 & 9 & 7 & 4 & 2 & 53248 & 28736 & 14140\\
T11.1323 & \texttt{01234332331} & 6 & 5 & 4 & 4 & 9 & 7 & 4 & 2 & 51200 & 29312 & 12448\\
T11.1324 & \texttt{01234332323} & 5 & 5 & 2 & 4 & 8 & 5 & 3 & 1 & 53248 & 23352 & 17428\\
T11.1325 & \texttt{01234332322} & 6 & 5 & 2 & 4 & 9 & 7 & 4 & 2 & 51200 & 29508 & 12908\\
T11.1326 & \texttt{01234332321} & 6 & 5 & 3 & 2 & 10 & 9 & 5 & 4 & 49152 & 31072 & 10400\\
T11.1327 & \texttt{01234332312} & 5 & 5 & 3 & 2 & 10 & 9 & 4 & 3 & 49152 & 21632 & 15820\\
T11.1328 & \texttt{01234332311} & 6 & 5 & 3 & 4 & 9 & 7 & 4 & 2 & 47104 & 29032 & 10352\\
T11.1329 & \texttt{01234332222} & 7 & 5 & 2 & 48 & 7 & 5 & 3 & 1 & 49152 & 31560 & 10908\\
T11.1330 & \texttt{01234332221} & 7 & 5 & 3 & 12 & 8 & 6 & 4 & 2 & 47104 & 34112 & 7456\\
T11.1331 & \texttt{01234332212} & 6 & 5 & 3 & 4 & 9 & 7 & 4 & 2 & 47104 & 27864 & 11092\\
T11.1332 & \texttt{01234332211} & 7 & 5 & 3 & 8 & 8 & 5 & 4 & 1 & 45056 & 33744 & 6208\\
T11.1333 & \texttt{01234332123} & 5 & 5 & 3 & 2 & 10 & 9 & 4 & 3 & 49152 & 21576 & 15260\\
T11.1334 & \texttt{01234332122} & 6 & 5 & 4 & 4 & 9 & 7 & 4 & 2 & 47104 & 26432 & 11696\\
T11.1335 & \texttt{01234332121} & 6 & 5 & 3 & 2 & 10 & 9 & 5 & 4 & 45056 & 28216 & 9516\\
T11.1336 & \texttt{01234332111} & 7 & 5 & 3 & 12 & 8 & 6 & 4 & 2 & 43008 & 31536 & 6312\\
T11.1337 & \texttt{01234331234} & 4 & 5 & 2 & 2 & 10 & 9 & 3 & 2 & 53248 & 13608 & 20592\\
T11.1338 & \texttt{01234331233} & 5 & 5 & 3 & 4 & 9 & 7 & 3 & 1 & 51200 & 18208 & 18224\\
T11.1339 & \texttt{01234331232} & 5 & 5 & 3 & 2 & 10 & 9 & 4 & 3 & 49152 & 20248 & 16268\\
T11.1340 & \texttt{01234331231} & 5 & 5 & 2 & 2 & 10 & 9 & 4 & 3 & 47104 & 21768 & 14048\\
T11.1341 & \texttt{01234331222} & 6 & 5 & 3 & 12 & 8 & 6 & 3 & 1 & 47104 & 23568 & 13492\\
T11.1342 & \texttt{01234331221} & 6 & 5 & 3 & 4 & 9 & 7 & 4 & 2 & 45056 & 25344 & 11068\\
T11.1343 & \texttt{01234331212} & 5 & 5 & 2 & 4 & 8 & 5 & 3 & 1 & 45056 & 19728 & 14364\\
T11.1344 & \texttt{01234331211} & 6 & 5 & 2 & 4 & 9 & 7 & 4 & 2 & 43008 & 25848 & 9728\\
T11.1345 & \texttt{01234331111} & 7 & 5 & 2 & 48 & 7 & 5 & 3 & 1 & 40960 & 27888 & 7364\\
T11.1346 & \texttt{01234323432} & 5 & 5 & 3 & 2 & 7 & 3 & 3 & 1 & 55296 & 25668 & 16456\\
T11.1347 & \texttt{01234323431} & 5 & 5 & 4 & 2 & 7 & 3 & 3 & 1 & 53248 & 25568 & 15236\\
T11.1348 & \texttt{01234323423} & 4 & 5 & 2 & 1 & 11 & 11 & 4 & 4 & 55296 & 17424 & 20484\\
T11.1349 & \texttt{01234323422} & 5 & 5 & 2 & 2 & 10 & 9 & 4 & 3 & 53248 & 25740 & 15244\\
T11.1350 & \texttt{01234323421} & 5 & 5 & 3 & 1 & 11 & 11 & 5 & 5 & 51200 & 26784 & 13248\\
T11.1351 & \texttt{01234323412} & 4 & 5 & 3 & 1 & 11 & 11 & 4 & 4 & 51200 & 15664 & 18892\\
T11.1352 & \texttt{01234323411} & 5 & 5 & 3 & 2 & 10 & 9 & 4 & 3 & 49152 & 23864 & 13536\\
T11.1353 & \texttt{01234323333} & 6 & 5 & 3 & 24 & 8 & 7 & 3 & 2 & 55296 & 25624 & 17588\\
T11.1354 & \texttt{01234323332} & 6 & 5 & 3 & 6 & 9 & 8 & 4 & 3 & 53248 & 28008 & 14452\\
T11.1355 & \texttt{01234323331} & 6 & 5 & 4 & 6 & 9 & 8 & 4 & 3 & 51200 & 28480 & 12840\\
T11.1356 & \texttt{01234323323} & 5 & 5 & 3 & 2 & 10 & 9 & 4 & 3 & 53248 & 22268 & 17412\\
T11.1357 & \texttt{01234323322} & 6 & 5 & 3 & 4 & 9 & 7 & 4 & 2 & 51200 & 28208 & 13108\\
T11.1358 & \texttt{01234323321} & 6 & 5 & 4 & 2 & 10 & 9 & 5 & 4 & 49152 & 29824 & 10632\\
T11.1359 & \texttt{01234323312} & 5 & 5 & 4 & 2 & 10 & 9 & 4 & 3 & 49152 & 20552 & 15820\\
T11.1360 & \texttt{01234323311} & 6 & 5 & 4 & 4 & 9 & 7 & 4 & 2 & 47104 & 27872 & 10520\\
T11.1361 & \texttt{01234323232} & 5 & 5 & 2 & 2 & 9 & 7 & 4 & 3 & 51200 & 24552 & 15340\\
T11.1362 & \texttt{01234323231} & 5 & 5 & 3 & 2 & 9 & 7 & 4 & 3 & 49152 & 25024 & 13616\\
T11.1363 & \texttt{01234323222} & 6 & 5 & 2 & 6 & 9 & 8 & 4 & 3 & 49152 & 28116 & 12500\\
T11.1364 & \texttt{01234323221} & 6 & 5 & 3 & 2 & 10 & 9 & 5 & 4 & 47104 & 30304 & 9496\\
T11.1365 & \texttt{01234323212} & 5 & 5 & 3 & 1 & 11 & 11 & 5 & 5 & 47104 & 22880 & 13636\\
T11.1366 & \texttt{01234323211} & 6 & 5 & 3 & 2 & 10 & 9 & 5 & 4 & 45056 & 29320 & 8696\\
T11.1367 & \texttt{01234323123} & 4 & 5 & 3 & 1 & 11 & 11 & 4 & 4 & 49152 & 14464 & 18644\\
T11.1368 & \texttt{01234323122} & 5 & 5 & 4 & 2 & 10 & 9 & 4 & 3 & 47104 & 19880 & 15024\\
T11.1369 & \texttt{01234323121} & 5 & 5 & 3 & 1 & 11 & 11 & 5 & 5 & 45056 & 22224 & 12676\\
T11.1370 & \texttt{01234323111} & 6 & 5 & 3 & 6 & 9 & 8 & 4 & 3 & 43008 & 26104 & 9304\\
T11.1371 & \texttt{01234322222} & 7 & 5 & 2 & 120 & 7 & 6 & 3 & 2 & 47104 & 29304 & 11124\\
T11.1372 & \texttt{01234322221} & 7 & 5 & 3 & 24 & 8 & 7 & 4 & 3 & 45056 & 32064 & 7464\\
T11.1373 & \texttt{01234322212} & 6 & 5 & 3 & 6 & 9 & 8 & 4 & 3 & 45056 & 26488 & 10620\\
T11.1374 & \texttt{01234322211} & 7 & 5 & 3 & 12 & 8 & 6 & 4 & 2 & 43008 & 32048 & 5912\\
T11.1375 & \texttt{01234322123} & 5 & 5 & 3 & 2 & 10 & 9 & 4 & 3 & 47104 & 21416 & 13996\\
T11.1376 & \texttt{01234322122} & 6 & 5 & 4 & 4 & 9 & 7 & 4 & 2 & 45056 & 25952 & 10560\\
T11.1377 & \texttt{01234322121} & 6 & 5 & 3 & 2 & 10 & 9 & 5 & 4 & 43008 & 27416 & 8572\\
T11.1378 & \texttt{01234322111} & 7 & 5 & 3 & 12 & 8 & 6 & 4 & 2 & 40960 & 30416 & 5560\\
T11.1379 & \texttt{01234321234} & 4 & 5 & 3 & 1 & 11 & 11 & 4 & 4 & 51200 & 15624 & 18116\\
T11.1380 & \texttt{01234321233} & 5 & 5 & 4 & 2 & 10 & 9 & 4 & 3 & 49152 & 19904 & 15788\\
T11.1381 & \texttt{01234321232} & 5 & 5 & 4 & 1 & 11 & 11 & 5 & 5 & 47104 & 21624 & 13968\\
T11.1382 & \texttt{01234321231} & 5 & 5 & 3 & 1 & 11 & 11 & 5 & 5 & 45056 & 22504 & 12116\\
T11.1383 & \texttt{01234321222} & 6 & 5 & 4 & 6 & 9 & 8 & 4 & 3 & 45056 & 24624 & 11296\\
T11.1384 & \texttt{01234321221} & 6 & 5 & 4 & 2 & 10 & 9 & 5 & 4 & 43008 & 25760 & 9272\\
T11.1385 & \texttt{01234321212} & 5 & 5 & 3 & 2 & 9 & 7 & 4 & 3 & 43008 & 20784 & 12232\\
T11.1386 & \texttt{01234321211} & 6 & 5 & 3 & 2 & 10 & 9 & 5 & 4 & 40960 & 25944 & 8132\\
T11.1387 & \texttt{01234321111} & 7 & 5 & 3 & 24 & 8 & 7 & 4 & 3 & 38912 & 27664 & 5968\\
T11.1388 & \texttt{01234312343} & 4 & 5 & 3 & 2 & 6 & 1 & 2 & 0 & 53248 & 14072 & 20468\\
T11.1389 & \texttt{01234312342} & 4 & 5 & 3 & 1 & 11 & 11 & 4 & 4 & 51200 & 15720 & 18444\\
T11.1390 & \texttt{01234312341} & 4 & 5 & 2 & 1 & 11 & 11 & 4 & 4 & 49152 & 16456 & 16632\\
T11.1391 & \texttt{01234312333} & 5 & 5 & 3 & 6 & 9 & 8 & 3 & 2 & 51200 & 17200 & 18568\\
T11.1392 & \texttt{01234312332} & 5 & 5 & 4 & 2 & 10 & 9 & 4 & 3 & 49152 & 18976 & 16180\\
T11.1393 & \texttt{01234312331} & 5 & 5 & 3 & 2 & 10 & 9 & 4 & 3 & 47104 & 20096 & 14408\\
T11.1394 & \texttt{01234312323} & 4 & 5 & 3 & 2 & 9 & 7 & 3 & 2 & 49152 & 13744 & 19012\\
T11.1395 & \texttt{01234312322} & 5 & 5 & 3 & 2 & 10 & 9 & 4 & 3 & 47104 & 19672 & 15020\\
T11.1396 & \texttt{01234312321} & 5 & 5 & 3 & 1 & 11 & 11 & 5 & 5 & 45056 & 21176 & 12884\\
T11.1397 & \texttt{01234312312} & 4 & 5 & 2 & 1 & 11 & 11 & 4 & 4 & 45056 & 13216 & 16772\\
T11.1398 & \texttt{01234312311} & 5 & 5 & 2 & 2 & 10 & 9 & 4 & 3 & 43008 & 20776 & 11816\\
T11.1399 & \texttt{01234312222} & 6 & 5 & 3 & 24 & 8 & 7 & 3 & 2 & 45056 & 21648 & 13052\\
T11.1400 & \texttt{01234312221} & 6 & 5 & 3 & 6 & 9 & 8 & 4 & 3 & 43008 & 23536 & 10444\\
T11.1401 & \texttt{01234312212} & 5 & 5 & 3 & 2 & 10 & 9 & 4 & 3 & 43008 & 17992 & 13440\\
T11.1402 & \texttt{01234312211} & 6 & 5 & 3 & 4 & 9 & 7 & 4 & 2 & 40960 & 23392 & 9268\\
T11.1403 & \texttt{01234312121} & 5 & 5 & 2 & 2 & 9 & 7 & 4 & 3 & 40960 & 19696 & 11508\\
T11.1404 & \texttt{01234312111} & 6 & 5 & 2 & 6 & 9 & 8 & 4 & 3 & 38912 & 22936 & 8552\\
T11.1405 & \texttt{01234311111} & 7 & 5 & 2 & 120 & 7 & 6 & 3 & 2 & 36864 & 24016 & 6812\\
T11.1406 & \texttt{01234234234} & 3 & 5 & 1 & 6 & 5 & 2 & 1 & 0 & 57344 & 9072 & 24296\\
T11.1407 & \texttt{01234234233} & 4 & 5 & 2 & 4 & 7 & 3 & 2 & 0 & 55296 & 15644 & 20912\\
T11.1408 & \texttt{01234234232} & 4 & 5 & 1 & 2 & 8 & 5 & 3 & 2 & 53248 & 19656 & 18144\\
T11.1409 & \texttt{01234234231} & 4 & 5 & 2 & 2 & 8 & 5 & 3 & 2 & 51200 & 18880 & 17116\\
T11.1410 & \texttt{01234234222} & 5 & 5 & 1 & 12 & 6 & 2 & 2 & 0 & 51200 & 24948 & 14608\\
T11.1411 & \texttt{01234234221} & 5 & 5 & 2 & 4 & 7 & 3 & 3 & 1 & 49152 & 26720 & 12092\\
T11.1412 & \texttt{01234234212} & 4 & 5 & 2 & 2 & 8 & 5 & 3 & 2 & 49152 & 17952 & 16512\\
T11.1413 & \texttt{01234234211} & 5 & 5 & 2 & 4 & 7 & 3 & 3 & 1 & 47104 & 25032 & 11660\\
T11.1414 & \texttt{01234234123} & 3 & 5 & 2 & 2 & 8 & 5 & 2 & 1 & 51200 & 7104 & 21936\\
T11.1415 & \texttt{01234234122} & 4 & 5 & 3 & 4 & 7 & 3 & 2 & 0 & 49152 & 13160 & 18372\\
T11.1416 & \texttt{01234234121} & 4 & 5 & 2 & 2 & 8 & 5 & 3 & 2 & 47104 & 16144 & 15952\\
T11.1417 & \texttt{01234234111} & 5 & 5 & 2 & 12 & 6 & 2 & 2 & 0 & 45056 & 20664 & 12508\\
T11.1418 & \texttt{01234233333} & 6 & 5 & 2 & 120 & 7 & 6 & 2 & 1 & 55296 & 23280 & 19488\\
T11.1419 & \texttt{01234233332} & 6 & 5 & 2 & 24 & 8 & 7 & 3 & 2 & 53248 & 26992 & 15296\\
T11.1420 & \texttt{01234233331} & 6 & 5 & 3 & 24 & 8 & 7 & 3 & 2 & 51200 & 26944 & 14012\\
T11.1421 & \texttt{01234233323} & 5 & 5 & 2 & 6 & 9 & 8 & 3 & 2 & 53248 & 20844 & 18216\\
T11.1422 & \texttt{01234233322} & 6 & 5 & 2 & 12 & 8 & 6 & 3 & 1 & 51200 & 27648 & 13408\\
T11.1423 & \texttt{01234233321} & 6 & 5 & 3 & 6 & 9 & 8 & 4 & 3 & 49152 & 29056 & 11068\\
T11.1424 & \texttt{01234233312} & 5 & 5 & 3 & 6 & 9 & 8 & 3 & 2 & 49152 & 19112 & 16632\\
T11.1425 & \texttt{01234233311} & 6 & 5 & 3 & 12 & 8 & 6 & 3 & 1 & 47104 & 26752 & 11164\\
T11.1426 & \texttt{01234233233} & 5 & 5 & 3 & 8 & 7 & 5 & 2 & 1 & 53248 & 20704 & 17792\\
T11.1427 & \texttt{01234233232} & 5 & 5 & 2 & 2 & 10 & 9 & 4 & 3 & 51200 & 23204 & 15624\\
T11.1428 & \texttt{01234233231} & 5 & 5 & 3 & 2 & 10 & 9 & 4 & 3 & 49152 & 23520 & 14020\\
T11.1429 & \texttt{01234233222} & 6 & 5 & 2 & 12 & 8 & 6 & 3 & 1 & 49152 & 26984 & 12688\\
T11.1430 & \texttt{01234233221} & 6 & 5 & 3 & 4 & 9 & 7 & 4 & 2 & 47104 & 29120 & 9772\\
T11.1431 & \texttt{01234233212} & 5 & 5 & 3 & 2 & 10 & 9 & 4 & 3 & 47104 & 21528 & 13944\\
T11.1432 & \texttt{01234233211} & 6 & 5 & 3 & 4 & 9 & 7 & 4 & 2 & 45056 & 28048 & 9036\\
T11.1433 & \texttt{01234233123} & 4 & 5 & 3 & 2 & 10 & 9 & 3 & 2 & 49152 & 12808 & 19016\\
T11.1434 & \texttt{01234233122} & 5 & 5 & 4 & 4 & 9 & 7 & 3 & 1 & 47104 & 18304 & 15412\\
T11.1435 & \texttt{01234233121} & 5 & 5 & 3 & 2 & 10 & 9 & 4 & 3 & 45056 & 20728 & 13064\\
T11.1436 & \texttt{01234233111} & 6 & 5 & 3 & 12 & 8 & 6 & 3 & 1 & 43008 & 24688 & 9692\\
T11.1437 & \texttt{01234232323} & 4 & 5 & 1 & 6 & 7 & 5 & 2 & 1 & 51200 & 17388 & 18864\\
T11.1438 & \texttt{01234232322} & 5 & 5 & 1 & 4 & 8 & 5 & 3 & 1 & 49152 & 24192 & 14392\\
T11.1439 & \texttt{01234232321} & 5 & 5 & 2 & 2 & 9 & 7 & 4 & 3 & 47104 & 25600 & 12004\\
T11.1440 & \texttt{01234232312} & 4 & 5 & 2 & 2 & 9 & 7 & 3 & 2 & 47104 & 15656 & 17200\\
T11.1441 & \texttt{01234232311} & 5 & 5 & 2 & 4 & 8 & 5 & 3 & 1 & 45056 & 23296 & 11908\\
T11.1442 & \texttt{01234232222} & 6 & 5 & 1 & 24 & 8 & 7 & 3 & 2 & 47104 & 26460 & 12392\\
T11.1443 & \texttt{01234232221} & 6 & 5 & 2 & 6 & 9 & 8 & 4 & 3 & 45056 & 28960 & 9092\\
T11.1444 & \texttt{01234232212} & 5 & 5 & 2 & 2 & 10 & 9 & 4 & 3 & 45056 & 22544 & 12616\\
T11.1445 & \texttt{01234232211} & 6 & 5 & 2 & 4 & 9 & 7 & 4 & 2 & 43008 & 28504 & 7892\\
T11.1446 & \texttt{01234232123} & 4 & 5 & 2 & 1 & 11 & 11 & 4 & 4 & 47104 & 15952 & 16632\\
T11.1447 & \texttt{01234232122} & 5 & 5 & 3 & 2 & 10 & 9 & 4 & 3 & 45056 & 20888 & 13164\\
T11.1448 & \texttt{01234232121} & 5 & 5 & 2 & 1 & 11 & 11 & 5 & 5 & 43008 & 22752 & 11064\\
T11.1449 & \texttt{01234232111} & 6 & 5 & 2 & 6 & 9 & 8 & 4 & 3 & 40960 & 26152 & 7940\\
T11.1450 & \texttt{01234231234} & 3 & 5 & 2 & 1 & 11 & 11 & 3 & 3 & 51200 & 7440 & 21624\\
T11.1451 & \texttt{01234231233} & 4 & 5 & 3 & 2 & 10 & 9 & 3 & 2 & 49152 & 12120 & 19336\\
T11.1452 & \texttt{01234231232} & 4 & 5 & 3 & 1 & 11 & 11 & 4 & 4 & 47104 & 14240 & 17436\\
T11.1453 & \texttt{01234231231} & 4 & 5 & 2 & 1 & 11 & 11 & 4 & 4 & 45056 & 15920 & 15304\\
T11.1454 & \texttt{01234231222} & 5 & 5 & 3 & 6 & 9 & 8 & 3 & 2 & 45056 & 17640 & 14724\\
T11.1455 & \texttt{01234231221} & 5 & 5 & 3 & 2 & 10 & 9 & 4 & 3 & 43008 & 19576 & 12380\\
T11.1456 & \texttt{01234231212} & 4 & 5 & 2 & 2 & 9 & 7 & 3 & 2 & 43008 & 13800 & 15580\\
T11.1457 & \texttt{01234231211} & 5 & 5 & 2 & 2 & 10 & 9 & 4 & 3 & 40960 & 20160 & 11080\\
T11.1458 & \texttt{01234231111} & 6 & 5 & 2 & 24 & 8 & 7 & 3 & 2 & 38912 & 22280 & 8756\\
T11.1459 & \texttt{01234222222} & 7 & 5 & 1 & 720 & 6 & 5 & 2 & 1 & 45056 & 27216 & 11328\\
T11.1460 & \texttt{01234222221} & 7 & 5 & 2 & 120 & 7 & 6 & 3 & 2 & 43008 & 30080 & 7516\\
T11.1461 & \texttt{01234222212} & 6 & 5 & 2 & 24 & 8 & 7 & 3 & 2 & 43008 & 24840 & 10456\\
T11.1462 & \texttt{01234222211} & 7 & 5 & 2 & 48 & 7 & 5 & 3 & 1 & 40960 & 30240 & 5788\\
T11.1463 & \texttt{01234222123} & 5 & 5 & 2 & 6 & 9 & 8 & 3 & 2 & 45056 & 20376 & 13448\\
T11.1464 & \texttt{01234222122} & 6 & 5 & 3 & 12 & 8 & 6 & 3 & 1 & 43008 & 24752 & 10052\\
T11.1465 & \texttt{01234222121} & 6 & 5 & 2 & 6 & 9 & 8 & 4 & 3 & 40960 & 26056 & 8136\\
T11.1466 & \texttt{01234222111} & 7 & 5 & 2 & 36 & 7 & 5 & 3 & 1 & 38912 & 28896 & 5196\\
T11.1467 & \texttt{01234221234} & 4 & 5 & 2 & 2 & 10 & 9 & 3 & 2 & 49152 & 15672 & 16968\\
T11.1468 & \texttt{01234221233} & 5 & 5 & 3 & 4 & 9 & 7 & 3 & 1 & 47104 & 19792 & 14648\\
T11.1469 & \texttt{01234221232} & 5 & 5 & 3 & 2 & 10 & 9 & 4 & 3 & 45056 & 21352 & 12884\\
T11.1470 & \texttt{01234221231} & 5 & 5 & 2 & 2 & 10 & 9 & 4 & 3 & 43008 & 21912 & 11192\\
T11.1471 & \texttt{01234221222} & 6 & 5 & 3 & 12 & 8 & 6 & 3 & 1 & 43008 & 24192 & 10252\\
T11.1472 & \texttt{01234221221} & 6 & 5 & 3 & 4 & 9 & 7 & 4 & 2 & 40960 & 25008 & 8404\\
T11.1473 & \texttt{01234221212} & 5 & 5 & 2 & 4 & 8 & 5 & 3 & 1 & 40960 & 20352 & 11220\\
T11.1474 & \texttt{01234221211} & 6 & 5 & 2 & 4 & 9 & 7 & 4 & 2 & 38912 & 25032 & 7352\\
T11.1475 & \texttt{01234221111} & 7 & 5 & 2 & 48 & 7 & 5 & 3 & 1 & 36864 & 26592 & 5276\\
T11.1476 & \texttt{01234212342} & 4 & 5 & 3 & 2 & 6 & 1 & 2 & 0 & 49152 & 17048 & 16548\\
T11.1477 & \texttt{01234212341} & 4 & 5 & 2 & 1 & 11 & 11 & 4 & 4 & 47104 & 17144 & 15040\\
T11.1478 & \texttt{01234212333} & 5 & 5 & 3 & 6 & 9 & 8 & 3 & 2 & 49152 & 18688 & 16512\\
T11.1479 & \texttt{01234212332} & 5 & 5 & 4 & 2 & 10 & 9 & 4 & 3 & 47104 & 20144 & 14268\\
T11.1480 & \texttt{01234212331} & 5 & 5 & 3 & 2 & 10 & 9 & 4 & 3 & 45056 & 20624 & 12832\\
T11.1481 & \texttt{01234212323} & 4 & 5 & 3 & 2 & 9 & 7 & 3 & 2 & 47104 & 15232 & 17004\\
T11.1482 & \texttt{01234212322} & 5 & 5 & 3 & 2 & 10 & 9 & 4 & 3 & 45056 & 20680 & 13156\\
T11.1483 & \texttt{01234212321} & 5 & 5 & 3 & 1 & 11 & 11 & 5 & 5 & 43008 & 21544 & 11372\\
T11.1484 & \texttt{01234212312} & 4 & 5 & 2 & 1 & 11 & 11 & 4 & 4 & 43008 & 14384 & 14924\\
T11.1485 & \texttt{01234212311} & 5 & 5 & 2 & 2 & 10 & 9 & 4 & 3 & 40960 & 20824 & 10480\\
T11.1486 & \texttt{01234212222} & 6 & 5 & 3 & 24 & 8 & 7 & 3 & 2 & 43008 & 22496 & 11204\\
T11.1487 & \texttt{01234212221} & 6 & 5 & 3 & 6 & 9 & 8 & 4 & 3 & 40960 & 23744 & 8980\\
T11.1488 & \texttt{01234212212} & 5 & 5 & 3 & 2 & 10 & 9 & 4 & 3 & 40960 & 18680 & 11752\\
T11.1489 & \texttt{01234212211} & 6 & 5 & 3 & 4 & 9 & 7 & 4 & 2 & 38912 & 23280 & 7996\\
T11.1490 & \texttt{01234212121} & 5 & 5 & 2 & 2 & 9 & 7 & 4 & 3 & 38912 & 19904 & 10076\\
T11.1491 & \texttt{01234212111} & 6 & 5 & 2 & 6 & 9 & 8 & 4 & 3 & 36864 & 22664 & 7376\\
T11.1492 & \texttt{01234211111} & 7 & 5 & 2 & 120 & 7 & 6 & 3 & 2 & 34816 & 23584 & 5732\\
T11.1493 & \texttt{01234123412} & 3 & 5 & 1 & 2 & 7 & 3 & 2 & 1 & 47104 & 6960 & 19280\\
T11.1494 & \texttt{01234123411} & 4 & 5 & 1 & 4 & 6 & 1 & 2 & 0 & 45056 & 15960 & 14236\\
T11.1495 & \texttt{01234123333} & 5 & 5 & 2 & 24 & 8 & 7 & 2 & 1 & 51200 & 15616 & 19364\\
T11.1496 & \texttt{01234123332} & 5 & 5 & 3 & 6 & 9 & 8 & 3 & 2 & 49152 & 17952 & 16256\\
T11.1497 & \texttt{01234123331} & 5 & 5 & 2 & 6 & 9 & 8 & 3 & 2 & 47104 & 19104 & 14660\\
T11.1498 & \texttt{01234123323} & 4 & 5 & 3 & 2 & 10 & 9 & 3 & 2 & 49152 & 12248 & 18992\\
T11.1499 & \texttt{01234123322} & 5 & 5 & 3 & 4 & 9 & 7 & 3 & 1 & 47104 & 18128 & 14920\\
T11.1500 & \texttt{01234123321} & 5 & 5 & 3 & 2 & 10 & 9 & 4 & 3 & 45056 & 19920 & 12752\\
T11.1501 & \texttt{01234123312} & 4 & 5 & 2 & 2 & 10 & 9 & 3 & 2 & 45056 & 11320 & 17040\\
T11.1502 & \texttt{01234123311} & 5 & 5 & 2 & 4 & 9 & 7 & 3 & 1 & 43008 & 19120 & 12084\\
T11.1503 & \texttt{01234123232} & 4 & 5 & 2 & 2 & 9 & 7 & 3 & 2 & 47104 & 14496 & 17040\\
T11.1504 & \texttt{01234123231} & 4 & 5 & 2 & 2 & 9 & 7 & 3 & 2 & 45056 & 15648 & 15248\\
T11.1505 & \texttt{01234123222} & 5 & 5 & 2 & 6 & 9 & 8 & 3 & 2 & 45056 & 18024 & 14304\\
T11.1506 & \texttt{01234123221} & 5 & 5 & 2 & 2 & 10 & 9 & 4 & 3 & 43008 & 20136 & 11760\\
T11.1507 & \texttt{01234123212} & 4 & 5 & 2 & 1 & 11 & 11 & 4 & 4 & 43008 & 13120 & 15284\\
T11.1508 & \texttt{01234123211} & 5 & 5 & 2 & 2 & 10 & 9 & 4 & 3 & 40960 & 19720 & 10776\\
T11.1509 & \texttt{01234123123} & 3 & 5 & 1 & 2 & 8 & 5 & 2 & 1 & 45056 & 5760 & 19200\\
T11.1510 & \texttt{01234123122} & 4 & 5 & 2 & 2 & 10 & 9 & 3 & 2 & 43008 & 11496 & 15764\\
T11.1511 & \texttt{01234123121} & 4 & 5 & 1 & 1 & 11 & 11 & 4 & 4 & 40960 & 14160 & 13536\\
T11.1512 & \texttt{01234123111} & 5 & 5 & 1 & 6 & 9 & 8 & 3 & 2 & 38912 & 18360 & 10284\\
T11.1513 & \texttt{01234122222} & 6 & 5 & 2 & 120 & 7 & 6 & 2 & 1 & 43008 & 19200 & 12904\\
T11.1514 & \texttt{01234122221} & 6 & 5 & 2 & 24 & 8 & 7 & 3 & 2 & 40960 & 21632 & 9864\\
T11.1515 & \texttt{01234122212} & 5 & 5 & 2 & 6 & 9 & 8 & 3 & 2 & 40960 & 16200 & 12684\\
T11.1516 & \texttt{01234122211} & 6 & 5 & 2 & 12 & 8 & 6 & 3 & 1 & 38912 & 21600 & 8504\\
T11.1517 & \texttt{01234122122} & 5 & 5 & 3 & 8 & 7 & 5 & 2 & 1 & 40960 & 16032 & 12536\\
T11.1518 & \texttt{01234122121} & 5 & 5 & 2 & 2 & 10 & 9 & 4 & 3 & 38912 & 17496 & 10756\\
T11.1519 & \texttt{01234122111} & 6 & 5 & 2 & 12 & 8 & 6 & 3 & 1 & 36864 & 20496 & 7952\\
T11.1520 & \texttt{01234121212} & 4 & 5 & 1 & 6 & 7 & 5 & 2 & 1 & 38912 & 12360 & 13604\\
T11.1521 & \texttt{01234121211} & 5 & 5 & 1 & 4 & 8 & 5 & 3 & 1 & 36864 & 17760 & 9664\\
T11.1522 & \texttt{01234121111} & 6 & 5 & 1 & 24 & 8 & 7 & 3 & 2 & 34816 & 19560 & 7548\\
T11.1523 & \texttt{01234111111} & 7 & 5 & 1 & 720 & 6 & 5 & 2 & 1 & 32768 & 20160 & 6120\\
T11.1524 & \texttt{01233333333} & 8 & 4 & 1 & 40320 & 4 & 3 & 1 & 0 & 55296 & 28448 & 20604\\
T11.1525 & \texttt{01233333332} & 8 & 4 & 2 & 5040 & 5 & 4 & 2 & 1 & 53248 & 33228 & 14512\\
T11.1526 & \texttt{01233333331} & 8 & 4 & 2 & 5040 & 5 & 4 & 2 & 1 & 51200 & 32920 & 13512\\
T11.1527 & \texttt{01233333323} & 7 & 4 & 2 & 720 & 6 & 5 & 2 & 1 & 53248 & 27480 & 17512\\
T11.1528 & \texttt{01233333322} & 8 & 4 & 2 & 1440 & 5 & 3 & 2 & 0 & 51200 & 34500 & 11528\\
T11.1529 & \texttt{01233333321} & 8 & 4 & 3 & 720 & 6 & 5 & 3 & 2 & 49152 & 35856 & 8988\\
T11.1530 & \texttt{01233333312} & 7 & 4 & 2 & 720 & 6 & 5 & 2 & 1 & 49152 & 25744 & 15988\\
T11.1531 & \texttt{01233333311} & 8 & 4 & 2 & 1440 & 5 & 3 & 2 & 0 & 47104 & 33464 & 9544\\
T11.1532 & \texttt{01233333233} & 7 & 4 & 3 & 240 & 6 & 4 & 2 & 0 & 53248 & 28124 & 15960\\
T11.1533 & \texttt{01233333232} & 7 & 4 & 2 & 120 & 7 & 6 & 3 & 2 & 51200 & 30408 & 13432\\
T11.1534 & \texttt{01233333231} & 7 & 4 & 3 & 120 & 7 & 6 & 3 & 2 & 49152 & 30880 & 11700\\
T11.1535 & \texttt{01233333222} & 8 & 4 & 2 & 720 & 5 & 3 & 2 & 0 & 49152 & 33972 & 10136\\
T11.1536 & \texttt{01233333221} & 8 & 4 & 3 & 240 & 6 & 4 & 3 & 1 & 47104 & 36160 & 6820\\
T11.1537 & \texttt{01233333212} & 7 & 4 & 3 & 120 & 7 & 6 & 3 & 2 & 47104 & 28736 & 11648\\
T11.1538 & \texttt{01233333211} & 8 & 4 & 3 & 240 & 6 & 4 & 3 & 1 & 45056 & 35176 & 6092\\
T11.1539 & \texttt{01233333123} & 6 & 4 & 2 & 120 & 7 & 6 & 2 & 1 & 49152 & 20320 & 17684\\
T11.1540 & \texttt{01233333122} & 7 & 4 & 3 & 240 & 6 & 4 & 2 & 0 & 47104 & 25736 & 13680\\
T11.1541 & \texttt{01233333121} & 7 & 4 & 2 & 120 & 7 & 6 & 3 & 2 & 45056 & 28080 & 10948\\
T11.1542 & \texttt{01233333111} & 8 & 4 & 2 & 720 & 5 & 3 & 2 & 0 & 43008 & 31960 & 7192\\
T11.1543 & \texttt{01233332333} & 7 & 4 & 3 & 144 & 6 & 4 & 2 & 0 & 53248 & 28612 & 15224\\
T11.1544 & \texttt{01233332332} & 7 & 4 & 3 & 48 & 7 & 5 & 3 & 1 & 51200 & 30108 & 12728\\
T11.1545 & \texttt{01233332331} & 7 & 4 & 4 & 48 & 7 & 5 & 3 & 1 & 49152 & 30944 & 10772\\
T11.1546 & \texttt{01233332323} & 6 & 4 & 2 & 48 & 6 & 3 & 2 & 0 & 51200 & 25156 & 15968\\
T11.1547 & \texttt{01233332322} & 7 & 4 & 2 & 48 & 7 & 5 & 3 & 1 & 49152 & 30664 & 11544\\
T11.1548 & \texttt{01233332321} & 7 & 4 & 3 & 24 & 8 & 7 & 4 & 3 & 47104 & 32384 & 8772\\
T11.1549 & \texttt{01233332312} & 6 & 4 & 3 & 24 & 8 & 7 & 3 & 2 & 47104 & 23448 & 14320\\
T11.1550 & \texttt{01233332311} & 7 & 4 & 3 & 48 & 7 & 5 & 3 & 1 & 45056 & 30608 & 8628\\
T11.1551 & \texttt{01233332222} & 8 & 4 & 2 & 576 & 5 & 3 & 2 & 0 & 47104 & 32500 & 9592\\
T11.1552 & \texttt{01233332221} & 8 & 4 & 3 & 144 & 6 & 4 & 3 & 1 & 45056 & 35104 & 5876\\
T11.1553 & \texttt{01233332212} & 7 & 4 & 3 & 48 & 7 & 5 & 3 & 1 & 45056 & 29024 & 9624\\
T11.1554 & \texttt{01233332211} & 8 & 4 & 3 & 96 & 6 & 3 & 3 & 0 & 43008 & 34824 & 4500\\
T11.1555 & \texttt{01233332123} & 6 & 4 & 3 & 24 & 8 & 7 & 3 & 2 & 47104 & 23040 & 13872\\
T11.1556 & \texttt{01233332122} & 7 & 4 & 4 & 48 & 7 & 5 & 3 & 1 & 45056 & 27816 & 10180\\
T11.1557 & \texttt{01233332121} & 7 & 4 & 3 & 24 & 8 & 7 & 4 & 3 & 43008 & 29520 & 7888\\
T11.1558 & \texttt{01233332111} & 8 & 4 & 3 & 144 & 6 & 4 & 3 & 1 & 40960 & 32760 & 4572\\
T11.1559 & \texttt{01233331233} & 6 & 4 & 3 & 48 & 7 & 5 & 2 & 0 & 49152 & 20136 & 16908\\
T11.1560 & \texttt{01233331232} & 6 & 4 & 3 & 24 & 8 & 7 & 3 & 2 & 47104 & 22096 & 14888\\
T11.1561 & \texttt{01233331231} & 6 & 4 & 2 & 24 & 8 & 7 & 3 & 2 & 45056 & 23456 & 12564\\
T11.1562 & \texttt{01233331222} & 7 & 4 & 3 & 144 & 6 & 4 & 2 & 0 & 45056 & 25336 & 12040\\
T11.1563 & \texttt{01233331221} & 7 & 4 & 3 & 48 & 7 & 5 & 3 & 1 & 43008 & 26952 & 9520\\
T11.1564 & \texttt{01233331212} & 6 & 4 & 2 & 48 & 6 & 3 & 2 & 0 & 43008 & 21496 & 12928\\
T11.1565 & \texttt{01233331211} & 7 & 4 & 2 & 48 & 7 & 5 & 3 & 1 & 40960 & 27376 & 8132\\
T11.1566 & \texttt{01233331111} & 8 & 4 & 2 & 576 & 5 & 3 & 2 & 0 & 38912 & 29336 & 5720\\
T11.1567 & \texttt{01233323332} & 7 & 4 & 3 & 72 & 5 & 3 & 2 & 1 & 51200 & 30132 & 12472\\
T11.1568 & \texttt{01233323331} & 7 & 4 & 4 & 72 & 5 & 3 & 2 & 1 & 49152 & 31072 & 10452\\
T11.1569 & \texttt{01233323323} & 6 & 4 & 3 & 12 & 8 & 6 & 3 & 1 & 51200 & 25040 & 15240\\
T11.1570 & \texttt{01233323322} & 7 & 4 & 3 & 24 & 7 & 4 & 3 & 0 & 49152 & 29900 & 11320\\
T11.1571 & \texttt{01233323321} & 7 & 4 & 4 & 12 & 8 & 6 & 4 & 2 & 47104 & 31776 & 8500\\
T11.1572 & \texttt{01233323312} & 6 & 4 & 4 & 12 & 8 & 6 & 3 & 1 & 47104 & 23344 & 13592\\
T11.1573 & \texttt{01233323311} & 7 & 4 & 4 & 24 & 7 & 4 & 3 & 0 & 45056 & 30264 & 8164\\
T11.1574 & \texttt{01233323232} & 6 & 4 & 2 & 12 & 7 & 4 & 3 & 1 & 49152 & 26676 & 13408\\
T11.1575 & \texttt{01233323231} & 6 & 4 & 3 & 12 & 7 & 4 & 3 & 1 & 47104 & 27616 & 11308\\
T11.1576 & \texttt{01233323222} & 7 & 4 & 2 & 36 & 7 & 5 & 3 & 1 & 47104 & 29592 & 10808\\
T11.1577 & \texttt{01233323221} & 7 & 4 & 3 & 12 & 8 & 6 & 4 & 2 & 45056 & 31936 & 7492\\
T11.1578 & \texttt{01233323212} & 6 & 4 & 3 & 6 & 9 & 8 & 4 & 3 & 45056 & 25016 & 11616\\
T11.1579 & \texttt{01233323211} & 7 & 4 & 3 & 12 & 8 & 6 & 4 & 2 & 43008 & 31216 & 6500\\
T11.1580 & \texttt{01233323123} & 5 & 4 & 3 & 6 & 9 & 8 & 3 & 2 & 47104 & 17512 & 16544\\
T11.1581 & \texttt{01233323122} & 6 & 4 & 4 & 12 & 8 & 6 & 3 & 1 & 45056 & 22688 & 12828\\
T11.1582 & \texttt{01233323121} & 6 & 4 & 3 & 6 & 9 & 8 & 4 & 3 & 43008 & 24792 & 10432\\
T11.1583 & \texttt{01233323111} & 7 & 4 & 3 & 36 & 7 & 5 & 3 & 1 & 40960 & 28432 & 7012\\
T11.1584 & \texttt{01233322222} & 8 & 4 & 2 & 720 & 5 & 3 & 2 & 0 & 45056 & 30564 & 9528\\
T11.1585 & \texttt{01233322221} & 8 & 4 & 3 & 144 & 6 & 4 & 3 & 1 & 43008 & 33376 & 5588\\
T11.1586 & \texttt{01233322212} & 7 & 4 & 3 & 36 & 7 & 5 & 3 & 1 & 43008 & 27968 & 8808\\
T11.1587 & \texttt{01233322211} & 8 & 4 & 3 & 72 & 6 & 3 & 3 & 0 & 40960 & 33448 & 3876\\
T11.1588 & \texttt{01233322123} & 6 & 4 & 3 & 12 & 8 & 6 & 3 & 1 & 45056 & 23200 & 12168\\
T11.1589 & \texttt{01233322122} & 7 & 4 & 4 & 24 & 7 & 4 & 3 & 0 & 43008 & 27656 & 8620\\
T11.1590 & \texttt{01233322121} & 7 & 4 & 3 & 12 & 8 & 6 & 4 & 2 & 40960 & 29040 & 6536\\
T11.1591 & \texttt{01233322111} & 8 & 4 & 3 & 72 & 6 & 3 & 3 & 0 & 38912 & 31960 & 3428\\
T11.1592 & \texttt{01233321233} & 6 & 4 & 4 & 12 & 8 & 6 & 3 & 1 & 47104 & 22152 & 13856\\
T11.1593 & \texttt{01233321232} & 6 & 4 & 4 & 6 & 9 & 8 & 4 & 3 & 45056 & 23792 & 11988\\
T11.1594 & \texttt{01233321231} & 6 & 4 & 3 & 6 & 9 & 8 & 4 & 3 & 43008 & 24512 & 10064\\
T11.1595 & \texttt{01233321222} & 7 & 4 & 4 & 36 & 7 & 5 & 3 & 1 & 43008 & 26712 & 9260\\
T11.1596 & \texttt{01233321221} & 7 & 4 & 4 & 12 & 8 & 6 & 4 & 2 & 40960 & 27688 & 7172\\
T11.1597 & \texttt{01233321212} & 6 & 4 & 3 & 12 & 7 & 4 & 3 & 1 & 40960 & 22872 & 10212\\
T11.1598 & \texttt{01233321211} & 7 & 4 & 3 & 12 & 8 & 6 & 4 & 2 & 38912 & 27792 & 6000\\
T11.1599 & \texttt{01233321111} & 8 & 4 & 3 & 144 & 6 & 4 & 3 & 1 & 36864 & 29432 & 3804\\
T11.1600 & \texttt{01233312333} & 6 & 4 & 3 & 72 & 4 & 1 & 1 & 0 & 49152 & 20248 & 16604\\
T11.1601 & \texttt{01233312332} & 6 & 4 & 4 & 12 & 8 & 6 & 3 & 1 & 47104 & 21864 & 14192\\
T11.1602 & \texttt{01233312331} & 6 & 4 & 3 & 12 & 8 & 6 & 3 & 1 & 45056 & 22664 & 12396\\
T11.1603 & \texttt{01233312323} & 5 & 4 & 3 & 12 & 7 & 4 & 2 & 0 & 47104 & 16792 & 17072\\
T11.1604 & \texttt{01233312322} & 6 & 4 & 3 & 12 & 8 & 6 & 3 & 1 & 45056 & 22480 & 13008\\
T11.1605 & \texttt{01233312321} & 6 & 4 & 3 & 6 & 9 & 8 & 4 & 3 & 43008 & 23664 & 10856\\
T11.1606 & \texttt{01233312312} & 5 & 4 & 2 & 6 & 9 & 8 & 3 & 2 & 43008 & 16104 & 14816\\
T11.1607 & \texttt{01233312311} & 6 & 4 & 2 & 12 & 8 & 6 & 3 & 1 & 40960 & 23104 & 9780\\
T11.1608 & \texttt{01233312222} & 7 & 4 & 3 & 144 & 6 & 4 & 2 & 0 & 43008 & 24376 & 11000\\
T11.1609 & \texttt{01233312221} & 7 & 4 & 3 & 36 & 7 & 5 & 3 & 1 & 40960 & 25944 & 8392\\
T11.1610 & \texttt{01233312212} & 6 & 4 & 3 & 12 & 8 & 6 & 3 & 1 & 40960 & 20640 & 11420\\
T11.1611 & \texttt{01233312211} & 7 & 4 & 3 & 24 & 7 & 4 & 3 & 0 & 38912 & 25640 & 7216\\
T11.1612 & \texttt{01233312121} & 6 & 4 & 2 & 12 & 7 & 4 & 3 & 1 & 38912 & 22104 & 9472\\
T11.1613 & \texttt{01233312111} & 7 & 4 & 2 & 36 & 7 & 5 & 3 & 1 & 36864 & 25104 & 6500\\
T11.1614 & \texttt{01233311111} & 8 & 4 & 2 & 720 & 5 & 3 & 2 & 0 & 34816 & 26104 & 4760\\
T11.1615 & \texttt{01233233233} & 6 & 4 & 4 & 48 & 4 & 2 & 1 & 0 & 51200 & 24684 & 14936\\
T11.1616 & \texttt{01233233232} & 6 & 4 & 3 & 8 & 7 & 5 & 3 & 2 & 49152 & 26104 & 13272\\
T11.1617 & \texttt{01233233231} & 6 & 4 & 4 & 8 & 7 & 5 & 3 & 2 & 47104 & 27200 & 11108\\
T11.1618 & \texttt{01233233222} & 7 & 4 & 3 & 48 & 5 & 2 & 2 & 0 & 47104 & 28804 & 10840\\
T11.1619 & \texttt{01233233221} & 7 & 4 & 4 & 16 & 6 & 3 & 3 & 1 & 45056 & 31200 & 7540\\
T11.1620 & \texttt{01233233212} & 6 & 4 & 4 & 8 & 7 & 5 & 3 & 2 & 45056 & 24448 & 11496\\
T11.1621 & \texttt{01233233211} & 7 & 4 & 4 & 16 & 6 & 3 & 3 & 1 & 43008 & 30568 & 6500\\
T11.1622 & \texttt{01233233123} & 5 & 4 & 4 & 8 & 7 & 5 & 2 & 1 & 47104 & 17248 & 16184\\
T11.1623 & \texttt{01233233122} & 6 & 4 & 5 & 16 & 6 & 3 & 2 & 0 & 45056 & 22344 & 12540\\
T11.1624 & \texttt{01233233121} & 6 & 4 & 4 & 8 & 7 & 5 & 3 & 2 & 43008 & 24368 & 10232\\
T11.1625 & \texttt{01233233111} & 7 & 4 & 4 & 48 & 5 & 2 & 2 & 0 & 40960 & 27928 & 6900\\
T11.1626 & \texttt{01233232323} & 5 & 4 & 2 & 12 & 6 & 3 & 2 & 0 & 49152 & 21584 & 16032\\
T11.1627 & \texttt{01233232322} & 6 & 4 & 2 & 8 & 7 & 3 & 3 & 0 & 47104 & 26444 & 12352\\
T11.1628 & \texttt{01233232321} & 6 & 4 & 3 & 4 & 8 & 5 & 4 & 2 & 45056 & 28320 & 9516\\
T11.1629 & \texttt{01233232312} & 5 & 4 & 3 & 4 & 8 & 5 & 3 & 1 & 45056 & 19888 & 14336\\
T11.1630 & \texttt{01233232311} & 6 & 4 & 3 & 8 & 7 & 3 & 3 & 0 & 43008 & 26808 & 9052\\
T11.1631 & \texttt{01233232222} & 7 & 4 & 2 & 48 & 7 & 5 & 3 & 1 & 45056 & 28064 & 10664\\
T11.1632 & \texttt{01233232221} & 7 & 4 & 3 & 12 & 8 & 6 & 4 & 2 & 43008 & 30720 & 7028\\
T11.1633 & \texttt{01233232212} & 6 & 4 & 3 & 4 & 9 & 7 & 4 & 2 & 43008 & 24808 & 10464\\
T11.1634 & \texttt{01233232211} & 7 & 4 & 3 & 8 & 8 & 5 & 4 & 1 & 40960 & 30528 & 5588\\
T11.1635 & \texttt{01233232123} & 5 & 4 & 3 & 2 & 10 & 9 & 4 & 3 & 45056 & 19128 & 14256\\
T11.1636 & \texttt{01233232122} & 6 & 4 & 4 & 4 & 9 & 7 & 4 & 2 & 43008 & 23824 & 10716\\
T11.1637 & \texttt{01233232121} & 6 & 4 & 3 & 2 & 10 & 9 & 5 & 4 & 40960 & 25448 & 8592\\
T11.1638 & \texttt{01233232111} & 7 & 4 & 3 & 12 & 8 & 6 & 4 & 2 & 38912 & 28608 & 5444\\
T11.1639 & \texttt{01233231233} & 5 & 4 & 4 & 4 & 9 & 7 & 3 & 1 & 47104 & 16688 & 16624\\
T11.1640 & \texttt{01233231232} & 5 & 4 & 4 & 2 & 10 & 9 & 4 & 3 & 45056 & 18568 & 14716\\
T11.1641 & \texttt{01233231231} & 5 & 4 & 3 & 2 & 10 & 9 & 4 & 3 & 43008 & 19768 & 12640\\
T11.1642 & \texttt{01233231222} & 6 & 4 & 4 & 12 & 8 & 6 & 3 & 1 & 43008 & 21728 & 11972\\
T11.1643 & \texttt{01233231221} & 6 & 4 & 4 & 4 & 9 & 7 & 4 & 2 & 40960 & 23184 & 9708\\
T11.1644 & \texttt{01233231212} & 5 & 4 & 3 & 4 & 8 & 5 & 3 & 1 & 40960 & 17888 & 12876\\
T11.1645 & \texttt{01233231211} & 6 & 4 & 3 & 4 & 9 & 7 & 4 & 2 & 38912 & 23528 & 8448\\
T11.1646 & \texttt{01233231111} & 7 & 4 & 3 & 48 & 7 & 5 & 3 & 1 & 36864 & 25408 & 6164\\
T11.1647 & \texttt{01233222222} & 8 & 4 & 2 & 1440 & 5 & 3 & 2 & 0 & 43008 & 28604 & 9720\\
T11.1648 & \texttt{01233222221} & 8 & 4 & 3 & 240 & 6 & 4 & 3 & 1 & 40960 & 31520 & 5620\\
T11.1649 & \texttt{01233222212} & 7 & 4 & 3 & 48 & 7 & 5 & 3 & 1 & 40960 & 26448 & 8600\\
T11.1650 & \texttt{01233222211} & 8 & 4 & 3 & 96 & 6 & 3 & 3 & 0 & 38912 & 31768 & 3716\\
T11.1651 & \texttt{01233222123} & 6 & 4 & 3 & 12 & 8 & 6 & 3 & 1 & 43008 & 22288 & 11528\\
T11.1652 & \texttt{01233222122} & 7 & 4 & 4 & 24 & 7 & 4 & 3 & 0 & 40960 & 26584 & 8028\\
T11.1653 & \texttt{01233222121} & 7 & 4 & 3 & 12 & 8 & 6 & 4 & 2 & 38912 & 27808 & 6024\\
T11.1654 & \texttt{01233222111} & 8 & 4 & 3 & 72 & 6 & 3 & 3 & 0 & 36864 & 30568 & 2996\\
T11.1655 & \texttt{01233221233} & 6 & 4 & 4 & 8 & 8 & 5 & 3 & 0 & 45056 & 22168 & 12536\\
T11.1656 & \texttt{01233221232} & 6 & 4 & 4 & 4 & 9 & 7 & 4 & 2 & 43008 & 23648 & 10732\\
T11.1657 & \texttt{01233221231} & 6 & 4 & 3 & 4 & 9 & 7 & 4 & 2 & 40960 & 24048 & 8984\\
T11.1658 & \texttt{01233221222} & 7 & 4 & 4 & 24 & 7 & 4 & 3 & 0 & 40960 & 26408 & 8052\\
T11.1659 & \texttt{01233221221} & 7 & 4 & 4 & 8 & 8 & 5 & 4 & 1 & 38912 & 27064 & 6156\\
T11.1660 & \texttt{01233221212} & 6 & 4 & 3 & 8 & 7 & 3 & 3 & 0 & 38912 & 22568 & 9036\\
T11.1661 & \texttt{01233221211} & 7 & 4 & 3 & 8 & 8 & 5 & 4 & 1 & 36864 & 27008 & 5080\\
T11.1662 & \texttt{01233221111} & 8 & 4 & 3 & 96 & 6 & 3 & 3 & 0 & 34816 & 28488 & 2980\\
T11.1663 & \texttt{01233212332} & 6 & 4 & 5 & 8 & 5 & 1 & 2 & 0 & 45056 & 23160 & 11940\\
T11.1664 & \texttt{01233212331} & 6 & 4 & 4 & 4 & 9 & 7 & 4 & 2 & 43008 & 23320 & 10512\\
T11.1665 & \texttt{01233212323} & 5 & 4 & 4 & 4 & 8 & 5 & 3 & 1 & 45056 & 18408 & 14708\\
T11.1666 & \texttt{01233212322} & 6 & 4 & 4 & 4 & 9 & 7 & 4 & 2 & 43008 & 23616 & 10812\\
T11.1667 & \texttt{01233212321} & 6 & 4 & 4 & 2 & 10 & 9 & 5 & 4 & 40960 & 24160 & 9044\\
T11.1668 & \texttt{01233212312} & 5 & 4 & 3 & 2 & 10 & 9 & 4 & 3 & 40960 & 17400 & 12628\\
T11.1669 & \texttt{01233212311} & 6 & 4 & 3 & 4 & 9 & 7 & 4 & 2 & 38912 & 23280 & 8160\\
T11.1670 & \texttt{01233212222} & 7 & 4 & 4 & 48 & 7 & 5 & 3 & 1 & 40960 & 25352 & 8828\\
T11.1671 & \texttt{01233212221} & 7 & 4 & 4 & 12 & 8 & 6 & 4 & 2 & 38912 & 26280 & 6636\\
T11.1672 & \texttt{01233212212} & 6 & 4 & 4 & 4 & 9 & 7 & 4 & 2 & 38912 & 21456 & 9416\\
T11.1673 & \texttt{01233212211} & 7 & 4 & 4 & 8 & 8 & 5 & 4 & 1 & 36864 & 25656 & 5668\\
T11.1674 & \texttt{01233212121} & 6 & 4 & 3 & 4 & 8 & 5 & 4 & 2 & 36864 & 22440 & 7748\\
T11.1675 & \texttt{01233212111} & 7 & 4 & 3 & 12 & 8 & 6 & 4 & 2 & 34816 & 24960 & 5056\\
T11.1676 & \texttt{01233211111} & 8 & 4 & 3 & 240 & 6 & 4 & 3 & 1 & 32768 & 25800 & 3420\\
T11.1677 & \texttt{01233123312} & 5 & 4 & 3 & 8 & 6 & 3 & 2 & 1 & 43008 & 15440 & 14736\\
T11.1678 & \texttt{01233123311} & 6 & 4 & 3 & 16 & 5 & 1 & 2 & 0 & 40960 & 22200 & 9916\\
T11.1679 & \texttt{01233123232} & 5 & 4 & 3 & 4 & 8 & 5 & 3 & 1 & 45056 & 18696 & 14672\\
T11.1680 & \texttt{01233123231} & 5 & 4 & 3 & 4 & 8 & 5 & 3 & 1 & 43008 & 19208 & 13008\\
T11.1681 & \texttt{01233123222} & 6 & 4 & 3 & 12 & 8 & 6 & 3 & 1 & 43008 & 21984 & 11904\\
T11.1682 & \texttt{01233123221} & 6 & 4 & 3 & 4 & 9 & 7 & 4 & 2 & 40960 & 23456 & 9520\\
T11.1683 & \texttt{01233123212} & 5 & 4 & 3 & 2 & 10 & 9 & 4 & 3 & 40960 & 17000 & 12988\\
T11.1684 & \texttt{01233123211} & 6 & 4 & 3 & 4 & 9 & 7 & 4 & 2 & 38912 & 22720 & 8616\\
T11.1685 & \texttt{01233123123} & 4 & 4 & 2 & 4 & 7 & 3 & 2 & 0 & 43008 & 10280 & 16928\\
T11.1686 & \texttt{01233123122} & 5 & 4 & 3 & 4 & 9 & 7 & 3 & 1 & 40960 & 15456 & 13468\\
T11.1687 & \texttt{01233123121} & 5 & 4 & 2 & 2 & 10 & 9 & 4 & 3 & 38912 & 17560 & 11328\\
T11.1688 & \texttt{01233123111} & 6 & 4 & 2 & 12 & 8 & 6 & 3 & 1 & 36864 & 21200 & 8164\\
T11.1689 & \texttt{01233122222} & 7 & 4 & 3 & 240 & 6 & 4 & 2 & 0 & 40960 & 23080 & 10456\\
T11.1690 & \texttt{01233122221} & 7 & 4 & 3 & 48 & 7 & 5 & 3 & 1 & 38912 & 24872 & 7608\\
T11.1691 & \texttt{01233122212} & 6 & 4 & 3 & 12 & 8 & 6 & 3 & 1 & 38912 & 19840 & 10372\\
T11.1692 & \texttt{01233122211} & 7 & 4 & 3 & 24 & 7 & 4 & 3 & 0 & 36864 & 24520 & 6344\\
T11.1693 & \texttt{01233122122} & 6 & 4 & 4 & 16 & 6 & 3 & 2 & 0 & 38912 & 19592 & 10248\\
T11.1694 & \texttt{01233122121} & 6 & 4 & 3 & 4 & 9 & 7 & 4 & 2 & 36864 & 20656 & 8556\\
T11.1695 & \texttt{01233122111} & 7 & 4 & 3 & 24 & 7 & 4 & 3 & 0 & 34816 & 23256 & 5840\\
T11.1696 & \texttt{01233121212} & 5 & 4 & 2 & 12 & 6 & 3 & 2 & 0 & 36864 & 16000 & 11340\\
T11.1697 & \texttt{01233121211} & 6 & 4 & 2 & 8 & 7 & 3 & 3 & 0 & 34816 & 20680 & 7520\\
T11.1698 & \texttt{01233121111} & 7 & 4 & 2 & 48 & 7 & 5 & 3 & 1 & 32768 & 22240 & 5460\\
T11.1699 & \texttt{01233111111} & 8 & 4 & 2 & 1440 & 5 & 3 & 2 & 0 & 30720 & 22760 & 4056\\
T11.1700 & \texttt{01232323232} & 5 & 4 & 1 & 24 & 5 & 3 & 2 & 1 & 47104 & 23220 & 14344\\
T11.1701 & \texttt{01232323231} & 5 & 4 & 2 & 24 & 5 & 3 & 2 & 1 & 45056 & 24160 & 12164\\
T11.1702 & \texttt{01232323222} & 6 & 4 & 1 & 36 & 5 & 2 & 2 & 0 & 45056 & 26136 & 11888\\
T11.1703 & \texttt{01232323221} & 6 & 4 & 2 & 12 & 6 & 3 & 3 & 1 & 43008 & 28480 & 8588\\
T11.1704 & \texttt{01232323212} & 5 & 4 & 2 & 6 & 7 & 5 & 3 & 2 & 43008 & 21560 & 12536\\
T11.1705 & \texttt{01232323211} & 6 & 4 & 2 & 12 & 6 & 3 & 3 & 1 & 40960 & 27760 & 7532\\
T11.1706 & \texttt{01232323123} & 4 & 4 & 2 & 6 & 7 & 5 & 2 & 1 & 45056 & 14056 & 17160\\
T11.1707 & \texttt{01232323122} & 5 & 4 & 3 & 12 & 6 & 3 & 2 & 0 & 43008 & 19232 & 13524\\
T11.1708 & \texttt{01232323121} & 5 & 4 & 2 & 6 & 7 & 5 & 3 & 2 & 40960 & 21336 & 11208\\
T11.1709 & \texttt{01232323111} & 6 & 4 & 2 & 36 & 5 & 2 & 2 & 0 & 38912 & 24976 & 7868\\
T11.1710 & \texttt{01232322222} & 7 & 4 & 1 & 240 & 5 & 2 & 2 & 0 & 43008 & 27108 & 10656\\
T11.1711 & \texttt{01232322221} & 7 & 4 & 2 & 48 & 6 & 3 & 3 & 1 & 40960 & 29920 & 6764\\
T11.1712 & \texttt{01232322212} & 6 & 4 & 2 & 12 & 7 & 4 & 3 & 1 & 40960 & 24512 & 9904\\
T11.1713 & \texttt{01232322211} & 7 & 4 & 2 & 24 & 6 & 2 & 3 & 0 & 38912 & 29992 & 5052\\
T11.1714 & \texttt{01232322123} & 5 & 4 & 2 & 4 & 8 & 5 & 3 & 1 & 43008 & 19744 & 13152\\
T11.1715 & \texttt{01232322122} & 6 & 4 & 3 & 8 & 7 & 3 & 3 & 0 & 40960 & 24200 & 9652\\
T11.1716 & \texttt{01232322121} & 6 & 4 & 2 & 4 & 8 & 5 & 4 & 2 & 38912 & 25584 & 7616\\
T11.1717 & \texttt{01232322111} & 7 & 4 & 2 & 24 & 6 & 2 & 3 & 0 & 36864 & 28504 & 4556\\
T11.1718 & \texttt{01232321232} & 5 & 4 & 3 & 2 & 9 & 7 & 4 & 3 & 43008 & 20336 & 12836\\
T11.1719 & \texttt{01232321231} & 5 & 4 & 2 & 2 & 9 & 7 & 4 & 3 & 40960 & 21056 & 10976\\
T11.1720 & \texttt{01232321222} & 6 & 4 & 3 & 12 & 7 & 4 & 3 & 1 & 40960 & 23256 & 10140\\
T11.1721 & \texttt{01232321221} & 6 & 4 & 3 & 4 & 8 & 5 & 4 & 2 & 38912 & 24232 & 8116\\
T11.1722 & \texttt{01232321212} & 5 & 4 & 2 & 4 & 7 & 3 & 3 & 1 & 38912 & 19416 & 11092\\
T11.1723 & \texttt{01232321211} & 6 & 4 & 2 & 4 & 8 & 5 & 4 & 2 & 36864 & 24336 & 6976\\
T11.1724 & \texttt{01232321111} & 7 & 4 & 2 & 48 & 6 & 3 & 3 & 1 & 34816 & 25976 & 4812\\
T11.1725 & \texttt{01232312323} & 4 & 4 & 3 & 8 & 4 & 1 & 1 & 0 & 45056 & 13336 & 17540\\
T11.1726 & \texttt{01232312322} & 5 & 4 & 3 & 4 & 8 & 5 & 3 & 1 & 43008 & 19024 & 13548\\
T11.1727 & \texttt{01232312321} & 5 & 4 & 3 & 2 & 9 & 7 & 4 & 3 & 40960 & 20208 & 11492\\
T11.1728 & \texttt{01232312312} & 4 & 4 & 2 & 2 & 9 & 7 & 3 & 2 & 40960 & 12648 & 15332\\
T11.1729 & \texttt{01232312311} & 5 & 4 & 2 & 4 & 8 & 5 & 3 & 1 & 38912 & 19648 & 10464\\
T11.1730 & \texttt{01232312222} & 6 & 4 & 3 & 48 & 6 & 3 & 2 & 0 & 40960 & 20920 & 11564\\
T11.1731 & \texttt{01232312221} & 6 & 4 & 3 & 12 & 7 & 4 & 3 & 1 & 38912 & 22488 & 9052\\
T11.1732 & \texttt{01232312212} & 5 & 4 & 3 & 4 & 8 & 5 & 3 & 1 & 38912 & 17184 & 12008\\
T11.1733 & \texttt{01232312211} & 6 & 4 & 3 & 8 & 7 & 3 & 3 & 0 & 36864 & 22184 & 7924\\
T11.1734 & \texttt{01232312121} & 5 & 4 & 2 & 4 & 7 & 3 & 3 & 1 & 36864 & 18648 & 10132\\
T11.1735 & \texttt{01232312111} & 6 & 4 & 2 & 12 & 7 & 4 & 3 & 1 & 34816 & 21648 & 7232\\
T11.1736 & \texttt{01232311111} & 7 & 4 & 2 & 240 & 5 & 2 & 2 & 0 & 32768 & 22648 & 5516\\
T11.1737 & \texttt{01232222222} & 8 & 4 & 1 & 5040 & 5 & 4 & 2 & 1 & 40960 & 27432 & 9880\\
T11.1738 & \texttt{01232222221} & 8 & 4 & 2 & 720 & 6 & 5 & 3 & 2 & 38912 & 30400 & 5604\\
T11.1739 & \texttt{01232222212} & 7 & 4 & 2 & 120 & 7 & 6 & 3 & 2 & 38912 & 25496 & 8512\\
T11.1740 & \texttt{01232222211} & 8 & 4 & 2 & 240 & 6 & 4 & 3 & 1 & 36864 & 30736 & 3556\\
T11.1741 & \texttt{01232222123} & 6 & 4 & 2 & 24 & 8 & 7 & 3 & 2 & 40960 & 21640 & 11248\\
T11.1742 & \texttt{01232222122} & 7 & 4 & 3 & 48 & 7 & 5 & 3 & 1 & 38912 & 25856 & 7724\\
T11.1743 & \texttt{01232222121} & 7 & 4 & 2 & 24 & 8 & 7 & 4 & 3 & 36864 & 27000 & 5712\\
T11.1744 & \texttt{01232222111} & 8 & 4 & 2 & 144 & 6 & 4 & 3 & 1 & 34816 & 29680 & 2676\\
T11.1745 & \texttt{01232221232} & 6 & 4 & 3 & 6 & 9 & 8 & 4 & 3 & 40960 & 23384 & 10108\\
T11.1746 & \texttt{01232221231} & 6 & 4 & 2 & 6 & 9 & 8 & 4 & 3 & 38912 & 23624 & 8400\\
T11.1747 & \texttt{01232221222} & 7 & 4 & 3 & 36 & 7 & 5 & 3 & 1 & 38912 & 26064 & 7428\\
T11.1748 & \texttt{01232221221} & 7 & 4 & 3 & 12 & 8 & 6 & 4 & 2 & 36864 & 26560 & 5580\\
T11.1749 & \texttt{01232221212} & 6 & 4 & 2 & 12 & 7 & 4 & 3 & 1 & 36864 & 22224 & 8428\\
T11.1750 & \texttt{01232221211} & 7 & 4 & 2 & 12 & 8 & 6 & 4 & 2 & 34816 & 26424 & 4528\\
T11.1751 & \texttt{01232221111} & 8 & 4 & 2 & 144 & 6 & 4 & 3 & 1 & 32768 & 27824 & 2452\\
T11.1752 & \texttt{01232212322} & 6 & 4 & 3 & 8 & 5 & 1 & 2 & 0 & 40960 & 23992 & 9700\\
T11.1753 & \texttt{01232212321} & 6 & 4 & 3 & 2 & 10 & 9 & 5 & 4 & 38912 & 24216 & 8076\\
T11.1754 & \texttt{01232212312} & 5 & 4 & 2 & 2 & 10 & 9 & 4 & 3 & 38912 & 17856 & 11532\\
T11.1755 & \texttt{01232212311} & 6 & 4 & 2 & 4 & 9 & 7 & 4 & 2 & 36864 & 23176 & 7264\\
T11.1756 & \texttt{01232212222} & 7 & 4 & 3 & 48 & 7 & 5 & 3 & 1 & 38912 & 25648 & 7716\\
T11.1757 & \texttt{01232212221} & 7 & 4 & 3 & 12 & 8 & 6 & 4 & 2 & 36864 & 26256 & 5684\\
T11.1758 & \texttt{01232212212} & 6 & 4 & 3 & 4 & 9 & 7 & 4 & 2 & 36864 & 21672 & 8376\\
T11.1759 & \texttt{01232212211} & 7 & 4 & 3 & 8 & 8 & 5 & 4 & 1 & 34816 & 25472 & 4796\\
T11.1760 & \texttt{01232212121} & 6 & 4 & 2 & 4 & 8 & 5 & 4 & 2 & 34816 & 22416 & 6812\\
T11.1761 & \texttt{01232212111} & 7 & 4 & 2 & 12 & 8 & 6 & 4 & 2 & 32768 & 24696 & 4224\\
T11.1762 & \texttt{01232211111} & 8 & 4 & 2 & 240 & 6 & 4 & 3 & 1 & 30720 & 25456 & 2628\\
T11.1763 & \texttt{01232123212} & 5 & 4 & 3 & 2 & 7 & 3 & 3 & 1 & 38912 & 18320 & 11320\\
T11.1764 & \texttt{01232123211} & 6 & 4 & 3 & 4 & 6 & 1 & 3 & 0 & 36864 & 23160 & 7348\\
T11.1765 & \texttt{01232123123} & 4 & 4 & 2 & 2 & 8 & 5 & 3 & 2 & 40960 & 12240 & 15092\\
T11.1766 & \texttt{01232123122} & 5 & 4 & 3 & 2 & 10 & 9 & 4 & 3 & 38912 & 16856 & 11776\\
T11.1767 & \texttt{01232123121} & 5 & 4 & 2 & 1 & 11 & 11 & 5 & 5 & 36864 & 18400 & 9892\\
T11.1768 & \texttt{01232123111} & 6 & 4 & 2 & 6 & 9 & 8 & 4 & 3 & 34816 & 21480 & 6984\\
T11.1769 & \texttt{01232122222} & 7 & 4 & 3 & 120 & 7 & 6 & 3 & 2 & 38912 & 24400 & 8636\\
T11.1770 & \texttt{01232122221} & 7 & 4 & 3 & 24 & 8 & 7 & 4 & 3 & 36864 & 25552 & 6172\\
T11.1771 & \texttt{01232122212} & 6 & 4 & 3 & 6 & 9 & 8 & 4 & 3 & 36864 & 20920 & 8760\\
T11.1772 & \texttt{01232122211} & 7 & 4 & 3 & 12 & 8 & 6 & 4 & 2 & 34816 & 24880 & 5100\\
T11.1773 & \texttt{01232122122} & 6 & 4 & 4 & 8 & 7 & 5 & 3 & 2 & 36864 & 20592 & 8684\\
T11.1774 & \texttt{01232122121} & 6 & 4 & 3 & 2 & 10 & 9 & 5 & 4 & 34816 & 21256 & 7200\\
T11.1775 & \texttt{01232122111} & 7 & 4 & 3 & 12 & 8 & 6 & 4 & 2 & 32768 & 23456 & 4692\\
T11.1776 & \texttt{01232121212} & 5 & 4 & 2 & 6 & 7 & 5 & 3 & 2 & 34816 & 17080 & 9776\\
T11.1777 & \texttt{01232121211} & 6 & 4 & 2 & 4 & 8 & 5 & 4 & 2 & 32768 & 21040 & 6292\\
T11.1778 & \texttt{01232121111} & 7 & 4 & 2 & 24 & 8 & 7 & 4 & 3 & 30720 & 22360 & 4360\\
T11.1779 & \texttt{01232111111} & 8 & 4 & 2 & 720 & 6 & 5 & 3 & 2 & 28672 & 22800 & 3004\\
T11.1780 & \texttt{01231231231} & 4 & 4 & 1 & 6 & 5 & 2 & 2 & 1 & 38912 & 13600 & 13136\\
T11.1781 & \texttt{01231231222} & 5 & 4 & 2 & 12 & 6 & 2 & 2 & 0 & 38912 & 15480 & 12428\\
T11.1782 & \texttt{01231231221} & 5 & 4 & 2 & 4 & 7 & 3 & 3 & 1 & 36864 & 17096 & 10276\\
T11.1783 & \texttt{01231231212} & 4 & 4 & 1 & 4 & 6 & 1 & 2 & 0 & 36864 & 11640 & 13316\\
T11.1784 & \texttt{01231231211} & 5 & 4 & 1 & 4 & 7 & 3 & 3 & 1 & 34816 & 17520 & 9072\\
T11.1785 & \texttt{01231231111} & 6 & 4 & 1 & 48 & 5 & 1 & 2 & 0 & 32768 & 19480 & 6844\\
T11.1786 & \texttt{01231222222} & 7 & 4 & 2 & 720 & 6 & 5 & 2 & 1 & 38912 & 22272 & 10168\\
T11.1787 & \texttt{01231222221} & 7 & 4 & 2 & 120 & 7 & 6 & 3 & 2 & 36864 & 24480 & 6904\\
T11.1788 & \texttt{01231222212} & 6 & 4 & 2 & 24 & 8 & 7 & 3 & 2 & 36864 & 19592 & 9580\\
T11.1789 & \texttt{01231222211} & 7 & 4 & 2 & 48 & 7 & 5 & 3 & 1 & 34816 & 24352 & 5400\\
T11.1790 & \texttt{01231222122} & 6 & 4 & 3 & 12 & 8 & 6 & 3 & 1 & 36864 & 19536 & 9200\\
T11.1791 & \texttt{01231222121} & 6 & 4 & 2 & 6 & 9 & 8 & 4 & 3 & 34816 & 20520 & 7516\\
T11.1792 & \texttt{01231222111} & 7 & 4 & 2 & 36 & 7 & 5 & 3 & 1 & 32768 & 23040 & 4808\\
T11.1793 & \texttt{01231221221} & 6 & 4 & 3 & 8 & 7 & 5 & 3 & 2 & 34816 & 20032 & 7560\\
T11.1794 & \texttt{01231221212} & 5 & 4 & 2 & 4 & 8 & 5 & 3 & 1 & 34816 & 15696 & 10200\\
T11.1795 & \texttt{01231221211} & 6 & 4 & 2 & 4 & 9 & 7 & 4 & 2 & 32768 & 19896 & 6612\\
T11.1796 & \texttt{01231221111} & 7 & 4 & 2 & 48 & 7 & 5 & 3 & 1 & 30720 & 21296 & 4640\\
T11.1797 & \texttt{01231212121} & 5 & 4 & 1 & 6 & 7 & 5 & 3 & 2 & 32768 & 16680 & 8628\\
T11.1798 & \texttt{01231212111} & 6 & 4 & 1 & 12 & 7 & 4 & 3 & 1 & 30720 & 19200 & 6032\\
T11.1799 & \texttt{01231211111} & 7 & 4 & 1 & 120 & 7 & 6 & 3 & 2 & 28672 & 20040 & 4428\\
T11.1800 & \texttt{01231111111} & 8 & 4 & 1 & 5040 & 5 & 4 & 2 & 1 & 26624 & 20320 & 3192\\
T11.1801 & \texttt{01222222222} & 9 & 3 & 1 & 362880 & 3 & 2 & 1 & 0 & 38912 & 27540 & 9272\\
T11.1802 & \texttt{01222222221} & 9 & 3 & 2 & 40320 & 4 & 3 & 2 & 1 & 36864 & 30560 & 4692\\
T11.1803 & \texttt{01222222212} & 8 & 3 & 2 & 5040 & 5 & 4 & 2 & 1 & 36864 & 25824 & 7592\\
T11.1804 & \texttt{01222222211} & 9 & 3 & 2 & 10080 & 4 & 2 & 2 & 0 & 34816 & 30984 & 2436\\
T11.1805 & \texttt{01222222122} & 8 & 3 & 3 & 1440 & 5 & 3 & 2 & 0 & 36864 & 26408 & 6556\\
T11.1806 & \texttt{01222222121} & 8 & 3 & 2 & 720 & 6 & 5 & 3 & 2 & 34816 & 27472 & 4472\\
T11.1807 & \texttt{01222222111} & 9 & 3 & 2 & 4320 & 4 & 2 & 2 & 0 & 32768 & 30072 & 1364\\
T11.1808 & \texttt{01222221222} & 8 & 3 & 3 & 720 & 5 & 3 & 2 & 0 & 36864 & 27000 & 5924\\
T11.1809 & \texttt{01222221221} & 8 & 3 & 3 & 240 & 6 & 4 & 3 & 1 & 34816 & 27336 & 4060\\
T11.1810 & \texttt{01222221212} & 7 & 3 & 2 & 240 & 5 & 2 & 2 & 0 & 34816 & 23160 & 6940\\
T11.1811 & \texttt{01222221211} & 8 & 3 & 2 & 240 & 6 & 4 & 3 & 1 & 32768 & 27120 & 3000\\
T11.1812 & \texttt{01222221111} & 9 & 3 & 2 & 2880 & 4 & 2 & 2 & 0 & 30720 & 28440 & 916\\
T11.1813 & \texttt{01222212222} & 8 & 3 & 3 & 1152 & 3 & 1 & 1 & 0 & 36864 & 27224 & 5716\\
T11.1814 & \texttt{01222212221} & 8 & 3 & 3 & 144 & 6 & 4 & 3 & 1 & 34816 & 27512 & 3780\\
T11.1815 & \texttt{01222212212} & 7 & 3 & 3 & 48 & 7 & 5 & 3 & 1 & 34816 & 23168 & 6432\\
T11.1816 & \texttt{01222212211} & 8 & 3 & 3 & 96 & 6 & 3 & 3 & 0 & 32768 & 26568 & 2940\\
T11.1817 & \texttt{01222212121} & 7 & 3 & 2 & 48 & 6 & 3 & 3 & 1 & 32768 & 23672 & 4924\\
T11.1818 & \texttt{01222212111} & 8 & 3 & 2 & 144 & 6 & 4 & 3 & 1 & 30720 & 25712 & 2392\\
T11.1819 & \texttt{01222211111} & 9 & 3 & 2 & 2880 & 4 & 2 & 2 & 0 & 28672 & 26392 & 820\\
T11.1820 & \texttt{01222122212} & 7 & 3 & 3 & 72 & 5 & 3 & 2 & 1 & 34816 & 23280 & 6232\\
T11.1821 & \texttt{01222122211} & 8 & 3 & 3 & 144 & 4 & 1 & 2 & 0 & 32768 & 26520 & 2868\\
T11.1822 & \texttt{01222122122} & 7 & 3 & 4 & 48 & 5 & 2 & 2 & 0 & 34816 & 22872 & 6196\\
T11.1823 & \texttt{01222122121} & 7 & 3 & 3 & 12 & 8 & 6 & 4 & 2 & 32768 & 23136 & 4880\\
T11.1824 & \texttt{01222122111} & 8 & 3 & 3 & 72 & 6 & 3 & 3 & 0 & 30720 & 24936 & 2540\\
T11.1825 & \texttt{01222121212} & 6 & 3 & 2 & 36 & 5 & 2 & 2 & 0 & 32768 & 19440 & 7296\\
T11.1826 & \texttt{01222121211} & 7 & 3 & 2 & 24 & 6 & 2 & 3 & 0 & 30720 & 22680 & 4076\\
T11.1827 & \texttt{01222121111} & 8 & 3 & 2 & 144 & 6 & 4 & 3 & 1 & 28672 & 23760 & 2248\\
T11.1828 & \texttt{01222111111} & 9 & 3 & 2 & 4320 & 4 & 2 & 2 & 0 & 26624 & 24120 & 932\\
T11.1829 & \texttt{01221221221} & 7 & 3 & 4 & 48 & 4 & 2 & 2 & 1 & 32768 & 22568 & 4972\\
T11.1830 & \texttt{01221221212} & 6 & 3 & 3 & 16 & 5 & 1 & 2 & 0 & 32768 & 19032 & 7276\\
T11.1831 & \texttt{01221221211} & 7 & 3 & 3 & 16 & 6 & 3 & 3 & 1 & 30720 & 22032 & 4232\\
T11.1832 & \texttt{01221221111} & 8 & 3 & 3 & 192 & 4 & 1 & 2 & 0 & 28672 & 23032 & 2468\\
T11.1833 & \texttt{01221212121} & 6 & 3 & 2 & 12 & 6 & 3 & 3 & 1 & 30720 & 19296 & 6040\\
T11.1834 & \texttt{01221212111} & 7 & 3 & 2 & 24 & 6 & 2 & 3 & 0 & 28672 & 21096 & 3780\\
T11.1835 & \texttt{01221211111} & 8 & 3 & 2 & 240 & 6 & 4 & 3 & 1 & 26624 & 21696 & 2304\\
T11.1836 & \texttt{01221111111} & 9 & 3 & 2 & 10080 & 4 & 2 & 2 & 0 & 24576 & 21896 & 1116\\
T11.1837 & \texttt{01212121212} & 5 & 3 & 1 & 120 & 3 & 1 & 1 & 0 & 30720 & 15600 & 8360\\
T11.1838 & \texttt{01212121211} & 6 & 3 & 1 & 48 & 4 & 1 & 2 & 0 & 28672 & 18840 & 5284\\
T11.1839 & \texttt{01212121111} & 7 & 3 & 1 & 144 & 4 & 1 & 2 & 0 & 26624 & 19920 & 3504\\
T11.1840 & \texttt{01212111111} & 8 & 3 & 1 & 1440 & 4 & 1 & 2 & 0 & 24576 & 20280 & 2204\\
T11.1841 & \texttt{01211111111} & 9 & 3 & 1 & 40320 & 4 & 3 & 2 & 1 & 22528 & 20400 & 1080\\
T11.1842 & \texttt{01111111111} & 10 & 2 & 1 & 3628800 & 2 & 1 & 1 & 0 & 20480 & 20440 & 20\\
\bottomrule
\end{longtable}}


\begin{thebibliography}{99}\small
\bibitem[Aho et~al.(1974)]{AhoHopcroftUllman1974} Aho, A.~V., Hopcroft, J.~E. and Ullman, J.~D. (1974). \emph{The Design and Analysis of Computer Algorithms}. Addison-Wesley.
\bibitem[Belnap(1992)]{Belnap1992} Belnap, N. (1992). Branching space-time. \emph{Synthese}, 92(3), 385--434.
\bibitem[Belnap and Green(1994)]{BelnapGreen1994} Belnap, N. and Green, M. (1994). Indeterminism and the thin red line. \emph{Philosophical Perspectives}, 8, 365--388.
\bibitem[Belnap et~al.(2001)]{Belnap2001} Belnap, N., Perloff, M. and Xu, M. (2001). \emph{Facing the Future: Agents and Choices in Our Indeterminist World}. Oxford University Press.
\bibitem[Beyer and Hedetniemi(1980)]{BeyerHedetniemi1980} Beyer, T. and Hedetniemi, S.~M. (1980). Constant time generation of rooted trees. \emph{SIAM Journal on Computing}, 9(4), 706--712.
\bibitem[Bombelli et~al.(1987)]{Bombelli1987} Bombelli, L., Lee, J., Meyer, D. and Sorkin, R.~D. (1987). Space-time as a causal set. \emph{Physical Review Letters}, 59(5), 521--524.
\bibitem[Bourne(2002)]{Bourne2002} Bourne, C. (2002). When am I? A tense time for some tense theorists? \emph{Australasian Journal of Philosophy}, 80(3), 359--371.
\bibitem[Braddon-Mitchell(2004)]{BraddonMitchell2004} Braddon-Mitchell, D. (2004). How do we know it is now now? \emph{Analysis}, 64(3), 199--203.
\bibitem[Broad(1923)]{Broad1923} Broad, C.~D. (1923). \emph{Scientific Thought}. Kegan Paul.
\bibitem[Broad(1938)]{Broad1938} Broad, C.~D. (1938). \emph{An Examination of McTaggart's Philosophy}, Vol.~II. Cambridge University Press.
\bibitem[Cameron(2015)]{Cameron2015} Cameron, R.~P. (2015). \emph{The Moving Spotlight: An Essay on Time and Ontology}. Oxford University Press.
\bibitem[Cayley(1857)]{Cayley1857} Cayley, A. (1857). On the theory of the analytical forms called trees. \emph{Philosophical Magazine}, 13, 172--176.
\bibitem[Correia and Rosenkranz(2018)]{CorreiaRosenkranz2018} Correia, F. and Rosenkranz, S. (2018). \emph{Nothing to Come: A Defence of the Growing Block Theory of Time}. Springer.
\bibitem[Dummett(1960)]{Dummett1960} Dummett, M. (1960). A defense of McTaggart's proof of the unreality of time. \emph{The Philosophical Review}, 69(4), 497--504.
\bibitem[Fine(2005)]{Fine2005} Fine, K. (2005). Tense and reality. In \emph{Modality and Tense: Philosophical Papers}, 261--320. Oxford University Press.
\bibitem[Forrest(2004)]{Forrest2004} Forrest, P. (2004). The real but dead past: A reply to Braddon-Mitchell. \emph{Analysis}, 64(4), 358--362.
\bibitem[Geach(1969)]{Geach1969} Geach, P.~T. (1969). \emph{God and the Soul}. Routledge \& Kegan Paul.
\bibitem[Gödel(1949)]{Godel1949} Gödel, K. (1949). A remark about the relationship between relativity theory and idealistic philosophy. In P.~A. Schilpp (ed.), \emph{Albert Einstein: Philosopher-Scientist}, 557--562. Library of Living Philosophers.
\bibitem[Hawking et~al.(1976)]{HawkingKingMcCarthy1976} Hawking, S.~W., King, A.~R. and McCarthy, P.~J. (1976). A new topology for curved space-time which incorporates the causal, differential, and conformal structures. \emph{Journal of Mathematical Physics}, 17(2), 174--181.
\bibitem[Le~Poidevin and Mellor(1987)]{LePoidevinMellor1987} Le~Poidevin, R. and Mellor, D.~H. (1987). Time, change, and the \enquote{indexical fallacy}. \emph{Mind}, 96(384), 534--538.
\bibitem[Lewis(1979)]{Lewis1979} Lewis, D. (1979). Attitudes de dicto and de se. \emph{The Philosophical Review}, 88(4), 513--543.
\bibitem[Lowe(1987)]{Lowe1987} Lowe, E.~J. (1987). The indexical fallacy in McTaggart's proof of the unreality of time. \emph{Mind}, 96(381), 62--70.
\bibitem[MacFarlane(2003)]{MacFarlane2003} MacFarlane, J. (2003). Future contingents and relative truth. \emph{The Philosophical Quarterly}, 53(212), 321--336.
\bibitem[Malament(1977)]{Malament1977} Malament, D.~B. (1977). The class of continuous timelike curves determines the topology of spacetime. \emph{Journal of Mathematical Physics}, 18(7), 1399--1404.
\bibitem[Markosian(2004)]{Markosian2004} Markosian, N. (2004). A defense of presentism. In D.~W. Zimmerman (ed.), \emph{Oxford Studies in Metaphysics}, Vol.~1, 47--82. Oxford University Press.
\bibitem[Markosian(1993)]{Markosian1993} Markosian, N. (1993). How fast does time pass? \emph{Philosophy and Phenomenological Research}, 53(4), 829--844.
\bibitem[Maudlin(2007)]{Maudlin2007} Maudlin, T. (2007). \emph{The Metaphysics Within Physics}. Oxford University Press.
\bibitem[McTaggart(1908)]{McTaggart1908} McTaggart, J.~M.~E. (1908). The unreality of time. \emph{Mind}, 17(68), 457--474.
\bibitem[McTaggart(1927)]{McTaggart1927} McTaggart, J.~M.~E. (1927). \emph{The Nature of Existence}, Vol.~II, ed. C.~D. Broad. Cambridge University Press.
\bibitem[Mellor(1981)]{Mellor1981} Mellor, D.~H. (1981). \emph{Real Time}. Cambridge University Press.
\bibitem[Mellor(1998)]{Mellor1998} Mellor, D.~H. (1998). \emph{Real Time II}. Routledge.
\bibitem[Otter(1948)]{Otter1948} Otter, R. (1948). The number of trees. \emph{Annals of Mathematics}, 49(3), 583--599.
\bibitem[Perry(1979)]{Perry1979} Perry, J. (1979). The problem of the essential indexical. \emph{Noûs}, 13(1), 3--21.
\bibitem[Price(1996)]{Price1996} Price, H. (1996). \emph{Time's Arrow and Archimedes' Point}. Oxford University Press.
\bibitem[Prior(1957)]{Prior1957} Prior, A.~N. (1957). \emph{Time and Modality}. Clarendon Press.
\bibitem[Prior(1959)]{Prior1959} Prior, A.~N. (1959). Thank goodness that's over. \emph{Philosophy}, 34(128), 12--17.
\bibitem[Prior(1967)]{Prior1967} Prior, A.~N. (1967). \emph{Past, Present and Future}. Clarendon Press.
\bibitem[Prior(1968)]{Prior1968} Prior, A.~N. (1968). \emph{Papers on Time and Tense}. Clarendon Press.
\bibitem[Putnam(1967)]{Putnam1967} Putnam, H. (1967). Time and physical geometry. \emph{The Journal of Philosophy}, 64(8), 240--247.
\bibitem[Robb(1914)]{Robb1914} Robb, A.~A. (1914). \emph{A Theory of Time and Space}. Cambridge University Press.
\bibitem[Russell(1903)]{Russell1903} Russell, B. (1903). \emph{The Principles of Mathematics}. Cambridge University Press.
\bibitem[Russell(1915)]{Russell1915} Russell, B. (1915). On the experience of time. \emph{The Monist}, 25(2), 212--233.
\bibitem[Shoemaker(1969)]{Shoemaker1969} Shoemaker, S. (1969). Time without change. \emph{The Journal of Philosophy}, 66(12), 363--381.
\bibitem[Skow(2015)]{Skow2015} Skow, B. (2015). \emph{Objective Becoming}. Oxford University Press.
\bibitem[Smart(1963)]{Smart1963} Smart, J.~J.~C. (1963). \emph{Philosophy and Scientific Realism}. Routledge \& Kegan Paul.
\bibitem[Stein(1968)]{Stein1968} Stein, H. (1968). On Einstein--Minkowski space-time. \emph{The Journal of Philosophy}, 65(1), 5--23.
\bibitem[Thomason(1970)]{Thomason1970} Thomason, R.~H. (1970). Indeterminist time and truth-value gaps. \emph{Theoria}, 36(3), 264--281.
\bibitem[Tooley(1997)]{Tooley1997} Tooley, M. (1997). \emph{Time, Tense, and Causation}. Clarendon Press.
\bibitem[Williams(1951)]{Williams1951} Williams, D.~C. (1951). The myth of passage. \emph{The Journal of Philosophy}, 48(15), 457--472.
\bibitem[Zimmerman(2005)]{Zimmerman2005} Zimmerman, D.~W. (2005). The A-theory of time, the B-theory of time, and \enquote{taking tense seriously}. \emph{Dialectica}, 59(4), 401--457.
\end{thebibliography}

\end{document}
